% Auto-generated by the bl-book typesetter (Bernstein Logic literate rule
% catalogue) from the mechanized rule sources, revision 5a296a5.
% Do not hand-edit: regenerate with `bl-book --emit-fragment` and re-copy.
% Lean statements are extracted from the sources at emission and hash-bound;
% verifier receipts were consulted at emission (their in-text badges
% record checker-layer coverage of the bound builds, not a live
% re-check by the host build).
% This file: the catalogue body (preface, stamp, coverage dashboard,
% rule synopsis, the rules grouped in ToC sections, references).
% \input it inside a chapter of a host document whose preamble loaded
% the companion preamble file.

\begingroup
\setlength{\parindent}{0pt}
\setlength{\parskip}{4pt}

\textbf{What this catalogue is.}\enspace The proof rules of the
certificate calculus, in two strata sharing one page discipline. The
base stratum is \textbf{Bernstein Logic} --- a small proof calculus for
certificates about polynomial facts over networks of
rational-tensor-valued wires: sequents $\Phi \vdash \Psi$ asserting that
every admissible valuation of the wires satisfying the premise $\Phi$
satisfies the conclusion $\Psi$, with all trusted arithmetic exact over
$\Q$. Its rules appear here in three tiers --- core rules (kernel
semantic content), admissible rules (verified unpackings into core
rules), and alias rules (verified single-application spellings) ---
together with the composition rules (\rulename{Promise},
\rulename{Syzygy}, \rulename{SmallGain}, \rulename{SmallGainN},
\rulename{ProductLattice}),
which are Bernstein Logic outright. Above the base stratum sit the
\textbf{QTSL rules}: rules concluding satisfaction judgments of the
quantitative temporal specification language --- a system satisfies a
specification built from the specification chapter\textquotesingle s
formers --- through the battery rules (safety, deadline, until, and
Streett acceptance), the judgment-structure and specification-surface
rules, the Morse-topology rules, composition and aggregation,
distance transport, and the stored-chain certificate species
(eventual-absorption brackets and generator-form contraction tails,
recomputed exactly over stored chain rows), supported by discharge
kits and arithmetic leaves
that remain ordinary Bernstein Logic sequents. QTSL premises bottom out
in Bernstein Logic sequents, and both strata are designed for one
exact-replay discipline: a \emph{certificate} is an acyclic derivation
graph (judgment premises may be shared) whose nodes carry enough
payload that a small checker can replay each rule\textquotesingle s
side conditions by exact rational computation --- no search, no
optimization, no floating point in the trusted base. Where a
rule\textquotesingle s row and checker arm are enacted, the checker
replays it; the coverage dashboard names, rule by rule, what is enacted
and what is not.
Producers (LP/SDP pipelines, abstract interpreters, hand
constructions) do whatever clever thing \emph{finds} the certificate;
the checker only replays. The catalogue opens with a coverage
dashboard --- one row per rule, covered or not, mapping the current
level of mechanized verification of soundness, honest absences
included --- then a synopsis parade of every rule, then the rules
grouped by subject.

The validity judgment itself is not free-floating: it is
kernel-checked to be the possibilistic face of the semantic
universe\textquotesingle s satisfaction judgment
(chapter~\ref{ch:qtsl-semantic-universe}), at the sequent\textquotesingle s
exact valuation carrier --- a sequent is semantically valid iff the
upper probability of its conclusion\textquotesingle s complement
vanishes under the credal compile of its premise
cylinder\textquotesingle s possibilistic wall
(\texttt{semanticallyValid\_iff\_face}, module
\texttt{SuFunctor.Plain.CheckerFace}) --- so each soundness theorem
in this catalogue is, through that identification, a statement about
the semantic universe\textquotesingle s own judgment.

Each rule of the calculus gets one section here, carrying five
ingredients. (i)~The rule itself, typeset in Steele\textquotesingle s
Computer Science Metanotation~\cite{steele2017}, with prose on what the
rule does, when a producer uses it, and why each hypothesis is there.
(ii)~The Lean~4 declarations that realize it --- the payload structure,
the soundness theorem, and (where a side condition\textquotesingle s
meaning must be visible) the checker definitions in full --- extracted
mechanically from the sources at every build of this catalogue and bound by
hash, so quoted text cannot drift from the tree silently. (iii)~Verifier
receipts: attestations that named checking layers ---
elaboration, kernel replay by the official external checker, a hardened
screen, and a decorrelated second kernel --- accepted the \emph{exact
build} being quoted, identified by the SHA-256 of the compiled module;
every layer is a program the reader can re-run over the quoted sources.
(iv)~The HOL4 statement where the independent HOL4 development carries
one, and a loud, reasoned absence where it does not. (v)~A proof at
textbook detail: ordinary mathematics for core rules, explicit rewriting
recipes for admissible rules, filled-in applications of the target rule
for alias rules. Where a rule\textquotesingle s certificate content is
carried by the baked theory, the section additionally quotes the theory
rows --- payload schema, side-condition demand relations,
well-formedness axioms --- mechanically from the theory source at every
build, so a renamed or removed row fails the build rather than drifting
silently.

\textbf{The theory as source.}\enspace The normative home of the
rules\textquotesingle{} certificate content is the baked theory pair,
shipped as data in \texttt{theories/} beside this book:
\texttt{theory BernsteinLogic} (the base calculus\textquotesingle s
proof tower, from ground dtypes through hypernets, formulas, and
judgments to the proof species) and \texttt{theory Qtsl} (one omnibus
extension carrying the specification language, the System stratum, the
flavored cell schemas, the certificate sections --- judgment kinds,
proof-node payloads, batteries, the verdict form --- and the
certificate-shape stratum). Where this catalogue\textquotesingle s
statements and the theory disagree, the theory is the authority. Not
every catalogued rule has a baked row home: entries whose certificate
content is out of scope of the current baked theory say so outright ---
the rule remains catalogued, tiered, and proved at its stated scope,
and its schema realization is named as a successor rather than
implied.

\textbf{The certificate-shape stratum.}\enspace One stratum of the
baked theory types certificates \emph{as artifacts} rather than as
derivations, and its consumers are the battery rules of this catalogue,
so it is documented here. A \emph{nondeterminism axis}
(\texttt{NDAxis}) is an entity --- identity is the choice source ---
carrying a total binding time and extent: \texttt{btStep} (re-chosen at
every step; the axis enters drift reads cellwise as a region axis, and
no envelope object is ever materialized) or \texttt{btRun} (committed
once per run), with extent \texttt{ndEnum} (an enumerated branch
count) or \texttt{ndParam} (a box enclosure of a parametric domain). A
\emph{certificate family} (\texttt{CertFam}) declares its own axis
roster (\texttt{cert\_axis}); its branch-key space is the product of
its enumerated run-bound axes in position order, keys riding the
ground \texttt{Shape} sort as natural tuples --- the scalar certificate
is the nullary case (one branch under the width-zero key, not a second
sort). Two refusals shape the roster: a family may not split over a
step-bound axis, and may not enumerate a parametric one. Per-branch
positional bodies (\texttt{cert\_body}) carry the certificate
functions; position semantics belongs to the consuming rule. The
existential face\textquotesingle s slot is the id-keyed
\emph{selection} (\texttt{cert\_selection}: a named branch plus a
properness witness; \texttt{cert\_sel\_cell} refines a parametric
domain, zero cell rows reading as the whole domain) --- selections are
mandatory by default, waived only when every consuming conclusion is
safety-only. The universal face\textquotesingle s economies are
branch-discharge rows: an emptiness certificate
(\texttt{cert\_branch\_empty}, the one optional slot),
grade domination (\texttt{cert\_dominates}, irreflexive and acyclic),
and progressive-revelation hand-off
(\texttt{cert\_handoff}, prefix agreement between branch certificates).
The stratum\textquotesingle s structural demand is counting and ground
comparison --- admission, never geometry --- and its deliberate
absences are testable negatives: no envelope slot (a demonic envelope
is a consuming rule\textquotesingle s cellwise reading of the step
axes, never an interface object), no verdict rows (those live with the
certificate sections that consume them), no checker-arm vocabulary,
and no branch-key structure sort.

\textbf{Why trust a green cell.}\enspace A green receipt row reduces to
three things a skeptical reader can check independently: the statement
printed on the page (hash-bound to the source file it was extracted
from); the SHA-256 of the compiled artefact the receipt names; and the
checker layers the receipt credits --- each layer an independently
runnable program (elaboration, official kernel replay, a hardened
screen, a decorrelated second kernel), so a green cell never asks for
trust in this book or its receipt store: rebuild the module, re-run the
checkers, and compare.
Staleness is always visible --- a receipt attesting an older build, a
statement that moved since its proof was written, a missing ingredient
each render as an \emph{amber} badge, never as silence. Red is reserved
for outright verification failure. The coverage dashboard below records,
rule by rule, which ingredients exist and which do not; the empty cells
are deliberate, and are meant to be read.

\textbf{Naming conventions.}\enspace Judgment kinds carry a \texttt{Q}
prefix (\texttt{QSat}, \texttt{QRegionSat}, \ldots): they are the
satisfaction-judgment ($\vDash$-side) sorts of the extension theory, and
the prefix is their namespace. Rules that conclude a satisfaction
judgment carry the same \texttt{Q}; derivation-side ($\vdash$) leaf rules
stay bare or carry \texttt{P}, and mathematical eponyms are never forced
into descriptive molds. Rule names end in a species suffix where one
applies: \emph{-Battery} --- a certificate kit concluding a region
judgment; \emph{-Discharge} --- evidence concluding a named judgment
kind; \emph{-Compose} --- transport across a composition axis, named by
the axis; \emph{-Alias} --- a verified one-application spelling;
\emph{-Rewrite} --- an equivalence-concluding rewrite. A few rules carry
a former name (marked ``n\'e'') from earlier editions of this catalogue.


\par\medskip
{\footnotesize\textcolor{NoteGray}{Rendered 2026-08-04 03:24 UTC \textperiodcentered{} receipts consulted from the store (checker-layer attestations) \textperiodcentered{} tree \texttt{5a296a5}}}

\clearpage
\pdfbookmark[1]{Coverage dashboard}{blcat-dashboard}%
\dashboardhead{One table over every rule of the catalogue, whether or not this edition
covers it yet: a coverage dashboard that maps the current level of
actual mechanized verification of soundness, rule by rule. Each covered
section joins its ingredients from
their single sources of truth: the rule cell from the hand-authored rules
source (the same file the two-column diff pages render), the Lean
statements from the sources by declaration name at every render, the
receipts from the store (cited at render,
never re-run), the HOL4 statement from the HOL4 development or a reasoned
absence, and the proof from its hand-authored file. Rules with no catalogue
entry yet render as all-absent rows, loudly. Some rows carry a Lean
column without a catalogue section yet: their soundness theorems exist in the
mechanized development (two sibling Lean packages --- the core calculus
and an auxiliary measure-theoretic library --- both statement-bound by
the same hash machinery), while their proof and HOL columns remain
marked absent until those sections are authored. A module whose build has
no attestation yet shows its receipt layers as none: statements
are source-bound, and receipt coverage arrives when the checking gate
next runs over that package.
}
\dashrow{\hyperref[blcatrule-Ax]{Ax}~\dashsrc{BL}}{core}{\dok}{\dok}{\lcur{b}\lstale{r}\lstale{h}\lcur{d}}{\dok}{\ddoc}{\dok}
\dashrow{\hyperref[blcatrule-Weaken]{Weaken}~\dashsrc{BL}}{core}{\dok}{\dok}{\lcur{b}\lstale{r}\lstale{h}\lcur{d}}{\dok}{\ddoc}{\dok}
\dashrow{\hyperref[blcatrule-Cut]{Cut}~\dashsrc{BL}}{core}{\dok}{\dok}{\lcur{b}\lstale{r}\lstale{h}\lcur{d}}{\dok}{\ddoc}{\dok}
\dashrow{\hyperref[blcatrule-AndI]{$\wedge$I}~\dashsrc{BL}}{core}{\dok}{\dok}{\lcur{b}\lstale{r}\lstale{h}\lcur{d}}{\dok}{\ddoc}{\dok}
\dashrow{\hyperref[blcatrule-AndE]{$\wedge$E}~\dashsrc{BL}}{core}{\dok}{\dok}{\lcur{b}\lstale{r}\lstale{h}\lcur{d}}{\dok}{\ddoc}{\dok}
\dashrow{\hyperref[blcatrule-OrI]{$\vee$I}~\dashsrc{BL}}{core}{\dok}{\dok}{\lcur{b}\lstale{r}\lstale{h}\lcur{d}}{\dok}{\ddoc}{\dok}
\dashrow{\hyperref[blcatrule-OrE]{$\vee$E}~\dashsrc{BL}}{core}{\dok}{\dok}{\lcur{b}\lcur{r}\lcur{h}\lcur{d}}{\dok}{\ddoc}{\dok}
\dashrow{\hyperref[blcatrule-EqElim]{EqElim}~\dashsrc{BL}}{core}{\dok}{\dok}{\lcur{b}\lstale{r}\lstale{h}\lcur{d}}{\dok}{\ddoc}{\dok}
\dashrow{\hyperref[blcatrule-EqRefl]{EqRefl}~\dashsrc{BL}}{core}{\dok}{\dok}{\lcur{b}\lstale{r}\lstale{h}\lcur{d}}{\dok}{\ddoc}{\dok}
\dashrow{\hyperref[blcatrule-EqSym]{EqSym}~\dashsrc{BL}}{core}{\dok}{\dok}{\lcur{b}\lstale{r}\lstale{h}\lcur{d}}{\dok}{\ddoc}{\dok}
\dashrow{\hyperref[blcatrule-Antisym]{Antisym}~\dashsrc{BL}}{core}{\dok}{\dok}{\lcur{b}\lstale{r}\lstale{h}\lcur{d}}{\dok}{\ddoc}{\dok}
\dashrow{\hyperref[blcatrule-ConvElim]{ConvElim}~\dashsrc{BL}}{core}{\dok}{\dok}{\lcur{b}\lstale{r}\lstale{h}\lcur{d}}{\dok}{\ddoc}{\dok}
\dashrow{\hyperref[blcatrule-ConvIntro]{ConvIntro}~\dashsrc{BL}}{core}{\dok}{\dok}{\lcur{b}\lstale{r}\lstale{h}\lcur{d}}{\dok}{\ddoc}{\dok}
\dashrow{\hyperref[blcatrule-ArithAtom]{ArithAtom}~\dashsrc{BL}}{core}{\dok}{\dok}{\lcur{b}\lstale{r}\lstale{h}\lcur{d}}{\dok}{\ddoc}{\dok}
\dashrow{\hyperref[blcatrule-BernCone]{BernCone}~\dashsrc{BL}}{core}{\dok}{\damber}{\lcur{b}\lstale{r}\lstale{h}\lcur{d}}{\dok}{\ddoc}{\dok}
\dashrow{\hyperref[blcatrule-Syzygy]{Syzygy}~\dashsrc{BL}}{core}{\dok}{\dok}{\lcur{b}\lstale{r}\lstale{h}\lcur{d}}{\dok}{\ddoc}{\dok}
\dashrow{\hyperref[blcatrule-QGuardIntro]{QGuardIntro}~\dashsrc{QTSL}}{core}{\dok}{\dok}{\lcur{b}\lnone{r}\lnone{h}\lcur{d}}{\dnone}{\dnone}{\dok}
\dashrow{\hyperref[blcatrule-QGuardElim]{QGuardElim}~\dashsrc{QTSL}}{core}{\dok}{\dok}{\lcur{b}\lnone{r}\lnone{h}\lcur{d}}{\dnone}{\dnone}{\dok}
\dashrow{\hyperref[blcatrule-QSafetyBattery]{QSafetyBattery}~\dashsrc{QTSL}}{core}{\dok}{\damber}{\lcur{b}\lcur{r}\lcur{h}\lcur{d}}{\dnone}{\dnone}{\dok}
\dashrow{\hyperref[blcatrule-QDeadlineBattery]{QDeadlineBattery}~\dashsrc{QTSL}}{core}{\dok}{\damber}{\lcur{b}\lstale{r}\lstale{h}\lcur{d}}{\dnone}{\dnone}{\dok}
\dashrow{\hyperref[blcatrule-QUnionBound]{QUnionBound}~\dashsrc{QTSL}}{core}{\dok}{\damber}{\lcur{b}\lstale{r}\lstale{h}\lcur{d}}{\dnone}{\dnone}{\dok}
\dashrow{\hyperref[blcatrule-QStreettBattery]{QStreettBattery}~\dashsrc{QTSL}}{core}{\dok}{\dok}{\lcur{b}\lcur{r}\lcur{h}\lcur{d}}{\dnone}{\dnone}{\dok}
\dashrow{\hyperref[blcatrule-QAvgBattery]{QAvgBattery}~\dashsrc{QTSL}}{core}{\dnone}{\dnone}{\lnone{b}\lnone{r}\lnone{h}\lnone{d}}{\dnone}{\dnone}{\dok}
\dashrow{\hyperref[blcatrule-QMorseDischarge]{QMorseDischarge}~\dashsrc{QTSL}}{core}{\dcond}{\damber}{\lcur{b}\lcur{r}\lcur{h}\lcur{d}}{\dnone}{\dnone}{\dok}
\dashrow{\hyperref[blcatrule-QMorseProduct]{QMorseProduct}~\dashsrc{QTSL}}{core}{\dok}{\dnone}{\lcur{b}\lcur{r}\lcur{h}\lcur{d}}{\dnone}{\dnone}{\dok}
\dashrow{\hyperref[blcatrule-QBlockInv]{QBlockInv}~\dashsrc{QTSL}}{core}{\dcond}{\damber}{\lcur{b}\lcur{r}\lcur{h}\lcur{d}}{\dok}{\ddoc}{\dok}
\dashrow{\hyperref[blcatrule-QBasin]{QBasin}~\dashsrc{QTSL}}{core}{\dok}{\damber}{\lcur{b}\lstale{r}\lstale{h}\lcur{d}}{\dnone}{\dnone}{\dok}
\dashrow{\hyperref[blcatrule-QZeroByMiranda]{QZeroByMiranda}~\dashsrc{QTSL}}{core}{\dnone}{\dnone}{\lnone{b}\lnone{r}\lnone{h}\lnone{d}}{\dnone}{\dnone}{\dok}
\dashrow{\hyperref[blcatrule-QCoverByMiranda]{QCoverByMiranda}~\dashsrc{QTSL}}{core}{\dnone}{\dnone}{\lnone{b}\lnone{r}\lnone{h}\lnone{d}}{\dnone}{\dnone}{\dok}
\dashrow{\hyperref[blcatrule-QWireCompose]{QWireCompose}~\dashsrc{QTSL}}{core}{\dcond}{\damber}{\lcur{b}\lcur{r}\lcur{h}\lcur{d}}{\dnone}{\dnone}{\dok}
\dashrow{\hyperref[blcatrule-QWireCertDischarge]{QWireCertDischarge}~\dashsrc{QTSL}}{core}{\dcond}{\dok}{\lcur{b}\lcur{r}\lcur{h}\lcur{d}}{\dnone}{\dnone}{\dok}
\dashrow{\hyperref[blcatrule-QWireComposeMarks]{QWireComposeMarks}~\dashsrc{QTSL}}{core}{\dcond}{\damber}{\lstale{b}\lnone{r}\lnone{h}\lstale{d}}{\dnone}{\dnone}{\dok}
\dashrow{\hyperref[blcatrule-QSimTransport]{QSimTransport}~\dashsrc{QTSL}}{core}{\dok}{\dok}{\lcur{b}\lcur{r}\lcur{h}\lcur{d}}{\dnone}{\dnone}{\dok}
\dashrow{\hyperref[blcatrule-QLinGrowthDischarge]{QLinGrowthDischarge}~\dashsrc{QTSL}}{core}{\dnone}{\dnone}{\lnone{b}\lnone{r}\lnone{h}\lnone{d}}{\dnone}{\dnone}{\dok}
\dashrow{\hyperref[blcatrule-QMarginalCompat]{QMarginalCompat}~\dashsrc{QTSL}}{core}{\dnone}{\dnone}{\lnone{b}\lnone{r}\lnone{h}\lnone{d}}{\dnone}{\dnone}{\dok}
\dashrow{\hyperref[blcatrule-QAggReduce]{QAggReduce}~\dashsrc{QTSL}}{core}{\dcond}{\dok}{\lcur{b}\lcur{r}\lcur{h}\lcur{d}}{\dnone}{\dnone}{\dok}
\dashrow{\hyperref[blcatrule-QPAlwaysFromFold]{QPAlwaysFromFold}~\dashsrc{QTSL}}{core}{\dnone}{\dnone}{\lnone{b}\lnone{r}\lnone{h}\lnone{d}}{\dnone}{\dnone}{\dok}
\dashrow{\hyperref[blcatrule-QFireBoundFromFold]{QFireBoundFromFold}~\dashsrc{QTSL}}{core}{\dnone}{\dnone}{\lnone{b}\lnone{r}\lnone{h}\lnone{d}}{\dnone}{\dnone}{\dok}
\dashrow{\hyperref[blcatrule-QPAlwaysTFromFold]{QPAlwaysTFromFold}~\dashsrc{QTSL}}{core}{\dnone}{\dnone}{\lnone{b}\lnone{r}\lnone{h}\lnone{d}}{\dnone}{\dnone}{\dok}
\dashrow{\hyperref[blcatrule-QRateFromClock]{QRateFromClock}~\dashsrc{QTSL}}{core}{\dnone}{\dnone}{\lnone{b}\lnone{r}\lnone{h}\lnone{d}}{\dnone}{\dnone}{\dok}
\dashrow{\hyperref[blcatrule-QPAlwaysTFromLineage]{QPAlwaysTFromLineage}~\dashsrc{QTSL}}{core}{\dnone}{\dnone}{\lnone{b}\lnone{r}\lnone{h}\lnone{d}}{\dnone}{\dnone}{\dok}
\dashrow{\hyperref[blcatrule-QCoupleDrift]{QCoupleDrift}~\dashsrc{QTSL}}{core}{\dcond[finite]}{\dok}{\lcur{b}\lcur{r}\lcur{h}\lcur{d}}{\dnone}{\dnone}{\dok}
\dashrow{\hyperref[blcatrule-QFixBall]{QFixBall}~\dashsrc{QTSL}}{core}{\dok}{\dok}{\lcur{b}\lcur{r}\lcur{h}\lcur{d}}{\dnone}{\dnone}{\dok}
\dashrow{\hyperref[blcatrule-QDistTransport]{QDistTransport}~\dashsrc{QTSL}}{core}{\dcond}{\dok}{\lcur{b}\lcur{r}\lcur{h}\lcur{d}}{\dnone}{\dnone}{\dok}
\dashrow{\hyperref[blcatrule-QInvRelativize]{QInvRelativize}~\dashsrc{QTSL}}{core}{\dok}{\dnone}{\lcur{b}\lstale{r}\lstale{h}\lcur{d}}{\dok}{\ddoc}{\dok}
\dashrow{\hyperref[blcatrule-QFirstService]{QFirstService}~\dashsrc{QTSL}}{core}{\dok}{\dnone}{\lcur{b}\lstale{r}\lstale{h}\lcur{d}}{\dok}{\ddoc}{\dok}
\dashrow{\hyperref[blcatrule-QFirstServiceSplit]{QFirstServiceSplit}~\dashsrc{QTSL}}{core}{\dok}{\dnone}{\lcur{b}\lstale{r}\lstale{h}\lcur{d}}{\dok}{\ddoc}{\dok}
\dashrow{\hyperref[blcatrule-DtmcAbsorbBracket]{DtmcAbsorbBracket}~\dashsrc{QTSL}}{core}{\dnone}{\dnone}{\lnone{b}\lnone{r}\lnone{h}\lnone{d}}{\dnone}{\dnone}{\dok}
\dashrow{\hyperref[blcatrule-DtmcPosteriorBracket]{DtmcPosteriorBracket}~\dashsrc{QTSL}}{core}{\dnone}{\dnone}{\lnone{b}\lnone{r}\lnone{h}\lnone{d}}{\dnone}{\dnone}{\dok}
\dashrow{\hyperref[blcatrule-QGenFormContraction]{QGenFormContraction}~\dashsrc{QTSL}}{core}{\dnone}{\dnone}{\lnone{b}\lnone{r}\lnone{h}\lnone{d}}{\dnone}{\dnone}{\dok}
\dashrow{\hyperref[blcatrule-Norm]{Norm}~\dashsrc{BL}}{adm.}{\dok}{\dok}{\lcur{b}\lstale{r}\lstale{h}\lcur{d}}{\dok}{\ddoc}{\dok}
\dashrow{\hyperref[blcatrule-Farkas]{Farkas}~\dashsrc{BL}}{adm.}{\dok}{\damber}{\lcur{b}\lstale{r}\lstale{h}\lcur{d}}{\dok}{\ddoc}{\dok}
\dashrow{\hyperref[blcatrule-Subdivide]{Subdivide}~\dashsrc{BL}}{adm.}{\dok}{\dok}{\lcur{b}\lstale{r}\lstale{h}\lcur{d}}{\dok}{\ddoc}{\dok}
\dashrow{\hyperref[blcatrule-LiRPA]{LiRPA}~\dashsrc{BL}}{adm.}{\dok}{\dok}{\lcur{b}\lcur{r}\lcur{h}\lcur{d}}{\dok}{\ddoc}{\dok}
\dashrow{\hyperref[blcatrule-ReLUPos]{ReLU-Pos}~\dashsrc{BL}}{adm.}{\dok}{\dok}{\lcur{b}\lstale{r}\lstale{h}\lcur{d}}{\dok}{\ddoc}{\dok}
\dashrow{\hyperref[blcatrule-ReLUNeg]{ReLU-Neg}~\dashsrc{BL}}{adm.}{\dok}{\dok}{\lcur{b}\lstale{r}\lstale{h}\lcur{d}}{\dok}{\ddoc}{\dok}
\dashrow{\hyperref[blcatrule-ReLUNN]{ReLU-NN}~\dashsrc{BL}}{adm.}{\dok}{\damber}{\lcur{b}\lstale{r}\lstale{h}\lcur{d}}{\dok}{\ddoc}{\dok}
\dashrow{\hyperref[blcatrule-ReLUAbove]{ReLU-Above}~\dashsrc{BL}}{adm.}{\dok}{\dok}{\lcur{b}\lstale{r}\lstale{h}\lcur{d}}{\dok}{\ddoc}{\dok}
\dashrow{\hyperref[blcatrule-CellInInterior]{Cell-In-Interior}~\dashsrc{BL}}{adm.}{\dok}{\dok}{\lcur{b}\lcur{r}\lcur{h}\lcur{d}}{\dok}{\ddoc}{\dok}
\dashrow{\hyperref[blcatrule-CellOut]{Cell-Out}~\dashsrc{BL}}{adm.}{\dok}{\dok}{\lcur{b}\lstale{r}\lstale{h}\lcur{d}}{\dok}{\ddoc}{\dok}
\dashrow{\hyperref[blcatrule-PConstExact]{PConstExact}~\dashsrc{BL}}{adm.}{\dok}{\dok}{\lcur{b}\lcur{r}\lcur{h}\lcur{d}}{\dok}{\ddoc}{\dok}
\dashrow{\hyperref[blcatrule-PsdCertTri]{PsdCertTri}~\dashsrc{BL}}{adm.}{\dok}{\dok}{\lcur{b}\lcur{r}\lcur{h}\lcur{d}}{\dok}{\ddoc}{\dok}
\dashrow{\hyperref[blcatrule-PsdCertDiag]{PsdCertDiag}~\dashsrc{BL}}{adm.}{\dok}{\dok}{\lcur{b}\lcur{r}\lcur{h}\lcur{d}}{\dok}{\ddoc}{\dok}
\dashrow{\hyperref[blcatrule-PsdCertKernel]{PsdCertKernel}~\dashsrc{BL}}{adm.}{\dok}{\dok}{\lcur{b}\lcur{r}\lcur{h}\lcur{d}}{\dok}{\ddoc}{\dok}
\dashrow{\hyperref[blcatrule-ReLUSplit]{ReLU-Split}~\dashsrc{BL}}{adm.}{\dok}{\dnone}{\lcur{b}\lcur{r}\lcur{h}\lcur{d}}{\dnone}{\dnone}{\dok}
\dashrow{\hyperref[blcatrule-Promise]{Promise}~\dashsrc{BL}}{adm.}{\dok}{\dok}{\lcur{b}\lstale{r}\lstale{h}\lcur{d}}{\dok}{\ddoc}{\dok}
\dashrow{\hyperref[blcatrule-ProductLattice]{ProductLattice}~\dashsrc{BL}}{adm.}{\dok}{\dok}{\lcur{b}\lcur{r}\lcur{h}\lcur{d}}{\dnone}{\dnone}{\dok}
\dashrow{\hyperref[blcatrule-QUntilBattery]{QUntilBattery}~\dashsrc{QTSL}}{adm.}{\dok}{\damber}{\lcur{b}\lstale{r}\lstale{h}\lcur{d}}{\dnone}{\dnone}{\dok}
\dashrow{\hyperref[blcatrule-QStreettCmplRequest]{QStreettCmplRequest}~\dashsrc{QTSL}}{adm.}{\dcond}{\dok}{\lcur{b}\lcur{r}\lcur{h}\lcur{d}}{\dnone}{\dnone}{\dok}
\dashrow{\hyperref[blcatrule-HypMultiplier]{HypMultiplier}~\dashsrc{QTSL}}{adm.}{\dok}{\dok}{\lcur{b}\lcur{r}\lcur{h}\lcur{d}}{\dnone}{\dnone}{\dok}
\dashrow{\hyperref[blcatrule-QVerifiedCover]{QVerifiedCover}~\dashsrc{QTSL}}{adm.}{\dok}{\dok}{\lcur{b}\lcur{r}\lcur{h}\lcur{d}}{\dnone}{\dnone}{\dok}
\dashrow{\hyperref[blcatrule-QPositiveCombinationRoot]{QPositiveCombinationRoot}~\dashsrc{QTSL}}{adm.}{\dok}{\dok}{\lcur{b}\lcur{r}\lcur{h}\lcur{d}}{\dnone}{\dnone}{\dok}
\dashrow{\hyperref[blcatrule-QClosureOutward]{QClosureOutward}~\dashsrc{QTSL}}{adm.}{\dok}{\dok}{\lcur{b}\lcur{r}\lcur{h}\lcur{d}}{\dnone}{\dnone}{\dok}
\dashrow{\hyperref[blcatrule-QClosureInwardEps]{QClosureInwardEps}~\dashsrc{QTSL}}{adm.}{\dok}{\dok}{\lcur{b}\lcur{r}\lcur{h}\lcur{d}}{\dnone}{\dnone}{\dok}
\dashrow{\hyperref[blcatrule-BoxBern]{BoxBern}~\dashsrc{QTSL}}{adm.}{\dok}{\dok}{\lcur{b}\lstale{r}\lstale{h}\lcur{d}}{\dnone}{\dnone}{\dok}
\dashrow{\hyperref[blcatrule-PExpTypeEnclosure]{PExpTypeEnclosure}~\dashsrc{QTSL}}{adm.}{\dok}{\dok}{\lcur{b}\lcur{r}\lcur{h}\lcur{d}}{\dok}{\ddoc}{\dok}
\dashrow{\hyperref[blcatrule-QEquilibriumFromInward]{QEquilibriumFromInward}~\dashsrc{QTSL}}{adm.}{\dnone}{\dnone}{\lnone{b}\lnone{r}\lnone{h}\lnone{d}}{\dnone}{\dnone}{\dok}
\dashrow{\hyperref[blcatrule-QFlatten]{QFlatten}~\dashsrc{QTSL}}{adm.}{\dok}{\dok}{\lcur{b}\lcur{r}\lcur{h}\lcur{d}}{\dnone}{\dnone}{\dok}
\dashrow{\hyperref[blcatrule-PBound]{PBound}~\dashsrc{BL}}{alias}{\dok}{\dok}{\lcur{b}\lstale{r}\lstale{h}\lcur{d}}{\dok}{\ddoc}{\dok}
\dashrow{\hyperref[blcatrule-SmallGain]{SmallGain}~\dashsrc{BL}}{alias}{\dok}{\dok}{\lcur{b}\lstale{r}\lstale{h}\lcur{d}}{\dscope}{\dnone}{\dok}
\dashrow{\hyperref[blcatrule-SmallGainN]{SmallGainN}~\dashsrc{BL}}{alias}{\dok}{\dnone}{\lcur{b}\lstale{r}\lstale{h}\lcur{d}}{\dscope}{\dnone}{\dok}
\dashrow{\hyperref[blcatrule-QSatFromRegion]{QSatFromRegion}~\dashsrc{QTSL}}{alias}{\dok}{\dok}{\lcur{b}\lcur{r}\lcur{h}\lcur{d}}{\dnone}{\dnone}{\dok}
\dashrow{\hyperref[blcatrule-QStreettDegenerate]{QStreettDegenerate}~\dashsrc{QTSL}}{alias}{\dok}{\dok}{\lcur{b}\lcur{r}\lcur{h}\lcur{d}}{\dscope}{\dnone}{\dok}
\dashrow{\hyperref[blcatrule-QStreettTargetCellCollapse]{QStreettTargetCellCollapse}~\dashsrc{QTSL}}{alias}{\dok}{\dok}{\lcur{b}\lcur{r}\lcur{h}\lcur{d}}{\dscope}{\dnone}{\dok}
\dashrow{\hyperref[blcatrule-QNnfRewrite]{QNnfRewrite}~\dashsrc{QTSL}}{alias}{\dok}{\dok}{\lcur{b}\lcur{r}\lcur{h}\lcur{d}}{\dscope}{\dnone}{\dok}
\dashrow{\hyperref[blcatrule-QStreettRequestNnf]{QStreettRequestNnf}~\dashsrc{QTSL}}{alias}{\dok}{\dok}{\lcur{b}\lcur{r}\lcur{h}\lcur{d}}{\dscope}{\dnone}{\dok}
\dashrow{\hyperref[blcatrule-QAvoid]{QAvoid}~\dashsrc{QTSL}}{alias}{\dok}{\dok}{\lcur{b}\lcur{r}\lcur{h}\lcur{d}}{\dscope}{\dnone}{\dok}
\dashrow{\hyperref[blcatrule-QStreettGF]{QStreettGF}~\dashsrc{QTSL}}{alias}{\dok}{\dok}{\lcur{b}\lcur{r}\lcur{h}\lcur{d}}{\dscope}{\dnone}{\dok}
\dashrow{\hyperref[blcatrule-QStreettFG]{QStreettFG}~\dashsrc{QTSL}}{alias}{\dok}{\dok}{\lcur{b}\lcur{r}\lcur{h}\lcur{d}}{\dscope}{\dnone}{\dok}
\dashrow{\hyperref[blcatrule-QNegFAlias]{QNegFAlias}~\dashsrc{QTSL}}{alias}{\dok}{\dok}{\lcur{b}\lcur{r}\lcur{h}\lcur{d}}{\dscope}{\dnone}{\dok}
\dashrow{\hyperref[blcatrule-QFullDomainWeaken]{QFullDomainWeaken}~\dashsrc{QTSL}}{alias}{\dok}{\dok}{\lcur{b}\lcur{r}\lcur{h}\lcur{d}}{\dnone}{\dnone}{\dok}
\dashrow{\hyperref[blcatrule-QClosureExact]{QClosureExact}~\dashsrc{QTSL}}{alias}{\dok}{\dok}{\lcur{b}\lcur{r}\lcur{h}\lcur{d}}{\dnone}{\dnone}{\dok}
\dashrow{\hyperref[blcatrule-TaylorBound]{TaylorBound}~\dashsrc{QTSL}}{alias}{\dok}{\dok}{\lcur{b}\lcur{r}\lcur{h}\lcur{d}}{\dscope}{\dnone}{\dok}
\dashrow{\hyperref[blcatrule-RootEval]{RootEval}~\dashsrc{QTSL}}{alias}{\dok}{\dok}{\lcur{b}\lcur{r}\lcur{h}\lcur{d}}{\dscope}{\dnone}{\dok}
\dashrow{\hyperref[blcatrule-QLatticeInv]{QLatticeInv}~\dashsrc{QTSL}}{alias}{\dok}{\dok}{\lcur{b}\lcur{r}\lcur{h}\lcur{d}}{\dnone}{\dnone}{\dok}
\dashrow{\hyperref[blcatrule-QCountAtomRounding]{QCountAtomRounding}~\dashsrc{QTSL}}{alias}{\dok}{\dok}{\lcur{b}\lcur{r}\lcur{h}\lcur{d}}{\dnone}{\dnone}{\dok}
\dashrow{\hyperref[blcatrule-ProofWellFormedness]{ProofWellFormedness}~\dashsrc{QTSL}}{glue}{\dcond}{\dnone}{\lcur{b}\lcur{r}\lcur{h}\lcur{d}}{\dscope}{\dnone}{\dok}
\dashboardfoot

\blcatsection{Rule Synopsis}
\sourcenote{Every rule of the catalogue at a glance: the rule as
realized in the mechanized development, its name, and the page of
its full entry. The closing entry, set off by a rule, is the
assembly-level discipline the numbered rules feed.}
\par\medskip
\paradeentry{\inferrule*{\;}{\Phi \vdash \Phi}}{Ax}{blcatrule-Ax}
\paradeentry{\inferrule*{\Phi \vdash \Psi}{\Phi \wedge \chi \vdash \Psi}}{Weaken}{blcatrule-Weaken}
\paradeentry{\inferrule*{\Phi \vdash \chi \\ \Phi \wedge \chi \vdash \Psi}{\Phi \vdash \Psi}}{Cut}{blcatrule-Cut}
\paradeentry{\inferrule*{\Phi \vdash \Psi_1 \\ \Phi \vdash \Psi_2}{\Phi \vdash \Psi_1 \wedge \Psi_2}}{$\wedge$I}{blcatrule-AndI}
\paradeentry{\inferrule*{\Phi \vdash \Psi_1 \wedge \Psi_2}{\Phi \vdash \Psi_i}}{$\wedge$E}{blcatrule-AndE}
\paradeentry{\inferrule*{\Phi \vdash \Psi_i}{\Phi \vdash \Psi_1 \vee \Psi_2}}{$\vee$I}{blcatrule-OrI}
\paradeentry{\inferrule*{\Phi_1 \vdash \Psi \\ \Phi_2 \vdash \Psi}{\Phi_1 \vee \Phi_2 \vdash \Psi}}{$\vee$E}{blcatrule-OrE}
\paradeentry{\inferrule*{\Phi \vdash t = u \\ \Phi \vdash \Psi\sub{a}{t} \\ \exists s.\; \bshape_H(t) = s = \bshape_H(u)}{\Phi \vdash \Psi\sub{a}{u}}}{EqElim}{blcatrule-EqElim}
\paradeentry{\inferrule*{\exists s.\; \bshape_H(t) = s \\ \wires(t) \text{ finitely typed in } H}{\Phi \vdash t = t}}{EqRefl}{blcatrule-EqRefl}
\paradeentry{\inferrule*{\Phi \vdash t = u}{\Phi \vdash u = t}}{EqSym}{blcatrule-EqSym}
\paradeentry{\inferrule*{\Phi \vdash t \leq u \\ \Phi \vdash t \geq u}{\Phi \vdash t = u}}{Antisym}{blcatrule-Antisym}
\paradeentry{\inferrule*{
  \Phi \vdash x \in \mathcal{R}(\Delta, S)
}{
  \Phi \vdash \textstyle\bigwedge_{i = 0}^{n} \mathrm{canon}(\epsilon_i\,\lambda_i(x)) \leq 0 \\ \epsilon_i \,=\, {+1} \text{ if } i \in S, \quad {-1} \text{ if } i \in \bar S
}}{ConvElim}{blcatrule-ConvElim}
\paradeentry{\inferrule*{
  \Phi \vdash \textstyle\bigwedge_{i = 0}^{n} \mathrm{canon}(\epsilon_i\,\lambda_i(x)) \leq 0 \\ \epsilon_i \,=\, {+1} \text{ if } i \in S, \ {-1} \text{ if } i \in \bar S \\ \bshape_H(x) = [\,n\,]
}{
  \Phi \vdash x \in \mathcal{R}(\Delta, S)
}}{ConvIntro}{blcatrule-ConvIntro}
\paradeblock{\begin{mathpar}
\inferrule*[Right=\rulename{ArithAtom}]{
  \Phi \vdash (t_1 - t_2) \bowtie 0
}{
  \Phi \vdash t_1 \bowtie t_2
}
\\
\inferrule*[Right=\rulename{ArithAtom}$^{-1}$]{
  \Phi \vdash t_1 \bowtie t_2
}{
  \Phi \vdash (t_1 - t_2) \bowtie 0
}
\end{mathpar}
with $\bowtie \in \{<, \leq, =, \geq, >\}$.}{ArithAtom}{blcatrule-ArithAtom}
\paradeentry{\inferrule*{
  \Delta = \mathrm{conv}(V_0, \ldots, V_n) \subset \R^n \text{ an } n\text{-simplex};
  \quad S \subsetneq \{0, 1, \ldots, n\}
  \\
  \Phi \vdash x \in \mathcal{R}(\Delta, S)
  \\
  p \colon \R^n \to W \text{ polynomial map of degree } d, \text{ with coefficient tensor } (b_\alpha)_{|\alpha|=d}
  \\
  \wires(p) = \emptyset \;\vee\; \bigl(x = w \,\wedge\, \wires(p) \subseteq \{w\}\bigr)
  \\
  S \neq \emptyset \Rightarrow P \text{ conic}; \qquad S = \emptyset \Rightarrow P \text{ convex suffices}
  \\
  \bigl(p \ \text{idx-free, with } p \ \text{constant} \vee n = 1\bigr) \;\vee\; \bigl(p \ \text{projection-indexed, indices} < n,\ S = \emptyset\bigr)
  \\
  \text{for every } j \in \N^S \text{ with } |j| \leq d, \text{ every degree-}d \text{ Bernstein coefficient of } P_j \text{ on } F_S \text{ satisfies } P
}{
  \Phi \vdash P(p(x))
}}{BernCone}{blcatrule-BernCone}
\paradeentry{\inferrule*{
  t_1, t_2 \ \text{polynomial terms in the same variables with equal monomial canonical forms over } \Q
  \\
  \bshape_H(t_1) = \bshape_H(t_2) \neq \bot \ \text{(a shared defined broadcast shape)}
  \\
  \wires(t_1), \wires(t_2) \ \text{all finitely typed in } H
}{
  \Phi \vdash t_1 = t_2
}}{Norm}{blcatrule-Norm}
\paradeentry{\inferrule*{
  \{\ell_j \bowtie_j 0\}_j \text{ a sub-conjunction of } \Phi \text{ with each } \ell_j \text{ linear, } \bowtie_j \in \{\leq, =\} \\
  \{\lambda_j \in \Q_{\geq 0}\}_j \text{ nonneg coefficients, with } \lambda_j \text{ free when } \bowtie_j \text{ is } {=} \\
  \bshape_H(\textstyle\sum_j \lambda_j \ell_j) \neq \bot \\
  \wires(\textstyle\sum_j \lambda_j \ell_j) \text{ all finitely typed in } H \\
  \wires(\Phi) \cup \wires(\textstyle\sum_j \lambda_j \ell_j) \text{ all typed at positive-size shapes in } H
}{
  \Phi \vdash \textstyle\sum_j \lambda_j \ell_j \leq 0
}}{Farkas}{blcatrule-Farkas}
\paradeentry{\inferrule*{
  \Phi \vdash x \in \mathcal{R}(\Delta, S)
  \\
  \{(\Delta_k, S_k)\}_{k=1}^N \text{ finite, with cover side-condition } \mathcal{R}(\Delta, S) \subseteq \textstyle\bigcup_k \mathcal{R}(\Delta_k, S_k)
  \\
  \forall k.\; \Phi \;\wedge\; (x \in \mathcal{R}(\Delta_k, S_k)) \;\vdash\; \Psi
}{
  \Phi \vdash \Psi
}}{Subdivide}{blcatrule-Subdivide}
\paradeentry{\inferrule*{
  \Phi \vdash a \leq h \\
  \Phi \vdash h \leq b \\
  \ReLU_{\{h\}\to\{y\}} \in H \\
  a < 0 \\
  0 < b \\\\
  \alpha a + \beta \leq 0 \\
  \beta \leq 0 \\
  \alpha b + \beta \leq b \\\\
  0 \leq \gamma a + \delta \\
  0 \leq \delta \\
  b \leq \gamma b + \delta
}{
  \Phi \vdash (\alpha h + \beta) \leq y \;\wedge\; y \leq (\gamma h + \delta)
}}{LiRPA}{blcatrule-LiRPA}
\paradeentry{\inferrule*{
  \Phi \vdash h \geq 0
  \\
  \ReLU_{\{h\}\to\{y\}} \in H
}{
  \Phi \vdash y = h
}}{ReLU-Pos}{blcatrule-ReLUPos}
\paradeentry{\inferrule*{
  \Phi \vdash h \leq 0
  \\
  \ReLU_{\{h\}\to\{y\}} \in H
}{
  \Phi \vdash y = 0
}}{ReLU-Neg}{blcatrule-ReLUNeg}
\paradeentry{\inferrule*{
  \ReLU_{\{h\}\to\{y\}} \in H
}{
  \Phi \vdash y \geq 0
}}{ReLU-NN}{blcatrule-ReLUNN}
\paradeentry{\inferrule*{
  \ReLU_{\{h\}\to\{y\}} \in H
}{
  \Phi \vdash y \geq h
}}{ReLU-Above}{blcatrule-ReLUAbove}
\paradeentry{\inferrule*{
  \Phi \vdash \bigwedge_{v} \lambda^{C_i}_v(x) \bowtie^{\circ}_v 0 \ \text{(strict barycentric interior atoms, flat vertex order; } < \text{ for extended } v\text{, } > \text{ for core } v\text{)}
  \\
  \Cell[\Sigma]_{\{x\}\to\{b_1,\ldots,b_N\}} \in H
}{
  \Phi \vdash b_i = 1 \wedge \bigwedge_{j \neq i} b_j = 0
}}{Cell-In-Interior}{blcatrule-CellInInterior}
\paradeentry{\inferrule*{
  \Phi \vdash \lambda^{C_j}_w(x) \bowtie^{\circ} 0 \ \text{(a single strict separating atom at a named witness vertex } w \text{ of } C_j\text{; } > \text{ for extended } w\text{, } < \text{ for core } w\text{)}
  \\
  \Cell[\Sigma]_{\{x\}\to\{b_1,\ldots,b_N\}} \in H
}{
  \Phi \vdash b_j = 0
}}{Cell-Out}{blcatrule-CellOut}
\paradeentry{\inferrule*{
  e_1, e_2 \text{ wire-free (closed) terms}
  \\
  \bowtie \in \{<, \leq, =, \geq, >\}
  \\
  e_1 \bowtie e_2 \ \text{holds under exact evaluation in } \Q
}{
  \Phi \vdash e_1 \bowtie e_2
}}{PConstExact}{blcatrule-PConstExact}
\paradeentry{\inferrule*{
  \bigl\{\Phi_i \vdash e_i \leq 0\bigr\}_{i=1}^{n}, \quad \Phi \supseteq \Phi_i \ \text{conjunct-wise (checked)}
  \\
  \bigl\{g_k = 0\bigr\}_k \ \text{equality sub-conjuncts of } \Phi
  \\
  \lambda_i \in \Q_{\geq 0}
  \\
  \nu_k \in \Q
  \\
  c \equiv \textstyle\sum_i \lambda_i e_i + \sum_k \nu_k g_k \ \text{(exact canonical polynomial)}
  \\
  \bshape_H(c) \neq \bot, \quad \wires(c) \ \text{all finitely typed in } H
  \\
  \wires(\Phi) \cup \wires(c) \cup \textstyle\bigcup_i \wires(e_i) \ \text{all typed at positive-size shapes in } H
  \\
  \Psi \in \bigl\{\, c \leq 0, \;\; e \geq 0 \text{ where } c = -e, \;\; l \leq r \text{ where } c = l - r \,\bigr\}
}{
  \Phi \vdash \Psi
}}{PBound}{blcatrule-PBound}
\paradeentry{\inferrule*{
  A \in \Q^{n \times n} \ \text{symmetric}
  \\
  Z \ \text{upper-triangular}, \ Z_{ii} > 0
  \\
  A \cdot Z \ \text{lower-triangular}, \ (A \cdot Z)_{ii} > 0 \ \text{(recomputed exactly)}
  \\
  p_1, \dots, p_n \ \text{payload terms}
  \\
  c \equiv -\textstyle\sum_{i,j} A_{ij}\, p_i\, p_j \ \text{(exact canonical polynomial)}
  \\
  \Psi \in \bigl\{\, c \leq 0, \;\; Q \geq 0 \text{ where } c = -Q, \;\; l \leq r \text{ where } c = l - r \,\bigr\}
}{
  \Phi \vdash \Psi
}}{PsdCertTri}{blcatrule-PsdCertTri}
\paradeentry{\inferrule*{
  A \in \Q^{n \times n} \ \text{symmetric}
  \\
  Z \ \text{upper-triangular}, \ Z_{ii} > 0
  \\
  B := Z^{\mathsf{T}} A Z \ \text{strictly row-diagonally dominant (recomputed exactly)}
  \\
  p_1, \dots, p_n \ \text{payload terms}
  \\
  c \equiv -\textstyle\sum_{i,j} A_{ij}\, p_i\, p_j \ \text{(exact canonical polynomial)}
  \\
  \Psi \in \bigl\{\, c \leq 0, \;\; Q \geq 0 \text{ where } c = -Q, \;\; l \leq r \text{ where } c = l - r \,\bigr\}
}{
  \Phi \vdash \Psi
}}{PsdCertDiag}{blcatrule-PsdCertDiag}
\paradeentry{\inferrule*{
  A \in \Q^{(m+k) \times (m+k)}, \ \text{kernel coordinates last}
  \\
  L := (I \mid N), \ N \in \Q^{m \times k}
  \\
  A = L^{\mathsf{T}} B L \ \text{entrywise exactly}
  \\
  B \in \Q^{m \times m} \ \text{with a /T-or-/DD inner witness } Z
  \\
  p_1, \dots, p_{m+k} \ \text{payload terms}
  \\
  c \equiv -\textstyle\sum_{i,j} A_{ij}\, p_i\, p_j \ \text{(exact canonical polynomial)}
  \\
  \Psi \in \bigl\{\, c \leq 0, \;\; Q \geq 0 \text{ where } c = -Q, \;\; l \leq r \text{ where } c = l - r \,\bigr\}
}{
  \Phi \vdash \Psi
}}{PsdCertKernel}{blcatrule-PsdCertKernel}
\paradeentry{\inferrule*{
  \Phi \wedge (h \geq 0) \vdash \Psi\sub{b}{h} \\
  \Phi \wedge (h \leq 0) \vdash \Psi\sub{b}{0} \\
  \ReLU_{\{h\}\to\{y\}} \in H \quad (y \text{ typed at rank } 0)
}{
  \Phi \vdash \Psi
}}{ReLU-Split}{blcatrule-ReLUSplit}
\paradeentry{\inferrule*{
  \Phi_1 \vdash Q \\
  \Phi_2 \wedge Q \vdash \Psi \\
  \wires(Q) \subseteq Y \\
  \wires(\Phi_1) \cap \bigl(\wires(\Phi_2) \cup \wires(\Psi)\bigr) \subseteq Y
}{
  \Phi_1 \wedge \Phi_2 \vdash \Psi
}}{Promise}{blcatrule-Promise}
\paradeentry{\inferrule*{
  \bigl\{\Phi_i \vdash e_i \leq 0\bigr\}_{i=1}^{n} \\
  \bigl\{g_k = 0\bigr\}_k \ \text{equality sub-conjuncts of } \Phi \text{ in wire variables} \\
  \lambda_i \in \Q_{\geq 0} \\
  \nu_k \in \Q \\
  c \equiv \textstyle\sum_i \lambda_i e_i + \sum_k \nu_k g_k \ \text{(exact canonical polynomial)} \\
  \Psi \in \bigl\{\, c \leq 0, \;\; e \geq 0 \text{ where } c = -e, \;\; l \leq r \text{ where } c = l - r \,\bigr\}
}{
  \Phi \vdash \Psi
}}{Syzygy}{blcatrule-Syzygy}
\paradeentry{\inferrule*{
  \Phi_1 \vdash e_1 \leq 0 \\
  \Phi_2 \vdash e_2 \leq 0 \\
  \Phi_w \vdash \lambda\sigma_2 - (1{-}\theta)\,\alpha_1 \leq 0 \\
  \Phi_w \vdash \sigma_1 - \lambda(1{-}\theta)\,\alpha_2 \leq 0 \\
  \Phi \supseteq \Phi_1,\ \Phi_2,\ \Phi_w \ \text{conjunct-wise (checked)} \\
  \lambda \in \Q_{>0} \\
  \theta \in \Q \cap (0,1)
}{
  \Phi \vdash \mathcal{L}(V_1 + \lambda V_2) + \theta\,(\alpha_1 + \lambda\alpha_2) \leq 0
}}{SmallGain}{blcatrule-SmallGain}
\paradeentry{\inferrule*{
  \bigl\{\Phi_i \vdash e_i \leq 0\bigr\}_{i=1}^{n} \\
  \bigl\{\Phi_w \vdash m_i \leq 0\bigr\}_{i=1}^{n} \\
  \Phi \supseteq \Phi_1, \ldots, \Phi_n, \Phi_w \ \text{conjunct-wise (checked)} \\
  \lambda \in \Q_{>0}^{n} \\
  B \in \Q_{\geq 0}^{n \times n} \\
  \textstyle\sum_i B_{ij} = \lambda_j \ (\forall j) \\
  \theta \in \Q \cap (0,1)
}{
  \Phi \vdash \mathcal{L}\bigl(\textstyle\sum_i \lambda_i V_i\bigr) + \theta \textstyle\sum_j \lambda_j \alpha_j \leq 0
}}{SmallGainN}{blcatrule-SmallGainN}
\paradeentry{\inferrule*{
  B_1 \ \text{a certified attracting-block lattice for } \varphi_1 \ \text{over } H_1
  \\
  B_2 \ \text{likewise for } \varphi_2 \ \text{over } H_2
  \\
  \mathrm{wires}(H_1) \cap \mathrm{wires}(H_2) = \emptyset
}{
  \{\, b_1 \times b_2 \mid b_1 \in B_1,\ b_2 \in B_2 \,\} \ \text{(the meet-generators) certified for } \varphi_1 \times \varphi_2 \ \text{over } H_1 \sqcup H_2
}}{ProductLattice}{blcatrule-ProductLattice}
\paradeentry{\inferrule*{
  \mathrm{child} : \mathsf{QRegionSat}(\varphi, R) \\
  \mathrm{eval} : \vdash x_0 \in R \ \text{(BL leaf; sublevel} \to \rulename{PConstExact}\text{, box} \to \text{registry interval check)} \\
  \mathsf{SInit}(\varphi) \in s.\forall
}{
  \mathcal{M} \vDash s
}}{QSatFromRegion}{blcatrule-QSatFromRegion}
\paradeentry{\inferrule*{
  \mathrm{child} : \mathsf{QRegionSat}(\varphi, R) \\
  \mathsf{SInit}(\varphi) \in s.\forall
}{
  \mathcal{M} \vDash \mathsf{Guard}(R,\ s)
}}{QGuardIntro}{blcatrule-QGuardIntro}
\paradeentry{\inferrule*{
  \mathcal{M} \vDash \mathsf{Guard}(R,\ s) \\
  \mathrm{eval} : \vdash X_0 \subseteq R
}{
  \mathcal{M} \vDash s
}}{QGuardElim}{blcatrule-QGuardElim}
\paradeentry{\inferrule*{
  V : \mathsf{QCertFam} \\
  \bigl\{D_m \vdash \mathcal{A}_m V_m \le 0\bigr\}_{m} \ \text{(interiors)} \\
  \bigl\{\mathrm{slab}_j \vdash g_j \fatsemi V_{m'_j} \le V_{m_j}\bigr\}_{j} \ \text{(guards, per jump $j$: source mode $m_j$, target $m'_j$, reset $g_j$)} \\
  \bigl\{U \cap \bar{D}_m \vdash V_m \ge 1\bigr\}_{m} \ \text{(floors)} \\
  \bigl\{\bar{D}_m \vdash V_m \ge 0\bigr\}_{m} \ \text{(nonnegs)} \\
  \mathrm{scope\_warrant} : \mathsf{QInvScope} \ \text{(optional — the dissolved ScopeRule)} \\
  p + y \le 1 \ \text{(threshold tie)}
}{
  \mathcal{M},\ x_0 \in \{V \le y\} \vDash \mathsf{SProbGE}\bigl(p,\ \mathsf{PAlways}\ \varphi\bigr)
}}{QSafetyBattery}{blcatrule-QSafetyBattery}
\paradeentry{\inferrule*{
  V : \mathsf{QCertFam} \\
  R \subseteq [\![\varphi]\!] \ \text{(target formula tie)} \\
  \bigl\{D_m \smallsetminus R \vdash \mathcal{A}_m V_m \le -\theta V_m\bigr\}_{m} \ \text{(decay)} \\
  \bigl\{\mathrm{slab}_j \vdash g_j \fatsemi V_{m'_j} \le V_{m_j}\bigr\}_{j} \ \text{(guards, per jump $j$: source mode $m_j$, target $m'_j$, reset $g_j$)} \\
  \bigl\{\bar{D}_m \smallsetminus R \vdash V_m \ge 1\bigr\}_{m} \ \text{(floors, off-target only)} \\
  \bigl\{\bar{D}_m \vdash V_m \ge 0\bigr\}_{m},\ \bigl\{\bar{D}_m \vdash V_m \le M\bigr\}_{m} \ \text{(nonnegs; bound)} \\
  \vdash q \le \textstyle\sum_{k \le n} (\theta T)^k / k! \ \text{(one \rulename{TaylorBound} leaf)} \\
  0 \le \beta,\ \beta \le e^{-\theta},\ \theta \ge 0,\ 0 \le q,\ p \le 1 \ \text{(drift tie; sign gates)}
}{
  \mathcal{M},\ x_0 \in \{V \le (1-p)\,q\} \vDash \mathsf{SProbGE}\bigl(p,\ \mathsf{PDeadline}(T, \varphi, b)\bigr)
}}{QDeadlineBattery}{blcatrule-QDeadlineBattery}
\paradeentry{\inferrule*{
  \mu(E_1) \le a \\ \mu(E_2) \le b \\ \mathrm{fail} \subseteq E_1 \cup E_2
}{
  \mu(\mathrm{fail}) \le a + b
}}{QUnionBound}{blcatrule-QUnionBound}
\paradeentry{\inferrule*{
  E_{\mathrm{bad}} := \{\text{the first exit from } D \smallsetminus (A \cup R) \ \text{lands in } A \cup D^{\mathrm{c}}\}, \quad E_{\mathrm{miss}} := \{A \cup R \ \text{not reached by } T\} \\
  [\![\neg\varphi_s]\!] \subseteq A, \qquad R \subseteq [\![\varphi]\!] \ \text{(formula ties)} \\
  \mu_x(E_{\mathrm{bad}}) \le V_a(x) \ \text{(the \rulename{QSafetyBattery} slice for } V_a \ \text{at stop set } (A \cup R) \cup D^{\mathrm{c}}\text{: floor on } A \ \text{only)} \\
  \mu_x(E_{\mathrm{miss}}) \le \beta^{T}\, V_d(x) \ \text{(the \rulename{QDeadlineBattery} slice for } V_d \ \text{at target } A \cup R\text{)} \\
  [\![\neg(\varphi_s\,\mathsf{U}_{\le T}\,\varphi)]\!] \subseteq E_{\mathrm{bad}} \cup E_{\mathrm{miss}} \ \text{(the d\'ebut decomposition; consumed by \rulename{QUnionBound})} \\
  0 \le \beta, \qquad 0 \le V_d(x), \qquad \beta^{T} \le c \ \text{(sign gates; rate tie)}
}{
  \mathcal{M},\ x_0 \in \{V_a + c\,V_d \le 1 - p\} \vDash \mathsf{SProbGE}\bigl(p,\ \mathsf{PUntil}(T, \varphi_s, \varphi)\bigr)
}}{QUntilBattery}{blcatrule-QUntilBattery}
\paradeentry{\inferrule*{
  V_0,\ (W_i)_{i \le k},\ (V_i)_{i \in A} : \mathsf{QCertFam}; \quad A \subseteq \{1, \ldots, k\} \ \text{(the pairs posting absorbing certificates)}; \quad \varepsilon_0 \in \Q_{>0}; \quad (\gamma_i)_{i \in A} \in \Q_{>0} \\
  \text{the \rulename{QSafetyBattery} rows for } V_0 \ \text{at } U := \mathrm{cmpl}(I) \ \text{(interior, guard, floor on } U \cap \bar{D}\text{, nonneg on } \bar{D}\text{)} \\
  \bigl\{S \smallsetminus [\![\rho_i]\!] \vdash \mathcal{A}_m V_{i,m} \le 0\bigr\}_{i \in A} \ \text{(absorbing drift: off-request, on all of } S\text{)} \\
  \bigl\{S \vdash 0 \le V_i \le C^{a}_i\bigr\}_{i \in A} \ \text{(absorbing bounds, on all of } S\text{)}, \quad \bigl\{[\![\rho_i]\!] \vdash V_i \ge 1\bigr\}_{i \in A} \ \text{(absorbing floors)} \\
  \bigl\{([\![\varphi_i]\!] \cup J_i)^{\complement} \vdash \mathcal{A}_m W_{i,m} \le -\varepsilon_0\bigr\}_{i \le k} \ \text{(ranking decay)}, \quad J_i := \{V_i \le 1 - \gamma_i\} \cap (I \smallsetminus [\![\rho_i]\!]) \ \text{for } i \in A; \quad J_i := \emptyset \ \text{for } i \notin A \\
  W_i \equiv 0 \ \text{on target cells (containment)}; \quad \bigl\{\bar{D} \vdash W_i \ge 0\bigr\}_{i \le k} \ \text{(ranking nonnegs)} \\
  \bigl\{\bar{D} \vdash W_i \le C_i\bigr\}_{i \le k}, \quad C_i \in \Q \\
  p + y \le 1 \ \text{(threshold tie)}
}{
  \mathcal{M},\ x_0 \in \{V_0 \le y\} \vDash \mathsf{SProbGE}\bigl(p,\ \mathsf{PDwellStreett}(\mathsf{PairBase}(\rho_1, \varphi_1), \ldots, \mathsf{PairBase}(\rho_k, \varphi_k))\bigr)
}}{QStreettBattery}{blcatrule-QStreettBattery}
\paradeentry{\inferrule*{
  \overline{\mathrm{drift}_m : R_m \vdash g - r_m - \mathcal{A}_m B_m \ \ge\ 0} \\
  \overline{\mathrm{guard}_j} \quad \overline{\mathrm{floor}_m} \quad \overline{\mathrm{nonneg}_m} \\
  \mathrm{tie} : \vdash g \bowtie c
}{
  \mathcal{M},\ x_0 \in R \vDash \mathsf{SReturn}{\bowtie}\bigl(c,\ \mathsf{poAvg}(o)\bigr)
}}{QAvgBattery}{blcatrule-QAvgBattery}
\paradeentry{\inferrule*{
  \mathsf{TrMeter}(t, d) \\
  \rho \cdot d \ge 1 \\
  \sigma \ge 1 \\
  \mathsf{qtsl\_rate\_clock}(n)
}{
  \mathcal{M} \vDash \mathsf{SRateLE}(\{t\},\ \rho,\ \sigma)
}}{QRateFromClock}{blcatrule-QRateFromClock}
\paradeentry{\inferrule*{
  \mathsf{BoxInv}(B, \nu, x)
}{
  \mathsf{AService}(B, T, \chi \cap \nu, x) \;\leftrightarrow\; \mathsf{AService}(B, T, \chi, x)
}}{QInvRelativize}{blcatrule-QInvRelativize}
\paradeentry{\inferrule*{
  B.\mathsf{FusionClosed} \\ 0 < n \\
  \forall i.\ \mathsf{FirstVisits}(B, \alpha_i) \\
  \forall i.\ \mathsf{AService}(B, T_i, \alpha_i, x)
}{
  \mathsf{AService}\big(B,\ {\textstyle\min_i T_i},\ \mathsf{fsTarget}(\alpha, T),\ x\big)
}}{QFirstService}{blcatrule-QFirstService}
\paradeentry{\inferrule*{
  B.\mathsf{FusionClosed} \\ 0 < n \\ \forall i.\ c < T_i \\
  \forall i.\ \mathsf{FirstVisits}(B, \alpha_i) \\
  \forall i.\ \mathsf{AService}(B, T_i, \alpha_i, x)
}{
  \mathsf{AService}\big(B,\ c,\ \mathsf{fsTarget}(\alpha, T) \cup \mathsf{residualCell}(\alpha, T, c),\ x\big)
}}{QFirstServiceSplit}{blcatrule-QFirstServiceSplit}
\paradeentry{\inferrule*{
  \text{one \rulename{QStreettBattery} node at } \rho_i := \mathsf{STrue} \ \text{for every pair} \\
  \text{(absorbing families and } J_i \ \text{syntactically empty; ranking decay off } [\![\varphi_i]\!] \ \text{alone)}
}{
  \mathcal{M},\ x_0 \in \{V_0 \le y\} \vDash \mathsf{SProbGE}\bigl(p,\ \mathsf{PDwellStreett}(\mathsf{PairBase}(\mathsf{STrue}, \varphi_1), \ldots, \mathsf{PairBase}(\mathsf{STrue}, \varphi_k))\bigr)
}}{QStreettDegenerate}{blcatrule-QStreettDegenerate}
\paradeentry{\inferrule*{
  \text{a \rulename{QStreettBattery} kit at certified regions } (R_i, G_i) \\
  \bigl\{ [\![\mathrm{cmpl}\,\rho_i]\!] \subseteq R_i \bigr\} \ \text{(requests over-approximated)} \\
  \bigl\{ G_i \subseteq [\![\varphi_i]\!] \bigr\} \ \text{(responses under-approximated)}
}{
  \mathcal{M},\ x_0 \in \{V_0 \le y\} \vDash \mathsf{SProbGE}\bigl(p,\ \mathsf{PDwellStreett}(\mathsf{PairCmpl}(\rho_1, \varphi_1), \ldots)\bigr)
}}{QStreettCmplRequest}{blcatrule-QStreettCmplRequest}
\paradeentry{\inferrule*{
  \mathcal{K},\ \mathrm{slab} \vdash W_i \ge 0
}{
  \mathcal{K},\ \mathrm{slab} \vdash g \fatsemi W_i \le W_i
}}{QStreettTargetCellCollapse}{blcatrule-QStreettTargetCellCollapse}
\paradeblock{\begin{mathpar}
\inferrule*{ }{ [\![\neg\neg\varphi]\!] = [\![\varphi]\!] }
\and
\inferrule*{ }{ [\![\neg(\varphi \wedge \psi)]\!] = [\![\neg\varphi \vee \neg\psi]\!] }
\and
\inferrule*{ }{ [\![\neg(\varphi \vee \psi)]\!] = [\![\neg\varphi \wedge \neg\psi]\!] }
\and
\inferrule*{ }{ [\![\neg\{f \le c\}]\!] = \{f > c\} }
\and
\inferrule*{ }{ [\![\neg\{f < c\}]\!] = \{f \ge c\} }
\and
\inferrule*{ }{ [\![\neg\{f \ge c\}]\!] = \{f < c\} }
\and
\inferrule*{ }{ [\![\neg\{f > c\}]\!] = \{f \le c\} }
\end{mathpar}}{QNnfRewrite}{blcatrule-QNnfRewrite}
\paradeblock{\begin{mathpar}
\inferrule*{
  \mathcal{M},\ x_0 \in R \vDash \mathsf{SProbGE}\bigl(p,\ \mathsf{PDwellStreett}(\ldots,\ \mathsf{PairCmpl}(\mathrm{elimNeg}(\mathrm{nnf}(\rho_i)),\ \varphi_i),\ \ldots)\bigr)
}{
  \mathcal{M},\ x_0 \in R \vDash \mathsf{SProbGE}\bigl(p,\ \mathsf{PDwellStreett}(\ldots,\ \mathsf{PairCmpl}(\rho_i,\ \varphi_i),\ \ldots)\bigr)
}
\and
\inferrule*{ }{ [\![\mathrm{elimNeg}(\mathrm{nnf}(\rho))]\!] = [\![\rho]\!] }
\end{mathpar}}{QStreettRequestNnf}{blcatrule-QStreettRequestNnf}
\paradeentry{\inferrule*{
  \mathcal{M},\ x_0 \in R \vDash \mathsf{SProbGE}\bigl(p,\ \mathsf{PAlways}(\mathsf{SNot}\,\beta)\bigr)
}{
  \mathcal{M},\ x_0 \in R \vDash \mathsf{P}_{\ge p}[\mathsf{avoid}\,\beta]
}}{QAvoid}{blcatrule-QAvoid}
\paradeentry{\inferrule*{
  \mathcal{M},\ x_0 \in R \vDash \mathsf{SProbGE}\bigl(p,\ \mathsf{PDwellStreett}(\mathsf{PairBase}(\mathsf{STrue}, \varphi))\bigr) \\
  \text{paths total over the state space}
}{
  \mathcal{M},\ x_0 \in R \vDash \mathsf{SProbGE}(p,\ \mathsf{GF}\,\varphi)
}}{QStreettGF}{blcatrule-QStreettGF}
\paradeentry{\inferrule*{
  \mathcal{M},\ x_0 \in R \vDash \mathsf{SProbGE}\bigl(p,\ \mathsf{PDwellStreett}(\mathsf{PairBase}(\beta, \mathsf{SFalse}))\bigr)
}{
  \mathcal{M},\ x_0 \in R \vDash \mathsf{SProbGE}(p,\ \mathsf{FG}\,\neg\beta)
}}{QStreettFG}{blcatrule-QStreettFG}
\paradeentry{\inferrule*{
  \mathrm{child} : \mathsf{QRegionSat}\bigl(\mathsf{SProbGE}(p, \mathsf{PAlways}(\mathsf{SNot}\,\varphi)),\ R\bigr)
}{
  \mathcal{M},\ x_0 \in R \vDash \mathsf{SProbGE}(p, \mathsf{PNotPath}(\mathsf{PDeadline}(T, \varphi, \mathrm{false})))
}}{QNegFAlias}{blcatrule-QNegFAlias}
\paradeentry{\inferrule*{
  \vdash \lambda \ge 0 \ \text{on } \mathrm{cell} \\
  \vdash N + \lambda\, g_q \le 0 \ \text{on } \mathrm{cell}
}{
  \vdash N \le 0 \ \text{on } \mathrm{cell} \cap \{g_q \ge 0\}
}}{HypMultiplier}{blcatrule-HypMultiplier}
\paradeentry{\inferrule*{
  \vdash P \ \text{on } D_m \\
  \mathrm{scope} \subseteq D_m \ \text{(compile fact, not a posted premise)}
}{
  \vdash P \ \text{on scope}
}}{QFullDomainWeaken}{blcatrule-QFullDomainWeaken}
\paradeentry{\inferrule*{
  \vdash P \ \text{on } C \ \text{(the cover row)} \\
  \bigl\{\vdash K_i \subseteq \mathrm{Off}\bigr\}_i \ \text{(per-component containments)} \\
  D_m \smallsetminus C \subseteq \textstyle\bigcup_i K_i \ \text{(residual decomposition)}
}{
  \vdash P \ \text{on } D_m \smallsetminus \mathrm{Off}
}}{QVerifiedCover}{blcatrule-QVerifiedCover}
\paradeblock{\begin{mathpar}
\inferrule*{
  \mathrm{child}_1 : \mathsf{QRegionSat}(\varphi_1, R_1) \\
  \mathrm{child}_2 : \mathsf{QRegionSat}(\varphi_2, R_2)
}{
  \mathcal{M},\ x_0 \in R_1 \cap R_2 \vDash \varphi_1 \wedge \varphi_2
}
\and
\inferrule*{
  \mathrm{child}_1 : \mathsf{QRegionSat}(\varphi_1, R_1) \\
  \mathrm{child}_2 : \mathsf{QRegionSat}(\varphi_2, R_2)
}{
  \mathcal{M},\ x_0 \in R_1 \cup R_2 \vDash \varphi_1 \vee \varphi_2
}
\end{mathpar}}{QPositiveCombinationRoot}{blcatrule-QPositiveCombinationRoot}
\paradeblock{\begin{mathpar}
\inferrule*{ }{ \mathrm{cmpl}\{p < c\} = \{c \le p\} }
\and
\inferrule*{ }{ \mathrm{cmpl}\{c < p\} = \{p \le c\} }
\end{mathpar}}{QClosureExact}{blcatrule-QClosureExact}
\paradeblock{\begin{mathpar}
\inferrule*{ }{ \mathrm{cmpl}\{p \le c\} \subseteq \{c \le p\} }
\and
\inferrule*{ }{ \mathrm{cmpl}\{c \le p\} \subseteq \{p \le c\} }
\and
\inferrule*{
  R \subseteq [\![\varphi]\!]
}{
  \mathrm{cmpl}\,[\![\varphi]\!] \subseteq \mathrm{cmpl}\,R
}
\end{mathpar}}{QClosureOutward}{blcatrule-QClosureOutward}
\paradeblock{\begin{mathpar}
\inferrule*{
  \varepsilon \in \Q_{>0}
}{
  \{p \le c - \varepsilon\} \subseteq \{p < c\}
}
\and
\inferrule*{
  \varepsilon \in \Q_{>0}
}{
  \{c + \varepsilon \le p\} \subseteq \{c < p\}
}
\end{mathpar}}{QClosureInwardEps}{blcatrule-QClosureInwardEps}
\paradeentry{\inferrule*{
  q \in \Q; \quad \theta \in \Q_{>0}; \quad T \in \N; \quad n \in \N \\
  q \le \textstyle\sum_{k \le n} (\theta T)^k / k! \ \text{holds under exact evaluation in } \Q
}{
  \Phi \vdash q \le \textstyle\sum_{k \le n} (\theta T)^k / k!
}}{TaylorBound}{blcatrule-TaylorBound}
\paradeblock{\begin{mathpar}
\inferrule*{
  R = \{x : V(x) \le y\} \ \text{(sublevel root atom)} \\
  \vdash V(x_0) \le y \ \text{(one \rulename{PConstExact} leaf: the substituted comparison is closed)}
}{
  \vdash x_0 \in R
}
\and
\inferrule*{
  R = \textstyle\prod_i [a_i, b_i] \ \text{(box root atom)} \\
  \forall i.\ a_i \le (x_0)_i \le b_i \ \text{(recomputed in exact } \Q\text{)}
}{
  \vdash x_0 \in R
}
\end{mathpar}}{RootEval}{blcatrule-RootEval}
\paradeentry{\inferrule*{
  N \in \Q[z_1, \ldots, z_n]; \quad D = \textstyle\prod_{v} [a_v, b_v] \subseteq \Q^n
  \\
  \varphi = (A,\ A^{-1},\ b): \quad A,\ A^{-1} : \Q^{[n] \times [n]},\ \ b : \Q^{[n]}; \qquad \varphi(w) = A\,w + b \ \text{(affine-rung witness; } \varphi \ \text{omitted reads as } \mathrm{id}\text{)}
  \\
  \vdash A^{-1} A = 1_{[n]} \ \text{(the one certified product, exact)}
  \\
  \hat N = \varphi \fatsemi N \ \text{(exact substitution)}
  \\
  \bigl\{b_\alpha \le 0\bigr\}_{\alpha} \ \text{(every tensor-product Bernstein coefficient of } \hat N \ \text{over } D\text{)}
}{
  D \vdash \hat N \le 0
}}{BoxBern}{blcatrule-BoxBern}
\paradeentry{\inferrule*{
  f \in F \quad n \in \mathbb{N} \quad r, R \in \Q
  \\
  a_k = a^f_k \ \text{for } k = 0..n \ \text{(exact coefficient recurrence)}
  \\
  0 \le r < n + 2
  \qquad
  R \;\ge\; \frac{r^{n+1}}{(n+1)!}\cdot\frac{n+2}{\,n+2-r\,}
}{
  \forall x,\ |x| \le r \;\Rightarrow\; \Bigl|f(x) - \textstyle\sum_{k \le n} a_k x^k\Bigr| \le R
}}{PExpTypeEnclosure}{blcatrule-PExpTypeEnclosure}
\paradeentry{\inferrule*{
  L : \mathsf{QBlockLattice} \ \text{(blocks over one rational grid)} \\
  \bigl\{F_f \vdash \sigma_f\, \mathcal{A}\ell_f > 0\bigr\}_{f \in \mathrm{facets}(L)} \ \text{(facet-sign rows: } \ell_f \ \text{the facet\textquotesingle s normal coordinate, } \sigma_f \in \{\pm 1\} \ \text{the declared crossing sign, } F_f \ \text{the facet)} \\
  h : \mathsf{QCellMap} \ (\text{cell} \to \text{vertex} \sqcup \{\text{transit}\}) \\
  \bigl\{\vdash T(b) \subseteq \rho(h(b))\bigr\}_{h(b) \ne \text{transit}} \ \text{(containments, per vertex cell)} \\
  \bigl\{T(b) \vdash \mathcal{A} V_b + \varepsilon_b \le 0\bigr\}_{h(b) = \text{transit}}, \quad \varepsilon_b \in \Q_{>0} \ \text{(transits, per transit cell)} \\
  \bigl\{\vdash \rho(u) \cap \rho(v) = \emptyset\bigr\}_{u \ne v \in G}, \quad \bigl\{\vdash \rho(v) \subseteq D\bigr\}_{v \in G} \ \text{(} \rho\text{-wf: rows where not box-aligned, recomputed where box-aligned)} \\
  \mathrm{CR}(\mathcal{M}|_D) \subseteq \textstyle\bigcup_{h(b) \ne \text{transit}} T(b) \ \text{(cell-level chain-recurrence cover)} \\
  \mathsf{MorseAssembly}(L, h, G, \rho) \ \text{(the assembly ledger, defined below)}
}{
  \mathcal{M} \vDash \mathsf{Morse}(D, G, \rho)
}}{QMorseDischarge}{blcatrule-QMorseDischarge}
\paradeentry{\inferrule*{
  \mathcal{M}_1 \vDash \mathsf{Morse}(D_1, G_1, \rho_1) \\
  \mathcal{M}_2 \vDash \mathsf{Morse}(D_2, G_2, \rho_2) \\
  L_1, L_2 : \mathsf{QBlockLattice} \ \text{(the factor lattices)}
}{
  \mathcal{M}_1 \times \mathcal{M}_2 \vDash \mathsf{Morse}(D_1 \times D_2,\ G_1 \times G_2,\ \rho_1 \times \rho_2)
}}{QMorseProduct}{blcatrule-QMorseProduct}
\paradeentry{\inferrule*{
  N : \mathsf{QRegionAtom} \ \text{(the block)} \\
  \bigl\{F_f \vdash \mathcal{A}\ell_f \ge \varepsilon_f\bigr\}_{f \in \mathrm{facets}(N)}, \quad \varepsilon_f \in \Q_{>0} \ \text{(} \ell_f \ \text{the facet\textquotesingle s inward affine coordinate; } F_f \ \text{its facet region)} \\
  \mathsf{FacetBridge}(\mathcal{M}, D, N, (F_f, \varepsilon_f)_f) \ \text{(the bridge package, defined below)} \\
  \vdash X_0 \subseteq N \ \text{(\texttt{init\_containment}, a single BL leaf)}
}{
  \mathcal{M} \vDash \mathsf{QInvScope}(N, X_0)
}}{QBlockInv}{blcatrule-QBlockInv}
\paradeentry{\inferrule*{
  L : \mathsf{QBlockLattice} \\
  C \ \text{(\texttt{downset}, a cell-set; membership and down-closedness recomputed)} \\
  \mathcal{M} \vDash \mathsf{QInvScope}\bigl(\textstyle\bigcup_{i \in C} \mathrm{blocks}(\mathrm{emb}\ i),\ X_0\bigr) \ \text{(the cited \rulename{QBlockInv} conclusion)} \\
  \vdash X_0 \subseteq \textstyle\bigcup\downarrow C \ \text{(\texttt{init\_containment})}
}{
  \mathcal{M} \vDash \mathsf{QInvScope}\bigl(\textstyle\bigcup\downarrow C,\ X_0\bigr)
}}{QLatticeInv}{blcatrule-QLatticeInv}
\paradeentry{\inferrule*{
  \mathcal{M} \vDash \mathsf{QInvScope}(S, X_0) \\
  \bigl\{\vdash S \cap [\![\beta_j]\!] = \emptyset\bigr\}_{j} \ \text{(\texttt{avoid\_disjoint} rows)}
}{
  \mathcal{M},\ x_0 \in X_0 \vDash \mathsf{SProbGE}\bigl(1,\ \mathsf{PAlways}(\mathsf{SNot}\,\beta_j)\bigr)
}}{QBasin}{blcatrule-QBasin}
\paradeentry{\inferrule*{
  b : \mathsf{QMirandaBattery} \\
  \mathsf{qtsl\_miranda\_battery}(b,\ C,\ g,\ \Pi)
}{
  \vDash \mathsf{QZeroExists}(g,\ C)
}}{QZeroByMiranda}{blcatrule-QZeroByMiranda}
\paradeentry{\inferrule*{
  \mathrm{inv} : \mathsf{QBlockInv}\ \text{(strictly-inward facets, box block)}
}{
  \vDash \mathsf{QEquilibriumExists}(f,\ C)
}}{QEquilibriumFromInward}{blcatrule-QEquilibriumFromInward}
\paradeentry{\inferrule*{
  b : \mathsf{QMirandaBattery} \\
  \mathsf{qtsl\_miranda\_battery}\bigl(b,\ W,\ \sigma - \mathrm{id}_{Z},\ Z\bigr)
}{
  \vDash Z \subseteq \sigma(W)
}}{QCoverByMiranda}{blcatrule-QCoverByMiranda}
\paradeentry{\inferrule*{
  \bigl\{\Sigma_i \vDash (I_i, \alpha_i, \gamma_i[, \kappa_i])\bigr\}_{i=1}^{n} \\
  \mathsf{QWireCert}\bigl(w,\ (\mathit{spec}_i)_{i} \Rightarrow \mathit{spec}_{\mathrm{out}},\ \kappa\bigr) \\
  \mathrm{rank} : \{1, \ldots, n\} \to \N; \quad \bigl\{\mathrm{rank}_i < \mathrm{rank}_j\bigr\}_{(i \to j) \in w} \ \text{(the causality certificate: one rank per component, strictly increasing along every wire)} \\
  \mathrm{W0}\!: \ \text{the composite martingale problem admits a c\`adl\`ag strong-Markov solution, unique in law, from every initial state (stochastic mechanisms; the martingale-problem sense --- distinct from the aggregation language\textquotesingle s own W0)} \\
  \mathrm{W1}\!: \ \text{transition data factor through per-place colour multisets (SDCPN mechanisms; typing-discharged in the aggregation language)} \\
  \mathrm{W2}\!: \ \text{the within-instant immediate-jump relation is well-founded, of uniformly bounded depth} \\
  \mathrm{W3}\!: \ \text{simultaneous enablings carry a declared serialization or a certified reset commutation (a.s.-null collision row for stochastic timers)} \\
  \kappa_{\mathrm{out}} := \textstyle\bigoplus_{\mathrm{legs}} \kappa_{\mathrm{leg}}, \quad \kappa_a \oplus \kappa_b := \kappa_a + \kappa_b \ \text{(slack addition along lens composition; checker-recomputed, never stored)}
}{
  \mathsf{SysWire}(w, [\Sigma_1, \ldots, \Sigma_n]) \vDash \mathit{spec}_{\mathrm{out}}
}}{QWireCompose}{blcatrule-QWireCompose}
\paradeentry{\inferrule*{
  \bigl\{R_\ell,\ \varphi(s) \wedge \alpha(v(s), a) \vdash \mathrm{Lift}_s\, \varphi(u(s, a))\bigr\}_{\ell},\ \ \bigl\{R_\ell,\ \varphi(s) \vdash \gamma(v(s))\bigr\}_{\ell} \ \text{(mechanism-certification legs; } s \ \text{a mechanism state, } a \ \text{an admissible input)} \\
  \bigl\{R_\ell \vdash \alpha_2(w(o_1), a_2) + \kappa(\gamma_1(o_1)) \ge \alpha_1(w^{\sharp}(o_1, a_2))\bigr\}_{\ell} \ \text{(assumption legs)} \\
  \bigl\{R_\ell \vdash \gamma_1(o_1) \ge \gamma_2(w(o_1))\bigr\}_{\ell} \ \text{(guarantee legs)} \\
  \vdash \kappa(0) = 0 \ \text{(closed evaluation)} \\
  [0, \infty) \vdash \kappa' \ge 0 \ \text{(derivative-nonnegativity row)} \\
  \vdash \deg \kappa > 0,\ \mathrm{lead}(\kappa) \ge 0 \ \text{(leading-coefficient row)}
}{
  \mathsf{QWireCert}\bigl(w,\ (\mathit{spec}_i) \Rightarrow \mathit{spec}_{\mathrm{out}},\ \kappa\bigr)
}}{QWireCertDischarge}{blcatrule-QWireCertDischarge}
\paradeentry{\inferrule*{
  \bigl\{\Sigma_i \vDash (I_i,\ A_i,\ E_i)\bigr\}_{i=1}^{n}
  \quad \text{where } A_i \equiv \textstyle\bigwedge_{j \in \mathrm{in}(i)} \mathrm{ev}(g_j)\ \mathsf{PathPred},
  \ \ E_i \equiv g_i\ \mathsf{ExistsSpec}
  \ \ \text{(conditional delivery: } \forall\text{-conforming-}\exists\text{-delivering, properness presupposed)} \\
  \bigl\{\ell_i\ \text{latency rows}\bigr\}_{i=1}^{n}
  \ \text{(each premise ships its delivery-stage arithmetic + seriality, at its own profile tier)} \\
  \mathsf{QWireCert}\bigl(w,\ (\mathit{spec}_i)_i \Rightarrow \mathit{spec}_{\mathrm{out}},\ -\bigr)
  \qquad
  \bigwedge_i g_i \vdash \mathrm{marks}(\mathit{spec}_{\mathrm{out}})
  \ \text{(mark-entailment cut)} \\
  \mathsf{MarkFlow}\bigl(w,\ (t_k)_k,\ (P_k)_{k \in \mathrm{cyc}_0(w)}\bigr)
  \ \text{(the schedule walk, two columns: stage rows --- difference-constraint feasibility at the labels; carve rows --- containment $P_k \subseteq \mathrm{mark}_k$ + carved-delivery citations + one species-tagged diagonal certificate per zero-lag cycle, exact values as the degenerate species)} \\
  \mathrm{W0}\ \text{(composite properness --- the existing family; the liveness construction\textquotesingle s tier discharge is stepwise, schedule-relative)}
}{
  \mathsf{SysWire}(w, [\Sigma_1, \ldots, \Sigma_n]) \vDash \mathit{spec}_{\mathrm{out}}
}}{QWireComposeMarks}{blcatrule-QWireComposeMarks}
\paradeentry{\inferrule*{
  S' \vDash \mathit{spec}' \\
  \mathsf{QSim}(\sigma, S, S')
}{
  S \vDash \sigma^{*}\mathit{spec}'
}}{QSimTransport}{blcatrule-QSimTransport}
\paradeentry{\inferrule*{
  \mathsf{QFlatEq}(\mathrm{sys}, \mathit{flat}) \\
  \mathcal{M}_{\mathit{flat}},\ x_0 \in R \vDash \varphi \ \text{(the closed-system pipeline\textquotesingle s judgment)}
}{
  \mathrm{sys},\ x_0 \in R \vDash \varphi
}}{QFlatten}{blcatrule-QFlatten}
\paradeentry{\inferrule*{
  \overline{\mathrm{env}_i : B \vdash\ \pm f_i(x) \ \le\ a + b\,\lVert x \rVert}
}{
  \vDash \mathsf{QLinGrowth}(f,\ a,\ b,\ B)
}}{QLinGrowthDischarge}{blcatrule-QLinGrowthDischarge}
\paradeentry{\inferrule*{
  \overline{\mathrm{dom}_k}\ \ \text{(per-contributor dominance rows)} \\
  \mathsf{qtsl\_stochwire\_shape}(\mathit{cl},\ \mathit{sp})
}{
  \vDash \mathsf{QMarginalCompat}(\mathit{sp},\ \mathit{cl})
}}{QMarginalCompat}{blcatrule-QMarginalCompat}
\paradeentry{\inferrule*{
  (\vec{s}, \vec{k}) \in K \vdash G(\vec{s}, \vec{k}) \le 0 \ \text{(one \rulename{BernCone}-family sequent over the aggregate coordinates)} \\
  s^P_j := \textstyle\sum_{t:P} \phi_j(c_t), \qquad k_P := \#P \ \text{(the aggregate coordinates, defined)} \\
  \bigl\{\vdash l_j\, k_P \le s^P_j \le u_j\, k_P\bigr\}_{j, P} \ \text{(coupling rows; } [l_j, u_j] \ \text{a Bernstein enclosure of } \phi_j \ \text{on } B_P\text{)} \\
  \bigl\{\vdash 0 \le k_P \le M_P\bigr\}_{P} \ \text{(count rows)} \\
  \bigl\{\vdash y^{\top} \vec{k} = y^{\top} \vec{k}(x_0)\bigr\}_{y \in Y} \ \text{(P-invariant conservation rows: each cited witness } y \ge 0,\ y^{\top} C = 0 \ \text{verified against the uncoloured skeleton\textquotesingle s incidence matrix } C\text{; } x_0 \ \text{the initial marking)} \\
  \bigl\{\vdash (s^P_\alpha)^2 \le k_P\, s^P_{2\alpha}\bigr\}_{\alpha, P} \ \text{(moment rows, optional tightening)} \\
  \forall x \in \mathrm{scope},\ \forall t{:}P \ \text{in } x:\ c_t \in B_P \ \text{(colour-box scope)}
}{
  x \in \mathrm{scope} \vdash G\bigl(\vec{s}(x), \vec{k}(x)\bigr) \le 0
}}{QAggReduce}{blcatrule-QAggReduce}
\paradeentry{\inferrule*{
  k_P \ \text{a count observable, integer-valued at every reachable marking} \\
  j \in \mathbb{Z}
}{
  \mathrm{cmpl}\,\{k_P \le j\} \ \rightsquigarrow \ \{k_P \ge j + 1\} \quad \text{(hypothesis positions, reachable scope)}
}}{QCountAtomRounding}{blcatrule-QCountAtomRounding}
\paradeentry{\inferrule*{
  \mathrm{upper} : \mathsf{qpaf\_upper}(n,\ c_{\mathrm{u}}) \\
  \mathrm{lower} : \mathsf{qpaf\_lower}(n,\ c_{\mathrm{l}}) \\
  \mathsf{qtsl\_palways\_fold}(n)
}{
  \mathcal{M} \vDash \mathsf{SProbGE}\bigl(p,\ \mathsf{PAlways}(\mathsf{SAnd}[\mathsf{SAtomLE}(c, V),\ \mathsf{SAtomGE}(c, V)])\bigr)
}}{QPAlwaysFromFold}{blcatrule-QPAlwaysFromFold}
\paradeentry{\inferrule*{
  \mathrm{cert} : \mathsf{HCert} \\
  d_{\min} > 0 \\
  \mathsf{qtsl\_firebound\_fold}(n)
}{
  \mathcal{M} \vDash \mathsf{SReturnLE}\bigl(B,\ \mathsf{retFireCount}(\mathit{trs})\bigr)
}}{QFireBoundFromFold}{blcatrule-QFireBoundFromFold}
\paradeentry{\inferrule*{
  n\,\mathsf{cert} = c \\ n\,\mathsf{goalfold} = f_g \\ n\,\mathsf{horizon} = T \\
  \mathsf{coverage\_lb}(c) \\ \mathsf{flow\_lb}(c) \\
  \forall r.\ u_r > 0 \Rightarrow \mathsf{TrMeter}(\mathit{tr}(r), d_r) \\
  \mathsf{qtsl\_palwayst\_fold}(n)
}{
  \mathcal{M} \vDash \mathsf{SProbGE}\bigl(p,\ \mathsf{PAlwaysT}(T,\ \mathrm{cmpl}(\mathit{goal}))\bigr)
}}{QPAlwaysTFromFold}{blcatrule-QPAlwaysTFromFold}
\paradeentry{\inferrule*{
  \mathsf{unit\_lineage}(c) \\
  \{\mathsf{hrpGap}(1), \mathsf{hrpDropAtMost}(1)\}\ \text{per advancing row on } \chi = -\psi\ \text{(pathwise species)} \\
  \mathsf{TrMeter}(\mathit{tr}(r), d)\ \text{per advancing row} \\
  \forall \tau \in m_0.\ \psi(\tau) \le 0 \\
  T < k \cdot d \\
  \mathsf{qtsl\_palwayst\_lineage}(n)
}{
  \mathcal{M} \vDash \mathsf{SProbGE}\bigl(p,\ \mathsf{PAlwaysT}(T,\ \textstyle\bigwedge_{q\,:\,\mathrm{depth}(q) \ge k} \mathrm{count}_q = 0)\bigr)
}}{QPAlwaysTFromLineage}{blcatrule-QPAlwaysTFromLineage}
\paradeentry{\inferrule*{
  \bigl\{\vdash K(p, \cdot) \ge 0,\ \textstyle\sum_q K(p, q) = 1,\ \pi_{i\#} K(p, \cdot) = \widehat{K}_i(p_i, \cdot)\bigr\} \ \text{(coupling rows, exact } \Q\text{; } \widehat{K}_i \ \text{the completed component kernels)} \\
  \bigl\{\vdash K(p, p) = 1 \ \text{for both-absorbed } p\bigr\} \ \text{(freeze rows)} \\
  \bigl\{\vdash \textstyle\sum_q K(p, q)\, V(q) \le V(p)\bigr\} \ \text{(relational supermartingale generator row, every product state)} \\
  \vdash V \ge 0 \\
  \bigl\{\vdash V \ge 1 \ \text{(disagreeing), resp.\ } V \ge |v_1 - v_2|, \ \text{on both-absorbed states}\bigr\} \ \text{(terminal domination)} \\
  \text{absorption: per-component reach criterion}; \quad \mu_0 \ \text{a probability law with } \pi_{i\#}\mu_0 = \text{the component initial laws (initial marginal rows)}; \quad \vdash \textstyle\sum_p \mu_0(p)\, V(p) \le b
}{
  \mathsf{QDist}[d](\mathcal{M}_1, \mathcal{M}_2) \le b \qquad d \in \{\mathrm{TV},\ \text{gap-pseudometrics}\}
}}{QCoupleDrift}{blcatrule-QCoupleDrift}
\paradeentry{\inferrule*{
  d : \mathsf{QDistRegistry}(X_d); \qquad (X_d, d) \ \text{complete} \\
  \vdash \rho < 1 \ \text{two-point on } \overline{B}(\hat\mu, r) \ \text{(contraction row)} \\
  \vdash d(T\hat\mu, \hat\mu) \le \delta \ \text{(residual row)} \\
  \vdash \delta + \rho\, r \le r \ \text{(radius row, division-free)}
}{
  \vdash \mathsf{QDist}[d](\hat\mu,\ \mathrm{fix}\,T) \le r
}}{QFixBall}{blcatrule-QFixBall}
\paradeentry{\inferrule*{
  J = \bigl(\mathcal{M}_c \vDash s\bigr) \ \text{(strict, law-level)} \quad \text{or} \quad J = \bigl(\mathsf{QDist}[\mathrm{TV}](\mu_c, \nu) \le \delta\bigr), \ \text{on the coarse side} \\
  \bigl\{\vdash \mathrm{tv}\bigl(\pi_{\#}\, K(s, \cdot),\ Q(\pi(s), \cdot)\bigr) \le \varepsilon\bigr\}_{s \in S_{\mathrm{f}}} \ \text{(} \sigma\text{-witness rows, one per fine state; } \varepsilon = 0 \ \text{is exact intertwining)} \\
  \vdash \mu^{0}_{c} = \pi_{\#}\, \mu^{0}_{\mathrm{f}} \ \text{(initial row)} \\
  n \in \N \ \text{(declared horizon; } \varepsilon\text{-cell only)}
}{
  \text{exact cell: } J \ \text{transports unchanged} \qquad \varepsilon\text{-cell: } \mathsf{QDist}[\mathrm{TV}]\bigl(\pi_{\#}(\text{fine law at } n),\ \text{coarse law at } n\bigr) \le n\,\varepsilon
}}{QDistTransport}{blcatrule-QDistTransport}
\paradeentry{\inferrule*{
  q : \mathsf{DAbsQuery} \\
  \mathsf{qtsl\_dabs\_wf}(q) \\
  \text{stored rows of } q.\mathsf{chain}\ (\mathsf{dprob}, \mathsf{dinit}, \mathsf{dclass}, \mathsf{dreject})
}{
  \mathsf{qtsl\_dabs\_absorb\_bracket}(q,\ c,\ lo,\ hi)
}}{DtmcAbsorbBracket}{blcatrule-DtmcAbsorbBracket}
\paradeentry{\inferrule*{
  q : \mathsf{DAbsQuery} \\
  \mathsf{qtsl\_dabs\_wf}(q) \\
  \text{stored rows of } q.\mathsf{chain}\ (\mathsf{dprob}, \mathsf{dinit}, \mathsf{dclass}, \mathsf{dreject}) \\
  S > 0\ \text{(recomputed; refusal otherwise)}
}{
  \mathsf{qtsl\_dabs\_posterior\_bracket}(q,\ c,\ lo,\ hi)
}}{DtmcPosteriorBracket}{blcatrule-DtmcPosteriorBracket}
\paradeentry{\inferrule*{
  w : \mathsf{DGenContraction} \\
  \text{stored rows: } \mathsf{dweight}\ (A = \mathrm{supp}\,V,\ \text{nonempty}),\ \mathsf{dprob},\ \mathsf{dinit} \\
  \textstyle\sum_{y \in A} K(x,y)\,V(y) \le c\,V(x)\ \ \forall x \in A \\
  K(x,y) = 0\ \text{for stored rows } x \notin A,\ y \in A \\
  \min_A V > 0 \\
  \mathsf{qtsl\_dtmc\_genform\_contraction}(w)
}{
  \text{each claimed row of } w:\ \mathsf{qtsl\_dgc\_partial\_lo},\ \mathsf{qtsl\_dgc\_ev\_bracket},\ \mathsf{qtsl\_dgc\_survival\_tail},\ \mathsf{qtsl\_dgc\_z\_one}
}}{QGenFormContraction}{blcatrule-QGenFormContraction}
\paradesep
\paradeblock{}{ProofWellFormedness}{blcatrule-ProofWellFormedness}

\blcatsection{Core Structural Rules}
\rulesection{1}{Ax}{core rule}
\label{blcatrule-Ax}
\ruleblock{The rule (Steele metanotation)}
\sourcenote{as realized in the mechanized development; source row: \texttt{\texttt{rules.toml}} --- \emph{Ax}. The metanotation statement is typeset by hand for notational quality; its mechanized ground truth is the Lean statement below, extracted from the sources at every build of this book and hash-bound. The rule \emph{roster} is machine-checked (every mechanized rule tag round-trips to exactly one catalogue row); the statement \emph{text} is editorial, checked by eye against the extracted Lean beside it.}
\begin{mathpar}
\inferrule*[Right=\rulename{Ax}]{\;}{\Phi \vdash \Phi}
\end{mathpar}

\par\medskip
\rulename{Ax} is the leaf of every derivation: it discharges a sequent
whose conclusion is literally its premise, $\Phi \vdash \Phi$. A
producer uses it wherever a derivation bottoms out at an assumption ---
the demonstration certificate cited under the comparator note below
begins with an \rulename{Ax} node discharging its composite premise
before a conjunction elimination extracts the region fact a
Bernstein-cone closure consumes.

What the checker recomputes is one thing: that the premise and
conclusion are the \emph{same formula tree}, compared node for node ---
no normalization, no reassociation, no matching. Formulas here are
ordered trees, so \rulename{Ax} does not close
$\Phi \vdash \Phi'$ for a logically equivalent but
differently shaped $\Phi'$; that is re-derived through
the conjunction rules when needed.

The single hypothesis is load-bearing in the obvious, total way: drop
the equality and the premise-free rule concludes any sequent at all ---
$\top \vdash \bot$ is the sharpest instance. Notably absent are
well-typedness side conditions, and their absence is correct: whatever
satisfaction amounts to for $\Phi$ --- even a $\Phi$ mentioning untyped
wires or atoms that never evaluate --- it sits on both sides of the
turnstile. The proof section contrasts this with \rulename{EqRefl},
which concludes a freshly minted \emph{atom} and therefore does need
typing fields.

\ruleblock{Lean statement}
\declhead{BernsteinLogic.Soundness}{sound\_ax}{\badge{BadgeGreen}{statement bound}}
\begin{leansrc}
theorem sound_ax {H : Hypernet} {seq : Sequent}
    (h : seq.pre = seq.post) : seq.semanticallyValid H
\end{leansrc}
\par\medskip
\textbf{How to read the payload.}\enspace A certificate node for this
rule carries an inhabitant of the following inductive: its
proof-obligation fields (here, one equality of formulas) are the
rule\textquotesingle s entire side-condition surface, and the soundness
lemma above takes exactly those fields as hypotheses. The six other
structural sections\textquotesingle{} payloads read the same way ---
sub-derivations plus structural equalities, nothing else.

\declhead{BernsteinLogic.Core.Catalogue}{PAxData}{\badge{BadgeGreen}{statement bound}}
\begin{leansrc}
inductive PAxData : Hypernet → Sequent → Type where
  | mk (H : Hypernet) (seq : Sequent)
       (preEqPost : seq.pre = seq.post) : PAxData H seq
\end{leansrc}
\ruleblock{Correspondence (Steele display $\leftrightarrow$ Lean)}
\badge{BadgeGreen}{reviewer-certified against the displayed texts}
\sourcenote{The dictionary below is a reviewer-certified statement of how the mechanized statements above interpret the displayed rule: display metavariable $\mapsto$ Lean binder, each display hypothesis $\mapsto$ the premise or side condition carrying it, the display conclusion $\mapsto$ the goal --- with any deliberate difference in width between display and mechanization named outright. Certification is recorded as a hash triple over (the displayed cell, the extracted statement texts, this dictionary\textquotesingle s own text); if any of the three changes, this block turns amber until a reviewer re-certifies. The correspondence claim itself remains editorial --- reviewer-grade, not machine-checked; its machine-checked successor types each soundness statement against the rule as a structured object, and supersedes this block per rule as it lands.}
\begin{itemize}
\item the judgment frame, shared by every structural dictionary below:
  a sequent is the pair of its premise and conclusion formulas, so a
  displayed judgment $\Phi \vdash \Psi$ $\mapsto$ \texttt{seq :
  Sequent} with \texttt{seq.pre}, \texttt{seq.post} the two sides, and
  asserting it $\mapsto$ the goal \texttt{seq.semanticallyValid H} ---
  at every admissible valuation of the ambient hypernet \texttt{H},
  satisfaction of \texttt{seq.pre} implies satisfaction of
  \texttt{seq.post}.
\item the display\textquotesingle s single metavariable $\Phi$,
  appearing on both sides of the turnstile, $\mapsto$ the sole
  hypothesis \texttt{h} $:$ \texttt{seq.pre = seq.post} --- sharing by
  field equation on one sequent binder, not by a shared metavariable:
  the statement quantifies over \emph{all} sequents and consumes the
  identity as its only input.
\item premise-freeness $\mapsto$ silence: no other hypothesis exists;
  \texttt{H} and \texttt{seq} are arbitrary.
\item the payload row \texttt{PAxData} (certified alongside) carries
  the same single tie as certificate data (\texttt{preEqPost}), so
  admission layer and soundness statement coincide exactly on this
  rule.
\end{itemize}
No width difference in either direction: one hypothesis, consumed by
rewriting; nothing else is stated or needed.

\ruleblock{Lean verifier receipts}
\sourcenote{module \texttt{BernsteinLogic.Soundness}; bound build \texttt{a20d90e66c13\ldots} (bound from receipt a9c0e34d3e3c9c8c, 2026-08-03T21:30:36Z). A receipt records that named checking layers accepted a build: (module, compiled-artifact SHA-256, environment SHA-256, checker layers), minted when the checking gate ran; the store receipts attesting the bound build are listed with the layers each run credited. Nothing here asks the reader to trust the store: every layer is an independently runnable program over the quoted sources --- rebuild the module and re-run the checkers to reproduce the verdict. Receipts attesting superseded builds credit no layer and are omitted from this rendering (the count of omissions, if any, is noted below).}
\begin{receipttable}
\receiptrow{a9c0e34d3e3c9c8c}{2026-08-03T21:30:36Z}{build+decorrelated}{\badge{BadgeGreen}{current}}
\end{receipttable}
\sourcenote{21 further receipt(s) in the store attest superseded builds of this module and are omitted from this rendering.}
\sourcenote{module \texttt{BernsteinLogic.Core.Catalogue}; bound build \texttt{cbc92e976a1d\ldots} (bound from receipt 5abdb3589707b5df, 2026-08-03T21:30:33Z). A receipt records that named checking layers accepted a build: (module, compiled-artifact SHA-256, environment SHA-256, checker layers), minted when the checking gate ran; the store receipts attesting the bound build are listed with the layers each run credited. Nothing here asks the reader to trust the store: every layer is an independently runnable program over the quoted sources --- rebuild the module and re-run the checkers to reproduce the verdict. Receipts attesting superseded builds credit no layer and are omitted from this rendering (the count of omissions, if any, is noted below).}
\begin{receipttable}
\receiptrow{24804e66edba0bc9}{2026-08-03T20:25:09Z}{build+decorrelated}{\badge{BadgeGreen}{current}}
\receiptrow{5abdb3589707b5df}{2026-08-03T21:30:33Z}{build+decorrelated}{\badge{BadgeGreen}{current}}
\end{receipttable}
\sourcenote{18 further receipt(s) in the store attest superseded builds of this module and are omitted from this rendering.}
\layerline{\layeritem{build}{lake build (elaboration)}{\badge{BadgeGreen}{current}}\layeritem{replay}{leanchecker (official kernel replay)}{\badge{BadgeAmber}{stale only}}\layeritem{hardened}{safe\_verify (hardened screen)}{\badge{BadgeAmber}{stale only}}\layeritem{decorrelated}{lean4export + nanoda (decorrelated kernel)}{\badge{BadgeGreen}{current}}}
\sourcenote{\badge{BadgeBlue}{comparator: exercised via the demonstration fixture}
the composite demonstration certificate (an axiom node, a conjunction
elimination, a Bernstein-cone closure) runs as a fixture through the
three-implementation comparator harness (compiled / kernel-interpreted /
reference-interpreted), so this rule\textquotesingle s node rides every
execution path of that harness.
}
\ruleblock{HOL4 statement}
\sourcenote{theory \texttt{bern\_fam\_ax}, HOL4 development script \texttt{\texttt{bern\_fam\_axScript.sml}} \badge{BadgeGreen}{script bound}}
\holhead{Definition checkFam\_ax\_def}
\begin{holsrc}
Definition checkFam_ax_def:
  checkFam_ax (cols:columns) (ct:ctx) (i:num) ⇔
    case conclRef cols ct i of
    | NONE => F
    | SOME cj =>
      case poolTagAt cols ct cj of
      | NONE => F
      | SOME t =>
        if t ≠ (48:num) then F else
        case poolField cols ct cj 0 of
        | NONE => F
        | SOME preR =>
          case poolField cols ct cj 1 of
          | NONE => F
          | SOME postR =>
            if preR = postR then T
            else
              case pv cols ct cj of
              | SOME (PJudg prew postw) =>
                  (expandSph FTt prew = expandSph FTt postw)
              | _ => F
End
\end{holsrc}
\holhead{Theorem famEq\_ax}
\begin{holsrc}
Theorem famEq_ax:
  wfColumns cols ct ∧ mkCtx cols = SOME ct ∧
  nodeCorr cols ct i (CertPf sq nd) ∧
  FST (THE (tagOf cols i)) ∈ {0}
  ⇒
  (checkFam_ax cols ct i ⇔ (sq.pre = sq.post))
\end{holsrc}
\ruleblock{HOL verifier receipts}
\badge{BadgeBlue}{documentary} Documentary citations. Spine census: \texttt{bern\_joins} GREEN under \texttt{hol-strict} (55 theorems, 0 offences, 208-theory ancestry including this theory), \texttt{checkJoins\_eq\_checkPf} verified verbatim in the built signature. The cited \texttt{famEq} lemma is the kernel-checked per-node reduction: under well-formed columns (\texttt{wfColumns}/\texttt{mkCtx}/\texttt{nodeCorr}) at wire tag 0, the columnar family check holds iff the node\textquotesingle s sequent satisfies the rule\textquotesingle s own side condition ($\mathtt{sq.pre} = \mathtt{sq.post}$). There is no signed HOL receipt store yet --- these citations are documentary.
Semantic stratum, also ported: the soundness twin \texttt{sound\_ax}
(theory \texttt{bern\_tr\_soundstruct}: \texttt{hol-strict} clean, 9
theorems, 0 offences, verification run of 2026-07-22 on an isolated
build host, content-addressed at the same sources), consumed by the
certified-fragment keystone (\texttt{checkPf\_sound\_frag}, theory
\texttt{bern\_tr\_splice}, 20 theorems, 0 offences, same run): a
reference-checker accept inside the proven wire fragment yields sequent
validity.

\sourcenote{HOL receipts are documentary: quoted results of \texttt{hol-strict} runs
and census rows from the HOL4 development\textquotesingle s build
evidence, not minted receipt artefacts. A HOL receipt store (the Lean
store\textquotesingle s twin) is named future work.
}
\ruleblock{Proof}

\sourcenote{Genre: core rule (tier 1) --- an ordinary mathematical proof
at Putnam-exam-solution detail. This section also fixes the semantic
setting shared by the seven structural sections (\rulename{Ax},
\rulename{Weaken}, \rulename{Cut}, $\wedge$I, $\wedge$E, $\vee$I,
$\vee$E); the later six abbreviate it as ``the \rulename{Ax} setting''.}

\textbf{Setting (shared by the structural sections).}\enspace Fix a
hypernet $H$ --- the ambient computation graph: typed wires, each
declaring a tensor shape and a dtype, and primitive applications
connecting them. A \emph{valuation} $\sigma$ assigns each typed wire a
tensor of values from its declared value set; $\sigma$ is
\emph{admissible} when every primitive application in $H$ satisfies its
semantic relation. A formula holds at $\sigma$ by structural recursion:
an atom $t_1 \bowtie t_2$ holds when --- and only when --- both terms
evaluate --- exact rational arithmetic, wires read whole, evaluation
failing on untyped wires, on shape-incompatible arithmetic, and on any
$\pm\infty$-valued wire mention --- to tensors compared entrywise under
$\bowtie$ (a rank-0 side broadcasts; incompatible shapes never hold), so
an atom either of whose terms fails to evaluate does not hold; a region
atom holds
when its term evaluates to a rational point of the named region;
$\Phi_1 \wedge \Phi_2$ holds when both conjuncts hold, $\Phi_1 \vee
\Phi_2$ when at least one disjunct holds, $\top$ always, $\bot$ never. A
sequent $\Phi \vdash \Psi$ is \emph{semantically valid} over $H$ when
every admissible $\sigma$ satisfying $\Phi$ satisfies $\Psi$.

Two facts frame all seven structural proofs. First, formulas are
\emph{ordered trees}: $\wedge$ and $\vee$ are binary constructors, and
every side condition of a structural rule is an identity of trees,
checked structurally --- no normalization, no reassociation, no
commuting. Wherever a proof below ``rewrites by'' a payload field, it is
replacing a formula by itself. Second, none of the seven proofs consumes
admissibility: each threads it, unexamined, to its sub-derivations, so
the structural fragment is sound over \emph{arbitrary} valuations ---
premise and conclusion validity both read over all valuations,
admissible or not --- and the admissibility hypothesis rides along only
because the validity statement quantifies over it. Where a counterexample is claimed in these
sections, read it over the one-wire hypernet ($w$ rational, rank-0, no
applications): there every valuation is admissible, so an unsound
conclusion genuinely falsifies.

\textbf{Claim (\texttt{sound\_ax}).}\enspace The payload certifies one
equality, $\mathrm{pre} = \mathrm{post}$ (\texttt{preEqPost}); under it
the premise-free rule concluding $\Phi \vdash \Phi$ is sound.

\begin{bthm}[soundness of \rulename{Ax}]\label{thm:ax-sound}
If $\mathrm{seq.pre} = \mathrm{seq.post}$, then $\mathrm{seq}$ is
semantically valid over every hypernet.
\end{bthm}
\begin{proof}
One step. Let $\sigma$ be an admissible valuation with $\sigma \models
\mathrm{seq.pre}$. The certified equality says $\mathrm{seq.post}$
\emph{is} $\mathrm{seq.pre}$ --- the same tree, not merely an equivalent
formula --- so $\sigma \models \mathrm{seq.post}$.
\end{proof}

\begin{bremark}[the one hypothesis is everything]\label{rem:ax-hyp}
\rulename{Ax} has no premises, so the payload equality is its entire
content. Drop \texttt{preEqPost} and the degenerate rule concludes an
arbitrary sequent; the sharpest instance is $\top \vdash \bot$ ---
every valuation satisfies the premise and none the conclusion.
\end{bremark}

\begin{bremark}[identity is free; reflexivity is not]\label{rem:ax-typing}
\rulename{Ax} carries no well-typedness side conditions, and this is
correct rather than an oversight: $\Phi \vdash \Phi$ is valid even when
$\Phi$ mentions untyped wires or atoms that can never evaluate, because
satisfaction --- whatever it turns out to be --- sits on both sides of
the turnstile. Contrast \rulename{EqRefl}, which concludes the
\emph{atom} $t = t$ under an arbitrary premise and does need shape and
finiteness fields: an atom is a claim about evaluation, and an ill-typed
$t$ never evaluates, so there the conclusion can fail at valuations the
premise never constrained. \rulename{Ax} never crosses that line because
it introduces no content.
\end{bremark}

\begin{bremark}[structural identity, not logical equivalence]
The checker compares the two formula trees node for node.
$\Phi \vdash \Phi'$ for a logically equivalent but differently shaped
$\Phi'$ --- $B \wedge A$ against $A \wedge B$, say --- is \emph{not} an
\rulename{Ax} instance; where needed it is re-derived through $\wedge$E
and $\wedge$I. In the lattice reading of the structural fragment
(developed at the $\vee$E section), \rulename{Ax} is reflexivity of the
entailment order --- the Gentzen-style leaf \cite{gentzen1935},
discharged semantically rather than by proof search.
\end{bremark}

\begin{anchors}
Theorem \ref{thm:ax-sound} & \texttt{sound\_ax}; payload \texttt{PAxData} (field \texttt{preEqPost}); wired as \texttt{RuleData.ax} into \texttt{ProofObject.sound}\\
ordered-tree formula grammar & \texttt{Formula} (constructors \texttt{and}/\texttt{or}); satisfaction \texttt{Formula.holds}\\
validity statement & \texttt{Sequent.semanticallyValid}\\
lattice reading & \texttt{M3.le\_refl} (the $\vee$E section's countermodel kit)\\
\end{anchors}


\rulesection{2}{Weaken}{core rule}
\label{blcatrule-Weaken}
\ruleblock{The rule (Steele metanotation)}
\sourcenote{as realized in the mechanized development; source row: \texttt{\texttt{rules.toml}} --- \emph{Weaken}. The metanotation statement is typeset by hand for notational quality; its mechanized ground truth is the Lean statement below, extracted from the sources at every build of this book and hash-bound. The rule \emph{roster} is machine-checked (every mechanized rule tag round-trips to exactly one catalogue row); the statement \emph{text} is editorial, checked by eye against the extracted Lean beside it.}
\begin{mathpar}
\inferrule*[Right=\rulename{Weaken}]{\Phi \vdash \Psi}{\Phi \wedge \chi \vdash \Psi}
\end{mathpar}

\par\medskip
\rulename{Weaken} moves a proven sequent under a larger set of
assumptions: from $\Phi \vdash \Psi$ to
$\Phi \wedge \chi \vdash \Psi$. A producer reaches for it when contexts
must be made to agree --- the structural rules with two premises
(\rulename{Cut}, $\wedge$I, $\vee$E) and the composition
rule\textquotesingle s containment check all compare premise trees
structurally, so a sequent proven from fewer assumptions is weakened
into the context the next node expects.

The checker recomputes two tree identities: the
conclusion\textquotesingle s premise is exactly the
sub-derivation\textquotesingle s premise conjoined \emph{on the right}
with the named conjunct $\chi$ (one new $\wedge$-node --- position is
data, since formulas are ordered trees), and the conclusion is carried
over unchanged. Naming $\chi$ in the payload keeps this a comparison,
never a search.

Each hypothesis earns its keep. Drop the premise identity and
weakening could silently \emph{replace} assumptions rather than extend
them --- the proof section\textquotesingle s counterexample derives $\top \vdash w \leq 0$ from the valid
$w \leq 0 \vdash w \leq 0$. Drop the conclusion identity and the
conclusion is free. Nothing at all is demanded of $\chi$ itself: an
unsatisfiable or ill-typed conjunct only makes the new premise harder to
satisfy, which costs applicability, never soundness --- the rule only
discards information.

\ruleblock{Lean statement}
\declhead{BernsteinLogic.Soundness}{sound\_weaken}{\badge{BadgeGreen}{statement bound}}
\begin{leansrc}
theorem sound_weaken {H : Hypernet} {seq premise : Sequent} {extra : Formula}
    (ih : premise.semanticallyValid H)
    (hpre : seq.pre = .and premise.pre extra)
    (hpost : seq.post = premise.post) : seq.semanticallyValid H
\end{leansrc}
\par\medskip
Payload, read as in the \rulename{Ax} section: a sub-derivation plus
structural equalities, nothing else.

\declhead{BernsteinLogic.Core.Catalogue}{PWeakenData}{\badge{BadgeGreen}{statement bound}}
\begin{leansrc}
inductive PWeakenData : Hypernet → Sequent → Type where
  | mk (H : Hypernet) (seq : Sequent)
       (premise : Sequent) (premiseProof : ProofObject H premise)
       (extra : Formula)
       (preEq  : seq.pre  = .and premise.pre extra)
       (postEq : seq.post = premise.post)
       : PWeakenData H seq
\end{leansrc}
\ruleblock{Correspondence (Steele display $\leftrightarrow$ Lean)}
\badge{BadgeGreen}{reviewer-certified against the displayed texts}
\sourcenote{The dictionary below is a reviewer-certified statement of how the mechanized statements above interpret the displayed rule: display metavariable $\mapsto$ Lean binder, each display hypothesis $\mapsto$ the premise or side condition carrying it, the display conclusion $\mapsto$ the goal --- with any deliberate difference in width between display and mechanization named outright. Certification is recorded as a hash triple over (the displayed cell, the extracted statement texts, this dictionary\textquotesingle s own text); if any of the three changes, this block turns amber until a reviewer re-certifies. The correspondence claim itself remains editorial --- reviewer-grade, not machine-checked; its machine-checked successor types each soundness statement against the rule as a structured object, and supersedes this block per rule as it lands.}
Judgment frame as in \rulename{Ax}.
\begin{itemize}
\item the premise judgment $\Phi \vdash \Psi$ $\mapsto$ \texttt{ih}
  $:$ \texttt{premise.semanticallyValid H} --- in the proof-object
  recursion this hypothesis is the sub-derivation\textquotesingle s
  induction hypothesis; the payload (\texttt{PWeakenData}, certified
  alongside) carries the sub-derivation itself as a
  \texttt{ProofObject} field beside the same two ties.
\item $\Phi \mapsto$ \texttt{premise.pre}; $\Psi \mapsto$
  \texttt{premise.post}; $\chi \mapsto$ the binder \texttt{extra :
  Formula}, of which nothing is demanded.
\item conclusion $\Phi \wedge \chi \vdash \Psi$ $\mapsto$ the goal
  under \texttt{hpre} $:$ \texttt{seq.pre = .and premise.pre extra}
  --- the displayed conjunction in the displayed order ($\Phi$ left,
  $\chi$ right; formulas are ordered trees, so the position is part of
  the certified statement) --- and \texttt{hpost} $:$ \texttt{seq.post
  = premise.post}.
\end{itemize}
No width difference in either direction: every hypothesis is consumed,
and the statement proves exactly the displayed rule at every hypernet
and sequent.

\ruleblock{Lean verifier receipts}
\sourcenote{module \texttt{BernsteinLogic.Soundness}; bound build \texttt{a20d90e66c13\ldots} (bound from receipt a9c0e34d3e3c9c8c, 2026-08-03T21:30:36Z). A receipt records that named checking layers accepted a build: (module, compiled-artifact SHA-256, environment SHA-256, checker layers), minted when the checking gate ran; the store receipts attesting the bound build are listed with the layers each run credited. Nothing here asks the reader to trust the store: every layer is an independently runnable program over the quoted sources --- rebuild the module and re-run the checkers to reproduce the verdict. Receipts attesting superseded builds credit no layer and are omitted from this rendering (the count of omissions, if any, is noted below).}
\begin{receipttable}
\receiptrow{a9c0e34d3e3c9c8c}{2026-08-03T21:30:36Z}{build+decorrelated}{\badge{BadgeGreen}{current}}
\end{receipttable}
\sourcenote{21 further receipt(s) in the store attest superseded builds of this module and are omitted from this rendering.}
\sourcenote{module \texttt{BernsteinLogic.Core.Catalogue}; bound build \texttt{cbc92e976a1d\ldots} (bound from receipt 5abdb3589707b5df, 2026-08-03T21:30:33Z). A receipt records that named checking layers accepted a build: (module, compiled-artifact SHA-256, environment SHA-256, checker layers), minted when the checking gate ran; the store receipts attesting the bound build are listed with the layers each run credited. Nothing here asks the reader to trust the store: every layer is an independently runnable program over the quoted sources --- rebuild the module and re-run the checkers to reproduce the verdict. Receipts attesting superseded builds credit no layer and are omitted from this rendering (the count of omissions, if any, is noted below).}
\begin{receipttable}
\receiptrow{24804e66edba0bc9}{2026-08-03T20:25:09Z}{build+decorrelated}{\badge{BadgeGreen}{current}}
\receiptrow{5abdb3589707b5df}{2026-08-03T21:30:33Z}{build+decorrelated}{\badge{BadgeGreen}{current}}
\end{receipttable}
\sourcenote{18 further receipt(s) in the store attest superseded builds of this module and are omitted from this rendering.}
\layerline{\layeritem{build}{lake build (elaboration)}{\badge{BadgeGreen}{current}}\layeritem{replay}{leanchecker (official kernel replay)}{\badge{BadgeAmber}{stale only}}\layeritem{hardened}{safe\_verify (hardened screen)}{\badge{BadgeAmber}{stale only}}\layeritem{decorrelated}{lean4export + nanoda (decorrelated kernel)}{\badge{BadgeGreen}{current}}}
\sourcenote{\badge{BadgeGray}{comparator: none} no fixture in the three-implementation comparator harness (compiled / kernel-interpreted / reference-interpreted) exercises this rule.}
\ruleblock{HOL4 statement}
\sourcenote{theory \texttt{bern\_fam\_resid}, HOL4 development script \texttt{\texttt{bern\_fam\_residScript.sml}} \badge{BadgeGreen}{script bound}}
\holhead{Definition checkFam\_weaken\_def}
\begin{holsrc}
Definition checkFam_weaken_def:
  checkFam_weaken (cols:columns) (ct:ctx) (i:num) ⇔
    case frSeq cols ct i of
    | NONE => F
    | SOME sq =>
      case frChildSeq cols ct i 0 of
      | NONE => F
      | SOME ps =>
        case frFormula cols ct i 0 of
        | NONE => F
        | SOME extra =>
            sq.pre = FAnd ps.pre extra ∧ sq.post = ps.post
End
\end{holsrc}
\holhead{Theorem famEq\_weaken}
\begin{holsrc}
Theorem famEq_weaken:
  wfColumns cols ct ∧ mkCtx cols = SOME ct ∧
  parseCert (cols.hdrImage ++ rest) = SOME (ws, pf) ∧
  nodeCorr cols ct i (CertPf sq nd) ∧
  FST (THE (tagOf cols i)) = 1 ⇒
  (checkFam_weaken cols ct i ⇔ armWeaken ws sq nd)
\end{holsrc}
\ruleblock{HOL verifier receipts}
\badge{BadgeBlue}{documentary} Documentary citations. Family validation: \texttt{bern\_fam\_resid} built GREEN with all ten \texttt{famEq} arms validated, \texttt{hol-strict} clean; spine census: \texttt{bern\_joins} GREEN under \texttt{hol-strict} (55 theorems, 0 offences, 208-theory ancestry including this theory), \texttt{checkJoins\_eq\_checkPf} verified verbatim in the built signature. The cited \texttt{famEq} lemma is the kernel-checked per-node equivalence: under well-formed columns (\texttt{wfColumns}/\texttt{mkCtx}/\texttt{nodeCorr}) at this rule\textquotesingle s wire tag, the columnar family check coincides with the flat reference walk\textquotesingle s arm for this rule. There is no signed HOL receipt store yet --- these citations are documentary.
Semantic stratum, also ported: the soundness twin \texttt{sound\_weaken}
(theory \texttt{bern\_tr\_soundstruct}: \texttt{hol-strict} clean, 9
theorems, 0 offences, verification run of 2026-07-22 on an isolated
build host, content-addressed at the same sources), consumed by the
certified-fragment keystone (\texttt{checkPf\_sound\_frag}, theory
\texttt{bern\_tr\_splice}, 20 theorems, 0 offences, same run): a
reference-checker accept inside the proven wire fragment yields sequent
validity.

\sourcenote{HOL receipts are documentary: quoted results of \texttt{hol-strict} runs
and census rows from the HOL4 development\textquotesingle s build
evidence, not minted receipt artefacts. A HOL receipt store (the Lean
store\textquotesingle s twin) is named future work.
}
\ruleblock{Proof}

\sourcenote{Genre: core rule (tier 1). Setting as in the \rulename{Ax}
section: valuations, admissibility, entrywise atom satisfaction,
semantic validity; counterexamples read over the one-wire hypernet.}

\textbf{Claim (\texttt{sound\_weaken}).}\enspace The payload supplies a
premise sequent $\Phi \vdash \Psi$ with its sub-derivation, the added
conjunct $\chi$ (\texttt{extra}), and two tree identities: the
conclusion's premise is $\Phi \wedge \chi$ (\texttt{preEq}) and its
conclusion is $\Psi$ (\texttt{postEq}). If $\Phi \vdash \Psi$ is
semantically valid, so is $\Phi \wedge \chi \vdash \Psi$.

\begin{bthm}[soundness of \rulename{Weaken}]\label{thm:weaken-sound}
Under \texttt{preEq} and \texttt{postEq}, semantic validity of the
premise sequent implies semantic validity of the conclusion sequent,
over every hypernet.
\end{bthm}
\begin{proof}
Let $\sigma$ be admissible with $\sigma \models \mathrm{seq.pre}$ ---
by \texttt{preEq}, the tree $\Phi \wedge \chi$.

\emph{Step 1 (project).} The $\wedge$-clause of satisfaction is
meta-conjunction, so $\sigma \models \Phi$ and $\sigma \models \chi$;
keep the left component and discard the right --- this rule's whole
activity is licensed forgetting.

\emph{Step 2 (apply the premise).} Premise validity at $\sigma$ gives
$\sigma \models \Psi$, which \texttt{postEq} pins as the conclusion's
post.
\end{proof}

\begin{bremark}[\texttt{preEq} forces extension, not replacement]
\label{rem:weaken-preeq}
The identity $\mathrm{seq.pre} = \Phi \wedge \chi$ makes the conclusion's
premise contain the sub-derivation's premise \emph{verbatim} as the left
conjunct of one new $\wedge$-node, with the named $\chi$ on the right.
Drop it and ``weakening'' could \emph{replace} the assumptions instead of
extending them: from $w \leq 0 \vdash w \leq 0$ (an \rulename{Ax}
instance) conclude $\top \vdash w \leq 0$ --- false at $\sigma(w) = 1$.
This context-swap family is the standard counterexample to every dropped
\texttt{preEq} in the structural sections; we cite it from here on
rather than restating it.
\end{bremark}

\begin{bremark}[positional discipline]
$\chi$ arrives as the right child of a single new $\wedge$-node:
$\Phi \wedge \chi$, never $\chi \wedge \Phi$, and never merged into
$\Phi$'s own tree. Any further placement is $\wedge$E/$\wedge$I work.
Naming $\chi$ in the payload keeps the checker's comparison pure
structure --- one tree identity, no matching or search.
\end{bremark}

\begin{bremark}[$\chi$ is unconstrained, harmlessly]
Like \rulename{Ax}, the rule carries no typing side conditions on what
it touches: $\chi$ may be unsatisfiable ($\bot$), ill-typed, or mention
untyped wires. At worst the conclusion's premise becomes harder or
impossible to satisfy --- vacuous validity, never unsoundness.
Weakening only discards information; rules that must \emph{evaluate}
are where typing fields live.
\end{bremark}

\begin{bremark}[monotonicity, locally and wholesale]
In the lattice reading (the $\vee$E section), meet is a lower bound:
the interpretation of $\Phi \wedge \chi$ lies below that of $\Phi$,
which lies below that of $\Psi$. \rulename{Weaken}
is the one-step premise-monotonicity of the consequence relation; the
composition rule's conjunct-wise premise-containment check
(\rulename{Syzygy}'s \texttt{preContains}) is the same monotonicity used
wholesale.
\end{bremark}

\begin{anchors}
Steps 1--2 & \texttt{sound\_weaken}; payload \texttt{PWeakenData} (fields \texttt{extra}, \texttt{preEq}, \texttt{postEq}); wired as \texttt{RuleData.weaken}\\
$\wedge$-clause of satisfaction & \texttt{Formula.holds} (case \texttt{and})\\
lattice reading & \texttt{M3.weaken\_kit}\\
\end{anchors}


\rulesection{3}{Cut}{core rule}
\label{blcatrule-Cut}
\ruleblock{The rule (Steele metanotation)}
\sourcenote{as realized in the mechanized development; source row: \texttt{\texttt{rules.toml}} --- \emph{Cut}. The metanotation statement is typeset by hand for notational quality; its mechanized ground truth is the Lean statement below, extracted from the sources at every build of this book and hash-bound. The rule \emph{roster} is machine-checked (every mechanized rule tag round-trips to exactly one catalogue row); the statement \emph{text} is editorial, checked by eye against the extracted Lean beside it.}
\begin{mathpar}
\inferrule*[Right=\rulename{Cut}]{\Phi \vdash \chi \\ \Phi \wedge \chi \vdash \Psi}{\Phi \vdash \Psi}
\end{mathpar}

\par\medskip
\rulename{Cut} is the lemma-staging rule: prove $\Phi \vdash \chi$
once, prove $\Phi \wedge \chi \vdash \Psi$ under the enlarged context,
conclude $\Phi \vdash \Psi$. A producer uses it to factor a
certificate the way a paper factors a proof --- a bound or region fact
is established once and then consumed as an assumption downstream;
iterated cuts stage a chain of such lemmas.

The checker recomputes four tree identities: the
conclusion\textquotesingle s premise and conclusion are inherited from
the left and right sub-derivations respectively; the left
sub-derivation concludes exactly the named cut formula $\chi$; and the
right sub-derivation\textquotesingle s premise is exactly the left
one\textquotesingle s premise conjoined with $\chi$ --- the same
context tree, verbatim. The payload must name $\chi$ explicitly
because the conclusion mentions it nowhere: \rulename{Cut} is the one
structural rule whose certified content is invisible in its
conclusion, so there is nothing to reconstruct it from.

The two interesting hypotheses fail independently, with near-twin
counterexamples walked in the proof section: unshare the context and a
fact proven under foreign assumptions leaks into the conclusion
($\top \vdash w \leq 0$ from two individually valid premises); unmoor
$\chi$ from the left conclusion and the right premise consumes an
assumption nobody earned (same false conclusion, different route). The
proof also records why no cut-elimination enters anywhere: soundness of
the rule is direct, and the checker never normalizes a proof.

\ruleblock{Lean statement}
\declhead{BernsteinLogic.Soundness}{sound\_cut}{\badge{BadgeGreen}{statement bound}}
\begin{leansrc}
theorem sound_cut {H : Hypernet} {seq leftPremise rightPremise : Sequent}
    {cutFormula : Formula}
    (ihl : leftPremise.semanticallyValid H)
    (ihr : rightPremise.semanticallyValid H)
    (hpre : seq.pre = leftPremise.pre)
    (hpost : seq.post = rightPremise.post)
    (hcutL : leftPremise.post = cutFormula)
    (hcutR : rightPremise.pre = .and leftPremise.pre cutFormula) :
    seq.semanticallyValid H
\end{leansrc}
\par\medskip
Payload, read as in the \rulename{Ax} section: sub-derivations plus
structural equalities, with the cut formula as explicit data.

\declhead{BernsteinLogic.Core.Catalogue}{PCutData}{\badge{BadgeGreen}{statement bound}}
\begin{leansrc}
inductive PCutData : Hypernet → Sequent → Type where
  | mk (H : Hypernet) (seq : Sequent)
       (cutFormula : Formula)
       (leftPremise  : Sequent) (leftProof  : ProofObject H leftPremise)
       (rightPremise : Sequent) (rightProof : ProofObject H rightPremise)
       (preEq    : seq.pre  = leftPremise.pre)
       (postEq   : seq.post = rightPremise.post)
       (cutLeft  : leftPremise.post = cutFormula)
       (cutRight : rightPremise.pre = .and leftPremise.pre cutFormula)
       : PCutData H seq
\end{leansrc}
\ruleblock{Correspondence (Steele display $\leftrightarrow$ Lean)}
\badge{BadgeGreen}{reviewer-certified against the displayed texts}
\sourcenote{The dictionary below is a reviewer-certified statement of how the mechanized statements above interpret the displayed rule: display metavariable $\mapsto$ Lean binder, each display hypothesis $\mapsto$ the premise or side condition carrying it, the display conclusion $\mapsto$ the goal --- with any deliberate difference in width between display and mechanization named outright. Certification is recorded as a hash triple over (the displayed cell, the extracted statement texts, this dictionary\textquotesingle s own text); if any of the three changes, this block turns amber until a reviewer re-certifies. The correspondence claim itself remains editorial --- reviewer-grade, not machine-checked; its machine-checked successor types each soundness statement against the rule as a structured object, and supersedes this block per rule as it lands.}
Judgment frame as in \rulename{Ax}; sub-derivation premises as at
\rulename{Weaken}.
\begin{itemize}
\item the two premise judgments $\Phi \vdash \chi$ and $\Phi \wedge
  \chi \vdash \Psi$ $\mapsto$ \texttt{ihl} $:$
  \texttt{leftPremise.semanticallyValid H} and \texttt{ihr} $:$
  \texttt{rightPremise.semanticallyValid H}.
\item $\Phi \mapsto$ \texttt{leftPremise.pre}; $\Psi \mapsto$
  \texttt{rightPremise.post}; $\chi \mapsto$ the binder
  \texttt{cutFormula}, tied to the left conclusion by \texttt{hcutL}
  $:$ \texttt{leftPremise.post = cutFormula} --- the explicit naming
  the commentary explains (the conclusion mentions $\chi$ nowhere).
\item the second premise\textquotesingle s displayed context $\Phi
  \wedge \chi$ $\mapsto$ \texttt{hcutR} $:$ \texttt{rightPremise.pre =
  .and leftPremise.pre cutFormula} --- the same context tree verbatim,
  cut formula conjoined on the right.
\item conclusion $\Phi \vdash \Psi$ $\mapsto$ the goal under
  \texttt{hpre} $:$ \texttt{seq.pre = leftPremise.pre} and
  \texttt{hpost} $:$ \texttt{seq.post = rightPremise.post}.
\end{itemize}
No width difference in either direction: the six hypotheses are
exactly the display\textquotesingle s two premises and four structural
identities, each consumed.

\ruleblock{Lean verifier receipts}
\sourcenote{module \texttt{BernsteinLogic.Soundness}; bound build \texttt{a20d90e66c13\ldots} (bound from receipt a9c0e34d3e3c9c8c, 2026-08-03T21:30:36Z). A receipt records that named checking layers accepted a build: (module, compiled-artifact SHA-256, environment SHA-256, checker layers), minted when the checking gate ran; the store receipts attesting the bound build are listed with the layers each run credited. Nothing here asks the reader to trust the store: every layer is an independently runnable program over the quoted sources --- rebuild the module and re-run the checkers to reproduce the verdict. Receipts attesting superseded builds credit no layer and are omitted from this rendering (the count of omissions, if any, is noted below).}
\begin{receipttable}
\receiptrow{a9c0e34d3e3c9c8c}{2026-08-03T21:30:36Z}{build+decorrelated}{\badge{BadgeGreen}{current}}
\end{receipttable}
\sourcenote{21 further receipt(s) in the store attest superseded builds of this module and are omitted from this rendering.}
\sourcenote{module \texttt{BernsteinLogic.Core.Catalogue}; bound build \texttt{cbc92e976a1d\ldots} (bound from receipt 5abdb3589707b5df, 2026-08-03T21:30:33Z). A receipt records that named checking layers accepted a build: (module, compiled-artifact SHA-256, environment SHA-256, checker layers), minted when the checking gate ran; the store receipts attesting the bound build are listed with the layers each run credited. Nothing here asks the reader to trust the store: every layer is an independently runnable program over the quoted sources --- rebuild the module and re-run the checkers to reproduce the verdict. Receipts attesting superseded builds credit no layer and are omitted from this rendering (the count of omissions, if any, is noted below).}
\begin{receipttable}
\receiptrow{24804e66edba0bc9}{2026-08-03T20:25:09Z}{build+decorrelated}{\badge{BadgeGreen}{current}}
\receiptrow{5abdb3589707b5df}{2026-08-03T21:30:33Z}{build+decorrelated}{\badge{BadgeGreen}{current}}
\end{receipttable}
\sourcenote{18 further receipt(s) in the store attest superseded builds of this module and are omitted from this rendering.}
\layerline{\layeritem{build}{lake build (elaboration)}{\badge{BadgeGreen}{current}}\layeritem{replay}{leanchecker (official kernel replay)}{\badge{BadgeAmber}{stale only}}\layeritem{hardened}{safe\_verify (hardened screen)}{\badge{BadgeAmber}{stale only}}\layeritem{decorrelated}{lean4export + nanoda (decorrelated kernel)}{\badge{BadgeGreen}{current}}}
\sourcenote{\badge{BadgeGray}{comparator: none} no fixture in the three-implementation comparator harness (compiled / kernel-interpreted / reference-interpreted) exercises this rule.}
\ruleblock{HOL4 statement}
\sourcenote{theory \texttt{bern\_fam\_resid}, HOL4 development script \texttt{\texttt{bern\_fam\_residScript.sml}} \badge{BadgeGreen}{script bound}}
\holhead{Definition checkFam\_cut\_def}
\begin{holsrc}
Definition checkFam_cut_def:
  checkFam_cut (cols:columns) (ct:ctx) (i:num) ⇔
    case frSeq cols ct i of
    | NONE => F
    | SOME sq =>
      case frChildSeq cols ct i 0 of
      | NONE => F
      | SOME ls =>
        case frChildSeq cols ct i 1 of
        | NONE => F
        | SOME rs =>
          case frFormula cols ct i 0 of
          | NONE => F
          | SOME cutf =>
              sq.pre = ls.pre ∧ sq.post = rs.post ∧
              ls.post = cutf ∧ rs.pre = FAnd ls.pre cutf
End
\end{holsrc}
\holhead{Theorem famEq\_cut}
\begin{holsrc}
Theorem famEq_cut:
  wfColumns cols ct ∧ mkCtx cols = SOME ct ∧
  parseCert (cols.hdrImage ++ rest) = SOME (ws, pf) ∧
  nodeCorr cols ct i (CertPf sq nd) ∧
  FST (THE (tagOf cols i)) = 2 ⇒
  (checkFam_cut cols ct i ⇔ armCut ws sq nd)
\end{holsrc}
\ruleblock{HOL verifier receipts}
\badge{BadgeBlue}{documentary} Documentary citations. Family validation: \texttt{bern\_fam\_resid} built GREEN with all ten \texttt{famEq} arms validated, \texttt{hol-strict} clean; spine census: \texttt{bern\_joins} GREEN under \texttt{hol-strict} (55 theorems, 0 offences, 208-theory ancestry including this theory), \texttt{checkJoins\_eq\_checkPf} verified verbatim in the built signature. The cited \texttt{famEq} lemma is the kernel-checked per-node equivalence: under well-formed columns (\texttt{wfColumns}/\texttt{mkCtx}/\texttt{nodeCorr}) at this rule\textquotesingle s wire tag, the columnar family check coincides with the flat reference walk\textquotesingle s arm for this rule. There is no signed HOL receipt store yet --- these citations are documentary.
Semantic stratum, also ported: the soundness twin \texttt{sound\_cut}
(theory \texttt{bern\_tr\_soundstruct}: \texttt{hol-strict} clean, 9
theorems, 0 offences, verification run of 2026-07-22 on an isolated
build host, content-addressed at the same sources), consumed by the
certified-fragment keystone (\texttt{checkPf\_sound\_frag}, theory
\texttt{bern\_tr\_splice}, 20 theorems, 0 offences, same run): a
reference-checker accept inside the proven wire fragment yields sequent
validity.

\sourcenote{HOL receipts are documentary: quoted results of \texttt{hol-strict} runs
and census rows from the HOL4 development\textquotesingle s build
evidence, not minted receipt artefacts. A HOL receipt store (the Lean
store\textquotesingle s twin) is named future work.
}
\ruleblock{Proof}

\sourcenote{Genre: core rule (tier 1). Setting as in the \rulename{Ax}
section; counterexamples read over the one-wire hypernet.}

\textbf{Claim (\texttt{sound\_cut}).}\enspace The payload supplies the
cut formula $\chi$ (\texttt{cutFormula}), a left premise
$\Phi \vdash \chi$ and a right premise $\Phi \wedge \chi \vdash \Psi$
with their sub-derivations, and four tree identities: the conclusion's
premise is the left premise's premise (\texttt{preEq}), its conclusion
is the right premise's conclusion (\texttt{postEq}), the left premise
concludes exactly $\chi$ (\texttt{cutLeft}), and the right premise's
premise is exactly $\Phi \wedge \chi$ over the \emph{same} $\Phi$
(\texttt{cutRight}). If both premise sequents are semantically valid,
so is $\Phi \vdash \Psi$.

\begin{bthm}[soundness of \rulename{Cut}]\label{thm:cut-sound}
Under the four payload identities, semantic validity of both premise
sequents implies semantic validity of the conclusion, over every
hypernet.
\end{bthm}
\begin{proof}
Let $\sigma$ be admissible with $\sigma \models \Phi$.

\emph{Step 1 (earn the cut formula).} By \texttt{preEq}, $\sigma$
satisfies the left premise's premise; left validity gives
$\sigma \models \chi$ (\texttt{cutLeft} identifies the left conclusion
with $\chi$).

\emph{Step 2 (assemble the enlarged context).} The $\wedge$-clause of
satisfaction is meta-conjunction, so the standing assumption and Step 1
give $\sigma \models \Phi \wedge \chi$ --- and by \texttt{cutRight} this tree
\emph{is} the right premise's premise.

\emph{Step 3 (apply the right premise).} Right validity gives
$\sigma \models \Psi$, pinned as the conclusion's post by
\texttt{postEq}.
\end{proof}

\begin{bremark}[context sharing is load-bearing]\label{rem:cut-sharing}
\texttt{cutRight} forces the right premise's context to be the left
premise's context --- the context-sharing (\emph{additive}) format
\cite{girard1987}. Allow the right context to be foreign, say
$\Phi' \wedge \chi$, and cut breaks: take the left premise
$\top \vdash \top$ (an \rulename{Ax} instance) and the right premise
$(w \leq 0) \wedge \top \vdash w \leq 0$ (valid: satisfaction of the
conjunction projects); the degenerate rule would conclude
$\top \vdash w \leq 0$ --- false at $\sigma(w) = 1$.
\end{bremark}

\begin{bremark}[\texttt{cutLeft} moors the cut formula]
Keep the sharing but unmoor $\chi$ from what the left premise actually
proved, and the right premise may consume an assumption nobody earned:
from the left premise $\top \vdash \top$ and the right premise
$\top \wedge (w \leq 0) \vdash w \leq 0$ (valid, again by projection),
the degenerate rule would conclude $\top \vdash w \leq 0$. The two
counterexamples are deliberately near-twins: \texttt{cutRight} guards
\emph{whose context}, \texttt{cutLeft} guards \emph{whose lemma}, and
each fails independently. \texttt{preEq}'s removal admits the
context-swap family (the \rulename{Weaken} section); \texttt{postEq}'s
removal frees the conclusion outright.
\end{bremark}

\begin{bremark}[the cut formula leaves no trace]
The conclusion $\Phi \vdash \Psi$ mentions $\chi$ nowhere --- the
conclusion does not determine the certified content (a hiddenness
\rulename{Cut} shares with $\wedge$E's discarded conjunct), so the
payload names \texttt{cutFormula} explicitly. That is a design choice
rather than a forced move --- the checker could equally read $\chi$
off the left premise's post --- and the explicit field keeps the check
local. The same invisibility is what the rule is \emph{for}:
certificate producers stage lemmas --- prove a bound or region fact
once, then continue under the enlarged context $\Phi \wedge \chi$ ---
and iterated cuts stage a chain of them.
\end{bremark}

\begin{bremark}[semantic cut costs nothing]
Gentzen's Hauptsatz \cite{gentzen1935} eliminates cut from derivations
at, in general, non-elementary blow-up in proof size
\cite{statman1979,orevkov1979}; nothing here needs it. The
soundness argument is the three steps above, the checker never
normalizes a proof, and the calculus keeps \rulename{Cut} as a
first-class node because producers want staged lemmas, not cut-free
normal forms. In the lattice reading: if $x \leq m$ then $x = x \sqcap
m$, so $x \sqcap m \leq y$ gives $x \leq y$ --- transitivity worn as
meet bookkeeping.
\end{bremark}

\begin{anchors}
Steps 1--3 & \texttt{sound\_cut}; payload \texttt{PCutData} (fields \texttt{cutFormula}, \texttt{preEq}, \texttt{postEq}, \texttt{cutLeft}, \texttt{cutRight}); wired as \texttt{RuleData.cut}\\
$\wedge$-clause of satisfaction & \texttt{Formula.holds} (case \texttt{and})\\
lattice reading & \texttt{M3.cut\_kit}\\
\end{anchors}


\rulesection{4}{$\wedge$I}{core rule}
\label{blcatrule-AndI}
\ruleblock{The rule (Steele metanotation)}
\sourcenote{as realized in the mechanized development; source row: \texttt{\texttt{rules.toml}} --- \emph{$\wedge$I}. The metanotation statement is typeset by hand for notational quality; its mechanized ground truth is the Lean statement below, extracted from the sources at every build of this book and hash-bound. The rule \emph{roster} is machine-checked (every mechanized rule tag round-trips to exactly one catalogue row); the statement \emph{text} is editorial, checked by eye against the extracted Lean beside it.}
\begin{mathpar}
\inferrule*[Right=$\wedge$I]{\Phi \vdash \Psi_1 \\ \Phi \vdash \Psi_2}{\Phi \vdash \Psi_1 \wedge \Psi_2}
\end{mathpar}

\par\medskip
$\wedge$I pairs two conclusions proven from the \emph{same} assumptions:
from $\Phi \vdash \Psi_1$ and $\Phi \vdash \Psi_2$, conclude
$\Phi \vdash \Psi_1 \wedge \Psi_2$. A producer uses it to assemble
composite claims from separately established facts --- two one-sided
bounds into a two-sided one, the conjuncts a later
node\textquotesingle s premise-form expects --- and, with $\wedge$E, to
rebuild conjunction trees in the shape a structural comparison demands
(the core never reassociates a conjunction on its own).

The checker recomputes three tree identities: the two
sub-derivations\textquotesingle{} premises are the same tree, the
conclusion inherits that premise, and the conclusion\textquotesingle s
post is one new $\wedge$-node with the left sub-derivation\textquotesingle s
conclusion on the left. Order is data: $\Psi_1 \wedge \Psi_2$ and
$\Psi_2 \wedge \Psi_1$ are different conclusions, obtained by swapping
which sub-derivation rides which payload slot.

The context-sharing identity is the load-bearing one. Drop it and the
rule pairs facts earned under different assumptions, keeping only one
set: from the valid $\top \vdash \top$ and $w \leq 0 \vdash w \leq 0$
it would conclude $\top \vdash \top \wedge (w \leq 0)$, false wherever
$w$ is positive. The proof section also records the structural
observation the $\vee$E section turns into a theorem: $\wedge$I and
$\wedge$E are the two directions of the universal property of meet, so
the conjunction fragment adds nothing beyond the lattice structure of
entailment.

\ruleblock{Lean statement}
\declhead{BernsteinLogic.Soundness}{sound\_andI}{\badge{BadgeGreen}{statement bound}}
\begin{leansrc}
theorem sound_andI {H : Hypernet} {seq leftPremise rightPremise : Sequent}
    (ihl : leftPremise.semanticallyValid H)
    (ihr : rightPremise.semanticallyValid H)
    (hpreMatch : leftPremise.pre = rightPremise.pre)
    (hpre : seq.pre = leftPremise.pre)
    (hpost : seq.post = .and leftPremise.post rightPremise.post) :
    seq.semanticallyValid H
\end{leansrc}
\par\medskip
Payload, read as in the \rulename{Ax} section: sub-derivations plus
structural equalities, nothing else.

\declhead{BernsteinLogic.Core.Catalogue}{PAndIData}{\badge{BadgeGreen}{statement bound}}
\begin{leansrc}
inductive PAndIData : Hypernet → Sequent → Type where
  | mk (H : Hypernet) (seq : Sequent)
       (leftPremise  : Sequent) (leftProof  : ProofObject H leftPremise)
       (rightPremise : Sequent) (rightProof : ProofObject H rightPremise)
       (preMatch : leftPremise.pre = rightPremise.pre)
       (preEq    : seq.pre  = leftPremise.pre)
       (postEq   : seq.post = .and leftPremise.post rightPremise.post)
       : PAndIData H seq
\end{leansrc}
\ruleblock{Correspondence (Steele display $\leftrightarrow$ Lean)}
\badge{BadgeGreen}{reviewer-certified against the displayed texts}
\sourcenote{The dictionary below is a reviewer-certified statement of how the mechanized statements above interpret the displayed rule: display metavariable $\mapsto$ Lean binder, each display hypothesis $\mapsto$ the premise or side condition carrying it, the display conclusion $\mapsto$ the goal --- with any deliberate difference in width between display and mechanization named outright. Certification is recorded as a hash triple over (the displayed cell, the extracted statement texts, this dictionary\textquotesingle s own text); if any of the three changes, this block turns amber until a reviewer re-certifies. The correspondence claim itself remains editorial --- reviewer-grade, not machine-checked; its machine-checked successor types each soundness statement against the rule as a structured object, and supersedes this block per rule as it lands.}
Judgment frame as in \rulename{Ax}; sub-derivation premises as at
\rulename{Weaken}.
\begin{itemize}
\item the two premise judgments $\mapsto$ \texttt{ihl}/\texttt{ihr}
  over the binders \texttt{leftPremise}/\texttt{rightPremise}.
\item the display writes one context $\Phi$ twice; the statement binds
  two sequents and demands \texttt{hpreMatch} $:$
  \texttt{leftPremise.pre = rightPremise.pre} --- sharing by equation
  rather than by a common binder, the pattern of every two-premise
  rule on this page.
\item $\Phi \mapsto$ \texttt{leftPremise.pre} (hence, through the
  match, \texttt{rightPremise.pre}); $\Psi_1 \mapsto$
  \texttt{leftPremise.post}; $\Psi_2 \mapsto$
  \texttt{rightPremise.post}.
\item conclusion $\Phi \vdash \Psi_1 \wedge \Psi_2$ $\mapsto$ the goal
  under \texttt{hpre} $:$ \texttt{seq.pre = leftPremise.pre} and
  \texttt{hpost} $:$ \texttt{seq.post = .and leftPremise.post
  rightPremise.post} --- left sub-derivation on the left, as
  displayed.
\end{itemize}
No width difference in either direction; every hypothesis is consumed.

\ruleblock{Lean verifier receipts}
\sourcenote{module \texttt{BernsteinLogic.Soundness}; bound build \texttt{a20d90e66c13\ldots} (bound from receipt a9c0e34d3e3c9c8c, 2026-08-03T21:30:36Z). A receipt records that named checking layers accepted a build: (module, compiled-artifact SHA-256, environment SHA-256, checker layers), minted when the checking gate ran; the store receipts attesting the bound build are listed with the layers each run credited. Nothing here asks the reader to trust the store: every layer is an independently runnable program over the quoted sources --- rebuild the module and re-run the checkers to reproduce the verdict. Receipts attesting superseded builds credit no layer and are omitted from this rendering (the count of omissions, if any, is noted below).}
\begin{receipttable}
\receiptrow{a9c0e34d3e3c9c8c}{2026-08-03T21:30:36Z}{build+decorrelated}{\badge{BadgeGreen}{current}}
\end{receipttable}
\sourcenote{21 further receipt(s) in the store attest superseded builds of this module and are omitted from this rendering.}
\sourcenote{module \texttt{BernsteinLogic.Core.Catalogue}; bound build \texttt{cbc92e976a1d\ldots} (bound from receipt 5abdb3589707b5df, 2026-08-03T21:30:33Z). A receipt records that named checking layers accepted a build: (module, compiled-artifact SHA-256, environment SHA-256, checker layers), minted when the checking gate ran; the store receipts attesting the bound build are listed with the layers each run credited. Nothing here asks the reader to trust the store: every layer is an independently runnable program over the quoted sources --- rebuild the module and re-run the checkers to reproduce the verdict. Receipts attesting superseded builds credit no layer and are omitted from this rendering (the count of omissions, if any, is noted below).}
\begin{receipttable}
\receiptrow{24804e66edba0bc9}{2026-08-03T20:25:09Z}{build+decorrelated}{\badge{BadgeGreen}{current}}
\receiptrow{5abdb3589707b5df}{2026-08-03T21:30:33Z}{build+decorrelated}{\badge{BadgeGreen}{current}}
\end{receipttable}
\sourcenote{18 further receipt(s) in the store attest superseded builds of this module and are omitted from this rendering.}
\layerline{\layeritem{build}{lake build (elaboration)}{\badge{BadgeGreen}{current}}\layeritem{replay}{leanchecker (official kernel replay)}{\badge{BadgeAmber}{stale only}}\layeritem{hardened}{safe\_verify (hardened screen)}{\badge{BadgeAmber}{stale only}}\layeritem{decorrelated}{lean4export + nanoda (decorrelated kernel)}{\badge{BadgeGreen}{current}}}
\sourcenote{\badge{BadgeGray}{comparator: none} no fixture in the three-implementation comparator harness (compiled / kernel-interpreted / reference-interpreted) exercises this rule.}
\ruleblock{HOL4 statement}
\sourcenote{theory \texttt{bern\_fam\_resid}, HOL4 development script \texttt{\texttt{bern\_fam\_residScript.sml}} \badge{BadgeGreen}{script bound}}
\holhead{Definition checkFam\_andi\_def}
\begin{holsrc}
Definition checkFam_andi_def:
  checkFam_andi (cols:columns) (ct:ctx) (i:num) ⇔
    case frSeq cols ct i of
    | NONE => F
    | SOME sq =>
      case frChildSeq cols ct i 0 of
      | NONE => F
      | SOME ls =>
        case frChildSeq cols ct i 1 of
        | NONE => F
        | SOME rs =>
            ls.pre = rs.pre ∧ sq.pre = ls.pre ∧
            sq.post = FAnd ls.post rs.post
End
\end{holsrc}
\holhead{Theorem famEq\_andi}
\begin{holsrc}
Theorem famEq_andi:
  wfColumns cols ct ∧ mkCtx cols = SOME ct ∧
  parseCert (cols.hdrImage ++ rest) = SOME (ws, pf) ∧
  nodeCorr cols ct i (CertPf sq nd) ∧
  FST (THE (tagOf cols i)) = 3 ⇒
  (checkFam_andi cols ct i ⇔ armAndI ws sq nd)
\end{holsrc}
\ruleblock{HOL verifier receipts}
\badge{BadgeBlue}{documentary} Documentary citations. Family validation: \texttt{bern\_fam\_resid} built GREEN with all ten \texttt{famEq} arms validated, \texttt{hol-strict} clean; spine census: \texttt{bern\_joins} GREEN under \texttt{hol-strict} (55 theorems, 0 offences, 208-theory ancestry including this theory), \texttt{checkJoins\_eq\_checkPf} verified verbatim in the built signature. The cited \texttt{famEq} lemma is the kernel-checked per-node equivalence: under well-formed columns (\texttt{wfColumns}/\texttt{mkCtx}/\texttt{nodeCorr}) at this rule\textquotesingle s wire tag, the columnar family check coincides with the flat reference walk\textquotesingle s arm for this rule. There is no signed HOL receipt store yet --- these citations are documentary.
Semantic stratum, also ported: the soundness twin \texttt{sound\_andI}
(theory \texttt{bern\_tr\_soundstruct}: \texttt{hol-strict} clean, 9
theorems, 0 offences, verification run of 2026-07-22 on an isolated
build host, content-addressed at the same sources), consumed by the
certified-fragment keystone (\texttt{checkPf\_sound\_frag}, theory
\texttt{bern\_tr\_splice}, 20 theorems, 0 offences, same run): a
reference-checker accept inside the proven wire fragment yields sequent
validity.

\sourcenote{HOL receipts are documentary: quoted results of \texttt{hol-strict} runs
and census rows from the HOL4 development\textquotesingle s build
evidence, not minted receipt artefacts. A HOL receipt store (the Lean
store\textquotesingle s twin) is named future work.
}
\ruleblock{Proof}

\sourcenote{Genre: core rule (tier 1). Setting as in the \rulename{Ax}
section; counterexamples read over the one-wire hypernet.}

\textbf{Claim (\texttt{sound\_andI}).}\enspace The payload supplies two
premise sequents $\Phi \vdash \Psi_1$ and $\Phi \vdash \Psi_2$ with
their sub-derivations and three tree identities: the two premises share
one context (\texttt{preMatch}), the conclusion's premise is that
context (\texttt{preEq}), and the conclusion's post is
$\Psi_1 \wedge \Psi_2$ with the left premise's conclusion as left
conjunct (\texttt{postEq}). If both premise sequents are semantically
valid, so is $\Phi \vdash \Psi_1 \wedge \Psi_2$.

\begin{bthm}[soundness of $\wedge$I]\label{thm:andi-sound}
Under the three payload identities, semantic validity of both premise
sequents implies semantic validity of the conclusion, over every
hypernet.
\end{bthm}
\begin{proof}
Let $\sigma$ be admissible with $\sigma \models \Phi$.

\emph{Step 1 (left component).} By \texttt{preEq}, $\sigma$ satisfies
the left premise's premise; left validity gives $\sigma \models \Psi_1$.

\emph{Step 2 (right component).} \texttt{preMatch} transports
$\sigma \models \Phi$ verbatim to the right premise's premise --- the
same tree, so the same satisfaction fact; right validity gives
$\sigma \models \Psi_2$.

\emph{Step 3 (pair).} The $\wedge$-clause of satisfaction is
meta-conjunction, so $\sigma \models \Psi_1 \wedge \Psi_2$, pinned by
\texttt{postEq}.
\end{proof}

\begin{bremark}[\texttt{preMatch} is load-bearing]\label{rem:andi-share}
The rule may only pair conclusions drawn from the \emph{same}
assumptions --- context sharing again, as at \rulename{Cut}
\cite{girard1987}. Drop \texttt{preMatch} and take the left premise
$\top \vdash \top$ and the right premise $w \leq 0 \vdash w \leq 0$
(both \rulename{Ax} instances): the degenerate rule concludes
$\top \vdash \top \wedge (w \leq 0)$ --- false at $\sigma(w) = 1$,
because the right conjunct was earned under an assumption the
conclusion no longer carries. \texttt{preEq}'s removal admits the
context-swap family (the \rulename{Weaken} section); \texttt{postEq}'s
removal frees the conclusion outright.
\end{bremark}

\begin{bremark}[order and association are data]
\texttt{postEq} places the left premise's conclusion on the left of one
new $\wedge$-node. The conjunction is an ordered tree: nothing in the
core commutes or reassociates it, so $\Psi_1 \wedge \Psi_2$ and
$\Psi_2 \wedge \Psi_1$ are different conclusions, each derivable from
the same premises by choosing which sub-derivation rides which payload
slot.
\end{bremark}

\begin{bremark}[the meet adjunction, split across two rules]
\label{rem:andi-meet}
In the lattice reading, $\wedge$I is one direction of the universal
property of meet --- $x \leq y$ and $x \leq z$ iff
$x \leq y \sqcap z$ --- and $\wedge$E is the other. Together the two
rules say the $\wedge$-fragment of the calculus adds nothing beyond the
lattice structure of entailment; the $\vee$E section makes this precise
(and load-bearing) through a countermodel over exactly that structure.
With $\wedge$E the pair is also computationally lossless at the
formula-tree level: the two projections undo pairing, and pairing the
two projections rebuilds the identical formula tree.
\end{bremark}

\begin{anchors}
Steps 1--3 & \texttt{sound\_andI}; payload \texttt{PAndIData} (fields \texttt{preMatch}, \texttt{preEq}, \texttt{postEq}); wired as \texttt{RuleData.andI}\\
$\wedge$-clause of satisfaction & \texttt{Formula.holds} (case \texttt{and})\\
lattice reading & \texttt{M3.andI\_kit}\\
\end{anchors}


\rulesection{5}{$\wedge$E}{core rule}
\label{blcatrule-AndE}
\ruleblock{The rule (Steele metanotation)}
\sourcenote{as realized in the mechanized development; source row: \texttt{\texttt{rules.toml}} --- \emph{$\wedge$E}. The metanotation statement is typeset by hand for notational quality; its mechanized ground truth is the Lean statement below, extracted from the sources at every build of this book and hash-bound. The rule \emph{roster} is machine-checked (every mechanized rule tag round-trips to exactly one catalogue row); the statement \emph{text} is editorial, checked by eye against the extracted Lean beside it.}
\begin{mathpar}
\inferrule*[Right=$\wedge$E]{\Phi \vdash \Psi_1 \wedge \Psi_2}{\Phi \vdash \Psi_i}
\end{mathpar}

\par\medskip
$\wedge$E projects a conjunction: from
$\Phi \vdash \Psi_1 \wedge \Psi_2$, conclude $\Phi \vdash \Psi_i$. A
producer uses it to take composite facts apart --- the demonstration
certificate applies it right after its axiom leaf, extracting from a
composite premise the region conjunct that the Bernstein-cone closure
needs --- and, with $\wedge$I, to reshape conjunction trees for the
structural comparisons other rules perform.

The rule is realized as two constructors, one per side; the
paper\textquotesingle s index $i \in \{1, 2\}$ is the constructor
choice itself, so there is no unchecked side flag. The checker
recomputes two tree identities: shared context, and that the
sub-derivation\textquotesingle s conclusion is exactly the
conclusion\textquotesingle s post conjoined with the \emph{named}
discarded conjunct --- on the right for the left constructor, on the
left for the right one. Naming the discarded conjunct in the payload
makes the side condition a single comparison; the checker never
searches a conjunction for the target, and a certificate claiming the
wrong side simply fails the identity.

Dropping the conclusion identity frees the conclusion outright;
dropping the context identity admits the context-swap counterexamples
walked at \rulename{Weaken}. The projection itself needs no defence:
discarding a conjunct only weakens a claim, and the proof section
notes the pairing with $\wedge$I under which the fragment is lossless
at the formula-tree level.

\ruleblock{Lean statement}
\declhead{BernsteinLogic.Soundness}{sound\_andE\_left}{\badge{BadgeGreen}{statement bound}}
\begin{leansrc}
theorem sound_andE_left {H : Hypernet} {seq premise : Sequent} {other : Formula}
    (ih : premise.semanticallyValid H)
    (hpre : seq.pre = premise.pre)
    (hpost : premise.post = .and seq.post other) :
    seq.semanticallyValid H
\end{leansrc}
\declhead{BernsteinLogic.Soundness}{sound\_andE\_right}{\badge{BadgeGreen}{statement bound}}
\begin{leansrc}
theorem sound_andE_right {H : Hypernet} {seq premise : Sequent} {other : Formula}
    (ih : premise.semanticallyValid H)
    (hpre : seq.pre = premise.pre)
    (hpost : premise.post = .and other seq.post) :
    seq.semanticallyValid H
\end{leansrc}
\par\medskip
Payload, read as in the \rulename{Ax} section: one constructor per
projection, each a sub-derivation plus structural equalities with the
discarded conjunct as explicit data.

\declhead{BernsteinLogic.Core.Catalogue}{PAndEData}{\badge{BadgeGreen}{statement bound}}
\begin{leansrc}
inductive PAndEData : Hypernet → Sequent → Type where
  | left  (H : Hypernet) (seq : Sequent)
          (premise : Sequent) (premiseProof : ProofObject H premise)
          (other : Formula)
          (preEq  : seq.pre = premise.pre)
          (postEq : premise.post = .and seq.post other)
          : PAndEData H seq
  | right (H : Hypernet) (seq : Sequent)
          (premise : Sequent) (premiseProof : ProofObject H premise)
          (other : Formula)
          (preEq  : seq.pre = premise.pre)
          (postEq : premise.post = .and other seq.post)
          : PAndEData H seq
\end{leansrc}
\ruleblock{Correspondence (Steele display $\leftrightarrow$ Lean)}
\badge{BadgeGreen}{reviewer-certified against the displayed texts}
\sourcenote{The dictionary below is a reviewer-certified statement of how the mechanized statements above interpret the displayed rule: display metavariable $\mapsto$ Lean binder, each display hypothesis $\mapsto$ the premise or side condition carrying it, the display conclusion $\mapsto$ the goal --- with any deliberate difference in width between display and mechanization named outright. Certification is recorded as a hash triple over (the displayed cell, the extracted statement texts, this dictionary\textquotesingle s own text); if any of the three changes, this block turns amber until a reviewer re-certifies. The correspondence claim itself remains editorial --- reviewer-grade, not machine-checked; its machine-checked successor types each soundness statement against the rule as a structured object, and supersedes this block per rule as it lands.}
Judgment frame as in \rulename{Ax}; sub-derivation premise as at
\rulename{Weaken}.
\begin{itemize}
\item the display\textquotesingle s index $i \in \{1, 2\}$ $\mapsto$
  the \emph{choice of lemma}: \texttt{sound\_andE\_left} is $i = 1$,
  \texttt{sound\_andE\_right} is $i = 2$, mirroring the
  payload\textquotesingle s two constructors --- the side is a
  constructor choice, never an unchecked flag.
\item premise judgment $\Phi \vdash \Psi_1 \wedge \Psi_2$ $\mapsto$
  \texttt{ih}, with the conjunction shape carried by \texttt{hpost}:
  in the left lemma \texttt{premise.post = .and seq.post other} ---
  the surviving conjunct $\Psi_i$ \emph{is} the
  conclusion\textquotesingle s own post field, and the discarded
  conjunct $\Psi_{3-i}$ $\mapsto$ the binder \texttt{other} (right
  lemma: mirrored, \texttt{.and other seq.post}).
\item $\Phi \mapsto$ \texttt{premise.pre}, shared via \texttt{hpre}
  $:$ \texttt{seq.pre = premise.pre}; conclusion $\Phi \vdash \Psi_i$
  $\mapsto$ the goal \texttt{seq.semanticallyValid H}.
\end{itemize}
No width difference in either direction: the lemma pair covers exactly
the displayed $i$-indexed family, and every hypothesis is consumed.

\ruleblock{Lean verifier receipts}
\sourcenote{module \texttt{BernsteinLogic.Soundness}; bound build \texttt{a20d90e66c13\ldots} (bound from receipt a9c0e34d3e3c9c8c, 2026-08-03T21:30:36Z). A receipt records that named checking layers accepted a build: (module, compiled-artifact SHA-256, environment SHA-256, checker layers), minted when the checking gate ran; the store receipts attesting the bound build are listed with the layers each run credited. Nothing here asks the reader to trust the store: every layer is an independently runnable program over the quoted sources --- rebuild the module and re-run the checkers to reproduce the verdict. Receipts attesting superseded builds credit no layer and are omitted from this rendering (the count of omissions, if any, is noted below).}
\begin{receipttable}
\receiptrow{a9c0e34d3e3c9c8c}{2026-08-03T21:30:36Z}{build+decorrelated}{\badge{BadgeGreen}{current}}
\end{receipttable}
\sourcenote{21 further receipt(s) in the store attest superseded builds of this module and are omitted from this rendering.}
\sourcenote{module \texttt{BernsteinLogic.Core.Catalogue}; bound build \texttt{cbc92e976a1d\ldots} (bound from receipt 5abdb3589707b5df, 2026-08-03T21:30:33Z). A receipt records that named checking layers accepted a build: (module, compiled-artifact SHA-256, environment SHA-256, checker layers), minted when the checking gate ran; the store receipts attesting the bound build are listed with the layers each run credited. Nothing here asks the reader to trust the store: every layer is an independently runnable program over the quoted sources --- rebuild the module and re-run the checkers to reproduce the verdict. Receipts attesting superseded builds credit no layer and are omitted from this rendering (the count of omissions, if any, is noted below).}
\begin{receipttable}
\receiptrow{24804e66edba0bc9}{2026-08-03T20:25:09Z}{build+decorrelated}{\badge{BadgeGreen}{current}}
\receiptrow{5abdb3589707b5df}{2026-08-03T21:30:33Z}{build+decorrelated}{\badge{BadgeGreen}{current}}
\end{receipttable}
\sourcenote{18 further receipt(s) in the store attest superseded builds of this module and are omitted from this rendering.}
\layerline{\layeritem{build}{lake build (elaboration)}{\badge{BadgeGreen}{current}}\layeritem{replay}{leanchecker (official kernel replay)}{\badge{BadgeAmber}{stale only}}\layeritem{hardened}{safe\_verify (hardened screen)}{\badge{BadgeAmber}{stale only}}\layeritem{decorrelated}{lean4export + nanoda (decorrelated kernel)}{\badge{BadgeGreen}{current}}}
\sourcenote{\badge{BadgeBlue}{comparator: exercised via the demonstration fixture}
the composite demonstration certificate (an axiom node, a conjunction
elimination, a Bernstein-cone closure) runs as a fixture through the
three-implementation comparator harness (compiled / kernel-interpreted /
reference-interpreted), so this rule\textquotesingle s node rides every
execution path of that harness.
}
\ruleblock{HOL4 statement}
\sourcenote{theory \texttt{bern\_fam\_ande}, HOL4 development script \texttt{\texttt{bern\_fam\_andeScript.sml}} \badge{BadgeGreen}{script bound}}
\holhead{Definition checkFam\_ande\_def}
\begin{holsrc}
Definition checkFam_ande_def:
  checkFam_ande (cols:columns) (ct:ctx) (i:num) ⇔
    case tagOf cols i of
    | NONE => F
    | SOME (r, v) =>
        r = 4 ∧
        (case childOrd ct i 0 of
         | NONE => F
         | SOME k =>
           case conclRef cols ct i of
           | NONE => F
           | SOME cj =>
           case conclRef cols ct k of
           | NONE => F
           | SOME ccj =>
           case pv cols ct cj of
           | SOME (PJudg prew postw) =>
             (case pv cols ct ccj of
              | SOME (PJudg cprew cpostw) =>
                (case prefField cols ct i 0 of
                 | NONE => F
                 | SOME ot =>
                   case pv cols ct ot of
                   | SOME (PFormula otw) =>
                       (* pre-chain, expanded level (design note N1) *)
                       expandSph FTt prew = expandSph FTt cprew ∧
                       (* child post = f.and; conclusion = taken side,
                          other = untaken side, by variant *)
                       expandSph FTt cpostw =
                         (if v = 0 then
                            FAnd (expandSph FTt postw) (expandSph FTt otw)
                          else
                            FAnd (expandSph FTt otw) (expandSph FTt postw))
                   | _ => F)
              | _ => F)
           | _ => F)
End
\end{holsrc}
\holhead{Theorem famEq\_ande}
\begin{holsrc}
Theorem famEq_ande:
  ∀cols ct i sq nd ws.
    wfColumns cols ct ∧ mkCtx cols = SOME ct ∧
    nodeCorr cols ct i (CertPf sq nd) ∧
    FST (THE (tagOf cols i)) ∈ {4} ⇒
    (checkFam_ande cols ct i ⇔ checkPf1_ande_arm ws sq nd)
\end{holsrc}
\ruleblock{HOL verifier receipts}
\badge{BadgeBlue}{documentary} Documentary citations. Spine census: \texttt{bern\_joins} GREEN under \texttt{hol-strict} (55 theorems, 0 offences, 208-theory ancestry including this theory), \texttt{checkJoins\_eq\_checkPf} verified verbatim in the built signature. The cited \texttt{famEq} lemma is the kernel-checked per-node equivalence: under well-formed columns (\texttt{wfColumns}/\texttt{mkCtx}/\texttt{nodeCorr}) at this rule\textquotesingle s wire tag, the columnar family check coincides with the flat reference walk\textquotesingle s arm for this rule. There is no signed HOL receipt store yet --- these citations are documentary.
Semantic stratum, also ported: the soundness twins
\texttt{sound\_andE\_left} / \texttt{sound\_andE\_right} (theory
\texttt{bern\_tr\_soundstruct}: \texttt{hol-strict} clean, 9 theorems, 0
offences, verification run of 2026-07-22 on an isolated build host,
content-addressed at the same sources), consumed by the certified-
fragment keystone (\texttt{checkPf\_sound\_frag}, theory
\texttt{bern\_tr\_splice}, 20 theorems, 0 offences, same run): a
reference-checker accept inside the proven wire fragment yields sequent
validity.

\sourcenote{HOL receipts are documentary: quoted results of \texttt{hol-strict} runs
and census rows from the HOL4 development\textquotesingle s build
evidence, not minted receipt artefacts. A HOL receipt store (the Lean
store\textquotesingle s twin) is named future work.
}
\ruleblock{Proof}

\sourcenote{Genre: core rule (tier 1). Setting as in the \rulename{Ax}
section; counterexamples read over the one-wire hypernet.}

\textbf{Claim (\texttt{sound\_andE\_left} /
\texttt{sound\_andE\_right}).}\enspace The rule is realized as two
constructors, one per projection; the payload of either supplies a
premise sequent with its sub-derivation, the \emph{discarded} conjunct
(\texttt{other}), and two tree identities: shared context
(\texttt{preEq}) and, for the left constructor, that the premise's
conclusion is $\mathrm{post} \wedge \mathrm{other}$ --- for the right,
$\mathrm{other} \wedge \mathrm{post}$ (\texttt{postEq}). If the premise
sequent $\Phi \vdash \Psi_1 \wedge \Psi_2$ is semantically valid, so is
$\Phi \vdash \Psi_i$.

\begin{bthm}[soundness of $\wedge$E]\label{thm:ande-sound}
Under the payload identities of either constructor, semantic validity of
the premise sequent implies semantic validity of the conclusion, over
every hypernet.
\end{bthm}
\begin{proof}
Left constructor; the right is its mirror image.

\emph{Step 1 (apply the premise).} Let $\sigma$ be admissible with
$\sigma \models \Phi$; by \texttt{preEq} and premise validity,
$\sigma$ satisfies the premise's conclusion.

\emph{Step 2 (read it as a pair and project).} \texttt{postEq} says that
conclusion \emph{is} the tree
$\mathrm{seq.post} \wedge \mathrm{other}$; the $\wedge$-clause of
satisfaction is meta-conjunction, so $\sigma$ satisfies both conjuncts.
Keep the left one: $\sigma \models \mathrm{seq.post}$.
\end{proof}

\begin{bremark}[the discarded conjunct is named, so nothing is searched]
\label{rem:ande-other}
Carrying \texttt{other} in the payload turns the side condition into one
structural identity: the checker never hunts through the premise's
conclusion for an occurrence of the target --- it compares two named
trees. Orientation is checked by the same identity: a left-constructor
certificate claiming the \emph{right} conjunct fails \texttt{postEq}
unless the two conjuncts happen to be the same tree, in which case the
distinction is empty anyway.
\end{bremark}

\begin{bremark}[the index is a constructor, not a flag]
The paper's $i \in \{1, 2\}$ is realized as the choice between two
constructors rather than an integer field --- there is no unchecked
side flag to corrupt, and each constructor's soundness lemma is its own
two-line proof. Dropping \texttt{postEq} in either frees the conclusion
outright; dropping \texttt{preEq} admits the context-swap family (the
\rulename{Weaken} section).
\end{bremark}

\begin{bremark}[one-way information loss, two-way fragment]
Projection discards its \texttt{other} irrecoverably --- the conclusion
carries no trace of it (compare \rulename{Cut}, whose conclusion hides
the cut formula; here the forgetting is the visible point of the rule).
Together with $\wedge$I the fragment is lossless at the formula-tree
level: see the meet-adjunction remark in the $\wedge$I section. In the lattice reading these are
\texttt{M3.andE\_l\_kit}/\texttt{M3.andE\_r\_kit}, the upper direction
of meet's universal property.
\end{bremark}

\begin{anchors}
Steps 1--2 (left) & \texttt{sound\_andE\_left}; payload \texttt{PAndEData.left} (fields \texttt{other}, \texttt{preEq}, \texttt{postEq})\\
mirror (right) & \texttt{sound\_andE\_right}; payload \texttt{PAndEData.right}\\
dispatch & wired as \texttt{RuleData.andE} into \texttt{ProofObject.sound}\\
lattice reading & \texttt{M3.andE\_l\_kit}, \texttt{M3.andE\_r\_kit}\\
\end{anchors}


\rulesection{6}{$\vee$I}{core rule}
\label{blcatrule-OrI}
\ruleblock{The rule (Steele metanotation)}
\sourcenote{as realized in the mechanized development; source row: \texttt{\texttt{rules.toml}} --- \emph{$\vee$I}. The metanotation statement is typeset by hand for notational quality; its mechanized ground truth is the Lean statement below, extracted from the sources at every build of this book and hash-bound. The rule \emph{roster} is machine-checked (every mechanized rule tag round-trips to exactly one catalogue row); the statement \emph{text} is editorial, checked by eye against the extracted Lean beside it.}
\begin{mathpar}
\inferrule*[Right=$\vee$I]{\Phi \vdash \Psi_i}{\Phi \vdash \Psi_1 \vee \Psi_2}
\end{mathpar}

\par\medskip
$\vee$I injects a proven conclusion into a disjunction: from
$\Phi \vdash \Psi_i$, conclude $\Phi \vdash \Psi_1 \vee \Psi_2$ ---
conclusion-side weakening, the mirror of \rulename{Weaken}\textquotesingle s
premise side. A producer mostly reaches for it to make the branches of
a disjunction elimination agree: each branch proves its own
branch-specific fact and $\vee$I-injects it into the common disjunction
the elimination\textquotesingle s conclusion-matching check demands.
The separation instance discussed at $\vee$E is the miniature --- both
of its branch derivations are an axiom node followed by one $\vee$I
node.

Like $\wedge$E, the rule is realized as two constructors, one per
side; the injected side is the constructor choice. The checker
recomputes two tree identities: shared context, and that the
conclusion\textquotesingle s post is one new $\vee$-node with the
sub-derivation\textquotesingle s conclusion on the constructor\textquotesingle s
side and the \emph{named} added disjunct on the other.

Nothing whatsoever is demanded of the added disjunct --- it may be
unsatisfiable, ill-typed, or mention untyped wires --- because
satisfaction of a disjunction demands only the witnessed side, and the
satisfaction relation carries no typing side conditions of its own;
typing enters the calculus only through rules that must evaluate.
Dropping the conclusion identity frees the conclusion outright;
dropping the context identity admits the context-swap family walked at
\rulename{Weaken}.

\ruleblock{Lean statement}
\declhead{BernsteinLogic.Soundness}{sound\_orI\_left}{\badge{BadgeGreen}{statement bound}}
\begin{leansrc}
theorem sound_orI_left {H : Hypernet} {seq premise : Sequent} {other : Formula}
    (ih : premise.semanticallyValid H)
    (hpre : seq.pre = premise.pre)
    (hpost : seq.post = .or premise.post other) :
    seq.semanticallyValid H
\end{leansrc}
\declhead{BernsteinLogic.Soundness}{sound\_orI\_right}{\badge{BadgeGreen}{statement bound}}
\begin{leansrc}
theorem sound_orI_right {H : Hypernet} {seq premise : Sequent} {other : Formula}
    (ih : premise.semanticallyValid H)
    (hpre : seq.pre = premise.pre)
    (hpost : seq.post = .or other premise.post) :
    seq.semanticallyValid H
\end{leansrc}
\par\medskip
Payload, read as in the \rulename{Ax} section: one constructor per
injection, each a sub-derivation plus structural equalities with the
added disjunct as explicit data.

\declhead{BernsteinLogic.Core.Catalogue}{POrIData}{\badge{BadgeGreen}{statement bound}}
\begin{leansrc}
inductive POrIData : Hypernet → Sequent → Type where
  | left  (H : Hypernet) (seq : Sequent)
          (premise : Sequent) (premiseProof : ProofObject H premise)
          (other : Formula)
          (preEq  : seq.pre = premise.pre)
          (postEq : seq.post = .or premise.post other)
          : POrIData H seq
  | right (H : Hypernet) (seq : Sequent)
          (premise : Sequent) (premiseProof : ProofObject H premise)
          (other : Formula)
          (preEq  : seq.pre = premise.pre)
          (postEq : seq.post = .or other premise.post)
          : POrIData H seq
\end{leansrc}
\ruleblock{Correspondence (Steele display $\leftrightarrow$ Lean)}
\badge{BadgeGreen}{reviewer-certified against the displayed texts}
\sourcenote{The dictionary below is a reviewer-certified statement of how the mechanized statements above interpret the displayed rule: display metavariable $\mapsto$ Lean binder, each display hypothesis $\mapsto$ the premise or side condition carrying it, the display conclusion $\mapsto$ the goal --- with any deliberate difference in width between display and mechanization named outright. Certification is recorded as a hash triple over (the displayed cell, the extracted statement texts, this dictionary\textquotesingle s own text); if any of the three changes, this block turns amber until a reviewer re-certifies. The correspondence claim itself remains editorial --- reviewer-grade, not machine-checked; its machine-checked successor types each soundness statement against the rule as a structured object, and supersedes this block per rule as it lands.}
Judgment frame as in \rulename{Ax}; sub-derivation premise as at
\rulename{Weaken}.
\begin{itemize}
\item the display\textquotesingle s index $i \in \{1, 2\}$ $\mapsto$
  the choice of lemma, exactly as at $\wedge$E:
  \texttt{sound\_orI\_left} is $i = 1$, \texttt{sound\_orI\_right} is
  $i = 2$, mirroring the payload\textquotesingle s two constructors.
\item premise judgment $\Phi \vdash \Psi_i$ $\mapsto$ \texttt{ih};
  the injected disjunct $\Psi_i$ \emph{is} \texttt{premise.post}, and
  the added disjunct $\Psi_{3-i}$ $\mapsto$ the binder \texttt{other},
  of which nothing is demanded.
\item conclusion $\Phi \vdash \Psi_1 \vee \Psi_2$ $\mapsto$ the goal
  under \texttt{hpost}: in the left lemma \texttt{seq.post = .or
  premise.post other} (right lemma mirrored); $\Phi \mapsto$
  \texttt{premise.pre}, shared via \texttt{hpre}.
\end{itemize}
No width difference in either direction: the lemma pair covers exactly
the displayed $i$-indexed family, and every hypothesis is consumed.

\ruleblock{Lean verifier receipts}
\sourcenote{module \texttt{BernsteinLogic.Soundness}; bound build \texttt{a20d90e66c13\ldots} (bound from receipt a9c0e34d3e3c9c8c, 2026-08-03T21:30:36Z). A receipt records that named checking layers accepted a build: (module, compiled-artifact SHA-256, environment SHA-256, checker layers), minted when the checking gate ran; the store receipts attesting the bound build are listed with the layers each run credited. Nothing here asks the reader to trust the store: every layer is an independently runnable program over the quoted sources --- rebuild the module and re-run the checkers to reproduce the verdict. Receipts attesting superseded builds credit no layer and are omitted from this rendering (the count of omissions, if any, is noted below).}
\begin{receipttable}
\receiptrow{a9c0e34d3e3c9c8c}{2026-08-03T21:30:36Z}{build+decorrelated}{\badge{BadgeGreen}{current}}
\end{receipttable}
\sourcenote{21 further receipt(s) in the store attest superseded builds of this module and are omitted from this rendering.}
\sourcenote{module \texttt{BernsteinLogic.Core.Catalogue}; bound build \texttt{cbc92e976a1d\ldots} (bound from receipt 5abdb3589707b5df, 2026-08-03T21:30:33Z). A receipt records that named checking layers accepted a build: (module, compiled-artifact SHA-256, environment SHA-256, checker layers), minted when the checking gate ran; the store receipts attesting the bound build are listed with the layers each run credited. Nothing here asks the reader to trust the store: every layer is an independently runnable program over the quoted sources --- rebuild the module and re-run the checkers to reproduce the verdict. Receipts attesting superseded builds credit no layer and are omitted from this rendering (the count of omissions, if any, is noted below).}
\begin{receipttable}
\receiptrow{24804e66edba0bc9}{2026-08-03T20:25:09Z}{build+decorrelated}{\badge{BadgeGreen}{current}}
\receiptrow{5abdb3589707b5df}{2026-08-03T21:30:33Z}{build+decorrelated}{\badge{BadgeGreen}{current}}
\end{receipttable}
\sourcenote{18 further receipt(s) in the store attest superseded builds of this module and are omitted from this rendering.}
\layerline{\layeritem{build}{lake build (elaboration)}{\badge{BadgeGreen}{current}}\layeritem{replay}{leanchecker (official kernel replay)}{\badge{BadgeAmber}{stale only}}\layeritem{hardened}{safe\_verify (hardened screen)}{\badge{BadgeAmber}{stale only}}\layeritem{decorrelated}{lean4export + nanoda (decorrelated kernel)}{\badge{BadgeGreen}{current}}}
\sourcenote{\badge{BadgeGray}{comparator: none} no fixture in the three-implementation comparator harness (compiled / kernel-interpreted / reference-interpreted) exercises this rule.}
\ruleblock{HOL4 statement}
\sourcenote{theory \texttt{bern\_fam\_resid}, HOL4 development script \texttt{\texttt{bern\_fam\_residScript.sml}} \badge{BadgeGreen}{script bound}}
\holhead{Definition checkFam\_ori\_def}
\begin{holsrc}
Definition checkFam_ori_def:
  checkFam_ori (cols:columns) (ct:ctx) (i:num) ⇔
    case tagOf cols i of
    | NONE => F
    | SOME (r, v) =>
      case frSeq cols ct i of
      | NONE => F
      | SOME sq =>
        case frChildSeq cols ct i 0 of
        | NONE => F
        | SOME ps =>
          case frFormula cols ct i 0 of
          | NONE => F
          | SOME other =>
              sq.pre = ps.pre ∧
              sq.post = (if v = 0 then FOr ps.post other
                         else FOr other ps.post)
End
\end{holsrc}
\holhead{Theorem famEq\_ori}
\begin{holsrc}
Theorem famEq_ori:
  wfColumns cols ct ∧ mkCtx cols = SOME ct ∧
  parseCert (cols.hdrImage ++ rest) = SOME (ws, pf) ∧
  nodeCorr cols ct i (CertPf sq nd) ∧
  FST (THE (tagOf cols i)) = 5 ⇒
  (checkFam_ori cols ct i ⇔ armOrI ws sq nd)
\end{holsrc}
\ruleblock{HOL verifier receipts}
\badge{BadgeBlue}{documentary} Documentary citations. Family validation: \texttt{bern\_fam\_resid} built GREEN with all ten \texttt{famEq} arms validated, \texttt{hol-strict} clean; spine census: \texttt{bern\_joins} GREEN under \texttt{hol-strict} (55 theorems, 0 offences, 208-theory ancestry including this theory), \texttt{checkJoins\_eq\_checkPf} verified verbatim in the built signature. The cited \texttt{famEq} lemma is the kernel-checked per-node equivalence: under well-formed columns (\texttt{wfColumns}/\texttt{mkCtx}/\texttt{nodeCorr}) at this rule\textquotesingle s wire tag, the columnar family check coincides with the flat reference walk\textquotesingle s arm for this rule. There is no signed HOL receipt store yet --- these citations are documentary.
Semantic stratum, also ported: the soundness twins
\texttt{sound\_orI\_left} / \texttt{sound\_orI\_right} (theory
\texttt{bern\_tr\_soundstruct}: \texttt{hol-strict} clean, 9 theorems, 0
offences, verification run of 2026-07-22 on an isolated build host,
content-addressed at the same sources), consumed by the certified-
fragment keystone (\texttt{checkPf\_sound\_frag}, theory
\texttt{bern\_tr\_splice}, 20 theorems, 0 offences, same run): a
reference-checker accept inside the proven wire fragment yields sequent
validity.

\sourcenote{HOL receipts are documentary: quoted results of \texttt{hol-strict} runs
and census rows from the HOL4 development\textquotesingle s build
evidence, not minted receipt artefacts. A HOL receipt store (the Lean
store\textquotesingle s twin) is named future work.
}
\ruleblock{Proof}

\sourcenote{Genre: core rule (tier 1). Setting as in the \rulename{Ax}
section; counterexamples read over the one-wire hypernet.}

\textbf{Claim (\texttt{sound\_orI\_left} /
\texttt{sound\_orI\_right}).}\enspace Two constructors, one per
injection; the payload of either supplies a premise sequent
$\Phi \vdash \Psi_i$ with its sub-derivation, the \emph{added} disjunct
(\texttt{other}), and two tree identities: shared context
(\texttt{preEq}) and, for the left constructor, that the conclusion's
post is $\Psi_i \vee \mathrm{other}$ --- for the right,
$\mathrm{other} \vee \Psi_i$ (\texttt{postEq}). If the premise sequent
is semantically valid, so is $\Phi \vdash \Psi_1 \vee \Psi_2$.

\begin{bthm}[soundness of $\vee$I]\label{thm:ori-sound}
Under the payload identities of either constructor, semantic validity of
the premise sequent implies semantic validity of the conclusion, over
every hypernet.
\end{bthm}
\begin{proof}
Left constructor; the right is its mirror image.

\emph{Step 1 (apply the premise).} Let $\sigma$ be admissible with
$\sigma \models \Phi$; by \texttt{preEq} and premise validity,
$\sigma \models \Psi_i$.

\emph{Step 2 (inject).} The $\vee$-clause of satisfaction is
meta-disjunction, and a disjunction is satisfied by satisfying either
side; the witnessed side is $\Psi_i$, and \texttt{postEq} pins the
conclusion's post as the tree $\Psi_i \vee \mathrm{other}$.
\end{proof}

\begin{bremark}[the added disjunct is wholly unconstrained]
\label{rem:ori-other}
\texttt{other} may be unsatisfiable, ill-typed, or mention untyped
wires; nothing is demanded of it, because satisfaction of a disjunction
demands only the witnessed side and the satisfaction relation carries no
typing side conditions of its own --- typing enters the calculus only
through rules that must evaluate. $\vee$I is conclusion-side weakening,
the mirror of \rulename{Weaken}'s premise side: both discard strength,
and discarding is always semantically safe. Dropping \texttt{postEq}
frees the conclusion outright; dropping \texttt{preEq} admits the
context-swap family (the \rulename{Weaken} section).
\end{bremark}

\begin{bremark}[where a producer actually reaches for it]
Mostly to make two branches agree: the branches of a disjunction
elimination must conclude a \emph{common} formula, so each branch
$\vee$I-injects its own branch-specific conclusion into the shared
disjunction. The separation instance quoted in the $\vee$E section is
the miniature: both of its branch derivations are an \rulename{Ax} node
followed by one $\vee$I node (left in one branch, right in the other).
\end{bremark}

\begin{bremark}[lattice reading]
$x \leq y$ implies $x \leq y \sqcup z$: the upper-bound half of join's
universal property (\texttt{M3.orI\_l\_kit}/\texttt{M3.orI\_r\_kit});
the other half is exactly $\vee$E. The four $\wedge$/$\vee$ rules are
the meet and join adjunctions in full, which is the observation the
$\vee$E section's countermodel turns into a theorem about what the core
\emph{cannot} derive.
\end{bremark}

\begin{anchors}
Steps 1--2 (left) & \texttt{sound\_orI\_left}; payload \texttt{POrIData.left} (fields \texttt{other}, \texttt{preEq}, \texttt{postEq})\\
mirror (right) & \texttt{sound\_orI\_right}; payload \texttt{POrIData.right}\\
dispatch & wired as \texttt{RuleData.orI} into \texttt{ProofObject.sound}\\
lattice reading & \texttt{M3.orI\_l\_kit}, \texttt{M3.orI\_r\_kit}\\
\end{anchors}


\rulesection{7}{$\vee$E}{core rule}
\label{blcatrule-OrE}
\ruleblock{The rule (Steele metanotation)}
\sourcenote{as realized in the mechanized development; source row: \texttt{\texttt{rules.toml}} --- \emph{$\vee$E}. The metanotation statement is typeset by hand for notational quality; its mechanized ground truth is the Lean statement below, extracted from the sources at every build of this book and hash-bound. The rule \emph{roster} is machine-checked (every mechanized rule tag round-trips to exactly one catalogue row); the statement \emph{text} is editorial, checked by eye against the extracted Lean beside it.}
\begin{mathpar}
\inferrule*[Right=$\vee$E]{\Phi_1 \vdash \Psi \\ \Phi_2 \vdash \Psi}{\Phi_1 \vee \Phi_2 \vdash \Psi}
\end{mathpar}

\par\medskip
$\vee$E is the calculus\textquotesingle s only branching rule --- case
analysis: from $\Phi_1 \vdash \Psi$ and $\Phi_2 \vdash \Psi$, conclude
$\Phi_1 \vee \Phi_2 \vdash \Psi$. The disjunctions worth eliminating
come from the hypernet\textquotesingle s synthesized primitive facts
--- a cell indicator\textquotesingle s finite case list, a
rectifier\textquotesingle s sign split --- and a producer eliminates
them by proving the goal once per case over the correspondingly
strengthened premise.

The checker recomputes three tree identities: the two branch
sub-derivations conclude the same formula, the
conclusion\textquotesingle s premise is exactly the $\vee$-node of the
two branch premises, and the conclusion\textquotesingle s post is the
common branch conclusion. The conclusion-matching identity is the
load-bearing one --- drop it and a fact proven only in one branch
escapes the case split (the proof section\textquotesingle s
counterexample); the premise identity\textquotesingle s removal admits
the context-swap family.

The essential caveat of this section, though, is the
rule\textquotesingle s \emph{reach}: the quoted form eliminates a
disjunction only when it is the \emph{whole} premise, while in practice
every eliminable disjunction arrives as a conjunct inside the premise
--- and the core provably cannot restructure its way there
(distribution of $\wedge$ over $\vee$ is underivable; the proof section
exhibits the mechanized lattice countermodel). The context-carrying
generalization --- same rule, plus a conjunct path addressing the
disjunction inside the premise\textquotesingle s $\wedge$-tree --- is
proven sound, proven to be a proper extension, and proven to collapse
onto the quoted form at the identity path; the accompanying development
treats it as the intended elimination rule, with the quoted plain form
destined to become its boundary alias.

\ruleblock{Lean statement}
\declhead{BernsteinLogic.Soundness}{sound\_orE}{\badge{BadgeGreen}{statement bound}}
\begin{leansrc}
theorem sound_orE {H : Hypernet} {seq leftPremise rightPremise : Sequent}
    (ihl : leftPremise.semanticallyValid H)
    (ihr : rightPremise.semanticallyValid H)
    (hpostMatch : leftPremise.post = rightPremise.post)
    (hpre : seq.pre = .or leftPremise.pre rightPremise.pre)
    (hpost : seq.post = leftPremise.post) :
    seq.semanticallyValid H
\end{leansrc}
\par\medskip
Payload of the plain form, read as in the \rulename{Ax} section: two
branch sub-derivations plus structural equalities.

\declhead{BernsteinLogic.Core.Catalogue}{POrEData}{\badge{BadgeGreen}{statement bound}}
\begin{leansrc}
inductive POrEData : Hypernet → Sequent → Type where
  | mk (H : Hypernet) (seq : Sequent)
       (leftPremise  : Sequent) (leftProof  : ProofObject H leftPremise)
       (rightPremise : Sequent) (rightProof : ProofObject H rightPremise)
       (postMatch : leftPremise.post = rightPremise.post)
       (preEq     : seq.pre  = .or leftPremise.pre rightPremise.pre)
       (postEq    : seq.post = leftPremise.post)
       : POrEData H seq
\end{leansrc}
\par\medskip
\textbf{The context-carrying satellite.}\enspace The four declarations
below locate the context-carrying generalization precisely: its
applicability data (a conjunct path with read/replace obligations; the
plain rule is the $\pi = \mathrm{here}$ case), its semantic soundness
over every hypernet, the concrete distribution instance no core
derivation concludes, and the resulting impossibility of elaborating
the generalization into the plain core. The proof section walks the
countermodel argument these rest on.

\declhead{BernsteinLogic.Admissible}{OrECtxStep}{\badge{BadgeGreen}{statement bound}}
\begin{leansrc}
/-- Applicability data of context-carrying disjunction elimination: a
conjunct path into the conclusion's premise addresses a disjunction;
each branch premise is the conclusion's premise with that conjunct
replaced by the respective disjunct; both branches conclude the common
conclusion. Field-for-field, these are the reference checker's
obligations for the rule (path resolves, path addresses a disjunction,
the two premise-replacement equations, the two conclusion equations).
The primitive elimination rule is the `path = here` case. -/
structure OrECtxStep where
  seq          : Sequent
  leftPremise  : Sequent
  rightPremise : Sequent
  path : ConjPath
  d₁ : Formula
  d₂ : Formula
  pathReads : seq.pre.conjGet? path = some (.or d₁ d₂)
  leftPre   : seq.pre.conjReplace? path d₁ = some leftPremise.pre
  rightPre  : seq.pre.conjReplace? path d₂ = some rightPremise.pre
  leftPost  : leftPremise.post = seq.post
  rightPost : rightPremise.post = seq.post
\end{leansrc}
\declhead{BernsteinLogic.Admissible}{sound\_orECtx}{\badge{BadgeGreen}{statement bound}}
\begin{leansrc}
/-- Context-carrying disjunction elimination is sound: whenever both
branch premises are semantically valid, so is the conclusion — over
every hypernet. This is the warrant for adding the rule as an
admissible extension: every sequent the extended calculus derives is
semantically valid (per-rule soundness of the core + this lemma).

The proof is the conjunct-path read lemma (extract the disjunction),
the case split of the primitive rule, and the conjunct-path replace
lemma (transport the branch's strengthened premise). -/
theorem sound_orECtx {H : Hypernet} (st : OrECtxStep)
    (ihl : st.leftPremise.semanticallyValid H)
    (ihr : st.rightPremise.semanticallyValid H) :
    st.seq.semanticallyValid H
\end{leansrc}
\declhead{BernsteinLogic.Admissible}{cx\_not\_derivable}{\badge{BadgeGreen}{statement bound}}
\begin{leansrc}
/-- The instance's conclusion is NOT derivable in the catalogue: its
interpretation drops from an atom to ⊥ across the turnstile, and the
metatheorem says no derivation does that. -/
theorem cx_not_derivable : IsEmpty (ProofObject cxH cxSeq)
\end{leansrc}
\declhead{BernsteinLogic.Admissible}{orECtx\_not\_elaborable}{\badge{BadgeGreen}{statement bound}}
\begin{leansrc}
/-- Context-carrying disjunction elimination does not elaborate into
the core calculus: no elaboration function can exist, because this
instance has derivable branch premises and an underivable conclusion.
(Contrast `OrECtxStep.elabHere`: at the boundary path the elaboration
exists; the failure is specific to genuinely deep paths.) -/
theorem orECtx_not_elaborable :
    ¬ (∀ (st : OrECtxStep) (H : Hypernet),
        Nonempty (ProofObject H st.leftPremise) →
        Nonempty (ProofObject H st.rightPremise) →
        Nonempty (ProofObject H st.seq))
\end{leansrc}
\declhead{BernsteinLogic.Admissible}{ProofObject.m3sound}{\badge{BadgeGreen}{statement bound}}
\begin{leansrc}
/-- Monotone replacement at a conjunct path: substituting a
lattice-smaller formula at any conjunct address shrinks the M3 value.
Meet monotonicity only — no distributivity — which is why Promise,
unlike ∨E-ctx, is lattice-sound and core-eligible. -/
theorem mval_conjReplace_mono :
    ∀ {π : ConjPath} {Γ Γ' x q : Formula},
    Γ.conjGet? π = some q → Γ.conjReplace? π x = some Γ' →
    M3.le (mval x) (mval q) = true →
    M3.le (mval Γ') (mval Γ) = true := by
  intro π
  induction π with
  | here =>
    intro Γ Γ' x q hget hrep hle
    simp only [Formula.conjGet?] at hget
    simp only [Formula.conjReplace?] at hrep
    injection hget with h1
    injection hrep with h2
    rw [← h2, h1]
    exact hle
  | left p ih =>
    intro Γ Γ' x q hget hrep hle
    cases Γ with
    | and l r =>
      simp only [Formula.conjGet?] at hget
      simp only [Formula.conjReplace?] at hrep
      cases hmap : l.conjReplace? p x with
      | none => rw [hmap] at hrep; nomatch hrep
      | some l' =>
        rw [hmap] at hrep
        injection hrep with h2
        rw [← h2]
        show M3.le ((mval l').inf (mval r)) ((mval l).inf (mval r)) = true
        exact M3.inf_mono_l _ _ _ (ih hget hmap hle)
    | atom b tl tr => nomatch hrep
    | inRegion t Δ S => nomatch hrep
    | or l r => nomatch hrep
    | tt => nomatch hrep
    | ff => nomatch hrep
  | right p ih =>
    intro Γ Γ' x q hget hrep hle
    cases Γ with
    | and l r =>
      simp only [Formula.conjGet?] at hget
      simp only [Formula.conjReplace?] at hrep
      cases hmap : r.conjReplace? p x with
      | none => rw [hmap] at hrep; nomatch hrep
      | some r' =>
        rw [hmap] at hrep
        injection hrep with h2
        rw [← h2]
        show M3.le ((mval l).inf (mval r')) ((mval l).inf (mval r)) = true
        exact M3.inf_mono_r _ _ _ (ih hget hmap hle)
    | atom b tl tr => nomatch hrep
    | inRegion t Δ S => nomatch hrep
    | or l r => nomatch hrep
    | tt => nomatch hrep
    | ff => nomatch hrep

/-! ## §A9 The metatheorem: derivability is lattice-sound

Over a hypernet with no typed wires, every derivable sequent satisfies
⟦pre⟧ ≤ ⟦post⟧ in M₃. One case per rule of the catalogue; the
propositional rules are the lattice axioms, equality elimination is the
substitution invariance (its typing field forces both terms closed over
an untyped hypernet), the region/arithmetic-closure rules are handled
by their conclusion shapes or payload absurdities, and the
empty-payload rules are vacuous. -/
theorem ProofObject.m3sound : ∀ {H : Hypernet} {seq : Sequent},
    (∀ w, H.wireType w = none) → ProofObject H seq →
    M3.le (mval seq.pre) (mval seq.post) = true
  | _, _, _, .mk _ _ _ (.ax _ _ (.mk _ _ hpp)) => by
      rw [hpp]; exact M3.le_refl _
  | _, _, hE, .mk _ _ _ (.weaken _ _ (.mk _ _ _ pproof _ hpre hpost)) => by
      rw [hpre, hpost]
      exact M3.weaken_kit _ _ _ (pproof.m3sound hE)
  | _, _, hE, .mk _ _ _ (.cut _ _
      (.mk _ _ _ _ lproof _ rproof hpre hpost hcl hcr)) => by
      have hl
\end{leansrc}
\declhead{BernsteinLogic.Admissible}{cxStep}{\badge{BadgeGreen}{statement bound}}
\begin{leansrc}
/-- The separation instance's applicability data: all five obligations
hold by computation. -/
def cxStep : OrECtxStep
\end{leansrc}
\ruleblock{Correspondence (Steele display $\leftrightarrow$ Lean)}
\badge{BadgeGreen}{reviewer-certified against the displayed texts}
\sourcenote{The dictionary below is a reviewer-certified statement of how the mechanized statements above interpret the displayed rule: display metavariable $\mapsto$ Lean binder, each display hypothesis $\mapsto$ the premise or side condition carrying it, the display conclusion $\mapsto$ the goal --- with any deliberate difference in width between display and mechanization named outright. Certification is recorded as a hash triple over (the displayed cell, the extracted statement texts, this dictionary\textquotesingle s own text); if any of the three changes, this block turns amber until a reviewer re-certifies. The correspondence claim itself remains editorial --- reviewer-grade, not machine-checked; its machine-checked successor types each soundness statement against the rule as a structured object, and supersedes this block per rule as it lands.}
Judgment frame as in \rulename{Ax}; sub-derivation premises as at
\rulename{Weaken}. The displayed rule corresponds to
\texttt{sound\_orE}:
\begin{itemize}
\item the branch judgments $\Phi_1 \vdash \Psi$ and $\Phi_2 \vdash
  \Psi$ $\mapsto$ \texttt{ihl}/\texttt{ihr}; the shared conclusion
  $\Psi$ is shared by equation, \texttt{hpostMatch} $:$
  \texttt{leftPremise.post = rightPremise.post} --- the load-bearing
  identity the commentary singles out.
\item $\Phi_1 \mapsto$ \texttt{leftPremise.pre}; $\Phi_2 \mapsto$
  \texttt{rightPremise.pre}; the conclusion\textquotesingle s premise
  $\Phi_1 \vee \Phi_2$ $\mapsto$ \texttt{hpre} $:$ \texttt{seq.pre =
  .or leftPremise.pre rightPremise.pre} --- the disjunction must be
  the \emph{whole} premise tree, which is precisely the reach
  limitation the commentary discusses.
\item conclusion post $\mapsto$ \texttt{hpost} $:$ \texttt{seq.post =
  leftPremise.post}; the judgment $\mapsto$ the goal
  \texttt{seq.semanticallyValid H}.
\end{itemize}
Width note, deliberate and one-sided: the certified statement set also
carries the context-carrying satellite (\texttt{OrECtxStep},
\texttt{sound\_orECtx}, \texttt{cx\_not\_derivable},
\texttt{orECtx\_not\_elaborable}). Those statements are \emph{about
the generalization}, not the displayed rule: \texttt{sound\_orECtx} is
strictly wider (a conjunct path addresses a disjunction \emph{inside}
the premise tree; the displayed rule is its identity-path case), and
the two negative statements certify that the extension is proper ---
its instance is underivable in, and not elaborable into, the plain
core. The displayed rule itself corresponds to \texttt{sound\_orE}
with no width gap in either direction.

\ruleblock{Lean verifier receipts}
\sourcenote{module \texttt{BernsteinLogic.Soundness}; bound build \texttt{a20d90e66c13\ldots} (bound from receipt a9c0e34d3e3c9c8c, 2026-08-03T21:30:36Z). A receipt records that named checking layers accepted a build: (module, compiled-artifact SHA-256, environment SHA-256, checker layers), minted when the checking gate ran; the store receipts attesting the bound build are listed with the layers each run credited. Nothing here asks the reader to trust the store: every layer is an independently runnable program over the quoted sources --- rebuild the module and re-run the checkers to reproduce the verdict. Receipts attesting superseded builds credit no layer and are omitted from this rendering (the count of omissions, if any, is noted below).}
\begin{receipttable}
\receiptrow{a9c0e34d3e3c9c8c}{2026-08-03T21:30:36Z}{build+decorrelated}{\badge{BadgeGreen}{current}}
\end{receipttable}
\sourcenote{21 further receipt(s) in the store attest superseded builds of this module and are omitted from this rendering.}
\sourcenote{module \texttt{BernsteinLogic.Core.Catalogue}; bound build \texttt{cbc92e976a1d\ldots} (bound from receipt 5abdb3589707b5df, 2026-08-03T21:30:33Z). A receipt records that named checking layers accepted a build: (module, compiled-artifact SHA-256, environment SHA-256, checker layers), minted when the checking gate ran; the store receipts attesting the bound build are listed with the layers each run credited. Nothing here asks the reader to trust the store: every layer is an independently runnable program over the quoted sources --- rebuild the module and re-run the checkers to reproduce the verdict. Receipts attesting superseded builds credit no layer and are omitted from this rendering (the count of omissions, if any, is noted below).}
\begin{receipttable}
\receiptrow{24804e66edba0bc9}{2026-08-03T20:25:09Z}{build+decorrelated}{\badge{BadgeGreen}{current}}
\receiptrow{5abdb3589707b5df}{2026-08-03T21:30:33Z}{build+decorrelated}{\badge{BadgeGreen}{current}}
\end{receipttable}
\sourcenote{18 further receipt(s) in the store attest superseded builds of this module and are omitted from this rendering.}
\sourcenote{module \texttt{BernsteinLogic.Admissible}; bound build \texttt{52e5099ba522\ldots} (bound from receipt c17abf2e103fe745, 2026-07-21T05:28:22Z). A receipt records that named checking layers accepted a build: (module, compiled-artifact SHA-256, environment SHA-256, checker layers), minted when the checking gate ran; the store receipts attesting the bound build are listed with the layers each run credited. Nothing here asks the reader to trust the store: every layer is an independently runnable program over the quoted sources --- rebuild the module and re-run the checkers to reproduce the verdict. Receipts attesting superseded builds credit no layer and are omitted from this rendering (the count of omissions, if any, is noted below).}
\begin{receipttable}
\receiptrow{7f3043c54f717472}{2026-07-20T07:43:38Z}{build+decorrelated}{\badge{BadgeGreen}{current}}
\receiptrow{acaf0e8871731343}{2026-07-20T20:44:28Z}{build+decorrelated}{\badge{BadgeGreen}{current}}
\receiptrow{62974466ea131d89}{2026-07-20T21:00:57Z}{build+replay+hardened+decorrelated}{\badge{BadgeGreen}{current}}
\receiptrow{c17abf2e103fe745}{2026-07-21T05:28:22Z}{build+replay+hardened+decorrelated}{\badge{BadgeGreen}{current}}
\end{receipttable}
\sourcenote{6 further receipt(s) in the store attest superseded builds of this module and are omitted from this rendering.}
\layerline{\layeritem{build}{lake build (elaboration)}{\badge{BadgeGreen}{current}}\layeritem{replay}{leanchecker (official kernel replay)}{\badge{BadgeGreen}{current}}\layeritem{hardened}{safe\_verify (hardened screen)}{\badge{BadgeGreen}{current}}\layeritem{decorrelated}{lean4export + nanoda (decorrelated kernel)}{\badge{BadgeGreen}{current}}}
\sourcenote{\badge{BadgeGray}{comparator: none} no fixture in the three-implementation comparator harness (compiled / kernel-interpreted / reference-interpreted) exercises this rule.}
\ruleblock{HOL4 statement}
\sourcenote{theory \texttt{bern\_fam\_ore}, HOL4 development script \texttt{\texttt{bern\_fam\_oreScript.sml}} \badge{BadgeGreen}{script bound}}
\holhead{Definition checkFam\_ore\_def}
\begin{holsrc}
Definition checkFam_ore_def:
  checkFam_ore (cols:columns) (ct:ctx) (i:num) ⇔
    case nodeSq cols ct i of
    | NONE => F
    | SOME sq =>
      case childOrd ct i 0 of
      | NONE => F
      | SOME kl =>
        case childOrd ct i 1 of
        | NONE => F
        | SOME kr =>
          case nodeSq cols ct kl of
          | NONE => F
          | SOME ls =>
            case nodeSq cols ct kr of
            | NONE => F
            | SOME rs =>
              case pathArgs cols ct i of
              | NONE => F
              | SOME (path, ifc) =>
                case conjGet sq.pre path of
                | SOME (FOr d1 d2) =>
                    (conjReplace sq.pre path d1 = SOME ls.pre ∧
                     conjReplace sq.pre path d2 = SOME rs.pre ∧
                     ls.post = sq.post ∧ rs.post = sq.post)
                | _ => F
End
\end{holsrc}
\holhead{Theorem famEq\_ore}
\begin{holsrc}
Theorem famEq_ore:
  ∀cols ct i sq nd ws.
    wfColumns cols ct ∧ mkCtx cols = SOME ct ∧
    nodeCorr cols ct i (CertPf sq nd) ∧
    FST (THE (tagOf cols i)) ∈ {6:num} ⇒
    (checkFam_ore cols ct i ⇔ oreArm ws sq nd)
\end{holsrc}
\ruleblock{HOL verifier receipts}
\badge{BadgeBlue}{documentary} Documentary citations. Spine census: \texttt{bern\_joins} GREEN under \texttt{hol-strict} (55 theorems, 0 offences, 208-theory ancestry including this theory), \texttt{checkJoins\_eq\_checkPf} verified verbatim in the built signature. The cited \texttt{famEq} lemma is the kernel-checked per-node equivalence: under well-formed columns (\texttt{wfColumns}/\texttt{mkCtx}/\texttt{nodeCorr}) at this rule\textquotesingle s wire tag, the columnar family check coincides with the flat reference walk\textquotesingle s arm for this rule. There is no signed HOL receipt store yet --- these citations are documentary.
Semantic stratum, also ported: the soundness twins \texttt{sound\_orE}
(theory \texttt{bern\_tr\_soundstruct}: \texttt{hol-strict} clean, 9
theorems, 0 offences, verification run of 2026-07-22 on an isolated
build host, content-addressed at the same sources) and
\texttt{sound\_orECtx} (theory \texttt{bern\_tr\_kitd}, 9 theorems, 0
offences, same run --- the context form the wire\textquotesingle s tag
carries), consumed by the certified-fragment keystone
(\texttt{checkPf\_sound\_frag}, theory \texttt{bern\_tr\_splice}, 20
theorems, 0 offences, same run): a reference-checker accept inside the
proven wire fragment yields sequent validity.

\sourcenote{HOL receipts are documentary: quoted results of \texttt{hol-strict} runs
and census rows from the HOL4 development\textquotesingle s build
evidence, not minted receipt artefacts. A HOL receipt store (the Lean
store\textquotesingle s twin) is named future work.
}
\ruleblock{Proof}

\sourcenote{Genre: core rule (tier 1). Setting as in the \rulename{Ax}
section; counterexamples read over the one-wire hypernet. This section
also carries the fork note on the rule's \emph{reach}: the plain form
quoted here, its mechanized limitation, and the context-carrying
generalization of which it is the boundary case.}

\textbf{Claim (\texttt{sound\_orE}).}\enspace The payload supplies two
branch sequents $\Phi_1 \vdash \Psi$ and $\Phi_2 \vdash \Psi$ with their
sub-derivations and three tree identities: the branches share one
conclusion (\texttt{postMatch}), the conclusion's premise is the
disjunction $\Phi_1 \vee \Phi_2$ of the branch premises (\texttt{preEq}),
and the conclusion's post is the left branch's post (\texttt{postEq}).
If both branch sequents are semantically valid, so is
$\Phi_1 \vee \Phi_2 \vdash \Psi$.

\begin{bthm}[soundness of $\vee$E, plain form]\label{thm:ore-sound}
Under the three payload identities, semantic validity of both branch
sequents implies semantic validity of the conclusion, over every
hypernet.
\end{bthm}
\begin{proof}
Let $\sigma$ be admissible with $\sigma$ satisfying the conclusion's
premise, which \texttt{preEq} pins as the tree $\Phi_1 \vee \Phi_2$.

\emph{Step 1 (case split).} The $\vee$-clause of satisfaction is
meta-disjunction: $\sigma \models \Phi_1$ or $\sigma \models \Phi_2$.

\emph{Step 2 (either branch delivers).} In the first case, left branch
validity gives $\sigma \models \Psi$ directly (\texttt{postEq}). In the
second, right branch validity gives satisfaction of the right branch's
conclusion, which \texttt{postMatch} identifies --- as a tree --- with
the left's, hence with the conclusion's post.
\end{proof}

\begin{bremark}[\texttt{postMatch} is load-bearing]\label{rem:ore-post}
The branches must agree on what they conclude \emph{before} the split is
trusted. Drop \texttt{postMatch} and take the left branch
$w \leq 0 \vdash w \leq 0$ and the right branch $\top \vdash \top$ (both
\rulename{Ax} instances): the degenerate rule, reading its conclusion
from the left branch, concludes $(w \leq 0) \vee \top \vdash w \leq 0$
--- false at $\sigma(w) = 1$, where the premise holds through the right
disjunct. \texttt{preEq}'s removal admits the context-swap family (the
\rulename{Weaken} section); \texttt{postEq}'s removal frees the
conclusion outright. In the lattice reading the rule is the bounding
half of join's universal property --- $x \leq z$ and $y \leq z$ iff
$x \sqcup y \leq z$ --- whose introducing half is $\vee$I.
\end{bremark}

\begin{bremark}[the reach problem, and why it is a theorem]
\label{rem:ore-reach}
\texttt{preEq} demands that the eliminated disjunction \emph{be} the
whole premise. But in this calculus the disjunctions worth eliminating
almost never are: a sequent's premise in practice has the shape
$\Gamma \wedge \mathsf{SP}$ --- a context conjoined with the synthesized
case facts of the hypernet's primitives (a cell indicator's finite case
list, a rectifier's sign split) --- so the disjunction arrives as a
\emph{conjunct} of the premise, one or more $\wedge$-nodes deep. The
plain rule cannot reach it, and the core cannot restructure its way
there, because distributing $\wedge$ over $\vee$ is underivable there.
That is a theorem, not a design suspicion. Interpret formulas into
$M_3$, the five-element diamond --- three pairwise-incomparable atoms
between $\bot$ and $\top$, the canonical modular non-distributive
lattice \cite{birkhoff1967} --- in full: $\wedge$ to meet, $\vee$ to
join, $\top$ to top, $\bot$ to bottom, three disequality atoms on
designated wires to the three lattice atoms, region atoms to bottom,
and every remaining atom to top. Work over a hypernet all of whose
wires are untyped; that hypothesis is what discharges every
non-structural rule (typing packs force closed terms there, region and
arithmetic payloads are handled by their conclusion shapes or turn
absurd, and the payload-free rules are vacuous). Then a derivable
sequent's premise always interprets $\leq$ its conclusion --- by
induction on derivations, one case per rule, the structural cases
being precisely the lattice axioms of order, meet, and join: the
sections \rulename{Ax} through $\vee$E in this book, one kit lemma
each (\texttt{ProofObject.m3sound}). No rule of the \emph{plain core}
--- the calculus as quoted in these sections, prior to the
context-carrying extension of \autoref{rem:ore-fork}, whose whole
point is otherwise --- consumes distributivity. Now take the
distribution instance
$P \wedge (Q \vee R) \;\vdash\; (P \wedge Q) \vee (P \wedge R)$:
its premise interprets to $a \sqcap (b \sqcup c) = a \sqcap \top = a$,
its conclusion to $(a \sqcap b) \sqcup (a \sqcap c) = \bot \sqcup \bot
= \bot$, and $a \not\leq \bot$ --- so no core derivation concludes it
(\texttt{cx\_not\_derivable}), even though both branch premises
$P \wedge Q \vdash \cdot$ and $P \wedge R \vdash \cdot$ are two-step
core derivations (an \rulename{Ax} node and a $\vee$I node each).
\end{bremark}

\begin{bremark}[the context-carrying form, and the fork discipline]
\label{rem:ore-fork}
The repair is not restructuring but a stronger elimination. The
context-carrying rule's payload adds a \emph{conjunct path} $\pi$ into
the premise's $\wedge$-tree, with read and replace obligations: the
conjunct at $\pi$ is a disjunction $\Phi_1 \vee \Phi_2$, each branch
premise is the conclusion's premise with that conjunct replaced by the
respective disjunct, and both branches conclude the common conclusion
(\texttt{OrECtxStep}). Three mechanized facts locate it precisely.
\emph{Sound}: whenever both branch premises are semantically valid, so
is the conclusion, over every hypernet (\texttt{sound\_orECtx} --- one
path-read lemma to extract the disjunction, the case split of
Theorem~\ref{thm:ore-sound}, one path-replace lemma to transport each
branch's strengthened premise). \emph{Proper}: no elaboration of the
context-carrying rule into the core exists
(\texttt{orECtx\_not\_elaborable}) --- concretely, at the conjunct path
addressing $Q \vee R$ as the premise's right conjunct, the read
obligation returns that disjunction and the replace obligations return
the branch premises $P \wedge Q$ and $P \wedge R$ of
\autoref{rem:ore-reach}'s instance (\texttt{cxStep}); its branches
are derivable there and its conclusion is not, so the extension proves
something true that the core cannot reach. \emph{Conservative at the boundary}: at $\pi = \mathrm{here}$
the rule elaborates into the plain constructor
(\texttt{OrECtxStep.elabHere}) --- the form proven sound above is
exactly the identity-path case. The cost side of the fork is small and
bounded: the certificate carries the path, the checker two small path
functions (read-at, replace-at), and the soundness proof one lemma
about conjunct-path replacement. The alternative --- adding
$\wedge$-restructuring rules to the core (rules whose premise and
conclusion rearrange the same conjuncts: reassociations, permutations)
--- cannot substitute: such rules are sound for every lattice
interpretation, so the induction above absorbs them --- as it absorbs
any rule that preserves $M_3$-order --- and the countermodel stands;
closing the gap requires a rule whose soundness genuinely uses
distributivity, which is the context-carrying elimination itself. The
accompanying development accordingly treats the context-carrying form
as the intended elimination rule; the constructor quoted in this
section's Lean cell is the plain form, destined to become its boundary
alias.
\end{bremark}

\begin{anchors}
Steps 1--2 & \texttt{sound\_orE}; payload \texttt{POrEData} (fields \texttt{postMatch}, \texttt{preEq}, \texttt{postEq}); wired as \texttt{RuleData.orE}\\
$\vee$-clause of satisfaction & \texttt{Formula.holds} (case \texttt{or})\\
lattice reading & \texttt{M3.orE\_kit}\\
Remark \ref{rem:ore-reach}: order-preservation metatheorem & \texttt{ProofObject.m3sound} (BernsteinLogic/Admissible.lean)\\
Remark \ref{rem:ore-reach}: underivable distribution instance & \texttt{cxStep}, \texttt{cx\_not\_derivable}\\
Remark \ref{rem:ore-fork}: extension payload + soundness & \texttt{OrECtxStep}, \texttt{sound\_orECtx}\\
Remark \ref{rem:ore-fork}: properness & \texttt{orECtx\_not\_elaborable}\\
Remark \ref{rem:ore-fork}: boundary elaboration & \texttt{OrECtxStep.elabHere}\\
\end{anchors}


\rulesection{8}{EqElim}{core rule}
\label{blcatrule-EqElim}
\ruleblock{The rule (Steele metanotation)}
\sourcenote{as realized in the mechanized development; source row: \texttt{\texttt{rules.toml}} --- \emph{EqElim}. The metanotation statement is typeset by hand for notational quality; its mechanized ground truth is the Lean statement below, extracted from the sources at every build of this book and hash-bound. The rule \emph{roster} is machine-checked (every mechanized rule tag round-trips to exactly one catalogue row); the statement \emph{text} is editorial, checked by eye against the extracted Lean beside it.}
\begin{mathpar}
\inferrule*[Right=\rulename{EqElim}]{\Phi \vdash t = u \\ \Phi \vdash \Psi\sub{a}{t} \\ \exists s.\; \bshape_H(t) = s = \bshape_H(u)}{\Phi \vdash \Psi\sub{a}{u}}
\end{mathpar}

\par\medskip
\rulename{EqElim} is Leibniz substitution as a proof rule: an established
equality $t = u$ licenses replacing $t$ by $u$ at a designated hole in an
established formula. The payload names the hole (a wire $a$), the body
formula $\Psi$, and the two terms; premises prove $t = u$ and
$\Psi\sub{a}{t}$ over the shared context, and the conclusion is
$\Psi\sub{a}{u}$. A producer reaches for it wherever a derived fact must
be rewritten along an equality: pushing a resolved value through a bound
(the ReLU rules conclude equalities such as $y = h$ exactly so that this
rule can consume them), transporting an atom across a
\rulename{Norm}-derived respelling, or assembling transitivity --- which
is not primitive in this calculus: a one-node derivation through this
rule (the proof section\textquotesingle s closing remark spells it out;
\rulename{EqRefl} and \rulename{EqSym} complete the kit).

What the checker recomputes: both substitution instances --- it applies
the (binder-free, hence capture-free) substitution function to
$(\Psi, a, t)$ and to $(\Psi, a, u)$ and demands syntactic equality with
the cited premise and conclusion formulas --- plus the premise-sharing
equations, and the typing side condition: $t$ and $u$ carry a
\emph{common defined} broadcast shape in the ambient hypernet. Nothing
semantic is consulted; every side condition is a decidable syntactic
equation.

The typing condition is the load-bearing hypothesis. Equality atoms hold
\emph{elementwise under broadcast}, so a scalar can stand equal to a
vector whose every entry it matches --- a genuine, derivable, usable
fact that is nonetheless not a license to interchange the two terms: in
a region atom, which constrains the shape of its term, the substitution
flips satisfiability. Dropping the condition admits a concrete derivation
of an unsatisfiable conclusion from a satisfiable premise (the proof
section constructs it); with the condition, that certificate is refused
because no common shape exists. Dropping either substitution equation
instead would let the conclusion be an arbitrary formula; dropping the
premise-match would combine facts proven under different contexts.

\cataloguenote{Producer-side, equational obligations are typically discharged by
congruence-closure engines \cite{nelsonoppen1980,nieuwenhuisoliveras2007};
the certificate carries not the closure computation but its residue ---
a chain of \rulename{EqRefl}/\rulename{EqSym}/\rulename{EqElim} nodes,
each checked by recomputing one substitution.
}
\ruleblock{Lean statement}
\declhead{BernsteinLogic.Soundness}{Term.evalQ\_subst\_congr}{\badge{BadgeGreen}{statement bound}}
\begin{leansrc}
theorem Term.evalQ_subst_congr {H : Hypernet} (σ : H.Valuation)
    {w : Wire} {u₁ u₂ : Term} (hu : Term.evalQ σ u₁ = Term.evalQ σ u₂) :
    ∀ (t : Term),
      Term.evalQ σ (Term.subst w u₁ t) = Term.evalQ σ (Term.subst w u₂ t)
\end{leansrc}
\declhead{BernsteinLogic.Soundness}{Formula.holds\_subst\_congr}{\badge{BadgeGreen}{statement bound}}
\begin{leansrc}
theorem Formula.holds_subst_congr {H : Hypernet} (σ : H.Valuation)
    {w : Wire} {u₁ u₂ : Term} (hu : Term.evalQ σ u₁ = Term.evalQ σ u₂) :
    ∀ (F : Formula), (Formula.holds σ (F.subst w u₁) ↔ Formula.holds σ (F.subst w u₂))
\end{leansrc}
\declhead{BernsteinLogic.Soundness}{Term.evalQ\_bshape}{\badge{BadgeGreen}{statement bound}}
\begin{leansrc}
/-- bshape soundness: a defined evaluation takes the broadcast shape. -/
theorem Term.evalQ_bshape {H : Hypernet} (σ : H.Valuation) :
    ∀ (t : Term) (s : Shape) (q : QTensor),
    t.bshape H = some s → Term.evalQ σ t = some q → q.shape = s
\end{leansrc}
\declhead{BernsteinLogic.Soundness}{QTensor.eq\_of\_rel2}{\badge{BadgeGreen}{statement bound}}
\begin{leansrc}
/-- rel2 with an equality-forcing relation at equal shapes forces tensor
equality (the PROPOSED-B13 EqElim extraction: without equal shapes the
broadcast branches of rel2 defeat this — the header's counterexample). -/
theorem QTensor.eq_of_rel2 {r : ℚ → ℚ → Prop} {a b : QTensor}
    (hr : ∀ x y, r x y → x = y) (hs : a.shape = b.shape)
    (h : QTensor.rel2 r a b) : a = b
\end{leansrc}
\declhead{BernsteinLogic.Soundness}{sound\_eqElim}{\badge{BadgeGreen}{statement bound}}
\begin{leansrc}
theorem sound_eqElim {H : Hypernet} {seq eqPremise bodyPremise : Sequent}
    {wire : Wire} {termT termU : Term} {bodyForm : Formula}
    (ihEq : eqPremise.semanticallyValid H)
    (ihBody : bodyPremise.semanticallyValid H)
    (hpre : seq.pre = eqPremise.pre)
    (hpreMatch : eqPremise.pre = bodyPremise.pre)
    (heqPost : eqPremise.post = .atom .eq termT termU)
    (htyEq : ∃ s, termT.bshape H = some s ∧ termU.bshape H = some s)
    (hbodySubT : substituteInFormula bodyForm wire termT bodyPremise.post)
    (hconclSubU : substituteInFormula bodyForm wire termU seq.post) :
    seq.semanticallyValid H
\end{leansrc}
\par\medskip
\textbf{The substitution the rule trusts, in full.}\enspace The payload
fields \texttt{bodySubT} and \texttt{conclSubU} are equations against the
substitution function displayed below: plain leaf replacement over a
grammar with no binders, so no capture or freshness side condition exists
anywhere. A skeptical reader should confirm by eye that
\texttt{substituteInFormula} is nothing but that equation --- the rule
manipulates formulas syntactically and consults no valuation; semantic
truth enters only through \texttt{sound\_eqElim}, from the premise
validity hypotheses.

\declhead{BernsteinLogic.Core.Catalogue}{Term.subst}{\badge{BadgeGreen}{definition bound}}
\begin{leansrc}
def Term.subst (w : Wire) (t : Term) : Term → Term
  | .var w' => if w' = w then t else .var w'
  | .const c => .const c
  | .add l r => .add (Term.subst w t l) (Term.subst w t r)
  | .mul l r => .mul (Term.subst w t l) (Term.subst w t r)
  | .sub l r => .sub (Term.subst w t l) (Term.subst w t r)
  | .neg u => .neg (Term.subst w t u)
  | .idx u i => .idx (Term.subst w t u) i
\end{leansrc}
\declhead{BernsteinLogic.Core.Catalogue}{Formula.subst}{\badge{BadgeGreen}{definition bound}}
\begin{leansrc}
def Formula.subst (w : Wire) (t : Term) : Formula → Formula
  | .atom b l r => .atom b (Term.subst w t l) (Term.subst w t r)
  | .inRegion u simplex S => .inRegion (Term.subst w t u) simplex S
  | .and l r => .and (Formula.subst w t l) (Formula.subst w t r)
  | .or l r => .or (Formula.subst w t l) (Formula.subst w t r)
  | .tt => .tt
  | .ff => .ff
\end{leansrc}
\declhead{BernsteinLogic.Core.Catalogue}{substituteInFormula}{\badge{BadgeGreen}{definition bound}}
\begin{leansrc}
def substituteInFormula (F : Formula) (w : Wire) (t : Term) (F' : Formula) : Prop :=
  F' = F.subst w t
\end{leansrc}
\declhead{BernsteinLogic.Core.Catalogue}{PEqElimData}{\badge{BadgeGreen}{definition bound}}
\begin{leansrc}
inductive PEqElimData : Hypernet → Sequent → Type where
  | mk (H : Hypernet) (seq : Sequent)
       (eqPremise   : Sequent) (eqProof   : ProofObject H eqPremise)
       (bodyPremise : Sequent) (bodyProof : ProofObject H bodyPremise)
       (wire : Wire) (termT termU : Term) (bodyForm : Formula)
       (preEq     : seq.pre = eqPremise.pre)
       (preMatch  : eqPremise.pre = bodyPremise.pre)
       (eqPost    : eqPremise.post = .atom .eq termT termU)
       (tyEq      : ∃ s, termT.bshape H = some s ∧ termU.bshape H = some s)
       (bodySubT  : substituteInFormula bodyForm wire termT bodyPremise.post)
       (conclSubU : substituteInFormula bodyForm wire termU seq.post)
       : PEqElimData H seq
\end{leansrc}
\ruleblock{Correspondence (Steele display $\leftrightarrow$ Lean)}
\badge{BadgeGreen}{reviewer-certified against the displayed texts}
\sourcenote{The dictionary below is a reviewer-certified statement of how the mechanized statements above interpret the displayed rule: display metavariable $\mapsto$ Lean binder, each display hypothesis $\mapsto$ the premise or side condition carrying it, the display conclusion $\mapsto$ the goal --- with any deliberate difference in width between display and mechanization named outright. Certification is recorded as a hash triple over (the displayed cell, the extracted statement texts, this dictionary\textquotesingle s own text); if any of the three changes, this block turns amber until a reviewer re-certifies. The correspondence claim itself remains editorial --- reviewer-grade, not machine-checked; its machine-checked successor types each soundness statement against the rule as a structured object, and supersedes this block per rule as it lands.}
Judgment frame as in \rulename{Ax}; sub-derivation premises as at
\rulename{Weaken}.
\begin{itemize}
\item $t, u \mapsto$ \texttt{termT termU : Term}; the hole $a \mapsto$
  \texttt{wire : Wire}; the body $\Psi \mapsto$ \texttt{bodyForm :
  Formula}; $\Phi \mapsto$ \texttt{seq.pre}, shared with both premise
  sequents through \texttt{hpre} and \texttt{hpreMatch}.
\item the equality premise $\Phi \vdash t = u$ $\mapsto$ \texttt{ihEq}
  with the shape tie \texttt{heqPost} $:$ \texttt{eqPremise.post =
  .atom .eq termT termU}.
\item the body premise $\Phi \vdash \Psi\sub{a}{t}$ $\mapsto$
  \texttt{ihBody} $:$ \texttt{bodyPremise.semanticallyValid H}, its
  post tied to $\Psi\sub{a}{t}$ by \texttt{hbodySubT}; the goal
  post is tied to $\Psi\sub{a}{u}$ by \texttt{hconclSubU}. Both ties
  run through \texttt{substituteInFormula}, quoted in full below: the
  one-line equation \texttt{F\textquotesingle{} = F.subst w t} over
  plain leaf replacement --- the grammar has no binders, so no capture
  or freshness side condition exists to omit.
\item the typing premise $\exists s.\; \bshape_H(t) = s = \bshape_H(u)$
  $\mapsto$ \texttt{htyEq}: existence of one shape \texttt{s} with
  \texttt{termT.bshape H = some s} and \texttt{termU.bshape H = some
  s} --- the \texttt{some} carries the definedness the
  display\textquotesingle s equation chain asserts by letting $s$
  range over shapes.
\item conclusion $\Phi \vdash \Psi\sub{a}{u}$ $\mapsto$ the goal
  \texttt{seq.semanticallyValid H}. Certified alongside as the
  semantic engine: \texttt{Term.evalQ\_subst\_congr} and
  \texttt{Formula.holds\_subst\_congr} --- substitution congruence
  along equal evaluations, the two lemmas the soundness proof routes
  through.
\end{itemize}
No width difference in either direction: statement and display carry
the same typing gate (the commentary\textquotesingle s counterexample
is what its absence admits), and every hypothesis is consumed.

\ruleblock{Lean verifier receipts}
\sourcenote{module \texttt{BernsteinLogic.Soundness}; bound build \texttt{a20d90e66c13\ldots} (bound from receipt a9c0e34d3e3c9c8c, 2026-08-03T21:30:36Z). A receipt records that named checking layers accepted a build: (module, compiled-artifact SHA-256, environment SHA-256, checker layers), minted when the checking gate ran; the store receipts attesting the bound build are listed with the layers each run credited. Nothing here asks the reader to trust the store: every layer is an independently runnable program over the quoted sources --- rebuild the module and re-run the checkers to reproduce the verdict. Receipts attesting superseded builds credit no layer and are omitted from this rendering (the count of omissions, if any, is noted below).}
\begin{receipttable}
\receiptrow{a9c0e34d3e3c9c8c}{2026-08-03T21:30:36Z}{build+decorrelated}{\badge{BadgeGreen}{current}}
\end{receipttable}
\sourcenote{21 further receipt(s) in the store attest superseded builds of this module and are omitted from this rendering.}
\sourcenote{module \texttt{BernsteinLogic.Core.Catalogue}; bound build \texttt{cbc92e976a1d\ldots} (bound from receipt 5abdb3589707b5df, 2026-08-03T21:30:33Z). A receipt records that named checking layers accepted a build: (module, compiled-artifact SHA-256, environment SHA-256, checker layers), minted when the checking gate ran; the store receipts attesting the bound build are listed with the layers each run credited. Nothing here asks the reader to trust the store: every layer is an independently runnable program over the quoted sources --- rebuild the module and re-run the checkers to reproduce the verdict. Receipts attesting superseded builds credit no layer and are omitted from this rendering (the count of omissions, if any, is noted below).}
\begin{receipttable}
\receiptrow{24804e66edba0bc9}{2026-08-03T20:25:09Z}{build+decorrelated}{\badge{BadgeGreen}{current}}
\receiptrow{5abdb3589707b5df}{2026-08-03T21:30:33Z}{build+decorrelated}{\badge{BadgeGreen}{current}}
\end{receipttable}
\sourcenote{18 further receipt(s) in the store attest superseded builds of this module and are omitted from this rendering.}
\layerline{\layeritem{build}{lake build (elaboration)}{\badge{BadgeGreen}{current}}\layeritem{replay}{leanchecker (official kernel replay)}{\badge{BadgeAmber}{stale only}}\layeritem{hardened}{safe\_verify (hardened screen)}{\badge{BadgeAmber}{stale only}}\layeritem{decorrelated}{lean4export + nanoda (decorrelated kernel)}{\badge{BadgeGreen}{current}}}
\sourcenote{\badge{BadgeGray}{comparator: none} no fixture in the three-implementation comparator harness (compiled / kernel-interpreted / reference-interpreted) exercises this rule.}
\ruleblock{HOL4 statement}
\sourcenote{theory \texttt{bern\_catalogue}, HOL4 development script \texttt{\texttt{bern\_catalogueScript.sml}} \badge{BadgeGreen}{script bound}}
\holhead{Definition termSubst\_def}
\begin{holsrc}
Definition termSubst_def:
  termSubst w t (TmVar w') = (if w' = w then t else TmVar w') ∧
  termSubst w t (TmConst c) = TmConst c ∧
  termSubst w t (TmAdd l r) = TmAdd (termSubst w t l) (termSubst w t r) ∧
  termSubst w t (TmMul l r) = TmMul (termSubst w t l) (termSubst w t r) ∧
  termSubst w t (TmSub l r) = TmSub (termSubst w t l) (termSubst w t r) ∧
  termSubst w t (TmNeg u) = TmNeg (termSubst w t u) ∧
  termSubst w t (TmIdx u i) = TmIdx (termSubst w t u) i
End
\end{holsrc}
\holhead{Definition formulaSubst\_def}
\begin{holsrc}
Definition formulaSubst_def:
  formulaSubst w t (FAtom b l r) =
    FAtom b (termSubst w t l) (termSubst w t r) ∧
  formulaSubst w t (FInRegion u sx es) = FInRegion (termSubst w t u) sx es ∧
  formulaSubst w t (FAnd l r) =
    FAnd (formulaSubst w t l) (formulaSubst w t r) ∧
  formulaSubst w t (FOr l r) =
    FOr (formulaSubst w t l) (formulaSubst w t r) ∧
  formulaSubst w t FTt = FTt ∧
  formulaSubst w t FFf = FFf
End
\end{holsrc}
\holhead{Definition substituteInFormulaB\_def}
\begin{holsrc}
Definition substituteInFormulaB_def:
  substituteInFormulaB fm w t fm' ⇔ fm' = formulaSubst w t fm
End
\end{holsrc}
\ruleblock{HOL verifier receipts}
\badge{BadgeBlue}{documentary} \texttt{hol-strict} CLEAN census over the ported checker theory:
\texttt{bern\_catalogue}, 472 theorems, 0 cheat-tagged, 0 nonstandard
axioms (whitelist: \texttt{DISK\_THM}; the four \texttt{bool} axioms) ---
verification run of 2026-07-22 on an isolated build host, against a
content-addressed checkout of the exact sources cited here --- the
earlier stale-run caveat (a 2026-07-18 script revision postdating the
prior census) is retired by this run, which attests the current text.
There is no signed HOL receipt store yet --- these citations are
documentary. Scope note, honest: the substitution core of
this rule\textquotesingle s check is what is ported; the broadcast-shape
side condition (\texttt{tyEq}) rides the hypernet typing island, which
the port deliberately defers to its proof-object stage.
Semantic stratum, also ported: the soundness twin \texttt{sound\_eqElim}
(theory \texttt{bern\_tr\_fragclose}: \texttt{hol-strict} clean, 26
theorems, 0 offences, verification run of 2026-07-22 on an isolated
build host, content-addressed at the same sources), consumed by the
certified-fragment keystone (\texttt{checkPf\_sound\_frag}, theory
\texttt{bern\_tr\_splice}, 20 theorems, 0 offences, same run): a
reference-checker accept inside the proven wire fragment yields sequent
validity.

\sourcenote{HOL receipts are documentary: quoted results of \texttt{hol-strict} runs
and census rows from the HOL4 development\textquotesingle s build
evidence, not minted receipt artefacts. A HOL receipt store (the Lean
store\textquotesingle s twin) is named future work.
}
\ruleblock{Proof}

\sourcenote{Genre: core rule (tier 1) --- an ordinary mathematical proof at
Putnam-exam-solution detail. Every hypothesis of the rule is used below and
remarked on; the formal-anchor table maps each step to the Lean declaration
that carries it. This section also fixes the semantic setting that the
other equality and conversion sections cite.}

\textbf{Setting.}\enspace Fix a hypernet $H$. A \emph{valuation} $\sigma$
assigns each \emph{typed} wire a tensor of its declared shape over its
declared dtype (an extended dtype admits infinite entries; untyped wires
are assigned nothing), and $\sigma$ is \emph{admissible} when every
primitive application in $H$ satisfies its semantic relation --- as at
the \rulename{Ax} setting. The rational term evaluation
$\mathrm{ev}_\sigma(t)$ is \emph{partial} and the evaluator is structural
and \emph{strict} --- a compound term evaluates exactly when its parts
evaluate and the constructor's own side condition is met: shapes joining
for the binary operations, nothing further for negation, a rank-1
operand with in-range index for projection --- so a mention of an
untyped wire, a wire holding an infinite value of an extended dtype, a
broadcast-incompatible operation, or a projection whose operand is not
rank-1 or whose index is out of range yields no value. A comparison atom
$t_1 \bowtie t_2$ \emph{holds} at $\sigma$
when both sides evaluate and the evaluated tensors stand in the relation
$\bowtie$ \emph{elementwise under broadcast}: equal shapes compare entry
by entry; a rank-0 (scalar) tensor against a tensor of any shape compares
against every entry; two unequal positive-rank shapes never compare. A
region atom holds when its term evaluates to a rank-1 tensor of extent
the region's ambient dimension whose entries form a rational point of
the region --- a rank-0 evaluation is never coerced to a one-dimensional
point. A sequent $\Phi \vdash \Psi$ is \emph{semantically valid} when
every admissible $\sigma$ satisfying $\Phi$ satisfies $\Psi$;
counterexamples in the equality and conversion sections read, as at the
\rulename{Ax} setting, over application-free hypernets, where every
valuation is admissible, so an unsound conclusion genuinely falsifies.
Note two
consequences that drive everything below: equality atoms are
\emph{evaluation-level} facts, not syntactic ones; and a holding equality
between tensors of \emph{different} broadcast shapes (a scalar against a
vector) does not make the two sides interchangeable --- it only says the
scalar matches every entry.

\textbf{Claim} (\texttt{sound\_eqElim}).\enspace The rule's payload
supplies a hole wire $a$, terms $t, u$, a body formula $\Psi$, and checked
side conditions: the two premises share the composite premise $\Phi$; the
first premise concludes the atom $t = u$; the second premise concludes
$\Psi\sub{a}{t}$; the conclusion is $\Psi\sub{a}{u}$; and (the typing
condition) $\exists s.\ \bshape_H(t) = s = \bshape_H(u)$ --- both terms
carry a \emph{common defined} broadcast shape in $H$. If both premises are
semantically valid, so is $\Phi \vdash \Psi\sub{a}{u}$.

\begin{blemma}[substitution congruence, term level]\label{lem:eqelim-term}
Fix $\sigma$, and suppose $\mathrm{ev}_\sigma(t) = \mathrm{ev}_\sigma(u)$
as partial values (both undefined, or both defined and equal). Then for
every term $e$,
$\mathrm{ev}_\sigma(e\sub{a}{t}) = \mathrm{ev}_\sigma(e\sub{a}{u})$.
\end{blemma}
\begin{proof}
Structural induction on $e$. The base case $e = a$ is the hypothesis; the
base cases $e$ a constant or a different wire are reflexivity; every
compound constructor is a congruence of the evaluator (the evaluator of a
sum, product, difference, negation, or projection is a function of the
evaluations of the parts). Two remarks. First, the hypothesis and
conclusion are equalities of \emph{partial} values --- the lemma never
manufactures definedness, it only transports whatever is there; in the
theorem, definedness will come out of the equality premise. Second, the
term grammar has no binders, so substitution is plain leaf replacement
and no capture condition is needed anywhere.
\end{proof}

\begin{blemma}[substitution congruence, formula level]\label{lem:eqelim-form}
Under the hypothesis of Lemma~\ref{lem:eqelim-term}: for every formula
$F$,\ $\sigma \models F\sub{a}{t}$ iff $\sigma \models F\sub{a}{u}$.
\end{blemma}
\begin{proof}
Induction on $F$. Comparison atoms substitute in both sides and satisfy
the same relation by Lemma~\ref{lem:eqelim-term}; a region atom
substitutes in its term slot only (its simplex and extension-set data
carry no wires), so its holding --- an assertion about the evaluated
tensor --- transports the same way; $\wedge$ and $\vee$ are the inductive
steps; $\top$ and $\bot$ are untouched by substitution.
\end{proof}

\begin{blemma}[the shape pin]\label{lem:eqelim-pin}
Suppose $\sigma \models t = u$, and suppose $t$ and $u$ carry a common
defined broadcast shape $s$ in $H$. Then
$\mathrm{ev}_\sigma(t) = \mathrm{ev}_\sigma(u)$ --- equality of
\emph{tensors}, not merely the elementwise relation.
\end{blemma}
\begin{proof}
The holding atom yields evaluations $A$ of $t$ and $B$ of $u$ in the
elementwise-equality relation. Broadcast-shape soundness
(\texttt{Term.evalQ\_bshape}) pins $A.\mathrm{shape} = s =
B.\mathrm{shape}$: a defined evaluation of a $\bshape$-typed term lands
exactly at the declared broadcast shape. Elementwise equality at
\emph{equal} shapes is equality of the (equal-length) entry lists,
position by position, hence equality of the tensors. This is the one step
that consumes the typing condition, and it is false without it: the
elementwise relation between a rank-0 tensor and a vector says only that
the scalar matches every entry (\autoref{rem:eqelim-cex}).
\end{proof}

\begin{bthm}[soundness of \rulename{EqElim}]\label{thm:eqelim-sound}
Under the payload's checked hypotheses, if $\Phi \vdash t = u$ and
$\Phi \vdash \Psi\sub{a}{t}$ are semantically valid, then so is
$\Phi \vdash \Psi\sub{a}{u}$.
\end{bthm}
\begin{proof}
Fix an admissible $\sigma$ with $\sigma \models \Phi$. By validity of the
first premise, $\sigma \models t = u$; with the typing condition,
Lemma~\ref{lem:eqelim-pin} upgrades this to
$\mathrm{ev}_\sigma(t) = \mathrm{ev}_\sigma(u)$ (both defined and equal,
so in particular equal as partial values). By validity of the second
premise --- its premise formula is $\Phi$ itself, by the checked premise
match --- $\sigma \models \Psi\sub{a}{t}$. Lemma~\ref{lem:eqelim-form}
transports this to $\sigma \models \Psi\sub{a}{u}$, which is the
conclusion formula by the checked substitution equation.
\end{proof}

\begin{bremark}[why the typing condition is not decoration]
\label{rem:eqelim-cex}
Drop the common-shape condition and the rule is \emph{unsound}
(orientation of the equality atom flipped here to spare an
\rulename{EqSym} step). Let $v$ be a wire of shape $[1]$ over $\Q$, let
$\Delta$ be the segment $\mathrm{conv}(1,2) \subset \Q^1$, and take
\[
\Phi \;:=\; \bigl(v = \tfrac{3}{2}\bigr) \;\wedge\; \bigl(v \in
\mathcal{R}(\Delta, \emptyset)\bigr).
\]
$\Phi$ is satisfiable: set $v := [\,3/2\,]$; the equality holds by rank-0
broadcast (the scalar $3/2$ matches the single entry), and the region
atom holds. \rulename{Ax} and $\wedge$E derive $\Phi \vdash v = 3/2$ and
$\Phi \vdash v \in \mathcal{R}(\Delta, \emptyset)$; the un-gated rule
with body $\Psi := (a \in \mathcal{R}(\Delta, \emptyset))$, $t := v$,
$u := 3/2$ then concludes
$\Phi \vdash (\tfrac{3}{2} : \mathrm{Term}) \in \mathcal{R}(\Delta,
\emptyset)$ --- which is \emph{unsatisfiable}: a constant term evaluates
to a rank-0 tensor, and a region atom demands a rank-1 tensor of extent
equal to the ambient dimension. A satisfiable premise, an unsatisfiable
conclusion. The root cause is precisely the gap between ``elementwise
equal under broadcast'' and ``interchangeable'': $3/2 = [\,3/2\,]$ holds
without the terms being exchangeable in region contexts. With the typing
condition the payload is uninhabitable here ---
$\bshape_H(\tfrac{3}{2}) = [\,]$ while $\bshape_H(v) = [1]$, and no
common $s$ exists --- so the certificate refuses. Note what the repair
does \emph{not} do: broadcast equalities remain derivable and usable as
facts; what is forbidden is exactly using one as a Leibniz license.
\end{bremark}

\begin{bremark}[the equational kit; relation to practice]
\label{rem:eqelim-kit}
\rulename{EqElim} with \rulename{EqRefl} and \rulename{EqSym} (next two
sections) completes a certificate-side equational kit: reflexivity and
symmetry are primitive, and transitivity is a one-node derivation ---
from $\Phi \vdash t = u$ and $\Phi \vdash u = w$, apply \rulename{EqElim}
with body $\Psi := (t = a)$, the hole wire $a$ chosen fresh for $t$
(substitution replaces \emph{every} occurrence of $a$, so a hole
occurring in $t$ would rewrite $t$ itself), rewriting the $u$ in the
\emph{first}
premise (taken as body premise, $\Psi\sub{a}{u}$) to $w$; the derived
node inherits the common-shape typing obligation on $(u, w)$ ---
broadcast equality is not transitive without it (two tensors of unequal
positive-rank shapes never compare, yet each can stand equal to one
scalar between them: $[\,0\,] = 0$ and $0 = [\,0\ \ 0\,]$ both hold
while $[\,0\,] = [\,0\ \ 0\,]$ does not). On the producer side,
equational goals are typically discharged by congruence-closure engines
\cite{nelsonoppen1980,nieuwenhuisoliveras2007}; what ships in a
certificate is not the closure computation but its residue: a chain of
\rulename{EqRefl}/\rulename{EqSym}/\rulename{EqElim} nodes, each checked
by syntactic recomputation of one substitution.
\end{bremark}

\begin{anchors}
Setting: valuation model, partial strict evaluation, atom and region semantics & \texttt{Hypernet.Valuation}, \texttt{Term.evalQ}, \texttt{QTensor.rel2}, \texttt{Formula.holds} (Layer 10)\\
Lemma \ref{lem:eqelim-term} & \texttt{Term.evalQ\_subst\_congr}\\
Lemma \ref{lem:eqelim-form} & \texttt{Formula.holds\_subst\_congr}\\
Lemma \ref{lem:eqelim-pin} & \texttt{Term.evalQ\_bshape}; \texttt{QTensor.eq\_of\_rel2}; payload field \texttt{tyEq}\\
substitution equations & \texttt{substituteInFormula} = \texttt{Formula.subst} (payload fields \texttt{bodySubT}, \texttt{conclSubU})\\
premise sharing & payload fields \texttt{preEq}, \texttt{preMatch}\\
Theorem \ref{thm:eqelim-sound} & \texttt{sound\_eqElim}; wired as core \texttt{RuleData.eqElim} into \texttt{ProofObject.sound}\\
Remark \ref{rem:eqelim-cex} & the counterexample inventory, \texttt{Soundness.lean} header\\
\end{anchors}


\rulesection{9}{EqRefl}{core rule}
\label{blcatrule-EqRefl}
\ruleblock{The rule (Steele metanotation)}
\sourcenote{as realized in the mechanized development; source row: \texttt{\texttt{rules.toml}} --- \emph{EqRefl}. The metanotation statement is typeset by hand for notational quality; its mechanized ground truth is the Lean statement below, extracted from the sources at every build of this book and hash-bound. The rule \emph{roster} is machine-checked (every mechanized rule tag round-trips to exactly one catalogue row); the statement \emph{text} is editorial, checked by eye against the extracted Lean beside it.}
\begin{mathpar}
\inferrule*[Right=\rulename{EqRefl}]{\exists s.\; \bshape_H(t) = s \\ \wires(t) \text{ finitely typed in } H}{\Phi \vdash t = t}
\end{mathpar}

\par\medskip
\rulename{EqRefl} derives $\Phi \vdash t = t$, premise-free, for an
arbitrary context $\Phi$ (a fact true in the empty context is true under
any premise; the proof never reads $\Phi$). A producer uses it to seed
equalities: the starting point of an \rulename{EqElim} rewrite chain
(transitivity, in particular, derives from these three rules and begins
here), and the $t = t$ obligations that alias unpacks and worked
derivations demand.

That the rule has side conditions at all is its content. Equality atoms
in this calculus are facts about \emph{evaluation}, and evaluation is
partial --- a wire of an extended dtype can hold $\pm\infty$ and fail
the rational read; a shape-incompatible sum never evaluates. So
``$t = t$'' costs something: the checker verifies the two typing
conditions the semantics actually needs --- the term\textquotesingle s
broadcast shape is \emph{defined} in the hypernet, and every free wire
is typed at a \emph{finite} dtype. Under exactly these, evaluation is
total at every valuation and reflexivity is free; drop either and the
un-gated rule asserts an invalid sequent under the empty premise --- in
one case a conclusion failing only at extended-valued assignments, in
the other one unsatisfiable outright (the proof section gives one
concrete refutation per field).
Nothing else is consulted: no premise, no admissibility, no evaluation
--- both conditions are static typing checks against the hypernet.

\cataloguenote{Where \rulename{Norm} applies, \rulename{EqRefl} is its diagonal; it is
kept primitive so the equational kit stands in fragments that omit the
polynomial engine entirely.
}
\ruleblock{Lean statement}
\declhead{BernsteinLogic.Soundness}{Term.evalQ\_defined\_of\_bshape\_fin}{\badge{BadgeGreen}{statement bound}}
\begin{leansrc}
/-- Totality of evaluation under the B13 typing pack: a `bshape`-typed
term over finitely-typed wires evaluates. Structural — no polynomial
view, so idx projections are covered (their rank-1 + range gates are
part of `bshape`). -/
theorem Term.evalQ_defined_of_bshape_fin {H : Hypernet} (σ : H.Valuation) :
    ∀ (t : Term) (s : Shape), t.bshape H = some s → t.finWires H →
    ∃ q, Term.evalQ σ t = some q ∧ q.shape = s
\end{leansrc}
\declhead{BernsteinLogic.Soundness}{QTensor.rel2\_refl}{\badge{BadgeGreen}{statement bound}}
\begin{leansrc}
/-- Reflexive relations lift to `rel2` on any tensor against itself. -/
theorem QTensor.rel2_refl {r : ℚ → ℚ → Prop} (hr : ∀ x, r x x)
    (a : QTensor) : QTensor.rel2 r a a
\end{leansrc}
\declhead{BernsteinLogic.Soundness}{sound\_eqRefl}{\badge{BadgeGreen}{statement bound}}
\begin{leansrc}
/-- EqRefl: under the B13 typing pack, t evaluates, and elementwise
equality against itself holds. -/
theorem sound_eqRefl {H : Hypernet} {seq : Sequent} {t : Term}
    (htyDef : ∃ s, t.bshape H = some s) (hfin : t.finWires H)
    (hpost : seq.post = .atom .eq t t) :
    seq.semanticallyValid H
\end{leansrc}
\par\medskip
\textbf{The two typing predicates, in full.}\enspace Both payload fields
are static checks against the hypernet\textquotesingle s wire typing ---
no valuation appears. \texttt{Term.bshape} is the broadcast-tolerant
shape synthesis (scalars join with any shape; the projection case carries
its own rank and range gates); \texttt{Term.finWires} demands a finite
(non-extended) dtype at every free wire.

\declhead{BernsteinLogic.Core.Catalogue}{Term.bshape}{\badge{BadgeGreen}{definition bound}}
\begin{leansrc}
/-- Broadcast-tolerant shape of a term: shapes must agree up to rank-0
broadcast (a scalar joins with any shape); dtypes are NOT consulted. -/
def Term.bshape (H : Hypernet) : Term → Option Shape
  | .var w => (H.wireType w).map (·.shape)
  | .const _ => some []
  | .add l r | .mul l r | .sub l r =>
    match l.bshape H, r.bshape H with
    | some s₁, some s₂ =>
      if s₁ = s₂ then some s₁
      else if s₁ = [] then some s₂
      else if s₂ = [] then some s₁
      else none
    | _, _ => none
  | .neg t => t.bshape H
  | .idx t i =>
    -- rank-1 argument, in-range index, scalar result (matches `evalQ`'s
    -- projection semantics — G1).
    match t.bshape H with
    | some [e] => if i < e then some [] else none
    | _ => none
\end{leansrc}
\declhead{BernsteinLogic.Core.Catalogue}{Term.finWires}{\badge{BadgeGreen}{definition bound}}
\begin{leansrc}
/-- Every wire of the term is typed at a finite (non-ext) dtype — the
carrier of such a wire contains no ±∞, so entry reads are exact
rationals. -/
def Term.finWires (H : Hypernet) (t : Term) : Prop :=
  ∀ w ∈ t.freeWires, ∃ τ, H.wireType w = some τ ∧ τ.dtype.ext = Ext.fin
\end{leansrc}
\declhead{BernsteinLogic.Core.Catalogue}{PEqReflData}{\badge{BadgeGreen}{definition bound}}
\begin{leansrc}
inductive PEqReflData : Hypernet → Sequent → Type where
  | mk (H : Hypernet) (seq : Sequent)
       (term : Term)
       (tyDef  : ∃ s, term.bshape H = some s)
       (fin    : term.finWires H)
       (postEq : seq.post = .atom .eq term term)
       : PEqReflData H seq
\end{leansrc}
\ruleblock{Correspondence (Steele display $\leftrightarrow$ Lean)}
\badge{BadgeGreen}{reviewer-certified against the displayed texts}
\sourcenote{The dictionary below is a reviewer-certified statement of how the mechanized statements above interpret the displayed rule: display metavariable $\mapsto$ Lean binder, each display hypothesis $\mapsto$ the premise or side condition carrying it, the display conclusion $\mapsto$ the goal --- with any deliberate difference in width between display and mechanization named outright. Certification is recorded as a hash triple over (the displayed cell, the extracted statement texts, this dictionary\textquotesingle s own text); if any of the three changes, this block turns amber until a reviewer re-certifies. The correspondence claim itself remains editorial --- reviewer-grade, not machine-checked; its machine-checked successor types each soundness statement against the rule as a structured object, and supersedes this block per rule as it lands.}
Judgment frame as in \rulename{Ax}.
\begin{itemize}
\item $t \mapsto$ the binder \texttt{t : Term}; the two displayed
  typing premises $\mapsto$ \texttt{htyDef}: existence of a shape
  \texttt{s} with \texttt{t.bshape H = some s} (defined broadcast
  shape --- the \texttt{some} is the definedness), and \texttt{hfin}
  $:$ \texttt{t.finWires H} (every free wire typed at a finite dtype,
  quoted below). Both are static checks against the hypernet; no
  valuation appears.
\item conclusion $\Phi \vdash t = t$ $\mapsto$ the goal under
  \texttt{hpost} $:$ \texttt{seq.post = .atom .eq t t}; the arbitrary
  context $\Phi$ is premise-freeness by silence --- no hypothesis
  mentions \texttt{seq.pre}, so the statement holds under every
  premise (the same representation the \rulename{PConstExact}
  dictionary describes).
\item certified alongside as the semantic engine:
  \texttt{Term.evalQ\_defined\_of\_bshape\_fin} --- under exactly the
  two typing conditions, evaluation is total at every valuation, at
  the synthesized shape --- and \texttt{QTensor.rel2\_refl}, the
  reflexivity of the elementwise comparison lifting.
\end{itemize}
No width difference in either direction: both typing hypotheses are
consumed, and the proof section refutes the rule with either one
dropped.

\ruleblock{Lean verifier receipts}
\sourcenote{module \texttt{BernsteinLogic.Soundness}; bound build \texttt{a20d90e66c13\ldots} (bound from receipt a9c0e34d3e3c9c8c, 2026-08-03T21:30:36Z). A receipt records that named checking layers accepted a build: (module, compiled-artifact SHA-256, environment SHA-256, checker layers), minted when the checking gate ran; the store receipts attesting the bound build are listed with the layers each run credited. Nothing here asks the reader to trust the store: every layer is an independently runnable program over the quoted sources --- rebuild the module and re-run the checkers to reproduce the verdict. Receipts attesting superseded builds credit no layer and are omitted from this rendering (the count of omissions, if any, is noted below).}
\begin{receipttable}
\receiptrow{a9c0e34d3e3c9c8c}{2026-08-03T21:30:36Z}{build+decorrelated}{\badge{BadgeGreen}{current}}
\end{receipttable}
\sourcenote{21 further receipt(s) in the store attest superseded builds of this module and are omitted from this rendering.}
\sourcenote{module \texttt{BernsteinLogic.Core.Catalogue}; bound build \texttt{cbc92e976a1d\ldots} (bound from receipt 5abdb3589707b5df, 2026-08-03T21:30:33Z). A receipt records that named checking layers accepted a build: (module, compiled-artifact SHA-256, environment SHA-256, checker layers), minted when the checking gate ran; the store receipts attesting the bound build are listed with the layers each run credited. Nothing here asks the reader to trust the store: every layer is an independently runnable program over the quoted sources --- rebuild the module and re-run the checkers to reproduce the verdict. Receipts attesting superseded builds credit no layer and are omitted from this rendering (the count of omissions, if any, is noted below).}
\begin{receipttable}
\receiptrow{24804e66edba0bc9}{2026-08-03T20:25:09Z}{build+decorrelated}{\badge{BadgeGreen}{current}}
\receiptrow{5abdb3589707b5df}{2026-08-03T21:30:33Z}{build+decorrelated}{\badge{BadgeGreen}{current}}
\end{receipttable}
\sourcenote{18 further receipt(s) in the store attest superseded builds of this module and are omitted from this rendering.}
\layerline{\layeritem{build}{lake build (elaboration)}{\badge{BadgeGreen}{current}}\layeritem{replay}{leanchecker (official kernel replay)}{\badge{BadgeAmber}{stale only}}\layeritem{hardened}{safe\_verify (hardened screen)}{\badge{BadgeAmber}{stale only}}\layeritem{decorrelated}{lean4export + nanoda (decorrelated kernel)}{\badge{BadgeGreen}{current}}}
\sourcenote{\badge{BadgeGray}{comparator: none} no fixture in the three-implementation comparator harness (compiled / kernel-interpreted / reference-interpreted) exercises this rule.}
\ruleblock{HOL4 statement}
\sourcenote{theory \texttt{bern\_fam\_resid}, HOL4 development script \texttt{\texttt{bern\_fam\_residScript.sml}} \badge{BadgeGreen}{script bound}}
\holhead{Definition checkFam\_eqrefl\_def}
\begin{holsrc}
Definition checkFam_eqrefl_def:
  checkFam_eqrefl (cols:columns) (ct:ctx) (i:num) ⇔
    case frSeq cols ct i of
    | NONE => F
    | SOME sq =>
      case frTerm cols ct i 0 of
      | NONE => F
      | SOME t =>
          IS_SOME (termBshape ct.ws t) ∧ finWiresB ct.ws t ∧
          sq.post = FAtom BowEq t t
End
\end{holsrc}
\holhead{Theorem famEq\_eqrefl}
\begin{holsrc}
Theorem famEq_eqrefl:
  wfColumns cols ct ∧ mkCtx cols = SOME ct ∧
  parseCert (cols.hdrImage ++ rest) = SOME (ws, pf) ∧
  nodeCorr cols ct i (CertPf sq nd) ∧
  FST (THE (tagOf cols i)) = 23 ⇒
  (checkFam_eqrefl cols ct i ⇔ armEqRefl ws sq nd)
\end{holsrc}
\ruleblock{HOL verifier receipts}
\badge{BadgeBlue}{documentary} Documentary citations. Family validation: \texttt{bern\_fam\_resid} built GREEN with all ten \texttt{famEq} arms validated, \texttt{hol-strict} clean; spine census: \texttt{bern\_joins} GREEN under \texttt{hol-strict} (55 theorems, 0 offences, 208-theory ancestry including this theory), \texttt{checkJoins\_eq\_checkPf} verified verbatim in the built signature. The cited \texttt{famEq} lemma is the kernel-checked per-node equivalence: under well-formed columns (\texttt{wfColumns}/\texttt{mkCtx}/\texttt{nodeCorr}) at this rule\textquotesingle s wire tag, the columnar family check coincides with the flat reference walk\textquotesingle s arm for this rule. There is no signed HOL receipt store yet --- these citations are documentary.
Semantic stratum, also ported: the soundness twin \texttt{sound\_eqRefl}
(theory \texttt{bern\_tr\_ready8}: \texttt{hol-strict} clean, 8
theorems, 0 offences, verification run of 2026-07-22 on an isolated
build host, content-addressed at the same sources), consumed by the
certified-fragment keystone (\texttt{checkPf\_sound\_frag}, theory
\texttt{bern\_tr\_splice}, 20 theorems, 0 offences, same run): a
reference-checker accept inside the proven wire fragment yields sequent
validity.

\sourcenote{HOL receipts are documentary: quoted results of \texttt{hol-strict} runs
and census rows from the HOL4 development\textquotesingle s build
evidence, not minted receipt artefacts. A HOL receipt store (the Lean
store\textquotesingle s twin) is named future work.
}
\ruleblock{Proof}

\sourcenote{Genre: core rule (tier 1). Setting (valuations, partial
evaluation, elementwise-under-broadcast atom semantics) as fixed in the
\rulename{EqElim} section.}

\textbf{Claim} (\texttt{sound\_eqRefl}).\enspace The payload supplies a
term $t$ and two checked typing conditions: $\bshape_H(t)$ is
\emph{defined} (some shape $s$), and every free wire of $t$ is typed at a
\emph{finite} (non-extended) dtype. Then $\Phi \vdash t = t$ is
semantically valid for an \emph{arbitrary} premise $\Phi$ --- the rule is
premise-free, and $\Phi$ rides along unread (a fact true in the empty
context is true under any premise). Indeed the proof below never reads
$\Phi$, the admissibility of $\sigma$, or the premise-satisfaction
hypothesis at all: the conclusion holds at \emph{every} valuation.

That anything must be checked here at all is the point of the section.
Equality atoms are evaluation-level facts (\rulename{EqElim} setting), so
``$t = t$'' is not free: it asserts that $t$ \emph{evaluates} --- and
evaluation is partial.

\begin{blemma}[totality under the typing pack]\label{lem:eqrefl-total}
If $\bshape_H(t) = s$ and every free wire of $t$ is finitely typed, then
at every valuation $\sigma$ the evaluation $\mathrm{ev}_\sigma(t)$ is
defined, with shape $s$.
\end{blemma}
\begin{proof}
Structural induction on $t$; both typing conditions descend to
subterms (a subterm's free wires are among $t$'s, and $\bshape$ is
synthesized bottom-up, so a defined $\bshape$ for $t$ is defined on
every subterm). A wire read succeeds because the wire is typed at a
finite dtype --- an untyped wire declares no dtype at all, so the
finite-wires condition already excludes it: the valuation's entries lie
in the dtype's value set, which for a finite dtype contains no
$\pm\infty$, so the whole-tensor rational read succeeds, at the
declared shape. A constant evaluates to its rank-0 tensor. For a
compound term, the induction hypotheses give defined operands whose
shapes are, by the definition of $\bshape$ (quoted in full beside this
section), either equal or joinable through rank-0 broadcast --- and on
such shapes the elementwise operations are total; the result lands at
the join. A projection's rank-1-and-in-range gates are part of
$\bshape$'s definedness, so the projected read is in range.
(\texttt{Term.evalQ\_defined\_of\_bshape\_fin}.)
\end{proof}

\begin{blemma}[reflexivity lifts]\label{lem:eqrefl-lift}
For any tensor $A$, the elementwise-under-broadcast equality relation
holds between $A$ and $A$.
\end{blemma}
\begin{proof}
The shapes are equal, so the relation reads the entry list zipped against
itself: every pair on the diagonal is $(x, x)$.
(\texttt{QTensor.rel2\_refl}.)
\end{proof}

\begin{bthm}[soundness of \rulename{EqRefl}]\label{thm:eqrefl-sound}
Under the payload's typing conditions, $\Phi \vdash t = t$ is
semantically valid for every $\Phi$.
\end{bthm}
\begin{proof}
Fix $\sigma$. Lemma~\ref{lem:eqrefl-total} gives
$\mathrm{ev}_\sigma(t) = A$ for some tensor $A$;
Lemma~\ref{lem:eqrefl-lift} closes the atom with both sides read as $A$.
\end{proof}

\begin{bremark}[why each typing condition is load-bearing]
\label{rem:eqrefl-cex}
Drop the finite-wires condition: let $w$ be a wire of an
\emph{extended} dtype (its value set admits $+\infty$), and let $\sigma$
assign it a tensor with an infinite entry --- admissible, since nothing
constrains a bare wire. The rational read of $w$ fails, $t := w$ does not
evaluate, and the atom $w = w$ is \emph{false} at $\sigma$; with
$\Phi := \top$ the sequent $\top \vdash w = w$ is refuted. Drop
shape-definedness: take $t := w_1 + w_2$ with $w_1$ of shape $[2]$ and
$w_2$ of shape $[3]$, both finite --- every wire reads, but the
elementwise sum refuses the incompatible shapes, $t$ never evaluates, and
$\vdash t = t$ is again false everywhere. In both cases the un-gated rule
would assert an invalid sequent under a satisfiable (indeed empty)
premise --- in the first case the conclusion fails at the exhibited
valuation (a finite assignment would satisfy it); in the second it is
unsatisfiable outright. The typing pack is exactly the semantic price of
equality-by-evaluation; it is the same field family that repairs
\rulename{Norm}, there compounded with canonical-form equality.
\end{bremark}

\begin{bremark}[taxonomy note]\label{rem:eqrefl-tax}
Wherever \rulename{Norm} applies, \rulename{EqRefl} is its diagonal
($t_1 = t_2 = t$, canonical forms trivially equal). It is nonetheless
kept primitive: the development's frozen minimal fragment (\texttt{v0} in
the sources) excludes the polynomial engine outright, and the equational
kit must stand without it.
\end{bremark}

\begin{anchors}
Lemma \ref{lem:eqrefl-total} & \texttt{Term.evalQ\_defined\_of\_bshape\_fin}; payload fields \texttt{tyDef}, \texttt{fin}\\
finite-wire reads succeed & \texttt{Hypernet.Valuation.wireQ\_isSome}\\
Lemma \ref{lem:eqrefl-lift} & \texttt{QTensor.rel2\_refl}\\
Theorem \ref{thm:eqrefl-sound} & \texttt{sound\_eqRefl}; wired as core \texttt{RuleData.eqRefl} into \texttt{ProofObject.sound}\\
typing predicates & \texttt{Term.bshape}, \texttt{Term.finWires} (Core/Catalogue)\\
\end{anchors}


\rulesection{10}{EqSym}{core rule}
\label{blcatrule-EqSym}
\ruleblock{The rule (Steele metanotation)}
\sourcenote{as realized in the mechanized development; source row: \texttt{\texttt{rules.toml}} --- \emph{EqSym}. The metanotation statement is typeset by hand for notational quality; its mechanized ground truth is the Lean statement below, extracted from the sources at every build of this book and hash-bound. The rule \emph{roster} is machine-checked (every mechanized rule tag round-trips to exactly one catalogue row); the statement \emph{text} is editorial, checked by eye against the extracted Lean beside it.}
\begin{mathpar}
\inferrule*[Right=\rulename{EqSym}]{\Phi \vdash t = u}{\Phi \vdash u = t}
\end{mathpar}

\par\medskip
\rulename{EqSym} flips an equality atom: from $\Phi \vdash t = u$
conclude $\Phi \vdash u = t$. The need is real because downstream
consumers of equalities match atoms \emph{syntactically} with a fixed
orientation --- \rulename{EqElim}\textquotesingle s equality premise, in
particular, names its two terms in order --- while only the combination
rules expose both polarities of an equality in their own matching. A
producer whose congruence engine emits $t = u$ but whose rewrite needs
$u = t$ inserts this one node.

What the checker recomputes: three syntactic form equations (premise
atom, conclusion atom, shared context) --- nothing else. The rule
carries \emph{no} typing fields, and that absence is principled twice
over: the premise\textquotesingle s holding atom already witnesses both
evaluations (nothing must be manufactured, contrast
\rulename{EqRefl}), and the elementwise-under-broadcast equality
relation is symmetric branch by branch, so no shape agreement is needed
(contrast \rulename{EqElim}, whose typing condition guards transport
into a \emph{different} formula --- exactly what a flip never does).

Every payload field is structural; there is no droppable hypothesis to
exhibit a counterexample against. The rule is core rather than
admissible because its would-be unpack (through
\rulename{EqRefl} + \rulename{EqElim}) demands typing side data the
premise atom does not syntactically carry --- the unpack is not
computable from the node it would replace.

\ruleblock{Lean statement}
\declhead{BernsteinLogic.Soundness}{QTensor.rel2\_symm}{\badge{BadgeGreen}{statement bound}}
\begin{leansrc}
/-- `rel2` of a symmetric relation is symmetric (all three broadcast
branches flip pointwise). -/
theorem QTensor.rel2_symm {r : ℚ → ℚ → Prop}
    (hr : ∀ x y, r x y → r y x) {a b : QTensor}
    (h : QTensor.rel2 r a b) : QTensor.rel2 r b a
\end{leansrc}
\declhead{BernsteinLogic.Soundness}{sound\_eqSym}{\badge{BadgeGreen}{statement bound}}
\begin{leansrc}
/-- EqSym: flip the premise atom's two evaluations and the symmetric
`rel2`. -/
theorem sound_eqSym {H : Hypernet} {seq premise : Sequent} {tT tU : Term}
    (ih : premise.semanticallyValid H)
    (hpre : seq.pre = premise.pre)
    (hpremPost : premise.post = .atom .eq tT tU)
    (hpost : seq.post = .atom .eq tU tT) :
    seq.semanticallyValid H
\end{leansrc}
\declhead{BernsteinLogic.Core.Catalogue}{PEqSymData}{\badge{BadgeGreen}{definition bound}}
\begin{leansrc}
inductive PEqSymData : Hypernet → Sequent → Type where
  | mk (H : Hypernet) (seq : Sequent)
       (premise : Sequent) (premiseProof : ProofObject H premise)
       (termT termU : Term)
       (preEq       : seq.pre = premise.pre)
       (premisePost : premise.post = .atom .eq termT termU)
       (postEq      : seq.post = .atom .eq termU termT)
       : PEqSymData H seq
\end{leansrc}
\ruleblock{Correspondence (Steele display $\leftrightarrow$ Lean)}
\badge{BadgeGreen}{reviewer-certified against the displayed texts}
\sourcenote{The dictionary below is a reviewer-certified statement of how the mechanized statements above interpret the displayed rule: display metavariable $\mapsto$ Lean binder, each display hypothesis $\mapsto$ the premise or side condition carrying it, the display conclusion $\mapsto$ the goal --- with any deliberate difference in width between display and mechanization named outright. Certification is recorded as a hash triple over (the displayed cell, the extracted statement texts, this dictionary\textquotesingle s own text); if any of the three changes, this block turns amber until a reviewer re-certifies. The correspondence claim itself remains editorial --- reviewer-grade, not machine-checked; its machine-checked successor types each soundness statement against the rule as a structured object, and supersedes this block per rule as it lands.}
Judgment frame as in \rulename{Ax}; sub-derivation premise as at
\rulename{Weaken}.
\begin{itemize}
\item $t, u \mapsto$ \texttt{tT tU : Term}; the premise judgment $\Phi
  \vdash t = u$ $\mapsto$ \texttt{ih} with the shape tie
  \texttt{hpremPost} $:$ \texttt{premise.post = .atom .eq tT tU};
  $\Phi \mapsto$ \texttt{premise.pre}, shared via \texttt{hpre}.
\item conclusion $\Phi \vdash u = t$ $\mapsto$ the goal under
  \texttt{hpost} $:$ \texttt{seq.post = .atom .eq tU tT} --- the same
  two term binders, transposed; nothing else changes.
\item certified alongside: \texttt{QTensor.rel2\_symm}, the branchwise
  symmetry of the elementwise comparison lifting --- the reason no
  typing field exists (the premise\textquotesingle s holding atom
  already witnesses both evaluations, as the commentary explains).
\end{itemize}
No width difference in either direction; every hypothesis is consumed.

\ruleblock{Lean verifier receipts}
\sourcenote{module \texttt{BernsteinLogic.Soundness}; bound build \texttt{a20d90e66c13\ldots} (bound from receipt a9c0e34d3e3c9c8c, 2026-08-03T21:30:36Z). A receipt records that named checking layers accepted a build: (module, compiled-artifact SHA-256, environment SHA-256, checker layers), minted when the checking gate ran; the store receipts attesting the bound build are listed with the layers each run credited. Nothing here asks the reader to trust the store: every layer is an independently runnable program over the quoted sources --- rebuild the module and re-run the checkers to reproduce the verdict. Receipts attesting superseded builds credit no layer and are omitted from this rendering (the count of omissions, if any, is noted below).}
\begin{receipttable}
\receiptrow{a9c0e34d3e3c9c8c}{2026-08-03T21:30:36Z}{build+decorrelated}{\badge{BadgeGreen}{current}}
\end{receipttable}
\sourcenote{21 further receipt(s) in the store attest superseded builds of this module and are omitted from this rendering.}
\sourcenote{module \texttt{BernsteinLogic.Core.Catalogue}; bound build \texttt{cbc92e976a1d\ldots} (bound from receipt 5abdb3589707b5df, 2026-08-03T21:30:33Z). A receipt records that named checking layers accepted a build: (module, compiled-artifact SHA-256, environment SHA-256, checker layers), minted when the checking gate ran; the store receipts attesting the bound build are listed with the layers each run credited. Nothing here asks the reader to trust the store: every layer is an independently runnable program over the quoted sources --- rebuild the module and re-run the checkers to reproduce the verdict. Receipts attesting superseded builds credit no layer and are omitted from this rendering (the count of omissions, if any, is noted below).}
\begin{receipttable}
\receiptrow{24804e66edba0bc9}{2026-08-03T20:25:09Z}{build+decorrelated}{\badge{BadgeGreen}{current}}
\receiptrow{5abdb3589707b5df}{2026-08-03T21:30:33Z}{build+decorrelated}{\badge{BadgeGreen}{current}}
\end{receipttable}
\sourcenote{18 further receipt(s) in the store attest superseded builds of this module and are omitted from this rendering.}
\layerline{\layeritem{build}{lake build (elaboration)}{\badge{BadgeGreen}{current}}\layeritem{replay}{leanchecker (official kernel replay)}{\badge{BadgeAmber}{stale only}}\layeritem{hardened}{safe\_verify (hardened screen)}{\badge{BadgeAmber}{stale only}}\layeritem{decorrelated}{lean4export + nanoda (decorrelated kernel)}{\badge{BadgeGreen}{current}}}
\sourcenote{\badge{BadgeGray}{comparator: none} no fixture in the three-implementation comparator harness (compiled / kernel-interpreted / reference-interpreted) exercises this rule.}
\ruleblock{HOL4 statement}
\sourcenote{theory \texttt{bern\_fam\_resid}, HOL4 development script \texttt{\texttt{bern\_fam\_residScript.sml}} \badge{BadgeGreen}{script bound}}
\holhead{Definition checkFam\_eqsym\_def}
\begin{holsrc}
Definition checkFam_eqsym_def:
  checkFam_eqsym (cols:columns) (ct:ctx) (i:num) ⇔
    case frSeq cols ct i of
    | NONE => F
    | SOME sq =>
      case frChildSeq cols ct i 0 of
      | NONE => F
      | SOME ps =>
        case frTerm cols ct i 0 of
        | NONE => F
        | SOME tT =>
          case frTerm cols ct i 1 of
          | NONE => F
          | SOME tU =>
              sq.pre = ps.pre ∧
              ps.post = FAtom BowEq tT tU ∧
              sq.post = FAtom BowEq tU tT
End
\end{holsrc}
\holhead{Theorem famEq\_eqsym}
\begin{holsrc}
Theorem famEq_eqsym:
  wfColumns cols ct ∧ mkCtx cols = SOME ct ∧
  parseCert (cols.hdrImage ++ rest) = SOME (ws, pf) ∧
  nodeCorr cols ct i (CertPf sq nd) ∧
  FST (THE (tagOf cols i)) = 24 ⇒
  (checkFam_eqsym cols ct i ⇔ armEqSym ws sq nd)
\end{holsrc}
\ruleblock{HOL verifier receipts}
\badge{BadgeBlue}{documentary} Documentary citations. Family validation: \texttt{bern\_fam\_resid} built GREEN with all ten \texttt{famEq} arms validated, \texttt{hol-strict} clean; spine census: \texttt{bern\_joins} GREEN under \texttt{hol-strict} (55 theorems, 0 offences, 208-theory ancestry including this theory), \texttt{checkJoins\_eq\_checkPf} verified verbatim in the built signature. The cited \texttt{famEq} lemma is the kernel-checked per-node equivalence: under well-formed columns (\texttt{wfColumns}/\texttt{mkCtx}/\texttt{nodeCorr}) at this rule\textquotesingle s wire tag, the columnar family check coincides with the flat reference walk\textquotesingle s arm for this rule. There is no signed HOL receipt store yet --- these citations are documentary.
Semantic stratum, also ported: the soundness twin \texttt{sound\_eqSym}
(theory \texttt{bern\_tr\_b14}: \texttt{hol-strict} clean, 2 theorems, 0
offences, verification run of 2026-07-22 on an isolated build host,
content-addressed at the same sources), consumed by the certified-
fragment keystone (\texttt{checkPf\_sound\_frag}, theory
\texttt{bern\_tr\_splice}, 20 theorems, 0 offences, same run): a
reference-checker accept inside the proven wire fragment yields sequent
validity.

\sourcenote{HOL receipts are documentary: quoted results of \texttt{hol-strict} runs
and census rows from the HOL4 development\textquotesingle s build
evidence, not minted receipt artefacts. A HOL receipt store (the Lean
store\textquotesingle s twin) is named future work.
}
\ruleblock{Proof}

\sourcenote{Genre: core rule (tier 1). Setting as fixed in the
\rulename{EqElim} section.}

\textbf{Claim} (\texttt{sound\_eqSym}).\enspace The payload supplies
terms $t, u$; the premise concludes the atom $t = u$ over premise formula
$\Phi$, shared with the conclusion sequent; the conclusion is the atom
$u = t$. If the premise is semantically valid, so is
$\Phi \vdash u = t$.

\begin{blemma}[the elementwise relation of a symmetric relation is
symmetric]\label{lem:eqsym-rel2}
Let $r$ be a symmetric binary relation on $\Q$, and suppose tensors $A$,
$B$ stand in the elementwise-under-broadcast lift of $r$. Then $B$, $A$
stand in it too.
\end{blemma}
\begin{proof}
The lift has three live branches, and each is separately symmetric ---
this per-branch analysis is the entire content, so we spell it out.
\emph{Equal shapes}: the relation reads the zipped entry lists; the
swapped claim reads the zip in the other order, and each swapped pair is
closed by symmetry of $r$. \emph{Rank-0 $A$ against positive-rank $B$}:
the hypothesis says $r(A_{\mathrm{scalar}}, y)$ for every entry $y$ of
$B$; the swapped claim is judged in the \emph{other} scalar-broadcast
branch (rank-0 on the right), which demands $r(y, A_{\mathrm{scalar}})$
for every $y$ --- symmetry of $r$, pointwise. \emph{Rank-0 $B$}: mirror
image. Two unequal positive-rank shapes never stand in the lift, so
there is nothing to prove there. (\texttt{QTensor.rel2\_symm}.)
\end{proof}

\begin{bthm}[soundness of \rulename{EqSym}]\label{thm:eqsym-sound}
If $\Phi \vdash t = u$ is semantically valid, so is $\Phi \vdash u = t$.
\end{bthm}
\begin{proof}
Fix an admissible $\sigma \models \Phi$. The premise yields evaluations
$A$ of $t$ and $B$ of $u$ in the elementwise equality lift. Present the
same two evaluations in the opposite roles: $u$ evaluates to $B$, $t$ to
$A$, and Lemma~\ref{lem:eqsym-rel2} (with $r$ the symmetric relation
$=$) closes the atom $u = t$.
\end{proof}

\begin{bremark}[the negative space: what this rule does \emph{not} need]
\label{rem:eqsym-neg}
The rule has no droppable side condition --- every payload field is
structural (the three form equations) --- so the honest remarks here
concern the two fields conspicuously absent. No typing pack
(\rulename{EqRefl}'s \texttt{tyDef}/\texttt{fin}): the premise's holding
atom already \emph{witnesses} both evaluations, and symmetry only re-uses
the premise's own evidence; nothing must be manufactured. No common-shape
condition (\rulename{EqElim}'s \texttt{tyEq}): Lemma~\ref{lem:eqsym-rel2}
is shape-agnostic --- a broadcast equality flips to the mirror broadcast
equality, same content, no interchangeability claim ever made. The
contrast triangulates \rulename{EqElim}'s repair: the typing condition
there guards \emph{transport into a different formula}, which is exactly
what \rulename{EqSym} never does.
\end{bremark}

\begin{bremark}[why a primitive, when a two-node unpack exists on paper]
\label{rem:eqsym-adm}
From $\Phi \vdash t = u$ one can derive $\Phi \vdash u = t$ through
\rulename{EqRefl} + \rulename{EqElim}: body $\Psi := (a = t)$ with the
hole wire $a$ fresh for $t$ --- substitution replaces \emph{every}
occurrence of $a$, and freshness is what makes $\Psi\sub{a}{t}$ the atom
$t = t$ and $\Psi\sub{a}{u}$ the atom $u = t$ --- rewriting
the left occurrence of $t$ in $t = t$ to $u$. But the derivation demands
\rulename{EqRefl}'s typing pack on $t$ \emph{and} \rulename{EqElim}'s
common-shape condition on $(t, u)$ --- side data the premise atom does
not syntactically carry, and which a producer cannot always reconstruct
at the flip site. The primitive constructor instead consumes the
premise's semantic evidence directly (the proof above), making the flip
unconditional. In the catalogue's three-tier taxonomy that obstruction is
what sorts \rulename{EqSym} core rather than admissible: the unpack is
not payload-computable from the node it would replace.
\end{bremark}

\begin{anchors}
Lemma \ref{lem:eqsym-rel2} & \texttt{QTensor.rel2\_symm}\\
Theorem \ref{thm:eqsym-sound} & \texttt{sound\_eqSym}; wired as core \texttt{RuleData.eqSym} into \texttt{ProofObject.sound}\\
form equations & payload fields \texttt{preEq}, \texttt{premisePost}, \texttt{postEq} (\texttt{PEqSymData})\\
\end{anchors}


\rulesection{11}{Antisym}{core rule}
\label{blcatrule-Antisym}
\ruleblock{The rule (Steele metanotation)}
\sourcenote{as realized in the mechanized development; source row: \texttt{\texttt{rules.toml}} --- \emph{Antisym}. The metanotation statement is typeset by hand for notational quality; its mechanized ground truth is the Lean statement below, extracted from the sources at every build of this book and hash-bound. The rule \emph{roster} is machine-checked (every mechanized rule tag round-trips to exactly one catalogue row); the statement \emph{text} is editorial, checked by eye against the extracted Lean beside it.}
\begin{mathpar}
\inferrule*[Right=\rulename{Antisym}]{\Phi \vdash t \leq u \\ \Phi \vdash t \geq u}{\Phi \vdash t = u}
\end{mathpar}

\par\medskip
\rulename{Antisym} squeezes an equality out of two opposite bounds: from
$\Phi \vdash t \leq u$ and $\Phi \vdash t \geq u$ --- the same two terms,
the same context --- conclude $\Phi \vdash t = u$. A producer reaches for
it wherever a pinching argument closes: an optimization pass that
certifies matching upper and lower bounds on the same quantity, a
fixed-point equality assembled from two directional inequalities, or any
pipeline whose native outputs are order facts but whose consumer
(\rulename{EqElim}\textquotesingle s rewrite chains, in particular) wants
an equality atom.

The rule earns its place in the core by a non-derivability fact, not by
convenience: equality-from-inequalities has no derivation in the
remainder of the calculus. Provably --- an induction over derivations
shows that every equality concluded from an equality-atom-free context
equates terms with identical polynomial normal forms --- so no chain of
existing rules can ever conclude $t = u$ from the two bounds when the
sides are genuinely different terms. Historically the calculus carried
this inference only inside two bespoke rules for the ReLU function,
which were thereby load-bearing for something much more general than
ReLU; with antisymmetry available as its own rule, that monopoly ends.

The side conditions are purely structural: two premise-atom form
equations, the conclusion-atom form equation, and the two context
equations tying each premise sequent to the conclusion\textquotesingle s
context. No typing fields --- and the absence is principled (contrast
\rulename{EqRefl}, which must \emph{manufacture} an evaluation and so
carries the typing pack): each premise atom holds only where both sides
already evaluate finitely at compatible shapes, and that is exactly the
evaluation evidence the equality atom\textquotesingle s semantics
consumes. Dropping either premise leaves a one-sided bound, and any
strictly-below evaluation pair refutes the un-gated rule.

One reading discipline, stated here because the conclusion invites
over-reading: the concluded equality holds in the calculus\textquotesingle
s elementwise-under-broadcast sense, which suffices for the truth of the
atom and for every consumer that reads it semantically --- but it is not
a license to interchange $t$ and $u$ inside an arbitrary formula.
Substitution rides \rulename{EqElim}, whose common-defined-shape typing
condition remains in force unchanged (that section constructs the
refutation its absence would admit).

\cataloguenote{This rule is \emph{staged, not enacted}: what exists today is the
statement register --- the soundness theorem below, proven in Lean with
a HOL twin --- while the certificate format knows nothing of it yet. No
payload constructor, no wire tag, no checker arm: a certificate cannot
currently carry an \rulename{Antisym} node, and this section will gain
its checker-side anchors when enactment lands. The mathematics is
settled; the plumbing is queued.
}
\ruleblock{Lean statement}
\declhead{BernsteinLogic.Soundness}{QTensor.rel2\_conj}{\badge{BadgeGreen}{statement bound}}
\begin{leansrc}
/-- Two `rel2` facts at the SAME tensor pair combine pointwise: the
broadcast branch taken is a function of the two shapes alone, so the
quantified domains coincide and the pointwise conjunction lifts.
(Stated for `sound_antisym`; any pointwise combination rides it.) -/
theorem QTensor.rel2_conj {r s t : ℚ → ℚ → Prop}
    (hrst : ∀ x y, r x y → s x y → t x y) {a b : QTensor}
    (hr : QTensor.rel2 r a b) (hs : QTensor.rel2 s a b) :
    QTensor.rel2 t a b
\end{leansrc}
\declhead{BernsteinLogic.Soundness}{sound\_antisym}{\badge{BadgeGreen}{statement bound}}
\begin{leansrc}
/-- Antisym: `≤` and `≥` premise validities at the same term pair force
the equality atom — entrywise `le_antisymm` lifted through `rel2`
(`QTensor.rel2_conj`: both premises constrain the same evaluated pair,
so the broadcast branches align). Like `sound_eqSym` and unlike
`sound_eqRefl`, no B13 typing pack is needed: the premise atoms only
hold where both sides already evaluate finitely, which is exactly the
evaluation evidence the conclusion atom needs. Core rather than
derived: the equality-from-inequalities inference has no catalogue
derivation — provably, since derivable equalities from eq-free contexts
have canonically-equal sides — so the rule must be primitive
(ReLU-Pos/Neg had been core solely as carriers of this inference). -/
theorem sound_antisym {H : Hypernet} {seq premiseLe premiseGe : Sequent}
    {tT tU : Term}
    (ihLe : premiseLe.semanticallyValid H)
    (ihGe : premiseGe.semanticallyValid H)
    (hpreLe : seq.pre = premiseLe.pre)
    (hpreGe : seq.pre = premiseGe.pre)
    (hpremLePost : premiseLe.post = .atom .le tT tU)
    (hpremGePost : premiseGe.post = .atom .ge tT tU)
    (hpost : seq.post = .atom .eq tT tU) :
    seq.semanticallyValid H
\end{leansrc}
\ruleblock{Correspondence (Steele display $\leftrightarrow$ Lean)}
\badge{BadgeGreen}{reviewer-certified against the displayed texts}
\sourcenote{The dictionary below is a reviewer-certified statement of how the mechanized statements above interpret the displayed rule: display metavariable $\mapsto$ Lean binder, each display hypothesis $\mapsto$ the premise or side condition carrying it, the display conclusion $\mapsto$ the goal --- with any deliberate difference in width between display and mechanization named outright. Certification is recorded as a hash triple over (the displayed cell, the extracted statement texts, this dictionary\textquotesingle s own text); if any of the three changes, this block turns amber until a reviewer re-certifies. The correspondence claim itself remains editorial --- reviewer-grade, not machine-checked; its machine-checked successor types each soundness statement against the rule as a structured object, and supersedes this block per rule as it lands.}
Judgment frame as in \rulename{Ax}; sub-derivation premises as at
\rulename{Weaken}.
\begin{itemize}
\item $t, u \mapsto$ \texttt{tT tU : Term}; the two premise judgments
  $\Phi \vdash t \leq u$ and $\Phi \vdash t \geq u$ $\mapsto$
  \texttt{ihLe}, \texttt{ihGe} with the shape ties \texttt{hpremLePost}
  $:$ \texttt{premiseLe.post = .atom .le tT tU} and \texttt{hpremGePost}
  $:$ \texttt{premiseGe.post = .atom .ge tT tU} --- note the $\geq$
  premise binds its own atom former at the \emph{same} binder order
  (\texttt{tT}, \texttt{tU}), not a flipped $\leq$ atom.
\item the shared context: the display writes one $\Phi$ across all
  three sequents; the statement renders it as \emph{two} equations,
  \texttt{hpreLe} $:$ \texttt{seq.pre = premiseLe.pre} and
  \texttt{hpreGe} $:$ \texttt{seq.pre = premiseGe.pre} --- each premise
  tied to the conclusion\textquotesingle s context, premise-to-premise
  equality following by transitivity. Same content, one equation more
  than the display suggests.
\item conclusion $\Phi \vdash t = u$ $\mapsto$ the goal under
  \texttt{hpost} $:$ \texttt{seq.post = .atom .eq tT tU}.
\item no typing hypotheses appear on either side --- principled, as the
  commentary explains (the premise atoms carry the evaluation evidence;
  contrast \rulename{EqRefl}\textquotesingle s \texttt{htyDef}/\texttt{hfin}).
\item certified alongside as the semantic engine:
  \texttt{QTensor.rel2\_conj} --- two elementwise-under-broadcast facts
  at the \emph{same} tensor pair combine pointwise, because the
  broadcast branch taken is a function of the two shapes alone, so the
  quantified domains coincide. The soundness proof feeds it entrywise
  antisymmetry of the rational order.
\end{itemize}
Width note, both directions: every hypothesis is consumed, and the one
place statement and display could be read apart is the conclusion
atom\textquotesingle s broadcast-sense semantics --- the proven equality
is the elementwise-under-broadcast lift (a scalar can stand equal to a
vector whose every entry it matches), which is the equality
atom\textquotesingle s meaning everywhere in this calculus; the
commentary\textquotesingle s reading discipline names what that does
\emph{not} license.

\ruleblock{Lean verifier receipts}
\sourcenote{module \texttt{BernsteinLogic.Soundness}; bound build \texttt{a20d90e66c13\ldots} (bound from receipt a9c0e34d3e3c9c8c, 2026-08-03T21:30:36Z). A receipt records that named checking layers accepted a build: (module, compiled-artifact SHA-256, environment SHA-256, checker layers), minted when the checking gate ran; the store receipts attesting the bound build are listed with the layers each run credited. Nothing here asks the reader to trust the store: every layer is an independently runnable program over the quoted sources --- rebuild the module and re-run the checkers to reproduce the verdict. Receipts attesting superseded builds credit no layer and are omitted from this rendering (the count of omissions, if any, is noted below).}
\begin{receipttable}
\receiptrow{a9c0e34d3e3c9c8c}{2026-08-03T21:30:36Z}{build+decorrelated}{\badge{BadgeGreen}{current}}
\end{receipttable}
\sourcenote{21 further receipt(s) in the store attest superseded builds of this module and are omitted from this rendering.}
\layerline{\layeritem{build}{lake build (elaboration)}{\badge{BadgeGreen}{current}}\layeritem{replay}{leanchecker (official kernel replay)}{\badge{BadgeAmber}{stale only}}\layeritem{hardened}{safe\_verify (hardened screen)}{\badge{BadgeAmber}{stale only}}\layeritem{decorrelated}{lean4export + nanoda (decorrelated kernel)}{\badge{BadgeGreen}{current}}}
\sourcenote{\badge{BadgeGray}{comparator: none} no fixture in the three-implementation comparator harness (compiled / kernel-interpreted / reference-interpreted) exercises this rule.}
\ruleblock{HOL4 statement}
\sourcenote{theory \texttt{bern\_tr\_antisym}, HOL4 development script \texttt{\texttt{bern\_tr\_antisymScript.sml}} \badge{BadgeGreen}{script bound}}
\holhead{Theorem qtRel2\_conj}
\begin{holsrc}
Theorem qtRel2_conj:
  ∀r s t a b.
    (∀x y. r x y ∧ s x y ⇒ t x y) ∧
    qtRel2 r a b ∧ qtRel2 s a b ⇒
    qtRel2 t a b
\end{holsrc}
\holhead{Theorem sound\_antisym}
\begin{holsrc}
Theorem sound_antisym:
  ∀H sq premiseLe premiseGe tT tU.
    seqSemanticallyValid H premiseLe ∧
    seqSemanticallyValid H premiseGe ∧
    sq.pre = premiseLe.pre ∧
    sq.pre = premiseGe.pre ∧
    premiseLe.post = FAtom BowLe tT tU ∧
    premiseGe.post = FAtom BowGe tT tU ∧
    sq.post = FAtom BowEq tT tU ⇒
    seqSemanticallyValid H sq
\end{holsrc}
\ruleblock{HOL verifier receipts}
\badge{BadgeBlue}{documentary} Documentary citations, one register-honest caveat: this theory is the
staged rule\textquotesingle s whole standalone footprint (both theorems
in one theory, deliberately, so the staged work moves as a unit), and
--- unlike the enacted siblings on this page --- no isolated-build-host
verification census is attached to it yet; the citation attests the
script text, not a checked build artifact. Statement fidelity follows
the ported-soundness register used throughout: the Lean statement\textquotesingle
s implicit arguments become leading universals and its curried named
hypotheses a conjunction in source order; \texttt{.atom .le/.ge/.eq}
renders as \texttt{FAtom BowLe/BowGe/BowEq}; the pointwise antisymmetry
witness closes by the rational order\textquotesingle s antisymmetry
lemma. Because the rule has no checker arm, there is no family-check or
certified-fragment claim to make here --- the semantic stratum is the
entire port, matching the statement-register staging the note describes.

\sourcenote{HOL receipts are documentary: quoted results of \texttt{hol-strict} runs
and census rows from the HOL4 development\textquotesingle s build
evidence, not minted receipt artefacts. A HOL receipt store (the Lean
store\textquotesingle s twin) is named future work.
}
\ruleblock{Proof}

\sourcenote{Genre: core rule (tier 1), staged --- statement register
only; see the section note. Setting as fixed in the \rulename{EqElim}
section.}

\textbf{Claim} (\texttt{sound\_antisym}).\enspace Two premise sequents
conclude the atoms $t \leq u$ and $t \geq u$, each over the premise
formula $\Phi$ shared with the conclusion sequent; the conclusion is the
atom $t = u$. If both premises are semantically valid, so is
$\Phi \vdash t = u$.

\begin{blemma}[elementwise lifts at a shared tensor pair combine
pointwise]\label{lem:antisym-rel2conj}
Let $r$, $s$, $w$ be binary relations on $\Q$ with
$r(x,y) \wedge s(x,y) \Rightarrow w(x,y)$ for all $x, y$, and suppose
tensors $A$, $B$ stand in the elementwise-under-broadcast lift of $r$
and also in the lift of $s$. Then they stand in the lift of $w$.
\end{blemma}
\begin{proof}
The branch the lift takes is a function of the two shapes alone --- not
of the relation --- so the two hypotheses and the goal all quantify over
the \emph{same} domain, and the pointwise implication closes each point.
Spelled per branch: \emph{equal shapes}: all three statements read the
same zipped entry lists, pair by pair. \emph{Rank-0 $A$ against
positive-rank $B$}: all three quantify $A$'s scalar against each entry
of $B$; mirror image for rank-0 $B$. \emph{Two unequal positive-rank
shapes}: no lift holds there at all, so the hypothesis for $r$ is
already false and there is nothing to prove.
(\texttt{QTensor.rel2\_conj}.)
\end{proof}

\begin{bthm}[soundness of \rulename{Antisym}]\label{thm:antisym-sound}
If $\Phi \vdash t \leq u$ and $\Phi \vdash t \geq u$ are semantically
valid, so is $\Phi \vdash t = u$.
\end{bthm}
\begin{proof}
Fix an admissible $\sigma \models \Phi$. The $\leq$ premise yields
evaluations $A$ of $t$ and $B$ of $u$ in the elementwise
$\leq$-lift; the $\geq$ premise yields evaluations $A'$ of $t$ and $B'$
of $u$ in the elementwise $\geq$-lift. Evaluation is a function of the
valuation and the term, so $A' = A$ and $B' = B$: both premises
constrain the \emph{same} evaluated pair, which is what makes
Lemma~\ref{lem:antisym-rel2conj} applicable at all. Apply it with
$r \mathbin{:=} {\leq}$, $s \mathbin{:=} {\geq}$,
$w \mathbin{:=} {=}$ and the pointwise fact
$x \leq y \wedge x \geq y \Rightarrow x = y$ --- antisymmetry of the
order on $\Q$ --- and the equality lift closes the atom $t = u$.
\end{proof}

\begin{bremark}[the negative space: no typing pack, and why both
premises]\label{rem:antisym-neg}
No typing fields (\rulename{EqRefl}'s \texttt{tyDef}/\texttt{fin}):
that rule must \emph{manufacture} an evaluation for its premise-free
conclusion, and pays for totality with static typing conditions; here
each premise atom holds only where both sides already evaluate finitely
at compatible shapes, and the proof above consumes exactly that
evidence --- nothing is manufactured (\rulename{EqSym} makes the same
economy for the same reason). Both orientations are load-bearing:
dropping either premise leaves a one-sided bound, and any evaluation
pair with $A$ strictly below $B$ entrywise satisfies the surviving
premise while refuting the conclusion --- the counterexample is
immediate at scalar shape already.
\end{bremark}

\begin{bremark}[what the broadcast-sense conclusion does not
license]\label{rem:antisym-broadcast}
The concluded atom holds in the elementwise-under-broadcast sense: a
scalar stands equal to a vector whose every entry it matches. That is
the equality atom's meaning everywhere in this calculus --- sufficient
for atom truth and for every consumer that reads the atom semantically
--- but it is not a license to interchange $t$ and $u$ inside an
arbitrary formula. Substitution rides \rulename{EqElim}, whose
common-defined-shape typing condition remains in force unchanged: that
section constructs a satisfiable-to-unsatisfiable flip from exactly the
scalar-equals-vector equalities this rule can legitimately conclude.
\end{bremark}

\begin{bremark}[why core, and what is staged]\label{rem:antisym-core}
Core because the inference is not derivable: an induction over
derivations shows every equality concluded from an equality-atom-free
context equates terms with identical polynomial normal forms, so no
chain of existing rules concludes $t = u$ from the two bounds when the
sides are genuinely different terms. Historically the calculus carried
antisymmetry only inside two bespoke rules for the ReLU function; a
core \rulename{Antisym} ends that monopoly and lets bespoke rules be
judged on their own content. Staged: the theorem above is the
statement register --- no payload constructor, wire tag, or checker arm
exists yet, so certificates cannot carry the rule today; the anchors
table below will gain its ``wired as core'' line at enactment.
\end{bremark}

\begin{anchors}
Lemma \ref{lem:antisym-rel2conj} & \texttt{QTensor.rel2\_conj}\\
Theorem \ref{thm:antisym-sound} & \texttt{sound\_antisym}; statement register only --- not yet wired (no \texttt{ProofObject} constructor exists)\\
HOL twin & theory \texttt{bern\_tr\_antisym}: \texttt{qtRel2\_conj}, \texttt{sound\_antisym}\\
\end{anchors}


\rulesection{12}{ConvElim}{core rule}
\label{blcatrule-ConvElim}
\ruleblock{The rule (Steele metanotation)}
\sourcenote{as realized in the mechanized development; source row: \texttt{\texttt{rules.toml}} --- \emph{ConvElim}. The metanotation statement is typeset by hand for notational quality; its mechanized ground truth is the Lean statement below, extracted from the sources at every build of this book and hash-bound. The rule \emph{roster} is machine-checked (every mechanized rule tag round-trips to exactly one catalogue row); the statement \emph{text} is editorial, checked by eye against the extracted Lean beside it.}
\begin{mathpar}
\inferrule*[Right=\rulename{ConvElim}]{
  \Phi \vdash x \in \mathcal{R}(\Delta, S)
}{
  \Phi \vdash \textstyle\bigwedge_{i = 0}^{n} \mathrm{canon}(\epsilon_i\,\lambda_i(x)) \leq 0 \\ \epsilon_i \,=\, {+1} \text{ if } i \in S, \quad {-1} \text{ if } i \in \bar S
}
\end{mathpar}

\par\medskip
\rulename{ConvElim} opens a region atom: from
$\Phi \vdash x \in \mathcal{R}(\Delta, S)$ it concludes the
region\textquotesingle s defining barycentric halfspace conjunction ---
one $(\cdot) \leq 0$ atom per support vertex, in a pinned canonical
spelling. It is the entry move of every argument that reasons
\emph{linearly} about a geometric premise: \rulename{BernCone} consumes
region atoms natively, but the combination rules (\rulename{Farkas},
\rulename{Syzygy}) match ordinary comparison atoms by exact canonical
polynomial, and this rule is what turns a region fact into atoms they
can see.

What the checker recomputes: the \emph{entire} conclusion, from
$(x, \Delta, S)$ alone --- the two barycentric affine forms of the 1-D
two-vertex support (weights are determined on a simplex, so nothing is
chosen and nothing rides the certificate), the sign-folding that turns
each core vertex\textquotesingle s wanted $\lambda \geq 0$ into a
$(-\lambda) \leq 0$ atom (the all-$\leq$ discipline the downstream
matchers consume), and the strictly positive primitive-integer rescale
that makes the spelling unique --- then demands syntactic equality with
the cited conclusion. No producer text is trusted; the side condition is
one decidable equation. Outside the checkable fragment (ambient
dimension 1, support exactly two vertices at distinct coordinates) the
build refuses and the rule cannot fire --- refusal is trivially sound,
and the distinct-coordinates gate is precisely where a degenerate
simplex\textquotesingle s division by zero is made unrepresentable.

The rule needs no typing or positivity fields: definedness and shape of
the region term are read out of the region atom\textquotesingle s own
semantics. Dropping the premise-form equation would sever the conclusion
from any premise at all; dropping the canonical-spelling pin would not
make the rule unsound but would strand its conclusions --- equivalent
spellings never match downstream --- which is why the spelling is
checked rather than suggested.

\cataloguenote{Geometric reading of extension sets, 1-D: $S = \emptyset$ is the segment
$[V_0, V_1]$; naming one vertex stretches the region to the closed ray
away from the other --- membership with a \emph{nonpositive} weight on
the extended vertex. The region atom is thus the calculus\textquotesingle
s interface between simplex geometry and sign arithmetic, and this rule
with \rulename{ConvIntro} is that interface\textquotesingle s two
directions, exact inverses on the fragment.
}
\ruleblock{Lean statement}
\declhead{BernsteinLogic.Soundness}{barycentricHalfspaceBuild\_spec}{\badge{BadgeGreen}{statement bound}}
\begin{leansrc}
/-- Structure of a successful barycentric-halfspace build (the 1-D case
`baryCoeffs1D` gates): the F4 all-≤ canonical conjunction, one
`f4VertexAtom` per vertex of the filter pair. -/
theorem barycentricHalfspaceBuild_spec {t : Term} {simplex : Simplex}
    {S : ExtensionSet} {F : Formula}
    (h : barycentricHalfspaceBuild t simplex S = some F) :
    S.simplex = simplex ∧ simplex.complex.ambDim = 1 ∧
    ∃ (a b : Fin simplex.complex.nVertices) (V₀ V₁ : ℚ),
      (List.finRange simplex.complex.nVertices).filter
        (fun v => simplex.contains v) = [a, b] ∧
      (∀ j, simplex.complex.coord a j = V₀) ∧
      (∀ j, simplex.complex.coord b j = V₁) ∧
      V₀ ≠ V₁ ∧
      F = .and
        (f4VertexAtom (S.containsNat a.val)
          (V₁ / (V₁ - V₀)) (-1 / (V₁ - V₀)) t)
        (f4VertexAtom (S.containsNat b.val)
          (-V₀ / (V₁ - V₀)) (1 / (V₁ - V₀)) t)
\end{leansrc}
\declhead{BernsteinLogic.Soundness}{f4VertexAtom\_holds\_iff}{\badge{BadgeGreen}{statement bound}}
\begin{leansrc}
/-- The F4 vertex atom holds iff the signed barycentric value is
nonpositive: `λ ≤ 0` for an extended vertex, `0 ≤ λ` for a core one
(the sign folded into the term, the rescale cancelled). -/
theorem f4VertexAtom_holds_iff {H : Hypernet} {σ : H.Valuation} {t : Term}
    {q : QTensor} (hq : Term.evalQ σ t = some q) (hsq : SingletonShape q.shape)
    (inS : Bool) (off slp : ℚ) :
    Formula.holds σ (f4VertexAtom inS off slp t) ↔
      (if inS then off + slp * q.scalar ≤ 0
       else 0 ≤ off + slp * q.scalar)
\end{leansrc}
\declhead{BernsteinLogic.Soundness}{maskRegion\_pair\_elim}{\badge{BadgeGreen}{statement bound}}
\begin{leansrc}
/-- Extraction: a masked region membership over a two-vertex 1-D support
pins the barycentric weights to their rational closed forms, and exports
the sign of each weight according to its extension flag. -/
theorem maskRegion_pair_elim {C : SimpC} {mask ext : Fin C.nVertices → Bool}
    {x : Fin C.ambDim → ℚ} {a b : Fin C.nVertices} {V₀ V₁ : ℚ}
    (hfil : (List.finRange C.nVertices).filter (fun v => mask v) = [a, b])
    (hVa : ∀ j, C.coord a j = V₀) (hVb : ∀ j, C.coord b j = V₁)
    (hne : V₀ ≠ V₁) (j0 : Fin C.ambDim)
    (h : MaskRegion C mask ext x) :
    (ext a = true → (V₁ - x j0) / (V₁ - V₀) ≤ 0) ∧
    (ext a = false → 0 ≤ (V₁ - x j0) / (V₁ - V₀)) ∧
    (ext b = true → (x j0 - V₀) / (V₁ - V₀) ≤ 0) ∧
    (ext b = false → 0 ≤ (x j0 - V₀) / (V₁ - V₀))
\end{leansrc}
\declhead{BernsteinLogic.Soundness}{sound\_convElim}{\badge{BadgeGreen}{statement bound}}
\begin{leansrc}
theorem sound_convElim {H : Hypernet} {seq premise : Sequent}
    {term : Term} {simplex : Simplex} {S : ExtensionSet}
    (ih : premise.semanticallyValid H)
    (hpre : seq.pre = premise.pre)
    (hpremiseForm : premise.post = .inRegion term simplex S)
    (hhalf : barycentricHalfspaceForm term simplex S seq.post) :
    seq.semanticallyValid H
\end{leansrc}
\par\medskip
\textbf{The canonical form, meaning bound by display.}\enspace The
conclusion formula is the value of the build displayed below ---
affine data (\texttt{baryCoeffs1D}), positive primitive-integer rescale
(\texttt{primAffine}), atom layout (\texttt{f4Lhs}), per-vertex
sign-folding (\texttt{f4VertexAtom}), assembly
(\texttt{barycentricHalfspaceBuild}) --- and the payload\textquotesingle
s side condition is the one-line equation
\texttt{barycentricHalfspaceForm}. A skeptical reader should confirm by
eye that every gate failure yields \texttt{none} (refusal), that the
rescale factor is strictly positive on every reachable branch, and that
no valuation is consulted anywhere in the chain.

\declhead{BernsteinLogic.Core.Catalogue}{baryCoeffs1D}{\badge{BadgeGreen}{definition bound}}
\begin{leansrc}
/-- The (vertex, offset, slope) affine data of the two barycentric
coordinate functions of a 1-dimensional simplex: `λ(x) = offset + slope·x`.
Same gates as `baryLambdas1D` (which builds the cell rules' pinned Term
shape from the same coordinates). -/
def baryCoeffs1D (cx : SimpC) (dContains : Fin cx.nVertices → Bool) :
    Option (List (Fin cx.nVertices × Rat × Rat)) :=
  if h : cx.ambDim = 1 then
    let j0 : Fin cx.ambDim := ⟨0, by omega⟩
    match (List.finRange cx.nVertices).filter (fun v => dContains v) with
    | [a, b] =>
      let V₀ := cx.coord a j0
      let V₁ := cx.coord b j0
      if V₀ = V₁ then none
      else
        let den := V₁ - V₀
        some [(a, V₁ / den, -1 / den), (b, -V₀ / den, 1 / den)]
    | _ => none
  else none
\end{leansrc}
\declhead{BernsteinLogic.Core.Catalogue}{primAffine}{\badge{BadgeGreen}{definition bound}}
\begin{leansrc}
/-- Primitive-integer rescale of the affine pair (slope `s`, offset `c`):
the positive multiple `(k·s, k·c)` whose entries are the coprime-integer
vector sympy's `canon_sympy` produces for `s·x + c` (producer finding F4).
`c = 0` degenerates to sign normalization (single-entry primitive vector
is ±1); `(0, 0)` is fixed. Computed rationally — `k·s` value-equals the
integer `n₁/g` because `s.den ∣ L` and `g ∣ n₁`. -/
def primAffine (s c : Rat) : Rat × Rat :=
  if c = 0 then
    (if s < 0 then -1 else if s = 0 then 0 else 1, 0)
  else
    let L : Nat := Nat.lcm s.den c.den
    let g : Nat := Int.gcd (s.num * (L / s.den : Nat)) (c.num * (L / c.den : Nat))
    let k : Rat := (L : Rat) / (g : Rat)
    (k * s, k * c)
\end{leansrc}
\declhead{BernsteinLogic.Core.Catalogue}{f4Lhs}{\badge{BadgeGreen}{definition bound}}
\begin{leansrc}
/-- The F4 atom body for the rescaled pair: degree-sorted monomial shape
`s'·t + c'`, the constant dropped when zero (sympy `_term_ast`). -/
def f4Lhs (s' c' : Rat) (t : Term) : Term :=
  if c' = 0 then .mul (.const s') t
  else .add (.mul (.const s') t) (.const c')
\end{leansrc}
\declhead{BernsteinLogic.Core.Catalogue}{f4VertexAtom}{\badge{BadgeGreen}{definition bound}}
\begin{leansrc}
/-- One vertex's F4 halfspace atom: sign-folded (complement vertex's λ
negated into the term), rescaled, `≤ 0`-shaped. -/
def f4VertexAtom (inS : Bool) (off slp : Rat) (t : Term) : Formula :=
  let sc := if inS then primAffine slp off else primAffine (-slp) (-off)
  .atom .le (f4Lhs sc.1 sc.2 t) (.const 0)
\end{leansrc}
\declhead{BernsteinLogic.Core.Catalogue}{barycentricHalfspaceBuild}{\badge{BadgeGreen}{definition bound}}
\begin{leansrc}
def barycentricHalfspaceBuild (t : Term) (simplex : Simplex)
    (S : ExtensionSet) : Option Formula := do
  guard (S.simplex = simplex)
  let cs ← baryCoeffs1D simplex.complex simplex.contains
  match cs.map (fun (v, off, slp) =>
    f4VertexAtom (S.containsNat v.val) off slp t) with
  | [x, y] => some (.and x y)
  | _ => none
\end{leansrc}
\declhead{BernsteinLogic.Core.Catalogue}{barycentricHalfspaceForm}{\badge{BadgeGreen}{definition bound}}
\begin{leansrc}
def barycentricHalfspaceForm (t : Term) (simplex : Simplex)
    (S : ExtensionSet) (F : Formula) : Prop :=
  barycentricHalfspaceBuild t simplex S = some F
\end{leansrc}
\declhead{BernsteinLogic.Core.Catalogue}{PConvElimData}{\badge{BadgeGreen}{definition bound}}
\begin{leansrc}
inductive PConvElimData : Hypernet → Sequent → Type where
  | mk (H : Hypernet) (seq : Sequent)
       (premise : Sequent) (premiseProof : ProofObject H premise)
       (term : Term) (simplex : Simplex) (S : ExtensionSet)
       (preEq       : seq.pre = premise.pre)
       (premiseForm : premise.post = .inRegion term simplex S)
       (halfspace   : barycentricHalfspaceForm term simplex S seq.post)
       : PConvElimData H seq
\end{leansrc}
\ruleblock{Correspondence (Steele display $\leftrightarrow$ Lean)}
\badge{BadgeGreen}{reviewer-certified against the displayed texts}
\sourcenote{The dictionary below is a reviewer-certified statement of how the mechanized statements above interpret the displayed rule: display metavariable $\mapsto$ Lean binder, each display hypothesis $\mapsto$ the premise or side condition carrying it, the display conclusion $\mapsto$ the goal --- with any deliberate difference in width between display and mechanization named outright. Certification is recorded as a hash triple over (the displayed cell, the extracted statement texts, this dictionary\textquotesingle s own text); if any of the three changes, this block turns amber until a reviewer re-certifies. The correspondence claim itself remains editorial --- reviewer-grade, not machine-checked; its machine-checked successor types each soundness statement against the rule as a structured object, and supersedes this block per rule as it lands.}
Judgment frame as in \rulename{Ax}; sub-derivation premise as at
\rulename{Weaken}.
\begin{itemize}
\item $x \mapsto$ \texttt{term : Term}; $\Delta \mapsto$
  \texttt{simplex : Simplex}; $S \mapsto$ \texttt{S : ExtensionSet};
  the premise judgment $\Phi \vdash x \in \mathcal{R}(\Delta, S)$
  $\mapsto$ \texttt{ih} with the shape tie \texttt{hpremiseForm} $:$
  \texttt{premise.post = .inRegion term simplex S}; $\Phi$ shared via
  \texttt{hpre}.
\item the \emph{entire} displayed conclusion --- the flat conjunction,
  its vertex order, the sign pattern $\epsilon_i$, the affine
  spellings $\lambda_i$ --- $\mapsto$ one hypothesis, \texttt{hhalf}
  $:$ \texttt{barycentricHalfspaceForm term simplex S seq.post}:
  \texttt{seq.post} is literally the value of the verified canonical
  build. The display\textquotesingle s $\lambda_i$ and $\epsilon_i$ are
  not statement binders but \emph{outputs} of
  \texttt{barycentricHalfspaceBuild}, quoted in full below.
\item conclusion judgment $\mapsto$ the goal
  \texttt{seq.semanticallyValid H}; certified alongside,
  \texttt{barycentricHalfspaceBuild\_spec} unpacks a successful build
  run into the geometric data the soundness proof consumes.
\end{itemize}
Width note, narrower and loud about it: a successful build run forces
ambient dimension $1$ with support exactly two vertices at distinct
coordinates (the first conjuncts of
\texttt{barycentricHalfspaceBuild\_spec}), so the mechanized statement
covers exactly the one-dimensional instance of the displayed
general-dimension rule. Outside that fragment the build returns
\texttt{none}, the hypothesis is unsatisfiable, and the rule cannot
fire --- refusal, never unverified acceptance; the displayed rule at
higher dimension is design intent, not yet certified statement.

Spelling alignment, formerly a live discrepancy: the displayed
conclusion and the certified build now agree --- both state the flat
sign-folded all-$\leq$ conjunction, each atom
$\mathrm{canon}(\epsilon_i\,\lambda_i(x)) \leq 0$ with the
extension/core orientation carried by the sign $\epsilon_i$
(\texttt{f4VertexAtom} negates the core vertex\textquotesingle s
affine form into the term and compares $\leq 0$ --- the commentary
describes this discipline). Earlier printings displayed a mixed
$\leq$/$\geq$ per-vertex comparator pattern; the roster cell has
since been reconciled to the realized form, and this certification
binds the reconciled texts.

\ruleblock{Lean verifier receipts}
\sourcenote{module \texttt{BernsteinLogic.Soundness}; bound build \texttt{a20d90e66c13\ldots} (bound from receipt a9c0e34d3e3c9c8c, 2026-08-03T21:30:36Z). A receipt records that named checking layers accepted a build: (module, compiled-artifact SHA-256, environment SHA-256, checker layers), minted when the checking gate ran; the store receipts attesting the bound build are listed with the layers each run credited. Nothing here asks the reader to trust the store: every layer is an independently runnable program over the quoted sources --- rebuild the module and re-run the checkers to reproduce the verdict. Receipts attesting superseded builds credit no layer and are omitted from this rendering (the count of omissions, if any, is noted below).}
\begin{receipttable}
\receiptrow{a9c0e34d3e3c9c8c}{2026-08-03T21:30:36Z}{build+decorrelated}{\badge{BadgeGreen}{current}}
\end{receipttable}
\sourcenote{21 further receipt(s) in the store attest superseded builds of this module and are omitted from this rendering.}
\sourcenote{module \texttt{BernsteinLogic.Core.Catalogue}; bound build \texttt{cbc92e976a1d\ldots} (bound from receipt 5abdb3589707b5df, 2026-08-03T21:30:33Z). A receipt records that named checking layers accepted a build: (module, compiled-artifact SHA-256, environment SHA-256, checker layers), minted when the checking gate ran; the store receipts attesting the bound build are listed with the layers each run credited. Nothing here asks the reader to trust the store: every layer is an independently runnable program over the quoted sources --- rebuild the module and re-run the checkers to reproduce the verdict. Receipts attesting superseded builds credit no layer and are omitted from this rendering (the count of omissions, if any, is noted below).}
\begin{receipttable}
\receiptrow{24804e66edba0bc9}{2026-08-03T20:25:09Z}{build+decorrelated}{\badge{BadgeGreen}{current}}
\receiptrow{5abdb3589707b5df}{2026-08-03T21:30:33Z}{build+decorrelated}{\badge{BadgeGreen}{current}}
\end{receipttable}
\sourcenote{18 further receipt(s) in the store attest superseded builds of this module and are omitted from this rendering.}
\layerline{\layeritem{build}{lake build (elaboration)}{\badge{BadgeGreen}{current}}\layeritem{replay}{leanchecker (official kernel replay)}{\badge{BadgeAmber}{stale only}}\layeritem{hardened}{safe\_verify (hardened screen)}{\badge{BadgeAmber}{stale only}}\layeritem{decorrelated}{lean4export + nanoda (decorrelated kernel)}{\badge{BadgeGreen}{current}}}
\sourcenote{\badge{BadgeGray}{comparator: none} no fixture in the three-implementation comparator harness (compiled / kernel-interpreted / reference-interpreted) exercises this rule.}
\ruleblock{HOL4 statement}
\sourcenote{theory \texttt{bern\_catalogue}, HOL4 development script \texttt{\texttt{bern\_catalogueScript.sml}} \badge{BadgeGreen}{script bound}}
\holhead{Definition baryCoeffs1D\_def}
\begin{holsrc}
Definition baryCoeffs1D_def:
  baryCoeffs1D cx dmask =
    if cx.ambDim = 1 then
      (case FILTER (λv. getD dmask v F) (rangeL cx.nVertices) of
       | [a; b] =>
           (let v0 = getD (getD cx.coord a []) 0 (0:rat) in
            let v1 = getD (getD cx.coord b []) 0 (0:rat) in
            let den = v1 - v0 in
              if den = 0 then NONE
              else SOME [(a, v1 / den, -(&1) / den);
                         (b, -v0 / den, &1 / den)])
       | _ => NONE)
    else NONE
End
\end{holsrc}
\holhead{Definition primAffine\_def}
\begin{holsrc}
Definition primAffine_def:
  primAffine (s:rat) (c:rat) =
    if c = 0 then
      ((if s < 0 then -1 else if s = 0 then 0 else 1) : rat, 0:rat)
    else
      let l = lcm (RATD s) (RATD c) in
      let g = gcd (intAbsNum (RATN s * &(l DIV RATD s)))
                  (intAbsNum (RATN c * &(l DIV RATD c))) in
      (* g = 0 is unreachable (c ≠ 0 ⇒ RATN c ≠ 0 ⇒ the second gcd
         argument is nonzero); the guard is the translator's division
         precondition, spelled on the divisor (Lean's Rat division by
         zero yields 0, so the two sides agree even on the dead branch). *)
      let k = if &g = (0:rat) then 0 else &l / &g in
        (k * s, k * c)
End
\end{holsrc}
\holhead{Definition f4Lhs\_def}
\begin{holsrc}
Definition f4Lhs_def:
  f4Lhs (s':rat) (c':rat) t =
    if c' = 0 then TmMul (TmConst s') t
    else TmAdd (TmMul (TmConst s') t) (TmConst c')
End
\end{holsrc}
\holhead{Definition f4VertexAtom\_def}
\begin{holsrc}
Definition f4VertexAtom_def:
  f4VertexAtom inS (off:rat) (slp:rat) t =
    let sc = if inS then primAffine slp off else primAffine (-slp) (-off) in
      FAtom BowLe (f4Lhs (FST sc) (SND sc) t) (TmConst 0)
End
\end{holsrc}
\holhead{Definition barycentricHalfspaceBuild\_def}
\begin{holsrc}
Definition barycentricHalfspaceBuild_def:
  barycentricHalfspaceBuild t sx es =
    if es.simplex = sx then
      (case baryCoeffs1D sx.complex sx.vmask of
       | NONE => NONE
       | SOME cs =>
         (case MAP (λ(v,off,slp). f4VertexAtom (extContainsNat es v) off slp t)
               cs of
          | [x; y] => SOME (FAnd x y)
          | _ => NONE))
    else NONE
End
\end{holsrc}
\holhead{Definition barycentricHalfspaceFormB\_def}
\begin{holsrc}
Definition barycentricHalfspaceFormB_def:
  barycentricHalfspaceFormB t sx es fm ⇔
    barycentricHalfspaceBuild t sx es = SOME fm
End
\end{holsrc}
\ruleblock{HOL verifier receipts}
\badge{BadgeBlue}{documentary} \texttt{hol-strict} CLEAN census over the ported checker theory:
\texttt{bern\_catalogue}, 472 theorems, 0 cheat-tagged, 0 nonstandard
axioms (whitelist: \texttt{DISK\_THM}; the four \texttt{bool} axioms) ---
verification run of 2026-07-22 on an isolated build host, against a
content-addressed checkout of the exact sources cited here. The earlier
loud caveat (the canonical-form revision of 2026-07-18 postdated the
prior census run) is retired: this run attests exactly the quoted text.
There is no signed HOL receipt store yet, and these citations are
documentary.
Semantic stratum, also ported: the soundness twin
\texttt{sound\_convElim} (theory \texttt{bern\_tr\_fragclose}:
\texttt{hol-strict} clean, 26 theorems, 0 offences, verification run of
2026-07-22 on an isolated build host, content-addressed at the same
sources), consumed by the certified-fragment keystone
(\texttt{checkPf\_sound\_frag}, theory \texttt{bern\_tr\_splice}, 20
theorems, 0 offences, same run): a reference-checker accept inside the
proven wire fragment yields sequent validity.

\sourcenote{HOL receipts are documentary: quoted results of \texttt{hol-strict} runs
and census rows from the HOL4 development\textquotesingle s build
evidence, not minted receipt artefacts. A HOL receipt store (the Lean
store\textquotesingle s twin) is named future work.
}
\ruleblock{Proof}

\sourcenote{Genre: core rule (tier 1). Setting (valuations, partial
evaluation, elementwise atom semantics) as fixed in the \rulename{EqElim}
section; this section additionally fixes the region-atom semantics and
the canonical halfspace form, which the \rulename{ConvIntro} section then
cites.}

\textbf{Region atoms, in the checkable fragment.}\enspace A region atom
$x \in \mathcal{R}(\Delta, S)$ holds at $\sigma$ when the extension set
$S$ is tied to the simplex $\Delta$, the region term evaluates to a
rank-1 tensor of extent equal to the ambient dimension, and the point it
holds lies in the \emph{masked region}: it is an affine combination of
$\Delta$'s vertices with weights summing to $1$, nonnegative on core
vertices and \emph{nonpositive} on the vertices named by $S$. The rules
below operate in the fragment their checker accepts --- ambient dimension
$1$, support exactly two vertices $a, b$ at distinct coordinates
$V_0 \neq V_1$ --- where the geometry is completely explicit: the two
barycentric weights are \emph{determined}, affine in the point,
\[
\lambda_a(x) = \frac{V_1 - x}{V_1 - V_0}, \qquad
\lambda_b(x) = \frac{x - V_0}{V_1 - V_0}, \qquad
\lambda_a + \lambda_b \equiv 1,
\]
and the region is: the segment $[V_0, V_1]$ for $S = \emptyset$; the
closed ray from $V_1$ away from $V_0$ for $S = \{a\}$ (weight of the
extended vertex $\leq 0$); the mirror ray for $S = \{b\}$. ($S$ is a
proper subset, so both-extended does not arise.) Determinacy of the
weights is why elimination below is complete --- no witness is chosen,
none rides the certificate.

\textbf{The canonical halfspace form.}\enspace The conclusion is not
``some conjunction equivalent to the sign facts'' but one \emph{pinned}
formula, recomputed by the checker from $(x, \Delta, S)$ alone: one atom
per support vertex, in support order, each in $(\cdot) \leq 0$ shape:
for $v \in S$ the atom $\mathrm{canon}(\lambda_v(x)) \leq 0$, for
$v \notin S$ the atom $\mathrm{canon}(-\lambda_v(x)) \leq 0$ (the wanted
$\lambda_v \geq 0$ \emph{sign-folded} into the term --- the all-$\leq$
discipline), where $\mathrm{canon}$ rescales the affine pair
(slope, offset) by a \emph{strictly positive} rational to its
primitive-integer representative and lays the atom body out in
degree-sorted shape. If the build's gates fail (dimension $\neq 1$,
support $\neq 2$ vertices, $V_0 = V_1$, or the atom's simplex not the
cited $\Delta$), the build refuses, the payload
is uninhabitable, and the rule cannot fire: refusal is trivially sound,
and in the degenerate-simplex case it is also where a division by zero
would otherwise live --- made unrepresentable rather than checked-then-divided.

\textbf{Claim} (\texttt{sound\_convElim}).\enspace If
$\Phi \vdash x \in \mathcal{R}(\Delta, S)$ is semantically valid and the
conclusion formula is the canonical halfspace form for
$(x, \Delta, S)$, then $\Phi \vdash$ (that conjunction) is semantically
valid.

\begin{blemma}[structure of a successful build]\label{lem:convelim-spec}
If the canonical build succeeds with output $F$, then the ambient
dimension is $1$, $S$ is tied to $\Delta$, the support filter returns
exactly $[a, b]$ with constant vertex coordinates $V_0 \neq V_1$, and
$F$ is the conjunction of the two vertex atoms described above.
\end{blemma}
\begin{proof}
Unfold the build: each claim is one of its gates, and the output shape
is its final assembly. (\texttt{barycentricHalfspaceBuild\_spec}.)
\end{proof}

\begin{blemma}[canonical atom $\leftrightarrow$ raw sign fact]
\label{lem:convelim-atom}
Let $q$ be a defined, singleton-shaped evaluation of $x$ at $\sigma$, and
let $(\mathit{off}, \mathit{slp})$ be an affine pair. The vertex atom for
membership flag $m$ holds at $\sigma$ iff
\[
m = \mathrm{true}:\quad \mathit{off} + \mathit{slp}\cdot q \;\leq\; 0,
\qquad\qquad
m = \mathrm{false}:\quad 0 \;\leq\; \mathit{off} + \mathit{slp}\cdot q.
\]
\end{blemma}
\begin{proof}
Two steps. \emph{Sign-folding}: for $m = \mathrm{false}$ the build
negates the pair before canonicalizing, so the emitted atom asserts
$\mathrm{canon}(-\lambda) \leq 0$, i.e.\ the displayed $\geq$. \emph{The
rescale is sign-transparent}: $\mathrm{canon}$ multiplies the pair by
$k > 0$ --- in the generic case $k = L/g$ with $L$ a positive lcm of
denominators and $g$ a positive gcd of the scaled numerators; in the
zero-offset case the output slope is the sign
$\mathrm{sgn}(\mathit{slp})$, i.e.\ the pair scaled by $1/|\mathit{slp}|$
--- well-defined on every pair in play: a successful barycentric
build has distinct vertices, so each emitted affine pair carries
slope $\pm 1/(V_1 - V_0) \neq 0$ (sign-folding only negates it), and
in the zero-offset case that slope is the whole pair; for the
degenerate all-zero pair, which no successful build emits, both sides
of the displayed equivalence hold trivially and the statement is
unaffected --- and a strictly positive rescale of an affine
function has the same sign locus. The atom's left side evaluates (constants and one already-
defined $q$), at singleton shape, so the elementwise reading is the
scalar comparison of the displayed value against $0$.
(\texttt{f4VertexAtom\_holds\_iff}; the rescale is
\texttt{primAffine}, the layout \texttt{f4Lhs}.)
\end{proof}

\begin{blemma}[region elimination, two-vertex case]
\label{lem:convelim-elim}
Suppose $\sigma \models x \in \mathcal{R}(\Delta, S)$ in the fragment of
Lemma~\ref{lem:convelim-spec}, with $q$ the evaluated coordinate. Then,
conditionally on membership: if $a \in S$ then $\lambda_a(q) \leq 0$, if
$a \notin S$ then $\lambda_a(q) \geq 0$; likewise for $b$.
\end{blemma}
\begin{proof}
The masked-region witness consists of weights on the support summing to
$1$ that reproduce $q$ and satisfy the sign discipline. On a two-vertex
support the reproduction and sum equations force the weights to be
exactly $\lambda_a(q), \lambda_b(q)$ (two linear equations, invertible
since $V_0 \neq V_1$), so the witness's sign discipline \emph{is} the
displayed conditional facts. (\texttt{maskRegion\_pair\_elim}.)
\end{proof}

\begin{bthm}[soundness of \rulename{ConvElim}]\label{thm:convelim-sound}
Under the payload's checked hypotheses, if
$\Phi \vdash x \in \mathcal{R}(\Delta, S)$ is semantically valid, so is
$\Phi \vdash F$ for the built conclusion $F$.
\end{bthm}
\begin{proof}
Fix an admissible $\sigma \models \Phi$. Premise validity gives the
region atom at $\sigma$: in particular $x$ \emph{evaluates}, to a rank-1
extent-1 tensor (the atom's own shape field --- no typing pack is needed
here; contrast \autoref{rem:convelim-asym}), whose scalar is $q$.
Lemma~\ref{lem:convelim-elim} delivers the four conditional sign facts;
Lemma~\ref{lem:convelim-atom} converts, per vertex, the applicable one
into the built vertex atom; Lemma~\ref{lem:convelim-spec} says the
conclusion is exactly the conjunction of those two atoms.
\end{proof}

\begin{bremark}[why a pinned canonical spelling]\label{rem:convelim-canon}
The combination rules downstream (\rulename{Farkas}, \rulename{Syzygy})
match premise atoms \emph{by exact canonical polynomial}; a conclusion
emitted in any equivalent-but-different spelling would simply never
match, stranding the geometric fact. Pinning the spelling in the checker
--- which recomputes the build from $(x, \Delta, S)$ and demands
syntactic equality with the cited conclusion, trusting no producer text
--- makes the emitted atoms the \emph{same} tokens every producer and
every downstream matcher agrees on. The all-$\leq$ polarity keeps each
atom in the $(\cdot) \leq 0$ normal form those rules consume, and the
positive primitive-integer rescale makes the representative unique. Note
the checked object is the \emph{formula}, not a semantic condition: this
rule's side condition is one syntactic equation, decidable by
construction.
\end{bremark}

\begin{bremark}[what is \emph{not} hypothesized]\label{rem:convelim-asym}
No shape or finiteness fields appear in this payload: elimination reads
definedness and shape out of the region atom itself. The mirror rule
cannot --- its premise is the halfspace conjunction, whose atoms
tolerate broadcast and are satisfied by a \emph{scalar} $x$ --- and that
asymmetry is exactly where \rulename{ConvIntro}'s one extra typing field
comes from (next section).
\end{bremark}

\begin{anchors}
Lemma \ref{lem:convelim-spec} & \texttt{barycentricHalfspaceBuild\_spec}\\
build, affine data, rescale, atom layout & \texttt{barycentricHalfspaceBuild}, \texttt{baryCoeffs1D}, \texttt{primAffine}, \texttt{f4Lhs}, \texttt{f4VertexAtom}\\
Lemma \ref{lem:convelim-atom} & \texttt{f4VertexAtom\_holds\_iff}\\
Lemma \ref{lem:convelim-elim} & \texttt{maskRegion\_pair\_elim}\\
Theorem \ref{thm:convelim-sound} & \texttt{sound\_convElim}; wired as core \texttt{RuleData.convElim} into \texttt{ProofObject.sound}\\
form predicate (the checked side condition) & \texttt{barycentricHalfspaceForm} (payload field \texttt{halfspace})\\
\end{anchors}


\rulesection{13}{ConvIntro}{core rule}
\label{blcatrule-ConvIntro}
\ruleblock{The rule (Steele metanotation)}
\sourcenote{as realized in the mechanized development; source row: \texttt{\texttt{rules.toml}} --- \emph{ConvIntro}. The metanotation statement is typeset by hand for notational quality; its mechanized ground truth is the Lean statement below, extracted from the sources at every build of this book and hash-bound. The rule \emph{roster} is machine-checked (every mechanized rule tag round-trips to exactly one catalogue row); the statement \emph{text} is editorial, checked by eye against the extracted Lean beside it.}
\begin{mathpar}
\inferrule*[Right=\rulename{ConvIntro}]{
  \Phi \vdash \textstyle\bigwedge_{i = 0}^{n} \mathrm{canon}(\epsilon_i\,\lambda_i(x)) \leq 0 \\ \epsilon_i \,=\, {+1} \text{ if } i \in S, \ {-1} \text{ if } i \in \bar S \\ \bshape_H(x) = [\,n\,]
}{
  \Phi \vdash x \in \mathcal{R}(\Delta, S)
}
\end{mathpar}

\par\medskip
\rulename{ConvIntro} closes a region atom: from the canonical barycentric
halfspace conjunction for $(x, \Delta, S)$ --- the same pinned formula
\rulename{ConvElim} emits, recomputed by the same build --- it concludes
$\Phi \vdash x \in \mathcal{R}(\Delta, S)$. A producer uses it to
re-enter the geometric world after arithmetic: a subdivision branch or a
combination argument establishes the two vertex atoms, and this rule
converts them back into the region fact that \rulename{BernCone} or a
further region rule consumes.

What the checker recomputes: the build (as in \rulename{ConvElim} ---
one syntactic equation, producer text untrusted, refusal outside the 1-D
two-vertex fragment), plus this rule\textquotesingle s one extra typing
field: $\bshape_H(x) = [\,d\,]$, the region term is statically rank-1 of
the ambient extent. The asymmetry with \rulename{ConvElim} is
principled: elimination reads the term\textquotesingle s shape out of
the region atom it consumes, while introduction is handed only halfspace
atoms --- which hold \emph{elementwise under broadcast} and are
perfectly satisfied by a scalar $x$ --- so the shape must arrive as a
checked hypothesis. Drop it and a closed scalar inside the segment
yields a satisfiable (indeed universally true) premise whose region
conclusion is unsatisfiable --- a rank-0 tensor never satisfies a
region atom; the proof section builds the refutation concretely. It is
the same broadcast liberality that breaks un-gated \rulename{EqElim},
repaired by the same typing-field pattern.

\cataloguenote{On the fragment, \rulename{ConvIntro} and \rulename{ConvElim} are exact
inverses --- membership, the four weight-sign facts, and the canonical
conjunction are three presentations of one determinate-weight geometry,
so the pair is a lossless interface between region atoms and the
arithmetic fragment.
}
\ruleblock{Lean statement}
\declhead{BernsteinLogic.Soundness}{Term.evalQ\_bshape}{\badge{BadgeGreen}{statement bound}}
\begin{leansrc}
/-- bshape soundness: a defined evaluation takes the broadcast shape. -/
theorem Term.evalQ_bshape {H : Hypernet} (σ : H.Valuation) :
    ∀ (t : Term) (s : Shape) (q : QTensor),
    t.bshape H = some s → Term.evalQ σ t = some q → q.shape = s
\end{leansrc}
\declhead{BernsteinLogic.Soundness}{f4VertexAtom\_holds\_inv}{\badge{BadgeGreen}{statement bound}}
\begin{leansrc}
/-- A held F4 vertex atom forces the core term to evaluate. -/
theorem f4VertexAtom_holds_inv {H : Hypernet} {σ : H.Valuation} {t : Term}
    {inS : Bool} {off slp : ℚ}
    (h : Formula.holds σ (f4VertexAtom inS off slp t)) :
    ∃ q, Term.evalQ σ t = some q
\end{leansrc}
\declhead{BernsteinLogic.Soundness}{maskRegion\_pair\_intro}{\badge{BadgeGreen}{statement bound}}
\begin{leansrc}
/-- Construction: rational weights of the right signs assemble a
`MaskRegion` witness over a two-vertex 1-D support. -/
theorem maskRegion_pair_intro {C : SimpC} {mask ext : Fin C.nVertices → Bool}
    {x : Fin C.ambDim → ℚ} {a b : Fin C.nVertices} {V₀ V₁ : ℚ}
    (hfil : (List.finRange C.nVertices).filter (fun v => mask v) = [a, b])
    (hVa : ∀ j, C.coord a j = V₀) (hVb : ∀ j, C.coord b j = V₁)
    (hne : V₀ ≠ V₁) (hdim : C.ambDim = 1) (j0 : Fin C.ambDim)
    (hw0t : ext a = true → (V₁ - x j0) / (V₁ - V₀) ≤ 0)
    (hw0f : ext a = false → 0 ≤ (V₁ - x j0) / (V₁ - V₀))
    (hw1t : ext b = true → (x j0 - V₀) / (V₁ - V₀) ≤ 0)
    (hw1f : ext b = false → 0 ≤ (x j0 - V₀) / (V₁ - V₀)) :
    MaskRegion C mask ext x
\end{leansrc}
\declhead{BernsteinLogic.Soundness}{sound\_convIntro}{\badge{BadgeGreen}{statement bound}}
\begin{leansrc}
theorem sound_convIntro {H : Hypernet} {seq premise : Sequent}
    {term : Term} {simplex : Simplex} {S : ExtensionSet}
    (ih : premise.semanticallyValid H)
    (hpre : seq.pre = premise.pre)
    (htermShape : term.bshape H = some [simplex.complex.ambDim])
    (hhalf : barycentricHalfspaceForm term simplex S premise.post)
    (hpost : seq.post = .inRegion term simplex S) :
    seq.semanticallyValid H
\end{leansrc}
\par\medskip
\textbf{Shared machinery.}\enspace The build chain is displayed in full
in the \rulename{ConvElim} section; here the book re-binds the assembly
and the form equation, plus this rule\textquotesingle s payload with its
one extra typing field (\texttt{termShape}).

\declhead{BernsteinLogic.Core.Catalogue}{barycentricHalfspaceBuild}{\badge{BadgeGreen}{definition bound}}
\begin{leansrc}
def barycentricHalfspaceBuild (t : Term) (simplex : Simplex)
    (S : ExtensionSet) : Option Formula := do
  guard (S.simplex = simplex)
  let cs ← baryCoeffs1D simplex.complex simplex.contains
  match cs.map (fun (v, off, slp) =>
    f4VertexAtom (S.containsNat v.val) off slp t) with
  | [x, y] => some (.and x y)
  | _ => none
\end{leansrc}
\declhead{BernsteinLogic.Core.Catalogue}{barycentricHalfspaceForm}{\badge{BadgeGreen}{definition bound}}
\begin{leansrc}
def barycentricHalfspaceForm (t : Term) (simplex : Simplex)
    (S : ExtensionSet) (F : Formula) : Prop :=
  barycentricHalfspaceBuild t simplex S = some F
\end{leansrc}
\declhead{BernsteinLogic.Core.Catalogue}{PConvIntroData}{\badge{BadgeGreen}{definition bound}}
\begin{leansrc}
inductive PConvIntroData : Hypernet → Sequent → Type where
  | mk (H : Hypernet) (seq : Sequent)
       (premise : Sequent) (premiseProof : ProofObject H premise)
       (term : Term) (simplex : Simplex) (S : ExtensionSet)
       (preEq     : seq.pre = premise.pre)
       (termShape : term.bshape H = some [simplex.complex.ambDim])
       (halfspace : barycentricHalfspaceForm term simplex S premise.post)
       (postEq    : seq.post = .inRegion term simplex S)
       : PConvIntroData H seq
\end{leansrc}
\ruleblock{Correspondence (Steele display $\leftrightarrow$ Lean)}
\badge{BadgeGreen}{reviewer-certified against the displayed texts}
\sourcenote{The dictionary below is a reviewer-certified statement of how the mechanized statements above interpret the displayed rule: display metavariable $\mapsto$ Lean binder, each display hypothesis $\mapsto$ the premise or side condition carrying it, the display conclusion $\mapsto$ the goal --- with any deliberate difference in width between display and mechanization named outright. Certification is recorded as a hash triple over (the displayed cell, the extracted statement texts, this dictionary\textquotesingle s own text); if any of the three changes, this block turns amber until a reviewer re-certifies. The correspondence claim itself remains editorial --- reviewer-grade, not machine-checked; its machine-checked successor types each soundness statement against the rule as a structured object, and supersedes this block per rule as it lands.}
Judgment frame as in \rulename{Ax}; sub-derivation premise as at
\rulename{Weaken}; binders as at \rulename{ConvElim}, mirrored.
\begin{itemize}
\item the displayed premise conjunction $\mapsto$ \texttt{hhalf} $:$
  \texttt{barycentricHalfspaceForm term simplex S premise.post} ---
  the same verified build, its value now equated with the
  \emph{premise\textquotesingle s} post; the premise judgment
  $\mapsto$ \texttt{ih}, context shared via \texttt{hpre}.
\item the display\textquotesingle s typing premise $\bshape_H(x) =
  [\,n\,]$ $\mapsto$ \texttt{htermShape} $:$ \texttt{term.bshape H =
  some [simplex.complex.ambDim]} --- statically rank-1 at the ambient
  extent, the one hypothesis \rulename{ConvElim} does not have; the
  commentary derives why introduction needs it and elimination does
  not.
\item conclusion $\Phi \vdash x \in \mathcal{R}(\Delta, S)$ $\mapsto$
  the goal under \texttt{hpost} $:$ \texttt{seq.post = .inRegion term
  simplex S}.
\end{itemize}
Width note: the same one-dimensional fragment confinement as
\rulename{ConvElim}, forced by the same spec conjuncts --- the
mechanized statement is the displayed rule at ambient dimension $1$;
outside it the build refuses and nothing is accepted unverified. The
\rulename{ConvElim} dictionary\textquotesingle s spelling alignment
carries over unchanged: the displayed premise conjunction and the
certified premise formula both state the build\textquotesingle s
sign-folded all-$\leq$ shape, the extension/core orientation carried
by the sign $\epsilon_i$ --- same halfspaces, same folded spelling,
display and mechanization reconciled.

\ruleblock{Lean verifier receipts}
\sourcenote{module \texttt{BernsteinLogic.Soundness}; bound build \texttt{a20d90e66c13\ldots} (bound from receipt a9c0e34d3e3c9c8c, 2026-08-03T21:30:36Z). A receipt records that named checking layers accepted a build: (module, compiled-artifact SHA-256, environment SHA-256, checker layers), minted when the checking gate ran; the store receipts attesting the bound build are listed with the layers each run credited. Nothing here asks the reader to trust the store: every layer is an independently runnable program over the quoted sources --- rebuild the module and re-run the checkers to reproduce the verdict. Receipts attesting superseded builds credit no layer and are omitted from this rendering (the count of omissions, if any, is noted below).}
\begin{receipttable}
\receiptrow{a9c0e34d3e3c9c8c}{2026-08-03T21:30:36Z}{build+decorrelated}{\badge{BadgeGreen}{current}}
\end{receipttable}
\sourcenote{21 further receipt(s) in the store attest superseded builds of this module and are omitted from this rendering.}
\sourcenote{module \texttt{BernsteinLogic.Core.Catalogue}; bound build \texttt{cbc92e976a1d\ldots} (bound from receipt 5abdb3589707b5df, 2026-08-03T21:30:33Z). A receipt records that named checking layers accepted a build: (module, compiled-artifact SHA-256, environment SHA-256, checker layers), minted when the checking gate ran; the store receipts attesting the bound build are listed with the layers each run credited. Nothing here asks the reader to trust the store: every layer is an independently runnable program over the quoted sources --- rebuild the module and re-run the checkers to reproduce the verdict. Receipts attesting superseded builds credit no layer and are omitted from this rendering (the count of omissions, if any, is noted below).}
\begin{receipttable}
\receiptrow{24804e66edba0bc9}{2026-08-03T20:25:09Z}{build+decorrelated}{\badge{BadgeGreen}{current}}
\receiptrow{5abdb3589707b5df}{2026-08-03T21:30:33Z}{build+decorrelated}{\badge{BadgeGreen}{current}}
\end{receipttable}
\sourcenote{18 further receipt(s) in the store attest superseded builds of this module and are omitted from this rendering.}
\layerline{\layeritem{build}{lake build (elaboration)}{\badge{BadgeGreen}{current}}\layeritem{replay}{leanchecker (official kernel replay)}{\badge{BadgeAmber}{stale only}}\layeritem{hardened}{safe\_verify (hardened screen)}{\badge{BadgeAmber}{stale only}}\layeritem{decorrelated}{lean4export + nanoda (decorrelated kernel)}{\badge{BadgeGreen}{current}}}
\sourcenote{\badge{BadgeGray}{comparator: none} no fixture in the three-implementation comparator harness (compiled / kernel-interpreted / reference-interpreted) exercises this rule.}
\ruleblock{HOL4 statement}
\sourcenote{theory \texttt{bern\_catalogue}, HOL4 development script \texttt{\texttt{bern\_catalogueScript.sml}} \badge{BadgeGreen}{script bound}}
\holhead{Definition barycentricHalfspaceBuild\_def}
\begin{holsrc}
Definition barycentricHalfspaceBuild_def:
  barycentricHalfspaceBuild t sx es =
    if es.simplex = sx then
      (case baryCoeffs1D sx.complex sx.vmask of
       | NONE => NONE
       | SOME cs =>
         (case MAP (λ(v,off,slp). f4VertexAtom (extContainsNat es v) off slp t)
               cs of
          | [x; y] => SOME (FAnd x y)
          | _ => NONE))
    else NONE
End
\end{holsrc}
\holhead{Definition barycentricHalfspaceFormB\_def}
\begin{holsrc}
Definition barycentricHalfspaceFormB_def:
  barycentricHalfspaceFormB t sx es fm ⇔
    barycentricHalfspaceBuild t sx es = SOME fm
End
\end{holsrc}
\ruleblock{HOL verifier receipts}
\badge{BadgeBlue}{documentary} \texttt{hol-strict} CLEAN census over the ported checker theory:
\texttt{bern\_catalogue}, 472 theorems, 0 cheat-tagged, 0 nonstandard
axioms (whitelist: \texttt{DISK\_THM}; the four \texttt{bool} axioms) ---
verification run of 2026-07-22 on an isolated build host, against a
content-addressed checkout of the exact sources cited here --- the
re-attestation the \rulename{ConvElim} cell\textquotesingle s earlier
caveat awaited (the canonical-form revision of 2026-07-18 postdated the
prior census) is this run (documentary citations --- no signed HOL
store yet).
Scope note: the shape side condition (\texttt{termShape}) rides the
un-ported hypernet typing island, deferred to the port\textquotesingle s
proof-object stage.
Semantic stratum, also ported: the soundness twin
\texttt{sound\_convIntro} (theory \texttt{bern\_tr\_fragclose}:
\texttt{hol-strict} clean, 26 theorems, 0 offences, verification run of
2026-07-22 on an isolated build host, content-addressed at the same
sources), consumed by the certified-fragment keystone
(\texttt{checkPf\_sound\_frag}, theory \texttt{bern\_tr\_splice}, 20
theorems, 0 offences, same run): a reference-checker accept inside the
proven wire fragment yields sequent validity.

\sourcenote{HOL receipts are documentary: quoted results of \texttt{hol-strict} runs
and census rows from the HOL4 development\textquotesingle s build
evidence, not minted receipt artefacts. A HOL receipt store (the Lean
store\textquotesingle s twin) is named future work.
}
\ruleblock{Proof}

\sourcenote{Genre: core rule (tier 1). Setting, region-atom semantics,
the checkable fragment, and the canonical halfspace form as fixed in the
\rulename{ConvElim} section.}

\textbf{Claim} (\texttt{sound\_convIntro}).\enspace The payload supplies
the region data $(x, \Delta, S)$ and two checked conditions: the premise
concludes the canonical halfspace form for $(x, \Delta, S)$, and (the
shape condition) $\bshape_H(x) = [\,d\,]$ where $d$ is the ambient
dimension --- the region term is statically a rank-1 tensor of the
ambient extent. If the premise is semantically valid, then
$\Phi \vdash x \in \mathcal{R}(\Delta, S)$ is semantically valid.

\begin{blemma}[evaluation recovered from the premise]
\label{lem:convintro-eval}
Suppose $\sigma$ satisfies the built conjunction. Then $x$ evaluates at
$\sigma$.
\end{blemma}
\begin{proof}
The first vertex atom holds, so its left side --- the canonicalized
affine image of $x$, built from constants and $x$ by product and sum ---
evaluates; evaluation of a compound is defined only if evaluation of
every part is --- the evaluator is strict, each node's value built
from its children's values, so a defined parent forces defined
children --- and inverting it at the product node (and sum
node, when the offset is present) yields a value for $x$ itself.
(\texttt{f4VertexAtom\_holds\_inv}.)
\end{proof}

\begin{blemma}[shape pinned by the payload field]
\label{lem:convintro-shape}
With the shape condition, the evaluation of Lemma~\ref{lem:convintro-eval}
is a rank-1 tensor of extent $1$ (the fragment has $d = 1$), hence
singleton-shaped, with scalar $q$.
\end{blemma}
\begin{proof}
Broadcast-shape soundness: a defined evaluation of a $\bshape$-typed term
lands at the declared shape (\texttt{Term.evalQ\_bshape}), here $[\,1\,]$
by the build's dimension gate (\rulename{ConvElim},
Lemma~\ref{lem:convelim-spec}).
\end{proof}

\begin{bthm}[soundness of \rulename{ConvIntro}]\label{thm:convintro-sound}
Under the payload's checked hypotheses, if $\Phi \vdash F$ is
semantically valid for the built conjunction $F$, so is
$\Phi \vdash x \in \mathcal{R}(\Delta, S)$.
\end{bthm}
\begin{proof}
Fix an admissible $\sigma \models \Phi$; premise validity gives both
vertex atoms at $\sigma$. Lemmas~\ref{lem:convintro-eval} and
\ref{lem:convintro-shape} produce the singleton evaluation $q$ of $x$.
Reading each held atom through the canonical-atom characterization
(\rulename{ConvElim}, Lemma~\ref{lem:convelim-atom}) yields the four
conditional sign facts on $\lambda_a(q), \lambda_b(q)$. The two-vertex
introduction assembles the masked-region witness from exactly these: the
weights $\lambda_a(q), \lambda_b(q)$ sum to $1$ and reproduce $q$ by
construction (affine identities, valid since $V_0 \neq V_1$), and the
sign discipline --- nonnegative on core vertices, nonpositive on
$S$-vertices --- is the four facts (\texttt{maskRegion\_pair\_intro}).
Together with the build's tie of $S$ to $\Delta$ and the shape of
Lemma~\ref{lem:convintro-shape}, this is the region atom's holding, term
slot included.
\end{proof}

\begin{bremark}[why the shape condition is not decoration]
\label{rem:convintro-cex}
Drop it and the rule is \emph{unsound}. Over $\Delta =
\mathrm{conv}(1, 2)$, $S = \emptyset$, instantiate the region term at the
\emph{closed scalar} $t := 3/2$. The built premise is then a conjunction
of two closed atoms asserting $\lambda_a(3/2) \geq 0$ and
$\lambda_b(3/2) \geq 0$ --- both true ($\lambda_a = \lambda_b = 1/2$), so
the conjunction holds at \emph{every} valuation, and taking $\Phi$ to be
that very conjunction, \rulename{Ax} proves the premise. The un-gated
introduction would conclude
$\Phi \vdash (\tfrac{3}{2} : \mathrm{Term}) \in \mathcal{R}(\Delta,
\emptyset)$ --- unsatisfiable: a constant evaluates to a rank-0 tensor,
and the region atom demands rank 1. Satisfiable premise, unsatisfiable
conclusion. The mechanism is the same broadcast liberality that broke the
un-gated \rulename{EqElim}: the halfspace atoms constrain the
\emph{values} of the affine images and are perfectly happy with a scalar
$x$; only the region atom cares about the tensor's shape. The shape field
restores, at the introduction rule, the well-formedness that region
atoms carry in the source calculus's typing discipline --- one repair
family (the \rulename{EqElim}/\rulename{Norm} typing pattern), one
failure family.
\end{bremark}

\begin{bremark}[elimination and introduction are exact inverses on the
fragment]\label{rem:convintro-inv}
On singleton-shaped, well-shaped region terms, Lemmas
\ref{lem:convelim-elim} (elimination) and the introduction step above are
mutually inverse readings of the same determinate-weight geometry:
membership $\leftrightarrow$ the four sign facts $\leftrightarrow$ the
canonical conjunction. The round trip loses nothing --- which is why the
pair can serve as the sole interface between region atoms and the
arithmetic fragment, with no completeness caveat inside the checkable
fragment.
\end{bremark}

\begin{anchors}
Lemma \ref{lem:convintro-eval} & \texttt{f4VertexAtom\_holds\_inv}\\
Lemma \ref{lem:convintro-shape} & \texttt{Term.evalQ\_bshape}; payload field \texttt{termShape}\\
canonical-atom reading & \texttt{f4VertexAtom\_holds\_iff} (\rulename{ConvElim} section)\\
witness assembly & \texttt{maskRegion\_pair\_intro}\\
Theorem \ref{thm:convintro-sound} & \texttt{sound\_convIntro}; wired as core \texttt{RuleData.convIntro} into \texttt{ProofObject.sound}\\
Remark \ref{rem:convintro-cex} & the counterexample inventory, \texttt{Soundness.lean} header\\
\end{anchors}


\rulesection{14}{ArithAtom}{core rule}
\label{blcatrule-ArithAtom}
\ruleblock{The rule (Steele metanotation)}
\sourcenote{as realized in the mechanized development; source row: \texttt{\texttt{rules.toml}} --- \emph{ArithAtom}. The metanotation statement is typeset by hand for notational quality; its mechanized ground truth is the Lean statement below, extracted from the sources at every build of this book and hash-bound. The rule \emph{roster} is machine-checked (every mechanized rule tag round-trips to exactly one catalogue row); the statement \emph{text} is editorial, checked by eye against the extracted Lean beside it.}
\begin{mathpar}
\inferrule*[Right=\rulename{ArithAtom}]{
  \Phi \vdash (t_1 - t_2) \bowtie 0
}{
  \Phi \vdash t_1 \bowtie t_2
}
\\
\inferrule*[Right=\rulename{ArithAtom}$^{-1}$]{
  \Phi \vdash t_1 \bowtie t_2
}{
  \Phi \vdash (t_1 - t_2) \bowtie 0
}
\end{mathpar}
with $\bowtie \in \{<, \leq, =, \geq, >\}$.

\par\medskip
\rulename{ArithAtom} toggles a comparison atom between its two
spellings: the \emph{direct form} $t_1 \bowtie t_2$ and the
\emph{zero-normal form} $(t_1 - t_2) \bowtie 0$, in both directions
(one payload constructor each), for
$\bowtie \in \{<, \leq, =, \geq, >\}$. It is the adapter at the
calculus\textquotesingle s normal-form boundary: the combination rules
(\rulename{Farkas}, \rulename{Syzygy}) and the canonical conclusions of
the geometry rules live in $(\cdot) \leq 0$-normal shapes, while
human-facing specifications, box premises of the relaxation rules, and
final theorem statements live in direct forms. A producer folds a
zero-form conclusion into the literal bound a downstream rule wants, or
unfolds a specification atom so its polynomial content can be matched.

What the checker recomputes: two syntactic form equations --- the
premise (respectively conclusion) must be \emph{literally} the atom
built from the same payload terms with the syntactic difference
$t_1 - t_2$ against the literal constant $0$ --- plus the comparator
gate $\bowtie \neq \mathord{\neq}$ and the shared-context equation.
Deliberately severe: no term evaluating to zero, no algebraic
rearrangement qualifies (a producer wanting those routes through
\rulename{Norm} first); the severity is what keeps the side conditions
decidable by construction and the trusted surface at the two form
equations, with the comparator gate and the context equation beside
them.

There is no droppable typing hypothesis, and that is the
rule\textquotesingle s quiet achievement: both forms mention the same
wires, and whenever either form holds, both evaluations exist and
transport --- in the
fold direction the difference\textquotesingle s definedness arrives with
the premise, and in the unfold direction it is manufactured from the
fact that a \emph{holding} elementwise comparison already certifies
broadcast compatibility, on exactly which subtraction is total. The
comparator gate is the one non-structural field: the rational semantics
would accept $\neq$ (the entrywise bridge is comparator-uniform), and
the gate instead pins the source calculus\textquotesingle s comparator
set and keeps the development\textquotesingle s auxiliary three-valued
screen uniform across the atom-shape rules (see the proof
section\textquotesingle s closing remark --- the screen itself would
survive $\neq$ here, unlike at the ground-fact rule).

\cataloguenote{The transported object is the per-entry rational comparison
$(p - q) \bowtie 0 \iff p \bowtie q$ --- translation invariance of the
order of $\Q$ --- lifted entrywise through the broadcast discipline by
an invertible totality argument. Exact arithmetic end to end: no
rounding, no tolerance, in either direction.
}
\ruleblock{Lean statement}
\declhead{BernsteinLogic.Soundness}{QTensor.map2\_defined\_of\_rel2}{\badge{BadgeGreen}{statement bound}}
\begin{leansrc}
/-- A holding rel2 forces broadcast-compatible shapes, so map2 succeeds. -/
theorem QTensor.map2_defined_of_rel2 {r : ℚ → ℚ → Prop} {a b : QTensor}
    (h : QTensor.rel2 r a b) (f : ℚ → ℚ → ℚ) :
    ∃ c, QTensor.map2 f a b = some c
\end{leansrc}
\declhead{BernsteinLogic.Soundness}{QTensor.rel2\_sub\_zero\_iff}{\badge{BadgeGreen}{statement bound}}
\begin{leansrc}
/-- The ArithAtom bridge, tensor-general: comparing a subtraction against
rank-0 zero is comparing the operands directly, in every broadcast case. -/
theorem QTensor.rel2_sub_zero_iff {b : BowTie} {al ar d : QTensor}
    (hd : QTensor.map2 (· - ·) al ar = some d) :
    QTensor.rel2 (BowTie.holdsQ b) d ⟨[], [0], rfl⟩
      ↔ QTensor.rel2 (BowTie.holdsQ b) al ar
\end{leansrc}
\declhead{BernsteinLogic.Soundness}{Term.evalQ\_sub\_inv}{\badge{BadgeGreen}{statement bound}}
\begin{leansrc}
theorem Term.evalQ_sub_inv {H : Hypernet} (σ : H.Valuation) {l r : Term}
    {q : QTensor} (h : Term.evalQ σ (.sub l r) = some q) :
    ∃ a b, Term.evalQ σ l = some a ∧ Term.evalQ σ r = some b ∧
      QTensor.map2 (· - ·) a b = some q
\end{leansrc}
\declhead{BernsteinLogic.Soundness}{sound\_arithAtom\_fold}{\badge{BadgeGreen}{statement bound}}
\begin{leansrc}
theorem sound_arithAtom_fold {H : Hypernet} {seq premise : Sequent}
    {termL termR : Term} {b : BowTie}
    (ih : premise.semanticallyValid H)
    (hpre : seq.pre = premise.pre)
    (hprem : atomicCompareZero b termL termR premise.post)
    (hconcl : atomicCompareForm b termL termR seq.post) :
    seq.semanticallyValid H
\end{leansrc}
\declhead{BernsteinLogic.Soundness}{sound\_arithAtom\_unfold}{\badge{BadgeGreen}{statement bound}}
\begin{leansrc}
theorem sound_arithAtom_unfold {H : Hypernet} {seq premise : Sequent}
    {termL termR : Term} {b : BowTie}
    (ih : premise.semanticallyValid H)
    (hpre : seq.pre = premise.pre)
    (hprem : atomicCompareForm b termL termR premise.post)
    (hconcl : atomicCompareZero b termL termR seq.post) :
    seq.semanticallyValid H
\end{leansrc}
\par\medskip
\textbf{The two form equations, in full.}\enspace Everything this rule
checks is displayed below: two one-line syntactic equations and the
payload\textquotesingle s two constructors (one per direction) with the
comparator gate \texttt{notNE}. No valuation is consulted; semantic
truth enters only through the soundness lemmas.

\declhead{BernsteinLogic.Core.Catalogue}{atomicCompareForm}{\badge{BadgeGreen}{definition bound}}
\begin{leansrc}
def atomicCompareForm (b : BowTie) (l r : Term) (F : Formula) : Prop :=
  F = .atom b l r
\end{leansrc}
\declhead{BernsteinLogic.Core.Catalogue}{atomicCompareZero}{\badge{BadgeGreen}{definition bound}}
\begin{leansrc}
def atomicCompareZero (b : BowTie) (l r : Term) (F : Formula) : Prop :=
  F = .atom b (.sub l r) (.const 0)
\end{leansrc}
\declhead{BernsteinLogic.Core.Catalogue}{PArithAtomData}{\badge{BadgeGreen}{definition bound}}
\begin{leansrc}
inductive PArithAtomData : Hypernet → Sequent → Type where
  | fold   (H : Hypernet) (seq : Sequent)
           (premise : Sequent) (premiseProof : ProofObject H premise)
           (termL termR : Term) (b : BowTie)
           (notNE          : b ≠ .ne)
           (preEq          : seq.pre = premise.pre)
           (premiseIsZero  : atomicCompareZero b termL termR premise.post)
           (conclIsDirect  : atomicCompareForm b termL termR seq.post)
           : PArithAtomData H seq
  | unfold (H : Hypernet) (seq : Sequent)
           (premise : Sequent) (premiseProof : ProofObject H premise)
           (termL termR : Term) (b : BowTie)
           (notNE           : b ≠ .ne)
           (preEq           : seq.pre = premise.pre)
           (premiseIsDirect : atomicCompareForm b termL termR premise.post)
           (conclIsZero     : atomicCompareZero b termL termR seq.post)
           : PArithAtomData H seq
\end{leansrc}
\ruleblock{Correspondence (Steele display $\leftrightarrow$ Lean)}
\badge{BadgeGreen}{reviewer-certified against the displayed texts}
\sourcenote{The dictionary below is a reviewer-certified statement of how the mechanized statements above interpret the displayed rule: display metavariable $\mapsto$ Lean binder, each display hypothesis $\mapsto$ the premise or side condition carrying it, the display conclusion $\mapsto$ the goal --- with any deliberate difference in width between display and mechanization named outright. Certification is recorded as a hash triple over (the displayed cell, the extracted statement texts, this dictionary\textquotesingle s own text); if any of the three changes, this block turns amber until a reviewer re-certifies. The correspondence claim itself remains editorial --- reviewer-grade, not machine-checked; its machine-checked successor types each soundness statement against the rule as a structured object, and supersedes this block per rule as it lands.}
Judgment frame as in \rulename{Ax}; sub-derivation premise as at
\rulename{Weaken}.
\begin{itemize}
\item the display\textquotesingle s two inference lines $\mapsto$ the
  lemma pair, mirroring the payload\textquotesingle s two
  constructors: \rulename{ArithAtom} (zero-form premise, direct
  conclusion) $\mapsto$ \texttt{sound\_arithAtom\_fold};
  \rulename{ArithAtom}$^{-1}$ $\mapsto$
  \texttt{sound\_arithAtom\_unfold}.
\item $t_1, t_2 \mapsto$ \texttt{termL termR : Term}; $\bowtie
  \mapsto$ \texttt{b : BowTie}; the two spellings $\mapsto$ the form
  predicates quoted below, each a one-line syntactic equation:
  \texttt{atomicCompareForm} (the atom is literally $t_1 \bowtie t_2$)
  and \texttt{atomicCompareZero} (literally $(t_1 - t_2) \bowtie 0$,
  syntactic difference against the literal constant).
\item $\Phi \mapsto$ the shared context via \texttt{hpre}; the
  premise judgment $\mapsto$ \texttt{ih}; the conclusion $\mapsto$ the
  goal \texttt{seq.semanticallyValid H}. Certified alongside:
  \texttt{QTensor.map2\_defined\_of\_rel2}, the totality bridge with
  which the unfold direction manufactures the
  difference\textquotesingle s definedness from the holding premise.
\end{itemize}
Width note, deliberate and load-bearing (the \rulename{PConstExact}
pattern): the display restricts $\bowtie \in \{<, \leq, =, \geq, >\}$,
while both lemmas are proven for \emph{all} constructors of
\texttt{BowTie}, including $\neq$ --- no $\bowtie \neq \mathord{\neq}$
hypothesis appears in either statement. The exclusion is an
admission-layer fact: the payload\textquotesingle s \texttt{notNE}
field refuses $\neq$ certificates to protect the three-valued screen
the commentary describes. The mechanized statements are therefore
strictly wider than the displayed rule, and the displayed rule is
exactly the admitted fragment of them.

\ruleblock{Lean verifier receipts}
\sourcenote{module \texttt{BernsteinLogic.Soundness}; bound build \texttt{a20d90e66c13\ldots} (bound from receipt a9c0e34d3e3c9c8c, 2026-08-03T21:30:36Z). A receipt records that named checking layers accepted a build: (module, compiled-artifact SHA-256, environment SHA-256, checker layers), minted when the checking gate ran; the store receipts attesting the bound build are listed with the layers each run credited. Nothing here asks the reader to trust the store: every layer is an independently runnable program over the quoted sources --- rebuild the module and re-run the checkers to reproduce the verdict. Receipts attesting superseded builds credit no layer and are omitted from this rendering (the count of omissions, if any, is noted below).}
\begin{receipttable}
\receiptrow{a9c0e34d3e3c9c8c}{2026-08-03T21:30:36Z}{build+decorrelated}{\badge{BadgeGreen}{current}}
\end{receipttable}
\sourcenote{21 further receipt(s) in the store attest superseded builds of this module and are omitted from this rendering.}
\sourcenote{module \texttt{BernsteinLogic.Core.Catalogue}; bound build \texttt{cbc92e976a1d\ldots} (bound from receipt 5abdb3589707b5df, 2026-08-03T21:30:33Z). A receipt records that named checking layers accepted a build: (module, compiled-artifact SHA-256, environment SHA-256, checker layers), minted when the checking gate ran; the store receipts attesting the bound build are listed with the layers each run credited. Nothing here asks the reader to trust the store: every layer is an independently runnable program over the quoted sources --- rebuild the module and re-run the checkers to reproduce the verdict. Receipts attesting superseded builds credit no layer and are omitted from this rendering (the count of omissions, if any, is noted below).}
\begin{receipttable}
\receiptrow{24804e66edba0bc9}{2026-08-03T20:25:09Z}{build+decorrelated}{\badge{BadgeGreen}{current}}
\receiptrow{5abdb3589707b5df}{2026-08-03T21:30:33Z}{build+decorrelated}{\badge{BadgeGreen}{current}}
\end{receipttable}
\sourcenote{18 further receipt(s) in the store attest superseded builds of this module and are omitted from this rendering.}
\layerline{\layeritem{build}{lake build (elaboration)}{\badge{BadgeGreen}{current}}\layeritem{replay}{leanchecker (official kernel replay)}{\badge{BadgeAmber}{stale only}}\layeritem{hardened}{safe\_verify (hardened screen)}{\badge{BadgeAmber}{stale only}}\layeritem{decorrelated}{lean4export + nanoda (decorrelated kernel)}{\badge{BadgeGreen}{current}}}
\sourcenote{\badge{BadgeGray}{comparator: none} no fixture in the three-implementation comparator harness (compiled / kernel-interpreted / reference-interpreted) exercises this rule.}
\ruleblock{HOL4 statement}
\sourcenote{theory \texttt{bern\_catalogue}, HOL4 development script \texttt{\texttt{bern\_catalogueScript.sml}} \badge{BadgeGreen}{script bound}}
\holhead{Definition atomicCompareFormB\_def}
\begin{holsrc}
Definition atomicCompareFormB_def:
  atomicCompareFormB b l r fm ⇔ fm = FAtom b l r
End
\end{holsrc}
\holhead{Definition atomicCompareZeroB\_def}
\begin{holsrc}
Definition atomicCompareZeroB_def:
  atomicCompareZeroB b l r fm ⇔ fm = FAtom b (TmSub l r) (TmConst 0)
End
\end{holsrc}
\ruleblock{HOL verifier receipts}
\badge{BadgeBlue}{documentary} \texttt{hol-strict} CLEAN census over the ported checker theory:
\texttt{bern\_catalogue}, 472 theorems, 0 cheat-tagged, 0 nonstandard
axioms (whitelist: \texttt{DISK\_THM}; the four \texttt{bool} axioms) ---
verification run of 2026-07-22 on an isolated build host, against a
content-addressed checkout of the exact sources cited here --- the
earlier stale-run caveat (a 2026-07-18 script revision postdating the
prior census) is retired by this run, which attests the current text.
There is no signed HOL receipt store yet --- these citations are
documentary.
Semantic stratum, also ported: the soundness twins
\texttt{sound\_arithAtom\_fold} (theory \texttt{bern\_tr\_ready8}:
\texttt{hol-strict} clean, 8 theorems, 0 offences, verification run of
2026-07-22 on an isolated build host, content-addressed at the same
sources) and \texttt{sound\_arithAtom\_unfold} (theory
\texttt{bern\_tr\_fragclose}, 26 theorems, 0 offences, same run),
consumed by the certified-fragment keystone
(\texttt{checkPf\_sound\_frag}, theory \texttt{bern\_tr\_splice}, 20
theorems, 0 offences, same run): a reference-checker accept inside the
proven wire fragment yields sequent validity.

\sourcenote{HOL receipts are documentary: quoted results of \texttt{hol-strict} runs
and census rows from the HOL4 development\textquotesingle s build
evidence, not minted receipt artefacts. A HOL receipt store (the Lean
store\textquotesingle s twin) is named future work.
}
\ruleblock{Proof}

\sourcenote{Genre: core rule (tier 1). Setting as fixed in the
\rulename{EqElim} section. The rule is a pair: one payload constructor
per direction, sharing the comparator gate; both directions are proven,
and the two proofs are deliberately \emph{not} mirror images --- the
asymmetry is the interesting content.}

\textbf{Claim} (\texttt{sound\_arithAtom\_fold},
\texttt{sound\_arithAtom\_unfold}).\enspace The payload supplies terms
$t_1, t_2$ and a comparator $\bowtie \in \{<, \leq, =, \geq, >\}$ (the
gate $\bowtie \neq \mathord{\neq}$ is a checked field). The two atom
spellings in play, both recomputed syntactically by the checker from the
same payload terms, are the \emph{direct form} $t_1 \bowtie t_2$ and the
\emph{zero-normal form} $(t_1 - t_2) \bowtie 0$, where the subtrahend
side is the literal constant $0$ and the left side the literal syntactic
difference. \emph{Fold}: premise concludes the zero-normal form,
conclusion is the direct form. \emph{Unfold}: the converse. In both
directions, validity of the premise sequent gives validity of the
conclusion sequent, over the shared $\Phi$.

\begin{blemma}[scalar bridge]\label{lem:arith-scalar}
For $p, q \in \Q$ and $\bowtie$ any of the five comparators:
$(p - q) \bowtie 0 \iff p \bowtie q$.
\end{blemma}
\begin{proof}
Add $q$ to both sides, respectively subtract it --- the order of $\Q$ is
translation-invariant. The argument is uniform in $\bowtie$; it would,
for that matter, be equally valid for $\neq$ --- the comparator gate is
\emph{not} forced at this level (\autoref{rem:arith-ne}).
\end{proof}

\begin{blemma}[tensor bridge]\label{lem:arith-tensor}
Let $A, B$ be tensors whose elementwise difference is defined, say $D$.
Then $D$ stands in the elementwise lift of $\bowtie$ against the rank-0
tensor $0$ iff $A, B$ stand in the elementwise lift of $\bowtie$.
\end{blemma}
\begin{proof}
$D$'s entries are the pointwise differences under the broadcast
discipline: at equal shapes, entry $i$ of $D$ is $A_i - B_i$; with
rank-0 $A$, entry $i$ is $A_{\mathrm{scalar}} - B_i$; with rank-0 $B$,
$A_i - B_{\mathrm{scalar}}$ --- and $D$ takes the positive-rank
operand's shape (the broadcast join). The left side of the claimed
equivalence reads every entry of $D$ against the scalar $0$; the right
side reads the operands in the corresponding branch of the elementwise
lift. In each of the three branches the two readings pair off the same
scalar comparisons entry by entry, and Lemma~\ref{lem:arith-scalar}
closes each pair. (\texttt{QTensor.rel2\_sub\_zero\_iff}.)
\end{proof}

\begin{blemma}[definedness, fold direction]\label{lem:arith-def-fold}
If the zero-normal atom holds at $\sigma$, then both $t_1$ and $t_2$
evaluate at $\sigma$, and their elementwise difference is the evaluation
of $t_1 - t_2$.
\end{blemma}
\begin{proof}
The holding atom includes a defined evaluation of the syntactic
difference; the evaluator is structural and strict (the Setting: a
compound evaluates exactly when its parts evaluate and their shapes
join), so inverting it at the
subtraction node yields defined evaluations of both operands and the
elementwise-difference equation. (\texttt{Term.evalQ\_sub\_inv}.)
\end{proof}

\begin{blemma}[definedness, unfold direction]\label{lem:arith-def-unfold}
If the direct atom holds at $\sigma$, then the elementwise difference of
the two (defined, by the atom's holding) evaluations is itself defined.
\end{blemma}
\begin{proof}
A \emph{holding} elementwise comparison forces broadcast-compatible
shapes --- the lift's fourth branch, unequal positive ranks, is empty ---
and elementwise subtraction is total on exactly the broadcast-compatible
pairs. A holding comparison is thereby itself a compatibility
certificate. (\texttt{QTensor.map2\_defined\_of\_rel2}.)
\end{proof}

\begin{bthm}[soundness of \rulename{ArithAtom}, both directions]
\label{thm:arith-sound}
Over shared $\Phi$: validity of $\Phi \vdash (t_1 - t_2) \bowtie 0$
gives validity of $\Phi \vdash t_1 \bowtie t_2$, and conversely.
\end{bthm}
\begin{proof}
Fix an admissible $\sigma \models \Phi$ in either direction.
\emph{Fold}: the premise's zero-normal atom holds; its constant side
evaluates to the rank-0 tensor $0$; Lemma~\ref{lem:arith-def-fold}
recovers evaluations $A, B$ of $t_1, t_2$ with difference $D$, and the
premise's relation is $D \bowtie 0$ elementwise;
Lemma~\ref{lem:arith-tensor} converts it to $A \bowtie B$ elementwise
--- the direct atom, on the operands' own evaluations. \emph{Unfold}:
the premise's direct atom holds with evaluations $A, B$;
Lemma~\ref{lem:arith-def-unfold} manufactures the difference $D$; the
evaluator assembles $\mathrm{ev}_\sigma(t_1 - t_2) = D$ from its
operands' evaluations, the constant side is again rank-0 $0$, and
Lemma~\ref{lem:arith-tensor}, read right to left, closes the zero-normal
atom.
\end{proof}

\begin{bremark}[the direction asymmetry, and payload minimalism]
\label{rem:arith-asym}
The fold direction \emph{consumes} definedness of the difference (the
premise evaluated it); the unfold direction must \emph{manufacture} it
--- and can, precisely because a holding comparison certifies broadcast
compatibility (Lemma~\ref{lem:arith-def-unfold}). The consequence is
worth stating plainly: whenever either atom \emph{holds}, both operands'
evaluations exist and transport to the other form. That is why this
payload carries \emph{no} typing or
positivity fields at all --- compare \rulename{EqRefl}'s typing pack,
\rulename{EqElim}'s shape agreement, and the positive-size fields of the
combination rules: every one of those exists to manufacture or protect
an evaluation this rule gets handed for free. (Definedness itself is not
mirrored form to form: at broadcast-incompatible operand shapes the
direct form's sides still evaluate while the zero-normal difference does
not --- but there neither atom can hold, each failing in its own way,
and holding is all the rule transports; at a compatible zero-size shape
the two forms hold together, vacuously.) The ``exact rational evaluation
transported both directions'' is exactly the pair
(Lemma~\ref{lem:arith-scalar}, Lemma~\ref{lem:arith-tensor}): a
translation-invariance fact of ordered $\Q$, lifted entrywise through an
invertible totality argument.
\end{bremark}

\begin{bremark}[the syntactic pins]\label{rem:arith-pins}
Both form conditions are single syntactic equations recomputed from the
payload --- the zero-normal form demands the \emph{literal} constant $0$
and the \emph{literal} difference term of the same $t_1, t_2$, not any
term evaluating to zero or any algebraic rearrangement. A producer
wanting $t_1 \bowtie t_2$ from, say, $(t_2 - t_1) \cdot (-1) \bowtie 0$
routes through \rulename{Norm} first; this rule does arithmetic on
\emph{atom shape} only. That severity is what makes the side conditions
decidable by construction and keeps the rule's checked surface at two
form equations, the shared-context equation, and one comparator gate.
\end{bremark}

\begin{bremark}[why $\neq$ is excluded]\label{rem:arith-ne}
Not for the reason one might guess --- and not for its sibling's reason
either. Lemmas \ref{lem:arith-scalar} and
\ref{lem:arith-tensor} are comparator-uniform, and the rational
semantics of this section would accept $\neq$ without change. Nor would
the development's auxiliary three-valued screen (the countermodel
construction that polices vacuity properties, developed at the $\vee$E
section) break here: both spellings of a $\neq$ atom mention the same
wires, and the screen's disequality clause reads only designated-wire
membership, so fold and unfold preserve its value --- unlike at the
premise-free ground-fact rule, where a wire-free $\neq$ atom's screen
value genuinely breaks the invariant (the \rulename{PConstExact}
section). The
comparator set $\{<, \leq, =, \geq, >\}$ is the source calculus's own
set for this rule, and the payload's gate keeps it checked: the screen's
per-rule argument stays a one-clause conclusion-shape case (the gate
vacates its $\neq$ branch), and the two atom-shape rules stay aligned.
$\neq$ remains in the
formula grammar --- it is simply not this rule's business to move it
across the zero-normal boundary.
\end{bremark}

\begin{anchors}
Lemma \ref{lem:arith-scalar} & inlined in \texttt{QTensor.rel2\_sub\_zero\_iff}'s entrywise closure\\
Lemma \ref{lem:arith-tensor} & \texttt{QTensor.rel2\_sub\_zero\_iff}\\
Lemma \ref{lem:arith-def-fold} & \texttt{Term.evalQ\_sub\_inv}\\
Lemma \ref{lem:arith-def-unfold} & \texttt{QTensor.map2\_defined\_of\_rel2}\\
Theorem \ref{thm:arith-sound}, fold & \texttt{sound\_arithAtom\_fold}\\
Theorem \ref{thm:arith-sound}, unfold & \texttt{sound\_arithAtom\_unfold}; both wired as core \texttt{RuleData.arithAtom} into \texttt{ProofObject.sound}\\
form equations, comparator gate & \texttt{atomicCompareForm}, \texttt{atomicCompareZero}; payload field \texttt{notNE} (\texttt{PArithAtomData})\\
\end{anchors}


\rulesection{15}{BernCone}{core rule}
\label{blcatrule-BernCone}
\ruleblock{The rule (Steele metanotation)}
\sourcenote{as realized in the mechanized development; source row: \texttt{\texttt{rules.toml}} --- \emph{BernCone}. The metanotation statement is typeset by hand for notational quality; its mechanized ground truth is the Lean statement below, extracted from the sources at every build of this book and hash-bound. The rule \emph{roster} is machine-checked (every mechanized rule tag round-trips to exactly one catalogue row); the statement \emph{text} is editorial, checked by eye against the extracted Lean beside it.}
\begin{mathpar}
\inferrule*[Right=\rulename{BernCone}]{
  \Delta = \mathrm{conv}(V_0, \ldots, V_n) \subset \R^n \text{ an } n\text{-simplex};
  \quad S \subsetneq \{0, 1, \ldots, n\}
  \\
  \Phi \vdash x \in \mathcal{R}(\Delta, S)
  \\
  p \colon \R^n \to W \text{ polynomial map of degree } d, \text{ with coefficient tensor } (b_\alpha)_{|\alpha|=d}
  \\
  \wires(p) = \emptyset \;\vee\; \bigl(x = w \,\wedge\, \wires(p) \subseteq \{w\}\bigr)
  \\
  S \neq \emptyset \Rightarrow P \text{ conic}; \qquad S = \emptyset \Rightarrow P \text{ convex suffices}
  \\
  \bigl(p \ \text{idx-free, with } p \ \text{constant} \vee n = 1\bigr) \;\vee\; \bigl(p \ \text{projection-indexed, indices} < n,\ S = \emptyset\bigr)
  \\
  \text{for every } j \in \N^S \text{ with } |j| \leq d, \text{ every degree-}d \text{ Bernstein coefficient of } P_j \text{ on } F_S \text{ satisfies } P
}{
  \Phi \vdash P(p(x))
}
\end{mathpar}

\ruleflavor{the disjunctive premise on $p$ is the checker\textquotesingle s scope guard: coefficient checks on coordinate polynomials over ambient dimension $\neq 1$ are refused rather than attempted outside the proven kernels.}
\par\medskip
\rulename{BernCone} is the calculus\textquotesingle s arithmetic-closure
primitive: from a \emph{geometric} fact (a wire value lies in a region)
it derives an \emph{analytic} one of unbounded degree (a polynomial of
that value satisfies a predicate) --- where the conversion rules exchange
region atoms with their defining affine sign conditions,
\rulename{BernCone} is the rule that concludes something new. From a proven region membership
$\Phi \vdash x \in \mathcal{R}(\Delta, S)$ and a degree-$d$ polynomial
term $p$, it concludes $\Phi \vdash P(p(x))$ for any predicate $P$ of the
accepted grammar, provided every certified Bernstein coefficient of $p$
over the region satisfies $P$. A producer reaches for it wherever a
Lyapunov decrease, an invariant bound, or an output enclosure must be
established over a whole region rather than at sample points: choose a
degree, optionally subdivide (\rulename{Subdivide} is the completeness
partner --- enclosures converge under degree elevation and subdivision),
and emit the payload.

What the checker recomputes, given the payload (simplex, extension set,
degree, polynomial term, predicate): that the extension set is tied to
the premise\textquotesingle s simplex and the simplex is
full-dimensional; that the polynomial term translates into the
region\textquotesingle s ambient coordinates (two routes: an idx-free
single-wire canonical form, or the projection-indexed coordinate form);
and then --- the heart --- \emph{every} forward-difference Bernstein
coefficient of the translated polynomial over the region, by exact
$\Q$ arithmetic, testing each against the predicate. No coefficient
tensors ride the wire: the certificate names the polynomial and the
region, and the checker rebuilds the entire coefficient table itself
(the recompute pattern), so a forged table is unrepresentable rather
than detected.

Each hypothesis is load-bearing. Drop the wire tie
(\texttt{polyTie}) and a certificate could prove its region fact about
one wire while running the coefficient check on a polynomial over a
different one --- concluding a bound about a value the premise never
constrained. Drop the predicate-shape clause and a convex-but-not-conic
predicate falsely transfers along an extension ray (the proof section
computes a two-line counterexample: certified coefficients $\{0, -1\}$
both satisfy $\square \geq -1$, yet the polynomial dips to $-2$ on the
ray). Drop the coefficient recomputation and the conclusion is unrelated
to the polynomial. Drop the simplex tie or full-dimensionality and the
barycentric reading of the premise breaks. The checker\textquotesingle s
own scope guard (constant canonical forms at any dimension; coordinate
polynomials only over 1-dimensional regions when the extension set is
nonempty) is part of the soundness story: everything beyond the proven
fragment is refused, never accepted on faith.

\cataloguenote{Predicates enter through a deliberately closed-world grammar: finite
conjunctions of scalar hole-atoms $s \cdot \square \bowtie c$ with
$s > 0$ and $\bowtie \in \{\leq, \geq, =\}$, classified for convexity
(always) and conicity ($c = 0$ throughout). Comparators outside the
grammar are refused --- conservatism that can lose completeness, never
soundness. The elementwise atom semantics batch the conclusion: a
tensor-shaped polynomial term gives per-entry predicate application for
free.
}
\ruleblock{Lean statement}
\declhead{BernsteinLogic.Soundness}{sound\_bernCone\_empty}{\badge{BadgeGreen}{statement bound}}
\begin{leansrc}
/-- BernCone at S = ∅ — the kernel assembly, at full D2 scope (leg D3:
the former `hD2` corner guard is gone — a polynomial whose wires cancel
in canonical form over a higher-dimensional region variable is handled
by the per-entry evaluation lemma `Term.evalQ_brRead`: every entry of
its semantic tensor equals the closed canonical value, which the kernel
certifies). -/
theorem sound_bernCone_empty {H : Hypernet} {seq regionPremise : Sequent}
    {termX : Term} {simplex : Simplex} {S : ExtensionSet} {degree : Nat}
    {polynomial : Term} {predicate : Formula}
    (ih : regionPremise.semanticallyValid H)
    (hpre : seq.pre = regionPremise.pre)
    (hregionPost : regionPremise.post = .inRegion termX simplex S)
    (hpolyTie : polynomial.freeWires = [] ∨
      ∃ w, termX = Term.var w ∧ polynomial.freeWires ⊆ [w])
    (hconvex : isExtensionSetEmpty S → isConvexPredicate predicate)
    (hSempty : isExtensionSetEmpty S)
    (hcoeff : allForwardDiffBernCoeffsSatisfyP simplex S degree polynomial
      predicate)
    (hpostApply : applyPredicateToTerm predicate polynomial seq.post) :
    seq.semanticallyValid H
\end{leansrc}
\declhead{BernsteinLogic.Soundness}{sound\_bernCone}{\badge{BadgeGreen}{statement bound}}
\begin{leansrc}
/-- BernCone, all extension sets (leg D3). At S = ∅ this is the
convex-predicate kernel assembly above; at S ≠ ∅ the check's own D2
scope guard confines reachable instances to constant canonical
polynomials (any dimension — `fdBernstein_sound_const`) and single-wire
polynomials over 1-dimensional regions, where a proper nonempty
extension set is exactly one of the two vertices and the region is the
extension ray (`fdBernstein_sound_dim1`, the univariate (★) identity).
The conic predicate-shape clause supplies zero/scaling/convexity
closure. -/
theorem sound_bernCone {H : Hypernet} {seq regionPremise : Sequent}
    {termX : Term} {simplex : Simplex} {S : ExtensionSet} {degree : Nat}
    {polynomial : Term} {predicate : Formula}
    (ih : regionPremise.semanticallyValid H)
    (hpre : seq.pre = regionPremise.pre)
    (hregionPost : regionPremise.post = .inRegion termX simplex S)
    (hpolyTie : polynomial.freeWires = [] ∨
      ∃ w, termX = Term.var w ∧ polynomial.freeWires ⊆ [w])
    (hpredShape : (¬ isExtensionSetEmpty S → isConicPredicate predicate)
      ∧ (isExtensionSetEmpty S → isConvexPredicate predicate))
    (hcoeff : allForwardDiffBernCoeffsSatisfyP simplex S degree polynomial
      predicate)
    (hpostApply : applyPredicateToTerm predicate polynomial seq.post) :
    seq.semanticallyValid H
\end{leansrc}
\par\medskip
\textbf{What the coefficient side-condition means --- and what it does
not.}\enspace The rule\textquotesingle s main premise appears in the
payload as the single name \texttt{allForwardDiffBernCoeffsSatisfyP}, so
this book binds its meaning by displaying the definition chain in full,
from the checker sources. Reading the chain: \texttt{satisfiesPredQ}
substitutes a candidate rational for the designated predicate-hole wire
and runs the \emph{closed} formula evaluator --- the predicate is judged
at a concrete value, by exact $\Q$ arithmetic, never at a wire.
\texttt{Term.asCoordPoly?} translates the polynomial term into the
region\textquotesingle s ambient coordinates, by one of two routes: an
idx-free term goes through its wire-keyed canonical polynomial (at most
one wire, remapped to coordinate 0), and a projection-indexed term maps
each projection $w[j]$ to coordinate variable $j$.
\texttt{bernConeCoeffCheck} then gates (simplex tie,
full-dimensionality, and the scope guard: non-constant idx-free
polynomials only over 1-dimensional regions; multivariate coordinate
forms only at $S = \emptyset$ with all projection indices in range),
computes the full forward-difference Bernstein coefficient table of the
translated polynomial over the masked vertex list, and sweeps every
grid entry through the predicate.

What a skeptical reader should confirm by eye in the definitions below:
the conclusion formula never enters the check --- only the polynomial,
the region data, and the predicate do (the conclusion is pinned
separately, as the syntactic substitution of the polynomial into the
predicate\textquotesingle s hole); nothing in the chain reads a
valuation, so the check cannot covertly assume anything about wire
values; and every \texttt{false} branch of the guards is a refusal ---
the fragment boundary (dimension, extension set, index range) fails
closed. Semantic truth is established only in \texttt{sound\_bernCone},
from the region premise\textquotesingle s validity.

\declhead{BernsteinLogic.Core.Catalogue}{satisfiesPredQ}{\badge{BadgeGreen}{definition bound}}
\begin{leansrc}
/-- Predicate satisfaction at a rational value. -/
def satisfiesPredQ (P : Formula) (c : Rat) : Bool :=
  (Formula.subst predicateHole (.const c) P).evalClosed = some true
\end{leansrc}
\declhead{BernsteinLogic.Core.Catalogue}{Term.asCoordPolyIdx?}{\badge{BadgeGreen}{definition bound}}
\begin{leansrc}
/-- The G1 coordinate-form translation: a Term whose every wire mention
is an idx projection `.idx (.var w) j` becomes a polynomial over the
ambient-coordinate variables j. Which wire is projected is deliberately
NOT recorded here — the BernCone payload's `polyTie` (B10/B7) pins every
free wire to the region variable. `none` on: a bare (vector-valued) wire
mention, and on idx of any non-var argument (v0.5 scope; einsum is the
planned v1+ extension). Index range is checked separately
(`Term.idxBounded`, conjoined by `bernConeCoeffCheck`). -/
def Term.asCoordPolyIdx? : Term → Option QArith.Poly
  | .var _ => none
  | .const c => some (QArith.Poly.const c)
  | .add l r =>
    do pure (QArith.Poly.add (← l.asCoordPolyIdx?) (← r.asCoordPolyIdx?))
  | .mul l r =>
    do pure (QArith.Poly.mul (← l.asCoordPolyIdx?) (← r.asCoordPolyIdx?))
  | .sub l r =>
    do pure (QArith.Poly.sub (← l.asCoordPolyIdx?) (← r.asCoordPolyIdx?))
  | .neg t => (Term.asCoordPolyIdx? t).map QArith.Poly.neg
  | .idx (.var _) i => some (QArith.Poly.var i)
  | .idx _ _ => none
\end{leansrc}
\declhead{BernsteinLogic.Core.Catalogue}{Term.asCoordPoly?}{\badge{BadgeGreen}{definition bound}}
\begin{leansrc}
def Term.asCoordPoly? (t : Term) : Option QArith.Poly :=
  if t.idxFree then
    let p := t.toPoly
    match QArith.Poly.vars p with
    | [] => some p
    | [_] => some (p.map fun (m, c) => (m.map fun (_, e) => (0, e), c))
    | _ => none
  else
    t.asCoordPolyIdx?
\end{leansrc}
\declhead{BernsteinLogic.Core.Catalogue}{bernConeCoeffCheck}{\badge{BadgeGreen}{definition bound}}
\begin{leansrc}
/-- The recomputation behind BernCone's main side-condition (rem:recompute-
vs-trust option (a), CORRESPONDENCE.md D-6; see original for the guard
rationale). -/
def bernConeCoeffCheck (simplex : Simplex) (S : ExtensionSet) (degree : Nat)
    (polynomial : Term) (predicate : Formula) : Bool :=
  if S.simplex = simplex then
    let cx := simplex.complex
    let verts := (List.finRange cx.nVertices).filter (fun v => simplex.contains v)
    -- full-dimensionality: an (ambDim)-simplex needs ambDim+1 vertices
    if verts.length = cx.ambDim + 1 then
      match Term.asCoordPoly? polynomial with
      | none => false
      | some p =>
        if polynomial.idxFree then
          -- legacy D2 scope: coordinate-referencing polynomials only
          -- over 1-D regions (any extension set)
          if QArith.Poly.vars p ≠ [] ∧ cx.ambDim ≠ 1 then false
          else
            let coords := verts.map fun v =>
              (List.finRange cx.ambDim).map fun j => cx.coord v j
            let smask := verts.map fun v => S.containsNat v.val
            match QArith.fdBernstein coords smask degree p with
            | some table => QArith.gridAll (satisfiesPredQ predicate) table
            | none => false
        else
          -- G1 idx coordinate form (multivariate): S = ∅ only — the
          -- general-S multivariate soundness is the v1 StarIdentity
          -- obligation — and every projection index below the ambient
          -- dimension (out-of-range reads make the conclusion
          -- unevaluable; see `Term.idxBounded`).
          if S.isEmptyB ∧ polynomial.idxBounded cx.ambDim then
            let coords := verts.map fun v =>
              (List.finRange cx.ambDim).map fun j => cx.coord v j
            let smask := verts.map fun v => S.containsNat v.val
            match QArith.fdBernstein coords smask degree p with
            | some table => QArith.gridAll (satisfiesPredQ predicate) table
            | none => false
          else false
    else false
  else false
\end{leansrc}
\declhead{BernsteinLogic.Core.Catalogue}{allForwardDiffBernCoeffsSatisfyP}{\badge{BadgeGreen}{definition bound}}
\begin{leansrc}
def allForwardDiffBernCoeffsSatisfyP
    (simplex : Simplex) (S : ExtensionSet) (degree : Nat)
    (polynomial : Term) (predicate : Formula) : Prop :=
  bernConeCoeffCheck simplex S degree polynomial predicate = true
\end{leansrc}
\ruleblock{Correspondence (Steele display $\leftrightarrow$ Lean)}
\badge{BadgeAmber}{a bound Lean statement CHANGED since certification --- re-review and re-certify}
\sourcenote{The dictionary below is a reviewer-certified statement of how the mechanized statements above interpret the displayed rule: display metavariable $\mapsto$ Lean binder, each display hypothesis $\mapsto$ the premise or side condition carrying it, the display conclusion $\mapsto$ the goal --- with any deliberate difference in width between display and mechanization named outright. Certification is recorded as a hash triple over (the displayed cell, the extracted statement texts, this dictionary\textquotesingle s own text); if any of the three changes, this block turns amber until a reviewer re-certifies. The correspondence claim itself remains editorial --- reviewer-grade, not machine-checked; its machine-checked successor types each soundness statement against the rule as a structured object, and supersedes this block per rule as it lands.}
Judgment frame as in \rulename{Ax}; sub-derivation premise as at
\rulename{Weaken}. The general statement is \texttt{sound\_bernCone};
\texttt{sound\_bernCone\_empty} (certified alongside) is its
$S = \emptyset$ specialization, through which the general proof routes
that case.
\begin{itemize}
\item $x \mapsto$ \texttt{termX : Term}; $\Delta \mapsto$
  \texttt{simplex}; $S \mapsto$ \texttt{S : ExtensionSet}; $d \mapsto$
  \texttt{degree : Nat}; $p \mapsto$ \texttt{polynomial : Term}; $P
  \mapsto$ \texttt{predicate : Formula} (the hole-atom grammar of the
  note, its hole a designated wire).
\item the region premise $\Phi \vdash x \in \mathcal{R}(\Delta, S)$
  $\mapsto$ \texttt{ih} with \texttt{hregionPost} $:$
  \texttt{regionPremise.post = .inRegion termX simplex S}; $\Phi$
  shared via \texttt{hpre}.
\item the wire-tie premise $\wires(p) = \emptyset \vee (x = w \wedge
  \wires(p) \subseteq \{w\})$ $\mapsto$ \texttt{hpolyTie}, the same
  disjunction verbatim --- note \texttt{termX = Term.var w}: the
  display\textquotesingle s ``$x = w$\textquotesingle\textquotesingle{}
  means the region term literally \emph{is} that wire.
\item the two-clause predicate premise $\mapsto$ \texttt{hpredShape}
  $:$ the conjunction of ($S$ nonempty $\to$ conic) and ($S$ empty
  $\to$ convex), by the decidable predicate classifiers.
\item the displayed coefficient sweep $\mapsto$ \texttt{hcoeff} $:$
  \texttt{allForwardDiffBernCoeffsSatisfyP simplex S degree polynomial
  predicate} --- one name whose meaning this book binds by displaying
  its definition chain in full (the intro above); the
  display\textquotesingle s first premise line (simplex tie,
  full-dimensionality) and its scope-guard premise have no Lean binder
  of their own because the checker\textquotesingle s gates certify
  them: \emph{absorbed}, as \rulename{PConstExact}\textquotesingle s
  closedness is.
\item conclusion $\Phi \vdash P(p(x))$ $\mapsto$ the goal under
  \texttt{hpostApply} $:$ \texttt{applyPredicateToTerm predicate
  polynomial seq.post} --- the conclusion is the syntactic
  substitution of the polynomial term into the
  predicate\textquotesingle s hole; the ``at
  $x$\textquotesingle\textquotesingle{} lives in the wire tie, not in
  a separate application node.
\end{itemize}
Width note, two-sided and honest: the statement is \emph{proven} at
every hypernet, sequent, and payload --- and \texttt{hcoeff} is
satisfiable only inside the displayed scope-guard premise (constant
canonical forms at any dimension; non-constant
idx-free forms over 1-dimensional regions; multivariate coordinate
forms at $S = \emptyset$ only), so the certified reach of the displayed
rule is exactly its guard-passing fragment, by refusal outside it ---
never by unverified acceptance. Conversely nothing in either statement
is wider than the display: every hypothesis is consumed, and the
display\textquotesingle s own two-clause predicate license is carried
exactly.

\ruleblock{Lean verifier receipts}
\sourcenote{module \texttt{BernsteinLogic.Soundness}; bound build \texttt{a20d90e66c13\ldots} (bound from receipt a9c0e34d3e3c9c8c, 2026-08-03T21:30:36Z). A receipt records that named checking layers accepted a build: (module, compiled-artifact SHA-256, environment SHA-256, checker layers), minted when the checking gate ran; the store receipts attesting the bound build are listed with the layers each run credited. Nothing here asks the reader to trust the store: every layer is an independently runnable program over the quoted sources --- rebuild the module and re-run the checkers to reproduce the verdict. Receipts attesting superseded builds credit no layer and are omitted from this rendering (the count of omissions, if any, is noted below).}
\begin{receipttable}
\receiptrow{a9c0e34d3e3c9c8c}{2026-08-03T21:30:36Z}{build+decorrelated}{\badge{BadgeGreen}{current}}
\end{receipttable}
\sourcenote{21 further receipt(s) in the store attest superseded builds of this module and are omitted from this rendering.}
\sourcenote{module \texttt{BernsteinLogic.Core.Catalogue}; bound build \texttt{cbc92e976a1d\ldots} (bound from receipt 5abdb3589707b5df, 2026-08-03T21:30:33Z). A receipt records that named checking layers accepted a build: (module, compiled-artifact SHA-256, environment SHA-256, checker layers), minted when the checking gate ran; the store receipts attesting the bound build are listed with the layers each run credited. Nothing here asks the reader to trust the store: every layer is an independently runnable program over the quoted sources --- rebuild the module and re-run the checkers to reproduce the verdict. Receipts attesting superseded builds credit no layer and are omitted from this rendering (the count of omissions, if any, is noted below).}
\begin{receipttable}
\receiptrow{24804e66edba0bc9}{2026-08-03T20:25:09Z}{build+decorrelated}{\badge{BadgeGreen}{current}}
\receiptrow{5abdb3589707b5df}{2026-08-03T21:30:33Z}{build+decorrelated}{\badge{BadgeGreen}{current}}
\end{receipttable}
\sourcenote{18 further receipt(s) in the store attest superseded builds of this module and are omitted from this rendering.}
\layerline{\layeritem{build}{lake build (elaboration)}{\badge{BadgeGreen}{current}}\layeritem{replay}{leanchecker (official kernel replay)}{\badge{BadgeAmber}{stale only}}\layeritem{hardened}{safe\_verify (hardened screen)}{\badge{BadgeAmber}{stale only}}\layeritem{decorrelated}{lean4export + nanoda (decorrelated kernel)}{\badge{BadgeGreen}{current}}}
\sourcenote{\badge{BadgeGreen}{comparator: fixture DemoBernCone} the
three-implementation comparator harness (compiled machine code,
kernel-reduction probes, and the compiler-IR interpreter, all executing
the same elaborated checker functions) carries the worked
example\textquotesingle s certificate --- an \rulename{Ax}
$\to$ $\wedge$E $\to$ \rulename{BernCone} chain at $S = \emptyset$ ---
as its fixture \texttt{DemoBernCone}. The suite passes only when all
three paths return identical verdicts on it, and the kernel path probes
two-sided (reporting its own verdict rather than confirming the compiled
one), so a false green would require the same wrong answer from
execution stacks that share nothing below the elaborated definitions.
}
\ruleblock{HOL4 statement}
\sourcenote{theory \texttt{bern\_catalogue}, HOL4 development script \texttt{\texttt{bern\_catalogueScript.sml}} \badge{BadgeGreen}{script bound}}
\holhead{Definition bernConeCoeffCheck\_def}
\begin{holsrc}
Definition bernConeCoeffCheck_def:
  bernConeCoeffCheck sx es dg pl pr =
    if es.simplex = sx then
      (let cx = sx.complex in
       let verts = FILTER (λv. simplexContainsNat sx v) (rangeL cx.nVertices) in
       if LENGTH verts = cx.ambDim + 1 then
         (case termAsCoordPoly pl of
          | NONE => F
          | SOME p =>
            if termIdxFree pl then
              (if polyVars p ≠ [] ∧ cx.ambDim ≠ 1 then F
               else
                 let coords = MAP (λv. MAP (λj. getD (getD cx.coord v []) j (0:rat))
                                           (rangeL cx.ambDim)) verts in
                 let smask = MAP (λv. extContainsNat es v) verts in
                 case fdBernstein coords smask dg p of
                 | SOME table => gridAll (satisfiesPredQ pr) table
                 | NONE => F)
            else
              (if extIsEmptyB es ∧ termIdxBounded cx.ambDim pl then
                 (let coords = MAP (λv. MAP (λj. getD (getD cx.coord v []) j (0:rat))
                                            (rangeL cx.ambDim)) verts in
                  let smask = MAP (λv. extContainsNat es v) verts in
                  case fdBernstein coords smask dg p of
                  | SOME table => gridAll (satisfiesPredQ pr) table
                  | NONE => F)
               else F))
       else F)
    else F
End
\end{holsrc}
\holhead{Definition bernConeIntvCheck\_def}
\begin{holsrc}
Definition bernConeIntvCheck_def:
  bernConeIntvCheck sx es dg pl pr dt =
    if es.simplex = sx then
      (let cx = sx.complex in
       let verts = FILTER (λv. simplexContainsNat sx v) (rangeL cx.nVertices) in
       if LENGTH verts = cx.ambDim + 1 then
         (case termAsCoordPoly pl of
          | NONE => F
          | SOME p =>
            if ¬ termIdxFree pl then F
            else if polyVars p ≠ [] ∧ cx.ambDim ≠ 1 then F
            else
              let coords = MAP (λv. MAP (λj. getD (getD cx.coord v []) j (0:rat))
                                        (rangeL cx.ambDim)) verts in
              let smask = MAP (λv. extContainsNat es v) verts in
              case fdBernstein coords smask dg p of
              | SOME table =>
                  gridAll (λc. case dtypeRoundE dt Rdn c of
                               | Fin lo => (lo ≤ c ∧ satisfiesPredQ pr lo)
                               | _ => F) table
              | NONE => F)
       else F)
    else F
End
\end{holsrc}
\ruleblock{HOL verifier receipts}
\badge{BadgeBlue}{documentary} \texttt{hol-strict} CLEAN over the ported checker chain: theory
\texttt{bern\_catalogue}, 472 theorems, 0 cheat-tagged, 0 nonstandard
axioms (whitelist: \texttt{DISK\_THM}; the four \texttt{bool} axioms) ---
verification run of 2026-07-22 on an isolated build host, scripts
sha-matched verbatim to the sources cited here (git blob
\texttt{e41b9ed0\ldots}). The ported receipt table replays the
\rulename{BernCone} rows of the Lean development\textquotesingle s own
receipt battery: the worked-example accept
(\texttt{bernConeCoeffCheck} at $S = \emptyset$, degree 1) plus four
mutant rejections (tightened goal, wrong dynamics sign, degree lie,
sub-dimensional simplex) and the nonconvex-predicate rejection.
Downstream, the same chain carries the checker to a CakeML binary with
the \texttt{machine\_code\_sound} tag (theory
\texttt{bernCertProofARM8}). \texttt{bernConeIntvCheck} is the ported
companion of a since-retired validated-numerics rule variant; its
round-down-then-check strategy is destined to fold into the BernCone
checker as an internal fast path (evaluation strategy, not proof
content), and the ported definition with its three receipt rows
(exact-dyadic accept, rounded-bound accept, non-representable-bound
reject) is retained as that branch\textquotesingle s reference.
Semantic stratum, also ported: the soundness twins
\texttt{sound\_bernCone} and its $S = \emptyset$ specialization
\texttt{sound\_bernCone\_empty} (theory \texttt{bern\_tr\_berncone}:
\texttt{hol-strict} clean, 14 theorems, 0 offences, verification run of
2026-07-22 on an isolated build host, content-addressed at the same
sources), consumed by the certified-fragment keystone
(\texttt{checkPf\_sound\_frag}, theory \texttt{bern\_tr\_splice}, 20
theorems, 0 offences, same run) under the payload\textquotesingle s
decode-time gates: a reference-checker accept inside the proven wire
fragment yields sequent validity.

\sourcenote{HOL receipts are documentary: quoted results of \texttt{hol-strict} runs
and census rows from the HOL4 development\textquotesingle s build
evidence, not minted receipt artefacts. A HOL receipt store (the Lean
store\textquotesingle s twin) is named future work.
}
\ruleblock{Proof}

\sourcenote{Genre: core rule (tier 1) --- an ordinary mathematical proof at
Putnam-exam-solution detail. Every hypothesis of the rule is used below and
remarked on, with the counterexample its removal admits where one exists;
the formal-anchor table maps each step to the Lean declaration that
carries it.}

\textbf{Setting and claim.}\enspace Fix a hypernet $H$. A valuation
$\sigma$ assigns each wire a rational tensor of its declared shape; a
sequent $\Phi \vdash \Psi$ is semantically valid when every admissible
$\sigma$ satisfying $\Phi$ satisfies $\Psi$, with comparison atoms read
per-entry. A region atom $x \in \mathcal{R}(\Delta, S)$, for
$\Delta = \mathrm{conv}(V_0, \dots, V_n)$ an $n$-simplex and
$S \subsetneq \{0, \dots, n\}$ a proper set of \emph{extended} vertices,
holds at $\sigma$ when the evaluated point admits barycentric weights
$\lambda_v$ with $\sum_v \lambda_v = 1$, $\lambda_v \geq 0$ at core
vertices $v \notin S$, and $\lambda_v \leq 0$ at $v \in S$ --- the closed
simplex when $S = \emptyset$, and the unbounded set extending beyond the
faces opposite $S$ otherwise. (Field bookkeeping, fixed once: the
vertices $V_0, \dots, V_n$ are exact rational data of the payload; the
weights $\lambda_v$ are quantified over $\R$; and the evaluated point is
rational, valuations assigning rational tensors --- Step 2 of the
theorem is where these three facts meet.) The rule's payload supplies the region
premise $\Phi \vdash x \in \mathcal{R}(\Delta, S)$, a polynomial term $p$
of degree $\leq d$, a predicate $P$ from the accepted grammar (a finite
conjunction of scalar hole-atoms $s \cdot \square \bowtie c$, $s > 0$,
$\bowtie \in \{\leq, \geq, =\}$), the checked wire tie \texttt{polyTie},
the predicate-shape clause ($P$ conic when $S \neq \emptyset$; convex
suffices when $S = \emptyset$), and the checked coefficient condition:
every forward-difference Bernstein coefficient of $p$ over
$\mathcal{R}(\Delta, S)$, recomputed exactly, satisfies $P$. For
$S \neq \emptyset$ the checker's scope guard is part of the rule's
acceptance: certificates are confined to constant canonical forms (any
dimension) or single-wire polynomials over a $1$-dimensional region
(\autoref{rem:bc-guard}) --- ``accepted'' below always means
guard-passing. The claim
(\texttt{sound\_bernCone}): the composite sequent
$\Phi \vdash P(p(x))$ is then semantically valid.

\begin{blemma}[Bernstein form on the closed simplex]\label{lem:bc-bern}
Let $p$ be a polynomial of degree $\leq d$ on $\R^n$ with degree-$d$
Bernstein coefficients $(b_\alpha)_{|\alpha| = d}$ on $\Delta$, so that
$p(x) = \sum_{|\alpha| = d} b_\alpha\, B^d_\alpha(\lambda)$ with
$B^d_\alpha(\lambda) = \binom{d}{\alpha} \lambda^\alpha$ in the
barycentric coordinates $\lambda$ of $x$. For $x \in \Delta$ (all
$\lambda_v \geq 0$),
\[
B^d_\alpha(\lambda) \geq 0
\qquad\text{and}\qquad
\sum_{|\alpha| = d} B^d_\alpha(\lambda)
  = \Bigl(\sum_v \lambda_v\Bigr)^{\!d} = 1,
\]
so $p(x)$ is a \emph{convex combination} of the coefficients $b_\alpha$.
\end{blemma}
\begin{proof}
Nonnegativity is termwise ($\lambda^\alpha \geq 0$); the partition of
unity is the multinomial theorem applied to
$(\lambda_0 + \dots + \lambda_n)^d = 1^d$. The representation itself is
the classical change to the Bernstein basis
\cite{bernstein1912,farouki2012}; the checker \emph{computes} the
$b_\alpha$ from $p$'s canonical monomial form by exact forward
differencing over $\Q$, so the representation is an identity of exact
rational data, not an approximation. That Bernstein coefficients
enclose the range of $p$ on $\Delta$ --- the reading of this lemma used
below --- is the sound (easy) direction of \cite{cargoshisha1966}.
\end{proof}

\begin{blemma}[accepted predicates are convex]\label{lem:bc-convex}
Let $P$ be accepted by the predicate classifier. Then its satisfaction
set $\{ v \in \Q : \mathrm{satisfiesPredQ}(P, v) \}$ is convex.
\end{blemma}
\begin{proof}
Each accepted atom orients (possibly by mirroring) to
$s \cdot v \bowtie c$ with $s > 0$ and
$\bowtie \in \{\leq, \geq, =\}$; its satisfaction set is a closed ray or
a point --- an interval either way --- and division by the positive $s$
does not disturb this. A conjunction intersects finitely many intervals.
Strict comparators and $\neq$ are \emph{refused} by the classifier ---
closed-world conservatism that can lose completeness but never
soundness; the refusal is a checker verdict, not a soundness case.
\end{proof}

\begin{blemma}[the univariate ray identity]\label{lem:bc-star}
Let $d \geq 0$ and let $b_0, \dots, b_d$ be the degree-$d$ Bernstein
coefficients of $p$ on the segment $[V_a, V_b]$, indexed by the
multiplicity of $V_b$ (labels chosen so that the extended vertex is
$V_a$ and the core vertex is $V_b$; the mirrored case relabels ---
concretely, with the other vertex extended, apply the identity to the
reversed sequence $c_i := b_{d-i}$, whose differences against the
un-reversed one are the signed
$(\Delta^{j} c)_{d-j} = (-1)^{j} (\Delta^{j} b)_{0}$; the checked table
is computed relative to the extension mask, so its rows \emph{are} the
$(\Delta^{j} c)_{d-j}$ of the relabeled orientation). Let
$x$ lie on the extension ray beyond $V_b$:
in barycentric coordinates $x = \lambda_a V_a + \lambda_b V_b$ with
$\lambda_a = -u \leq 0$, $\lambda_b = 1 + u$, $u \geq 0$. Then, with
$(\Delta b)_i := b_{i+1} - b_i$ the forward difference,
\[
p(x) \;=\; \sum_{j=0}^{d} \binom{d}{j}\, u^{j}\, (\Delta^{j} b)_{d-j},
\tag{$\star$}
\]
a combination of the certified table values $(\Delta^{j} b)_{d-j}$ with
\emph{nonnegative} weights $\binom{d}{j} u^{j}$, the $j = 0$ weight
being $1$.
\end{blemma}
\begin{proof}
Substitute $\lambda_a = -u$, $\lambda_b = 1 + u$ into the Bernstein
form and expand $(1+u)^i$ binomially:
\[
p(x) = \sum_{i=0}^{d} \binom{d}{i} b_i\, (-u)^{d-i} (1+u)^{i}
     = \sum_{i=0}^{d} \sum_{k=0}^{i} \binom{d}{i} \binom{i}{k}
       (-1)^{d-i}\, b_i\, u^{d-i+k}.
\]
Collect the coefficient of $u^{j}$: setting $i = d - j + m$ with
$m = 0, \dots, j$, it is
$\sum_{m=0}^{j} \binom{d}{d-j+m} \binom{d-j+m}{m} (-1)^{j-m}
b_{d-j+m}$. The product of binomials telescopes,
\[
\binom{d}{d-j+m} \binom{d-j+m}{m}
  = \frac{d!}{(j-m)!\, m!\, (d-j)!}
  = \binom{d}{j} \binom{j}{m},
\]
so the coefficient of $u^{j}$ is
$\binom{d}{j} \sum_{m=0}^{j} \binom{j}{m} (-1)^{j-m} b_{d-j+m}
 = \binom{d}{j} (\Delta^{j} b)_{d-j}$, which is $(\star)$. (Degree-$1$
sanity check: $p(x) = b_1 + u\,(b_1 - b_0)$, matching
$(1+u)\,b_1 - u\,b_0$ directly.) The values $(\Delta^{j} b)_{d-j}$ are
exactly the rows the checker's forward-difference table certifies on
this fragment: the $j$-th difference of the coefficient sequence,
anchored at the core vertex --- in the rule display's notation, the
degree-$d$ Bernstein coefficients of the difference polynomials $P_j$
on the core face $F_S$, specialized to the one-dimensional and constant
shapes the scope guard admits.
\end{proof}

\begin{blemma}[conic transfer]\label{lem:bc-conic}
Let $P$ be accepted with the conic classification (every atom's
constant is $0$). Then its satisfaction set contains $0$, and is closed
under nonnegative scaling and under finite nonnegative combinations.
\end{blemma}
\begin{proof}
Each conic atom reads $s \cdot v \bowtie 0$, satisfied at $v = 0$ and
preserved under $v \mapsto t v$ for $t \geq 0$; conjunction preserves
both facts. Closure under addition follows from convexity
(Lemma~\ref{lem:bc-convex}) plus scaling:
$v + w = 2 \cdot \tfrac{1}{2}(v + w)$; induction extends to finite
nonnegative combinations.
\end{proof}

\begin{bthm}[soundness of \rulename{BernCone}]\label{thm:bc-sound}
Under the payload's checked hypotheses, if the region premise is
semantically valid, then so is $\Phi \vdash P(p(x))$.
\end{bthm}
\begin{proof}
Fix an admissible $\sigma \models \Phi$; the region premise yields
barycentric weights for the evaluated point of $x$ as in the Setting.
The conclusion formula is, by the checked \texttt{postApply}, the
predicate $P$ with the polynomial term substituted for its hole, so it
suffices to establish $\mathrm{satisfiesPredQ}(P, \cdot)$ at the
evaluated value(s) of $p$ --- per-entry, when $p$ is tensor-shaped.

\emph{Step 1 (the polynomial reads one point).} The wire tie
\texttt{polyTie} says $p$ is closed, or mentions exactly the region
wire $w$ with $x = w$. In the second case the region premise pins
$\sigma$'s read of $w$ to a point of $\mathcal{R}(\Delta, S)$; in the
first, no wire is read at all and $p$'s value is its closed constant.

\emph{Step 2 (rational weights).} The model's barycentric weights are
\emph{real}; the kernel computes over $\Q$. The weights solve a linear
system with rational data (the vertex coordinates and the evaluated
point) under the region atom's sign pattern --- $\lambda_v \ge 0$ at
core vertices, $\lambda_v \le 0$ at extended ones. Substituting
$\mu_v := -\lambda_v$ at the extended vertices turns the pattern into a
nonnegative-orthant condition on a rational polyhedron, so a real
solution can be traded for a rational one with the same signs, and the
exact kernel consumes the rational weights. This
step is invisible in informal treatments and load-bearing in the
mechanization.

\emph{Step 3 (case $S = \emptyset$: the convex kernel).} By
Lemma~\ref{lem:bc-bern}, $p$'s value is a convex combination of the
certified coefficients; every coefficient satisfies $P$ (the checked
table sweep); by Lemma~\ref{lem:bc-convex} satisfaction transfers to
the combination. On this case the checker admits both translation
routes --- the idx-free single-wire form and the multivariate
projection-indexed form with all indices in range --- and the per-entry
evaluation lemma carries the scalar argument to every entry of a
tensor-shaped value.

\emph{Step 4 (case $S \neq \emptyset$).} The checker's scope guard
confines accepted certificates to two shapes. \emph{(4a: constant
canonical form, any dimension.)} If $p$'s canonical polynomial has no
variables, its value everywhere is the constant --- the zero-order
forward-difference row is that constant, certified directly.
\emph{(4b: single-wire polynomial over a 1-dimensional region.)} A
proper nonempty extension set on a 1-simplex extends exactly one of the
two vertices, and $\mathcal{R}(\Delta, S)$ is the closed ray from the
core vertex; Lemma~\ref{lem:bc-star} writes $p$'s value as a
nonnegative combination of the certified table values with the $j = 0$
weight equal to $1$, and Lemma~\ref{lem:bc-conic} transfers
satisfaction through the combination. Everything else is
\emph{refused} by the guard, so no further case exists for soundness
to cover.

\emph{Step 5 (read off).} In each case
$\mathrm{satisfiesPredQ}(P, \cdot)$ holds at the evaluated value; the
substitution reading turns this into
$\sigma \models P(p(x))$.
\end{proof}

\begin{bremark}[why \texttt{polyTie} is not decoration]
\label{rem:bc-polytie}
Nothing else ties the checked polynomial to the wire the premise
constrains. Without the field, take the satisfiable premise
$w_0 \in \mathcal{R}(\mathrm{conv}(1,2), \emptyset)$ and the polynomial
$p = w_1$ over a \emph{different} wire: the translation remaps the
single-variable canonical form onto the region coordinate, the
certified coefficients are $\{1, 2\}$ --- both satisfying
$\square \geq 1$ --- and the rule would conclude
$\Phi \vdash w_1 \geq 1$, false at any valuation reading
$\sigma(w_1) = 0$. Closed polynomials need no tie: their value reads no
wire.
\end{bremark}

\begin{bremark}[why conic, not merely convex, on rays]
\label{rem:bc-conic-needed}
The predicate-shape clause is two-sided because the two region shapes
transfer differently: convex combinations
(Lemma~\ref{lem:bc-bern}) on the closed simplex, conic combinations
with unbounded weights ($\star$) on a ray. A convex-but-not-conic
predicate fails on the ray, by hand: take $d = 1$,
$\Delta = \mathrm{conv}(0, 1)$ with the vertex at $0$ extended (the
region is the ray $[1, \infty)$), $p(t) = 1 - t$, and
$P = (\square \geq -1)$. The certified table is
$\{\,b_1,\ (\Delta b)_0\,\} = \{\,0,\ -1\,\}$, every value satisfying
$P$; yet at $u = 2$ (the point $t = 3$), $(\star)$ gives
$p = b_1 + 2(b_1 - b_0) = -2 < -1$. The classifier refuses exactly
this: an atom with nonzero constant classifies non-conic, and the
$S \neq \emptyset$ clause then rejects the certificate. Conversely, at
$S = \emptyset$ demanding conic would be needless strength --- the rule
carries the license (any accepted convex predicate) in the premise
itself rather than in surrounding prose.
\end{bremark}

\begin{bremark}[the scope guard is part of the theorem]
\label{rem:bc-guard}
The checker refuses non-constant idx-free polynomials over regions of
dimension $\neq 1$, and multivariate projection-indexed polynomials
whenever $S \neq \emptyset$. Soundness is stated and proven for every
certificate the checker \emph{accepts}; the guard makes the unproven
territory unrepresentable rather than discouraged. The fully general
multivariate ray identity (the arbitrary-shape analogue of
Lemma~\ref{lem:bc-star}) remains open in the development and becomes
load-bearing only if the guard lifts --- the two must move together.
\end{bremark}

\begin{bremark}[relation to the literature]
The rule is the certificate form of Bernstein-coefficient range
enclosure: coefficients bound the range on the simplex
\cite{cargoshisha1966}, enclosures converge under degree elevation and
subdivision \cite{garloff1986} --- which is why \rulename{Subdivide} is
this rule's completeness partner --- and the basis itself goes back to
\cite{bernstein1912}, surveyed in \cite{farouki2012}. With polynomial
atoms and positivity certificates the rule sits beside the
Krivine--Handelman family \cite{krivine1964,handelman1988}: as
everywhere in this calculus, only the trivial (sound) direction is
claimed or needed --- no Positivstellensatz completeness enters the
kernel. The forward-difference computation of the coefficient table is
the classical finite-difference reading of the Bernstein form; the
subdivision algorithms of \cite{decasteljau1959,laneriesenfeld1980}
operate on the same data one rule over.
\end{bremark}

\begin{anchors}
Setting: region atoms as masked barycentric membership & \texttt{MaskRegion}; \texttt{Formula.holds} (\texttt{inRegion} case)\\
Lemma \ref{lem:bc-bern} + Step 3 kernel & the $S = \emptyset$ convex kernel consumed by \texttt{sound\_bernCone\_empty}; per-entry carry \texttt{Term.evalQ\_brRead}\\
Step 2 rational weights & \texttt{maskRegion\_rat\_weights}; \texttt{ConvexRat.exists\_nonneg\_rat\_solution}\\
Lemma \ref{lem:bc-convex} & \texttt{classifyPred}, \texttt{satisfiesPredQ\_convex}\\
Lemma \ref{lem:bc-star} (univariate $(\star)$) & \texttt{QArith.fdBernstein\_sound\_dim1}\\
Step 4a constant corner & \texttt{QArith.fdBernstein\_sound\_const}\\
Lemma \ref{lem:bc-conic} & \texttt{satisfiesPredQ\_conic\_zero}, \texttt{satisfiesPredQ\_conic\_scale}\\
coefficient recomputation (the checked table) & \texttt{bernConeCoeffCheck}, \texttt{allForwardDiffBernCoeffsSatisfyP}, \texttt{QArith.fdBernstein}, \texttt{QArith.gridAll}\\
coordinate translation + index gate & \texttt{Term.asCoordPoly?}, \texttt{Term.asCoordPolyIdx?}, \texttt{Term.idxBounded}\\
Step 1 wire tie & payload field \texttt{polyTie}\\
Step 5 conclusion reading & \texttt{applyPredicateToTerm}, \texttt{holds\_substHole\_of\_satisfies}\\
Theorem \ref{thm:bc-sound} & \texttt{sound\_bernCone} (case split onto \texttt{sound\_bernCone\_empty}); wired as core \texttt{RuleData.bernCone} into \texttt{ProofObject.sound}\\
open residue (Remark \ref{rem:bc-guard}) & \texttt{ConeLemma.StarIdentity} (general shapes; not consulted by the current guard)\\
\end{anchors}


\blcatsection{Admissible and Alias Rules}
\rulesection{16}{Norm}{admissible rule}
\label{blcatrule-Norm}
\ruleblock{The rule (Steele metanotation)}
\sourcenote{as realized in the mechanized development; source row: \texttt{\texttt{rules.toml}} --- \emph{Norm}. The metanotation statement is typeset by hand for notational quality; its mechanized ground truth is the Lean statement below, extracted from the sources at every build of this book and hash-bound. The rule \emph{roster} is machine-checked (every mechanized rule tag round-trips to exactly one catalogue row); the statement \emph{text} is editorial, checked by eye against the extracted Lean beside it.}
\begin{mathpar}
\inferrule*[Right=\rulename{Norm}]{
  t_1, t_2 \ \text{polynomial terms in the same variables with equal monomial canonical forms over } \Q
  \\
  \bshape_H(t_1) = \bshape_H(t_2) \neq \bot \ \text{(a shared defined broadcast shape)}
  \\
  \wires(t_1), \wires(t_2) \ \text{all finitely typed in } H
}{
  \Phi \vdash t_1 = t_2
}
\end{mathpar}

\par\medskip
\rulename{Norm} is the calculus\textquotesingle s computation rule: it
states, as a premise-free sequent, every equality of polynomial terms
that holds by \emph{canonical form} --- $w + w = 2 \cdot w$,
distributivity, collection of like monomials, exact $\Q$ coefficient
arithmetic. A producer uses it constantly and invisibly: whenever an
emitted derivation needs two spellings of the same polynomial to be
interchangeable (before an \rulename{EqElim}, to massage a term into
the pinned shape another rule\textquotesingle s payload demands), a
single \rulename{Norm} node discharges the rewriting that an
equational-logic presentation would smear across many congruence steps.

What the checker recomputes: that both sides are idx-free (projection
terms have no canonical form in the wire-keyed polynomial space --- the
comparison would not carry the rule\textquotesingle s meaning), and
that the two canonical polynomials are \emph{identical} ---
coefficient-for-coefficient equality over $\Q$, never evaluation at
test points, and never rescaling: $2 \cdot w$ and $w + w$ are equal as
canonical forms; $2 \cdot w$ and $w$ are not, proportional though they
are. The remaining payload fields are well-typedness, not computation:
both sides share a defined broadcast shape and mention only finitely
typed wires.

Each hypothesis is load-bearing. Drop the canonical-form identity and
the rule asserts arbitrary equalities. Drop the shared-shape and
finiteness fields and the rule is false \emph{as staged}: at a wire of
extended dtype, even $\vdash t = t$ is invalid --- the atom fails to
hold at every valuation giving the wire an infinite entry, while
holding at the rational ones --- and since the premise is
arbitrary, the staged rule concluded failing sequents from satisfiable
premises; a cancelling sub-term $w_2 - w_2$ over a wire
whose shape does not join with the other side\textquotesingle s leaves
the same canonical form but an atom that never evaluates at any
valuation --- unsatisfiable outright. The typing
fields are exactly what makes ``equal canonical forms\textquotesingle
\textquotesingle{} imply ``equal evaluated tensors\textquotesingle
\textquotesingle{} at every valuation.

\cataloguenote{The rule also names a convention the premise-free rules of this
catalogue share (\emph{the \rulename{Norm} convention}): the premise
$\Phi$ is arbitrary and never consulted --- a fact that holds at every
valuation holds under any premise. The catalogued stratum records an
unpacking sketch through \rulename{BernCone} and \rulename{ArithAtom};
the shipped warrant is direct and stronger (one verified semantic
lemma, all tensor shapes), with the well-typedness fields carrying what
the sketch left implicit.
}
\ruleblock{Lean statement}
\declhead{BernsteinLogic.Soundness}{sound\_norm}{\badge{BadgeGreen}{statement bound}}
\begin{leansrc}
/-- Norm, sound under the PROPOSED-B13-Norm fields: both sides evaluate
(bshape + finWires), at the same shape (`tyEq` + `evalQ_bshape`), with
per-entry values equal because the canonical polynomials are literally
equal. -/
theorem sound_norm {H : Hypernet} {seq : Sequent} {termL termR : Term}
    (hcanon : polynomialCanonicalEq termL termR)
    (htyEq : termL.bshape H = termR.bshape H)
    (htyDef : ∃ s, termL.bshape H = some s)
    (hfinL : termL.finWires H) (hfinR : termR.finWires H)
    (hpost : seq.post = .atom .eq termL termR) :
    seq.semanticallyValid H
\end{leansrc}
\par\medskip
\textbf{The whole computational content, displayed.}\enspace The
check is two gates and one equality: both sides idx-free, and the
canonical polynomials literally equal. Everything else the rule needs
is well-typedness data carried in the payload and consumed by the
soundness lemma, not by the checker.

\declhead{BernsteinLogic.Core.Catalogue}{Term.toPoly}{\badge{BadgeGreen}{statement bound}}
\begin{leansrc}
/-- Exact sparse-polynomial canonical form of a term, variables = wires.
The ONLY canonicalization in leg C (F7 discipline; see original). -/
def Term.toPoly : Term → QArith.Poly
  | .var w => QArith.Poly.var w
  | .const c => QArith.Poly.const c
  | .add l r => QArith.Poly.add l.toPoly r.toPoly
  | .mul l r => QArith.Poly.mul l.toPoly r.toPoly
  | .sub l r => QArith.Poly.sub l.toPoly r.toPoly
  | .neg t => QArith.Poly.neg t.toPoly
  -- G1 INVARIANT (load-bearing): the wire-keyed variable space cannot
  -- express per-entry projections, so idx has NO canonical form here —
  -- the placeholder below is semantically inert because every consumer
  -- of `toPoly` first requires `Term.idxFree` (the canonical-form
  -- relations `polynomialCanonicalEq`/`farkasCheck` carry the guard).
  -- idx terms enter polynomial land ONLY through `asCoordPolyIdx?`.
  | .idx _ _ => QArith.Poly.const 0

/-- Ground-term evaluation: `none` iff a wire occurs — and on any idx
node: a well-typed idx (rank-1 argument) is never closed. -/
def Term.evalQ? : Term → Option Rat
  | .var _ => none
  | .const c => some c
  | .add l r => do pure ((← l.evalQ?) + (← r.evalQ?))
  | .mul l r => do pure ((← l.evalQ?) * (← r.evalQ?))
  | .sub l r => do pure ((← l.evalQ?) - (← r.evalQ?))
  | .neg t => do pure (-(← t.evalQ?))
  | .idx _ _ => none

/-- Comparison evaluation over ℚ. -/
def BowTie.evalQ : BowTie → Rat → Rat → Bool
  | .lt, x, y => decide (x < y)
  | .le, x, y => decide (x ≤ y)
  | .eq, x, y => decide (x = y)
  | .ge, x, y => decide (y ≤ x)
  | .gt, x, y => decide (y < x)
  | .ne, x, y => decide (¬ x = y)

/-- Closed-formula evaluation: `none` iff a wire or region atom occurs
(strict; see original for the closed-world reading). -/
def Formula.evalClosed : Formula → Option Bool
  | .atom b l r => do pure (b.evalQ (← l.evalQ?) (← r.evalQ?))
  | .inRegion _ _ _ => none
  | .and l r => do pure ((← l.evalClosed) && (← r.evalClosed))
  | .or l r => do pure ((← l.evalClosed) || (← r.evalClosed))
  | .tt => some true
  | .ff => some false

/-- The barycentric coordinate functions of a 1-dimensional simplex, as
canonical affine Terms in `t` (pinned shape `offset + slope·t`; see the
original for the full derivation notes and D2 scope). -/
def baryLambdas1D (cx : SimpC) (dContains : Fin cx.nVertices → Bool)
    (t : Term) : Option (List (Fin cx.nVertices × Term))
\end{leansrc}
\declhead{BernsteinLogic.Core.Catalogue}{PNormData}{\badge{BadgeGreen}{statement bound}}
\begin{leansrc}
inductive PNormData : Hypernet → Sequent → Type where
  | mk (H : Hypernet) (seq : Sequent)
       (termL termR : Term)
       (canonEq : polynomialCanonicalEq termL termR)
       -- PROPOSED-B13-Norm (register B13 rider; realized provisionally
       -- 2026-07-16 on the repair hint "asserting wires have finite
       -- type" — the same field family that repaired EqElim/ConvIntro):
       -- both sides are shape-stable at a common broadcast shape and
       -- finitely typed. Without these, an ∞-valued ext wire makes even
       -- ⊢ t = t unsatisfiable (pre is arbitrary), and a cancelling
       -- wire of incompatible shape (w₂ − w₂ at a shape unjoinable with
       -- termL's) makes the conclusion atom never evaluate.
       (tyEq   : termL.bshape H = termR.bshape H)
       (tyDef  : ∃ s, termL.bshape H = some s)
       (finL   : termL.finWires H)
       (finR   : termR.finWires H)
       (postEq  : seq.post = .atom .eq termL termR)
       : PNormData H seq
\end{leansrc}
\declhead{BernsteinLogic.Core.Catalogue}{Term.idxFree}{\badge{BadgeGreen}{definition bound}}
\begin{leansrc}
/-- Structural absence of the idx constructor. The polynomial
canonical-form fragment (D2): `toPoly`'s wire-keyed variable space cannot
express per-entry projections, so every canonical-form consumer
(`polynomialCanonicalEq`, `farkasCheck`, the legacy `asCoordPoly?`
remap) requires this; idx coordinate polynomials take the
`asCoordPolyIdx?` route instead (G1). -/
def Term.idxFree : Term → Bool
  | .var _ | .const _ => true
  | .add l r | .mul l r | .sub l r => l.idxFree && r.idxFree
  | .neg t => t.idxFree
  | .idx _ _ => false
\end{leansrc}
\declhead{BernsteinLogic.Core.Catalogue}{polynomialCanonicalEq}{\badge{BadgeGreen}{definition bound}}
\begin{leansrc}
def polynomialCanonicalEq (l r : Term) : Prop :=
  l.idxFree = true ∧ r.idxFree = true ∧ l.toPoly = r.toPoly
\end{leansrc}
\ruleblock{Correspondence (Steele display $\leftrightarrow$ Lean)}
\badge{BadgeGreen}{reviewer-certified against the displayed texts}
\sourcenote{The dictionary below is a reviewer-certified statement of how the mechanized statements above interpret the displayed rule: display metavariable $\mapsto$ Lean binder, each display hypothesis $\mapsto$ the premise or side condition carrying it, the display conclusion $\mapsto$ the goal --- with any deliberate difference in width between display and mechanization named outright. Certification is recorded as a hash triple over (the displayed cell, the extracted statement texts, this dictionary\textquotesingle s own text); if any of the three changes, this block turns amber until a reviewer re-certifies. The correspondence claim itself remains editorial --- reviewer-grade, not machine-checked; its machine-checked successor types each soundness statement against the rule as a structured object, and supersedes this block per rule as it lands.}
Judgment frame as in \rulename{Ax}.
\begin{itemize}
\item $t_1, t_2 \mapsto$ \texttt{termL termR : Term}; the displayed
  side condition (equal monomial canonical forms) $\mapsto$
  \texttt{hcanon} $:$ \texttt{polynomialCanonicalEq termL termR},
  quoted in full below --- both sides idx-free and the canonical
  polynomials literally equal. The display\textquotesingle s ``same
  variables\textquotesingle\textquotesingle{} and ``polynomial
  terms\textquotesingle\textquotesingle{} have no separate binders:
  canonical forms are wire-keyed (equality forces the same canonical
  support --- a mentioned-but-cancelled wire is no constraint on the
  polynomial; evaluability polices it through the typing fields
  instead) and the term grammar is polynomial by construction.
\item conclusion $\Phi \vdash t_1 = t_2$ $\mapsto$ the goal under
  \texttt{hpost} $:$ \texttt{seq.post = .atom .eq termL termR};
  the arbitrary $\Phi$ is premise-freeness by silence --- the
  precondition component of \texttt{seq : Sequent} is unconstrained
  by any hypothesis (the \rulename{PConstExact} representation).
\end{itemize}
Width note: the display carries the
typing pack inline, matching the mechanized statement\textquotesingle s
four well-typedness hypotheses one for one ---
\texttt{htyEq} (both broadcast shapes equal), \texttt{htyDef} (and
defined --- the \texttt{some}), \texttt{hfinL}/\texttt{hfinR} (all
free wires at finite dtypes). The bound \texttt{PNormData} payload
carries the same conditions under its own field names ---
\texttt{hcanon} $\leftrightarrow$ \texttt{canonEq}, \texttt{htyEq}
$\leftrightarrow$ \texttt{tyEq}, \texttt{htyDef} $\leftrightarrow$
\texttt{tyDef}, \texttt{hfinL}/\texttt{hfinR} $\leftrightarrow$
\texttt{finL}/\texttt{finR}, \texttt{hpost} $\leftrightarrow$
\texttt{postEq} --- and they are not optional: without them the
bare canonical-form rule is refuted (an $\infty$-valued wire makes
even $t = t$ fail at every extended-valued assignment --- the roster
row\textquotesingle s
counterexample). The one-line canonical-form sentence is the
rule\textquotesingle s
\emph{intent}; the certified statement is that sentence plus exactly
this typing pack.

\ruleblock{Lean verifier receipts}
\sourcenote{module \texttt{BernsteinLogic.Soundness}; bound build \texttt{a20d90e66c13\ldots} (bound from receipt a9c0e34d3e3c9c8c, 2026-08-03T21:30:36Z). A receipt records that named checking layers accepted a build: (module, compiled-artifact SHA-256, environment SHA-256, checker layers), minted when the checking gate ran; the store receipts attesting the bound build are listed with the layers each run credited. Nothing here asks the reader to trust the store: every layer is an independently runnable program over the quoted sources --- rebuild the module and re-run the checkers to reproduce the verdict. Receipts attesting superseded builds credit no layer and are omitted from this rendering (the count of omissions, if any, is noted below).}
\begin{receipttable}
\receiptrow{a9c0e34d3e3c9c8c}{2026-08-03T21:30:36Z}{build+decorrelated}{\badge{BadgeGreen}{current}}
\end{receipttable}
\sourcenote{21 further receipt(s) in the store attest superseded builds of this module and are omitted from this rendering.}
\sourcenote{module \texttt{BernsteinLogic.Core.Catalogue}; bound build \texttt{cbc92e976a1d\ldots} (bound from receipt 5abdb3589707b5df, 2026-08-03T21:30:33Z). A receipt records that named checking layers accepted a build: (module, compiled-artifact SHA-256, environment SHA-256, checker layers), minted when the checking gate ran; the store receipts attesting the bound build are listed with the layers each run credited. Nothing here asks the reader to trust the store: every layer is an independently runnable program over the quoted sources --- rebuild the module and re-run the checkers to reproduce the verdict. Receipts attesting superseded builds credit no layer and are omitted from this rendering (the count of omissions, if any, is noted below).}
\begin{receipttable}
\receiptrow{24804e66edba0bc9}{2026-08-03T20:25:09Z}{build+decorrelated}{\badge{BadgeGreen}{current}}
\receiptrow{5abdb3589707b5df}{2026-08-03T21:30:33Z}{build+decorrelated}{\badge{BadgeGreen}{current}}
\end{receipttable}
\sourcenote{18 further receipt(s) in the store attest superseded builds of this module and are omitted from this rendering.}
\layerline{\layeritem{build}{lake build (elaboration)}{\badge{BadgeGreen}{current}}\layeritem{replay}{leanchecker (official kernel replay)}{\badge{BadgeAmber}{stale only}}\layeritem{hardened}{safe\_verify (hardened screen)}{\badge{BadgeAmber}{stale only}}\layeritem{decorrelated}{lean4export + nanoda (decorrelated kernel)}{\badge{BadgeGreen}{current}}}
\sourcenote{\badge{BadgeGreen}{comparator: fixture NormOutsideFragment} the
three-implementation comparator harness (compiled machine code,
kernel-reduction probes, and the compiler-IR interpreter, all executing
the same elaborated checker functions) carries a well-typed single
\rulename{Norm} node as its fixture \texttt{NormOutsideFragment},
expected verdicts \texttt{check = true}, \texttt{v0frag = false}. The
fixture is the suite\textquotesingle s false-verdict control:
\rulename{Norm} sits outside the restricted fragment recognizer by
design, so the fixture pins the reject lane end to end --- a comparator
that could only confirm acceptance would be blind to a path that
wrongly accepts everything.
}
\ruleblock{HOL4 statement}
\sourcenote{theory \texttt{bern\_catalogue}, HOL4 development script \texttt{\texttt{bern\_catalogueScript.sml}} \badge{BadgeGreen}{script bound}}
\holhead{Definition polynomialCanonicalEqB\_def}
\begin{holsrc}
Definition polynomialCanonicalEqB_def:
  polynomialCanonicalEqB l r ⇔
    termIdxFree l ∧ termIdxFree r ∧ termToPoly l = termToPoly r
End
\end{holsrc}
\ruleblock{HOL verifier receipts}
\badge{BadgeBlue}{documentary} \texttt{hol-strict} CLEAN over the ported checker chain: theory
\texttt{bern\_catalogue}, 472 theorems, 0 cheat-tagged, 0 nonstandard
axioms (whitelist: \texttt{DISK\_THM}; the four \texttt{bool} axioms) ---
verification run of 2026-07-22 on an isolated build host, scripts
sha-matched verbatim to the sources cited here (git blob
\texttt{e41b9ed0\ldots}). The ported receipt table pins the
rule\textquotesingle s exactness discipline in both directions: the
canonical-form accept ($w_0 + w_0 = 2 \cdot w_0$) and the
scale-preserving reject ($2 \cdot w_0 \neq w_0$ as canonical forms).
Scope note, honest by the port\textquotesingle s own seam inventory:
the ported surface is the first-order check function; the
payload\textquotesingle s well-typedness fields (broadcast shapes,
finite typing) live in the dependently-typed proof-object layer, which
is a separately tracked porting stage.
Semantic stratum, also ported: the soundness twin \texttt{sound\_norm}
(theory \texttt{bern\_tr\_fragclose}: \texttt{hol-strict} clean, 26
theorems, 0 offences, verification run of 2026-07-22 on an isolated
build host, content-addressed at the same sources), consumed by the
certified-fragment keystone (\texttt{checkPf\_sound\_frag}, theory
\texttt{bern\_tr\_splice}, 20 theorems, 0 offences, same run): a
reference-checker accept inside the proven wire fragment yields sequent
validity.

\sourcenote{HOL receipts are documentary: quoted results of \texttt{hol-strict} runs
and census rows from the HOL4 development\textquotesingle s build
evidence, not minted receipt artefacts. A HOL receipt store (the Lean
store\textquotesingle s twin) is named future work.
}
\ruleblock{Proof}

\sourcenote{Genre: admissible rule (tier 2) --- a proof-tree-rewriting
recipe. The tier taxonomy defines tier 2 as \emph{verified unpacking
into core rules} (one node $\to$ many), with rules admissible by
\emph{reduction to verified machinery} as the recorded boundary case;
\rulename{Norm} is of the latter kind, and this section --- the book's
first admissible rule --- states the general template that the later
recipes instantiate.}

\begin{brecipe}[the admissible template, general form]\label{rec:norm-template}
An admissible rule $R$ with premises $P_1, \dots, P_n$ and conclusion $C$
is warranted by a \emph{rewrite}: a function taking any proof-tree node
\[
\inferrule*[Right=$R$]{\;\mathcal{T}_1 \rhd P_1 \\ \cdots \\ \mathcal{T}_n \rhd P_n\;}{C}
\qquad\rightsquigarrow\qquad
\mathcal{D}\bigl[\mathcal{T}_1, \dots, \mathcal{T}_n\bigr] \rhd C
\]
where $\mathcal{D}$ is a derivation fragment over core (and
previously-justified admissible) rules with the same conclusion,
consuming the premise subtrees $\mathcal{T}_i$ at its leaves. Where such
a fragment exists, applying
it leaf-first eliminates every $R$-node from a finite tree in
finitely many steps, so any judgment derivable with $R$ is derivable
without it --- admissibility in the derivability sense. In the tier's
recorded boundary case --- rules admissible \emph{by reduction to
verified machinery}, this one among them --- no fragment is applied:
where no catalogue derivation of $C$ can exist, the node is instead
discharged in place against a named verified theorem, no derivability
claim is made, and the conserved property --- weaker, and sufficient ---
is the calculus's soundness theorem itself. Either way, no new kernel
semantic content enters. The verified object
shipped per admissible rule is a soundness
lemma for $R$ that factors through already-verified machinery: for a
fragment-backed rule, provable by replaying $\mathcal{D}$; for a
boundary-case rule, the discharge theorem itself.
\end{brecipe}

\begin{brecipe}[the \rulename{Norm} instance --- the degenerate rewrite]
\label{rec:norm-instance}
\rulename{Norm} has no proof premises ($n = 0$): the node
\[
\inferrule*[Right=\rulename{Norm}]{
  \mathrm{polynomialCanonicalEq}(t_1, t_2) \\
  \text{\upshape both sides shape-stable and finitely typed}
}{\Phi \vdash t_1 = t_2}
\qquad\rightsquigarrow\qquad
\text{\upshape the checker replays the canonical forms;}\;
\text{\texttt{sound\_norm} stands as warrant.}
\]
The side condition is \emph{computational}: both sides idx-free, and
their canonical monomial polynomials \emph{identical} over $\Q$ ---
coefficient-for-coefficient, decidable, replayed by the checker. The
catalogued stratum records an unpacking sketch through
\rulename{BernCone} and \rulename{ArithAtom} (an equality as a pair of
zero-polynomial bounds); the shipped warrant is direct and stronger ---
one verified semantic lemma at all tensor shapes --- with the payload's
well-typedness fields carrying exactly what the sketch left implicit.
\end{brecipe}

\begin{bprop}[semantic warrant]\label{prop:norm-sound}
Suppose $t_1, t_2$ are idx-free with $t_1.\mathrm{toPoly} =
t_2.\mathrm{toPoly}$, the two sides carry the same defined broadcast
shape in $H$, and every wire of either side is finitely typed. Then
$\Phi \vdash t_1 = t_2$ is semantically valid over $H$ for
\emph{arbitrary} $\Phi$.
\end{bprop}
\begin{proof}
Fix any admissible valuation $\sigma$; the premise is not consulted ---
this is \emph{the \rulename{Norm} convention}, shared by every
premise-free rule of the catalogue: a fact holding at every valuation
holds under any premise. The shape and finiteness fields give both
sides defined evaluations, as tensors of the common broadcast shape
$s$: a term's broadcast shape is computed structurally, so its
definedness already certifies that every subterm's join exists;
finitely typed wires always read; and idx-freedom leaves no projection
to fail. If $s$ has positive rank and zero total size, both sides
evaluate to the unique empty tensor of shape $s$ and the equality atom
holds outright --- there is no entry to compare. Otherwise fix an entry
index of $s$ (for $s = [\,]$, the single scalar
position) and read each wire at it: a rank-0 wire contributes its
scalar, and a positive-rank wire --- necessarily of shape $s$ itself,
the only positive rank that joins with $s$ under this broadcast
discipline --- contributes its entry at the index. Evaluation at that
read-out is a
ring homomorphism into $\Q$ from the polynomial ring over the union of
the wires the two sides mention (a wire absent from one side, or
cancelled from its canonical form, simply carries coefficient zero
there); a syntactic identity of canonical
forms is an identity of polynomials, and ring homomorphisms preserve
polynomial identities --- so the two evaluated tensors agree entry by
entry, and their shapes agree by hypothesis. Hence the equality atom
holds at $\sigma$. That is the whole proof; \texttt{sound\_norm}
packages it once, per-entry, at every tensor shape.
\end{proof}

\begin{bremark}[what broke without the typing fields]
\label{rem:norm-b13}
As first staged --- canonical-form equality alone --- the rule was
\emph{false}, twice over, and the counterexamples are why the fields
exist. (i) At a wire of extended dtype, the staged rule concluded
\emph{invalid} sequents: at any valuation giving the wire an infinite
entry, even $t = t$ fails to hold --- evaluation is over exact
rationals, and an infinite entry refuses --- so $\vdash t = t$ fails
exactly at the extended-valued assignments while holding at the
rational ones. (ii) A cancelling
sub-term leaves the canonical form unchanged but not the evaluation:
$t_2 = t_1 + (w_2 - w_2)$ has $t_2.\mathrm{toPoly} =
t_1.\mathrm{toPoly}$, yet if $w_2$'s shape does not join with $t_1$'s,
the right-hand atom side never evaluates at \emph{any} valuation ---
there the concluded atom is unsatisfiable outright. In either case a
satisfiable (indeed arbitrary) premise concluded a failing sequent; the
shared-shape and
finiteness fields exclude both. A hypothesis whose necessity is
semantic (evaluability), not algebraic, is exactly the kind a recipe
must state, or the next porter re-learns it the hard way.
\end{bremark}

\begin{bremark}[exactness discipline, and the idx-freedom gate]
\label{rem:norm-f7}
Canonical-form equality is never up-to-rescaling: $2 \cdot w$ and
$w + w$ are equal as canonical forms; $2 \cdot w$ and $w$ are not,
proportional though they are --- the matching discipline every
combination rule of this catalogue (\rulename{Farkas},
\rulename{Syzygy}) inherits when it compares atom polynomials. The
idx-freedom gate is the meaning-preservation condition of the
translation: the wire-keyed polynomial space cannot express per-entry
projections, so a projection-bearing term has no canonical form here
and the comparison is refused rather than performed on the wrong
object. (The combination rules' widened replays translate projections
into their own variable spaces; those translations belong to those
rules, not to \rulename{Norm}.)
\end{bremark}

\begin{anchors}
Recipe \ref{rec:norm-instance} warrant & \texttt{sound\_norm} (per-entry evaluation at the common broadcast shape)\\
the computational check & \texttt{polynomialCanonicalEq} $=$ idx-freedom gates $+$ \texttt{Term.toPoly} equality\\
idx-freedom gate (Remark \ref{rem:norm-f7}) & \texttt{Term.idxFree}; \texttt{Term.toPoly}'s wire-keyed variable space\\
typing fields (Remark \ref{rem:norm-b13}) & payload fields \texttt{tyEq}, \texttt{tyDef}, \texttt{finL}, \texttt{finR} (\texttt{PNormData})\\
premise-freedom ($\Phi$ arbitrary) & the \rulename{Norm} convention; \texttt{pre} unused in \texttt{sound\_norm}\\
wire placement & \texttt{RuleTag.norm}\\
\end{anchors}


\rulesection{17}{Farkas}{admissible rule}
\label{blcatrule-Farkas}
\ruleblock{The rule (Steele metanotation)}
\sourcenote{as realized in the mechanized development; source row: \texttt{\texttt{rules.toml}} --- \emph{Farkas}. The metanotation statement is typeset by hand for notational quality; its mechanized ground truth is the Lean statement below, extracted from the sources at every build of this book and hash-bound. The rule \emph{roster} is machine-checked (every mechanized rule tag round-trips to exactly one catalogue row); the statement \emph{text} is editorial, checked by eye against the extracted Lean beside it.}
\begin{mathpar}
\inferrule*[Right=\rulename{Farkas}]{
  \{\ell_j \bowtie_j 0\}_j \text{ a sub-conjunction of } \Phi \text{ with each } \ell_j \text{ linear, } \bowtie_j \in \{\leq, =\} \\
  \{\lambda_j \in \Q_{\geq 0}\}_j \text{ nonneg coefficients, with } \lambda_j \text{ free when } \bowtie_j \text{ is } {=} \\
  \bshape_H(\textstyle\sum_j \lambda_j \ell_j) \neq \bot \\
  \wires(\textstyle\sum_j \lambda_j \ell_j) \text{ all finitely typed in } H \\
  \wires(\Phi) \cup \wires(\textstyle\sum_j \lambda_j \ell_j) \text{ all typed at positive-size shapes in } H
}{
  \Phi \vdash \textstyle\sum_j \lambda_j \ell_j \leq 0
}
\end{mathpar}

\par\medskip
\rulename{Farkas} is the affine-combination rule over the premise: from
atoms $\{\ell_j \bowtie_j 0\}$ occurring as sub-conjuncts of $\Phi$ and
nonnegative multipliers $\lambda_j$ (free-signed on equality atoms), it
concludes $\Phi \vdash \sum_j \lambda_j \ell_j \leq 0$. It is the
one-rule form of what \rulename{Syzygy} does with sub-\emph{proofs}:
here every combined atom must already sit in the premise conjunction, so
no component derivations enter. A producer finds the multipliers by
whatever method it likes (an LP feasibility call, typically) and ships
only the combination term; nothing about the search survives into the
certificate.

What the checker recomputes, given the payload (combination term,
conclusion term): decode the combination into weighted entries --- a
fixed canonical encoding (a right-nested sum of
$\lambda_j \cdot e_j$ products with a zero terminator), anything else
rejected; require each weight nonnegative; require each
entry\textquotesingle s polynomial to occur among the
$({\cdot} \leq 0)$-normal polynomials of a premise conjunct --- equality
atoms expose \emph{both} polarities to the matcher, which is exactly how
a free-sign multiplier on an equality is realized as a nonnegative
weight on the appropriate polarity; and require the weighted fold of the
entry polynomials to equal the conclusion\textquotesingle s polynomial
\emph{identically}, as canonical forms over $\Q$. The check is a
disjunction of three branches of one replay algorithm, differing only in
the variable space where polynomials are compared: the legacy idx-free
canonical-form branch, a wire-aware coordinate branch for
projection-indexed terms, and a mixed tagged-union branch for
combinations mentioning both bare wires and projections.

Each hypothesis is load-bearing. Drop the sign condition and a producer
multiplies a bound by $-1$ and concludes its negation. Drop the
atom-sourcing check and the combination smuggles in atoms nobody
assumed. Drop the exact fold identity and the conclusion is unrelated to
the premise. Drop the positive-size wire condition and the rule is
unsound at degenerate shapes: comparison atoms over a zero-size
(empty-tensor) wire hold \emph{vacuously}, and a combination of two
vacuous contradictory atoms --- $5 \leq w$ and $w \leq 3$ at a wire of
shape $[0]$ --- derives the false scalar conclusion
$\vdash 2 \leq 0$ from a satisfiable premise. Drop the
conclusion\textquotesingle s shape/finiteness fields and a conclusion
mentioning an untyped wire has a correct canonical polynomial but never
evaluates, so the concluded atom never holds. The proof section walks
these in order.

\cataloguenote{The catalogued statement ranges over \emph{linear} $\ell_j$; the replay
itself is degree-blind, and soundness never consults linearity --- the
restriction is the rule\textquotesingle s published register (the shape
producers emit and the wire format prices), not a checked side
condition. \rulename{Syzygy} lifts the same fold to component
conclusions with degrees unrestricted; \rulename{PBound} re-spells that
lift as the compact affine bound-node.
}
\ruleblock{Lean statement}
\declhead{BernsteinLogic.Soundness}{farkas\_fold\_eval\_nonpos}{\badge{BadgeGreen}{statement bound}}
\begin{leansrc}
/-- Evaluation of the Farkas combination fold: a nonnegative combination
of nonpositive values is nonpositive (valuation-general). -/
theorem farkas_fold_eval_nonpos (ρ : Nat → ℚ) :
    ∀ (entries : List (ℚ × Term)),
    (∀ e ∈ entries, 0 ≤ e.1 ∧ QArith.Poly.eval ρ e.2.toPoly ≤ 0) →
    QArith.Poly.eval ρ
      (entries.foldr
        (fun (le : ℚ × Term) acc =>
          QArith.Poly.add (QArith.Poly.scale le.1 le.2.toPoly) acc) []) ≤ 0
\end{leansrc}
\declhead{BernsteinLogic.Soundness}{farkas\_fold\_eval\_nonpos\_polys}{\badge{BadgeGreen}{statement bound}}
\begin{leansrc}
/-- The nonneg fold in the paired space (valuation-general; entries
carry their polynomials directly — the `farkasEntriesW`-translated
form). -/
theorem farkas_fold_eval_nonpos_polys (ρ : Nat → ℚ) :
    ∀ (pentries : List (ℚ × QArith.Poly)),
    (∀ e ∈ pentries, 0 ≤ e.1 ∧ QArith.Poly.eval ρ e.2 ≤ 0) →
    QArith.Poly.eval ρ
      (pentries.foldr
        (fun (le : ℚ × QArith.Poly) acc =>
          QArith.Poly.add (QArith.Poly.scale le.1 le.2) acc) []) ≤ 0
\end{leansrc}
\declhead{BernsteinLogic.Soundness}{sound\_farkas}{\badge{BadgeGreen}{statement bound}}
\begin{leansrc}
/-- Farkas, general tensor shapes (supersedes the v0 singleton form):
the conclusion's B13 shape/finiteness fields give it a defined
evaluation; the PROPOSED-B14 positivity field rules out vacuously-held
zero-size atoms; the per-entry claim runs through the CLAMPED valuation,
which serves heterogeneous premise-atom shapes at every conclusion
index. -/
theorem sound_farkas {H : Hypernet} {seq : Sequent}
    {combinationCoeffs conclusion : Term}
    (hshape : ∃ s, conclusion.bshape H = some s)
    (hfin : conclusion.finWires H)
    (hwpos : ∀ w ∈ seq.pre.freeWires ++ conclusion.freeWires,
      ∃ τ, H.wireType w = some τ ∧ 0 < Shape.totalSize τ.shape)
    (hvalid : farkasCombinationValid combinationCoeffs seq.pre conclusion)
    (hpost : seq.post = .atom .le conclusion (.const 0)) :
    seq.semanticallyValid H
\end{leansrc}
\par\medskip
\textbf{What the combination check means --- and what it does
not.}\enspace The side condition appears in the payload as the single
name \texttt{farkasCombinationValid}, so this book binds its meaning by
displaying the decoding chain from the checker sources.
\texttt{farkasEntries} decodes the combination term against a
\emph{pinned} canonical encoding --- anything else is rejected, so the
wire format cannot express a malformed combination.
\texttt{atomLeZeroPolys} is the $({\cdot} \leq 0)$-normalization of one
premise atom, exact and never rescaled; note that equality atoms
contribute both polarities, which is the entire realization of the
classical free-sign multiplier discipline for equality constraints.
\texttt{farkasCheckCore} strings the gates together in the idx-free
canonical-form space; \texttt{farkasCheck} is the three-way branch
disjunction --- the two widened branches replay the same gates
verbatim in the wire-aware coordinate space (projection terms
$w[j]$ keyed by an injective pairing of wire and coordinate, so
distinct wires can never collide) and in the mixed tagged-union space
(bare wires and projections as disjoint variable classes; the variant
identifying a bare wire with its zeroth projection admits a false
sequent and is deliberately unrepresentable).

What a skeptical reader should confirm by eye: the conclusion term
enters only as the target of the exact fold identity; the
step\textquotesingle s posted formula appears nowhere in the check; and
nothing in the chain evaluates or consults the truth of anything ---
the check manipulates polynomials, never valuations. Semantic truth is
established only in \texttt{sound\_farkas}, which case-splits on the
accepted branch.

\declhead{BernsteinLogic.Core.Catalogue}{atomLeZeroPolysW}{\badge{BadgeGreen}{statement bound}}
\begin{leansrc}
/-- `atomLeZeroPolys` in the paired coordinate space: the
`(e ≤ 0)`-normal-form polynomials of one atom under the wire-aware idx
translation. An atom with an untranslatable side entails nothing (like
non-atoms — it can then never source a combination entry). -/
def atomLeZeroPolysW : Formula → List QArith.Poly
  | .atom op l r =>
    match l.asCoordPolyIdxW? with
    | none => []
    | some pl =>
      match r.asCoordPolyIdxW? with
      | none => []
      | some pr =>
        match op with
        | .le | .lt => [QArith.Poly.sub pl pr]
        | .ge | .gt => [QArith.Poly.sub pr pl]
        | .eq => [QArith.Poly.sub pl pr, QArith.Poly.sub pr pl]
        | .ne => []
  | _ => []

/-- Translate every decoded entry's vehicle term; `none` if any vehicle
is outside the idx coordinate fragment. (Explicit recursion rather than
`mapM` so soundness inverts it by plain `cases`.) -/
def farkasEntriesW : List (Rat × Term) → Option (List (Rat × QArith.Poly))
  | [] => some []
  | (l, e) :: rest =>
    match e.asCoordPolyIdxW?, farkasEntriesW rest with
    | some pe, some prest => some ((l, pe) :: prest)
    | _, _ => none

/-- The idx-form branch of the Farkas check: every polynomial — entry
vehicles, matched premise atoms, and the conclusion — rides the
wire-aware coordinate translation; entry encoding, nonneg weights,
atom matching, and the exact combination fold mirror the idx-free
branch verbatim. -/
def farkasCheckW (cc : Term) (premise : Formula) (conclusion : Term) : Bool
\end{leansrc}
\declhead{BernsteinLogic.Core.Catalogue}{atomLeZeroPolysM}{\badge{BadgeGreen}{statement bound}}
\begin{leansrc}
/-- `atomLeZeroPolys` in the mixed tagged space: the `(e ≤ 0)`-normal-
form polynomials of one atom under the mixed translation. An atom with
an untranslatable side entails nothing (the W-branch discipline). -/
def atomLeZeroPolysM : Formula → List QArith.Poly
  | .atom op l r =>
    match l.asCoordPolyMixed? with
    | none => []
    | some pl =>
      match r.asCoordPolyMixed? with
      | none => []
      | some pr =>
        match op with
        | .le | .lt => [QArith.Poly.sub pl pr]
        | .ge | .gt => [QArith.Poly.sub pr pl]
        | .eq => [QArith.Poly.sub pl pr, QArith.Poly.sub pr pl]
        | .ne => []
  | _ => []

/-- Translate every decoded entry's vehicle term into the mixed space;
`none` if any vehicle is outside the mixed fragment. -/
def farkasEntriesM : List (Rat × Term) → Option (List (Rat × QArith.Poly))
  | [] => some []
  | (l, e) :: rest =>
    match e.asCoordPolyMixed?, farkasEntriesM rest with
    | some pe, some prest => some ((l, pe) :: prest)
    | _, _ => none

/-- The mixed-variable branch of the Farkas check: every polynomial —
entry vehicles, matched premise atoms, and the conclusion — rides the
mixed tagged translation; entry encoding, nonneg weights, atom matching,
and the exact combination fold mirror the sibling branches verbatim. -/
def farkasCheckM (cc : Term) (premise : Formula) (conclusion : Term) : Bool
\end{leansrc}
\declhead{BernsteinLogic.Core.Catalogue}{farkasCheckW}{\badge{BadgeGreen}{statement bound}}
\begin{leansrc}
/-- The idx-form branch of the Farkas check: every polynomial — entry
vehicles, matched premise atoms, and the conclusion — rides the
wire-aware coordinate translation; entry encoding, nonneg weights,
atom matching, and the exact combination fold mirror the idx-free
branch verbatim. -/
def farkasCheckW (cc : Term) (premise : Formula) (conclusion : Term) : Bool
\end{leansrc}
\declhead{BernsteinLogic.Core.Catalogue}{farkasCheckM}{\badge{BadgeGreen}{statement bound}}
\begin{leansrc}
/-- The mixed-variable branch of the Farkas check: every polynomial —
entry vehicles, matched premise atoms, and the conclusion — rides the
mixed tagged translation; entry encoding, nonneg weights, atom matching,
and the exact combination fold mirror the sibling branches verbatim. -/
def farkasCheckM (cc : Term) (premise : Formula) (conclusion : Term) : Bool
\end{leansrc}
\declhead{BernsteinLogic.Core.Catalogue}{farkasCombinationValid}{\badge{BadgeGreen}{statement bound}}
\begin{leansrc}
def farkasCombinationValid (cc : Term) (premise : Formula)
    (conclusion : Term) : Prop
\end{leansrc}
\declhead{BernsteinLogic.Core.Catalogue}{wireCoordKey}{\badge{BadgeGreen}{statement bound}}
\begin{leansrc}
/-- Injective key for the (wire, coordinate) variable space of the
multi-wire idx coordinate translation below. This is `Nat.pair`'s
definition inlined (the core is mathlib-free; Soundness.lean proves the
`Nat.unpair` round-trip through the definitional equality). Distinct
wires MUST get distinct polynomial variables here: a Syzygy composite
mixes projections of several region wires in one combination, and
conflating `w0[j]` with `w2[j]` would let semantically different
polynomials collide (contrast `asCoordPolyIdx?`, whose BernCone scope
has a single region wire pinned by `polyTie`, so dropping the wire is
sound THERE and only there). -/
def wireCoordKey (w i : Nat) : Nat
\end{leansrc}
\declhead{BernsteinLogic.Core.Catalogue}{Term.asCoordPolyIdxW?}{\badge{BadgeGreen}{statement bound}}
\begin{leansrc}
/-- The wire-aware idx coordinate translation: a Term whose every wire
mention is a projection `.idx (.var w) j` becomes a polynomial over the
paired variables `wireCoordKey w j`. `none` on bare wire mentions and
on idx of any non-var argument (v0.5 scope, matching
`asCoordPolyIdx?`). Constants translate, so idx-free closed terms ride
either route. -/
def Term.asCoordPolyIdxW? : Term → Option QArith.Poly
  | .var _ => none
  | .const c => some (QArith.Poly.const c)
  | .add l r =>
    do pure (QArith.Poly.add (← l.asCoordPolyIdxW?) (← r.asCoordPolyIdxW?))
  | .mul l r =>
    do pure (QArith.Poly.mul (← l.asCoordPolyIdxW?) (← r.asCoordPolyIdxW?))
  | .sub l r =>
    do pure (QArith.Poly.sub (← l.asCoordPolyIdxW?) (← r.asCoordPolyIdxW?))
  | .neg t => (Term.asCoordPolyIdxW? t).map QArith.Poly.neg
  | .idx (.var w) i => some (QArith.Poly.var (wireCoordKey w i))
  | .idx _ _ => none

/-- `atomLeZeroPolys` in the paired coordinate space: the
`(e ≤ 0)`-normal-form polynomials of one atom under the wire-aware idx
translation. An atom with an untranslatable side entails nothing (like
non-atoms — it can then never source a combination entry). -/
def atomLeZeroPolysW : Formula → List QArith.Poly
  | .atom op l r =>
    match l.asCoordPolyIdxW? with
    | none => []
    | some pl =>
      match r.asCoordPolyIdxW? with
      | none => []
      | some pr =>
        match op with
        | .le | .lt => [QArith.Poly.sub pl pr]
        | .ge | .gt => [QArith.Poly.sub pr pl]
        | .eq => [QArith.Poly.sub pl pr, QArith.Poly.sub pr pl]
        | .ne => []
  | _ => []

/-- Translate every decoded entry's vehicle term; `none` if any vehicle
is outside the idx coordinate fragment. (Explicit recursion rather than
`mapM` so soundness inverts it by plain `cases`.) -/
def farkasEntriesW : List (Rat × Term) → Option (List (Rat × QArith.Poly))
  | [] => some []
  | (l, e) :: rest =>
    match e.asCoordPolyIdxW?, farkasEntriesW rest with
    | some pe, some prest => some ((l, pe) :: prest)
    | _, _ => none

/-- The idx-form branch of the Farkas check: every polynomial — entry
vehicles, matched premise atoms, and the conclusion — rides the
wire-aware coordinate translation; entry encoding, nonneg weights,
atom matching, and the exact combination fold mirror the idx-free
branch verbatim. -/
def farkasCheckW (cc : Term) (premise : Formula) (conclusion : Term) : Bool
\end{leansrc}
\declhead{BernsteinLogic.Core.Catalogue}{Term.asCoordPolyMixed?}{\badge{BadgeGreen}{statement bound}}
\begin{leansrc}
/-- The mixed-variable translation: bare wires and idx projections are
DISJOINT opaque variables of one combination algebra — bare `.var w`
keyed even (`2·w`), projections `.idx (.var w) j` keyed odd
(`2·wireCoordKey w j + 1`) — so the two classes can never collide. In
particular bare `w` is NEVER identified with `w[0]`: the identification
variant admits a false sequent (a countermodel witnesses the failure,
alongside the REJECT specimen in EvalDemo). This matches the python reference
(`bl_checker` keys bare and projected mentions as distinct symbols) and
the geolog/fastcheck mixed evaluators. `none` only on idx-of-non-var
(outside every reading's fragment). -/
def Term.asCoordPolyMixed? : Term → Option QArith.Poly
  | .var w => some (QArith.Poly.var (2 * w))
  | .const c => some (QArith.Poly.const c)
  | .add l r =>
    do pure (QArith.Poly.add (← l.asCoordPolyMixed?) (← r.asCoordPolyMixed?))
  | .mul l r =>
    do pure (QArith.Poly.mul (← l.asCoordPolyMixed?) (← r.asCoordPolyMixed?))
  | .sub l r =>
    do pure (QArith.Poly.sub (← l.asCoordPolyMixed?) (← r.asCoordPolyMixed?))
  | .neg t => (Term.asCoordPolyMixed? t).map QArith.Poly.neg
  | .idx (.var w) j => some (QArith.Poly.var (2 * wireCoordKey w j + 1))
  | .idx _ _ => none

/-- `atomLeZeroPolys` in the mixed tagged space: the `(e ≤ 0)`-normal-
form polynomials of one atom under the mixed translation. An atom with
an untranslatable side entails nothing (the W-branch discipline). -/
def atomLeZeroPolysM : Formula → List QArith.Poly
  | .atom op l r =>
    match l.asCoordPolyMixed? with
    | none => []
    | some pl =>
      match r.asCoordPolyMixed? with
      | none => []
      | some pr =>
        match op with
        | .le | .lt => [QArith.Poly.sub pl pr]
        | .ge | .gt => [QArith.Poly.sub pr pl]
        | .eq => [QArith.Poly.sub pl pr, QArith.Poly.sub pr pl]
        | .ne => []
  | _ => []

/-- Translate every decoded entry's vehicle term into the mixed space;
`none` if any vehicle is outside the mixed fragment. -/
def farkasEntriesM : List (Rat × Term) → Option (List (Rat × QArith.Poly))
  | [] => some []
  | (l, e) :: rest =>
    match e.asCoordPolyMixed?, farkasEntriesM rest with
    | some pe, some prest => some ((l, pe) :: prest)
    | _, _ => none

/-- The mixed-variable branch of the Farkas check: every polynomial —
entry vehicles, matched premise atoms, and the conclusion — rides the
mixed tagged translation; entry encoding, nonneg weights, atom matching,
and the exact combination fold mirror the sibling branches verbatim. -/
def farkasCheckM (cc : Term) (premise : Formula) (conclusion : Term) : Bool
\end{leansrc}
\declhead{BernsteinLogic.Core.Catalogue}{PFarkasData}{\badge{BadgeGreen}{statement bound}}
\begin{leansrc}
inductive PFarkasData : Hypernet → Sequent → Type where
  | mk (H : Hypernet) (seq : Sequent)
       (combinationCoeffs : Term)
       (conclusion : Term)
       -- The per-witness-node denominator bound, log₂ — claim semantics
       -- den(q) ≤ 2^denBoundLog2 over this node's witness rationals.
       -- ZERO proof-theoretic weight in v1 (no applicability axiom
       -- consumes it). Wire: u64 den_bound_log2 on tag 14 (BINARY.md §4).
       (denBoundLog2 : Nat)
       (conclShape : ∃ s, conclusion.bshape H = some s)
       (conclFin   : conclusion.finWires H)
       -- PROPOSED-B14 (approval-gated; realized provisionally per the
       -- B10/B13 pattern, 2026-07-16): every premise and conclusion wire
       -- is typed at a positive-size shape. Comparison atoms over
       -- zero-size (empty-tensor) wires hold vacuously, and a Farkas
       -- combination of two vacuous contradictory atoms — e.g. 5 ≤ w and
       -- w ≤ 3 at w : ⟨[0], Q⟩ — derives the false scalar conclusion
       -- ⊢ 2 ≤ 0 from a satisfiable premise. (The v0 singleton wire
       -- discipline excluded this by accident of strength; this is the
       -- honest general-shape guard.)
       (wiresPos : ∀ w ∈ seq.pre.freeWires ++ conclusion.freeWires,
                     ∃ τ, H.wireType w = some τ ∧
                       0 < Shape.totalSize τ.shape)
       (valid  : farkasCombinationValid combinationCoeffs seq.pre conclusion)
       (postEq : seq.post = .atom .le conclusion (.const 0))
       : PFarkasData H seq
\end{leansrc}
\declhead{BernsteinLogic.Core.Catalogue}{Formula.conjunctList}{\badge{BadgeGreen}{definition bound}}
\begin{leansrc}
/-- Flatten a conjunction tree into its leaf list (left-to-right). -/
def Formula.conjunctList : Formula → List Formula
  | .and l r => l.conjunctList ++ r.conjunctList
  | f => [f]
\end{leansrc}
\declhead{BernsteinLogic.Core.Catalogue}{atomLeZeroPolys}{\badge{BadgeGreen}{definition bound}}
\begin{leansrc}
/-- The `(e ≤ 0)`-normal-form polynomials entailed by one premise atom
(exact polynomial forms — never rescaled; F7 discipline). -/
def atomLeZeroPolys : Formula → List QArith.Poly
  | .atom .le l r => [QArith.Poly.sub l.toPoly r.toPoly]
  | .atom .lt l r => [QArith.Poly.sub l.toPoly r.toPoly]
  | .atom .ge l r => [QArith.Poly.sub r.toPoly l.toPoly]
  | .atom .gt l r => [QArith.Poly.sub r.toPoly l.toPoly]
  | .atom .eq l r => [QArith.Poly.sub l.toPoly r.toPoly,
                      QArith.Poly.sub r.toPoly l.toPoly]
  | _ => []
\end{leansrc}
\declhead{BernsteinLogic.Core.Catalogue}{farkasEntries}{\badge{BadgeGreen}{definition bound}}
\begin{leansrc}
/-- Decode the combination-coefficient Term (pinned canonical encoding:
right-nested `.add` chain of `.mul (.const λⱼ) eⱼ`, `.const 0` end). -/
def farkasEntries : Term → Option (List (Rat × Term))
  | .const c => if c = 0 then some [] else none
  | .add (.mul (.const l) e) rest =>
    (farkasEntries rest).map ((l, e) :: ·)
  | _ => none
\end{leansrc}
\declhead{BernsteinLogic.Core.Catalogue}{farkasCheckCore}{\badge{BadgeGreen}{definition bound}}
\begin{leansrc}
def farkasCheckCore (cc : Term) (premise : Formula) (conclusion : Term) : Bool :=
  match farkasEntries cc with
  | none => false
  | some entries =>
    let atoms := premise.conjunctList
    conclusion.idxFree &&
    entries.all (fun (l, e) =>
      decide (0 ≤ l) &&
      atoms.any fun a =>
        a.atomSidesIdxFree &&
        (atomLeZeroPolys a).any fun pa => decide (pa = e.toPoly)) &&
    decide ((entries.foldr
        (fun (l, e) acc => QArith.Poly.add (QArith.Poly.scale l e.toPoly) acc)
        []) = conclusion.toPoly)
\end{leansrc}
\declhead{BernsteinLogic.Core.Catalogue}{farkasCheck}{\badge{BadgeGreen}{definition bound}}
\begin{leansrc}
/-- The Farkas check: the legacy idx-free branch (exact `toPoly`
replay), OR the wire-aware idx coordinate branch (`farkasCheckW`),
OR the mixed tagged branch (`farkasCheckM`). Soundness
case-splits on the disjuncts; each branch carries its own argument
(`sound_farkas`). -/
def farkasCheck (cc : Term) (premise : Formula) (conclusion : Term) : Bool :=
  farkasCheckCore cc premise conclusion || farkasCheckW cc premise conclusion ||
  farkasCheckM cc premise conclusion
\end{leansrc}
\ruleblock{Correspondence (Steele display $\leftrightarrow$ Lean)}
\badge{BadgeAmber}{a bound Lean statement CHANGED since certification --- re-review and re-certify}
\sourcenote{The dictionary below is a reviewer-certified statement of how the mechanized statements above interpret the displayed rule: display metavariable $\mapsto$ Lean binder, each display hypothesis $\mapsto$ the premise or side condition carrying it, the display conclusion $\mapsto$ the goal --- with any deliberate difference in width between display and mechanization named outright. Certification is recorded as a hash triple over (the displayed cell, the extracted statement texts, this dictionary\textquotesingle s own text); if any of the three changes, this block turns amber until a reviewer re-certifies. The correspondence claim itself remains editorial --- reviewer-grade, not machine-checked; its machine-checked successor types each soundness statement against the rule as a structured object, and supersedes this block per rule as it lands.}
Judgment frame as in \rulename{Ax}. Distinctively, there is no
sub-derivation hypothesis at all: the displayed premise row describes
the syntactic conjunct structure of $\Phi$ itself, so the statement
consumes \texttt{seq.pre} directly and no \texttt{ih} exists.
\begin{itemize}
\item the combination $\sum_j \lambda_j \ell_j \mapsto$ the binder
  \texttt{conclusion : Term}; the multipliers $\lambda_j \mapsto$
  encoded in \texttt{combinationCoeffs : Term} under the pinned
  canonical encoding the intro describes.
\item the displayed premise row --- sub-conjunction of $\Phi$, each
  $\ell_j$ linear, $\bowtie_j \in \{\leq, =\}$, multipliers
  nonnegative and free exactly on equalities, and the identity of
  \texttt{conclusion} with the combination --- $\mapsto$ one
  hypothesis, \texttt{hvalid} $:$ \texttt{farkasCombinationValid
  combinationCoeffs seq.pre conclusion}, whose meaning this book binds
  by displaying the decoding chain in full (absorbed hypotheses, the
  \rulename{PConstExact} pattern: a successful check certifies the
  structure; the free-sign discipline is realized as both-polarity
  contribution of equality atoms). The displayed row is the
  rule\textquotesingle s published register; the realized matcher is
  wider --- see the width note below.
\item the displayed typing premises $\bshape_H(\sum \lambda_j \ell_j)
  \neq \bot$ and finite typing $\mapsto$ \texttt{hshape} (existence of
  a defined broadcast shape --- the \texttt{some}) and \texttt{hfin}
  $:$ \texttt{conclusion.finWires H}.
\item conclusion $\Phi \vdash \sum_j \lambda_j \ell_j \leq 0$
  $\mapsto$ the goal under \texttt{hpost} $:$ \texttt{seq.post = .atom
  .le conclusion (.const 0)}.
\end{itemize}
Width notes. On the side-condition surface the statement and the
display agree: \texttt{hwpos} --- every free wire of \texttt{seq.pre}
\emph{and} of the conclusion typed at a shape of strictly positive
total size --- appears inline as the display\textquotesingle s
positive-size row, and it is not optional: without it, zero-size wires
make atoms vacuously true and a combination of two vacuous atoms
derives a false scalar bound (the zero-size counterexample beside the
warrant proof). One direction is genuinely \emph{wider}: the matcher
behind \texttt{hvalid} normalizes arbitrary
two-sided comparison atoms over all five admitted comparators (strict
forms contribute their non-strict normal polynomials; equality atoms
both polarities) and the replay is degree-blind --- the displayed
linear zero-form $\{\leq, =\}$ row is the rule\textquotesingle s
published register (the note above), not the checked boundary.
The displayed cell is the rule\textquotesingle s intent; the certified
statement is that intent with the register widened exactly as stated.

\ruleblock{Lean verifier receipts}
\sourcenote{module \texttt{BernsteinLogic.Soundness}; bound build \texttt{a20d90e66c13\ldots} (bound from receipt a9c0e34d3e3c9c8c, 2026-08-03T21:30:36Z). A receipt records that named checking layers accepted a build: (module, compiled-artifact SHA-256, environment SHA-256, checker layers), minted when the checking gate ran; the store receipts attesting the bound build are listed with the layers each run credited. Nothing here asks the reader to trust the store: every layer is an independently runnable program over the quoted sources --- rebuild the module and re-run the checkers to reproduce the verdict. Receipts attesting superseded builds credit no layer and are omitted from this rendering (the count of omissions, if any, is noted below).}
\begin{receipttable}
\receiptrow{a9c0e34d3e3c9c8c}{2026-08-03T21:30:36Z}{build+decorrelated}{\badge{BadgeGreen}{current}}
\end{receipttable}
\sourcenote{21 further receipt(s) in the store attest superseded builds of this module and are omitted from this rendering.}
\sourcenote{module \texttt{BernsteinLogic.Core.Catalogue}; bound build \texttt{cbc92e976a1d\ldots} (bound from receipt 5abdb3589707b5df, 2026-08-03T21:30:33Z). A receipt records that named checking layers accepted a build: (module, compiled-artifact SHA-256, environment SHA-256, checker layers), minted when the checking gate ran; the store receipts attesting the bound build are listed with the layers each run credited. Nothing here asks the reader to trust the store: every layer is an independently runnable program over the quoted sources --- rebuild the module and re-run the checkers to reproduce the verdict. Receipts attesting superseded builds credit no layer and are omitted from this rendering (the count of omissions, if any, is noted below).}
\begin{receipttable}
\receiptrow{24804e66edba0bc9}{2026-08-03T20:25:09Z}{build+decorrelated}{\badge{BadgeGreen}{current}}
\receiptrow{5abdb3589707b5df}{2026-08-03T21:30:33Z}{build+decorrelated}{\badge{BadgeGreen}{current}}
\end{receipttable}
\sourcenote{18 further receipt(s) in the store attest superseded builds of this module and are omitted from this rendering.}
\layerline{\layeritem{build}{lake build (elaboration)}{\badge{BadgeGreen}{current}}\layeritem{replay}{leanchecker (official kernel replay)}{\badge{BadgeAmber}{stale only}}\layeritem{hardened}{safe\_verify (hardened screen)}{\badge{BadgeAmber}{stale only}}\layeritem{decorrelated}{lean4export + nanoda (decorrelated kernel)}{\badge{BadgeGreen}{current}}}
\sourcenote{\badge{BadgeGray}{comparator: none} no fixture in the three-implementation comparator harness (compiled / kernel-interpreted / reference-interpreted) exercises this rule.}
\ruleblock{HOL4 statement}
\sourcenote{theory \texttt{bern\_catalogue}, HOL4 development script \texttt{\texttt{bern\_catalogueScript.sml}} \badge{BadgeGreen}{script bound}}
\holhead{Definition farkasEntries\_def}
\begin{holsrc}
Definition farkasEntries_def:
  farkasEntries (TmConst c) = (if c = 0 then SOME [] else NONE) ∧
  farkasEntries (TmAdd (TmMul (TmConst l) e) rest) =
    (case farkasEntries rest of
     | NONE => NONE
     | SOME es => SOME ((l,e)::es)) ∧
  farkasEntries _ = NONE
End
\end{holsrc}
\holhead{Definition atomLeZeroPolys\_def}
\begin{holsrc}
Definition atomLeZeroPolys_def:
  atomLeZeroPolys (FAtom BowLe l r) = [polySub (termToPoly l) (termToPoly r)] ∧
  atomLeZeroPolys (FAtom BowLt l r) = [polySub (termToPoly l) (termToPoly r)] ∧
  atomLeZeroPolys (FAtom BowGe l r) = [polySub (termToPoly r) (termToPoly l)] ∧
  atomLeZeroPolys (FAtom BowGt l r) = [polySub (termToPoly r) (termToPoly l)] ∧
  atomLeZeroPolys (FAtom BowEq l r) = [polySub (termToPoly l) (termToPoly r);
                                       polySub (termToPoly r) (termToPoly l)] ∧
  atomLeZeroPolys _ = []
End
\end{holsrc}
\holhead{Definition farkasCheckCore\_def}
\begin{holsrc}
Definition farkasCheckCore_def:
  farkasCheckCore cc premise conclusion =
    case farkasEntries cc of
    | NONE => F
    | SOME entries =>
        (let atoms = conjunctList premise in
         termIdxFree conclusion ∧
         EVERY (λ(l,e).
           0 ≤ l ∧
           EXISTS (λa. atomSidesIdxFree a ∧
                       EXISTS (λpa. pa = termToPoly e) (atomLeZeroPolys a))
                  atoms) entries ∧
         (FOLDR (λ(l,e) acc. polyAdd (polyScale l (termToPoly e)) acc)
                [] entries =
          termToPoly conclusion))
End
\end{holsrc}
\holhead{Definition farkasCheckW\_def}
\begin{holsrc}
Definition farkasCheckW_def:
  farkasCheckW cc premise conclusion =
    case farkasEntries cc of NONE => F | SOME entries =>
    case farkasEntriesW entries of NONE => F | SOME pentries =>
    case asCoordPolyIdxW conclusion of NONE => F | SOME pc =>
      (let atoms = conjunctList premise in
       EVERY (λ(l,pe).
         0 ≤ l ∧
         EXISTS (λa. EXISTS (λpa. pa = pe) (atomLeZeroPolysW a)) atoms)
             pentries ∧
       (FOLDR (λ(l,pe) acc. polyAdd (polyScale l pe) acc) [] pentries =
        pc))
End
\end{holsrc}
\holhead{Definition farkasCheck\_def}
\begin{holsrc}
Definition farkasCheck_def:
  farkasCheck cc premise conclusion ⇔
    farkasCheckCore cc premise conclusion ∨
    farkasCheckW cc premise conclusion
End
\end{holsrc}
\holhead{Definition farkasCheckFull\_def}
\begin{holsrc}
Definition farkasCheckFull_def:
  farkasCheckFull cc premise conclusion ⇔
    farkasCheckCore cc premise conclusion ∨
    farkasCheckW cc premise conclusion ∨
    farkasCheckM cc premise conclusion
End
\end{holsrc}
\ruleblock{HOL verifier receipts}
\badge{BadgeBlue}{documentary} \texttt{hol-strict} CLEAN over the ported checker chain: theory
\texttt{bern\_catalogue}, 472 theorems, 0 cheat-tagged, 0 nonstandard
axioms (whitelist: \texttt{DISK\_THM}; the four \texttt{bool} axioms) ---
verification run of 2026-07-22 on an isolated build host, scripts
sha-matched verbatim to the sources cited here (git blob
\texttt{e41b9ed0\ldots}). The ported receipt table replays the
$\lambda = 1$ re-canonicalization accept. Downstream, the same chain
carries the checker to a CakeML binary with the
\texttt{machine\_code\_sound} tag (theory \texttt{bernCertProofARM8}).
Divergence repaired, dated: an earlier port run predated the Lean
checker\textquotesingle s third branch (the mixed tagged-union replay,
added 2026-07-19) and carried the two-branch \texttt{farkasCheck}
with a refusal-direction-only caveat; \texttt{farkasCheckFull} ---
displayed above --- is now the honest three-branch disjunction, and
the reference checker\textquotesingle s Farkas arm consumes it
verbatim (verified in the sources at the 2026-07-22 run cited here).
Semantic stratum, also ported: the soundness twin \texttt{sound\_farkas}
(theory \texttt{bern\_tr\_ready8}: \texttt{hol-strict} clean, 8
theorems, 0 offences, verification run of 2026-07-22 on an isolated
build host, content-addressed at the same sources), stated at
\texttt{farkasCheckFull} --- matching the widened checker arm verbatim
--- and consumed by the certified-fragment keystone
(\texttt{checkPf\_sound\_frag}, theory \texttt{bern\_tr\_splice}, 20
theorems, 0 offences, same run): a reference-checker accept inside the
proven wire fragment yields sequent validity.

\sourcenote{HOL receipts are documentary: quoted results of \texttt{hol-strict} runs
and census rows from the HOL4 development\textquotesingle s build
evidence, not minted receipt artefacts. A HOL receipt store (the Lean
store\textquotesingle s twin) is named future work.
}
\ruleblock{Proof}

\sourcenote{Genre: admissible rule (tier 2) --- a proof-tree-rewriting
recipe; the general admissible template is stated in the \rulename{Norm}
section and instantiated here. \rulename{Farkas} is admissible by
reduction to verified machinery: the affine-combination replay is
checked computation, and one verified lemma warrants every accepted
node.}

\begin{brecipe}[the \rulename{Farkas} instance]\label{rec:fk-instance}
\rulename{Farkas} has no proof premises --- its inputs are
sub-conjuncts of $\Phi$ itself, not sub-derivations:
\[
\inferrule*[Right=\rulename{Farkas}]{
  \{a_j : \ell_j \bowtie_j r_j\}_j \subseteq \Phi \text{ conjunct-wise,}
  \ p_j \ \text{\upshape their } ({\cdot}) \leq 0\text{\upshape-normal
  forms} \\
  \lambda_j \geq 0 \text{ (free-signed on equalities)} \\\\
  \textstyle\sum_j \lambda_j p_j \equiv c
  \ \text{\upshape(exact canonical fold)} \\
  c \ \text{\upshape shape-stable, finitely typed;}\;
  \wires(\Phi) \cup \wires(c) \ \text{\upshape positive-size}
}{\Phi \vdash c \leq 0}
\qquad\rightsquigarrow\qquad
\text{\upshape the checker replays the fold;}\;
\text{\texttt{sound\_farkas} stands as warrant.}
\]
The rewrite target contains zero catalogue nodes: the entire content is
the checked combination (decoded, sign-gated, atom-sourced, folded ---
the fragment display quoted above walks the chain) plus the semantic
argument of Proposition~\ref{prop:fk-sound}. The matcher is wider than
the catalogued cell's published register (linear zero-form
$\{\leq, =\}$ rows): it normalizes arbitrary two-sided comparison atoms
of the five admitted comparators --- strict forms contribute their
non-strict normal forms, semantically (Proposition~\ref{prop:fk-comb})
--- and the replay is degree-blind (the book note). The catalogued
stratum
records an unpacking sketch through \rulename{BernCone} at degree $1$;
the shipped warrant is direct, at all tensor shapes, heterogeneous
premise-atom shapes included.
\end{brecipe}

\begin{bprop}[nonnegative combination, per entry]\label{prop:fk-comb}
Work in an ordered field. Suppose values $a_j$ satisfy $a_j \leq 0$,
and weights $\lambda_j \geq 0$. Then
$\sum_j \lambda_j a_j \leq 0$. Moreover the $({\cdot} \leq 0)$-normal
reading of each premise atom delivers such values: an atom
$l \bowtie r$ with $\bowtie \in \{\leq, <\}$ yields
$l - r \leq 0$; with $\bowtie \in \{\geq, >\}$ it yields
$r - l \leq 0$; and an equality yields \emph{both}
$l - r \leq 0$ and $r - l \leq 0$ --- so a free-sign multiplier
$\nu$ on an equality is realized as the nonnegative weight $|\nu|$ on
the polarity matching $\nu$'s sign.
\end{bprop}
\begin{proof}
Each $\lambda_j a_j \leq \lambda_j \cdot 0 = 0$ (multiplication by a
nonnegative element preserves $\leq$), and a finite sum of nonpositive
elements is nonpositive, by induction. The polarity readings are the
definitions of the normal forms; a holding strict atom \emph{implies}
its non-strict normal form --- the matcher emits the same polynomial
for $<$ as for $\leq$, the weakening being a fact about holding, not a
checker step. The
equality discipline is the textbook sign convention for multipliers of
equality constraints \cite[\S7.3]{schrijver1986}, realized here
syntactically: the matcher sees both polarities, so a negative
multiplier is unrepresentable rather than checked for.
\end{proof}

\begin{bprop}[semantic warrant]\label{prop:fk-sound}
Suppose the payload's combination check accepts (any branch), the
conclusion term $c$ is shape-stable and finitely typed, and every wire
of $\Phi$ and $c$ is typed at a positive-size shape. If
$\sigma \models \Phi$ for an admissible $\sigma$, then
$\sigma \models c \leq 0$.
\end{bprop}
\begin{proof}
The shape and finiteness fields give $c$ a defined evaluation; fix a
conclusion entry $i$. The accepted branch fixes a variable space and a
per-entry valuation for it: the idx-free branch reads each bare wire
\emph{clamped} at entry $i$ --- a rank-0 wire reads its scalar; a
positive-rank wire of $m$ entries reads entry $i$ when $i < m$ and
entry $0$ otherwise, so the read is always an honest in-range entry of
the same tensor, and is the plain positional read wherever $i$ is in
range (a read that serves heterogeneous premise
atom shapes at every conclusion index); the wire-aware branch reads
each projection $w[j]$ at its fixed coordinate; the mixed branch does
both at once, on disjoint variable classes. Two facts from the
Setting's broadcast discipline make the clamped reads honest and exact.
First, a holding atom never compares unequal positive ranks, and a
term's defined broadcast shape forces every positive-rank wire in it to
carry that very shape (only rank-0 joins across shapes) --- so all
positive-rank wires of one premise atom share one shape, and the clamp
lands them at one \emph{common} honest entry of that atom, rank-0 wires
reading their scalar under any index; here positive size is
load-bearing: an atom over positively-sized wires has at least one
honest entry for the clamp to land on, so its per-entry normal form is
genuinely nonpositive at
the branch's valuation (\autoref{rem:fk-b14} shows what happens
otherwise). Second, the same discipline pins every positive-rank wire
of $c$ at $c$'s own shape, so at conclusion entry $i$ the clamp is
exact and the fold's value at the branch's valuation \emph{is} entry
$i$. Thus each sourced entry polynomial is the normal form of a premise
conjunct that holds at $\sigma$, read per-entry;
Proposition~\ref{prop:fk-comb} bounds the weighted fold by
$0$; the certified fold identity says the fold \emph{is} $c$'s
polynomial in the branch's variable space; evaluating at the branch's
valuation gives entry $i$ of $c$'s tensor $\leq 0$. Entries were
arbitrary, so the atom $c \leq 0$ holds. \texttt{sound\_farkas}
case-splits on the branch disjunct and runs exactly this argument
three times.
\end{proof}

\begin{bremark}[why positive-size wires are not decoration]
\label{rem:fk-b14}
Comparison atoms hold per-entry, and a zero-size wire has no entries:
at $w$ of shape $[0]$, the atoms $5 \leq w$ and $w \leq 3$ both hold
\emph{vacuously} --- yet their $({\cdot} \leq 0)$-normal forms
$5 - w$ and $w - 3$ are perfectly good polynomials, and with
$\lambda = 1$ each they fold to the constant $2$. Without the
positive-size condition the rule derives $\vdash 2 \leq 0$ from a
satisfiable premise: the classic mismatch between empty-product
semantics and nonempty algebra, and the reason the field exists. (A
narrower wire discipline once masked the failure by accident of
strength; the field is the honest general-shape guard.)
\end{bremark}

\begin{bremark}[why the conclusion's shape fields are not decoration]
\label{rem:fk-b13}
A conclusion mentioning an untyped wire --- say
$c' = 0 \cdot w_{999} + c$ with $w_{999}$ untyped --- has a correct
canonical polynomial: the dead monomial vanishes in canonical form, so
the fold identity holds exactly as it does for $c$. But evaluation is
structural on the \emph{term}, not on its canonical form --- it still
reads $w_{999}$, which is untyped --- so the conclusion never
evaluates, and the concluded atom never holds at any valuation. The shape-stability and finiteness
fields make ``the fold is $c$'s polynomial'' imply ``$c$
evaluates'', closing the gap between syntax and semantics.
\end{bremark}

\begin{bremark}[the three branches are one algorithm]
\label{rem:fk-branches}
The branches differ only in the variable space where polynomials are
compared: canonical wire-keyed forms; projection coordinates keyed by
an injective (wire, coordinate) pairing --- distinct wires \emph{must}
get distinct variables, since a composite may mix projections of
several wires and conflating $w_0[j]$ with $w_2[j]$ would let
semantically different polynomials collide (a one-line witness: from
the holding premise $w_0[0] \leq 0$ at $w_0[0] = -1$, $w_2[0] = 1$,
the conflated space accepts the false conclusion $w_2[0] \leq 0$); and
the mixed tagged union,
where bare wires and projections are disjoint classes --- in
particular a bare $w$ is never identified with $w[0]$: the
identification variant admits a false sequent (witness: $w = (-1, 1)$
satisfies $w[0] \leq 0$, and the identified space would conclude the
whole-wire bound $w \leq 0$, false at the second entry), so the
tagging is a
soundness fact, not a convenience. Entry encoding, sign gates, atom
sourcing, and the exact fold are the same gates, replayed branch by
branch in its own variable space (the branch definitions cited at the
anchors), and each
branch carries its own per-entry valuation in the soundness proof.
\end{bremark}

\begin{bremark}[relation to the literature]
The linear ancestor is Farkas' lemma \cite{farkas1902}; the multiplier
sign conventions ($\lambda \geq 0$ on inequalities, free on equalities)
are \cite[\S7]{schrijver1986}. As everywhere in this calculus, only
the sound (trivial) direction is claimed: the classical lemma's
content --- that every valid affine consequence \emph{has} such a
combination --- is a completeness statement about producers, never
consulted by the kernel. Degree-blindness of the replay places the
rule's polynomial generalizations in the Krivine--Handelman family
\cite{krivine1964,handelman1988}, reached in this catalogue through
\rulename{Syzygy} and \rulename{BernCone} rather than by widening this
rule.
\end{bremark}

\begin{anchors}
Recipe \ref{rec:fk-instance} warrant & \texttt{sound\_farkas} (branch case split); \texttt{farkasCombinationValid}\\
entry decoding (pinned encoding) & \texttt{farkasEntries}\\
atom normal forms, both polarities on equality & \texttt{atomLeZeroPolys} (idx-free); \texttt{atomLeZeroPolysW}, \texttt{atomLeZeroPolysM} (widened branches)\\
the three branches & \texttt{farkasCheckCore}, \texttt{farkasCheckW}, \texttt{farkasCheckM}; disjunction \texttt{farkasCheck}\\
projection variable spaces (Remark \ref{rem:fk-branches}) & \texttt{wireCoordKey}, \texttt{Term.asCoordPolyIdxW?}, \texttt{Term.asCoordPolyMixed?}\\
Proposition \ref{prop:fk-comb} fold & \texttt{farkas\_fold\_eval\_nonpos}, \texttt{farkas\_fold\_eval\_nonpos\_polys}\\
per-entry valuations & the clamped, coordinate, and mixed entry valuations consumed by \texttt{sound\_farkas}'s three cases\\
positive-size guard (Remark \ref{rem:fk-b14}) & payload field \texttt{wiresPos} (\texttt{PFarkasData})\\
conclusion evaluability (Remark \ref{rem:fk-b13}) & payload fields \texttt{conclShape}, \texttt{conclFin}\\
wire placement & \texttt{RuleTag.farkas}\\
\end{anchors}


\rulesection{18}{Subdivide}{admissible rule}
\label{blcatrule-Subdivide}
\ruleblock{The rule (Steele metanotation)}
\sourcenote{as realized in the mechanized development; source row: \texttt{\texttt{rules.toml}} --- \emph{Subdivide}. The metanotation statement is typeset by hand for notational quality; its mechanized ground truth is the Lean statement below, extracted from the sources at every build of this book and hash-bound. The rule \emph{roster} is machine-checked (every mechanized rule tag round-trips to exactly one catalogue row); the statement \emph{text} is editorial, checked by eye against the extracted Lean beside it.}
\begin{mathpar}
\inferrule*[Right=\rulename{Subdivide}]{
  \Phi \vdash x \in \mathcal{R}(\Delta, S)
  \\
  \{(\Delta_k, S_k)\}_{k=1}^N \text{ finite, with cover side-condition } \mathcal{R}(\Delta, S) \subseteq \textstyle\bigcup_k \mathcal{R}(\Delta_k, S_k)
  \\
  \forall k.\; \Phi \;\wedge\; (x \in \mathcal{R}(\Delta_k, S_k)) \;\vdash\; \Psi
}{
  \Phi \vdash \Psi
}
\end{mathpar}

\par\medskip
\rulename{Subdivide} is the case-split rule of region reasoning: from a
proven membership $\Phi \vdash x \in \mathcal{R}(\Delta, S)$, a finite
family of branch regions covering the parent, and one sub-proof of the
goal per branch (each under the strengthened premise
$\Phi \wedge x \in \mathcal{R}(\Delta_k, S_k)$), it concludes
$\Phi \vdash \Psi$ on the whole parent region. It is degree
escalation\textquotesingle s partner and load-bearing for every non-toy
proof: where a single \rulename{BernCone} certificate at feasible degree
is too coarse, a producer splits the region, certifies each piece at low
degree, and lets \rulename{Subdivide} reassemble the global conclusion
--- the branch-and-bound shape that makes Bernstein enclosures converge.

What the checker recomputes, given the payload (parent region, branch
family, per-branch sub-proofs): that the parent region reads as a closed
interval --- ambient dimension 1, support exactly two vertices at
distinct coordinates, empty extension set; that every branch region
reads as a closed interval and is tied to its own extension set; and
that the branch intervals \emph{chain} across the parent interval --- a
greedy reach, iterated to its fixpoint, that extends past any interval
whose left endpoint has already been reached and must arrive at the
parent\textquotesingle s right endpoint. Everything beyond this
fragment (extension-ray parents or branches, regions of dimension
$\geq 2$, supports of more than two vertices) is refused, never
accepted on faith. The repair for that narrowness is
\rulename{ConvIntro}: $n$-dimensional box grids arise by
\emph{nesting} per-coordinate splits across several rule
applications, not inside one --- each branch establishes its
coordinate facts as order atoms, and \rulename{ConvIntro} converts
them back into the region atom the next coordinate\textquotesingle s
split consumes, so the one-dimensional rule reaches every box grid
without widening its checker.

Each hypothesis is load-bearing. Drop the cover check and a producer
omits the branch where the goal fails --- a gap family such as
$[1, \frac{4}{3}] \cup [\frac{3}{2}, 2]$ over $[1, 2]$ proves the goal
everywhere except where it is false. Drop the branch-region tie and the
chained intervals belong to different regions than the ones the
sub-proofs assume. Drop the branch-premise shape and a sub-proof could
be consulted outside the piece it was proven on. Drop the
branch-conclusion tie and the branches may prove any formulas whatever
--- $\Psi$ never enters their obligations while the rule concludes it
anyway. Drop the parent region premise and a valuation satisfying
$\Phi$ with the point outside the parent falls through every branch:
the cover says nothing about it, yet the rule still concludes $\Psi$
there. The certificate slot
for cover evidence rides the wire but is deliberately \emph{not
consulted}: the checker recomputes the cover from the branch regions
themselves, so ignoring the slot can only make checking stricter, never
unsound.

\cataloguenote{The region variable itself plays no role in the cover geometry: the
soundness argument feeds the \emph{same} evaluated point to the parent
and to the selected branch, so no tie between the term and the cover is
needed --- the cover is a fact about intervals, not about wires. A
future interpretation may read per-gap infeasibility certificates out of
the wire slot; the catalogued stratum records the corresponding
unpacking sketch (per-branch cuts discharged through
\rulename{Farkas} and $\vee$-elimination).
}
\ruleblock{Lean statement}
\declhead{BernsteinLogic.Soundness}{sound\_subdivide}{\badge{BadgeGreen}{statement bound}}
\begin{leansrc}
/-- Soundness of Subdivide (the covering argument): if the region
premise puts the evaluated point in the parent region, the checker's
cover chains the branch intervals across the parent interval, and every
branch sequent is valid on its region, then the conclusion holds on the
parent. -/
theorem sound_subdivide {H : Hypernet} {seq regionPremise : Sequent}
    {nBranches : Nat} {branchPremise : Fin nBranches → Sequent}
    {termX : Term} {simplex : Simplex} {S : ExtensionSet}
    {branchSimplex : Fin nBranches → Simplex}
    {branchS : Fin nBranches → ExtensionSet}
    {coverCertificate : Term}
    (ihRegion : regionPremise.semanticallyValid H)
    (ihBranch : ∀ k, (branchPremise k).semanticallyValid H)
    (hpre : seq.pre = regionPremise.pre)
    (hregionPost : regionPremise.post = .inRegion termX simplex S)
    (hbranchPre : ∀ k, (branchPremise k).pre
                   = .and regionPremise.pre
                          (.inRegion termX (branchSimplex k) (branchS k)))
    (hbranchPost : ∀ k, (branchPremise k).post = seq.post)
    (hcover : coverInfeasibilityValid termX simplex S nBranches
                branchSimplex branchS coverCertificate) :
    seq.semanticallyValid H
\end{leansrc}
\declhead{BernsteinLogic.Soundness}{interval1D\_elim}{\badge{BadgeGreen}{statement bound}}
\begin{leansrc}
/-- Elim: a `MaskRegion` witness against a successful interval reading
lands the point's coordinate in `[lo, hi]`, and the interval is proper. -/
theorem interval1D_elim {S : ExtensionSet} {lo hi : ℚ}
    (hint : S.interval1D? = some (lo, hi))
    {x : Fin S.simplex.complex.ambDim → ℚ}
    (hmask : MaskRegion S.simplex.complex S.simplex.contains S.contains x)
    (j0 : Fin S.simplex.complex.ambDim) :
    lo ≤ x j0 ∧ x j0 ≤ hi ∧ lo < hi
\end{leansrc}
\declhead{BernsteinLogic.Soundness}{interval1D\_intro}{\badge{BadgeGreen}{statement bound}}
\begin{leansrc}
/-- Intro: a coordinate inside a successful interval reading assembles a
`MaskRegion` witness (the extension set is empty in this fragment, so
both barycentric weights come out nonnegative). -/
theorem interval1D_intro {S : ExtensionSet} {lo hi : ℚ}
    (hint : S.interval1D? = some (lo, hi))
    {x : Fin S.simplex.complex.ambDim → ℚ}
    (j0 : Fin S.simplex.complex.ambDim)
    (hxlo : lo ≤ x j0) (hxhi : x j0 ≤ hi) :
    MaskRegion S.simplex.complex S.simplex.contains S.contains x
\end{leansrc}
\declhead{BernsteinLogic.Soundness}{ChainInv}{\badge{BadgeGreen}{statement bound}}
\begin{leansrc}
/-- The greedy chain invariant: the reach never sits below `lo`, and
once it has moved past `lo`, every point of `[lo, reach]` lies in some
interval of the FULL list. -/
def ChainInv (ivs : List (ℚ × ℚ)) (lo r : ℚ) : Prop
\end{leansrc}
\declhead{BernsteinLogic.Soundness}{chainInv\_extendPass}{\badge{BadgeGreen}{statement bound}}
\begin{leansrc}
/-- One pass preserves the invariant (stated over a sublist so the foldl
induction goes through; the covering intervals always come from the full
list). -/
theorem chainInv_extendPass (ivs : List (ℚ × ℚ)) (lo : ℚ) :
    ∀ (sub : List (ℚ × ℚ)), (∀ iv ∈ sub, iv ∈ ivs) →
    ∀ (r : ℚ), ChainInv ivs lo r →
    ChainInv ivs lo (coverExtendPass sub r)
\end{leansrc}
\declhead{BernsteinLogic.Soundness}{chainInv\_reach}{\badge{BadgeGreen}{statement bound}}
\begin{leansrc}
theorem chainInv_reach (ivs : List (ℚ × ℚ)) (lo : ℚ) :
    ∀ (n : Nat) (r : ℚ), ChainInv ivs lo r →
    ChainInv ivs lo (coverReach n ivs r)
\end{leansrc}
\declhead{BernsteinLogic.Soundness}{coverChain\_covers}{\badge{BadgeGreen}{statement bound}}
\begin{leansrc}
/-- Soundness of the greedy chain check: on success, every point of
`[lo, hi]` lies in some interval of the list. -/
theorem coverChain_covers {ivs : List (ℚ × ℚ)} {lo hi : ℚ} (hlt : lo < hi)
    (hcheck : coverChainCheck ivs lo hi = true) :
    ∀ x : ℚ, lo ≤ x → x ≤ hi → ∃ iv ∈ ivs, iv.1 ≤ x ∧ x ≤ iv.2
\end{leansrc}
\par\medskip
\textbf{What the cover check means --- and what it does not.}\enspace
The side condition appears in the payload as the single name
\texttt{coverInfeasibilityValid}, so this book binds its meaning by
displaying the definition chain in full, from the checker sources.
Reading the chain: \texttt{ExtensionSet.interval1D?} is the
closed-interval reading of a region --- it succeeds only on the checkable
fragment (dimension 1, two-vertex support at distinct coordinates,
empty extension set), returning sorted endpoints, and everything else
reads \texttt{none}, which the cover check turns into refusal.
\texttt{coverExtendPass} is one greedy pass: extend the reach past any
interval whose left endpoint has been reached;
\texttt{coverReach} iterates it, and as many passes as there are
intervals suffice for the fixpoint (a productive pass strictly
increases the reach onto a new right endpoint, of which there are only
finitely many). \texttt{coverChainCheck} asks whether the reach from
the parent\textquotesingle s left endpoint arrives at its right
endpoint; \texttt{coverCheck} assembles the whole condition, tying each
branch\textquotesingle s extension set to its simplex on the way.

What a skeptical reader should confirm by eye: the region term and the
wire-borne cover certificate are \emph{ignored} --- the relation
discards both arguments, and the cover is recomputed from the branch
regions alone (the recompute pattern; a forged certificate cannot
influence the verdict, and a missing one costs nothing); no valuation
is consulted anywhere; and the branch sub-proofs enter soundness, not
the cover computation --- the checker decides pure interval geometry.

\declhead{BernsteinLogic.Core.Catalogue}{ExtensionSet.interval1D?}{\badge{BadgeGreen}{definition bound}}
\begin{leansrc}
/-- The closed-interval reading of a region `R(Δ, S)`: succeeds iff the
ambient dimension is 1, the simplex's support is exactly two vertices at
distinct coordinates, and the extension set is empty — returning the
region's endpoints, sorted (`lo < hi` by construction). Everything else
reads `none` (and the cover check then rejects). -/
def ExtensionSet.interval1D? (S : ExtensionSet) : Option (Rat × Rat) :=
  if h : S.simplex.complex.ambDim = 1 then
    if S.isEmptyB then
      match (List.finRange S.simplex.complex.nVertices).filter
          (fun v => S.simplex.contains v) with
      | [a, b] =>
        let j0 : Fin S.simplex.complex.ambDim := ⟨0, by omega⟩
        let V₀ := S.simplex.complex.coord a j0
        let V₁ := S.simplex.complex.coord b j0
        if V₀ < V₁ then some (V₀, V₁)
        else if V₁ < V₀ then some (V₁, V₀)
        else none
      | _ => none
    else none
  else none
\end{leansrc}
\declhead{BernsteinLogic.Core.Catalogue}{coverExtendPass}{\badge{BadgeGreen}{definition bound}}
\begin{leansrc}
/-- One greedy pass of the 1-D chain cover: extend the reach past any
interval whose left endpoint has already been reached. -/
def coverExtendPass (ivs : List (Rat × Rat)) (r : Rat) : Rat :=
  ivs.foldl (fun r iv => if iv.1 ≤ r ∧ r < iv.2 then iv.2 else r) r
\end{leansrc}
\declhead{BernsteinLogic.Core.Catalogue}{coverReach}{\badge{BadgeGreen}{definition bound}}
\begin{leansrc}
/-- Iterated greedy reach. `ivs.length` passes suffice for the fixpoint:
a productive pass strictly increases the reach onto a new interval's
right endpoint (of which there are at most `ivs.length`), and an
unproductive pass is already the fixpoint. -/
def coverReach : Nat → List (Rat × Rat) → Rat → Rat
  | 0, _, r => r
  | n + 1, ivs, r => coverReach n ivs (coverExtendPass ivs r)
\end{leansrc}
\declhead{BernsteinLogic.Core.Catalogue}{coverChainCheck}{\badge{BadgeGreen}{definition bound}}
\begin{leansrc}
/-- Greedy 1-D chain check: `[lo, hi] ⊆ ⋃ ivs`, all intervals closed.
Soundness: `coverChain_covers` (Soundness §S5′). -/
def coverChainCheck (ivs : List (Rat × Rat)) (lo hi : Rat) : Bool :=
  decide (hi ≤ coverReach ivs.length ivs lo)
\end{leansrc}
\declhead{BernsteinLogic.Core.Catalogue}{coverCheck}{\badge{BadgeGreen}{definition bound}}
\begin{leansrc}
/-- The Subdivide cover computation (fragment documented at the section
header): the parent region is an interval, every branch region is an
interval tied to its own extension set, and the branch intervals chain
across the parent's. -/
def coverCheck (simplex : Simplex) (S : ExtensionSet) (nBranches : Nat)
    (branchSimplex : Fin nBranches → Simplex)
    (branchS : Fin nBranches → ExtensionSet) : Bool :=
  match S.interval1D? with
  | none => false
  | some (lo, hi) =>
    decide (S.simplex = simplex) &&
    ((List.finRange nBranches).all fun k =>
      decide ((branchS k).simplex = branchSimplex k) &&
      ((branchS k).interval1D?).isSome) &&
    coverChainCheck
      ((List.finRange nBranches).filterMap fun k => (branchS k).interval1D?)
      lo hi
\end{leansrc}
\declhead{BernsteinLogic.Core.Catalogue}{coverInfeasibilityValid}{\badge{BadgeGreen}{definition bound}}
\begin{leansrc}
def coverInfeasibilityValid
    (_termX : Term) (simplex : Simplex) (S : ExtensionSet)
    (nBranches : Nat)
    (branchSimplex : Fin nBranches → Simplex)
    (branchS : Fin nBranches → ExtensionSet)
    (_coverCert : Term) : Prop :=
  coverCheck simplex S nBranches branchSimplex branchS = true
\end{leansrc}
\ruleblock{Correspondence (Steele display $\leftrightarrow$ Lean)}
\badge{BadgeGreen}{reviewer-certified against the displayed texts}
\sourcenote{The dictionary below is a reviewer-certified statement of how the mechanized statements above interpret the displayed rule: display metavariable $\mapsto$ Lean binder, each display hypothesis $\mapsto$ the premise or side condition carrying it, the display conclusion $\mapsto$ the goal --- with any deliberate difference in width between display and mechanization named outright. Certification is recorded as a hash triple over (the displayed cell, the extracted statement texts, this dictionary\textquotesingle s own text); if any of the three changes, this block turns amber until a reviewer re-certifies. The correspondence claim itself remains editorial --- reviewer-grade, not machine-checked; its machine-checked successor types each soundness statement against the rule as a structured object, and supersedes this block per rule as it lands.}
Judgment frame as in \rulename{Ax}; sub-derivation premises as at
\rulename{Weaken}.
\begin{itemize}
\item $x \mapsto$ \texttt{termX : Term} (a free term --- unlike
  \rulename{BernCone}, no wire tie is needed here: the note above
  explains why); $(\Delta, S) \mapsto$ \texttt{simplex},
  \texttt{S}; the branch family $\{(\Delta_k, S_k)\}_{k=1}^N
  \mapsto$ \texttt{branchSimplex}, \texttt{branchS} $:$ \texttt{Fin
  nBranches → \_}.
\item region premise $\Phi \vdash x \in \mathcal{R}(\Delta, S)$
  $\mapsto$ \texttt{ihRegion} with \texttt{hregionPost} $:$
  \texttt{regionPremise.post = .inRegion termX simplex S}; $\Phi$
  shared via \texttt{hpre}.
\item the branch judgments $\forall k.\ \Phi \wedge (x \in
  \mathcal{R}(\Delta_k, S_k)) \vdash \Psi$ $\mapsto$
  \texttt{ihBranch} with \texttt{hbranchPre} pinning each branch
  premise to the displayed conjunction in the displayed order ($\Phi$
  left) and \texttt{hbranchPost} tying every branch conclusion to
  \texttt{seq.post} --- $\Psi$ is whatever formula the conclusion
  sequent carries; nothing constrains its shape.
\item the cover side-condition $\mathcal{R}(\Delta, S) \subseteq
  \bigcup_k \mathcal{R}(\Delta_k, S_k)$ $\mapsto$ \texttt{hcover} $:$
  \texttt{coverInfeasibilityValid}, i.e.\ \texttt{coverCheck = true}:
  the parent extension set must present as a 1-dimensional interval
  (\texttt{ExtensionSet.interval1D?} --- outside that fragment the
  check refuses, so the certified reach of the displayed
  set-inclusion is the interval fragment, the
  \rulename{ConvElim} class), each branch\textquotesingle s
  simplex-tie is checked, and the greedy chain
  (\texttt{coverExtendPass}/\texttt{coverReach}/\texttt{coverChainCheck},
  certified alongside) verifies the branch intervals chain across the
  parent interval --- a decidable \emph{sufficient} witness of the
  displayed inclusion, conservative like every refusal boundary.
\item consumed by nothing: the payload\textquotesingle s cover
  certificate term rides into \texttt{coverInfeasibilityValid} as an
  ignored argument (the binder is underscore-named in the displayed
  definition) --- retained for the per-gap-infeasibility
  interpretation the note records; the realized checker recomputes the
  cover itself and reads no producer-supplied cover data.
\end{itemize}
Width note: within the interval fragment neither side is wider ---
every other hypothesis is consumed, and the displayed rule\textquotesingle
s ``finite family\textquotesingle\textquotesingle{} is carried by the
\texttt{Fin nBranches} index type itself.

\ruleblock{Lean verifier receipts}
\sourcenote{module \texttt{BernsteinLogic.Soundness}; bound build \texttt{a20d90e66c13\ldots} (bound from receipt a9c0e34d3e3c9c8c, 2026-08-03T21:30:36Z). A receipt records that named checking layers accepted a build: (module, compiled-artifact SHA-256, environment SHA-256, checker layers), minted when the checking gate ran; the store receipts attesting the bound build are listed with the layers each run credited. Nothing here asks the reader to trust the store: every layer is an independently runnable program over the quoted sources --- rebuild the module and re-run the checkers to reproduce the verdict. Receipts attesting superseded builds credit no layer and are omitted from this rendering (the count of omissions, if any, is noted below).}
\begin{receipttable}
\receiptrow{a9c0e34d3e3c9c8c}{2026-08-03T21:30:36Z}{build+decorrelated}{\badge{BadgeGreen}{current}}
\end{receipttable}
\sourcenote{21 further receipt(s) in the store attest superseded builds of this module and are omitted from this rendering.}
\sourcenote{module \texttt{BernsteinLogic.Core.Catalogue}; bound build \texttt{cbc92e976a1d\ldots} (bound from receipt 5abdb3589707b5df, 2026-08-03T21:30:33Z). A receipt records that named checking layers accepted a build: (module, compiled-artifact SHA-256, environment SHA-256, checker layers), minted when the checking gate ran; the store receipts attesting the bound build are listed with the layers each run credited. Nothing here asks the reader to trust the store: every layer is an independently runnable program over the quoted sources --- rebuild the module and re-run the checkers to reproduce the verdict. Receipts attesting superseded builds credit no layer and are omitted from this rendering (the count of omissions, if any, is noted below).}
\begin{receipttable}
\receiptrow{24804e66edba0bc9}{2026-08-03T20:25:09Z}{build+decorrelated}{\badge{BadgeGreen}{current}}
\receiptrow{5abdb3589707b5df}{2026-08-03T21:30:33Z}{build+decorrelated}{\badge{BadgeGreen}{current}}
\end{receipttable}
\sourcenote{18 further receipt(s) in the store attest superseded builds of this module and are omitted from this rendering.}
\layerline{\layeritem{build}{lake build (elaboration)}{\badge{BadgeGreen}{current}}\layeritem{replay}{leanchecker (official kernel replay)}{\badge{BadgeAmber}{stale only}}\layeritem{hardened}{safe\_verify (hardened screen)}{\badge{BadgeAmber}{stale only}}\layeritem{decorrelated}{lean4export + nanoda (decorrelated kernel)}{\badge{BadgeGreen}{current}}}
\sourcenote{\badge{BadgeGray}{comparator: none} no fixture in the three-implementation comparator harness (compiled / kernel-interpreted / reference-interpreted) exercises this rule.}
\ruleblock{HOL4 statement}
\sourcenote{theory \texttt{bern\_catalogue}, HOL4 development script \texttt{\texttt{bern\_catalogueScript.sml}} \badge{BadgeGreen}{script bound}}
\holhead{Definition extInterval1D\_def}
\begin{holsrc}
Definition extInterval1D_def:
  extInterval1D es =
    if es.simplex.complex.ambDim = 1 then
      (if extIsEmptyB es then
         (case FILTER (λv. simplexContainsNat es.simplex v)
                      (rangeL es.simplex.complex.nVertices) of
          | [a; b] =>
              (let v0 = getD (getD es.simplex.complex.coord a []) 0 (0:rat) in
               let v1 = getD (getD es.simplex.complex.coord b []) 0 (0:rat) in
               if v0 < v1 then SOME (v0, v1)
               else if v1 < v0 then SOME (v1, v0)
               else NONE)
          | _ => NONE)
       else NONE)
    else NONE
End
\end{holsrc}
\holhead{Definition coverExtendPass\_def}
\begin{holsrc}
Definition coverExtendPass_def:
  coverExtendPass ivs (r:rat) =
    FOLDL (λr iv. if FST iv ≤ r ∧ r < SND iv then SND iv else r) r ivs
End
\end{holsrc}
\holhead{Definition coverReach\_def}
\begin{holsrc}
Definition coverReach_def:
  coverReach 0 ivs (r:rat) = r ∧
  coverReach (SUC n) ivs r = coverReach n ivs (coverExtendPass ivs r)
End
\end{holsrc}
\holhead{Definition coverChainCheck\_def}
\begin{holsrc}
Definition coverChainCheck_def:
  coverChainCheck ivs (lo:rat) hi = (hi ≤ coverReach (LENGTH ivs) ivs lo)
End
\end{holsrc}
\holhead{Definition coverCheck\_def}
\begin{holsrc}
Definition coverCheck_def:
  coverCheck sx es branches =
    case extInterval1D es of
    | NONE => F
    | SOME (lo, hi) =>
        (es.simplex = sx ∧
         EVERY (λ(bsx,bes). bes.simplex = bsx ∧ IS_SOME (extInterval1D bes))
               branches ∧
         coverChainCheck
           (FLAT (MAP (λ(bsx,bes). case extInterval1D bes of
                                   | NONE => []
                                   | SOME iv => [iv]) branches))
           lo hi)
End
\end{holsrc}
\ruleblock{HOL verifier receipts}
\badge{BadgeBlue}{documentary} \texttt{hol-strict} CLEAN over the ported checker chain: theory
\texttt{bern\_catalogue}, 472 theorems, 0 cheat-tagged, 0 nonstandard
axioms (whitelist: \texttt{DISK\_THM}; the four \texttt{bool} axioms) ---
verification run of 2026-07-22 on an isolated build host, scripts
sha-matched verbatim to the sources cited here (git blob
\texttt{e41b9ed0\ldots}). The ported receipt table replays the cover
battery: the positive chain
($[1, \frac{3}{2}] \cup [\frac{3}{2}, 2]$ over $[1, 2]$) and four
negative controls --- the gap mutant, a point-supported branch, an
extension-ray parent, and untied branch simplices. Representation seam,
inventoried in the script: the indexed branch family becomes one list
of (simplex, extension-set) pairs in index order, and the Lean
relation\textquotesingle s two inert arguments (region term, cover
certificate) are not reproduced --- both value-preserving deltas under
the port\textquotesingle s seam discipline.

\sourcenote{HOL receipts are documentary: quoted results of \texttt{hol-strict} runs
and census rows from the HOL4 development\textquotesingle s build
evidence, not minted receipt artefacts. A HOL receipt store (the Lean
store\textquotesingle s twin) is named future work.
}
\ruleblock{Proof}

\sourcenote{Genre: admissible rule (tier 2) --- a proof-tree-rewriting
recipe; the general admissible template is stated in the \rulename{Norm}
section and instantiated here. \rulename{Subdivide} is the genuinely
one-node-to-many member of this batch: its rewrite is a case analysis
over the branch sub-proofs, and the checked cover condition is what
makes the case analysis exhaustive.}

\begin{brecipe}[the \rulename{Subdivide} instance]\label{rec:sd-instance}
The node carries a region premise and one sub-proof per branch:
\[
\inferrule*[Right=\rulename{Subdivide}]{
  \mathcal{T}_0 \rhd \bigl(\Phi \vdash x \in \mathcal{R}(\Delta, S)\bigr) \\\\
  \mathcal{T}_k \rhd \bigl(\Phi \wedge (x \in \mathcal{R}(\Delta_k, S_k))
    \vdash \Psi\bigr) \quad (k = 1, \dots, N) \\\\
  \mathcal{R}(\Delta, S) \subseteq \textstyle\bigcup_k
    \mathcal{R}(\Delta_k, S_k) \ \text{\upshape(checked: the interval
    chain)}
}{\Phi \vdash \Psi}
\qquad\rightsquigarrow\qquad
\text{\upshape select the covering branch; consult its subtree.}
\]
Conceptually the rewrite is an $N$-way case analysis: at any valuation
satisfying $\Phi$, the premise puts the evaluated point in the parent
region, the cover places it in some branch region $k$, and
$\mathcal{T}_k$'s strengthened premise then holds --- so $\mathcal{T}_k$
delivers $\Psi$. The shipped warrant packages this selection argument as
one verified lemma rather than as catalogue surgery (the catalogued
stratum records the corresponding derivation sketch through
\rulename{Farkas}-discharged infeasibility cuts and
$\vee$-elimination); the checked side condition is
Proposition~\ref{prop:sd-chain}, the exhaustiveness of the case
analysis.
\end{brecipe}

\begin{bprop}[the greedy chain covers]\label{prop:sd-chain}
Let $\mathrm{ivs}$ be a finite list of closed intervals and
$\mathrm{lo} < \mathrm{hi}$ rationals. Define one greedy pass
$r \mapsto \mathrm{extend}(r)$: fold over the list, replacing $r$ by an
interval's right endpoint whenever that interval's left endpoint is
$\leq r$ and its right endpoint is $> r$. If after
$|\mathrm{ivs}|$ passes from $\mathrm{lo}$ the reach is
$\geq \mathrm{hi}$, then
$[\mathrm{lo}, \mathrm{hi}] \subseteq \bigcup \mathrm{ivs}$.
\end{bprop}
\begin{proof}
Maintain the invariant: the reach $r$ satisfies $\mathrm{lo} \leq r$,
and either $r \leq \mathrm{lo}$ or every point of
$[\mathrm{lo}, r]$ lies in some interval of the \emph{full} list. It
holds initially ($r = \mathrm{lo}$, the first disjunct). One pass
preserves it: when the fold takes an interval $[a, b]$ with
$a \leq r < b$, any point $y \in [\mathrm{lo}, b]$ is either
$\leq r$ --- covered by the invariant, or by $[a, b]$ itself in the
$r \leq \mathrm{lo}$ base case since then $a \leq \mathrm{lo} \leq y$
--- or lies in $(r, b] \subseteq [a, b]$; and when the fold declines an
interval the state is unchanged. As many passes as there are intervals
reach the fixpoint: a productive pass strictly increases $r$ onto a
\emph{new} interval's right endpoint, of which there are at most
$|\mathrm{ivs}|$, and an unproductive pass is already the fixpoint. If
the final reach is $\geq \mathrm{hi} > \mathrm{lo}$, the invariant's
second disjunct is in force and covers all of
$[\mathrm{lo}, \mathrm{hi}]$.
\end{proof}

\begin{bprop}[semantic warrant]\label{prop:subdiv-sound}
Suppose the region premise and every branch premise-sequent are
semantically valid, the branch premises have the checked shape
($(\mathrm{branch}\ k)$'s premise is
$\Phi \wedge (x \in \mathcal{R}(\Delta_k, S_k))$ and its conclusion is
$\Psi$), and the cover check accepts. Then $\Phi \vdash \Psi$ is
semantically valid.
\end{bprop}
\begin{proof}
Fix an admissible $\sigma \models \Phi$. The region premise puts the
evaluated point $q$ of $x$ in $\mathcal{R}(\Delta, S)$; the accepted
cover check says the parent reads as a closed interval
$[\mathrm{lo}, \mathrm{hi}]$ (dimension 1, two-vertex support,
$S = \emptyset$), so $q$'s coordinate lies in
$[\mathrm{lo}, \mathrm{hi}]$ --- eliminating the parent's barycentric
weights to endpoint bounds. By Proposition~\ref{prop:sd-chain} some
branch interval contains the coordinate. Each branch interval was
itself produced by running the closed-interval reading on that
branch's \emph{own} simplex and extension set (the cover check
refuses any branch it cannot so read, and checks the tie between the
two), so ``the coordinate lies in branch $k$'s interval'' converts
back --- re-introducing barycentric weights over branch $k$'s own
support from the two endpoint bounds --- into membership of the
\emph{same} evaluated point in $\mathcal{R}(\Delta_k, S_k)$: this is
where nothing about the region term is needed beyond its one
evaluation, used twice.
Hence $\sigma$ satisfies branch $k$'s premise
$\Phi \wedge (x \in \mathcal{R}(\Delta_k, S_k))$; its validity delivers
its conclusion, which is $\Psi$ by the checked branch-conclusion
shape.
\end{proof}

\begin{bremark}[what breaks without the cover check]
\label{rem:sd-gap}
Everything: the rule without the cover is ``prove $\Psi$ wherever
convenient.'' The canonical mutant is the gap family
$\{[1, \tfrac{4}{3}],\ [\tfrac{3}{2}, 2]\}$ over the parent
$[1, 2]$ --- each branch sub-proof can be perfectly honest, and the goal
simply fails on $(\tfrac{4}{3}, \tfrac{3}{2})$, where no branch
constrains it. The greedy reach from $1$ arrives at $\tfrac{4}{3}$ and
stops: no remaining interval's left endpoint is $\leq \tfrac{4}{3}$,
so the check refuses. The checker's replayed battery keeps this mutant,
together with a point-supported branch (a degenerate interval that
cannot read), an extension-ray parent, and untied branch simplices, as
standing negative controls.
\end{bremark}

\begin{bremark}[the fragment boundary, and why refusal is the right
default]\label{rem:sd-fragment}
The checkable fragment is deliberate: 1-dimensional closed intervals,
chained. Extension-ray parents or branches, higher-dimensional regions,
and richer supports are refused --- not because the rule is false there,
but because no verified covering argument exists there yet, and the
fail-toward-refusal posture makes the unproven territory
unrepresentable. Box domains lose no expressive reach:
$n$-dimensional grids are expressible by nesting per-coordinate splits
across several applications, each within the fragment, the inner
split's parent membership entering through the outer split's branch
premise. The cover-certificate
wire slot is likewise part of the posture: recomputing the cover from
the branch regions means a forged certificate cannot influence the
verdict --- ignoring the slot cannot affect it either --- and a future
interpretation that additionally \emph{required} per-gap infeasibility
certificates from the slot would only tighten the check (one that
instead accepted new covers on certificate evidence would be a new
checking route, owing its own soundness argument).
\end{bremark}

\begin{bremark}[relation to the literature]
The rule is the reassembly half of branch-and-bound over Bernstein
enclosures: subdivision sharpens coefficient bounds quadratically in the
subdivision width and
drives them to the range in the limit \cite{garloff1986}, with the
subdivision algorithms themselves classical
\cite{decasteljau1959,laneriesenfeld1980}. In that literature the
splitting is a completeness device --- \emph{some} refinement certifies
any strict global fact; here only the soundness direction enters the
kernel: whatever refinement the producer chose, the chain check makes
its case analysis exhaustive. None of the producer-side effectivity
data (piece counts, degrees, tolerances) is consulted.
\end{bremark}

\begin{anchors}
Recipe \ref{rec:sd-instance} warrant & \texttt{sound\_subdivide}\\
interval reading of the parent and branches & \texttt{ExtensionSet.interval1D?}; elim/intro pair \texttt{interval1D\_elim}, \texttt{interval1D\_intro}\\
Proposition \ref{prop:sd-chain} & \texttt{coverExtendPass}, \texttt{coverReach}, \texttt{coverChainCheck}; invariant \texttt{ChainInv}, \texttt{chainInv\_extendPass}, \texttt{chainInv\_reach}; covering fact \texttt{coverChain\_covers}\\
the assembled check (ties included) & \texttt{coverCheck}, \texttt{coverInfeasibilityValid}\\
branch premise/conclusion shape & payload fields \texttt{branchPre}, \texttt{branchPost} (\texttt{PSubdivideData})\\
same-point re-entry (Prop.\ \ref{prop:subdiv-sound}) & the region read used by both parent elimination and branch introduction in \texttt{sound\_subdivide}'s body\\
wire placement & \texttt{RuleTag.subdivide}; the unconsulted certificate slot \texttt{coverCertificate}\\
\end{anchors}


\rulesection{19}{LiRPA}{admissible rule}
\label{blcatrule-LiRPA}
\ruleblock{The rule (Steele metanotation)}
\sourcenote{as realized in the mechanized development; source row: \texttt{\texttt{rules.toml}} --- \emph{LiRPA}. The metanotation statement is typeset by hand for notational quality; its mechanized ground truth is the Lean statement below, extracted from the sources at every build of this book and hash-bound. The rule \emph{roster} is machine-checked (every mechanized rule tag round-trips to exactly one catalogue row); the statement \emph{text} is editorial, checked by eye against the extracted Lean beside it.}
\begin{mathpar}
\inferrule*[Right=\rulename{LiRPA}]{
  \Phi \vdash a \leq h \\
  \Phi \vdash h \leq b \\
  \ReLU_{\{h\}\to\{y\}} \in H \\
  a < 0 \\
  0 < b \\\\
  \alpha a + \beta \leq 0 \\
  \beta \leq 0 \\
  \alpha b + \beta \leq b \\\\
  0 \leq \gamma a + \delta \\
  0 \leq \delta \\
  b \leq \gamma b + \delta
}{
  \Phi \vdash (\alpha h + \beta) \leq y \;\wedge\; y \leq (\gamma h + \delta)
}
\end{mathpar}

\par\medskip
\rulename{LiRPA} is the calculus\textquotesingle s neural-network
bound-propagation rule — the certificate form of the linear relaxation
of the ReLU activation that bound-propagation verifiers compute (the
CROWN/LiRPA family, the triangle relaxation, DeepPoly). Given
already-proven box bounds $\Phi \vdash a \leq h$ and $\Phi \vdash h
\leq b$ on the preactivation wire of a cited ReLU application, with
$a < 0 < b$ (the box straddles the kink — the \emph{unstable} case, in
the verification literature\textquotesingle s vocabulary), it concludes
the two-sided affine envelope $\alpha h + \beta \leq y \;\wedge\; y
\leq \gamma h + \delta$ on the application' s output
wire $y$. The producer finds the envelope however it likes —
CROWN\textquotesingle s adaptive slopes, the triangle
relaxation\textquotesingle s chord, any heuristic — and ships only the
box endpoints and the four coefficients; no relaxation machinery is
re-run or trusted.

What the checker recomputes, given the payload (two premise proofs, the
two wires, six rationals): the \emph{linkage} — the hypernet really
contains a packaged ReLU application at exactly these ports, of
positive total size (a one-line decidable port lookup); that the
premises are literally the atoms $a \leq h$ and $h \leq b$ over one
shared premise formula; the sign conditions $a < 0$ and $0 < b$; the
six envelope inequalities, exactly over $\Q$ — the lower line under the
function at the three checkpoints $a$, $0$, $b$ ($\alpha a + \beta \leq
0$, $\beta \leq 0$, $\alpha b + \beta \leq b$) and the upper line above
it at the same three ($0 \leq \gamma a + \delta$, $0 \leq \delta$, $b
\leq \gamma b + \delta$); and that the conclusion is \emph{exactly} the
pinned conjunction of the two affine comparison atoms, built by the
checker from the coefficients — the producer proposes no conclusion
shape at all.

Each hypothesis but one is load-bearing. Drop the linkage and $y$ is an
arbitrary wire with no semantic tie to $h$ (the rule as first staged
had exactly this gap and was kept vacuous — rejecting every certificate
— until the port-lookup premise landed). Drop the lower kink check
($\beta \leq 0$) and the chord through the box endpoints, which
passes both endpoint checks, would certify as a lower envelope while
lying strictly above the function at $h = 0$. (The upper kink check
$0 \leq \delta$ is the one redundancy: under the strict straddle the
two upper endpoint checks already force $\delta \geq -ab/(b-a) > 0$ —
evaluate the affine interpolation of the endpoint bounds at $0$ — so
it rides along because the certificate's check set is the uniform
six.) Drop an endpoint check
and a steep line passing the remaining two checkpoints escapes the
function at the dropped end (the proof section constructs both
failures). The box premises themselves are what make an affine
\emph{upper} envelope possible at all — no affine function dominates
$\max(0, h)$ on all of $\Q$ — and what license lower-line slopes
outside $[0, 1]$.

\cataloguenote{The family\textquotesingle s shared machinery — the packaged-node
linkage, the elementwise two-case semantic contract, and the
conservativity warrant behind it — is stated once, in this
section\textquotesingle s recipe preamble, and referenced by the four
sibling rules (\rulename{ReLU-Pos}, \rulename{ReLU-Neg},
\rulename{ReLU-NN}, \rulename{ReLU-Above}). One envelope serves the
whole node: the rule is elementwise over the cited
application\textquotesingle s tensor shape, with a single $(a, b,
\alpha, \beta, \gamma, \delta)$ for every entry — per-entry envelopes
are expressed with singleton-shaped applications. The payload
additionally declares a denominator bound on its witness rationals (a
size claim for resource accounting; it carries no proof-theoretic
weight in the soundness argument). The rule\textquotesingle s name is
the literature\textquotesingle s: linear-relaxation-based perturbation
analysis~\cite{xu2020lirpa}.
}
\ruleblock{Lean statement}
\par\medskip
\textbf{The linkage and the pinned conclusion, from the checker
sources.}\enspace \texttt{reluLinked} is the whole of the
family\textquotesingle s side condition: the hypernet\textquotesingle s
application list contains a packaged ReLU application at exactly the
cited preactivation/output ports, of positive total size — one
decidable lookup (\texttt{PrimApp.isReluAt}), no pattern-matching over
composites. The positive-size gate refuses vacuous surface at the
source: a size-zero node\textquotesingle s comparison atoms hold
vacuously, and a fact no valuation can falsify must not be citable
(the same lesson the \rulename{Farkas} rule\textquotesingle s
zero-size counterexample taught that payload).
\texttt{affineWireTerm} and \texttt{lirpaConclusionForm} pin the
conclusion: the rule concludes exactly the conjunction of the two
affine comparison atoms, built from the payload coefficients — there
is no producer-supplied conclusion shape to audit.

\declhead{BernsteinLogic.Core.Catalogue}{PrimApp.isReluAt}{\badge{BadgeGreen}{definition bound}}
\begin{leansrc}
/-- Port-lookup matcher for the ReLU linkage: a packaged ReLU
application at exactly these ports, of positive total size. The size
gate is the shape-aware rider (conservative-extension memo register 2,
the `PFarkasData.wiresPos` lesson generalized): a shape-[0] node's
atoms hold vacuously, and vacuous surface must not be citable — the
rules refuse it at the linkage, before any soundness question arises. -/
def PrimApp.isReluAt (preact post : Wire) : PrimApp → Bool
  | .relu s i o =>
      decide (i = preact) && decide (o = post) &&
      decide (0 < Shape.totalSize s)
  | _ => false
\end{leansrc}
\declhead{BernsteinLogic.Core.Catalogue}{reluLinked}{\badge{BadgeGreen}{definition bound}}
\begin{leansrc}
def reluLinked (H : Hypernet) (preact post : Wire) : Prop :=
  H.applications.any (PrimApp.isReluAt preact post) = true
\end{leansrc}
\declhead{BernsteinLogic.Core.Catalogue}{affineWireTerm}{\badge{BadgeGreen}{definition bound}}
\begin{leansrc}
def affineWireTerm (offset slope : Rat) (w : Wire) : Term :=
  .add (.const offset) (.mul (.const slope) (.var w))
\end{leansrc}
\declhead{BernsteinLogic.Core.Catalogue}{lirpaConclusionForm}{\badge{BadgeGreen}{definition bound}}
\begin{leansrc}
def lirpaConclusionForm
    (preact post : Wire) (α β γ δ : Rat) (F : Formula) : Prop :=
  F = .and (.atom .le (affineWireTerm β α preact) (.var post))
           (.atom .le (.var post) (affineWireTerm δ γ preact))
\end{leansrc}
\par\medskip
\textbf{The verified chain.}\enspace \texttt{reluLinked\_facts} turns
the linkage into semantics: under any admissible valuation the two
cited wires carry tensors of one shared shape and positive entry
count, related entrywise by the two-case ReLU relation. The four
interpolation lemmas are the checkpoint arithmetic — an affine
function suitably bounded at an interval\textquotesingle s endpoints
is bounded on the interval, one lemma per envelope side per sign
branch — and \texttt{sound\_lirpa} assembles them entrywise into the
rule\textquotesingle s warrant.

\declhead{BernsteinLogic.Soundness}{reluLinked\_facts}{\badge{BadgeGreen}{statement bound}}
\begin{leansrc}
/-- Extraction for the realized B12 linkage: under any admissible
valuation, a cited ReLU application yields both wire tensors — one
shared shape, positive entry count — and the elementwise two-case
relation. Typing rides `H.typeConsistent` at the cited application (no
payload typing fields needed — the wire format is untouched). -/
theorem reluLinked_facts {H : Hypernet} {preact post : Wire}
    (hlink : reluLinked H preact post) (σ : H.Valuation)
    (hadm : σ.Admissible) :
    ∃ qp qy : QTensor,
      σ.wireQ preact = some qp ∧ σ.wireQ post = some qy ∧
      qp.shape = qy.shape ∧ 0 < qp.entries.length ∧
      qp.entries.length = qy.entries.length ∧
      ∀ i, i < qp.entries.length →
        (qp.entries.getD i 0 ≤ 0 ∧ qy.entries.getD i 0 = 0) ∨
        (0 ≤ qp.entries.getD i 0 ∧
          qy.entries.getD i 0 = qp.entries.getD i 0)
\end{leansrc}
\declhead{BernsteinLogic.Soundness}{lirpa\_lower\_neg}{\badge{BadgeGreen}{statement bound}}
\begin{leansrc}
theorem lirpa_lower_neg {a α β p : ℚ} (ha : a < 0)
    (h1 : α * a + β ≤ 0) (h2 : β ≤ 0) (hap : a ≤ p) (hp0 : p ≤ 0) :
    α * p + β ≤ 0
\end{leansrc}
\declhead{BernsteinLogic.Soundness}{lirpa\_lower\_pos}{\badge{BadgeGreen}{statement bound}}
\begin{leansrc}
theorem lirpa_lower_pos {b α β p : ℚ} (hb : 0 < b)
    (h2 : β ≤ 0) (h3 : α * b + β ≤ b) (hp0 : 0 ≤ p) (hpb : p ≤ b) :
    α * p + β ≤ p
\end{leansrc}
\declhead{BernsteinLogic.Soundness}{lirpa\_upper\_neg}{\badge{BadgeGreen}{statement bound}}
\begin{leansrc}
theorem lirpa_upper_neg {a γ δ p : ℚ} (ha : a < 0)
    (h1 : 0 ≤ γ * a + δ) (h2 : 0 ≤ δ) (hap : a ≤ p) (hp0 : p ≤ 0) :
    0 ≤ γ * p + δ
\end{leansrc}
\declhead{BernsteinLogic.Soundness}{lirpa\_upper\_pos}{\badge{BadgeGreen}{statement bound}}
\begin{leansrc}
theorem lirpa_upper_pos {b γ δ p : ℚ} (hb : 0 < b)
    (h2 : 0 ≤ δ) (h3 : b ≤ γ * b + δ) (hp0 : 0 ≤ p) (hpb : p ≤ b) :
    p ≤ γ * p + δ
\end{leansrc}
\declhead{BernsteinLogic.Soundness}{sound\_lirpa}{\badge{BadgeGreen}{statement bound}}
\begin{leansrc}
/-- LiRPA: the packaged degree-1 relaxation — from box premises
`a ≤ h ≤ b` and the eight endpoint side conditions, the affine sandwich
`αh + β ≤ y ∧ y ≤ γh + δ` holds entrywise, by the four interpolation
lemmas applied per two-case branch. -/
theorem sound_lirpa {H : Hypernet}
    {seq lowerPremise upperPremise : Sequent}
    {preact post : Wire} {a b α β γ δ : ℚ}
    (ihl : lowerPremise.semanticallyValid H)
    (ihu : upperPremise.semanticallyValid H)
    (hlink : reluLinked H preact post)
    (hAneg : a < 0) (hBpos : 0 < b)
    (hLower1 : α * a + β ≤ 0) (hLower2 : β ≤ 0) (hLower3 : α * b + β ≤ b)
    (hUpper1 : 0 ≤ γ * a + δ) (hUpper2 : 0 ≤ δ) (hUpper3 : b ≤ γ * b + δ)
    (hpre : seq.pre = lowerPremise.pre)
    (hpreMatch : lowerPremise.pre = upperPremise.pre)
    (hlowerPost : lowerPremise.post = .atom .le (.const a) (.var preact))
    (hupperPost : upperPremise.post = .atom .le (.var preact) (.const b))
    (hpostForm : lirpaConclusionForm preact post α β γ δ seq.post) :
    seq.semanticallyValid H
\end{leansrc}
\par\medskip
\textbf{Why consuming the packaged two-case fact is not new kernel
content.}\enspace The two-case relation is the packaged
symbol\textquotesingle s defined semantics, consumed by the family as
model fact at every positive size; what this pair of declarations adds
is the conservativity warrant --- that the definition mints no
genuinely new primitive: \texttt{reluExpH} is the definitional
expansion (a copy, a sign-split cell complex, an indicator cast, and
an exact rank-0 contraction), and \texttt{relu\_projection} shows —
both directions — that the two-case relation is exactly the
expansion\textquotesingle s model-level facts existentially projected
onto the external ports, proven at the canonical single-entry
instance. The boundary point $h = 0$ rides the one-hot cell semantics:
both cells admit it, and every admissible choice forces $y = 0$.
Applications of total size other than one currently have no core
expansion (a recorded packing gap in the development\textquotesingle s
notes): at those sizes the node stands as a genuine primitive with the
stated semantics — the rules remain sound there, while the
not-a-new-primitive claim is certified only at size one.

\declhead{BernsteinLogic.ReluExpansion}{relu\_projection}{\badge{BadgeGreen}{statement bound}}
\begin{leansrc}
/-- **Obligation P (the projection lemma), both directions**: the
packaged ReLU's derived sp-fact — the two-case relation — EQUALS the
expansion's sp-facts existentially projected onto the external ports.
The internal wires (the copies, the indicators, the cast image, the
rank-0 legs) are exactly the existentially-quantified content of the
valuation; the boundary case p = 0 is covered because both one-hot
choices are admissible there and both force y = 0. -/
theorem relu_projection (p y : ℚ) :
    (∃ σ : reluExpH.Valuation, σ.Admissible ∧
        σ.entryFlat 0 0 = some (.fin p) ∧
        σ.entryFlat 8 0 = some (.fin y)) ↔
    ((p ≤ 0 ∧ y = 0) ∨ (0 ≤ p ∧ y = p))
\end{leansrc}
\declhead{BernsteinLogic.ReluExpansion}{reluExpH}{\badge{BadgeGreen}{definition bound}}
\begin{leansrc}
/-- The expansion: copy;(id ⊗ Cell[Σ±]);Contract, with the typing
plumbing (Cast 𝔹→ℚ on the indicator, reshapes on the rank-1 ports). -/
def reluExpH : Hypernet where
  wireType := expTy
  applications :=
    [.copy 2 0 [1, 2],
     .cell sigmaPM 2 [3, 4],
     .cast Dtype.B Dtype.Q .exact [] 4 5,
     .reshape [] 1 6,
     .contract Dtype.Q .exact [6, 5] 7 [0, 0] 0 pat0,
     .reshape [1] 7 8]
  linearity := by decide
  acyclic := by decide
  typeConsistent := by decide +kernel
\end{leansrc}
\ruleblock{Correspondence (Steele display $\leftrightarrow$ Lean)}
\badge{BadgeGreen}{reviewer-certified against the displayed texts}
\sourcenote{The dictionary below is a reviewer-certified statement of how the mechanized statements above interpret the displayed rule: display metavariable $\mapsto$ Lean binder, each display hypothesis $\mapsto$ the premise or side condition carrying it, the display conclusion $\mapsto$ the goal --- with any deliberate difference in width between display and mechanization named outright. Certification is recorded as a hash triple over (the displayed cell, the extracted statement texts, this dictionary\textquotesingle s own text); if any of the three changes, this block turns amber until a reviewer re-certifies. The correspondence claim itself remains editorial --- reviewer-grade, not machine-checked; its machine-checked successor types each soundness statement against the rule as a structured object, and supersedes this block per rule as it lands.}
Judgment frame as in \rulename{Ax}; sub-derivation premises as at
\rulename{Weaken}. This dictionary doubles as the family preamble: its
linkage row is shared verbatim by the four sibling rules.
\begin{itemize}
\item $h \mapsto$ \texttt{preact : Wire}; $y \mapsto$ \texttt{post :
  Wire}; $a, b, \alpha, \beta, \gamma, \delta \mapsto$ the six
  rational binders of the same names in \texttt{sound\_lirpa}.
\item the linkage premise $\ReLU_{\{h\}\to\{y\}} \in H$ $\mapsto$
  \texttt{hlink} $:$ \texttt{reluLinked H preact post} ---
  membership \emph{plus} the positive-size gate: unfolding (the
  definition above), the certified premise demands a packaged ReLU
  application at exactly these ports \emph{of positive total size}.
  The display\textquotesingle s bare $\in$ shows no size condition;
  the certified linkage is strictly narrower, by design (vacuous
  surface refused at the source --- the intro above), and the
  narrowing is family-wide.
\item box premises $\Phi \vdash a \leq h$ and $\Phi \vdash h \leq b$
  $\mapsto$ \texttt{ihl}, \texttt{ihu} with posts pinned by
  \texttt{hlowerPost} $:$ \texttt{lowerPremise.post = .atom .le
  (.const a) (.var preact)} and \texttt{hupperPost} $:$
  \texttt{upperPremise.post = .atom .le (.var preact) (.const b)} ---
  displayed comparators and operand orders carried exactly; $\Phi$
  shared three ways via \texttt{hpre} and \texttt{hpreMatch}.
\item the eight arithmetic side conditions $\mapsto$ \texttt{hAneg},
  \texttt{hBpos}, \texttt{hLower1}--\texttt{hLower3},
  \texttt{hUpper1}--\texttt{hUpper3}, verbatim over $\Q$ --- scalar
  hypotheses on the payload rationals, not judgments.
\item conclusion $\Phi \vdash (\alpha h + \beta) \leq y \wedge y \leq
  (\gamma h + \delta)$ $\mapsto$ the goal under \texttt{hpostForm}
  $:$ \texttt{lirpaConclusionForm preact post} applied to the four
  coefficients and \texttt{seq.post}, pinning \texttt{seq.post} to
  the conjunction of the two affine atoms in the displayed order
  (lower envelope left). Spelling drift, disclosed: the pinned
  affine term is \texttt{affineWireTerm}\textquotesingle s
  offset-first $\beta + \alpha h$, where the display writes $\alpha h
  + \beta$ --- semantically identical, but formulas are ordered
  trees, so the certified rule accepts exactly the offset-first
  spelling (the \rulename{ConvElim}/\rulename{ConvIntro}
  build-output class).
\end{itemize}
Width notes beyond the linkage gate: the certified atoms are
elementwise over the cited application\textquotesingle s tensor shape
--- one envelope for every entry, per the note above --- where the
display reads scalar; that is the batch semantics of the atom
grammar, not a strengthening. \texttt{reluLinked\_facts}, the four
interpolation lemmas, \texttt{relu\_projection}, and
\texttt{reluExpH} are certified alongside as machinery: the first
extracts the elementwise two-case relation the family consumes, the
last two ground the packaged symbol in the core expansion --- they
are \emph{about} the linkage, in the satellite sense of
\rulename{OrE}\textquotesingle s ctx-rule. No hypothesis of
\texttt{sound\_lirpa} is unused, and nothing certified is wider than
the display.

\ruleblock{Lean verifier receipts}
\sourcenote{module \texttt{BernsteinLogic.Core.Catalogue}; bound build \texttt{cbc92e976a1d\ldots} (bound from receipt 5abdb3589707b5df, 2026-08-03T21:30:33Z). A receipt records that named checking layers accepted a build: (module, compiled-artifact SHA-256, environment SHA-256, checker layers), minted when the checking gate ran; the store receipts attesting the bound build are listed with the layers each run credited. Nothing here asks the reader to trust the store: every layer is an independently runnable program over the quoted sources --- rebuild the module and re-run the checkers to reproduce the verdict. Receipts attesting superseded builds credit no layer and are omitted from this rendering (the count of omissions, if any, is noted below).}
\begin{receipttable}
\receiptrow{24804e66edba0bc9}{2026-08-03T20:25:09Z}{build+decorrelated}{\badge{BadgeGreen}{current}}
\receiptrow{5abdb3589707b5df}{2026-08-03T21:30:33Z}{build+decorrelated}{\badge{BadgeGreen}{current}}
\end{receipttable}
\sourcenote{18 further receipt(s) in the store attest superseded builds of this module and are omitted from this rendering.}
\sourcenote{module \texttt{BernsteinLogic.Soundness}; bound build \texttt{a20d90e66c13\ldots} (bound from receipt a9c0e34d3e3c9c8c, 2026-08-03T21:30:36Z). A receipt records that named checking layers accepted a build: (module, compiled-artifact SHA-256, environment SHA-256, checker layers), minted when the checking gate ran; the store receipts attesting the bound build are listed with the layers each run credited. Nothing here asks the reader to trust the store: every layer is an independently runnable program over the quoted sources --- rebuild the module and re-run the checkers to reproduce the verdict. Receipts attesting superseded builds credit no layer and are omitted from this rendering (the count of omissions, if any, is noted below).}
\begin{receipttable}
\receiptrow{a9c0e34d3e3c9c8c}{2026-08-03T21:30:36Z}{build+decorrelated}{\badge{BadgeGreen}{current}}
\end{receipttable}
\sourcenote{21 further receipt(s) in the store attest superseded builds of this module and are omitted from this rendering.}
\sourcenote{module \texttt{BernsteinLogic.ReluExpansion}; bound build \texttt{68395e3f44af\ldots} (bound from receipt b1a5795a8e57ce30, 2026-07-18T00:07:06Z). A receipt records that named checking layers accepted a build: (module, compiled-artifact SHA-256, environment SHA-256, checker layers), minted when the checking gate ran; the store receipts attesting the bound build are listed with the layers each run credited. Nothing here asks the reader to trust the store: every layer is an independently runnable program over the quoted sources --- rebuild the module and re-run the checkers to reproduce the verdict. Receipts attesting superseded builds credit no layer and are omitted from this rendering (the count of omissions, if any, is noted below).}
\begin{receipttable}
\receiptrow{7d0d53fd422f509e}{2026-07-17T20:41:28Z}{build+replay+hardened+decorrelated}{\badge{BadgeGreen}{current}}
\receiptrow{3219a0974f9b52a3}{2026-07-17T21:00:42Z}{build+replay+hardened+decorrelated}{\badge{BadgeGreen}{current}}
\receiptrow{b1a5795a8e57ce30}{2026-07-18T00:07:06Z}{build+decorrelated}{\badge{BadgeGreen}{current}}
\end{receipttable}
\sourcenote{4 further receipt(s) in the store attest superseded builds of this module and are omitted from this rendering.}
\layerline{\layeritem{build}{lake build (elaboration)}{\badge{BadgeGreen}{current}}\layeritem{replay}{leanchecker (official kernel replay)}{\badge{BadgeGreen}{current}}\layeritem{hardened}{safe\_verify (hardened screen)}{\badge{BadgeGreen}{current}}\layeritem{decorrelated}{lean4export + nanoda (decorrelated kernel)}{\badge{BadgeGreen}{current}}}
\sourcenote{\badge{BadgeGray}{comparator: none} no fixture in the three-implementation comparator harness (compiled / kernel-interpreted / reference-interpreted) exercises this rule.}
\ruleblock{HOL4 statement}
\sourcenote{theory \texttt{bern\_tr\_catdata2}, HOL4 development script \texttt{\texttt{bern\_tr\_catdata2Script.sml}} \badge{BadgeGreen}{script bound}}
\holhead{Theorem plirpadata\_receipt}
\begin{holsrc}
Theorem plirpadata_receipt:
  ruledataOk [] cd2Hrelu
    <| pre := FAnd (FAtom BowLe (TmConst (-&1)) (TmVar 3))
                   (FAtom BowLe (TmVar 3) (TmConst (&1))) ;
       post := FAnd (FAtom BowLe (affineWireTerm 0 0 3) (TmVar 4))
                    (FAtom BowLe (TmVar 4) (affineWireTerm 1 0 3)) |>
    (RdLirpa (PLirpaD
       <| pre := FAnd (FAtom BowLe (TmConst (-&1)) (TmVar 3))
                      (FAtom BowLe (TmVar 3) (TmConst (&1))) ;
          post := FAtom BowLe (TmConst (-&1)) (TmVar 3) |>
       (ProofObj (RdAndE (PAndEDL
          <| pre := FAnd (FAtom BowLe (TmConst (-&1)) (TmVar 3))
                         (FAtom BowLe (TmVar 3) (TmConst (&1))) ;
             post := FAnd (FAtom BowLe (TmConst (-&1)) (TmVar 3))
                          (FAtom BowLe (TmVar 3) (TmConst (&1))) |>
          (ProofObj RdAx)
          (FAtom BowLe (TmVar 3) (TmConst (&1))))))
       <| pre := FAnd (FAtom BowLe (TmConst (-&1)) (TmVar 3))
                      (FAtom BowLe (TmVar 3) (TmConst (&1))) ;
          post := FAtom BowLe (TmVar 3) (TmConst (&1)) |>
       (ProofObj (RdAndE (PAndEDR
          <| pre := FAnd (FAtom BowLe (TmConst (-&1)) (TmVar 3))
                         (FAtom BowLe (TmVar 3) (TmConst (&1))) ;
             post := FAnd (FAtom BowLe (TmConst (-&1)) (TmVar 3))
                          (FAtom BowLe (TmVar 3) (TmConst (&1))) |>
          (ProofObj RdAx)
          (FAtom BowLe (TmConst (-&1)) (TmVar 3)))))
       3 4 (-&1) (&1) 0 0 0 (&1) 0)) ∧
  ¬ ruledataOk [] hnEmpty
      <| pre := FAnd (FAtom BowLe (TmConst (-&1)) (TmVar 3))
                     (FAtom BowLe (TmVar 3) (TmConst (&1))) ;
         post := FAnd (FAtom BowLe (affineWireTerm 0 0 3) (TmVar 4))
                      (FAtom BowLe (TmVar 4) (affineWireTerm 1 0 3)) |>
      (RdLirpa (PLirpaD
         <| pre := FTt ; post := FTt |> (ProofObj RdAx)
         <| pre := FTt ; post := FTt |> (ProofObj RdAx)
         3 4 (-&1) (&1) 0 0 0 (&1) 0))
\end{holsrc}
\ruleblock{HOL verifier receipts}
\badge{BadgeBlue}{documentary} \texttt{hol-strict} CLEAN: theory \texttt{bern\_tr\_catdata2}, 272
theorems, 0 cheat-tagged, 0 nonstandard axioms (whitelist:
\texttt{DISK\_THM}; the four \texttt{bool} axioms) --- verification
run of 2026-07-22 on an isolated build host, against a
content-addressed checkout of the exact sources cited here (git blob
\texttt{cd3c1e95\ldots}). The proof-object stratum is ported as one
mutual datatype family with a mutual wellformedness predicate,
exactly as the Lean mutual inductive block; this
rule\textquotesingle s arm carries the two premise sub-proofs, the
six sandwich inequalities verbatim, the application-list linkage
(\texttt{reluLinkedB}, theory \texttt{bern\_tr\_catdata}, 252
theorems, 0 offences, same run --- the lookup the earlier port
excluded wholesale now has a ported realization), and the conclusion
form (\texttt{lirpaConclusionFormB}, theory
\texttt{bern\_tr\_catforms}, 23 theorems, 0 offences). The receipt
theorem cited here is the kernel-evaluated accept/reject pair over a
hypernet carrying the packaged ReLU application. The four bound
lemmas behind the sandwich soundness are ported as well
(\texttt{lirpa\_lower\_neg} / \texttt{lirpa\_lower\_pos} /
\texttt{lirpa\_upper\_neg} / \texttt{lirpa\_upper\_pos}, theory
\texttt{bern\_tr\_soundmid}, 16 theorems, 0 offences, run of
2026-07-22). Honest residue: the apex semantic theorem assembling
them (the Lean development\textquotesingle s \texttt{sound\_lirpa})
is not yet ported; its ingredients above are.

\sourcenote{HOL receipts are documentary: quoted results of \texttt{hol-strict} runs
and census rows from the HOL4 development\textquotesingle s build
evidence, not minted receipt artefacts. A HOL receipt store (the Lean
store\textquotesingle s twin) is named future work.
}
\ruleblock{Proof}

\sourcenote{Genre: admissible rule (tier 2) --- a proof-tree-rewriting
recipe. This section additionally carries the \emph{family template}
for the five ReLU rules (\rulename{LiRPA}, \rulename{ReLU-Pos},
\rulename{ReLU-Neg}, \rulename{ReLU-NN}, \rulename{ReLU-Above}): the
packaged-node machinery they share is stated once here, and the four
sibling sections instantiate it by reference.}

\begin{brecipe}[the ReLU family template]\label{rec:relu-template}
The general admissible template --- one rule node rewrites to a
derivation fragment over core and previously-justified rules with the
same conclusion --- is stated in full in the \rulename{PConstExact}
section; the five ReLU rules are its instances over one shared
semantic fact. The realized term grammar has no $\ReLU$ symbol:
$\ReLU$ is a \emph{packaged hypernet primitive}, an application
$\ReLU_{\{h\}\to\{y\}} \in H$ with preactivation (input) wire $h$ and
output wire $y$. Every family rule cites it through one side
condition, the \emph{linkage} $\mathrm{reluLinked}(H, h, y)$: the
hypernet's application list contains a packaged ReLU at exactly these
ports, of positive total size --- a single decidable port lookup, not a
pattern-match over composites. The packaged node's semantic contract
is the elementwise two-case fact: under every admissible valuation,
at every entry $i$ of the node's (shared) tensor shape,
\[
\bigl(h_i \leq 0 \;\wedge\; y_i = 0\bigr)
\;\vee\;
\bigl(0 \leq h_i \;\wedge\; y_i = h_i\bigr).
\]
Three remarks fix what the template means and what it does not.

\emph{(i) The rewrite shape.} Each family node rewrites to a case
analysis over this disjunction: the sign-split fact enters the branch
contexts (context-carrying $\vee$E --- the disjunction arrives as a
sub-conjunct of the premise; see the $\vee$E section's note), equality
transport (\rulename{EqElim}) carries the branch equality $y = 0$ or
$y = h$ through the goal, and \rulename{ArithAtom}/\rulename{Norm}
plumbing normalizes atom orientations. Where a branch discharges fully
into catalogue nodes the instance recipe displays the derivation
(\rulename{ReLU-NN} and \rulename{ReLU-Above}); where it does not ---
the boundary branches of \rulename{ReLU-Pos}/\rulename{ReLU-Neg} need
an antisymmetry step no catalogue rule provides --- the instance says
so and the rule is admissible by \emph{reduction to verified
machinery}, the honest tier note the \rulename{PConstExact} section
established.

\emph{(ii) The semantic contract and the conservativity warrant, kept
distinct.} The two-case relation is the packaged symbol's
\emph{defined} semantics --- part of what an admissible valuation of a
hypernet containing the node means, at every entry of every
positive-size application --- and the family's soundness lemmas consume
it as model fact at any size (that is (iii)'s extraction lemma). What
the conservativity warrant adds is that the definition mints no
genuinely new primitive: a definitional expansion into core primitives
(a copy, the sign-split cell complex $\Sigma^{\pm}$, an indicator
cast, and an exact rank-0 contraction --- \texttt{reluExpH}), a typing
rule derivable from the expansion's, and the projection theorem
\texttt{relu\_projection}: \emph{both directions}, the two-case
relation is exactly the expansion's model-level facts existentially
projected onto the external ports --- proven at the canonical
single-entry instance. The boundary point $h = 0$ rides the one-hot
cell semantics: both cells of $\Sigma^{\pm}$ admit it, and every
admissible choice forces $y = 0$. Scope, stated loudly: applications
of total size other than one currently have \emph{no} core expansion
(a recorded packing gap), so at those sizes the packaged node stands
as a genuine primitive with the stated semantics --- the five rules
remain sound there, consuming the defined relation, while the
not-a-new-primitive claim is certified only at size one.

\emph{(iii) The shipped verified objects.} Per the catalogue's
convention, each family rule ships a direct soundness lemma factoring
through already-verified machinery rather than a derivation object:
\texttt{reluLinked\_facts} converts the linkage, under any admissible
valuation, into (a) tensors at both ports, (b) one shared shape of
positive entry count, and (c) the per-entry two-case relation; each
rule's lemma is then an entrywise case analysis. The positive-size
gate inside the linkage is deliberate hygiene: a size-zero node's
comparison atoms hold vacuously, and vacuous surface must not be
citable --- the general lesson of the \rulename{Farkas} rule's
zero-size counterexample, enforced here at the source.
\end{brecipe}

\begin{brecipe}[the \rulename{LiRPA} instance --- subdivide at the kink]
\label{rec:lirpa-instance}
The node, with its premise subtrees and side conditions:
\[
\inferrule*[Right=\rulename{LiRPA}]{
  \mathcal{T}_a \rhd \Phi \vdash a \leq h \\
  \mathcal{T}_b \rhd \Phi \vdash h \leq b \\
  \ReLU_{\{h\}\to\{y\}} \in H \\
  a < 0 \\ 0 < b \\\\
  \alpha a + \beta \leq 0 \\ \beta \leq 0 \\ \alpha b + \beta \leq b \\
  0 \leq \gamma a + \delta \\ 0 \leq \delta \\ b \leq \gamma b + \delta
}{
  \Phi \vdash (\alpha h + \beta) \leq y \;\wedge\; y \leq (\gamma h + \delta)
}
\]
rewrites to a two-branch case analysis at the kink,
\[
\inferrule*[Right=cases]{
  \mathcal{D}_{-} \rhd \Phi \wedge (h \leq 0) \vdash
    (\alpha h + \beta) \leq y \,\wedge\, y \leq (\gamma h + \delta) \\
  \mathcal{D}_{+} \rhd \Phi \wedge (0 \leq h) \vdash
    (\alpha h + \beta) \leq y \,\wedge\, y \leq (\gamma h + \delta)
}{
  \Phi \vdash (\alpha h + \beta) \leq y \;\wedge\; y \leq (\gamma h + \delta)
}
\]
where the case split is the box subdivision $[a, b] = [a, 0] \cup
[0, b]$ (region atoms enter from the order atoms via
\rulename{ConvIntro}; the cover is the one \rulename{Subdivide}
checks), and the branches discharge as follows. In $\mathcal{D}_{-}$,
the branch context proves $h \leq 0$, so \rulename{ReLU-Neg} yields
$y = 0$ (straight from the packaged contract: on $h \leq 0$ both
disjuncts force $y = 0$); the two envelope atoms then reduce, through \rulename{EqElim}
along $y = 0$, to degree-$1$ bound facts on $[a, 0]$ --- $\alpha h +
\beta \leq 0$ and $0 \leq \gamma h + \delta$ --- whose degree-$1$
Bernstein certificates on that cell are \emph{exactly} the checkpoint
conditions at $a$ and at $0$ (\rulename{BernCone} at $d = 1$; on a
$1$-simplex the degree-$1$ Bernstein coefficients of an affine gap
polynomial are its two endpoint values). Dually in $\mathcal{D}_{+}$:
\rulename{ReLU-Pos} yields $y = h$ (from the contract's active
disjunct; at the boundary $h = 0$ the inactive disjunct gives the same
$y = 0 = h$), and the envelope atoms reduce to
degree-$1$ facts on $[0, b]$ --- $\alpha h + \beta \leq h$ and $h \leq
\gamma h + \delta$ --- whose Bernstein certificates are the checkpoint
conditions at $0$ and at $b$. The six side conditions are thus the six
distinct endpoint values of the four affine gap polynomials, the kink
checkpoints $\beta \leq 0$ and $0 \leq \delta$ each shared between
their line's two cells. \rulename{ArithAtom}, \rulename{Norm}, and
\rulename{Cut} supply the atom-normalization and stitching plumbing
throughout (the \rulename{ReLU-Above} section walks a chain of this
kind in full). Scope of the rewrite, matching template remark (ii):
this two-branch derivation is the canonical \emph{single-entry} story
--- at larger tensor shapes the sign split is per-entry, so the
displayed branch contexts $h \le 0$ and $0 \le h$ do not jointly cover
a mixed-sign tensor and the derivation does not lift; the rule's
warrant at every positive size is the direct semantic lemma
(Proposition~\ref{prop:lirpa-sound}), per the family convention of
remark (iii).
\end{brecipe}

\begin{bprop}[semantic warrant]\label{prop:lirpa-sound}
Suppose the linkage holds, the two premise sequents are semantically
valid with a common premise formula, and the eight arithmetic side
conditions hold (the strict straddle pair $a < 0$, $0 < b$, plus the
six envelope checkpoints --- ``the six side conditions'' elsewhere in
this section always means the envelope six). Then the conclusion sequent --- the pinned conjunction
$(\alpha h + \beta) \leq y \wedge y \leq (\gamma h + \delta)$ over the
same premise --- is semantically valid.
\end{bprop}
\begin{proof}
Fix an admissible valuation $\sigma$ satisfying $\Phi$. The premises,
instantiated at $\sigma$, eliminate to entrywise bounds: every entry
$p$ of the preactivation tensor satisfies $a \leq p \leq b$ (a
rank-$0$ constant against a tensor compares against every entry). The
linkage extraction (\texttt{reluLinked\_facts}) gives the output
tensor at the same shape and, per entry, either $p \leq 0 \wedge y =
0$ or $0 \leq p \wedge y = p$. Fix an entry and take the cases.

\emph{Inactive case} ($p \leq 0$, $y = 0$; here $p \in [a, 0]$).
Lower line: $\alpha p + \beta \leq 0 = y$ by two-endpoint
interpolation on $[a, 0]$ --- the affine function $\alpha({\cdot}) +
\beta$ is $\leq 0$ at $a$ (checkpoint $\alpha a + \beta \leq 0$) and
at $0$ (checkpoint $\beta \leq 0$), hence on all of $[a, 0]$
(\texttt{lirpa\_lower\_neg}; the verified proof is an explicit product
certificate riding the identity
$(-a)(\alpha p + \beta) + (-p)\bigl(-(\alpha a + \beta)\bigr)
+ (p - a)(-\beta) = 0$: on $[a, 0]$ the last two products are
nonnegative, so the assumed violation $\alpha p + \beta > 0$ --- which,
with $-a > 0$, makes the first term strictly positive --- is impossible
in a zero sum). Upper
line: $y = 0 \leq \gamma p + \delta$ likewise from the checkpoints at
$a$ and $0$ (\texttt{lirpa\_upper\_neg}).

\emph{Active case} ($0 \leq p$, $y = p$; here $p \in [0, b]$). Lower
line: $\alpha p + \beta \leq p = y$ by interpolation on $[0, b]$ from
the checkpoints $\beta \leq 0$ and $\alpha b + \beta \leq b$
(\texttt{lirpa\_lower\_pos}). Upper line: $y = p \leq \gamma p +
\delta$ from $0 \leq \delta$ and $b \leq \gamma b + \delta$
(\texttt{lirpa\_upper\_pos}).

The conclusion's two atoms evaluate entrywise at the shared shape (the
affine terms evaluate to $\beta + \alpha\,p$ and $\delta + \gamma\,p$
per entry --- \texttt{affineWireTerm\_evalQ}), so the four cases
assemble into the conjunction. That is \texttt{sound\_lirpa}, whose
structure the paragraph mirrors step for step.
\end{proof}

\begin{bremark}[three checkpoints suffice]\label{rem:lirpa-checkpoints}
Because $\ReLU$ is affine on each side of $0$ and $a < 0 < b$, an
affine line bounds it on $[a, b]$ \emph{iff} it bounds it at the three
points $a$, $0$, $b$ --- the six side conditions are precisely these
six facts (two lines $\times$ three checkpoints), and equivalently the
degree-$1$ Bernstein certificates of Recipe~\ref{rec:lirpa-instance}.
A producer can therefore emit the tightest certifiable envelope for
its chosen slopes by solving three scalar inequalities per line ---
nothing about the rule rewards loose coefficients.
\end{bremark}

\begin{bremark}[what the strict straddle buys]\label{rem:lirpa-strict}
With $a < 0 < b$ the checkpoint conditions read, verbatim, ``the line
is on the correct side of $\ReLU$ at $a$, $0$, and $b$'' ($\ReLU(a) =
0$ needs $a \leq 0$; $\ReLU(b) = b$ needs $0 \leq b$). Semantic
soundness would in fact survive relaxing the strictness: whatever the
signs of $a \leq b$, the six conditions pin the lower line beneath
$\ReLU$ on $[\min(a,0), \max(b,0)] \supseteq [a,b]$ --- at $a$ because
$\alpha a + \beta \leq 0 \leq \ReLU(a)$ holds unconditionally, and on
each of $[\min(a,0), 0]$, $[0, \max(b,0)]$ by the same endpoint
interpolation; for the upper line the interval containment does the
work (the checks at $0$ and $b$ dominate $\ReLU$ on $[0, \max(b,0)]$,
the checks at $a$ and $0$ dominate it on $[\min(a,0), 0]$). The
rule as adopted nevertheless demands the strict straddle, and reasonably:
it pins \rulename{LiRPA} to the genuinely unstable case --- a stable
neuron rides \rulename{ReLU-Pos} or \rulename{ReLU-Neg}
\emph{exactly}, losing nothing to a relaxation --- it keeps both
interpolation cells nondegenerate, and it is consumed as stated by the
verified interpolation lemmas, whose contradiction products multiply
by $-a$ and $b$.
\end{bremark}

\begin{bremark}[the conclusion is built, not audited]
The conclusion formula is pinned syntactically
(\texttt{lirpaConclusionForm}): it \emph{is} the conjunction of the
two comparison atoms over the affine terms
$\texttt{affineWireTerm}\,\beta\,\alpha\,h$ and
$\texttt{affineWireTerm}\,\delta\,\gamma\,h$, constructed from the
payload coefficients --- the producer supplies no conclusion shape.
A payload field that the checker rebuilds rather than reads is one
the payload cannot mis-declare (the \rulename{PConstExact} section's
closedness observation, in conclusion-shape form): no ``equivalent''
respelling can smuggle in a different wire, comparator, or shape.
\end{bremark}

\begin{anchors}
family linkage extraction & \texttt{reluLinked\_facts}\\
checkpoint interpolation, lower line & \texttt{lirpa\_lower\_neg}, \texttt{lirpa\_lower\_pos}\\
checkpoint interpolation, upper line & \texttt{lirpa\_upper\_neg}, \texttt{lirpa\_upper\_pos}\\
affine conclusion evaluation & \texttt{affineWireTerm\_evalQ}\\
rule warrant (Proposition \ref{prop:lirpa-sound}) & \texttt{sound\_lirpa}\\
linkage side condition & \texttt{reluLinked}; \texttt{PrimApp.isReluAt} (the positive-size gate)\\
pinned conclusion & \texttt{lirpaConclusionForm}; \texttt{affineWireTerm}\\
packaged-node conservativity & \texttt{relu\_projection} over \texttt{reluExpH} (canonical single-entry instance)\\
\end{anchors}

\sourcenote{Literature: the two-sided affine relaxation of ReLU is the
verification community's workhorse --- the triangle
relaxation~\cite{ehlers2017}, its dual/convex-outer
form~\cite{wongkolter2018}, CROWN's adaptive slope
choice~\cite{zhang2018crown}, DeepPoly's abstract-domain
rendering~\cite{singh2019deeppoly}, and the LiRPA framework that named
the family~\cite{xu2020lirpa}. This rule is the checker-side
certificate form: the cited works \emph{find} envelopes; the rule only
ever \emph{checks} one, at six rational inequalities.}


\rulesection{20}{ReLU-Pos}{admissible rule}
\label{blcatrule-ReLUPos}
\ruleblock{The rule (Steele metanotation)}
\sourcenote{as realized in the mechanized development; source row: \texttt{\texttt{rules.toml}} --- \emph{ReLU-Pos}. The metanotation statement is typeset by hand for notational quality; its mechanized ground truth is the Lean statement below, extracted from the sources at every build of this book and hash-bound. The rule \emph{roster} is machine-checked (every mechanized rule tag round-trips to exactly one catalogue row); the statement \emph{text} is editorial, checked by eye against the extracted Lean beside it.}
\begin{mathpar}
\inferrule*[Right=\rulename{ReLU-Pos}]{
  \Phi \vdash h \geq 0
  \\
  \ReLU_{\{h\}\to\{y\}} \in H
}{
  \Phi \vdash y = h
}
\end{mathpar}

\par\medskip
\rulename{ReLU-Pos} is the exactness rule for the active side of a
ReLU: once $\Phi \vdash h \geq 0$ is proven for the preactivation wire
of a cited ReLU application, the output wire is not merely bounded but
\emph{equal} — $\Phi \vdash y = h$. A producer reaches for it when
earlier bound propagation has shown a neuron stably active (a proven
lower box bound at or above zero): the lossy affine envelope of
\rulename{LiRPA} is replaced by an identity, and the equality feeds
\rulename{EqElim}, so every fact known about $h$ transports to $y$
verbatim.

What the checker recomputes, given the payload (premise proof, the two
wires): the linkage — the hypernet really contains a packaged ReLU
application at exactly these ports, of positive total size (the
family\textquotesingle s port-lookup side condition, displayed in the
\rulename{LiRPA} section); that the premise\textquotesingle s
conclusion is literally the atom $h \geq 0$ over the same premise
formula; and that the conclusion is literally $y = h$. Both atoms are
elementwise over the node\textquotesingle s tensor shape.

Drop the premise and the rule is false at any entry with $h < 0$
(there $y = 0 \neq h$). Drop the linkage and $y$ is an arbitrary wire
with no semantic tie to $h$ — the linkage-free rule is unsound over
unlinked wires, and a realization without the port lookup could only
be kept sound by rejecting every certificate.
The interesting hypothesis is the one that is \emph{not}
there: no side condition excludes the boundary $h = 0$, because the
node\textquotesingle s two semantic cases agree at the kink — the
proof section makes this the centre of the argument.

\cataloguenote{TIER NOTE: the unpacking by case analysis over the sign-split fact
discharges the active branch cleanly, but its boundary branch needs
the antisymmetry step $h \leq 0 \wedge h \geq 0 \Rightarrow h = 0$,
and no catalogue rule concludes an equality from a pair of opposite
inequalities — so, like \rulename{PConstExact}, this rule is
admissible by \emph{reduction to verified machinery} rather than by a
complete catalogue derivation: the shipped warrant is the direct lemma
\texttt{sound\_reluPos}, whose boundary case is one antisymmetry step
in the model. The proof section states the enumeration behind that
claim.
}
\ruleblock{Lean statement}
\par\medskip
The warrant consumes the family linkage extraction
(\texttt{reluLinked\_facts}, displayed in the \rulename{LiRPA}
section) and the premise, eliminated entrywise (the elimination lemma
is displayed alongside); the payload structure is
\texttt{PReLUPosData}, quoted in full below.

\declhead{BernsteinLogic.Soundness}{sound\_reluPos}{\badge{BadgeGreen}{statement bound}}
\begin{leansrc}
/-- ReLU-Pos: from `Φ ⊢ h ≥ 0` conclude `Φ ⊢ y = h` — on each entry the
premise forces the nonpositive case's input to 0, where both cases
agree. -/
theorem sound_reluPos {H : Hypernet} {seq premise : Sequent}
    {preact post : Wire}
    (ih : premise.semanticallyValid H)
    (hlink : reluLinked H preact post)
    (hpre : seq.pre = premise.pre)
    (hpremPost : premise.post = .atom .ge (.var preact) (.const 0))
    (hpost : seq.post = .atom .eq (.var post) (.var preact)) :
    seq.semanticallyValid H
\end{leansrc}
\declhead{BernsteinLogic.Soundness}{QTensor.rel2\_const\_right\_elim}{\badge{BadgeGreen}{statement bound}}
\begin{leansrc}
/-- rel2 against a rank-0 constant on the right, eliminated to the
all-entries relation. -/
theorem QTensor.rel2_const_right_elim {r : ℚ → ℚ → Prop} {q : QTensor}
    {c : ℚ} (h : QTensor.rel2 r q ⟨[], [c], rfl⟩) :
    ∀ x ∈ q.entries, r x c
\end{leansrc}
\declhead{BernsteinLogic.Core.Catalogue}{PReLUPosData}{\badge{BadgeGreen}{definition bound}}
\begin{leansrc}
inductive PReLUPosData : Hypernet → Sequent → Type where
  | mk (H : Hypernet) (seq : Sequent)
       (premise : Sequent) (premiseProof : ProofObject H premise)
       (preact post : Wire)
       (link : reluLinked H preact post)
       (preEq       : seq.pre = premise.pre)
       (premisePost : premise.post = .atom .ge (.var preact) (.const 0))
       (conclPost   : seq.post     = .atom .eq (.var post) (.var preact))
       : PReLUPosData H seq
\end{leansrc}
\ruleblock{Correspondence (Steele display $\leftrightarrow$ Lean)}
\badge{BadgeGreen}{reviewer-certified against the displayed texts}
\sourcenote{The dictionary below is a reviewer-certified statement of how the mechanized statements above interpret the displayed rule: display metavariable $\mapsto$ Lean binder, each display hypothesis $\mapsto$ the premise or side condition carrying it, the display conclusion $\mapsto$ the goal --- with any deliberate difference in width between display and mechanization named outright. Certification is recorded as a hash triple over (the displayed cell, the extracted statement texts, this dictionary\textquotesingle s own text); if any of the three changes, this block turns amber until a reviewer re-certifies. The correspondence claim itself remains editorial --- reviewer-grade, not machine-checked; its machine-checked successor types each soundness statement against the rule as a structured object, and supersedes this block per rule as it lands.}
Judgment frame as in \rulename{Ax}; sub-derivation premise as at
\rulename{Weaken}; linkage premise --- including the positive-size
narrowing hidden under the displayed $\in$ --- exactly as in the
\rulename{LiRPA} dictionary, the family preamble.
\begin{itemize}
\item $h \mapsto$ \texttt{preact}; $y \mapsto$ \texttt{post};
  $\ReLU_{\{h\}\to\{y\}} \in H$ $\mapsto$ \texttt{hlink} $:$
  \texttt{reluLinked H preact post} (payload alias:
  \texttt{PReLUPosData.link}, same proposition).
\item premise $\Phi \vdash h \geq 0$ $\mapsto$ \texttt{ih} with
  \texttt{hpremPost} $:$ \texttt{premise.post = .atom .ge (.var
  preact) (.const 0)} --- displayed comparator and operand order
  carried exactly, with no sign-folding (in contrast to the
  \rulename{ConvElim} family\textquotesingle s all-$\leq$ build
  output); $\Phi$ shared via \texttt{hpre}. Payload aliases, same
  propositions: \texttt{ih} $\leftrightarrow$
  \texttt{premiseProof} (over the shared \texttt{premise}),
  \texttt{hpre} $\leftrightarrow$ \texttt{preEq},
  \texttt{hpremPost} $\leftrightarrow$ \texttt{premisePost}.
\item conclusion $\Phi \vdash y = h$ $\mapsto$ the goal under
  \texttt{hpost} $:$ \texttt{seq.post = .atom .eq (.var post) (.var
  preact)} (payload alias: \texttt{conclPost}, same equation) ---
  output wire left, preactivation right, as displayed.
\end{itemize}
Width note: the certified atoms are elementwise over the cited
application\textquotesingle s tensor shape (batch semantics, as at
\rulename{LiRPA}); the boundary $h = 0$ needs no side condition ---
the node\textquotesingle s two semantic cases agree at the kink, by
one antisymmetry step in the model (the commentary and tier note).
Beyond the family\textquotesingle s linkage gate, neither side is
wider: every hypothesis of \texttt{sound\_reluPos} is consumed.

\ruleblock{Lean verifier receipts}
\sourcenote{module \texttt{BernsteinLogic.Soundness}; bound build \texttt{a20d90e66c13\ldots} (bound from receipt a9c0e34d3e3c9c8c, 2026-08-03T21:30:36Z). A receipt records that named checking layers accepted a build: (module, compiled-artifact SHA-256, environment SHA-256, checker layers), minted when the checking gate ran; the store receipts attesting the bound build are listed with the layers each run credited. Nothing here asks the reader to trust the store: every layer is an independently runnable program over the quoted sources --- rebuild the module and re-run the checkers to reproduce the verdict. Receipts attesting superseded builds credit no layer and are omitted from this rendering (the count of omissions, if any, is noted below).}
\begin{receipttable}
\receiptrow{a9c0e34d3e3c9c8c}{2026-08-03T21:30:36Z}{build+decorrelated}{\badge{BadgeGreen}{current}}
\end{receipttable}
\sourcenote{21 further receipt(s) in the store attest superseded builds of this module and are omitted from this rendering.}
\sourcenote{module \texttt{BernsteinLogic.Core.Catalogue}; bound build \texttt{cbc92e976a1d\ldots} (bound from receipt 5abdb3589707b5df, 2026-08-03T21:30:33Z). A receipt records that named checking layers accepted a build: (module, compiled-artifact SHA-256, environment SHA-256, checker layers), minted when the checking gate ran; the store receipts attesting the bound build are listed with the layers each run credited. Nothing here asks the reader to trust the store: every layer is an independently runnable program over the quoted sources --- rebuild the module and re-run the checkers to reproduce the verdict. Receipts attesting superseded builds credit no layer and are omitted from this rendering (the count of omissions, if any, is noted below).}
\begin{receipttable}
\receiptrow{24804e66edba0bc9}{2026-08-03T20:25:09Z}{build+decorrelated}{\badge{BadgeGreen}{current}}
\receiptrow{5abdb3589707b5df}{2026-08-03T21:30:33Z}{build+decorrelated}{\badge{BadgeGreen}{current}}
\end{receipttable}
\sourcenote{18 further receipt(s) in the store attest superseded builds of this module and are omitted from this rendering.}
\layerline{\layeritem{build}{lake build (elaboration)}{\badge{BadgeGreen}{current}}\layeritem{replay}{leanchecker (official kernel replay)}{\badge{BadgeAmber}{stale only}}\layeritem{hardened}{safe\_verify (hardened screen)}{\badge{BadgeAmber}{stale only}}\layeritem{decorrelated}{lean4export + nanoda (decorrelated kernel)}{\badge{BadgeGreen}{current}}}
\sourcenote{\badge{BadgeGray}{comparator: none} no fixture in the three-implementation comparator harness (compiled / kernel-interpreted / reference-interpreted) exercises this rule.}
\ruleblock{HOL4 statement}
\sourcenote{theory \texttt{bern\_tr\_catdata2}, HOL4 development script \texttt{\texttt{bern\_tr\_catdata2Script.sml}} \badge{BadgeGreen}{script bound}}
\holhead{Theorem preluposdata\_receipt}
\begin{holsrc}
Theorem preluposdata_receipt:
  ruledataOk [] cd2Hrelu
    <| pre := FAtom BowGe (TmVar 3) (TmConst 0) ;
       post := FAtom BowEq (TmVar 4) (TmVar 3) |>
    (RdReluPos (PReluPosD
       <| pre := FAtom BowGe (TmVar 3) (TmConst 0) ;
          post := FAtom BowGe (TmVar 3) (TmConst 0) |> (ProofObj RdAx)
       3 4)) ∧
  ¬ ruledataOk [] hnEmpty
      <| pre := FAtom BowGe (TmVar 3) (TmConst 0) ;
         post := FAtom BowEq (TmVar 4) (TmVar 3) |>
      (RdReluPos (PReluPosD
         <| pre := FAtom BowGe (TmVar 3) (TmConst 0) ;
            post := FAtom BowGe (TmVar 3) (TmConst 0) |> (ProofObj RdAx)
         3 4))
\end{holsrc}
\ruleblock{HOL verifier receipts}
\badge{BadgeBlue}{documentary} \texttt{hol-strict} CLEAN: theory \texttt{bern\_tr\_catdata2}, 272
theorems, 0 cheat-tagged, 0 nonstandard axioms (whitelist:
\texttt{DISK\_THM}; the four \texttt{bool} axioms) --- verification
run of 2026-07-22 on an isolated build host, against a
content-addressed checkout of the exact sources cited here (git blob
\texttt{cd3c1e95\ldots}). The proof-object stratum is ported as one
mutual datatype family with a mutual wellformedness predicate,
exactly as the Lean mutual inductive block; this
rule\textquotesingle s arm states its entire checkable content ---
the premise post-form, the family linkage, and the conclusion
equation --- and the linkage lookup into the
hypernet\textquotesingle s application list has a ported realization
(\texttt{reluLinkedB}, theory \texttt{bern\_tr\_catdata}, 252
theorems, 0 offences, same run). The receipt theorem cited here is
the kernel-evaluated pair: an accept over a hypernet carrying the
packaged ReLU application, and a refusal over one without it ---
linkage doing exactly its job. Honest residue: the semantic soundness
twin (the Lean development\textquotesingle s \texttt{sound\_reluPos})
is not yet ported; the payload/checker layer above is.

\sourcenote{HOL receipts are documentary: quoted results of \texttt{hol-strict} runs
and census rows from the HOL4 development\textquotesingle s build
evidence, not minted receipt artefacts. A HOL receipt store (the Lean
store\textquotesingle s twin) is named future work.
}
\ruleblock{Proof}

\sourcenote{Genre: admissible rule (tier 2). The family machinery ---
the packaged node, the linkage, the two-case semantic fact, and its
conservativity warrant --- is stated in the \rulename{LiRPA} section's
family template and is not restated here.}

\begin{brecipe}[the \rulename{ReLU-Pos} instance]\label{rec:relupos-instance}
The node, and its intended rewrite over the family's sign-split fact:
\[
\inferrule*[Right=\rulename{ReLU-Pos}]{
  \mathcal{T} \rhd \Phi \vdash h \geq 0 \\
  \ReLU_{\{h\}\to\{y\}} \in H
}{
  \Phi \vdash y = h
}
\qquad\rightsquigarrow\qquad
\text{case analysis on }
(h \leq 0 \wedge y = 0) \vee (0 \leq h \wedge y = h).
\]
The \emph{active} branch discharges cleanly: its context carries the
conjunct $y = h$ (the case disjunct enters the branch premise as one
ordered subtree, so two $\wedge$E steps reach it), which \rulename{Ax}
and $\wedge$E extract --- the goal verbatim. The \emph{boundary} branch does not: its context
carries $h \leq 0$ and $y = 0$, the premise supplies $h \geq 0$, and
the goal $y = h$ now needs the antisymmetry fact $h = 0$ --- which, as
\autoref{rem:relupos-antisym} shows, no catalogue derivation that
withholds this rule itself provides. The rewrite therefore does \emph{not} fully discharge into
catalogue nodes; like \rulename{PConstExact}, the rule is admissible
by reduction to verified machinery: the shipped warrant is the direct
lemma \texttt{sound\_reluPos}, in which the boundary case is a single
antisymmetry step in the model ($p \leq 0$ from the case, $0 \leq p$
from the premise, so $p = 0$ and the two cases agree).
\end{brecipe}

\begin{bprop}[semantic warrant]\label{prop:relupos-sound}
Suppose the linkage holds and the premise sequent --- same premise
formula, conclusion the atom $h \geq 0$ --- is semantically valid.
Then the sequent concluding $y = h$ is semantically valid.
\end{bprop}
\begin{proof}
Fix an admissible valuation $\sigma$ satisfying $\Phi$. The premise
eliminates entrywise: every entry $p$ of the preactivation tensor
satisfies $0 \leq p$. The linkage extraction
(\texttt{reluLinked\_facts}) gives the output tensor at the same
shape and, per entry, either $p \leq 0 \wedge y = 0$ or $0 \leq p
\wedge y = p$. In the second case $y = p$ is the goal entry. In the
first, $p \leq 0$ meets the premise's $0 \leq p$, so $p = 0$
(antisymmetry --- the model step the catalogue lacks) and $y = 0 = p$.
The conclusion atom is the entrywise equality of the two tensors at
their shared shape. That is \texttt{sound\_reluPos}, step for step.
\end{proof}

\begin{bremark}[the kink case has no catalogue derivation]
\label{rem:relupos-antisym}
The claim in Recipe~\ref{rec:relupos-instance} is an enumeration over
the catalogue's equality-capable conclusions. \rulename{Norm} and
\rulename{EqRefl} conclude only canonical-form identities ($h$ and $0$
are not canonically equal); \rulename{EqSym}, \rulename{EqElim}, and
\rulename{ArithAtom} at the $=$ comparator only \emph{transport}
equalities already proven; this family and the cell rules conclude
equalities only over their packaged applications' ports; the
combination rules (\rulename{Farkas}, \rulename{Syzygy},
\rulename{PBound}, \rulename{BernCone} with an order predicate)
conclude order atoms, never $=$; and \rulename{BernCone}'s degenerate
conic-predicate route would need a region atom pinning $h$ to the
single point $0$, which no rule derives --- region atoms denote
regions of positive-dimensional simplices. The catalogue at large \emph{does} reach the fact, and the route is
instructive: over the two-wire hypernet carrying only the linked ReLU
application, \rulename{ReLU-Neg} (premise $h \leq 0$, an
\rulename{Ax} extraction) concludes $y = 0$, this rule (premise
$h \geq 0$) concludes $y = h$, and one \rulename{EqElim} along
$y = h$ rewrites $y = 0$ into $h = 0$ --- the linked application is
itself the antisymmetry oracle, its output forced to agree with both
branches at the kink. But the elaboration needs the fact
\emph{without the rule being eliminated} --- the leaf-first species
of the admissible template consumes every rule except that one ---
and with either envelope rule withheld the supply dries up. Write
$d(t = u)$ for the canonical form of $t - u$, and consider a smallest
derivation, withholding one envelope rule, that concludes an equality
atom with $d$ \emph{not} a polynomial multiple of the surviving port
difference ($y$ if \rulename{ReLU-Pos} is withheld with
\rulename{ReLU-Neg}'s $y = 0$ surviving; $y - h$ mirrored). Its final
step cannot be the surviving port conclusion or a canonical-form
identity ($d$ a multiple of the port difference, respectively $0$);
not \rulename{EqSym}, the zero-form respelling, or \rulename{EqElim}
transport (each changes $d$ by a sign or by polynomial multiples of a
previously concluded equality's own $d$ --- a smaller derivation,
hence a permitted multiple); not a closed ground fact ($d = 0$, the
constants being exactly evaluated); not a combination or region rule
(order and region conclusions only); and not a structural or
case-analysis step, which only re-exposes an equality already
concluded, or sitting in a premise set --- the root premises
($h \leq 0 \wedge h \geq 0$, at either spelling) contain no equality,
premise sets grow only by \rulename{Cut} formulas and case disjuncts
concluded or contained in smaller derivations, and a case split whose
smuggled disjunct carries the target is idle: its sibling branch must
conclude the same target without it, a smaller offender. The target
$d(h = 0) = h$ is a multiple of neither $y$ nor $y - h$, so $h \leq 0
\wedge h \geq 0 \vdash h = 0$ has no \rulename{ReLU-Pos}-free
derivation (nor, mirrored, a \rulename{ReLU-Neg}-free one), and the
boundary branch cannot discharge into nodes. This is not a defect to
repair but a fact about where the calculus keeps its arithmetic: an
elaboration can source order antisymmetry only from the verified
model step, one line of \texttt{sound\_reluPos}, rather than from a
rule every certificate would pay wire format for.
\end{bremark}

\begin{bremark}[what the premise is for]
Without $\Phi \vdash h \geq 0$ the rule is false at any entry with $h
< 0$: there the node's semantics forces $y = 0 \neq h$. The premise is
consumed \emph{entrywise} --- a tensor-shaped preactivation must be
proven nonnegative in every entry, which is what a rank-$0$ constant
against a tensor atom asserts. No side condition excludes the boundary
$h = 0$, and none is needed: the two semantic cases agree there, which
is exactly the antisymmetry step above.
\end{bremark}

\begin{anchors}
Recipe \ref{rec:relupos-instance} warrant & \texttt{sound\_reluPos}\\
family linkage extraction & \texttt{reluLinked\_facts} (see the \rulename{LiRPA} section)\\
entrywise premise elimination & \texttt{QTensor.rel2\_const\_right\_elim}\\
the kink antisymmetry & the \texttt{le\_antisymm} step inside \texttt{sound\_reluPos}\\
payload structure & \texttt{PReLUPosData} (premise proof, two wires, linkage)\\
\end{anchors}


\rulesection{21}{ReLU-Neg}{admissible rule}
\label{blcatrule-ReLUNeg}
\ruleblock{The rule (Steele metanotation)}
\sourcenote{as realized in the mechanized development; source row: \texttt{\texttt{rules.toml}} --- \emph{ReLU-Neg}. The metanotation statement is typeset by hand for notational quality; its mechanized ground truth is the Lean statement below, extracted from the sources at every build of this book and hash-bound. The rule \emph{roster} is machine-checked (every mechanized rule tag round-trips to exactly one catalogue row); the statement \emph{text} is editorial, checked by eye against the extracted Lean beside it.}
\begin{mathpar}
\inferrule*[Right=\rulename{ReLU-Neg}]{
  \Phi \vdash h \leq 0
  \\
  \ReLU_{\{h\}\to\{y\}} \in H
}{
  \Phi \vdash y = 0
}
\end{mathpar}

\par\medskip
\rulename{ReLU-Neg} is the exactness rule for the inactive side of a
ReLU, \rulename{ReLU-Pos}\textquotesingle s mirror: once $\Phi \vdash
h \leq 0$ is proven for the preactivation wire of a cited ReLU
application, the output is pinned to a constant — $\Phi \vdash y = 0$.
A producer reaches for it when bound propagation has shown a neuron
stably inactive (a proven upper box bound at or below zero): the
neuron is dead on the region of interest, and the equality lets
\rulename{EqElim} rewrite $y$ to the literal constant $0$ downstream
— dead-neuron elimination in certificate form, at the rank-0
instances the constant-sided transport gate admits (the
family\textquotesingle s shape-scope remark).

What the checker recomputes, given the payload (premise proof, the two
wires): the linkage — a packaged ReLU application at exactly these
ports, of positive total size (the family\textquotesingle s
port-lookup side condition, displayed in the \rulename{LiRPA}
section); that the premise\textquotesingle s conclusion is literally
the atom $h \leq 0$ over the same premise formula; and that the
conclusion is literally $y = 0$. Both atoms are elementwise over the
node\textquotesingle s tensor shape.

Drop the premise and the rule is false at any entry with $h > 0$
(there $y = h \neq 0$). Drop the linkage and an arbitrary wire is
asserted zero. As with \rulename{ReLU-Pos}, the boundary $h = 0$ needs
no exclusion: both semantic cases of the node force $y = 0$ there.

\cataloguenote{TIER NOTE: as for \rulename{ReLU-Pos}, and for the mirrored reason —
the unpacking\textquotesingle s active branch would need $h = 0$ from
$h \geq 0 \wedge h \leq 0$, an antisymmetry step no catalogue rule
provides — the rule is admissible by reduction to verified machinery:
the shipped warrant is the direct lemma \texttt{sound\_reluNeg}. The
enumeration argument is stated once, in the \rulename{ReLU-Pos}
proof section.
}
\ruleblock{Lean statement}
\par\medskip
The warrant consumes the family linkage extraction
(\texttt{reluLinked\_facts}, displayed in the \rulename{LiRPA}
section) and the premise, eliminated entrywise (the elimination lemma
is displayed alongside); the payload structure is
\texttt{PReLUNegData}, quoted in full below.

\declhead{BernsteinLogic.Soundness}{sound\_reluNeg}{\badge{BadgeGreen}{statement bound}}
\begin{leansrc}
/-- ReLU-Neg: from `Φ ⊢ h ≤ 0` conclude `Φ ⊢ y = 0` — on each entry the
premise forces the nonnegative case's input to 0, where both cases
agree. -/
theorem sound_reluNeg {H : Hypernet} {seq premise : Sequent}
    {preact post : Wire}
    (ih : premise.semanticallyValid H)
    (hlink : reluLinked H preact post)
    (hpre : seq.pre = premise.pre)
    (hpremPost : premise.post = .atom .le (.var preact) (.const 0))
    (hpost : seq.post = .atom .eq (.var post) (.const 0)) :
    seq.semanticallyValid H
\end{leansrc}
\declhead{BernsteinLogic.Soundness}{QTensor.rel2\_const\_right\_elim}{\badge{BadgeGreen}{statement bound}}
\begin{leansrc}
/-- rel2 against a rank-0 constant on the right, eliminated to the
all-entries relation. -/
theorem QTensor.rel2_const_right_elim {r : ℚ → ℚ → Prop} {q : QTensor}
    {c : ℚ} (h : QTensor.rel2 r q ⟨[], [c], rfl⟩) :
    ∀ x ∈ q.entries, r x c
\end{leansrc}
\declhead{BernsteinLogic.Core.Catalogue}{PReLUNegData}{\badge{BadgeGreen}{definition bound}}
\begin{leansrc}
inductive PReLUNegData : Hypernet → Sequent → Type where
  | mk (H : Hypernet) (seq : Sequent)
       (premise : Sequent) (premiseProof : ProofObject H premise)
       (preact post : Wire)
       (link : reluLinked H preact post)
       (preEq       : seq.pre = premise.pre)
       (premisePost : premise.post = .atom .le (.var preact) (.const 0))
       (conclPost   : seq.post     = .atom .eq (.var post) (.const 0))
       : PReLUNegData H seq
\end{leansrc}
\ruleblock{Correspondence (Steele display $\leftrightarrow$ Lean)}
\badge{BadgeGreen}{reviewer-certified against the displayed texts}
\sourcenote{The dictionary below is a reviewer-certified statement of how the mechanized statements above interpret the displayed rule: display metavariable $\mapsto$ Lean binder, each display hypothesis $\mapsto$ the premise or side condition carrying it, the display conclusion $\mapsto$ the goal --- with any deliberate difference in width between display and mechanization named outright. Certification is recorded as a hash triple over (the displayed cell, the extracted statement texts, this dictionary\textquotesingle s own text); if any of the three changes, this block turns amber until a reviewer re-certifies. The correspondence claim itself remains editorial --- reviewer-grade, not machine-checked; its machine-checked successor types each soundness statement against the rule as a structured object, and supersedes this block per rule as it lands.}
Judgment frame as in \rulename{Ax}; sub-derivation premise as at
\rulename{Weaken}; linkage premise --- including the positive-size
narrowing hidden under the displayed $\in$ --- exactly as in the
\rulename{LiRPA} dictionary, the family preamble.
\begin{itemize}
\item $h \mapsto$ \texttt{preact}; $y \mapsto$ \texttt{post};
  $\ReLU_{\{h\}\to\{y\}} \in H$ $\mapsto$ \texttt{hlink} $:$
  \texttt{reluLinked H preact post} (payload alias:
  \texttt{PReLUNegData.link}, same proposition).
\item premise $\Phi \vdash h \leq 0$ $\mapsto$ \texttt{ih} with
  \texttt{hpremPost} $:$ \texttt{premise.post = .atom .le (.var
  preact) (.const 0)} --- displayed comparator and operand order
  carried exactly; $\Phi$ shared via \texttt{hpre}. Payload aliases, same
  propositions: \texttt{ih} $\leftrightarrow$
  \texttt{premiseProof} (over the shared \texttt{premise}),
  \texttt{hpre} $\leftrightarrow$ \texttt{preEq},
  \texttt{hpremPost} $\leftrightarrow$ \texttt{premisePost}.
\item conclusion $\Phi \vdash y = 0$ $\mapsto$ the goal under
  \texttt{hpost} $:$ \texttt{seq.post = .atom .eq (.var post)
  (.const 0)} (payload alias: \texttt{conclPost}, same equation) --- output
  wire left, constant right, as displayed.
\end{itemize}
Width note: certified atoms elementwise over the cited
application\textquotesingle s tensor shape (batch semantics, as at
\rulename{LiRPA}); the boundary $h = 0$ is again condition-free ---
the mirrored antisymmetry step of the tier note. Beyond the
family\textquotesingle s linkage gate, neither side is wider: every
hypothesis of \texttt{sound\_reluNeg} is consumed.

\ruleblock{Lean verifier receipts}
\sourcenote{module \texttt{BernsteinLogic.Soundness}; bound build \texttt{a20d90e66c13\ldots} (bound from receipt a9c0e34d3e3c9c8c, 2026-08-03T21:30:36Z). A receipt records that named checking layers accepted a build: (module, compiled-artifact SHA-256, environment SHA-256, checker layers), minted when the checking gate ran; the store receipts attesting the bound build are listed with the layers each run credited. Nothing here asks the reader to trust the store: every layer is an independently runnable program over the quoted sources --- rebuild the module and re-run the checkers to reproduce the verdict. Receipts attesting superseded builds credit no layer and are omitted from this rendering (the count of omissions, if any, is noted below).}
\begin{receipttable}
\receiptrow{a9c0e34d3e3c9c8c}{2026-08-03T21:30:36Z}{build+decorrelated}{\badge{BadgeGreen}{current}}
\end{receipttable}
\sourcenote{21 further receipt(s) in the store attest superseded builds of this module and are omitted from this rendering.}
\sourcenote{module \texttt{BernsteinLogic.Core.Catalogue}; bound build \texttt{cbc92e976a1d\ldots} (bound from receipt 5abdb3589707b5df, 2026-08-03T21:30:33Z). A receipt records that named checking layers accepted a build: (module, compiled-artifact SHA-256, environment SHA-256, checker layers), minted when the checking gate ran; the store receipts attesting the bound build are listed with the layers each run credited. Nothing here asks the reader to trust the store: every layer is an independently runnable program over the quoted sources --- rebuild the module and re-run the checkers to reproduce the verdict. Receipts attesting superseded builds credit no layer and are omitted from this rendering (the count of omissions, if any, is noted below).}
\begin{receipttable}
\receiptrow{24804e66edba0bc9}{2026-08-03T20:25:09Z}{build+decorrelated}{\badge{BadgeGreen}{current}}
\receiptrow{5abdb3589707b5df}{2026-08-03T21:30:33Z}{build+decorrelated}{\badge{BadgeGreen}{current}}
\end{receipttable}
\sourcenote{18 further receipt(s) in the store attest superseded builds of this module and are omitted from this rendering.}
\layerline{\layeritem{build}{lake build (elaboration)}{\badge{BadgeGreen}{current}}\layeritem{replay}{leanchecker (official kernel replay)}{\badge{BadgeAmber}{stale only}}\layeritem{hardened}{safe\_verify (hardened screen)}{\badge{BadgeAmber}{stale only}}\layeritem{decorrelated}{lean4export + nanoda (decorrelated kernel)}{\badge{BadgeGreen}{current}}}
\sourcenote{\badge{BadgeGray}{comparator: none} no fixture in the three-implementation comparator harness (compiled / kernel-interpreted / reference-interpreted) exercises this rule.}
\ruleblock{HOL4 statement}
\sourcenote{theory \texttt{bern\_tr\_catdata2}, HOL4 development script \texttt{\texttt{bern\_tr\_catdata2Script.sml}} \badge{BadgeGreen}{script bound}}
\holhead{Theorem prelunegdata\_receipt}
\begin{holsrc}
Theorem prelunegdata_receipt:
  ruledataOk [] cd2Hrelu
    <| pre := FAtom BowLe (TmVar 3) (TmConst 0) ;
       post := FAtom BowEq (TmVar 4) (TmConst 0) |>
    (RdReluNeg (PReluNegD
       <| pre := FAtom BowLe (TmVar 3) (TmConst 0) ;
          post := FAtom BowLe (TmVar 3) (TmConst 0) |> (ProofObj RdAx)
       3 4)) ∧
  ¬ ruledataOk [] cd2Hrelu
      <| pre := FAtom BowLe (TmVar 3) (TmConst 0) ;
         post := FAtom BowEq (TmVar 4) (TmVar 3) |>
      (RdReluNeg (PReluNegD
         <| pre := FAtom BowLe (TmVar 3) (TmConst 0) ;
            post := FAtom BowLe (TmVar 3) (TmConst 0) |> (ProofObj RdAx)
         3 4))
\end{holsrc}
\ruleblock{HOL verifier receipts}
\badge{BadgeBlue}{documentary} \texttt{hol-strict} CLEAN: theory \texttt{bern\_tr\_catdata2}, 272
theorems, 0 cheat-tagged, 0 nonstandard axioms (whitelist:
\texttt{DISK\_THM}; the four \texttt{bool} axioms) --- verification
run of 2026-07-22 on an isolated build host, against a
content-addressed checkout of the exact sources cited here (git blob
\texttt{cd3c1e95\ldots}). The proof-object stratum is ported as one
mutual datatype family with a mutual wellformedness predicate,
exactly as the Lean mutual inductive block; this
rule\textquotesingle s arm states its entire checkable content ---
the premise post-form, the family linkage, and the zero-clamp
conclusion equation --- and the linkage lookup into the
hypernet\textquotesingle s application list has a ported realization
(\texttt{reluLinkedB}, theory \texttt{bern\_tr\_catdata}, 252
theorems, 0 offences, same run). The receipt theorem cited here is
the kernel-evaluated pair: an accept over a hypernet carrying the
packaged ReLU application, and a wrong-conclusion refusal. Honest
residue: the semantic soundness twin (the Lean
development\textquotesingle s \texttt{sound\_reluNeg}) is not yet
ported; the payload/checker layer above is.

\sourcenote{HOL receipts are documentary: quoted results of \texttt{hol-strict} runs
and census rows from the HOL4 development\textquotesingle s build
evidence, not minted receipt artefacts. A HOL receipt store (the Lean
store\textquotesingle s twin) is named future work.
}
\ruleblock{Proof}

\sourcenote{Genre: admissible rule (tier 2). The family machinery is
stated in the \rulename{LiRPA} section's family template;
\rulename{ReLU-Pos} is this rule's mirror, and the shared boundary
argument is stated there.}

\begin{brecipe}[the \rulename{ReLU-Neg} instance]\label{rec:reluneg-instance}
The node, and its intended rewrite over the family's sign-split fact:
\[
\inferrule*[Right=\rulename{ReLU-Neg}]{
  \mathcal{T} \rhd \Phi \vdash h \leq 0 \\
  \ReLU_{\{h\}\to\{y\}} \in H
}{
  \Phi \vdash y = 0
}
\qquad\rightsquigarrow\qquad
\text{case analysis on }
(h \leq 0 \wedge y = 0) \vee (0 \leq h \wedge y = h).
\]
The \emph{inactive} branch discharges cleanly: its context carries the
conjunct $y = 0$ (one ordered subtree of the branch premise; two
$\wedge$E steps), extracted by \rulename{Ax} and $\wedge$E --- the goal
verbatim. The \emph{boundary} branch is the mirror of
\rulename{ReLU-Pos}'s: its context carries $0 \leq h$ and $y = h$, the
premise supplies $h \leq 0$, and the goal $y = 0$ needs $h = 0$ --- the
antisymmetry fact that no catalogue derivation withholding this rule
itself provides (the enumeration argument, covering both withheld
orientations, is stated once, as
\autoref{rem:relupos-antisym} in the \rulename{ReLU-Pos} section).
The rule is therefore admissible by reduction to verified machinery:
the shipped warrant is the direct lemma \texttt{sound\_reluNeg}, whose
boundary case is the same one-line antisymmetry step in the model.
\end{brecipe}

\begin{bprop}[semantic warrant]\label{prop:reluneg-sound}
Suppose the linkage holds and the premise sequent --- same premise
formula, conclusion the atom $h \leq 0$ --- is semantically valid.
Then the sequent concluding $y = 0$ is semantically valid.
\end{bprop}
\begin{proof}
Fix an admissible valuation $\sigma$ satisfying $\Phi$. The premise
eliminates entrywise: every entry $p$ of the preactivation tensor
satisfies $p \leq 0$. The linkage extraction
(\texttt{reluLinked\_facts}) gives, per entry, either $p \leq 0 \wedge
y = 0$ --- the goal entry directly --- or $0 \leq p \wedge y = p$, in
which case $p = 0$ by antisymmetry against the premise, so $y = p =
0$. The conclusion atom compares the output tensor entrywise against
the rank-$0$ constant $0$. That is \texttt{sound\_reluNeg}, step for
step.
\end{proof}

\begin{bremark}[what the premise is for, and what the rule feeds]
Without $\Phi \vdash h \leq 0$ the rule is false at any entry with
$h > 0$ (there $y = h \neq 0$). Downstream, the conclusion differs
from \rulename{ReLU-Pos}'s in kind: $y = 0$ pins the output to a
\emph{literal constant}, so \rulename{EqElim} can rewrite $y$ away ---
dead-neuron elimination inside a certificate --- at the rank-$0$
instances the transport gate admits (a constant-sided equality rides
\rulename{EqElim}'s shared-broadcast-shape condition; the family's
shape-scope remark, \rulename{ReLU-NN} section), after which ground
arithmetic on the constant rides \rulename{PConstExact}.
\end{bremark}

\begin{anchors}
Recipe \ref{rec:reluneg-instance} warrant & \texttt{sound\_reluNeg}\\
family linkage extraction & \texttt{reluLinked\_facts} (see the \rulename{LiRPA} section)\\
entrywise premise elimination & \texttt{QTensor.rel2\_const\_right\_elim}\\
the kink antisymmetry & the \texttt{le\_antisymm} step inside \texttt{sound\_reluNeg}\\
payload structure & \texttt{PReLUNegData} (premise proof, two wires, linkage)\\
\end{anchors}


\rulesection{22}{ReLU-NN}{admissible rule}
\label{blcatrule-ReLUNN}
\ruleblock{The rule (Steele metanotation)}
\sourcenote{as realized in the mechanized development; source row: \texttt{\texttt{rules.toml}} --- \emph{ReLU-NN}. The metanotation statement is typeset by hand for notational quality; its mechanized ground truth is the Lean statement below, extracted from the sources at every build of this book and hash-bound. The rule \emph{roster} is machine-checked (every mechanized rule tag round-trips to exactly one catalogue row); the statement \emph{text} is editorial, checked by eye against the extracted Lean beside it.}
\begin{mathpar}
\inferrule*[Right=\rulename{ReLU-NN}]{
  \ReLU_{\{h\}\to\{y\}} \in H
}{
  \Phi \vdash y \geq 0
}
\end{mathpar}

\par\medskip
\rulename{ReLU-NN} is the range fact: the output wire of a cited ReLU
application is nonnegative, $\Phi \vdash y \geq 0$, with no proof
premises at all — $\Phi$ is arbitrary (the premise-free convention
shared with \rulename{Norm} and \rulename{PConstExact}). It is the
first thing a producer knows about every ReLU output: the lower half
of every output box starts here, before any information about the
preactivation exists, and together with \rulename{ReLU-Above} it
yields the exact pointwise lower envelope of the activation
($y \geq \max(0, h)$, as two atoms).

What the checker recomputes is nearly everything, because there is
nearly nothing: the linkage — the hypernet really contains a packaged
ReLU application at exactly these ports, of positive total size (the
family\textquotesingle s port-lookup side condition, displayed in the
\rulename{LiRPA} section) — and that the conclusion is literally the
atom $y \geq 0$, elementwise over the node\textquotesingle s tensor
shape.

The linkage is the only hypothesis, and dropping it is the whole
failure mode: an arbitrary wire asserted nonnegative, against an
arbitrary premise — the linkage-free rule is unsound over unlinked
wires, and a realization without the port lookup could only be kept
sound by rejecting every certificate. The positive-size gate inside the
linkage keeps even the vacuous-surface corner unciteable (a size-zero
node satisfies every elementwise atom trivially; the linkage refuses
it at the source).

\cataloguenote{This is the family\textquotesingle s cleanest unpacking: the case
analysis over the sign-split fact discharges \emph{both} branches into
catalogue nodes — the inactive branch\textquotesingle s leaf is the
ground fact $0 \geq 0$, which rides \rulename{PConstExact} (the one
rule that can state a closed ground fact). The
proof section displays the full intended derivation (its case-closing
disjunction elimination is the context-carrying form --- a satellite
theorem today, so the display is an elaboration blueprint rather than
an emittable certificate), and states its shape
scope honestly: the constant-sided equality transport type-checks at
rank-0 instances, and the shape-general warrant is the direct lemma.
}
\ruleblock{Lean statement}
\par\medskip
The warrant consumes the family linkage extraction
(\texttt{reluLinked\_facts}, displayed in the \rulename{LiRPA}
section): each output entry is either $0$ or a copy of a nonnegative
input entry. The payload structure, \texttt{PReLUNNData}
--- quoted in full below --- carries no premise field.

\declhead{BernsteinLogic.Soundness}{sound\_reluNN}{\badge{BadgeGreen}{statement bound}}
\begin{leansrc}
/-- ReLU-NN: `Φ ⊢ y ≥ 0`, premise-free — every entry of the packaged
output is 0 or a nonnegative copy of the input entry. (A review finding
exhibited the rule as unsound over unbound wires; the linkage is the repair.) -/
theorem sound_reluNN {H : Hypernet} {seq : Sequent} {preact post : Wire}
    (hlink : reluLinked H preact post)
    (hpost : seq.post = .atom .ge (.var post) (.const 0)) :
    seq.semanticallyValid H
\end{leansrc}
\declhead{BernsteinLogic.Core.Catalogue}{PReLUNNData}{\badge{BadgeGreen}{definition bound}}
\begin{leansrc}
inductive PReLUNNData : Hypernet → Sequent → Type where
  | mk (H : Hypernet) (seq : Sequent)
       (preact post : Wire)
       (link : reluLinked H preact post)
       (conclPost : seq.post = .atom .ge (.var post) (.const 0))
       : PReLUNNData H seq
\end{leansrc}
\ruleblock{Correspondence (Steele display $\leftrightarrow$ Lean)}
\badge{BadgeAmber}{a bound Lean statement CHANGED since certification --- re-review and re-certify}
\sourcenote{The dictionary below is a reviewer-certified statement of how the mechanized statements above interpret the displayed rule: display metavariable $\mapsto$ Lean binder, each display hypothesis $\mapsto$ the premise or side condition carrying it, the display conclusion $\mapsto$ the goal --- with any deliberate difference in width between display and mechanization named outright. Certification is recorded as a hash triple over (the displayed cell, the extracted statement texts, this dictionary\textquotesingle s own text); if any of the three changes, this block turns amber until a reviewer re-certifies. The correspondence claim itself remains editorial --- reviewer-grade, not machine-checked; its machine-checked successor types each soundness statement against the rule as a structured object, and supersedes this block per rule as it lands.}
Judgment frame as in \rulename{Ax}; linkage premise --- including the
positive-size narrowing hidden under the displayed $\in$ --- exactly
as in the \rulename{LiRPA} dictionary, the family preamble.
\begin{itemize}
\item $h \mapsto$ \texttt{preact}; $y \mapsto$ \texttt{post};
  $\ReLU_{\{h\}\to\{y\}} \in H$ $\mapsto$ \texttt{hlink} $:$
  \texttt{reluLinked H preact post} (payload alias:
  \texttt{PReLUNNData.link}, same proposition). The
  display\textquotesingle s
  universality over the preactivation --- no premise constrains $h$
  --- is carried by the
  linkage alone: \texttt{preact} occurs in no other hypothesis (the
  roster note records that the appendix payload named only the output
  wire --- the certified linkage names both ports).
\item premise-freeness $\mapsto$ silence, the \rulename{PConstExact}
  representation: the precondition component of \texttt{seq :
  Sequent} is unconstrained by any hypothesis, so the statement
  holds under arbitrary premises $\Phi$.
\item conclusion $\Phi \vdash y \geq 0$ $\mapsto$ the goal under
  \texttt{hpost} $:$ \texttt{seq.post = .atom .ge (.var post)
  (.const 0)} (payload alias: \texttt{conclPost}, same equation) ---
  displayed comparator and operand order carried exactly.
\end{itemize}
Width note: certified atoms elementwise over the cited
application\textquotesingle s tensor shape (batch semantics, as at
\rulename{LiRPA}). Beyond the family\textquotesingle s linkage gate,
neither side is wider: both hypotheses of \texttt{sound\_reluNN} are
consumed, and nothing else is stated.

\ruleblock{Lean verifier receipts}
\sourcenote{module \texttt{BernsteinLogic.Soundness}; bound build \texttt{a20d90e66c13\ldots} (bound from receipt a9c0e34d3e3c9c8c, 2026-08-03T21:30:36Z). A receipt records that named checking layers accepted a build: (module, compiled-artifact SHA-256, environment SHA-256, checker layers), minted when the checking gate ran; the store receipts attesting the bound build are listed with the layers each run credited. Nothing here asks the reader to trust the store: every layer is an independently runnable program over the quoted sources --- rebuild the module and re-run the checkers to reproduce the verdict. Receipts attesting superseded builds credit no layer and are omitted from this rendering (the count of omissions, if any, is noted below).}
\begin{receipttable}
\receiptrow{a9c0e34d3e3c9c8c}{2026-08-03T21:30:36Z}{build+decorrelated}{\badge{BadgeGreen}{current}}
\end{receipttable}
\sourcenote{21 further receipt(s) in the store attest superseded builds of this module and are omitted from this rendering.}
\sourcenote{module \texttt{BernsteinLogic.Core.Catalogue}; bound build \texttt{cbc92e976a1d\ldots} (bound from receipt 5abdb3589707b5df, 2026-08-03T21:30:33Z). A receipt records that named checking layers accepted a build: (module, compiled-artifact SHA-256, environment SHA-256, checker layers), minted when the checking gate ran; the store receipts attesting the bound build are listed with the layers each run credited. Nothing here asks the reader to trust the store: every layer is an independently runnable program over the quoted sources --- rebuild the module and re-run the checkers to reproduce the verdict. Receipts attesting superseded builds credit no layer and are omitted from this rendering (the count of omissions, if any, is noted below).}
\begin{receipttable}
\receiptrow{24804e66edba0bc9}{2026-08-03T20:25:09Z}{build+decorrelated}{\badge{BadgeGreen}{current}}
\receiptrow{5abdb3589707b5df}{2026-08-03T21:30:33Z}{build+decorrelated}{\badge{BadgeGreen}{current}}
\end{receipttable}
\sourcenote{18 further receipt(s) in the store attest superseded builds of this module and are omitted from this rendering.}
\layerline{\layeritem{build}{lake build (elaboration)}{\badge{BadgeGreen}{current}}\layeritem{replay}{leanchecker (official kernel replay)}{\badge{BadgeAmber}{stale only}}\layeritem{hardened}{safe\_verify (hardened screen)}{\badge{BadgeAmber}{stale only}}\layeritem{decorrelated}{lean4export + nanoda (decorrelated kernel)}{\badge{BadgeGreen}{current}}}
\sourcenote{\badge{BadgeGray}{comparator: none} no fixture in the three-implementation comparator harness (compiled / kernel-interpreted / reference-interpreted) exercises this rule.}
\ruleblock{HOL4 statement}
\sourcenote{theory \texttt{bern\_tr\_catdata2}, HOL4 development script \texttt{\texttt{bern\_tr\_catdata2Script.sml}} \badge{BadgeGreen}{script bound}}
\holhead{Theorem prelunndata\_receipt}
\begin{holsrc}
Theorem prelunndata_receipt:
  ruledataOk [] cd2Hrelu
    <| pre := FTt ; post := FAtom BowGe (TmVar 4) (TmConst 0) |>
    (RdReluNN (PReluNND 3 4)) ∧
  ¬ ruledataOk [] hnEmpty
      <| pre := FTt ; post := FAtom BowGe (TmVar 4) (TmConst 0) |>
      (RdReluNN (PReluNND 3 4))
\end{holsrc}
\ruleblock{HOL verifier receipts}
\badge{BadgeBlue}{documentary} \texttt{hol-strict} CLEAN: theory \texttt{bern\_tr\_catdata2}, 272
theorems, 0 cheat-tagged, 0 nonstandard axioms (whitelist:
\texttt{DISK\_THM}; the four \texttt{bool} axioms) --- verification
run of 2026-07-22 on an isolated build host, against a
content-addressed checkout of the exact sources cited here (git blob
\texttt{cd3c1e95\ldots}). The proof-object stratum is ported as one
mutual datatype family with a mutual wellformedness predicate,
exactly as the Lean mutual inductive block; this premise-free
rule\textquotesingle s arm states its entire checkable content ---
the family linkage plus the nonnegativity conclusion equation ---
and the linkage lookup into the hypernet\textquotesingle s
application list has a ported realization (\texttt{reluLinkedB},
theory \texttt{bern\_tr\_catdata}, 252 theorems, 0 offences, same
run). The receipt theorem cited here is the kernel-evaluated
accept/refusal pair for the arm. Honest residue: the semantic
soundness twin (the Lean development\textquotesingle s
\texttt{sound\_reluNN}) is not yet ported; the payload/checker layer
above is.

\sourcenote{HOL receipts are documentary: quoted results of \texttt{hol-strict} runs
and census rows from the HOL4 development\textquotesingle s build
evidence, not minted receipt artefacts. A HOL receipt store (the Lean
store\textquotesingle s twin) is named future work.
}
\ruleblock{Proof}

\sourcenote{Genre: admissible rule (tier 2). The family machinery is
stated in the \rulename{LiRPA} section's family template. This rule is
premise-free: $\Phi$ is arbitrary and never consulted (the
\rulename{Norm} convention).}

\begin{brecipe}[the \rulename{ReLU-NN} instance --- the full unpacking]
\label{rec:relunn-instance}
The node:
\[
\inferrule*[Right=\rulename{ReLU-NN}]{
  \ReLU_{\{h\}\to\{y\}} \in H
}{
  \Phi \vdash y \geq 0
}
\qquad\rightsquigarrow\qquad
\text{case analysis on }
(h \leq 0 \wedge y = 0) \vee (0 \leq h \wedge y = h),
\]
and both branches discharge into catalogue nodes --- the family's
cleanest unpacking. \emph{Inactive branch} (context carries $y = 0$):
\begin{enumerate}
\item $\vdash 0 \geq 0$ --- \rulename{PConstExact}, the closed exact
  ground fact (the leaf that rule exists to state: no other catalogue
  rule concludes a closed ground fact);
\item $\Phi_{-} \vdash y = 0$ --- \rulename{Ax} + $\wedge$E from the
  branch context (the case disjunct enters as one ordered subtree; two
  $\wedge$E steps); then $\Phi_{-} \vdash 0 = y$ by \rulename{EqSym};
\item $\Phi_{-} \vdash y \geq 0$ --- \rulename{EqElim} along $0 = y$,
  transporting $\Psi[a \coloneqq 0] = (0 \geq 0)$ (step 1, weakened)
  to $\Psi[a \coloneqq y]$, where $\Psi = (a \geq 0)$.
\end{enumerate}
\emph{Active branch} (context carries $0 \leq h$ and $y = h$): the
context's sign atom is extracted (\rulename{Ax} + $\wedge$E) and
re-oriented to $h \geq 0$ by the bound node's polarity-flip reading
(\rulename{PBound}, reading (G): the single component $0 \leq h$ with
weight $1$, canonical polynomial $c = -h$, read as $e \geq 0$ with
$e = h$ --- one node of the toolkit whose longer chains the
\rulename{ReLU-Above} section walks); \rulename{EqSym} flips the
context's $y = h$ to $h = y$, and \rulename{EqElim} along $h = y$ (an
equality between two wires of the node's one shared shape) transports
it to $y \geq 0$. $\vee$E in its context-carrying form --- today a
satellite theorem, not a core constructor (the split rule's section
prices exactly this), so what follows is the intended elaboration,
emittable when that constructor lands --- closes the case
analysis, and \rulename{Cut} stitches the sign-split fact into place.
\end{brecipe}

\begin{bremark}[the shape scope of the displayed derivation]
\label{rem:relunn-shapescope}
The realized \rulename{EqElim} carries a broadcast-shape side
condition (its section documents the counterexample that forced it):
both sides of the transported equality must share a broadcast shape.
Step 3's equality $0 = y$ has a rank-$0$ constant on one side, so the
displayed transport type-checks where the output wire is rank-$0$ ---
at tensor-shaped instances that step is stopped by the gate, and the family's
realized warrant does not route through the derivation at all: it is
the direct lemma \texttt{sound\_reluNN}, shape-general, via the
linkage extraction. (Wire-to-wire transports such as the active
branch's $h = y$ are unaffected: the linkage forces one shared shape.)
The derivation above is thus the intended admissibility story at
scalar instances --- intended, since its case-closing step awaits the
constructor named at its step --- and the
reduction-to-verified-machinery story, the
\rulename{PConstExact} pattern stated in the family template, is
the standing warrant at every shape either way.
\end{bremark}

\begin{bprop}[semantic warrant]\label{prop:relunn-sound}
If the linkage holds, the sequent concluding $y \geq 0$ is
semantically valid over any premise formula $\Phi$.
\end{bprop}
\begin{proof}
Fix an admissible valuation $\sigma$ ($\sigma \models \Phi$ is taken
but never consulted --- the rule is premise-free). The linkage
extraction (\texttt{reluLinked\_facts}) gives the output tensor and,
per entry, either $y = 0$ or $y = p$ with $0 \leq p$; in both cases
$0 \leq y$. The conclusion atom compares the output tensor entrywise
against the rank-$0$ constant $0$. That is \texttt{sound\_reluNN},
step for step.
\end{proof}

\begin{bremark}[the only hypothesis]
The linkage is the rule's sole side condition --- the conclusion-shape
check beside it is a reconstruction pin, not a droppable hypothesis:
the checker rebuilds the atom $y \geq 0$ rather than trusting it, so
``dropping'' the pin merely re-addresses the rule at a conclusion it
never concluded. Dropping the linkage is the whole failure mode: an
arbitrary wire asserted nonnegative against an arbitrary premise. Its positive-size gate additionally keeps the
vacuous corner unciteable --- a size-zero node satisfies every
elementwise atom trivially, and a fact no valuation can falsify must
not enter a derivation (the family template's hygiene remark).
\end{bremark}

\begin{anchors}
Recipe \ref{rec:relunn-instance} warrant & \texttt{sound\_reluNN}\\
family linkage extraction & \texttt{reluLinked\_facts} (see the \rulename{LiRPA} section)\\
ground-fact leaf (step 1) & \texttt{sound\_pConstExact} (the \rulename{PConstExact} section)\\
payload structure & \texttt{PReLUNNData} (two wires, linkage --- no premise field)\\
\end{anchors}


\rulesection{23}{ReLU-Above}{admissible rule}
\label{blcatrule-ReLUAbove}
\ruleblock{The rule (Steele metanotation)}
\sourcenote{as realized in the mechanized development; source row: \texttt{\texttt{rules.toml}} --- \emph{ReLU-Above}. The metanotation statement is typeset by hand for notational quality; its mechanized ground truth is the Lean statement below, extracted from the sources at every build of this book and hash-bound. The rule \emph{roster} is machine-checked (every mechanized rule tag round-trips to exactly one catalogue row); the statement \emph{text} is editorial, checked by eye against the extracted Lean beside it.}
\begin{mathpar}
\inferrule*[Right=\rulename{ReLU-Above}]{
  \ReLU_{\{h\}\to\{y\}} \in H
}{
  \Phi \vdash y \geq h
}
\end{mathpar}

\par\medskip
\rulename{ReLU-Above} is the identity minorant: the output wire of a
cited ReLU application dominates its preactivation, $\Phi \vdash y
\geq h$, with no proof premises — $\Phi$ arbitrary, the premise-free
convention. Together with \rulename{ReLU-NN} it gives the exact
pointwise lower envelope of the activation as two always-available
atoms ($y \geq 0$ and $y \geq h$); a producer needs no box knowledge
at all to emit either. Upper bounds are different in kind — no affine
upper bound on the activation exists without a box, which is why the
upper side of the family lives in \rulename{LiRPA} behind two box
premises.

What the checker recomputes: the linkage — the hypernet really
contains a packaged ReLU application at exactly these ports, of
positive total size (the family\textquotesingle s port-lookup side
condition, displayed in the \rulename{LiRPA} section) — and that the
conclusion is literally the atom $y \geq h$, elementwise over the
node\textquotesingle s tensor shape.

The linkage is the only hypothesis; without it, two arbitrary wires
are asserted ordered. Entrywise the fact is immediate from the
node\textquotesingle s two-case semantics: on the inactive branch $y =
0 \geq h$ (there $h \leq 0$), on the active branch $y = h$.

\cataloguenote{The proof section uses this rule\textquotesingle s derivation as the
family\textquotesingle s worked example of atom-orientation plumbing:
its inactive branch discharges into catalogue nodes through a chain of
a bound-node polarity flip, a \rulename{Norm} identity, an
\rulename{ArithAtom} fold, and two \rulename{EqElim} transports — a
faithful specimen of what ``one node unpacks to many\textquotesingle
\textquotesingle{} means in practice.
}
\ruleblock{Lean statement}
\par\medskip
The warrant consumes the family linkage extraction
(\texttt{reluLinked\_facts}, displayed in the \rulename{LiRPA}
section): entrywise, $y = 0 \geq h$ on the inactive case and $y = h$
on the active one. The payload structure, \texttt{PReLUAboveData}
--- quoted in full below --- carries no premise field.

\declhead{BernsteinLogic.Soundness}{sound\_reluAbove}{\badge{BadgeGreen}{statement bound}}
\begin{leansrc}
/-- ReLU-Above: `Φ ⊢ y ≥ h`, premise-free — entrywise, `y = 0 ≥ h` on
the nonpositive case and `y = h` on the nonnegative one. -/
theorem sound_reluAbove {H : Hypernet} {seq : Sequent} {preact post : Wire}
    (hlink : reluLinked H preact post)
    (hpost : seq.post = .atom .ge (.var post) (.var preact)) :
    seq.semanticallyValid H
\end{leansrc}
\declhead{BernsteinLogic.Core.Catalogue}{PReLUAboveData}{\badge{BadgeGreen}{definition bound}}
\begin{leansrc}
inductive PReLUAboveData : Hypernet → Sequent → Type where
  | mk (H : Hypernet) (seq : Sequent)
       (preact post : Wire)
       (link : reluLinked H preact post)
       (conclPost : seq.post = .atom .ge (.var post) (.var preact))
       : PReLUAboveData H seq
\end{leansrc}
\ruleblock{Correspondence (Steele display $\leftrightarrow$ Lean)}
\badge{BadgeGreen}{reviewer-certified against the displayed texts}
\sourcenote{The dictionary below is a reviewer-certified statement of how the mechanized statements above interpret the displayed rule: display metavariable $\mapsto$ Lean binder, each display hypothesis $\mapsto$ the premise or side condition carrying it, the display conclusion $\mapsto$ the goal --- with any deliberate difference in width between display and mechanization named outright. Certification is recorded as a hash triple over (the displayed cell, the extracted statement texts, this dictionary\textquotesingle s own text); if any of the three changes, this block turns amber until a reviewer re-certifies. The correspondence claim itself remains editorial --- reviewer-grade, not machine-checked; its machine-checked successor types each soundness statement against the rule as a structured object, and supersedes this block per rule as it lands.}
Judgment frame as in \rulename{Ax}; linkage premise --- including the
positive-size narrowing hidden under the displayed $\in$ --- exactly
as in the \rulename{LiRPA} dictionary, the family preamble.
\begin{itemize}
\item $h \mapsto$ \texttt{preact}; $y \mapsto$ \texttt{post};
  $\ReLU_{\{h\}\to\{y\}} \in H$ $\mapsto$ \texttt{hlink} $:$
  \texttt{reluLinked H preact post} (payload alias:
  \texttt{PReLUAboveData.link}, same proposition).
\item premise-freeness $\mapsto$ silence, the \rulename{PConstExact}
  representation: the precondition component of \texttt{seq :
  Sequent} is unconstrained by any hypothesis, so the
  statement holds under arbitrary premises $\Phi$ --- the
  display\textquotesingle s universality over the preactivation
  rides the linkage, as at
  \rulename{ReLU-NN}.
\item conclusion $\Phi \vdash y \geq h$ $\mapsto$ the goal under
  \texttt{hpost} $:$ \texttt{seq.post = .atom .ge (.var post) (.var
  preact)} (payload alias: \texttt{conclPost}, same equation) ---
  output wire left, preactivation right, as displayed;
  the only two-wire atom among the family\textquotesingle s
  premise-free conclusions, so this dictionary is where the shared
  shape of the two tensors (from the linkage extraction) is
  load-bearing rather than convenient.
\end{itemize}
Width note: certified atoms elementwise over the cited
application\textquotesingle s tensor shape (batch semantics, as at
\rulename{LiRPA}). Beyond the family\textquotesingle s linkage gate,
neither side is wider: both hypotheses of \texttt{sound\_reluAbove}
are consumed, and nothing else is stated.

\ruleblock{Lean verifier receipts}
\sourcenote{module \texttt{BernsteinLogic.Soundness}; bound build \texttt{a20d90e66c13\ldots} (bound from receipt a9c0e34d3e3c9c8c, 2026-08-03T21:30:36Z). A receipt records that named checking layers accepted a build: (module, compiled-artifact SHA-256, environment SHA-256, checker layers), minted when the checking gate ran; the store receipts attesting the bound build are listed with the layers each run credited. Nothing here asks the reader to trust the store: every layer is an independently runnable program over the quoted sources --- rebuild the module and re-run the checkers to reproduce the verdict. Receipts attesting superseded builds credit no layer and are omitted from this rendering (the count of omissions, if any, is noted below).}
\begin{receipttable}
\receiptrow{a9c0e34d3e3c9c8c}{2026-08-03T21:30:36Z}{build+decorrelated}{\badge{BadgeGreen}{current}}
\end{receipttable}
\sourcenote{21 further receipt(s) in the store attest superseded builds of this module and are omitted from this rendering.}
\sourcenote{module \texttt{BernsteinLogic.Core.Catalogue}; bound build \texttt{cbc92e976a1d\ldots} (bound from receipt 5abdb3589707b5df, 2026-08-03T21:30:33Z). A receipt records that named checking layers accepted a build: (module, compiled-artifact SHA-256, environment SHA-256, checker layers), minted when the checking gate ran; the store receipts attesting the bound build are listed with the layers each run credited. Nothing here asks the reader to trust the store: every layer is an independently runnable program over the quoted sources --- rebuild the module and re-run the checkers to reproduce the verdict. Receipts attesting superseded builds credit no layer and are omitted from this rendering (the count of omissions, if any, is noted below).}
\begin{receipttable}
\receiptrow{24804e66edba0bc9}{2026-08-03T20:25:09Z}{build+decorrelated}{\badge{BadgeGreen}{current}}
\receiptrow{5abdb3589707b5df}{2026-08-03T21:30:33Z}{build+decorrelated}{\badge{BadgeGreen}{current}}
\end{receipttable}
\sourcenote{18 further receipt(s) in the store attest superseded builds of this module and are omitted from this rendering.}
\layerline{\layeritem{build}{lake build (elaboration)}{\badge{BadgeGreen}{current}}\layeritem{replay}{leanchecker (official kernel replay)}{\badge{BadgeAmber}{stale only}}\layeritem{hardened}{safe\_verify (hardened screen)}{\badge{BadgeAmber}{stale only}}\layeritem{decorrelated}{lean4export + nanoda (decorrelated kernel)}{\badge{BadgeGreen}{current}}}
\sourcenote{\badge{BadgeGray}{comparator: none} no fixture in the three-implementation comparator harness (compiled / kernel-interpreted / reference-interpreted) exercises this rule.}
\ruleblock{HOL4 statement}
\sourcenote{theory \texttt{bern\_tr\_catdata2}, HOL4 development script \texttt{\texttt{bern\_tr\_catdata2Script.sml}} \badge{BadgeGreen}{script bound}}
\holhead{Theorem preluabovedata\_receipt}
\begin{holsrc}
Theorem preluabovedata_receipt:
  ruledataOk [] cd2Hrelu
    <| pre := FTt ; post := FAtom BowGe (TmVar 4) (TmVar 3) |>
    (RdReluAbove (PReluAboveD 3 4)) ∧
  ¬ ruledataOk [] cd2Hrelu
      <| pre := FTt ; post := FAtom BowGe (TmVar 3) (TmVar 4) |>
      (RdReluAbove (PReluAboveD 3 4))
\end{holsrc}
\ruleblock{HOL verifier receipts}
\badge{BadgeBlue}{documentary} \texttt{hol-strict} CLEAN: theory \texttt{bern\_tr\_catdata2}, 272
theorems, 0 cheat-tagged, 0 nonstandard axioms (whitelist:
\texttt{DISK\_THM}; the four \texttt{bool} axioms) --- verification
run of 2026-07-22 on an isolated build host, against a
content-addressed checkout of the exact sources cited here (git blob
\texttt{cd3c1e95\ldots}). The proof-object stratum is ported as one
mutual datatype family with a mutual wellformedness predicate,
exactly as the Lean mutual inductive block; this premise-free
rule\textquotesingle s arm states its entire checkable content ---
the family linkage plus the output-dominates-preactivation
conclusion equation --- and the linkage lookup into the
hypernet\textquotesingle s application list has a ported realization
(\texttt{reluLinkedB}, theory \texttt{bern\_tr\_catdata}, 252
theorems, 0 offences, same run). The receipt theorem cited here is
the kernel-evaluated accept/refusal pair for the arm. Honest
residue: the semantic soundness twin (the Lean
development\textquotesingle s \texttt{sound\_reluAbove}) is not yet
ported; the payload/checker layer above is.

\sourcenote{HOL receipts are documentary: quoted results of \texttt{hol-strict} runs
and census rows from the HOL4 development\textquotesingle s build
evidence, not minted receipt artefacts. A HOL receipt store (the Lean
store\textquotesingle s twin) is named future work.
}
\ruleblock{Proof}

\sourcenote{Genre: admissible rule (tier 2). The family machinery is
stated in the \rulename{LiRPA} section's family template. This rule is
premise-free: $\Phi$ is arbitrary and never consulted.}

\begin{brecipe}[the \rulename{ReLU-Above} instance --- worked plumbing]
\label{rec:reluabove-instance}
The node:
\[
\inferrule*[Right=\rulename{ReLU-Above}]{
  \ReLU_{\{h\}\to\{y\}} \in H
}{
  \Phi \vdash y \geq h
}
\qquad\rightsquigarrow\qquad
\text{case analysis on }
(h \leq 0 \wedge y = 0) \vee (0 \leq h \wedge y = h).
\]
We walk the \emph{inactive} branch in full --- the family's honest
specimen of what ``one node unpacks to many'' means in catalogue
practice. Context $\Phi_{-}$ carries $h \leq 0$ and $y = 0$; the goal
is $y \geq h$.
\begin{enumerate}
\item $\Phi_{-} \vdash h \leq 0$ --- \rulename{Ax} + $\wedge$E (the
  case disjunct enters as one ordered subtree; two $\wedge$E steps).
\item $\Phi_{-} \vdash -h \geq 0$ --- the bound node's polarity-flip
  reading (\rulename{PBound}, reading (G)): the combination is the
  single component $h \leq 0$ with weight $1$, canonical polynomial
  $c = h$, read as $e \geq 0$ with $e = -h$.
\item $\Phi_{-} \vdash (-h) = (0 - h)$ --- \rulename{Norm}: equal
  monomial canonical forms over $\Q$, same wires, shared broadcast
  shape.
\item $\Phi_{-} \vdash (0 - h) \geq 0$ --- \rulename{EqElim} along
  step 3's equality (both sides shaped alike --- the transport gate is
  met), rewriting step 2's atom.
\item $\Phi_{-} \vdash 0 \geq h$ --- \rulename{ArithAtom} (fold): from
  $(t_1 - t_2) \bowtie 0$ to $t_1 \bowtie t_2$ at $t_1 = 0$, $t_2 =
  h$, $\bowtie = {\geq}$.
\item $\Phi_{-} \vdash 0 = y$ --- \rulename{EqSym} on the context's
  $y = 0$; then $\Phi_{-} \vdash y \geq h$ --- \rulename{EqElim} along
  $0 = y$ at $\Psi = (a \geq h)$.
\end{enumerate}
The \emph{active} branch (context $0 \leq h$, $y = h$) derives $h \geq
h$ by the same pattern, spelled at the same grain: \rulename{Norm}
gives $(h - h) = 0$, whence $0 = (h - h)$ by \rulename{EqSym};
\rulename{PConstExact} supplies the ground fact $0 \geq 0$, weakened
into the branch context; \rulename{EqElim} along $0 = (h - h)$
transports it to $(h - h) \geq 0$, and \rulename{ArithAtom} folds that
to $h \geq h$ --- then \rulename{EqSym} on the context's $y = h$ and
\rulename{EqElim} along $h = y$ (wire-to-wire, one shared shape) close
the goal. Six arithmetic-plumbing nodes in the inactive branch ---
\rulename{PBound}, \rulename{Norm}, \rulename{EqElim},
\rulename{ArithAtom}, \rulename{EqSym}, \rulename{EqElim}; the
extraction's \rulename{Ax} and $\wedge$E steps aside --- to move one
minus sign across a comparator: this is why the family rules exist as
single nodes, and why the wire format prices them.
\end{brecipe}

\begin{bremark}[shape scope]
Step 6 of the inactive branch and the active branch's route through
$(h - h) = 0$ transport along equalities with a rank-$0$ constant
side, so the displayed derivation type-checks at rank-$0$ instances
(step 4's transport along $(-h) = (0 - h)$ is wire-shaped on both
sides --- the constant $0$ sits \emph{inside} the term $0 - h$ --- and
is unaffected); the shape-general warrant is the direct lemma, exactly
as \autoref{rem:relunn-shapescope} in the \rulename{ReLU-NN}
section states for the whole family.
\end{bremark}

\begin{bprop}[semantic warrant]\label{prop:reluabove-sound}
If the linkage holds, the sequent concluding $y \geq h$ is
semantically valid over any premise formula $\Phi$.
\end{bprop}
\begin{proof}
Fix an admissible valuation $\sigma$. The linkage extraction
(\texttt{reluLinked\_facts}) gives both tensors at one shared shape
and, per entry, either $p \leq 0 \wedge y = 0$ --- whence $y = 0 \geq
p$ --- or $y = p$ --- whence $y \geq p$ reflexively. The conclusion
atom compares the two tensors entrywise at their shared shape. That is
\texttt{sound\_reluAbove}, step for step.
\end{proof}

\begin{bremark}[the exact lower envelope, premise-free]
Together with \rulename{ReLU-NN} this rule states the exact pointwise
lower envelope of the activation, $y \geq \max(0, h)$, as two atoms a
producer can emit with \emph{no} box knowledge whatsoever --- exact in
that the bound is attained: $\max(0, h)$ \emph{is} the activation's
value, so no sound lower bound can improve on the pair. The
asymmetry with the upper side is structural, not an artefact: no
affine function dominates $\max(0, h)$ on all of $\Q$ (a candidate
with positive slope dives below $0$ as $h \to -\infty$; one with
slope $\leq 0$ falls below $h$ as $h \to +\infty$) --- so upper
bounds require a box and live in \rulename{LiRPA}.
\end{bremark}

\begin{anchors}
Recipe \ref{rec:reluabove-instance} warrant & \texttt{sound\_reluAbove}\\
family linkage extraction & \texttt{reluLinked\_facts} (see the \rulename{LiRPA} section)\\
polarity-flip reading (step 2) & \texttt{sound\_pBound} (the \rulename{PBound} section)\\
ground-fact leaf (active branch) & \texttt{sound\_pConstExact} (the \rulename{PConstExact} section)\\
payload structure & \texttt{PReLUAboveData} (two wires, linkage --- no premise field)\\
\end{anchors}


\rulesection{24}{Cell-In-Interior}{admissible rule}
\label{blcatrule-CellInInterior}
\ruleblock{The rule (Steele metanotation)}
\sourcenote{as realized in the mechanized development; source row: \texttt{\texttt{rules.toml}} --- \emph{Cell-In-Interior}. The metanotation statement is typeset by hand for notational quality; its mechanized ground truth is the Lean statement below, extracted from the sources at every build of this book and hash-bound. The rule \emph{roster} is machine-checked (every mechanized rule tag round-trips to exactly one catalogue row); the statement \emph{text} is editorial, checked by eye against the extracted Lean beside it.}
\begin{mathpar}
\inferrule*[Right=\rulename{Cell-In-Interior}]{
  \Phi \vdash \bigwedge_{v} \lambda^{C_i}_v(x) \bowtie^{\circ}_v 0 \ \text{(strict barycentric interior atoms, flat vertex order; } < \text{ for extended } v\text{, } > \text{ for core } v\text{)}
  \\
  \Cell[\Sigma]_{\{x\}\to\{b_1,\ldots,b_N\}} \in H
}{
  \Phi \vdash b_i = 1 \wedge \bigwedge_{j \neq i} b_j = 0
}
\end{mathpar}

\par\medskip
\rulename{Cell-In-Interior} is one of the two rules that read the cell
primitive --- the calculus\textquotesingle s only source of disjunction.
A cell application $\Cell[\Sigma]_{\{x\}\to\{b_1,\ldots,b_N\}}$ tabulates
which cell of a simplicial decomposition its input point lies in, as
one-hot indicator outputs; piecewise definitions (a piecewise-linear
activation, a case-defined dynamics) enter the hypernet through it. This
rule discharges the case a proof is actually in: from a premise placing
$x$ \emph{strictly inside} cell $C_i$ --- a conjunction of strict
barycentric inequalities in the pinned flat vertex order --- it concludes
the full one-hot vector, $b_i = 1$ and $b_j = 0$ for every $j \neq i$. A
producer reaches for it after region reasoning (the region introduction
rules, \rulename{Subdivide}) has confined the input to one
cell\textquotesingle s interior, and needs the indicator values as
arithmetic facts downstream --- typically to substitute a branch\textquotesingle s
closed form for the primitive\textquotesingle s output.

What the checker recomputes, given the payload (schema, cell index, cited
application, premise proof): that the cited cell application is present
in the hypernet with the input wire typed at shape $[\mathrm{ambDim}]$
and every output at total size one; that the premise formula is
byte-for-byte the schema-derived strict-interior build for cell $i$
(barycentric atoms recomputed from the vertex coordinates --- no formula
text trusted from the wire); and that the conclusion is the recomputed
one-hot form. Soundness then rides the model: admissibility of the cell
application forces some cell choice $k$ with weak membership and the
one-hot outputs at $k$; the strict premise pins $k = i$ through the
schema\textquotesingle s interior-disjointness well-formedness.

Each hypothesis is load-bearing. Weaken strictness to weak membership
and a boundary point of two cells admits a valuation one-hotting the
\emph{other} cell --- the conclusion fails. Drop interior-disjointness
(distinct cells of the enumeration have disjoint interiors --- which the
complex\textquotesingle s simplicial conformity does NOT imply: conformity never
sees the extension masks and permits duplicate cell entries) and the
model may one-hot an overlapping cell. Drop the output typing and the
model\textquotesingle s entry-level one-hot fact decouples from the whole-wire
equality atoms (a zero-size output satisfies every atom vacuously; a
multi-entry output can falsify the atom away from entry zero). Drop the
input-shape tie and the premise\textquotesingle s constraint on the input tensor
no longer speaks about the model\textquotesingle s evaluation point. The proof
section walks these in order.

\cataloguenote{The checkable fragment is 1-dimensional (two-vertex cell supports on
distinct coordinates) --- the same fragment as the region machinery;
everything beyond is refused, never approximated. The intended
catalogue elaboration (case analysis over the cell axiom\textquotesingle s
disjunction) awaits the context-carrying case-analysis constructor; the
shipped warrant is a direct semantic lemma, the degenerate-rewrite
shape the \rulename{PConstExact} section defines.
}
\ruleblock{Lean statement}
\declhead{BernsteinLogic.Soundness}{sound\_cellInInterior}{\badge{BadgeGreen}{statement bound}}
\begin{leansrc}
/-- Cell-In-Interior: the strict-interior premise pins the model's cell
choice to whichI (via `cellIntDisjoint`), and the sp-fact's one-hot
valuation IS the conclusion. -/
theorem sound_cellInInterior {H : Hypernet} {seq premise : Sequent}
    {schema : SimpC} {input : Wire} {outputs : List Wire}
    {whichI : Fin schema.cells.length}
    (ih : premise.semanticallyValid H)
    (happ : PrimApp.cell schema input outputs ∈ H.applications)
    (hinshape : ∃ τ, H.wireType input = some τ ∧
      τ.shape = [schema.ambDim])
    (houts : ∀ o ∈ outputs, ∃ τ, H.wireType o = some τ ∧
      Shape.totalSize τ.shape = 1)
    (hpre : seq.pre = premise.pre)
    (hint : cellStrictInteriorForm schema (.var input) whichI premise.post)
    (hpost : cellOneHotForm schema outputs whichI seq.post) :
    seq.semanticallyValid H
\end{leansrc}
\declhead{BernsteinLogic.Soundness}{cellStrictInteriorBuild\_spec}{\badge{BadgeGreen}{statement bound}}
\begin{leansrc}
/-- Extraction for the strict-interior premise: the build succeeds only
over a 1-D complex with a two-vertex cell mask on distinct coordinates,
and the formula is the pinned pair of strict barycentric atoms. -/
theorem cellStrictInteriorBuild_spec {schema : SimpC} {t : Term}
    {k : Fin schema.cells.length} {F : Formula}
    (h : cellStrictInteriorForm schema t k F) :
    schema.ambDim = 1 ∧
    ∃ (a b : Fin schema.nVertices) (V₀ V₁ : ℚ),
      (List.finRange schema.nVertices).filter
        (fun v => (schema.cells.get k).1 v) = [a, b] ∧
      (∀ j, schema.coord a j = V₀) ∧
      (∀ j, schema.coord b j = V₁) ∧
      V₀ ≠ V₁ ∧
      F = .and
        (.atom (if (schema.cells.get k).2 a then .lt else .gt)
          (.add (.const (V₁ / (V₁ - V₀))) (.mul (.const (-1 / (V₁ - V₀))) t))
          (.const 0))
        (.atom (if (schema.cells.get k).2 b then .lt else .gt)
          (.add (.const (-V₀ / (V₁ - V₀))) (.mul (.const (1 / (V₁ - V₀))) t))
          (.const 0))
\end{leansrc}
\declhead{BernsteinLogic.Soundness}{maskRegionStrict\_pair\_intro}{\badge{BadgeGreen}{statement bound}}
\begin{leansrc}
/-- Construction: rational weights of strict signs assemble a
`MaskRegionStrict` witness over a two-vertex 1-D support (the strict
twin of `maskRegion_pair_intro`). Needs `hscope : ext ⊆ mask` — the
cellsWF invariant — because the strict region demands lam < 0 at every
ext vertex, and support vertices are the only ones with nonzero lam
(the weak twin's 0 ≤ 0 escape is unavailable at strict signs). -/
theorem maskRegionStrict_pair_intro {C : SimpC}
    {mask ext : Fin C.nVertices → Bool}
    {x : Fin C.ambDim → ℚ} {a b : Fin C.nVertices} {V₀ V₁ : ℚ}
    (hfil : (List.finRange C.nVertices).filter (fun v => mask v) = [a, b])
    (hscope : ∀ v, ext v = true → mask v = true)
    (hVa : ∀ j, C.coord a j = V₀) (hVb : ∀ j, C.coord b j = V₁)
    (hne : V₀ ≠ V₁) (hdim : C.ambDim = 1) (j0 : Fin C.ambDim)
    (hw0t : ext a = true → (V₁ - x j0) / (V₁ - V₀) < 0)
    (hw0f : ext a = false → 0 < (V₁ - x j0) / (V₁ - V₀))
    (hw1t : ext b = true → (x j0 - V₀) / (V₁ - V₀) < 0)
    (hw1f : ext b = false → 0 < (x j0 - V₀) / (V₁ - V₀)) :
    MaskRegionStrict C mask ext x
\end{leansrc}
\declhead{BernsteinLogic.Soundness}{holds\_cellOneHot}{\badge{BadgeGreen}{statement bound}}
\begin{leansrc}
/-- The one-hot conclusion holds at any valuation whose output entries
realize the one-hot pattern at k (unit-size output typing carries the
entry facts to full wire reads). -/
theorem holds_cellOneHot {H : Hypernet} (σ : H.Valuation) {schema : SimpC}
    {outputs : List Wire} {k : Fin schema.cells.length} {F : Formula}
    (hbuild : cellOneHotForm schema outputs k F)
    (houts : ∀ o ∈ outputs, ∃ τ, H.wireType o = some τ ∧
      Shape.totalSize τ.shape = 1)
    (hsem : ∀ j : Fin outputs.length,
      σ.entryFlat (outputs.get j) 0
        = some (.fin (if j.val = k.val then 1 else 0))) :
    Formula.holds σ F
\end{leansrc}
\par\medskip
\textbf{The recomputed forms and the payload.}\enspace The premise and
conclusion formulas are never read off the certificate as text: the
checker rebuilds both from the schema --- the strict-interior
conjunction from the 1-D barycentric coordinate terms
(\texttt{baryLambdas1D}, offset-first affine normal forms over the
two-vertex support), the conclusion from the one-hot pattern --- and the
payload\textquotesingle s form fields assert equality with the rebuilt
formulas. A skeptical reader should confirm by eye: the builds succeed
only over a 1-dimensional complex with exactly two support vertices at
distinct coordinates (everything else reads \texttt{none} and the rule
refuses); the strict comparators are $<$ at extended vertices and $>$
at core vertices; and the payload carries the port-typing obligations
(input at $[\mathrm{ambDim}]$, outputs at total size one) explicitly.

\declhead{BernsteinLogic.Core.Catalogue}{baryLambdas1D}{\badge{BadgeGreen}{definition bound}}
\begin{leansrc}
/-- The barycentric coordinate functions of a 1-dimensional simplex, as
canonical affine Terms in `t` (pinned shape `offset + slope·t`; see the
original for the full derivation notes and D2 scope). -/
def baryLambdas1D (cx : SimpC) (dContains : Fin cx.nVertices → Bool)
    (t : Term) : Option (List (Fin cx.nVertices × Term)) :=
  if h : cx.ambDim = 1 then
    let j0 : Fin cx.ambDim := ⟨0, by omega⟩
    match (List.finRange cx.nVertices).filter (fun v => dContains v) with
    | [a, b] =>
      let V₀ := cx.coord a j0
      let V₁ := cx.coord b j0
      if V₀ = V₁ then none
      else
        let den := V₁ - V₀
        let lam := fun (c s : Rat) => Term.add (.const c) (.mul (.const s) t)
        some [(a, lam (V₁ / den) (-1 / den)), (b, lam (-V₀ / den) (1 / den))]
    | _ => none
  else none
\end{leansrc}
\declhead{BernsteinLogic.Core.Catalogue}{cellStrictInteriorBuild}{\badge{BadgeGreen}{definition bound}}
\begin{leansrc}
/-- Strict-interior conjunction (flat vertex-order convention, strict
bowties). -/
def cellStrictInteriorBuild (schema : SimpC) (t : Term)
    (k : Fin schema.cells.length) : Option Formula := do
  let c := schema.cells[k.val]'(k.isLt)
  let lams ← baryLambdas1D schema c.1 t
  match lams.map (fun (v, lt) =>
    Formula.atom (if c.2 v then .lt else .gt) lt (.const 0)) with
  | [x, y] => some (.and x y)
  | _ => none
\end{leansrc}
\declhead{BernsteinLogic.Core.Catalogue}{cellStrictInteriorForm}{\badge{BadgeGreen}{definition bound}}
\begin{leansrc}
def cellStrictInteriorForm (schema : SimpC) (t : Term)
    (k : Fin schema.cells.length) (F : Formula) : Prop :=
  cellStrictInteriorBuild schema t k = some F
\end{leansrc}
\declhead{BernsteinLogic.Core.Catalogue}{cellOneHotBuild}{\badge{BadgeGreen}{definition bound}}
\begin{leansrc}
/-- One-hot conclusion: `(b_k = 1) ∧ chainAnd [(b_j = 0) | j ≠ k]`. -/
def cellOneHotBuild (schema : SimpC) (outputs : List Wire)
    (k : Fin schema.cells.length) : Option Formula :=
  if outputs.length = schema.cells.length then
    let bk := Formula.atom .eq (.var (outputs.getD k.val 0)) (.const 1)
    let others := ((List.range outputs.length).filter (· ≠ k.val)).map
      fun j => Formula.atom .eq (.var (outputs.getD j 0)) (.const 0)
    some (.and bk (Formula.chainAnd others))
  else none
\end{leansrc}
\declhead{BernsteinLogic.Core.Catalogue}{cellOneHotForm}{\badge{BadgeGreen}{definition bound}}
\begin{leansrc}
def cellOneHotForm (schema : SimpC) (outputs : List Wire)
    (k : Fin schema.cells.length) (F : Formula) : Prop :=
  cellOneHotBuild schema outputs k = some F
\end{leansrc}
\declhead{BernsteinLogic.Core.Catalogue}{PCellInIntData}{\badge{BadgeGreen}{definition bound}}
\begin{leansrc}
inductive PCellInIntData : Hypernet → Sequent → Type where
  | mk (H : Hypernet) (seq : Sequent)
       (premise : Sequent) (premiseProof : ProofObject H premise)
       (schema : SimpC) (input : Wire) (outputs : List Wire)
       (whichI : Fin schema.cells.length)
       (appInH : PrimApp.cell schema input outputs ∈ H.applications)
       (inputShape : ∃ τ, H.wireType input = some τ ∧
                       τ.shape = [schema.ambDim])
       (outputsScalar : ∀ o ∈ outputs, ∃ τ, H.wireType o = some τ ∧
                          Shape.totalSize τ.shape = 1)
       (preEq  : seq.pre = premise.pre)
       (interiorPremise : cellStrictInteriorForm schema (.var input) whichI
                            premise.post)
       (postOneHot : cellOneHotForm schema outputs whichI seq.post)
       : PCellInIntData H seq
\end{leansrc}
\par\medskip
\textbf{The semantic side: strict regions and interior
disjointness.}\enspace \texttt{MaskRegionStrict} is the strict-sign twin
of the weak region membership the model\textquotesingle s cell fact
carries; \texttt{SimpC.CellIntDisjoint} is the well-formedness the
pinning argument consumes, stated over cell \emph{indices} (so duplicate
entries are excluded) at rational points (the model\textquotesingle s
cell inputs are rational). It is consumed as a conjunct of the cell
application\textquotesingle s model-level fact: an application over a
non-disjoint enumeration admits no valuation --- hypothesis-side
vacuity, sound but weak --- rather than an unsound conclusion. The
predicate is a proposed schema amendment realized provisionally; the
roster page records its status.

\declhead{BernsteinLogic.Catalogue}{MaskRegionStrict}{\badge{BadgeGreen}{definition bound}}
\begin{leansrc}
/-- Strict-interior membership: the same weights with STRICT signs on
every mask vertex (< 0 extended, > 0 core) — the region the
Cell-In-Interior premise certifies. -/
def MaskRegionStrict (C : SimpC) (mask ext : Fin C.nVertices → Bool)
    (x : Fin C.ambDim → ℚ) : Prop :=
  ∃ lam : Fin C.nVertices → ℝ,
    (∀ v, mask v = false → lam v = 0) ∧
    finSum C.nVertices lam = 1 ∧
    (∀ j : Fin C.ambDim,
      finSum C.nVertices (fun v => lam v * (C.coord v j : ℝ)) = (x j : ℝ)) ∧
    (∀ v, ext v = true → lam v < 0) ∧
    (∀ v, mask v = true → ext v = false → 0 < lam v)
\end{leansrc}
\declhead{BernsteinLogic.Catalogue}{SimpC.CellIntDisjoint}{\badge{BadgeGreen}{definition bound}}
\begin{leansrc}
/-- PROPOSED-B2-DISJOINT (provisional): distinct cells of a schema's
enumeration have disjoint interiors — no rational point lies strictly
inside cell i and in the CLOSED region of a different cell k. This is
what the cell rules' admissibility derivation cites
("conformity: distinct cells have disjoint interiors"), and the
schema's `conformity` field does NOT deliver it: that field constrains
simplex-mask intersections only — it never sees the (Δ,S) extension
masks, and it permits literally duplicate cell entries, either of which
makes Cell-In-Interior UNSOUND (the model may one-hot a different
overlapping cell). Stated over INDICES so duplicates are excluded, at
ℚ points because the model's Cell inputs are rational. Σ± satisfies it
(the two rays share only the origin, strictly inside neither).

Lives here (not as a `SimpC` field) because the quantification is over
REAL weights — mathlib-side of the core/mathlib module split — and is
consumed as a
conjunct of the Cell sp-fact below: a cell application over a
non-disjoint schema admits NO valuation, the same "soundness errors
are excluded from every value set" reading that rejects NaN
executions (a hypothesis-side vacuity: sound, merely weak). -/
def SimpC.CellIntDisjoint (C : SimpC) : Prop :=
  ∀ i k : Fin C.cells.length, i ≠ k →
    ∀ x : Fin C.ambDim → ℚ,
      MaskRegionStrict C (C.cells.get i).1 (C.cells.get i).2 x →
      ¬ MaskRegion C (C.cells.get k).1 (C.cells.get k).2 x
\end{leansrc}
\ruleblock{Correspondence (Steele display $\leftrightarrow$ Lean)}
\badge{BadgeGreen}{reviewer-certified against the displayed texts}
\sourcenote{The dictionary below is a reviewer-certified statement of how the mechanized statements above interpret the displayed rule: display metavariable $\mapsto$ Lean binder, each display hypothesis $\mapsto$ the premise or side condition carrying it, the display conclusion $\mapsto$ the goal --- with any deliberate difference in width between display and mechanization named outright. Certification is recorded as a hash triple over (the displayed cell, the extracted statement texts, this dictionary\textquotesingle s own text); if any of the three changes, this block turns amber until a reviewer re-certifies. The correspondence claim itself remains editorial --- reviewer-grade, not machine-checked; its machine-checked successor types each soundness statement against the rule as a structured object, and supersedes this block per rule as it lands.}
Judgment frame as in \rulename{Ax}; sub-derivation premise as at
\rulename{Weaken}.
\begin{itemize}
\item $\Sigma \mapsto$ \texttt{schema : SimpC}; $i \mapsto$
  \texttt{whichI : Fin schema.cells.length}; $x \mapsto$ the term
  \texttt{.var input} --- the region term literally \emph{is} the
  input wire, as at \rulename{BernCone}; $b_1, \ldots, b_N \mapsto$
  \texttt{outputs : List Wire}.
\item the linkage $\Cell[\Sigma]_{\{x\}\to\{b_1,\ldots,b_N\}} \in H$
  $\mapsto$ \texttt{happ}, plain application-list membership ---
  contrast the \rulename{LiRPA} family\textquotesingle s size-gated
  lookup: here the port discipline rides two explicit typing
  hypotheses instead, \texttt{hinshape} (the input wire\textquotesingle
  s type has shape \texttt{[schema.ambDim]}) and \texttt{houts}
  (every output wire singleton-sized). Neither has a displayed row;
  the payload carries both verbatim (\texttt{PCellInIntData}\textquotesingle
  s port-typing fields, displayed above); they are
  load-bearing twice: the model\textquotesingle s
  sp-fact pins entry 0 of each output while the conclusion atom
  evaluates the whole wire, and the premise atoms constrain the input
  tensor, which ties to the model\textquotesingle s rational point
  only at the declared shape. The same undisplayed typing pack as at
  \rulename{Norm}.
\item premise $\Phi \vdash \bigwedge_v \lambda^{C_i}_v(x)
  \bowtie^{\circ}_v 0$ $\mapsto$ \texttt{ih} with \texttt{hint} $:$
  \texttt{cellStrictInteriorForm schema (.var input) whichI
  premise.post} --- the premise formula is pinned as the
  \emph{build\textquotesingle s} output: the two-atom conjunction in
  flat vertex order, comparators $<$ on extended vertices and $>$ on
  core ones exactly as displayed, each atom comparing the barycentric
  entry to \texttt{.const 0}. The build succeeds only at
  \texttt{schema.ambDim = 1}, so the certified reach of the displayed
  rule is the 1-dimensional fragment, by refusal (the
  \rulename{ConvElim}/\rulename{ConvIntro} class; the note above says
  it, this triple certifies it). $\Phi$ shared via \texttt{hpre}.
\item conclusion $\Phi \vdash b_i = 1 \wedge \bigwedge_{j \neq i} b_j
  = 0$ $\mapsto$ the goal under \texttt{hpost} $:$
  \texttt{cellOneHotForm schema outputs whichI seq.post} --- pinned
  as the one-hot build\textquotesingle s output: the $b_i = 1$ atom
  first, then the ascending-$j$ chain of zeros; the displayed
  big-wedge has exactly one accepted spelling.
\item nowhere in the certificate: the interior-disjointness
  well-formedness. It is consumed from the \emph{model
  side} --- admissible valuations carry
  \texttt{SimpC.CellIntDisjoint} in the cell sp-fact (certified
  alongside), and the proof pins the model\textquotesingle s cell
  choice with it --- so the display is right to show no such premise,
  and the rule\textquotesingle s soundness is honestly conditional on
  that provisional schema conjunct (the commentary\textquotesingle s
  PROPOSED-B2-DISJOINT story, visible in the statement chain).
\end{itemize}
Width note: within the 1-dimensional fragment, and modulo the
undisplayed port-typing pack flagged above, nothing certified is
narrower or wider than the display --- every hypothesis is consumed,
and the two builds pin both formulas to single spellings.

\ruleblock{Lean verifier receipts}
\sourcenote{module \texttt{BernsteinLogic.Soundness}; bound build \texttt{a20d90e66c13\ldots} (bound from receipt a9c0e34d3e3c9c8c, 2026-08-03T21:30:36Z). A receipt records that named checking layers accepted a build: (module, compiled-artifact SHA-256, environment SHA-256, checker layers), minted when the checking gate ran; the store receipts attesting the bound build are listed with the layers each run credited. Nothing here asks the reader to trust the store: every layer is an independently runnable program over the quoted sources --- rebuild the module and re-run the checkers to reproduce the verdict. Receipts attesting superseded builds credit no layer and are omitted from this rendering (the count of omissions, if any, is noted below).}
\begin{receipttable}
\receiptrow{a9c0e34d3e3c9c8c}{2026-08-03T21:30:36Z}{build+decorrelated}{\badge{BadgeGreen}{current}}
\end{receipttable}
\sourcenote{21 further receipt(s) in the store attest superseded builds of this module and are omitted from this rendering.}
\sourcenote{module \texttt{BernsteinLogic.Core.Catalogue}; bound build \texttt{cbc92e976a1d\ldots} (bound from receipt 5abdb3589707b5df, 2026-08-03T21:30:33Z). A receipt records that named checking layers accepted a build: (module, compiled-artifact SHA-256, environment SHA-256, checker layers), minted when the checking gate ran; the store receipts attesting the bound build are listed with the layers each run credited. Nothing here asks the reader to trust the store: every layer is an independently runnable program over the quoted sources --- rebuild the module and re-run the checkers to reproduce the verdict. Receipts attesting superseded builds credit no layer and are omitted from this rendering (the count of omissions, if any, is noted below).}
\begin{receipttable}
\receiptrow{24804e66edba0bc9}{2026-08-03T20:25:09Z}{build+decorrelated}{\badge{BadgeGreen}{current}}
\receiptrow{5abdb3589707b5df}{2026-08-03T21:30:33Z}{build+decorrelated}{\badge{BadgeGreen}{current}}
\end{receipttable}
\sourcenote{18 further receipt(s) in the store attest superseded builds of this module and are omitted from this rendering.}
\sourcenote{module \texttt{BernsteinLogic.Catalogue}; bound build \texttt{c477b6015133\ldots} (bound from receipt d87d31fb92f86564, 2026-07-20T07:43:40Z). A receipt records that named checking layers accepted a build: (module, compiled-artifact SHA-256, environment SHA-256, checker layers), minted when the checking gate ran; the store receipts attesting the bound build are listed with the layers each run credited. Nothing here asks the reader to trust the store: every layer is an independently runnable program over the quoted sources --- rebuild the module and re-run the checkers to reproduce the verdict. Receipts attesting superseded builds credit no layer and are omitted from this rendering (the count of omissions, if any, is noted below).}
\begin{receipttable}
\receiptrow{d87d31fb92f86564}{2026-07-20T07:43:40Z}{build+decorrelated}{\badge{BadgeGreen}{current}}
\receiptrow{85a439e3f753d74e}{2026-07-20T20:44:30Z}{build+decorrelated}{\badge{BadgeGreen}{current}}
\receiptrow{1e680b8d5ca91dae}{2026-07-20T21:00:59Z}{build+replay+hardened+decorrelated}{\badge{BadgeGreen}{current}}
\receiptrow{c707aa2f41b0f58f}{2026-07-21T05:28:24Z}{build+replay+hardened+decorrelated}{\badge{BadgeGreen}{current}}
\end{receipttable}
\sourcenote{4 further receipt(s) in the store attest superseded builds of this module and are omitted from this rendering.}
\layerline{\layeritem{build}{lake build (elaboration)}{\badge{BadgeGreen}{current}}\layeritem{replay}{leanchecker (official kernel replay)}{\badge{BadgeGreen}{current}}\layeritem{hardened}{safe\_verify (hardened screen)}{\badge{BadgeGreen}{current}}\layeritem{decorrelated}{lean4export + nanoda (decorrelated kernel)}{\badge{BadgeGreen}{current}}}
\sourcenote{\badge{BadgeGray}{comparator: none} no fixture in the three-implementation comparator harness (compiled / kernel-interpreted / reference-interpreted) exercises this rule.}
\ruleblock{HOL4 statement}
\sourcenote{theory \texttt{bern\_tr\_catdata2}, HOL4 development script \texttt{\texttt{bern\_tr\_catdata2Script.sml}} \badge{BadgeGreen}{script bound}}
\holhead{Theorem pcellinintdata\_receipt}
\begin{holsrc}
Theorem pcellinintdata_receipt:
  ruledataOk cd2wsCell cd2Hcell
    <| pre := FAnd (FAtom BowLt (TmAdd (TmConst (&1))
                                       (TmMul (TmConst (-&1)) (TmVar 9)))
                          (TmConst 0))
                   (FAtom BowGt (TmAdd (TmConst (&0))
                                       (TmMul (TmConst (&1)) (TmVar 9)))
                          (TmConst 0)) ;
       post := FAnd (FAtom BowEq (TmVar 5) (TmConst 1))
                    (FAtom BowEq (TmVar 7) (TmConst 0)) |>
    (RdCellInInterior (PCellInIntD
       <| pre := FAnd (FAtom BowLt (TmAdd (TmConst (&1))
                                          (TmMul (TmConst (-&1))
                                                 (TmVar 9)))
                             (TmConst 0))
                      (FAtom BowGt (TmAdd (TmConst (&0))
                                          (TmMul (TmConst (&1))
                                                 (TmVar 9)))
                             (TmConst 0)) ;
          post := FAnd (FAtom BowLt (TmAdd (TmConst (&1))
                                           (TmMul (TmConst (-&1))
                                                  (TmVar 9)))
                              (TmConst 0))
                       (FAtom BowGt (TmAdd (TmConst (&0))
                                           (TmMul (TmConst (&1))
                                                  (TmVar 9)))
                              (TmConst 0)) |> (ProofObj RdAx)
       checkfnDemoSchema 9 [5; 7] 0)) ∧
  ¬ ruledataOk cd2wsCell hnEmpty
      <| pre := FTt ; post := FTt |>
      (RdCellInInterior (PCellInIntD
         <| pre := FTt ; post := FTt |> (ProofObj RdAx)
         checkfnDemoSchema 9 [5; 7] 0))
\end{holsrc}
\ruleblock{HOL verifier receipts}
\badge{BadgeBlue}{documentary} \texttt{hol-strict} CLEAN: theory \texttt{bern\_tr\_catdata2}, 272
theorems, 0 cheat-tagged, 0 nonstandard axioms (whitelist:
\texttt{DISK\_THM}; the four \texttt{bool} axioms) --- verification
run of 2026-07-22 on an isolated build host, against a
content-addressed checkout of the exact sources cited here (git blob
\texttt{cd3c1e95\ldots}). The formerly dropped cell-payload island is
ported: this rule\textquotesingle s arm of the mutual wellformedness
predicate states the application membership, the input-shape gate,
the strict-interior premise and the one-hot conclusion (the form
faces as boolean Build-equations, the characterization convention),
and the form characterizations themselves are ported in theory
\texttt{bern\_tr\_catforms}
(\texttt{cellStrictInteriorForm\_char} /
\texttt{cellOneHotForm\_char}, 23 theorems, 0 offences, same run).
The receipt theorem cited above is the arm\textquotesingle s
kernel-evaluated accept/refusal pair; the shared barycentric
ingredient (\texttt{baryCoeffs1D}) was already ported and receipted.
Honest residue: the cell soundness argument (the Lean
development\textquotesingle s \texttt{sound\_cellInInterior}) is not
yet ported; the payload/builder layer above is.

\sourcenote{HOL receipts are documentary: quoted results of \texttt{hol-strict} runs
and census rows from the HOL4 development\textquotesingle s build
evidence, not minted receipt artefacts. A HOL receipt store (the Lean
store\textquotesingle s twin) is named future work.
}
\ruleblock{Proof}

\sourcenote{Genre: admissible rule (tier 2) --- a proof-tree-rewriting
recipe. The intended elaboration is a case analysis over the cell
primitive's one-hot disjunction (the $\sp[\Cell[\Sigma]]$ axiom ---
the calculus's only source of disjunction); the
\emph{realized} warrant is the degenerate rewrite in the
\rulename{PConstExact} sense, discharging against a named verified
semantic lemma. Both are presented, and the honest reason the first is
not yet a checkable catalogue derivation is stated.}

\begin{brecipe}[the intended rewrite --- case analysis over the cell disjunction]
\label{rec:cii-intended}
A cell application $\Cell[\Sigma]_{\{x\}\to\{b_1,\ldots,b_N\}}$
contributes to the ambient premise its one-hot disjunction
\[
\sp[\Cell[\Sigma]] \;=\; \bigvee_{k=1}^{N}\Bigl( x \in \overline{C_k}
\;\wedge\; b_k = 1 \;\wedge\; \textstyle\bigwedge_{j \neq k} b_j = 0 \Bigr).
\]
The intended unpacking of one \rulename{Cell-In-Interior} node is a
case split on $k$: in branch $k = i$ the conclusion
$b_i = 1 \wedge \bigwedge_{j\neq i} b_j = 0$ is reached from the case
hypothesis by conjunction plumbing ($\wedge$E and $\wedge$I under
whatever association the axiom's serialization fixes); in a branch
$k \neq i$ the case hypothesis
$x \in \overline{C_k}$ contradicts the rule's strict-interior premise
$x \in \mathring{C_i}$ --- \emph{provided no point lies strictly
inside one cell of the enumeration and weakly inside another}, which
is exactly the schema well-formedness the soundness theorem consumes
(\autoref{rem:cii-disjoint}). Two honest gaps keep this a template
rather than a shipped derivation: (i) the disjunction sits \emph{inside}
the premise conjunction, and the calculus's disjunction elimination as
currently realized acts only on a disjunction at the root of the premise
--- the context-carrying form that reaches it is proven as a satellite
theorem but its core constructor swap is still scheduled; (ii) the
contradiction in the off-branches is semantic (disjoint interiors), not
yet a formula-level refutation in the negation-free grammar.
\end{brecipe}

\begin{brecipe}[the realized rewrite --- discharge against the semantic lemma]
\label{rec:cii-realized}
As shipped, the rule is admissible by \emph{reduction to verified
machinery} (the \rulename{PConstExact} pattern): the node
\[
\inferrule*[Right=\rulename{Cell-In-Interior}]
  {\;\mathcal{T} \rhd \Phi \vdash \textstyle\bigwedge_{v}
    \lambda^{C_i}_v(x) \bowtie^{\circ}_v 0\;}
  {\Phi \vdash b_i = 1 \wedge \textstyle\bigwedge_{j \neq i} b_j = 0}
\qquad\rightsquigarrow\qquad
\text{\upshape discharged by } \texttt{sound\_cellInInterior},
\]
where the premise is the \emph{strict} barycentric-interior conjunction
for cell $C_i$ (one strict atom per support vertex, $<$ for extended
vertices, $>$ for core vertices, in the pinned flat vertex order) and
the checker recomputes both formula shapes from the schema --- the
premise via \texttt{cellStrictInteriorBuild}, the one-hot conclusion via
\texttt{cellOneHotBuild} --- so no formula text is trusted from the
certificate.
\end{brecipe}

\begin{bprop}[semantic warrant]\label{prop:cii-sound}
Suppose the payload carries a proof of the strict-interior premise, the
cited cell application is in the hypernet, the input wire is typed at
shape $[\mathrm{ambDim}]$, every output wire at total size $1$, and the
premise and conclusion formulas match the recomputed builds. Then the
conclusion sequent is semantically valid.
\end{bprop}
\begin{proof}[Proof sketch (the lemma's actual route)]
Fix an admissible valuation $\sigma$ satisfying $\Phi$. Admissibility of
the cited cell application yields the model-level fact: the schema's
cells have disjoint interiors, the input wire holds a rational point
$x$, and there is a cell choice $k$ with $x \in
\mathcal{R}(\Delta_k, S_k)$ (weak membership) whose one-hot valuation
the outputs realize --- entry $0$ of output $j$ reads $1$ if $j = k$,
else $0$. The premise, transported through the input wire's declared
shape (the tensor-to-point tie), gives \emph{strict} barycentric signs
at $x$ for cell $i$, i.e.\ $x$ lies in the strict interior of $C_i$
(\texttt{maskRegionStrict\_pair\_intro} assembles the region witness
from the two strict atoms). If $k \neq i$, the interior-exclusion
well-formedness is contradicted: $x$ is strictly inside $C_i$ and
weakly inside $C_k$, the exact configuration
\autoref{rem:cii-disjoint}'s predicate forbids (strict-versus-weak
exclusion --- deliberately stronger than bare disjointness of open
interiors, which a strictly-interior-to-$C_i$ point on $C_k$'s
boundary would survive). So $k = i$, and the one-hot output pattern at $k$ \emph{is} the conclusion:
each conjunct is an equality between a singleton-shaped output wire and
the constant its entry holds (\texttt{holds\_cellOneHot}).
\end{proof}

\begin{bremark}[why the premise must be strict]\label{rem:cii-strict}
At a shared boundary face the one-hot semantics deliberately admits
\emph{either} adjacent cell as the model's choice: a boundary point has
the one-hot set of every cell whose closure contains it. A weak-interior
premise therefore cannot pin $k$ --- for $x$ on the face between $C_i$
and $C_k$ the valuation one-hotting $C_k$ satisfies the weak premise
and falsifies $b_i = 1$. Strictness excludes exactly the boundary, where
the pinning genuinely fails; the companion rule \rulename{Cell-Out}
covers boundary points from the other side.
\end{bremark}

\begin{bremark}[the disjointness hypothesis, and why the complex's
conformity does not deliver it]\label{rem:cii-disjoint}
The step $k = i$ uses: no rational point lies strictly inside cell $i$
and weakly inside a \emph{different} cell $k$
(\texttt{SimpC.CellIntDisjoint}, stated over cell \emph{indices} so
duplicate entries are excluded). The schema's simplicial conformity
field does not imply this: conformity constrains simplex-mask
intersections only, never sees the extension masks that carve cells out
of simplices, and permits literally duplicate cell entries --- either
defect, were the predicate not consumed, would let the model one-hot an
overlapping cell and falsify the conclusion. The predicate is consumed as a conjunct of the cell
application's model-level fact: an application over a non-disjoint
enumeration admits no valuation at all (hypothesis-side vacuity ---
sound, merely weak), the same reading that excludes NaN executions.
It is a schema amendment adopted provisionally, pending
ratification.
\end{bremark}

\begin{bremark}[the port-typing hypotheses are load-bearing twice]
The model's one-hot fact pins \emph{entry $0$} of each output wire,
while the conclusion's equality atoms evaluate the \emph{whole} wire;
the total-size-$1$ typing is what makes the two coincide (a multi-entry
output could satisfy the entry fact while its remaining entries falsify
the atom; a zero-size output would satisfy every atom vacuously).
Dually, the premise's barycentric atoms constrain the input
\emph{tensor}, and the declared shape $[\mathrm{ambDim}]$ (with
$\mathrm{ambDim} = 1$ in the checkable fragment) ties that tensor to the
model's rational point.
\end{bremark}

\begin{bremark}[scope of the checkable fragment]
\texttt{cellStrictInteriorBuild} succeeds only over a $1$-dimensional
complex with a two-vertex cell support on distinct coordinates --- the
same fragment as the region machinery elsewhere in the catalogue;
everything beyond it is refused, not approximated (fail-toward-refusal).
Higher-dimensional cell complexes await their own barycentric builders;
the rule's statement is dimension-general and only the builder gates
narrow it.
\end{bremark}

\begin{anchors}
Recipe \ref{rec:cii-realized} warrant & \texttt{sound\_cellInInterior} (Soundness)\\
premise shape extraction & \texttt{cellStrictInteriorBuild\_spec}; builder \texttt{cellStrictInteriorBuild} / \texttt{baryLambdas1D}\\
strict-region assembly & \texttt{maskRegionStrict\_pair\_intro}; \texttt{MaskRegionStrict}\\
the disjointness pin (Rem.\ \ref{rem:cii-disjoint}) & \texttt{SimpC.CellIntDisjoint}; the cell case of \texttt{PrimApp.SemOk}\\
one-hot conclusion & \texttt{cellOneHotBuild}; \texttt{holds\_cellOneHot}\\
port typing ties & \texttt{PCellInIntData} fields \texttt{inputShape}, \texttt{outputsScalar}\\
\end{anchors}


\rulesection{25}{Cell-Out}{admissible rule}
\label{blcatrule-CellOut}
\ruleblock{The rule (Steele metanotation)}
\sourcenote{as realized in the mechanized development; source row: \texttt{\texttt{rules.toml}} --- \emph{Cell-Out}. The metanotation statement is typeset by hand for notational quality; its mechanized ground truth is the Lean statement below, extracted from the sources at every build of this book and hash-bound. The rule \emph{roster} is machine-checked (every mechanized rule tag round-trips to exactly one catalogue row); the statement \emph{text} is editorial, checked by eye against the extracted Lean beside it.}
\begin{mathpar}
\inferrule*[Right=\rulename{Cell-Out}]{
  \Phi \vdash \lambda^{C_j}_w(x) \bowtie^{\circ} 0 \ \text{(a single strict separating atom at a named witness vertex } w \text{ of } C_j\text{; } > \text{ for extended } w\text{, } < \text{ for core } w\text{)}
  \\
  \Cell[\Sigma]_{\{x\}\to\{b_1,\ldots,b_N\}} \in H
}{
  \Phi \vdash b_j = 0
}
\end{mathpar}

\ruleflavor{applies even at boundary points.}
\par\medskip
\rulename{Cell-Out} is \rulename{Cell-In-Interior}\textquotesingle s
refutation companion: from a premise strictly separating the input
point from cell $C_j$, it zeroes that cell\textquotesingle s indicator
--- $\Phi \vdash b_j = 0$ --- and it does so \emph{even at boundary
points of other cells}, where the pinning rule cannot fire. The paper
states the premise as $x \notin \overline{C_j}$, a disjunction of
strict halfspace violations; the negation-free formula grammar has no
rendering of that, so the realized rule serializes a witness instead:
a single strict barycentric atom at a \emph{named witness vertex} $w$
of $C_j$, with the polarity that contradicts membership ($>$ at an
extended vertex, $<$ at a core vertex). One atom replaces the
disjunction --- no loss for producers, who know which halfspace their
point violates, and a strict gain for the checker, which does no
disjunctive reasoning. A producer uses the rule to kill off-branches
of a piecewise definition: once every indicator but one is zero,
downstream arithmetic (a weighted combination over branches, a
substitution) simplifies to the live branch.

What the checker recomputes: presence of the cited cell application
with the same port typing as the companion rule (input at
$[\mathrm{ambDim}]$, outputs at total size one); that the witness
vertex lies in cell $j$\textquotesingle s support mask (the build
refuses otherwise); that the premise is byte-for-byte the recomputed
separating atom for $(j, w)$; and that the conclusion is the
recomputed $b_j = 0$ atom.

Each hypothesis is load-bearing. The support-mask guard: an atom at a
vertex outside $C_j$\textquotesingle s support constrains a weight the
membership predicate never mentions --- no contradiction, no soundness.
Strictness of the atom: the model\textquotesingle s membership fact
delivers weak signs, and only a strict opposite sign contradicts them
(a weak atom is compatible with a boundary weight of zero, where the
model may still choose cell $j$). The port typing: same double duty as
in the companion rule --- the tensor-to-point tie on the input, and
the entry-to-whole-wire lift on the output. Notably ABSENT here: the
interior-disjointness well-formedness --- the refutation is by weight
uniqueness on the two-vertex support against the strict witness atom,
within cell $j$ alone, so this rule stands even where only the
companion needs the provisional schema amendment. The proof section
walks these in order.

\ruleblock{Lean statement}
\declhead{BernsteinLogic.Soundness}{sound\_cellOut}{\badge{BadgeGreen}{statement bound}}
\begin{leansrc}
/-- Cell-Out: the strict separating witness refutes the model choosing
cell whichJ (weight uniqueness on the two-vertex support), so the
one-hot sp-fact zeroes its indicator. -/
theorem sound_cellOut {H : Hypernet} {seq premise : Sequent}
    {schema : SimpC} {input : Wire} {outputs : List Wire}
    {whichJ : Fin schema.cells.length} {wv : Fin schema.nVertices}
    (ih : premise.semanticallyValid H)
    (happ : PrimApp.cell schema input outputs ∈ H.applications)
    (hinshape : ∃ τ, H.wireType input = some τ ∧
      τ.shape = [schema.ambDim])
    (houts : ∀ o ∈ outputs, ∃ τ, H.wireType o = some τ ∧
      Shape.totalSize τ.shape = 1)
    (hpre : seq.pre = premise.pre)
    (hwit : cellOutWitnessForm schema (.var input) whichJ wv premise.post)
    (hpost : cellOutZeroForm schema outputs whichJ seq.post) :
    seq.semanticallyValid H
\end{leansrc}
\declhead{BernsteinLogic.Soundness}{cellOutWitnessBuild\_spec}{\badge{BadgeGreen}{statement bound}}
\begin{leansrc}
/-- Extraction for the Cell-Out witness premise: additionally pins the
witness vertex to the cell's mask and identifies it with one of the two
support vertices, selecting the corresponding strict atom. -/
theorem cellOutWitnessBuild_spec {schema : SimpC} {t : Term}
    {k : Fin schema.cells.length} {w : Fin schema.nVertices} {F : Formula}
    (h : cellOutWitnessForm schema t k w F) :
    schema.ambDim = 1 ∧
    (schema.cells.get k).1 w = true ∧
    ∃ (a b : Fin schema.nVertices) (V₀ V₁ : ℚ),
      (List.finRange schema.nVertices).filter
        (fun v => (schema.cells.get k).1 v) = [a, b] ∧
      (∀ j, schema.coord a j = V₀) ∧
      (∀ j, schema.coord b j = V₁) ∧
      V₀ ≠ V₁ ∧
      ((w = a ∧ F = .atom (if (schema.cells.get k).2 a then .gt else .lt)
          (.add (.const (V₁ / (V₁ - V₀))) (.mul (.const (-1 / (V₁ - V₀))) t))
          (.const 0)) ∨
       (w = b ∧ F = .atom (if (schema.cells.get k).2 b then .gt else .lt)
          (.add (.const (-V₀ / (V₁ - V₀))) (.mul (.const (1 / (V₁ - V₀))) t))
          (.const 0)))
\end{leansrc}
\declhead{BernsteinLogic.Soundness}{maskRegion\_pair\_elim}{\badge{BadgeGreen}{statement bound}}
\begin{leansrc}
/-- Extraction: a masked region membership over a two-vertex 1-D support
pins the barycentric weights to their rational closed forms, and exports
the sign of each weight according to its extension flag. -/
theorem maskRegion_pair_elim {C : SimpC} {mask ext : Fin C.nVertices → Bool}
    {x : Fin C.ambDim → ℚ} {a b : Fin C.nVertices} {V₀ V₁ : ℚ}
    (hfil : (List.finRange C.nVertices).filter (fun v => mask v) = [a, b])
    (hVa : ∀ j, C.coord a j = V₀) (hVb : ∀ j, C.coord b j = V₁)
    (hne : V₀ ≠ V₁) (j0 : Fin C.ambDim)
    (h : MaskRegion C mask ext x) :
    (ext a = true → (V₁ - x j0) / (V₁ - V₀) ≤ 0) ∧
    (ext a = false → 0 ≤ (V₁ - x j0) / (V₁ - V₀)) ∧
    (ext b = true → (x j0 - V₀) / (V₁ - V₀) ≤ 0) ∧
    (ext b = false → 0 ≤ (x j0 - V₀) / (V₁ - V₀))
\end{leansrc}
\par\medskip
\textbf{The recomputed forms and the payload.}\enspace As for the
companion rule, no formula text is trusted from the certificate: the
witness atom is rebuilt from the schema\textquotesingle s barycentric
data at the named vertex (note the \texttt{guard} --- the build refuses
a witness vertex outside the cell\textquotesingle s support mask), the
conclusion atom from the output list, and the payload\textquotesingle s
form fields assert equality with the rebuilt formulas. The same
two-vertex 1-D fragment gates apply (\texttt{baryLambdas1D}, displayed
in the \rulename{Cell-In-Interior} section); everything beyond the
fragment is refused.

\declhead{BernsteinLogic.Core.Catalogue}{cellOutWitnessBuild}{\badge{BadgeGreen}{definition bound}}
\begin{leansrc}
/-- Single strict separating atom at witness vertex w. -/
def cellOutWitnessBuild (schema : SimpC) (t : Term)
    (k : Fin schema.cells.length) (w : Fin schema.nVertices) :
    Option Formula := do
  let c := schema.cells[k.val]'(k.isLt)
  guard (c.1 w = true)
  let lams ← baryLambdas1D schema c.1 t
  let entry ← lams.find? fun p => decide (p.1 = w)
  some (.atom (if c.2 w then .gt else .lt) entry.2 (.const 0))
\end{leansrc}
\declhead{BernsteinLogic.Core.Catalogue}{cellOutWitnessForm}{\badge{BadgeGreen}{definition bound}}
\begin{leansrc}
def cellOutWitnessForm (schema : SimpC) (t : Term)
    (k : Fin schema.cells.length) (w : Fin schema.nVertices)
    (F : Formula) : Prop :=
  cellOutWitnessBuild schema t k w = some F
\end{leansrc}
\declhead{BernsteinLogic.Core.Catalogue}{cellOutZeroBuild}{\badge{BadgeGreen}{definition bound}}
\begin{leansrc}
def cellOutZeroBuild (schema : SimpC) (outputs : List Wire)
    (k : Fin schema.cells.length) : Option Formula :=
  if outputs.length = schema.cells.length then
    some (.atom .eq (.var (outputs.getD k.val 0)) (.const 0))
  else none
\end{leansrc}
\declhead{BernsteinLogic.Core.Catalogue}{cellOutZeroForm}{\badge{BadgeGreen}{definition bound}}
\begin{leansrc}
def cellOutZeroForm (schema : SimpC) (outputs : List Wire)
    (k : Fin schema.cells.length) (F : Formula) : Prop :=
  cellOutZeroBuild schema outputs k = some F
\end{leansrc}
\declhead{BernsteinLogic.Core.Catalogue}{PCellOutData}{\badge{BadgeGreen}{definition bound}}
\begin{leansrc}
inductive PCellOutData : Hypernet → Sequent → Type where
  | mk (H : Hypernet) (seq : Sequent)
       (premise : Sequent) (premiseProof : ProofObject H premise)
       (schema : SimpC) (input : Wire) (outputs : List Wire)
       (whichJ : Fin schema.cells.length)
       (witnessVertex : Fin schema.nVertices)
       (appInH : PrimApp.cell schema input outputs ∈ H.applications)
       (inputShape : ∃ τ, H.wireType input = some τ ∧
                       τ.shape = [schema.ambDim])
       (outputsScalar : ∀ o ∈ outputs, ∃ τ, H.wireType o = some τ ∧
                          Shape.totalSize τ.shape = 1)
       (preEq  : seq.pre = premise.pre)
       (outWitness : cellOutWitnessForm schema (.var input) whichJ
                       witnessVertex premise.post)
       (postZero : cellOutZeroForm schema outputs whichJ seq.post)
       : PCellOutData H seq
\end{leansrc}
\ruleblock{Correspondence (Steele display $\leftrightarrow$ Lean)}
\badge{BadgeGreen}{reviewer-certified against the displayed texts}
\sourcenote{The dictionary below is a reviewer-certified statement of how the mechanized statements above interpret the displayed rule: display metavariable $\mapsto$ Lean binder, each display hypothesis $\mapsto$ the premise or side condition carrying it, the display conclusion $\mapsto$ the goal --- with any deliberate difference in width between display and mechanization named outright. Certification is recorded as a hash triple over (the displayed cell, the extracted statement texts, this dictionary\textquotesingle s own text); if any of the three changes, this block turns amber until a reviewer re-certifies. The correspondence claim itself remains editorial --- reviewer-grade, not machine-checked; its machine-checked successor types each soundness statement against the rule as a structured object, and supersedes this block per rule as it lands.}
Judgment frame as in \rulename{Ax}; sub-derivation premise as at
\rulename{Weaken}; binders, linkage row, and the undisplayed
port-typing pack (\texttt{happ}, \texttt{hinshape}, \texttt{houts} ---
payload fields \texttt{PCellOutData}) exactly as in the
\rulename{Cell-In-Interior} dictionary, with $j \mapsto$
\texttt{whichJ} and the witness vertex $w \mapsto$ \texttt{wv : Fin
schema.nVertices}.
\begin{itemize}
\item premise $\Phi \vdash \lambda^{C_j}_w(x) \bowtie^{\circ} 0$
  $\mapsto$ \texttt{ih} with \texttt{hwit} $:$
  \texttt{cellOutWitnessForm schema (.var input) whichJ wv
  premise.post} --- pinned as the witness build\textquotesingle s
  output: a single strict atom at the flat-order entry of the named
  vertex, comparator $>$ on an extended witness and $<$ on a core
  one, exactly the display\textquotesingle s polarity (note the
  deliberate flip against \rulename{Cell-In-Interior}\textquotesingle
  s interior comparators); the build\textquotesingle s guard demands
  the witness vertex lie in $C_j$\textquotesingle s support --- the
  display\textquotesingle s ``of $C_j$\textquotesingle\textquotesingle{}
  --- and, as at the companion, succeeds only at
  \texttt{schema.ambDim = 1}: the certified reach is the
  1-dimensional fragment, by refusal. $\Phi$ shared via
  \texttt{hpre}.
\item conclusion $\Phi \vdash b_j = 0$ $\mapsto$ the goal under
  \texttt{hpost} $:$ \texttt{cellOutZeroForm schema outputs whichJ
  seq.post} --- the single equality atom, indicator wire left,
  \texttt{.const 0} right.
\item the disjointness conjunct of the model\textquotesingle s
  sp-fact is \emph{discarded} by this rule\textquotesingle s proof
  (the statement chain consumes the one-hot fact only) --- certifying
  the commentary\textquotesingle s claim that \rulename{Cell-Out}
  stands independent of the provisional schema amendment its
  companion needs.
\end{itemize}
Width note: ``applies even at boundary points\textquotesingle\textquotesingle{}
is carried structurally --- the witness atom is strict and nothing
demands interiority of $x$; within the 1-dimensional fragment, and
modulo the undisplayed port-typing pack inherited from the companion
dictionary, neither side is wider, every hypothesis consumed.

\ruleblock{Lean verifier receipts}
\sourcenote{module \texttt{BernsteinLogic.Soundness}; bound build \texttt{a20d90e66c13\ldots} (bound from receipt a9c0e34d3e3c9c8c, 2026-08-03T21:30:36Z). A receipt records that named checking layers accepted a build: (module, compiled-artifact SHA-256, environment SHA-256, checker layers), minted when the checking gate ran; the store receipts attesting the bound build are listed with the layers each run credited. Nothing here asks the reader to trust the store: every layer is an independently runnable program over the quoted sources --- rebuild the module and re-run the checkers to reproduce the verdict. Receipts attesting superseded builds credit no layer and are omitted from this rendering (the count of omissions, if any, is noted below).}
\begin{receipttable}
\receiptrow{a9c0e34d3e3c9c8c}{2026-08-03T21:30:36Z}{build+decorrelated}{\badge{BadgeGreen}{current}}
\end{receipttable}
\sourcenote{21 further receipt(s) in the store attest superseded builds of this module and are omitted from this rendering.}
\sourcenote{module \texttt{BernsteinLogic.Core.Catalogue}; bound build \texttt{cbc92e976a1d\ldots} (bound from receipt 5abdb3589707b5df, 2026-08-03T21:30:33Z). A receipt records that named checking layers accepted a build: (module, compiled-artifact SHA-256, environment SHA-256, checker layers), minted when the checking gate ran; the store receipts attesting the bound build are listed with the layers each run credited. Nothing here asks the reader to trust the store: every layer is an independently runnable program over the quoted sources --- rebuild the module and re-run the checkers to reproduce the verdict. Receipts attesting superseded builds credit no layer and are omitted from this rendering (the count of omissions, if any, is noted below).}
\begin{receipttable}
\receiptrow{24804e66edba0bc9}{2026-08-03T20:25:09Z}{build+decorrelated}{\badge{BadgeGreen}{current}}
\receiptrow{5abdb3589707b5df}{2026-08-03T21:30:33Z}{build+decorrelated}{\badge{BadgeGreen}{current}}
\end{receipttable}
\sourcenote{18 further receipt(s) in the store attest superseded builds of this module and are omitted from this rendering.}
\layerline{\layeritem{build}{lake build (elaboration)}{\badge{BadgeGreen}{current}}\layeritem{replay}{leanchecker (official kernel replay)}{\badge{BadgeAmber}{stale only}}\layeritem{hardened}{safe\_verify (hardened screen)}{\badge{BadgeAmber}{stale only}}\layeritem{decorrelated}{lean4export + nanoda (decorrelated kernel)}{\badge{BadgeGreen}{current}}}
\sourcenote{\badge{BadgeGray}{comparator: none} no fixture in the three-implementation comparator harness (compiled / kernel-interpreted / reference-interpreted) exercises this rule.}
\ruleblock{HOL4 statement}
\sourcenote{theory \texttt{bern\_tr\_catdata2}, HOL4 development script \texttt{\texttt{bern\_tr\_catdata2Script.sml}} \badge{BadgeGreen}{script bound}}
\holhead{Definition cellOutWitnessBuild\_def}
\begin{holsrc}
Definition cellOutWitnessBuild_def:
  cellOutWitnessBuild (schema:simpc) t (k:num) (w:num) =
    if k < LENGTH schema.cells then
      (let c = EL k schema.cells in
       if getD (FST c) w F then
         (case baryLambdas1D schema (FST c) t of
          | NONE => NONE
          | SOME lams =>
            (case FILTER (λp. FST p = w) lams of
             | [] => NONE
             | (entry::rest) =>
                 SOME (FAtom (if getD (SND c) w F then BowGt else BowLt)
                             (SND entry) (TmConst 0))))
       else NONE)
    else NONE
End
\end{holsrc}
\holhead{Definition cellOutWitnessFormB\_def}
\begin{holsrc}
Definition cellOutWitnessFormB_def:
  cellOutWitnessFormB schema t k w fm ⇔
    cellOutWitnessBuild schema t k w = SOME fm
End
\end{holsrc}
\holhead{Theorem pcelloutdata\_receipt}
\begin{holsrc}
Theorem pcelloutdata_receipt:
  ruledataOk cd2wsCell cd2Hcell
    <| pre := FAtom BowGt (TmAdd (TmConst (&1))
                                 (TmMul (TmConst (-&1)) (TmVar 9)))
                    (TmConst 0) ;
       post := FAtom BowEq (TmVar 5) (TmConst 0) |>
    (RdCellOut (PCellOutD
       <| pre := FAtom BowGt (TmAdd (TmConst (&1))
                                    (TmMul (TmConst (-&1)) (TmVar 9)))
                       (TmConst 0) ;
          post := FAtom BowGt (TmAdd (TmConst (&1))
                                     (TmMul (TmConst (-&1)) (TmVar 9)))
                        (TmConst 0) |> (ProofObj RdAx)
       checkfnDemoSchema 9 [5; 7] 0 0)) ∧
  ¬ ruledataOk cd2wsCell hnEmpty
      <| pre := FTt ; post := FTt |>
      (RdCellOut (PCellOutD
         <| pre := FTt ; post := FTt |> (ProofObj RdAx)
         checkfnDemoSchema 9 [5; 7] 0 0))
\end{holsrc}
\ruleblock{HOL verifier receipts}
\badge{BadgeBlue}{documentary} \texttt{hol-strict} CLEAN: theory \texttt{bern\_tr\_catdata2}, 272
theorems, 0 cheat-tagged, 0 nonstandard axioms (whitelist:
\texttt{DISK\_THM}; the four \texttt{bool} axioms) --- verification
run of 2026-07-22 on an isolated build host, against a
content-addressed checkout of the exact sources cited here (git blob
\texttt{cd3c1e95\ldots}). The formerly dropped cell-builder island is
ported: the witness build/form pair displayed here verbatim (the Form
face as a boolean Build-equation, the characterization convention),
with kernel-evaluated receipts (\texttt{cellOutWitnessBuild\_receipt}
/ \texttt{cellOutWitnessFormB\_receipt}), the payload arm of the
mutual wellformedness predicate with its own accept/refusal receipt
(cited above), the zero-build twin in theory
\texttt{bern\_tr\_checkfn} (\texttt{cellOutZeroBuild}, 8 theorems, 0
offences, same run), and its form characterization in theory
\texttt{bern\_tr\_catforms} (\texttt{cellOutZeroForm\_char}, 23
theorems, 0 offences). The shared barycentric ingredient
(\texttt{baryCoeffs1D}) was already ported and receipted. Honest
residue: the cell soundness argument (the Lean
development\textquotesingle s \texttt{sound\_cellOut}) is not yet
ported; the payload/builder layer above is.

\sourcenote{HOL receipts are documentary: quoted results of \texttt{hol-strict} runs
and census rows from the HOL4 development\textquotesingle s build
evidence, not minted receipt artefacts. A HOL receipt store (the Lean
store\textquotesingle s twin) is named future work.
}
\ruleblock{Proof}

\sourcenote{Genre: admissible rule (tier 2) --- a proof-tree-rewriting
recipe. Same family as \rulename{Cell-In-Interior}: the intended
elaboration is a case analysis over the cell primitive's one-hot
disjunction; the realized warrant is the degenerate rewrite discharging
against a verified semantic lemma, for the same two reasons stated in
that rule's first recipe (the context-carrying case analysis is not yet
a core constructor, and the off-branch contradiction is semantic). Only
the branch-closing argument differs --- and it is \emph{lighter}: no
disjointness hypothesis is consumed.}

\begin{brecipe}[the \rulename{Cell-Out} instance --- the realized rewrite]
\label{rec:co-realized}
\[
\inferrule*[Right=\rulename{Cell-Out}]
  {\;\mathcal{T} \rhd \Phi \vdash \lambda^{C_j}_w(x) \bowtie^{\circ} 0\;}
  {\Phi \vdash b_j = 0}
\qquad\rightsquigarrow\qquad
\text{\upshape discharged by } \texttt{sound\_cellOut},
\]
where the premise is a \emph{single strict separating atom} at a named
witness vertex $w$ of cell $C_j$ ($>$ if $w$ is an extended vertex, $<$
if a core vertex --- the sign that contradicts membership), and the
checker recomputes premise and conclusion shapes from the schema
(\texttt{cellOutWitnessBuild}, \texttt{cellOutZeroBuild}). The payload
names $w$ explicitly; the build refuses any $w$ outside $C_j$'s support
mask.
\end{brecipe}

\begin{bprop}[semantic warrant]\label{prop:co-sound}
Under the payload hypotheses (premise proof; cited cell application in
the hypernet; input typed at shape $[\mathrm{ambDim}]$; outputs at total
size $1$; premise and conclusion matching the recomputed builds), the
conclusion sequent is semantically valid.
\end{bprop}
\begin{proof}[Proof sketch (the lemma's actual route)]
Fix an admissible valuation $\sigma$ satisfying $\Phi$. The cell
application's model-level fact yields the rational input point $x$, a
cell choice $k$ with $x \in \mathcal{R}(\Delta_k, S_k)$, and the
one-hot output valuation at $k$. Suppose $k = j$. On the two-vertex
support the barycentric weights of $x$ are \emph{unique} --- they are
determined affinely by the point --- so weak membership in cell $j$
pins their signs: $\leq 0$ at extended vertices, $\geq 0$ at core
vertices (the pair-elimination lemma). The premise delivers the
\emph{strict opposite} sign for the weight at the witness vertex $w$
(transported to the model's point through the input wire's declared
shape), a direct contradiction. Hence $k \neq j$, and the one-hot fact
reads entry $0$ of output $j$ as $0$; the total-size-$1$ typing lifts
that entry to the whole-wire equality $b_j = 0$.
\end{proof}

\begin{bremark}[why a witness vertex replaces the paper's set complement]
The source statement's premise is $x \notin \overline{C_j}$ --- a
\emph{disjunction} of strict halfspace violations, one per bounding
halfspace. The formula grammar is negation-free and the premise of a
rule application must be a formula the calculus can derive, so the
amendment serializes a choice of disjunct instead: \emph{which}
halfspace is strictly violated, named by \emph{which} vertex carries
the violated barycentric sign. One atom replaces the disjunction with
no loss for producers --- within the fragment the builders admit,
where the cell's closure is exactly the conjunction of its per-vertex
sign constraints, a point outside the closure violates at least one of
them, and the producer knows which --- and a strict gain for the
checker (no disjunctive reasoning at check time).
\end{bremark}

\begin{bremark}[the witness must be a support vertex of $C_j$]
The build guards $w \in \mathrm{supp}(C_j)$. Without the guard the
premise atom would constrain a barycentric coordinate that cell $j$'s
membership predicate never mentions --- the sign contradiction in
Proposition~\ref{prop:co-sound} needs the weight at $w$ to be one of
the two the membership fact actually constrains. (On the two-vertex
support this is also what makes the atom well-defined: the builder
computes $\lambda^{C_j}_w$ from the support pair's coordinates.)
\end{bremark}

\begin{bremark}[boundary points, and the division of labor with
\rulename{Cell-In-Interior}]
The rule applies even at boundary points of \emph{other} cells: the
contradiction is between a strict violated sign and weak membership of
the \emph{same} cell $j$, so no interior-disjointness hypothesis is
consumed --- a point on the shared face of $C_i$ and $C_k$ can still be
strictly separated from a third cell $C_j$, and the rule concludes
$b_j = 0$ regardless of which of the two adjacent cells the model
one-hots. This is the refutation companion of
\rulename{Cell-In-Interior},
whose pinning argument fails precisely on boundaries and needs
disjointness; between them, interior points get the full one-hot
vector, and a boundary point gets $b_j = 0$ for every cell $j$ whose
closure excludes it --- the cells actually sharing the face stay
unpinned, necessarily: the model may one-hot either, so no sound rule
could zero theirs.
\end{bremark}

\begin{bremark}[scope of the checkable fragment]
As for \rulename{Cell-In-Interior}: the builders succeed over
$1$-dimensional complexes with two-vertex cell supports on distinct
coordinates, and refuse everything else. The statement is
dimension-general; only the builder gates narrow it.
\end{bremark}

\begin{anchors}
Recipe \ref{rec:co-realized} warrant & \texttt{sound\_cellOut} (Soundness)\\
witness-atom extraction & \texttt{cellOutWitnessBuild\_spec}; builder \texttt{cellOutWitnessBuild} / \texttt{baryLambdas1D}\\
weight uniqueness on the support pair & the pair-elimination lemma cited inside \texttt{sound\_cellOut} (\texttt{maskRegion\_pair\_elim})\\
conclusion form & \texttt{cellOutZeroBuild}\\
support-mask guard & the \texttt{guard} in \texttt{cellOutWitnessBuild}; payload field \texttt{witnessVertex}\\
port typing ties & \texttt{PCellOutData} fields \texttt{inputShape}, \texttt{outputsScalar}\\
\end{anchors}


\rulesection{26}{PConstExact}{admissible rule}
\label{blcatrule-PConstExact}
\ruleblock{The rule (Steele metanotation)}
\sourcenote{as realized in the mechanized development; source row: \texttt{\texttt{rules.toml}} --- \emph{PConstExact}. The metanotation statement is typeset by hand for notational quality; its mechanized ground truth is the Lean statement below, extracted from the sources at every build of this book and hash-bound. The rule \emph{roster} is machine-checked (every mechanized rule tag round-trips to exactly one catalogue row); the statement \emph{text} is editorial, checked by eye against the extracted Lean beside it.}
\begin{mathpar}
\inferrule*[Right=\rulename{PConstExact}]{
  e_1, e_2 \text{ wire-free (closed) terms}
  \\
  \bowtie \in \{<, \leq, =, \geq, >\}
  \\
  e_1 \bowtie e_2 \ \text{holds under exact evaluation in } \Q
}{
  \Phi \vdash e_1 \bowtie e_2
}
\end{mathpar}

\ruleflavor{the third premise is decided, not assumed: both terms are closed, so the comparison evaluates in exact rational arithmetic (the checker\textquotesingle s realization is the verified closed evaluator, which returns no verdict on any term containing a wire --- closedness is thereby implied, never a trusted annotation).}
\par\medskip
\rulename{PConstExact} is the calculus\textquotesingle s ground-fact
rule: it lets a derivation state an exact closed arithmetic fact ---
$\vdash -3/4 \leq 0$, or $\vdash (2/3)^2 = 4/9$ --- with no premises.
Before it, ground facts were nearly out of reach: \rulename{Norm},
premise-free, already states closed equalities of identical canonical
forms ($\vdash (2/3)^2 = 4/9$ falls to it), but no rule stated a ground
\emph{order} comparison --- $\vdash -3/4 \leq 0$ had no derivation. A
producer uses it
wherever a numeric side condition enters a derivation --- a constant
bound, an exact constant comparison feeding a later combination step.

What the checker recomputes is the entire side condition:
\texttt{evalClosed}, the verified closed-term evaluator, exact over $\Q$,
must return \texttt{true} on the comparison. There is no separate
``closedness'' hypothesis to check: the
evaluator returns \texttt{none} the moment it meets a wire, so a
\texttt{true} result certifies closedness as a byproduct --- a
hypothesis the payload cannot mis-declare. One comparator, $\neq$, is
excluded from the admissible set even though the underlying semantic
lemma covers it: the calculus maintains a second, three-valued soundness
screen, and a premise-free $\neq$ atom would break it (the proof section
unpacks why). Without the exclusion, the rule would silently invalidate
an invariant the rest of the calculus is entitled to assume.

\cataloguenote{The rule generalizes the earlier scalar-constant-fact rule, retained as
the corner case $e_2 = 0$. Its tier is the taxonomy\textquotesingle s
recorded boundary case: admissible by \emph{reduction to verified
machinery} --- the soundness lemma is one step to the
evaluator-correctness theorem --- rather than by unpacking into a
catalogue derivation, because for its order conclusions no such
derivation exists; that
impossibility is the rule\textquotesingle s point. It is not core
because it adds no new kernel semantic content.
}
\ruleblock{Lean statement}
\declhead{BernsteinLogic.Soundness}{Formula.holds\_of\_evalClosed\_true}{\badge{BadgeGreen}{statement bound}}
\begin{leansrc}
/-- Closed formulas that evaluate true hold at every valuation. -/
theorem Formula.holds_of_evalClosed_true {H : Hypernet} (σ : H.Valuation) :
    ∀ (F : Formula), F.evalClosed = some true → F.holds σ
\end{leansrc}
\declhead{BernsteinLogic.Soundness}{sound\_pConstExact}{\badge{BadgeGreen}{statement bound}}
\begin{leansrc}
/-- PConstExact soundness: a closed exact comparison whose closed evaluator
`Formula.evalClosed` decides `some true` holds at every valuation. One step to
the verified `Formula.holds_of_evalClosed_true` — closedness of `e₁`, `e₂` is
IMPLIED by `evalTrue` succeeding (`evalClosed` is `none` on any wire), and the
comparator is decided by exact ℚ evaluation. `pre` is arbitrary (premise-free). -/
theorem sound_pConstExact {H : Hypernet} {seq : Sequent} {e₁ e₂ : Term}
    {b : BowTie}
    (hEval : (Formula.atom b e₁ e₂).evalClosed = some true)
    (hpost : seq.post = .atom b e₁ e₂) :
    seq.semanticallyValid H
\end{leansrc}
\declhead{BernsteinLogic.Core.Catalogue}{Formula.evalClosed}{\badge{BadgeGreen}{statement bound}}
\begin{leansrc}
/-- Closed-formula evaluation: `none` iff a wire or region atom occurs
(strict; see original for the closed-world reading). -/
def Formula.evalClosed : Formula → Option Bool
  | .atom b l r => do pure (b.evalQ (← l.evalQ?) (← r.evalQ?))
  | .inRegion _ _ _ => none
  | .and l r => do pure ((← l.evalClosed) && (← r.evalClosed))
  | .or l r => do pure ((← l.evalClosed) || (← r.evalClosed))
  | .tt => some true
  | .ff => some false

/-- The barycentric coordinate functions of a 1-dimensional simplex, as
canonical affine Terms in `t` (pinned shape `offset + slope·t`; see the
original for the full derivation notes and D2 scope). -/
def baryLambdas1D (cx : SimpC) (dContains : Fin cx.nVertices → Bool)
    (t : Term) : Option (List (Fin cx.nVertices × Term))
\end{leansrc}
\declhead{BernsteinLogic.Core.Catalogue}{PConstExactData}{\badge{BadgeGreen}{statement bound}}
\begin{leansrc}
inductive PConstExactData : Hypernet → Sequent → Type where
  | mk (H : Hypernet) (seq : Sequent)
       (e₁ e₂ : Term) (b : BowTie)
       (notNE : b ≠ .ne)
       (evalTrue : (Formula.atom b e₁ e₂).evalClosed = some true)
       (postEq : seq.post = .atom b e₁ e₂)
       : PConstExactData H seq
\end{leansrc}
\declhead{BernsteinLogic.Admissible}{ProofObject.m3sound}{\badge{BadgeGreen}{statement bound}}
\begin{leansrc}
/-- Monotone replacement at a conjunct path: substituting a
lattice-smaller formula at any conjunct address shrinks the M3 value.
Meet monotonicity only — no distributivity — which is why Promise,
unlike ∨E-ctx, is lattice-sound and core-eligible. -/
theorem mval_conjReplace_mono :
    ∀ {π : ConjPath} {Γ Γ' x q : Formula},
    Γ.conjGet? π = some q → Γ.conjReplace? π x = some Γ' →
    M3.le (mval x) (mval q) = true →
    M3.le (mval Γ') (mval Γ) = true := by
  intro π
  induction π with
  | here =>
    intro Γ Γ' x q hget hrep hle
    simp only [Formula.conjGet?] at hget
    simp only [Formula.conjReplace?] at hrep
    injection hget with h1
    injection hrep with h2
    rw [← h2, h1]
    exact hle
  | left p ih =>
    intro Γ Γ' x q hget hrep hle
    cases Γ with
    | and l r =>
      simp only [Formula.conjGet?] at hget
      simp only [Formula.conjReplace?] at hrep
      cases hmap : l.conjReplace? p x with
      | none => rw [hmap] at hrep; nomatch hrep
      | some l' =>
        rw [hmap] at hrep
        injection hrep with h2
        rw [← h2]
        show M3.le ((mval l').inf (mval r)) ((mval l).inf (mval r)) = true
        exact M3.inf_mono_l _ _ _ (ih hget hmap hle)
    | atom b tl tr => nomatch hrep
    | inRegion t Δ S => nomatch hrep
    | or l r => nomatch hrep
    | tt => nomatch hrep
    | ff => nomatch hrep
  | right p ih =>
    intro Γ Γ' x q hget hrep hle
    cases Γ with
    | and l r =>
      simp only [Formula.conjGet?] at hget
      simp only [Formula.conjReplace?] at hrep
      cases hmap : r.conjReplace? p x with
      | none => rw [hmap] at hrep; nomatch hrep
      | some r' =>
        rw [hmap] at hrep
        injection hrep with h2
        rw [← h2]
        show M3.le ((mval l).inf (mval r')) ((mval l).inf (mval r)) = true
        exact M3.inf_mono_r _ _ _ (ih hget hmap hle)
    | atom b tl tr => nomatch hrep
    | inRegion t Δ S => nomatch hrep
    | or l r => nomatch hrep
    | tt => nomatch hrep
    | ff => nomatch hrep

/-! ## §A9 The metatheorem: derivability is lattice-sound

Over a hypernet with no typed wires, every derivable sequent satisfies
⟦pre⟧ ≤ ⟦post⟧ in M₃. One case per rule of the catalogue; the
propositional rules are the lattice axioms, equality elimination is the
substitution invariance (its typing field forces both terms closed over
an untyped hypernet), the region/arithmetic-closure rules are handled
by their conclusion shapes or payload absurdities, and the
empty-payload rules are vacuous. -/
theorem ProofObject.m3sound : ∀ {H : Hypernet} {seq : Sequent},
    (∀ w, H.wireType w = none) → ProofObject H seq →
    M3.le (mval seq.pre) (mval seq.post) = true
  | _, _, _, .mk _ _ _ (.ax _ _ (.mk _ _ hpp)) => by
      rw [hpp]; exact M3.le_refl _
  | _, _, hE, .mk _ _ _ (.weaken _ _ (.mk _ _ _ pproof _ hpre hpost)) => by
      rw [hpre, hpost]
      exact M3.weaken_kit _ _ _ (pproof.m3sound hE)
  | _, _, hE, .mk _ _ _ (.cut _ _
      (.mk _ _ _ _ lproof _ rproof hpre hpost hcl hcr)) => by
      have hl
\end{leansrc}
\ruleblock{Correspondence (Steele display $\leftrightarrow$ Lean)}
\badge{BadgeGreen}{reviewer-certified against the displayed texts}
\sourcenote{The dictionary below is a reviewer-certified statement of how the mechanized statements above interpret the displayed rule: display metavariable $\mapsto$ Lean binder, each display hypothesis $\mapsto$ the premise or side condition carrying it, the display conclusion $\mapsto$ the goal --- with any deliberate difference in width between display and mechanization named outright. Certification is recorded as a hash triple over (the displayed cell, the extracted statement texts, this dictionary\textquotesingle s own text); if any of the three changes, this block turns amber until a reviewer re-certifies. The correspondence claim itself remains editorial --- reviewer-grade, not machine-checked; its machine-checked successor types each soundness statement against the rule as a structured object, and supersedes this block per rule as it lands.}
\begin{itemize}
\item $e_1$, $e_2$ $\mapsto$ \texttt{e₁ e₂ : Term}; $\bowtie$
  $\mapsto$ \texttt{b : BowTie}; the judgment context $\Phi$ and the
  concluded sequent $\mapsto$ \texttt{seq : Sequent} with the
  conclusion tie \texttt{hpost} $:$ \texttt{seq.post = .atom b e₁ e₂}.
  Premise-freeness is represented by silence: no hypothesis mentions
  \texttt{seq.pre}, so the statement holds under arbitrary premises.
\item the displayed premise ($e_1 \bowtie e_2$ holds under exact
  evaluation in $\Q$) $\mapsto$ \texttt{hEval} $:$
  \texttt{(Formula.atom b e₁ e₂).evalClosed = some true} --- note the
  \texttt{some}: the verified evaluator is partial and returns
  \texttt{none} on any wire, so the display\textquotesingle s separate
  ``wire-free (closed) terms\textquotesingle\textquotesingle{}
  hypothesis is \emph{absorbed}: it has no Lean binder of its own
  because a successful evaluation certifies it.
\item conclusion $\Phi \vdash e_1 \bowtie e_2$ $\mapsto$ the goal
  \texttt{seq.semanticallyValid H}, validity at every valuation of the
  ambient hypernet.
\end{itemize}
Width note, deliberate and load-bearing: the display restricts
$\bowtie \in \{<, \leq, =, \geq, >\}$, while the Lean statement is
proven for \emph{all} constructors of \texttt{BowTie}, including
$\neq$. The comparator exclusion is an admission-layer fact --- the
checker refuses a premise-free $\neq$ atom to protect the three-valued
screen, as the commentary explains --- not a semantic one. The
mechanized statement is therefore strictly wider than the displayed
rule, and the displayed rule is exactly the admitted fragment of it.

\ruleblock{Lean verifier receipts}
\sourcenote{module \texttt{BernsteinLogic.Soundness}; bound build \texttt{a20d90e66c13\ldots} (bound from receipt a9c0e34d3e3c9c8c, 2026-08-03T21:30:36Z). A receipt records that named checking layers accepted a build: (module, compiled-artifact SHA-256, environment SHA-256, checker layers), minted when the checking gate ran; the store receipts attesting the bound build are listed with the layers each run credited. Nothing here asks the reader to trust the store: every layer is an independently runnable program over the quoted sources --- rebuild the module and re-run the checkers to reproduce the verdict. Receipts attesting superseded builds credit no layer and are omitted from this rendering (the count of omissions, if any, is noted below).}
\begin{receipttable}
\receiptrow{a9c0e34d3e3c9c8c}{2026-08-03T21:30:36Z}{build+decorrelated}{\badge{BadgeGreen}{current}}
\end{receipttable}
\sourcenote{21 further receipt(s) in the store attest superseded builds of this module and are omitted from this rendering.}
\sourcenote{module \texttt{BernsteinLogic.Core.Catalogue}; bound build \texttt{cbc92e976a1d\ldots} (bound from receipt 5abdb3589707b5df, 2026-08-03T21:30:33Z). A receipt records that named checking layers accepted a build: (module, compiled-artifact SHA-256, environment SHA-256, checker layers), minted when the checking gate ran; the store receipts attesting the bound build are listed with the layers each run credited. Nothing here asks the reader to trust the store: every layer is an independently runnable program over the quoted sources --- rebuild the module and re-run the checkers to reproduce the verdict. Receipts attesting superseded builds credit no layer and are omitted from this rendering (the count of omissions, if any, is noted below).}
\begin{receipttable}
\receiptrow{24804e66edba0bc9}{2026-08-03T20:25:09Z}{build+decorrelated}{\badge{BadgeGreen}{current}}
\receiptrow{5abdb3589707b5df}{2026-08-03T21:30:33Z}{build+decorrelated}{\badge{BadgeGreen}{current}}
\end{receipttable}
\sourcenote{18 further receipt(s) in the store attest superseded builds of this module and are omitted from this rendering.}
\sourcenote{module \texttt{BernsteinLogic.Admissible}; bound build \texttt{52e5099ba522\ldots} (bound from receipt c17abf2e103fe745, 2026-07-21T05:28:22Z). A receipt records that named checking layers accepted a build: (module, compiled-artifact SHA-256, environment SHA-256, checker layers), minted when the checking gate ran; the store receipts attesting the bound build are listed with the layers each run credited. Nothing here asks the reader to trust the store: every layer is an independently runnable program over the quoted sources --- rebuild the module and re-run the checkers to reproduce the verdict. Receipts attesting superseded builds credit no layer and are omitted from this rendering (the count of omissions, if any, is noted below).}
\begin{receipttable}
\receiptrow{7f3043c54f717472}{2026-07-20T07:43:38Z}{build+decorrelated}{\badge{BadgeGreen}{current}}
\receiptrow{acaf0e8871731343}{2026-07-20T20:44:28Z}{build+decorrelated}{\badge{BadgeGreen}{current}}
\receiptrow{62974466ea131d89}{2026-07-20T21:00:57Z}{build+replay+hardened+decorrelated}{\badge{BadgeGreen}{current}}
\receiptrow{c17abf2e103fe745}{2026-07-21T05:28:22Z}{build+replay+hardened+decorrelated}{\badge{BadgeGreen}{current}}
\end{receipttable}
\sourcenote{6 further receipt(s) in the store attest superseded builds of this module and are omitted from this rendering.}
\layerline{\layeritem{build}{lake build (elaboration)}{\badge{BadgeGreen}{current}}\layeritem{replay}{leanchecker (official kernel replay)}{\badge{BadgeGreen}{current}}\layeritem{hardened}{safe\_verify (hardened screen)}{\badge{BadgeGreen}{current}}\layeritem{decorrelated}{lean4export + nanoda (decorrelated kernel)}{\badge{BadgeGreen}{current}}}
\sourcenote{\badge{BadgeGray}{comparator: none} no fixture in the three-implementation comparator harness (compiled / kernel-interpreted / reference-interpreted) exercises this rule.}
\ruleblock{HOL4 statement}
\sourcenote{theory \texttt{bern\_tr\_catdata}, HOL4 development script \texttt{\texttt{bern\_tr\_catdataScript.sml}} \badge{BadgeGreen}{script bound}}
\holhead{Definition pconstexactdataOk\_def}
\begin{holsrc}
Definition pconstexactdataOk_def:
  pconstexactdataOk (sq:sequent) (d:pconstexactdata) ⇔
    d.ceBow ≠ BowNe ∧
    formulaEvalClosed (FAtom d.ceBow d.ceL d.ceR) = SOME T ∧
    sq.post = FAtom d.ceBow d.ceL d.ceR
End
\end{holsrc}
\holhead{Theorem pconstexactdataOk\_receipt}
\begin{holsrc}
Theorem pconstexactdataOk_receipt:
  pconstexactdataOk
    <| pre := FTt ; post := FAtom BowGt (TmConst 1) (TmConst 0) |>
    <| ceBow := BowGt ; ceL := TmConst 1 ; ceR := TmConst 0 |> ∧
  ¬ pconstexactdataOk
    <| pre := FTt ; post := FAtom BowNe (TmConst 1) (TmConst 0) |>
    <| ceBow := BowNe ; ceL := TmConst 1 ; ceR := TmConst 0 |>
\end{holsrc}
\ruleblock{HOL verifier receipts}
\badge{BadgeBlue}{documentary} \texttt{hol-strict} CLEAN: theory \texttt{bern\_tr\_catdata}, 252
theorems, 0 cheat-tagged, 0 nonstandard axioms (whitelist:
\texttt{DISK\_THM}; the four \texttt{bool} axioms) --- verification
run of 2026-07-22 on an isolated build host, against a
content-addressed checkout of the exact sources cited here (git blob
\texttt{f8fceac5\ldots}). The closed-comparison payload predicate
displayed above is the rule\textquotesingle s entire checkable
content --- comparator confined to the admissible five, both terms
evaluated by the closed exact-$\Q$ evaluator, conclusion pinned to
the payload --- with a kernel-evaluated accept/refusal receipt. The
proof-object stratum consumes it verbatim as this
rule\textquotesingle s arm of the mutual wellformedness predicate
(theory \texttt{bern\_tr\_catdata2}, 272 theorems, 0 offences, same
run), and the rule-tag roster carries the rule. As on the Lean side,
the rule rides the proof-object family, not the first-order wire
format: the wire tag remains deferred on both sides (the emit path
refuses until allocated), so there is deliberately no column-side
check to mirror. Honest residue: the semantic reduction (the Lean
development\textquotesingle s \texttt{sound\_pConstExact}, one step
to the verified closed evaluator) is not yet ported.

\sourcenote{HOL receipts are documentary: quoted results of \texttt{hol-strict} runs
and census rows from the HOL4 development\textquotesingle s build
evidence, not minted receipt artefacts. A HOL receipt store (the Lean
store\textquotesingle s twin) is named future work.
}
\ruleblock{Proof}

\sourcenote{Genre: admissible rule (tier 2) --- a proof-tree-rewriting
recipe. The tier taxonomy defines tier 2 as \emph{verified unpacking
into core rules} (one node $\to$ many); this rule is the taxonomy's
recorded boundary case --- admissible by \emph{reduction to verified
machinery}. The general template, both species included, is stated in
the \rulename{Norm} section; the recipe below typesets the honest
degenerate instance.}

\begin{brecipe}[the \rulename{PConstExact} instance --- the degenerate rewrite]
\label{rec:pce-instance}
\rulename{PConstExact} has \emph{no proof premises} ($n = 0$), and its
rewrite target contains \emph{zero catalogue nodes}: no catalogue
derivation states a ground \emph{order} comparison --- closed equalities
of identical canonical forms fall to \rulename{Norm} already, but
$\vdash -3/4 \leq 0$ had no derivation before this rule --- which is the
rule's entire point. The rewrite therefore discharges the node against a
named \emph{verified theorem} rather than against catalogue structure:
\[
\inferrule*[Right=\rulename{PConstExact}]{\;e_1 \bowtie e_2 \ \text{holds under exact evaluation in } \Q\;}{\Phi \vdash e_1 \bowtie e_2}
\qquad\rightsquigarrow\qquad
\text{\upshape the checker replays } \mathrm{evalClosed};\;
\text{\texttt{holds\_of\_evalClosed\_true} stands as warrant.}
\]
The side condition is \emph{computational} (exact $\Q$ evaluation of a
closed comparison --- decidable, replayed by the checker), and the warrant
theorem is kernel-verified once, globally: soundness of the rule is
literally ``the exact evaluator is correct.''
\end{brecipe}

\begin{bprop}[semantic warrant]\label{prop:pce-sound}
If $(\mathrm{atom}\ \bowtie\ e_1\ e_2).\mathrm{evalClosed} =
\mathsf{some\ true}$ and the sequent's conclusion is that atom, the
sequent is semantically valid over any hypernet, with arbitrary premise
$\Phi$.
\end{bprop}
\begin{proof}
Fix any valuation $\sigma$ (admissibility and $\sigma \models \Phi$ are
not consulted --- the rule is premise-free, the \rulename{Norm}
convention). By the verified evaluator theorem
\texttt{Formula.holds\_of\_evalClosed\_true}, a formula whose closed
evaluation returns $\mathsf{true}$ holds at \emph{every} valuation: a
closed term's evaluation consults no wire, so $\sigma$ is never read, and
the comparison was decided by exact arithmetic over $\Q$. That is the
whole proof --- \texttt{sound\_pConstExact} is one \texttt{exact} step.
\end{proof}

\begin{bremark}[closedness is implied, not declared]
There is no ``$e_1, e_2$ closed'' payload field to check:
$\mathrm{evalClosed}$ returns $\mathsf{none}$ on any wire occurrence, so
$\mathsf{some\ true}$ \emph{certifies} closedness as a byproduct. A
well-formedness hypothesis that the evaluator itself enforces is one the
payload cannot mis-declare --- unrepresentable, not discouraged.
\end{bremark}

\begin{bremark}[why $\neq$ is excluded --- the $M_3$ countermodel]
The comparator set is confined to $\{<, \leq, =, \geq, >\}$: five of the
six. The semantic lemma covers all six, but the calculus also maintains a
three-valued screen (\texttt{ProofObject.m3sound}), whose disequality
clause reads designated free-wire membership --- so a wire-free $\neq$
atom, its wire list empty, interprets to $\bot$ under the $M_3$
countermodel. Admitting
$\neq$ against arbitrary $\Phi$ would then let a premise-free rule
conclude screen value $\bot$ under a premise valued $\top$ ($\vdash 1
\neq 2$ is true over $\Q$, and breaks the screen's
premise-$\leq$-conclusion invariant), breaking \texttt{m3sound}. The
\texttt{notNE} payload field is shared with \rulename{ArithAtom}, where
the same gate carries different weight: both of that rule's spellings
mention the same wires, so its screen value is preserved and the gate
serves alignment (that section's closing remark); here the breakage is
real. A
hypothesis whose necessity comes not from $\Q$-semantics but from a
\emph{second} soundness invariant is exactly the kind a recipe must state,
or the next porter re-learns it the hard way.
\end{bremark}

\begin{bremark}[the subsumed corner, and coverage]
\rulename{PRatFact} (the scalar constant fact $c \bowtie 0$) is the corner
$e_2 = 0$ of this rule and was subsumed when the roster was consolidated. Coverage today:
the \texttt{Term} grammar makes every closed term rank-0, so the rule
states every closed \emph{scalar} exact comparison; the matrix corner
(exact closed matrix identities, MIX's use case) needs closed-tensor
constructors and a \texttt{QTensor}-valued closed evaluator --- the named
next leg, deliberately not hand-faked here.
\end{bremark}

\begin{anchors}
Recipe \ref{rec:pce-instance} warrant & \texttt{sound\_pConstExact} (one step to \texttt{Formula.holds\_of\_evalClosed\_true})\\
closedness-is-implied & \texttt{Formula.evalClosed} ($\mathsf{none}$ on any wire)\\
the $\neq$ exclusion & \texttt{notNE} payload field; \texttt{ProofObject.m3sound}\\
premise-freedom ($\Phi$ arbitrary) & the \rulename{Norm} convention; \texttt{pre} unused in \texttt{sound\_pConstExact}\\
wire placement & \texttt{RuleTag.pConstExact}; wire tag not yet allocated --- the emit path refuses the rule until allocation\\
\end{anchors}


\rulesection{27}{PBound}{alias for \textcolor{DiffBlue}{\rulename{Syzygy}}}
\label{blcatrule-PBound}
\ruleblock{The rule (Steele metanotation)}
\sourcenote{as realized in the mechanized development; source row: \texttt{\texttt{rules.toml}} --- \emph{PBound}. The metanotation statement is typeset by hand for notational quality; its mechanized ground truth is the Lean statement below, extracted from the sources at every build of this book and hash-bound. The rule \emph{roster} is machine-checked (every mechanized rule tag round-trips to exactly one catalogue row); the statement \emph{text} is editorial, checked by eye against the extracted Lean beside it.}
\begin{mathpar}
\inferrule*[Right=\rulename{PBound}]{
  \bigl\{\Phi_i \vdash e_i \leq 0\bigr\}_{i=1}^{n}, \quad \Phi \supseteq \Phi_i \ \text{conjunct-wise (checked)}
  \\
  \bigl\{g_k = 0\bigr\}_k \ \text{equality sub-conjuncts of } \Phi
  \\
  \lambda_i \in \Q_{\geq 0}
  \\
  \nu_k \in \Q
  \\
  c \equiv \textstyle\sum_i \lambda_i e_i + \sum_k \nu_k g_k \ \text{(exact canonical polynomial)}
  \\
  \bshape_H(c) \neq \bot, \quad \wires(c) \ \text{all finitely typed in } H
  \\
  \wires(\Phi) \cup \wires(c) \cup \textstyle\bigcup_i \wires(e_i) \ \text{all typed at positive-size shapes in } H
  \\
  \Psi \in \bigl\{\, c \leq 0, \;\; e \geq 0 \text{ where } c = -e, \;\; l \leq r \text{ where } c = l - r \,\bigr\}
}{
  \Phi \vdash \Psi
}
\end{mathpar}

\par\medskip
\rulename{PBound} is the workhorse bound-node: a compact spelling of a
\rulename{Syzygy} application whose entries the register prices as
affine, with the
conclusion exposed through the three-way reading selector --- the raw
zero-form $c \leq 0$, the sign-flipped $e \geq 0$ where $c = -e$, or
the literal atom $l \leq r$ where $c = l - r$. It is the
taxonomy\textquotesingle s canonical alias: the payload is
\rulename{Syzygy}\textquotesingle s field-for-field, the unpack is the
identity transport, and the wire format prices the affine register
compactly. A producer emits it constantly: it is how an affine
consequence of proven bounds becomes a first-class node whose
conclusion can feed the rules that demand a particular atom shape ---
the $\geq$-reading makes a nonnegativity premise reachable from bound
facts without a detour, and the literal-atom reading feeds
interval-shaped consumers with no separate re-arrangement step.

What the checker recomputes is exactly
\rulename{Syzygy}\textquotesingle s obligation set --- entry decoding
against the pinned combination encoding, nonnegative weights,
entry-sourcing from component conclusions and premise equality
sub-conjuncts (both polarities on equalities), and the exact
canonical-polynomial fold identity --- plus the same bookkeeping
(conjunct-wise premise containment, conclusion shape and finiteness,
positive-size wires on premise, conclusion, and component conclusions).
That section\textquotesingle s commentary on what breaks without each
hypothesis applies verbatim here, through the identity transport.

The affine restriction is the
rule\textquotesingle s published register (the roster note\textquotesingle s
pricing vocabulary --- wire-format compactness
for coefficient rows); it appears in neither the displayed premises nor
the statement, and it is not a checked side condition: the replay
is degree-blind and soundness never consults linearity. The three
readings are one exact $\Q$ identity re-presented under a sign flip;
their semantic transport is proven once, in the target
rule\textquotesingle s section, and inherited here as a pure
re-spelling.

\cataloguenote{The wire tag for this rule is deferred to a payload-window allocation;
the emit path refuses the rule until the tag is allocated. The
compact-register pricing is the reason the alias exists as a distinct
spelling at all --- semantically it adds nothing to its target, which
is the alias genre working as intended.
}
\ruleblock{Lean statement}
\declhead{BernsteinLogic.Soundness}{holds\_of\_syzygyConclusionForm}{\badge{BadgeGreen}{statement bound}}
\begin{leansrc}
/-- The ge-bridge ring lemma (the both-polarity core change): every
reading admitted by
`syzygyConclusionForm conclusion` holds whenever the zero-normal atom
`conclusion ≤ 0` holds — the readings are the same exact ℚ identity
re-presented under a sign flip. The zero-form is the identity transport,
the ge-polarity flip is `-e ≤ 0 → e ≥ 0`, and the literal atom is
`l - r ≤ 0 → l ≤ r` (ArithAtom's tensor-general sub-zero bridge).
Applied once inside `sound_syzygy`, so PSyzygy (tier-1), PBound's
compact spelling (tier-3), and any future ≤-0-shaped fact rule inherit
all readings with no separate polarity justification. -/
theorem holds_of_syzygyConclusionForm {H : Hypernet} {σ : H.Valuation}
    {conclusion : Term} {F : Formula}
    (hle : Formula.holds σ (.atom .le conclusion (.const 0)))
    (hform : syzygyConclusionForm conclusion F) :
    Formula.holds σ F
\end{leansrc}
\declhead{BernsteinLogic.Soundness}{sound\_syzygy}{\badge{BadgeGreen}{statement bound}}
\begin{leansrc}
/-- Soundness of Syzygy (= Lemma B, dialect D): valid components give a
valid composite. The argument is `sound_farkas`'s, with the per-entry
source discharge generalized: a component-sourced entry's atom holds by
component validity (through `premiseContains`), a wiring-sourced entry's
atom holds as a premise sub-conjunct; either way
`atomLeZeroPolys_sound_entryC` delivers per-entry nonpositivity at every
clamped per-entry valuation, and the nonnegative fold closes. The
zero-normal bound then reaches `seq.post` through whichever
`syzygyConclusionForm` reading the payload declared (both-polarity core
change, 2026-07-18). -/
theorem sound_syzygy {H : Hypernet} (st : SyzygyStep H)
    (hcomps : ∀ c ∈ st.components, c.semanticallyValid H) :
    st.seq.semanticallyValid H
\end{leansrc}
\declhead{BernsteinLogic.Soundness}{sound\_pBound}{\badge{BadgeGreen}{statement bound}}
\begin{leansrc}
/-- PBound soundness: INHERITED from `sound_syzygy`, literally — the payload's
fields are `PSyzygyData`'s verbatim (tier-3 compact spelling, 1 node → 1 node),
so the proof is the same `SyzygyStep` handed to the same theorem. Polarity is
not justified here: since the both-polarity core change,
`sound_syzygy` itself carries every
`syzygyConclusionForm` reading through `holds_of_syzygyConclusionForm`, and
PBound inherits them as a pure re-spelling. -/
theorem sound_pBound {H : Hypernet} {seq : Sequent} {components : List Sequent}
    {combination conclusion : Term}
    (hcomps : ∀ c ∈ components, c.semanticallyValid H)
    (hpreC : ∀ c ∈ components, premiseContains seq.pre c.pre)
    (hshape : ∃ s, conclusion.bshape H = some s)
    (hfin : conclusion.finWires H)
    (hwpos : ∀ w ∈ seq.pre.freeWires ++ conclusion.freeWires,
      ∃ τ, H.wireType w = some τ ∧ 0 < Shape.totalSize τ.shape)
    (hcompPos : ∀ c ∈ components, ∀ w ∈ c.post.freeWires,
      ∃ τ, H.wireType w = some τ ∧ 0 < Shape.totalSize τ.shape)
    (hvalid : syzygyCombinationValid combination components seq.pre conclusion)
    (hform : syzygyConclusionForm conclusion seq.post) :
    seq.semanticallyValid H
\end{leansrc}
\declhead{BernsteinLogic.Core.Catalogue}{PBoundData}{\badge{BadgeGreen}{statement bound}}
\begin{leansrc}
inductive PBoundData : Hypernet → Sequent → Type where
  | mk (H : Hypernet) (seq : Sequent)
       (components : List Sequent)
       (componentProofs : ProofSeqList H components)
       (combination : Term)
       (conclusion : Term)
       (denBoundLog2 : Nat)
       (preContains : ∀ c ∈ components, premiseContains seq.pre c.pre)
       (conclShape : ∃ s, conclusion.bshape H = some s)
       (conclFin   : conclusion.finWires H)
       (wiresPos : ∀ w ∈ seq.pre.freeWires ++ conclusion.freeWires,
                     ∃ τ, H.wireType w = some τ ∧
                       0 < Shape.totalSize τ.shape)
       (compPostWiresPos : ∀ c ∈ components, ∀ w ∈ c.post.freeWires,
                     ∃ τ, H.wireType w = some τ ∧
                       0 < Shape.totalSize τ.shape)
       (valid : syzygyCombinationValid combination components seq.pre conclusion)
       (postForm : syzygyConclusionForm conclusion seq.post)
       : PBoundData H seq
\end{leansrc}
\declhead{BernsteinLogic.Core.Catalogue}{syzygyConclusionForm}{\badge{BadgeGreen}{definition bound}}
\begin{leansrc}
def syzygyConclusionForm (conclusion : Term) (F : Formula) : Prop :=
  F = .atom .le conclusion (.const 0)
  ∨ (∃ e, conclusion = .neg e ∧ F = .atom .ge e (.const 0))
  ∨ (∃ l r, conclusion = .sub l r ∧ F = .atom .le l r)
\end{leansrc}
\declhead{BernsteinLogic.Core.Catalogue}{syzygyCombinationValid}{\badge{BadgeGreen}{definition bound}}
\begin{leansrc}
def syzygyCombinationValid (cc : Term) (components : List Sequent)
    (premise : Formula) (conclusion : Term) : Prop :=
  syzygyCheck cc components premise conclusion = true
\end{leansrc}
\ruleblock{Correspondence (Steele display $\leftrightarrow$ Lean)}
\badge{BadgeGreen}{reviewer-certified against the displayed texts}
\sourcenote{The dictionary below is a reviewer-certified statement of how the mechanized statements above interpret the displayed rule: display metavariable $\mapsto$ Lean binder, each display hypothesis $\mapsto$ the premise or side condition carrying it, the display conclusion $\mapsto$ the goal --- with any deliberate difference in width between display and mechanization named outright. Certification is recorded as a hash triple over (the displayed cell, the extracted statement texts, this dictionary\textquotesingle s own text); if any of the three changes, this block turns amber until a reviewer re-certifies. The correspondence claim itself remains editorial --- reviewer-grade, not machine-checked; its machine-checked successor types each soundness statement against the rule as a structured object, and supersedes this block per rule as it lands.}
Verbatim the \rulename{Syzygy} dictionary, unbundled: where
\rulename{Syzygy}\textquotesingle s statement packs its obligations
into a \texttt{SyzygyStep}, \texttt{sound\_pBound} takes the same
eleven fields as flat binders --- \texttt{hpreC}, \texttt{hshape},
\texttt{hfin}, \texttt{hwpos}, \texttt{hcompPos}, \texttt{hvalid},
\texttt{hform}, over the same \texttt{components}/\texttt{combination}/
\texttt{conclusion} data with \texttt{hcomps} the component
validities --- and its entire proof is one application of
\texttt{sound\_syzygy} to the repacked Step, visible in the displayed
statement. The alias relation is thereby itself certified: 1 $\to$ 1
at statement level, and (since the both-polarity core change) 1 $\to$
1 at reading level too --- the displayed $\Psi$ selector is the same
\texttt{syzygyConclusionForm} the target consumes. This cell displays
the evaluability and positive-size rows inline; their mapping is the
\rulename{Syzygy} dictionary\textquotesingle s, unchanged. No hypothesis differs, none
is unused, and neither side is wider --- the affine restriction the
roster note mentions is wire-format compactness, absent from both the
display\textquotesingle s premises and the statement.

\ruleblock{Lean verifier receipts}
\sourcenote{module \texttt{BernsteinLogic.Soundness}; bound build \texttt{a20d90e66c13\ldots} (bound from receipt a9c0e34d3e3c9c8c, 2026-08-03T21:30:36Z). A receipt records that named checking layers accepted a build: (module, compiled-artifact SHA-256, environment SHA-256, checker layers), minted when the checking gate ran; the store receipts attesting the bound build are listed with the layers each run credited. Nothing here asks the reader to trust the store: every layer is an independently runnable program over the quoted sources --- rebuild the module and re-run the checkers to reproduce the verdict. Receipts attesting superseded builds credit no layer and are omitted from this rendering (the count of omissions, if any, is noted below).}
\begin{receipttable}
\receiptrow{a9c0e34d3e3c9c8c}{2026-08-03T21:30:36Z}{build+decorrelated}{\badge{BadgeGreen}{current}}
\end{receipttable}
\sourcenote{21 further receipt(s) in the store attest superseded builds of this module and are omitted from this rendering.}
\sourcenote{module \texttt{BernsteinLogic.Core.Catalogue}; bound build \texttt{cbc92e976a1d\ldots} (bound from receipt 5abdb3589707b5df, 2026-08-03T21:30:33Z). A receipt records that named checking layers accepted a build: (module, compiled-artifact SHA-256, environment SHA-256, checker layers), minted when the checking gate ran; the store receipts attesting the bound build are listed with the layers each run credited. Nothing here asks the reader to trust the store: every layer is an independently runnable program over the quoted sources --- rebuild the module and re-run the checkers to reproduce the verdict. Receipts attesting superseded builds credit no layer and are omitted from this rendering (the count of omissions, if any, is noted below).}
\begin{receipttable}
\receiptrow{24804e66edba0bc9}{2026-08-03T20:25:09Z}{build+decorrelated}{\badge{BadgeGreen}{current}}
\receiptrow{5abdb3589707b5df}{2026-08-03T21:30:33Z}{build+decorrelated}{\badge{BadgeGreen}{current}}
\end{receipttable}
\sourcenote{18 further receipt(s) in the store attest superseded builds of this module and are omitted from this rendering.}
\layerline{\layeritem{build}{lake build (elaboration)}{\badge{BadgeGreen}{current}}\layeritem{replay}{leanchecker (official kernel replay)}{\badge{BadgeAmber}{stale only}}\layeritem{hardened}{safe\_verify (hardened screen)}{\badge{BadgeAmber}{stale only}}\layeritem{decorrelated}{lean4export + nanoda (decorrelated kernel)}{\badge{BadgeGreen}{current}}}
\sourcenote{\badge{BadgeGray}{comparator: none} no fixture in the three-implementation comparator harness (compiled / kernel-interpreted / reference-interpreted) exercises this rule.}
\ruleblock{HOL4 statement}
\sourcenote{theory \texttt{bern\_catalogue}, HOL4 development script \texttt{\texttt{bern\_catalogueScript.sml}} \badge{BadgeGreen}{script bound}}
\holhead{Definition syzygyCheck\_def}
\begin{holsrc}
Definition syzygyCheck_def:
  syzygyCheck cc components premise conclusion ⇔
    syzygyCheckCore cc components premise conclusion ∨
    syzygyCheckW cc components premise conclusion
End
\end{holsrc}
\ruleblock{HOL verifier receipts}
\badge{BadgeBlue}{documentary} \texttt{hol-strict} CLEAN over the ported checker chain: theory
\texttt{bern\_catalogue}, 472 theorems, 0 cheat-tagged, 0 nonstandard
axioms (whitelist: \texttt{DISK\_THM}; the four \texttt{bool} axioms) ---
verification run of 2026-07-22 on an isolated build host, scripts
sha-matched verbatim to the sources cited here (git blob
\texttt{e41b9ed0\ldots}). Scope, stated honestly: the
alias\textquotesingle s entire combination check \emph{is} the target
rule\textquotesingle s, so the ported \texttt{syzygyCheck} (with its
accept and two reject receipt rows, quoted in the target
rule\textquotesingle s section) covers this rule\textquotesingle s
computational content; the conclusion-form reading selector and the
payload layer ride the dependently-typed proof-object surface, which is
a separately tracked porting stage. As on the \rulename{Farkas} page:
the Lean check has since grown a third (mixed tagged-union) branch,
added 2026-07-19 after this port run --- the divergence is
refusal-direction only, and its port is named residue.
Semantic stratum, also ported: the soundness twin \texttt{sound\_pBound}
(theory \texttt{bern\_tr\_soundsyz}: \texttt{hol-strict} clean, 7
theorems, 0 offences, verification run of 2026-07-22 on an isolated
build host, content-addressed at the same sources); the certified-
fragment keystone (\texttt{checkPf\_sound\_frag}, theory
\texttt{bern\_tr\_splice}, 20 theorems, 0 offences, same run) consumes
it at the \rulename{Syzygy} wire arm, this rule\textquotesingle s target
relation.

\sourcenote{HOL receipts are documentary: quoted results of \texttt{hol-strict} runs
and census rows from the HOL4 development\textquotesingle s build
evidence, not minted receipt artefacts. A HOL receipt store (the Lean
store\textquotesingle s twin) is named future work.
}
\ruleblock{Proof}

\sourcenote{Genre: alias rule (tier 3) --- a respelling template. An alias
owes the book five artefacts: the verified 1$\to$1 unpack, the filled-in
application of the target rule, the payload-transport table, the exact
algebraic identity that makes the respelling correct, and the soundness
corollary through the target rule's lemma. All five follow --- and this
alias is the tier taxonomy's \emph{canonical} case, the one the
tier was named for: the unpack is the identity.}

\begin{brecipe}[the alias template]\label{rec:pb-template}
An alias rule $A$ is warranted by a verified function
$\mathrm{unpack} : \mathrm{Step}_A \to \mathrm{Step}_T$ into the payload
of a \emph{single} target node $T$ (core or admissible), such that the
$A$-node and its image conclude the same sequent from the same premises:
one node $\to$ one node, payload transported, nothing derived. Soundness
is then a corollary, $\mathrm{sound}_A = \mathrm{sound}_T \circ
\mathrm{unpack}$. The wire format may carry $A$ as a compact register of
its own or never carry it at all (producers elaborating before
emission, as the \rulename{SmallGain} section records); either way the
checker's obligations are the target's, reached through the unpack.
This rule takes the second branch in v0: no producer emits the
once-contemplated affine coefficient-row register, so no tag is
allocated --- and the identity unpack makes a register pure wire
compactness, purchasable later without semantic motion.
\end{brecipe}

\begin{brecipe}[the verified unpack --- the identity transport]
\label{rec:pb-unpack}
\rulename{PBound}'s payload carries components, combination term,
conclusion term, and the bookkeeping fields \emph{verbatim} as the
target's payload does, including the conclusion-form selector. The
unpack is therefore not a computation but a re-labelling: the verified
object is the direct specialization
\[
\mathrm{sound}_{\rulename{PBound}}
\;=\;
\mathrm{sound}_{\rulename{Syzygy}}
\ \text{\upshape applied to the same fields, re-assembled as the
target's step}
\]
--- in the sources, the body of \texttt{sound\_pBound} \emph{is} one
application of \texttt{sound\_syzygy} to the step built from
\rulename{PBound}'s own fields. Where \rulename{SmallGain}'s unpack
computes a weighting, this one computes nothing; that is the tier-3
boundary case made flesh, and it is why the alias relation here is
$1 \to 1$ at the level of \emph{readings} too
(Recipe~\ref{rec:pb-transport}).
\end{brecipe}

\begin{brecipe}[the filled-in application]\label{rec:pb-filled}
Under the identity transport the \rulename{PBound} node \emph{is} the
following application of the target rule (compare the rule display in
the \rulename{Syzygy} section, field for field):
\begin{mathpar}
\inferrule*[Right=\rulename{Syzygy}]{
  \bigl\{\Phi_i \vdash e_i \leq 0\bigr\}_{i=1}^{n}, \quad
  \Phi \supseteq \Phi_i \ \text{conjunct-wise (checked)} \\
  \bigl\{g_k = 0\bigr\}_k \ \text{sub-conjuncts of } \Phi
  \ \text{in wire variables} \\\\
  \lambda_i \in \Q_{\geq 0} \\
  \nu_k \in \Q \\
  c \;\equiv\; \textstyle\sum_i \lambda_i e_i + \sum_k \nu_k g_k
  \ \text{(exact canonical polynomial)} \\\\
  \Psi \in \bigl\{\, c \leq 0, \;\; e \geq 0 \text{ where } c = -e,
  \;\; l \leq r \text{ where } c = l - r \,\bigr\}
}{
  \Phi \vdash \Psi
}
\end{mathpar}
where $\Psi$ is whichever reading of $c \leq 0$ the shared selector
declares: \emph{(Z$\leq$)} $c \leq 0$ itself; \emph{(G)} $e \geq 0$
where $c = -e$; \emph{(L$\leq$)} $l \leq r$ where $c = l - r$.
Discharging the target's obligations at this instance is immediate ---
they are \emph{the same obligations on the same data}: entry decoding,
nonnegative weights, entry-sourcing (component conclusions and premise
equality sub-conjuncts, both polarities on equalities), the exact fold
identity, premise containment, conclusion shape/finiteness, and the
positive-size wire conditions on premise, conclusion, and component
conclusions --- the last three families displayed inline in
\rulename{PBound}'s own cell, while in \rulename{Syzygy}'s cell the
shape/finiteness and positive-size conditions ride the payload (its
display carries containment as a trailing line); the obligations are
identical either way, only the display convention differs. \rulename{PBound}'s published register prices the
$e_i$ as \emph{affine} entries (the roster's register note); the
displayed premises carry no such restriction, and nothing above
consults one
(\autoref{rem:pb-affine}).
\end{brecipe}

\begin{brecipe}[payload transport --- field for field]
\label{rec:pb-transport}
\begin{center}\small
\begin{tabular}{@{}ll@{}}
\textbf{\rulename{PBound} payload} & \textbf{target (\rulename{Syzygy}) payload}\\
\midrule
\texttt{components} & \texttt{components} (verbatim)\\
\texttt{combination} & \texttt{combination} (verbatim)\\
\texttt{conclusion} & \texttt{conclusion} (verbatim)\\
\texttt{denBoundLog2} & \texttt{denBoundLog2} (verbatim; witness-side claim, zero proof-theoretic weight)\\
\texttt{preContains} & \texttt{preContains} (verbatim)\\
\texttt{conclShape, conclFin} & \texttt{conclShape, conclFin} (verbatim)\\
\texttt{wiresPos, compPostWiresPos} & \texttt{wiresPos, compPostWiresPos} (verbatim)\\
\texttt{valid} & \texttt{valid} (the same exact-combination proposition)\\
\texttt{postForm} & \texttt{postForm} (the same three-reading selector)\\
\end{tabular}
\end{center}
Every row is the identity --- including the last: since the target rule
itself carries the three-reading conclusion selector, the alias
inherits all readings with no residue, and the transport is $1 \to 1$
at reading level, not merely at node level.
\end{brecipe}

\begin{bprop}[the exactness identity]\label{prop:pb-exact}
The respelling is exact definitionally: the unpacked node's combination
identity is \emph{the same} exact canonical-polynomial identity
$c \equiv \sum_i \lambda_i e_i + \sum_k \nu_k g_k$, unchanged. What
remains to make the three readings correct is the fact that they are
one exact $\Q$ identity re-presented under a sign flip --- per entry,
$-e(\sigma) \leq 0 \iff e(\sigma) \geq 0$ and
$l(\sigma) - r(\sigma) \leq 0 \iff l(\sigma) \leq r(\sigma)$ --- which
is the target rule's reading-transport lemma, proven once in the
\rulename{Syzygy} section and \emph{not} re-justified here.
\end{bprop}
\begin{proof}
The first claim is reflexivity of the transported field. The second is
the cited lemma (the ordered-field axioms negate-and-flip and
add-to-both-sides, packaged as the ge-bridge ring lemma); the alias
contributes no polarity argument of its own, which is what keeps the
$1 \to 1$ relation clean.
\end{proof}

\begin{bprop}[soundness corollary]\label{prop:pb-sound}
If every component sequent is semantically valid, then under the
payload's checked hypotheses the sequent $\Phi \vdash \Psi$ is
semantically valid, for every reading $\Psi$ the selector admits.
\end{bprop}
\begin{proof}
\texttt{sound\_pBound} $=$ \texttt{sound\_syzygy} at the transported
(identical) payload data, fed the same component validities. There is
no other content --- the alias genre working as intended.
\end{proof}

\begin{bremark}[why the rule exists at all]\label{rem:pb-why}
Semantically the alias adds nothing to its target; its existence is
pragmatic, twice over. \emph{Register}: affine entries admit a compact
coefficient-row wire spelling, and bound nodes are the bulk commodity
of produced certificates --- the compact register is worth a tag.
\emph{Readings at the seams}: the (G) reading makes a nonnegativity
premise (as \rulename{ReLU-Pos} demands, $h \geq 0$) reachable
directly from bound facts --- the once-contemplated standalone
$\geq$-bridge rule dissolved into the conclusion selector instead ---
and the (L$\leq$) reading hands interval-shaped consumers
(\rulename{LiRPA}'s endpoint premises) their literal atom with no
\rulename{ArithAtom} re-arrangement node in between. Producers get the
convenient conclusion shape; the kernel keeps exactly one combination
argument.
\end{bremark}

\begin{bremark}[the affine restriction is register, not soundness]
\label{rem:pb-affine}
The register's ``affine entries'' --- the roster note's pricing
vocabulary, appearing in neither the displayed premises nor the
statement --- names the wire spelling
the rule is priced for, and is not a checked side condition: the
combination replay is degree-blind, and the soundness argument
(Proposition~\ref{prop:pb-sound}) never consults linearity. A
higher-degree entry simply belongs to the target rule's own spelling
--- what would refuse it is the published register\textquotesingle s
compact coefficient-row spelling, an encoding fact of that format
(asserted at its register, not displayed here), never the checker.
Publishing the restriction beside the rule, while the checker enforces
only what soundness needs, mirrors the \rulename{Farkas} page's
linearity note --- the register tells producers what the format prices;
the checker tells them what truth requires.
\end{bremark}

\begin{anchors}
Recipe \ref{rec:pb-unpack} unpack (identity transport) & \texttt{sound\_pBound}'s body: one application of \texttt{sound\_syzygy} to the step assembled from \rulename{PBound}'s fields (\texttt{PBoundData} $=$ the target's fields verbatim)\\
Recipe \ref{rec:pb-filled} obligations & \texttt{syzygyCombinationValid} and the target's bookkeeping fields, unchanged\\
Recipe \ref{rec:pb-transport} reading selector & \texttt{syzygyConclusionForm} (shared, not duplicated)\\
Proposition \ref{prop:pb-exact} reading transport & \texttt{holds\_of\_syzygyConclusionForm} (proven in the \rulename{Syzygy} section)\\
Proposition \ref{prop:pb-sound} & \texttt{sound\_pBound} (a \texttt{sound\_syzygy} specialization)\\
wire placement & \texttt{RuleTag.pBound}; no wire tag allocated in v0 --- producers elaborate before emission (the template\textquotesingle s second branch), and the emit path refuses the rule accordingly\\
\end{anchors}


\rulesection{28}{PsdCertTri}{admissible rule}
\label{blcatrule-PsdCertTri}
\ruleblock{The rule (Steele metanotation)}
\sourcenote{as realized in the mechanized development; source row: \texttt{\texttt{rules.toml}} --- \emph{PsdCertTri}. The metanotation statement is typeset by hand for notational quality; its mechanized ground truth is the Lean statement below, extracted from the sources at every build of this book and hash-bound. The rule \emph{roster} is machine-checked (every mechanized rule tag round-trips to exactly one catalogue row); the statement \emph{text} is editorial, checked by eye against the extracted Lean beside it.}
\begin{mathpar}
\inferrule*[Right=\rulename{PsdCertTri}]{
  A \in \Q^{n \times n} \ \text{symmetric}
  \\
  Z \ \text{upper-triangular}, \ Z_{ii} > 0
  \\
  A \cdot Z \ \text{lower-triangular}, \ (A \cdot Z)_{ii} > 0 \ \text{(recomputed exactly)}
  \\
  p_1, \dots, p_n \ \text{payload terms}
  \\
  c \equiv -\textstyle\sum_{i,j} A_{ij}\, p_i\, p_j \ \text{(exact canonical polynomial)}
  \\
  \Psi \in \bigl\{\, c \leq 0, \;\; Q \geq 0 \text{ where } c = -Q, \;\; l \leq r \text{ where } c = l - r \,\bigr\}
}{
  \Phi \vdash \Psi
}
\end{mathpar}

\par\medskip
\rulename{PsdCertTri} is the exact tier of the positive-definiteness
certificate family: a premise-free rule concluding nonnegativity of the
quadratic form $\sum_{i,j} A_{ij}\,p_i\,p_j$ of a symmetric rational
matrix $A$ at a payload term vector, warranted by an untrusted witness
$Z$ passing three decidable syntactic checks --- $A$ symmetric, $Z$
upper-triangular with positive diagonal, and the product $A \cdot Z$
(recomputed exactly, one sparse multiply) lower-triangular with
positive diagonal. Those premises are a theorem-strength certificate of
$A \succ 0$: the congruence $Z^{\mathsf T}(AZ)$ collapses to a positive
diagonal, and positivity transports back through the invertible $Z$.
The matrix-witness check is division-free, $O(\alpha n^2)$
multiplications ($\alpha$
average nonzeros per row), and runs in $\mathbb{Z}$ on integer inputs.

A producer reaches for it wherever a proof obligation is
\textquotedblleft this explicit matrix is positive
(semi)definite\textquotedblright\ --- Lyapunov and storage-function
certificates, energy quadratic forms, elimination-free replacements for
exact rational LDL factorization, whose $O(n^3)$
divisions-with-fill-in it displaces. The witness is produced off-stage
(a float Cholesky, exact elimination, anything); soundness never
consults its accuracy --- a forged $Z$ can only refuse. The concluded
fact enters the logic through the same conclusion-form selector as
\rulename{Syzygy} (raw zero-form, the flipped $Q \geq 0$ reading, or a
literal $l \leq r$ split), so it slots directly into
\rulename{Syzygy}/\rulename{Farkas} combinations as a component
conclusion; strictness is deliberately dropped at the formula boundary
(a premise-free atom carries no $p \neq 0$), so the in-logic fact is
$\succcurlyeq$-shaped even though the matrix-side theorem is strict.

What the checker recomputes: the three matrix checks entrywise over
exact $\Q$, and --- the recompute pattern --- that the payload
conclusion term\textquotesingle s canonical polynomial IS the negated
quadratic form $-\sum_{i,j} A_{ij}\,p_i\,p_j$ built from the payload
terms (in either the wire-keyed or the wire-coordinate polynomial
space). What breaks without each hypothesis: without symmetry the
congruence argument dies and the premises say nothing about the
quadratic form; without the $Z$ diagonal signs the product argument
loses invertibility or positivity (the two sign families are
determinant-linked --- for an indefinite $A$, flipping a column to
repair one violates the
other); the $AZ$ pattern and sign checks are the content itself, and
for an indefinite $A$ no witness whatsoever passes; without the
conclusion identity the posted atom could be unrelated to $A$; the
shape/finiteness/positive-size fields on the conclusion are the
standard evaluability-and-vacuity discipline (a zero-size wire
satisfies atoms vacuously, and the positive-size field refuses that
vacuous surface). The proof section carries the family
preamble and walks these in order.

\cataloguenote{Provenance: the triangular-inverse method is due to David
\textquotedblleft davidad\textquotedblright\ Dalrymple; Geoffrey Irving
derived and published the rounding-tolerant integer analysis and
formalized that variant at scale (a certified 65536-node spectral
expander) --- see the proof section\textquotesingle s citations. The
kernel development here is an independent proof of the exact tier.
Wire tag: deferred to the payload-window allocation; the emit path
refuses the rule until allocated.
}
\ruleblock{Lean statement}
\declhead{BernsteinLogic.Soundness}{sound\_psdCertOf}{\badge{BadgeGreen}{statement bound}}
\begin{leansrc}
/-- The shared PsdCert semantic lemma: from the abstract rational-points
quadratic-form fact (supplied per-variant by the PsdCert bridge
theorems), through the checked conclusion identity, to sequent validity
at every reading of `syzygyConclusionForm`. -/
theorem sound_psdCertOf {H : Hypernet} {seq : Sequent} {a : Nat}
    {M : Fin a → Fin a → ℚ} {ts : List Term} {conclusion : Term}
    (hpsd : ∀ x : Fin a → ℚ, 0 ≤ PsdCert.finQSum a fun i =>
      PsdCert.finQSum a fun j => M i j * x i * x j)
    (hconcl : PsdCert.conclCheck a M ts conclusion = true)
    (hshape : ∃ s, conclusion.bshape H = some s)
    (hfin : conclusion.finWires H)
    (hwpos : ∀ w ∈ conclusion.freeWires,
      ∃ τ, H.wireType w = some τ ∧ 0 < Shape.totalSize τ.shape)
    (hform : syzygyConclusionForm conclusion seq.post) :
    seq.semanticallyValid H
\end{leansrc}
\declhead{BernsteinLogic.Soundness}{sound\_psdCertTri}{\badge{BadgeGreen}{statement bound}}
\begin{leansrc}
/-- PsdCertTri soundness: the /T witness warrants A ≻ 0 over ℝ
(`PsdCert.checkTri_posDef`, davidad's triangular-inverse theorem), whose
quadratic form is nonnegative at every rational point. -/
theorem sound_psdCertTri {H : Hypernet} {seq : Sequent} {a : Nat}
    {A Z : Fin a → Fin a → ℚ} {ts : List Term} {conclusion : Term}
    (hwit : PsdCert.checkTri a A Z = true)
    (hconcl : PsdCert.conclCheck a A ts conclusion = true)
    (hshape : ∃ s, conclusion.bshape H = some s)
    (hfin : conclusion.finWires H)
    (hwpos : ∀ w ∈ conclusion.freeWires,
      ∃ τ, H.wireType w = some τ ∧ 0 < Shape.totalSize τ.shape)
    (hform : syzygyConclusionForm conclusion seq.post) :
    seq.semanticallyValid H
\end{leansrc}
\par\medskip
\textbf{The matrix theorems (mathlib side).}\enspace The chain from
witness check to consumable fact: the congruence collapse and transport
lemmas prove the triangular-witness theorem; \texttt{checkTri\_posDef}
bridges the executable Bool check to positive definiteness of the
$\R$-cast; and the rational-points corollary is the exact shape the
rule soundness lemma consumes. Soundness never consults witness
accuracy --- every hypothesis of every theorem below is a checked
premise, so a forged witness can only fail a decidable check.

\declhead{BernsteinLogic.PsdCert}{congruence\_eq\_diagonal}{\badge{BadgeGreen}{statement bound}}
\begin{leansrc}
/-- The congruence `Zᴴ * (A * Z)` collapses to a diagonal matrix when `Z` is
upper-triangular and `A * Z` is lower-triangular, and its diagonal entries
are the products of the two matrices' diagonals. -/
theorem congruence_eq_diagonal {A Z : Matrix n n ℝ}
    (hA : ∀ i j, A i j = A j i)
    (hZu : ∀ ⦃i j⦄, j < i → Z i j = 0)
    (hMl : ∀ ⦃i j⦄, i < j → (A * Z) i j = 0) :
    Zᴴ * (A * Z) = diagonal fun i => Z i i * (A * Z) i i
\end{leansrc}
\declhead{BernsteinLogic.PsdCert}{posDef\_of\_congruence\_posDef}{\badge{BadgeGreen}{statement bound}}
\begin{leansrc}
/-- Congruence transport (factored from the /T proof for reuse by /DD): if
the congruence `Zᴴ * (A * Z)` is positive definite and `Z` is
upper-triangular with positive diagonal (hence invertible), then `A` is
positive definite. -/
theorem posDef_of_congruence_posDef {A Z : Matrix n n ℝ}
    (hZu : ∀ ⦃i j⦄, j < i → Z i j = 0)
    (hZd : ∀ i, 0 < Z i i)
    (hBd : (Zᴴ * (A * Z)).PosDef) :
    A.PosDef
\end{leansrc}
\declhead{BernsteinLogic.PsdCert}{posDef\_of\_triangular\_witness}{\badge{BadgeGreen}{statement bound}}
\begin{leansrc}
/-- **PsdCertTri soundness** (davidad's theorem). A symmetric matrix with an
upper-triangular positive-diagonal witness `Z` such that `A * Z` is
lower-triangular with positive diagonal is positive definite. -/
theorem posDef_of_triangular_witness {A Z : Matrix n n ℝ}
    (hA : ∀ i j, A i j = A j i)
    (hZu : ∀ ⦃i j⦄, j < i → Z i j = 0)
    (hZd : ∀ i, 0 < Z i i)
    (hMl : ∀ ⦃i j⦄, i < j → (A * Z) i j = 0)
    (hMd : ∀ i, 0 < (A * Z) i i) :
    A.PosDef
\end{leansrc}
\declhead{BernsteinLogic.PsdCert}{checkTri\_posDef}{\badge{BadgeGreen}{statement bound}}
\begin{leansrc}
/-- `checkTri` warrants positive definiteness of the ℝ-cast. -/
theorem checkTri_posDef {a : Nat} {A Z : Fin a → Fin a → ℚ}
    (h : checkTri a A Z = true) : (matR A).PosDef
\end{leansrc}
\declhead{BernsteinLogic.PsdCert}{quadForm\_nonneg\_of\_posSemidef}{\badge{BadgeGreen}{statement bound}}
\begin{leansrc}
/-- Nonnegativity of the ℚ quadratic form at every rational point, from
positive semidefiniteness of the ℝ-cast — the fact the rule soundness
lemmas consume. -/
theorem quadForm_nonneg_of_posSemidef {a : Nat} {M : Fin a → Fin a → ℚ}
    (h : (matR M).PosSemidef) (x : Fin a → ℚ) :
    0 ≤ finQSum a (fun i => finQSum a (fun j => M i j * x i * x j))
\end{leansrc}
\declhead{BernsteinLogic.PsdCert}{checkTri\_quadForm\_nonneg}{\badge{BadgeGreen}{statement bound}}
\begin{leansrc}
/-- /T's consumable corollary. -/
theorem checkTri_quadForm_nonneg {a : Nat} {A Z : Fin a → Fin a → ℚ}
    (h : checkTri a A Z = true) (x : Fin a → ℚ) :
    0 ≤ finQSum a (fun i => finQSum a (fun j => A i j * x i * x j))
\end{leansrc}
\par\medskip
\textbf{The executable check (zero-mathlib closure).}\enspace The
witness checks are List-based exact rational computations on the
executable checker\textquotesingle s import closure --- displayed in
full so the reader can confirm by eye that they are sign, zero-pattern,
and entrywise-equality inspections over recomputed products, with no
divisions and no trusted certificate data beyond the witness itself;
and that \texttt{conclCheck} pins the payload conclusion\textquotesingle s
canonical polynomial to the negated quadratic form (both polynomial
branches mirror the \rulename{Syzygy} check\textquotesingle s split).
These definitions are shared by all three family members; the /DD and
/K sections display only their own additional entries.

\declhead{BernsteinLogic.Core.Catalogue}{qlsum}{\badge{BadgeGreen}{definition bound}}
\begin{leansrc}
/-- List sum over Rat (zero-mathlib; `List.sum`'s instance path avoided
deliberately on the core closure). -/
def qlsum (l : List Rat) : Rat := l.foldr (· + ·) 0
\end{leansrc}
\declhead{BernsteinLogic.Core.Catalogue}{qabs}{\badge{BadgeGreen}{definition bound}}
\begin{leansrc}
/-- `|q|` without lattice instances. -/
def qabs (q : Rat) : Rat := if q < 0 then -q else q
\end{leansrc}
\declhead{BernsteinLogic.Core.Catalogue}{finQSum}{\badge{BadgeGreen}{definition bound}}
\begin{leansrc}
/-- Σ_{i < n} f i, List-based. -/
def finQSum (n : Nat) (f : Fin n → Rat) : Rat := qlsum ((List.finRange n).map f)
\end{leansrc}
\declhead{BernsteinLogic.Core.Catalogue}{mulEntry}{\badge{BadgeGreen}{definition bound}}
\begin{leansrc}
/-- Entry (i,j) of the product A·B, recomputed exactly. -/
def mulEntry (n : Nat) (A B : Fin n → Fin n → Rat) (i j : Fin n) : Rat :=
  finQSum n (fun k => A i k * B k j)
\end{leansrc}
\declhead{BernsteinLogic.Core.Catalogue}{symB}{\badge{BadgeGreen}{definition bound}}
\begin{leansrc}
/-- Symmetry, checked entrywise. -/
def symB (n : Nat) (A : Fin n → Fin n → Rat) : Bool :=
  (List.finRange n).all fun i => (List.finRange n).all fun j =>
    decide (A i j = A j i)
\end{leansrc}
\declhead{BernsteinLogic.Core.Catalogue}{upperTriPosDiagB}{\badge{BadgeGreen}{definition bound}}
\begin{leansrc}
/-- Upper-triangular with positive diagonal, checked entrywise. -/
def upperTriPosDiagB (n : Nat) (Z : Fin n → Fin n → Rat) : Bool :=
  ((List.finRange n).all fun i => (List.finRange n).all fun j =>
    decide (j < i → Z i j = 0)) &&
  ((List.finRange n).all fun i => decide (0 < Z i i))
\end{leansrc}
\declhead{BernsteinLogic.Core.Catalogue}{checkTri}{\badge{BadgeGreen}{definition bound}}
\begin{leansrc}
/-- PsdCertTri witness check (davidad's theorem, all-syntactic premises):
A symmetric; Z upper-triangular with positive diagonal; A·Z
lower-triangular with positive diagonal. Concludes A ≻ 0 (matrix side);
division-free, O(αn²) multiplications. -/
def checkTri (n : Nat) (A Z : Fin n → Fin n → Rat) : Bool :=
  symB n A && upperTriPosDiagB n Z &&
  ((List.finRange n).all fun i => (List.finRange n).all fun j =>
    decide (i < j → mulEntry n A Z i j = 0)) &&
  ((List.finRange n).all fun i => decide (0 < mulEntry n A Z i i))
\end{leansrc}
\declhead{BernsteinLogic.Core.Catalogue}{negQuadFormPoly}{\badge{BadgeGreen}{definition bound}}
\begin{leansrc}
/-- The negated quadratic-form canonical polynomial
−Σᵢⱼ Mᵢⱼ·pᵢ·pⱼ (zero-normal orientation: the rule's primitive fact is
`conclusion ≤ 0` with conclusion ~ −Q, exposed through
`syzygyConclusionForm` — a producer typically emits the ge-flip reading
`Q ≥ 0` with the conclusion term literally `.neg Q`). -/
def negQuadFormPoly (n : Nat) (M : Fin n → Fin n → Rat)
    (ps : List QArith.Poly) : QArith.Poly :=
  QArith.Poly.neg ((List.finRange n).foldr (fun i acc =>
    (List.finRange n).foldr (fun j acc2 =>
      QArith.Poly.add (QArith.Poly.scale (M i j)
        (QArith.Poly.mul (ps.getD i.val []) (ps.getD j.val []))) acc2) acc) [])
\end{leansrc}
\declhead{BernsteinLogic.Core.Catalogue}{conclCheckCore}{\badge{BadgeGreen}{definition bound}}
\begin{leansrc}
/-- Conclusion check, idx-free branch (syzygyCheckCore's discipline):
the payload conclusion's canonical polynomial IS the negated quadratic
form of M at the payload terms' canonical polynomials — the recompute
pattern, no coefficient tensors trusted. -/
def conclCheckCore (n : Nat) (M : Fin n → Fin n → Rat) (ts : List Term)
    (conclusion : Term) : Bool :=
  decide (ts.length = n) &&
  conclusion.idxFree && ts.all Term.idxFree &&
  decide (conclusion.toPoly = negQuadFormPoly n M (ts.map Term.toPoly))
\end{leansrc}
\declhead{BernsteinLogic.Core.Catalogue}{conclCheckW}{\badge{BadgeGreen}{definition bound}}
\begin{leansrc}
/-- Conclusion check, wire-coordinate branch (syzygyCheckW's G1 route:
multi-wire idx atoms — the vector-wire quadratic forms the E-4a /
Lyapunov customers emit). -/
def conclCheckW (n : Nat) (M : Fin n → Fin n → Rat) (ts : List Term)
    (conclusion : Term) : Bool :=
  decide (ts.length = n) &&
  (match conclusion.asCoordPolyIdxW?, ts.mapM Term.asCoordPolyIdxW? with
   | some pc, some pts => decide (pc = negQuadFormPoly n M pts)
   | _, _ => false)
\end{leansrc}
\declhead{BernsteinLogic.Core.Catalogue}{conclCheck}{\badge{BadgeGreen}{definition bound}}
\begin{leansrc}
/-- Either conclusion branch (mirrors `syzygyCheck`'s split). -/
def conclCheck (n : Nat) (M : Fin n → Fin n → Rat) (ts : List Term)
    (conclusion : Term) : Bool :=
  conclCheckCore n M ts conclusion || conclCheckW n M ts conclusion
\end{leansrc}
\ruleblock{Correspondence (Steele display $\leftrightarrow$ Lean)}
\badge{BadgeGreen}{reviewer-certified against the displayed texts}
\sourcenote{The dictionary below is a reviewer-certified statement of how the mechanized statements above interpret the displayed rule: display metavariable $\mapsto$ Lean binder, each display hypothesis $\mapsto$ the premise or side condition carrying it, the display conclusion $\mapsto$ the goal --- with any deliberate difference in width between display and mechanization named outright. Certification is recorded as a hash triple over (the displayed cell, the extracted statement texts, this dictionary\textquotesingle s own text); if any of the three changes, this block turns amber until a reviewer re-certifies. The correspondence claim itself remains editorial --- reviewer-grade, not machine-checked; its machine-checked successor types each soundness statement against the rule as a structured object, and supersedes this block per rule as it lands.}
Judgment frame as in \rulename{Ax}. This dictionary is the
trio\textquotesingle s template; /DD and /K differ in the witness row
alone.
\begin{itemize}
\item $n \mapsto$ \texttt{a : Nat}; $A \mapsto$ \texttt{A}; $Z
  \mapsto$ \texttt{Z} (both \texttt{Fin a → Fin a → ℚ}); $p_1,
  \ldots, p_n \mapsto$ \texttt{ts : List Term}; $c \mapsto$
  \texttt{conclusion : Term}.
\item premise-freeness $\mapsto$ silence, the \rulename{PConstExact}
  representation: no hypothesis mentions \texttt{seq.pre}, so $\Phi$
  is arbitrary --- the \rulename{Norm} convention of the roster note,
  carried structurally.
\item the three displayed matrix rows (symmetry; upper-triangular $Z$
  with positive diagonal; lower-triangular $A \cdot Z$ with positive
  diagonal, recomputed exactly) $\mapsto$ the single Boolean witness
  \texttt{hwit} $:$ \texttt{PsdCert.checkTri a A Z = true}, whose
  definition (displayed above) recomputes all three. The mathematics
  enters through the certified-alongside bridge:
  \texttt{checkTri\_posDef} warrants $A \succ 0$ over $\R$, and
  \texttt{sound\_psdCertTri} consumes only its rational-points
  quadratic-form corollary \texttt{checkTri\_quadForm\_nonneg} --- the
  displayed strictness drop at the formula boundary, visible in the
  statement\textquotesingle s type.
\item the canonical-polynomial row $c \equiv -\sum_{i,j} A_{ij} p_i
  p_j$ $\mapsto$ \texttt{hconcl} $:$ \texttt{conclCheck a A ts
  conclusion = true} --- the identity is checked on canonical forms by
  either of two certified routes (the idx-free \texttt{toPoly} replay
  of \texttt{conclCheckCore}, or the coordinate-indexed
  \texttt{conclCheckW}), and \texttt{sound\_psdCertOf} discharges
  both.
\item $\Psi \in \{c \leq 0;\ Q \geq 0$ where $c = -Q;\ l \leq r$
  where $c = l - r\}$ $\mapsto$ \texttt{hform} $:$
  \texttt{syzygyConclusionForm conclusion seq.post} --- the certified
  selector is literally that three-membered disjunction, reading the
  conclusion term\textquotesingle s head syntactically (shared with
  \rulename{Syzygy} and \rulename{PBound}).
\item with no displayed row: the evaluability pack ---
  \texttt{hshape} (the conclusion has a broadcast shape in
  \texttt{H}), \texttt{hfin} (its wires are finite-typed), and
  \texttt{hwpos} (every free wire positively sized). The display
  omits all three; the certified statement demands them, and
  \texttt{hwpos} is the same vacuous-surface refusal the
  \rulename{Farkas} display carries inline and the
  \rulename{LiRPA} linkage builds in.
\end{itemize}
Width note: elementwise atom semantics batch the conclusion (the
roster note), so the display\textquotesingle s scalar reading is the
singleton case of the certified one; conversely the certified
statement never uses witness accuracy beyond the Boolean check ---
a forged $Z$ refuses, never falsely accepts --- and no hypothesis is
unused.

\ruleblock{Lean verifier receipts}
\sourcenote{module \texttt{BernsteinLogic.Soundness}; bound build \texttt{a20d90e66c13\ldots} (bound from receipt a9c0e34d3e3c9c8c, 2026-08-03T21:30:36Z). A receipt records that named checking layers accepted a build: (module, compiled-artifact SHA-256, environment SHA-256, checker layers), minted when the checking gate ran; the store receipts attesting the bound build are listed with the layers each run credited. Nothing here asks the reader to trust the store: every layer is an independently runnable program over the quoted sources --- rebuild the module and re-run the checkers to reproduce the verdict. Receipts attesting superseded builds credit no layer and are omitted from this rendering (the count of omissions, if any, is noted below).}
\begin{receipttable}
\receiptrow{a9c0e34d3e3c9c8c}{2026-08-03T21:30:36Z}{build+decorrelated}{\badge{BadgeGreen}{current}}
\end{receipttable}
\sourcenote{21 further receipt(s) in the store attest superseded builds of this module and are omitted from this rendering.}
\sourcenote{module \texttt{BernsteinLogic.PsdCert}; bound build \texttt{be59faf6aff0\ldots} (bound from receipt 872a60504989acef, 2026-07-29T07:57:01Z). A receipt records that named checking layers accepted a build: (module, compiled-artifact SHA-256, environment SHA-256, checker layers), minted when the checking gate ran; the store receipts attesting the bound build are listed with the layers each run credited. Nothing here asks the reader to trust the store: every layer is an independently runnable program over the quoted sources --- rebuild the module and re-run the checkers to reproduce the verdict. Receipts attesting superseded builds credit no layer and are omitted from this rendering (the count of omissions, if any, is noted below).}
\begin{receipttable}
\receiptrow{13403ea3a4e42f77}{2026-07-28T23:44:21Z}{build+decorrelated}{\badge{BadgeGreen}{current}}
\receiptrow{6ad2a620d75bd04e}{2026-07-28T23:56:32Z}{build+decorrelated}{\badge{BadgeGreen}{current}}
\receiptrow{d37ce088bc4caf69}{2026-07-29T00:55:46Z}{build+replay+hardened+decorrelated}{\badge{BadgeGreen}{current}}
\receiptrow{872a60504989acef}{2026-07-29T07:57:01Z}{build+replay+hardened+decorrelated}{\badge{BadgeGreen}{current}}
\end{receipttable}
\sourcenote{5 further receipt(s) in the store attest superseded builds of this module and are omitted from this rendering.}
\sourcenote{module \texttt{BernsteinLogic.Core.Catalogue}; bound build \texttt{cbc92e976a1d\ldots} (bound from receipt 5abdb3589707b5df, 2026-08-03T21:30:33Z). A receipt records that named checking layers accepted a build: (module, compiled-artifact SHA-256, environment SHA-256, checker layers), minted when the checking gate ran; the store receipts attesting the bound build are listed with the layers each run credited. Nothing here asks the reader to trust the store: every layer is an independently runnable program over the quoted sources --- rebuild the module and re-run the checkers to reproduce the verdict. Receipts attesting superseded builds credit no layer and are omitted from this rendering (the count of omissions, if any, is noted below).}
\begin{receipttable}
\receiptrow{24804e66edba0bc9}{2026-08-03T20:25:09Z}{build+decorrelated}{\badge{BadgeGreen}{current}}
\receiptrow{5abdb3589707b5df}{2026-08-03T21:30:33Z}{build+decorrelated}{\badge{BadgeGreen}{current}}
\end{receipttable}
\sourcenote{18 further receipt(s) in the store attest superseded builds of this module and are omitted from this rendering.}
\layerline{\layeritem{build}{lake build (elaboration)}{\badge{BadgeGreen}{current}}\layeritem{replay}{leanchecker (official kernel replay)}{\badge{BadgeGreen}{current}}\layeritem{hardened}{safe\_verify (hardened screen)}{\badge{BadgeGreen}{current}}\layeritem{decorrelated}{lean4export + nanoda (decorrelated kernel)}{\badge{BadgeGreen}{current}}}
\sourcenote{\badge{BadgeGray}{comparator: none} no fixture in the three-implementation comparator harness (compiled / kernel-interpreted / reference-interpreted) exercises this rule.}
\ruleblock{HOL4 statement}
\sourcenote{theory \texttt{bern\_tr\_soundpsd}, HOL4 development script \texttt{\texttt{bern\_tr\_soundpsdScript.sml}} \badge{BadgeGreen}{script bound}}
\holhead{Theorem sound\_psdCertOf}
\begin{holsrc}
Theorem sound_psdCertOf:
  ∀H sq n M ts conclusion.
    (∀x. 0 ≤ quadForm n M x) ∧
    psdConclCheck n M ts conclusion ∧
    (∃s. tmBshape H conclusion = SOME s) ∧
    tmFinWires H conclusion ∧
    (∀w. MEM w (termFreeWires conclusion) ⇒
         ∃τ. H.wireType w = SOME τ ∧ 0 < shapeTotalSize τ.shape) ∧
    syzygyConclusionFormB conclusion sq.post ⇒
    seqSemanticallyValid H sq
\end{holsrc}
\holhead{Theorem sound\_psdCertTri}
\begin{holsrc}
Theorem sound_psdCertTri:
  ∀H sq n A Z ts conclusion.
    psdCheckT n A Z ∧
    psdConclCheck n A ts conclusion ∧
    (∃s. tmBshape H conclusion = SOME s) ∧
    tmFinWires H conclusion ∧
    (∀w. MEM w (termFreeWires conclusion) ⇒
         ∃τ. H.wireType w = SOME τ ∧ 0 < shapeTotalSize τ.shape) ∧
    syzygyConclusionFormB conclusion sq.post ⇒
    seqSemanticallyValid H sq
\end{holsrc}
\ruleblock{HOL verifier receipts}
\badge{BadgeBlue}{documentary} \texttt{hol-strict} CLEAN: theory \texttt{bern\_tr\_soundpsd}, 4
theorems, 0 cheat-tagged, 0 nonstandard axioms (whitelist:
\texttt{DISK\_THM}; the four \texttt{bool} axioms), 249-theory
ancestry --- verification run of 2026-07-22 on an isolated build host,
against a content-addressed checkout of the exact sources cited here
(git blob \texttt{c01e0b0c\ldots}). The twins mirror the Lean
statements clause-for-clause per the script\textquotesingle s
statement-fidelity inventory: the shared semantic spine
(\texttt{sound\_psdCertOf}) runs from the abstract rational-points
quadratic-form fact through the checked conclusion identity to sequent
validity at every reading of the shared conclusion-form selector, and
the per-shape corollary discharges that fact via the ported
triangular-witness bridge (\texttt{psdCheckT\_quadForm}, theory
\texttt{bern\_psd}, where the witness check itself ---
\texttt{psdCheckT} --- is ported with kernel-EVAL receipts). The
payload record and its wellformedness arm ride the ported proof-object
stratum (theory \texttt{bern\_tr\_catdata2}, 272 theorems, 0 offences,
same run).

\sourcenote{HOL receipts are documentary: quoted results of \texttt{hol-strict} runs
and census rows from the HOL4 development\textquotesingle s build
evidence, not minted receipt artefacts. A HOL receipt store (the Lean
store\textquotesingle s twin) is named future work.
}
\ruleblock{Proof}

\sourcenote{Genre: admissible rule (tier 2) --- a recipe by
\emph{reduction to verified machinery} (the \rulename{PConstExact}
pattern): no catalogue derivation of an exact matrix-positivity fact
exists, which is the trio's point; each rule discharges against a named
kernel-verified matrix theorem. This section also states the
certificate-family preamble the \rulename{PsdCertDiag} and
\rulename{PsdCertKernel} sections reference.}

\begin{brecipe}[the PSD-certificate family --- shared shape]
\label{rec:psd-preamble}
All three rules certify positivity of a symmetric matrix
$A \in \Q^{n\times n}$ from an \emph{untrusted witness}, by decidable
syntactic checks in exact rational arithmetic --- no divisions, no
pivoting, no floating point on the trusted side. The shared discipline:

\emph{(i) Soundness never consults the witness's accuracy.} The
conclusion is a theorem of the checked premises; a forged or degraded
witness can only cause a false REJECT, never a false ACCEPT. Producers
may therefore run arbitrary float factorizations off-stage.

\emph{(ii) The in-logic conclusion is the quadratic form, not the
matrix.} The formula language has no matrix atoms, so the rule's
conclusion is nonnegativity of $Q(p) = \sum_{i,j} A_{ij}\,p_i\,p_j$ at
a payload term vector $p_1,\ldots,p_n$: the payload conclusion term $c$
is checked --- by exact canonical-polynomial identity, the recompute
pattern --- to \emph{be} $-\!\sum_{i,j} A_{ij}\,p_i\,p_j$, and the
sequent's posted formula is one of three readings of $c \leq 0$ (the
raw zero-form; the flipped $Q \geq 0$; the literal $l \leq r$ where
$c = l - r$) --- the same conclusion-form selector displayed in the
\rulename{Syzygy} section, so the fact slots directly into
\rulename{Syzygy}/\rulename{Farkas} combinations as a component
conclusion.

\emph{(iii) Strictness is deliberately dropped at the formula
boundary.} The matrix theorems certify $A \succ 0$ (for /T and /DD),
but the premise-free atom ships the $\succcurlyeq$-shaped fact: strict
positivity of $Q(p)$ would need $p \neq 0$, which no premise-free atom
carries.

\emph{(iv) The rules are premise-free}; $\Phi$ is arbitrary (a fact
holding in the empty context holds under any premise). Elementwise atom
semantics batch the conclusion: at tensor-shaped payload terms, entry
$e$ of the conclusion is the quadratic form at the terms' entry-$e$
values, so one node states a family of instances.

The rewrite, per rule, is degenerate in the \rulename{PConstExact}
sense: one node, zero catalogue premises, discharged by the rule's
matrix theorem composed with the shared transport lemma
(\texttt{sound\_psdCertOf}) that carries ``PSD over $\R$'' through
``nonnegative at every rational point'' and the checked polynomial
identity to the posted formula.
\end{brecipe}

\begin{brecipe}[the \rulename{PsdCertTri} instance]\label{rec:psdt-instance}
Witness: $Z \in \Q^{n\times n}$. Checked premises (all decidable, all
exact): $A$ symmetric; $Z$ upper-triangular with positive diagonal;
$M := A\,Z$ (recomputed exactly, one sparse multiply) lower-triangular
with positive diagonal. Then
\[
\inferrule*[Right=\rulename{PsdCertTri}]
  {\;\mathrm{checkTri}(A, Z) = \mathsf{true} \\
    c \equiv -\textstyle\sum_{i,j} A_{ij}\,p_i\,p_j\;}
  {\Phi \vdash \Psi}
\qquad\rightsquigarrow\qquad
\text{\upshape discharged by } \texttt{sound\_psdCertTri},
\]
with $\Psi$ any reading of $c \leq 0$ per the preamble. The
matrix-witness check costs $O(\alpha n^2)$ multiplications ($\alpha$ =
average nonzeros per row of $A$) plus $O(n^2)$ sign/zero inspections,
division-free; on integer inputs it runs entirely in $\mathbb{Z}$ with
$\{*, +, <\}$. (The conclusion-identity recompute is additional, and
its cost tracks the payload terms' polynomial sizes, not $n$ alone.)
\end{brecipe}

\begin{bprop}[matrix warrant]\label{prop:psdt-matrix}
If $A$ is symmetric, $Z$ upper-triangular with $Z_{ii} > 0$, and $AZ$
lower-triangular with $(AZ)_{ii} > 0$, then $A \succ 0$ over $\R$.
\end{bprop}
\begin{proof}
Let $B := Z^{\mathsf T}(AZ)$. As a product of two lower-triangular
matrices ($Z^{\mathsf T}$ and $AZ$) it is lower-triangular; since
$Z^{\mathsf T} A Z$ is a congruence of the symmetric $A$ it is also
symmetric; hence $B$ is \emph{diagonal}, and the sum
$B_{ii} = \sum_k Z_{ki} (AZ)_{ki}$ collapses to the single term
$Z_{ii}(AZ)_{ii} > 0$. So $B \succ 0$. Moreover
$\det Z = \prod_i Z_{ii} > 0$ makes $Z$ invertible, and
$A = Z^{-\mathsf T} B\, Z^{-1} \succ 0$ by congruence. (This is the
route the kernel proof takes --- \texttt{congruence\_eq\_diagonal}
then \texttt{posDef\_of\_congruence\_posDef} --- avoiding leading
principal minors; the classical alternative reads $A = M Z^{-1}$ as an
LU factorization with positive-diagonal factors and applies Sylvester's
criterion \cite{hornjohnson2013}.)
\end{proof}

\begin{bremark}[each hypothesis, and what removing it admits]
\emph{Symmetry of $A$}: without it $Z^{\mathsf T}(AZ)$ need not be
symmetric, the diagonal-collapse argument dies, and the quadratic form
reads only the symmetric part of $A$ --- the premises would say nothing
about it. \emph{Upper-triangularity and positive diagonal of $Z$}:
triangularity is the cheapest syntactic invertibility certificate
($\det Z = \prod Z_{ii}$), and the diagonal signs feed
$B_{ii} > 0$; the two sign families are determinant-linked: for an
indefinite $A$ --- a genuinely negative pivot --- ``repairing'' the
$AZ$ diagonal sign by flipping the offending column of $Z$ only moves
the violation into $Z$'s own diagonal check ($A = \mathrm{diag}(1,-1)$,
$Z = \mathrm{diag}(1,z)$: making $-z$ positive forces $z < 0$; when $A$
is in fact positive definite, a flip can repair both checks at once ---
acceptance working as intended, not a leak). \emph{The pattern and sign
checks on $AZ$}: this is the entire content --- the diagonal of $AZ$
carries the pivot signs of $A$, and for an indefinite $A$ \emph{no}
witness whatsoever passes (by Proposition~\ref{prop:psdt-matrix} an
accepting witness would prove $A \succ 0$). \emph{The conclusion
identity}: without the recomputed polynomial check the posted atom
could be an arbitrary formula unrelated to $A$.
\end{bremark}

\begin{bremark}[/T needs an exact witness --- the honest split with /DD]
The natural producer emits $Z \approx$ an inverse transposed Cholesky
factor. A \emph{rounded} such $Z$ makes the strictly-upper entries of
$AZ$ small but nonzero, and /T's pattern check rightly refuses. An
exact witness always exists when $A \succ 0$, and rationally: LDL$^{\mathsf T}$
elimination gives $A = L D L^{\mathsf T}$ with unit lower-triangular
rational $L$ and positive rational diagonal $D$, so
$Z := L^{-\mathsf T}$ is rational upper-triangular with positive
diagonal and $A Z = L D$ is lower-triangular with positive diagonal
--- clear denominators columnwise if integer entries are wanted (each
column is determined up to positive scale) --- but the entries grow
like scaled minors, i.e.\ LDL-style bit growth. (The Cholesky factor
itself is generally irrational; the LDL route is the rational
construction.) At sizes where that growth
bites, the rounding-tolerant companion \rulename{PsdCertDiag} is the
production tier; at moderate sizes /T's $O(\alpha n^2)$ check wins the
clock. Neither tier ever risks soundness in the trade --- only
completeness margin moves.
\end{bremark}

\begin{bremark}[provenance]
The triangular-inverse method is due to David ``davidad'' Dalrymple
\cite{dalrymple2026psd}; Geoffrey Irving derived the rounding-tolerant
integer variant and published the certificate analysis
\cite{irving2026certificate}, then formalized that variant in Lean at
scale (a certified $65536$-node spectral expander) in the course of a
sorting-network formalization \cite{irvingaks}. The kernel development
cited here is an independent in-house proof of the exact tier, grown to
the /DD and /K companions.
\end{bremark}

\begin{anchors}
Recipe \ref{rec:psdt-instance} warrant & \texttt{sound\_psdCertTri} (Soundness); shared transport \texttt{sound\_psdCertOf}\\
matrix theorem (Prop.\ \ref{prop:psdt-matrix}) & \texttt{PsdCert.posDef\_of\_triangular\_witness} via \texttt{congruence\_eq\_diagonal}, \texttt{posDef\_of\_congruence\_posDef}\\
check-to-theorem bridge & \texttt{PsdCert.checkTri\_posDef}; rational-points corollary \texttt{checkTri\_quadForm\_nonneg}, \texttt{quadForm\_nonneg\_of\_posSemidef}\\
the executable check & \texttt{PsdCert.checkTri} (with \texttt{symB}, \texttt{upperTriPosDiagB}, \texttt{mulEntry}) --- zero-mathlib, exact $\Q$\\
conclusion identity & \texttt{PsdCert.conclCheck} against \texttt{negQuadFormPoly} (idx-free and wire-coordinate branches)\\
conclusion readings & the \rulename{Syzygy} section's conclusion-form selector and its transport lemma\\
\end{anchors}


\rulesection{29}{PsdCertDiag}{admissible rule}
\label{blcatrule-PsdCertDiag}
\ruleblock{The rule (Steele metanotation)}
\sourcenote{as realized in the mechanized development; source row: \texttt{\texttt{rules.toml}} --- \emph{PsdCertDiag}. The metanotation statement is typeset by hand for notational quality; its mechanized ground truth is the Lean statement below, extracted from the sources at every build of this book and hash-bound. The rule \emph{roster} is machine-checked (every mechanized rule tag round-trips to exactly one catalogue row); the statement \emph{text} is editorial, checked by eye against the extracted Lean beside it.}
\begin{mathpar}
\inferrule*[Right=\rulename{PsdCertDiag}]{
  A \in \Q^{n \times n} \ \text{symmetric}
  \\
  Z \ \text{upper-triangular}, \ Z_{ii} > 0
  \\
  B := Z^{\mathsf{T}} A Z \ \text{strictly row-diagonally dominant (recomputed exactly)}
  \\
  p_1, \dots, p_n \ \text{payload terms}
  \\
  c \equiv -\textstyle\sum_{i,j} A_{ij}\, p_i\, p_j \ \text{(exact canonical polynomial)}
  \\
  \Psi \in \bigl\{\, c \leq 0, \;\; Q \geq 0 \text{ where } c = -Q, \;\; l \leq r \text{ where } c = l - r \,\bigr\}
}{
  \Phi \vdash \Psi
}
\end{mathpar}

\par\medskip
\rulename{PsdCertDiag} is the rounding-tolerant tier of the
positive-definiteness certificate family (the family discipline ---
untrusted witness, quadratic-form conclusion through
\rulename{Syzygy}\textquotesingle s conclusion-form selector,
strictness dropped at the formula boundary, premise-freedom --- is
stated once in the \rulename{PsdCertTri} section). Here the checked
premises are: $A$ symmetric; $Z$ upper-triangular with positive
diagonal; and the congruence $B := Z^{\mathsf T} A Z$, recomputed
exactly by the checker, strictly row-diagonally dominant (positive
diagonal implied, since the off-row absolute sum is nonnegative). A
symmetric strictly dominant matrix is positive definite --- the
Gershgorin-circle reading; the kernel proof is a self-contained
arithmetic-geometric-mean regrouping of the quadratic form --- and
positivity transports back through the invertible $Z$.

This is the tier a real producer feeds: $Z := \lceil 2^s R^{-1}
\rfloor$ rounded from an off-stage float Cholesky $A \approx
R^{\mathsf T} R$ works directly, with the required scale $s$ growing
only logarithmically in dimension and input scale rather than tracking
any accuracy target --- witness entries encode in $s +
O(\log^{+} \|R^{-1}\|_{\max})$ bits (scale exponent plus the
factorization\textquotesingle s own magnitude), dimension
entering only through the logarithm, where the exact /T witness
inherits minor-scale bit growth.
What a too-coarse witness costs is only completeness margin: the check
fails closed and the producer retries finer, with soundness never in
play. The price relative to /T is the $O(n^3)$ congruence product ---
on bounded-width integers, embarrassingly parallel, still
division-free.

What breaks without each hypothesis: without symmetry of $A$ the
regrouping step (which folds the dominance envelope using
$B_{ij} = B_{ji}$) fails and row dominance alone bounds nothing about
the form; without strictness the matrix warrant $A \succ 0$ would overclaim on
singular matrices
(a path-graph Laplacian is weakly dominant with the constants in its
kernel --- the $\succcurlyeq$-shaped posted atom would survive, but
the strict theorem /K consumes on its inner block would not); without exact recomputation of $B$ a wire-supplied congruence
could be dominant while unrelated to $A$; the $Z$ checks are, as in
/T, the cheapest syntactic invertibility certificate. The conclusion
identity and evaluability fields carry the same discipline as /T. The
proof section walks these in order.

\ruleblock{Lean statement}
\declhead{BernsteinLogic.Soundness}{sound\_psdCertDiag}{\badge{BadgeGreen}{statement bound}}
\begin{leansrc}
/-- PsdCertDiag soundness: the /DD witness warrants A ≻ 0 over ℝ
(`PsdCert.checkDiag_posDef`, the diagonal-dominance congruence). -/
theorem sound_psdCertDiag {H : Hypernet} {seq : Sequent} {a : Nat}
    {A Z : Fin a → Fin a → ℚ} {ts : List Term} {conclusion : Term}
    (hwit : PsdCert.checkDiag a A Z = true)
    (hconcl : PsdCert.conclCheck a A ts conclusion = true)
    (hshape : ∃ s, conclusion.bshape H = some s)
    (hfin : conclusion.finWires H)
    (hwpos : ∀ w ∈ conclusion.freeWires,
      ∃ τ, H.wireType w = some τ ∧ 0 < Shape.totalSize τ.shape)
    (hform : syzygyConclusionForm conclusion seq.post) :
    seq.semanticallyValid H
\end{leansrc}
\par\medskip
\textbf{The dominance core.}\enspace \texttt{posDef\_of\_dominant} is
the self-contained finite-sum argument (no spectral machinery):
off-diagonal mass retreats into its arithmetic-geometric-mean envelope
and symmetry folds the envelope onto the diagonal, leaving
$\sum_i (B_{ii} - \sum_{j\neq i} |B_{ij}|)\,x_i^2 > 0$.
\texttt{posDef\_of\_dd\_witness} composes it with the congruence
transport displayed in the \rulename{PsdCertTri} section;
\texttt{checkDiag\_posDef} bridges the executable check, and the
rational-points corollary feeds the rule lemma.

\declhead{BernsteinLogic.PsdCert}{posDef\_of\_dominant}{\badge{BadgeGreen}{statement bound}}
\begin{leansrc}
/-- **The /DD dominance core**: a symmetric strictly row-diagonally
dominant real matrix is positive definite. Self-contained Finset
algebra — the AM–GM regrouping: off-diagonal mass retreats into
`Σᵢ (Σ_{j≠i} |Bᵢⱼ|)·xᵢ²` by symmetry, leaving `Σᵢ (Bᵢᵢ − Σ_{j≠i}|Bᵢⱼ|)·xᵢ²`,
positive at any `x ≠ 0`. (Positivity of the diagonal is implied:
`Bᵢᵢ > Σ_{j≠i}|Bᵢⱼ| ≥ 0`.) -/
theorem posDef_of_dominant [DecidableEq n] {B : Matrix n n ℝ}
    (hsym : ∀ i j, B i j = B j i)
    (hdom : ∀ i, (∑ j ∈ Finset.univ.erase i, |B i j|) < B i i) :
    B.PosDef
\end{leansrc}
\declhead{BernsteinLogic.PsdCert}{posDef\_of\_dd\_witness}{\badge{BadgeGreen}{statement bound}}
\begin{leansrc}
/-- **PsdCertDiag soundness** (Irving's formalized shape): a symmetric
matrix whose congruence `Zᴴ(AZ)` against an upper-triangular
positive-diagonal witness is strictly row-diagonally dominant is
positive definite. -/
theorem posDef_of_dd_witness [DecidableEq n] {A Z : Matrix n n ℝ}
    (hA : ∀ i j, A i j = A j i)
    (hZu : ∀ ⦃i j⦄, j < i → Z i j = 0)
    (hZd : ∀ i, 0 < Z i i)
    (hdom : ∀ i, (∑ j ∈ Finset.univ.erase i, |(Zᴴ * (A * Z)) i j|)
              < (Zᴴ * (A * Z)) i i) :
    A.PosDef
\end{leansrc}
\declhead{BernsteinLogic.PsdCert}{checkDiag\_posDef}{\badge{BadgeGreen}{statement bound}}
\begin{leansrc}
/-- `checkDiag` warrants positive definiteness of the ℝ-cast. -/
theorem checkDiag_posDef {a : Nat} {A Z : Fin a → Fin a → ℚ}
    (h : checkDiag a A Z = true) : (matR A).PosDef
\end{leansrc}
\declhead{BernsteinLogic.PsdCert}{checkDiag\_quadForm\_nonneg}{\badge{BadgeGreen}{statement bound}}
\begin{leansrc}
/-- /DD's consumable corollary. -/
theorem checkDiag_quadForm_nonneg {a : Nat} {A Z : Fin a → Fin a → ℚ}
    (h : checkDiag a A Z = true) (x : Fin a → ℚ) :
    0 ≤ finQSum a (fun i => finQSum a (fun j => A i j * x i * x j))
\end{leansrc}
\par\medskip
\textbf{The executable check --- what /DD adds.}\enspace Only the
exact congruence entries and the dominance inspection are new; the
shared helpers (\texttt{finQSum}, \texttt{qabs},
\texttt{upperTriPosDiagB}, the conclusion check) are displayed in the
\rulename{PsdCertTri} section.

\declhead{BernsteinLogic.Core.Catalogue}{congEntry}{\badge{BadgeGreen}{definition bound}}
\begin{leansrc}
/-- Entry (i,j) of the congruence B := Zᵀ·(A·Z), recomputed exactly. -/
def congEntry (n : Nat) (A Z : Fin n → Fin n → Rat) (i j : Fin n) : Rat :=
  finQSum n (fun k => Z k i * mulEntry n A Z k j)
\end{leansrc}
\declhead{BernsteinLogic.Core.Catalogue}{checkDiag}{\badge{BadgeGreen}{definition bound}}
\begin{leansrc}
/-- PsdCertDiag witness check (Irving's formalized shape — the tier a
float-LDL ROUNDED witness actually feeds): A symmetric; Z
upper-triangular with positive diagonal; B := ZᵀAZ strictly
row-diagonally dominant (whose positive diagonal is implied: the
off-row sum is ≥ 0). Concludes A ≻ 0. -/
def checkDiag (n : Nat) (A Z : Fin n → Fin n → Rat) : Bool :=
  symB n A && upperTriPosDiagB n Z &&
  ((List.finRange n).all fun i =>
    decide (finQSum n (fun j => if j = i then 0 else qabs (congEntry n A Z i j))
            < congEntry n A Z i i))
\end{leansrc}
\ruleblock{Correspondence (Steele display $\leftrightarrow$ Lean)}
\badge{BadgeGreen}{reviewer-certified against the displayed texts}
\sourcenote{The dictionary below is a reviewer-certified statement of how the mechanized statements above interpret the displayed rule: display metavariable $\mapsto$ Lean binder, each display hypothesis $\mapsto$ the premise or side condition carrying it, the display conclusion $\mapsto$ the goal --- with any deliberate difference in width between display and mechanization named outright. Certification is recorded as a hash triple over (the displayed cell, the extracted statement texts, this dictionary\textquotesingle s own text); if any of the three changes, this block turns amber until a reviewer re-certifies. The correspondence claim itself remains editorial --- reviewer-grade, not machine-checked; its machine-checked successor types each soundness statement against the rule as a structured object, and supersedes this block per rule as it lands.}
Verbatim the \rulename{PsdCertTri} dictionary --- premise-free silence,
the same \texttt{ts}/\texttt{conclusion} binders, the same
\texttt{conclCheck} and \texttt{syzygyConclusionForm} rows, the same
undisplayed evaluability pack (\texttt{hshape}, \texttt{hfin},
\texttt{hwpos}) --- except the witness row: the displayed congruence
$B := Z^{\mathsf{T}} A Z$ strictly row-diagonally dominant $\mapsto$
\texttt{hwit} $:$ \texttt{PsdCert.checkDiag a A Z = true}, entering
through \texttt{checkDiag\_posDef} and consumed as its rational-points
corollary \texttt{checkDiag\_quadForm\_nonneg}. The shared assembly is
one certified lemma, \texttt{sound\_psdCertOf}, visible above ---
the dictionary\textquotesingle s claim that the trio differs only in
the witness is itself a fact of the statements, not editorial.

\ruleblock{Lean verifier receipts}
\sourcenote{module \texttt{BernsteinLogic.Soundness}; bound build \texttt{a20d90e66c13\ldots} (bound from receipt a9c0e34d3e3c9c8c, 2026-08-03T21:30:36Z). A receipt records that named checking layers accepted a build: (module, compiled-artifact SHA-256, environment SHA-256, checker layers), minted when the checking gate ran; the store receipts attesting the bound build are listed with the layers each run credited. Nothing here asks the reader to trust the store: every layer is an independently runnable program over the quoted sources --- rebuild the module and re-run the checkers to reproduce the verdict. Receipts attesting superseded builds credit no layer and are omitted from this rendering (the count of omissions, if any, is noted below).}
\begin{receipttable}
\receiptrow{a9c0e34d3e3c9c8c}{2026-08-03T21:30:36Z}{build+decorrelated}{\badge{BadgeGreen}{current}}
\end{receipttable}
\sourcenote{21 further receipt(s) in the store attest superseded builds of this module and are omitted from this rendering.}
\sourcenote{module \texttt{BernsteinLogic.PsdCert}; bound build \texttt{be59faf6aff0\ldots} (bound from receipt 872a60504989acef, 2026-07-29T07:57:01Z). A receipt records that named checking layers accepted a build: (module, compiled-artifact SHA-256, environment SHA-256, checker layers), minted when the checking gate ran; the store receipts attesting the bound build are listed with the layers each run credited. Nothing here asks the reader to trust the store: every layer is an independently runnable program over the quoted sources --- rebuild the module and re-run the checkers to reproduce the verdict. Receipts attesting superseded builds credit no layer and are omitted from this rendering (the count of omissions, if any, is noted below).}
\begin{receipttable}
\receiptrow{13403ea3a4e42f77}{2026-07-28T23:44:21Z}{build+decorrelated}{\badge{BadgeGreen}{current}}
\receiptrow{6ad2a620d75bd04e}{2026-07-28T23:56:32Z}{build+decorrelated}{\badge{BadgeGreen}{current}}
\receiptrow{d37ce088bc4caf69}{2026-07-29T00:55:46Z}{build+replay+hardened+decorrelated}{\badge{BadgeGreen}{current}}
\receiptrow{872a60504989acef}{2026-07-29T07:57:01Z}{build+replay+hardened+decorrelated}{\badge{BadgeGreen}{current}}
\end{receipttable}
\sourcenote{5 further receipt(s) in the store attest superseded builds of this module and are omitted from this rendering.}
\sourcenote{module \texttt{BernsteinLogic.Core.Catalogue}; bound build \texttt{cbc92e976a1d\ldots} (bound from receipt 5abdb3589707b5df, 2026-08-03T21:30:33Z). A receipt records that named checking layers accepted a build: (module, compiled-artifact SHA-256, environment SHA-256, checker layers), minted when the checking gate ran; the store receipts attesting the bound build are listed with the layers each run credited. Nothing here asks the reader to trust the store: every layer is an independently runnable program over the quoted sources --- rebuild the module and re-run the checkers to reproduce the verdict. Receipts attesting superseded builds credit no layer and are omitted from this rendering (the count of omissions, if any, is noted below).}
\begin{receipttable}
\receiptrow{24804e66edba0bc9}{2026-08-03T20:25:09Z}{build+decorrelated}{\badge{BadgeGreen}{current}}
\receiptrow{5abdb3589707b5df}{2026-08-03T21:30:33Z}{build+decorrelated}{\badge{BadgeGreen}{current}}
\end{receipttable}
\sourcenote{18 further receipt(s) in the store attest superseded builds of this module and are omitted from this rendering.}
\layerline{\layeritem{build}{lake build (elaboration)}{\badge{BadgeGreen}{current}}\layeritem{replay}{leanchecker (official kernel replay)}{\badge{BadgeGreen}{current}}\layeritem{hardened}{safe\_verify (hardened screen)}{\badge{BadgeGreen}{current}}\layeritem{decorrelated}{lean4export + nanoda (decorrelated kernel)}{\badge{BadgeGreen}{current}}}
\sourcenote{\badge{BadgeGray}{comparator: none} no fixture in the three-implementation comparator harness (compiled / kernel-interpreted / reference-interpreted) exercises this rule.}
\ruleblock{HOL4 statement}
\sourcenote{theory \texttt{bern\_tr\_soundpsd}, HOL4 development script \texttt{\texttt{bern\_tr\_soundpsdScript.sml}} \badge{BadgeGreen}{script bound}}
\holhead{Theorem sound\_psdCertDiag}
\begin{holsrc}
Theorem sound_psdCertDiag:
  ∀H sq n A Z ts conclusion.
    psdCheckDD n A Z ∧
    psdConclCheck n A ts conclusion ∧
    (∃s. tmBshape H conclusion = SOME s) ∧
    tmFinWires H conclusion ∧
    (∀w. MEM w (termFreeWires conclusion) ⇒
         ∃τ. H.wireType w = SOME τ ∧ 0 < shapeTotalSize τ.shape) ∧
    syzygyConclusionFormB conclusion sq.post ⇒
    seqSemanticallyValid H sq
\end{holsrc}
\ruleblock{HOL verifier receipts}
\badge{BadgeBlue}{documentary} \texttt{hol-strict} CLEAN: theory \texttt{bern\_tr\_soundpsd}, 4
theorems, 0 cheat-tagged, 0 nonstandard axioms (whitelist:
\texttt{DISK\_THM}; the four \texttt{bool} axioms), 249-theory
ancestry --- verification run of 2026-07-22 on an isolated build host,
against a content-addressed checkout of the exact sources cited here
(git blob \texttt{c01e0b0c\ldots}). The shared semantic spine
(\texttt{sound\_psdCertOf}) is displayed in the /Tri section per the
family-preamble convention; this corollary discharges the
quadratic-form fact via the ported congruence-plus-row-dominance
bridge (\texttt{psdCheckDD\_quadForm}, theory \texttt{bern\_psd},
where the witness check itself --- \texttt{psdCheckDD} --- is ported
with kernel-EVAL receipts). The payload record and its wellformedness
arm ride the ported proof-object stratum (theory
\texttt{bern\_tr\_catdata2}, 272 theorems, 0 offences, same run).

\sourcenote{HOL receipts are documentary: quoted results of \texttt{hol-strict} runs
and census rows from the HOL4 development\textquotesingle s build
evidence, not minted receipt artefacts. A HOL receipt store (the Lean
store\textquotesingle s twin) is named future work.
}
\ruleblock{Proof}

\sourcenote{Genre: admissible rule (tier 2) --- reduction to verified
machinery. The certificate-family preamble --- untrusted witness,
quadratic-form conclusion through the \rulename{Syzygy}
conclusion-form selector, strictness dropped at the formula boundary,
premise-freedom, elementwise batching --- is stated in the
\rulename{PsdCertTri} section (its first recipe) and applies verbatim;
only the witness check differs here.}

\begin{brecipe}[the \rulename{PsdCertDiag} instance]\label{rec:psddd-instance}
Witness: $Z \in \Q^{n\times n}$. Checked premises: $A$ symmetric; $Z$
upper-triangular with positive diagonal; and the congruence
$B := Z^{\mathsf T} A Z$ --- \emph{recomputed exactly by the checker} ---
strictly row-diagonally dominant:
$\sum_{j \neq i} |B_{ij}| < B_{ii}$ for every $i$. (Positivity of the
diagonal is implied: the off-row sum is $\geq 0$.) Then
\[
\inferrule*[Right=\rulename{PsdCertDiag}]
  {\;\mathrm{checkDiag}(A, Z) = \mathsf{true} \\
    c \equiv -\textstyle\sum_{i,j} A_{ij}\,p_i\,p_j\;}
  {\Phi \vdash \Psi}
\qquad\rightsquigarrow\qquad
\text{\upshape discharged by } \texttt{sound\_psdCertDiag}.
\]
The matrix-witness check costs $O(\alpha n^2 + n^3)$ multiplications
($\alpha$ the average nonzeros per row of $A$, as at /T)
with the product $AZ$ materialized once (the $Z^{\mathsf T}(AZ)$
congruence is the $n^3$ term; the reference check displayed below
inlines the product entrywise for readability --- it is the
specification, not the cost model) on integers of \emph{bounded}
width --- input bits plus twice the witness bits plus $O(\log n)$ for
the nested sums --- embarrassingly parallel, still division-free.
\end{brecipe}

\begin{bprop}[matrix warrant]\label{prop:psddd-matrix}
A symmetric real matrix $B$ with
$\sum_{j\neq i}|B_{ij}| < B_{ii}$ for all $i$ is positive definite;
hence if additionally $Z$ is upper-triangular with positive diagonal
and $B = Z^{\mathsf T} A Z$ for symmetric $A$, then $A \succ 0$.
\end{bprop}
\begin{proof}
For $x \neq 0$, split the quadratic form and bound each off-diagonal
term by its arithmetic--geometric-mean envelope,
$x_i B_{ij} x_j \geq -|B_{ij}|\,(x_i^2 + x_j^2)/2$; folding the
envelope's two symmetric halves by $B_{ij} = B_{ji}$ gives
\[
x^{\mathsf T} B\, x \;\geq\; \sum_i \Bigl(B_{ii} -
\textstyle\sum_{j \neq i} |B_{ij}|\Bigr)\, x_i^2 \;>\; 0 .
\]
(The kernel proof is exactly this self-contained finite-sum argument
--- \texttt{posDef\_of\_dominant} --- with no spectral machinery. The
classical reading: strict dominance with positive diagonal places
every Gershgorin disc in the open right half-plane, and symmetry makes
the spectrum real, hence positive
\cite{gershgorin1931,taussky1949,hornjohnson2013}.)
Transport through the invertible $Z$ as in the /T proof
(\texttt{posDef\_of\_congruence\_posDef}) yields $A \succ 0$.
\end{proof}

\begin{bremark}[what each hypothesis is doing]
\emph{Symmetry of $A$} makes $B$ symmetric by construction, which the
envelope-folding step consumes ($B_{ij} = B_{ji}$) --- row dominance of
a nonsymmetric matrix bounds neither the symmetric part's rows nor the
form. \emph{Triangularity + positive diagonal of $Z$} is, as in /T,
only the cheapest syntactic invertibility certificate --- invertibility
is all the transport step needs, and triangularity is what the natural
producer emits anyway. \emph{Strict dominance} is the content of the strict tier:
weakening it to non-strict admits singular $B$ (e.g.\ the graph
Laplacian of a path: weakly dominant, kernel = constants), and the
matrix warrant $A \succ 0$ would overclaim --- the posted
$\succcurlyeq$-shaped atom would in fact survive, but the strict
theorem is what \rulename{PsdCertKernel} consumes on its inner block to
confine kernels, and that consumer admits no slack. \emph{Exact recomputation of $B$}: the
congruence is never taken from the certificate --- a wire-supplied $B$
could be dominant while unrelated to $A$.
\end{bremark}

\begin{bremark}[rounding tolerance --- why this is the production tier]
Feed $Z := \lceil 2^s R^{-1} \rfloor$ from an off-stage float Cholesky
$A \approx R^{\mathsf T} R$: writing $Z = 2^s R^{-1} + E$ with rounding
$|E_{ij}| \leq \tfrac12$, exact expansion (against an exact
factorization) gives $B = Z^{\mathsf T} A Z = 4^s I + 2^s\bigl(RE +
(RE)^{\mathsf T}\bigr) + E^{\mathsf T} A E$ --- entrywise contamination
$O(2^s n \|R\|_{\max} + n^2 \|A\|_{\max})$, off-row dominance sums one
factor of $n$ larger --- plus the image of the float residual
$\Delta := A - R^{\mathsf T} R$ --- the term $Z^{\mathsf T} \Delta Z$,
of normalized leading size $R^{-\mathsf T} \Delta R^{-1}$ --- whose
control is a mantissa requirement of
order $\log_2 \kappa(A) + O(\log n)$ bits on the factorization (the
dimension guard bits of the standard backward-error bound), not on the
witness. Strict dominance therefore holds once $2^s$ clears a
polynomial in $n$ times the input scales: a scale exponent $s =
O(\log n +
\log \|R\|_{\max})$ suffices (sharp constants in the cited
certificate analysis), so the witness \emph{scale} grows only
\emph{logarithmically} in dimension and scale and never tracks an
accuracy target; the witness \emph{entries} encode in $s +
O(\log^{+} \|R^{-1}\|_{\max})$ bits --- the second term is the
factorization's own magnitude ($Z$ rounds $2^s R^{-1}$), present at
any scale, and no accuracy target either. Contrast /T,
whose exact witness inherits minor-scale bit growth. What a too-coarse
$Z$ costs is only completeness margin: the check fails closed, and the
producer retries at a finer scale with soundness never in play. This
rounded-integer-witness shape is the one Irving analysed and formalized
at scale \cite{irving2026certificate,irvingaks}; the underlying
tolerance phenomenon is the factor-width/scaled-diagonal-dominance
circle of ideas \cite{boman2005} --- context, not a characterization:
that theorem equates factor-width two with scaled diagonal dominance
of $A$ \emph{itself} (diagonal scaling only), while /DD checks
dominance of the congruence target $B = Z^{\mathsf T} A Z$ for
arbitrary triangular $Z$, and so accepts matrices far outside the SDD
class --- with an exact LDL witness the congruence target is diagonal,
so /DD's exact-witness reach is \emph{every} rational SPD matrix,
while e.g.\ the $3 \times 3$ equicorrelation matrix with off-diagonal
$\rho \in (\tfrac12, 1)$ is SPD but not SDD under any positive
diagonal scaling.
\end{bremark}

\begin{bremark}[the /T---/DD complementarity, stated once]
Same witness idea, dual regimes: /T is exact-witness,
$O(\alpha n^2)$, and wins at moderate sizes and exact-by-construction
witnesses; /DD tolerates any rounding at $O(n^3)$ bounded-width
parallel arithmetic and wins at scale. A producer holding a float
factorization reaches for /DD by default and /T when it can clear
denominators cheaply. Both refuse indefinite $A$ against every witness.
\end{bremark}

\begin{anchors}
Recipe \ref{rec:psddd-instance} warrant & \texttt{sound\_psdCertDiag} (Soundness); shared transport \texttt{sound\_psdCertOf} (see \rulename{PsdCertTri})\\
dominance core (Prop.\ \ref{prop:psddd-matrix}) & \texttt{PsdCert.posDef\_of\_dominant}; assembled in \texttt{posDef\_of\_dd\_witness}\\
check-to-theorem bridge & \texttt{PsdCert.checkDiag\_posDef}; rational-points corollary \texttt{checkDiag\_quadForm\_nonneg}\\
the executable check & \texttt{PsdCert.checkDiag} with the exact congruence entries \texttt{congEntry}\\
conclusion identity + readings & \texttt{PsdCert.conclCheck}; the \rulename{Syzygy} section's conclusion-form selector\\
\end{anchors}


\rulesection{30}{PsdCertKernel}{admissible rule}
\label{blcatrule-PsdCertKernel}
\ruleblock{The rule (Steele metanotation)}
\sourcenote{as realized in the mechanized development; source row: \texttt{\texttt{rules.toml}} --- \emph{PsdCertKernel}. The metanotation statement is typeset by hand for notational quality; its mechanized ground truth is the Lean statement below, extracted from the sources at every build of this book and hash-bound. The rule \emph{roster} is machine-checked (every mechanized rule tag round-trips to exactly one catalogue row); the statement \emph{text} is editorial, checked by eye against the extracted Lean beside it.}
\begin{mathpar}
\inferrule*[Right=\rulename{PsdCertKernel}]{
  A \in \Q^{(m+k) \times (m+k)}, \ \text{kernel coordinates last}
  \\
  L := (I \mid N), \ N \in \Q^{m \times k}
  \\
  A = L^{\mathsf{T}} B L \ \text{entrywise exactly}
  \\
  B \in \Q^{m \times m} \ \text{with a /T-or-/DD inner witness } Z
  \\
  p_1, \dots, p_{m+k} \ \text{payload terms}
  \\
  c \equiv -\textstyle\sum_{i,j} A_{ij}\, p_i\, p_j \ \text{(exact canonical polynomial)}
  \\
  \Psi \in \bigl\{\, c \leq 0, \;\; Q \geq 0 \text{ where } c = -Q, \;\; l \leq r \text{ where } c = l - r \,\bigr\}
}{
  \Phi \vdash \Psi
}
\end{mathpar}

\par\medskip
\rulename{PsdCertKernel} is the semidefinite member of the certificate
family (family discipline in the \rulename{PsdCertTri} section): it
certifies $A \succcurlyeq 0$ \emph{with an exactly-checked kernel},
for the boundary cases the two strict tiers cannot touch --- energy
forms with conservation laws, graph Laplacians with their all-ones
kernel, sharp certificates whose slack vanishes identically. Data:
$A$ over $\Q^{(m+k)\times(m+k)}$ with the $k$ kernel coordinates last
(the producer applies its permutation before emission --- conjugation
is index renaming, so the identity-permutation form loses no
generality); a rational kernel description $N$; and a leading block
$B$ with a /T-or-/DD inner witness. The checked premise is a single
entrywise-exact factorization: $A = L^{\mathsf T} B\, L$ for the
kernel-completion rows $L = (I_m \mid N)$ --- equivalent, given the
inner check\textquotesingle s own symmetry verification of $B$, to
($A$ symmetric $\wedge$ $A \cdot K = 0$ exactly for
$K = (-N;\, I_k)$ $\wedge$ leading block $= B$), the cross and
trailing blocks being forced either way, so a forged $N$ or $B$
refuses --- plus the strict inner check on $B$. Rectangular congruence
then gives $A \succcurlyeq 0$ with the form vanishing exactly on the
span of $K$: the certificate fixes kernel and rank, not just
semidefiniteness.

The kernel data must be exact by necessity, not by preference: the
semidefinite-with-kernel matrices form a boundary set with empty
interior, so no rounding-tolerant witness for kernel directions can
exist --- any check tolerant to kernel perturbation also accepts
slightly-indefinite matrices. Exact kernel plus rounding-tolerant bulk
(the /DD inner witness) is therefore the forced split --- tolerance
can live nowhere else --- and it
matches the producers: kernels are semantic (theorems supply $N$),
floats only ever touch the bulk witness $Z$. Corners are structural:
$k = 0$ degenerates to the inner strict tier; $m = 0$ forces $A = 0$,
where the factorization premise is the entire content.

What breaks without each hypothesis: an approximate factorization
certifies nothing (the boundary argument above); a merely semidefinite
inner block would hide undeclared kernel directions inside $B$; the
kernel-coordinates-last convention is pure bookkeeping. What reaches
the formula layer is the quadratic-form conclusion --- honest
for this rule, since it is $\succcurlyeq$-shaped --- while the checked
kernel/rank content becomes formula-expressible only with closed
tensor terms (the extension the \rulename{PConstExact} section
records); it is certificate-level assurance meanwhile, deliberately
not hand-faked into the formula grammar. The proof section walks these
in order.

\ruleblock{Lean statement}
\declhead{BernsteinLogic.Soundness}{sound\_psdCertKernel}{\badge{BadgeGreen}{statement bound}}
\begin{leansrc}
/-- PsdCertKernel soundness: the /K witness warrants A ⪰ 0 over ℝ
(`PsdCert.checkKernel_posSemidef`, the kernel-completion congruence). -/
theorem sound_psdCertKernel {H : Hypernet} {seq : Sequent} {m k : Nat}
    {A : Fin (m + k) → Fin (m + k) → ℚ} {N : Fin m → Fin k → ℚ}
    {B : Fin m → Fin m → ℚ} {innerT : Bool} {Z : Fin m → Fin m → ℚ}
    {ts : List Term} {conclusion : Term}
    (hwit : PsdCert.checkKernel m k A N B innerT Z = true)
    (hconcl : PsdCert.conclCheck (m + k) A ts conclusion = true)
    (hshape : ∃ s, conclusion.bshape H = some s)
    (hfin : conclusion.finWires H)
    (hwpos : ∀ w ∈ conclusion.freeWires,
      ∃ τ, H.wireType w = some τ ∧ 0 < Shape.totalSize τ.shape)
    (hform : syzygyConclusionForm conclusion seq.post) :
    seq.semanticallyValid H
\end{leansrc}
\par\medskip
\textbf{The rectangular congruence.}\enspace
\texttt{posSemidef\_of\_factor} is one line --- semidefiniteness
transports through ANY rectangular congruence, no invertibility
anywhere --- which is exactly why the factorization premise can carry
the whole rule. \texttt{checkKernel\_posSemidef} bridges the executable
check (selecting the /T or /DD theorem for the inner block), and the
rational-points corollary feeds the rule lemma.

\declhead{BernsteinLogic.PsdCert}{posSemidef\_of\_factor}{\badge{BadgeGreen}{statement bound}}
\begin{leansrc}
/-- **PsdCertKernel soundness core**: a matrix exhibited as the two-sided
congruence `LᴴBL` of a positive-semidefinite block is positive
semidefinite (rectangular congruence; no invertibility anywhere). -/
theorem posSemidef_of_factor {m' n' : Type*} [Fintype m'] [Fintype n']
    {A : Matrix n' n' ℝ} {B : Matrix m' m' ℝ} {L : Matrix m' n' ℝ}
    (hB : B.PosSemidef) (hA : A = Lᴴ * B * L) : A.PosSemidef
\end{leansrc}
\declhead{BernsteinLogic.PsdCert}{checkKernel\_posSemidef}{\badge{BadgeGreen}{statement bound}}
\begin{leansrc}
/-- `checkKernel` warrants positive semidefiniteness of the ℝ-cast. -/
theorem checkKernel_posSemidef {m k : Nat} {A : Fin (m + k) → Fin (m + k) → ℚ}
    {N : Fin m → Fin k → ℚ} {B : Fin m → Fin m → ℚ} {innerT : Bool}
    {Z : Fin m → Fin m → ℚ}
    (h : checkKernel m k A N B innerT Z = true) : (matR A).PosSemidef
\end{leansrc}
\declhead{BernsteinLogic.PsdCert}{checkKernel\_quadForm\_nonneg}{\badge{BadgeGreen}{statement bound}}
\begin{leansrc}
/-- /K's consumable corollary. -/
theorem checkKernel_quadForm_nonneg {m k : Nat} {A : Fin (m + k) → Fin (m + k) → ℚ}
    {N : Fin m → Fin k → ℚ} {B : Fin m → Fin m → ℚ} {innerT : Bool}
    {Z : Fin m → Fin m → ℚ}
    (h : checkKernel m k A N B innerT Z = true) (x : Fin (m + k) → ℚ) :
    0 ≤ finQSum (m + k) (fun i => finQSum (m + k) (fun j => A i j * x i * x j))
\end{leansrc}
\par\medskip
\textbf{The executable check --- what /K adds.}\enspace The
kernel-completion rows and the entrywise factorization check; the
inner witness reuses \texttt{checkTri}/\texttt{checkDiag} and the
conclusion check is shared (both displayed in the
\rulename{PsdCertTri} section).

\declhead{BernsteinLogic.Core.Catalogue}{kernelExt}{\badge{BadgeGreen}{definition bound}}
\begin{leansrc}
/-- The kernel-completion row matrix L = (I | N) : m × (m+k). -/
def kernelExt (m k : Nat) (N : Fin m → Fin k → Rat) (a : Fin m)
    (j : Fin (m + k)) : Rat :=
  if h : j.val < m then (if a.val = j.val then 1 else 0)
  else N a ⟨j.val - m, by have := j.isLt; omega⟩
\end{leansrc}
\declhead{BernsteinLogic.Core.Catalogue}{checkKernel}{\badge{BadgeGreen}{definition bound}}
\begin{leansrc}
/-- PsdCertKernel witness check (PSD-with-kernel companion): A — over
`Fin (m+k)`, the kernel coordinates LAST — IS the two-sided
kernel-completion congruence LᵀBL of its leading block B, entrywise
exactly, plus a T-or-DD inner witness for B. Equivalent to the note's §1
premises (A symmetric ∧ A·K = 0 exactly for K = (−N; Iₖ) ∧ leading
block = B): C = BN and D = NᵀBN are forced either way, and a forged N
or B refuses. The producer applies its permutation P before emission —
P-conjugation is index renaming on matrix and term vector alike, so the
catalogue rule is the P = id core. Corners are structural: k = 0 gives
L = I (A = B); m = 0 forces A = 0 (empty B, no inner witness content). -/
def checkKernel (m k : Nat) (A : Fin (m + k) → Fin (m + k) → Rat)
    (N : Fin m → Fin k → Rat) (B : Fin m → Fin m → Rat)
    (innerT : Bool) (Z : Fin m → Fin m → Rat) : Bool :=
  ((List.finRange (m + k)).all fun i => (List.finRange (m + k)).all fun j =>
    decide (A i j = qlsum ((List.finRange m).map fun a =>
      qlsum ((List.finRange m).map fun b =>
        kernelExt m k N a i * B a b * kernelExt m k N b j)))) &&
  (if innerT then checkTri m B Z else checkDiag m B Z)
\end{leansrc}
\ruleblock{Correspondence (Steele display $\leftrightarrow$ Lean)}
\badge{BadgeGreen}{reviewer-certified against the displayed texts}
\sourcenote{The dictionary below is a reviewer-certified statement of how the mechanized statements above interpret the displayed rule: display metavariable $\mapsto$ Lean binder, each display hypothesis $\mapsto$ the premise or side condition carrying it, the display conclusion $\mapsto$ the goal --- with any deliberate difference in width between display and mechanization named outright. Certification is recorded as a hash triple over (the displayed cell, the extracted statement texts, this dictionary\textquotesingle s own text); if any of the three changes, this block turns amber until a reviewer re-certifies. The correspondence claim itself remains editorial --- reviewer-grade, not machine-checked; its machine-checked successor types each soundness statement against the rule as a structured object, and supersedes this block per rule as it lands.}
As at \rulename{PsdCertTri} --- premise-free silence, the shared
\texttt{conclCheck} row at size $m + k$, the same
\texttt{syzygyConclusionForm} selector and undisplayed evaluability
pack --- with the witness row exchanged for the kernel-completion
data:
\begin{itemize}
\item the block sizes $\mapsto$ \texttt{m k : Nat} with $A$ $:$
  \texttt{Fin (m + k) → Fin (m + k) → ℚ} (kernel coordinates last,
  as displayed --- the producer\textquotesingle s permutation is
  applied before emission, so the certified rule is the $P =
  \mathrm{id}$ core of the roster note); $N \mapsto$ \texttt{N : Fin
  m → Fin k → ℚ}; $B \mapsto$ \texttt{B}; the ``/T-or-/DD inner
  witness\textquotesingle\textquotesingle{} $\mapsto$ the pair
  \texttt{innerT : Bool} (which inner check) and \texttt{Z : Fin m →
  Fin m → ℚ} (its witness).
\item the displayed factorization row $A = L^{\mathsf{T}} B L$
  entrywise exactly, $L = (I \mid N)$, together with the inner
  witness $\mapsto$ the single Boolean \texttt{hwit} $:$
  \texttt{PsdCert.checkKernel m k A N B innerT Z = true}, entering
  through \texttt{checkKernel\_posSemidef} ($A \succcurlyeq 0$ --- the
  trio\textquotesingle s one genuinely $\succcurlyeq$-shaped member,
  so here the formula boundary loses nothing) and consumed as
  \texttt{checkKernel\_quadForm\_nonneg}.
\end{itemize}
Width note, beyond /T\textquotesingle s: the checked kernel and rank
content (the roster note\textquotesingle s \S1 equivalents) is
consumed by the matrix-side bridge but never surfaces in the sequent
--- the certified formula-level claim is exactly the quadratic-form
conclusion, neither more nor less, which is the honest reading the
commentary states and this triple pins.

\ruleblock{Lean verifier receipts}
\sourcenote{module \texttt{BernsteinLogic.Soundness}; bound build \texttt{a20d90e66c13\ldots} (bound from receipt a9c0e34d3e3c9c8c, 2026-08-03T21:30:36Z). A receipt records that named checking layers accepted a build: (module, compiled-artifact SHA-256, environment SHA-256, checker layers), minted when the checking gate ran; the store receipts attesting the bound build are listed with the layers each run credited. Nothing here asks the reader to trust the store: every layer is an independently runnable program over the quoted sources --- rebuild the module and re-run the checkers to reproduce the verdict. Receipts attesting superseded builds credit no layer and are omitted from this rendering (the count of omissions, if any, is noted below).}
\begin{receipttable}
\receiptrow{a9c0e34d3e3c9c8c}{2026-08-03T21:30:36Z}{build+decorrelated}{\badge{BadgeGreen}{current}}
\end{receipttable}
\sourcenote{21 further receipt(s) in the store attest superseded builds of this module and are omitted from this rendering.}
\sourcenote{module \texttt{BernsteinLogic.PsdCert}; bound build \texttt{be59faf6aff0\ldots} (bound from receipt 872a60504989acef, 2026-07-29T07:57:01Z). A receipt records that named checking layers accepted a build: (module, compiled-artifact SHA-256, environment SHA-256, checker layers), minted when the checking gate ran; the store receipts attesting the bound build are listed with the layers each run credited. Nothing here asks the reader to trust the store: every layer is an independently runnable program over the quoted sources --- rebuild the module and re-run the checkers to reproduce the verdict. Receipts attesting superseded builds credit no layer and are omitted from this rendering (the count of omissions, if any, is noted below).}
\begin{receipttable}
\receiptrow{13403ea3a4e42f77}{2026-07-28T23:44:21Z}{build+decorrelated}{\badge{BadgeGreen}{current}}
\receiptrow{6ad2a620d75bd04e}{2026-07-28T23:56:32Z}{build+decorrelated}{\badge{BadgeGreen}{current}}
\receiptrow{d37ce088bc4caf69}{2026-07-29T00:55:46Z}{build+replay+hardened+decorrelated}{\badge{BadgeGreen}{current}}
\receiptrow{872a60504989acef}{2026-07-29T07:57:01Z}{build+replay+hardened+decorrelated}{\badge{BadgeGreen}{current}}
\end{receipttable}
\sourcenote{5 further receipt(s) in the store attest superseded builds of this module and are omitted from this rendering.}
\sourcenote{module \texttt{BernsteinLogic.Core.Catalogue}; bound build \texttt{cbc92e976a1d\ldots} (bound from receipt 5abdb3589707b5df, 2026-08-03T21:30:33Z). A receipt records that named checking layers accepted a build: (module, compiled-artifact SHA-256, environment SHA-256, checker layers), minted when the checking gate ran; the store receipts attesting the bound build are listed with the layers each run credited. Nothing here asks the reader to trust the store: every layer is an independently runnable program over the quoted sources --- rebuild the module and re-run the checkers to reproduce the verdict. Receipts attesting superseded builds credit no layer and are omitted from this rendering (the count of omissions, if any, is noted below).}
\begin{receipttable}
\receiptrow{24804e66edba0bc9}{2026-08-03T20:25:09Z}{build+decorrelated}{\badge{BadgeGreen}{current}}
\receiptrow{5abdb3589707b5df}{2026-08-03T21:30:33Z}{build+decorrelated}{\badge{BadgeGreen}{current}}
\end{receipttable}
\sourcenote{18 further receipt(s) in the store attest superseded builds of this module and are omitted from this rendering.}
\layerline{\layeritem{build}{lake build (elaboration)}{\badge{BadgeGreen}{current}}\layeritem{replay}{leanchecker (official kernel replay)}{\badge{BadgeGreen}{current}}\layeritem{hardened}{safe\_verify (hardened screen)}{\badge{BadgeGreen}{current}}\layeritem{decorrelated}{lean4export + nanoda (decorrelated kernel)}{\badge{BadgeGreen}{current}}}
\sourcenote{\badge{BadgeGray}{comparator: none} no fixture in the three-implementation comparator harness (compiled / kernel-interpreted / reference-interpreted) exercises this rule.}
\ruleblock{HOL4 statement}
\sourcenote{theory \texttt{bern\_tr\_soundpsd}, HOL4 development script \texttt{\texttt{bern\_tr\_soundpsdScript.sml}} \badge{BadgeGreen}{script bound}}
\holhead{Theorem sound\_psdCertKernel}
\begin{holsrc}
Theorem sound_psdCertKernel:
  ∀H sq m k A N B innerT Z ts conclusion.
    psdCheckK m k A N B innerT Z ∧
    psdConclCheck (m + k) A ts conclusion ∧
    (∃s. tmBshape H conclusion = SOME s) ∧
    tmFinWires H conclusion ∧
    (∀w. MEM w (termFreeWires conclusion) ⇒
         ∃τ. H.wireType w = SOME τ ∧ 0 < shapeTotalSize τ.shape) ∧
    syzygyConclusionFormB conclusion sq.post ⇒
    seqSemanticallyValid H sq
\end{holsrc}
\ruleblock{HOL verifier receipts}
\badge{BadgeBlue}{documentary} \texttt{hol-strict} CLEAN: theory \texttt{bern\_tr\_soundpsd}, 4
theorems, 0 cheat-tagged, 0 nonstandard axioms (whitelist:
\texttt{DISK\_THM}; the four \texttt{bool} axioms), 249-theory
ancestry --- verification run of 2026-07-22 on an isolated build host,
against a content-addressed checkout of the exact sources cited here
(git blob \texttt{c01e0b0c\ldots}). The shared semantic spine
(\texttt{sound\_psdCertOf}) is displayed in the /Tri section per the
family-preamble convention; this corollary discharges the
quadratic-form fact via the ported kernel-extension bridge
(\texttt{psdCheckK\_quadForm}, theory \texttt{bern\_psd}, where the
witness check itself --- \texttt{psdCheckK}, over the extension
\texttt{kernelExt} --- is ported with kernel-EVAL receipts). The
payload record and its wellformedness arm ride the ported proof-object
stratum (theory \texttt{bern\_tr\_catdata2}, 272 theorems, 0 offences,
same run).

\sourcenote{HOL receipts are documentary: quoted results of \texttt{hol-strict} runs
and census rows from the HOL4 development\textquotesingle s build
evidence, not minted receipt artefacts. A HOL receipt store (the Lean
store\textquotesingle s twin) is named future work.
}
\ruleblock{Proof}

\sourcenote{Genre: admissible rule (tier 2) --- reduction to verified
machinery. The certificate-family preamble is stated in the
\rulename{PsdCertTri} section and applies verbatim. This is the
\emph{semidefinite} member: it certifies $A \succcurlyeq 0$ together
with an exactly-checked kernel, for the boundary cases the two strict
tiers cannot touch.}

\begin{brecipe}[the \rulename{PsdCertKernel} instance]\label{rec:psdk-instance}
Data: $A \in \Q^{(m+k)\times(m+k)}$ with the $k$ kernel coordinates
\emph{last} (the producer applies its permutation before emission ---
conjugation by a permutation is index renaming on matrix and term
vector alike, so the identity-permutation form loses no generality);
$N \in \Q^{m \times k}$; the kernel-completion row matrix
$L := (\,I_m \mid N\,)$; a leading block $B \in \Q^{m\times m}$ with a
/T-or-/DD inner witness $Z$. Checked premises: the factorization
$A = L^{\mathsf T} B\, L$ holds \emph{entrywise exactly} (recomputed),
and $B$ passes the selected strict inner check. Then
\[
\inferrule*[Right=\rulename{PsdCertKernel}]
  {\;\mathrm{checkKernel}(A, N, B, Z) = \mathsf{true} \\
    c \equiv -\textstyle\sum_{i,j} A_{ij}\,p_i\,p_j\;}
  {\Phi \vdash \Psi}
\qquad\rightsquigarrow\qquad
\text{\upshape discharged by } \texttt{sound\_psdCertKernel}.
\]
The factorization check is --- given the symmetry of $B$, which the
selected inner check itself verifies --- equivalent to the three
separate conditions ($A$ symmetric; $A \cdot K = 0$ exactly for
$K = \binom{-N}{\,I_k}$; leading block of $A$ equal to $B$): forward,
$B$'s symmetry makes $A = L^{\mathsf T} B L$ symmetric and $LK = 0$
kills $AK$; backward, $AK = 0$ forces the cross block $BN$ and then
the trailing block $N^{\mathsf T} B N$. Either way the off-leading
blocks are forced, so a forged $N$ or $B$ refuses.
\end{brecipe}

\begin{bprop}[matrix warrant]\label{prop:psdk-matrix}
If $B \succ 0$ (hence $\succcurlyeq 0$) and
$A = L^{\mathsf T} B L$ for the rectangular $L = (I \mid N)$, then
$A \succcurlyeq 0$. Moreover $x^{\mathsf T} A x = 0$ exactly on
$\operatorname{span} K$, so the certificate fixes
$\ker A = \operatorname{span} K$ and
$\operatorname{rank} A = m$.
\end{bprop}
\begin{proof}
Positive semidefiniteness is rectangular congruence, with no
invertibility anywhere: $x^{\mathsf T} A x = (Lx)^{\mathsf T} B\,(Lx)
\geq 0$ (the kernel proof is this one line,
\texttt{posSemidef\_of\_factor}). For the kernel content, write
$x = (x_1, x_2)$ against the block split and note
$Lx = x_1 + N x_2$, while $L K = -N + N = 0$; with $B \succ 0$,
$x^{\mathsf T} A x = 0$ iff $Lx = 0$ iff $x = K x_2$ --- i.e.\ the
zero set of the form is exactly $\operatorname{span} K$. That zero set
\emph{is} the kernel: if $Lx = 0$ then $Ax = L^{\mathsf T} B (Lx) =
0$, and conversely $Ax = 0$ forces $x^{\mathsf T} A x = 0$, hence
$Lx = 0$. The columns of $K$ are independent (its lower block is
$I_k$), so $\ker A = \operatorname{span} K$ has dimension $k$ and
$\operatorname{rank} A = m$ by rank--nullity.
\end{proof}

\begin{bremark}[exact kernels are forced, not a method limitation]
A rounding-tolerant witness for kernel \emph{directions} cannot exist
for a residual-style check --- one attesting the kernel claim only up
to an $\varepsilon$-residual, with no exact kernel equation:
$\{A \succcurlyeq 0,\ \det A = 0\}$ is a boundary set with empty
interior, so a check whose acceptance region contains an
entrywise-$\varepsilon$ neighbourhood of a true instance also accepts
some $A - \varepsilon' I \not\succcurlyeq 0$. The split is therefore
forced, not preferred: exact kernel data plus a rounding-tolerant bulk
(the /DD inner witness), with tolerance able to live in the bulk and
nowhere else --- and it matches the intended
producers, whose kernels are \emph{semantic} (a conservation law, a
syzygy relation, the all-ones vector of a graph Laplacian): $N$ comes
from theorems, and floats only ever touch $Z$.
\end{bremark}

\begin{bremark}[each hypothesis, and the corners]
\emph{Entrywise-exact factorization}: the whole kernel claim rides on
$A K = 0$ holding exactly; an $\varepsilon$-approximate factorization
certifies nothing (previous remark). \emph{The strict inner witness}:
$B \succ 0$ is what confines the zero set to
$\operatorname{span} K$ --- a merely semidefinite $B$ would leave
undeclared kernel directions inside the leading block. \emph{Kernel
coordinates last}: bookkeeping only, by the permutation argument in
the recipe. Corners are structural, not special-cased: $k = 0$ gives
$L = I$ and the rule degenerates to its inner strict tier ($A = B$);
$m = 0$ forces $A = 0$ (empty $B$, no inner witness content --- the
factorization premise \emph{is} the entire content).
\end{bremark}

\begin{bremark}[what reaches the formula layer]
The checker verifies the kernel and rank content (a forged $N$
refuses), but the formula language states only the
quadratic-form conclusion $\sum_{i,j} A_{ij}\,p_i\,p_j \geq 0$ ---
which for this rule is honest ($\succcurlyeq$-shaped, per the
preamble's strictness discipline) and loses only the kernel/rank
\emph{statement}, not its checking. Making rank and kernel
formula-expressible needs closed-tensor terms and a tensor-valued
closed evaluator --- the same named extension the
\rulename{PConstExact} section records for exact matrix identities;
short of that extension, the kernel content is certificate-level
assurance
surfaced here, deliberately not hand-faked into the formula grammar.
\end{bremark}

\begin{anchors}
Recipe \ref{rec:psdk-instance} warrant & \texttt{sound\_psdCertKernel} (Soundness); shared transport \texttt{sound\_psdCertOf} (see \rulename{PsdCertTri})\\
rectangular congruence (Prop.\ \ref{prop:psdk-matrix}) & \texttt{PsdCert.posSemidef\_of\_factor}\\
check-to-theorem bridge & \texttt{PsdCert.checkKernel\_posSemidef} (inner tier via \texttt{checkTri\_posDef} / \texttt{checkDiag\_posDef}); corollary \texttt{checkKernel\_quadForm\_nonneg}\\
the executable check & \texttt{PsdCert.checkKernel} with the completion rows \texttt{kernelExt}\\
conclusion identity + readings & \texttt{PsdCert.conclCheck}; the \rulename{Syzygy} section's conclusion-form selector\\
\end{anchors}


\rulesection{31}{ReLU-Split}{admissible rule}
\label{blcatrule-ReLUSplit}
\ruleblock{The rule (Steele metanotation)}
\sourcenote{as realized in the mechanized development; source row: \texttt{\texttt{rules.toml}} --- \emph{ReLU-Split}. The metanotation statement is typeset by hand for notational quality; its mechanized ground truth is the Lean statement below, extracted from the sources at every build of this book and hash-bound. The rule \emph{roster} is machine-checked (every mechanized rule tag round-trips to exactly one catalogue row); the statement \emph{text} is editorial, checked by eye against the extracted Lean beside it.}
\begin{mathpar}
\inferrule*[Right=\rulename{ReLU-Split}]{
  \Phi \wedge (h \geq 0) \vdash \Psi\sub{b}{h} \\
  \Phi \wedge (h \leq 0) \vdash \Psi\sub{b}{0} \\
  \ReLU_{\{h\}\to\{y\}} \in H \quad (y \text{ typed at rank } 0)
}{
  \Phi \vdash \Psi
}
\end{mathpar}

\par\medskip
\rulename{ReLU-Split} is the case-split rule of the ReLU family: from
two branch proofs — the goal with the output position replaced by $h$,
under $\Phi$ strengthened by $h \geq 0$; and the goal with the output
position replaced by $0$, under $\Phi$ strengthened by $h \leq 0$ —
conclude the goal outright, per rank-0 node: the tensor-shaped reading
of the rule is \emph{false}, not merely unrealized (the proof
section\textquotesingle s scope remark carries the mixed-sign
countermodel), so a producer splits scalar by scalar. It is the \emph{complete}-verification
move of the neural-network literature (the Reluplex and
branch-and-bound lineage): where the affine envelope of
\rulename{LiRPA} loses too much, a producer splits on the
neuron\textquotesingle s sign and pays with proof size — each split
doubles the branch count, which is exactly the branch-and-bound
tradeoff.

Unlike every other rule in this book so far, \rulename{ReLU-Split} has
\emph{no certificate node}: no payload constructor, no wire tag — a
certificate stream never contains one, so there is nothing for this
rule\textquotesingle s checker column to recompute. A producer
elaborates it away. The elaboration\textquotesingle s target is the
\emph{derivation} the rule abbreviates (case analysis over the
sign-split fact, equality transport in each branch, and structural
plumbing), whose nodes the checker checks by their own rules — a
certificate fragment that becomes emittable exactly when the
context-carrying case split it opens with lands as a core constructor
(today that form is a satellite soundness theorem only; the proof
section states the dependency). What stands unconditionally is the
rule-level semantic theorem of the proof
section\textquotesingle s mechanization remark.

Each ingredient of the two premises is load-bearing for the
elaboration. The branch strengthenings ($h \geq 0$, $h \leq 0$) are
what the case fact discharges into each branch context — remove one
and that branch proof has no purchase on the sign information. The
substitution alignment ($\Psi$ with the output position filled by $h$,
respectively $0$) is what the equality transport rewrites through:
inside the active branch the node\textquotesingle s fact $y = h$
carries $\Psi[b \coloneqq h]$ to $\Psi$, and inside the inactive
branch $y = 0$ carries $\Psi[b \coloneqq 0]$ to $\Psi$.

\cataloguenote{The checker column remains empty for the structural reason: no
certificate stream ever contains a \rulename{ReLU-Split} node, so
there is nothing for the checker to recompute. The rule-level
statement exists:
\texttt{sound\_reluSplit} takes the two branch validities to the
conclusion\textquotesingle s validity, at the rank-$0$ scope the
recipe\textquotesingle s own per-scalar route always meant — and the
companion theorem \texttt{reluSplit\_rank0\_necessary} mechanizes the
mixed-sign countermodel showing the tensor-shaped reading is false,
upgrading the family\textquotesingle s shape-scope remark from
derivation-obstruction to semantic necessity for this rule. A
verified elaboration function producing the derivation does not yet
exist. The derivation requires the context-carrying form
of disjunction elimination — the sign-split disjunction arrives as a
sub-conjunct of the premise, which the plain sequent form of
$\vee$E cannot reach; see the $\vee$E section\textquotesingle s note.
}
\ruleblock{Lean statement}
\par\medskip
The rule-level semantic lemma, at the rank-$0$ scope: the template
$\Psi$ is an arbitrary formula with an arbitrary marker wire for $b$;
the three substitution instances are hypotheses via
\texttt{substituteInFormula} (\rulename{EqElim}\textquotesingle s
substitution predicate); the branch strengthenings are
\rulename{ReLU-Pos}/\rulename{ReLU-Neg}\textquotesingle s premise
atoms verbatim; the linkage and its extraction are the
family\textquotesingle s (\rulename{LiRPA} section). The second
theorem is the scope\textquotesingle s necessity witness: a
shape-$[2]$ fixture with mixed-sign input on which both branch
premises are valid and the conclusion fails.

\declhead{BernsteinLogic.ReluSplit}{sound\_reluSplit}{\badge{BadgeGreen}{statement bound}}
\begin{leansrc}
/-- ReLU-Split (rule-level semantic lemma; rank-0 scope, which
`reluSplit_rank0_necessary` shows is mandatory): from the validity of
the two branch sequents — the conclusion template at the preactivation
under `Φ ∧ (h ≥ 0)`, and at the constant 0 under `Φ ∧ (h ≤ 0)` — the
conclusion sequent `Φ ⊢ Ψ` is valid, where `Ψ` is the template at the
cited ReLU application's output wire. -/
theorem sound_reluSplit {H : Hypernet}
    {seq activePremise inactivePremise : Sequent}
    {preact post marker : Wire} {bodyForm : Formula}
    (ihActive : activePremise.semanticallyValid H)
    (ihInactive : inactivePremise.semanticallyValid H)
    (hlink : reluLinked H preact post)
    (hrank0 : ∃ τ, H.wireType post = some τ ∧ τ.shape = [])
    (hActPre : activePremise.pre
      = .and seq.pre (.atom .ge (.var preact) (.const 0)))
    (hInactPre : inactivePremise.pre
      = .and seq.pre (.atom .le (.var preact) (.const 0)))
    (hActPost : substituteInFormula bodyForm marker (.var preact)
      activePremise.post)
    (hInactPost : substituteInFormula bodyForm marker (.const 0)
      inactivePremise.post)
    (hconcl : substituteInFormula bodyForm marker (.var post) seq.post) :
    seq.semanticallyValid H
\end{leansrc}
\declhead{BernsteinLogic.ReluSplit}{reluSplit\_rank0\_necessary}{\badge{BadgeGreen}{statement bound}}
\begin{leansrc}
/-- The rank-0 hypothesis of `sound_reluSplit` is NECESSARY: the same
rule schema with the shape scope deleted is refuted — at a shape-[2]
ReLU node with mixed-sign input, both branch premises are valid while
the conclusion fails. (The header records the shape-[1] countermodel,
via broadcast-arity mixing, in prose; this is the mechanized mixed-sign
witness.) -/
theorem reluSplit_rank0_necessary :
    ¬ (∀ (H : Hypernet) (seq activePremise inactivePremise : Sequent)
        (preact post marker : Wire) (bodyForm : Formula),
        activePremise.semanticallyValid H →
        inactivePremise.semanticallyValid H →
        reluLinked H preact post →
        activePremise.pre
          = .and seq.pre (.atom .ge (.var preact) (.const 0)) →
        inactivePremise.pre
          = .and seq.pre (.atom .le (.var preact) (.const 0)) →
        substituteInFormula bodyForm marker (.var preact)
          activePremise.post →
        substituteInFormula bodyForm marker (.const 0)
          inactivePremise.post →
        substituteInFormula bodyForm marker (.var post) seq.post →
        seq.semanticallyValid H)
\end{leansrc}
\ruleblock{Correspondence (Steele display $\leftrightarrow$ Lean)}
\absentbox{NO CORRESPONDENCE DICTIONARY YET}{no reviewer-certified dictionary maps the displayed rule onto the mechanized statements above; until one is written and certified, that mapping rests on this section\textquotesingle s editorial care alone.}
\ruleblock{Lean verifier receipts}
\sourcenote{module \texttt{BernsteinLogic.ReluSplit}; bound build \texttt{01611bc8a2c3\ldots} (bound from receipt 6dd1e30ab7bf0a19, 2026-07-20T05:56:22Z). A receipt records that named checking layers accepted a build: (module, compiled-artifact SHA-256, environment SHA-256, checker layers), minted when the checking gate ran; the store receipts attesting the bound build are listed with the layers each run credited. Nothing here asks the reader to trust the store: every layer is an independently runnable program over the quoted sources --- rebuild the module and re-run the checkers to reproduce the verdict. Receipts attesting superseded builds credit no layer and are omitted from this rendering (the count of omissions, if any, is noted below).}
\begin{receipttable}
\receiptrow{6dd1e30ab7bf0a19}{2026-07-20T05:56:22Z}{build+decorrelated}{\badge{BadgeGreen}{current}}
\receiptrow{f5e25a81aaa84445}{2026-07-20T20:44:58Z}{build+decorrelated}{\badge{BadgeGreen}{current}}
\receiptrow{ca5158835179141d}{2026-07-20T21:01:26Z}{build+replay+hardened+decorrelated}{\badge{BadgeGreen}{current}}
\receiptrow{1a900df2daf383f4}{2026-07-21T05:28:56Z}{build+replay+hardened+decorrelated}{\badge{BadgeGreen}{current}}
\end{receipttable}
\layerline{\layeritem{build}{lake build (elaboration)}{\badge{BadgeGreen}{current}}\layeritem{replay}{leanchecker (official kernel replay)}{\badge{BadgeGreen}{current}}\layeritem{hardened}{safe\_verify (hardened screen)}{\badge{BadgeGreen}{current}}\layeritem{decorrelated}{lean4export + nanoda (decorrelated kernel)}{\badge{BadgeGreen}{current}}}
\sourcenote{\badge{BadgeGray}{comparator: none} no fixture in the three-implementation comparator harness (compiled / kernel-interpreted / reference-interpreted) exercises this rule.}
\ruleblock{HOL4 statement}
\absentbox{NO HOL STATEMENT YET}{\rulename{ReLU-Split} has no certificate representation on either side
of the seam: no payload constructor, no wire tag, no check function —
a producer elaborates the rule into its derivation, so a certificate
stream never contains a \rulename{ReLU-Split} node and there is
nothing for a HOL port to mirror. Any future HOL treatment would be a
formalization of the rule-level statement or the derivation itself
(the Lean side of the former exists — \texttt{sound\_reluSplit}), not
a port of checker code.
}
\ruleblock{HOL verifier receipts}
\absentbox{NONE}{no HOL verification receipt exists for this rule.}
\ruleblock{Proof}

\sourcenote{Genre: admissible rule (tier 2) --- but of the family's
third kind. \rulename{PConstExact} is admissible by reduction to
verified machinery; its siblings here are admissible with shipped
direct lemmas; \rulename{ReLU-Split} is admissible in the purest
sense: it has \emph{no proof-term constructor and no wire tag} --- a
certificate stream never contains a \rulename{ReLU-Split} node. The
derivation below is the rule's intended elaboration, and it becomes a
shippable certificate fragment exactly when the context-carrying case
split it opens with becomes a core constructor
(\autoref{rem:relusplit-ctxve}); the warrant that stands on its own
is the rule-level theorem: \texttt{sound\_reluSplit} states the rule
itself, at the rank-$0$ scope \autoref{rem:relusplit-scope} proves
is the rule's true extent.}

\begin{brecipe}[the \rulename{ReLU-Split} elaboration]
\label{rec:relusplit-instance}
The rule, as the source catalogue states it:
\[
\inferrule*[Right=\rulename{ReLU-Split}]{
  \mathcal{T}_{+} \rhd \Phi \wedge (h \geq 0) \vdash \Psi\sub{b}{h} \\
  \mathcal{T}_{-} \rhd \Phi \wedge (h \leq 0) \vdash \Psi\sub{b}{0}
}{
  \Phi \vdash \Psi
}
\]
where $b$ marks the occurrences of the output wire $y$ of a ReLU
application $\ReLU_{\{h\}\to\{y\}} \in H$, so that $\Psi =
\Psi\sub{b}{y}$. Two side conditions ride under that ``where'', made
explicit by the realized statement exactly as the siblings' displays
make theirs: the linkage itself (the warrant consumes the family's
port lookup), and --- semantic necessity, not caution --- the output
wire typed at rank $0$ (\autoref{rem:relusplit-scope}). The
intended elaboration of one such node is the derivation:
\begin{enumerate}
\item the node's sign-split fact --- $(h \leq 0 \wedge y = 0) \vee
  (0 \leq h \wedge y = h)$, the packaged semantic contract stated in
  the \rulename{LiRPA} section's family template --- enters as a
  sub-conjunct of the premise (the sp-axiom convention: for an
  application in $H$, its packaged contract is available premise-side
  as an axiom, true at every admissible valuation --- the same
  convention under which the cell rules' one-hot disjunction enters);
\item $\vee$E \emph{in its context-carrying form} splits on it (see
  \autoref{rem:relusplit-ctxve});
\item \emph{active branch} (context gains $0 \leq h$ and $y = h$):
  the branch context is aligned with $\mathcal{T}_{+}$'s premise $\Phi
  \wedge (h \geq 0)$ and consulted, in five named moves --- $\wedge$E
  extracts $\Phi$ and the case conjunct $0 \leq h$; the bound node's
  polarity-flip reading re-orients the latter to $h \geq 0$
  (\rulename{PBound}, reading (G), the one-node move displayed at the
  \rulename{ReLU-NN} section's active branch; the \rulename{ReLU-Above}
  section walks the toolkit's longer chains); $\wedge$I rebuilds $\Phi
  \wedge (h \geq 0)$; \rulename{Cut} against $\mathcal{T}_{+}$ yields
  $\Psi\sub{b}{h}$ in the branch context; and \rulename{EqSym} on the
  context's $y = h$ followed by \rulename{EqElim} along $h = y$
  (wire-to-wire, the node's one shared shape) transports it to
  $\Psi\sub{b}{y} = \Psi$;
\item \emph{inactive branch} (context gains $h \leq 0$ and $y = 0$):
  the same route with the mirrored atoms, four grouped moves
  ($h \leq 0$ needs no
  reorientation, so the \rulename{PBound} step drops out);
  $\mathcal{T}_{-}$ yields $\Psi\sub{b}{0}$, and \rulename{EqSym} +
  \rulename{EqElim} along $0 = y$ transport it to $\Psi$;
\item \rulename{Cut} and \rulename{Weaken} stitch the case analysis
  back under the original $\Phi$.
\end{enumerate}
Step 4's transport has a rank-$0$ constant side, so at tensor-shaped
instances it is stopped by \rulename{EqElim}'s broadcast-shape gate,
exactly as the family's shape-scope remark records
(\autoref{rem:relunn-shapescope}, \rulename{ReLU-NN} section) ---
at such instances a producer splits per scalar node
(\rulename{ReLU-Neg}'s direct node supplies the shape-general fact
$y = 0$, but transporting it into an arbitrary $\Psi$ rides the same
gate, so the per-scalar split is the general route --- and
necessarily so: \autoref{rem:relusplit-scope} proves the
tensor-shaped reading admits a countermodel).
\end{brecipe}

\begin{bremark}[why the context-carrying $\vee$E]
\label{rem:relusplit-ctxve}
The disjunction being eliminated does not stand alone as the premise:
it arrives \emph{inside} the premise conjunction, next to everything
else $\Phi$ carries. The plain sequent form of $\vee$E ($\Phi_1
\vdash \Psi$ and $\Phi_2 \vdash \Psi$ give $\Phi_1 \vee \Phi_2 \vdash
\Psi$) cannot reach it --- and this is not a stylistic gap: the
$\vee$E section records the countermodel showing the context-carrying
form is strictly stronger than the printed one (no elaboration of the
former into the latter exists). This derivation is the load-bearing consumer of that form. The form is
proven sound as a satellite lemma (\texttt{sound\_orECtx}), and that
is currently \emph{all} it is: it has no core constructor and no wire
tag, so the derivation above is checkable as a certificate fragment
only once the constructor lands --- until then the rule's standing
warrant is \autoref{rem:relusplit-mech}'s theorem, and the
elaboration is the design the constructor exists to serve.
\end{bremark}

\begin{bremark}[mechanization status --- the rule-level statement]
\label{rem:relusplit-mech}
The checker column remains empty for the structural reason it always
was: no certificate ever contains the rule, so there is no payload to
validate and no soundness case in the checker's recursion. The
rule-level statement exists: \texttt{sound\_reluSplit} is
the semantic lemma --- from the validity of the two branch sequents
and the linkage, the conclusion sequent is valid --- built from
exactly the vocabulary the rest of the catalogue already uses (the linkage extraction
\texttt{reluLinked\_facts}, the substitution predicate
\texttt{substituteInFormula} that \rulename{EqElim}'s payload already
uses, and the semantic transport \texttt{Formula.holds\_subst\_congr}
that \rulename{EqElim}'s soundness already rides). The template
$\Psi$ with its marked occurrences $b$ is an arbitrary formula over an
arbitrary marker wire, with the three substitution instances
($\Psi\sub{b}{h}$, $\Psi\sub{b}{0}$, $\Psi = \Psi\sub{b}{y}$) as
hypotheses --- no freshness condition, since the proof transports
along value equality, never syntax. ``Derivable, by inspection of
this recipe'' is thereby backed by a theorem at every
instance the recipe itself reaches (rank-$0$ nodes ---
\autoref{rem:relusplit-scope}); a
verified elaboration function producing
Recipe~\ref{rec:relusplit-instance}'s derivation does not yet exist.
\end{bremark}

\begin{bremark}[the rank-$0$ scope is semantic necessity]
\label{rem:relusplit-scope}
\texttt{sound\_reluSplit} hypothesizes the output wire typed at
rank~$0$, and this is the rule's true extent rather than proof
convenience. The family's shape-scope remark
(\autoref{rem:relunn-shapescope}) records that the
\emph{derivation}'s constant-sided transport is stopped by
\rulename{EqElim}'s broadcast-shape gate at tensor shapes; for the
siblings the direct lemma is nonetheless shape-general, so the gate
reads there as a technicality of the derivation route. Not here:
\texttt{reluSplit\_rank0\_necessary} (same module) mechanizes a
countermodel showing the shape-general reading of
\rulename{ReLU-Split} is \emph{false}. Template $\Psi(b) = (b = 0)
\vee (b = h)$ over a shape-$[2]$ node, $\Phi = \top$: both premises
are valid ($\mathcal{T}_{+}$ through $h = h$, $\mathcal{T}_{-}$
through $0 = 0$), yet at the mixed-sign valuation $h = [1, -1]$,
$y = [1, 0]$ neither disjunct of the conclusion holds. Mixed signs
are exactly what the rule's case analysis cannot see from
sign-definite branch hypotheses --- the semantic content of ``each
split doubles the branch count'': a tensor node is $2^{n}$ scalar
cases, not $2$. Even shape $[1]$ fails, by a different leak
(broadcast-arity mixing): take $\Psi(b) = (b = h) \vee (b + w \geq 0)$
with $w$ typed at $[2]$ and premise $\Phi = (w \geq 0)$. Both branch
sequents are valid --- $\mathcal{T}_{+}$ through $h = h$,
$\mathcal{T}_{-}$ through $0 + w \geq 0$, the rank-$0$ constant
broadcasting to $w$'s shape --- yet at $h = [-1]$, $y = [0]$,
$w = [1, 1]$ the conclusion fails: $y = h$ fails entrywise, and the
$[1]$-shaped $y + w$ atom is unsatisfiable outright by the
strict-shape arithmetic, which broadcasts rank-$0$ sides only (this
leak is stated as prose beside the mechanized $[2]$ witness, not
itself mechanized). The recipe's step-4 note
(``the per-scalar split is the general route'') is thus not a
workaround for a gate but the only sound reading; producers split per
rank-$0$ node, and \texttt{sound\_reluSplit} is the rule those splits
instantiate.
\end{bremark}

\begin{bremark}[when a producer splits]
Splitting is the \emph{complete}-verification move of the
neural-network literature --- Reluplex's case
analysis~\cite{katz2017reluplex}, systematized as branch and
bound~\cite{bunel2020branch}: exact where \rulename{LiRPA}'s envelope
is lossy, at the price of one branch doubling per split neuron, so
proof size is exponential in the split count. The practical frontier
--- which neurons to split, which to relax --- is the branch-and-bound
literature's whole subject; the calculus prices the choice honestly
(a split is two subtrees, an envelope is one node) and takes no
position on it.
\end{bremark}

\begin{anchors}
the sign-split fact (step 1) & \texttt{PrimApp.SemOk} --- the packaged contract; conservativity via \texttt{relu\_projection} (\rulename{LiRPA} section)\\
context-carrying case split (step 2) & \texttt{sound\_orECtx} (satellite lemma, Admissible.lean)\\
equality transport (steps 3--4) & \texttt{sound\_eqElim}; \texttt{substituteInFormula}\\
rule-level statement (rank-$0$ scope) & \texttt{sound\_reluSplit} (ReluSplit.lean) --- \autoref{rem:relusplit-mech}\\
scope necessity (mixed-sign countermodel) & \texttt{reluSplit\_rank0\_necessary}; fixture \texttt{ReluSplitScope.HSplit} --- \autoref{rem:relusplit-scope}\\
verified elaboration function & \emph{none exists} --- the remaining gap, with the core constructor for the context-carrying case split\\
\end{anchors}


\blcatsection{Assume--Guarantee Composition}
\rulesection{32}{Promise}{admissible rule}
\label{blcatrule-Promise}
\ruleblock{The rule (Steele metanotation)}
\sourcenote{as realized in the mechanized development; source row: \texttt{\texttt{rules.toml}} --- \emph{Promise}. The metanotation statement is typeset by hand for notational quality; its mechanized ground truth is the Lean statement below, extracted from the sources at every build of this book and hash-bound. The rule \emph{roster} is machine-checked (every mechanized rule tag round-trips to exactly one catalogue row); the statement \emph{text} is editorial, checked by eye against the extracted Lean beside it.}
\begin{mathpar}
\inferrule*[Right=\rulename{Promise}]{
  \Phi_1 \vdash Q \\
  \Phi_2 \wedge Q \vdash \Psi \\
  \wires(Q) \subseteq Y \\
  \wires(\Phi_1) \cap \bigl(\wires(\Phi_2) \cup \wires(\Psi)\bigr) \subseteq Y
}{
  \Phi_1 \wedge \Phi_2 \vdash \Psi
}
\end{mathpar}
$Y$ = the declared interface (payload wire list; bundle entries addressed as $\mathrm{TIdx}(w,i)$). The two confinement checks are the locality warrant, not soundness needs: the contract $Q$ is the \emph{datasheet} --- the public artefact whose vocabulary is the separator algebra --- and the cost identities (LP size, ZK disclosure set) quantify over certificates carrying these checks.

\ruleflavor{a Bernstein Logic rule outright: the baked base theory lists it in its core roster (grouped there under an assume--guarantee comment), the certified checker decodes its wire node, and its certificate relations ride the pin table --- there is no intermediate composition logic between this rule and the base calculus. The tier word above classifies proof theory (the unpack is \rulename{Weaken} + $\wedge$-restructuring + \rulename{Cut}); core-roster membership is wire placement.}
\par\medskip
\rulename{Promise} is the calculus\textquotesingle s acyclic composition
rule: it discharges an assumption one component makes about another. The
supplier proves the contract $Q$ from its own premises
($\Phi_1 \vdash Q$); the consumer proves the goal \emph{using} $Q$
($\Phi_2 \wedge Q \vdash \Psi$); the composite concludes
$\Phi_1 \wedge \Phi_2 \vdash \Psi$ with the contract discharged
internally --- cut, wearing interface bookkeeping. A producer reaches for
it wherever a system decomposes along a directed interface: certify the
downstream component against a \emph{datasheet} (the contract, written
over the declared interface wires only), certify the upstream component
as meeting the datasheet, and compose without ever re-deriving either
side inside the other.

What the checker recomputes: the splice plumbing and the locality
warrant. The plumbing --- the contract sits at a payload-declared
conjunct path in the consumer\textquotesingle s premise; the composite
premise is the consumer\textquotesingle s with the contract replaced by
the supplier\textquotesingle s premise at that path --- is verified by
exact read/replace on the conjunction tree. The locality warrant is two
wire-confinement checks: the contract\textquotesingle s free wires lie
inside the declared interface, and the supplier\textquotesingle s wires
meet the consumer\textquotesingle s and the goal\textquotesingle s only
inside it.

The hypotheses split by role, and the split is the
rule\textquotesingle s point. Soundness consumes only the plumbing:
drop the path checks and the \textquotedblleft composite\textquotedblright{}
premise could be unrelated to the two sub-proofs. The confinement checks
are \emph{warrant, not soundness} --- cut is sound regardless --- but
dropping them destroys the datasheet reading: a contract mentioning
supplier-private wires is no longer a public interface artefact, and
every quantitative consequence hung on the interface (certificate size
scaling with the separator rather than the components, bounded
disclosure of component internals) quantifies over certificates that
carry these checks. The proof section separates the two roles
explicitly.

\cataloguenote{Tier reading: proof-theoretically the rule is admissible --- one node
unpacks into \rulename{Cut} against premise-transported copies of the
cited sub-derivations, a global rewrite of those sub-proofs rather
than a local chain (the proof section\textquotesingle s recipe walks
why: \rulename{Weaken} and \rulename{Cut} are head-positional) ---
while in the realized
calculus it is a native constructor, kept for cost and for the
locality-warrant payload. The tier column records the proof-theoretic
status; constructor membership is wire placement.
}
\ruleblock{Lean statement}
\declhead{BernsteinLogic.Soundness}{PromiseStep}{\badge{BadgeGreen}{statement bound}}
\begin{leansrc}
/-- Applicability data of the Promise rule, ctx-style — field-for-field
the reference checker's obligations, and field-for-field `PPromiseData`
minus the two sub-proofs (which enter `sound_promise` as validity
hypotheses). The two confinement fields are checked side-conditions that
soundness does NOT consume (AG draft §1.2: cut is sound regardless) —
they are the locality warrant. -/
structure PromiseStep where
  seq : Sequent
  supplier : Sequent
  consumer : Sequent
  contract : Formula
  path : ConjPath
  iface : List Wire
  ctx : Formula
  supplierPost : supplier.post = contract
  pathReads : consumer.pre.conjGet? path = some contract
  prePlumb : consumer.pre.conjReplace? path supplier.pre = some seq.pre
  postEq : seq.post = consumer.post
  ctxReads : consumer.pre.conjReplace? path .tt = some ctx
  contractConfined : formulaWiresWithin contract iface
  separation : wiresSeparated supplier.pre ctx seq.post iface
\end{leansrc}
\declhead{BernsteinLogic.Soundness}{sound\_promise}{\badge{BadgeGreen}{statement bound}}
\begin{leansrc}
/-- Soundness of Promise: valid supplier + valid consumer give a valid
composite. Read the supplier's premise at the path, obtain the contract
from supplier validity, replace it back to recover the consumer's
premise, conclude by consumer validity. Pure predicate logic over the
semantics; the confinement fields are untouched. -/
theorem sound_promise {H : Hypernet} (st : PromiseStep)
    (ihs : st.supplier.semanticallyValid H)
    (ihc : st.consumer.semanticallyValid H) :
    st.seq.semanticallyValid H
\end{leansrc}
\declhead{BernsteinLogic.Soundness}{Formula.holds\_of\_conjGet}{\badge{BadgeGreen}{statement bound}}
\begin{leansrc}
/-- The conjunct addressed by a path is entailed: if Γ holds, Γ|π holds.
(Promoted from Admissible.lean — the ∨E-ctx kit; the Promise contract
path reuses it.) -/
theorem Formula.holds_of_conjGet {H : Hypernet} {σ : H.Valuation} :
    ∀ {π : ConjPath} {Γ d : Formula},
    Γ.conjGet? π = some d → Γ.holds σ → d.holds σ
\end{leansrc}
\declhead{BernsteinLogic.Soundness}{Formula.holds\_conjReplace}{\badge{BadgeGreen}{statement bound}}
\begin{leansrc}
/-- Replacing the addressed conjunct by a formula that holds preserves
holding: if Γ holds and g holds, Γ[π := g] holds. (Promoted from
Admissible.lean.) -/
theorem Formula.holds_conjReplace {H : Hypernet} {σ : H.Valuation} :
    ∀ {π : ConjPath} {Γ Γ' g : Formula},
    Γ.conjReplace? π g = some Γ' → Γ.holds σ → g.holds σ → Γ'.holds σ
\end{leansrc}
\declhead{BernsteinLogic.Soundness}{Formula.conjGet\_of\_conjReplace}{\badge{BadgeGreen}{statement bound}}
\begin{leansrc}
/-- Reading at the path of a successful replacement returns the
replacement. (Promoted from Composition.lean §G2.) -/
theorem Formula.conjGet_of_conjReplace :
    ∀ {π : ConjPath} {Γ Γ' g : Formula},
    Γ.conjReplace? π g = some Γ' → Γ'.conjGet? π = some g
\end{leansrc}
\declhead{BernsteinLogic.Soundness}{Formula.conjReplace\_back}{\badge{BadgeGreen}{statement bound}}
\begin{leansrc}
/-- Replacing back the conjunct a read found recovers the original: the
round trip Γ ⟶[π := g] Γ' ⟶[π := Γ|π] Γ. (Promoted from
Composition.lean §G2.) -/
theorem Formula.conjReplace_back :
    ∀ {π : ConjPath} {Γ Γ' q g : Formula},
    Γ.conjGet? π = some q → Γ.conjReplace? π g = some Γ' →
    Γ'.conjReplace? π q = some Γ
\end{leansrc}
\par\medskip
\textbf{The confinement side conditions, as checked.}\enspace The
proof\textquotesingle s premise-transport induction leans on the
wire-discipline predicates the checker recomputes; their definitions
are displayed so the side conditions\textquotesingle{} meaning is on
the page.

\declhead{BernsteinLogic.Core.Catalogue}{premiseContains}{\badge{BadgeGreen}{definition bound}}
\begin{leansrc}
/-- Conjunct-wise premise containment (geolog `premiseContains`): every
conjunct of `sub` is a conjunct of `pre`. Decouples component premises
from the composite premise — the core has no ∧-restructuring rules, so
without this a component proved over its own box could not enter a
composite whose premise merely CONTAINS that box. -/
def premiseContains (pre sub : Formula) : Prop :=
  ∀ a ∈ sub.conjunctList, a ∈ pre.conjunctList
\end{leansrc}
\declhead{BernsteinLogic.Core.Catalogue}{formulaWiresWithin}{\badge{BadgeGreen}{definition bound}}
\begin{leansrc}
/-- Wire confinement (geolog `formulaWiresWithin`): every free wire of
the formula lies in the declared interface. Prover-side warrant that
demanding this costs no completeness: Lemma A, `farkas_interpolant_vars`
(contracts are mined confined). -/
def formulaWiresWithin (f : Formula) (iface : List Wire) : Prop :=
  ∀ w ∈ f.freeWires, w ∈ iface
\end{leansrc}
\declhead{BernsteinLogic.Core.Catalogue}{wiresSeparated}{\badge{BadgeGreen}{definition bound}}
\begin{leansrc}
/-- Separation (geolog `wiresSeparated`, ctx-form): every wire the
supplier's premise shares with the consumer's context or the conclusion
routes through the interface. The context formula is the consumer's
premise with the contract slot blanked (`conjReplace?` at `.tt`). -/
def wiresSeparated (left ctx post : Formula) (iface : List Wire) : Prop :=
  ∀ w ∈ left.freeWires,
    (w ∈ ctx.freeWires ∨ w ∈ post.freeWires) → w ∈ iface
\end{leansrc}
\ruleblock{Correspondence (Steele display $\leftrightarrow$ Lean)}
\badge{BadgeGreen}{reviewer-certified against the displayed texts}
\sourcenote{The dictionary below is a reviewer-certified statement of how the mechanized statements above interpret the displayed rule: display metavariable $\mapsto$ Lean binder, each display hypothesis $\mapsto$ the premise or side condition carrying it, the display conclusion $\mapsto$ the goal --- with any deliberate difference in width between display and mechanization named outright. Certification is recorded as a hash triple over (the displayed cell, the extracted statement texts, this dictionary\textquotesingle s own text); if any of the three changes, this block turns amber until a reviewer re-certifies. The correspondence claim itself remains editorial --- reviewer-grade, not machine-checked; its machine-checked successor types each soundness statement against the rule as a structured object, and supersedes this block per rule as it lands.}
Judgment frame as in \rulename{Ax}; the statement bundles its data as
\texttt{PromiseStep} (field-for-field the payload minus the two
sub-proofs, which enter as \texttt{ihs}/\texttt{ihc}).
\begin{itemize}
\item $\Phi_1 \vdash Q$ $\mapsto$ \texttt{ihs} with
  \texttt{supplierPost} $:$ \texttt{supplier.post = contract}; $Q
  \mapsto$ \texttt{contract : Formula}; $\Phi_1 \mapsto$
  \texttt{supplier.pre}.
\item $\Phi_2 \wedge Q \vdash \Psi$ $\mapsto$ \texttt{ihc}, with the
  certified statement taking the path-general form: $Q$ sits at a
  \emph{payload-declared conjunct path} (\texttt{pathReads} $:$
  \texttt{consumer.pre.conjGet? path = some contract}), not at a fixed
  position --- the ctx-$\vee$E style. $\Phi_2 \mapsto$ \texttt{ctx},
  the consumer premise with the contract deleted at the path
  (\texttt{ctxReads}).
\item conclusion $\Phi_1 \wedge \Phi_2 \vdash \Psi$ $\mapsto$ the
  goal for \texttt{st.seq}, whose premise is pinned by
  \texttt{prePlumb}: the consumer\textquotesingle s premise with the
  supplier\textquotesingle s premise spliced in \emph{at the same
  path} where $Q$ sat. Order honesty: at the displayed root-right
  instance ($Q$ the right conjunct of $\Phi_2 \wedge Q$), the splice
  therefore pins \texttt{seq.pre} to $\Phi_2 \wedge \Phi_1$ --- the
  displayed $\Phi_1 \wedge \Phi_2$ is obtained by semantic
  conjunction commutativity, not by direct instantiation. Path
  generality remains strictly wider than the display, in the safe
  direction; the conjunct \emph{order} is the one representational
  difference, named here. \texttt{postEq} carries $\Psi$ across
  unchanged.
\item the two confinement premises ($\wires(Q) \subseteq Y$; the
  interface-separation check) $\mapsto$ the Step fields
  \texttt{contractConfined} and \texttt{separation} over
  \texttt{iface} $\mapsto Y$ --- carried, checked,
  and consumed by \emph{no} step of \texttt{sound\_promise}: the
  displayed caption\textquotesingle s claim that they are the
  locality warrant rather than soundness needs is visible in the
  statement itself: the proof never touches them, as on the QTSL
  pages.
\end{itemize}
Width note: beyond the deliberate path-generality and the
conclusion-order transposition named above (the splice yields
$\Phi_2 \wedge \Phi_1$; the display\textquotesingle s order needs
conjunction commutativity, semantically free), neither side is
wider; the soundness content is pure predicate-logic cut over the
semantics, exactly as the witness column says.

\ruleblock{Lean verifier receipts}
\sourcenote{module \texttt{BernsteinLogic.Soundness}; bound build \texttt{a20d90e66c13\ldots} (bound from receipt a9c0e34d3e3c9c8c, 2026-08-03T21:30:36Z). A receipt records that named checking layers accepted a build: (module, compiled-artifact SHA-256, environment SHA-256, checker layers), minted when the checking gate ran; the store receipts attesting the bound build are listed with the layers each run credited. Nothing here asks the reader to trust the store: every layer is an independently runnable program over the quoted sources --- rebuild the module and re-run the checkers to reproduce the verdict. Receipts attesting superseded builds credit no layer and are omitted from this rendering (the count of omissions, if any, is noted below).}
\begin{receipttable}
\receiptrow{a9c0e34d3e3c9c8c}{2026-08-03T21:30:36Z}{build+decorrelated}{\badge{BadgeGreen}{current}}
\end{receipttable}
\sourcenote{21 further receipt(s) in the store attest superseded builds of this module and are omitted from this rendering.}
\sourcenote{module \texttt{BernsteinLogic.Core.Catalogue}; bound build \texttt{cbc92e976a1d\ldots} (bound from receipt 5abdb3589707b5df, 2026-08-03T21:30:33Z). A receipt records that named checking layers accepted a build: (module, compiled-artifact SHA-256, environment SHA-256, checker layers), minted when the checking gate ran; the store receipts attesting the bound build are listed with the layers each run credited. Nothing here asks the reader to trust the store: every layer is an independently runnable program over the quoted sources --- rebuild the module and re-run the checkers to reproduce the verdict. Receipts attesting superseded builds credit no layer and are omitted from this rendering (the count of omissions, if any, is noted below).}
\begin{receipttable}
\receiptrow{24804e66edba0bc9}{2026-08-03T20:25:09Z}{build+decorrelated}{\badge{BadgeGreen}{current}}
\receiptrow{5abdb3589707b5df}{2026-08-03T21:30:33Z}{build+decorrelated}{\badge{BadgeGreen}{current}}
\end{receipttable}
\sourcenote{18 further receipt(s) in the store attest superseded builds of this module and are omitted from this rendering.}
\layerline{\layeritem{build}{lake build (elaboration)}{\badge{BadgeGreen}{current}}\layeritem{replay}{leanchecker (official kernel replay)}{\badge{BadgeAmber}{stale only}}\layeritem{hardened}{safe\_verify (hardened screen)}{\badge{BadgeAmber}{stale only}}\layeritem{decorrelated}{lean4export + nanoda (decorrelated kernel)}{\badge{BadgeGreen}{current}}}
\sourcenote{\badge{BadgeGray}{comparator: none} no fixture in the three-implementation comparator harness (compiled / kernel-interpreted / reference-interpreted) exercises this rule.}
\ruleblock{HOL4 statement}
\sourcenote{theory \texttt{bern\_fam\_resid}, HOL4 development script \texttt{\texttt{bern\_fam\_residScript.sml}} \badge{BadgeGreen}{script bound}}
\holhead{Definition checkFam\_promise\_def}
\begin{holsrc}
Definition checkFam_promise_def:
  checkFam_promise (cols:columns) (ct:ctx) (i:num) ⇔
    case frSeq cols ct i of
    | NONE => F
    | SOME sq =>
      case frChildSeq cols ct i 0 of
      | NONE => F
      | SOME ss =>
        case frChildSeq cols ct i 1 of
        | NONE => F
        | SOME cs =>
          case frFormula cols ct i 0 of
          | NONE => F
          | SOME ctr =>
            case pathOf cols ct i of
            | NONE => F
            | SOME (stps, ifc) =>
              let pth = stepsToPath stps in
                ss.post = ctr ∧
                conjGet cs.pre pth = SOME ctr ∧
                conjReplace cs.pre pth ss.pre = SOME sq.pre ∧
                sq.post = cs.post ∧
                (case conjReplace cs.pre pth FTt of
                 | SOME cx0 =>
                     (EVERY (λw. MEM w ifc) (formulaFreeWires ctr) ∧
                      EVERY (λw. (MEM w (formulaFreeWires cx0) ∨
                                  MEM w (formulaFreeWires sq.post)) ⇒
                                 MEM w ifc)
                            (formulaFreeWires ss.pre))
                 | NONE => F)
End
\end{holsrc}
\holhead{Theorem famEq\_promise}
\begin{holsrc}
Theorem famEq_promise:
  wfColumns cols ct ∧ mkCtx cols = SOME ct ∧
  parseCert (cols.hdrImage ++ rest) = SOME (ws, pf) ∧
  nodeCorr cols ct i (CertPf sq nd) ∧
  FST (THE (tagOf cols i)) = 25 ⇒
  (checkFam_promise cols ct i ⇔ armPromise ws sq nd)
\end{holsrc}
\ruleblock{HOL verifier receipts}
\badge{BadgeBlue}{documentary} Documentary citations. Family validation: \texttt{bern\_fam\_resid} built GREEN with all ten \texttt{famEq} arms validated, \texttt{hol-strict} clean; spine census: \texttt{bern\_joins} GREEN under \texttt{hol-strict} (55 theorems, 0 offences, 208-theory ancestry including this theory), \texttt{checkJoins\_eq\_checkPf} verified verbatim in the built signature. The cited \texttt{famEq} lemma is the kernel-checked per-node equivalence: under well-formed columns (\texttt{wfColumns}/\texttt{mkCtx}/\texttt{nodeCorr}) at this rule\textquotesingle s wire tag, the columnar family check coincides with the flat reference walk\textquotesingle s arm for this rule. There is no signed HOL receipt store yet --- these citations are documentary.
Semantic stratum, also ported: the soundness twin
\texttt{sound\_promise} (theory \texttt{bern\_tr\_b14}: \texttt{hol-
strict} clean, 2 theorems, 0 offences, verification run of 2026-07-22 on
an isolated build host, content-addressed at the same sources), consumed
by the certified-fragment keystone (\texttt{checkPf\_sound\_frag},
theory \texttt{bern\_tr\_splice}, 20 theorems, 0 offences, same run): a
reference-checker accept inside the proven wire fragment yields sequent
validity.

\sourcenote{HOL receipts are documentary: quoted results of \texttt{hol-strict} runs
and census rows from the HOL4 development\textquotesingle s build
evidence, not minted receipt artefacts. A HOL receipt store (the Lean
store\textquotesingle s twin) is named future work.
}
\ruleblock{Proof}

\sourcenote{Genre: admissible rule (tier 2) --- a proof-tree-rewriting
recipe. Proof-theoretically \rulename{Promise} unpacks into core rules
(conjunct plumbing + \rulename{Cut}); in the realized calculus it is a
native constructor, kept for cost and for the locality-warrant payload.
The recipe shows the unpack; the semantic warrant is the kernel lemma
\texttt{sound\_promise}, which is the unpack's semantic shadow run in
one pass.}

\textbf{Setting.}\enspace The payload (the \texttt{PromiseStep}
structure quoted above) carries a supplier sequent $\Phi_1 \vdash Q$, a
consumer sequent whose premise contains the contract $Q$ at a declared
conjunct path $\pi$ (write it $\Phi_2 \wedge_{\pi} Q \vdash \Psi$ ---
the $\pi$-notation meaning $Q$ sits at path $\pi$ in the premise's
conjunction tree, $\Phi_2$ being the rest), a declared interface wire
list $Y$, and the composite sequent, whose premise is the consumer's
with $Q$ \emph{replaced} by $\Phi_1$ at $\pi$ and whose conclusion is
$\Psi$. Four checked hypotheses: \emph{(P1)} the path reads $Q$ in the
consumer's premise; \emph{(P2)} replacing $Q$ by $\Phi_1$ at $\pi$
yields exactly the composite premise; \emph{(P3)}
$\wires(Q) \subseteq Y$; \emph{(P4)} the supplier's wires meet the
consumer's and the goal's only inside $Y$. (P1)--(P2) are the splice
plumbing; (P3)--(P4) are the locality warrant.

\begin{brecipe}[the \rulename{Promise} instance --- cut through a
conjunct path]\label{rec:promise-instance}
The general admissible template (one rule node $\rightsquigarrow$ a
derivation fragment over core rules with the same conclusion, applied
leaf-first until no such node remains) instantiates as:
\[
\inferrule*[Right=\rulename{Promise}]{\;\mathcal{T}_s \rhd \Phi_1 \vdash Q \\ \mathcal{T}_c \rhd \Phi_2 \wedge_{\pi} Q \vdash \Psi\;}{\Phi_2 \wedge_{\pi} \Phi_1 \vdash \Psi}
\qquad\rightsquigarrow\qquad
\inferrule*[Right=\rulename{Cut}]{
  \inferrule*{\mathcal{T}_s}{\Phi_2 \wedge_{\pi} \Phi_1 \vdash Q}
  \\
  \inferrule*{\mathcal{T}_c}{\Phi_2 \wedge_{\pi} Q \vdash \Psi}
}{\Phi_2 \wedge_{\pi} \Phi_1 \vdash \Psi}
\]
where, writing $C$ for the composite premise
$\Phi_2 \wedge_{\pi} \Phi_1$: the left branch derives $C \vdash Q$ by
\rulename{Cut} from $C \vdash \Phi_1$ (an \rulename{Ax} leaf
$C \vdash C$, then a conclusion-side $\wedge$E chain along $\pi$)
against a \emph{premise-transported} copy of $\mathcal{T}_s$, and
\rulename{Cut}'s right premise proper is $C \wedge Q \vdash \Psi$, a
premise-transported copy of $\mathcal{T}_c$ --- the display elides both
transports. Premise transport (re-rooting a derivation $X \vdash \Psi$
at a premise $Y$ whose addressed conjuncts entail $X$) is \emph{not} a
local chain in this calculus: \rulename{Weaken} and \rulename{Cut} are
deliberately head-positional --- each grafts one conjunct at the
premise root's right --- so the premises reachable from an opaque
sub-proof $X \vdash \Psi$ by insertion alone are exactly $X$ with
right-appended conjuncts (\rulename{Cut} can also \emph{consume} an
already-present root-right conjunct when a proof of it is supplied ---
removal, never interior insertion or permutation), and even
$Y \wedge X \vdash \Psi$ is locally out of reach.
The transport is instead an induction over the sub-derivation itself;
its principal cases: rules that copy their premise tree verbatim from
their sub-proofs re-instantiate on the transported children;
premise-free leaves conclude at an arbitrary premise by the
\rulename{Norm} convention (an \rulename{Ax} leaf $\Phi \vdash \Phi$
re-derives at $Y$ as $Y \vdash Y$ plus a conclusion-side
$\wedge$E/$\wedge$I rebuild); and the two head-positional rules
are their own inductive cases, because re-rooting commutes with the
head graft: a transported \rulename{Weaken} grafts its conjunct onto
the new root, and a transported \rulename{Cut}'s right premise
re-forms as the new root with the cut formula appended. This list is
not exhaustive over the catalogue: the premise-\emph{forming} rules
($\vee$E's conclusion premise is the disjunction of its branch
premises; \rulename{Subdivide}'s branch premises conjoin the region
data) and the context-carrying rules (\rulename{Promise},
\rulename{Syzygy}, whose path plumbing must compose with the splice
--- the $\vee$E section's machinery) are neither verbatim copies nor
leaves; each needs its own commutation case, and those cases belong to
the same transport metatheorem flagged as unmechanized below.
Rewriting the premise field node by node under that case analysis
yields the transported copies of $\mathcal{T}_s$ and $\mathcal{T}_c$. Elimination costs
the sizes of $\mathcal{T}_s$ and $\mathcal{T}_c$ plus the
$\wedge$E chain along $\pi$ at the auxiliary derivation plus the
per-leaf \rulename{Ax} rebuilds --- and each transported leaf pays its
own conclusion-side rebuild chain (the proof-object grammar has no
sharing), so the total is bounded by the certificate sizes plus the
number of transported leaves times the addressing depth, linear in the
certificate data only when that depth is bounded ---
and consumes \rulename{Ax} leaves throughout; conservativity holds
with no new kernel semantic content, but as a global rewrite warranted
by the transport induction (stated here as prose --- the transport
metatheorem is not yet mechanized), not as insertion of a
constant-size fragment.
\end{brecipe}

\begin{bprop}[semantic warrant]\label{prop:promise-sound}
If the supplier and consumer sequents are semantically valid and
(P1)--(P2) hold, the composite sequent is semantically valid.
\end{bprop}
\begin{proof}
Fix an admissible valuation $\sigma$ satisfying the composite premise.
Reading the composite premise at $\pi$ returns $\Phi_1$: (P1) says
$\pi$ addresses $Q$ in the consumer's premise, so the replacement (P2)
performs succeeded at $\pi$, and the read of a successful replacement
returns the replacement (the round-trip lemma
\texttt{Formula.conjGet\_of\_conjReplace}); an addressed conjunct
of a holding conjunction holds (\texttt{Formula.holds\_of\_conjGet});
so $\sigma \models \Phi_1$, and supplier validity gives
$\sigma \models Q$. Replacing $\Phi_1$ back by $Q$ at $\pi$ recovers
the consumer's premise exactly
(\texttt{Formula.conjReplace\_back}, using (P1)--(P2)), and replacing
the addressed conjunct of a holding conjunction by a formula that
itself holds at $\sigma$ preserves holding
(\texttt{Formula.holds\_conjReplace}, applied at $\pi$ with $Q$ for
$\Phi_1$); so
$\sigma \models \Phi_2 \wedge_{\pi} Q$, and consumer validity gives
$\sigma \models \Psi$. This is \texttt{sound\_promise} verbatim: pure
predicate logic over the semantics --- admissibility of $\sigma$ is
passed through untouched, and neither confinement hypothesis is
consulted.
\end{proof}

\begin{bremark}[the confinement premises are warrant, not soundness]
\label{rem:promise-warrant}
(P3)--(P4) appear in no step of Proposition~\ref{prop:promise-sound},
and dropping them alone --- (P1)--(P2) retained --- admits no false
sequent: cut is sound regardless of where the contract's wires live. What their removal \emph{does} admit
is a degenerate reading of the rule --- a ``contract'' mentioning
supplier-private wires, i.e.\ a composition whose interface is as wide
as the components themselves. Everything the rule is \emph{for} then
fails: the contract stops being a datasheet (a public artefact whose
vocabulary is the declared separator), certificate size stops scaling
with the interface, and any bounded-disclosure reading of ``what the
consumer learns about the supplier'' becomes vacuous. Those
consequences are theorems proved elsewhere, \emph{about certificates
that carry the checks} --- they quantify over exactly the certificates
these hypotheses admit --- which is why the checks are checked --- a hypothesis whose
necessity is economic and epistemic rather than semantic is exactly
the kind a recipe must label, or the next reader deletes it as
redundant.
\end{bremark}

\begin{bremark}[relation to the literature]
Assume--guarantee composition descends from composing-specifications
calculi \cite{abadilamport1993}; the subtlety those calculi negotiate
is \emph{circularity} --- each component assuming what the other
guarantees --- which needs an induction over time or a fixpoint
argument \cite{mcmillan1999}. \rulename{Promise} deliberately takes
the acyclic case, where plain cut suffices and the whole rule is
bookkeeping around it; the circular case enters this calculus through
the feedback rules --- \rulename{SmallGain} and \rulename{SmallGainN}
at the sequent layer and the feedback cells of the satisfaction-layer
composition rule, where the guarded-fixpoint structure is explicit,
and \rulename{QWireComposeMarks} for liveness marks, where cyclic
wiring needs no well-foundedness device at all and the wiring
condition is a feasibility walk (schedules and pinned values must
close around cycles) rather than a causality walk. Separating the
cases --- cut where cut suffices, guarded fixpoints where feedback
threatens a safety face, a schedule walk where a cycle only has to
deliver a mark --- is the design's answer to the classical observation
that \emph{unguarded} circular assume--guarantee reasoning about
safety is unsound --- there, some well-foundedness device is required.
\end{bremark}

\begin{anchors}
Recipe \ref{rec:promise-instance} payload \& plumbing & \texttt{PromiseStep} (fields \texttt{path}, \texttt{pathReads}, \texttt{prePlumb}, \texttt{postEq})\\
path read/replace round trips & \texttt{Formula.conjGet\_of\_conjReplace}, \texttt{Formula.conjReplace\_back}\\
addressed conjunct entailed; replacement preserves holding & \texttt{Formula.holds\_of\_conjGet}, \texttt{Formula.holds\_conjReplace}\\
Proposition \ref{prop:promise-sound} & \texttt{sound\_promise}; wired as core \texttt{RuleData.promise} into \texttt{ProofObject.sound}\\
confinement checks (warrant only) & \texttt{formulaWiresWithin}, \texttt{wiresSeparated} (payload fields \texttt{contractConfined}, \texttt{separation})\\
\end{anchors}


\rulesection{33}{Syzygy}{core rule}
\label{blcatrule-Syzygy}
\ruleblock{The rule (Steele metanotation)}
\sourcenote{as realized in the mechanized development; source row: \texttt{\texttt{rules.toml}} --- \emph{Syzygy}. The metanotation statement is typeset by hand for notational quality; its mechanized ground truth is the Lean statement below, extracted from the sources at every build of this book and hash-bound. The rule \emph{roster} is machine-checked (every mechanized rule tag round-trips to exactly one catalogue row); the statement \emph{text} is editorial, checked by eye against the extracted Lean beside it.}
\begin{mathpar}
\inferrule*[Right=\rulename{Syzygy}]{
  \bigl\{\Phi_i \vdash e_i \leq 0\bigr\}_{i=1}^{n} \\
  \bigl\{g_k = 0\bigr\}_k \ \text{equality sub-conjuncts of } \Phi \text{ in wire variables} \\
  \lambda_i \in \Q_{\geq 0} \\
  \nu_k \in \Q \\
  c \equiv \textstyle\sum_i \lambda_i e_i + \sum_k \nu_k g_k \ \text{(exact canonical polynomial)} \\
  \Psi \in \bigl\{\, c \leq 0, \;\; e \geq 0 \text{ where } c = -e, \;\; l \leq r \text{ where } c = l - r \,\bigr\}
}{
  \Phi \vdash \Psi
}
\end{mathpar}
$\Phi \supseteq$ each $\Phi_i$ conjunct-wise (checked); the $\nu_k$ are free-sign exactly as \rulename{Farkas} frees $\lambda_j$ on $=$ atoms.

\ruleflavor{free-sign $\nu_k$ is Lemma A\textquotesingle s CommRing sign-blindness, one theorem for both dialects. Lifts \rulename{Farkas} in exactly two ways: atoms from sub-\emph{proofs} (not sub-conjuncts), atoms of arbitrary polynomial degree. The conclusion is exposed in any of the three \texttt{syzygyConclusionForm} readings --- one exact $\Q$ identity under a sign flip; the transport is the ge-bridge ring lemma, applied once inside \texttt{sound\_syzygy}. A Bernstein Logic rule outright: the baked base theory lists it core (its wire node decoded by the certified checker, its combination relation pinned); \textquotedblleft assume--guarantee\textquotedblright{} names its purpose, not a separate layer.}
\par\medskip
\rulename{Syzygy} is the calculus\textquotesingle s composition rule: it
combines already-proven component bounds into a bound on a weighted sum.
Each component sequent $\Phi_i \vdash e_i \leq 0$ is a sub-proof ---
typically ``this subsystem dissipates its storage function'' --- and the
composite premise $\Phi$ carries, among its conjuncts,
\emph{equality atoms} $g_k = 0$ describing how the subsystems are wired
together (power-preserving interconnection: what one port supplies,
another consumes). A producer reaches for the rule wherever a
system-level argument says ``add the component certificates with weights
$\lambda_i$ and let the internal supplies
cancel'': the weights are found by whatever
method the producer likes (an LP, typically), and only the resulting
exact identity ships in the certificate.

What the checker recomputes, given the payload (combination term,
component list, conclusion term): that every entry weight is
${}\geq 0$ --- premise \emph{equality} atoms expose both polarities to
the matcher, so a free-sign multiplier on an equality is realized as a
nonnegative weight on the appropriate polarity, and that freedom is what
makes them syzygies; that every entry
of the combination really is a component conclusion or a premise equality
sub-conjunct, matched by exact canonical polynomial; and that the
weighted sum of those polynomials is \emph{identically} the conclusion
polynomial as canonical forms over $\Q$ --- coefficient-for-coefficient
equality, not evaluation at test points. It further checks the
bookkeeping hypotheses: each component\textquotesingle s premise is
contained conjunct-wise in $\Phi$ (\texttt{preContains}), the conclusion
term has a determinate finite shape (\texttt{conclShape},
\texttt{conclFin}), and every wire in sight has positive size
(\texttt{wiresPos}, \texttt{compPostWiresPos}).

Each hypothesis is load-bearing. Drop the sign condition on the
inequality weights and a producer could multiply a bound by $-1$ and
conclude its negation. Drop the entry-sourcing check and the
``combination'' could smuggle in atoms
nobody proved. Drop the exact-identity check and the conclusion would be
unrelated to the premises. Drop premise containment and a component
proved over its own operating box could be applied outside it. Drop the
positive-size conditions and a component conclusion over a zero-size wire
--- vacuously true, since it quantifies over no tensor entries --- could
enter the combination with full algebraic weight and push the composite
to a falsehood. The proof section walks these in order.

\cataloguenote{Port-Hamiltonian reading: with $e_i = \mathcal{L}_i V_i -
s_i^{\mathrm{ext}} - s_i^{\mathrm{int}}$, the equality atoms $g_k = 0$
cancel the internal supplies and the conclusion is a dissipation
inequality for $\sum_i \lambda_i V_i$ --- Tellegen\textquotesingle s
theorem in certificate form. The free-body diagram is the junction tree;
internal forces at the cuts are the separator algebra.
}
\ruleblock{Lean statement}
\declhead{BernsteinLogic.Soundness}{SyzygyStep}{\badge{BadgeGreen}{statement bound}}
\begin{leansrc}
/-- Applicability data of the Syzygy rule over a hypernet — field-for-field
the reference checker's obligations, and field-for-field `PSyzygyData`
minus the component-proof list (validity hypotheses of `sound_syzygy`)
and `denBoundLog2` (the Q7 family; zero proof-theoretic weight in v1).
B13-shaped conclusion fields and B14 (`wiresPos`) are inherited verbatim
from `PFarkasData`; `compPostWiresPos` is the honest EXTENSION of B14 to
component-conclusion wires (Composition.lean header: without it a
vacuously-held component conclusion over a zero-size wire poisons the
combination — counterexample kernel-checkable). -/
structure SyzygyStep (H : Hypernet) where
  seq : Sequent
  components : List Sequent
  combination : Term
  conclusion : Term
  preContains : ∀ c ∈ components, premiseContains seq.pre c.pre
  conclShape : ∃ s, conclusion.bshape H = some s
  conclFin : conclusion.finWires H
  wiresPos : ∀ w ∈ seq.pre.freeWires ++ conclusion.freeWires,
      ∃ τ, H.wireType w = some τ ∧ 0 < Shape.totalSize τ.shape
  compPostWiresPos : ∀ c ∈ components, ∀ w ∈ c.post.freeWires,
      ∃ τ, H.wireType w = some τ ∧ 0 < Shape.totalSize τ.shape
  valid : syzygyCombinationValid combination components seq.pre conclusion
  postForm : syzygyConclusionForm conclusion seq.post
\end{leansrc}
\declhead{BernsteinLogic.Soundness}{holds\_of\_syzygyConclusionForm}{\badge{BadgeGreen}{statement bound}}
\begin{leansrc}
/-- The ge-bridge ring lemma (the both-polarity core change): every
reading admitted by
`syzygyConclusionForm conclusion` holds whenever the zero-normal atom
`conclusion ≤ 0` holds — the readings are the same exact ℚ identity
re-presented under a sign flip. The zero-form is the identity transport,
the ge-polarity flip is `-e ≤ 0 → e ≥ 0`, and the literal atom is
`l - r ≤ 0 → l ≤ r` (ArithAtom's tensor-general sub-zero bridge).
Applied once inside `sound_syzygy`, so PSyzygy (tier-1), PBound's
compact spelling (tier-3), and any future ≤-0-shaped fact rule inherit
all readings with no separate polarity justification. -/
theorem holds_of_syzygyConclusionForm {H : Hypernet} {σ : H.Valuation}
    {conclusion : Term} {F : Formula}
    (hle : Formula.holds σ (.atom .le conclusion (.const 0)))
    (hform : syzygyConclusionForm conclusion F) :
    Formula.holds σ F
\end{leansrc}
\declhead{BernsteinLogic.Soundness}{sound\_syzygy}{\badge{BadgeGreen}{statement bound}}
\begin{leansrc}
/-- Soundness of Syzygy (= Lemma B, dialect D): valid components give a
valid composite. The argument is `sound_farkas`'s, with the per-entry
source discharge generalized: a component-sourced entry's atom holds by
component validity (through `premiseContains`), a wiring-sourced entry's
atom holds as a premise sub-conjunct; either way
`atomLeZeroPolys_sound_entryC` delivers per-entry nonpositivity at every
clamped per-entry valuation, and the nonnegative fold closes. The
zero-normal bound then reaches `seq.post` through whichever
`syzygyConclusionForm` reading the payload declared (both-polarity core
change, 2026-07-18). -/
theorem sound_syzygy {H : Hypernet} (st : SyzygyStep H)
    (hcomps : ∀ c ∈ st.components, c.semanticallyValid H) :
    st.seq.semanticallyValid H
\end{leansrc}
\declhead{BernsteinLogic.Soundness}{atomLeZeroPolys\_sound\_entryC}{\badge{BadgeGreen}{statement bound}}
\begin{leansrc}
/-- Every `(e ≤ 0)`-normal-form polynomial a held atom entails evaluates
nonpositively at EVERY clamped per-entry valuation, given positive-size
wire types (general shapes — heterogeneous atoms read through the
clamp; the positivity is PROPOSED-B14's honest content: zero-size atoms
hold vacuously and entail nothing). -/
theorem atomLeZeroPolys_sound_entryC {H : Hypernet} (σ : H.Valuation)
    {F : Formula} (hF : Formula.holds σ F)
    (hw : ∀ w ∈ F.freeWires, ∃ τ, H.wireType w = some τ ∧
      0 < Shape.totalSize τ.shape)
    (hix : F.atomSidesIdxFree = true) (i : ℕ) :
    ∀ pa ∈ atomLeZeroPolys F, QArith.Poly.eval (σ.entryValC i) pa ≤ 0
\end{leansrc}
\declhead{BernsteinLogic.Soundness}{wire\_discipline\_of\_bshape\_fin}{\badge{BadgeGreen}{statement bound}}
\begin{leansrc}
/-- The PROPOSED-B13 field pair (`bshape` + `finWires`) yields the wire
discipline that per-entry evaluation consumes: every free wire reads a
finite tensor of shape `[]` or `s`. -/
theorem wire_discipline_of_bshape_fin {H : Hypernet} (σ : H.Valuation)
    {t : Term} {s : Shape} (hb : t.bshape H = some s)
    (hf : t.finWires H) (hix : t.idxFree = true) :
    ∀ w ∈ t.freeWires, ∃ qw, σ.wireQ w = some qw ∧
      (qw.shape = [] ∨ qw.shape = s)
\end{leansrc}
\declhead{BernsteinLogic.Soundness}{Term.evalQ\_total\_singleton}{\badge{BadgeGreen}{statement bound}}
\begin{leansrc}
/-- Total singleton evaluation: if every free wire reads as a finite
singleton tensor, the term evaluates to a singleton tensor whose scalar
is the polynomial evaluation at the valuation read-out. (Definedness half
of the PROPOSED-B13 Farkas fields, and the workhorse behind every scalar
extraction below.) -/
theorem Term.evalQ_total_singleton {H : Hypernet} (σ : H.Valuation) :
    ∀ (t : Term),
    (∀ w ∈ t.freeWires, ∃ qw, σ.wireQ w = some qw ∧ SingletonShape qw.shape) →
    t.idxFree = true →
    ∃ q, Term.evalQ σ t = some q ∧ SingletonShape q.shape ∧
      q.scalar = QArith.Poly.eval σ.rho t.toPoly
\end{leansrc}
\par\medskip
\textbf{What \texttt{syzygyCombinationValid} means --- and what it does
not.}\enspace The rule\textquotesingle s side condition appears in the
payload structure above as a single opaque name, so this book binds its
meaning by displaying the definition chain in full, from the checker
sources. Reading the chain: \texttt{syzygyCombinationValid} is precisely
the proposition ``\texttt{syzygyCheck} returns
\texttt{true}'', and \texttt{syzygyCheck}
is the disjunction of three branches of one algorithm (a branch for
index-free terms, a wire-coordinate branch for indexed terms, and a
mixed tagged-union branch for heterogeneous
combinations of both kinds). All three
branches do the same three things: decode the combination term into
weighted entries (\texttt{farkasEntries} --- a fixed canonical encoding,
anything else is rejected); require each weight nonnegative and each
entry\textquotesingle s polynomial to occur among the
$({\cdot} \leq 0)$-normal polynomials of a \emph{component conclusion} or
of an \emph{equality sub-conjunct of the premise}
(\texttt{atomLeZeroPolys} --- note that equality atoms contribute both
polarities, which is exactly the free-sign multiplier discipline); and
require the weighted fold of those polynomials to equal the conclusion
term\textquotesingle s canonical polynomial.

What a skeptical reader should confirm by eye in the definitions below:
the conclusion term enters the check \emph{only} as the target of the
exact-identity comparison (and through a syntactic index-freedom gate);
the step\textquotesingle s posted formula \texttt{seq.post} does not
appear in the check at all; and nothing in the chain evaluates, assumes,
or otherwise consults the \emph{truth} of the conclusion --- the check
manipulates polynomials, never valuations, so it could not covertly
``assert the conclusion is
true''. \texttt{seq.post} connects to the
conclusion term only through the separate, purely syntactic field
\texttt{postForm} --- the three readings of $c \leq 0$ enumerated by
\texttt{syzygyConclusionForm} below --- whose semantic transport is the
separately proven ring lemma \texttt{holds\_of\_syzygyConclusionForm}
quoted above. Semantic truth is established only in
\texttt{sound\_syzygy}, from the component-validity hypotheses.

\declhead{BernsteinLogic.Core.Catalogue}{syzygyConclusionForm}{\badge{BadgeGreen}{definition bound}}
\begin{leansrc}
def syzygyConclusionForm (conclusion : Term) (F : Formula) : Prop :=
  F = .atom .le conclusion (.const 0)
  ∨ (∃ e, conclusion = .neg e ∧ F = .atom .ge e (.const 0))
  ∨ (∃ l r, conclusion = .sub l r ∧ F = .atom .le l r)
\end{leansrc}
\declhead{BernsteinLogic.Core.Catalogue}{atomLeZeroPolys}{\badge{BadgeGreen}{definition bound}}
\begin{leansrc}
/-- The `(e ≤ 0)`-normal-form polynomials entailed by one premise atom
(exact polynomial forms — never rescaled; F7 discipline). -/
def atomLeZeroPolys : Formula → List QArith.Poly
  | .atom .le l r => [QArith.Poly.sub l.toPoly r.toPoly]
  | .atom .lt l r => [QArith.Poly.sub l.toPoly r.toPoly]
  | .atom .ge l r => [QArith.Poly.sub r.toPoly l.toPoly]
  | .atom .gt l r => [QArith.Poly.sub r.toPoly l.toPoly]
  | .atom .eq l r => [QArith.Poly.sub l.toPoly r.toPoly,
                      QArith.Poly.sub r.toPoly l.toPoly]
  | _ => []
\end{leansrc}
\declhead{BernsteinLogic.Core.Catalogue}{farkasEntries}{\badge{BadgeGreen}{definition bound}}
\begin{leansrc}
/-- Decode the combination-coefficient Term (pinned canonical encoding:
right-nested `.add` chain of `.mul (.const λⱼ) eⱼ`, `.const 0` end). -/
def farkasEntries : Term → Option (List (Rat × Term))
  | .const c => if c = 0 then some [] else none
  | .add (.mul (.const l) e) rest =>
    (farkasEntries rest).map ((l, e) :: ·)
  | _ => none
\end{leansrc}
\declhead{BernsteinLogic.Core.Catalogue}{syzygyCheckCore}{\badge{BadgeGreen}{definition bound}}
\begin{leansrc}
/-- The Syzygy combination check — `farkasCheck`'s shape, lifted the
AG draft's two ways: entry atoms may come from COMPONENT CONCLUSIONS
(the sub-proofs) as well as from equality sub-conjuncts of the composite
premise; atom polynomials are unrestricted in degree. Entry encoding,
fold, exactness discipline, and the G1 idx-freedom gates are inherited
verbatim from `farkasCheck`. -/
def syzygyCheckCore (cc : Term) (components : List Sequent) (premise : Formula)
    (conclusion : Term) : Bool :=
  match farkasEntries cc with
  | none => false
  | some entries =>
    conclusion.idxFree &&
    entries.all (fun (l, e) =>
      decide (0 ≤ l) &&
      ((components.any fun c =>
          c.post.atomSidesIdxFree &&
          ((atomLeZeroPolys c.post).any fun pa => decide (pa = e.toPoly))) ||
       (premise.conjunctList.any fun a =>
          a.isEqAtom && a.atomSidesIdxFree &&
          ((atomLeZeroPolys a).any fun pa => decide (pa = e.toPoly))))) &&
    decide ((entries.foldr
        (fun (l, e) acc => QArith.Poly.add (QArith.Poly.scale l e.toPoly) acc)
        []) = conclusion.toPoly)
\end{leansrc}
\declhead{BernsteinLogic.Core.Catalogue}{syzygyCheckW}{\badge{BadgeGreen}{definition bound}}
\begin{leansrc}
/-- The idx-form branch of the Syzygy check: entries
match component conclusions or premise equality sub-conjuncts by their
wire-aware coordinate polynomials — the multi-wire composites (DP
Syzygy/SmallGain, distinct region wires per component) are exactly why
`wireCoordKey` keys variables by (wire, coordinate). -/
def syzygyCheckW (cc : Term) (components : List Sequent) (premise : Formula)
    (conclusion : Term) : Bool :=
  match farkasEntries cc with
  | none => false
  | some entries =>
    match farkasEntriesW entries with
    | none => false
    | some pentries =>
      match conclusion.asCoordPolyIdxW? with
      | none => false
      | some pc =>
        pentries.all (fun le =>
          decide (0 ≤ le.1) &&
          ((components.any fun c =>
              (atomLeZeroPolysW c.post).any fun pa => decide (pa = le.2)) ||
           (premise.conjunctList.any fun a =>
              a.isEqAtom &&
              ((atomLeZeroPolysW a).any fun pa => decide (pa = le.2))))) &&
        decide ((pentries.foldr
            (fun le acc => QArith.Poly.add (QArith.Poly.scale le.1 le.2) acc)
            []) = pc)
\end{leansrc}
\declhead{BernsteinLogic.Core.Catalogue}{syzygyCheckM}{\badge{BadgeGreen}{definition bound}}
\begin{leansrc}
/-- The mixed-variable branch of the Syzygy check: entries match
component conclusions or premise equality sub-conjuncts by their MIXED
tagged polynomials — heterogeneous composites (a chart wire's TIdx
coordinates and bare scalar interface wires in one combination) are
exactly why the tagged union space exists. -/
def syzygyCheckM (cc : Term) (components : List Sequent) (premise : Formula)
    (conclusion : Term) : Bool :=
  match farkasEntries cc with
  | none => false
  | some entries =>
    match farkasEntriesM entries with
    | none => false
    | some pentries =>
      match conclusion.asCoordPolyMixed? with
      | none => false
      | some pc =>
        pentries.all (fun le =>
          decide (0 ≤ le.1) &&
          ((components.any fun c =>
              (atomLeZeroPolysM c.post).any fun pa => decide (pa = le.2)) ||
           (premise.conjunctList.any fun a =>
              a.isEqAtom &&
              ((atomLeZeroPolysM a).any fun pa => decide (pa = le.2))))) &&
        decide ((pentries.foldr
            (fun le acc => QArith.Poly.add (QArith.Poly.scale le.1 le.2) acc)
            []) = pc)
\end{leansrc}
\declhead{BernsteinLogic.Core.Catalogue}{syzygyCheck}{\badge{BadgeGreen}{definition bound}}
\begin{leansrc}
/-- The Syzygy combination check: legacy idx-free branch OR the
wire-aware idx coordinate branch OR the mixed tagged branch (mirrors
`farkasCheck`'s three-way split). -/
def syzygyCheck (cc : Term) (components : List Sequent) (premise : Formula)
    (conclusion : Term) : Bool :=
  syzygyCheckCore cc components premise conclusion ||
  syzygyCheckW cc components premise conclusion ||
  syzygyCheckM cc components premise conclusion
\end{leansrc}
\declhead{BernsteinLogic.Core.Catalogue}{syzygyCombinationValid}{\badge{BadgeGreen}{definition bound}}
\begin{leansrc}
def syzygyCombinationValid (cc : Term) (components : List Sequent)
    (premise : Formula) (conclusion : Term) : Prop :=
  syzygyCheck cc components premise conclusion = true
\end{leansrc}
\ruleblock{Correspondence (Steele display $\leftrightarrow$ Lean)}
\badge{BadgeGreen}{reviewer-certified against the displayed texts}
\sourcenote{The dictionary below is a reviewer-certified statement of how the mechanized statements above interpret the displayed rule: display metavariable $\mapsto$ Lean binder, each display hypothesis $\mapsto$ the premise or side condition carrying it, the display conclusion $\mapsto$ the goal --- with any deliberate difference in width between display and mechanization named outright. Certification is recorded as a hash triple over (the displayed cell, the extracted statement texts, this dictionary\textquotesingle s own text); if any of the three changes, this block turns amber until a reviewer re-certifies. The correspondence claim itself remains editorial --- reviewer-grade, not machine-checked; its machine-checked successor types each soundness statement against the rule as a structured object, and supersedes this block per rule as it lands.}
Judgment frame as in \rulename{Ax}. The statement bundles its data as
\texttt{SyzygyStep} --- field-for-field the reference
checker\textquotesingle s obligations, and field-for-field the payload
minus the component-proof list (which enters \texttt{sound\_syzygy} as
the validity hypothesis \texttt{hcomps}) and minus the
denominator-bound declaration (resource accounting; zero
proof-theoretic weight, so rightly absent from the soundness
statement).
\begin{itemize}
\item the component judgments $\{\Phi_i \vdash e_i \leq 0\}_{i=1}^n$
  $\mapsto$ \texttt{st.components : List Sequent} with
  \texttt{hcomps}; the atoms $e_i \leq 0$ have no binders of their own
  --- each combination entry tagged \textsc{from-component} is checked
  against a component\textquotesingle s conclusion inside the
  combination check, which is where the $e_i$ live.
\item $\Phi \supseteq \Phi_i$ conjunct-wise $\mapsto$
  \texttt{st.preContains} $:$ \texttt{premiseContains seq.pre c.pre},
  per component.
\item the equality sub-conjuncts $\{g_k = 0\}$ of $\Phi$ $\mapsto$ the
  \textsc{from-wiring} entry tag: such an entry\textquotesingle s atom
  is checked to be a sub-conjunct of \texttt{seq.pre} directly.
\item $\lambda_i \in \Q_{\geq 0}$, $\nu_k \in \Q$, and $c \equiv
  \sum_i \lambda_i e_i + \sum_k \nu_k g_k$ $\mapsto$
  \texttt{st.combination} and \texttt{st.conclusion} under
  \texttt{st.valid} $:$ \texttt{syzygyCombinationValid} --- the
  canonical-polynomial identity is recomputed by one of \emph{three}
  certified routes (the idx-free \texttt{syzygyCheckCore} replay, the
  wire-coordinate \texttt{syzygyCheckW}, and the mixed
  \texttt{syzygyCheckM}), each discharged separately in
  \texttt{sound\_syzygy}.
\item the evaluability and positive-size packs $\mapsto$
  \texttt{conclShape}, \texttt{conclFin}, \texttt{wiresPos},
  \texttt{compPostWiresPos} --- displayed inline at
  \rulename{PBound}\textquotesingle s cell (this page\textquotesingle
  s cell leaves them to the caption); \texttt{compPostWiresPos} is
  the honest extension of the positive-size discipline to component
  conclusions, whose absence has a kernel-checkable
  vacuous-conclusion counterexample (the source docstring).
\item $\Psi$\textquotesingle s three readings $\mapsto$
  \texttt{st.postForm} $:$ \texttt{syzygyConclusionForm}, literally
  the displayed three-membered disjunction; the transport of the
  zero-normal bound to every reading is the certified-alongside
  ge-bridge lemma \texttt{holds\_of\_syzygyConclusionForm}, applied
  once here, inherited by \rulename{PBound} and the
  \rulename{PsdCert} trio.
\end{itemize}
Width note: nothing in the statement consumes the port-Hamiltonian
reading --- supply rates, storage functions, and the wiring identity
are producer-side interpretations of which the certified content is
exactly the polynomial identity and its component sourcing; and the
statement is proven at every hypernet and component list, with every
hypothesis consumed.

\ruleblock{Lean verifier receipts}
\sourcenote{module \texttt{BernsteinLogic.Soundness}; bound build \texttt{a20d90e66c13\ldots} (bound from receipt a9c0e34d3e3c9c8c, 2026-08-03T21:30:36Z). A receipt records that named checking layers accepted a build: (module, compiled-artifact SHA-256, environment SHA-256, checker layers), minted when the checking gate ran; the store receipts attesting the bound build are listed with the layers each run credited. Nothing here asks the reader to trust the store: every layer is an independently runnable program over the quoted sources --- rebuild the module and re-run the checkers to reproduce the verdict. Receipts attesting superseded builds credit no layer and are omitted from this rendering (the count of omissions, if any, is noted below).}
\begin{receipttable}
\receiptrow{a9c0e34d3e3c9c8c}{2026-08-03T21:30:36Z}{build+decorrelated}{\badge{BadgeGreen}{current}}
\end{receipttable}
\sourcenote{21 further receipt(s) in the store attest superseded builds of this module and are omitted from this rendering.}
\sourcenote{module \texttt{BernsteinLogic.Core.Catalogue}; bound build \texttt{cbc92e976a1d\ldots} (bound from receipt 5abdb3589707b5df, 2026-08-03T21:30:33Z). A receipt records that named checking layers accepted a build: (module, compiled-artifact SHA-256, environment SHA-256, checker layers), minted when the checking gate ran; the store receipts attesting the bound build are listed with the layers each run credited. Nothing here asks the reader to trust the store: every layer is an independently runnable program over the quoted sources --- rebuild the module and re-run the checkers to reproduce the verdict. Receipts attesting superseded builds credit no layer and are omitted from this rendering (the count of omissions, if any, is noted below).}
\begin{receipttable}
\receiptrow{24804e66edba0bc9}{2026-08-03T20:25:09Z}{build+decorrelated}{\badge{BadgeGreen}{current}}
\receiptrow{5abdb3589707b5df}{2026-08-03T21:30:33Z}{build+decorrelated}{\badge{BadgeGreen}{current}}
\end{receipttable}
\sourcenote{18 further receipt(s) in the store attest superseded builds of this module and are omitted from this rendering.}
\layerline{\layeritem{build}{lake build (elaboration)}{\badge{BadgeGreen}{current}}\layeritem{replay}{leanchecker (official kernel replay)}{\badge{BadgeAmber}{stale only}}\layeritem{hardened}{safe\_verify (hardened screen)}{\badge{BadgeAmber}{stale only}}\layeritem{decorrelated}{lean4export + nanoda (decorrelated kernel)}{\badge{BadgeGreen}{current}}}
\sourcenote{\badge{BadgeGray}{comparator: none for this rule} no fixture in the
three-implementation comparator harness (compiled / kernel-interpreted /
reference-interpreted) exercises a \rulename{Syzygy} node yet; current
fixtures cover \rulename{BernCone} and \rulename{Norm} shapes. The
end-to-end \rulename{Syzygy} certificate
\texttt{DpCompositeSyzygy.lean} is a kernel-checked instance but rides
the handshake harness, not the comparator.
}
\ruleblock{HOL4 statement}
\sourcenote{theory \texttt{bern\_catalogue}, HOL4 development script \texttt{\texttt{bern\_catalogueScript.sml}} \badge{BadgeGreen}{script bound}}
\holhead{Definition syzygyCheckCore\_def}
\begin{holsrc}
Definition syzygyCheckCore_def:
  syzygyCheckCore cc components premise conclusion =
    case farkasEntries cc of
    | NONE => F
    | SOME entries =>
        termIdxFree conclusion ∧
        EVERY (λ(l,e).
          0 ≤ l ∧
          (EXISTS (λc. atomSidesIdxFree c.post ∧
                       EXISTS (λpa. pa = termToPoly e) (atomLeZeroPolys c.post))
                  components ∨
           EXISTS (λa. isEqAtom a ∧ atomSidesIdxFree a ∧
                       EXISTS (λpa. pa = termToPoly e) (atomLeZeroPolys a))
                  (conjunctList premise))) entries ∧
        (FOLDR (λ(l,e) acc. polyAdd (polyScale l (termToPoly e)) acc)
               [] entries =
         termToPoly conclusion)
End
\end{holsrc}
\holhead{Definition syzygyCheck\_def}
\begin{holsrc}
Definition syzygyCheck_def:
  syzygyCheck cc components premise conclusion ⇔
    syzygyCheckCore cc components premise conclusion ∨
    syzygyCheckW cc components premise conclusion
End
\end{holsrc}
\ruleblock{HOL verifier receipts}
\badge{BadgeBlue}{documentary} \texttt{hol-strict} CLEAN over the ported checker chain: theory
\texttt{bern\_catalogue}, 472 theorems, 0 cheat-tagged, 0 nonstandard
axioms (whitelist: \texttt{DISK\_THM}; the four \texttt{bool} axioms) ---
verification run of 2026-07-22 on an isolated build host, scripts
sha-matched verbatim to the sources cited here (git blob
\texttt{e41b9ed0\ldots}). Downstream, the same chain carries the checker
to a CakeML binary with the \texttt{machine\_code\_sound} tag
\texttt{(["DISK\_THM"], [])} (theory \texttt{bernCertProofARM8}).
Family-level bridge lemmas for the \rulename{Syzygy} node
(\texttt{famEq\_syz}, \texttt{bern\_fam\_syz}) ride an in-flight assembly
and are not yet in the census table. As on the \rulename{Farkas} and
\rulename{PBound} pages: the Lean check has since grown a third (mixed
tagged-union) branch, added 2026-07-19 after this port run --- the
divergence is refusal-direction only, and its port is named residue.
Semantic stratum, also ported: the soundness twins
\texttt{sound\_syzygy} (the proof-object form) and
\texttt{sound\_pBound} (theory \texttt{bern\_tr\_soundsyz}: \texttt{hol-
strict} clean, 7 theorems, 0 offences, verification run of 2026-07-22 on
an isolated build host, content-addressed at the same sources); the
certified-fragment keystone (\texttt{checkPf\_sound\_frag}, theory
\texttt{bern\_tr\_splice}, 20 theorems, 0 offences, same run) discharges
this rule\textquotesingle s wire arm through the latter.

\sourcenote{HOL receipts are documentary: quoted results of \texttt{hol-strict} runs
and census rows from the HOL4 development\textquotesingle s build
evidence, not minted receipt artefacts. A HOL receipt store (the Lean
store\textquotesingle s twin) is named future work.
}
\ruleblock{Proof}

\sourcenote{Genre: core rule (tier 1) --- an ordinary mathematical proof at
Putnam-exam-solution detail. Every hypothesis of the rule is used below and
remarked on; the formal-anchor table maps each step to the Lean declaration
that carries it.}

\textbf{Setting and claim.}\enspace Fix a hypernet $H$. A \emph{valuation}
$\sigma$ assigns each wire a rational tensor of its declared shape; a
sequent $\Phi \vdash \Psi$ is \emph{semantically valid} when every
admissible $\sigma$ satisfying $\Phi$ satisfies $\Psi$. Atoms
$e \leq 0$ over tensor-shaped terms hold per-entry: every entry of the
evaluated tensor is $\leq 0$. The rule's data (the \texttt{SyzygyStep}
payload quoted above) supplies component sequents
$\Phi_i \vdash e_i \leq 0$ for $i = 1,\dots,n$, a combination certificate
with weights $\lambda_i \in \Q_{\geq 0}$ on the components and free
weights $\nu_k \in \Q$ on equality sub-conjuncts $g_k = 0$ of $\Phi$, a
conclusion term $c$, and a checked \emph{syntactic} identity
\begin{equation}
c \;\equiv\; \sum_{i} \lambda_i\, e_i \;+\; \sum_{k} \nu_k\, g_k
\qquad\text{(equality of exact canonical polynomials over } \Q\text{).}
\tag{$\ast$}
\end{equation}
The claim (\texttt{sound\_syzygy}): if every component sequent is
semantically valid, the composite sequent $\Phi \vdash \Psi$ is
semantically valid, where $\Psi$ is any reading of $c \leq 0$ admitted by
\texttt{syzygyConclusionForm} (Lemma~\ref{lem:syz-transport}).

\begin{blemma}[nonnegative combination]\label{lem:syz-comb}
In an ordered field, if $a_i \leq 0$ and $\lambda_i \geq 0$ for
$i = 1,\dots,n$, and $b_k = 0$ for $k = 1,\dots,m$ with $\nu_k$ arbitrary,
then $\sum_i \lambda_i a_i + \sum_k \nu_k b_k \leq 0$.
\end{blemma}
\begin{proof}
Each $\lambda_i a_i \leq 0$: multiplication by a nonnegative element
preserves $\leq$ (applied to $a_i \leq 0$), so
$\lambda_i a_i \leq \lambda_i \cdot 0 = 0$. Each $\nu_k b_k = \nu_k \cdot 0
= 0$ \emph{regardless of the sign of $\nu_k$} --- this sign-blindness on
equality atoms is precisely the freedom classical Farkas variants grant
multipliers of equality constraints \cite[\S7.3]{schrijver1986}. A finite
sum of nonpositive elements is nonpositive, by induction on $n + m$.
\end{proof}

\begin{blemma}[the certified identity evaluates]\label{lem:syz-eval}
Let $(\ast)$ hold as an identity of canonical polynomials, and let
$\sigma$ be a valuation at which $c$ and every $e_i$, $g_k$ evaluate
(per-entry, through the clamped per-entry valuation for tensor shapes).
Then the rational identity
$c(\sigma) = \sum_i \lambda_i\, e_i(\sigma) + \sum_k \nu_k\, g_k(\sigma)$
holds at every entry.
\end{blemma}
\begin{proof}
First fix what ``per-entry'' means. For an entry index $j$ of the
conclusion's shape, the \emph{clamped per-entry valuation} reads each
wire at position $j$ when $j$ is in range for that wire's shape and at
its entry $0$ otherwise --- total on every wire in sight, because the
B14 positivity pack guarantees each carries at least one entry, and
agreeing with the plain positional read wherever $j$ is in range. No
shape-agreement hypothesis is needed beyond that: a sourced
polynomial's wires share its own atom's shape (rank-0 broadcast
aside), so they clamp in step, and its clamped value at $j$ is an
honest entry of its own evaluated tensor --- which is where Step 4's
per-entry case split picks it up. Evaluation at this scalar assignment
is a ring homomorphism $\Q[\text{wires}] \to \Q$; a syntactic
identity of canonical forms is an identity of polynomials, and ring
homomorphisms preserve polynomial identities. Two remarks make this step deserve its own lemma. First,
$(\ast)$ is checked by \emph{exact canonical-form equality} over $\Q$
(\texttt{syzygyCombinationValid}, replayed through \texttt{toPoly}), not by
evaluation at test points --- so no completeness, degree bound, or
interpolation argument is needed anywhere in this proof; soundness is pure
arithmetic. Second, the checker admits three certified routes to the same
fact (the idx-free branch, the wire-coordinate branch, and the mixed
tagged-union branch:
\texttt{syzygyCheck = Core $\vee$ W $\vee$ M}); whichever branch
fires, the certified fact consumed downstream is the same --- the
identity $(\ast)$, under the same per-entry reading. The branches
differ in how terms are decoded, never in what is certified.
\end{proof}

\begin{blemma}[reading transport]\label{lem:syz-transport}
Suppose $\sigma$ satisfies the zero-normal atom $c \leq 0$. Then $\sigma$
satisfies every formula $F$ admitted by
$\mathrm{syzygyConclusionForm}(c, F)$, namely:
\emph{(Z$\leq$)} $F = (c \leq 0)$ itself;
\emph{(G)} $F = (e \geq 0)$ where $c$ is \emph{literally} the
negation term $-e$;
\emph{(L$\leq$)} $F = (l \leq r)$ where $c$ is literally the
difference term $l - r$ --- structural identities of terms, not merely
canonical-polynomial ones.
\end{blemma}
\begin{proof}
(Z$\leq$) is the identity transport. (G): per entry,
$-e(\sigma) \leq 0 \iff e(\sigma) \geq 0$ (negate and flip, an ordered-field
axiom). (L$\leq$): per entry, $l(\sigma) - r(\sigma) \leq 0 \iff l(\sigma)
\leq r(\sigma)$ (add $r(\sigma)$ to both sides). Structurality is load-bearing:
because $e$, $l$, $r$ are subterms of $c$, the evaluability that the
zero-normal atom's holding certifies for $c$ passes to them --- a
merely canonical identity would not carry evaluation (a cancelled
ill-typed subterm vanishes from the polynomial yet still poisons the
term), which is why the admitted readings are stated on term
structure. This is the ge-bridge ring
lemma \texttt{holds\_of\_syzygyConclusionForm} (the both-polarity
conclusion-form design): the readings are one exact $\Q$ identity
re-presented under a sign flip, so the polarity bridge lives \emph{once},
here, and every $\leq 0$-shaped rule (\rulename{Syzygy} itself,
\rulename{PBound}'s compact spelling, the PsdCert family) inherits all
readings with no separate justification.
\end{proof}

\begin{bthm}[soundness of \rulename{Syzygy} = Lemma B, dialect D]
\label{thm:syz-sound}
Under the payload's checked hypotheses, if every component sequent
$\Phi_i \vdash e_i \leq 0$ is semantically valid, then so is
$\Phi \vdash \Psi$.
\end{bthm}
\begin{proof}
Fix an admissible valuation $\sigma$ with $\sigma \models \Phi$; we must
show $\sigma \models \Psi$.

\emph{Step 1 (components hold at $\sigma$).} For each $i$, the payload
checks $\Phi \supseteq \Phi_i$ conjunct-wise (\texttt{premiseContains}
--- an ordered-tree containment certificate, and satisfaction descends
along it by one $\wedge$-projection per step); hence
$\sigma \models \Phi_i$, and by the assumed semantic validity of the
$i$-th component, $\sigma \models e_i \leq 0$: every entry of
$e_i(\sigma)$ is $\leq 0$.

\emph{Step 2 (wiring identities hold at $\sigma$).} Each $g_k = 0$ is an
equality sub-conjunct of $\Phi$, so $\sigma \models g_k = 0$ directly:
$g_k(\sigma) = 0$ at every entry.

\emph{Step 3 (the conclusion evaluates).} The B13 shape fields
(\texttt{conclShape}, \texttt{conclFin}) guarantee $c$ has a determinate
finite Bernstein shape over $H$, and the B14 positivity fields
(\autoref{rem:syz-b14}) guarantee every wire in sight has positive total
size; together these give a total per-entry evaluation of $c$ at $\sigma$
--- the two-case shape analysis (\,singleton vs.\ indexed\,) in
\texttt{sound\_syzygy}'s body.

\emph{Step 4 (combine).} At each entry, Lemma~\ref{lem:syz-eval} turns the
certified identity $(\ast)$ into the rational identity
$c(\sigma) = \sum_i \lambda_i e_i(\sigma) + \sum_k \nu_k g_k(\sigma)$; the
per-entry source discharge is the case split of Steps 1--2 (a
component-sourced entry's atom holds by component validity, a
wiring-sourced entry's atom holds as a premise sub-conjunct ---
\texttt{atomLeZeroPolys\_sound\_entryC} delivers per-entry nonpositivity in
both cases), and Lemma~\ref{lem:syz-comb} closes the fold:
$c(\sigma) \leq 0$ at every entry.

\emph{Step 5 (read off).} Lemma~\ref{lem:syz-transport} transports the
zero-normal bound $c \leq 0$ to whichever reading
$\mathrm{syzygyConclusionForm}(c, \Psi)$ the payload declared
(\texttt{postForm}). Hence $\sigma \models \Psi$.
\end{proof}

\begin{bremark}[why the B14 hypotheses are not decoration]
\label{rem:syz-b14}
\texttt{wiresPos} and its honest extension \texttt{compPostWiresPos}
require every wire appearing in $\Phi$'s free wires, the conclusion, and
each component conclusion to have positive total size. Drop
\texttt{compPostWiresPos} and the rule is \emph{unsound}: a component
conclusion over a zero-size wire holds \emph{vacuously} (its evaluated
tensor has no entries, so ``every entry $\leq 0$'' is empty), yet its
polynomial enters the combination $(\ast)$ with full algebraic weight ---
a vacuously-held component can then push the composite conclusion to
something false. Concretely: over a hypernet with one zero-size wire $z$, take the
single component $\top \vdash z - z + 1 \leq 0$ --- vacuously valid,
its evaluated tensor having no entries --- with weight $1$ and empty
$\nu$-row; the combination's canonical polynomial is the constant
$1$, so the checked identity admits the conclusion $1 \leq 0$, a
rank-0 falsehood concluded from a vacuous truth. (The counterexample
is also kernel-checkable, recorded at the field's introduction.) The
B13 shape fields have the matching witness one step earlier: with $c$
the term $w - w$ for an untyped wire $w$, the canonical polynomial is
$0$ --- identical to the empty combination's --- yet the conclusion
atom never evaluates and so never holds; \texttt{conclShape} and
\texttt{conclFin} exclude exactly this. This is the classic mismatch between
empty-product semantics and nonempty algebra, caught by making the
degenerate case unrepresentable rather than discouraged.
\end{bremark}

\begin{bremark}[relation to the literature]
The rule is the \emph{sound} (easy) direction of an affine Farkas-type
combination, lifted twice: atoms come from sub-\emph{proofs} rather than
sub-conjuncts, and atoms may have arbitrary polynomial degree. The linear
ancestor is Farkas' lemma \cite{farkas1902}; the textbook treatment of
multiplier sign conventions ($\lambda \geq 0$ on inequalities, free on
equalities) is \cite[\S7]{schrijver1986}. With polynomial atoms the
combination sits in the Krivine--Handelman family of positivity
certificates \cite{krivine1964,handelman1988} --- of which, again, only
the trivial direction is claimed or needed: no Positivstellensatz
completeness enters the kernel. The wiring reading --- components as
open subsystems, $g_k = 0$ as power-preserving interconnection, the
composite as dissipation of $\sum_i \lambda_i V_i$ --- is Tellegen's
theorem \cite{tellegen1952} in certificate form, i.e.\ compositional
dissipativity in the sense of Willems \cite{willems1972} and
Moylan--Hill \cite{moylanhill1978}; the port-Hamiltonian formulation is
surveyed in \cite{vanderschaftjeltsema2014}. Within the accompanying
development the rule is known as Lemma~B (dialect~D); a quantitative
effectivity result for tree-structured interconnections --- sparse
Handelman certificates with running-intersection support --- is developed
in a companion note.
\end{bremark}

\begin{anchors}
Step 1: conjunct-wise premise containment & \texttt{premiseContains} (payload field \texttt{preContains})\\
Steps 1--2, per-entry atom discharge & \texttt{atomLeZeroPolys\_sound\_entryC}; source tags \textsc{from-component}/\textsc{from-wiring}\\
Step 3: conclusion evaluates & \texttt{conclShape}, \texttt{conclFin} (B13); \texttt{wiresPos}, \texttt{compPostWiresPos} (B14); \texttt{wire\_discipline\_of\_bshape\_fin}, \texttt{Term.evalQ\_total\_singleton}\\
Step 4: certified identity, all three checker branches & \texttt{syzygyCombinationValid}, \texttt{syzygyCheck} = \texttt{syzygyCheckCore} $\vee$ \texttt{syzygyCheckW} $\vee$ \texttt{syzygyCheckM}\\
Lemma \ref{lem:syz-comb} fold & the nonnegative fold in \texttt{sound\_syzygy}'s body\\
Lemma \ref{lem:syz-transport} & \texttt{holds\_of\_syzygyConclusionForm}\\
Theorem \ref{thm:syz-sound} & \texttt{sound\_syzygy}; wired as core \texttt{RuleData.syzygy} into \texttt{ProofObject.sound}\\
\end{anchors}


\rulesection{34}{SmallGain}{alias for \textcolor{DiffBlue}{\rulename{Syzygy}}}
\label{blcatrule-SmallGain}
\ruleblock{The rule (Steele metanotation)}
\sourcenote{as realized in the mechanized development; source row: \texttt{\texttt{rules.toml}} --- \emph{SmallGain}. The metanotation statement is typeset by hand for notational quality; its mechanized ground truth is the Lean statement below, extracted from the sources at every build of this book and hash-bound. The rule \emph{roster} is machine-checked (every mechanized rule tag round-trips to exactly one catalogue row); the statement \emph{text} is editorial, checked by eye against the extracted Lean beside it.}
\begin{mathpar}
\inferrule*[Right=\rulename{SmallGain}]{
  \Phi_1 \vdash e_1 \leq 0 \\
  \Phi_2 \vdash e_2 \leq 0 \\
  \Phi_w \vdash \lambda\sigma_2 - (1{-}\theta)\,\alpha_1 \leq 0 \\
  \Phi_w \vdash \sigma_1 - \lambda(1{-}\theta)\,\alpha_2 \leq 0 \\
  \Phi \supseteq \Phi_1,\ \Phi_2,\ \Phi_w \ \text{conjunct-wise (checked)} \\
  \lambda \in \Q_{>0} \\
  \theta \in \Q \cap (0,1)
}{
  \Phi \vdash \mathcal{L}(V_1 + \lambda V_2) + \theta\,(\alpha_1 + \lambda\alpha_2) \leq 0
}
\end{mathpar}
$e_i \equiv \mathcal{L}_i V_i + \alpha_i - \sigma_i$; \quad $m_1 \equiv \lambda\sigma_2 - (1-\theta)\,\alpha_1$, \quad $m_2 \equiv \sigma_1 - \lambda(1-\theta)\,\alpha_2$; \quad conclusion identity, checked as exact canonical polynomial: $e_1 + \lambda e_2 + m_1 + m_2 = \mathcal{L}(V_1 + \lambda V_2) + \theta(\alpha_1 + \lambda\alpha_2)$.

\ruleflavor{the two membership sequents are ordinary BL sequents over the declared operating box $\Phi_w$, their \rulename{BernCone} multipliers \emph{are} the LP witness; in the conclusion identity the cross-supplies cancel by construction --- the checker never sees a fixpoint. In the baked base theory this rule sits in the admissible roster --- no proof-term sort, no wire tag: it elaborates to one \rulename{Syzygy} node at weights $(1, \lambda, 1, 1)$, which this catalogue renders as the alias tier.}
\par\medskip
\rulename{SmallGain} is the calculus\textquotesingle s
feedback-composition rule in its classical costume: two subsystems, each
dissipating its own storage function except for a supply fed by the
other, compose into a single dissipation certificate for
$V_1 + \lambda V_2$ when the loop gain is strictly less than one
($\theta$ is the margin). It is an \emph{alias}: the wire format never
carries a \rulename{SmallGain} node. A producer that finds
$(\lambda, \theta)$ --- an LP feasibility problem over the membership
certificates --- elaborates the step to the single \rulename{Syzygy}
node displayed in the recipe below before emission, so what the checker
recomputes is exactly \rulename{Syzygy}\textquotesingle s obligation
set: the exact conclusion identity at weights $(1, \lambda, 1, 1)$, in
which the cross-supplies cancel by construction. Nothing fixpoint-shaped
survives to checking time; the classical contraction argument lives
entirely on the producer\textquotesingle s side, in the choice of
$(\lambda, \theta)$.

What would break without each hypothesis is inherited from
\rulename{Syzygy} through the unpack --- that
section\textquotesingle s commentary applies verbatim to the unpacked
node. The two conditions native to this rule: $\lambda > 0$, of which
soundness uses only $\lambda \geq 0$ (the strictness backs the
rule\textquotesingle s published meaning), and $\theta \in (0,1)$, which
is well-formedness for $\theta$\textquotesingle s published meaning as a
robustness margin rather than a soundness need --- the membership
sequents self-police, and an infeasible pair simply yields unprovable
memberships.

\cataloguenote{The four component sequents are ordinary sequents over the declared
operating box; their cone-membership multipliers \emph{are} the LP
witness. Signed polynomial contracts (rather than norm gains) are
constitutive, not stylistic: swing-up analyses in the
{\AA}str\"om--Furuta family exploit the \emph{sign} of
$\dot\vartheta\cos\vartheta$, which no norm-gain contract can express.
$(\lambda,\theta)$-existence stays on the producer side, deliberately
outside the checker; the $n$-component generalization is the
\rulename{SmallGainN} section --- the same alias discipline over an
arbitrary interconnection graph, elaborating to one \rulename{Syzygy}
node, of which this rule is the $n = 2$ instance.
}
\ruleblock{Lean statement}
\declhead{BernsteinLogic.Composition}{SmallGainStep}{\badge{BadgeGreen}{statement bound}}
\begin{leansrc}
/-- Applicability data of the SmallGain rule — field-for-field the
draft's `smallGainShapes` obligations (§3.3):

* (a) the dissipation conclusions are `(eᵢ ≤ 0)` atoms — the residual
  reading eᵢ = 𝓛ᵢVᵢ + αᵢ − σᵢ is definitional at the checker level
  (`𝓛ᵢVᵢ := eᵢ − αᵢ + σᵢ`), so no separate field pins it;
* (b) the membership conclusions are `(mᵢ ≤ 0)` atoms whose polynomials
  equal their schema instances at (λ, θ) — exact canonical polynomial;
* (c) λ > 0 and 0 < θ < 1 — three rational comparisons. The θ range is
  NOT consumed by soundness (the memberships self-police; the range backs
  the published guard-modulus semantics — draft §3.2's honesty note);
* (d) the conclusion identity: the (1, λ, 1, 1) fold of the four posted
  polynomials equals the conclusion schema 𝓛(V₁+λV₂) + θ(α₁+λα₂) — the
  checker replays the cancellation (`smallGain_is_syzygy` proves the
  schema coherent: given (b) and the residual reading, (d) is a theorem,
  not a search).

B13/B14 fields as in `SyzygyStep`, component-post coverage included. -/
structure SmallGainStep (H : Hypernet) where
  seq : Sequent
  diss1 : Sequent
  diss2 : Sequent
  mem1 : Sequent
  mem2 : Sequent
  lambda : ℚ
  theta : ℚ
  alpha1 : Term
  alpha2 : Term
  sigma1 : Term
  sigma2 : Term
  lyapTerm : Term
  e1 : Term
  e2 : Term
  m1 : Term
  m2 : Term
  concl : Term
  -- (a) dissipation shapes
  diss1Post : diss1.post = .atom .le e1 (.const 0)
  diss2Post : diss2.post = .atom .le e2 (.const 0)
  -- (b) membership shapes and schemas
  mem1Post : mem1.post = .atom .le m1 (.const 0)
  mem2Post : mem2.post = .atom .le m2 (.const 0)
  mem1Poly : m1.toPoly = Poly.sub (Poly.scale lambda sigma2.toPoly)
      (Poly.scale (1 - theta) alpha1.toPoly)
  mem2Poly : m2.toPoly = Poly.sub sigma1.toPoly
      (Poly.scale (lambda * (1 - theta)) alpha2.toPoly)
  -- G1 rider: the posted residuals and the conclusion are idx-free —
  -- their canonical polynomials carry the rule's meaning (`Term.toPoly`'s
  -- G1 invariant; mirrors `farkasCheck`'s gates)
  e1IdxFree : e1.idxFree = true
  e2IdxFree : e2.idxFree = true
  m1IdxFree : m1.idxFree = true
  m2IdxFree : m2.idxFree = true
  conclIdxFree : concl.idxFree = true
  -- (c) the three rational comparisons
  lambdaPos : 0 < lambda
  thetaLo : 0 < theta
  thetaHi : theta < 1
  -- (d) the conclusion identity (checker-replayed cancellation), and the
  -- conclusion term's tie to the schema
  foldEq : Poly.add e1.toPoly (Poly.add (Poly.scale lambda e2.toPoly)
        (Poly.add m1.toPoly m2.toPoly))
      = Poly.add lyapTerm.toPoly
          (Poly.scale theta
            (Poly.add alpha1.toPoly (Poly.scale lambda alpha2.toPoly)))
  conclPoly : concl.toPoly = Poly.add lyapTerm.toPoly
      (Poly.scale theta
        (Poly.add alpha1.toPoly (Poly.scale lambda alpha2.toPoly)))
  postEq : seq.post = .atom .le concl (.const 0)
  -- premise containment for all four components
  pre1 : premiseContains seq.pre diss1.pre
  pre2 : premiseContains seq.pre diss2.pre
  preM1 : premiseContains seq.pre mem1.pre
  preM2 : premiseContains seq.pre mem2.pre
  -- B13/B14 (component-post coverage per the §R2 finding)
  conclShape : ∃ s, concl.bshape H = some s
  conclFin : concl.finWires H
  wiresPos : ∀ w ∈ seq.pre.freeWires ++ concl.freeWires,
      ∃ τ, H.wireType w = some τ ∧ 0 < Shape.totalSize τ.shape
  compPostWiresPos : ∀ c ∈ [diss1, diss2, mem1, mem2],
      ∀ w ∈ c.post.freeWires,
      ∃ τ, H.wireType w = some τ ∧ 0 < Shape.totalSize τ.shape
\end{leansrc}
\declhead{BernsteinLogic.Composition}{toSyzygy}{\badge{BadgeGreen}{statement bound}}
\begin{leansrc}
/-- The SmallGain node, AS the Syzygy instance it is (weights
(1, λ, 1, 1), components in combination order). -/
def toSyzygy (st : SmallGainStep H) : SyzygyStep H
\end{leansrc}
\declhead{BernsteinLogic.Composition}{sound\_smallGain}{\badge{BadgeGreen}{statement bound}}
\begin{leansrc}
/-- Soundness of SmallGain (= Lemma B, dialect G), as the Syzygy
corollary: the four component validities feed `sound_syzygy` on the
collapsed instance. The three rational comparisons (λ > 0, 0 < θ < 1)
enter only through λ ≥ 0 inside the combination check; θ's range is
well-formedness, exactly as the draft's honesty note states. -/
theorem sound_smallGain {H : Hypernet} (st : SmallGainStep H)
    (h1 : st.diss1.semanticallyValid H)
    (h2 : st.diss2.semanticallyValid H)
    (hm1 : st.mem1.semanticallyValid H)
    (hm2 : st.mem2.semanticallyValid H) :
    st.seq.semanticallyValid H
\end{leansrc}
\declhead{BernsteinLogic.Composition}{toSyzygy\_components}{\badge{BadgeGreen}{statement bound}}
\begin{leansrc}
theorem toSyzygy_components (st : SmallGainStep H) :
    st.toSyzygy.components = st.components
\end{leansrc}
\ruleblock{Correspondence (Steele display $\leftrightarrow$ Lean)}
\badge{BadgeGreen}{reviewer-certified against the displayed texts}
\sourcenote{The dictionary below is a reviewer-certified statement of how the mechanized statements above interpret the displayed rule: display metavariable $\mapsto$ Lean binder, each display hypothesis $\mapsto$ the premise or side condition carrying it, the display conclusion $\mapsto$ the goal --- with any deliberate difference in width between display and mechanization named outright. Certification is recorded as a hash triple over (the displayed cell, the extracted statement texts, this dictionary\textquotesingle s own text); if any of the three changes, this block turns amber until a reviewer re-certifies. The correspondence claim itself remains editorial --- reviewer-grade, not machine-checked; its machine-checked successor types each soundness statement against the rule as a structured object, and supersedes this block per rule as it lands.}
The alias\textquotesingle s Lean vehicle is the payload structure
\texttt{SmallGainStep H} (quoted in full above --- its fields carry the
display\textquotesingle s metavariables), the verified unpack
\texttt{toSyzygy}, and the soundness corollary
\texttt{sound\_smallGain}.
\begin{itemize}
\item $\lambda \mapsto$ \texttt{lambda} $: \Q$ with \texttt{lambdaPos}
  $: 0 < \lambda$; $\theta \mapsto$ \texttt{theta} with
  \texttt{thetaLo}, \texttt{thetaHi} $: 0 < \theta < 1$; $\alpha_i,
  \sigma_i \mapsto$ \texttt{alpha1}/\texttt{alpha2},
  \texttt{sigma1}/\texttt{sigma2}; the composite storage $V_1 + \lambda
  V_2 \mapsto$ \texttt{lyapTerm}.
\item the dissipation premises $\Phi_i \vdash e_i \leq 0$ $\mapsto$
  \texttt{diss1}/\texttt{diss2} with the shape ties
  \texttt{diss1Post}/\texttt{diss2Post} (each $e_i$ is a field; its
  published spelling $\mathcal{L}_i V_i + \alpha_i - \sigma_i$ lives on
  the producer\textquotesingle s side). The operating-box membership
  premises $\mapsto$ \texttt{mem1}/\texttt{mem2} with
  \texttt{mem1Poly}/\texttt{mem2Poly} tying their residuals to the
  displayed schemas $\lambda\sigma_2 - (1{-}\theta)\alpha_1$ and
  $\sigma_1 - \lambda(1{-}\theta)\alpha_2$ as canonical polynomials.
\item the display\textquotesingle s side text --- the exact conclusion
  identity (the cross-supply-cancellation gloss is the flavor
  text\textquotesingle s) --- $\mapsto$ \texttt{foldEq}
  and \texttt{conclPoly} (the checker-replayed cancellation, exact over
  canonical polynomials) plus \texttt{postEq} (the conclusion tie);
  the displayed containment premise ($\Phi \supseteq \Phi_1, \Phi_2,
  \Phi_w$ conjunct-wise) $\mapsto$
  \texttt{pre1}/\texttt{pre2}/\texttt{preM1}/\texttt{preM2}.
\item conclusion $\Phi \vdash \mathcal{L}(V_1 + \lambda V_2) +
  \theta\,(\alpha_1 + \lambda\alpha_2) \leq 0$ $\mapsto$ the goal
  \texttt{st.seq.semanticallyValid H} of \texttt{sound\_smallGain},
  whose four hypotheses are exactly the component validities.
\item the alias artefacts: the unpack function is \texttt{toSyzygy}
  (weights $(1, \lambda, 1, 1)$, components in combination order); the
  soundness corollary routes through the target rule\textquotesingle s
  lemma on the collapsed instance --- the displayed rule never reaches
  the wire.
\end{itemize}
Width notes, both directions: soundness uses only $\lambda \geq 0$ of
\texttt{lambdaPos} (the strictness backs the rule\textquotesingle s
published meaning), and $\theta$\textquotesingle s range is
well-formedness rather than a soundness need --- both as the
commentary states. Conversely the structure carries idx-freeness and
wire-shape fields (\texttt{e1IdxFree} through
\texttt{compPostWiresPos}) that the display leaves to the ambient
calculus\textquotesingle s conventions: the mechanization is narrower
there, naming every side condition the checker replays.

\ruleblock{Lean verifier receipts}
\sourcenote{module \texttt{BernsteinLogic.Composition}; bound build \texttt{9116f454beb0\ldots} (bound from receipt 548ea3dda8ed73a2, 2026-08-03T21:30:29Z). A receipt records that named checking layers accepted a build: (module, compiled-artifact SHA-256, environment SHA-256, checker layers), minted when the checking gate ran; the store receipts attesting the bound build are listed with the layers each run credited. Nothing here asks the reader to trust the store: every layer is an independently runnable program over the quoted sources --- rebuild the module and re-run the checkers to reproduce the verdict. Receipts attesting superseded builds credit no layer and are omitted from this rendering (the count of omissions, if any, is noted below).}
\begin{receipttable}
\receiptrow{548ea3dda8ed73a2}{2026-08-03T21:30:29Z}{build+decorrelated}{\badge{BadgeGreen}{current}}
\end{receipttable}
\sourcenote{10 further receipt(s) in the store attest superseded builds of this module and are omitted from this rendering.}
\layerline{\layeritem{build}{lake build (elaboration)}{\badge{BadgeGreen}{current}}\layeritem{replay}{leanchecker (official kernel replay)}{\badge{BadgeAmber}{stale only}}\layeritem{hardened}{safe\_verify (hardened screen)}{\badge{BadgeAmber}{stale only}}\layeritem{decorrelated}{lean4export + nanoda (decorrelated kernel)}{\badge{BadgeGreen}{current}}}
\sourcenote{\badge{BadgeGray}{comparator: structurally n/a} the alias unpacks
\emph{before} emission (no rule tag; producers elaborate a
\rulename{SmallGain} step to its \rulename{Syzygy} node), so no wire
certificate --- and hence no comparator fixture --- can ever contain a
\rulename{SmallGain} node. Comparator coverage for this rule \emph{is}
\rulename{Syzygy}\textquotesingle s. The end-to-end instance
\texttt{DpCompositeSmallGain.lean} is kernel-checked upstream of the
unpack.
}
\ruleblock{HOL4 statement}
\absentbox{OUT OF SCOPE BY DESIGN}{Out of scope for the HOL side \emph{by design}, not by backlog: the wire
format never carries a \rulename{SmallGain} node (no tag --- producers
elaborate to the $(1,\lambda,1,1)$ \rulename{Syzygy} node before
emission, warrant \texttt{SmallGainStep.toSyzygy}), and the HOL4
development\textquotesingle s object is the wire checker. The HOL
statement that verifies a composed small-gain certificate is therefore
\rulename{Syzygy}\textquotesingle s (\texttt{syzygyCheck}, previous
section); an independent HOL port of the alias\textquotesingle s unpack
function would be new assurance surface only if HOL-side producers ever
emit \rulename{SmallGain} steps directly.
}
\ruleblock{HOL verifier receipts}
\absentbox{NONE}{no HOL verification receipt exists for this rule.}
\ruleblock{Proof}

\sourcenote{Genre: alias rule (tier 3) --- a respelling template. An alias
owes the book five artefacts: the verified 1$\to$1 unpack function, the
filled-in application of the target rule, the payload-transport table, the
exact algebraic identity that makes the respelling correct, and the
soundness corollary through the target rule's lemma. All five follow.
(The tier criterion counts derivation nodes, not instantiation work: the
unpack \emph{computes} the weighting, but its image is a single
\rulename{Syzygy} node, so the rule is an alias.)}

\begin{brecipe}[the alias template]\label{rec:sg-template}
An alias rule $A$ is warranted by a verified function
$\mathrm{unpack} : \mathrm{Step}_A \to \mathrm{Step}_T$ into the payload
of a \emph{single} target node $T$ (core or admissible), such that the
$A$-node and its image conclude the same sequent from the same premises:
one node $\to$ one node, payload transported, nothing derived. Soundness
is then a corollary, $\mathrm{sound}_A = \mathrm{sound}_T \circ
\mathrm{unpack}$, and the wire format need never carry $A$ at all ---
producers elaborate before emission, which is why \rulename{SmallGain} has
no \texttt{RuleTag}.
\end{brecipe}

\begin{brecipe}[the filled-in application --- \rulename{SmallGain}'s
premises rewritten \emph{into} the \rulename{Syzygy} instance]
\label{rec:sg-filled}
A \rulename{SmallGain} step carries four component proofs: the two
dissipation sequents (over their own boxes $\Phi_1, \Phi_2$) and the two
gain-membership sequents (over the declared operating box $\Phi_w$).
Define the four atom terms and the substitution into
\rulename{Syzygy}'s metavariables:
\[
e_1 \equiv \mathcal{L}_1 V_1 + \alpha_1 - \sigma_1, \qquad
e_2 \equiv \mathcal{L}_2 V_2 + \alpha_2 - \sigma_2, \qquad
m_1 \equiv \lambda\sigma_2 - (1{-}\theta)\,\alpha_1, \qquad
m_2 \equiv \sigma_1 - \lambda(1{-}\theta)\,\alpha_2,
\]
\[
n := 4, \qquad
(e_1, e_2, e_3, e_4) := (e_1, e_2, m_1, m_2), \qquad
(\Phi_1, \Phi_2, \Phi_3, \Phi_4) := (\Phi_1, \Phi_2, \Phi_w, \Phi_w),
\]
\[
(\lambda_1, \lambda_2, \lambda_3, \lambda_4) := (1, \lambda, 1, 1),
\qquad
\{g_k\}_k := \varnothing \ \text{(no equality entries)}, \qquad
c := \mathcal{L}(V_1 + \lambda V_2) + \theta\,(\alpha_1 + \lambda\alpha_2),
\]
with $\lambda \in \Q_{>0}$ and $\theta \in \Q \cap (0,1)$ from the
payload. Under this substitution the \rulename{SmallGain} node \emph{is}
the following \rulename{Syzygy} application (compare the rule display in
the \rulename{Syzygy} section, field for field):
\begin{mathpar}
\inferrule*[Right=\rulename{Syzygy}]{
  \Phi_1 \vdash e_1 \leq 0 \\
  \Phi_2 \vdash e_2 \leq 0 \\
  \Phi_w \vdash m_1 \leq 0 \\
  \Phi_w \vdash m_2 \leq 0 \\\\
  (1, \lambda, 1, 1) \in \Q_{\geq 0}^{\,4} \\
  c \;\equiv\; 1\cdot e_1 + \lambda\cdot e_2 + 1\cdot m_1 + 1\cdot m_2
  \ \text{(exact canonical polynomial)}
}{
  \Phi \vdash c \leq 0
}
\end{mathpar}
with $\Phi \supseteq \Phi_1, \Phi_2, \Phi_w$ conjunct-wise (the
four-way \texttt{preContains} split). Discharging \rulename{Syzygy}'s
obligations at this instance: the nonnegativity checks are $1 \geq 0$
(three times) and $\lambda \geq 0$ (from $\lambda > 0$ --- the
\emph{only} use soundness makes of $\lambda$'s sign); the sourcing check
is immediate (each of the four entries is the corresponding component's
posted atom, verbatim); and the exact-identity check is
Proposition~\ref{prop:sg-identity} below; the five index-freeness
gates (the four posted residuals and the composite conclusion carry no
free indices --- recomputed syntactic checks, the G1 rider that lets
canonical polynomials carry the rule's meaning) discharge by
recomputation on the five terms; the remaining obligations
--- the B13 shape pack, the B14 positivity pack, and the zero-normal
\texttt{postForm} declaration --- transport verbatim, per
Recipe~\ref{rec:sg-transport}. The dead premise slots of
\rulename{Syzygy} are honestly dead: no equality sub-conjuncts of $\Phi$
enter (the $\nu$-row is empty), because the cross-supply cancellation is
performed by the membership atoms themselves. The verified function
computing this rewriting is \texttt{SmallGainStep.toSyzygy}, quoted
above.
\end{brecipe}

\begin{bprop}[the conclusion identity]\label{prop:sg-identity}
As an identity of exact canonical polynomials,
\[
e_1 + \lambda e_2 + m_1 + m_2
\;=\;
\mathcal{L}(V_1 + \lambda V_2) \;+\; \theta\,(\alpha_1 + \lambda\alpha_2).
\]
\end{bprop}
\begin{proof}
Expand and collect --- two lines, which is the point:
\begin{align*}
e_1 + \lambda e_2 + m_1 + m_2
&= \bigl(\mathcal{L}_1 V_1 + \lambda\,\mathcal{L}_2 V_2\bigr)
 + \alpha_1 + \lambda\alpha_2
 \underbrace{-\,\sigma_1 - \lambda\sigma_2
 + \lambda\sigma_2 + \sigma_1}_{=\,0\ \text{(cross-supplies cancel)}}
 - (1{-}\theta)\alpha_1 - \lambda(1{-}\theta)\alpha_2\\
&= \mathcal{L}(V_1 + \lambda V_2)
 + \bigl(1 - (1{-}\theta)\bigr)\alpha_1
 + \lambda\bigl(1 - (1{-}\theta)\bigr)\alpha_2
\;=\; \mathcal{L}(V_1 + \lambda V_2) + \theta\,(\alpha_1 + \lambda\alpha_2).
\end{align*}
The final rewrite reads $\mathcal{L}(V_1 + \lambda V_2)$ for
$\mathcal{L}_1 V_1 + \lambda\,\mathcal{L}_2 V_2$, and that is a
definition, not a step: over disjoint variable sets the composite
generator acts coordinate-wise, so the composite storage-rate term
--- the payload's \texttt{lyapTerm}, pinned to exactly this sum by
the exact canonical ties (\texttt{foldEq} $\cdot$
\texttt{conclPoly}) --- has $\mathcal{L}(V_1 + \lambda V_2)$ as its
published spelling. The cross-supply cancellation is \emph{by
construction of the multipliers}: the membership sequents were chosen so that their
polynomials supply exactly $+\lambda\sigma_2$ and $+\sigma_1$. The checker
replays this as \rulename{Syzygy}'s exact canonical-polynomial identity
$(\ast)$ and \textbf{never sees a fixpoint} --- the classical small-gain
contraction argument is dissolved into two lines of ring arithmetic at
certificate-checking time.
\end{proof}

\begin{brecipe}[payload transport --- \texttt{toSyzygy}, field for field]
\label{rec:sg-transport}
\begin{center}\small
\begin{tabular}{@{}ll@{}}
\textbf{\texttt{SmallGainStep}} & \textbf{\texttt{SyzygyStep}}\\
\midrule
\texttt{seq} & \texttt{seq}\\
\texttt{diss1, diss2, mem1, mem2} & \texttt{components} (that order = combination order)\\
\texttt{combTerm} \textnormal{(weights $(1,\lambda,1,1)$)} & \texttt{combination}\\
\texttt{concl} & \texttt{conclusion}\\
\texttt{pre1, pre2, preM1, preM2} & \texttt{preContains} (four-way case split)\\
\texttt{conclShape, conclFin} & \texttt{conclShape, conclFin} (verbatim)\\
\texttt{wiresPos, compPostWiresPos} & \texttt{wiresPos, compPostWiresPos} (verbatim)\\
\texttt{e1/e2/m1/m2/conclIdxFree} & \textnormal{consumed by} \texttt{combValid} $\mapsto$ \texttt{valid} \textnormal{(G1 gates)}\\
\texttt{combValid} & \texttt{valid}\\
\texttt{postEq} & \texttt{postForm} $=$ \textnormal{(Z$\leq$), the zero-normal reading}\\
\end{tabular}
\end{center}
The alias declares the \emph{zero-normal} reading of the both-polarity
conclusion form: \rulename{SmallGain}'s own contract stays
$\leq$-only, and any polarity flexibility is inherited from
\rulename{Syzygy}'s side of the fence (the reading-transport lemma
\texttt{holds\_of\_syzygyConclusionForm}, proved in the \rulename{Syzygy}
section), not re-justified here.
\end{brecipe}

\begin{bprop}[soundness corollary]\label{prop:sg-sound}
If the four component sequents are semantically valid, the composite
sequent
$\Phi \vdash \mathcal{L}(V_1 + \lambda V_2) + \theta(\alpha_1 +
\lambda\alpha_2) \leq 0$ is semantically valid.
\end{bprop}
\begin{proof}
\texttt{sound\_smallGain} $=$ \texttt{sound\_syzygy} applied to
\texttt{st.toSyzygy}, feeding the four validities through the four-way
membership case split. Beyond that routing there is no further
argument --- the alias genre working as intended.
\end{proof}

\begin{bremark}[$\theta$'s range is well-formedness, not soundness]
The identity of Proposition~\ref{prop:sg-identity} holds at \emph{any}
$\theta$; the memberships self-police: an infeasible $(\lambda, \theta)$
makes the required membership inequalities false over the operating
box, so no sound certificate can discharge them. The $(0,1)$ range check
backs the \emph{published meaning} of $\theta$ as a guard modulus --- a
first-class robustness radius travelling in the proof term, per-node in
the payload, with the certificate header publishing the composite
modulus as the min over nodes --- and $\lambda > 0$ enters soundness
only as $\lambda \geq 0$.
\end{bremark}

\begin{bremark}[relation to the literature]
The rule certificate-izes the Lyapunov small-gain construction: two ISS-style
dissipation inequalities with gain interconnection conditions yield a
composite Lyapunov function $V_1 + \lambda V_2$ --- the nonlinear
small-gain theorem of Jiang--Teel--Praly \cite{jiangteelpraly1994} in its
Lyapunov formulation \cite{jiangmareelswang1996}, sitting inside the
dissipativity-composition tradition \cite{willems1972,moylanhill1978}.
Where the classical theorems close a fixpoint/contraction argument, the
certificate form closes a ring identity: the analytic content is moved
into the \emph{choice} of $(\lambda, \theta)$ and the membership
certificates, and the kernel checks only arithmetic. The contracts are
\emph{signed polynomial} inequalities rather than norm gains for
cause: {\AA}str\"om--Furuta's swing-up analysis \cite{astromfuruta2000}
exploits the sign of $\dot\vartheta\cos\vartheta$, information a
pure norm-gain contract does not record. $(\lambda,\theta)$-existence
results are deliberately outside the checker; the $n$-component
generalization over arbitrary interconnection graphs
\cite{dashkovskiyrufferwirth2007} is catalogued as
\rulename{SmallGainN} --- the same alias discipline, one
\rulename{Syzygy} node, with this rule as its $n = 2$ instance.
\end{bremark}

\begin{anchors}
Recipe \ref{rec:sg-transport} unpack & \texttt{SmallGainStep.toSyzygy} (Composition.lean); \texttt{toSyzygy\_components}\\
Proposition \ref{prop:sg-identity} & \texttt{combValid}'s fold identity (\texttt{foldEq} $\cdot$ \texttt{conclPoly}), replayed as \rulename{Syzygy}'s $(\ast)$\\
Proposition \ref{prop:sg-sound} & \texttt{sound\_smallGain} (a \texttt{sound\_syzygy} corollary)\\
no wire tag, by design & \texttt{RuleTag} comment block (Core/Catalogue.lean): producers elaborate before emission\\
end-to-end instance & \texttt{Handshake/DpCompositeSmallGain.lean} (kernel-checked upstream of the unpack)\\
\end{anchors}


\rulesection{35}{SmallGainN}{alias for \textcolor{DiffBlue}{\rulename{Syzygy}}}
\label{blcatrule-SmallGainN}
\ruleblock{The rule (Steele metanotation)}
\sourcenote{as realized in the mechanized development; source row: \texttt{\texttt{rules.toml}} --- \emph{SmallGainN}. The metanotation statement is typeset by hand for notational quality; its mechanized ground truth is the Lean statement below, extracted from the sources at every build of this book and hash-bound. The rule \emph{roster} is machine-checked (every mechanized rule tag round-trips to exactly one catalogue row); the statement \emph{text} is editorial, checked by eye against the extracted Lean beside it.}
\begin{mathpar}
\inferrule*[Right=\rulename{SmallGainN}]{
  \bigl\{\Phi_i \vdash e_i \leq 0\bigr\}_{i=1}^{n} \\
  \bigl\{\Phi_w \vdash m_i \leq 0\bigr\}_{i=1}^{n} \\
  \Phi \supseteq \Phi_1, \ldots, \Phi_n, \Phi_w \ \text{conjunct-wise (checked)} \\
  \lambda \in \Q_{>0}^{n} \\
  B \in \Q_{\geq 0}^{n \times n} \\
  \textstyle\sum_i B_{ij} = \lambda_j \ (\forall j) \\
  \theta \in \Q \cap (0,1)
}{
  \Phi \vdash \mathcal{L}\bigl(\textstyle\sum_i \lambda_i V_i\bigr) + \theta \textstyle\sum_j \lambda_j \alpha_j \leq 0
}
\end{mathpar}
$e_i \equiv \mathcal{L}_i V_i + \alpha_i - \sigma_i$; \quad $m_i \equiv \lambda_i \sigma_i - (1-\theta) \textstyle\sum_j B_{ij}\, \alpha_j$ (the repayment matrix\textquotesingle s membership rows); \quad conclusion identity, checked as exact canonical polynomial: $\textstyle\sum_i \lambda_i e_i + \sum_i m_i = \mathcal{L}(\sum_i \lambda_i V_i) + \theta \sum_j \lambda_j \alpha_j$ under the column condition.

\ruleflavor{the $n$ membership sequents are ordinary BL sequents over the declared operating box $\Phi_w$, their \rulename{BernCone} multipliers \emph{are} the LP witness; the repayment matrix $B$ routes each component\textquotesingle s supply through its neighbours, and the column condition $\sum_i B_{ij} = \lambda_j$ makes every cross-supply cancel in the conclusion identity --- the checker never sees a fixpoint, a spectral radius, or the interconnection graph. Like the binary rule it is an alias: no proof-term sort, no wire tag; it elaborates to one \rulename{Syzygy} node at weights $(\lambda_1, \ldots, \lambda_n, 1, \ldots, 1)$, and $B$ itself never reaches the wire (the membership rows carry it; the exact identity is the whole obligation). The binary rule is the $n = 2$ instance on the nose: $\lambda = (1, \lambda)$, $B_{21} = 1$, $B_{12} = \lambda$ reproduce \rulename{SmallGain}\textquotesingle s membership rows and its $(1, \lambda, 1, 1)$ weights.}
\par\medskip
\rulename{SmallGainN} is the $n$-component quantitative small-gain
rule: $n$ subsystems on an arbitrary directed interconnection graph,
each dissipating its own storage function except for supplies fed by
its neighbours, compose into a single dissipation certificate for the
weighted sum $\sum_i \lambda_i V_i$ when a weighted loop-gain condition
holds with margin $\theta$. The graph structure enters through the
\emph{repayment matrix} $B$: the membership row
$m_i \equiv \lambda_i \sigma_i - (1{-}\theta)\sum_j B_{ij}\alpha_j$
charges component $i$\textquotesingle s incoming supply against a
$(1{-}\theta)$-discounted draw on its neighbours\textquotesingle{}
decay budgets, and the column condition $\sum_i B_{ij} = \lambda_j$
says every component\textquotesingle s budget is drawn on at exactly
its weight in the composite --- which is what makes the cross-supplies
cancel identically in the conclusion.

Like the binary rule, it is an \emph{alias}: the wire format never
carries a \rulename{SmallGainN} node, and the checker gains no new
obligation kind. A producer that finds $(\lambda, B, \theta)$ --- an LP
feasibility problem over the membership certificates; when the gain
matrix has spectral radius below $1 - \theta$, a Neumann-series weight
vector witnesses feasibility, so the classical spectral small-gain
condition is exactly an existence theory for this rule\textquotesingle s
side data --- elaborates the step to a single \rulename{Syzygy} node at
weights $(\lambda_1, \ldots, \lambda_n, 1, \ldots, 1)$ before emission.
What the checker recomputes is exactly
\rulename{Syzygy}\textquotesingle s obligation set: the exact
conclusion identity over canonical polynomials. Neither the graph, nor
$B$, nor anything fixpoint- or spectrum-shaped survives to checking
time.

What would break without each hypothesis is inherited from
\rulename{Syzygy} through the unpack --- that
section\textquotesingle s commentary applies verbatim to the unpacked
node. The conditions native to this rule: $\lambda_i > 0$, of which
soundness uses only $\lambda_i \geq 0$ (strictness backs the published
meaning: every component genuinely participates); $\theta \in (0,1)$,
well-formedness for $\theta$\textquotesingle s published meaning as a
robustness margin (the identity holds at any $\theta$; the memberships
self-police); and the column condition, which is well-formedness at the
\emph{schema} level --- it is what makes the membership rows deliver
the stated conclusion for every decay budget $\alpha$, while at any
concrete instance the exact identity subsumes it: an emitted node whose
identity checks is sound regardless of how the producer bookkept $B$.

The binary rule is the $n = 2$ instance on the nose: $\lambda = (1,
\lambda)$, $B_{21} = 1$, $B_{12} = \lambda$ reproduce
\rulename{SmallGain}\textquotesingle s membership rows (as a set ---
the two rows arrive in the opposite order, a bookkeeping difference in
combination order) and its $(1, \lambda, 1, 1)$ weights.

\cataloguenote{The $n$-component small-gain theorem for general interconnection
topologies is Dashkovskiy--R\"uffer--Wirth
\cite{dashkovskiyrufferwirth2007}; this rule is its certificate form,
with the cycle/spectral conditions of that tradition dissolved into the
producer\textquotesingle s choice of $(\lambda, B, \theta)$ and the
checker replaying only ring arithmetic. Signed polynomial contracts
rather than norm gains, for the same cause as the binary section.
The dedicated $n$-ary payload vehicle is mechanized as the binary
rule\textquotesingle s pattern at arbitrary $n$: the exact identity
(\texttt{smallGainN\_is\_syzygy} --- the cross-supplies cancel under the
column condition, by canonical-polynomial extensionality), the payload
structure with its verified collapse
(\texttt{SmallGainNStep.toSyzygy}, weights
$(\lambda_1, \ldots, \lambda_n, 1, \ldots, 1)$), the soundness
corollary (\texttt{sound\_smallGainN}), and the column
condition\textquotesingle s necessity as a machine-checked mutant
(\texttt{smallGainN\_colSum\_mutant\_refuted}: delete it and a
one-component instance separates fold from conclusion, by
\texttt{decide}). The semantic face lives beside the binary
one\textquotesingle s in the discrete-kernel module:
\texttt{smallGainN\_drift} with its assembled avoid bound, and the
binary recovery on the nose (\texttt{smallGain\_drift\_of\_N}: $\lambda
= (1, \lambda)$, $B_{21} = 1$, $B_{12} = \lambda$ reproduce the binary
rows --- order-swapped, set-equal, as the commentary states). Every
\emph{emitted} instance remains checked and warranted as
\rulename{Syzygy}\textquotesingle s node: the alias unpacks before
emission.
}
\ruleblock{Lean statement}
\declhead{BernsteinLogic.Composition}{SmallGainNStep}{\badge{BadgeGreen}{statement bound}}
\begin{leansrc}
/-- The n-ary payload vehicle: `SmallGainStep`'s pattern at arbitrary `kk`.
Every field is the checker's obligation for the rule; the collapse target is
one `Syzygy` node (`toSyzygy`). `lambdaPos`, `thetaLo/Hi`, and `colSum` are
rule-level well-formedness exactly as in the binary structure (soundness
uses only `0 ≤ λᵢ` inside the combination check; an emitted node whose
`foldEq` checks is sound however the producer bookkept `B` — the identity
`smallGainN_is_syzygy` is what guarantees the row shapes always discharge
`foldEq` under the column condition). -/
structure SmallGainNStep (H : Hypernet) (kk : ℕ) where
  seq : Sequent
  diss : Fin kk → Sequent
  mem : Fin kk → Sequent
  lambda : Fin kk → ℚ
  B : Fin kk → Fin kk → ℚ
  theta : ℚ
  e : Fin kk → Term
  m : Fin kk → Term
  concl : Term

  dissPost : ∀ i, (diss i).post = .atom .le (e i) (.const 0)
  memPost : ∀ i, (mem i).post = .atom .le (m i) (.const 0)

  eIdxFree : ∀ i, (e i).idxFree = true
  mIdxFree : ∀ i, (m i).idxFree = true
  conclIdxFree : concl.idxFree = true

  lambdaPos : ∀ i, 0 < lambda i
  thetaLo : 0 < theta
  thetaHi : theta < 1
  colSum : ∀ j, ((List.finRange kk).map (fun i => B i j)).sum = lambda j

  foldEq : (((List.finRange kk).map (fun i => ((lambda i : ℚ), e i))
      ++ (List.finRange kk).map (fun i => ((1 : ℚ), m i))).foldr
      (fun pr acc => Poly.add (Poly.scale pr.1 pr.2.toPoly) acc) [])
    = concl.toPoly
  postEq : seq.post = .atom .le concl (.const 0)

  preDiss : ∀ i, premiseContains seq.pre (diss i).pre
  preMem : ∀ i, premiseContains seq.pre (mem i).pre

  conclShape : ∃ s, concl.bshape H = some s
  conclFin : concl.finWires H
  wiresPos : ∀ w ∈ seq.pre.freeWires ++ concl.freeWires,
      ∃ τ, H.wireType w = some τ ∧ 0 < Shape.totalSize τ.shape
  compPostWiresPos : ∀ c ∈ (List.finRange kk).map diss
        ++ (List.finRange kk).map mem,
      ∀ w ∈ c.post.freeWires,
      ∃ τ, H.wireType w = some τ ∧ 0 < Shape.totalSize τ.shape
\end{leansrc}
\declhead{BernsteinLogic.Composition}{smallGainN\_is\_syzygy}{\badge{BadgeGreen}{statement bound}}
\begin{leansrc}
/-- **The n-ary exact identity** (the producer warrant that `foldEq` is
always dischargeable from the row shapes): under the column condition, the
`(λ₁, …, λ_kk, 1, …, 1)`-weighted entry fold over the dissipation and
membership polynomials equals the conclusion polynomial
`∑ᵢ λᵢ·LVᵢ + θ·∑ⱼ λⱼ·Aⱼ` — the cross-supplies cancel identically. Reduces
by canonical-polynomial extensionality to per-key rational arithmetic:
distribute, exchange the double sum, apply the column condition. -/
theorem smallGainN_is_syzygy (kk : ℕ) (LV A Sg : Fin kk → Poly)
    (lam : Fin kk → ℚ) (B : Fin kk → Fin kk → ℚ) (th : ℚ)
    (hcol : ∀ j, ((List.finRange kk).map (fun i => B i j)).sum = lam j) :
    (((List.finRange kk).map
        (fun i => (lam i, Poly.sub (Poly.add (LV i) (A i)) (Sg i)))
      ++ (List.finRange kk).map
        (fun i => ((1 : ℚ), smallGainNMemPoly kk A Sg lam B th i))).foldr
      (fun pr acc => Poly.add (Poly.scale pr.1 pr.2) acc) [])
    = Poly.add
        ((List.finRange kk).foldr
          (fun i acc => Poly.add (Poly.scale (lam i) (LV i)) acc) [])
        (Poly.scale th
          ((List.finRange kk).foldr
            (fun i acc => Poly.add (Poly.scale (lam i) (A i)) acc) []))
\end{leansrc}
\declhead{BernsteinLogic.Composition}{sound\_smallGainN}{\badge{BadgeGreen}{statement bound}}
\begin{leansrc}
/-- **Soundness of SmallGainN**, as the Syzygy corollary: the `2·kk`
component validities feed `sound_syzygy` on the collapsed instance. The
rational comparisons (`λᵢ > 0`, `0 < θ < 1`) enter only through `λᵢ ≥ 0`
inside the combination check; `θ`'s range and the column condition are
well-formedness, exactly as the catalogue commentary states. -/
theorem sound_smallGainN {H : Hypernet} {kk : ℕ} (st : SmallGainNStep H kk)
    (hdiss : ∀ i, (st.diss i).semanticallyValid H)
    (hmem : ∀ i, (st.mem i).semanticallyValid H) :
    st.seq.semanticallyValid H
\end{leansrc}
\declhead{BernsteinLogic.Composition}{smallGainN\_colSum\_mutant\_refuted}{\badge{BadgeGreen}{statement bound}}
\begin{leansrc}
/-- **M-N2 (the column condition deleted): the identity is REFUTED** — by
the B-swap, computation-free: the mutant at `B = 0` and the PROVED identity
at `B = 1` (column sums check: `1 = λ`) share every other datum
(`kk = 1`, `LV = Sg = 0`, `A = const 1`, `λ = 1`, `θ = 0`), so their entry
folds would be equal; reading coefficients through the fold lemmas, the two
membership polynomials differ by exactly the decay budget's coefficient —
forcing `coeffSum (const 1) = 0` everywhere, against canonicity. The column
condition is what makes the membership rows deliver the stated conclusion
for every decay budget. -/
theorem smallGainN_colSum_mutant_refuted :
    ¬ (∀ (kk : ℕ) (LV A Sg : Fin kk → Poly) (lam : Fin kk → ℚ)
        (B : Fin kk → Fin kk → ℚ) (th : ℚ),
        (((List.finRange kk).map
            (fun i => (lam i, Poly.sub (Poly.add (LV i) (A i)) (Sg i)))
          ++ (List.finRange kk).map
            (fun i => ((1 : ℚ), smallGainNMemPoly kk A Sg lam B th i))).foldr
          (fun pr acc => Poly.add (Poly.scale pr.1 pr.2) acc) [])
        = Poly.add
            ((List.finRange kk).foldr
              (fun i acc => Poly.add (Poly.scale (lam i) (LV i)) acc) [])
            (Poly.scale th
              ((List.finRange kk).foldr
                (fun i acc => Poly.add (Poly.scale (lam i) (A i)) acc) [])))
\end{leansrc}
\declhead{BernsteinLogicAux.QtslCompose.AGComposeDiscrete}{smallGainN\_drift}{\badge{BadgeGreen}{statement bound}}
\begin{leansrc}
/-- **B3′, the cancellation row at any n** (the semantic face of the
checker-level `(λ₁, …, λₙ, 1, …, 1)` syzygy fold): per-component dissipation
rows with cross-supplies on the joint closed-loop kernel, under the n-ary
membership rows and the exact column condition, yield the `θ`-margin drift
row for `W = ∑ i, lam i * V i`. The cross-supplies cancel by construction:
`∑ i, lam i * sigma i ≤ (1-θ) ∑ i ∑ j, B i j * alpha j = (1-θ) ∑ j, lam j * alpha j`
(swap the double sum, then the column condition). -/
theorem smallGainN_drift (hlam : ∀ i, 0 ≤ lam i)
    (hcol : ∀ j, ∑ i, B i j = lam j)
    (hint : ∀ i, ∀ y, y ∉ A → Integrable (V i) (K y))
    (hd : ∀ i, ∀ y, y ∉ A →
        ∫ z, V i z ∂K y ≤ V i y - alpha i y + sigma i y)
    (hm : ∀ i, ∀ y, y ∉ A →
        lam i * sigma i y ≤ (1 - theta) * ∑ j, B i j * alpha j y) :
    ∀ y, y ∉ A → ∫ z, (∑ i, lam i * V i z) ∂K y
      ≤ (∑ i, lam i * V i y) - theta * ∑ j, lam j * alpha j y
\end{leansrc}
\declhead{BernsteinLogicAux.QtslCompose.AGComposeDiscrete}{smallGain\_drift\_of\_N}{\badge{BadgeGreen}{statement bound}}
\begin{leansrc}
/-- **Binary recovery, on the nose**: `smallGain_drift`'s exact statement
re-derived from `smallGainN_drift` at `m = 2`, `lam = ![1, λ]`,
`B = ![![0, λ], ![1, 0]]` (column sums `0 + 1 = 1` and `λ + 0 = λ`). One
bookkeeping disanalogy, inherited from the catalogue commentary: the n-ary
membership row at `i = 0` reads `1 * σ₁ ≤ (1-θ)·λ·α₂` — the binary `hm₂`
with its factors reassociated — and at `i = 1` reads `λ * σ₂ ≤ (1-θ)·α₁` =
the binary `hm₁`; the two rows arrive order-swapped vs the binary labels,
set-equal. -/
theorem smallGain_drift_of_N {V₁ V₂ alpha₁ alpha₂ sigma₁ sigma₂ : S → ℝ}
    {lam theta : ℝ} (hlam : 0 ≤ lam)
    (hint₁ : ∀ y, y ∉ A → Integrable V₁ (K y))
    (hint₂ : ∀ y, y ∉ A → Integrable V₂ (K y))
    (hd₁ : ∀ y, y ∉ A → ∫ z, V₁ z ∂K y ≤ V₁ y - alpha₁ y + sigma₁ y)
    (hd₂ : ∀ y, y ∉ A → ∫ z, V₂ z ∂K y ≤ V₂ y - alpha₂ y + sigma₂ y)
    (hm₁ : ∀ y, y ∉ A → lam * sigma₂ y ≤ (1 - theta) * alpha₁ y)
    (hm₂ : ∀ y, y ∉ A → sigma₁ y ≤ lam * ((1 - theta) * alpha₂ y)) :
    ∀ y, y ∉ A → ∫ z, (V₁ z + lam * V₂ z) ∂K y
      ≤ (V₁ y + lam * V₂ y) - theta * (alpha₁ y + lam * alpha₂ y)
\end{leansrc}
\declhead{BernsteinLogicAux.QtslCompose.AGComposeDiscrete}{avoid\_smallGainN}{\badge{BadgeGreen}{statement bound}}
\begin{leansrc}
/-- **B3′, assembled**: on the closed loop, the n-ary small-gain data yields
the composite avoid bound for the weighted-sum certificate (the discrete face
of the catalogue rule end-to-end; Foster–Lyapunov metatheory through
`Battery.measure_reach_le`, exactly as the binary `avoid_smallGain`). -/
theorem avoid_smallGainN (hlam : ∀ i, 0 ≤ lam i) (htheta : 0 ≤ theta)
    {C : Fin m → ℝ} (hA : MeasurableSet A)
    (hcol : ∀ j, ∑ i, B i j = lam j)
    (hV : ∀ i, Measurable (V i))
    (hV0 : ∀ i y, 0 ≤ V i y) (hVC : ∀ i y, V i y ≤ C i)
    (hbar : ∀ y ∈ A, 1 ≤ ∑ i, lam i * V i y)
    (halpha : ∀ i, ∀ y, y ∉ A → 0 ≤ alpha i y)
    (hd : ∀ i, ∀ y, y ∉ A →
        ∫ z, V i z ∂K y ≤ V i y - alpha i y + sigma i y)
    (hm : ∀ i, ∀ y, y ∉ A →
        lam i * sigma i y ≤ (1 - theta) * ∑ j, B i j * alpha j y) (x : S) :
    pathMeasure K x {ω | ∃ n, ω n ∈ A}
      ≤ ENNReal.ofReal (∑ i, lam i * V i x)
\end{leansrc}
\ruleblock{Correspondence (Steele display $\leftrightarrow$ Lean)}
\absentbox{NO CORRESPONDENCE DICTIONARY YET}{no reviewer-certified dictionary maps the displayed rule onto the mechanized statements above; until one is written and certified, that mapping rests on this section\textquotesingle s editorial care alone.}
\ruleblock{Lean verifier receipts}
\sourcenote{module \texttt{BernsteinLogic.Composition}; bound build \texttt{9116f454beb0\ldots} (bound from receipt 548ea3dda8ed73a2, 2026-08-03T21:30:29Z). A receipt records that named checking layers accepted a build: (module, compiled-artifact SHA-256, environment SHA-256, checker layers), minted when the checking gate ran; the store receipts attesting the bound build are listed with the layers each run credited. Nothing here asks the reader to trust the store: every layer is an independently runnable program over the quoted sources --- rebuild the module and re-run the checkers to reproduce the verdict. Receipts attesting superseded builds credit no layer and are omitted from this rendering (the count of omissions, if any, is noted below).}
\begin{receipttable}
\receiptrow{548ea3dda8ed73a2}{2026-08-03T21:30:29Z}{build+decorrelated}{\badge{BadgeGreen}{current}}
\end{receipttable}
\sourcenote{10 further receipt(s) in the store attest superseded builds of this module and are omitted from this rendering.}
\sourcenote{module \texttt{BernsteinLogicAux.QtslCompose.AGComposeDiscrete}; bound build \texttt{d5233632f95a\ldots} (bound from receipt 88ff74180adf5b95, 2026-08-03T20:13:53Z). A receipt records that named checking layers accepted a build: (module, compiled-artifact SHA-256, environment SHA-256, checker layers), minted when the checking gate ran; the store receipts attesting the bound build are listed with the layers each run credited. Nothing here asks the reader to trust the store: every layer is an independently runnable program over the quoted sources --- rebuild the module and re-run the checkers to reproduce the verdict. Receipts attesting superseded builds credit no layer and are omitted from this rendering (the count of omissions, if any, is noted below).}
\begin{receipttable}
\receiptrow{88ff74180adf5b95}{2026-08-03T20:13:53Z}{build+decorrelated}{\badge{BadgeGreen}{current}}
\end{receipttable}
\sourcenote{11 further receipt(s) in the store attest superseded builds of this module and are omitted from this rendering.}
\layerline{\layeritem{build}{lake build (elaboration)}{\badge{BadgeGreen}{current}}\layeritem{replay}{leanchecker (official kernel replay)}{\badge{BadgeAmber}{stale only}}\layeritem{hardened}{safe\_verify (hardened screen)}{\badge{BadgeAmber}{stale only}}\layeritem{decorrelated}{lean4export + nanoda (decorrelated kernel)}{\badge{BadgeGreen}{current}}}
\sourcenote{\badge{BadgeGray}{comparator: structurally n/a} the alias unpacks
\emph{before} emission (no rule tag; producers elaborate a
\rulename{SmallGainN} step to its \rulename{Syzygy} node), so no wire
certificate --- and hence no comparator fixture --- can ever contain a
\rulename{SmallGainN} node. Comparator coverage for this rule \emph{is}
\rulename{Syzygy}\textquotesingle s.
}
\ruleblock{HOL4 statement}
\absentbox{OUT OF SCOPE BY DESIGN}{Out of scope for the HOL side \emph{by design}, not by backlog, and for
exactly the binary rule\textquotesingle s reason: the wire format never
carries a \rulename{SmallGainN} node --- producers elaborate to the
$(\lambda_1, \ldots, \lambda_n, 1, \ldots, 1)$ \rulename{Syzygy} node
before emission --- and the HOL4 development\textquotesingle s object
is the wire checker. The HOL statement that verifies a composed
$n$-component small-gain certificate is therefore
\rulename{Syzygy}\textquotesingle s (\texttt{syzygyCheck}).
}
\ruleblock{HOL verifier receipts}
\absentbox{NONE}{no HOL verification receipt exists for this rule.}
\ruleblock{Proof}

\sourcenote{Genre: alias rule (tier 3) --- a respelling template. The
five artefacts an alias owes the book: the 1$\to$1 unpack, the
filled-in application of the target rule, the payload-transport table,
the exact algebraic identity that makes the respelling correct, and the
soundness corollary through the target rule's lemma. All five follow;
the unpack and transport table are stated here as the producer's
contract, their dedicated mechanized vehicle being the one artefact not
yet in the Lean tree (stated loudly below and in the section's Lean
cell) --- soundness of every emitted node meanwhile rests on the target
rule's already-mechanized lemma, because the alias unpacks before
emission.}

\begin{brecipe}[the unpack --- one node, any $n$]\label{rec:sgn-unpack}
A \rulename{SmallGainN} step carries $2n$ component proofs: the $n$
dissipation sequents (each over its own box $\Phi_i$) and the $n$
membership sequents (over the declared operating box $\Phi_w$), with
side data $\lambda \in \Q_{>0}^{n}$, $B \in \Q_{\geq 0}^{n \times n}$
with $\sum_i B_{ij} = \lambda_j$ for every $j$, and $\theta \in \Q \cap
(0,1)$. Define the atom terms
\[
e_i \equiv \mathcal{L}_i V_i + \alpha_i - \sigma_i,
\qquad
m_i \equiv \lambda_i \sigma_i - (1{-}\theta) \sum_j B_{ij}\, \alpha_j ,
\]
and the substitution into \rulename{Syzygy}'s metavariables:
\[
N := 2n, \qquad
(e_1, \ldots, e_N) := (e_1, \ldots, e_n, m_1, \ldots, m_n), \qquad
(\Phi_1, \ldots, \Phi_N) := (\Phi_1, \ldots, \Phi_n,
  \underbrace{\Phi_w, \ldots, \Phi_w}_{n}),
\]
\[
(\lambda_1, \ldots, \lambda_N) := (\lambda_1, \ldots, \lambda_n,
  1, \ldots, 1), \qquad
\{g_k\}_k := \varnothing, \qquad
c := \mathcal{L}\bigl(\textstyle\sum_i \lambda_i V_i\bigr)
  + \theta \textstyle\sum_j \lambda_j \alpha_j .
\]
The repayment matrix $B$ is consumed here and does not appear in the
image: it is the producer's bookkeeping for \emph{finding} membership
rows that make the identity of
Proposition~\ref{prop:sgn-identity} true, and the identity is the only
trace of it the checker ever replays.
\end{brecipe}

\begin{brecipe}[the filled-in application]\label{rec:sgn-filled}
Under the substitution of Recipe~\ref{rec:sgn-unpack} the
\rulename{SmallGainN} node \emph{is} the following \rulename{Syzygy}
application:
\begin{mathpar}
\inferrule*[Right=\rulename{Syzygy}]{
  \bigl\{\Phi_i \vdash e_i \leq 0\bigr\}_{i=1}^{n} \\
  \bigl\{\Phi_w \vdash m_i \leq 0\bigr\}_{i=1}^{n} \\\\
  (\lambda_1, \ldots, \lambda_n, 1, \ldots, 1) \in \Q_{\geq 0}^{\,2n} \\
  c \;\equiv\; \textstyle\sum_i \lambda_i\, e_i + \sum_i 1 \cdot m_i
  \ \text{(exact canonical polynomial)}
}{
  \Phi \vdash c \leq 0
}
\end{mathpar}
with $\Phi \supseteq \Phi_1, \ldots, \Phi_n, \Phi_w$ conjunct-wise (the
$2n$-way containment split). Discharging \rulename{Syzygy}'s
obligations at this instance: the nonnegativity checks are $\lambda_i
\geq 0$ (from $\lambda_i > 0$ --- the only use soundness makes of the
signs) and $1 \geq 0$, $n$ times each; the sourcing check is immediate
(each of the $2n$ entries is the corresponding component's posted atom,
verbatim); the exact-identity check is
Proposition~\ref{prop:sgn-identity}; the index-freeness gates discharge
by recomputation on the $2n + 1$ terms; and the shape, positivity, and
zero-normal \texttt{postForm} obligations transport exactly as in the
binary section. The $\nu$-row is empty for the same reason as there:
the cross-supply cancellation is performed by the membership atoms
themselves, so no equality sub-conjuncts of $\Phi$ enter.
\end{brecipe}

\begin{bprop}[the conclusion identity]\label{prop:sgn-identity}
Under the column condition $\sum_i B_{ij} = \lambda_j$ (all $j$), as an
identity of exact canonical polynomials,
\[
\sum_i \lambda_i e_i + \sum_i m_i
\;=\;
\mathcal{L}\bigl(\textstyle\sum_i \lambda_i V_i\bigr)
\;+\; \theta \sum_j \lambda_j \alpha_j .
\]
\end{bprop}
\begin{proof}
Expand and collect --- still two lines at any $n$, which is the point:
\begin{align*}
\sum_i \lambda_i e_i + \sum_i m_i
&= \sum_i \lambda_i \mathcal{L}_i V_i
 + \sum_i \lambda_i \alpha_i
 \underbrace{- \sum_i \lambda_i \sigma_i
 + \sum_i \lambda_i \sigma_i}_{=\,0\ \text{(cross-supplies cancel)}}
 - (1{-}\theta) \sum_j \Bigl(\sum_i B_{ij}\Bigr) \alpha_j \\
&= \mathcal{L}\bigl(\textstyle\sum_i \lambda_i V_i\bigr)
 + \sum_j \bigl(1 - (1{-}\theta)\bigr) \lambda_j \alpha_j
\;=\; \mathcal{L}\bigl(\textstyle\sum_i \lambda_i V_i\bigr)
 + \theta \sum_j \lambda_j \alpha_j ,
\end{align*}
where the second line swaps the summation order in the $B$-term and
applies the column condition $\sum_i B_{ij} = \lambda_j$ --- the one
place it is used, and exactly the discount that makes the neighbour
draws rejoin the decay budgets at weight $\lambda_j$. As in the binary
section, reading $\mathcal{L}(\sum_i \lambda_i V_i)$ for $\sum_i
\lambda_i \mathcal{L}_i V_i$ is a definition, not a step: over disjoint
variable sets the composite generator acts coordinate-wise. The checker
replays this as \rulename{Syzygy}'s exact canonical-polynomial identity
and never sees a fixpoint, a spectral radius, or the interconnection
graph --- the $n$-component small-gain argument is dissolved into ring
arithmetic at certificate-checking time.
\end{proof}

\begin{bprop}[binary recovery, on the nose]\label{prop:sgn-binary}
At $n = 2$ with $\lambda = (1, \lambda)$, $B_{21} = 1$, $B_{12} =
\lambda$ (and $B_{11} = B_{22} = 0$), the column condition holds, the
membership rows are exactly the binary section's
$\{\,\lambda\sigma_2 - (1{-}\theta)\alpha_1,\;
\sigma_1 - \lambda(1{-}\theta)\alpha_2\,\}$, the conclusion is exactly
$\mathcal{L}(V_1 + \lambda V_2) + \theta(\alpha_1 + \lambda\alpha_2)
\leq 0$, and the \rulename{Syzygy} weights are $(1, \lambda, 1, 1)$ ---
i.e.\ \rulename{SmallGain} is the $n = 2$ instance of this rule.
\end{bprop}
\begin{proof}
Column sums: $B_{11} + B_{21} = 1 = \lambda_1$ and $B_{12} + B_{22} =
\lambda = \lambda_2$. Membership rows: $m_1 = \sigma_1 -
\lambda(1{-}\theta)\alpha_2$ and $m_2 = \lambda\sigma_2 -
(1{-}\theta)\alpha_1$ --- the binary section's rows as a \emph{set},
arriving in the opposite order (its $m_1$ is this rule's $m_2$): a
bookkeeping difference in combination order, visible in the node bytes
and irrelevant to validity, since \rulename{Syzygy}'s identity is a
sum. Conclusion and weights read off directly.
\end{proof}

\begin{brecipe}[payload transport --- the producer's contract]
\label{rec:sgn-transport}
The $n$-ary payload transports into \texttt{SyzygyStep} exactly as the
binary payload does, with the four fixed slots generalized to $n$-vectors:
\begin{center}\small
\begin{tabular}{@{}ll@{}}
\textbf{\rulename{SmallGainN} step (contract)} & \textbf{\texttt{SyzygyStep}}\\
\midrule
\texttt{seq} & \texttt{seq}\\
$\texttt{diss}_1, \ldots, \texttt{diss}_n, \texttt{mem}_1, \ldots, \texttt{mem}_n$ & \texttt{components} (that order = combination order)\\
weights $(\lambda_1, \ldots, \lambda_n, 1, \ldots, 1)$ & \texttt{combination}\\
\texttt{concl} & \texttt{conclusion}\\
$2n$-way containment split & \texttt{preContains}\\
shape/positivity/idx-freeness packs & verbatim, as in the binary table\\
$(\lambda, B, \theta)$ & consumed by the unpack; $B$ leaves no image field\\
\texttt{postEq} & \texttt{postForm} $=$ the zero-normal reading\\
\end{tabular}
\end{center}
This table is the contract of the dedicated mechanized vehicle --- the
$n$-ary payload structure with its verified unpack and soundness
corollary, the binary section's \texttt{SmallGainStep} apparatus at
arbitrary $n$ --- which is not yet in the Lean tree (the section's Lean
cell states the absence). Nothing checked on the wire waits on it: the
unpack happens producer-side, and the emitted node is
\rulename{Syzygy}'s.
\end{brecipe}

\begin{bprop}[soundness corollary]\label{prop:sgn-sound}
If the $2n$ component sequents are semantically valid, the composite
sequent $\Phi \vdash \mathcal{L}(\sum_i \lambda_i V_i) + \theta \sum_j
\lambda_j \alpha_j \leq 0$ is semantically valid.
\end{bprop}
\begin{proof}
The target rule's kernel lemma applied to the unpacked instance of
Recipe~\ref{rec:sgn-filled}, feeding the $2n$ validities through the
containment split; the exact-identity obligation is
Proposition~\ref{prop:sgn-identity}. Beyond that routing there is no
further argument --- the alias genre working as intended, at every $n$
at once.
\end{proof}

\begin{bremark}[what stays outside the checker]
$(\lambda, B, \theta)$-\emph{existence} is the producer's problem, and
the classical spectral conditions are its solution theory: when the
gain matrix $\Gamma$ has spectral radius $\rho(\Gamma) < 1 - \theta$,
the Neumann-series vector $\lambda := \sum_k (1-\theta)^{-k}
(\Gamma^{\mathsf{T}})^k \mathbb{1}$ converges and witnesses
$\Gamma^{\mathsf{T}} \lambda \leq (1-\theta) \lambda$, from which a
feasible $B$ is assembled column by column (under-allocation is
repaired by diagonal padding, which only weakens the membership row it
feeds). The $n$-component ISS small-gain theorem for general networks
is Dashkovskiy--R\"uffer--Wirth \cite{dashkovskiyrufferwirth2007};
where that tradition closes a monotone-operator fixpoint argument over
the interconnection graph, the certificate form closes the ring
identity of Proposition~\ref{prop:sgn-identity}, and the graph
combinatorics survive only as the producer's search space. The signed
polynomial contract discipline is the binary section's, unchanged.
\end{bremark}

\begin{anchors}
Recipe \ref{rec:sgn-unpack} unpack & producer-side elaboration; dedicated mechanized vehicle = named successor (binary vehicle: \texttt{SmallGainStep.toSyzygy})\\
Proposition \ref{prop:sgn-identity} & \texttt{combValid}'s fold identity, replayed as \rulename{Syzygy}'s $(\ast)$ on the emitted node\\
Proposition \ref{prop:sgn-binary} & the binary section's \texttt{SmallGainStep} apparatus at $n = 2$ (weights $(1, \lambda, 1, 1)$)\\
Proposition \ref{prop:sgn-sound} & \texttt{sound\_syzygy} on the unpacked instance ($n$-ary corollary = named successor)\\
no wire tag, by design & \texttt{RuleTag} comment block (Core/Catalogue.lean): producers elaborate before emission\\
\end{anchors}


\rulesection{36}{ProductLattice}{admissible rule}
\label{blcatrule-ProductLattice}
\ruleblock{The rule (Steele metanotation)}
\sourcenote{as realized in the mechanized development; source row: \texttt{\texttt{rules.toml}} --- \emph{ProductLattice}. The metanotation statement is typeset by hand for notational quality; its mechanized ground truth is the Lean statement below, extracted from the sources at every build of this book and hash-bound. The rule \emph{roster} is machine-checked (every mechanized rule tag round-trips to exactly one catalogue row); the statement \emph{text} is editorial, checked by eye against the extracted Lean beside it.}
\begin{mathpar}
\inferrule*[Right=\rulename{ProductLattice}]{
  B_1 \ \text{a certified attracting-block lattice for } \varphi_1 \ \text{over } H_1
  \\
  B_2 \ \text{likewise for } \varphi_2 \ \text{over } H_2
  \\
  \mathrm{wires}(H_1) \cap \mathrm{wires}(H_2) = \emptyset
}{
  \{\, b_1 \times b_2 \mid b_1 \in B_1,\ b_2 \in B_2 \,\} \ \text{(the meet-generators) certified for } \varphi_1 \times \varphi_2 \ \text{over } H_1 \sqcup H_2
}
\end{mathpar}
the displayed family is the list of pairwise-product meet-generators.

\ruleflavor{the full lattice $\langle B_1 \times B_2 \rangle_{\vee}$ (join-closure of the pairwise products) is the paper-strength reading --- the bound theorem quantifies over the meet-generators, and the finite-union/boundary-decomposition step remains paper-strength, exactly as the correspondence concedes. Every composite facet sequent is discharged BY CITATION: a composite facet is normal to one factor\textquotesingle s axis, its inwardness polynomial IS that factor\textquotesingle s facet polynomial with the other factor\textquotesingle s coordinates riding along as free box variables (the rider, conjoined onto the premise). Bernstein-wise the composite facet grid is the factor grid tensored with all-ones --- a coefficient identity; nothing for a checker to recompute beyond the ordinary \rulename{BernCone} recompute pattern. Declared in the baked base theory itself (its first certificate-family rule), not in any separate composition layer.}
\par\medskip
\rulename{ProductLattice} composes at the certificate-\emph{family}
level, one judgment up from the sequent rules: given certified
attracting-block lattices for two systems over wire-disjoint hypernets,
it certifies the join-closure of pairwise products as an
attracting-block lattice for the product system over the disjoint
union. The geometric fact that powers it: every internal facet of an
axis-aligned product block is normal to one factor\textquotesingle s
axis, so its inwardness polynomial \emph{is} that
factor\textquotesingle s facet polynomial, with the other
factor\textquotesingle s coordinates riding along as free box
variables. Each composite facet sequent is therefore discharged
\emph{by citation}: conclusion kept verbatim, the other
factor\textquotesingle s box conjoined onto the premise as a rider.

What the checker recomputes: no new Bernstein arithmetic ---
Bernstein-wise the composite facet grid is the factor grid tensored
with all-ones, a coefficient identity, so the per-facet discharge is
syntactic identity of the cited conclusion. What remains checked is
the citation plumbing itself: wire-disjointness of the two hypernets,
factor scopedness, and the rider conjunction on each cited premise.

Wire-disjointness is load-bearing twice. Structurally, the disjoint
union of the two hypernets is only well-formed without shared wires
(linearity and acyclicity of the union are proven \emph{from}
disjointness). Semantically, no single transported facet sequent\textquotesingle s
validity fails under a shared wire --- the weaken lemma conjoins any
rider --- but the assembled object stops being a product of the
factors: the identification of every composite facet with a cited
factor facet, the other factor entering only as a rider over fresh
coordinates, is proven from disjointness, and with a shared wire the
recomputed readout misdescribes the composite. A single shared wire is
precisely an interface, and composition along an interface is
\rulename{Promise}/small-gain territory, not this
rule\textquotesingle s.

\cataloguenote{The soundness statement is dynamics-free --- sound for flows and maps
alike, because it certifies a \emph{lower bound} family (\textquotedblleft
these blocks are certified\textquotedblright), never \textquotedblleft
these are all the attractors\textquotedblright. The completeness
reading --- that the generated lattice is the whole attractor lattice
of the product --- is a flows-only theorem of the source literature
and lives paper-side; the proof section exhibits the map-dynamics
counterexample that forces the fence. No wire tag is allocated yet:
soundness lands first, executable checking follows the family-object
wire-format work, and the emit path refuses the rule until then.
}
\ruleblock{Lean statement}
\par\medskip
\textbf{What \textquotedblleft discharge by
citation\textquotedblright{} means, stated in the flavor text and
bound by the declarations below.}\enspace \texttt{citeFacet} keeps the cited
factor facet\textquotesingle s conclusion verbatim and conjoins the
rider onto the premise; \texttt{prodBlock} assembles a product
block\textquotesingle s facet list as exactly the two citation
families; \texttt{prodBlocks} enumerates only the meet-generators of
the join-closure --- the exponentially many join blocks are never
materialized; their sub-facets discharge through the paper-strength
boundary decomposition (which selects a factor facet and rider), whose
final transport-and-weaken step is the same citation lemmas with
smaller riders. A skeptical reader should confirm by eye
that no definition below evaluates dynamics anywhere: the rule\textquotesingle s
semantic content is entirely the transport-and-weaken pair quoted
above.

\declhead{BernsteinLogic.ProductLattice}{BlockCert}{\badge{BadgeGreen}{statement bound}}
\begin{leansrc}
/-- One certified attracting block, certificate view (the recon's Route-2
shape): its region — the compact axis-aligned box(-union), as the formula
the producer wrote it with — and its internal-facet inwardness sequents
(strict inwardness via non-strict predicates with explicit positive
rational margins; block-hood is elementary for strictly inward box unions,
no combinatorial index data claimed — note §1.1). -/
structure BlockCert where
  region : Formula
  facets : List Sequent
\end{leansrc}
\declhead{BernsteinLogic.ProductLattice}{CertifiedBlockLattice}{\badge{BadgeGreen}{statement bound}}
\begin{leansrc}
/-- A certified attracting-block lattice over a hypernet — the catalogue's
first certificate-FAMILY object (premise object of ProductLattice).
The family carries its lattice structure explicitly at the index level:
`join`/`meet` with the six standard lattice laws (idempotence is
derivable). These are the certificate-level shadow of the note §1.1
closure ("closed under ∪ and regularized ∩"); that a join index denotes
the set-union of the denoted blocks is producer-side geometry in v1 —
kernel teeth for the denotation arrive with the family-object wire format
(C-4/format cycle). Wire confinement makes the family transportable into
any extension (§3) — it is what `sound_productLattice` consumes. -/
structure CertifiedBlockLattice (H : Hypernet) where
  blocks : List BlockCert
  join : Fin blocks.length → Fin blocks.length → Fin blocks.length
  meet : Fin blocks.length → Fin blocks.length → Fin blocks.length
  join_comm : ∀ i j, join i j = join j i
  join_assoc : ∀ i j k, join (join i j) k = join i (join j k)
  meet_comm : ∀ i j, meet i j = meet j i
  meet_assoc : ∀ i j k, meet (meet i j) k = meet i (meet j k)
  absorb_join : ∀ i j, join i (meet i j) = i
  absorb_meet : ∀ i j, meet i (join i j) = i
  regionConfined : ∀ b ∈ blocks, ∀ w ∈ b.region.freeWires, (H.wireType w).isSome
  facetsScoped : ∀ b ∈ blocks, ∀ s ∈ b.facets, s.wellScoped H
\end{leansrc}
\declhead{BernsteinLogic.ProductLattice}{PProductLatticeData}{\badge{BadgeGreen}{statement bound}}
\begin{leansrc}
/-- Applicability data of the ProductLattice rule — the two factor
families and the side condition. Field-for-field what a checker would
verify: the factor certificates (whose facet-sequent validity enters
`sound_productLattice` as hypotheses, Step-style) and wire-disjointness.
No first-class `RuleTag`/wire constructor in v1 (note §1.4: emission-layer
posture; BINARY.md §6.2 reserves tag 27 for the family-object format
cycle); correspondingly NOT in `v0frag` — like the ReLU family
pre-enactment, soundness lands first and executable checking follows the
format work. -/
structure PProductLatticeData (H1 H2 : Hypernet) where
  disjoint : H1.WireDisjoint H2
  left : CertifiedBlockLattice H1
  right : CertifiedBlockLattice H2
\end{leansrc}
\declhead{BernsteinLogic.ProductLattice}{Sequent.semanticallyValid\_extend}{\badge{BadgeGreen}{statement bound}}
\begin{leansrc}
/-- TRANSPORT: semantic validity is monotone under hypernet extension —
a sequent valid over H, well-scoped in H, is valid over any H' that
extends H's wire typing and contains H's applications. (For ProductLattice
this is "the factor's certified sign is the composite's inwardness": the
facet polynomial reads only the factor's wires, so the other factor's
coordinates ride along invisibly.) -/
theorem Sequent.semanticallyValid_extend {H H' : Hypernet} {seq : Sequent}
    (hext : H.WireExtends H')
    (happs : ∀ app ∈ H.applications, app ∈ H'.applications)
    (hws : seq.wellScoped H) (hval : seq.semanticallyValid H) :
    seq.semanticallyValid H'
\end{leansrc}
\declhead{BernsteinLogic.ProductLattice}{Sequent.semanticallyValid\_weaken}{\badge{BadgeGreen}{statement bound}}
\begin{leansrc}
/-- WEAKEN, semantic and standalone: any conjunct may join a valid
sequent's premise (PWeaken's semantic content — the rider needs no
side condition). -/
theorem Sequent.semanticallyValid_weaken {H : Hypernet}
    {pre extra post : Formula}
    (h : Sequent.semanticallyValid H ⟨pre, post⟩) :
    Sequent.semanticallyValid H ⟨.and pre extra, post⟩
\end{leansrc}
\declhead{BernsteinLogic.ProductLattice}{sound\_productLattice\_citeLeft}{\badge{BadgeGreen}{statement bound}}
\begin{leansrc}
/-- The standing citation discharge, left factor: ANY sequent citing a
valid, well-scoped left-factor facet — conclusion verbatim, premise
= factor premise ∧ an arbitrary rider — is valid over the union.
Soundness is transport (§3) then weaken: the note §1.2's one-line
argument, in two lines. A join block's sub-facet, once the paper-strength
boundary decomposition has selected its factor facet and rider, is
discharged by this lemma (rider = the sub-facet's free-coordinate cell);
the decomposition itself is not mechanized here. -/
theorem sound_productLattice_citeLeft {H1 H2 : Hypernet}
    (hd : H1.WireDisjoint H2) {s0 : Sequent} (rider : Formula)
    (hws : s0.wellScoped H1) (hval : s0.semanticallyValid H1) :
    (citeFacet s0 rider).semanticallyValid (H1.disjUnion H2 hd)
\end{leansrc}
\declhead{BernsteinLogic.ProductLattice}{sound\_productLattice\_citeRight}{\badge{BadgeGreen}{statement bound}}
\begin{leansrc}
/-- The standing citation discharge, right factor. -/
theorem sound_productLattice_citeRight {H1 H2 : Hypernet}
    (hd : H1.WireDisjoint H2) {s0 : Sequent} (rider : Formula)
    (hws : s0.wellScoped H2) (hval : s0.semanticallyValid H2) :
    (citeFacet s0 rider).semanticallyValid (H1.disjUnion H2 hd)
\end{leansrc}
\declhead{BernsteinLogic.ProductLattice}{sound\_productLattice}{\badge{BadgeGreen}{statement bound}}
\begin{leansrc}
/-- SOUNDNESS of ProductLattice (note §1.2, written properly): given two
certified block lattices over wire-disjoint hypernets, every facet
sequent of every product block in the conclusion family is semantically
valid over the disjoint union. Each facet is a factor facet discharged by
citation; validity = transport of the factor's validity along the union
injection (§3), then weakening in the rider. Dynamics-free — sound for
flows AND maps; the flows-only fence (header) concerns completeness,
which this statement deliberately does not express: the conclusion is a
certified LOWER BOUND family, and "the whole attractor lattice" is a
paper-side, flows-only reading (KW `cor_flows_att_iso`). -/
theorem sound_productLattice {H1 H2 : Hypernet}
    (d : PProductLatticeData H1 H2)
    (hL : d.left.Valid) (hR : d.right.Valid) :
    ∀ b ∈ d.prodBlocks, ∀ s ∈ b.facets,
      s.semanticallyValid (H1.disjUnion H2 d.disjoint)
\end{leansrc}
\declhead{BernsteinLogic.ProductLattice}{Hypernet.disjUnion}{\badge{BadgeGreen}{statement bound}}
\begin{leansrc}
/-- The disjoint union of two wire-disjoint hypernets: union typing,
concatenated applications; the three well-formedness predicates are
proven from the factors' plus disjointness. This is the composite system
φ₁ × φ₂'s hypernet, H₁ ⊔ H₂ (note §1.1). -/
def Hypernet.disjUnion (H1 H2 : Hypernet) (hd : H1.WireDisjoint H2) : Hypernet
\end{leansrc}
\declhead{BernsteinLogic.ProductLattice}{Hypernet.Valuation.restrict}{\badge{BadgeGreen}{statement bound}}
\begin{leansrc}
/-- Restrict a valuation of an extending hypernet to the extended one. -/
def Hypernet.Valuation.restrict {H H' : Hypernet} (σ : H'.Valuation)
    (hext : H.WireExtends H') : H.Valuation
\end{leansrc}
\declhead{BernsteinLogic.ProductLattice}{Formula.holds\_restrict}{\badge{BadgeGreen}{statement bound}}
\begin{leansrc}
/-- Formulas scoped to the extended hypernet's wires hold identically
under restriction. -/
theorem Formula.holds_restrict {H H' : Hypernet} (σ : H'.Valuation)
    (hext : H.WireExtends H') :
    ∀ (F : Formula), (∀ w ∈ F.freeWires, (H.wireType w).isSome) →
      (Formula.holds (σ.restrict hext) F ↔ Formula.holds σ F)
\end{leansrc}
\declhead{BernsteinLogic.ProductLattice}{PrimApp.semOk\_restrict}{\badge{BadgeGreen}{statement bound}}
\begin{leansrc}
/-- Per-application semantics survive restriction: every entry SemOk reads
is at a port wire, typed in the extended hypernet by type consistency. -/
theorem PrimApp.semOk_restrict {H H' : Hypernet} (σ : H'.Valuation)
    (hext : H.WireExtends H') {app : PrimApp}
    (hTC : app.TypeConsistent H.wireType) (h : app.SemOk σ) :
    app.SemOk (σ.restrict hext)
\end{leansrc}
\declhead{BernsteinLogic.ProductLattice}{prodBlocks\_confined}{\badge{BadgeGreen}{statement bound}}
\begin{leansrc}
/-- The conclusion family is itself wire-confined over the union — the
confinement half of what a NESTED application (product of products,
AG-inside/product-outside, note §2) consumes from its premise object. -/
theorem prodBlocks_confined {H1 H2 : Hypernet}
    (d : PProductLatticeData H1 H2) :
    ∀ b ∈ d.prodBlocks,
      (∀ w ∈ b.region.freeWires,
        ((H1.disjUnion H2 d.disjoint).wireType w).isSome) ∧
      ∀ s ∈ b.facets, s.wellScoped (H1.disjUnion H2 d.disjoint)
\end{leansrc}
\declhead{BernsteinLogic.ProductLattice}{Hypernet.WireDisjoint}{\badge{BadgeGreen}{definition bound}}
\begin{leansrc}
/-- No wire is typed in both hypernets — ProductLattice's side condition
(note §1.1: wires(H₁) ∩ wires(H₂) = ∅). A single shared wire demotes the
composition to the interface axis (Promise/SmallGain — note §2). -/
def Hypernet.WireDisjoint (H1 H2 : Hypernet) : Prop :=
  ∀ w : Wire, (H1.wireType w).isSome → (H2.wireType w).isSome → False
\end{leansrc}
\declhead{BernsteinLogic.ProductLattice}{CertifiedBlockLattice.Valid}{\badge{BadgeGreen}{definition bound}}
\begin{leansrc}
/-- The family is certified when every block's every facet sequent is
semantically valid. Validity enters `sound_productLattice` as a
hypothesis — the Step pattern: sub-proof warrants are supplied at the
use site (by `sound_v0` on the factor certificates, or manually). -/
def CertifiedBlockLattice.Valid {H : Hypernet}
    (L : CertifiedBlockLattice H) : Prop :=
  ∀ b ∈ L.blocks, ∀ s ∈ b.facets, s.semanticallyValid H
\end{leansrc}
\declhead{BernsteinLogic.ProductLattice}{citeFacet}{\badge{BadgeGreen}{definition bound}}
\begin{leansrc}
/-- Discharge BY CITATION (note §1.1): the composite facet sequent keeps
the cited factor facet's conclusion VERBATIM — polynomial and predicate
unchanged; the Bernstein grid of the composite facet is the factor's grid
replicated along the free coordinates, a coefficient identity, so
syntactic identity of the conclusion is the whole discharge — and
conjoins the rider (the other factor's box, riding along as free box
variables) onto the premise. -/
def citeFacet (s : Sequent) (rider : Formula) : Sequent :=
  ⟨.and s.pre rider, s.post⟩
\end{leansrc}
\declhead{BernsteinLogic.ProductLattice}{prodBlock}{\badge{BadgeGreen}{definition bound}}
\begin{leansrc}
/-- The product of two blocks: region = the conjunction of boxes; facets =
each factor's facets, cited with the other factor's region as rider.
(Every internal facet of an axis-aligned product is normal to one
factor's axis — note §1.2.) -/
def prodBlock (b1 b2 : BlockCert) : BlockCert :=
  ⟨.and b1.region b2.region,
   b1.facets.map (fun s => citeFacet s b2.region)
     ++ b2.facets.map (fun s => citeFacet s b1.region)⟩
\end{leansrc}
\declhead{BernsteinLogic.ProductLattice}{PProductLatticeData.prodBlocks}{\badge{BadgeGreen}{definition bound}}
\begin{leansrc}
/-- The conclusion family's meet-generators: all pairwise products.
⟨B₁ × B₂⟩∨ = the join-closure of THIS list (KW Prop
`prop:Lwedgevee-generated`); the joins — exponentially many — are never
enumerated: their sub-facets discharge via the paper-strength boundary
decomposition (unmechanized — it selects the factor facet and rider),
whose final step is `sound_productLattice_citeLeft/Right` below with
smaller riders. -/
def PProductLatticeData.prodBlocks {H1 H2 : Hypernet}
    (d : PProductLatticeData H1 H2) : List BlockCert :=
  d.left.blocks.flatMap fun b1 => d.right.blocks.map fun b2 => prodBlock b1 b2
\end{leansrc}
\ruleblock{Correspondence (Steele display $\leftrightarrow$ Lean)}
\badge{BadgeGreen}{reviewer-certified against the displayed texts}
\sourcenote{The dictionary below is a reviewer-certified statement of how the mechanized statements above interpret the displayed rule: display metavariable $\mapsto$ Lean binder, each display hypothesis $\mapsto$ the premise or side condition carrying it, the display conclusion $\mapsto$ the goal --- with any deliberate difference in width between display and mechanization named outright. Certification is recorded as a hash triple over (the displayed cell, the extracted statement texts, this dictionary\textquotesingle s own text); if any of the three changes, this block turns amber until a reviewer re-certifies. The correspondence claim itself remains editorial --- reviewer-grade, not machine-checked; its machine-checked successor types each soundness statement against the rule as a structured object, and supersedes this block per rule as it lands.}
The one rule of this book whose judgment sits a layer up: the
conclusion is not a sequent but a certified \emph{family}, so the
\rulename{Ax} frame applies one level in --- per facet sequent.
\begin{itemize}
\item $B_1, B_2 \mapsto$ \texttt{d.left}, \texttt{d.right} $:$
  \texttt{CertifiedBlockLattice} over \texttt{H1}, \texttt{H2};
  ``certified for $\varphi_i$\textquotesingle\textquotesingle{}
  $\mapsto$ \texttt{hL}, \texttt{hR} $:$ \texttt{Valid} --- every
  block\textquotesingle s every facet sequent semantically valid.
  Note what does \emph{not} appear: any semiflow. The certified
  statement is dynamics-free (the note above) --- the
  display\textquotesingle s $\varphi_1 \times \varphi_2$ is the
  intended reading of the facet sequents, not a consumed hypothesis.
\item $\mathrm{wires}(H_1) \cap \mathrm{wires}(H_2) = \emptyset$
  $\mapsto$ \texttt{d.disjoint} $:$ \texttt{H1.WireDisjoint H2}, and
  the composite carrier $H_1 \sqcup H_2 \mapsto$
  \texttt{Hypernet.disjUnion}.
\item the displayed conclusion (the meet-generator family
  certified) $\mapsto$ \texttt{sound\_productLattice}, which
  discharges every facet
  of every \emph{pairwise product} block (\texttt{prodBlocks} --- the
  join-closure\textquotesingle s meet-generators) --- the displayed
  family and the bound quantification now coincide. The full
  join-closure $\langle B_1 \times B_2 \rangle_{\vee}$ is the flavor
  text\textquotesingle s paper-strength reading: the joins,
  never enumerated, discharge facet-by-facet through the standing
  citation lemmas \texttt{sound\_productLattice\_citeLeft/Right}
  (certified alongside, at \emph{arbitrary} rider), with the boundary
  decomposition itself paper-strength --- no separate join statement
  is bound.
\item with no displayed row: well-scopedness --- the factor
  lattices\textquotesingle{} facets and regions are confined to their
  own hypernets (\texttt{facetsScoped}/\texttt{regionConfined}
  fields of \texttt{CertifiedBlockLattice}), which is what makes
  transport along the union injection sound; the undisplayed-pack
  class of the \rulename{PsdCert} evaluability rows.
\end{itemize}
Width note: the citation lemmas accept \emph{any} rider formula where
the rule\textquotesingle s conclusion only ever conjoins the other
factor\textquotesingle s region --- wider machinery in the safe
direction, and exactly the generality the join-block sub-facets
consume.

\ruleblock{Lean verifier receipts}
\sourcenote{module \texttt{BernsteinLogic.ProductLattice}; bound build \texttt{84b7018f79da\ldots} (bound from receipt 9ebc3d9af324d1f7, 2026-07-21T04:48:38Z). A receipt records that named checking layers accepted a build: (module, compiled-artifact SHA-256, environment SHA-256, checker layers), minted when the checking gate ran; the store receipts attesting the bound build are listed with the layers each run credited. Nothing here asks the reader to trust the store: every layer is an independently runnable program over the quoted sources --- rebuild the module and re-run the checkers to reproduce the verdict. Receipts attesting superseded builds credit no layer and are omitted from this rendering (the count of omissions, if any, is noted below).}
\begin{receipttable}
\receiptrow{224af3a8595b550f}{2026-07-21T01:49:09Z}{build+decorrelated}{\badge{BadgeGreen}{current}}
\receiptrow{7187ff3aa085cd58}{2026-07-21T01:58:41Z}{build+decorrelated}{\badge{BadgeGreen}{current}}
\receiptrow{9ebc3d9af324d1f7}{2026-07-21T04:48:38Z}{build+decorrelated}{\badge{BadgeGreen}{current}}
\receiptrow{fd5ec354b89369c0}{2026-07-21T05:28:29Z}{build+replay+hardened+decorrelated}{\badge{BadgeGreen}{current}}
\end{receipttable}
\sourcenote{7 further receipt(s) in the store attest superseded builds of this module and are omitted from this rendering.}
\layerline{\layeritem{build}{lake build (elaboration)}{\badge{BadgeGreen}{current}}\layeritem{replay}{leanchecker (official kernel replay)}{\badge{BadgeGreen}{current}}\layeritem{hardened}{safe\_verify (hardened screen)}{\badge{BadgeGreen}{current}}\layeritem{decorrelated}{lean4export + nanoda (decorrelated kernel)}{\badge{BadgeGreen}{current}}}
\sourcenote{\badge{BadgeGray}{comparator: none} no fixture in the three-implementation comparator harness (compiled / kernel-interpreted / reference-interpreted) exercises this rule.}
\ruleblock{HOL4 statement}
\absentbox{NO HOL STATEMENT YET}{The certificate-object stratum of this rule now has HOL twins, with
kernel-evaluated transcription receipts, across five ported theories:
the certified-block record and the citation constructor, the
product-block assembly, the certified block lattice with its
wellformedness and validity predicates, the applicability payload
record with its wellformedness predicate, the meet-generator
enumeration, and the valuation-restriction lemmas the transport
argument reads (verified against the scripts this writing;
\texttt{hol-strict} clean at the same 2026-07-22 run as the ported
entries). What has no HOL twin yet is exactly the
rule\textquotesingle s semantic content --- the transport-and-weaken
citation lemmas and the soundness theorem they assemble into (the
commentary\textquotesingle s own point: nothing in the certificate
stratum evaluates dynamics anywhere, so the semantic layer IS the
rule). The rule also has no arm in the ported proof-object family and
no wire tag on either side. Absent until the citation/soundness
stratum ports; the certificate-object groundwork above is in place
for that seating.
}
\ruleblock{HOL verifier receipts}
\absentbox{NONE}{no HOL verification receipt exists for this rule.}
\ruleblock{Proof}

\sourcenote{Genre: admissible rule (tier 2) --- a proof-tree-rewriting
recipe, at the certificate-family layer: one lattice node unpacks into
a family of cited facet sequents. Like \rulename{PConstExact}, part of
the discharge is by reduction to verified machinery rather than to
catalogue nodes --- the transport step is a semantic monotonicity
lemma with no catalogue counterpart, because hypernet extension is
invisible to the sequent calculus --- and the recipe says so.}

\textbf{Setting.}\enspace A \emph{certified attracting-block lattice}
over a hypernet $H$ (the \texttt{CertifiedBlockLattice} structure
quoted above) is a finite family of blocks --- each a region box
formula plus its internal-facet inwardness sequents --- carrying
explicit index-level lattice structure and wire confinement; the family
is \emph{valid} when every facet sequent is semantically valid over
$H$. The rule's premises: valid families $B_1$ over $H_1$ and $B_2$
over $H_2$, and wire-disjointness
$\wires(H_1) \cap \wires(H_2) = \emptyset$. Conclusion: the
join-closure $\langle B_1 \times B_2 \rangle_{\vee}$ is certified over
the disjoint union $H_1 \sqcup H_2$.

\begin{brecipe}[the \rulename{ProductLattice} instance --- one node,
a citation family]\label{rec:pl-instance}
The rewrite replaces the lattice node by one derivation fragment
\emph{per facet of per product block}, each of the same two-move shape:
\[
\inferrule*[Right=cite]{\;s_0 \rhd \text{a factor facet sequent (cited sub-proof)}\;}{\bigl\langle \mathrm{pre}(s_0) \wedge \mathrm{rider},\ \mathrm{post}(s_0) \bigr\rangle \ \text{over } H_1 \sqcup H_2}
\]
where the rider is the \emph{other} factor's box. The conclusion is
kept verbatim --- the polynomial and predicate of the cited facet,
unchanged --- because a composite facet of an axis-aligned product is
normal to one factor's axis: its inwardness polynomial \emph{is} that
factor's, the other factor's coordinates riding along as free box
variables. Only the meet-generators $\{b_1 \times b_2\}$ are
enumerated ($O(|B_1| \cdot |B_2|)$ blocks, $O(|B_1| + |B_2|)$
distinct cited factor sub-proofs, each reused across its products
with varying riders); the join-closure's remaining elements ---
exponentially many --- are never materialized. They discharge by the
same mechanism: a boundary point of a finite union of axis-aligned
blocks lies on the boundary of some member containing it (member
interiors are interior to the union), so a union block's boundary
decomposes into sub-facets, each a subset of some generator's facet
and hence again normal to one factor's axis; inwardness for the member
implies inwardness for the union (monotonicity under $\subseteq$), so
the sub-facet's sequent is the same factor citation with its premise
further constrained to the exposed portion --- and the citation
lemmas accept these further-constrained premises because they accept
an \emph{arbitrary} rider.
\end{brecipe}

\begin{blemma}[transport: validity is monotone under hypernet
extension]\label{lem:pl-transport}
Let $H'$ extend $H$'s wire typing and contain $H$'s applications, and
let $s$ be semantically valid over $H$ and well-scoped in $H$. Then
$s$ is semantically valid over $H'$.
\end{blemma}
\begin{proof}
An admissible valuation $\sigma$ of $H'$ restricts to a valuation of
$H$ (read each $H$-typed wire through the extension), and the
restriction is admissible: every entry a per-application semantic
predicate reads is at a port wire, typed in $H$ by type consistency,
where the restriction agrees with $\sigma$ by construction. Formulas
whose free wires are $H$-typed therefore evaluate identically under
$\sigma$ and its restriction --- by induction on terms and formulas,
the only base case being a wire read, which is where well-scopedness
is consumed. So from $\sigma \models \mathrm{pre}(s)$: the restriction
satisfies $\mathrm{pre}(s)$, validity over $H$ gives the restriction
satisfies $\mathrm{post}(s)$, and evaluation-agreement carries it back
to $\sigma$. This is \texttt{Sequent.semanticallyValid\_extend};
the well-scopedness hypothesis is not decoration --- a sequent reading
a wire \emph{untyped} in $H$ evaluates that read differently in an
extension that types it, and the induction's base case fails exactly
there.
\end{proof}

\begin{blemma}[weaken, standalone]\label{lem:pl-weaken}
If $\langle \mathrm{pre}, \mathrm{post} \rangle$ is semantically valid
over $H$, so is
$\langle \mathrm{pre} \wedge \mathrm{extra}, \mathrm{post} \rangle$,
for any formula $\mathrm{extra}$.
\end{blemma}
\begin{proof}
Project the premise (\texttt{Sequent.semanticallyValid\_weaken}). The
rider needs no side condition --- in particular it may mention wires
the cited factor has never heard of, which is the whole point.
\end{proof}

\begin{bprop}[soundness of \rulename{ProductLattice}]
\label{prop:pl-sound}
Given valid families over wire-disjoint hypernets, every facet sequent
of every product block is semantically valid over
$H_1 \sqcup H_2$.
\end{bprop}
\begin{proof}
A product block's facet list is, by construction (\texttt{prodBlock}),
the left factor's facets cited with the right factor's region as
rider, concatenated with the symmetric family. For a left-cited facet:
the disjoint union extends $H_1$'s wire typing and contains its
applications (disjointness makes the union's linearity, acyclicity,
and type consistency inherit from the factors --- this is where the
side condition does structural work), so
Lemma~\ref{lem:pl-transport} transports the factor's validity to the
union, and Lemma~\ref{lem:pl-weaken} conjoins the rider. Symmetrically
on the right. That is the entire proof
(\texttt{sound\_productLattice}, through the two standing citation
lemmas \texttt{sound\_productLattice\_citeLeft}/\texttt{citeRight}).
For the un-enumerated join blocks' sub-facets the same two lemmas
supply the per-generator transport and arbitrary-rider steps, with the
exposure constraints as riders --- but the join-closure extension
itself (the boundary decomposition of Recipe~\ref{rec:pl-instance}'s
argument) is paper-strength: no bound declaration represents finite
unions, their boundary decomposition, or the union-inwardness step,
and the mechanized theorem quantifies over the generator blocks only.
\end{proof}

\begin{bremark}[why wire-disjointness is not decoration]
\label{rem:pl-disjoint}
With a shared wire the recipe fails twice. Structurally, the disjoint
union is no longer well-formed as built (linearity of the concatenated
application lists is proven \emph{from} disjointness). Semantically, no
single transported sequent's validity would fail ---
Lemma~\ref{lem:pl-weaken} conjoins \emph{any} rider, shared wires
included --- but the assembled \emph{readout} would misdescribe the
composite: the identification of every product-block facet with a
cited factor facet (the other factor entering only as a rider) is
proven from disjointness, and with a shared wire the built object is
simply not a product of the factors.
A single shared wire is precisely an \emph{interface}, and composition
along an interface is \rulename{Promise}/small-gain territory --- the
rule's applicability condition is the boundary between the two
composition regimes, not a technicality.
\end{bremark}

\begin{bremark}[the completeness fence: flows only, and the
map counterexample]\label{rem:pl-fence}
Proposition~\ref{prop:pl-sound} is dynamics-free --- no premise
mentions a flow or a map, so the rule is sound for both. What is
\emph{flows-only} is the completeness reading from the source
literature \cite{kalieswehbe2026}: for flows, the attractor lattice of
a product is the coproduct of the factors' lattices, so the generated
family is everything. For maps it is not: with
$f(x) = -\tanh(2x)$ on $[-p, p]$ (where $\tanh(2p) = p$), the factor
attractor lattice is $\{\emptyset, P, X\}$ with $P = \{p, -p\}$, yet
$f \times f$ has the phase-locked period-two attractors
$\{(p,p), (-p,-p)\}$ and $\{(p,-p), (-p,p)\}$ --- neither generated by
factor attractors; the coproduct is a proper sublattice. (The
arithmetic: $\tanh(2x)$ leaves the diagonal with slope $2$, so a
fixed point $p \approx 0.96$ of $\tanh(2p) = p$ exists; $0$ repels,
$|f'(0)| = 2$; each product two-cycle attracts,
$(f^{2})'(\pm p) = 4(1 - p^{2})^{2} < 1$; and that the factor lattice
has nothing beyond $\emptyset, P, X$ is the source example's own
analysis \cite{kalieswehbe2026}.) The Lean
statement is shaped so this cannot overreach: it certifies a lower
bound family, and any decomposition-level readout from a map instance
must be stated as ``at least this structure,'' never ``exactly this.''
\end{bremark}

\begin{anchors}
Recipe \ref{rec:pl-instance}: citation \& product blocks & \texttt{citeFacet}, \texttt{prodBlock}, \texttt{PProductLatticeData.prodBlocks}\\
Lemma \ref{lem:pl-transport} & \texttt{Sequent.semanticallyValid\_extend}; restriction machinery \texttt{Hypernet.Valuation.restrict}, \texttt{Formula.holds\_restrict}, \texttt{PrimApp.semOk\_restrict}\\
Lemma \ref{lem:pl-weaken} & \texttt{Sequent.semanticallyValid\_weaken}\\
union well-formedness from disjointness & \texttt{Hypernet.disjUnion} (fields \texttt{linearity}, \texttt{acyclic}, \texttt{typeConsistent}); \texttt{Hypernet.WireDisjoint}\\
Proposition \ref{prop:pl-sound} & \texttt{sound\_productLattice}; \texttt{sound\_productLattice\_citeLeft}/\texttt{citeRight}\\
nested-application confinement & \texttt{prodBlocks\_confined}\\
family validity (premise form) & \texttt{CertifiedBlockLattice.Valid}\\
\end{anchors}


\blcatsection{Judgment Structure and Batteries}
\rulesection{37}{QSatFromRegion}{alias for \textcolor{DiffBlue}{\rulename{PConstExact}}}
\label{blcatrule-QSatFromRegion}
\ruleblock{The rule (Steele metanotation)}
\sourcenote{as realized in the mechanized development; source row: \texttt{\texttt{QTSL-rules.toml}} --- \emph{QSatFromRegion}. The metanotation statement is typeset by hand for notational quality; its mechanized ground truth is the Lean statement below, extracted from the sources at every build of this book and hash-bound. The rule \emph{roster} is machine-checked (every mechanized rule tag round-trips to exactly one catalogue row); the statement \emph{text} is editorial, checked by eye against the extracted Lean beside it.}
\begin{mathpar}
\inferrule*{
  \mathrm{child} : \mathsf{QRegionSat}(\varphi, R) \\
  \mathrm{eval} : \vdash x_0 \in R \ \text{(BL leaf; sublevel} \to \rulename{PConstExact}\text{, box} \to \text{registry interval check)} \\
  \mathsf{SInit}(\varphi) \in s.\forall
}{
  \mathcal{M} \vDash s
}
\end{mathpar}

\par\medskip
\rulename{QSatFromRegion} is the extension theory\textquotesingle s
closing move: it turns a region-level certificate into the root
satisfaction verdict. The workhorse judgment of the extension page
concludes region-level facts --- every point of a region $R$ satisfies
the formula $\varphi$, with $R$ under-approximating the
formula\textquotesingle s denotation --- and this rule is the single
bridge from that shape to the top-level claim
$\mathcal{M} \vDash s$: cite the child, check that the
specification\textquotesingle s initial point $x_0$ lies in $R$ by
exact rational arithmetic, and require that the warranted safety
component of the specification is the initial-state injection
$\mathsf{SInit}(\varphi)$ of the child\textquotesingle s formula. A producer uses
it exactly once per initialized certificate, as the last row.

What the checker recomputes: the membership leaf --- for a sublevel
atom $\{x : V(x) \le y\}$, exact substitution of $x_0$ into
$V$ and one closed comparison, which is precisely the ground-fact rule
\rulename{PConstExact}\textquotesingle s own check; for a box atom,
per-coordinate containment recomputed by the checker\textquotesingle s
side-condition registry --- and the component tie: the warrant row names which safety
component it warrants, anything other than the $\mathsf{SInit}$ image
of the child\textquotesingle s formula refuses at load (a menu of
one), and that the named component is among the
specification\textquotesingle s consumed components is admission
conformance.

What would break without each hypothesis: drop the child and $R$ is an
arbitrary set nobody certified --- the conclusion is unmoored from any
proof; drop the eval leaf and $x_0$ may lie outside $R$, so the
child\textquotesingle s universal never instantiates (a certificate
could certify a region the system never starts in); drop the component tie
and a certificate about some subformula --- or an unrelated formula
--- could claim a component it never certified.

\cataloguenote{The soundness content is one line --- instantiation of the
child\textquotesingle s universal at $x_0$ --- which is why the rule
prices as an alias even though it has two premise fields: the child is
a judgment \emph{citation} (an edge under the tower discipline, the
\rulename{ProofWellFormedness} section), and the sole derivation node the
respelling produces is the \rulename{PConstExact} leaf.
Mechanization state, plainly: the bound statements below are that one
line and its assembly. \texttt{sound\_QSatFromRegion} (module
\texttt{BernsteinLogicAux.QtslAliases}) is the instantiation itself,
generic in the satisfaction predicate and profile-free --- axiom-free,
even: the kernel certifies it depends on no axioms at all ---
and \texttt{sound\_QSatFromRegion\_sublevel} composes it over the
exact-$\Q$ membership leaf (the \rulename{RootEval} identities, this
page), which is the sole derivation node named above. The
payload-level elaborator is bound beside them (module
\texttt{BernsteinLogic.QtslElaborators}, the batch
template\textquotesingle s unpack made concrete):
\texttt{qsatFromRegionElab} takes the rule payload to the one
\rulename{PConstExact} payload --- an inhabitant of the
target\textquotesingle s own payload type, comparator confinement
discharged at elaboration --- with
\texttt{qsatFromRegionElab\_sound} the factoring through
\texttt{sound\_pConstExact} and \texttt{qsatFromRegionElab\_sat} the
judgment-level composition over the elaborated leaf; the semantic
content carries no profile residue, so the entry carries no
conditional flag. Inherited flag, travelling with the alias: the target
excludes the $\neq$ comparator from its admissible set (a
second-soundness-invariant exclusion, unpacked in the
target\textquotesingle s section); the membership leaf uses $\le$, and
any port must preserve the exclusion.
}
\ruleblock{Lean statement}
\declhead{BernsteinLogicAux.QtslAliases}{sound\_QSatFromRegion}{\badge{BadgeGreen}{statement bound}}
\begin{leansrc}
/-- **`sound_QSatFromRegion`** (QTSL-rules.toml row `QSatFromRegion`; the
catalogue's "Soundness: one line — instantiation of the universal"): the
child certificate concludes region-level satisfaction (`R` under-approximates
the truth set: every point of `R` satisfies), the eval leaf places the
specification's initial point in `R`, and the root claim follows by
instantiation. Generic in the satisfaction predicate: for the probabilistic
judgments the extension page concludes, `sat x` reads as the concluded
bound at start `x` (e.g. `sound_QSafetyBattery`'s conclusion), and the rule
is profile-free. The root tie `s.root = φ` is a syntactic side condition
recomputed by the checker; it selects *which* predicate this lemma is
instantiated at, and carries no semantic content of its own. -/
theorem sound_QSatFromRegion {X : Type*} {R : Set X} {sat : X → Prop} {x₀ : X}
    (hchild : ∀ x ∈ R, sat x) (heval : x₀ ∈ R) : sat x₀
\end{leansrc}
\declhead{BernsteinLogicAux.QtslAliases}{sound\_QSatFromRegion\_sublevel}{\badge{BadgeGreen}{statement bound}}
\begin{leansrc}
/-- **`QSatFromRegion` assembled over its `RootEval` leaf, sublevel shape**:
the child's region-level universal on the cast sublevel atom, instantiated
at the cast initial point through the exact-ℚ evaluation — the sole
derivation node the respelling produces, composed. -/
theorem sound_QSatFromRegion_sublevel {σ : Type*} {sat : (σ → ℝ) → Prop}
    {P : MvPolynomial σ ℚ} {thr : ℚ} {xq : σ → ℚ}
    (hchild : ∀ x ∈ {x : σ → ℝ |
      MvPolynomial.eval x (MvPolynomial.map (Rat.castHom ℝ) P) ≤ (thr : ℝ)}, sat x)
    (heval : MvPolynomial.eval xq P ≤ thr) :
    sat (fun i => (xq i : ℝ))
\end{leansrc}
\par\medskip
\textbf{The payload-level elaborator, displayed whole.}\enspace The
unpack builds the target\textquotesingle s payload from the
rule\textquotesingle s wire data --- the eval field \emph{is} the
\rulename{RootEval} leaf (that section binds the closed-term builder
and the wire-level exactness identity it rides).

\declhead{BernsteinLogic.QtslElaborators}{qsatFromRegionElab\_sound}{\badge{BadgeGreen}{statement bound}}
\begin{leansrc}
/-- The factoring, BL face — through `sound_pConstExact`, as the amber
box demanded. -/
theorem qsatFromRegionElab_sound {H : Hypernet} {seq : Sequent}
    {x0 : Nat → Rat} {V : Poly} {thr : Rat}
    (hEval : (rootEvalFormula x0 V thr).evalClosed = some true)
    (hpost : seq.post = rootEvalFormula x0 V thr) :
    seq.semanticallyValid H
\end{leansrc}
\declhead{BernsteinLogic.QtslElaborators}{qsatFromRegionElab\_sat}{\badge{BadgeGreen}{statement bound}}
\begin{leansrc}
/-- The judgement-level composition over the elaborated leaf
(`prop:qsr-sound` at the wire): instantiation of the child's universal
at x₀, the membership hypothesis supplied by the checked leaf through
the ℚ-face exactness identity. `sound_QSatFromRegion` is generic in the
satisfaction predicate and the carrier, so the composition consumes the
bound lemma verbatim — sound_A = (instantiation) ∘ (elaborated leaf). -/
theorem qsatFromRegionElab_sat {x0 : Nat → Rat} {V : Poly} {thr : Rat}
    {sat : (Nat → Rat) → Prop}
    (hchild : ∀ x ∈ {x : Nat → Rat | Poly.eval x V ≤ thr}, sat x)
    (hEval : (rootEvalFormula x0 V thr).evalClosed = some true) :
    sat x0
\end{leansrc}
\declhead{BernsteinLogic.QtslElaborators}{qsatFromRegionElab}{\badge{BadgeGreen}{definition bound}}
\begin{leansrc}
/-- **The QSatFromRegion payload elaborator**, sublevel case: the
rule's wire payload elaborates to exactly the RootEval leaf's payload —
the leaf-level respelling has its own row (`RootEval`), and this rule's
eval field IS that leaf (the consumption-site anchor of the RootEval
section). -/
def qsatFromRegionElab (H : Hypernet) (seq : Sequent) (x0 : Nat → Rat)
    (V : Poly) (thr : Rat)
    (hEval : (rootEvalFormula x0 V thr).evalClosed = some true)
    (hpost : seq.post = rootEvalFormula x0 V thr) :
    PConstExactData H seq :=
  rootEvalElab H seq x0 V thr hEval hpost
\end{leansrc}
\ruleblock{Correspondence (Steele display $\leftrightarrow$ Lean)}
\badge{BadgeGreen}{reviewer-certified against the displayed texts}
\sourcenote{The dictionary below is a reviewer-certified statement of how the mechanized statements above interpret the displayed rule: display metavariable $\mapsto$ Lean binder, each display hypothesis $\mapsto$ the premise or side condition carrying it, the display conclusion $\mapsto$ the goal --- with any deliberate difference in width between display and mechanization named outright. Certification is recorded as a hash triple over (the displayed cell, the extracted statement texts, this dictionary\textquotesingle s own text); if any of the three changes, this block turns amber until a reviewer re-certifies. The correspondence claim itself remains editorial --- reviewer-grade, not machine-checked; its machine-checked successor types each soundness statement against the rule as a structured object, and supersedes this block per rule as it lands.}
The displayed premises map onto \texttt{sound\_QSatFromRegion}
verbatim: the child $\mathsf{QRegionSat}(\varphi, R)$ $\mapsto$
\texttt{hchild : ∀ x ∈ R, sat x} (region-level satisfaction,
compiled); the eval leaf $\vdash x_0 \in R$ $\mapsto$ \texttt{heval :
x₀ ∈ R}; the conclusion $\mathcal{M} \vDash s$ $\mapsto$ \texttt{sat
x₀} --- universal instantiation, one line, axiom-free. The component tie
($\mathsf{SInit}(\varphi) \in s.\forall$, the refuse-mode warrant
naming) has deliberately no Lean counterpart: it
is the syntactic side condition that selects \emph{which} predicate
\texttt{sat} the lemma is instantiated at, recomputed by the checker,
carrying no semantic content --- the statement\textquotesingle s
docstring says exactly this. \texttt{sound\_QSatFromRegion\_sublevel}
is the assembled form over the exact-$\Q$ \rulename{RootEval} leaf:
child universal on the cast sublevel atom + exact-$\Q$ evaluation
\texttt{heval : MvPolynomial.eval xq P ≤ thr} $\Rightarrow$
satisfaction at the cast initial point --- the sole derivation node
the respelling produces, composed. Width note: \texttt{sat : X →
Prop} is generic --- probabilistic conclusions read \texttt{sat x} as
the concluded bound at start \texttt{x}, so the alias is
profile-free by binder shape; wider than any single displayed
instantiation, safe direction.

\ruleblock{Lean verifier receipts}
\sourcenote{module \texttt{BernsteinLogicAux.QtslAliases}; bound build \texttt{570599779c29\ldots} (bound from receipt 5a60f4d01fd81ac5, 2026-07-21T05:39:02Z). A receipt records that named checking layers accepted a build: (module, compiled-artifact SHA-256, environment SHA-256, checker layers), minted when the checking gate ran; the store receipts attesting the bound build are listed with the layers each run credited. Nothing here asks the reader to trust the store: every layer is an independently runnable program over the quoted sources --- rebuild the module and re-run the checkers to reproduce the verdict. Receipts attesting superseded builds credit no layer and are omitted from this rendering (the count of omissions, if any, is noted below).}
\begin{receipttable}
\receiptrow{0a3819a978e31127}{2026-07-21T04:51:36Z}{build+decorrelated}{\badge{BadgeGreen}{current}}
\receiptrow{0bb1074df705d863}{2026-07-21T05:21:58Z}{build+decorrelated}{\badge{BadgeGreen}{current}}
\receiptrow{5a60f4d01fd81ac5}{2026-07-21T05:39:02Z}{build+replay+hardened+decorrelated}{\badge{BadgeGreen}{current}}
\end{receipttable}
\sourcenote{11 further receipt(s) in the store attest superseded builds of this module and are omitted from this rendering.}
\sourcenote{module \texttt{BernsteinLogic.QtslElaborators}; bound build \texttt{0205ec213b91\ldots} (hashed from the built olean at bind time). A receipt records that named checking layers accepted a build: (module, compiled-artifact SHA-256, environment SHA-256, checker layers), minted when the checking gate ran; the store receipts attesting the bound build are listed with the layers each run credited. Nothing here asks the reader to trust the store: every layer is an independently runnable program over the quoted sources --- rebuild the module and re-run the checkers to reproduce the verdict. Receipts attesting superseded builds credit no layer and are omitted from this rendering (the count of omissions, if any, is noted below).}
\begin{receipttable}
\receiptrow{a237f4cff27cb15a}{2026-07-22T10:06:47Z}{build+decorrelated}{\badge{BadgeGreen}{current}}
\end{receipttable}
\layerline{\layeritem{build}{lake build (elaboration)}{\badge{BadgeGreen}{current}}\layeritem{replay}{leanchecker (official kernel replay)}{\badge{BadgeGreen}{current}}\layeritem{hardened}{safe\_verify (hardened screen)}{\badge{BadgeGreen}{current}}\layeritem{decorrelated}{lean4export + nanoda (decorrelated kernel)}{\badge{BadgeGreen}{current}}}
\sourcenote{\badge{BadgeGray}{comparator: none} no fixture in the three-implementation comparator harness (compiled / kernel-interpreted / reference-interpreted) exercises this rule.}
\ruleblock{HOL4 statement}
\absentbox{NO HOL STATEMENT YET}{The HOL4 development\textquotesingle s object is the base
calculus\textquotesingle s wire-format checker (theory
\texttt{bern\_catalogue} ports the base catalogue\textquotesingle s
check-function cone); no HOL development of the extension theory
exists, so there is no HOL statement for any extension rule to bind.
Absent rather than out-of-scope: the bridge\textquotesingle s baked
realization is the extension theory\textquotesingle s warrant seam
(\texttt{qsat\_warrant}), so a future HOL port of an
extension-theory checker would state it. The one leaf the alias adds
checks through the base ground-fact rule \rulename{PConstExact}, whose
own HOL statement is not yet writable either (it postdates the HOL4
port and its wire tag is unallocated --- see that
rule\textquotesingle s section).
}
\ruleblock{HOL verifier receipts}
\absentbox{NONE}{no HOL verification receipt exists for this rule.}
\ruleblock{Proof}

\sourcenote{Genre: alias rule (tier 3) --- a respelling template: the
1$\to$1 unpack, the filled-in target application, the payload-transport
table, the exactness identity, and the soundness corollary through the
target\textquotesingle s lemma. Two riders. First, the tier criterion
counts derivation \emph{nodes}, not premise fields: this rule has two
premise fields, but the child is a judgment \emph{citation} --- an edge
under the tower discipline (the \rulename{ProofWellFormedness} section), not
a node any unpack must produce --- so the rule\textquotesingle s proof
content per instance is at most one node: in the sublevel branch, the
single \rulename{PConstExact} application discharging its membership
leaf; in the box branch, no node at all (the leaf is registry
recomputation). The child instantiation and the component tie are
definitional transports along edges, not nodes any unpack produces ---
that, and nothing stronger, is the sense in which the rule is an
alias.
Second, the semantic composition is mechanized and bound in the entry,
and so --- now that the extension certificate schema exists --- is the
payload-level elaborator (the batch template\textquotesingle s unpack,
at the wire level; the remark below names the bound functions).}

\begin{brecipe}[the alias template, one extends-edge up]\label{rec:qsr-template}
An alias at the satisfaction level is warranted by a 1$\to$1 respelling
into a single application of a target rule, with judgment citations
transported as \emph{edges} and the concluded judgment recovered
definitionally from the cited child. Here the target is the base
calculus\textquotesingle s ground-fact rule \rulename{PConstExact}; the
one new node is the membership leaf; and the $\vDash$-level conclusion
is an instantiation of the child\textquotesingle s universal
(Proposition~\ref{prop:qsr-inst}), not a derived fact needing catalogue
structure.
\end{brecipe}

\begin{brecipe}[the filled-in application]\label{rec:qsr-filled}
The rule\textquotesingle s fields: a child judgment $\mathrm{child} :
\mathsf{QRegionSat}(\varphi, R)$ --- every point of $R$ satisfies
$\varphi$, with $R$ under-approximating the formula\textquotesingle s
denotation --- a membership leaf $\mathrm{eval}$, and the component tie
$\mathsf{SInit}(\varphi) \in s.\forall$ (the warranted safety
component is the initial-state injection of the child\textquotesingle s
formula --- the record carries component lists, not a single root, and
the warrant row names its component). For a sublevel region atom
$R = \{x : V(x) \le y\}$, with $V$ a polynomial over $\Q$
and $x_0$ the specification\textquotesingle s rational initial point,
the eval leaf \emph{is} the \rulename{PConstExact} node
\begin{mathpar}
\inferrule*[Right=\rulename{PConstExact}]{
  V(x_0) \le y \ \text{holds under exact evaluation in } \Q
}{
  \Phi \vdash V(x_0) \le y
}
\end{mathpar}
with arbitrary premise $\Phi$, as the target licenses. $V(x_0)$ is the
closed term obtained by exact substitution of $x_0$\textquotesingle s
coordinates --- wire-free, so the target\textquotesingle s verified
evaluator decides it, with comparator $\le$ inside the
target\textquotesingle s admissible set. The leaf-level respelling has
its own row on this page (\rulename{RootEval}), whose section carries
the transport and exactness detail; this section carries the
judgment-level composition. For a \emph{box} atom the leaf is not a
\rulename{PConstExact} node: per-coordinate interval containment is
recomputed directly by the checker\textquotesingle s side-condition
registry --- the alias reading covers the sublevel branch, and the box
branch is stated here so the respelling never overclaims.
\end{brecipe}

\begin{brecipe}[payload transport --- field for field]\label{rec:qsr-transport}
\begin{center}\small
\begin{tabular}{@{}ll@{}}
\textbf{\rulename{QSatFromRegion} field} & \textbf{destination}\\
\midrule
$\mathrm{child}$ & not transported --- a judgment citation (edge; the tower discipline)\\
$\mathrm{eval}$ & the \rulename{PConstExact} payload: closed comparison $V(x_0) \le y$\\
component tie $\mathsf{SInit}(\varphi) \in s.\forall$ & the concluded judgment\textquotesingle s formula\\
concluded $\mathcal{M} \vDash s$ & instantiation of the child\textquotesingle s universal at $x_0$\\
\end{tabular}
\end{center}
\end{brecipe}

\begin{bprop}[the instantiation identity]\label{prop:qsr-inst}
If $R \subseteq \{x : \mathcal{M}, x \vDash \varphi\}$ and
$x_0 \in R$, then $\mathcal{M}, x_0 \vDash \varphi$; with
the component tie, this is $\mathcal{M} \vDash s$. Moreover,
for the sublevel atom the membership hypothesis is \emph{exactly} the
checked leaf, as an identity of exact rationals:
\[
x_0 \in R \iff V(x_0) \le y,
\]
because substituting a rational point into a rational polynomial
evaluates with no rounding --- the checker\textquotesingle s
$\mathsf{true}$ verdict \emph{is} membership, not an approximation to
it (the \rulename{RootEval} section proves this as its exactness
identity).
\end{bprop}
\begin{proof}
Instantiate the universal at $x_0$; with the component tie
the instance \emph{is} $\mathcal{M} \vDash s$. For the moreover clause,
two equivalences compose: $x_0 \in R \iff V(x_0) \le y$ is
the sublevel region\textquotesingle s membership predicate read at
$x_0$ --- a definition, not a computation --- and the identification of
the checker\textquotesingle s $\mathsf{true}$ verdict with that exact
inequality is the \rulename{RootEval} exactness identity, cited rather
than re-proved. The instantiation is the rule\textquotesingle s entire
own content --- one node\textquotesingle s worth, which is why it
prices as an alias.
\end{proof}

\begin{bprop}[soundness corollary]
\label{prop:qsr-sound}
In the sublevel branch,
$\mathrm{sound}_{\mathsf{QSatFromRegion}} = $ \texttt{sound\_pConstExact}
(applied to the eval leaf) $\fatsemi$ (instantiation at $x_0$); in the
box branch the recomputed containment supplies the membership fact and
the same instantiation finishes. In both branches the
child\textquotesingle s validity is supplied by the induction hypothesis
of the tower discipline\textquotesingle s induction
(\rulename{ProofWellFormedness} section).
\end{bprop}
\begin{proof}
Sublevel branch: \texttt{sound\_pConstExact} (the target\textquotesingle s
kernel-verified lemma, quoted in the \rulename{PConstExact} section)
turns the $\mathsf{true}$ evaluation into the validity of
$V(x_0) \le y$; Proposition~\ref{prop:qsr-inst} finishes.
Box branch: the recomputed containment plays the leaf\textquotesingle s
role --- per-coordinate comparison of exact rationals, the box face of
the same no-rounding exactness the \rulename{RootEval} section proves
for sublevel leaves, so the recompute\textquotesingle s verdict is
membership itself; Proposition~\ref{prop:qsr-inst}\textquotesingle s
instantiation, which never asked how membership was obtained, finishes
identically.
\end{proof}

\begin{bremark}[the payload-level elaborator, mechanized]
\label{rem:qsr-elab}
The \emph{semantic} composition above is bound in the entry:
\texttt{sound\_QSatFromRegion} is the instantiation (axiom-free), and
\texttt{sound\_QSatFromRegion\_sublevel} is
Proposition~\ref{prop:qsr-sound}\textquotesingle s statement over the
exact-$\Q$ leaf (the \rulename{RootEval} identities). The payload-level
elaborator --- previously this page\textquotesingle s one absent
artefact, held back until the extension certificate schema existed ---
is bound with it: \texttt{qsatFromRegionElab} (module
\texttt{BernsteinLogic.QtslElaborators}) takes the rule\textquotesingle
s wire payload to (in the sublevel case) exactly the one
\rulename{PConstExact} payload of Recipe~\ref{rec:qsr-transport}, as an
inhabitant of the target\textquotesingle s own payload type --- the
comparator confinement discharged at elaboration, once ---
with \texttt{qsatFromRegionElab\_sound} the factoring through
\texttt{sound\_pConstExact} and \texttt{qsatFromRegionElab\_sat} the
judgment-level composition over the elaborated leaf, the
checker\textquotesingle s verdict entering as membership through the
wire-level exactness identity \texttt{rootEvalFormula\_true\_iff} (the
$\Q$ face of the \rulename{RootEval} identities). The box branch
remains checker-side registry recompute, by design --- nothing to
elaborate.
\end{bremark}

\begin{bremark}[the inherited comparator confinement]
\rulename{PConstExact} excludes one comparator ($\neq$) from its
admissible set, for reasons of a second soundness invariant unpacked in
that rule\textquotesingle s section. The membership leaf uses $\le$
only --- comfortably inside --- but any port or mechanization of this
rule must \emph{preserve} the exclusion rather than re-derive the
comparator set from the semantic lemma alone; the flag travels with
the alias.
\end{bremark}

\begin{anchors}
target lemma (Prop.~\ref{prop:qsr-sound}) & \texttt{sound\_pConstExact}; \texttt{Formula.holds\_of\_evalClosed\_true} (quoted in the \rulename{PConstExact} section)\\
leaf-level transport + exactness & the \rulename{RootEval} section (this page)\\
child validity & the tower discipline\textquotesingle s induction (\rulename{ProofWellFormedness} section)\\
alias-level corollary & \texttt{sound\_QSatFromRegion} + \texttt{sound\_QSatFromRegion\_sublevel} (bound in the entry)\\
payload-level elaborator (\autoref{rem:qsr-elab}) & \texttt{qsatFromRegionElab} + \texttt{qsatFromRegionElab\_sound} + \texttt{qsatFromRegionElab\_sat} (module \texttt{BernsteinLogic.QtslElaborators}, bound in the entry)\\
box-atom branch & checker-side registry recompute, by design --- no rule anchor\\
\end{anchors}


\rulesection{38}{QGuardIntro}{core rule}
\label{blcatrule-QGuardIntro}
\ruleblock{The rule (Steele metanotation)}
\sourcenote{as realized in the mechanized development; source row: \texttt{\texttt{QTSL-rules.toml}} --- \emph{QGuardIntro}. The metanotation statement is typeset by hand for notational quality; its mechanized ground truth is the Lean statement below, extracted from the sources at every build of this book and hash-bound. The rule \emph{roster} is machine-checked (every mechanized rule tag round-trips to exactly one catalogue row); the statement \emph{text} is editorial, checked by eye against the extracted Lean beside it.}
\begin{mathpar}
\inferrule*{
  \mathrm{child} : \mathsf{QRegionSat}(\varphi, R) \\
  \mathsf{SInit}(\varphi) \in s.\forall
}{
  \mathcal{M} \vDash \mathsf{Guard}(R,\ s)
}
\end{mathpar}

\par\medskip
\rulename{QGuardIntro} is implication introduction at the judgment
stratum: from a region-level certificate alone --- every state of $R$
satisfies $\varphi$ --- it concludes the \emph{guarded} claim
$\mathcal{M} \vDash \mathsf{Guard}(R, s)$, with no initial-membership
premise anywhere. The hypothesis that the unguarded route discharges
at the ground seam ($X_0 \subseteq R$, checked by exact rational
arithmetic) is instead discharged \emph{into the specification}: a
guarded specification demands its body\textquotesingle s guarantee
only from initial states inside $R$ and promises nothing off-region.
The guard region imports the certificate-side region grammar
(sublevel, box, affine) into the specification record at exactly one
slot; the record still binds no model and fixes no initial condition
--- the guard is a predicate \emph{on} the subject\textquotesingle s
initial data, not initial data.

What the checker recomputes: nothing beyond the child\textquotesingle s
own obligations --- that is the point. The two ties are load-time
conformance, not computation: the guarded warrant row names its
safety component, and anything other than the initial-state injection
of the warrant\textquotesingle s formula refuses at load (a menu of
one); the guard region stored on the record is derived from the
warrant row itself (store-and-check: a disagreeing producer row is a
second region on a one-region slot, refused at conformance).

What would break without each hypothesis: drop the child and the
guarded claim is unmoored from any proof; drop the component tie and
a certificate about an unrelated formula could claim the guarded
specification; drop the disjointness from the unguarded warrant
relation and the initial-membership demand would fire on a guarded
introduction --- exactly the hypothesis the former exists to
internalize.

\cataloguenote{The semantics is \emph{pointwise} over the credal stack: a behavior
object satisfies $\mathsf{Guard}(R, s)$ iff every admitted law whose
initial marginal is supported in $\llbracket R \rrbracket$ meets
$s$\textquotesingle s demand --- at demonic-box initial data this is
exactly conditioning the belief evaluation on initial membership.
Pointwise is \emph{strictly stronger} than the blanket conditional
(``if all of $X_0$ lies in $R$ then\ldots\textquotedbl), and the
strictness is a theorem, not a taste: the blanket reading would
vacuously discharge any guard the initial set merely straddles. The
class algebra stays a complement-free distributive lattice --- the
guard\textquotesingle s antecedent is a ground condition on system
data, not a behavior class, which is why the conditional lives at the
record stratum and no formula-level implication former exists or is
introduced.
}
\ruleblock{Theory rows (geolog)}
\sourcenote{The rule\textquotesingle s certificate content as it is declared in the geolog theory the checker loads --- payload schema, side-condition demands, and axiom rows, quoted from the theory source at every build of this catalogue (a stale quotation cannot exist). Axiom modes: \texttt{chase} = the engine derives; \texttt{report} = a violation is a loud judgment; \texttt{refuse} = a violation gates admission. Side-condition relations are discharged by bounded evaluation at the engine\textquotesingle s ground seam, never by derivation.}
\par\medskip
The baked realization: the guard slot on the specification record,
the guarded-route warrant relation (disjoint-by-relation from the
unguarded one, so the initial-membership axiom structurally never
fires on a guarded introduction), and the ties.

\declhead{theory Qtsl}{spec\_guard}{\badge{BadgeGreen}{quoted from the theory source at render}}
\begin{leansrc}
spec_guard : {spec: Spec, region: QRegionAtom} -> Prop;
\end{leansrc}
\declhead{theory Qtsl}{qsat\_guard\_warrant}{\badge{BadgeGreen}{quoted from the theory source at render}}
\begin{leansrc}
qsat_guard_warrant : {sat: QSat, warrant: QForallSat} -> Prop;
\end{leansrc}
\declhead{theory Qtsl}{ax/qguard/region\_tie}{\badge{BadgeGreen}{quoted from the theory source at render}}
\begin{leansrc}
chase ax/qguard/region_tie : forall r : QSat, c : QForallSat,
    sp : Spec.
  {sat: r, warrant: c} qsat_guard_warrant,
  r QSat/spec = sp
  |- {spec: sp, region: c QForallSat/region} spec_guard;
\end{leansrc}
\declhead{theory Qtsl}{ax/wf/qguard\_warrant\_member}{\badge{BadgeGreen}{quoted from the theory source at render}}
\begin{leansrc}
report ax/wf/qguard_warrant_member : forall r : QSat, sp : Spec,
    c : QForallSat.
  {sat: r, warrant: c} qsat_guard_warrant,
  r QSat/spec = sp
  |- {spec: sp, f: c QForallSat/component} qtsl_warrant_member;
\end{leansrc}
\ruleblock{Lean statement}
\par\medskip
\textbf{The guarded denotation and the introduction\textquotesingle s
soundness, displayed whole.}\enspace \texttt{GuardSat} is the
pointwise reading; the strictness witness against the blanket
conditional and the wrong-region countermodel bind beside the
soundness statements.

\declhead{QtslAux.QGuard}{sound\_QGuardIntro}{\badge{BadgeGreen}{statement bound}}
\begin{leansrc}
/-- Soundness of the introduction rule: a region warrant — the body's
demand holds from every state of `R` — yields the guarded claim with no
initial-membership premise, for any initial set. The membership hypothesis
that region-warranted satisfaction discharges at the ground seam is here
discharged into the former (implication introduction). -/
theorem sound_QGuardIntro {init R : Set X} {sat : X → Prop}
    (hchild : ∀ x ∈ R, sat x) : GuardSat init R sat
\end{leansrc}
\declhead{QtslAux.QGuard}{sound\_QGuardIntro\_sublevel}{\badge{BadgeGreen}{statement bound}}
\begin{leansrc}
/-- The rational-sublevel guard (a certificate family's sublevel set at an
exact-rational threshold, single-piece form): a warrant on the real
sublevel set of the coefficient-cast polynomial yields the guarded claim
with that sublevel set in the guard slot, for any initial set. -/
theorem sound_QGuardIntro_sublevel {σ : Type*} {init : Set (σ → ℝ)}
    {sat : (σ → ℝ) → Prop} {P : MvPolynomial σ ℚ} {thr : ℚ}
    (hchild : ∀ x ∈ {x : σ → ℝ |
        MvPolynomial.eval x (MvPolynomial.map (Rat.castHom ℝ) P) ≤ (thr : ℝ)},
      sat x) :
    GuardSat init
      {x : σ → ℝ |
        MvPolynomial.eval x (MvPolynomial.map (Rat.castHom ℝ) P) ≤ (thr : ℝ)}
      sat
\end{leansrc}
\declhead{QtslAux.QGuard}{sound\_QGuardIntro\_prob}{\badge{BadgeGreen}{statement bound}}
\begin{leansrc}
/-- Kernel-form corollary in the battery idiom: a per-initial-state
probability floor warranted on the region yields the guarded floor. -/
theorem sound_QGuardIntro_prob [MeasurableSpace X] {init R : Set X}
    {μ : X → Measure (ℕ → X)} {p : ENNReal} {A : Set (ℕ → X)}
    (hchild : ∀ x ∈ R, p ≤ μ x A) :
    GuardSat init R fun x => p ≤ μ x A
\end{leansrc}
\declhead{QtslAux.QGuard}{exists\_blanket\_not\_guardSat}{\badge{BadgeGreen}{statement bound}}
\begin{leansrc}
/-- Strictness witness for the semantics adjudication, and the permanent
twin of the dropped-membership mutant: the blanket reading does not imply
the pointwise one. On `Bool` with the full initial set, guard region
`{true}`, and an everywhere-false demand, the blanket conditional holds
vacuously (the initial set is not contained in the region) while the
pointwise guard is refuted at `true`. A blanket semantics would license
claims the pointwise semantics rejects; the former adopts the pointwise
reading. -/
theorem exists_blanket_not_guardSat :
    ∃ (init R : Set Bool) (sat : Bool → Prop),
      GuardSatBlanket init R sat ∧ ¬ GuardSat init R sat
\end{leansrc}
\declhead{QtslAux.QGuard}{exists\_not\_guardSat\_wrong\_region}{\badge{BadgeGreen}{statement bound}}
\begin{leansrc}
/-- Wrong-region countermodel (the permanent twin of the region-substitution
mutant): a warrant on one region does not license a guard on an unrelated
one — an everywhere-false demand is vacuously warranted on the empty
region, and the guard on the full region is refuted. -/
theorem exists_not_guardSat_wrong_region :
    ∃ (init R R' : Set Bool) (sat : Bool → Prop),
      (∀ x ∈ R, sat x) ∧ ¬ GuardSat init R' sat
\end{leansrc}
\declhead{QtslAux.QGuard}{GuardSat}{\badge{BadgeGreen}{definition bound}}
\begin{leansrc}
/-- The guarded-satisfaction denotation, pointwise (demonic-initial-set)
form: every admitted initial point inside the guard region meets the body's
demand; initial points outside the region are unconstrained. -/
def GuardSat (init R : Set X) (sat : X → Prop) : Prop :=
  ∀ x₀ ∈ init, x₀ ∈ R → sat x₀
\end{leansrc}
\ruleblock{Correspondence (Steele display $\leftrightarrow$ Lean)}
\badge{BadgeGreen}{reviewer-certified against the displayed texts}
\sourcenote{The dictionary below is a reviewer-certified statement of how the mechanized statements above interpret the displayed rule: display metavariable $\mapsto$ Lean binder, each display hypothesis $\mapsto$ the premise or side condition carrying it, the display conclusion $\mapsto$ the goal --- with any deliberate difference in width between display and mechanization named outright. Certification is recorded as a hash triple over (the displayed cell, the extracted statement texts, this dictionary\textquotesingle s own text); if any of the three changes, this block turns amber until a reviewer re-certifies. The correspondence claim itself remains editorial --- reviewer-grade, not machine-checked; its machine-checked successor types each soundness statement against the rule as a structured object, and supersedes this block per rule as it lands.}
The display maps onto \texttt{sound\_QGuardIntro} verbatim: the child
$\mathsf{QRegionSat}(\varphi, R)$ $\mapsto$ \texttt{h : ∀ x ∈ R, sat
x}; the conclusion $\mathcal{M} \vDash \mathsf{Guard}(R, s)$ $\mapsto$
\texttt{GuardSat init R sat} (pointwise: every admitted initial point
inside $R$ meets the body\textquotesingle s demand), generic in the
initial set \texttt{init} --- the rule is sound for \emph{any}
initial data, which is what ``no membership premise\textquotedbl{}
means formally. The component tie ($\mathsf{SInit}(\varphi) \in
s.\forall$) has deliberately no Lean counterpart: it is the load-time
refusal selecting \emph{which} predicate \texttt{sat} the lemma is
instantiated at. Width note: \texttt{sat : X → Prop} is generic, so
probabilistic bodies read \texttt{sat x} as the concluded bound at
start \texttt{x} --- wider than any single displayed instantiation,
safe direction. The specializations bind beside it: sublevel-atom
guards over exact-$\Q$ polynomial evaluation
(\texttt{sound\_QGuardIntro\_sublevel}, \texttt{guard\_eval\_sublevel})
and the probability-threshold body form
(\texttt{sound\_QGuardIntro\_prob}).

\ruleblock{Lean verifier receipts}
\sourcenote{module \texttt{QtslAux.QGuard}; bound build \texttt{3eb84fb9b559\ldots} (bound from receipt bd99b3a9109c4daf, 2026-07-28T18:45:47Z). A receipt records that named checking layers accepted a build: (module, compiled-artifact SHA-256, environment SHA-256, checker layers), minted when the checking gate ran; the store receipts attesting the bound build are listed with the layers each run credited. Nothing here asks the reader to trust the store: every layer is an independently runnable program over the quoted sources --- rebuild the module and re-run the checkers to reproduce the verdict. Receipts attesting superseded builds credit no layer and are omitted from this rendering (the count of omissions, if any, is noted below).}
\begin{receipttable}
\receiptrow{bd99b3a9109c4daf}{2026-07-28T18:45:47Z}{build+decorrelated}{\badge{BadgeGreen}{current}}
\end{receipttable}
\layerline{\layeritem{build}{lake build (elaboration)}{\badge{BadgeGreen}{current}}\layeritem{replay}{leanchecker (official kernel replay)}{\badge{BadgeGray}{none}}\layeritem{hardened}{safe\_verify (hardened screen)}{\badge{BadgeGray}{none}}\layeritem{decorrelated}{lean4export + nanoda (decorrelated kernel)}{\badge{BadgeGreen}{current}}}
\sourcenote{\badge{BadgeGray}{comparator: none} no fixture in the three-implementation comparator harness (compiled / kernel-interpreted / reference-interpreted) exercises this rule.}
\ruleblock{HOL4 statement}
\absentbox{NO HOL STATEMENT YET}{No HOL development of the extension theory exists (see the
region-to-satisfaction bridge\textquotesingle s section); the guarded
former\textquotesingle s realization is extension-theory record
structure, so a future HOL port of an extension-theory checker would
state it.
}
\ruleblock{HOL verifier receipts}
\absentbox{NONE}{no HOL verification receipt exists for this rule.}
\ruleblock{Proof}

\sourcenote{Genre: core rule (tier 1), structural group. The proof is
deliberately one step --- the rule exists to \emph{relocate} a
hypothesis, not to add content --- and the substantive work is the
semantics adjudication in the second remark, certified by the
strictness witness and the wrong-region countermodel bound at this
section.}

\textbf{Claim (\texttt{sound\_QGuardIntro}).}\enspace Let $R$ be the
warrant's region, $s$ the specification whose warranted safety
component is $\mathsf{SInit}(\varphi)$, and suppose the region
warrant: every state of $R$ satisfies $\varphi$. Then for \emph{any}
initial set $X_0$, the guarded claim holds: every admitted initial
point inside $R$ meets the body's demand.

\begin{proof}
Let $x \in X_0 \cap R$. The warrant instantiates at $x$: $x \in R$,
so $x \vDash \varphi$, which is the body's demand at $x$. Since $x$
was an arbitrary admitted initial point inside the guard region, the
pointwise guarded denotation holds. No hypothesis on $X_0$ was used
--- the rule is sound for arbitrary initial data, which is exactly
what licenses dropping the membership premise.
\end{proof}

\begin{bremark}[why pointwise, not blanket]
The blanket conditional --- \emph{if all of $X_0$ lies in $R$, then
the body holds from all of $X_0$} --- is strictly weaker: on an
initial set that merely straddles $R$, the blanket antecedent is
false and the blanket conditional vacuously true, while the pointwise
reading still demands the body from the inside points. The
strictness is mechanized (\texttt{exists\_blanket\_not\_guardSat});
the pointwise reading is the adopted semantics, and it is the one the
introduction rule is sound for.
\end{bremark}

\begin{bremark}[wrong region refuted]
Substituting any region other than the warrant's own breaks
soundness, by countermodel
(\texttt{exists\_not\_guardSat\_wrong\_region}); the load-time tie
(the stored guard region is derived from the warrant row,
store-and-check) is what makes the substitution unrepresentable
rather than detected.
\end{bremark}


\rulesection{39}{QGuardElim}{core rule}
\label{blcatrule-QGuardElim}
\ruleblock{The rule (Steele metanotation)}
\sourcenote{as realized in the mechanized development; source row: \texttt{\texttt{QTSL-rules.toml}} --- \emph{QGuardElim}. The metanotation statement is typeset by hand for notational quality; its mechanized ground truth is the Lean statement below, extracted from the sources at every build of this book and hash-bound. The rule \emph{roster} is machine-checked (every mechanized rule tag round-trips to exactly one catalogue row); the statement \emph{text} is editorial, checked by eye against the extracted Lean beside it.}
\begin{mathpar}
\inferrule*{
  \mathcal{M} \vDash \mathsf{Guard}(R,\ s) \\
  \mathrm{eval} : \vdash X_0 \subseteq R
}{
  \mathcal{M} \vDash s
}
\end{mathpar}

\par\medskip
\rulename{QGuardElim} is the elimination: a guarded claim plus initial
containment yields the body\textquotesingle s claim. The membership
premise is the same ground computation the unguarded route uses ---
per region-constructor pair, a degenerate box against a sublevel atom
is the exact evaluation leaf, a fat box against a sublevel atom is an
enclosure family, and box against box is interval containment ---
re-entering at exactly one place. Introduction followed by
elimination recovers the dissolved region-warranted route whole, so
the guarded former adds expressive structure without adding any new
soundness surface: everything the pair can conclude, the old route
could conclude given the same membership evidence.

What would break without the containment premise is the consumption
fence\textquotesingle s whole reason: a guarded claim consumed as if
unguarded is an unsound \emph{strengthening} --- the guarded claim
promises nothing off-region, so treating it as the body\textquotesingle
s claim smuggles in the very hypothesis the former internalized.

\cataloguenote{Realized \emph{at consumption}, deliberately --- there is no
standalone elimination node. Guard-presence triggers the membership
demand on any consumer of the claim: the composition rule walks its
premise list and demands initial membership for every guarded premise
(the quoted consumer-fence axiom below), and the verdict-compile and
equivalence evaluators are guard-aware. A consumer that cannot
discharge the membership defers loudly rather than assuming it.
}
\ruleblock{Theory rows (geolog)}
\sourcenote{The rule\textquotesingle s certificate content as it is declared in the geolog theory the checker loads --- payload schema, side-condition demands, and axiom rows, quoted from the theory source at every build of this catalogue (a stale quotation cannot exist). Axiom modes: \texttt{chase} = the engine derives; \texttt{report} = a violation is a loud judgment; \texttt{refuse} = a violation gates admission. Side-condition relations are discharged by bounded evaluation at the engine\textquotesingle s ground seam, never by derivation.}
\par\medskip
The baked realization of elimination-at-consumption: list membership
over the composition rule\textquotesingle s premises, the
consumer-fence axiom demanding initial membership for guarded
premises, and the ground membership relation it demands.

\declhead{theory Qtsl}{qsat\_list\_member}{\badge{BadgeGreen}{quoted from the theory source at render}}
\begin{leansrc}
qsat_list_member : {list: QSatList, sat: QSat} -> Prop;
\end{leansrc}
\declhead{theory Qtsl}{ax/wf/qwirecompose\_guard\_elim}{\badge{BadgeGreen}{quoted from the theory source at render}}
\begin{leansrc}
report ax/wf/qwirecompose_guard_elim : forall n : QWireCompose,
    l : QSatList, q : QSat, y : System, sp : Spec,
    g : QRegionAtom.
  n QWireCompose/components = l,
  {list: l, sat: q} qsat_list_member,
  q QSat/sys = y,
  q QSat/spec = sp,
  {spec: sp, region: g} spec_guard
  |- {region: g, sys: y} qtsl_init_member;
\end{leansrc}
\declhead{theory Qtsl}{qtsl\_init\_member}{\badge{BadgeGreen}{quoted from the theory source at render}}
\begin{leansrc}
qtsl_init_member : {region: QRegionAtom, sys: System} -> Prop;
\end{leansrc}
\ruleblock{Lean statement}
\declhead{QtslAux.QGuard}{sound\_QGuardElim}{\badge{BadgeGreen}{statement bound}}
\begin{leansrc}
/-- Soundness of the elimination rule: the initial-membership evaluation
re-enters exactly here — a guarded claim together with containment of the
initial set in the guard region yields the unguarded claim. -/
theorem sound_QGuardElim {init R : Set X} {sat : X → Prop}
    (hg : GuardSat init R sat) (heval : init ⊆ R) : ∀ x₀ ∈ init, sat x₀
\end{leansrc}
\declhead{QtslAux.QGuard}{qsatFromRegion\_derived}{\badge{BadgeGreen}{statement bound}}
\begin{leansrc}
/-- Coherence: introduction followed by elimination recovers the
region-warranted record route (warrant plus initial containment yields
satisfaction from every initial point). -/
theorem qsatFromRegion_derived {init R : Set X} {sat : X → Prop}
    (hchild : ∀ x ∈ R, sat x) (heval : init ⊆ R) : ∀ x₀ ∈ init, sat x₀
\end{leansrc}
\declhead{QtslAux.QGuard}{exists\_guardSat\_not\_unguarded}{\badge{BadgeGreen}{statement bound}}
\begin{leansrc}
/-- The permanent twin of the dropped-membership mutant, stated directly:
elimination WITHOUT its containment premise is false — a guarded claim
does not yield the unguarded one. On `Bool` with the full initial set and
guard region `{true}`, the demand "the state is `true`" is guarded-satisfied
(vacuously off-region, directly on it) yet fails from the initial point
`false`. -/
theorem exists_guardSat_not_unguarded :
    ∃ (init R : Set Bool) (sat : Bool → Prop),
      GuardSat init R sat ∧ ¬ ∀ x₀ ∈ init, sat x₀
\end{leansrc}
\declhead{QtslAux.QGuard}{guardSat\_singleton\_iff}{\badge{BadgeGreen}{statement bound}}
\begin{leansrc}
/-- At a singleton initial set the guarded denotation is the pointwise
conditional — the shape `Qtsl.Aliases.sound_QSatFromRegion` consumes. -/
theorem guardSat_singleton_iff {R : Set X} {sat : X → Prop} {x₀ : X} :
    GuardSat {x₀} R sat ↔ (x₀ ∈ R → sat x₀)
\end{leansrc}
\ruleblock{Correspondence (Steele display $\leftrightarrow$ Lean)}
\badge{BadgeGreen}{reviewer-certified against the displayed texts}
\sourcenote{The dictionary below is a reviewer-certified statement of how the mechanized statements above interpret the displayed rule: display metavariable $\mapsto$ Lean binder, each display hypothesis $\mapsto$ the premise or side condition carrying it, the display conclusion $\mapsto$ the goal --- with any deliberate difference in width between display and mechanization named outright. Certification is recorded as a hash triple over (the displayed cell, the extracted statement texts, this dictionary\textquotesingle s own text); if any of the three changes, this block turns amber until a reviewer re-certifies. The correspondence claim itself remains editorial --- reviewer-grade, not machine-checked; its machine-checked successor types each soundness statement against the rule as a structured object, and supersedes this block per rule as it lands.}
The display maps onto \texttt{sound\_QGuardElim} verbatim: the guarded
premise $\mapsto$ \texttt{hg : GuardSat init R sat}; the containment
premise $\vdash X_0 \subseteq R$ $\mapsto$ \texttt{hsub : init ⊆ R};
the conclusion $\mathcal{M} \vDash s$ $\mapsto$ \texttt{∀ x ∈ init,
sat x} (the unguarded pointwise demand). The dropped-membership
variant is refuted by countermodel
(\texttt{exists\_guardSat\_not\_unguarded}); the coherence statement
\texttt{qsatFromRegion\_derived} is the factoring of the dissolved
region-warranted route through the pair.

\ruleblock{Lean verifier receipts}
\sourcenote{module \texttt{QtslAux.QGuard}; bound build \texttt{3eb84fb9b559\ldots} (bound from receipt bd99b3a9109c4daf, 2026-07-28T18:45:47Z). A receipt records that named checking layers accepted a build: (module, compiled-artifact SHA-256, environment SHA-256, checker layers), minted when the checking gate ran; the store receipts attesting the bound build are listed with the layers each run credited. Nothing here asks the reader to trust the store: every layer is an independently runnable program over the quoted sources --- rebuild the module and re-run the checkers to reproduce the verdict. Receipts attesting superseded builds credit no layer and are omitted from this rendering (the count of omissions, if any, is noted below).}
\begin{receipttable}
\receiptrow{bd99b3a9109c4daf}{2026-07-28T18:45:47Z}{build+decorrelated}{\badge{BadgeGreen}{current}}
\end{receipttable}
\layerline{\layeritem{build}{lake build (elaboration)}{\badge{BadgeGreen}{current}}\layeritem{replay}{leanchecker (official kernel replay)}{\badge{BadgeGray}{none}}\layeritem{hardened}{safe\_verify (hardened screen)}{\badge{BadgeGray}{none}}\layeritem{decorrelated}{lean4export + nanoda (decorrelated kernel)}{\badge{BadgeGreen}{current}}}
\sourcenote{\badge{BadgeGray}{comparator: none} no fixture in the three-implementation comparator harness (compiled / kernel-interpreted / reference-interpreted) exercises this rule.}
\ruleblock{HOL4 statement}
\absentbox{NO HOL STATEMENT YET}{As at the introduction rule: no HOL development of the extension
theory exists; the consumption fence is extension-theory walk
structure.
}
\ruleblock{HOL verifier receipts}
\absentbox{NONE}{no HOL verification receipt exists for this rule.}
\ruleblock{Proof}

\sourcenote{Genre: core rule (tier 1), structural group. One-step
proof; the load-bearing content is the consumption fence (a guarded
claim consumed as if unguarded is an unsound strengthening), certified
by the dropped-membership countermodel.}

\textbf{Claim (\texttt{sound\_QGuardElim}).}\enspace If the guarded
claim holds (every admitted initial point inside $R$ meets the body's
demand) and $X_0 \subseteq R$, then the body's claim holds from every
admitted initial point.

\begin{proof}
Let $x \in X_0$. By containment $x \in R$, so $x \in X_0 \cap R$, and
the guarded claim yields the body's demand at $x$.
\end{proof}

\begin{bremark}[the fence, and its countermodel]
Dropping the containment premise is refuted by countermodel
(\texttt{exists\_guardSat\_not\_unguarded}): a guarded claim promises
nothing off-region, so an initial point outside $R$ can violate the
body while the guarded claim holds. This is why the elimination is
realized \emph{at consumption} --- guard-presence triggers the
membership demand on any consumer --- rather than trusting producers
to remember it.
\end{bremark}

\begin{bremark}[coherence]
Introduction followed by elimination recovers the dissolved
region-warranted route whole
(\texttt{qsatFromRegion\_derived}): warrant in, guarded claim out,
membership in, body's claim out --- the same conclusion the old
single rule drew from the same two pieces of evidence. The former
adds a place to \emph{stand} between them, not a new way to conclude.
\end{bremark}


\rulesection{40}{QSafetyBattery}{core rule}
\label{blcatrule-QSafetyBattery}
\ruleblock{The rule (Steele metanotation)}
\sourcenote{as realized in the mechanized development; source row: \texttt{\texttt{QTSL-rules.toml}} --- \emph{QSafetyBattery}. The metanotation statement is typeset by hand for notational quality; its mechanized ground truth is the Lean statement below, extracted from the sources at every build of this book and hash-bound. The rule \emph{roster} is machine-checked (every mechanized rule tag round-trips to exactly one catalogue row); the statement \emph{text} is editorial, checked by eye against the extracted Lean beside it.}
\begin{mathpar}
\inferrule*{
  V : \mathsf{QCertFam} \\
  \bigl\{D_m \vdash \mathcal{A}_m V_m \le 0\bigr\}_{m} \ \text{(interiors)} \\
  \bigl\{\mathrm{slab}_j \vdash g_j \fatsemi V_{m'_j} \le V_{m_j}\bigr\}_{j} \ \text{(guards, per jump $j$: source mode $m_j$, target $m'_j$, reset $g_j$)} \\
  \bigl\{U \cap \bar{D}_m \vdash V_m \ge 1\bigr\}_{m} \ \text{(floors)} \\
  \bigl\{\bar{D}_m \vdash V_m \ge 0\bigr\}_{m} \ \text{(nonnegs)} \\
  \mathrm{scope\_warrant} : \mathsf{QInvScope} \ \text{(optional — the dissolved ScopeRule)} \\
  p + y \le 1 \ \text{(threshold tie)}
}{
  \mathcal{M},\ x_0 \in \{V \le y\} \vDash \mathsf{SProbGE}\bigl(p,\ \mathsf{PAlways}\ \varphi\bigr)
}
\end{mathpar}

\ruleflavor{a region written left of the turnstile scopes the row: $R \vdash P$ asserts $P$ at every point of $R$ (each such row is an inherited BL \rulename{BernCone}-family sequent). The certificate family $V$ enters as certificate data --- a shared reference with no claim content of its own; the rows impose every claim --- never as a derivation edge. The threshold tie $p + y \le 1$ is the soundness theorem\textquotesingle s own hypothesis, displayed as a side condition: unlike the deadline and until nodes (whose baked sorts carry a \texttt{tie} proof field), the baked safety node carries no tie field yet --- a named schema divergence, not a checked leaf. The conclusion is the theory row\textquotesingle s own judgment, in constructor vocabulary: the region-conditioned probability formula ($\mathsf{SProbGE}(p, \mathsf{PAlways}\ \varphi)$ on the region $\mathsf{qraSublevel}(V, y)$). The record\textquotesingle s licensed alias tie is $\mathsf{SProbGE}(p, \psi) = \mathsf{SReturnGE}(p, \mathsf{retPath}\ \psi)$ --- a return respelling, which does not touch the state-region antecedent; the initial-set hypothesis $x_0 \in \{V \le y\}$ left of the judgment is the satisfaction binding\textquotesingle s own guard (the theory demands the subject\textquotesingle s initial data lie in the concluded region), rendered in the same hypothesis convention as every \rulename{BernCone}-family sequent. The internalized variant LANDED at the grid re-sort revision: the guarded-spec former ($\mathsf{Guard}$, a certificate-stratum record slot carrying the initial-region antecedent) with its introduction rule \rulename{QGuardIntro} (this catalogue, the structural group). The hypothesis form remains this catalogue\textquotesingle s rendering convention (one convention catalogue-wide); the guarded former is its internalized spelling, one \rulename{QGuardIntro} application away, and no belief-gauge reading is displayed.}
\par\medskip
\rulename{QSafetyBattery} is the probabilistic safety battery: from a
family of purely polynomial premise rows about a certificate $V$, it
concludes that with probability at least $p$ the system satisfies
$\mathsf{PAlways}\ \varphi$ (surface syntax $\mathsf{G}\,\varphi$) ---
\emph{forever} --- from every start state in
the sublevel region $\{V \le y\}$. The rule concludes a
satisfaction judgment ($\vDash$), one theory-extension edge above the
base calculus, and every premise row is an inherited base-calculus
sequent: \emph{interior} rows say the certificate does not increase
along the in-mode dynamics, \emph{guard} rows say jumps do not increase
it, \emph{floor} rows say $V \ge 1$ on the avoid set (the compiled
over-approximation of where $\varphi$ fails), and \emph{nonneg} rows say
$V \ge 0$ on the domain hull. The battery adds zero kernel arithmetic;
its entire content is the supermartingale assembly one level up: $V$
read along the process and stopped at the first failure is a nonnegative
supermartingale, and Ville\textquotesingle s inequality converts its
initial value into a bound on the probability of ever failing.

A producer reaches for the rule wherever a quantitative safety guarantee
over a hybrid stochastic system is wanted. The certificate is found by
whatever synthesis method the producer likes (typically linear or
sum-of-squares programming over the same polynomial cones the premise
rows are checked in); only the rows ship in the certificate.

What the checker recomputes are the side conditions: the domain-hull
condition (one step from any in-scope state stays inside the hull
$\bar{D}$, which is exactly where the floor and nonneg rows were
checked), the fail-closed containment (one-step mass stays in
$D \cup U$: leaving the modelled domain is counted as \emph{failure},
never as silent success), and the floor-scope resolution (the floor
rows\textquotesingle{} scopes resolve to an exact cover of the avoid set
within the hull). The threshold tie $p + y \le 1$ is displayed as a
side condition, not a checked leaf: it is the soundness
theorem\textquotesingle s own hypothesis, and the baked safety node
carries no tie proof field yet (contrast the deadline and until nodes)
--- a named schema divergence, as the rule\textquotesingle s flavor
text states.

Each hypothesis is load-bearing. Drop the floors and failure is free of
charge: a certificate identically $0$ satisfies every other row and
would certify almost-sure safety of a chain that fails at its first
step. Drop the nonnegs and the maximal inequality itself fails --- a
signed process can pay for a barrier visit with mass borrowed from
below. Drop the fail-closed containment and probability mass could exit
the modelled domain unaccounted. Drop the threshold tie and the
conclusion would assert success probability $p$ alongside a failure
budget $y$ with $p + y > 1$. The proof section
walks each of these, with counterexamples.

The dashboard row\textquotesingle s mark states the status honestly:
the hash-bound theorems warrant the rule at the
\emph{discrete-kernel profile} --- system semantics a one-step Markov
kernel on the flattened hybrid state space, interior rows read as
one-step expected decrease (the Foster--Lyapunov reading) --- with the
PDMP sub-profile of the continuous-time semantics additionally
machine-checked end-to-end. The general continuous-time profile
(jump--diffusion) is \emph{proven}, at the martingale-problem
presentation: the system enters through its generator drift data,
$\mathcal{A}V$ as data with the martingale-problem hypothesis for the
pair, and the battery\textquotesingle s conclusion follows end-to-end
(\texttt{sound\_GSHS\_safety\_mp},
\texttt{sound\_GSHS\_safety\_battery\_mp}, module
\texttt{QtslAux.LayerD.MPDischarge} of the semantic-universe
strand\textquotesingle s auxiliary package) --- the same epistemic
shape as the staged-row theorem\textquotesingle s (\texttt{sound\_QSafetyBattery\_staged}) mixture bridge: an
explicit presentation hypothesis, quantifying over exactly the
systems whose presentation satisfies it. The martingale-problem
theorems are hash-bound below like every other block --- their
package, the semantic-universe strand\textquotesingle s auxiliary
package in the public tree, is registered with this
catalogue\textquotesingle s binder --- so the dashboard row carries no
conditional mark: each named profile is proven, and its theorems are
bound.

At the general stratum the battery\textquotesingle s shape is also
stated over the semantic universe itself, in the semantic-universe
functor package (hash-bound below): trajectories live as closed record
sets under the Fell topology, and the undisplayed regularity pack ---
the measurability the profile proofs \emph{assume} of their region
events --- restates there as a theorem of the site\textquotesingle s
topology: the confinement event of a closed safe region is Borel over
every trajectory window, with no compactness and no
$\sigma$-compactness anywhere
(\texttt{measurableSet\_confinedEvent\_cylinder}). On that carrier the
safety conclusion is a \emph{wall} --- a $\forall$-face spec:
per-window confinement classes are closed under record restriction,
their persistence follows from window-calculus transitivity alone, and
the judgment bridge is an equivalence --- satisfaction \emph{equals}
almost-sure confinement at the window
(\texttt{satAtM\_safetySpecF\_iff}), with the indicator calculus
admitted exactly by the measurability above (\texttt{pull\_indPenM}).
Two further site-level facts sharpen the story. Record confinement of
a trajectory\textquotesingle s closed-up graph to a closed wall is
\emph{pointwise} safety at every stage, unconditional in the
trajectory (\texttt{graphCl\_mem\_confinedEvent\_iff}) --- the
segment-constancy side condition a skeleton-based route would carry is
absent by construction. And composing two wall demands along a window
compositor yields exactly the wall demand of the glued record class
(\texttt{mboxplusWF\_indPenM}): wall specs are closed under
composition, and persistence survives it. These statements bind the
site-general stratum --- the acceptance events and spec classes that
every profile\textquotesingle s conclusion instantiates; the displayed
judgment\textquotesingle s warrant remains with the profile theorems
above.

\cataloguenote{Fail-closed as a design reading: the underlying theorem\textquotesingle s
concluded event is \emph{stronger} than the safety
formula\textquotesingle s truth set --- never fail \emph{and} never
leave the modelled domain --- and the formula tie then weakens it. The
node presents its premises per mode and per jump (mode-indexed
certificate families); the bound theorem for that shape assembles the
flattened generator row through a sub-convex mixture decomposition of the
one-step kernel --- the mode\textquotesingle s flow step plus the
enabled guard jumps --- which is the presentation-to-semantics bridge,
carried as an explicit hypothesis and quoted in the
proof\textquotesingle s final remark.
}
\ruleblock{Theory rows (geolog)}
\sourcenote{The rule\textquotesingle s certificate content as it is declared in the geolog theory the checker loads --- payload schema, side-condition demands, and axiom rows, quoted from the theory source at every build of this catalogue (a stale quotation cannot exist). Axiom modes: \texttt{chase} = the engine derives; \texttt{report} = a violation is a loud judgment; \texttt{refuse} = a violation gates admission. Side-condition relations are discharged by bounded evaluation at the engine\textquotesingle s ground seam, never by derivation.}
\par\medskip
In the geolog theory the checker loads, every object this rule touches
is a declared row, and the display\textquotesingle s remaining ambient
notation resolves as follows. The concluded formula
$\mathsf{SProbGE}(p, \mathsf{PAlways}\ \varphi)$ (surface syntax
$\mathsf{P}_{\ge p}[\mathsf{G}\,\varphi]$) \emph{is} the safety-face
cap former \texttt{SProbGE} (concluding into the record\textquotesingle s
$\forall$-component sort) applied at the path former \texttt{PAlways}
--- the probability-floor former whose threshold canonicity is the
range axiom quoted below. The displayed $\mathcal{M}$ is the paired
subject: a \texttt{System} element, bound to the claim by the
\texttt{QSat} pairing (the specification consumes this formula among
its safety components; the warrant row names it). The generator reads $\mathcal{A}_m V_m$ are not
free-floating operator applications: \texttt{qtsl\_gen\_form} is the
derived generator-form seam --- given the subject system and a
certificate component (a hypernet from the state space to the
nonnegative extended reals), the row names the hypernet saying how the
certificate can change with duration, recomputed from the
subject\textquotesingle s content rows and never trusted as stored
data. The battery\textquotesingle s E2 payload is the
\texttt{QSafetyBattery} entity below, field for field the
display\textquotesingle s premise kit; the conclusion region
$\mathsf{qraSublevel}(V, y)$ is the region-atom constructor
of the same name.

\declhead{theory Qtsl}{SProbGE}{\badge{BadgeGreen}{quoted from the theory source at render}}
\begin{leansrc}
SProbGE      : {p: Q, path: PathPred} -> ForallSpec;
\end{leansrc}
\declhead{theory Qtsl}{ax/wf/prob\_range}{\badge{BadgeGreen}{quoted from the theory source at render}}
\begin{leansrc}
report ax/wf/prob_range : forall p : Q, ps : PathPred, f : ForallSpec.
  f = {p: p, path: ps} SProbGE
  |- {l: 0, r: p} q_lt, {l: p, r: 1} q_le;
\end{leansrc}
\declhead{theory Qtsl}{PAlways}{\badge{BadgeGreen}{quoted from the theory source at render}}
\begin{leansrc}
PAlways : StatePred -> PathPred;                          // G phi
\end{leansrc}
\declhead{theory Qtsl}{QSafetyBattery}{\badge{BadgeGreen}{quoted from the theory source at render}}
\begin{leansrc}
QSafetyBattery : Set;
\end{leansrc}
\declhead{theory Qtsl}{QSafetyBattery/*}{\badge{BadgeGreen}{quoted from the theory source at render}}
\begin{leansrc}
QSafetyBattery/proof : QSafetyBattery -> QProof;
QSafetyBattery/p : QSafetyBattery -> Q;
QSafetyBattery/arg : QSafetyBattery -> StatePred;
QSafetyBattery/cert : QSafetyBattery -> QCertFam;
QSafetyBattery/interiors : QSafetyBattery -> QProofList;
QSafetyBattery/guards : QSafetyBattery -> QProofList;
QSafetyBattery/lowbounds : QSafetyBattery -> QProofList;
QSafetyBattery/nonnegs : QSafetyBattery -> QProofList;
QSafetyBattery/avoid : QSafetyBattery -> QRegionAtom;
QSafetyBattery/scope_warrant : QSafetyBattery -> QInvScope;
\end{leansrc}
\declhead{theory Qtsl}{qraSublevel}{\badge{BadgeGreen}{quoted from the theory source at render}}
\begin{leansrc}
qraSublevel : {cert: QCertFam, thr: Q} -> QRegionAtom;
\end{leansrc}
\declhead{theory Qtsl}{qtsl\_gen\_form}{\badge{BadgeGreen}{quoted from the theory source at render}}
\begin{leansrc}
qtsl_gen_form : {sys: System, cert: HypernetRef, gen: HypernetRef}
  -> Prop;   // side-condition (oracle discharge / checker
\end{leansrc}
\ruleblock{Lean statement}
\declhead{BernsteinLogicAux.Battery}{sound\_QSafetyBattery}{\badge{BadgeGreen}{statement bound}}
\begin{leansrc}
/-- **`sound_QSafetyBattery`, discrete profile** (QTSL-rules.toml row `QSafetyBattery`
at the discrete-kernel profile; the concluded judgment reads as region
under-approximation — every start in the region satisfies the formula). Over the
fail-closed form, the formula tie (F) — the avoid region over-approximates the falsity
set of the `𝖦`-argument `φ`, compiled as `U := cmpl(φ)` — weakens the concluded event to
the truth of `𝖦 φ` along the path:
`Pr_{x}[∀ t, X t ⊨ φ] ≥ p` for every start `x` in the concluded region
`D ∩ {V ≤ thr}`. No measurability is asked of the satisfaction predicate: the `𝖦 φ`
event is bounded below through `measure_mono` on the (measurable) fail-closed event. -/
theorem sound_QSafetyBattery {sat : S → Prop}
    (hD : MeasurableSet D) (hDbar : MeasurableSet Dbar) (hU : MeasurableSet U)
    (hV : Measurable V)
    (hF : ∀ y, ¬ sat y → y ∈ U)
    (hΘ : p + thr ≤ 1)
    (hDsub : D ⊆ Dbar)
    (hR3 : ∀ y ∈ D, y ∉ U → K y Dbar = 1)
    (hR4 : ∀ y ∈ D, y ∉ U → K y (D ∪ U) = 1)
    (hB12int : ∀ y ∈ D, y ∉ U → Integrable V (K y))
    (hB12 : ∀ y ∈ D, y ∉ U → ∫ z, V z ∂K y ≤ V y)
    (hB3 : ∀ y ∈ U, y ∈ Dbar → 1 ≤ V y)
    (hB4 : ∀ y ∈ Dbar, 0 ≤ V y)
    (hx : x ∈ D) (hthr : V x ≤ thr) :
    ENNReal.ofReal p ≤ pathMeasure K x {ω | ∀ n, sat (ω n)}
\end{leansrc}
\declhead{BernsteinLogicAux.Battery}{sound\_QSafetyBattery\_failClosed}{\badge{BadgeGreen}{statement bound}}
\begin{leansrc}
/-- **`sound_QSafetyBattery`, discrete profile, fail-closed form** (the discrete-kernel
profile: the hybrid model flattened to one measurable state space with a single one-step
Markov kernel `K`). Hypotheses
are the checked rows (B1)∧(B2) as the kernel drift row, (B3) floors, (B4) nonnegs, the
side conditions (R3) (hull: `D ⊆ D̄` and one-step mass in `D̄`) and (R4) (fail-closed:
one-step mass in `D ∪ U`), the threshold arithmetic (Θ) `p + thr ≤ 1` (the staged
`q_add`/`q_le` vocabulary, in exact ℚ), and a start state in the concluded region
`D ∩ {V ≤ thr}` (`qraSublevel`). Conclusion: with probability at least `p` the chain
remains in `D \ U` **forever** — never reaching the avoid region and never leaving the
modeled domain. This is the strong (fail-closed) reading of the concluded `QRegionSat`;
`sound_QSafetyBattery` weakens it through the formula tie (F).

The `Integrable V (K y)` hypothesis is the discrete reading of
"domain membership of `V` in the extended generator's domain"; for the intended
instances (`V` polynomial, one-step images in the compact
`D̄` by (R3)) it is automatic. -/
theorem sound_QSafetyBattery_failClosed
    (hD : MeasurableSet D) (hDbar : MeasurableSet Dbar) (hU : MeasurableSet U)
    (hV : Measurable V)
    (hΘ : p + thr ≤ 1)
    (hDsub : D ⊆ Dbar)
    (hR3 : ∀ y ∈ D, y ∉ U → K y Dbar = 1)
    (hR4 : ∀ y ∈ D, y ∉ U → K y (D ∪ U) = 1)
    (hB12int : ∀ y ∈ D, y ∉ U → Integrable V (K y))
    (hB12 : ∀ y ∈ D, y ∉ U → ∫ z, V z ∂K y ≤ V y)
    (hB3 : ∀ y ∈ U, y ∈ Dbar → 1 ≤ V y)
    (hB4 : ∀ y ∈ Dbar, 0 ≤ V y)
    (hx : x ∈ D) (hthr : V x ≤ thr) :
    ENNReal.ofReal p ≤ pathMeasure K x {ω | ∀ n, ω n ∈ D \ U}
\end{leansrc}
\declhead{BernsteinLogicAux.LayerD.Skeleton}{sound\_GSHS\_safety\_battery}{\badge{BadgeGreen}{statement bound}}
\begin{leansrc}
/-- **`sound_QSafetyBattery`, GSHS instance** (T5; lifetime-relative in the
value-presentation reading — and unconditional for W-Z-free models,
those a.s. free of the two stall phenomena, Zeno accumulation and
vanishing into a never-terminating segment): with the threshold arithmetic `p + thr ≤ 1` and a
start whose certificate is under the threshold, the chain remains outside
the avoid region — at every alive entry *and* every halt value — with
probability at least `p`. -/
theorem sound_GSHS_safety_battery [IsMarkovKernel M.seg] [IsMarkovKernel M.jump]
    {V : State → ℝ} {U : Set State} {p thr : ℝ}
    (hV : Measurable V) (hU : MeasurableSet U)
    (hSegD : SegmentDynkin M 0 V) (hForce : ForcedNonexpansion M V)
    (hNonneg : NonnegV V) {C : ℝ} (hC0 : 0 ≤ C) (hVC : ∀ y, V y ≤ C)
    (hFloor : ∀ y ∈ U, 1 ≤ V y)
    (hΘ : p + thr ≤ 1) (x : State) (hthr : V x ≤ thr) :
    ENNReal.ofReal p ≤ Battery.pathMeasure (skel M) (SkelState.alive x)
      {ω | ∀ n, ω n ∉ skelBad U}
\end{leansrc}
\par\medskip
\textbf{Discharging the segment field, corner-immune route.}\enspace
The per-model segment-Dynkin field the statement above consumes is
itself a theorem for piecewise-deterministic segments: the statement
below derives it from upper-Dini rate rows alone --- no
differentiability of the certificate along the flow is assumed
anywhere, so corner points, absolutely-continuous solutions, and
Filippov-style trajectories are in scope exactly as the
module\textquotesingle s header fences state.

\declhead{BernsteinLogicAux.LayerD.DiniDischarge}{segmentDynkin\_of\_pdmp\_dini}{\badge{BadgeGreen}{statement bound}}
\begin{leansrc}
/-- **[SegD] discharged at Dini strength** for models whose entry states
carry terminating Dini flow data. Mixing with the landed smooth data is
free — the field is a conjunction over states, so discharge is pointwise;
this theorem states the all-Dini case, and mixed cases compose by
`rcases` at the call site exactly as `segmentDynkin_of_pdmp` does. -/
theorem segmentDynkin_of_pdmp_dini (M : SegmentModel State) (θ : ℝ)
    (V : State → ℝ) [IsSFiniteKernel M.jump] (hV : Measurable V)
    (h : ∀ x : State, Nonempty (PDMPFlowFinDini M θ V x)) :
    SegmentDynkin M θ V
\end{leansrc}
\declhead{BernsteinLogicAux.LayerD.DiniDischarge}{pdmp\_dynkin\_core\_dini}{\badge{BadgeGreen}{statement bound}}
\begin{leansrc}
/-- **The corner-immune [SegD] core**: the B1 row at upper-right-Dini
strength — stated in the eventual/liminf-slope form the mean-value fence
consumes, so it survives switch corners, Zeno accumulation of corners,
and (through the AC/Filippov reading) sliding — still bounds the
θ-discounted, survival-weighted terminal read plus the accumulated jump
read by the entry value. No derivative of `v` is hypothesized anywhere. -/
theorem pdmp_dynkin_core_dini {v r j : ℝ → ℝ} {θ T : ℝ} (hT : 0 ≤ T)
    (hr : Continuous r) (hj : Continuous j)
    (hv_cont : ContinuousOn v (Icc 0 T))
    (hrow : ∀ t ∈ Ico 0 T, ∀ ρ : ℝ, (r t - θ) * v t - j t < ρ →
      ∃ᶠ z in 𝓝[>] t, slope v t z < ρ) :
    jint r j θ T + survWeight r θ T * v T ≤ v 0
\end{leansrc}
\par\medskip
\textbf{Discharging the segment field, diffusion route.}\enspace For
jump-diffusion segments the same field follows from a single quarantined
per-model input --- the sampled-supermartingale inequality --- stated
one stratum below as an interface structure; the theorem below converts
any inhabitant into the segment-Dynkin field, with a deterministic
positive control included in the module.

\declhead{BernsteinLogicAux.LayerD.ItoInterface}{segmentDynkin\_of\_ito}{\badge{BadgeGreen}{statement bound}}
\begin{leansrc}
/-- **[SegD] discharged for the diffusion profile against the stated
interface** (assault doc §3 row 1(a)): a model all of whose entry states
carry Itô-interface data satisfies the `SegmentDynkin` hypothesis field.
Mixing with PDMP/Dini data is free at call sites, exactly as for
`segmentDynkin_of_pdmp_dini` — the field is a conjunction over states,
discharged pointwise. -/
theorem segmentDynkin_of_ito (M : SegmentModel State) (θ : ℝ)
    (V : State → ℝ) [IsSFiniteKernel M.jump] (hV : Measurable V)
    (h : ∀ x : State, Nonempty (ItoInterface M θ V x)) :
    SegmentDynkin M θ V
\end{leansrc}
\par\medskip
\textbf{The interface\textquotesingle s own discharge.}\enspace The
quarantined inequality is in turn delivered by optional sampling: from a
supermartingale property of the certificate process (plus the
d\'ebut-measurability family), the assembly below inhabits the full
interface --- leaving, per diffusion entry state, exactly the
supermartingale property and the d\'ebut family. The blocks that follow
discharge each in turn.

\declhead{BernsteinLogicAux.LayerD.ItoSampling}{ItoInterface.ofSupermartingale}{\badge{BadgeGreen}{statement bound}}
\begin{leansrc}
/-- **The interface, discharged down to the supermartingale property**:
a full `ItoInterface` instance from the non-analytic pathwise data plus
`Supermartingale` + the début family. Composing with
`segmentDynkin_of_ito`, the chain
`Supermartingale + début family ⟹ [SegD]` is Lean end-to-end; what a
concrete diffusion model still owes is exactly the supermartingale
property (theorem 6's shape / upstream Itô). -/
noncomputable def ItoInterface.ofSupermartingale
    (M : SegmentModel State) [IsProbabilityMeasure P]
    (hT : ∀ ω, 0 ≤ T ω)
    (hk : ∀ ω, k ω = ExitKind.guard ∨ k ω = ExitKind.halt
      ∨ k ω = ExitKind.live)
    (hγ : ∀ ω, Measurable (γ ω))
    (hvc : ∀ ω, Continuous fun u => V (γ ω u))
    (hVγ : ∀ ω u, 0 ≤ V (γ ω u))
    (hr : ∀ ω, Continuous (r ω)) (hr0 : ∀ ω u, 0 ≤ r ω u)
    (hjr : ∀ ω, Continuous (jr ω)) (hjr0 : ∀ ω u, 0 ≤ jr ω u)
    (hjr_couple : ∀ ω, ∀ u ∈ segDom (T ω) (k ω),
      ENNReal.ofReal (jr ω u)
        = ENNReal.ofReal (r ω u)
          * ∫⁻ z, ENNReal.ofReal (V z) ∂(M.jump (γ ω u, ExitKind.alarm)))
    (hread : ∫⁻ o, outcomeRead M θ V o ∂(M.seg x)
      ≤ ∫⁻ ω, ∫⁻ o, outcomeRead M θ V o
          ∂(segMeasureIto (γ ω) (r ω) (T ω) (k ω)) ∂P)
    (hγ0 : ∀ ω, γ ω 0 = x)
    (hZ : Supermartingale (certProcess θ V γ r jr) ℱ P)
    (hτ : ∀ s : ℝ≥0, MeasurableSet[ℱ s]
      {ω | exitDebut T k ω ≤ (s : ℝ≥0∞)}) :
    ItoInterface M θ V x
\end{leansrc}
\declhead{BernsteinLogicAux.LayerD.ItoSampling}{sampled\_le\_of\_supermartingale}{\badge{BadgeGreen}{statement bound}}
\begin{leansrc}
/-- **Partial discharge of (A.1)** from the banked optional-sampling
suite: the supermartingale property of the certificate process (the
[WR-2]/[WR-3] residue) plus the exit début's measurability family imply
the `hsampled` field, with the start pinned at `x` and `P` a probability
measure supplying `Y₀ = V x`. -/
theorem sampled_le_of_supermartingale [IsProbabilityMeasure P]
    (hT : ∀ ω, 0 ≤ T ω)
    (hvc : ∀ ω, Continuous fun u => V (γ ω u))
    (hVγ : ∀ ω u, 0 ≤ V (γ ω u))
    (hr : ∀ ω, Continuous (r ω)) (hjr : ∀ ω, Continuous (jr ω))
    (hjr0 : ∀ ω u, 0 ≤ jr ω u)
    (hγ0 : ∀ ω, γ ω 0 = x)
    (hZ : Supermartingale (certProcess θ V γ r jr) ℱ P)
    (hτ : ∀ s : ℝ≥0, MeasurableSet[ℱ s]
      {ω | exitDebut T k ω ≤ (s : ℝ≥0∞)}) :
    ∀ t : ℝ, 0 ≤ t →
      ∫⁻ ω, sampledRead θ V (γ ω) (r ω) (jr ω) (T ω) (k ω) t ∂P
        ≤ ENNReal.ofReal (V x)
\end{leansrc}
\par\medskip
\textbf{The residue, bundled.}\enspace The per-model inputs the two
routes leave open are packaged as a structure whose only analytic
fields are the certificate-process supermartingale property and the
exit d\'ebut\textquotesingle s measurability family; the theorems below
convert per-entry-state families of such bundles (all-diffusion, or a
per-state diffusion/Dini mixture) into the segment field, and the class
constructor builds the battery\textquotesingle s Ville class with that
field proven rather than assumed.

\declhead{BernsteinLogicAux.GshsItoDischarge}{ItoSupermartingaleData}{\badge{BadgeGreen}{statement bound}}
\begin{leansrc}
/-- **The supermartingale-floor bundle**: exactly
`ItoInterface.ofSupermartingale`'s inputs as a structure, so that
per-entry-state families can be quantified. The two analytic fields are
`hZ` and `hτ`; everything else is the pathwise bookkeeping the interface
already carried. -/
structure ItoSupermartingaleData (M : SegmentModel State) (θ : ℝ)
    (V : State → ℝ) (x : State) where
  /-- the driving path-data space -/
  Ω : Type*
  /-- its σ-algebra -/
  [mΩ : MeasurableSpace Ω]
  /-- the law of the driving data issued from `x` -/
  P : Measure Ω
  /-- `P` is a probability measure (`Y₀ = V x` is the model's business —
  the optional-sampling route consumes this) -/
  hprob : IsProbabilityMeasure P
  /-- the filtration the certificate process rides -/
  ℱ : Filtration ℝ≥0 mΩ
  /-- exit time of each path (unread on live paths) -/
  T : Ω → ℝ
  hT : ∀ ω, 0 ≤ T ω
  /-- terminator species of each path -/
  k : Ω → ExitKind
  hk : ∀ ω, k ω = ExitKind.guard ∨ k ω = ExitKind.halt
    ∨ k ω = ExitKind.live
  /-- the trajectory of each path (clamped extension off the segment) -/
  γ : Ω → ℝ → State
  hγ : ∀ ω, Measurable (γ ω)
  hvc : ∀ ω, Continuous fun u => V (γ ω u)
  hVγ : ∀ ω u, 0 ≤ V (γ ω u)
  /-- paths start at the entry state -/
  hγ0 : ∀ ω, γ ω 0 = x
  /-- alarm intensity along each path -/
  r : Ω → ℝ → ℝ
  hr : ∀ ω, Continuous (r ω)
  hr0 : ∀ ω u, 0 ≤ r ω u
  /-- λ-weighted post-jump read along each path -/
  jr : Ω → ℝ → ℝ
  hjr : ∀ ω, Continuous (jr ω)
  hjr0 : ∀ ω u, 0 ≤ jr ω u
  hjr_couple : ∀ ω, ∀ u ∈ segDom (T ω) (k ω),
    ENNReal.ofReal (jr ω u)
      = ENNReal.ofReal (r ω u)
        * ∫⁻ z, ENNReal.ofReal (V z) ∂(M.jump (γ ω u, ExitKind.alarm))
  /-- the read-level disintegration of `M.seg x` along paths
  (definitional for `Rep`-constructed models) -/
  hread : ∫⁻ o, outcomeRead M θ V o ∂(M.seg x)
    ≤ ∫⁻ ω, ∫⁻ o, outcomeRead M θ V o
        ∂(segMeasureIto (γ ω) (r ω) (T ω) (k ω)) ∂P
  /-- **analytic field 1**: the certificate process is a supermartingale
  — the exact [WR-2]/[WR-3] residue -/
  hZ : Supermartingale (certProcess θ V γ r jr) ℱ P
  /-- **analytic field 2**: the exit début's measurability family -/
  hτ : ∀ s : ℝ≥0, MeasurableSet[ℱ s]
    {ω | exitDebut T k ω ≤ (s : ℝ≥0∞)}
\end{leansrc}
\declhead{BernsteinLogicAux.GshsItoDischarge}{segmentDynkin\_of\_supermartingale}{\badge{BadgeGreen}{statement bound}}
\begin{leansrc}
/-- **[SegD] from the supermartingale floor, all-diffusion form**: a
model all of whose entry states carry supermartingale-floor bundles
satisfies the `SegmentDynkin` hypothesis field. -/
theorem segmentDynkin_of_supermartingale (M : SegmentModel State) (θ : ℝ)
    (V : State → ℝ) [IsSFiniteKernel M.jump] (hV : Measurable V)
    (h : ∀ x : State, Nonempty (ItoSupermartingaleData M θ V x)) :
    SegmentDynkin M θ V
\end{leansrc}
\declhead{BernsteinLogicAux.GshsItoDischarge}{segmentDynkin\_of\_mixed}{\badge{BadgeGreen}{statement bound}}
\begin{leansrc}
/-- **The pointwise diffusion/Dini mixture, enacted** (both discharge
docstrings promised this "free at call sites"; here it is as one
theorem): each entry state carries either a supermartingale-floor bundle
(diffusion segments) or terminating Dini flow data (deterministic/PDMP
segments, switch corners included) — mixed freely, [SegD] follows. Two
disjuncts only: Dini is the deterministic floor (smooth data restates at
it, or lifts by `dini_row_of_deriv_row`). -/
theorem segmentDynkin_of_mixed (M : SegmentModel State) (θ : ℝ)
    (V : State → ℝ) [IsSFiniteKernel M.jump] (hV : Measurable V)
    (h : ∀ x : State, Nonempty (ItoSupermartingaleData M θ V x)
      ∨ Nonempty (PDMPFlowFinDini M θ V x)) :
    SegmentDynkin M θ V
\end{leansrc}
\declhead{BernsteinLogicAux.GshsItoDischarge}{GshsSegmentVille.ofDischarge}{\badge{BadgeGreen}{statement bound}}
\begin{leansrc}
/-- **The discharge constructor — `gshsGen:segment-dynkin`'s diffusion
sub-case retired to the floor**: the rung's class with its quarantined
field PROVEN from row-conditional per-state discharge data, never
assumed. `hdata`'s conditional shape mirrors the field's own: the
certified θ-discounted generator rows are where a concrete model's
analytic side takes its input (Itô's formula consumes `𝒜V ≤ 0` on the
cells to prove `hZ`; the Dini row data is the compiled rows themselves).
What a model owes per diffusion entry state is exactly `hZ` +`hτ` — the
`sdeGen:dynkin-supermartingale` route-(i) floor. -/
noncomputable def GshsSegmentVille.ofDischarge
    [MeasurableSpace Mode] [DecidableEq Mode]
    {P : GshsProfile n m T Mode} {V : Mode → State n → ℝ}
    {gradV : Mode → State n → State n}
    {hessV : Mode → State n → Fin n → Fin n → ℝ} {θ : ℝ}
    (M : LayerD.SegmentModel (Flat n Mode)) [IsSFiniteKernel M.jump]
    (rd : Flat n Mode → Flat n Mode)
    (jump_guard_det : ∀ p, M.jump (p, LayerD.ExitKind.guard)
      = Measure.dirac (rd p))
    (dispatch_cover : ∀ p : Flat n Mode,
      (∃ e : Fin T, (P.trans e).forced = true ∧ (P.trans e).src = p.1
        ∧ p.2 ∈ (P.trans e).guard
        ∧ rd p = ((P.trans e).dst, (P.trans e).reset p.2))
      ∨ rd p = p)
    (hV : Measurable (flatCert V))
    (hdata : (∀ mo : Mode, ∀ z ∈ P.cell mo,
        θ * V mo z + generator P V gradV hessV mo z ≤ 0) →
      ∀ p : Flat n Mode,
        Nonempty (LayerD.ItoSupermartingaleData M θ (flatCert V) p)
        ∨ Nonempty (LayerD.PDMPFlowFinDini M θ (flatCert V) p)) :
    GshsSegmentVille P V gradV hessV θ
\end{leansrc}
\par\medskip
\textbf{Discharging the d\'ebut family.}\enspace The second analytic
field is a library application: exit from a closed certified cell is
the d\'ebut of the open complement, and the d\'ebut of an open set
under right-continuous paths is a stopping time of the
right-continuous filtration --- no completion, no metric, no per-model
measure theory.

\declhead{BernsteinLogicAux.Debut}{isStoppingTime\_debut}{\badge{BadgeGreen}{statement bound}}
\begin{leansrc}
/-- The début of an open set by a right-continuous adapted process is a
bona-fide mathlib **stopping time** for any right-continuous filtration
(`ℝ≥0∞ := WithTop ℝ≥0` definitionally, so no translation layer): the pin's
`isStoppingTime_of_measurableSet_lt_of_isRightContinuous` upgrades the
`<`-family. This is the wide-sense-to-strict step the census's G3 entry
priced; it costs an instance, never a per-model proof — models on a raw
filtration `𝓕` use `𝓕.rightCont` (`𝓕₊`), whose `IsRightContinuous`
instance mathlib provides. -/
theorem isStoppingTime_debut [ℱ.IsRightContinuous]
    (hX : ∀ t : ℝ≥0, Measurable[ℱ t] (X t))
    (hrc : ∀ ω, ∀ s : ℝ≥0, ContinuousWithinAt (fun u => X u ω) (Set.Ici s) s)
    (hU : IsOpen U) :
    IsStoppingTime ℱ (debut X U)
\end{leansrc}
\declhead{BernsteinLogicAux.Debut}{exitDebut\_le\_family\_of\_isClosed\_cell}{\badge{BadgeGreen}{statement bound}}
\begin{leansrc}
/-- The certified-cell instantiation: exit from a CLOSED cell (the
box-union geometry — finite unions of closed boxes are closed by
`isClosed_Icc`/`isClosed_set_pi` + `Set.Finite.isClosed_biUnion`, one line
per model) is the début of the open complement, so the `hτ` field is a
library application. -/
theorem exitDebut_le_family_of_isClosed_cell [ℱ.IsRightContinuous]
    {T : Ω → ℝ} {k : Ω → ExitKind} {X : ℝ≥0 → Ω → E} {cell : Set E}
    (hcell : IsClosed cell)
    (hEq : ∀ ω, exitDebut T k ω = debut X cellᶜ ω)
    (hX : ∀ t : ℝ≥0, Measurable[ℱ t] (X t))
    (hrc : ∀ ω, ∀ s : ℝ≥0, ContinuousWithinAt (fun u => X u ω) (Set.Ici s) s) :
    ∀ s : ℝ≥0, MeasurableSet[ℱ s] {ω | exitDebut T k ω ≤ (s : ℝ≥0∞)}
\end{leansrc}
\par\medskip
\textbf{Discharging the supermartingale property, generic half.}\enspace
The first analytic field\textquotesingle s martingale theory is also a
library application: the exponential-discounting theorem produces the
required supermartingale from a martingale-problem hypothesis (the
compensated certificate process is a martingale --- what It\^o
\textquotesingle s formula and the interior rows establish per model,
with no stochastic integration in the certified statement), the
bounded-horizon and horizon-free optional-sampling inequalities are the
theorems the interface discharge above consumes, and the continuous-time
Doob/Ville maximal inequality converts supermartingale level crossings
into probability bounds. What remains per model is exactly the
martingale-problem hypothesis and the uniform bounds it consumes.

\declhead{BernsteinLogicAux.OptionalStoppingRC}{supermartingale\_exp\_discount}{\badge{BadgeGreen}{statement bound}}
\begin{leansrc}
/-- **Exponential discounting of a drift-dominated process.** Suppose
`W : ℝ≥0 → Ω → ℝ` is adapted with continuous paths, `g : ℝ≥0 → Ω → ℝ` is
jointly measurable, both are uniformly bounded, the compensated process
`fun t ω => W t ω - ∫ s in (0 : ℝ)..(t : ℝ), g s.toNNReal ω` is a
martingale (the martingale-problem formulation of `g` being the generator
drift of `W`), and `g t ω ≤ c * W t ω` pointwise. Then

`fun t ω => Real.exp (-c * t) * W t ω` is a supermartingale.

This is the discounting step of Dynkin-formula arguments, carried out with
pathwise Riemann–Stieltjes integration by parts against the `C¹` function
`s ↦ exp (-c * s)` and a conditional Fubini argument — no stochastic
integral is involved. -/
theorem supermartingale_exp_discount
    [IsProbabilityMeasure μ]
    (W : ℝ≥0 → Ω → ℝ) (g : ℝ≥0 → Ω → ℝ) (c : ℝ) (C : ℝ)
    (hWadp : Adapted ℱ W)
    (hWcont : ∀ ω, Continuous fun t => W t ω)
    (hWb : ∀ t ω, |W t ω| ≤ C)
    (hgmeas : Measurable (Function.uncurry g))
    (hgb : ∀ t ω, |g t ω| ≤ C)
    (hmart : Martingale
      (fun (t : ℝ≥0) (ω : Ω) =>
        W t ω - ∫ s in (0 : ℝ)..(t : ℝ), g s.toNNReal ω) ℱ μ)
    (hdom : ∀ t ω, g t ω ≤ c * W t ω) :
    Supermartingale
      (fun (t : ℝ≥0) (ω : Ω) => Real.exp (-c * (t : ℝ)) * W t ω) ℱ μ
\end{leansrc}
\declhead{BernsteinLogicAux.OptionalStoppingRC}{lintegral\_stoppedValue\_min\_le}{\badge{BadgeGreen}{statement bound}}
\begin{leansrc}
/-- The upper dyadic approximation to a nonnegative real. -/
noncomputable def dyadicCeil (n : ℕ) (x : ℝ≥0) : ℕ :=
  ⌈x * (2 ^ n : ℝ≥0)⌉₊

lemma dyadicCeil_time_tendsto (x : ℝ≥0) :
    Filter.Tendsto (fun n => dyadicTime n (dyadicCeil n x)) Filter.atTop (nhds x) := by
  -- Use the fact that $x \leq \lceil x \cdot 2^n \rceil / 2^n \leq x + 1 / 2^n$.
  have h_bounds : ∀ n : ℕ, x ≤ (dyadicCeil n x : ℝ≥0) / 2^n ∧ (dyadicCeil n x : ℝ≥0) / 2^n ≤ x + 1 / 2^n := by
    intro n
    simp [dyadicCeil, dyadicTime];
    rw [ le_div_iff₀, div_le_iff₀ ] <;> norm_num;
    exact ⟨ Nat.le_ceil _, by rw [ add_mul, inv_mul_cancel₀ ( by positivity ) ] ; exact_mod_cast Nat.ceil_lt_add_one ( by positivity ) |> le_of_lt ⟩;
  refine' tendsto_of_tendsto_of_tendsto_of_le_of_le' tendsto_const_nhds _ ( Filter.Eventually.of_forall fun n => h_bounds n |>.1 ) ( Filter.Eventually.of_forall fun n => h_bounds n |>.2 );
  convert tendsto_const_nhds.add ( tendsto_inv_atTop_zero.comp ( tendsto_pow_atTop_atTop_of_one_lt ( show ( 1 : ℝ≥0 ) < 2 by norm_num ) ) ) using 2 ; norm_num;
  exacts [ rfl, by rw [ add_zero ] ]

lemma dyadicCeil_time_mem_Ici (n : ℕ) (x : ℝ≥0) :
    dyadicTime n (dyadicCeil n x) ∈ Set.Ici x := by
  unfold dyadicTime dyadicCeil;
  exact Set.mem_Ici.mpr ( by rw [ le_div_iff₀ ( by positivity ) ] ; exact Nat.le_ceil _ )

lemma dyadicCeil_le_iff (n k : ℕ) (x : ℝ≥0) :
    dyadicCeil n x ≤ k ↔ x ≤ dyadicTime n k := by
  rw [dyadicCeil, Nat.ceil_le, dyadicTime, le_div_iff₀ (by positivity)]

lemma dyadicCeil_stoppingTime
    (τ : Ω → ℝ≥0) (hτ : ∀ t : ℝ≥0, MeasurableSet[ℱ t] {ω | τ ω ≤ t})
    (T : ℝ≥0) (n : ℕ) :
    IsStoppingTime (dyadicFiltration (ℱ := ℱ) n)
      (fun ω => (dyadicCeil n (min (τ ω) T) : ℕ∞)) := by
  intro k
  by_cases hT : T ≤ dyadicTime n k
  · refine MeasurableSet.univ.congr ?_
    refine (Set.eq_univ_of_forall fun ω => ?_).symm
    show (dyadicCeil n (min (τ ω) T) : ℕ∞) ≤ (k : ℕ∞)
    exact_mod_cast (dyadicCeil_le_iff n k (min (τ ω) T)).mpr ((min_le_right _ _).trans hT)
  · refine (hτ (dyadicTime n k)).congr ?_
    ext ω
    simp only [Set.mem_setOf_eq]
    constructor
    · intro h
      show (dyadicCeil n (min (τ ω) T) : ℕ∞) ≤ (k : ℕ∞)
      exact_mod_cast (dyadicCeil_le_iff n k (min (τ ω) T)).mpr ((min_le_left _ _).trans h)
    · intro h
      have h' : (dyadicCeil n (min (τ ω) T) : ℕ∞) ≤ (k : ℕ∞) := h
      have hnk : dyadicCeil n (min (τ ω) T) ≤ k := by exact_mod_cast h'
      have := (dyadicCeil_le_iff n k (min (τ ω) T)).mp hnk
      exact ((min_le_iff.mp this).resolve_right hT)

lemma dyadicCeil_bounded (τ : Ω → ℝ≥0) (T : ℝ≥0) (n : ℕ) (ω : Ω) :
    (dyadicCeil n (min (τ ω) T) : ℕ∞) ≤ dyadicCeil n T := by
  exact_mod_cast Nat.ceil_mono ( mul_le_mul_of_nonneg_right ( min_le_right _ _ ) ( by positivity ) )

lemma integral_dyadic_stopped_le
    [IsProbabilityMeasure μ]
    (Z : ℝ≥0 → Ω → ℝ) (hZ : Supermartingale Z ℱ μ)
    (τ : Ω → ℝ≥0)
    (hτ : ∀ t : ℝ≥0, MeasurableSet[ℱ t] {ω | τ ω ≤ t})
    (T : ℝ≥0) (n : ℕ) :
    ∫ ω, Z (dyadicTime n (dyadicCeil n (min (τ ω) T))) ω ∂μ
      ≤ ∫ ω, Z 0 ω ∂μ := by
  have h_submartingale :
      Submartingale (fun k => -Z (dyadicTime n k)) (dyadicFiltration (ℱ := ℱ) n) μ :=
    (supermartingale_dyadic Z hZ n).neg
  have h_stopping_time : IsStoppingTime (dyadicFiltration (ℱ := ℱ) n)
      (fun ω => (dyadicCeil n (min (τ ω) T) : ℕ∞)) := dyadicCeil_stoppingTime τ hτ T n
  have h0 : IsStoppingTime (dyadicFiltration (ℱ := ℱ) n) (fun _ : Ω => ((0 : ℕ) : ℕ∞)) := by
    intro k
    refine MeasurableSet.univ.congr ?_
    refine (Set.eq_univ_of_forall fun ω => ?_).symm
    show ((0 : ℕ) : ℕ∞) ≤ (k : ℕ∞)
    exact_mod_cast Nat.zero_le k
  have key := h_submartingale.expected_stoppedValue_mono h0 h_stopping_time
    (fun ω => by
      show ((0 : ℕ) : ℕ∞) ≤ (dyadicCeil n (min (τ ω) T) : ℕ∞)
      exact_mod_cast Nat.zero_le _)
    (N := dyadicCeil n T) (fun ω => by exact_mod_cast dyadicCeil_bounded τ T n ω)
  have hL : stoppedValue (fun k => -Z (dyadicTime n k)) (fun _ : Ω => ((0 : ℕ) : ℕ∞))
      = fun ω => -Z (dyadicTime n 0) ω := rfl
  have hd0 : dyadicTime n 0 = 0 := by simp [dyadicTime]
  have hR : stoppedValue (fun k => -Z (dyadicTime n k))
      (fun ω => (dyadicCeil n (min (τ ω) T) : ℕ∞))
      = fun ω => -Z (dyadicTime n (dyadicCeil n (min (τ ω) T))) ω := rfl
  rw [hL, hR, hd0] at key
  simp only [integral_neg] at key
  exact neg_le_neg_iff.mp key

lemma integrable_dyadic_stopped
    [IsProbabilityMeasure μ]
    (Z : ℝ≥0 → Ω → ℝ) (hZ : Supermartingale Z ℱ μ)
    (τ : Ω → ℝ≥0)
    (hτ : ∀ t : ℝ≥0, MeasurableSet[ℱ t] {ω | τ ω ≤ t})
    (T : ℝ≥0) (n : ℕ) :
    Integrable (fun ω => Z (dyadicTime n (dyadicCeil n (min (τ ω) T))) ω) μ := by
  have h_sub : Submartingale (fun k => -Z (dyadicTime n k)) (dyadicFiltration (ℱ := ℱ) n) μ :=
    (supermartingale_dyadic Z hZ n).neg
  have h_int := h_sub.integrable_stoppedValue (dyadicCeil_stoppingTime τ hτ T n)
    (N := dyadicCeil n T) (fun ω => by exact_mod_cast dyadicCeil_bounded τ T n ω)
  have hR : stoppedValue (fun k => -Z (dyadicTime n k))
      (fun ω => (dyadicCeil n (min (τ ω) T) : ℕ∞))
      = fun ω => -Z (dyadicTime n (dyadicCeil n (min (τ ω) T))) ω := rfl
  rw [hR] at h_int
  simpa using h_int.neg

lemma dyadic_stopped_tendsto
    (Z : ℝ≥0 → Ω → ℝ)
    (hZrc : ∀ ω, ∀ t : ℝ≥0,
      ContinuousWithinAt (fun s => Z s ω) (Set.Ici t) t)
    (τ : Ω → ℝ≥0) (T : ℝ≥0) (ω : Ω) :
    Filter.Tendsto
      (fun n => Z (dyadicTime n (dyadicCeil n (min (τ ω) T))) ω)
      Filter.atTop (nhds (Z (min (τ ω) T) ω)) := by
  refine' Filter.Tendsto.comp ( hZrc ω ( min ( τ ω ) T ) ) _;
  rw [ tendsto_nhdsWithin_iff ];
  exact ⟨ dyadicCeil_time_tendsto _, Filter.Eventually.of_forall fun n => dyadicCeil_time_mem_Ici _ _ ⟩

/-
**Optional sampling inequality at a bounded horizon.** For a nonnegative
supermartingale `Z` over `ℝ≥0` with right-continuous paths, a stopping time
`τ` (given by its measurability family) and a deterministic horizon `T`,

`∫⁻ ω, ENNReal.ofReal (Z (min (τ ω) T) ω) ∂μ ≤ ∫⁻ ω, ENNReal.ofReal (Z 0 ω) ∂μ`.

Stated with the lower Lebesgue integral so that no integrability of the
stopped value needs to be assumed.
-/
theorem lintegral_stoppedValue_min_le
    [IsProbabilityMeasure μ]
    (Z : ℝ≥0 → Ω → ℝ) (hZ : Supermartingale Z ℱ μ)
    (hZnn : ∀ t ω, 0 ≤ Z t ω)
    (hZrc : ∀ ω, ∀ t : ℝ≥0,
      ContinuousWithinAt (fun s => Z s ω) (Set.Ici t) t)
    (τ : Ω → ℝ≥0)
    (hτ : ∀ t : ℝ≥0, MeasurableSet[ℱ t] {ω | τ ω ≤ t})
    (T : ℝ≥0) :
    ∫⁻ ω, ENNReal.ofReal (Z (min (τ ω) T) ω) ∂μ
      ≤ ∫⁻ ω, ENNReal.ofReal (Z 0 ω) ∂μ
\end{leansrc}
\declhead{BernsteinLogicAux.OptionalStoppingRC}{lintegral\_stoppedValue\_le}{\badge{BadgeGreen}{statement bound}}
\begin{leansrc}
/-- Truncate a wide-sense nonnegative stopping time at a finite horizon. -/
noncomputable def ennrealTrunc (τ : Ω → ℝ≥0∞) (T : ℝ≥0) (ω : Ω) : ℝ≥0 :=
  if τ ω ≤ (T : ℝ≥0∞) then (τ ω).toNNReal else T

lemma ennrealTrunc_le_iff (τ : Ω → ℝ≥0∞) (T t : ℝ≥0) (ht : t < T) :
    {ω | ennrealTrunc τ T ω ≤ t} = {ω | τ ω ≤ (t : ℝ≥0∞)} := by
  ext ω;
  by_cases h : τ ω ≤ T <;> simp +decide [ h, ennrealTrunc ];
  · cases h' : τ ω <;> aesop;
  · exact iff_of_false ( not_le_of_gt ht ) ( fun h' => h <| le_trans h' <| mod_cast ht.le )

lemma ennrealTrunc_eq_of_le (τ : Ω → ℝ≥0∞) (T : ℝ≥0) (ω : Ω)
    (h : τ ω ≤ (T : ℝ≥0∞)) :
    ennrealTrunc τ T ω = (τ ω).toNNReal := by
  exact if_pos h

/-
**Optional sampling inequality, horizon-free form.** For a nonnegative
right-continuous supermartingale `Z` and a wide-sense stopping time `τ` with
values in `ℝ≥0∞` (given by its measurability family), the contribution of
the event `{τ < ∞}` satisfies

`∫⁻ ω, (if τ ω = ∞ then 0 else ENNReal.ofReal (Z (τ ω).toNNReal ω)) ∂μ
   ≤ ∫⁻ ω, ENNReal.ofReal (Z 0 ω) ∂μ`.

Follows from `lintegral_stoppedValue_min_le` by monotone convergence along
an increasing sequence of horizons, using nonnegativity to discard the
`{τ > T}` remainder.
-/
theorem lintegral_stoppedValue_le
    [IsProbabilityMeasure μ]
    (Z : ℝ≥0 → Ω → ℝ) (hZ : Supermartingale Z ℱ μ)
    (hZnn : ∀ t ω, 0 ≤ Z t ω)
    (hZrc : ∀ ω, ∀ t : ℝ≥0,
      ContinuousWithinAt (fun s => Z s ω) (Set.Ici t) t)
    (τ : Ω → ℝ≥0∞)
    (hτ : ∀ t : ℝ≥0, MeasurableSet[ℱ t] {ω | τ ω ≤ (t : ℝ≥0∞)}) :
    ∫⁻ ω, (if τ ω = ∞ then 0
           else ENNReal.ofReal (Z (τ ω).toNNReal ω)) ∂μ
      ≤ ∫⁻ ω, ENNReal.ofReal (Z 0 ω) ∂μ
\end{leansrc}
\declhead{BernsteinLogicAux.OptionalStoppingRC}{measure\_sup\_ge\_le\_of\_supermartingale}{\badge{BadgeGreen}{statement bound}}
\begin{leansrc}
/-- Restrict a supermartingale along a monotone map of time indices. -/
private lemma Supermartingale.comp_monotone
    {ι κ : Type*} [Preorder ι] [Preorder κ]
    {ℱι : Filtration ι m0} {Z : ι → Ω → ℝ}
    (hZ : Supermartingale Z ℱι μ) (f : κ → ι) (hf : Monotone f) :
    Supermartingale (fun k ↦ Z (f k))
      ({ seq := fun k ↦ ℱι (f k)
         mono' := fun _ _ h ↦ ℱι.mono (hf h)
         le' := fun k ↦ ℱι.le (f k) } : Filtration κ m0) μ := by
  exact ⟨fun k ↦ hZ.stronglyAdapted (f k),
    fun i j hij ↦ hZ.condExp_ae_le (hf hij), fun k ↦ hZ.integrable (f k)⟩

/-
A continuous path hitting `c` on `[0,T]` is approximated on uniform grids
at every strictly smaller level.
-/
private lemma exists_grid_hit_of_continuous
    (z : ℝ≥0 → ℝ) (hz : Continuous z) (T : ℝ≥0) (c d : ℝ) (hdc : d < c)
    (h : ∃ s : ℝ≥0, s ≤ T ∧ c ≤ z s) :
    ∃ m : ℕ, 0 < m ∧ ∃ k ≤ m,
      d ≤ z (⟨(k : ℝ) * T / m, by positivity⟩ : ℝ≥0) := by
  by_contra h_contra;
  obtain ⟨ s, hs₁, hs₂ ⟩ := h;
  -- By the continuity of $z$, for any $\epsilon > 0$, there exists a $\delta > 0$ such that $|z(t) - z(s)| < \epsilon$ whenever $|t - s| < \delta$.
  obtain ⟨δ, hδ_pos, hδ⟩ : ∃ δ > 0, ∀ t : ℝ≥0, |(t : ℝ) - s| < δ → |z t - z s| < (c - d) / 2 := by
    obtain ⟨δ, hδ, hδ'⟩ :=
      Metric.continuous_iff.mp hz s ((c - d) / 2) (half_pos (sub_pos.mpr hdc))
    refine ⟨δ, hδ, fun t ht ↦ ?_⟩
    simpa only [Real.dist_eq] using hδ' t (by simpa only [NNReal.dist_eq] using ht);
  -- Choose $m$ large enough such that $T/m < \delta$.
  obtain ⟨m, hm_pos, hm⟩ : ∃ m : ℕ, 0 < m ∧ (T : ℝ) / m < δ := by
    exact ⟨ ⌊T / δ⌋₊ + 1, Nat.succ_pos _, by rw [ div_lt_iff₀ ] <;> push_cast <;> nlinarith [ Nat.lt_floor_add_one ( T / δ ), mul_div_cancel₀ ( T : ℝ ) hδ_pos.ne' ] ⟩;
  -- Choose $k$ such that $k * T / m$ is within $\delta$ of $s$.
  obtain ⟨k, hk⟩ : ∃ k : ℕ, k ≤ m ∧ |(k * T / m : ℝ) - s| < δ := by
    refine' ⟨ ⌊ ( s : ℝ ) * m / T⌋₊, _, _ ⟩;
    · exact Nat.floor_le_of_le ( div_le_of_le_mul₀ ( by positivity ) ( by positivity ) ( by nlinarith [ show ( s : ℝ ) ≤ T from hs₁, show ( m : ℝ ) ≥ 1 by norm_cast ] ) );
    · by_cases hT : T = 0 <;> simp_all +decide [div_lt_iff₀];
      rw [ abs_lt ] ; constructor <;> nlinarith [ Nat.floor_le ( show 0 ≤ ( s : ℝ ) * m / T by positivity ), Nat.lt_floor_add_one ( ( s : ℝ ) * m / T ), show ( T : ℝ ) > 0 by positivity, mul_div_cancel₀ ( ( s : ℝ ) * m ) ( show ( T : ℝ ) ≠ 0 by positivity ), mul_div_cancel₀ ( ( ⌊ ( s : ℝ ) * m / T⌋₊ : ℝ ) * T ) ( show ( m : ℝ ) ≠ 0 by positivity ) ];
  exact h_contra ⟨ m, hm_pos, k, hk.1, by linarith [ abs_lt.mp ( hδ ⟨ k * T / m, by positivity ⟩ hk.2 ) ] ⟩

/-
**Doob-type maximal inequality, continuous time.** For a nonnegative
supermartingale `Z` over the index set `ℝ≥0` with continuous paths, the
probability that the running supremum over `[0, T]` reaches a level `c > 0`
is at most `(∫ ω, Z 0 ω ∂μ) / c`.
-/
theorem measure_sup_ge_le_of_supermartingale
    [IsProbabilityMeasure μ]
    (Z : ℝ≥0 → Ω → ℝ) (hZ : Supermartingale Z ℱ μ)
    (hZnn : ∀ t ω, 0 ≤ Z t ω)
    (hZcont : ∀ ω, Continuous fun t => Z t ω)
    (T : ℝ≥0) (c : ℝ) (hc : 0 < c) :
    μ {ω | ∃ s : ℝ≥0, s ≤ T ∧ c ≤ Z s ω}
      ≤ ENNReal.ofReal ((∫ ω, Z 0 ω ∂μ) / c)
\end{leansrc}
\par\medskip
\textbf{The martingale-problem profile, end to end.}\enspace The
general continuous-time theorems: the inter-jump motion enters through
its generator --- $\mathcal{A}V$ as data, tied to the process only by
the martingale hypothesis for the discounted compensated certificate
process --- the segment field is derived rather than assumed, and the
battery conclusion follows for every system whose presentation
satisfies the martingale-problem hypothesis.

\declhead{QtslAux.LayerD.MPDischarge}{sound\_GSHS\_safety\_mp}{\badge{BadgeGreen}{statement bound}}
\begin{leansrc}
/-- Safety (avoidance) at the martingale-problem profile: the probability
that the segment skeleton ever enters the bad set is at most `V x`. -/
theorem sound_GSHS_safety_mp [IsMarkovKernel M.seg] [IsMarkovKernel M.jump]
    {V : State → ℝ} {U : Set State} (hV : Measurable V) (hU : MeasurableSet U)
    (hmp : ∀ x : State, Nonempty (MPFlow M 0 V x))
    (hForce : ForcedNonexpansion M V)
    (hNonneg : NonnegV V) {C : ℝ} (hC0 : 0 ≤ C) (hVC : ∀ y, V y ≤ C)
    (hFloor : ∀ y ∈ U, 1 ≤ V y) (x : State) :
    Battery.pathMeasure (skel M) (SkelState.alive x)
        {ω | ∃ n, ω n ∈ skelBad U}
      ≤ ENNReal.ofReal (V x)
\end{leansrc}
\declhead{QtslAux.LayerD.MPDischarge}{sound\_GSHS\_safety\_battery\_mp}{\badge{BadgeGreen}{statement bound}}
\begin{leansrc}
/-- The safety battery at the martingale-problem profile. -/
theorem sound_GSHS_safety_battery_mp
    [IsMarkovKernel M.seg] [IsMarkovKernel M.jump]
    {V : State → ℝ} {U : Set State} {p thr : ℝ}
    (hV : Measurable V) (hU : MeasurableSet U)
    (hmp : ∀ x : State, Nonempty (MPFlow M 0 V x))
    (hForce : ForcedNonexpansion M V)
    (hNonneg : NonnegV V) {C : ℝ} (hC0 : 0 ≤ C) (hVC : ∀ y, V y ≤ C)
    (hFloor : ∀ y ∈ U, 1 ≤ V y)
    (hΘ : p + thr ≤ 1) (x : State) (hthr : V x ≤ thr) :
    ENNReal.ofReal p ≤ Battery.pathMeasure (skel M) (SkelState.alive x)
      {ω | ∀ n, ω n ∉ skelBad U}
\end{leansrc}
\par\medskip
The staged-row form binds the rule\textquotesingle s premise shapes as
the schema realizes them: per-mode interior rows against the
mode\textquotesingle s flow mechanism, per-jump guard rows, per-mode
floors and nonnegs, and the exact-$\Q$ threshold row cast once ---
assembled onto the flattened battery through the mixture bridge
discussed in the proof section\textquotesingle s final remark.

\declhead{BernsteinLogicAux.QtslBattery}{sound\_QSafetyBattery\_staged}{\badge{BadgeGreen}{statement bound}}
\begin{leansrc}
/-- **`sound_QSafetyBattery_staged`** — the `QSafetyBattery` node's staged
right column (QTSL-rules.toml row `QSafetyBattery`,
amended side) as a bound theorem, discrete-kernel profile. Hypothesis
families mirror the staged rosters:

* `hB1` **interiors** — per mode `m`, the flow chart's one-step drift
  `∫ V_m d(flow_m z) ≤ V_m z` on `D_m \ U_m` (the refined admissible scope;
  the B1 row), with `hint` its domain-membership integrability;
* `hB2` **guards** — per jump `j`, the boundary row
  `V_{dst j}(g_j z) ≤ V_{src j} z` on the guard slab (the B2 row);
* `hB3` **floors** and `hB4` **nonnegs** — per mode on `U ∩ D̄_m` / `D̄_m`
  (the B3/B4 rows);
* `hΘ` — the threshold row `p + thr ≤ 1` in **exact ℚ**, the staged
  `q_add`/`q_le` vocabulary, cast once here;
* `hR3`/`hR4` — the registry side conditions at the compiled kernel;
* `hmix` — **the System bridge**: the presentation's
  charts compile, at this profile, to the step kernel `K` — pointwise, the
  one-step measure from any in-scope state is a sub-convex mixture of the
  mode's flow step and the jumps enabled there (weights supported on
  `src j = m ∧ z ∈ G j`). A System whose compilation satisfies `hmix` is
  exactly what stands left of the concluded `⊨`.

Conclusion: region-level, as the concluded `QRegionSat` with the
`qraSublevel(V, thr)` atom reads — from every start in
`D_m ∩ {V_m ≤ thr}`, the chain satisfies `𝖦 φ` with probability at least
`p`. Assembled onto `sound_QSafetyBattery` through `mixture_drift` /
`mixture_integrable`; the GSHS-profile transport is the standing
residue. -/
theorem sound_QSafetyBattery_staged
    {K : Kernel (Md × X) (Md × X)} [IsMarkovKernel K]
    {Dm Dbarm Um : Md → Set X} {p thr : ℚ} {sat : Md × X → Prop}
    (hDm : ∀ m, MeasurableSet (Dm m)) (hDbarm : ∀ m, MeasurableSet (Dbarm m))
    (hUm : ∀ m, MeasurableSet (Um m)) (hVmeas : ∀ m, Measurable (V m))
    (hF : ∀ m z, ¬ sat (m, z) → z ∈ Um m)
    (hΘ : p + thr ≤ 1)
    (hDsub : ∀ m, Dm m ⊆ Dbarm m)
    (hR3 : ∀ m z, z ∈ Dm m → z ∉ Um m → K (m, z) {y | y.2 ∈ Dbarm y.1} = 1)
    (hR4 : ∀ m z, z ∈ Dm m → z ∉ Um m →
      K (m, z) ({y | y.2 ∈ Dm y.1} ∪ {y | y.2 ∈ Um y.1}) = 1)
    (hint : ∀ m z, z ∈ Dm m → z ∉ Um m → Integrable (V m) ((flow m) z))
    (hB1 : ∀ m z, z ∈ Dm m → z ∉ Um m → ∫ ζ, V m ζ ∂(flow m) z ≤ V m z)
    (hB2 : ∀ j, ∀ z ∈ G j, V (dst j) (g j z) ≤ V (src j) z)
    (hB3 : ∀ m z, z ∈ Um m → z ∈ Dbarm m → 1 ≤ V m z)
    (hB4 : ∀ m z, z ∈ Dbarm m → 0 ≤ V m z)
    (hmix : ∀ m z, z ∈ Dm m → z ∉ Um m →
      ∃ (w₀ : ℝ≥0) (w : J → ℝ≥0), w₀ + ∑ j, w j = 1 ∧
        (∀ j, w j ≠ 0 → src j = m ∧ z ∈ G j) ∧
        K (m, z) = (w₀ : ℝ≥0∞) • ((flow m) z).map (Prod.mk m)
          + ∑ j, (w j : ℝ≥0∞) • Measure.dirac (dst j, g j z)) :
    ∀ m z, z ∈ Dm m → V m z ≤ (thr : ℝ) →
      ENNReal.ofReal (p : ℝ)
        ≤ pathMeasure K (m, z) {ω | ∀ n, sat (ω n)}
\end{leansrc}
\par\medskip
\textbf{The general stratum: acceptance events on the record
carrier.}\enspace Over the semantic universe\textquotesingle s Fell
carrier the safety acceptance event is measurable by the topology of
the site rather than by assumption: confinement events of closed
walls are Fell-closed, hence Borel under the Effros structure ---
over every trajectory window, with no compactness hypothesis
anywhere.

\declhead{SuFunctor.Site.FellAcceptance}{measurableSet\_confinedEvent}{\badge{BadgeGreen}{statement bound}}
\begin{leansrc}
/-- The safety acceptance event of a closed wall is Borel — no
compactness, no σ-compactness, no Hausdorffness. -/
theorem measurableSet_confinedEvent {A : Set α} (hA : IsClosed A) :
    MeasurableSet (confinedEvent A)
\end{leansrc}
\declhead{SuFunctor.Site.FellAcceptance}{measurableSet\_confinedEvent\_cylinder}{\badge{BadgeGreen}{statement bound}}
\begin{leansrc}
/-- **The safety acceptance event is Borel over every window**: the
confinement event of a closed safe region's cylinder, with no
compactness and no σ-compactness — the A8 discharge of the safety
battery's measurability pack. -/
theorem measurableSet_confinedEvent_cylinder {D : Set X}
    (hD : IsClosed D) :
    MeasurableSet
      (confinedEvent ((Set.univ : Set ↥W) ×ˢ D) : Set (FTraj W X))
\end{leansrc}
\par\medskip
\textbf{Walls over the record carrier.}\enspace The
battery\textquotesingle s conclusion shape at the general stratum:
per-window confinement classes with the full judgment bridge
(satisfaction equals almost-sure confinement), the indicator
transport law that admits them, the closed-graph strengthening
(record confinement is pointwise safety, unconditional in the
trajectory), and the composition seam (the window compositor of two
wall demands is the wall demand of the glued class --- persistence
survives composition).

\declhead{SuFunctor.Site.FellBehaviorSpec}{indPenM}{\badge{BadgeGreen}{statement bound}}
\begin{leansrc}
/-- The indicator penalty of a class at the general stratum: permitted
outright iff the joint's mass avoids the class's complement, forbidden
outright otherwise.  Total in the class (mathlib's outer-measure
evaluation); its transport law below is where measurability binds. -/
def indPenM (S : Set K) : MeasPen K
\end{leansrc}
\declhead{SuFunctor.Site.FellBehaviorSpec}{pull\_indPenM}{\badge{BadgeGreen}{statement bound}}
\begin{leansrc}
/-- **Transport of the indicator along a measurable map is the
indicator of the preimage** — the general-stratum form of
`pull_indPen`, with the two measurability hypotheses that stratum
charges: the map's, and the class's.  The latter is the A8 admission
ticket: only Borel acceptance events enter the indicator calculus. -/
theorem pull_indPenM (h : K → L) (hm : Measurable h) {S : Set L}
    (hS : MeasurableSet S) :
    MeasPen.pull h (indPenM S) = indPenM (h ⁻¹' S)
\end{leansrc}
\declhead{SuFunctor.Site.FellBehaviorSpec}{isPersistentMF\_clipDemandsF}{\badge{BadgeGreen}{statement bound}}
\begin{leansrc}
/-- **Persistence of clip demands**: window-calculus transitivity
alone (`penExtMF_penExtMF`), with openness of the site's windows as
the only hypothesis.  No condition on the bases whatsoever. -/
theorem isPersistentMF_clipDemandsF (U : ι → Set T)
    (hU : ∀ i, IsOpen (U i))
    (base : ∀ i, MeasPen (FTraj (U i) X)) :
    IsPersistentMF U (clipDemandsF U base)
\end{leansrc}
\declhead{SuFunctor.Site.FellBehaviorSpec}{satAtM\_specOfBehaviorF\_iff}{\badge{BadgeGreen}{statement bound}}
\begin{leansrc}
/-- **The judgment bridge at the general stratum**: a joint satisfies
the spec of a Fell behavior subtype at a window iff its mass lies in
the class there.  Forward consumes the reflexive demand; backward
consumes the subpresheaf condition, clip measurability (openness of
the site's windows), and class measurability (the A8 ticket) — each
exactly once. -/
theorem satAtM_specOfBehaviorF_iff (U : ι → Set T)
    (hU : ∀ i, IsOpen (U i)) (B : WBehaviorF U X)
    (hB : ∀ i, MeasurableSet (B.cls i)) (j : ι)
    (μ : MeasJoint (FTraj (U j) X)) :
    SatAtM μ (specOfBehaviorF U B j) ↔
      (μ.1 : Measure (FTraj (U j) X)) (B.cls j)ᶜ = 0
\end{leansrc}
\declhead{SuFunctor.Site.FellBehaviorSpec}{satAtM\_safetySpecF\_iff}{\badge{BadgeGreen}{statement bound}}
\begin{leansrc}
/-- The safety wall's judgment bridge, with its measurability paid by
A8: for a closed safe region, satisfaction of the safety spec at a
window is almost-sure confinement there — the SU-general form of the
safety battery's acceptance reading, with the measurability pack
discharged by the topology of the site. -/
theorem satAtM_safetySpecF_iff (U : ι → Set T)
    (hU : ∀ i, IsOpen (U i)) {D : Set X} (hD : IsClosed D) (j : ι)
    (μ : MeasJoint (FTraj (U j) X)) :
    SatAtM μ (specOfBehaviorF U (safetyBehaviorF U D) j) ↔
      (μ.1 : Measure (FTraj (U j) X))
        (confinedEvent ((Set.univ : Set ↥(U j)) ×ˢ D))ᶜ = 0
\end{leansrc}
\declhead{SuFunctor.Site.FellBehaviorSpec}{graphCl\_mem\_confinedEvent\_iff}{\badge{BadgeGreen}{statement bound}}
\begin{leansrc}
/-- **The A5 strengthening**: record confinement of a trajectory's
closed-up graph to a closed wall is pointwise confinement at every
stage — unconditional in the trajectory.  The skeleton route's
segment-constancy side condition is absent by construction: the record
constrains every `(t, x)` of the closed graph, mid-segment stages
included, and the closure completion costs nothing against a closed
wall. -/
theorem graphCl_mem_confinedEvent_iff {D : Set X} (hD : IsClosed D)
    (f : Traj W X) :
    graphCl f ∈ confinedEvent ((Set.univ : Set ↥W) ×ˢ D) ↔
      ∀ w, f w ∈ D
\end{leansrc}
\declhead{SuFunctor.Site.FellBehaviorSpec}{mboxplusWF\_indPenM}{\badge{BadgeGreen}{statement bound}}
\begin{leansrc}
/-- **The record-carrier composition seam on indicators** (the
general-stratum sibling of `boxplusW_indPen`): the window compositor
of two wall demands is the wall demand of the glued record class —
the pullback intersection along the two state transports.  Properness
feeds the transports' measurability, and the classes' measurability
is the A8 ticket again.  This is the A4 transport core: composing
wall specs along the compositor never leaves the wall fragment, so
the composed conclusion is again a per-window class demand, and
persistence survives by `IsPersistentMF.image2_mboxplusWF`.  The
per-profile grid cells of the wire-composition rule (ODE, sampled
data, gauge-form feedback) are drift-assembly content at the
presentation profiles, deliberately not restated here. -/
theorem mboxplusWF_indPenM {W : Set T} {π₁ : Z → U'} {π₂ : Z → V'}
    (hp₁ : IsProperMap π₁) (hp₂ : IsProperMap π₂)
    {S₁ : Set (FTraj W U')} {S₂ : Set (FTraj W V')}
    (hS₁ : MeasurableSet S₁) (hS₂ : MeasurableSet S₂) :
    mboxplusWF π₁ π₂ (indPenM S₁) (indPenM S₂)
      = indPenM (fellMap (W
\end{leansrc}
\ruleblock{Correspondence (Steele display $\leftrightarrow$ Lean)}
\badge{BadgeAmber}{a bound Lean statement CHANGED since certification --- re-review and re-certify}
\sourcenote{The dictionary below is a reviewer-certified statement of how the mechanized statements above interpret the displayed rule: display metavariable $\mapsto$ Lean binder, each display hypothesis $\mapsto$ the premise or side condition carrying it, the display conclusion $\mapsto$ the goal --- with any deliberate difference in width between display and mechanization named outright. Certification is recorded as a hash triple over (the displayed cell, the extracted statement texts, this dictionary\textquotesingle s own text); if any of the three changes, this block turns amber until a reviewer re-certifies. The correspondence claim itself remains editorial --- reviewer-grade, not machine-checked; its machine-checked successor types each soundness statement against the rule as a structured object, and supersedes this block per rule as it lands.}
Model frame as at \rulename{QBasin}, probabilistically read: the
displayed $\mathcal{M},\ x_0 \in \{V \le y\} \vDash
\mathsf{SProbGE}\bigl(p,\ \mathsf{PAlways}\ \varphi\bigr)$ $\mapsto$
the goal $\mathrm{ofReal}\ p \le$
\texttt{pathMeasure K x \{ω | ∀ n, sat (ω n)\}} for every start in the
concluded region --- \texttt{sat} the $\varphi$-satisfaction
predicate, of which \emph{no measurability is asked} (the fail-closed
event is the measurable one; the tie weakens by
\texttt{measure\_mono}). The staged-row theorem
\texttt{sound\_QSafetyBattery\_staged} is the amended
display\textquotesingle s binding; the flattened
\texttt{sound\_QSafetyBattery} is its engine.
\begin{itemize}
\item profile, stated plainly: the map splits in two. The bullets
  below marked \emph{discrete map} bind the rule at the
  \emph{discrete-kernel profile} --- $\mathcal{M}$\textquotesingle s
  step semantics is a Markov kernel chain, its trajectory measure
  rendered as \texttt{pathMeasure} in the bound statements --- via
  \texttt{sound\_QSafetyBattery\_staged}, with the PDMP sub-profile
  carried end-to-end by \texttt{sound\_GSHS\_safety\_battery}. The
  martingale-problem theorems (\texttt{sound\_GSHS\_safety\_mp},
  \texttt{sound\_GSHS\_safety\_battery\_mp} --- hash-bound below) are
  the \emph{continuous-time avoidance engine}, mapped separately in
  the \emph{MP map} bullet; they do not themselves bind the full
  displayed judgment.
\item \emph{MP map}, its own binders and its own goal:
  the generator presentation and the interior rows $\mapsto$
  \texttt{hmp} $:$ $\forall x$, \texttt{Nonempty (MPFlow M 0 V x)}
  (the martingale-problem presentation hypothesis, carried as data);
  guards $\mapsto$ \texttt{hForce} $:$ \texttt{ForcedNonexpansion M
  V}; nonnegs $\mapsto$ \texttt{hNonneg} $:$ \texttt{NonnegV V};
  floors $\mapsto$ \texttt{hFloor} $:$ $\forall y \in U$, $1 \le V\,
  y$; the threshold row $\mapsto$ \texttt{hΘ}, and the start
  $\mapsto$ \texttt{x}/\texttt{hthr}. Undisplayed MP regularity, named
  outright: \texttt{hV}/\texttt{hU} (measurability),
  \texttt{hC0}/\texttt{hVC} (a global certificate bound $C$), and the
  \texttt{IsMarkovKernel} instances on \texttt{M.seg}/\texttt{M.jump}.
  Honesty about the goal: the MP battery\textquotesingle s literal
  conclusion is skeleton avoidance ---
  $\mathrm{ofReal}\ p \le$ \texttt{Battery.pathMeasure (skel M)
  (SkelState.alive x) \{ω | ∀ n, ω n ∉ skelBad U\}} --- avoidance of
  the $U$-bad set along the embedded skeleton, not the displayed
  $\mathsf{SProbGE}(p, \mathsf{PAlways}\ \varphi)$ judgment verbatim; the
  formula tie from $\varphi$ to $U$ is the discrete
  map\textquotesingle s \texttt{hF} business, and only the discrete
  staged theorem binds the displayed judgment itself.
\item \emph{discrete map} --- what stands left of the $\vDash$
  $\mapsto$ \texttt{hmix}, the
  System bridge: the presentation\textquotesingle s charts compile to
  the kernel as a sub-convex mixture of the mode\textquotesingle s
  flow step and the enabled jumps --- an explicit hypothesis, so the
  certified statement quantifies over exactly the Systems whose
  compilation satisfies it.
\item \emph{discrete map} --- the four premise families map
  one-to-one: interiors $\mapsto$
  \texttt{hB1} (per-mode one-step flow drift on $D_m \smallsetminus
  U_m$) with the undisplayed integrability rider \texttt{hint};
  guards $\mapsto$ \texttt{hB2} (per-jump boundary rows on the guard
  slab); floors $\mapsto$ \texttt{hB3} (on $U \cap \bar{D}_m$, the
  ruled fail-closed scope); nonnegs $\mapsto$ \texttt{hB4}.
\item \emph{discrete map} --- the threshold tie $p +
  y \le 1$ (a side condition in the display) $\mapsto$ \texttt{hΘ},
  in exact $\Q$, cast once inside the statement.
\item \emph{discrete map} --- side conditions R3/R4 $\mapsto$
  \texttt{hR3}/\texttt{hR4} at
  the compiled kernel (hull containment; fail-closed reset
  containment); the formula tie $U := \mathrm{cmpl}(\varphi)$
  $\mapsto$ \texttt{hF} $:$ $\neg\,\texttt{sat}\ y \to y \in U$ ---
  note the certified hypothesis asks only that $U$
  \emph{over-approximate} the falsity set: strictly wider than the
  displayed equation, in the safe direction.
\item \emph{discrete map} --- conclusion region
  $\mathsf{qraSublevel}(V, y)$
  $\mapsto$ the start conditions \texttt{z ∈ Dm m} and \texttt{V m z
  ≤ thr} (the Lean source\textquotesingle s own identifier for the
  display\textquotesingle s $y$), region-level as the judgment reads.
\end{itemize}
Width note: measurability of the region/certificate data
(\texttt{hDm}/\texttt{hDbarm}/\texttt{hUm}/\texttt{hVmeas}) and the
interior integrability rider are profile-level regularity with no
displayed rows --- the probabilistic sibling of the
\rulename{PsdCert} evaluability pack; every mandatory displayed row is
consumed --- the display\textquotesingle s \emph{optional}
$\mathrm{scope\_warrant}$ slot (the dissolved ScopeRule\textquotesingle
s interface) is taken by neither certified statement and has no entry
in this map --- and the two widenings (over-approximating $U$,
sub-convex \texttt{hmix}) both enlarge the certified reach beyond the
display\textquotesingle s letter.

\ruleblock{Lean verifier receipts}
\sourcenote{module \texttt{BernsteinLogicAux.Battery}; bound build \texttt{2fd1c031c80b\ldots} (bound from receipt 8d9ccb68993ea0ec, 2026-08-04T01:56:55Z). A receipt records that named checking layers accepted a build: (module, compiled-artifact SHA-256, environment SHA-256, checker layers), minted when the checking gate ran; the store receipts attesting the bound build are listed with the layers each run credited. Nothing here asks the reader to trust the store: every layer is an independently runnable program over the quoted sources --- rebuild the module and re-run the checkers to reproduce the verdict. Receipts attesting superseded builds credit no layer and are omitted from this rendering (the count of omissions, if any, is noted below).}
\begin{receipttable}
\receiptrow{8d9ccb68993ea0ec}{2026-08-04T01:56:55Z}{build+decorrelated}{\badge{BadgeGreen}{current}}
\receiptrow{5075586e4813a7c5}{2026-08-04T03:11:30Z}{build+decorrelated}{\badge{BadgeGreen}{current}}
\end{receipttable}
\sourcenote{13 further receipt(s) in the store attest superseded builds of this module and are omitted from this rendering.}
\sourcenote{module \texttt{BernsteinLogicAux.LayerD.Skeleton}; bound build \texttt{9a4f5ca8bc26\ldots} (bound from receipt f3d40eda96c4513b, 2026-08-04T01:57:01Z). A receipt records that named checking layers accepted a build: (module, compiled-artifact SHA-256, environment SHA-256, checker layers), minted when the checking gate ran; the store receipts attesting the bound build are listed with the layers each run credited. Nothing here asks the reader to trust the store: every layer is an independently runnable program over the quoted sources --- rebuild the module and re-run the checkers to reproduce the verdict. Receipts attesting superseded builds credit no layer and are omitted from this rendering (the count of omissions, if any, is noted below).}
\begin{receipttable}
\receiptrow{f3d40eda96c4513b}{2026-08-04T01:57:01Z}{build+decorrelated}{\badge{BadgeGreen}{current}}
\receiptrow{5966aa764026dde9}{2026-08-04T03:11:31Z}{build+decorrelated}{\badge{BadgeGreen}{current}}
\end{receipttable}
\sourcenote{13 further receipt(s) in the store attest superseded builds of this module and are omitted from this rendering.}
\sourcenote{module \texttt{BernsteinLogicAux.LayerD.DiniDischarge}; bound build \texttt{db41370b5e77\ldots} (bound from receipt 2a294d6f4ffc9a7b, 2026-07-21T04:51:51Z). A receipt records that named checking layers accepted a build: (module, compiled-artifact SHA-256, environment SHA-256, checker layers), minted when the checking gate ran; the store receipts attesting the bound build are listed with the layers each run credited. Nothing here asks the reader to trust the store: every layer is an independently runnable program over the quoted sources --- rebuild the module and re-run the checkers to reproduce the verdict. Receipts attesting superseded builds credit no layer and are omitted from this rendering (the count of omissions, if any, is noted below).}
\begin{receipttable}
\receiptrow{073080d556f0cf08}{2026-07-20T22:02:26Z}{build+decorrelated}{\badge{BadgeGreen}{current}}
\receiptrow{6f4929afcf219d67}{2026-07-20T22:52:48Z}{build+decorrelated}{\badge{BadgeGreen}{current}}
\receiptrow{f258f966d21ec502}{2026-07-21T00:03:10Z}{build+decorrelated}{\badge{BadgeGreen}{current}}
\receiptrow{2a294d6f4ffc9a7b}{2026-07-21T04:51:51Z}{build+decorrelated}{\badge{BadgeGreen}{current}}
\receiptrow{277ec2c0fa8227a9}{2026-07-21T05:38:41Z}{build+replay+hardened+decorrelated}{\badge{BadgeGreen}{current}}
\receiptrow{a142e9bdad29ed20}{2026-07-22T18:27:03Z}{build+decorrelated}{\badge{BadgeGreen}{current}}
\receiptrow{6aeb017a3939a647}{2026-07-31T21:59:58Z}{build+decorrelated}{\badge{BadgeGreen}{current}}
\receiptrow{85b98a8efabbfc19}{2026-07-31T22:28:52Z}{build+decorrelated}{\badge{BadgeGreen}{current}}
\receiptrow{fcd8973bd3d8aa0b}{2026-08-02T23:31:02Z}{build+decorrelated}{\badge{BadgeGreen}{current}}
\receiptrow{45fcba5d9746e88b}{2026-08-03T16:27:06Z}{build+decorrelated}{\badge{BadgeGreen}{current}}
\receiptrow{5a035508867866dd}{2026-08-03T22:45:49Z}{build+decorrelated}{\badge{BadgeGreen}{current}}
\end{receipttable}
\sourcenote{module \texttt{BernsteinLogicAux.LayerD.ItoInterface}; bound build \texttt{9b18e4e710f1\ldots} (bound from receipt 8b9184a6df93a83e, 2026-07-21T04:51:52Z). A receipt records that named checking layers accepted a build: (module, compiled-artifact SHA-256, environment SHA-256, checker layers), minted when the checking gate ran; the store receipts attesting the bound build are listed with the layers each run credited. Nothing here asks the reader to trust the store: every layer is an independently runnable program over the quoted sources --- rebuild the module and re-run the checkers to reproduce the verdict. Receipts attesting superseded builds credit no layer and are omitted from this rendering (the count of omissions, if any, is noted below).}
\begin{receipttable}
\receiptrow{4e3a32599ecb1605}{2026-07-21T00:03:12Z}{build+decorrelated}{\badge{BadgeGreen}{current}}
\receiptrow{8b9184a6df93a83e}{2026-07-21T04:51:52Z}{build+decorrelated}{\badge{BadgeGreen}{current}}
\receiptrow{04aa3a28f81ea9f9}{2026-07-21T05:38:43Z}{build+replay+hardened+decorrelated}{\badge{BadgeGreen}{current}}
\receiptrow{d95924f2dd7f318f}{2026-07-22T18:27:04Z}{build+decorrelated}{\badge{BadgeGreen}{current}}
\receiptrow{c78dba97e2eee2ba}{2026-07-31T21:59:59Z}{build+decorrelated}{\badge{BadgeGreen}{current}}
\receiptrow{ca54011b021b532e}{2026-07-31T22:28:52Z}{build+decorrelated}{\badge{BadgeGreen}{current}}
\receiptrow{bdf792a36abf7da8}{2026-08-02T23:31:02Z}{build+decorrelated}{\badge{BadgeGreen}{current}}
\receiptrow{e8478cde1b56df65}{2026-08-03T16:27:06Z}{build+decorrelated}{\badge{BadgeGreen}{current}}
\receiptrow{e987ba886a91fc22}{2026-08-03T22:45:50Z}{build+decorrelated}{\badge{BadgeGreen}{current}}
\end{receipttable}
\sourcenote{module \texttt{BernsteinLogicAux.LayerD.ItoSampling}; bound build \texttt{fb8de8b804d7\ldots} (bound from receipt 467820f314b0750c, 2026-07-21T08:16:24Z). A receipt records that named checking layers accepted a build: (module, compiled-artifact SHA-256, environment SHA-256, checker layers), minted when the checking gate ran; the store receipts attesting the bound build are listed with the layers each run credited. Nothing here asks the reader to trust the store: every layer is an independently runnable program over the quoted sources --- rebuild the module and re-run the checkers to reproduce the verdict. Receipts attesting superseded builds credit no layer and are omitted from this rendering (the count of omissions, if any, is noted below).}
\begin{receipttable}
\receiptrow{467820f314b0750c}{2026-07-21T08:16:24Z}{build+decorrelated}{\badge{BadgeGreen}{current}}
\end{receipttable}
\sourcenote{9 further receipt(s) in the store attest superseded builds of this module and are omitted from this rendering.}
\sourcenote{module \texttt{BernsteinLogicAux.GshsItoDischarge}; bound build \texttt{858d3234fc4f\ldots} (bound from receipt d6088699caa41da2, 2026-08-03T20:13:46Z). A receipt records that named checking layers accepted a build: (module, compiled-artifact SHA-256, environment SHA-256, checker layers), minted when the checking gate ran; the store receipts attesting the bound build are listed with the layers each run credited. Nothing here asks the reader to trust the store: every layer is an independently runnable program over the quoted sources --- rebuild the module and re-run the checkers to reproduce the verdict. Receipts attesting superseded builds credit no layer and are omitted from this rendering (the count of omissions, if any, is noted below).}
\begin{receipttable}
\receiptrow{1d9003db41305809}{2026-08-02T12:53:42Z}{build+decorrelated}{\badge{BadgeGreen}{current}}
\receiptrow{615e37d28395ee85}{2026-08-02T23:30:51Z}{build+decorrelated}{\badge{BadgeGreen}{current}}
\receiptrow{306da0b697e1ba52}{2026-08-03T16:26:54Z}{build+decorrelated}{\badge{BadgeGreen}{current}}
\receiptrow{d6088699caa41da2}{2026-08-03T20:13:46Z}{build+decorrelated}{\badge{BadgeGreen}{current}}
\receiptrow{b13efcc7c507c534}{2026-08-03T22:45:38Z}{build+decorrelated}{\badge{BadgeGreen}{current}}
\end{receipttable}
\sourcenote{7 further receipt(s) in the store attest superseded builds of this module and are omitted from this rendering.}
\sourcenote{module \texttt{BernsteinLogicAux.Debut}; bound build \texttt{cc670db616db\ldots} (bound from receipt 38571a839f9da48c, 2026-07-24T09:28:37Z). A receipt records that named checking layers accepted a build: (module, compiled-artifact SHA-256, environment SHA-256, checker layers), minted when the checking gate ran; the store receipts attesting the bound build are listed with the layers each run credited. Nothing here asks the reader to trust the store: every layer is an independently runnable program over the quoted sources --- rebuild the module and re-run the checkers to reproduce the verdict. Receipts attesting superseded builds credit no layer and are omitted from this rendering (the count of omissions, if any, is noted below).}
\begin{receipttable}
\receiptrow{a7c2ce15b3ea5ce7}{2026-07-22T16:23:09Z}{build+decorrelated}{\badge{BadgeGreen}{current}}
\receiptrow{e8af72cc3335497f}{2026-07-22T18:26:50Z}{build+decorrelated}{\badge{BadgeGreen}{current}}
\receiptrow{38571a839f9da48c}{2026-07-24T09:28:37Z}{build+decorrelated}{\badge{BadgeGreen}{current}}
\receiptrow{d0e0aa26ac8f8d34}{2026-07-31T21:59:41Z}{build+decorrelated}{\badge{BadgeGreen}{current}}
\receiptrow{b5cc2fcbabda9723}{2026-07-31T22:28:35Z}{build+decorrelated}{\badge{BadgeGreen}{current}}
\receiptrow{f494a0acd9a534b4}{2026-08-02T23:30:43Z}{build+decorrelated}{\badge{BadgeGreen}{current}}
\receiptrow{72ed3556543f0ada}{2026-08-03T16:26:47Z}{build+decorrelated}{\badge{BadgeGreen}{current}}
\receiptrow{aea75952b3723556}{2026-08-03T22:45:30Z}{build+decorrelated}{\badge{BadgeGreen}{current}}
\end{receipttable}
\sourcenote{module \texttt{BernsteinLogicAux.OptionalStoppingRC}; bound build \texttt{3f566c19afe8\ldots} (bound from receipt 704016ef1d4bd23d, 2026-07-24T09:28:43Z). A receipt records that named checking layers accepted a build: (module, compiled-artifact SHA-256, environment SHA-256, checker layers), minted when the checking gate ran; the store receipts attesting the bound build are listed with the layers each run credited. Nothing here asks the reader to trust the store: every layer is an independently runnable program over the quoted sources --- rebuild the module and re-run the checkers to reproduce the verdict. Receipts attesting superseded builds credit no layer and are omitted from this rendering (the count of omissions, if any, is noted below).}
\begin{receipttable}
\receiptrow{704016ef1d4bd23d}{2026-07-24T09:28:43Z}{build+decorrelated}{\badge{BadgeGreen}{current}}
\end{receipttable}
\sourcenote{6 further receipt(s) in the store attest superseded builds of this module and are omitted from this rendering.}
\sourcenote{module \texttt{QtslAux.LayerD.MPDischarge}; bound build \texttt{baded74dd943\ldots} (bound from receipt 7228834dea10e5f9, 2026-07-24T11:04:04Z). A receipt records that named checking layers accepted a build: (module, compiled-artifact SHA-256, environment SHA-256, checker layers), minted when the checking gate ran; the store receipts attesting the bound build are listed with the layers each run credited. Nothing here asks the reader to trust the store: every layer is an independently runnable program over the quoted sources --- rebuild the module and re-run the checkers to reproduce the verdict. Receipts attesting superseded builds credit no layer and are omitted from this rendering (the count of omissions, if any, is noted below).}
\begin{receipttable}
\receiptrow{7228834dea10e5f9}{2026-07-24T11:04:04Z}{build+decorrelated}{\badge{BadgeGreen}{current}}
\end{receipttable}
\sourcenote{module \texttt{BernsteinLogicAux.QtslBattery}; bound build \texttt{5c75f9923f27\ldots} (bound from receipt 3fd830b4914bdeda, 2026-08-04T01:56:58Z). A receipt records that named checking layers accepted a build: (module, compiled-artifact SHA-256, environment SHA-256, checker layers), minted when the checking gate ran; the store receipts attesting the bound build are listed with the layers each run credited. Nothing here asks the reader to trust the store: every layer is an independently runnable program over the quoted sources --- rebuild the module and re-run the checkers to reproduce the verdict. Receipts attesting superseded builds credit no layer and are omitted from this rendering (the count of omissions, if any, is noted below).}
\begin{receipttable}
\receiptrow{3fd830b4914bdeda}{2026-08-04T01:56:58Z}{build+decorrelated}{\badge{BadgeGreen}{current}}
\receiptrow{96b18d8572451f65}{2026-08-04T03:11:31Z}{build+decorrelated}{\badge{BadgeGreen}{current}}
\end{receipttable}
\sourcenote{13 further receipt(s) in the store attest superseded builds of this module and are omitted from this rendering.}
\sourcenote{module \texttt{SuFunctor.Site.FellAcceptance}; bound build \texttt{651fb557bbfe\ldots} (bound from receipt aa2169560f8f571a, 2026-08-03T20:19:04Z). A receipt records that named checking layers accepted a build: (module, compiled-artifact SHA-256, environment SHA-256, checker layers), minted when the checking gate ran; the store receipts attesting the bound build are listed with the layers each run credited. Nothing here asks the reader to trust the store: every layer is an independently runnable program over the quoted sources --- rebuild the module and re-run the checkers to reproduce the verdict. Receipts attesting superseded builds credit no layer and are omitted from this rendering (the count of omissions, if any, is noted below).}
\begin{receipttable}
\receiptrow{7ad85720de6e419e}{2026-08-03T10:13:50Z}{build+decorrelated}{\badge{BadgeGreen}{current}}
\receiptrow{94a408109996d1bf}{2026-08-03T13:41:21Z}{build+decorrelated}{\badge{BadgeGreen}{current}}
\receiptrow{580a7e551fab03b6}{2026-08-03T15:39:44Z}{build+decorrelated}{\badge{BadgeGreen}{current}}
\receiptrow{aa2169560f8f571a}{2026-08-03T20:19:04Z}{build+decorrelated}{\badge{BadgeGreen}{current}}
\end{receipttable}
\sourcenote{2 further receipt(s) in the store attest superseded builds of this module and are omitted from this rendering.}
\sourcenote{module \texttt{SuFunctor.Site.FellBehaviorSpec}; bound build \texttt{345f9caf9453\ldots} (bound from receipt 72f2c3b62b6771fc, 2026-08-03T20:19:08Z). A receipt records that named checking layers accepted a build: (module, compiled-artifact SHA-256, environment SHA-256, checker layers), minted when the checking gate ran; the store receipts attesting the bound build are listed with the layers each run credited. Nothing here asks the reader to trust the store: every layer is an independently runnable program over the quoted sources --- rebuild the module and re-run the checkers to reproduce the verdict. Receipts attesting superseded builds credit no layer and are omitted from this rendering (the count of omissions, if any, is noted below).}
\begin{receipttable}
\receiptrow{d7b846e785c1d7fb}{2026-08-03T10:13:53Z}{build+decorrelated}{\badge{BadgeGreen}{current}}
\receiptrow{c0e3dc15419765aa}{2026-08-03T13:41:24Z}{build+decorrelated}{\badge{BadgeGreen}{current}}
\receiptrow{9a62b29c08b63b4f}{2026-08-03T15:39:47Z}{build+decorrelated}{\badge{BadgeGreen}{current}}
\receiptrow{72f2c3b62b6771fc}{2026-08-03T20:19:08Z}{build+decorrelated}{\badge{BadgeGreen}{current}}
\end{receipttable}
\sourcenote{2 further receipt(s) in the store attest superseded builds of this module and are omitted from this rendering.}
\layerline{\layeritem{build}{lake build (elaboration)}{\badge{BadgeGreen}{current}}\layeritem{replay}{leanchecker (official kernel replay)}{\badge{BadgeGreen}{current}}\layeritem{hardened}{safe\_verify (hardened screen)}{\badge{BadgeGreen}{current}}\layeritem{decorrelated}{lean4export + nanoda (decorrelated kernel)}{\badge{BadgeGreen}{current}}}
\sourcenote{\badge{BadgeGray}{comparator: none} no fixture in the three-implementation comparator harness (compiled / kernel-interpreted / reference-interpreted) exercises this rule.}
\ruleblock{HOL4 statement}
\absentbox{NO HOL STATEMENT YET}{The HOL4 development ports the base calculus\textquotesingle s checker
chain (theory \texttt{bern\_catalogue}, developed side by side with the
base catalogue sources, stopping below the proof-object layer); it
contains no probability or measure theory.
\rulename{QSafetyBattery} lives one theory-extension edge above the base
calculus --- it concludes a satisfaction judgment, not a sequent --- and
its soundness content is path-space measure theory (trajectory kernels,
stopped supermartingales, Ville\textquotesingle s inequality), for which
no counterpart exists anywhere in the ported development. A HOL
statement of this rule would be a new measure-theoretic campaign, not an
extension of the ported checker chain; until one exists, this box
records the current state of the edifice.
}
\ruleblock{HOL verifier receipts}
\absentbox{NONE}{no HOL verification receipt exists for this rule.}
\ruleblock{Proof}

\sourcenote{Genre: core rule (tier 1) --- an ordinary mathematical proof
at Putnam-exam-solution detail. This rule concludes a \emph{satisfaction
judgment} ($\vDash$), one theory-extension edge above the base calculus:
every premise row is an inherited base-calculus sequent, and soundness
is proved \emph{given} base soundness (the extension-edge convention).
The theorems quoted above warrant the rule at the
\textbf{discrete-kernel profile}; the profile remark at the end of this
section states exactly what remains a hypothesis beyond it. Lemmas
\ref{lem:qsb-stopped}--\ref{lem:qsb-ville} and
Proposition~\ref{prop:qsb-pattern} are the battery family\textquotesingle s
shared probabilistic preamble.}

\textbf{Setting and claim.}\enspace Fix a measurable state space $S$ ---
the flattened hybrid state space, i.e.\ the disjoint sum of the per-mode
domains --- and a Markov one-step kernel $K$ on $S$, carrying both
discretized flow steps and guarded jumps. For $x \in S$ write $\Pr_x$
for the path measure on $S^{\N}$: the trajectory law of the
time-homogeneous chain driven by $K$, started deterministically at $x$
(the Ionescu--Tulcea construction, with the constant kernel family).
$(\mathcal{F}_n)_n$ is the canonical filtration, $\mathcal{F}_n$
generated by the coordinates $\le n$.

The rule\textquotesingle s data: a modelled domain $D$, a hull
$\bar{D} \supseteq D$ (the region on which the polynomial rows were
actually checked), an avoid set $U$ compiled to over-approximate the
falsity set of the $\mathsf{G}$-argument --- $D$, $\bar{D}$, $U$ all
measurable subsets of $S$, as is every region a row speaks on, here
and in the sibling rules riding this substrate --- a measurable
certificate $V : S \to \R$, and rationals $p, y$. The checked rows and
side conditions, read semantically at this profile:
\begin{itemize}
\item[(B1$\wedge$B2)] \emph{generator}: for every $z \in D \smallsetminus U$,
  $V$ is integrable against $K\,z$ and $\int V \,dK\,z \le V(z)$. (At
  this profile the interior and guard rows land \emph{together}: the
  one-step kernel carries flow and enabled jumps in a single step, so
  the two row families compile to one kernel-level generator row.)
\item[(B3)] \emph{floors}: $V \ge 1$ on $U \cap \bar{D}$.
\item[(B4)] \emph{nonnegs}: $V \ge 0$ on $\bar{D}$.
\item[(R3)] \emph{hull}: for every $z \in D \smallsetminus U$,
  $K\,z\,(\bar{D}) = 1$ --- one-step mass stays where the rows were
  checked.
\item[(R4)] \emph{fail-closed containment}: for every
  $z \in D \smallsetminus U$, $K\,z\,(D \cup U) = 1$.
\item[($\Theta$)] \emph{threshold}: $p + y \le 1$, checked in
  exact rational arithmetic.
\item[(F)] \emph{formula tie}: $\lnot \mathrm{sat}(z) \Rightarrow z \in U$,
  where $\mathrm{sat}$ is the state-satisfaction predicate of the
  $\mathsf{G}$-argument $\varphi$.
\end{itemize}
The claim (\texttt{sound\_QSafetyBattery}): for every $x \in D$ with
$V(x) \le y$,
\[
\Pr_x^{\,*}\bigl[\forall n,\ \mathrm{sat}(\omega_n)\bigr] \;\ge\; p
\]
--- outer measure, because no measurability is ever asked of the
satisfaction predicate (Lemma~\ref{lem:qsb-ville} states the
convention); it coincides with the ordinary probability whenever the
satisfaction event is measurable.

\begin{blemma}[the stopped certificate]\label{lem:qsb-stopped}
Let $f : S \to \R$ and let $A \subseteq S$ be a stop set. Define the
stopped certificate process $X_n(\omega) := f(\omega_{\,n \wedge \tau_A})$,
where $\tau_A$ is the first visit to $A$ --- equivalently, by the
primitive recursion that freezes the value once $A$ has been visited.
Then: (i) if $A$ is measurable and $f$ is measurable, $X$ is adapted to
$(\mathcal{F}_n)$; (ii) any pointwise property of $f$\textquotesingle s
values holds of $X$\textquotesingle s values; (iii) \emph{barrier
transport}: if $f \ge 1$ on $A$, then $X_n \ge 1$ on the event
$\{\tau_A \le n\}$; (iv) $X_n$ depends only on the coordinates $\le n$.
\end{blemma}
\begin{proof}
All four are elementary inductions on $n$ over the two branches of the
recursion (surviving: $X_{n+1} = f(\omega_{n+1})$; frozen:
$X_{n+1} = X_n$). No stopping-time API is needed anywhere --- the
recursion \emph{is} the stopped process, and (iii) is proved by chasing
the first index at which $A$ was visited.
\end{proof}

\begin{blemma}[supermartingale bridge]\label{lem:qsb-bridge}
Let $f$ be measurable with $0 \le f \le C$, let $A$ be measurable, and
suppose the one-step generator row holds \emph{off} $A$: for all
$z \notin A$, $\int f\,dK\,z \le f(z)$. Then for every start $x$, the
stopped certificate $X$ of Lemma~\ref{lem:qsb-stopped} is a
supermartingale with respect to $(\mathcal{F}_n, \Pr_x)$.
\end{blemma}
\begin{proof}
Integrability is the bound $0 \le X_n \le C$ (Lemma
\ref{lem:qsb-stopped}(ii)). For the supermartingale inequality,
condition on $\mathcal{F}_n$: by the Markov property of the trajectory
construction, the conditional expectation of a bounded
$\mathcal{F}_{n+1}$-measurable function given $\mathcal{F}_n$ is its
integral against the trajectory kernel continued from the $n$-prefix.
Split on the survival event
$S_n := \{\omega : \omega_k \notin A \ \text{for all } k \le n\}$ ---
$\mathcal{F}_n$-measurable, $A$ being measurable, and determined by the
coordinates $\le n$, so paths patched to agree with $\omega$ up to $n$
have the same survival status. On the surviving branch the integrand is
$f(\text{next coordinate})$, whose one-step marginal is $K(\omega_n)$;
since survival means $\omega_n \notin A$, the generator row applies and the
conditional expectation is $\le f(\omega_n) = X_n(\omega)$. On the
frozen branch the integrand is the constant $X_n(\omega)$. Note that
\emph{every} hypothesis here is pointwise --- there is no
almost-everywhere invariance argument anywhere in the bridge; this is
what the fail-closed design buys, and the reachability sibling
(\rulename{QDeadlineBattery}) will have to pay precisely this cost when
it can no longer enlarge its stop set.
\end{proof}

\begin{blemma}[Ville\textquotesingle s inequality, barrier form]
\label{lem:qsb-ville}
Let $X$ be a nonnegative supermartingale with $X_n \ge 1$ on the event
$\{$some coordinate $\le n$ lies in $A\}$. Then
$\Pr_x^{\,*}[\exists n,\ \omega_n \in A] \le \mathbb{E}\,X_0$, where
$\Pr_x^{\,*}$ is the outer measure --- equal to the ordinary
probability whenever the visit event is measurable, as it is at every
consumption in this family ($A$ is measurable there).
\end{blemma}
\begin{proof}
This is the maximal inequality for nonnegative supermartingales at
barrier $1$ \cite{ville1939} (textbook treatment: \cite{doob1953}): the
probability that a nonnegative supermartingale ever reaches level $a$ is
at most $\mathbb{E}\,X_0 / a$. The mechanized form deliberately asks no
measurability of the visit event itself --- the bad event is bounded by
inclusion in the \emph{measurable} event $\{\exists n,\ X_n \ge 1\}$,
whose probability the maximal inequality bounds, and only the process
needs to be adapted; the outer-measure phrasing of the conclusion is
exactly this inclusion, spelled.
\end{proof}

\begin{bprop}[\textbf{the battery pattern}]\label{prop:qsb-pattern}
Let $A$ be measurable and let $f$ be measurable with
$0 \le f \le C$ globally, $f \ge 1$ on $A$, and one-step drift off $A$.
Then for every $x$,
\[
\Pr_x\bigl[\exists n,\ \omega_n \in A\bigr] \;\le\; f(x).
\]
\end{bprop}
\begin{proof}
Lemmas \ref{lem:qsb-stopped}--\ref{lem:qsb-ville} assembled; the initial
value is $\mathbb{E}\,X_0 = f(x)$ because coordinate $0$ is
deterministically $x$. This proposition --- \emph{nonnegative
supermartingale $+$ barrier $+$ initial value $\Rightarrow$ probability
bound} --- is the shape every rule in the battery family consumes; the
siblings cite it as ``the battery pattern
(\rulename{QSafetyBattery}, Proposition~\ref{prop:qsb-pattern})''.
\end{proof}

\begin{blemma}[the clamp: from checked rows to global hypotheses]
\label{lem:qsb-clamp}
The rows (B3), (B4) speak only on $\bar{D}$, while
Proposition~\ref{prop:qsb-pattern} wants its bounds globally. Set
$A := U \cup D^{\mathrm{c}}$ (the fail-closed stop set),
$C := \max(y, 1)$, and let $g$ be the clamped certificate:
on $D$, clamp $V$ into $[0, C]$; off $D$, clamp $\max(V, 1)$ into
$[0,C]$ (so the floor is \emph{forced} off the modelled domain). Then:
(i) $0 \le g \le C$ globally; (ii) $g \ge 1$ on $A$; (iii) on
$\bar{D} \cap (D \cup U)$ the clamp is invisible:
$g = \min(V, C)$ there; (iv) under (R3) $+$ (R4), $g$ satisfies the
one-step generator row off $A$; (v) $g(x) = V(x)$ at any
$x \in D$ with $V(x) \le y$. Throughout, $g$ is measurable:
a case split between measurable pieces ($V$, $D$ measurable) composed
with clamps --- so $(g, A)$ is an admissible input to
Proposition~\ref{prop:qsb-pattern}.
\end{blemma}
\begin{proof}
(i), (ii), (v) are case checks from (B3), (B4) and the definition ---
(ii) uses (B3) on $U \cap \bar D$ and the forced floor off $D$. (iii):
on $D$ the lower clamp is inert by (B4); on
$(\bar{D} \cap U) \smallsetminus D$ the forced floor is inert by (B3).
For (iv), fix $z \in D \smallsetminus U$ (i.e.\ $z \notin A$). By (R3)
and (R4) the one-step measure $K\,z$ is carried by
$\bar{D} \cap (D \cup U)$ --- the complement has measure
$\le K z (\bar D^{\mathrm c}) + K z ((D \cup U)^{\mathrm c}) = 0$ ---
so by (iii) the integrand may be replaced by $\min(V, C)$. Then
$\int \min(V, C)\,dK\,z \le \min\bigl(\int V\,dK\,z,\; C\bigr)
\le \min(V(z), C) = g(z)$, using monotonicity of the integral, the
generator row (B1$\wedge$B2), and (iii) at $z$ itself.
\end{proof}

\begin{bthm}[fail-closed form]\label{thm:qsb-failclosed}
Under (B1)--(B4), (R3), (R4), ($\Theta$): for every $x \in D$ with
$V(x) \le y$,
\[
\Pr_x\bigl[\forall n,\ \omega_n \in D \smallsetminus U\bigr] \;\ge\; p .
\]
\end{bthm}
\begin{proof}
Apply Proposition~\ref{prop:qsb-pattern} to $(g, A)$ of
Lemma~\ref{lem:qsb-clamp}:
$\Pr_x[\text{reach } A] \le g(x) = V(x) \le y$. The
complement of ever visiting $A = U \cup D^{\mathrm{c}}$ is exactly
$\{\forall n,\ \omega_n \in D \smallsetminus U\}$, so its probability is
at least $1 - y \ge p$ by ($\Theta$). Note the concluded
event is \emph{stronger} than the safety formula\textquotesingle s truth
set: the chain never fails \emph{and never leaves the modelled domain}.
Leaving counts as failure --- that is the fail-closed reading, and it is
what (R4) is recomputed for.
\end{proof}

\begin{bprop}[soundness of \rulename{QSafetyBattery}]\label{prop:qsb-sound}
Under the hypotheses above including the formula tie (F): for every
$x \in D$ with $V(x) \le y$,
$\Pr_x^{\,*}[\forall n,\ \mathrm{sat}(\omega_n)] \ge p$ --- the concluded
satisfaction judgment
$\mathcal{M},\ x_0 \in \{V \le y\} \vDash
\mathsf{SProbGE}\bigl(p,\ \mathsf{PAlways}\ \varphi\bigr)$, read on the
sublevel region $\{V \le y\}$.
\end{bprop}
\begin{proof}
By the contrapositive of (F), $\omega_n \in D \smallsetminus U$ implies
$\mathrm{sat}(\omega_n)$; so the event of
Theorem~\ref{thm:qsb-failclosed} is contained in
$\{\forall n,\ \mathrm{sat}(\omega_n)\}$, and monotonicity of the
(outer) measure finishes. No measurability is ever asked of the
satisfaction predicate --- the smaller, fail-closed event is the
measurable one.
\end{proof}

\begin{bremark}[why the floors are not decoration]\label{rem:qsb-floors}
Drop (B3) and take $V \equiv 0$: drift, nonneg, and threshold rows all
hold, with $V(x) = 0 \le y$ for any $y \ge 0$, so
with $y = 0$, ($\Theta$) permits $p = 1$: the rule would
certify \emph{almost-sure} safety --- of any chain at all, including the
two-state chain $D = \{d\}$, $U = \{u\}$, $K\,d = \delta_u$, which fails
at its first step. With the floor present, this certificate dies
immediately: (B3) forces $V(u) \ge 1$, and then the generator row at $d$
demands $\int V\,dK\,d = V(u) \le V(d) = 0$ --- contradiction. Failure
must cost at least one unit of certificate value; the floor is the
price list.
\end{bremark}

\begin{bremark}[why the nonnegs are not decoration]\label{rem:qsb-nonnegs}
The maximal inequality of Lemma~\ref{lem:qsb-ville} is \emph{false} for
signed supermartingales, and a two-branch chain isolates the sign
hypothesis exactly --- every other row intact, boundedness included.
Take $D = \{s, b\}$, $U = \{a\}$, $\bar{D} = \{s, a, b\}$, with
$K\,s = \tfrac{1}{2}\delta_a + \tfrac{1}{2}\delta_b$,
$K\,b = \delta_b$, and the signed certificate $V(s) = 0$, $V(a) = 1$,
$V(b) = -1$. The generator row holds on $D \smallsetminus U$
($\int V\,dK\,s = 0 \le 0$; $\int V\,dK\,b = -1 \le -1$), the floor
holds on $U \cap \bar{D}$, the hull and fail-closed rows hold, and
$|V| \le 1$. Were Proposition~\ref{prop:qsb-pattern} available without
nonnegativity, the certificate would conclude
$\Pr_s[\text{reach } a] \le V(s) = 0$ --- but the chain reaches $a$
with probability $\tfrac{1}{2}$: the $a$-branch's barrier visit is
subsidized by the mass parked at $V = -1$ on the other branch, and
(B4) is what forbids the loan. The classical unbounded picture is the
simple symmetric random walk on $\mathbb{Z}$ with $V(z) = z$ and
$A = \{1\}$ --- same phenomenon, though as a drop-one witness it is
impure: it violates the boundedness hypothesis along with the sign,
which is why the two-branch chain is the recorded counterexample.
\end{bremark}

\begin{bremark}[the hull and fail-closed conditions (R3), (R4)]
\label{rem:qsb-r34}
Both are recomputed by the checker, never taken from the producer, and
each blocks a distinct hole. Without (R3), the one-step integral in the
generator row could consult states outside $\bar{D}$ --- states where the
floor and nonneg rows were \emph{never checked} --- and
Lemma~\ref{lem:qsb-clamp}(iv)\textquotesingle s clamp-invisibility
computation collapses. Without (R4), one-step mass could leave
$D \cup U$ entirely: the concluded event counts leaving the modelled
domain as failure, and (R4) is exactly what makes the clamped
certificate\textquotesingle s drift hold at the \emph{enlarged} stop set
$U \cup D^{\mathrm{c}}$ --- the design that let every hypothesis in
Lemma~\ref{lem:qsb-bridge} stay pointwise. A companion condition
(floor-scope resolution) is discharged at row-compilation time: it is
what makes the checked floor rows imply (B3) as stated here, and it has
no semantic residue at this layer.
\end{bremark}

\begin{bremark}[threshold and integrability]\label{rem:qsb-theta}
($\Theta$) is bookkeeping, and it is sharp: the certificate buys
$\Pr[\mathrm{fail}] \le y$, so asserting success
$\ge p$ requires $p \le 1 - y$, and
$y = 1 - p$ is the tightest admissible choice --- producer
weakening ($y < 1 - p$) is sound for free. The integrability
half of (B1$\wedge$B2) is the discrete-profile reading of
``the certificate lies in the generator\textquotesingle s
domain''; for the intended instances --- polynomial $V$, one-step
images inside the compact hull by (R3) --- it is automatic.
\end{bremark}

\begin{bremark}[\textbf{profile} --- what is proven, at which
presentations]
\label{rem:qsb-profile}
Proven, and quoted above this proof: the discrete-kernel profile ---
one-step Markov kernel semantics over the trajectory construction, the
theorems \texttt{sound\_QSafetyBattery} and
\texttt{sound\_QSafetyBattery\_failClosed}, sorry-free. Also proven,
end-to-end: the \emph{PDMP sub-profile} of the continuous-time semantics
--- piecewise-deterministic flow with guarded jumps, no diffusion
(\texttt{sound\_GSHS\_safety\_battery}). And proven, at the profile
that carries the general continuous-time semantics, jump--diffusions
included: the \emph{martingale-problem profile} --- the system
presented by its generator drift data, $\mathcal{A}V$ carried as data
beside the certificate with the martingale-problem hypothesis for the
pair, c\`adl\`ag paths per the presentation --- the theorems
\texttt{sound\_GSHS\_safety\_mp} and
\texttt{sound\_GSHS\_safety\_battery\_mp}, machine-checked sorry-free
in the semantic-universe strand\textquotesingle s auxiliary package
(module \texttt{QtslAux.LayerD.MPDischarge}). The three ingredients an
earlier edition of this remark filed as missing are each discharged
there: continuous-time generator domains ride the presentation ---
for the calculus\textquotesingle s own denoted class, compositional
symbolic piecewise-polynomial hypernets, $\mathcal{A}V$ is again
polynomial data and domain bookkeeping dissolves; fine path
properties ride the presentation\textquotesingle s c\`adl\`ag
clause; and optional stopping for right-continuous nonnegative
supermartingales is a library theorem, consumed at the exit
d\'ebut (\texttt{OptionalSamplingRC}, with
\texttt{supermartingale\_of\_compensator\_nonpos} and the
exponential-discount construction). The epistemic shape is the
discrete profile\textquotesingle s own: exactly as the staged theorem
quantifies over systems whose compilation satisfies the mixture
bridge, the martingale-problem theorems quantify over systems whose
presentation satisfies the martingale-problem hypothesis --- an
explicit bridge carried as data, never a silent assumption, so the
argument this section walks is now the theorem\textquotesingle s own
at every profile. The martingale-problem theorems live in the
semantic-universe strand\textquotesingle s auxiliary package, which is
registered with this catalogue\textquotesingle s binder: their
statements are extracted and hash-bound like every other block in
this section, so the dashboard row carries no conditional mark ---
each named profile is proven, and its theorems are bound.
\end{bremark}

\begin{bremark}[the staged premise shapes and the system bridge]
\label{rem:qsb-staged}
The rule\textquotesingle s node presents its premises per mode and per
jump: mode-indexed certificate family $(V_m)$, interior rows against
each mode\textquotesingle s flow mechanism, guard rows per jump, floors
and nonnegs per mode, and the threshold row in exact rational
arithmetic, cast once. The bound theorem for that shape
(\texttt{sound\_QSafetyBattery\_staged}) closes the gap to the flattened
hypotheses used above through an explicit bridge hypothesis: at every
in-scope state, the one-step kernel decomposes as a sub-convex mixture
of the mode\textquotesingle s flow step and the enabled guard-reset
jumps (weights supported on enabled jumps). Under that decomposition the
mode\textquotesingle s interior row and the enabled guard rows
\emph{are} the flattened generator row (\texttt{mixture\_drift}): the
one-step integral splits over the mixture, each jump term is a guard row
evaluation, and the weights sum to one. What stands left of the
concluded $\vDash$ is any system presentation whose compilation at this
profile satisfies the bridge.
\end{bremark}

\begin{bremark}[relation to the literature]\label{rem:qsb-lit}
Supermartingale safety certificates for stochastic systems go back to
Kushner \cite{kushner1967}; the modern certificate-form idiom ---
synthesize a barrier function, check polynomial rows, conclude a
probability bound --- is Prajna--Jadbabaie--Pappas
\cite{prajnajadbabaiepappas2007}. The probability engine is
Ville\textquotesingle s inequality \cite{ville1939}, with the systematic
supermartingale treatment in Doob \cite{doob1953}. The reading of the
interior rows as one-step expected decrease is the Foster--Lyapunov
drift idiom \cite{meyntweedie2009}, whose geometric form powers the
deadline sibling. What this rule adds to that literature is not
mathematics but \emph{trust structure}: every analytic ingredient lives
in a verified lemma, every polynomial ingredient is an inherited
base-calculus row replayed by the kernel, and the checker recomputes the
side conditions rather than trusting the producer.
\end{bremark}

\begin{anchors}
Lemma \ref{lem:qsb-stopped}: stopped certificate, adaptedness, barrier transport & \texttt{stopped}, \texttt{survive}, \texttt{stronglyAdapted\_stopped}, \texttt{one\_le\_stopped}, \texttt{stopped\_congr}\\
Lemma \ref{lem:qsb-bridge}: supermartingale bridge & \texttt{supermartingale\_stopped} (Markov property via \texttt{Kernel.condExp\_traj}; prefix patching via \texttt{Kernel.integral\_traj}; one-step marginal via \texttt{Kernel.map\_traj\_succ\_self})\\
Lemma \ref{lem:qsb-ville}: maximal inequality at barrier $1$ & \texttt{BernsteinLogic.Ville.reach\_avoid}, consumed at $a = 1$\\
Proposition \ref{prop:qsb-pattern}: the battery pattern & \texttt{measure\_reach\_le}\\
Lemma \ref{lem:qsb-clamp}: the clamp & the construction $g$, \texttt{hclamp}, \texttt{hgdrift} inside \texttt{sound\_QSafetyBattery\_failClosed}\\
Theorem \ref{thm:qsb-failclosed} & \texttt{sound\_QSafetyBattery\_failClosed}\\
Proposition \ref{prop:qsb-sound}: formula tie, monotonicity & \texttt{sound\_QSafetyBattery}\\
Remark \ref{rem:qsb-staged}: staged shapes, mixture bridge & \texttt{mixture\_integrable}, \texttt{mixture\_drift}, \texttt{sound\_QSafetyBattery\_staged}\\
Remark \ref{rem:qsb-profile}: PDMP sub-profile & \texttt{sound\_GSHS\_safety\_battery}\\
Remark \ref{rem:qsb-profile}: martingale-problem profile & \texttt{sound\_GSHS\_safety\_mp}, \texttt{sound\_GSHS\_safety\_battery\_mp} (\texttt{QtslAux.LayerD.MPDischarge}; hash-bound via the qtsl-aux package registration)\\
\end{anchors}


\rulesection{41}{QDeadlineBattery}{core rule}
\label{blcatrule-QDeadlineBattery}
\ruleblock{The rule (Steele metanotation)}
\sourcenote{as realized in the mechanized development; source row: \texttt{\texttt{QTSL-rules.toml}} --- \emph{QDeadlineBattery}. The metanotation statement is typeset by hand for notational quality; its mechanized ground truth is the Lean statement below, extracted from the sources at every build of this book and hash-bound. The rule \emph{roster} is machine-checked (every mechanized rule tag round-trips to exactly one catalogue row); the statement \emph{text} is editorial, checked by eye against the extracted Lean beside it.}
\begin{mathpar}
\inferrule*{
  V : \mathsf{QCertFam} \\
  R \subseteq [\![\varphi]\!] \ \text{(target formula tie)} \\
  \bigl\{D_m \smallsetminus R \vdash \mathcal{A}_m V_m \le -\theta V_m\bigr\}_{m} \ \text{(decay)} \\
  \bigl\{\mathrm{slab}_j \vdash g_j \fatsemi V_{m'_j} \le V_{m_j}\bigr\}_{j} \ \text{(guards, per jump $j$: source mode $m_j$, target $m'_j$, reset $g_j$)} \\
  \bigl\{\bar{D}_m \smallsetminus R \vdash V_m \ge 1\bigr\}_{m} \ \text{(floors, off-target only)} \\
  \bigl\{\bar{D}_m \vdash V_m \ge 0\bigr\}_{m},\ \bigl\{\bar{D}_m \vdash V_m \le M\bigr\}_{m} \ \text{(nonnegs; bound)} \\
  \vdash q \le \textstyle\sum_{k \le n} (\theta T)^k / k! \ \text{(one \rulename{TaylorBound} leaf)} \\
  0 \le \beta,\ \beta \le e^{-\theta},\ \theta \ge 0,\ 0 \le q,\ p \le 1 \ \text{(drift tie; sign gates)}
}{
  \mathcal{M},\ x_0 \in \{V \le (1-p)\,q\} \vDash \mathsf{SProbGE}\bigl(p,\ \mathsf{PDeadline}(T, \varphi, b)\bigr)
}
\end{mathpar}

\ruleflavor{rule sort baked in \texttt{theory Qtsl} (\texttt{QDeadlineBattery}; the tabular tier\textquotesingle s \texttt{QDeadlineTable} beside it) --- the display is the realized form, as the discrete-profile theorems bind it. A region left of the turnstile scopes the row ($R \vdash P$: $P$ at every point of $R$); the certificate family $V$ enters as certificate data --- a shared reference with no claim content of its own; the rows impose every claim --- never as a derivation edge. In the drift tie and sign gates, both transcendental sides ($\beta \le e^{-\theta}$ and the Taylor leaf\textquotesingle s $e^{\theta T}$ enclosure) live in the soundness lemma, never in the checker. The conclusion is the region-conditioned judgment in constructor vocabulary, as at \rulename{QSafetyBattery}; $b$ is \texttt{PDeadline}\textquotesingle s producer-chosen-horizon mark.}
\par\medskip
\rulename{QDeadlineBattery} is the bounded-reachability battery: it
concludes that with probability at least $p$ the system reaches the
target --- satisfies $\mathsf{PDeadline}(T, \varphi, b)$ (surface
syntax $\mathsf{F}_{\le T}\,\varphi$) --- from every start
in the sublevel region $\{V \le (1-p)\,q\}$. Its engine is exponential
decay: decay rows say the certificate shrinks by a factor $\beta$ per
step off the target, so the discounted process $\beta^{-t} V$ is a
supermartingale off the target, while the off-target floor $V \ge 1$
makes \emph{missing} the deadline cost at least one unit of discounted
certificate value; Markov\textquotesingle s inequality at the horizon
gives $\Pr[\mathrm{miss}] \le \beta^{T} V(x)$.

The rule\textquotesingle s division of labour between checker and lemma
is its most instructive feature. The checker replays only exact rational
comparisons: the one arithmetic leaf is the Taylor comparison
$q \le \sum_{k \le n} (\theta T)^k / k!$ --- a closed, wire-free
inequality discharged by the ground-fact rule --- while both
transcendental facts ($\sum_{k \le n} (\theta T)^k/k! \le e^{\theta T}$
and $e^{-\theta} e^{\theta} = 1$) live inside the soundness lemma. The
checker never needs to know $e$ exists.

What the checker recomputes beyond that leaf: the domain-hull and
fail-closed side conditions (as in the safety battery), and the
threshold arithmetic in the exact rational vocabulary
($\mathrm{thr}_{\mathsf F} = \mathrm{thr} \cdot q$ atop
$p + \mathrm{thr} \le 1$). Soundness is independent of the producer
knobs $\theta$ and $q$: any discrete rate $0 \le \beta \le e^{-\theta}$
and any $q$ below the Taylor sum are admissible.

Each hypothesis is load-bearing. Drop the off-target floor and a
certificate identically $0$ would certify certain reachability of an
unreachable target. Drop the boundedness row and the decay row no longer transfers to the
clamped certificate --- the discount multiplies the certificate by
$\beta^{-t}$, so the ceiling must be invisible wherever one-step mass
reaches, which is exactly what $V \le M$ buys (and, in the mechanized
engine, the integrability of the discounted certificate rides the same
bound). Drop the
fail-closed containment and escape from the modelled domain would count
as success: a reach guarantee must show the escape event \emph{null},
where the safety battery could simply enlarge its stop set (leaving
counts as failing is conservative for safety, but would be unsound
generosity for reachability). Drop the threshold arithmetic and the
concluded region over-budgets. The proof section walks each.

The dashboard mark states the status: the hash-bound theorems warrant
the rule at the discrete-kernel profile (geometric drift,
$\beta = e^{-\theta}$ discretized); the continuous-time profile is
proven at the martingale-problem presentation --- the discounted
supermartingale built from the undiscounted martingale-problem
hypothesis by the exponential-discount construction, stopped at the
d\'ebut by the right-continuous optional-sampling library theorem
(\texttt{sound\_GSHS\_deadline\_mp},
\texttt{sound\_GSHS\_deadline\_battery\_mp},
\texttt{QtslAux.LayerD.MPDischarge}) --- hash-bound below via the
qtsl-aux package registration, so the dashboard row carries no
conditional mark: each named profile is proven, and its theorems are
bound.

At the general stratum, over the semantic universe\textquotesingle s
record carrier (the semantic-universe functor package, hash-bound
below), the deadline conclusion is a \emph{mark} --- an
$\exists$-face spec: the per-window family of acceptance-event
demands is persistent, and satisfaction forces almost-sure acceptance
at the window (\texttt{satAtM\_acceptanceSpecF\_ae}). The backward
direction is deliberately absent, and the absence is the semantics: a
record that reaches its target only in a superwindow is invisible to
the subwindow, so a subwindow\textquotesingle s reach demand is a
genuinely stronger demand --- a deadline --- and extension is exactly
how a met deadline persists to every superwindow. The measurability
the mark face charges is issued by the site\textquotesingle s
topology rather than assumed: hit-events commute with countable
unions in the set argument (\texttt{hitEvent\_iUnion}), turning grid
decompositions into event decompositions, and a compact target
cell\textquotesingle s visit event over a $\sigma$-compact stage
window is Borel
(\texttt{measurableSet\_hitEvent\_univ\_prod\_of\_isCompact}). One
typing note rides the transposition: the site-general grids are
indexed by bare naturals; at clocked continuum strata the legs
re-type to rational-duration stage cuts --- the tick-indexed reading
of a windowed horizon is a forcing of the finitary record site, not
of the grammar.

\cataloguenote{The checked floor is \emph{off-target only}: on the target the
certificate may take any value, since success needs no budget. This is
deliberately weaker than a global floor, and the proven statement binds
the weaker form; a catalogued rendering that states a global floor is
sound but relaxable to match. The concluded threshold $(1-p)\,q$ is the
compile\textquotesingle s choice; the underlying theorem consumes only
the single inequality $\beta^{T} \cdot \mathrm{thr} \le 1 - p$. At this
profile, guard rows land inside the one-step decay row (the kernel
carries flow and jumps together); the instantaneous-jump reading, where
the discount factor does not tick during a jump, is carried at the
continuous-time profile by the martingale-problem
discharge\textquotesingle s no-alarm segment assembly.
}
\ruleblock{Lean statement}
\declhead{BernsteinLogicAux.Battery}{sound\_QDeadline}{\badge{BadgeGreen}{statement bound}}
\begin{leansrc}
/-- **`sound_QDeadline`, discrete profile** (T6): over the target-hit form, the deadline
formula tie — the target region *under*-approximates the truth set of the
`𝖥_{≤T}`-argument — weakens the concluded event to the truth of the reach formula:
`Pr_x[∃ t ≤ T, X t ⊨ φ] ≥ p` on the concluded region `D ∩ {V ≤ thr}`. No measurability
is asked of the satisfaction predicate (`measure_mono` on the measurable hit event). -/
theorem sound_QDeadline {T : ℕ} {M : ℝ} {sat : S → Prop}
    (hD : MeasurableSet D) (hDbar : MeasurableSet Dbar) (hB : MeasurableSet B)
    (hV : Measurable V)
    (hF : ∀ y ∈ B, sat y)
    (hβ : 0 ≤ β)
    (hΘF : β ^ T * thr ≤ 1 - p)
    (hDsub : D ⊆ Dbar)
    (hR3 : ∀ y ∈ D, y ∉ B → K y Dbar = 1)
    (hR4 : ∀ y ∈ D, y ∉ B → K y (D ∪ B) = 1)
    (hint : ∀ y ∈ D, y ∉ B → Integrable V (K y))
    (hdecay : ∀ y ∈ D, y ∉ B → ∫ z, V z ∂K y ≤ β * V y)
    (hfloor : ∀ y ∈ Dbar, y ∉ B → 1 ≤ V y)
    (hnonneg : ∀ y ∈ Dbar, 0 ≤ V y)
    (hM : ∀ y ∈ Dbar, V y ≤ M)
    (hx : x ∈ D) (hthr : V x ≤ thr) :
    ENNReal.ofReal p ≤ pathMeasure K x {ω | ∃ k ≤ T, sat (ω k)}
\end{leansrc}
\declhead{BernsteinLogicAux.Battery}{sound\_QDeadline\_hit}{\badge{BadgeGreen}{statement bound}}
\begin{leansrc}
/-- **`sound_QDeadline`, discrete profile, target-hit form** (T6 — the safety battery's
deadline successor — at the discrete-kernel profile). Hypotheses are the
decay rows (`∫ V ∂K y ≤ β·V y` off the target `B`, the discrete `𝓛V ≤ -θV` with
`β = e^{-θ}`), floors **off** the target, nonnegs, a certificate bound `M` on the cover,
side conditions (R3)/(R4), and the deadline threshold arithmetic
`(Θ_F): β ^ T * thr ≤ 1 - p` — the inequality form of the compile's `thr = (1-p)·q` with
`q ≤ e^{θT}`. Conclusion: from any start in `D ∩ {V ≤ thr}`, the chain visits `B` by time
`T` with probability at least `p`. Consumes the landed geometric-drift engine
`BernsteinLogic.Deadline.deadline_tail` at the enlarged target `B ∪ Dᶜ`,
then removes the escape branch with `ae_no_escape`. -/
theorem sound_QDeadline_hit {T : ℕ} {M : ℝ}
    (hD : MeasurableSet D) (hDbar : MeasurableSet Dbar) (hB : MeasurableSet B)
    (hV : Measurable V)
    (hβ : 0 ≤ β)
    (hΘF : β ^ T * thr ≤ 1 - p)
    (hDsub : D ⊆ Dbar)
    (hR3 : ∀ y ∈ D, y ∉ B → K y Dbar = 1)
    (hR4 : ∀ y ∈ D, y ∉ B → K y (D ∪ B) = 1)
    (hint : ∀ y ∈ D, y ∉ B → Integrable V (K y))
    (hdecay : ∀ y ∈ D, y ∉ B → ∫ z, V z ∂K y ≤ β * V y)
    (hfloor : ∀ y ∈ Dbar, y ∉ B → 1 ≤ V y)
    (hnonneg : ∀ y ∈ Dbar, 0 ≤ V y)
    (hM : ∀ y ∈ Dbar, V y ≤ M)
    (hx : x ∈ D) (hthr : V x ≤ thr) :
    ENNReal.ofReal p ≤ pathMeasure K x {ω | ∃ k ≤ T, ω k ∈ B}
\end{leansrc}
\par\medskip
The catalogue form: the $(\theta, q)$ parameterization, concluding on
$\{V \le (1-p)\,q\}$ from the Taylor leaf and the drift tie
$\beta \le e^{-\theta}$, with both transcendental facts inside the
lemma. It is a wrapper over \texttt{sound\_QDeadline}, which carries the
general inequality form $\beta^{T} \cdot \mathrm{thr} \le 1 - p$.

\declhead{BernsteinLogicAux.QtslBattery}{sound\_QDeadlineBattery}{\badge{BadgeGreen}{statement bound}}
\begin{leansrc}
/-- **`sound_QDeadlineBattery`** (QTSL-rules.toml row `QDeadlineBattery` —
the catalogue's `(θ, q)` parameterization of T6, discrete-kernel profile).
The concluded region is `{V ≤ (1 − p)·q}`: the deadline threshold
`thr_F = thr · q` in the staged `q_mul` vocabulary atop `thr = 1 − p`.
Its arithmetic premises are

* the **T7 Taylor leaf** `q ≤ ∑_{k ≤ n} (θT)^k / k!` — the one closed
  wire-free comparison the checker replays in exact ℚ (rule row
  `TaylorBound (T7)`, alias → `PConstExact`), and
* the **drift tie** `β ≤ e^{−θ}` — the producer's discrete geometric rate
  matching the catalogue's `β = e^{−θ}` (any smaller rate is sound).

The transcendental facts — `∑_{k ≤ n} (θT)^k/k! ≤ e^{θT}`
(`Real.sum_le_exp_of_nonneg`) and `e^{−θ}·e^{θ} = 1` — live in this lemma,
never in the checker (the catalogue's "the checker never needs to know `e`
exists"). Wrapper over `sound_QDeadline`, which carries the general
inequality form `β ^ T · thr ≤ 1 − p`.

System bridge: for any System whose chart compilation at
the discrete-kernel profile yields `K`, this is the concluded
`QRegionSat`'s reading on that System's records. -/
theorem sound_QDeadlineBattery {n : ℕ} {sat : S → Prop}
    (hD : MeasurableSet D) (hDbar : MeasurableSet Dbar) (hB : MeasurableSet B)
    (hV : Measurable V)
    (hF : ∀ y ∈ B, sat y)
    (hβ : 0 ≤ β) (hβθ : β ≤ Real.exp (-θ)) (hθ : 0 ≤ θ)
    (hq0 : 0 ≤ q)
    (hq : q ≤ ∑ k ∈ range (n + 1), (θ * T) ^ k / (Nat.factorial k : ℝ))
    (hp : p ≤ 1)
    (hDsub : D ⊆ Dbar)
    (hR3 : ∀ y ∈ D, y ∉ B → K y Dbar = 1)
    (hR4 : ∀ y ∈ D, y ∉ B → K y (D ∪ B) = 1)
    (hint : ∀ y ∈ D, y ∉ B → Integrable V (K y))
    (hdecay : ∀ y ∈ D, y ∉ B → ∫ z, V z ∂K y ≤ β * V y)
    (hfloor : ∀ y ∈ Dbar, y ∉ B → 1 ≤ V y)
    (hnonneg : ∀ y ∈ Dbar, 0 ≤ V y)
    (hM : ∀ y ∈ Dbar, V y ≤ M)
    (hx : x ∈ D) (hthr : V x ≤ (1 - p) * q) :
    ENNReal.ofReal p ≤ pathMeasure K x {ω | ∃ k ≤ T, sat (ω k)}
\end{leansrc}
\par\medskip
\textbf{The general-stochastic instance, assembled.}\enspace The
deadline bound on the timed skeleton of the segment model: the
$\theta$-discounted generator rows enter per segment, the clock is a
coordinate of the timed chain, and the conclusion is the displayed
deadline-overrun bound on that chain\textquotesingle s events ---
headline, off-target, and battery forms.

\declhead{BernsteinLogicAux.LayerD.Deadline}{sound\_GSHS\_deadline}{\badge{BadgeGreen}{statement bound}}
\begin{leansrc}
/-- **The GSHS deadline tail bound** (T6, headline form). Inherited
premises, stated (the leg's discipline — never silent):

* `hSegD : SegmentDynkin M θ V` at the discount `θ` — **the θ > 0 premise**:
  the B1 row compiled at positive discount; the
  statement is typed at `0 ≤ θ` and degenerates to the trivial floor bound
  at `θ = 0`.
* `hFloor : ∀ y, ε ≤ V y` — **the F7 floor premise**: the
  positive floor on the working region, conclusion scaled by `1/ε`.
  Without it the conclusion is unsound (the review's two-point witness: `V ≡ 0`
  off-target satisfies every remaining hypothesis while the deadline fails
  surely).

Conclusion: the probability that the timed skeleton chain ever presents a
value stamped past the deadline `T` is at most `e^{−θT} V(x)/ε` — the T6
display `Pr[τ_B > T] ≤ e^{−θT} V(x₀)/ε` in the value-presentation reading:
instantiated at the stopped presentation `Rep(P^{†B})` with halts into the
target ("the stop IS the target" — a design choice
made here and visible), a path outside this event either reached the
target by `T` or stalled (Zeno / vanished — the residues emptied a.s. for
in-scope models by the prose W-Z-freedom theorem).
Consumes `Battery.measure_reach_le` (= `Ville.reach_avoid`) over the
timed chain: the deadline battery rides literally the same engine as T5,
with the capped certificate `min(e^{θs}V, εe^{θT})/εe^{θT}` — and no
global bound `V ≤ C` is demanded (the cap manufactures its own). -/
theorem sound_GSHS_deadline [IsMarkovKernel M.seg] [IsMarkovKernel M.jump]
    {V : State → ℝ} {θ ε : ℝ} {T : ℝ≥0}
    (hV : Measurable V) (hθ : 0 ≤ θ)
    (hSegD : SegmentDynkin M θ V) (hForce : ForcedNonexpansion M V)
    (hNonneg : NonnegV V) (hε : 0 < ε) (hFloor : ∀ y, ε ≤ V y) (x : State) :
    Battery.pathMeasure (tskel M) (TSkelState.alive (0, x))
        {ω | ∃ n, T < tclock (ω n)}
      ≤ ENNReal.ofReal (Real.exp (-(θ * T)) * V x / ε)
\end{leansrc}
\declhead{BernsteinLogicAux.LayerD.Deadline}{sound\_GSHS\_deadline\_offTarget}{\badge{BadgeGreen}{statement bound}}
\begin{leansrc}
/-- **T6, off-target-floor variant** (the sharper premise, the narrower
event): the floor is demanded off the target only — the F7 floor premise
read narrowly — and the conclusion event asks for a past-deadline
presentation *outside* `B`. What this variant gives up is the crossing-jump
corner: a path whose first past-`T` presentation is the target hit itself
(the dispatch that crossed the deadline landed in `B`) misses the deadline
without entering this event — `exceed_cases` makes the dichotomy explicit.
The headline `sound_GSHS_deadline` closes the corner by flooring the whole
working region instead. -/
theorem sound_GSHS_deadline_offTarget [IsMarkovKernel M.seg] [IsMarkovKernel M.jump]
    {V : State → ℝ} {θ ε : ℝ} {T : ℝ≥0} {B : Set State}
    (hV : Measurable V) (hθ : 0 ≤ θ) (hB : MeasurableSet B)
    (hSegD : SegmentDynkin M θ V) (hForce : ForcedNonexpansion M V)
    (hNonneg : NonnegV V) (hε : 0 < ε) (hFloorB : ∀ y ∉ B, ε ≤ V y)
    (x : State) :
    Battery.pathMeasure (tskel M) (TSkelState.alive (0, x))
        {ω | ∃ n, T < tclock (ω n) ∧ ω n ∉ tbadIn B}
      ≤ ENNReal.ofReal (Real.exp (-(θ * T)) * V x / ε)
\end{leansrc}
\declhead{BernsteinLogicAux.LayerD.Deadline}{sound\_GSHS\_deadline\_battery}{\badge{BadgeGreen}{statement bound}}
\begin{leansrc}
/-- **The deadline battery form** (T6 as `sound_QDeadline`'s GSHS instance,
complement arithmetic): with the threshold accounting
`p + e^{−θT}·thr/ε ≤ 1` and a start under the threshold, the chain's
presented clock stays `≤ T` at every index with probability at least `p`.
On a stopped-at-target instantiation with halts into the target, that event
reads *reached the target by the deadline* — up to the two stall residues
(forced-jump Zeno accumulation; vanishing into a never-terminating
segment), which are inside the event, named, and emptied a.s. for in-scope
models by the prose W-Z-freedom theorem: the honest
one-sided residue, visible rather than erased. -/
theorem sound_GSHS_deadline_battery [IsMarkovKernel M.seg] [IsMarkovKernel M.jump]
    {V : State → ℝ} {θ ε p thr : ℝ} {T : ℝ≥0}
    (hV : Measurable V) (hθ : 0 ≤ θ)
    (hSegD : SegmentDynkin M θ V) (hForce : ForcedNonexpansion M V)
    (hNonneg : NonnegV V) (hε : 0 < ε) (hFloor : ∀ y, ε ≤ V y)
    (hΘ : p + Real.exp (-(θ * T)) * thr / ε ≤ 1) (x : State) (hthr : V x ≤ thr) :
    ENNReal.ofReal p ≤ Battery.pathMeasure (tskel M) (TSkelState.alive (0, x))
      {ω | ∀ n, tclock (ω n) ≤ T}
\end{leansrc}
\par\medskip
\textbf{The martingale-problem profile, end to end.}\enspace The
general continuous-time theorems: the discounted supermartingale is
built from the undiscounted martingale-problem hypothesis by the
exponential-discount construction, stopped at the d\'ebut by the
right-continuous optional-sampling library theorem, and the deadline
conclusion follows for every system whose presentation satisfies the
martingale-problem hypothesis.

\declhead{QtslAux.LayerD.MPDischarge}{sound\_GSHS\_deadline\_mp}{\badge{BadgeGreen}{statement bound}}
\begin{leansrc}
/-- The deadline reach bound at the martingale-problem profile. -/
theorem sound_GSHS_deadline_mp [IsMarkovKernel M.seg] [IsMarkovKernel M.jump]
    {V : State → ℝ} {θ ε : ℝ} {T : ℝ≥0}
    (hV : Measurable V) (hθ : 0 ≤ θ)
    (hmp : ∀ x : State, Nonempty (MPFlow M θ V x))
    (hForce : ForcedNonexpansion M V)
    (hNonneg : NonnegV V) (hε : 0 < ε) (hFloor : ∀ y, ε ≤ V y) (x : State) :
    Battery.pathMeasure (tskel M) (TSkelState.alive (0, x))
        {ω | ∃ n, T < tclock (ω n)}
      ≤ ENNReal.ofReal (Real.exp (-(θ * T)) * V x / ε)
\end{leansrc}
\declhead{QtslAux.LayerD.MPDischarge}{sound\_GSHS\_deadline\_battery\_mp}{\badge{BadgeGreen}{statement bound}}
\begin{leansrc}
/-- The deadline battery at the martingale-problem profile. -/
theorem sound_GSHS_deadline_battery_mp
    [IsMarkovKernel M.seg] [IsMarkovKernel M.jump]
    {V : State → ℝ} {θ ε p thr : ℝ} {T : ℝ≥0}
    (hV : Measurable V) (hθ : 0 ≤ θ)
    (hmp : ∀ x : State, Nonempty (MPFlow M θ V x))
    (hForce : ForcedNonexpansion M V)
    (hNonneg : NonnegV V) (hε : 0 < ε) (hFloor : ∀ y, ε ≤ V y)
    (hΘ : p + Real.exp (-(θ * T)) * thr / ε ≤ 1) (x : State)
    (hthr : V x ≤ thr) :
    ENNReal.ofReal p ≤ Battery.pathMeasure (tskel M) (TSkelState.alive (0, x))
      {ω | ∀ n, tclock (ω n) ≤ T}
\end{leansrc}
\par\medskip
\textbf{The general stratum: grid tickets by topology.}\enspace The
countable-combination engine and the reach leg\textquotesingle s
measurability ticket over the record carrier: hitting commutes with
unions in the set argument, and a compact target cell\textquotesingle s
visit event over a $\sigma$-compact stage window is Borel --- issued
by the site\textquotesingle s topology, not assumed.

\declhead{SuFunctor.Site.FellAcceptance}{hitEvent\_iUnion}{\badge{BadgeGreen}{statement bound}}
\begin{leansrc}
/-- **Hitting commutes with unions in the set argument** — the engine
that turns countable set decompositions into countable event
decompositions.  Arbitrary index, pure set algebra. -/
theorem hitEvent_iUnion {ι : Sort*} (A : ι → Set α) :
    hitEvent (⋃ i, A i) = ⋃ i, hitEvent (A i)
\end{leansrc}
\declhead{SuFunctor.Site.FellAcceptance}{measurableSet\_hitEvent\_iUnion}{\badge{BadgeGreen}{statement bound}}
\begin{leansrc}
/-- Hit events pass countable unions of measurable pieces. -/
theorem measurableSet_hitEvent_iUnion {ι : Sort*} [Countable ι]
    {A : ι → Set α} (h : ∀ i, MeasurableSet (hitEvent (A i))) :
    MeasurableSet (hitEvent (⋃ i, A i))
\end{leansrc}
\declhead{SuFunctor.Site.FellAcceptance}{measurableSet\_hitEvent\_univ\_prod\_of\_isCompact}{\badge{BadgeGreen}{statement bound}}
\begin{leansrc}
/-- **Hit-events of certified cells are Effros-measurable for free**
(the memo's headline, cashed): a compact cell's visit event over a
σ-compact stage window is a countable union of Fell-closed hit
events. -/
theorem measurableSet_hitEvent_univ_prod_of_isCompact
    [SigmaCompactSpace ↥W] {R : Set X} (hR : IsCompact R) :
    MeasurableSet
      (hitEvent ((Set.univ : Set ↥W) ×ˢ R) : Set (FTraj W X))
\end{leansrc}
\par\medskip
\textbf{Marks over the record carrier.}\enspace The deadline
conclusion\textquotesingle s general-stratum shape: acceptance-event
spec families, persistent by window-calculus transitivity, with the
forward judgment lemma; the backward direction is deliberately absent
--- deadline semantics, as the commentary states.

\declhead{SuFunctor.Site.FellBehaviorSpec}{acceptanceSpecF}{\badge{BadgeGreen}{statement bound}}
\begin{leansrc}
/-- The spec family of a per-window acceptance event (deadline, until,
Streett): clip demands of the events' indicators. -/
def acceptanceSpecF (U : ι → Set T)
    (E : ∀ i, Set (FTraj (U i) X)) (j : ι) :
    Set (MeasPen (FTraj (U j) X))
\end{leansrc}
\declhead{SuFunctor.Site.FellBehaviorSpec}{isPersistentMF\_acceptanceSpecF}{\badge{BadgeGreen}{statement bound}}
\begin{leansrc}
theorem isPersistentMF_acceptanceSpecF (U : ι → Set T)
    (hU : ∀ i, IsOpen (U i)) (E : ∀ i, Set (FTraj (U i) X)) :
    IsPersistentMF U (acceptanceSpecF U E)
\end{leansrc}
\declhead{SuFunctor.Site.FellBehaviorSpec}{satAtM\_acceptanceSpecF\_ae}{\badge{BadgeGreen}{statement bound}}
\begin{leansrc}
/-- The forward judgment: satisfying an acceptance spec at a window
forces almost-sure acceptance there.  The backward direction is
deliberately absent — liveness events do not clip-restrict, so
subwindow demands are genuinely stronger (a reach demand within a
subwindow is a deadline); the module docstring carries the reading. -/
theorem satAtM_acceptanceSpecF_ae (U : ι → Set T)
    (E : ∀ i, Set (FTraj (U i) X)) (j : ι)
    {μ : MeasJoint (FTraj (U j) X)}
    (h : SatAtM μ (acceptanceSpecF U E j)) :
    (μ.1 : Measure (FTraj (U j) X)) (E j)ᶜ = 0
\end{leansrc}
\ruleblock{Correspondence (Steele display $\leftrightarrow$ Lean)}
\badge{BadgeAmber}{a bound Lean statement CHANGED since certification --- re-review and re-certify}
\sourcenote{The dictionary below is a reviewer-certified statement of how the mechanized statements above interpret the displayed rule: display metavariable $\mapsto$ Lean binder, each display hypothesis $\mapsto$ the premise or side condition carrying it, the display conclusion $\mapsto$ the goal --- with any deliberate difference in width between display and mechanization named outright. Certification is recorded as a hash triple over (the displayed cell, the extracted statement texts, this dictionary\textquotesingle s own text); if any of the three changes, this block turns amber until a reviewer re-certifies. The correspondence claim itself remains editorial --- reviewer-grade, not machine-checked; its machine-checked successor types each soundness statement against the rule as a structured object, and supersedes this block per rule as it lands.}
Model frame and profile as at \rulename{QSafetyBattery}: hash-bound at
the discrete-kernel profile (the amended cell\textquotesingle s own
caveat), with the continuous-time profile proven at the
martingale-problem presentation (hash-bound below, per the profile
remark). The catalogue form binds
\texttt{sound\_QDeadlineBattery}; \texttt{sound\_QDeadline} carries
the general inequality it wraps.
\begin{itemize}
\item the displayed judgment $\mathcal{M}
  ,\ x_0 \in \{V \le (1-p)\,q\} \vDash
  \mathsf{SProbGE}\bigl(p,\ \mathsf{PDeadline}(T, \varphi, b)\bigr)$
  $\mapsto$ the goal
  $\mathrm{ofReal}\ p \le$ \texttt{pathMeasure K x \{ω | ∃ k ≤ T, sat
  (ω k)\}} from every start with \texttt{x ∈ D} and \texttt{V x ≤ (1
  − p)·q}; the target tie $\mapsto$ \texttt{hF} $:$ $\forall y \in B,
  \texttt{sat}\ y$ --- $B$ \emph{under}-approximates the target truth
  set (the polarity opposite to \rulename{QSafetyBattery}\textquotesingle
  s avoid tie, both safe directions), and again no measurability is
  asked of \texttt{sat}.
\item decay rows $\mapsto$ \texttt{hdecay} (one-step $\int V \le
  \beta V$ on $D \smallsetminus B$; at this profile guard rows land
  inside it, per the note) with integrability rider \texttt{hint};
  floors \emph{off-target only} $\mapsto$ \texttt{hfloor} (on
  $\bar{D} \smallsetminus B$ --- the display\textquotesingle s
  weaker, deliberately relaxable form is exactly what is bound);
  nonnegs $\mapsto$ \texttt{hnonneg}; the bound row $V \le M$
  $\mapsto$ \texttt{hM}.
\item the TaylorBound leaf $q \le \sum_{k \le n} (\theta T)^k / k!$
  $\mapsto$ \texttt{hq}, verbatim --- the one closed comparison a
  checker replays in exact $\Q$ --- and the drift tie $\beta \le
  e^{-\theta}$ $\mapsto$ \texttt{hβθ}, with both transcendental
  facts (\texttt{Real.sum\_le\_exp\_of\_nonneg}; $e^{-\theta}
  e^{\theta} = 1$) discharged \emph{inside} the lemma: the certified
  statement is the witness that the checker never needs to know $e$
  exists.
\item sign gates $0 \le \beta$, $\theta \ge 0$, $0 \le q$, $p \le 1$
  $\mapsto$ \texttt{hβ}/\texttt{hθ}/\texttt{hq0}/\texttt{hp}, verbatim
  --- the base sign gate \texttt{hβ} is consumed by the power-order
  step; side conditions and region
  measurability as at
  \rulename{QSafetyBattery} (\texttt{hR3}/\texttt{hR4};
  \texttt{hD}/\texttt{hDbar}/\texttt{hB}/\texttt{hV}, the
  undisplayed regularity pack).
\end{itemize}
Width note: the wrapped \texttt{sound\_QDeadline} consumes only
$\beta^T \cdot \mathrm{thr} \le 1 - p$ --- every $(\theta, q)$
satisfying the displayed leaves yields it, so the catalogue
parameterization is an instance of a strictly wider certified core;
producer knobs, exactly as the roster note says.

\ruleblock{Lean verifier receipts}
\sourcenote{module \texttt{BernsteinLogicAux.Battery}; bound build \texttt{2fd1c031c80b\ldots} (bound from receipt 8d9ccb68993ea0ec, 2026-08-04T01:56:55Z). A receipt records that named checking layers accepted a build: (module, compiled-artifact SHA-256, environment SHA-256, checker layers), minted when the checking gate ran; the store receipts attesting the bound build are listed with the layers each run credited. Nothing here asks the reader to trust the store: every layer is an independently runnable program over the quoted sources --- rebuild the module and re-run the checkers to reproduce the verdict. Receipts attesting superseded builds credit no layer and are omitted from this rendering (the count of omissions, if any, is noted below).}
\begin{receipttable}
\receiptrow{8d9ccb68993ea0ec}{2026-08-04T01:56:55Z}{build+decorrelated}{\badge{BadgeGreen}{current}}
\receiptrow{5075586e4813a7c5}{2026-08-04T03:11:30Z}{build+decorrelated}{\badge{BadgeGreen}{current}}
\end{receipttable}
\sourcenote{13 further receipt(s) in the store attest superseded builds of this module and are omitted from this rendering.}
\sourcenote{module \texttt{BernsteinLogicAux.QtslBattery}; bound build \texttt{5c75f9923f27\ldots} (bound from receipt 3fd830b4914bdeda, 2026-08-04T01:56:58Z). A receipt records that named checking layers accepted a build: (module, compiled-artifact SHA-256, environment SHA-256, checker layers), minted when the checking gate ran; the store receipts attesting the bound build are listed with the layers each run credited. Nothing here asks the reader to trust the store: every layer is an independently runnable program over the quoted sources --- rebuild the module and re-run the checkers to reproduce the verdict. Receipts attesting superseded builds credit no layer and are omitted from this rendering (the count of omissions, if any, is noted below).}
\begin{receipttable}
\receiptrow{3fd830b4914bdeda}{2026-08-04T01:56:58Z}{build+decorrelated}{\badge{BadgeGreen}{current}}
\receiptrow{96b18d8572451f65}{2026-08-04T03:11:31Z}{build+decorrelated}{\badge{BadgeGreen}{current}}
\end{receipttable}
\sourcenote{13 further receipt(s) in the store attest superseded builds of this module and are omitted from this rendering.}
\sourcenote{module \texttt{BernsteinLogicAux.LayerD.Deadline}; bound build \texttt{ae6767f0b59d\ldots} (bound from receipt 1c87bcb62534ed03, 2026-08-04T01:57:04Z). A receipt records that named checking layers accepted a build: (module, compiled-artifact SHA-256, environment SHA-256, checker layers), minted when the checking gate ran; the store receipts attesting the bound build are listed with the layers each run credited. Nothing here asks the reader to trust the store: every layer is an independently runnable program over the quoted sources --- rebuild the module and re-run the checkers to reproduce the verdict. Receipts attesting superseded builds credit no layer and are omitted from this rendering (the count of omissions, if any, is noted below).}
\begin{receipttable}
\receiptrow{1c87bcb62534ed03}{2026-08-04T01:57:04Z}{build+decorrelated}{\badge{BadgeGreen}{current}}
\receiptrow{ea6411e23e04ba3c}{2026-08-04T03:11:30Z}{build+decorrelated}{\badge{BadgeGreen}{current}}
\end{receipttable}
\sourcenote{14 further receipt(s) in the store attest superseded builds of this module and are omitted from this rendering.}
\sourcenote{module \texttt{QtslAux.LayerD.MPDischarge}; bound build \texttt{baded74dd943\ldots} (bound from receipt 7228834dea10e5f9, 2026-07-24T11:04:04Z). A receipt records that named checking layers accepted a build: (module, compiled-artifact SHA-256, environment SHA-256, checker layers), minted when the checking gate ran; the store receipts attesting the bound build are listed with the layers each run credited. Nothing here asks the reader to trust the store: every layer is an independently runnable program over the quoted sources --- rebuild the module and re-run the checkers to reproduce the verdict. Receipts attesting superseded builds credit no layer and are omitted from this rendering (the count of omissions, if any, is noted below).}
\begin{receipttable}
\receiptrow{7228834dea10e5f9}{2026-07-24T11:04:04Z}{build+decorrelated}{\badge{BadgeGreen}{current}}
\end{receipttable}
\sourcenote{module \texttt{SuFunctor.Site.FellAcceptance}; bound build \texttt{651fb557bbfe\ldots} (bound from receipt aa2169560f8f571a, 2026-08-03T20:19:04Z). A receipt records that named checking layers accepted a build: (module, compiled-artifact SHA-256, environment SHA-256, checker layers), minted when the checking gate ran; the store receipts attesting the bound build are listed with the layers each run credited. Nothing here asks the reader to trust the store: every layer is an independently runnable program over the quoted sources --- rebuild the module and re-run the checkers to reproduce the verdict. Receipts attesting superseded builds credit no layer and are omitted from this rendering (the count of omissions, if any, is noted below).}
\begin{receipttable}
\receiptrow{7ad85720de6e419e}{2026-08-03T10:13:50Z}{build+decorrelated}{\badge{BadgeGreen}{current}}
\receiptrow{94a408109996d1bf}{2026-08-03T13:41:21Z}{build+decorrelated}{\badge{BadgeGreen}{current}}
\receiptrow{580a7e551fab03b6}{2026-08-03T15:39:44Z}{build+decorrelated}{\badge{BadgeGreen}{current}}
\receiptrow{aa2169560f8f571a}{2026-08-03T20:19:04Z}{build+decorrelated}{\badge{BadgeGreen}{current}}
\end{receipttable}
\sourcenote{2 further receipt(s) in the store attest superseded builds of this module and are omitted from this rendering.}
\sourcenote{module \texttt{SuFunctor.Site.FellBehaviorSpec}; bound build \texttt{345f9caf9453\ldots} (bound from receipt 72f2c3b62b6771fc, 2026-08-03T20:19:08Z). A receipt records that named checking layers accepted a build: (module, compiled-artifact SHA-256, environment SHA-256, checker layers), minted when the checking gate ran; the store receipts attesting the bound build are listed with the layers each run credited. Nothing here asks the reader to trust the store: every layer is an independently runnable program over the quoted sources --- rebuild the module and re-run the checkers to reproduce the verdict. Receipts attesting superseded builds credit no layer and are omitted from this rendering (the count of omissions, if any, is noted below).}
\begin{receipttable}
\receiptrow{d7b846e785c1d7fb}{2026-08-03T10:13:53Z}{build+decorrelated}{\badge{BadgeGreen}{current}}
\receiptrow{c0e3dc15419765aa}{2026-08-03T13:41:24Z}{build+decorrelated}{\badge{BadgeGreen}{current}}
\receiptrow{9a62b29c08b63b4f}{2026-08-03T15:39:47Z}{build+decorrelated}{\badge{BadgeGreen}{current}}
\receiptrow{72f2c3b62b6771fc}{2026-08-03T20:19:08Z}{build+decorrelated}{\badge{BadgeGreen}{current}}
\end{receipttable}
\sourcenote{2 further receipt(s) in the store attest superseded builds of this module and are omitted from this rendering.}
\layerline{\layeritem{build}{lake build (elaboration)}{\badge{BadgeGreen}{current}}\layeritem{replay}{leanchecker (official kernel replay)}{\badge{BadgeAmber}{stale only}}\layeritem{hardened}{safe\_verify (hardened screen)}{\badge{BadgeAmber}{stale only}}\layeritem{decorrelated}{lean4export + nanoda (decorrelated kernel)}{\badge{BadgeGreen}{current}}}
\sourcenote{\badge{BadgeGray}{comparator: none} no fixture in the three-implementation comparator harness (compiled / kernel-interpreted / reference-interpreted) exercises this rule.}
\ruleblock{HOL4 statement}
\absentbox{NO HOL STATEMENT YET}{The HOL4 development ports the base calculus\textquotesingle s checker
chain (theory \texttt{bern\_catalogue}, developed side by side with the
base catalogue sources, stopping below the proof-object layer); it
contains no probability or measure theory.
\rulename{QDeadlineBattery} lives one theory-extension edge above the
base calculus and its soundness content is path-space measure theory
(discounted supermartingales, optional stopping, an escape-null
argument), for which no counterpart exists anywhere in the ported
development. A HOL statement of this rule would be a new
measure-theoretic campaign, not an extension of the ported checker
chain; until one exists, this box is the honest state of the edifice.
}
\ruleblock{HOL verifier receipts}
\absentbox{NONE}{no HOL verification receipt exists for this rule.}
\ruleblock{Proof}

\sourcenote{Genre: core rule (tier 1) --- an ordinary mathematical proof
at Putnam-exam-solution detail. Substrate, stopped-process machinery,
and the battery pattern are stated with \rulename{QSafetyBattery}
(Lemmas \ref{lem:qsb-stopped}--\ref{lem:qsb-ville},
Proposition~\ref{prop:qsb-pattern}) and cited here by reference. The
theorems quoted above warrant the rule at the \textbf{discrete-kernel
profile}; the profile remark at the end states what remains a
hypothesis beyond it.}

\textbf{Setting and claim.}\enspace The substrate of
\rulename{QSafetyBattery}: state space $S$, one-step Markov kernel $K$,
path measure $\Pr_x$, canonical filtration. The rule\textquotesingle s
data: modelled domain $D$, hull $\bar{D} \supseteq D$, a \emph{target}
$B$ compiled to \emph{under}-approximate the truth set of the
$\mathsf{F}_{\le T}$-argument --- note the orientation flip against the
safety battery\textquotesingle s avoid set: the formula tie here is
(F$'$)\ $\forall y \in B,\ \mathrm{sat}(y)$ --- a measurable certificate
$V$, a discrete decay factor $\beta \ge 0$, a horizon $T \in \N$, and
$p, \mathrm{thr}$. The checked rows, read semantically:
\begin{itemize}
\item[(D1)] \emph{decay}: for $y \in D \smallsetminus B$, $V$ integrable
  against $K\,y$ and $\int V \,dK\,y \le \beta\, V(y)$ --- the discrete
  reading of $\mathcal{A} V \le -\theta V$ with $\beta = e^{-\theta}$.
  (At this profile guard rows land inside this one row; see
  \autoref{rem:qdb-guards}.)
\item[(D2)] \emph{floor, off-target}: $V \ge 1$ on
  $\bar{D} \smallsetminus B$.
\item[(D3)] \emph{nonneg}: $V \ge 0$ on $\bar{D}$; \emph{bound}:
  $V \le M$ on $\bar{D}$.
\item[(R3)/(R4)] hull and fail-closed containment at target $B$: from
  $y \in D \smallsetminus B$, one-step mass stays in $\bar{D}$ and in
  $D \cup B$.
\item[($\Theta_{\mathsf F}$)] \emph{deadline threshold}:
  $\beta^{T} \cdot \mathrm{thr} \le 1 - p$.
\end{itemize}
The claim (\texttt{sound\_QDeadline}): for every $x \in D$ with
$V(x) \le \mathrm{thr}$,
\[
\Pr_x^{\,*}\bigl[\exists k \le T,\ \mathrm{sat}(\omega_k)\bigr] \;\ge\; p
\]
--- outer measure, as everywhere in the family: no measurability is
asked of the satisfaction predicate, and the mark is inert whenever its
truth set is measurable.

\begin{blemma}[discounted tail bound at the enlarged target]
\label{lem:qdb-tail}
Let $B' := B \cup D^{\mathrm{c}}$. Under (D1)--(D3), (R3), (R4): for
every $x \in D \smallsetminus B$,
\[
\Pr_x\bigl[\forall k \le T,\ \omega_k \notin B'\bigr]
\;\le\; \beta^{T}\, V(x).
\]
\end{blemma}
\begin{proof}
Clamp as in the safety battery, but with the floor forced on
$D \smallsetminus B$ rather than off $D$: let
$C := \max(M, 1)$ and let $g$ clamp $\max(V,1)$ into $[0,C]$ on
$D \smallsetminus B$ and $\max(V,0)$ into $[0,C]$ elsewhere. As in
Lemma~\ref{lem:qsb-clamp}: $0 \le g \le C$; $g \ge 1$ \emph{off} $B'$
(here the floor (D2) does the work on
$\bar{D} \smallsetminus B$, and the forced branch covers the rest); on
$\bar{D} \cap (D \cup B)$ --- which carries all one-step mass from
$D \smallsetminus B$, by (R3)$+$(R4) --- the clamp is invisible,
$g = \min(V, C)$. And by the bound row (D3), $V \le M \le C$ on the
hull, so the ceiling never truncates there: $g = V$ on the whole
mass-carrying set and at $y$ itself. The decay row therefore transfers,
$\int g \,dK\,y = \int V \,dK\,y \le \beta\, V(y) = \beta\, g(y)$ for
$y \in D \smallsetminus B$ --- and this is where the bound row is
load-bearing, not an integrability convenience: without it the transfer
would need $\min(\beta V(y), C) \le \beta \min(V(y), C)$, false for
$0 < \beta < 1$ wherever $V > C$ (at $\beta = 0$ both sides vanish;
that branch is handled separately below). Likewise $g(x) = V(x)$. Assume first
$\beta > 0$. Now the
discounted process
$M_n := \beta^{-(n \wedge \tau_{B'})}\, g(\omega_{n \wedge \tau_{B'}})$
is a nonnegative supermartingale (the decay row supplies exactly the
extra factor $\beta$ the discount consumes at each surviving step), and
on the miss event $\{\forall k \le T,\ \omega_k \notin B'\}$ we have
$\tau_{B'} > T$ and $g(\omega_T) \ge 1$, so $M_T \ge \beta^{-T}$.
Markov\textquotesingle s inequality at time $T$ gives
$\Pr[\text{miss}] \le \beta^{T}\, \mathbb{E} M_0
= \beta^{T} g(x) = \beta^{T} V(x)$. For $\beta = 0$ the inverse
discount is unavailable and the bound holds directly: at $T = 0$ the
right side is $V(x) = g(x) \ge 1 \ge \Pr[\text{miss}]$ (the floor off
$B'$), while for $T \ge 1$ the decay row at $\beta = 0$ forces
$\int g \,dK_y = 0$ for every $y$ off $B'$, so --- $g \ge 1$ off $B'$
--- the kernel from any off-$B'$ point charges $B'$ with full mass and
the miss event is null, by the one-step marginal of the trajectory
law --- the same prefix--step disintegration
Lemma~\ref{lem:qdb-noescape} rides. The mechanized route packages this
computation as the geometric-drift engine (the discounted
supermartingale and its hitting-time moment bound), consumed at the
enlarged target $B'$.
\end{proof}

\begin{blemma}[no escape]\label{lem:qdb-noescape}
Under (R4) at target $B$: for $x \in D$,
\[
\Pr_x\Bigl[\,\exists n,\ \bigl(\forall k \le n,\ \omega_k \notin B\bigr)
\wedge \omega_n \notin D \Bigr] \;=\; 0 .
\]
\end{blemma}
\begin{proof}
Write $E_n$ for the event that some escape as displayed occurs by time
$n$. Induction on $n$ gives $\Pr_x[E_n] = 0$ for every $n$. Base:
coordinate $0$ is deterministically $x \in D$. Step: an escape at time
$n+1$ that has not yet visited $B$ either already escaped by time $n$
(null by the inductive hypothesis) or steps from a state in
$D \smallsetminus B$ to a state outside $D \cup B$ --- and that step is
null, because the joint law of the $n$-prefix and the next step
disintegrates through the one-step kernel, at which (R4) says the
section $(D \cup B)^{\mathrm{c}}$ has measure zero from every surviving
prefix endpoint. The displayed event is $\bigcup_{n} E_n$, null as a
countable union of null events. This lemma is the almost-everywhere
invariance content of (R4) that a \emph{reach} conclusion cannot avoid:
the safety battery got fail-closure for free by enlarging its stop set
--- leaving counts as failing is conservative there --- but a reach
guarantee cannot count escape as success, so the escape event must be
shown null outright.
\end{proof}

\begin{bthm}[target-hit form]\label{thm:qdb-hit}
Under (D1)--(D3), (R3), (R4), ($\Theta_{\mathsf F}$): for every
$x \in D$ with $V(x) \le \mathrm{thr}$,
$\Pr_x[\exists k \le T,\ \omega_k \in B] \ge p$.
\end{bthm}
\begin{proof}
If $x \in B$, the deadline is met at time $0$ almost surely (coordinate
$0$ is deterministic) and $p \le 1$ suffices --- the degenerate corner,
handled first so the clamp may assume $x \notin B$. Otherwise decompose
the miss event: $\{$miss $B$ through $T\} \subseteq \{$miss $B'$ through
$T\} \cup \{$escape before visiting $B\}$. The second is null
(Lemma~\ref{lem:qdb-noescape}); the first has probability
$\le \beta^{T} V(x) \le \beta^{T}\,\mathrm{thr} \le 1 - p$
(Lemma~\ref{lem:qdb-tail} and ($\Theta_{\mathsf F}$)). Complement.
\end{proof}

\begin{bprop}[soundness of \rulename{QDeadlineBattery}]
\label{prop:qdb-sound}
With the formula tie (F$'$): $\Pr_x^{\,*}[\exists k \le T,\
\mathrm{sat}(\omega_k)] \ge p$ for every $x \in D$ with
$V(x) \le \mathrm{thr}$ --- the concluded judgment read on the sublevel
region.
\end{bprop}
\begin{proof}
$B$ under-approximates the truth set, so the hit event of
Theorem~\ref{thm:qdb-hit} is contained in the reach-formula event;
monotonicity of the (outer) measure. As always, no measurability is
asked of the satisfaction predicate.
\end{proof}

\begin{bprop}[catalogue form: the $(\theta, q)$ parameterization]
\label{prop:qdb-catalogue}
Suppose additionally $\beta \le e^{-\theta}$ with $\theta \ge 0$ (the
drift tie), $0 \le q \le \sum_{k \le n} (\theta T)^k / k!$ for some $n$
(the Taylor leaf --- the one closed, wire-free comparison the checker
replays in exact rational arithmetic, discharged as a ground-fact-rule
application), and $p \le 1$. Then the conclusion of
Proposition~\ref{prop:qdb-sound} holds on the region
$\{V \le (1-p)\,q\}$.
\end{bprop}
\begin{proof}
$q \le \sum_{k \le n} (\theta T)^k/k! \le e^{\theta T}$ --- the first
transcendental fact, living here and not in the checker, applied at the
nonnegative exponent $\theta T$ --- and $\beta^{T} \le e^{-\theta T}$,
powers preserving order between nonnegative bases ($0 \le \beta \le
e^{-\theta}$, $T \in \N$); so
$\beta^{T} q \le e^{-\theta T} e^{\theta T} = 1$ --- the second. Hence
$\beta^{T} \cdot (1-p)\,q \le 1 - p$: hypothesis
($\Theta_{\mathsf F}$) holds at $\mathrm{thr} = (1-p)\,q$, and
Proposition~\ref{prop:qdb-sound} applies. The checker never needs to
know $e$ exists; soundness is indifferent to the producer knobs
$\theta$ and $q$ --- any smaller rate or budget is admissible.
\end{proof}

\begin{bremark}[why the off-target floor is not decoration]
\label{rem:qdb-floor}
Drop (D2) and take $V \equiv 0$ with the identity kernel
$K\,y = \delta_y$ at some $x \in D \smallsetminus B$: decay holds
($0 \le \beta \cdot 0$), nonneg and bound hold, and
($\Theta_{\mathsf F}$) permits any $p \le 1$ since
$\beta^T \cdot 0 = 0$. The rule would then certify that a chain frozen
forever outside the target reaches it with probability one. With the
floor present the certificate dies: $V \ge 1$ on
$\bar{D} \smallsetminus B \ni x$ contradicts $V(x) = 0$. Missing the
deadline must cost at least one unit of discounted certificate value ---
that is precisely the step $M_T \ge \beta^{-T}$ in
Lemma~\ref{lem:qdb-tail}. And the floor is asked \emph{off the target
only}: on $B$ the certificate may take any value, because success needs
no budget. This checked form is deliberately weaker than a global
floor, and the proven statement binds the weaker form.
\end{bremark}

\begin{bremark}[the two intrinsic asymmetries against the safety battery]
\label{rem:qdb-asym}
Both are consequences of concluding \emph{reach} rather than
\emph{avoid}, and neither is an artefact of the mechanization. First,
the boundedness row $V \le M$: the safety battery's \emph{additive}
drift survives its ceiling for free ---
$\int \min(V, C)\,dK\,y \le \min(\int V\,dK\,y,\, C) \le
\min(V(y), C)$ --- so a ceiling at
$\max(\mathrm{thr}, 1)$ cost nothing; here the discount multiplies the
certificate by $\beta^{-n}$, and the transfer of the decay row to the
clamped certificate needs the ceiling to be invisible wherever one-step
mass reaches --- exactly (D3)'s bound, consumed inside
Lemma~\ref{lem:qdb-tail} (and in the mechanized geometric-drift engine
it additionally supplies integrability of the discounted certificate)
--- for polynomial certificates over a compact hull it is free. Second, escape: the safety battery counts
leaving the modelled domain as failure by enlarging its stop set, which
is conservative for an avoid conclusion but would be unsound generosity
for a reach conclusion --- hence Lemma~\ref{lem:qdb-noescape}, whose
almost-everywhere argument is the one place this
section\textquotesingle s hypotheses stop being pointwise.
\end{bremark}

\begin{bremark}[guards at this profile]\label{rem:qdb-guards}
The catalogued rule lists guard rows (jumps instantaneous). At the
discrete-kernel profile the one-step kernel carries flow and enabled
jumps together, so the guard rows land inside the single decay row (D1)
--- exactly as the safety battery\textquotesingle s interior and guard
rows land together. The instantaneous-jump reading --- the discount
factor does not tick during a jump, so the $e^{\theta t}$ factor cancels
across it --- belongs to the continuous-time profile, where it is
carried by the martingale-problem discharge of
\autoref{rem:qdb-profile}.
\end{bremark}

\begin{bremark}[\textbf{profile} --- what is proven, at which
presentations]
\label{rem:qdb-profile}
Proven: the discrete-kernel profile, sorry-free
(\texttt{sound\_QDeadline\_hit}, \texttt{sound\_QDeadline},
\texttt{sound\_QDeadlineBattery}). Also proven: the continuous-time
profile at the martingale-problem presentation
(\texttt{sound\_GSHS\_deadline\_mp},
\texttt{sound\_GSHS\_deadline\_battery\_mp}, module
\texttt{QtslAux.LayerD.MPDischarge} of the semantic-universe
strand\textquotesingle s auxiliary package) --- and the three
ingredients this remark once filed are exactly what that discharge
supplies: the $e^{\theta t} V$ supermartingale over continuous time
is the exponential-discount construction from the undiscounted
martingale-problem hypothesis; optional stopping for right-continuous
nonnegative supermartingales is the library theorem, consumed at the
d\'ebut; and the instantaneous-jump reading rides the
presentation\textquotesingle s no-alarm segment assembly. The
martingale-problem theorems are hash-bound by this
catalogue\textquotesingle s binder via the qtsl-aux package
registration, as for the safety sibling, so the dashboard row carries
no conditional mark --- each named profile is proven, and its
theorems are bound.
\end{bremark}

\begin{bremark}[relation to the literature]\label{rem:qdb-lit}
The engine is the geometric drift condition of the Foster--Lyapunov
family \cite{meyntweedie2009}: expected decrease by a factor per step
off the target, yielding exponential hitting-time tails. The discounted
supermartingale and its optional-stopping consumption are classical
\cite{doob1953}; the barrier mechanics are the same
Ville-type argument \cite{ville1939} the whole family rides. The
certificate-form packaging --- exact rational rows, transcendental facts
quarantined inside a verified lemma --- is this
calculus\textquotesingle s contribution rather than the
literature\textquotesingle s.
\end{bremark}

\begin{anchors}
Lemma \ref{lem:qdb-tail}: discounted engine, clamp for the deadline reading & \texttt{BernsteinLogic.Deadline.deadline\_tail} (via \texttt{supermartingale\_inv\_pow\_mul}, \texttt{expMoment\_hitting}); the clamp $g$ inside \texttt{sound\_QDeadline\_hit}\\
Lemma \ref{lem:qdb-noescape}: escape-null induction & \texttt{ae\_no\_escape} (file-local; joint prefix--step law via \texttt{Kernel.partialTraj\_compProd\_eq\_map\_traj}); base case \texttt{ae\_coord\_zero\_mem}\\
Theorem \ref{thm:qdb-hit}: assembly, degenerate corner & \texttt{sound\_QDeadline\_hit}\\
Proposition \ref{prop:qdb-sound}: formula tie & \texttt{sound\_QDeadline}\\
Proposition \ref{prop:qdb-catalogue}: Taylor leaf, drift tie, transcendental quarantine & \texttt{sound\_QDeadlineBattery} (via \texttt{Real.sum\_le\_exp\_of\_nonneg})\\
\end{anchors}


\rulesection{42}{QUnionBound}{core rule}
\label{blcatrule-QUnionBound}
\ruleblock{The rule (Steele metanotation)}
\sourcenote{as realized in the mechanized development; source row: \texttt{\texttt{QTSL-rules.toml}} --- \emph{QUnionBound}. The metanotation statement is typeset by hand for notational quality; its mechanized ground truth is the Lean statement below, extracted from the sources at every build of this book and hash-bound. The rule \emph{roster} is machine-checked (every mechanized rule tag round-trips to exactly one catalogue row); the statement \emph{text} is editorial, checked by eye against the extracted Lean beside it.}
\begin{mathpar}
\inferrule*{
  \mu(E_1) \le a \\ \mu(E_2) \le b \\ \mathrm{fail} \subseteq E_1 \cup E_2
}{
  \mu(\mathrm{fail}) \le a + b
}
\end{mathpar}

\ruleflavor{rule sort baked in \texttt{theory Qtsl} --- realized measure-generically (arbitrary measure; no profile residue). Packaged instance (the d\'ebut decomposition, pure path-set algebra over sequence space): failure of $(\varphi_s\,\mathsf{U}_{\le T}\,\varphi)$, at stop set $\mathit{Stop}$ and good-exit set $R$, is contained in $\{\text{some \emph{first} } \mathit{Stop}\text{-visit presents a bad state}\} \cup \{\text{no } \mathit{Stop}\text{-visit through } T\}$ --- the containment the until battery cites as its $\mathrm{fail} \subseteq E_1 \cup E_2$ premise.}
\par\medskip
\rulename{QUnionBound} is the battery family\textquotesingle s glue:
from two probability bounds and an event decomposition it concludes that
the sum of the budgets bounds the failure event ---
$\Pr[E_1] \le a$, $\Pr[E_2] \le b$, and
$\mathrm{fail} \subseteq E_1 \cup E_2$ give
$\Pr[\mathrm{fail}] \le a + b$. Its semantic content is monotonicity
plus subadditivity (Boole\textquotesingle s inequality), stated over an
\emph{arbitrary} measure: the rule is profile-generic, and its instance
at any system-semantics profile is the same lemma applied at that
profile\textquotesingle s path measure. Its dashboard row accordingly
carries an unconditional check with no profile discharge even needed
--- there was never a profile residue to discharge (the battery
siblings\textquotesingle{} rows earned theirs through the
martingale-problem discharge; this one is measure-generic from
birth).

The hypothesis $\mathrm{fail} \subseteq E_1 \cup E_2$ is where all the
content concentrates. The rule is sound for \emph{any} decomposition;
the citing rule must justify the containment. The packaged instance ---
the one the bounded-until battery cites --- is the d\'ebut
decomposition: failure of ``stay safe until the target, by the
deadline'' is contained in $\{$some \emph{first} stop-set visit presents
a bad state$\}$ $\cup$ $\{$no stop-set visit by the horizon$\}$. That
containment is pure path-set algebra over sequence space: no measure, no
measurability, only the well-ordering of the time axis.

A producer rarely writes this node by hand; it appears where a composite
temporal property splits its failure modes across two certificates, and
what the checker recomputes at such a node is precisely the containment
side of the decomposition (or it discharges it by citing the packaged
path-set lemma). The same additivity doubles as the slack-composition
device in assume--guarantee composition, where interface slack budgets
add across a wiring --- one lemma, two consumers.

Why core rather than admissible: there is nothing to unpack into.
Subadditivity of measure is not derivable from the other battery rules;
it is new --- deliberately small --- semantic content at the measure
level, and keeping it a separate node keeps every consumer\textquotesingle s
soundness proof honest about where its additivity comes from.

\cataloguenote{Stated in two forms: an extended-nonneg form (the primitive), and a
real-valued form in the shape battery bounds actually arrive
(certificate initial values), whose nonnegativity hypotheses are free at
every use site.
}
\ruleblock{Lean statement}
\declhead{BernsteinLogicAux.QtslBattery}{sound\_QUnionBound}{\badge{BadgeGreen}{statement bound}}
\begin{leansrc}
/-- **`sound_QUnionBound`** (ℝ≥0∞ form). Subadditivity through an event
decomposition — the entire semantic content of the `QUnionBound` rule
(QTSL-rules.toml row `QUnionBound`). -/
theorem sound_QUnionBound {Ω : Type*} [MeasurableSpace Ω] (μ : Measure Ω)
    {E₁ E₂ F : Set Ω} {a b : ℝ≥0∞}
    (h₁ : μ E₁ ≤ a) (h₂ : μ E₂ ≤ b) (hsub : F ⊆ E₁ ∪ E₂) :
    μ F ≤ a + b
\end{leansrc}
\declhead{BernsteinLogicAux.QtslBattery}{sound\_QUnionBound\_ofReal}{\badge{BadgeGreen}{statement bound}}
\begin{leansrc}
/-- **`sound_QUnionBound`** (`ENNReal.ofReal` form), in the shape the battery
bounds arrive: real-valued certificate bounds `a`, `b` combine to `a + b`.
The nonnegativity hypotheses are free at every use site (certificate initial
values, nonneg rows). -/
theorem sound_QUnionBound_ofReal {Ω : Type*} [MeasurableSpace Ω] (μ : Measure Ω)
    {E₁ E₂ F : Set Ω} {a b : ℝ} (ha : 0 ≤ a) (hb : 0 ≤ b)
    (h₁ : μ E₁ ≤ ENNReal.ofReal a) (h₂ : μ E₂ ≤ ENNReal.ofReal b)
    (hsub : F ⊆ E₁ ∪ E₂) :
    μ F ≤ ENNReal.ofReal (a + b)
\end{leansrc}
\declhead{BernsteinLogicAux.QtslBattery}{compl\_firstExit\_subset}{\badge{BadgeGreen}{statement bound}}
\begin{leansrc}
/-- **The first-exit event decomposition** (the `QUnionBound` premise's
début decomposition): over
sequence space, the complement of the first-exit-success event — some time
`n ≤ T` whose state is in the good exit set `Good` with no earlier visit to
the stop set `Stop` — is contained in the union of {some `Stop`-visit is a
*first* visit presenting a bad state (`Stop \ Good`)} and {no `Stop`-visit
through `T`}. Pure path-set algebra, no measure and no measurability; the
until battery instantiates `Stop := (A ∪ R) ∪ Dᶜ`, `Good := R`. This is the
lemma the `QUnionBound` premise `fail ⊆ E₁ ∪ E₂` cites for T8. -/
theorem compl_firstExit_subset {S : Type*} (Stop Good : Set S) (T : ℕ) :
    {ω : ℕ → S | ∃ n ≤ T, ω n ∈ Good ∧ ∀ j < n, ω j ∉ Stop}ᶜ ⊆
      {ω : ℕ → S | ∃ n, ω n ∈ Stop \ Good ∧ ∀ j < n, ω j ∉ Stop} ∪
        {ω : ℕ → S | ∀ k ≤ T, ω k ∉ Stop}
\end{leansrc}
\ruleblock{Correspondence (Steele display $\leftrightarrow$ Lean)}
\badge{BadgeAmber}{a bound Lean statement CHANGED since certification --- re-review and re-certify}
\sourcenote{The dictionary below is a reviewer-certified statement of how the mechanized statements above interpret the displayed rule: display metavariable $\mapsto$ Lean binder, each display hypothesis $\mapsto$ the premise or side condition carrying it, the display conclusion $\mapsto$ the goal --- with any deliberate difference in width between display and mechanization named outright. Certification is recorded as a hash triple over (the displayed cell, the extracted statement texts, this dictionary\textquotesingle s own text); if any of the three changes, this block turns amber until a reviewer re-certifies. The correspondence claim itself remains editorial --- reviewer-grade, not machine-checked; its machine-checked successor types each soundness statement against the rule as a structured object, and supersedes this block per rule as it lands.}
The cleanest dictionary in the Q-family --- the certified statement is
the displayed rule at full measure-generality, with no profile
residue.
\begin{itemize}
\item $\mu \mapsto$ an arbitrary \texttt{Measure μ} on an arbitrary
  measurable space --- the amended display\textquotesingle s own
  generality (\emph{wider} than the original display\textquotesingle s
  path-probability reading, in the safe direction): the rule\textquotesingle
  s GSHS instance and the AG slack-composition use are the same lemma
  at other measures (the roster\textquotesingle s ``one item, two
  consumers\textquotesingle\textquotesingle).
\item the three premises $\mu(E_1) \le a$, $\mu(E_2) \le b$,
  $\mathrm{fail} \subseteq E_1 \cup E_2$ $\mapsto$ \texttt{h₁},
  \texttt{h₂}, \texttt{hsub}, verbatim; conclusion $\mu(\mathrm{fail})
  \le a + b$ $\mapsto$ the goal, verbatim. No hypothesis is unused
  and nothing undisplayed is demanded --- not even measurability of
  the events (monotonicity and subadditivity hold for outer measure).
\item the \texttt{ofReal} companion carries the same rule in the
  shape battery bounds arrive (real-valued $a$, $b$), its two
  nonnegativity riders free at every use site, per the note.
\item the packaged instance --- the flavor text\textquotesingle s
  d\'ebut decomposition
  --- $\mapsto$ \texttt{compl\_firstExit\_subset}, certified
  alongside: pure path-set algebra over sequence space (no measure,
  no measurability), delivering exactly the $\mathrm{fail} \subseteq
  E_1 \cup E_2$ premise the until battery cites at $\mathit{Stop} :=
  (A \cup R) \cup D^{\mathrm{c}}$, $\mathit{Good} := R$.
\end{itemize}

\ruleblock{Lean verifier receipts}
\sourcenote{module \texttt{BernsteinLogicAux.QtslBattery}; bound build \texttt{5c75f9923f27\ldots} (bound from receipt 3fd830b4914bdeda, 2026-08-04T01:56:58Z). A receipt records that named checking layers accepted a build: (module, compiled-artifact SHA-256, environment SHA-256, checker layers), minted when the checking gate ran; the store receipts attesting the bound build are listed with the layers each run credited. Nothing here asks the reader to trust the store: every layer is an independently runnable program over the quoted sources --- rebuild the module and re-run the checkers to reproduce the verdict. Receipts attesting superseded builds credit no layer and are omitted from this rendering (the count of omissions, if any, is noted below).}
\begin{receipttable}
\receiptrow{3fd830b4914bdeda}{2026-08-04T01:56:58Z}{build+decorrelated}{\badge{BadgeGreen}{current}}
\receiptrow{96b18d8572451f65}{2026-08-04T03:11:31Z}{build+decorrelated}{\badge{BadgeGreen}{current}}
\end{receipttable}
\sourcenote{13 further receipt(s) in the store attest superseded builds of this module and are omitted from this rendering.}
\layerline{\layeritem{build}{lake build (elaboration)}{\badge{BadgeGreen}{current}}\layeritem{replay}{leanchecker (official kernel replay)}{\badge{BadgeAmber}{stale only}}\layeritem{hardened}{safe\_verify (hardened screen)}{\badge{BadgeAmber}{stale only}}\layeritem{decorrelated}{lean4export + nanoda (decorrelated kernel)}{\badge{BadgeGreen}{current}}}
\sourcenote{\badge{BadgeGray}{comparator: none} no fixture in the three-implementation comparator harness (compiled / kernel-interpreted / reference-interpreted) exercises this rule.}
\ruleblock{HOL4 statement}
\absentbox{NO HOL STATEMENT YET}{The HOL4 development ports the base calculus\textquotesingle s checker
chain (theory \texttt{bern\_catalogue}, developed side by side with the
base catalogue sources, stopping below the proof-object layer); it
contains no probability or measure theory. \rulename{QUnionBound} lives
one theory-extension edge above the base calculus, and even its small
soundness content (measure subadditivity, a path-set decomposition over
sequence space) has no counterpart anywhere in the ported development. A
HOL statement of this rule would open a measure-theory development, not
extend the ported checker chain; until one exists, this box is the
honest state of the edifice.
}
\ruleblock{HOL verifier receipts}
\absentbox{NONE}{no HOL verification receipt exists for this rule.}
\ruleblock{Proof}

\sourcenote{Genre: core rule (tier 1) --- the battery
family\textquotesingle s smallest member, and its only
\emph{measure-generic} one: the theorems below are stated over an
arbitrary measure, so the rule carries no profile conditionality at all
--- its instance at any system-semantics profile is the same lemma
applied at that profile\textquotesingle s path measure. A Putnam grader
would accept the one-line proof and then ask where the content went; the
honest answer --- into the decomposition hypothesis --- is
\autoref{rem:qub-content}.}

\textbf{Setting and claim.}\enspace Let $(\Omega, \mu)$ be any
measurable space with a measure, read throughout as the induced outer
measure on arbitrary subsets (coinciding with $\mu$ on measurable sets)
--- no measurability is asked of any event below, and both monotonicity
and finite subadditivity hold in that reading. The rule: from
$\mu(E_1) \le a$, $\mu(E_2) \le b$, and a decomposition
$\mathrm{fail} \subseteq E_1 \cup E_2$, conclude
$\mu(\mathrm{fail}) \le a + b$. In its battery consumption, $\mu$ is the
path measure $\Pr_x$ of the profile at hand, $\mathrm{fail}$ is the
failure event of a composite temporal property, and $a, b$ are the
initial values of two certificates.

\begin{bthm}[soundness of \rulename{QUnionBound}]\label{thm:qub-sound}
$\mu(E_1) \le a$, $\mu(E_2) \le b$, and
$\mathrm{fail} \subseteq E_1 \cup E_2$ imply
$\mu(\mathrm{fail}) \le a + b$.
\end{bthm}
\begin{proof}
Monotonicity, then subadditivity (Boole\textquotesingle s inequality):
$\mu(\mathrm{fail}) \le \mu(E_1 \cup E_2) \le \mu(E_1) + \mu(E_2)
\le a + b$. Note what is \emph{not} used: no measurability of
\emph{any} of the three events (the Setting\textquotesingle s
outer-measure reading supplies monotonicity and finite subadditivity on
arbitrary subsets), no
disjointness of $E_1, E_2$, no probability normalization --- the
statement holds for any measure, which is exactly what makes the rule
profile-generic.
\end{proof}

\begin{bprop}[real-valued form]\label{prop:qub-ofreal}
With $a, b \ge 0$ real and the bounds given as
$\mu(E_i) \le \mathrm{ofReal}$-embedded budgets, the conclusion is
$\mu(\mathrm{fail}) \le a + b$ under the same embedding.
\end{bprop}
\begin{proof}
The embedding of nonnegative reals into extended nonnegative reals is
additive on nonnegatives; apply Theorem~\ref{thm:qub-sound}. This is the
shape battery bounds actually arrive in --- certificate initial values
--- and its nonnegativity hypotheses are free at every use site
(certificate nonneg rows).
\end{proof}

\begin{blemma}[the d\'ebut decomposition --- the packaged instance]
\label{lem:qub-debut}
Over sequence space $S^{\N}$, fix a stop set
$\mathit{Stop} \subseteq S$, a good-exit set $\mathit{Good} \subseteq S$,
and a horizon $T$. Then the complement of the first-exit-success event
--- some $n \le T$ with $\omega_n \in \mathit{Good}$ and no earlier
visit to $\mathit{Stop}$ --- is contained in
\[
\bigl\{\exists n,\ \omega_n \in \mathit{Stop} \smallsetminus \mathit{Good}
\ \wedge\ \forall j < n,\ \omega_j \notin \mathit{Stop}\bigr\}
\;\cup\;
\bigl\{\forall k \le T,\ \omega_k \notin \mathit{Stop}\bigr\},
\]
i.e.\ in $\{$some \emph{first} $\mathit{Stop}$-visit presents a bad
state$\}$ $\cup$ $\{$no $\mathit{Stop}$-visit through $T\}$.
\end{blemma}
\begin{proof}
Suppose the path is not a first-exit success. If it never visits
$\mathit{Stop}$ through $T$, the second disjunct holds. Otherwise some
$k \le T$ has $\omega_k \in \mathit{Stop}$; let $n$ be the \emph{least}
index with $\omega_n \in \mathit{Stop}$ (the well-ordering of $\N$ ---
this is the entire ``d\'ebut'' content; at continuous time the
first-entry construction is a genuine theorem, while over $\N$ it is a
minimum). Then $n \le k \le T$ and no $j < n$ visits $\mathit{Stop}$;
were $\omega_n \in \mathit{Good}$, the path would be a first-exit
success at $n$ --- contradiction. So $\omega_n \in \mathit{Stop}
\smallsetminus \mathit{Good}$, and $n$ is a first visit: the first
disjunct holds. No measure, no measurability, no topology was consulted
--- the lemma is pure path-set algebra.
\end{proof}

\begin{bremark}[where the content lives]\label{rem:qub-content}
The rule is sound for \emph{any} events $E_1, E_2$; every gram of
content sits in the hypothesis
$\mathrm{fail} \subseteq E_1 \cup E_2$, which the citing rule must
justify and the checker must see justified. Drop it and the conclusion
is unrelated to the premises --- one could ``bound'' any event by any
two budgets. In the bounded-until battery\textquotesingle s use,
Lemma~\ref{lem:qub-debut} \emph{is} the justification, at
$\mathit{Stop} := (A \cup R) \cup D^{\mathrm{c}}$ and
$\mathit{Good} := R$: failure of ``stay safe until the target, by the
deadline'' splits into a bad first exit (bounded by the avoid
certificate) and no exit by the horizon (bounded by the deadline
certificate). The split covers every failure: a failing path either
never visits $\mathit{Stop}$ through the horizon (the second disjunct)
or has a first $\mathit{Stop}$-visit, which either is or is not a good
landing --- and a good landing, since $R \subseteq \mathit{Stop}$ in
this instance, is a success, so no third thing remains.
\end{bremark}

\begin{bremark}[measure-genericity, stated sibling-aware]
\label{rem:qub-generic}
Within this book\textquotesingle s battery family --- safety, deadline,
until, and this rule --- the other three carry conditional marks (now
recording, for safety, deadline, and until, that their
continuous-time martingale-problem theorems are cited rather than
hash-bound; see each profile remark); this rule alone is warranted
unconditionally, because its statement never mentions a kernel, a path
construction, or a filtration. The prediction this remark has carried
since its first edition held: the continuous-time discharges of the
siblings consume these same two theorems verbatim --- nothing here
needed to be touched. (The d\'ebut
lemma is the sequence-space instance of its decomposition; a
continuous-time consumer would supply its own first-entry construction
--- a genuine theorem there, as the lemma\textquotesingle s proof notes
--- through those same transports.)
\end{bremark}

\begin{bremark}[the second consumer]\label{rem:qub-second}
The same additivity is the slack-composition device in assume--guarantee
composition: when interface contracts compose across a wiring, their
slack budgets add, and the addition is this rule\textquotesingle s
theorem applied to the composed failure events. One lemma, two consumers
--- which is the reason the rule is kept as a separate node rather than
inlined into the until battery: every consumer\textquotesingle s
soundness proof stays honest about where its additivity comes from.
\end{bremark}

\begin{bremark}[why core, and why this small]\label{rem:qub-tier}
There is nothing to unpack into: subadditivity of measure is not
derivable from the other battery rules --- each of them concludes a
bound on a single certificate-shaped event from kernel rows, and none
mentions two arbitrary events, let alone their union --- so the rule
cannot be admissible: it is new semantic content, deliberately minimal. Its smallness is a
design decision inherited from the base
calculus\textquotesingle s taste: keep kernel-level content atomic, let
composite rules be admissible unpackings over the atoms.
\end{bremark}

\begin{anchors}
Theorem \ref{thm:qub-sound} & \texttt{sound\_QUnionBound}\\
Proposition \ref{prop:qub-ofreal} & \texttt{sound\_QUnionBound\_ofReal}\\
Lemma \ref{lem:qub-debut} & \texttt{compl\_firstExit\_subset}\\
the until-battery consumption & the \texttt{sound\_QUnionBound\_ofReal} application inside \texttt{sound\_QUntilBattery\_firstHit}\\
\end{anchors}


\rulesection{43}{QUntilBattery}{admissible rule}
\label{blcatrule-QUntilBattery}
\ruleblock{The rule (Steele metanotation)}
\sourcenote{as realized in the mechanized development; source row: \texttt{\texttt{QTSL-rules.toml}} --- \emph{QUntilBattery}. The metanotation statement is typeset by hand for notational quality; its mechanized ground truth is the Lean statement below, extracted from the sources at every build of this book and hash-bound. The rule \emph{roster} is machine-checked (every mechanized rule tag round-trips to exactly one catalogue row); the statement \emph{text} is editorial, checked by eye against the extracted Lean beside it.}
\begin{mathpar}
\inferrule*{
  E_{\mathrm{bad}} := \{\text{the first exit from } D \smallsetminus (A \cup R) \ \text{lands in } A \cup D^{\mathrm{c}}\}, \quad E_{\mathrm{miss}} := \{A \cup R \ \text{not reached by } T\} \\
  [\![\neg\varphi_s]\!] \subseteq A, \qquad R \subseteq [\![\varphi]\!] \ \text{(formula ties)} \\
  \mu_x(E_{\mathrm{bad}}) \le V_a(x) \ \text{(the \rulename{QSafetyBattery} slice for } V_a \ \text{at stop set } (A \cup R) \cup D^{\mathrm{c}}\text{: floor on } A \ \text{only)} \\
  \mu_x(E_{\mathrm{miss}}) \le \beta^{T}\, V_d(x) \ \text{(the \rulename{QDeadlineBattery} slice for } V_d \ \text{at target } A \cup R\text{)} \\
  [\![\neg(\varphi_s\,\mathsf{U}_{\le T}\,\varphi)]\!] \subseteq E_{\mathrm{bad}} \cup E_{\mathrm{miss}} \ \text{(the d\'ebut decomposition; consumed by \rulename{QUnionBound})} \\
  0 \le \beta, \qquad 0 \le V_d(x), \qquad \beta^{T} \le c \ \text{(sign gates; rate tie)}
}{
  \mathcal{M},\ x_0 \in \{V_a + c\,V_d \le 1 - p\} \vDash \mathsf{SProbGE}\bigl(p,\ \mathsf{PUntil}(T, \varphi_s, \varphi)\bigr)
}
\end{mathpar}

\ruleflavor{rule sort baked in \texttt{theory Qtsl} (\texttt{QUntilBattery}, avoidance and target certificates as separate family fields), the concluded affine region atom (\texttt{qraAffine}) baked beside it --- the display is the realized form, as the discrete-profile assembly binds it. The two slice premises are the sibling batteries\textquotesingle{} kits re-aimed --- the avoid slice is the safety kit at the composite stop set with its floor asked on the bad exit $A$ only, the deadline slice drives into $A \cup R$ by the horizon --- and each premise is the measure bound that slice\textquotesingle s rows conclude; the rows themselves are the sibling entries\textquotesingle{} premises. Catalogue instance of the rate tie: $c = q^{-1}$ with $q \le e^{\theta T}$, $\beta = e^{-\theta}$ --- fully rational at the discrete profile, no transcendental leaf. The conclusion is the region-conditioned judgment in constructor vocabulary, its composite region the baked \texttt{qraAffine} atom over $V_a + c\,V_d$ at threshold $1 - p$.}
\par\medskip
\rulename{QUntilBattery} concludes the bounded-until judgment
$\mathsf{SProbGE}\bigl(p,\ \mathsf{PUntil}(T, \varphi_s, \varphi)\bigr)$
(surface syntax $\mathsf{P}_{\ge p}[\varphi_s\, \mathsf{U}_{\le T}\,
\varphi]$) on a composite \emph{affine} region
$\mathsf{qraAffine}([(1,\, V_a),\, (c,\, V_d)],\ 1 - p)$ --- the
sublevel set $\{V_a + c\, V_d \le 1 - p\}$ --- from two certificates doing different
work: an avoid certificate $V_a$ guarding the safe-side formula
$\varphi_s$ until the target is reached, and a deadline certificate
$V_d$ driving the process into the target by the horizon. It is this
catalogue\textquotesingle s exemplar admissible rule at the satisfaction
stratum: one node unpacks into two battery slices and one
\rulename{QUnionBound} application, and the mechanized soundness proof
\emph{is} that unpacking --- it consumes the deadline
battery\textquotesingle s theorem as a black box, applies the union
bound once, and never re-proves either battery.

The two floors do different work, and the asymmetry is the
rule\textquotesingle s design. $V_a$\textquotesingle s floor is asked on
the bad-exit set $A$ \emph{only} --- the avoid certificate is silent on
$R \smallsetminus A$, so a landing there costs no certificate
budget: the rule does not charge for winning (a landing in
$A \cap R$, when the formula ties permit one, succeeds but is still
priced as avoid-territory). $V_d$\textquotesingle s
floor is on all of the hull off the target, exactly as its parent
deadline battery
demands. Fail-closed: leaving the modelled domain joins the stop set on
the \emph{bad} side, while an $R$-landing succeeds wherever it lands
(the target\textquotesingle s formula tie is global).

What the checker recomputes: the side conditions inherited from both
battery slices, plus the one arithmetic hypothesis binding the slices
together --- the rate tie $\beta^{T} \le c$ on the affine
atom\textquotesingle s second coefficient. The catalogued instance
$c = q^{-1}$ with $q \le e^{\theta T}$ and $\beta = e^{-\theta}$
satisfies the tie; at this profile the whole certificate is exact
rational arithmetic, with no transcendental leaf (one becomes available
by the deadline battery\textquotesingle s Taylor-leaf pattern when
$\theta$-rates are wanted). Drop the tie and the affine atom
under-budgets the miss mass; drop either slice\textquotesingle s rows
and that slice\textquotesingle s own counterexamples apply.

Tier, argued: admissible, not alias --- three premise sub-derivations,
not a respelling of one. And deliberately not monolithic core: a third
self-contained lemma would duplicate two batteries\textquotesingle{}
soundness; the mechanized assembly vindicates the choice by factoring
through the named machinery instead. The dashboard mark states the
status: hash-bound at the discrete-kernel profile; proven at the
continuous-time martingale-problem profile
(\texttt{sound\_GSHS\_until\_mp},
\texttt{sound\_GSHS\_until\_battery\_mp},
\texttt{QtslAux.LayerD.MPDischarge}) --- hash-bound below via the
qtsl-aux package registration, so the row carries no conditional
mark: each named profile is proven, and its theorems are bound.

At the general stratum the until acceptance atom is stated over the
semantic universe\textquotesingle s record carrier directly (the
semantic-universe functor package, hash-bound below): at some grid
index the reach leg is hit and the avoid leg is missed
(\texttt{untilEvent}) --- a countable Boolean combination of
hit-events, measurable from per-component tickets alone
(\texttt{measurableSet\_untilEvent}). The clocked instantiation shape
is recorded at the corollary: reach legs closed (a $\le$-cut of
stages times a closed target), avoid legs open (a $<$-cut times an
open unsafe region) --- and miss-events of open sets are measurable
with no compactness anywhere. As at the deadline battery, the until
conclusion is a mark ($\exists$-face): persistent by the same
window-calculus engine, with the forward judgment lemma and the same
deliberate absence of a backward direction --- the deadline semantics
the two batteries share.

\cataloguenote{The concluded region is the affine region atom with rows
$[(1, V_a), (c, V_d)]$ and threshold $1 - p$ --- an affine sublevel
region over certificate values, the region-vocabulary generalization
whose $k = 1$ case is the ordinary sublevel atom.
}
\ruleblock{Lean statement}
\declhead{BernsteinLogicAux.QtslBattery}{sound\_QUntilBattery}{\badge{BadgeGreen}{statement bound}}
\begin{leansrc}
/-- **`sound_QUntilBattery`, discrete profile** (T8): over the first-hit
form, the two formula ties — `A` over-approximates the falsity set of the
safe-side formula (`¬ sat_s y → y ∈ A`, the T13 OUTWARD direction), `R`
under-approximates the truth set of the target formula (`∀ y ∈ R, sat y`) —
weaken the concluded event to the truth of `φ_s U_{≤T} φ` along the path.
No measurability is asked of either satisfaction predicate (`measure_mono`
on the measurable first-hit event). -/
theorem sound_QUntilBattery {sat_s sat : S → Prop}
    (hD : MeasurableSet D) (hDbar : MeasurableSet Dbar)
    (hA : MeasurableSet A) (hR : MeasurableSet R)
    (hVa : Measurable Va) (hVd : Measurable Vd)
    (hFs : ∀ y, ¬ sat_s y → y ∈ A) (hFt : ∀ y ∈ R, sat y)
    (hβ : 0 ≤ β) (hc : β ^ T ≤ c)
    (hDsub : D ⊆ Dbar)
    (hR3 : ∀ y ∈ D, y ∉ A ∪ R → K y Dbar = 1)
    (hR4 : ∀ y ∈ D, y ∉ A ∪ R → K y (D ∪ (A ∪ R)) = 1)
    (hVaint : ∀ y ∈ D, y ∉ A ∪ R → Integrable Va (K y))
    (hVadrift : ∀ y ∈ D, y ∉ A ∪ R → ∫ z, Va z ∂K y ≤ Va y)
    (hVafloor : ∀ y ∈ A, y ∈ Dbar → 1 ≤ Va y)
    (hVanonneg : ∀ y ∈ Dbar, 0 ≤ Va y)
    (hVdint : ∀ y ∈ D, y ∉ A ∪ R → Integrable Vd (K y))
    (hVddecay : ∀ y ∈ D, y ∉ A ∪ R → ∫ z, Vd z ∂K y ≤ β * Vd y)
    (hVdfloor : ∀ y ∈ Dbar, y ∉ A ∪ R → 1 ≤ Vd y)
    (hVdnonneg : ∀ y ∈ Dbar, 0 ≤ Vd y)
    (hVdM : ∀ y ∈ Dbar, Vd y ≤ M)
    (hx : x ∈ D) (hthr : Va x + c * Vd x ≤ 1 - p) :
    ENNReal.ofReal p ≤ pathMeasure K x
      {ω | ∃ n ≤ T, sat (ω n) ∧ ∀ j < n, sat_s (ω j)}
\end{leansrc}
\declhead{BernsteinLogicAux.QtslBattery}{sound\_QUntilBattery\_firstHit}{\badge{BadgeGreen}{statement bound}}
\begin{leansrc}
/-- **`sound_QUntilBattery`, discrete profile, first-hit form** (the
catalogue's T8 row, fail-closed reading). The rule is
admissible over {T5, T6, `QUnionBound`}, and this proof *is* that unpacking:

* the T5 slice is `measure_firstBad_le` on the composite stop set
  `Stop := (A ∪ R) ∪ Dᶜ` with the clamped avoid certificate `V_a` — floor
  rows on `A` only, drift rows off `A ∪ R` (the process is stopped at the
  target, so `V_a` is silent on `R`);
* the T6 slice is `measure_missTarget_le` at target `A ∪ R` with the
  deadline certificate `V_d` — global floor off the target;
* the glue is one `sound_QUnionBound_ofReal` application whose event
  decomposition is `compl_firstExit_subset`.

The concluded region `{V_a + c·V_d ≤ 1 − p}` with `β ^ T ≤ c` is the
`qraAffine` atom (row `[(1, V_a), (c, V_d)]`, threshold `1 − p`);
the catalogue's `c = q⁻¹` with `q ≤ e^{θT}`, `β = e^{−θ}` satisfies
`β ^ T ≤ c` — the tie is dischargeable by the same T7 leaf pattern as
`sound_QDeadlineBattery`. Fail-closed: leaving the modeled domain `D`
before reaching `R` counts as failure (`Dᶜ` joins the stop set on the bad
side); reaching `R` counts as success wherever it lands, since the formula
tie for `R` is global.

System bridge: for any System whose chart compilation at
the discrete-kernel profile yields `K`, this is the concluded
`QRegionSat`'s reading on that System's records, fail-closed form. -/
theorem sound_QUntilBattery_firstHit
    (hD : MeasurableSet D) (hDbar : MeasurableSet Dbar)
    (hA : MeasurableSet A) (hR : MeasurableSet R)
    (hVa : Measurable Va) (hVd : Measurable Vd)
    (hβ : 0 ≤ β) (hc : β ^ T ≤ c)
    (hDsub : D ⊆ Dbar)
    (hR3 : ∀ y ∈ D, y ∉ A ∪ R → K y Dbar = 1)
    (hR4 : ∀ y ∈ D, y ∉ A ∪ R → K y (D ∪ (A ∪ R)) = 1)
    (hVaint : ∀ y ∈ D, y ∉ A ∪ R → Integrable Va (K y))
    (hVadrift : ∀ y ∈ D, y ∉ A ∪ R → ∫ z, Va z ∂K y ≤ Va y)
    (hVafloor : ∀ y ∈ A, y ∈ Dbar → 1 ≤ Va y)
    (hVanonneg : ∀ y ∈ Dbar, 0 ≤ Va y)
    (hVdint : ∀ y ∈ D, y ∉ A ∪ R → Integrable Vd (K y))
    (hVddecay : ∀ y ∈ D, y ∉ A ∪ R → ∫ z, Vd z ∂K y ≤ β * Vd y)
    (hVdfloor : ∀ y ∈ Dbar, y ∉ A ∪ R → 1 ≤ Vd y)
    (hVdnonneg : ∀ y ∈ Dbar, 0 ≤ Vd y)
    (hVdM : ∀ y ∈ Dbar, Vd y ≤ M)
    (hx : x ∈ D) (hthr : Va x + c * Vd x ≤ 1 - p) :
    ENNReal.ofReal p ≤ pathMeasure K x
      {ω | ∃ n ≤ T, ω n ∈ R ∧ ∀ j < n, ω j ∉ (A ∪ R) ∪ Dᶜ}
\end{leansrc}
\declhead{BernsteinLogicAux.QtslBattery}{measure\_firstBad\_le}{\badge{BadgeGreen}{statement bound}}
\begin{leansrc}
/-- **The stopped avoid bound** (the until battery's T5 slice, stopped at
`τ_{A∪R}`): a certificate `f`, bounded in `[0, C]`,
one-step non-increasing off the stop set `Stop`, and `≥ 1` on a bad-exit set
`Bad`, bounds the probability that some *first* `Stop`-visit presents a
`Bad` state: `μ {∃ n, ω n ∈ Bad ∧ ∀ j < n, ω j ∉ Stop} ≤ f x`.
`Supermartingale.ville_ineq` consumed on the `Battery.stopped` process at
barrier `1`. Unlike `measure_reach_le` (whose floor is on the stop set
itself), the floor here is asked only on `Bad` — which is what lets the
until battery's `V_a` be silent on the target region; `Bad ⊆ Stop` is not
required by the proof. -/
theorem measure_firstBad_le {K : Kernel S S} [IsMarkovKernel K]
    {f : S → ℝ} {Stop Bad : Set S} {C : ℝ}
    (hStop : MeasurableSet Stop) (hf : Measurable f)
    (hf0 : ∀ y, 0 ≤ f y) (hfC : ∀ y, f y ≤ C)
    (hf1 : ∀ y ∈ Bad, 1 ≤ f y)
    (hdrift : ∀ y, y ∉ Stop → ∫ z, f z ∂K y ≤ f y) (x : S) :
    pathMeasure K x {ω | ∃ n, ω n ∈ Bad ∧ ∀ j < n, ω j ∉ Stop}
      ≤ ENNReal.ofReal (f x)
\end{leansrc}
\declhead{BernsteinLogicAux.QtslBattery}{measure\_missTarget\_le}{\badge{BadgeGreen}{statement bound}}
\begin{leansrc}
/-- **The deadline miss bound** (the until battery's T6 slice, at target
`A ∪ R`): under the deadline battery's rows at a target
`B`, the probability of missing `B` through time `T` is at most
`β ^ T · V x`. This is `sound_QDeadline_hit` consumed as a black box and
complemented once — the Lean proof structure mirroring the until rule's
admissible-tier unpacking, which cites the T6 battery rather than re-proving
it. The deliberately weak conclusion shape (a bound on the miss event) is
what the `QUnionBound` application consumes. -/
theorem measure_missTarget_le {K : Kernel S S} [IsMarkovKernel K]
    {D Dbar B : Set S} {V : S → ℝ} {β : ℝ} {x : S} {T : ℕ} {M : ℝ}
    (hD : MeasurableSet D) (hDbar : MeasurableSet Dbar) (hB : MeasurableSet B)
    (hV : Measurable V) (hβ : 0 ≤ β)
    (hDsub : D ⊆ Dbar)
    (hR3 : ∀ y ∈ D, y ∉ B → K y Dbar = 1)
    (hR4 : ∀ y ∈ D, y ∉ B → K y (D ∪ B) = 1)
    (hint : ∀ y ∈ D, y ∉ B → Integrable V (K y))
    (hdecay : ∀ y ∈ D, y ∉ B → ∫ z, V z ∂K y ≤ β * V y)
    (hfloor : ∀ y ∈ Dbar, y ∉ B → 1 ≤ V y)
    (hnonneg : ∀ y ∈ Dbar, 0 ≤ V y)
    (hM : ∀ y ∈ Dbar, V y ≤ M)
    (hx : x ∈ D) :
    pathMeasure K x {ω | ∀ k ≤ T, ω k ∉ B} ≤ ENNReal.ofReal (β ^ T * V x)
\end{leansrc}
\par\medskip
\textbf{The general-stochastic instance, assembled.}\enspace The
until-failure bound on the timed skeleton of the stopped presentation:
failure decomposes into an unsafe presentation or a deadline overrun,
both events of the one timed chain --- the safety half runs the
avoidance engine at zero discount (the clock drops out pointwise), the
deadline half is the deadline instance below, and the glue is measure
algebra. Headline and battery forms.

\declhead{BernsteinLogicAux.LayerD.Until}{sound\_GSHS\_until}{\badge{BadgeGreen}{statement bound}}
\begin{leansrc}
/-- **The GSHS until failure bound** (T8, headline form — the K4 event
decomposition on one chain). Hypotheses are exactly the union of the two
halves' (each inherited premise stated, never silent — the safety half's
global bound `C`, the deadline half's `θ`-row and F7 floor `ε ≤ V_d`):
the probability that the timed skeleton of the stopped presentation ever
presents an unsafe value **or** overruns the deadline is at most
`V_a(x) + e^{−θT}·V_d(x)/ε`. -/
theorem sound_GSHS_until [IsMarkovKernel M.seg] [IsMarkovKernel M.jump]
    {Va Vd : State → ℝ} {Us : Set State} {θ ε : ℝ} {T : ℝ≥0}
    (hVa : Measurable Va) (hUs : MeasurableSet Us)
    (hSegDa : SegmentDynkin M 0 Va) (hForcea : ForcedNonexpansion M Va)
    (hNonnega : NonnegV Va) {C : ℝ} (hC0 : 0 ≤ C) (hVC : ∀ y, Va y ≤ C)
    (hFloora : ∀ y ∈ Us, 1 ≤ Va y)
    (hVd : Measurable Vd) (hθ : 0 ≤ θ)
    (hSegDd : SegmentDynkin M θ Vd) (hForced : ForcedNonexpansion M Vd)
    (hNonnegd : NonnegV Vd) (hε : 0 < ε) (hFloord : ∀ y, ε ≤ Vd y)
    (x : State) :
    Battery.pathMeasure (tskel M) (TSkelState.alive (0, x))
        {ω | ∃ n, ω n ∈ tbadIn Us ∨ T < tclock (ω n)}
      ≤ ENNReal.ofReal (Va x)
        + ENNReal.ofReal (Real.exp (-(θ * T)) * Vd x / ε)
\end{leansrc}
\declhead{BernsteinLogicAux.LayerD.Until}{sound\_GSHS\_until\_battery}{\badge{BadgeGreen}{statement bound}}
\begin{leansrc}
/-- **The until battery form** (T8 as `QUntilBattery`'s GSHS instance,
complement arithmetic): with the additive threshold accounting
`p + thr_a + e^{−θT}·thr_d/ε ≤ 1` and a start under both thresholds, the
chain presents only safe values with the clock never past `T`, with
probability at least `p`. On the stop-IS-target instantiation that event
reads *`φ_s` held until the target, reached by the deadline* — up to the
stall residues (inside the event, named, emptied a.s. in-scope by
W-Z-freedom). The composite-atom rendering `{V_a + (1/(qε))·V_d ≤ 1−p}`
of the catalogue's T8 entry maps onto (`thr_a`, `thr_d`) with
`q ≤ e^{θT}`; the `qraAffine` staging decision is
schema-side and does not gate this lemma. -/
theorem sound_GSHS_until_battery [IsMarkovKernel M.seg] [IsMarkovKernel M.jump]
    {Va Vd : State → ℝ} {Us : Set State} {θ ε p thra thrd : ℝ} {T : ℝ≥0}
    (hVa : Measurable Va) (hUs : MeasurableSet Us)
    (hSegDa : SegmentDynkin M 0 Va) (hForcea : ForcedNonexpansion M Va)
    (hNonnega : NonnegV Va) {C : ℝ} (hC0 : 0 ≤ C) (hVC : ∀ y, Va y ≤ C)
    (hFloora : ∀ y ∈ Us, 1 ≤ Va y)
    (hVd : Measurable Vd) (hθ : 0 ≤ θ)
    (hSegDd : SegmentDynkin M θ Vd) (hForced : ForcedNonexpansion M Vd)
    (hNonnegd : NonnegV Vd) (hε : 0 < ε) (hFloord : ∀ y, ε ≤ Vd y)
    (hΘ : p + (thra + Real.exp (-(θ * T)) * thrd / ε) ≤ 1)
    (x : State) (hthra : Va x ≤ thra) (hthrd : Vd x ≤ thrd) :
    ENNReal.ofReal p ≤ Battery.pathMeasure (tskel M) (TSkelState.alive (0, x))
      {ω | ∀ n, ω n ∉ tbadIn Us ∧ tclock (ω n) ≤ T}
\end{leansrc}
\declhead{BernsteinLogicAux.LayerD.Until}{sound\_GSHS\_avoid\_timed}{\badge{BadgeGreen}{statement bound}}
\begin{leansrc}
/-- **The safety half on the timed chain**: the probability that the timed
skeleton ever presents a value in `Us` (alive entries *and* value-carrying
halts — kills into `Us` included) is bounded by the initial safety
certificate. This is `sound_GSHS_safety` transported to the timed chain —
at `θ = 0` the clock drops out of the certificate pointwise, so the same
engine (`Battery.measure_reach_le` = K1) runs unchanged. -/
theorem sound_GSHS_avoid_timed [IsMarkovKernel M.seg] [IsMarkovKernel M.jump]
    {Va : State → ℝ} {Us : Set State}
    (hVa : Measurable Va) (hUs : MeasurableSet Us)
    (hSegD : SegmentDynkin M 0 Va) (hForce : ForcedNonexpansion M Va)
    (hNonneg : NonnegV Va) {C : ℝ} (hC0 : 0 ≤ C) (hVC : ∀ y, Va y ≤ C)
    (hFloor : ∀ y ∈ Us, 1 ≤ Va y) (x : State) :
    Battery.pathMeasure (tskel M) (TSkelState.alive (0, x))
        {ω | ∃ n, ω n ∈ tbadIn Us}
      ≤ ENNReal.ofReal (Va x)
\end{leansrc}
\par\medskip
\textbf{The martingale-problem profile, end to end.}\enspace The
general continuous-time theorems: both slices\textquotesingle{}
formerly per-model constituents --- the safety slice\textquotesingle s
stopping theory, the deadline slice\textquotesingle s discounted
supermartingale --- are supplied by the same martingale-problem
discharge that serves the parents, composite stop set included, and
the until conclusion follows for every system whose presentation
satisfies the martingale-problem hypothesis.

\declhead{QtslAux.LayerD.MPDischarge}{sound\_GSHS\_until\_mp}{\badge{BadgeGreen}{statement bound}}
\begin{leansrc}
/-- The timed until reach bound at the martingale-problem profile: both the
avoidance certificate (rate `0`) and the deadline certificate (rate `θ`)
are presented by martingale problems. -/
theorem sound_GSHS_until_mp [IsMarkovKernel M.seg] [IsMarkovKernel M.jump]
    {Va Vd : State → ℝ} {Us : Set State} {θ ε : ℝ} {T : ℝ≥0}
    (hVa : Measurable Va) (hUs : MeasurableSet Us)
    (hmpa : ∀ x : State, Nonempty (MPFlow M 0 Va x))
    (hForcea : ForcedNonexpansion M Va)
    (hNonnega : NonnegV Va) {C : ℝ} (hC0 : 0 ≤ C) (hVC : ∀ y, Va y ≤ C)
    (hFloora : ∀ y ∈ Us, 1 ≤ Va y)
    (hVd : Measurable Vd) (hθ : 0 ≤ θ)
    (hmpd : ∀ x : State, Nonempty (MPFlow M θ Vd x))
    (hForced : ForcedNonexpansion M Vd)
    (hNonnegd : NonnegV Vd) (hε : 0 < ε) (hFloord : ∀ y, ε ≤ Vd y)
    (x : State) :
    Battery.pathMeasure (tskel M) (TSkelState.alive (0, x))
        {ω | ∃ n, ω n ∈ tbadIn Us ∨ T < tclock (ω n)}
      ≤ ENNReal.ofReal (Va x)
        + ENNReal.ofReal (Real.exp (-(θ * T)) * Vd x / ε)
\end{leansrc}
\declhead{QtslAux.LayerD.MPDischarge}{sound\_GSHS\_until\_battery\_mp}{\badge{BadgeGreen}{statement bound}}
\begin{leansrc}
/-- The until battery at the martingale-problem profile. -/
theorem sound_GSHS_until_battery_mp
    [IsMarkovKernel M.seg] [IsMarkovKernel M.jump]
    {Va Vd : State → ℝ} {Us : Set State} {θ ε p thra thrd : ℝ} {T : ℝ≥0}
    (hVa : Measurable Va) (hUs : MeasurableSet Us)
    (hmpa : ∀ x : State, Nonempty (MPFlow M 0 Va x))
    (hForcea : ForcedNonexpansion M Va)
    (hNonnega : NonnegV Va) {C : ℝ} (hC0 : 0 ≤ C) (hVC : ∀ y, Va y ≤ C)
    (hFloora : ∀ y ∈ Us, 1 ≤ Va y)
    (hVd : Measurable Vd) (hθ : 0 ≤ θ)
    (hmpd : ∀ x : State, Nonempty (MPFlow M θ Vd x))
    (hForced : ForcedNonexpansion M Vd)
    (hNonnegd : NonnegV Vd) (hε : 0 < ε) (hFloord : ∀ y, ε ≤ Vd y)
    (hΘ : p + (thra + Real.exp (-(θ * T)) * thrd / ε) ≤ 1)
    (x : State) (hthra : Va x ≤ thra) (hthrd : Vd x ≤ thrd) :
    ENNReal.ofReal p ≤ Battery.pathMeasure (tskel M) (TSkelState.alive (0, x))
      {ω | ∀ n, ω n ∉ tbadIn Us ∧ tclock (ω n) ≤ T}
\end{leansrc}
\par\medskip
\textbf{The general stratum: the until atom on the record
carrier.}\enspace The deadline-grid until event --- some grid index
at which the reach leg is hit and the avoid leg missed --- as a
countable Boolean combination of hit-events, measurable from
per-component tickets; the clocked corollary takes closed reach legs
and open avoid legs, with no compactness on the avoid side.

\declhead{SuFunctor.Site.FellAcceptance}{untilEvent}{\badge{BadgeGreen}{statement bound}}
\begin{leansrc}
/-- **The deadline-grid until event**: at some grid index, the reach
leg is hit and the avoid leg is missed.  The countable-combination
skeleton of safety-until-a-deadline at the record tier: the clocked
instantiation takes `Reach n` = (stages ≤ nth cut) × target and
`Avoid n` = (stages < nth cut) × unsafe. -/
def untilEvent (Reach Avoid : ℕ → Set α) : Set (FellCloseds α)
\end{leansrc}
\declhead{SuFunctor.Site.FellAcceptance}{mem\_untilEvent\_iff}{\badge{BadgeGreen}{statement bound}}
\begin{leansrc}
theorem mem_untilEvent_iff {Reach Avoid : ℕ → Set α}
    {F : FellCloseds α} :
    F ∈ untilEvent Reach Avoid ↔
      ∃ n, ((F : Set α) ∩ Reach n).Nonempty ∧
        Disjoint (F : Set α) (Avoid n)
\end{leansrc}
\declhead{SuFunctor.Site.FellAcceptance}{measurableSet\_untilEvent}{\badge{BadgeGreen}{statement bound}}
\begin{leansrc}
theorem measurableSet_untilEvent {Reach Avoid : ℕ → Set α}
    (h₁ : ∀ n, MeasurableSet (hitEvent (Reach n)))
    (h₂ : ∀ n, MeasurableSet (missEvent (Avoid n))) :
    MeasurableSet (untilEvent Reach Avoid)
\end{leansrc}
\declhead{SuFunctor.Site.FellAcceptance}{measurableSet\_untilEvent\_of\_isClosed\_isOpen}{\badge{BadgeGreen}{statement bound}}
\begin{leansrc}
/-- The clocked until shape: closed reach legs (a `≤`-cut of stages
times a closed target), open avoid legs (a `<`-cut times an open
unsafe region).  The avoid legs' measurability costs nothing at all —
miss events of opens are Fell-closed at every stratum. -/
theorem measurableSet_untilEvent_of_isClosed_isOpen [SigmaCompactSpace α]
    {Reach Avoid : ℕ → Set α} (h₁ : ∀ n, IsClosed (Reach n))
    (h₂ : ∀ n, IsOpen (Avoid n)) :
    MeasurableSet (untilEvent Reach Avoid)
\end{leansrc}
\par\medskip
\textbf{Marks over the record carrier.}\enspace As at the deadline
battery: the until conclusion\textquotesingle s general-stratum spec
family is persistent with the forward judgment lemma, and the
deliberately absent backward direction is the shared deadline
semantics.

\declhead{SuFunctor.Site.FellBehaviorSpec}{acceptanceSpecF}{\badge{BadgeGreen}{statement bound}}
\begin{leansrc}
/-- The spec family of a per-window acceptance event (deadline, until,
Streett): clip demands of the events' indicators. -/
def acceptanceSpecF (U : ι → Set T)
    (E : ∀ i, Set (FTraj (U i) X)) (j : ι) :
    Set (MeasPen (FTraj (U j) X))
\end{leansrc}
\declhead{SuFunctor.Site.FellBehaviorSpec}{isPersistentMF\_acceptanceSpecF}{\badge{BadgeGreen}{statement bound}}
\begin{leansrc}
theorem isPersistentMF_acceptanceSpecF (U : ι → Set T)
    (hU : ∀ i, IsOpen (U i)) (E : ∀ i, Set (FTraj (U i) X)) :
    IsPersistentMF U (acceptanceSpecF U E)
\end{leansrc}
\declhead{SuFunctor.Site.FellBehaviorSpec}{satAtM\_acceptanceSpecF\_ae}{\badge{BadgeGreen}{statement bound}}
\begin{leansrc}
/-- The forward judgment: satisfying an acceptance spec at a window
forces almost-sure acceptance there.  The backward direction is
deliberately absent — liveness events do not clip-restrict, so
subwindow demands are genuinely stronger (a reach demand within a
subwindow is a deadline); the module docstring carries the reading. -/
theorem satAtM_acceptanceSpecF_ae (U : ι → Set T)
    (E : ∀ i, Set (FTraj (U i) X)) (j : ι)
    {μ : MeasJoint (FTraj (U j) X)}
    (h : SatAtM μ (acceptanceSpecF U E j)) :
    (μ.1 : Measure (FTraj (U j) X)) (E j)ᶜ = 0
\end{leansrc}
\ruleblock{Correspondence (Steele display $\leftrightarrow$ Lean)}
\badge{BadgeAmber}{a bound Lean statement CHANGED since certification --- re-review and re-certify}
\sourcenote{The dictionary below is a reviewer-certified statement of how the mechanized statements above interpret the displayed rule: display metavariable $\mapsto$ Lean binder, each display hypothesis $\mapsto$ the premise or side condition carrying it, the display conclusion $\mapsto$ the goal --- with any deliberate difference in width between display and mechanization named outright. Certification is recorded as a hash triple over (the displayed cell, the extracted statement texts, this dictionary\textquotesingle s own text); if any of the three changes, this block turns amber until a reviewer re-certifies. The correspondence claim itself remains editorial --- reviewer-grade, not machine-checked; its machine-checked successor types each soundness statement against the rule as a structured object, and supersedes this block per rule as it lands.}
Model frame and profile as at \rulename{QSafetyBattery}
(discrete-kernel; continuous time at the martingale-problem
presentation, hash-bound below per the profile remark). One
structural honesty first:
the display presents the two battery slices by their \emph{conclusions}
--- the measure bounds $\mu_x(E_{\mathrm{bad}}) \le V_a(x)$ and
$\mu_x(E_{\mathrm{miss}}) \le \beta^T V_d(x)$ --- plus the d\'ebut
containment that \rulename{QUnionBound} consumes, while the certified
statement \texttt{sound\_QUntilBattery} takes the two slices\textquotesingle{}
row families \emph{inline} (flattened hypothesis lists over one
kernel) and derives those bounds internally (the witness names the
waypoints: \texttt{measure\_firstBad\_le},
\texttt{measure\_missTarget\_le}, with the containment bound as
\texttt{compl\_firstExit\_subset}); the admissible-tier
unpacking is realized by the \emph{proof}, which consumes the QDeadlineBattery
theorem as a black box, not by the statement\textquotesingle s shape.
\begin{itemize}
\item the concluded judgment $\mathcal{M}
  ,\ x_0 \in \{V_a + c\,V_d \le 1 - p\} \vDash
  \mathsf{SProbGE}\bigl(p,\ \mathsf{PUntil}(T, \varphi_s,
  \varphi)\bigr)$ $\mapsto$ the goal
  $\mathrm{ofReal}\ p \le$ \texttt{pathMeasure K x \{ω | ∃ n ≤ T, sat
  (ω n) ∧ ∀ j < n, sat\_s (ω j)\}} --- the until event verbatim, both
  satisfaction predicates measurability-free.
\item the two formula ties, opposite polarities (the QClosure directions),
  displayed as the inclusions $[\![\neg\varphi_s]\!] \subseteq A$ and
  $R \subseteq [\![\varphi]\!]$: $A$ \emph{over}-approximates the
  falsity set of $\varphi_s$
  (\texttt{hFs}) and $R$ \emph{under}-approximates the truth set of
  $\varphi$ (\texttt{hFt}) --- the displayed inclusions verbatim, safe
  direction on both sides.
\item the \rulename{QSafetyBattery} slice for $V_a$ (rows behind the
  displayed bound $\mu_x(E_{\mathrm{bad}}) \le V_a(x)$) $\mapsto$
  \texttt{hVadrift}/\texttt{hVaint}
  (drift at the composite stop set $(A \cup R) \cup D^{\mathrm{c}}$,
  encoded as $y \in D,\ y \notin A \cup R$),
  \texttt{hVafloor} --- on $A$ \emph{only}, the display\textquotesingle
  s ``floor on $A$ only\textquotesingle\textquotesingle{}
  --- and \texttt{hVanonneg}.
\item the \rulename{QDeadlineBattery} slice for $V_d$ (rows behind the
  displayed bound $\mu_x(E_{\mathrm{miss}}) \le \beta^T V_d(x)$) $\mapsto$
  \texttt{hVddecay}/\texttt{hVdint} ($\beta$-decay off $A \cup R$),
  \texttt{hVdfloor} (global off-target), \texttt{hVdnonneg},
  \texttt{hVdM}.
\item the rate tie $\beta^T \le c$ and its displayed sign gates
  $0 \le \beta$, $0 \le V_d(x)$ $\mapsto$ \texttt{hc}, \texttt{hβ},
  and the nonneg row read at the start point
  (\texttt{hVdnonneg} with \texttt{hx}/\texttt{hDsub}) --- the
  $V_d(x)$ gate is what lets the outer assembly pass from the
  $\beta^T V_d(x)$ miss bound to the affine row\textquotesingle s
  $c\,V_d(x)$; fully
  rational at this profile; the catalogue\textquotesingle s
  $c = q^{-1}$ is one instance.
\item the concluded \texttt{qraAffine} atom $[(1, V_a), (c, V_d)]$ at
  threshold $1 - p$ $\mapsto$ the start condition \texttt{hthr} $:$
  \texttt{Va x + c·Vd x ≤ 1 − p}, verbatim affine row.
\item side conditions and regularity as at the siblings
  (\texttt{hR3}/\texttt{hR4} at the composite stop set; the
  measurability pack undisplayed).
\end{itemize}
Width note: beyond the two safe-direction ties, neither side is wider
--- every displayed row lands on a named hypothesis, and the
statement\textquotesingle s flattening loses nothing the display
promises (each inline family is exactly its slice\textquotesingle s
row family, stopped where the display stops it).

\ruleblock{Lean verifier receipts}
\sourcenote{module \texttt{BernsteinLogicAux.QtslBattery}; bound build \texttt{5c75f9923f27\ldots} (bound from receipt 3fd830b4914bdeda, 2026-08-04T01:56:58Z). A receipt records that named checking layers accepted a build: (module, compiled-artifact SHA-256, environment SHA-256, checker layers), minted when the checking gate ran; the store receipts attesting the bound build are listed with the layers each run credited. Nothing here asks the reader to trust the store: every layer is an independently runnable program over the quoted sources --- rebuild the module and re-run the checkers to reproduce the verdict. Receipts attesting superseded builds credit no layer and are omitted from this rendering (the count of omissions, if any, is noted below).}
\begin{receipttable}
\receiptrow{3fd830b4914bdeda}{2026-08-04T01:56:58Z}{build+decorrelated}{\badge{BadgeGreen}{current}}
\receiptrow{96b18d8572451f65}{2026-08-04T03:11:31Z}{build+decorrelated}{\badge{BadgeGreen}{current}}
\end{receipttable}
\sourcenote{13 further receipt(s) in the store attest superseded builds of this module and are omitted from this rendering.}
\sourcenote{module \texttt{BernsteinLogicAux.LayerD.Until}; bound build \texttt{5abcb1cb41f6\ldots} (bound from receipt dbbf5ae370de0748, 2026-08-04T01:57:07Z). A receipt records that named checking layers accepted a build: (module, compiled-artifact SHA-256, environment SHA-256, checker layers), minted when the checking gate ran; the store receipts attesting the bound build are listed with the layers each run credited. Nothing here asks the reader to trust the store: every layer is an independently runnable program over the quoted sources --- rebuild the module and re-run the checkers to reproduce the verdict. Receipts attesting superseded builds credit no layer and are omitted from this rendering (the count of omissions, if any, is noted below).}
\begin{receipttable}
\receiptrow{dbbf5ae370de0748}{2026-08-04T01:57:07Z}{build+decorrelated}{\badge{BadgeGreen}{current}}
\receiptrow{e2dab1e692fdab41}{2026-08-04T03:11:31Z}{build+decorrelated}{\badge{BadgeGreen}{current}}
\end{receipttable}
\sourcenote{11 further receipt(s) in the store attest superseded builds of this module and are omitted from this rendering.}
\sourcenote{module \texttt{QtslAux.LayerD.MPDischarge}; bound build \texttt{baded74dd943\ldots} (bound from receipt 7228834dea10e5f9, 2026-07-24T11:04:04Z). A receipt records that named checking layers accepted a build: (module, compiled-artifact SHA-256, environment SHA-256, checker layers), minted when the checking gate ran; the store receipts attesting the bound build are listed with the layers each run credited. Nothing here asks the reader to trust the store: every layer is an independently runnable program over the quoted sources --- rebuild the module and re-run the checkers to reproduce the verdict. Receipts attesting superseded builds credit no layer and are omitted from this rendering (the count of omissions, if any, is noted below).}
\begin{receipttable}
\receiptrow{7228834dea10e5f9}{2026-07-24T11:04:04Z}{build+decorrelated}{\badge{BadgeGreen}{current}}
\end{receipttable}
\sourcenote{module \texttt{SuFunctor.Site.FellAcceptance}; bound build \texttt{651fb557bbfe\ldots} (bound from receipt aa2169560f8f571a, 2026-08-03T20:19:04Z). A receipt records that named checking layers accepted a build: (module, compiled-artifact SHA-256, environment SHA-256, checker layers), minted when the checking gate ran; the store receipts attesting the bound build are listed with the layers each run credited. Nothing here asks the reader to trust the store: every layer is an independently runnable program over the quoted sources --- rebuild the module and re-run the checkers to reproduce the verdict. Receipts attesting superseded builds credit no layer and are omitted from this rendering (the count of omissions, if any, is noted below).}
\begin{receipttable}
\receiptrow{7ad85720de6e419e}{2026-08-03T10:13:50Z}{build+decorrelated}{\badge{BadgeGreen}{current}}
\receiptrow{94a408109996d1bf}{2026-08-03T13:41:21Z}{build+decorrelated}{\badge{BadgeGreen}{current}}
\receiptrow{580a7e551fab03b6}{2026-08-03T15:39:44Z}{build+decorrelated}{\badge{BadgeGreen}{current}}
\receiptrow{aa2169560f8f571a}{2026-08-03T20:19:04Z}{build+decorrelated}{\badge{BadgeGreen}{current}}
\end{receipttable}
\sourcenote{2 further receipt(s) in the store attest superseded builds of this module and are omitted from this rendering.}
\sourcenote{module \texttt{SuFunctor.Site.FellBehaviorSpec}; bound build \texttt{345f9caf9453\ldots} (bound from receipt 72f2c3b62b6771fc, 2026-08-03T20:19:08Z). A receipt records that named checking layers accepted a build: (module, compiled-artifact SHA-256, environment SHA-256, checker layers), minted when the checking gate ran; the store receipts attesting the bound build are listed with the layers each run credited. Nothing here asks the reader to trust the store: every layer is an independently runnable program over the quoted sources --- rebuild the module and re-run the checkers to reproduce the verdict. Receipts attesting superseded builds credit no layer and are omitted from this rendering (the count of omissions, if any, is noted below).}
\begin{receipttable}
\receiptrow{d7b846e785c1d7fb}{2026-08-03T10:13:53Z}{build+decorrelated}{\badge{BadgeGreen}{current}}
\receiptrow{c0e3dc15419765aa}{2026-08-03T13:41:24Z}{build+decorrelated}{\badge{BadgeGreen}{current}}
\receiptrow{9a62b29c08b63b4f}{2026-08-03T15:39:47Z}{build+decorrelated}{\badge{BadgeGreen}{current}}
\receiptrow{72f2c3b62b6771fc}{2026-08-03T20:19:08Z}{build+decorrelated}{\badge{BadgeGreen}{current}}
\end{receipttable}
\sourcenote{2 further receipt(s) in the store attest superseded builds of this module and are omitted from this rendering.}
\layerline{\layeritem{build}{lake build (elaboration)}{\badge{BadgeGreen}{current}}\layeritem{replay}{leanchecker (official kernel replay)}{\badge{BadgeAmber}{stale only}}\layeritem{hardened}{safe\_verify (hardened screen)}{\badge{BadgeAmber}{stale only}}\layeritem{decorrelated}{lean4export + nanoda (decorrelated kernel)}{\badge{BadgeGreen}{current}}}
\sourcenote{\badge{BadgeGray}{comparator: none} no fixture in the three-implementation comparator harness (compiled / kernel-interpreted / reference-interpreted) exercises this rule.}
\ruleblock{HOL4 statement}
\absentbox{NO HOL STATEMENT YET}{The HOL4 development ports the base calculus\textquotesingle s checker
chain (theory \texttt{bern\_catalogue}, developed side by side with the
base catalogue sources, stopping below the proof-object layer); it
contains no probability or measure theory. \rulename{QUntilBattery}
lives one theory-extension edge above the base calculus and its
soundness content is path-space measure theory (a stopped-supermartingale
slice, a discounted-supermartingale slice, and a union bound over a
d\'ebut decomposition), for which no counterpart exists anywhere in the
ported development. A HOL statement of this rule would be a new
measure-theoretic campaign, not an extension of the ported checker
chain; until one exists, this box is the honest state of the edifice.
}
\ruleblock{HOL verifier receipts}
\absentbox{NONE}{no HOL verification receipt exists for this rule.}
\ruleblock{Proof}

\sourcenote{Genre: admissible rule (tier 2) --- a proof-tree-rewriting
recipe. The general admissible template (one node $\to$ a derivation
fragment over core and previously-justified rules, applied leaf-first)
is stated in the \rulename{Norm} section, Recipe~\ref{rec:norm-template},
and not restated here; this section gives the rule\textquotesingle s
instance. The verified
theorems quoted above warrant the rule at the \textbf{discrete-kernel
profile}; the profile remark at the end states what remains a
hypothesis beyond it.}

\begin{brecipe}[the \rulename{QUntilBattery} rewrite]\label{rec:qutb}
A \rulename{QUntilBattery} node carries: an avoid certificate $V_a$ with
its rows (drift off $A \cup R$, floor \emph{on $A$ only}, nonneg), a
deadline certificate $V_d$ with its rows (decay factor $\beta$ off
$A \cup R$, floor on all of the hull off the target, nonneg, bound),
the shared side
conditions (hull; fail-closed at the composite exit set $A \cup R$), the
sign gate $\beta \ge 0$ with the rate tie $\beta^{T} \le c$, and the
threshold row
$V_a(x) + c\,V_d(x) \le 1 - p$ on the concluded affine region. Here $A$
over-approximates the falsity set of the safe-side formula $\varphi_s$
and $R$ under-approximates the truth set of the target $\varphi$. The
node rewrites to a derivation fragment over three constituents:
\[
\inferrule*[Right=\rulename{QUntilBattery}]{
  \text{$V_a$ rows} \\ \text{$V_d$ rows} \\ \beta^{T} \le c \\
  V_a + c\,V_d \le 1 - p \ \text{on the region}
}{\mathcal{M},\ x_0 \in \{V_a + c\,V_d \le 1 - p\} \vDash
  \mathsf{SProbGE}\bigl(p,\ \mathsf{PUntil}(T, \varphi_s, \varphi)\bigr)}
\;\rightsquigarrow\;
\begin{array}{l}
\text{(i) a \emph{QSafetyBattery slice}: the battery pattern
  (\rulename{QSafetyBattery}, Prop.~\ref{prop:qsb-pattern})}\\
\quad\text{at the composite stop set
  $\mathit{Stop} := (A \cup R) \cup D^{\mathrm{c}}$, certificate $V_a$;}\\
\text{(ii) a \emph{QDeadlineBattery slice}: the deadline miss bound
  (\rulename{QDeadlineBattery}, Lem.~\ref{lem:qdb-tail} via
  Thm.~\ref{thm:qdb-hit})}\\
\quad\text{at target $A \cup R$, certificate $V_d$;}\\
\text{(iii) one \rulename{QUnionBound} node on the d\'ebut
  decomposition (Lem.~\ref{lem:qub-debut}).}
\end{array}
\]
Honesty note, in the \rulename{PConstExact} tradition: the two slices
are consumed through their \emph{verified slice lemmas}, not by literal
re-application of full battery nodes --- the QSafetyBattery
slice\textquotesingle s floor family is deliberately weaker than a full
safety node\textquotesingle s (floor on the bad-exit set only; see
\autoref{rem:qutb-floors}), and the QDeadlineBattery slice is consumed in
complemented (miss-bound) form. The shipped verified object is therefore
the assembled soundness lemma, whose Lean structure \emph{is} this
rewrite: it consumes the deadline theorem as a black box, applies the
union bound once, and re-proves neither battery.
\end{brecipe}

\begin{bprop}[semantic warrant, first-hit form]\label{prop:qutb-firsthit}
On the substrate of the shared preamble, suppose: the $V_a$ rows ---
drift off $A \cup R$ within $D$, floor $V_a \ge 1$ on $A \cap \bar{D}$,
nonneg on $\bar{D}$; the $V_d$ rows --- decay
$\int V_d\,dK\,y \le \beta V_d(y)$ off $A \cup R$ within $D$, floor
$V_d \ge 1$ on $\bar{D} \smallsetminus (A \cup R)$, nonneg and bound $M$
on $\bar{D}$; hull and fail-closed side conditions at exit set
$A \cup R$; $\beta \ge 0$ and the rate tie $\beta^{T} \le c$. Then for
every $x \in D$ with $V_a(x) + c\, V_d(x) \le 1 - p$,
\[
\Pr_x\bigl[\exists n \le T,\ \omega_n \in R \ \wedge\
\forall j < n,\ \omega_j \notin \mathit{Stop}\bigr] \;\ge\; p,
\qquad \mathit{Stop} = (A \cup R) \cup D^{\mathrm{c}} .
\]
\end{bprop}
\begin{proof}
Three bounds, then arithmetic.
\emph{(i, the QSafetyBattery slice.)} Clamp $V_a$ as in the safety battery, with the
floor forced only off $D \cup R$ (the certificate is silent on the
target away from $A$). The stopped process at $\mathit{Stop}$ is a
nonnegative
supermartingale, $\ge 1$ whenever a \emph{first} $\mathit{Stop}$-visit
presents a bad state: a bad state lies in
$\mathit{Stop} \smallsetminus R = (A \smallsetminus R) \cup
(D^{\mathrm{c}} \smallsetminus R)$, and every reachable one is covered
--- in $A \smallsetminus R$ it lies in $\bar D$ (a first visit is either
the start, in $D$, or one kernel step from $D \smallsetminus
\mathit{Stop}$, confined to $\bar D$ by the hull row), where the floor
row on $A \cap \bar D$ applies; in $D^{\mathrm{c}} \smallsetminus R$,
which \emph{is} the region off $D \cup R$, the forced clamp applies.
Ville\textquotesingle s inequality gives
\[
E_1 := \bigl\{\text{some first } \mathit{Stop}\text{-visit is bad}\bigr\},
\qquad \Pr_x[E_1] \le V_a(x).
\]
The floor is asked on $A$ and off $D \cup R$ only --- in particular
never on $R \smallsetminus A$, which is what lets a landing there cost
no $V_a$-budget (a landing in $A \cap R$, if the ties permit one, still
succeeds, but the avoid floor prices it as avoid-territory; the
no-charge-for-winning story is the $R \smallsetminus A$ story ---
\autoref{rem:qutb-floors}).
\emph{(ii, the QDeadlineBattery slice.)} The deadline battery\textquotesingle s miss
bound at target $A \cup R$, consumed as a black box and weakened from
missing $A \cup R$ to missing the larger $\mathit{Stop}$ --- avoiding
$\mathit{Stop} \supseteq A \cup R$ through $T$ entails avoiding
$A \cup R$ through $T$, so $E_2$ below is contained in the
black box\textquotesingle s miss event:
\[
E_2 := \bigl\{\forall k \le T,\ \omega_k \notin \mathit{Stop}\bigr\},
\qquad \Pr_x[E_2] \le \beta^{T} V_d(x) \le c\, V_d(x),
\]
the last step by the rate tie and $V_d(x) \ge 0$.
\emph{(iii, the glue.)} The d\'ebut decomposition
(Lemma~\ref{lem:qub-debut}) at $(\mathit{Stop}, R, T)$ says the
complement of the concluded first-hit event is contained in
$E_1 \cup E_2$; one \rulename{QUnionBound} application in real-valued
form (Prop.~\ref{prop:qub-ofreal}, its nonnegativity riders supplied by
the nonneg rows and by $c \ge \beta^{T} \ge 0$) gives
$\Pr_x[\text{fail}] \le V_a(x) + c\,V_d(x) \le 1 - p$. Complement ---
licit because the first-hit event is a region event over measurable
regions (the profile\textquotesingle s regularity pack), so its measure
and its complement\textquotesingle s sum to one; the weakening to
possibly-nonmeasurable satisfaction events happens only in the next
proposition, in the monotone direction.
\end{proof}

\begin{bprop}[soundness of \rulename{QUntilBattery}]\label{prop:qutb-sound}
With the two formula ties --- $A$ over-approximating where $\varphi_s$
fails, $R$ under-approximating where $\varphi$ holds --- the first-hit
event weakens to the truth of the bounded-until formula along the path:
for $x$ in the affine region,
$\Pr_x[\exists n \le T,\ \mathrm{sat}(\omega_n) \wedge
\forall j < n,\ \mathrm{sat}_s(\omega_j)] \ge p$ --- the concluded
judgment
$\mathcal{M},\ x_0 \in \{V_a + c\,V_d \le 1 - p\} \vDash
\mathsf{SProbGE}\bigl(p,\ \mathsf{PUntil}(T, \varphi_s, \varphi)\bigr)$,
its region the sublevel set $\{V_a + c\,V_d \le 1-p\}$.
\end{bprop}
\begin{proof}
On the first-hit event: the landing state is in $R$, hence satisfies
$\varphi$; every earlier state avoided $\mathit{Stop} \supseteq A$,
hence satisfies $\varphi_s$ by the contrapositive of the safe-side tie.
Monotonicity of the (outer) measure; neither satisfaction predicate is
required measurable.
\end{proof}

\begin{bremark}[two floors doing different work]\label{rem:qutb-floors}
$V_a$\textquotesingle s floor is on $A$ only: the avoid certificate is
silent on $R \smallsetminus A$, so a landing there costs no
certificate budget --- the rule does not charge for winning (a state of
$A \cap R$, when the ties permit one, is a successful landing that the
avoid floor nonetheless prices as avoid-territory; the free wins are
exactly the $R \smallsetminus A$ ones). Demand the
floor on all of $R$ too and the concluded region shrinks pointlessly:
every
reachable success state would need $V_a \ge 1$, pushing
$V_a(x) + c V_d(x)$ past any useful threshold. $V_d$\textquotesingle s
floor is on all of $\bar D$ off the target, exactly as its parent
deadline battery
requires --- weaken \emph{it} and \autoref{rem:qdb-floor}\textquotesingle s
counterexample (the frozen chain certified to arrive) imports after one
padding: take $A = \varnothing$ and $V_a \equiv 0$ (all avoid rows
trivially satisfied), and the until node \emph{is} its deadline slice,
frozen chain and all.
The asymmetry is not an optimization; it is the two
slices\textquotesingle{} different jobs made visible in their price
lists.
\end{bremark}

\begin{bremark}[fail-closed, and the rate tie]\label{rem:qutb-failclosed}
$D^{\mathrm{c}}$ joins the stop set on the \emph{bad} side: leaving the
modelled domain before reaching the target counts as failure, while an
$R$-landing succeeds wherever it lands (the target tie is global). Drop
the rate tie $\beta^{T} \le c$ and step (ii)\textquotesingle s budget
transfer fails --- the affine atom would under-budget the
\emph{certified} miss bound by the gap $(\beta^{T} - c)\, V_d(x)$
(the actual miss mass sits somewhere below the bound; it is the
transfer, not the mass, that the tie protects). The catalogued instance
$c = q^{-1}$ with $0 < q \le e^{\theta T}$, $\beta = e^{-\theta}$
satisfies the tie, and at this profile the whole certificate is exact rational
arithmetic --- no transcendental leaf anywhere; one becomes available
through the deadline battery\textquotesingle s Taylor-leaf pattern
(Prop.~\ref{prop:qdb-catalogue}) when $\theta$-rates are wanted.
\end{bremark}

\begin{bremark}[tier, argued]\label{rem:qutb-tier}
Admissible, not alias: three premise sub-derivations (two battery
slices and a union bound), not a respelling of one. And deliberately
not monolithic core: a third self-contained soundness lemma would
duplicate two batteries\textquotesingle{} arguments line by line ---
the assembled proof instead consumes the deadline theorem as a black
box and the union bound as one application, which is this
catalogue\textquotesingle s admissibility discipline enforced at the
Lean level. The concluded region is the affine atom with rows
$[(1, V_a), (c, V_d)]$ and threshold $1 - p$ --- the
region-vocabulary generalization whose $k = 1$ case is the ordinary
sublevel atom.
\end{bremark}

\begin{bremark}[\textbf{profile} --- what is proven, at which
presentations]
\label{rem:qutb-profile}
Proven: the discrete-kernel profile, sorry-free
(\texttt{sound\_QUntilBattery\_firstHit},
\texttt{sound\_QUntilBattery}), with both slice lemmas and the union
bound as named constituents. Also proven: the continuous-time profile
at the martingale-problem presentation
(\texttt{sound\_GSHS\_until\_mp},
\texttt{sound\_GSHS\_until\_battery\_mp}, module
\texttt{QtslAux.LayerD.MPDischarge} of the semantic-universe
strand\textquotesingle s auxiliary package) --- the recipe above is
profile-independent in \emph{shape}, and both
parents\textquotesingle{} formerly filed constituents (the safety
slice\textquotesingle s stopping theory, the deadline
slice\textquotesingle s discounted supermartingale) are supplied by
the same martingale-problem discharge that serves the parents,
composite stop set included. The martingale-problem theorems are
hash-bound by this catalogue\textquotesingle s binder via the
qtsl-aux package registration, family-wide, so the dashboard row
carries no conditional mark --- each named profile is proven, and
its theorems are bound.
\end{bremark}

\begin{anchors}
Recipe \ref{rec:qutb}, slice (i) & \texttt{measure\_firstBad\_le} (Ville via \texttt{Supermartingale.ville\_ineq} on the stopped process; floor asked on the bad-exit set only)\\
Recipe \ref{rec:qutb}, slice (ii) & \texttt{measure\_missTarget\_le} (black-boxing \texttt{sound\_QDeadline\_hit}, complemented once)\\
Recipe \ref{rec:qutb}, glue (iii) & \texttt{sound\_QUnionBound\_ofReal} on \texttt{compl\_firstExit\_subset}\\
Proposition \ref{prop:qutb-firsthit}: the assembly & \texttt{sound\_QUntilBattery\_firstHit}\\
Proposition \ref{prop:qutb-sound}: formula ties & \texttt{sound\_QUntilBattery}\\
the concluded affine atom & region rows $[(1, V_a), (c, V_d)]$, threshold $1-p$; rate-tie hypothesis \texttt{hc} : $\beta^{T} \le c$\\
\end{anchors}


\rulesection{44}{QStreettBattery}{core rule}
\label{blcatrule-QStreettBattery}
\ruleblock{The rule (Steele metanotation)}
\sourcenote{as realized in the mechanized development; source row: \texttt{\texttt{QTSL-rules.toml}} --- \emph{QStreettBattery}. The metanotation statement is typeset by hand for notational quality; its mechanized ground truth is the Lean statement below, extracted from the sources at every build of this book and hash-bound. The rule \emph{roster} is machine-checked (every mechanized rule tag round-trips to exactly one catalogue row); the statement \emph{text} is editorial, checked by eye against the extracted Lean beside it.}
\begin{mathpar}
\inferrule*{
  V_0,\ (W_i)_{i \le k},\ (V_i)_{i \in A} : \mathsf{QCertFam}; \quad A \subseteq \{1, \ldots, k\} \ \text{(the pairs posting absorbing certificates)}; \quad \varepsilon_0 \in \Q_{>0}; \quad (\gamma_i)_{i \in A} \in \Q_{>0} \\
  \text{the \rulename{QSafetyBattery} rows for } V_0 \ \text{at } U := \mathrm{cmpl}(I) \ \text{(interior, guard, floor on } U \cap \bar{D}\text{, nonneg on } \bar{D}\text{)} \\
  \bigl\{S \smallsetminus [\![\rho_i]\!] \vdash \mathcal{A}_m V_{i,m} \le 0\bigr\}_{i \in A} \ \text{(absorbing drift: off-request, on all of } S\text{)} \\
  \bigl\{S \vdash 0 \le V_i \le C^{a}_i\bigr\}_{i \in A} \ \text{(absorbing bounds, on all of } S\text{)}, \quad \bigl\{[\![\rho_i]\!] \vdash V_i \ge 1\bigr\}_{i \in A} \ \text{(absorbing floors)} \\
  \bigl\{([\![\varphi_i]\!] \cup J_i)^{\complement} \vdash \mathcal{A}_m W_{i,m} \le -\varepsilon_0\bigr\}_{i \le k} \ \text{(ranking decay)}, \quad J_i := \{V_i \le 1 - \gamma_i\} \cap (I \smallsetminus [\![\rho_i]\!]) \ \text{for } i \in A; \quad J_i := \emptyset \ \text{for } i \notin A \\
  W_i \equiv 0 \ \text{on target cells (containment)}; \quad \bigl\{\bar{D} \vdash W_i \ge 0\bigr\}_{i \le k} \ \text{(ranking nonnegs)} \\
  \bigl\{\bar{D} \vdash W_i \le C_i\bigr\}_{i \le k}, \quad C_i \in \Q \\
  p + y \le 1 \ \text{(threshold tie)}
}{
  \mathcal{M},\ x_0 \in \{V_0 \le y\} \vDash \mathsf{SProbGE}\bigl(p,\ \mathsf{PDwellStreett}(\mathsf{PairBase}(\rho_1, \varphi_1), \ldots, \mathsf{PairBase}(\rho_k, \varphi_k))\bigr)
}
\end{mathpar}

\ruleflavor{a region written left of the turnstile scopes the row ($R \vdash P$: $P$ at every point of $R$); the certificate families enter as certificate data --- shared references with no claim content; the rows impose every claim. $J_i$ is \emph{defined, not supplied}: the mixed soundness theorem defines it from the posted absorbing data, with the $J_i := \emptyset$ case (absorbing family omitted, $i \notin A$) the conservative scope --- the baked node carries no absorbing-family fields yet, so this optional face is not yet wire-representable (a named schema divergence). $V_0$\textquotesingle s own nonnegativity rides the \rulename{QSafetyBattery} kit row, stated once there. The margin premise $\gamma_i > 0$ together with the absorbing floor forces $J_i \cap [\![\rho_i]\!] = \emptyset$; no disjointness row is asked. The target-cell containment $W_i \equiv 0$ is recomputed by the checker. The threshold tie $p + y \le 1$ is the soundness theorem\textquotesingle s own hypothesis, displayed as a side condition: the baked Streett node carries no tie proof field yet (contrast the deadline and until nodes) --- a named schema divergence, not a checked leaf. The conclusion is the region-conditioned judgment in constructor vocabulary, as at \rulename{QSafetyBattery}; the base-exact request consumption is the $\mathsf{PairBase}$ constructor.}
\par\medskip
\rulename{QStreettBattery} concludes a probabilistic Streett-acceptance
judgment: with probability at least $p$, a path of the modelled system
satisfies every pair of the Streett condition
$\mathsf{PDwellStreett}(\mathsf{PairBase}(\rho_1, \varphi_1), \ldots,
\mathsf{PairBase}(\rho_k, \varphi_k))$ (surface syntax
$\mathsf{Streett}((\rho_1, \varphi_1), \ldots, (\rho_k, \varphi_k))$) ---
for each pair $i$, if \emph{request} states (those satisfying $\rho_i$)
occur infinitely often along the path, then \emph{response} states
(those satisfying $\varphi_i$) also occur infinitely often. This is the
shape of liveness-under-fairness obligations, and it subsumes recurrence
and stability claims as degenerate pairs (the sub-rule sections that
follow this one). A producer discharges such a specification by
synthesizing the certificate family
$(V_0;\ (V_i, \gamma_i, W_i)_{i \le k};\ \varepsilon_0)$: one safety
certificate $V_0$ confining probability mass to the operating region,
and, per pair, a ranking certificate $W_i$ that pays out a fixed rational
$\varepsilon_0 > 0$ at every un-responded request --- plus, optionally,
absorbing certificates $(V_i, \gamma_i)$ licensing the sharper ranking
scope discussed in the proof. Every premise row is an inherited
base-calculus sequent (a Bernstein-cone-family polynomial inequality the
base kernel replays), so the bundle adds \emph{zero} kernel arithmetic:
its entire content is the supermartingale assembly one judgment level
up, which is where its soundness lemma lives.

What the checker recomputes rather than trusts: the target-cell
containment behind the $W_i \equiv 0$ convention (the collapse
side-condition of the \rulename{QStreettTargetCellCollapse} section), and
--- when the absorbing face becomes wire-representable --- the carved-out
set $J_i := \{V_i \le 1 - \gamma_i\} \cap
(I \smallsetminus [\![\rho_i]\!])$, $I$ the display's confinement
region $\{V_0 < 1\}$: $J_i$ is \emph{defined}, never supplied by the
producer, but today it is defined by the soundness theorem from posted
theorem data --- the baked node carries no absorbing-family fields, so no
enacted recomputation exists yet. Also recomputed: the
region-hull mass side-conditions that make the certificate scopes
fail-closed. The threshold tie $p + y \le 1$ is displayed as a side
condition, not a checked leaf: it is the soundness
theorem\textquotesingle s own hypothesis, and the baked Streett node
carries no tie proof field yet (contrast the deadline and until
nodes) --- a named schema divergence, as the rule\textquotesingle s
flavor text states. What would break without the hypotheses: the nonnegativity
rows for all three certificate families are load-bearing (a signed
certificate can subsidize probability-mass leakage --- the recorded
two-branch counterexample at the quantitative-safety rule); the
strictness $\varepsilon_0 > 0$ is load-bearing ($W \equiv 1$ satisfies
every row at $\varepsilon_0 = 0$ while no response is ever forced); the
floor and hull rows pair up fail-closed. Each is remarked, with its
counterexample, in the proof.

Scope, stated plainly: the statements bound and receipt-credited in
this edition are the five
excursion-cluster theorems quoted below --- the analytic engine of the
typeset proof --- together with the end-to-end
assembly \texttt{sound\_QStreettBattery} (the conservative-scope bundle
theorem over \texttt{Kernel.traj}) and its discharged carved-scope
corollary \texttt{sound\_QStreettBattery\_carved\_discharged} (absorbing
certificates supplied as kernel-level rows, composed through the
absorption discharge), both in the auxiliary module
\texttt{BernsteinLogicAux.QStreett}. These cover the rule at the
\emph{discrete-kernel profile} (a time-homogeneous Markov chain
presented by a single step kernel); the complement-request tie
\texttt{sound\_QStreettBattery\_tie} is likewise bound in this edition,
under the \rulename{QStreettCmplRequest} entry that carries its
event-level kernel; the parameterized
\texttt{sound\_QStreettBattery\_carved} and the absorption discharge
\texttt{absorptionBound\_of\_kernel} remain machine-checked and cited,
not bound. The continuous-time transport that used to sit beyond that
--- the residue this rule shared with its safety and deadline
siblings, and this entry\textquotesingle s former conditional mark ---
is discharged and bound: at the martingale-problem profile the whole
bundle is proven, confinement and zero-decrement ranking rows
(\texttt{sound\_GSHS\_streett\_bundle\_mp}) together with the
strict-decrement rows from strict presentations
(\texttt{sound\_GSHS\_streett\_bundle\_mp\_strict}, whose ranking-row
input is produced by \texttt{rankingSegmentDynkin\_of\_mp\_strict}
from the pathwise progress row of \texttt{MPFlowStrict}), leaving no
analytic per-model hypothesis in the bundle. All four are hash-bound
below via the registered qtsl-aux package, so the dashboard row
carries no conditional mark: each named profile is proven, and its
theorems are bound.

\cataloguenote{Two design notes travel with the rule. First, schema: the bundle
rule\textquotesingle s rows are baked in the extension theory
(per-pair rank certificates and nonneg proof lists); the source
catalogue\textquotesingle s old MISSING-ROW finding is discharged. Second, why Streett is the
\emph{primitive} acceptance condition here: a B\"uchi condition is
exactly a one-pair Streett condition with the trivially-true request,
and the degenerate-pair specialization erases the resulting tautological
rows --- so a Streett-primitive compile pays no premium on B\"uchi-shaped
workloads, while a B\"uchi-primitive compile would pay an
exponential-in-$k$ automaton re-encoding on Streett-shaped ones. The
sub-rules of the following sections (degenerate request, complement
request, target-cell collapse) are the compile-level faces of this
design.
}
\ruleblock{Lean statement}
\declhead{BernsteinLogicAux.Excursions}{Supermartingale.exists\_ae\_tendsto\_of\_nonneg}{\badge{BadgeGreen}{statement bound}}
\begin{leansrc}
/-- **Almost-everywhere convergence of nonnegative supermartingales** (K5a): a
nonnegative supermartingale is automatically bounded in `L¹` (its expectation is
nonincreasing), so Doob's convergence theorem applies with no further hypothesis. Dual
of `MeasureTheory.Submartingale.exists_ae_tendsto_of_bdd` via `Supermartingale.neg`. -/
theorem Supermartingale.exists_ae_tendsto_of_nonneg [IsFiniteMeasure μ] {f : ℕ → Ω → ℝ}
    (hf : Supermartingale f 𝒢 μ) (hnonneg : 0 ≤ f) :
    ∀ᵐ ω ∂μ, ∃ c, Tendsto (fun n => f n ω) atTop (𝓝 c)
\end{leansrc}
\declhead{BernsteinLogicAux.Excursions}{ranking\_hitting\_sum\_le}{\badge{BadgeGreen}{statement bound}}
\begin{leansrc}
/-- **Ranking hitting-time bound** (K5b, layer form): `ε₀ · 𝔼[τ_B] ≤ 𝔼[V 0]`, the
expected hitting time decomposed by the layer-cake formula as
`𝔼[τ_B] = ∑' k, μ {τ_B > k}` over the strict-survival events
`{τ_B > k} = {∀ j < k + 1, X j ∉ B}`. This is the certificate-facing form: the ranking
rows certify `hdrift`, the nonneg rows certify `hnonneg`, and the conclusion is a bound
on the expected number of steps before reaching `B`. -/
theorem ranking_hitting_sum_le (hε : 0 ≤ ε₀)
    (hint : ∀ n, Integrable (V n) μ) (hnonneg : 0 ≤ V)
    (hmB : ∀ n, MeasurableSet[𝒢 n] {ω | X n ω ∈ B})
    (hdrift : ∀ n, ∀ᵐ ω ∂μ, X n ω ∉ B → (μ[V (n + 1) | 𝒢 n]) ω ≤ V n ω - ε₀) :
    ENNReal.ofReal ε₀ * ∑' k, μ {ω | ∀ j < k + 1, X j ω ∉ B}
      ≤ ENNReal.ofReal (∫ ω, V 0 ω ∂μ)
\end{leansrc}
\declhead{BernsteinLogicAux.Excursions}{ae\_hits\_of\_ranking}{\badge{BadgeGreen}{statement bound}}
\begin{leansrc}
/-- **Almost-sure hitting from a ranking certificate** (K5b, qualitative form): a
strictly positive ranking drift off `B` forces the state process to reach `B` almost
surely — the summability of the survival probabilities from `ranking_hitting_sum_le`
sends them to zero. -/
theorem ae_hits_of_ranking (hε : 0 < ε₀)
    (hint : ∀ n, Integrable (V n) μ) (hnonneg : 0 ≤ V)
    (hmB : ∀ n, MeasurableSet[𝒢 n] {ω | X n ω ∈ B})
    (hdrift : ∀ n, ∀ᵐ ω ∂μ, X n ω ∉ B → (μ[V (n + 1) | 𝒢 n]) ω ≤ V n ω - ε₀) :
    ∀ᵐ ω ∂μ, ∃ n, X n ω ∈ B
\end{leansrc}
\declhead{BernsteinLogicAux.Excursions}{streett\_window\_sum\_le}{\badge{BadgeGreen}{statement bound}}
\begin{leansrc}
/-- **Summability of un-responded requests** (K5c, quantitative core): for each window
start `N`, `ε₀ · ∑_{j} μ(no response in [N, N+j] ∧ request at N+j) ≤ 𝔼[V N]` — the form
the Borel–Cantelli lemma consumes. -/
theorem streett_window_sum_le (hε : 0 ≤ ε₀)
    (hint : ∀ n, Integrable (V n) μ) (hnonneg : 0 ≤ V)
    (hmR : ∀ n, MeasurableSet[𝒢 n] {ω | X n ω ∈ R})
    (hmG : ∀ n, MeasurableSet[𝒢 n] {ω | X n ω ∈ G})
    (hdrift : ∀ n, ∀ᵐ ω ∂μ, X n ω ∉ G → (μ[V (n + 1) | 𝒢 n]) ω ≤ V n ω)
    (hdriftR : ∀ n, ∀ᵐ ω ∂μ,
      X n ω ∈ R → X n ω ∉ G → (μ[V (n + 1) | 𝒢 n]) ω ≤ V n ω - ε₀)
    (N : ℕ) :
    ENNReal.ofReal ε₀ * ∑' j : ℕ,
        μ ({ω | ∀ k, N ≤ k → k < N + j + 1 → X k ω ∉ G} ∩ {ω | X (N + j) ω ∈ R})
      ≤ ENNReal.ofReal (∫ ω, V N ω ∂μ)
\end{leansrc}
\declhead{BernsteinLogicAux.Excursions}{streett\_ae\_frequently\_response}{\badge{BadgeGreen}{statement bound}}
\begin{leansrc}
/-- **Streett excursion assembly** (K5c): if `V` is nonincreasing in conditional
expectation off the response set `G` and decreases by `ε₀ > 0` on requests (`R`, off
`G`), then almost surely infinitely many requests imply infinitely many responses —
the probabilistic content of a Streett pair `(R, G)`. Proof: by
`streett_window_sum_le` the un-responded-request events of each window are summable, so
their limsup is null (`MeasureTheory.measure_limsup_atTop_eq_zero`, Borel–Cantelli);
a path with infinitely many requests and no response after `N` lies in window `N`'s
limsup, and the union over `N` is null. -/
theorem streett_ae_frequently_response (hε : 0 < ε₀)
    (hint : ∀ n, Integrable (V n) μ) (hnonneg : 0 ≤ V)
    (hmR : ∀ n, MeasurableSet[𝒢 n] {ω | X n ω ∈ R})
    (hmG : ∀ n, MeasurableSet[𝒢 n] {ω | X n ω ∈ G})
    (hdrift : ∀ n, ∀ᵐ ω ∂μ, X n ω ∉ G → (μ[V (n + 1) | 𝒢 n]) ω ≤ V n ω)
    (hdriftR : ∀ n, ∀ᵐ ω ∂μ,
      X n ω ∈ R → X n ω ∉ G → (μ[V (n + 1) | 𝒢 n]) ω ≤ V n ω - ε₀) :
    ∀ᵐ ω ∂μ, (∃ᶠ n in atTop, X n ω ∈ R) → ∃ᶠ n in atTop, X n ω ∈ G
\end{leansrc}
\declhead{BernsteinLogicAux.QStreett}{sound\_QStreettBattery}{\badge{BadgeGreen}{statement bound}}
\begin{leansrc}
/-- **`sound_QStreettBattery`, discrete profile** (catalogue T9, rev 3
`thm:bundle-sound` transported; the assembly the receipt ledger carried as
FILED). Hypotheses: the `V₀` battery rows — quantitative safety of the
confinement region `D \ U` (the catalogue's `I := {V₀ < 1}` is the instance
`U := {y | 1 ≤ V₀ y}`, whose floor row is tautological) — and, per pair
`i : ι` (countable; a Streett bundle's finite `k` is an instance), the ranking
rows for `Wᵢ` at the **conservative scope** of the design note §4.1's ranking
row ("conservative on `I \ ⟦φᵢ⟧`"): nonincrease off the response region,
`εᵢ`-decrease at requests, on the confinement region. Conclusion:
`Pr[⋂ᵢ (Fin(Rqᵢ) ∪ Inf(Gᵢ))] ≥ p` from the start region `D ∩ {V₀ ≤ thr}` with
`p + thr ≤ 1` — the compile's atom `{V₀ ≤ 1 − p}` is `thr = 1 − p`.

The absorbing certificates `(Vᵢ, γᵢ)` of the full bundle datum do not appear:
at the conservative ranking scope they are not consumed (they license only the
`Jᵢ` carve-out — see `sound_QStreettBattery_carved`). Presentation→semantics
bridge and GSHS residue: module docstring. -/
theorem sound_QStreettBattery {Rq G : ι → Set S} {W : ι → S → ℝ} {ε C : ι → ℝ}
    (hD : MeasurableSet D) (hDbar : MeasurableSet Dbar) (hU : MeasurableSet U)
    (hV₀ : Measurable V₀)
    (hRq : ∀ i, MeasurableSet (Rq i)) (hG : ∀ i, MeasurableSet (G i))
    (hW : ∀ i, Measurable (W i))
    (hΘ : p + thr ≤ 1)
    (hDsub : D ⊆ Dbar)
    (hR3 : ∀ y ∈ D, y ∉ U → K y Dbar = 1)
    (hR4 : ∀ y ∈ D, y ∉ U → K y (D ∪ U) = 1)
    (hV₀int : ∀ y ∈ D, y ∉ U → Integrable V₀ (K y))
    (hV₀drift : ∀ y ∈ D, y ∉ U → ∫ z, V₀ z ∂K y ≤ V₀ y)
    (hV₀floor : ∀ y ∈ U, y ∈ Dbar → 1 ≤ V₀ y)
    (hV₀nonneg : ∀ y ∈ Dbar, 0 ≤ V₀ y)
    (hε : ∀ i, 0 < ε i)
    (hW0 : ∀ i, ∀ y ∈ Dbar, 0 ≤ W i y) (hWC : ∀ i, ∀ y ∈ Dbar, W i y ≤ C i)
    (hWint : ∀ i, ∀ y ∈ D, y ∉ U → y ∉ G i → Integrable (W i) (K y))
    (hWdrift : ∀ i, ∀ y ∈ D, y ∉ U → y ∉ G i → ∫ z, W i z ∂K y ≤ W i y)
    (hWdriftR : ∀ i, ∀ y ∈ D, y ∉ U → y ∈ Rq i → y ∉ G i →
      ∫ z, W i z ∂K y ≤ W i y - ε i)
    (hx : x ∈ D) (hthr : V₀ x ≤ thr) :
    ENNReal.ofReal p ≤ pathMeasure K x
      {ω | ∀ i, (∃ᶠ n in atTop, ω n ∈ Rq i) → ∃ᶠ n in atTop, ω n ∈ G i}
\end{leansrc}
\declhead{BernsteinLogicAux.QStreett}{sound\_QStreettBattery\_carved\_discharged}{\badge{BadgeGreen}{statement bound}}
\begin{leansrc}
/-- **`sound_QStreettBattery_carved`, discharged form** — the unconditional
corollary: the full bundle datum with the absorbing certificates `(Vᵢ, γᵢ)`
supplied as kernel-level rows (global scope, per the section docstring)
instead of the named hypothesis. The parameterized
`sound_QStreettBattery_carved` remains the citable assembly; this composes it
with `absorptionBound_of_kernel` per pair. -/
theorem sound_QStreettBattery_carved_discharged {D Dbar U : Set S} {V₀ : S → ℝ}
    {p thr : ℝ} {x : S} {ι : Type*} [Countable ι] {Rq G J : ι → Set S}
    {W Va : ι → S → ℝ} {ε C γ Ca : ι → ℝ}
    (hD : MeasurableSet D) (hDbar : MeasurableSet Dbar) (hU : MeasurableSet U)
    (hV₀ : Measurable V₀)
    (hRq : ∀ i, MeasurableSet (Rq i)) (hG : ∀ i, MeasurableSet (G i))
    (hJ : ∀ i, MeasurableSet (J i)) (hW : ∀ i, Measurable (W i))
    (hΘ : p + thr ≤ 1)
    (hDsub : D ⊆ Dbar)
    (hR3 : ∀ y ∈ D, y ∉ U → K y Dbar = 1)
    (hR4 : ∀ y ∈ D, y ∉ U → K y (D ∪ U) = 1)
    (hV₀int : ∀ y ∈ D, y ∉ U → Integrable V₀ (K y))
    (hV₀drift : ∀ y ∈ D, y ∉ U → ∫ z, V₀ z ∂K y ≤ V₀ y)
    (hV₀floor : ∀ y ∈ U, y ∈ Dbar → 1 ≤ V₀ y)
    (hV₀nonneg : ∀ y ∈ Dbar, 0 ≤ V₀ y)
    (hε : ∀ i, 0 < ε i) (hγ : ∀ i, 0 < γ i)
    (hW0 : ∀ i, ∀ y ∈ Dbar, 0 ≤ W i y) (hWC : ∀ i, ∀ y ∈ Dbar, W i y ≤ C i)
    (hWint : ∀ i, ∀ y ∈ D, y ∉ U → y ∉ G i → y ∉ J i → Integrable (W i) (K y))
    (hWdrift : ∀ i, ∀ y ∈ D, y ∉ U → y ∉ G i → y ∉ J i →
      ∫ z, W i z ∂K y ≤ W i y)
    (hWdriftR : ∀ i, ∀ y ∈ D, y ∉ U → y ∈ Rq i → y ∉ G i → y ∉ J i →
      ∫ z, W i z ∂K y ≤ W i y - ε i)
    (hVa : ∀ i, Measurable (Va i))
    (hVa0 : ∀ i y, 0 ≤ Va i y) (hVaC : ∀ i y, Va i y ≤ Ca i)
    (hVafloor : ∀ i, ∀ y ∈ Rq i, 1 ≤ Va i y)
    (hVadrift : ∀ i y, y ∉ Rq i → ∫ z, Va i z ∂K y ≤ Va i y)
    (hVaJ : ∀ i, ∀ y ∈ J i, Va i y ≤ 1 - γ i)
    (hx : x ∈ D) (hthr : V₀ x ≤ thr) :
    ENNReal.ofReal p ≤ pathMeasure K x
      {ω | ∀ i, (∃ᶠ n in atTop, ω n ∈ Rq i) → ∃ᶠ n in atTop, ω n ∈ G i}
\end{leansrc}
\par\medskip
\textbf{The general-stochastic instance, assembled.}\enspace The
bundle on the segment model\textquotesingle s plain skeleton chain,
obtained by instantiating the discrete bundle above at the skeleton
kernel: alive-only region transport, scoped per-pair ranking fields
with additive decrement (subtraction-free extended-nonnegative form),
and the three fail-closed containment rows enter as per-model fields;
the discrete engine\textquotesingle s localization and windowed
telescoping do the rest, restart-free. Plain and confined forms.

\declhead{BernsteinLogicAux.LayerD.Streett}{sound\_GSHS\_streett\_bundle}{\badge{BadgeGreen}{statement bound}}
\begin{leansrc}
/-- **`sound_QStreettBattery`, GSHS instance** (T9, conservative ranking
scope) — the discrete-statement-shaped corollary: the chain-level Streett
event alone. The `_confined` form is the semantically loaded one; this one
matches `QStreett.sound_QStreettBattery` line for line. -/
theorem sound_GSHS_streett_bundle
    [IsMarkovKernel M.seg] [IsMarkovKernel M.jump]
    (hD : MeasurableSet D) (hDbar : MeasurableSet Dbar) (hU : MeasurableSet U)
    (hV₀ : Measurable V₀)
    (hRq : ∀ i, MeasurableSet (Rq i)) (hG : ∀ i, MeasurableSet (G i))
    (hW : ∀ i, Measurable (W i))
    (hΘ : p + thr ≤ 1) (hDsub : D ⊆ Dbar)
    (hLive : NoVanish M (D \ U)) (hHalt : HaltsCaught M (D \ U) (U ∩ Dbar))
    (hReset : ResetsInto M (D \ U) ((D ∪ U) ∩ Dbar))
    (hSegD₀ : RankingSegmentDynkin M (D \ U) V₀ 0)
    (hForce₀ : ForcedNonexpansion M V₀) (hNonneg₀ : NonnegV V₀)
    (hC₀0 : 0 ≤ C₀) (hV₀C : ∀ y, V₀ y ≤ C₀)
    (hFloor₀ : ∀ y ∈ U, y ∈ Dbar → 1 ≤ V₀ y)
    (hε : ∀ i, 0 < ε i)
    (hSegDW : ∀ i, RankingSegmentDynkin M ((D \ U) \ G i) (W i) 0)
    (hSegDWR : ∀ i, RankingSegmentDynkin M (((D \ U) ∩ Rq i) \ G i) (W i) (ε i))
    (hForceW : ∀ i, ForcedNonexpansion M (W i)) (hNonnegW : ∀ i, NonnegV (W i))
    (hC0 : ∀ i, 0 ≤ C i) (hWC : ∀ i, ∀ y, W i y ≤ C i)
    (hx : x ∈ D) (hthr : V₀ x ≤ thr) :
    ENNReal.ofReal p ≤ Battery.pathMeasure (skel M) (SkelState.alive x)
      {ω | ∀ i, (∃ᶠ n in atTop, ω n ∈ skelIn (Rq i)) →
        ∃ᶠ n in atTop, ω n ∈ skelIn (G i)}
\end{leansrc}
\declhead{BernsteinLogicAux.LayerD.Streett}{sound\_GSHS\_streett\_bundle\_confined}{\badge{BadgeGreen}{statement bound}}
\begin{leansrc}
/-- **`sound_QStreettBattery`, GSHS instance, confined form** (T9; assault
doc §2.7 Theorem G at the conservative ranking scope, mechanized
restart-free — module docstring for why G1/G2's epoch machinery is not
consumed). Hypothesis fields are segment-level and per-model: the three
containment rows, the θ = 0 confinement field for `V₀`, and per pair the
scoped ranking fields for `W i` (nonincrease off the response region,
`ε i`-decrease at request entries — per-step shape, module docstring for
the rate-form conversion). Conclusion: with probability at least `p` the
skeleton chain of `Rep(P)` stays alive in `D \ U` forever **and** every
certified pair satisfies the chain-level Streett condition. Semantic
transport of the chain events: module docstring (T9b monotonicity for
requests, entry reads for responses, W-Z for time-cofinality). -/
theorem sound_GSHS_streett_bundle_confined
    [IsMarkovKernel M.seg] [IsMarkovKernel M.jump]
    (hD : MeasurableSet D) (hDbar : MeasurableSet Dbar) (hU : MeasurableSet U)
    (hV₀ : Measurable V₀)
    (hRq : ∀ i, MeasurableSet (Rq i)) (hG : ∀ i, MeasurableSet (G i))
    (hW : ∀ i, Measurable (W i))
    (hΘ : p + thr ≤ 1) (hDsub : D ⊆ Dbar)
    (hLive : NoVanish M (D \ U)) (hHalt : HaltsCaught M (D \ U) (U ∩ Dbar))
    (hReset : ResetsInto M (D \ U) ((D ∪ U) ∩ Dbar))
    (hSegD₀ : RankingSegmentDynkin M (D \ U) V₀ 0)
    (hForce₀ : ForcedNonexpansion M V₀) (hNonneg₀ : NonnegV V₀)
    (hC₀0 : 0 ≤ C₀) (hV₀C : ∀ y, V₀ y ≤ C₀)
    (hFloor₀ : ∀ y ∈ U, y ∈ Dbar → 1 ≤ V₀ y)
    (hε : ∀ i, 0 < ε i)
    (hSegDW : ∀ i, RankingSegmentDynkin M ((D \ U) \ G i) (W i) 0)
    (hSegDWR : ∀ i, RankingSegmentDynkin M (((D \ U) ∩ Rq i) \ G i) (W i) (ε i))
    (hForceW : ∀ i, ForcedNonexpansion M (W i)) (hNonnegW : ∀ i, NonnegV (W i))
    (hC0 : ∀ i, 0 ≤ C i) (hWC : ∀ i, ∀ y, W i y ≤ C i)
    (hx : x ∈ D) (hthr : V₀ x ≤ thr) :
    ENNReal.ofReal p ≤ Battery.pathMeasure (skel M) (SkelState.alive x)
      {ω | (∀ n, ω n ∈ skelIn D \ skelBad U) ∧
        ∀ i, (∃ᶠ n in atTop, ω n ∈ skelIn (Rq i)) →
          ∃ᶠ n in atTop, ω n ∈ skelIn (G i)}
\end{leansrc}
\par\medskip
\textbf{The martingale-problem profile, end to end.}\enspace The
general continuous-time theorems: the inter-jump motion enters through
its generator --- $\mathcal{A}V$ as data, tied to the process only by
the martingale hypothesis for the discounted compensated certificate
process. The mixed-profile bundle discharges the confinement half and
every zero-decrement ranking row; \texttt{MPFlowStrict} adds the one
pathwise progress row (accumulated generator drift reaching
$-\varepsilon$ by a finite sample time), from which the strict
ranking rows are derived at exactly the requested-but-ungranted
states --- and the strict bundle theorem then carries no analytic
per-model hypothesis at all.

\declhead{QtslAux.LayerD.MPDischarge}{sound\_GSHS\_streett\_bundle\_mp}{\badge{BadgeGreen}{statement bound}}
\begin{leansrc}
/-- The Streett/recurrence bundle at the martingale-problem profile,
mixed ranking scope: the confinement certificate `V₀` and every pair's
ranking certificate `W i` are presented by martingale problems, which
discharges all the `ε = 0` ranking rows. The strict-decrement rows
`hSegDWR` (`ε i > 0`) stay hypotheses *of this theorem*, for models that
discharge them by other means; `sound_GSHS_streett_bundle_mp_strict` below
discharges them too, from strict presentations (`MPFlowStrict`) on the
requested-but-ungranted region. -/
theorem sound_GSHS_streett_bundle_mp
    {ι : Type*} [Countable ι]
    {D Dbar U : Set State} {V₀ : State → ℝ} {Rq G : ι → Set State}
    {W : ι → State → ℝ} {p thr C₀ : ℝ} {ε C : ι → ℝ} {x : State}
    [IsMarkovKernel M.seg] [IsMarkovKernel M.jump]
    (hD : MeasurableSet D) (hDbar : MeasurableSet Dbar) (hU : MeasurableSet U)
    (hV₀ : Measurable V₀)
    (hRq : ∀ i, MeasurableSet (Rq i)) (hG : ∀ i, MeasurableSet (G i))
    (hW : ∀ i, Measurable (W i))
    (hΘ : p + thr ≤ 1) (hDsub : D ⊆ Dbar)
    (hLive : NoVanish M (D \ U)) (hHalt : HaltsCaught M (D \ U) (U ∩ Dbar))
    (hReset : ResetsInto M (D \ U) ((D ∪ U) ∩ Dbar))
    (hmp₀ : ∀ x : State, Nonempty (MPFlow M 0 V₀ x))
    (hForce₀ : ForcedNonexpansion M V₀) (hNonneg₀ : NonnegV V₀)
    (hC₀0 : 0 ≤ C₀) (hV₀C : ∀ y, V₀ y ≤ C₀)
    (hFloor₀ : ∀ y ∈ U, y ∈ Dbar → 1 ≤ V₀ y)
    (hε : ∀ i, 0 < ε i)
    (hmpW : ∀ i, ∀ x : State, Nonempty (MPFlow M 0 (W i) x))
    (hSegDWR : ∀ i, RankingSegmentDynkin M (((D \ U) ∩ Rq i) \ G i) (W i) (ε i))
    (hForceW : ∀ i, ForcedNonexpansion M (W i)) (hNonnegW : ∀ i, NonnegV (W i))
    (hC0 : ∀ i, 0 ≤ C i) (hWC : ∀ i, ∀ y, W i y ≤ C i)
    (hx : x ∈ D) (hthr : V₀ x ≤ thr) :
    ENNReal.ofReal p ≤ Battery.pathMeasure (skel M) (SkelState.alive x)
      {ω | ∀ i, (∃ᶠ n in atTop, ω n ∈ skelIn (Rq i)) →
        ∃ᶠ n in atTop, ω n ∈ skelIn (G i)}
\end{leansrc}
\declhead{QtslAux.LayerD.MPDischarge}{MPFlowStrict}{\badge{BadgeGreen}{statement bound}}
\begin{leansrc}
/-- **The strict martingale-problem presentation of a segment.** An
`MPFlow` together with the pathwise progress row: along every path, the
accumulated generator drift of the discounted certificate reaches `≤ -ε`
by some finite sample time. For exit paths the natural witness is the exit
time itself (`stopClip` caps the accumulation there), so the row reads
`∫₀^{T ω} mpDrift ≤ -ε`; the intended certificate-side derivation is a
strict drift bound `≤ -δ` on the segment together with a minimum dwell
time `t_min`, giving `ε = δ · t_min`. -/
structure MPFlowStrict (M : SegmentModel State) (θ : ℝ) (V : State → ℝ)
    (x : State) (ε : ℝ) extends MPFlow M θ V x where
  /-- the pathwise progress row -/
  hprog : ∀ ω, ∃ t : ℝ, 0 ≤ t ∧
    ∫ s in (0:ℝ)..(stopClip (T ω) (k ω) t),
      mpDrift θ V γ r jr AV s ω ≤ -ε
\end{leansrc}
\declhead{QtslAux.LayerD.MPDischarge}{rankingSegmentDynkin\_of\_mp\_strict}{\badge{BadgeGreen}{statement bound}}
\begin{leansrc}
/-- **The strict-decrement ranking rows at the martingale-problem
profile.** Entry states of the ranking region carrying strict
presentations at level `ε` satisfy `RankingSegmentDynkin M Sc V ε` — the
row family that `rankingSegmentDynkin_of_mp` could not produce. Note the
quantifier: only states in `Sc` need presentations (the `ε = 0` route
quantified over all states because it factored through `SegmentDynkin`). -/
theorem rankingSegmentDynkin_of_mp_strict (M : SegmentModel State)
    (V : State → ℝ) [IsSFiniteKernel M.jump] (hV : Measurable V)
    {Sc : Set State} {ε : ℝ} (hε : 0 ≤ ε)
    (h : ∀ x ∈ Sc, Nonempty (MPFlowStrict M 0 V x ε)) :
    RankingSegmentDynkin M Sc V ε
\end{leansrc}
\declhead{QtslAux.LayerD.MPDischarge}{sound\_GSHS\_streett\_bundle\_mp\_strict}{\badge{BadgeGreen}{statement bound}}
\begin{leansrc}
/-- The Streett/recurrence bundle at the **fully martingale-problem
profile**: the confinement certificate `V₀` and every pair's ranking
certificate `W i` are presented by martingale problems, with the
requested-but-ungranted region additionally carrying *strict*
presentations (`MPFlowStrict` at level `ε i`). This discharges the
`hSegDWR` hypothesis of `sound_GSHS_streett_bundle_mp` — at this profile
no analytic per-model hypothesis remains. -/
theorem sound_GSHS_streett_bundle_mp_strict
    {ι : Type*} [Countable ι]
    {D Dbar U : Set State} {V₀ : State → ℝ} {Rq G : ι → Set State}
    {W : ι → State → ℝ} {p thr C₀ : ℝ} {ε C : ι → ℝ} {x : State}
    [IsMarkovKernel M.seg] [IsMarkovKernel M.jump]
    (hD : MeasurableSet D) (hDbar : MeasurableSet Dbar) (hU : MeasurableSet U)
    (hV₀ : Measurable V₀)
    (hRq : ∀ i, MeasurableSet (Rq i)) (hG : ∀ i, MeasurableSet (G i))
    (hW : ∀ i, Measurable (W i))
    (hΘ : p + thr ≤ 1) (hDsub : D ⊆ Dbar)
    (hLive : NoVanish M (D \ U)) (hHalt : HaltsCaught M (D \ U) (U ∩ Dbar))
    (hReset : ResetsInto M (D \ U) ((D ∪ U) ∩ Dbar))
    (hmp₀ : ∀ x : State, Nonempty (MPFlow M 0 V₀ x))
    (hForce₀ : ForcedNonexpansion M V₀) (hNonneg₀ : NonnegV V₀)
    (hC₀0 : 0 ≤ C₀) (hV₀C : ∀ y, V₀ y ≤ C₀)
    (hFloor₀ : ∀ y ∈ U, y ∈ Dbar → 1 ≤ V₀ y)
    (hε : ∀ i, 0 < ε i)
    (hmpW : ∀ i, ∀ x : State, Nonempty (MPFlow M 0 (W i) x))
    (hmpWR : ∀ i, ∀ x ∈ ((D \ U) ∩ Rq i) \ G i,
      Nonempty (MPFlowStrict M 0 (W i) x (ε i)))
    (hForceW : ∀ i, ForcedNonexpansion M (W i)) (hNonnegW : ∀ i, NonnegV (W i))
    (hC0 : ∀ i, 0 ≤ C i) (hWC : ∀ i, ∀ y, W i y ≤ C i)
    (hx : x ∈ D) (hthr : V₀ x ≤ thr) :
    ENNReal.ofReal p ≤ Battery.pathMeasure (skel M) (SkelState.alive x)
      {ω | ∀ i, (∃ᶠ n in atTop, ω n ∈ skelIn (Rq i)) →
        ∃ᶠ n in atTop, ω n ∈ skelIn (G i)}
\end{leansrc}
\ruleblock{Correspondence (Steele display $\leftrightarrow$ Lean)}
\badge{BadgeGreen}{reviewer-certified against the displayed texts}
\sourcenote{The dictionary below is a reviewer-certified statement of how the mechanized statements above interpret the displayed rule: display metavariable $\mapsto$ Lean binder, each display hypothesis $\mapsto$ the premise or side condition carrying it, the display conclusion $\mapsto$ the goal --- with any deliberate difference in width between display and mechanization named outright. Certification is recorded as a hash triple over (the displayed cell, the extracted statement texts, this dictionary\textquotesingle s own text); if any of the three changes, this block turns amber until a reviewer re-certifies. The correspondence claim itself remains editorial --- reviewer-grade, not machine-checked; its machine-checked successor types each soundness statement against the rule as a structured object, and supersedes this block per rule as it lands.}
The displayed judgment is concluded by
\texttt{sound\_QStreettBattery} (second block): from the $V_0$ battery
rows on the confinement region and, per pair, the conservative-scope
ranking rows (nonincrease off the response region, $\varepsilon_i$-decay
at requests), it concludes
$\Pr[\bigcap_i (\mathrm{Fin}(Rq_i) \cup \mathrm{Inf}(G_i))] \ge p$
from the start region at \texttt{p + thr ≤ 1} --- the
display\textquotesingle s conclusion at the discrete-kernel profile
over \texttt{Kernel.traj}; \texttt{\_carved\_discharged} adds the
carved scope with the absorbing certificates consumed as kernel-level
rows. (Notation: $\mathrm{Fin}(A)$ / $\mathrm{Inf}(A)$ denote the path
events that a trajectory visits $A$ finitely / infinitely often ---
the acceptance atoms of the Hanoi Omega-Automata
format~\cite{babiak2015hoa} and of Emerson--Lei
conditions~\cite{emersonlei1987}; equivalently $\mathrm{Inf}(A) =
\{\rho \mid \mathrm{Inf}(\rho) \cap A \neq \emptyset\}$ in the
infinity-set convention, $\mathrm{Inf}(\rho)$ being the set of states
visited infinitely often. The upright $\mathrm{Fin}$ is distinct from
the sans-serif type former $\mathsf{Fin}$ of the mechanization.) The five bound excursion theorems are the K5 analytic engine the
assembly consumes, stated over an arbitrary filtered probability space
\texttt{(Ω, 𝒢, μ)} with an adapted state process \texttt{X} ---
\emph{wider} than the discrete-kernel chain. What the
display\textquotesingle s rows reach in the engine\textquotesingle s
bound set, per pair:
\begin{itemize}
\item ranking interior rows $\mathcal{A}_m W_{i,m} \le -\varepsilon_0$
  off $[\![\varphi_i]\!] \cup J_i$ $\mapsto$ the conditional-shadow
  pair in \texttt{streett\_ae\_frequently\_response}: \texttt{hdrift}
  (conditional non-increase off the response set \texttt{G}) and
  \texttt{hdriftR} (decay by \texttt{ε₀} at requests off \texttt{G})
  --- per-time a.e.\ conditional-expectation clauses, the
  probabilistic sibling of the integrated semantic-shadow forms on
  the deterministic pages; the BernCone rows reach them through the
  kernel bridge in the unbound assembly. Conclusion: \texttt{∀ᵐ ω,
  (∃ᶠ n, X n ω ∈ R) → ∃ᶠ n, X n ω ∈ G} --- the probabilistic content
  of one Streett pair, exactly.
\item nonneg rows $\mapsto$ \texttt{hnonneg : 0 ≤ V} --- and the
  roster\textquotesingle s two-branch signed-certificate
  counterexample is why the binder is there.
\item $\varepsilon_0 \in \Q_{>0}$ $\mapsto$ \texttt{hε : 0 < ε₀}
  (real, the $\Q$ cast upstream); the quantitative window/hitting
  forms (\texttt{streett\_window\_sum\_le},
  \texttt{ranking\_hitting\_sum\_le}) take the nonstrict \texttt{0 ≤
  ε₀} --- strictness is consumed exactly where Borel--Cantelli or
  summability-to-hitting needs it (\texttt{ae\_hits\_of\_ranking},
  the frequently-response theorem).
\item the measurability pack \texttt{hmR}/\texttt{hmG}
  (request/response events adapted at each time) is undisplayed
  regularity --- the probabilistic sibling of the PsdCert
  evaluability class, as on the battery pages.
\item \texttt{Supermartingale.exists\_ae\_tendsto\_of\_nonneg} (K5a)
  is engine stock: Doob convergence with the $L^1$ bound derived
  from nonnegativity, no display row corresponds.
\end{itemize}
Width note: the $V_0$ safety kit, the $W_i \le C_i$
upper-bound rows, the threshold tie $p + y \le 1$, and the
carved $J_i$ with the absorbing families $(V_i, \gamma_i)$ as
kernel-level rows are all consumed by the two bound assembly
statements; the G5 target-cell containment binds in the
\rulename{QStreettTargetCellCollapse} entry, and the complement-request
tie remains cited, not bound. The martingale-problem profile, stated
at its true extent: the confinement half and all zero-decrement
ranking rows are discharged by
\texttt{sound\_GSHS\_streett\_bundle\_mp}, and the strict-decrement
rows ($\varepsilon_i > 0$) by strict presentations ---
\texttt{MPFlowStrict} extends \texttt{MPFlow} by one pathwise progress
row (along every path the accumulated generator drift of the
discounted certificate reaches $-\varepsilon$ by a finite sample time;
the certificate-side derivation is $\varepsilon = \delta \cdot
t_{\min}$ from a strict segment drift bound and a minimum dwell time),
consumed by \texttt{rankingSegmentDynkin\_of\_mp\_strict}, which needs
presentations only at requested-but-ungranted states --- so
\texttt{sound\_GSHS\_streett\_bundle\_mp\_strict} leaves no analytic
per-model hypothesis in the bundle; all four hash-bound below. The
honest scope note that replaces the retired residue sentence: the
strict route\textquotesingle s progress row is \emph{pathwise}
(per-path accumulation, not in-expectation), the per-state bound
$\varepsilon \le V(x)$ is derived by the theorem rather than assumed
as a row, and the strict theorem is a companion to the mixed-profile
\texttt{\_mp} form, not an in-place replacement --- the landed
$\varepsilon = 0$ signature survives for citation stability. Width
note, deliberate strengthening
stated: the displayed ranking row asks strict $-\varepsilon_0$ decay
everywhere off $[\![\varphi_i]\!] \cup J_i$, where the bound
statements need nonincrease off the response and strict decrease only
at requests --- the displayed single-row form is stronger than
consumed, safe direction, and the producer emits it as one family.

\ruleblock{Lean verifier receipts}
\sourcenote{module \texttt{BernsteinLogicAux.Excursions}; bound build \texttt{ad62a846f92e\ldots} (hashed from the built olean at bind time). A receipt records that named checking layers accepted a build: (module, compiled-artifact SHA-256, environment SHA-256, checker layers), minted when the checking gate ran; the store receipts attesting the bound build are listed with the layers each run credited. Nothing here asks the reader to trust the store: every layer is an independently runnable program over the quoted sources --- rebuild the module and re-run the checkers to reproduce the verdict. Receipts attesting superseded builds credit no layer and are omitted from this rendering (the count of omissions, if any, is noted below).}
\begin{receipttable}
\receiptrow{495e3739c82a865d}{2026-07-18T13:23:35Z}{build+decorrelated}{\badge{BadgeGreen}{current}}
\receiptrow{fd91ed9c49139bb3}{2026-07-19T18:23:53Z}{build+decorrelated}{\badge{BadgeGreen}{current}}
\receiptrow{40087fade64fc483}{2026-07-20T01:21:05Z}{build+decorrelated}{\badge{BadgeGreen}{current}}
\receiptrow{483b08e5a60a003d}{2026-07-20T20:46:29Z}{build+decorrelated}{\badge{BadgeGreen}{current}}
\receiptrow{7142f9529d463930}{2026-07-20T20:58:02Z}{build+replay+hardened+decorrelated}{\badge{BadgeGreen}{current}}
\receiptrow{a477dcd09a788edc}{2026-07-21T00:02:42Z}{build+decorrelated}{\badge{BadgeGreen}{current}}
\receiptrow{30b972191d8805c2}{2026-07-21T05:38:35Z}{build+replay+hardened+decorrelated}{\badge{BadgeGreen}{current}}
\receiptrow{c6df511b240a81db}{2026-07-22T18:26:52Z}{build+decorrelated}{\badge{BadgeGreen}{current}}
\receiptrow{7d4ae5654cfa6cc8}{2026-07-31T21:59:44Z}{build+decorrelated}{\badge{BadgeGreen}{current}}
\receiptrow{50671021df35b6cc}{2026-07-31T22:28:38Z}{build+decorrelated}{\badge{BadgeGreen}{current}}
\receiptrow{153a2e7124bfd8d7}{2026-08-02T23:30:45Z}{build+decorrelated}{\badge{BadgeGreen}{current}}
\receiptrow{069303d7f7ac0945}{2026-08-03T16:26:49Z}{build+decorrelated}{\badge{BadgeGreen}{current}}
\receiptrow{b0dd60fee0daa2eb}{2026-08-03T22:45:33Z}{build+decorrelated}{\badge{BadgeGreen}{current}}
\end{receipttable}
\sourcenote{module \texttt{BernsteinLogicAux.QStreett}; bound build \texttt{469881beb770\ldots} (bound from receipt 9d1d848e0ba1718a, 2026-07-29T01:03:31Z). A receipt records that named checking layers accepted a build: (module, compiled-artifact SHA-256, environment SHA-256, checker layers), minted when the checking gate ran; the store receipts attesting the bound build are listed with the layers each run credited. Nothing here asks the reader to trust the store: every layer is an independently runnable program over the quoted sources --- rebuild the module and re-run the checkers to reproduce the verdict. Receipts attesting superseded builds credit no layer and are omitted from this rendering (the count of omissions, if any, is noted below).}
\begin{receipttable}
\receiptrow{e3990fe544d4708d}{2026-07-28T23:48:35Z}{build+decorrelated}{\badge{BadgeGreen}{current}}
\receiptrow{68e9374fa1b418fe}{2026-07-29T00:00:34Z}{build+decorrelated}{\badge{BadgeGreen}{current}}
\receiptrow{e5895c9706854cd3}{2026-07-29T00:49:04Z}{build+replay+hardened+decorrelated}{\badge{BadgeGreen}{current}}
\receiptrow{9d1d848e0ba1718a}{2026-07-29T01:03:31Z}{build+replay+hardened+decorrelated}{\badge{BadgeGreen}{current}}
\receiptrow{4a8d377b721b8523}{2026-07-31T22:00:09Z}{build+decorrelated}{\badge{BadgeGreen}{current}}
\receiptrow{d251a122540ab741}{2026-07-31T22:29:03Z}{build+decorrelated}{\badge{BadgeGreen}{current}}
\receiptrow{333635b8e047113a}{2026-08-02T23:31:14Z}{build+decorrelated}{\badge{BadgeGreen}{current}}
\receiptrow{9e96f80d63e4ad82}{2026-08-03T16:27:18Z}{build+decorrelated}{\badge{BadgeGreen}{current}}
\receiptrow{932b4b9e0b0fd323}{2026-08-03T22:46:02Z}{build+decorrelated}{\badge{BadgeGreen}{current}}
\end{receipttable}
\sourcenote{9 further receipt(s) in the store attest superseded builds of this module and are omitted from this rendering.}
\sourcenote{module \texttt{BernsteinLogicAux.LayerD.Streett}; bound build \texttt{7393065bbc7f\ldots} (bound from receipt abac471866cfcfd1, 2026-07-21T05:38:48Z). A receipt records that named checking layers accepted a build: (module, compiled-artifact SHA-256, environment SHA-256, checker layers), minted when the checking gate ran; the store receipts attesting the bound build are listed with the layers each run credited. Nothing here asks the reader to trust the store: every layer is an independently runnable program over the quoted sources --- rebuild the module and re-run the checkers to reproduce the verdict. Receipts attesting superseded builds credit no layer and are omitted from this rendering (the count of omissions, if any, is noted below).}
\begin{receipttable}
\receiptrow{0f3b50b4d3f0d33c}{2026-07-20T22:52:50Z}{build+decorrelated}{\badge{BadgeGreen}{current}}
\receiptrow{109669f848453ec2}{2026-07-21T00:03:19Z}{build+decorrelated}{\badge{BadgeGreen}{current}}
\receiptrow{7e17aa32b276623b}{2026-07-21T04:51:49Z}{build+decorrelated}{\badge{BadgeGreen}{current}}
\receiptrow{abac471866cfcfd1}{2026-07-21T05:38:48Z}{build+replay+hardened+decorrelated}{\badge{BadgeGreen}{current}}
\end{receipttable}
\sourcenote{6 further receipt(s) in the store attest superseded builds of this module and are omitted from this rendering.}
\sourcenote{module \texttt{QtslAux.LayerD.MPDischarge}; bound build \texttt{baded74dd943\ldots} (bound from receipt 7228834dea10e5f9, 2026-07-24T11:04:04Z). A receipt records that named checking layers accepted a build: (module, compiled-artifact SHA-256, environment SHA-256, checker layers), minted when the checking gate ran; the store receipts attesting the bound build are listed with the layers each run credited. Nothing here asks the reader to trust the store: every layer is an independently runnable program over the quoted sources --- rebuild the module and re-run the checkers to reproduce the verdict. Receipts attesting superseded builds credit no layer and are omitted from this rendering (the count of omissions, if any, is noted below).}
\begin{receipttable}
\receiptrow{7228834dea10e5f9}{2026-07-24T11:04:04Z}{build+decorrelated}{\badge{BadgeGreen}{current}}
\end{receipttable}
\layerline{\layeritem{build}{lake build (elaboration)}{\badge{BadgeGreen}{current}}\layeritem{replay}{leanchecker (official kernel replay)}{\badge{BadgeGreen}{current}}\layeritem{hardened}{safe\_verify (hardened screen)}{\badge{BadgeGreen}{current}}\layeritem{decorrelated}{lean4export + nanoda (decorrelated kernel)}{\badge{BadgeGreen}{current}}}
\sourcenote{\badge{BadgeGray}{comparator: none} no fixture in the three-implementation comparator harness (compiled / kernel-interpreted / reference-interpreted) exercises this rule.}
\ruleblock{HOL4 statement}
\absentbox{NO HOL STATEMENT YET}{Absent, with the gap located precisely: the HOL4 development ports the
\emph{base-calculus} checker cone (the \texttt{bern\_catalogue} theory,
built side-by-side with the Lean catalogue module) --- its object is the
wire-format certificate checker of the sequent stratum. The extension
theory\textquotesingle s satisfaction stratum has no HOL artefact at
all, so there is no HOL-side wire object a checker could
consume today. A HOL statement becomes meaningful
once the extension stratum is ported; it would then join the proof-object
layer the HOL port\textquotesingle s own coordination notes name as
successor territory. (Checked against the HOL sources: no Streett-,
satisfaction-, or NNF-shaped definitions exist there.)
}
\ruleblock{HOL verifier receipts}
\absentbox{NONE}{no HOL verification receipt exists for this rule.}
\ruleblock{Proof}

\sourcenote{Genre: core rule (tier 1) --- an ordinary mathematical proof at
Putnam-exam-solution detail, at the rule's \emph{proven scope}: discrete-kernel
profile, conservative ranking scope (\autoref{rem:qstb-scope} states the
scope loudly, including what is and is not hash-bound in this edition). Every
hypothesis is used and remarked on; the anchor table maps steps to the Lean
declarations that carry them.}

\textbf{Setting and claim.}\enspace The semantic object at the discrete-kernel
profile is a time-homogeneous Markov chain on a measurable state space $S$,
presented by a single step kernel $K$ (modes, guards, and resets of the
underlying hybrid presentation flattened into $K$), with path measure $\mu_x$
on $S^{\mathbb{N}}$ from a start state $x$ (the Ionescu--Tulcea trajectory
construction). Region data: the operating region $D$, its checker-recomputed
hull $\bar{D} \supseteq D$, and the exit set $U$ for the confinement
certificate; per Streett pair $i$ (finitely or countably many), a certified
request region $R_i$ and response region $G_i$. Certificate data: the
confinement certificate $V_0 : S \to \R$ with threshold
$y$, and per pair a ranking certificate $W_i : S \to \R$ with decay
rate $\varepsilon_i \in \Q_{>0}$ and bound $C_i$. The premise rows, each an
inherited base-calculus sequent:

\begin{itemize}
\item[(H1)] threshold arithmetic: $p + y \le 1$ (exact rational).
\item[(H2)] hull and mass rows, for every $z \in D \smallsetminus U$:
  $K(z)(\bar{D}) = 1$ and $K(z)(D \cup U) = 1$ --- one-step mass stays in the
  hull, and stays in operating-or-exit.
\item[(H3)] confinement rows: one-step generator row
  $\int V_0 \, dK(z) \le V_0(z)$ on $D \smallsetminus U$ (with one-step
  integrability); floor $V_0 \ge 1$ on $U \cap \bar{D}$; nonnegativity
  $V_0 \ge 0$ on $\bar{D}$.
\item[(H4)] ranking rows, per pair, on the confinement region
  $D \smallsetminus U$: bounds $0 \le W_i \le C_i$ on $\bar{D}$; nonincrease
  $\int W_i \, dK(z) \le W_i(z)$ off the response region $G_i$; request decay
  $\int W_i \, dK(z) \le W_i(z) - \varepsilon_i$ on $R_i \smallsetminus G_i$.
\item[(H5)] initialization: $x \in D$ and $V_0(x) \le y$ --- i.e.\
  $x$ lies in the concluded region atom
  $\mathsf{qraSublevel}(V_0, y)$.
\end{itemize}

The claim: under (H1)--(H5),
\[
\mu_x\Bigl\{\omega \;\Bigm|\; \forall i,\ \bigl(\exists^\infty n,\ \omega_n
\in R_i\bigr) \Rightarrow \bigl(\exists^\infty n,\ \omega_n \in G_i\bigr)
\Bigr\} \;\ge\; p ,
\]
which is the Streett-acceptance reading of the concluded judgment
$\mathcal{M},\ x_0 \in \{V_0 \le y\} \vDash
\mathsf{SProbGE}\bigl(p,\ \mathsf{PDwellStreett}(\mathsf{PairBase}(\rho_1,
\varphi_1), \ldots, \mathsf{PairBase}(\rho_k, \varphi_k))\bigr)$ once the
certified regions are tied to the formulas (\autoref{rem:qstb-tie}).

\begin{blemma}[confinement]\label{lem:qstb-confine}
$\mu_x\{\omega \mid \forall n,\ \omega_n \in D \smallsetminus U\} \ge
1 - V_0(x) \ge 1 - y \ge p$.
\end{blemma}
\begin{proof}
This is the quantitative-safety rule consumed whole: by (H2)--(H3), $V_0$
read along the path and stopped at the first $U$-visit is a nonnegative
supermartingale with barrier $\ge 1$ on the stopped set, so Ville's maximal
inequality bounds the exit probability by $V_0(x)$; the mass rows (H2) keep
every reachable state inside $\bar{D}$, where the nonnegativity and floor
rows actually hold --- the fail-closed pairing. The full argument, with the
two-branch counterexample showing the nonnegativity row is not decoration,
is the \rulename{QSafetyBattery} section's; the displayed intermediate
bound is the chain inside that argument --- Proposition~\ref{prop:qsb-pattern}
applied to the clamped certificate of Lemma~\ref{lem:qsb-clamp} bounds
the exit probability by $g(x) = V_0(x)$ --- consumed here through the
same end-to-end statement (anchor table), whose packaged form
Theorem~\ref{thm:qsb-failclosed} then states as the $\ge p$ complement.
The last two inequalities are (H5) and (H1).
\end{proof}

\begin{blemma}[windowed telescoping invariant]\label{lem:qstb-window}
Fix one pair $(R, G)$ with certificate $W$, rate $\varepsilon$, and suppose
the generator rows hold \emph{globally}: $\mathbb{E}[W_{n+1} \mid \mathcal{F}_n]
\le W_n$ a.e.\ off $G$, and $\le W_n - \varepsilon$ a.e.\ on requests off
$G$, for the process $W_n := W(X_n)$ adapted to a filtration
$(\mathcal{F}_n)$, with $W_n \ge 0$ and integrable. For a window start $N$
let $A_{N,n} := \{\text{no response in } [N, n)\}$. Then for all $n \ge N$,
\[
\int_{A_{N,n}} W_n \, d\mu \;+\; \varepsilon \sum_{k = N}^{n-1}
\mu\bigl(A_{N,k+1} \cap \{X_k \in R\}\bigr) \;\le\; \int W_N \, d\mu .
\]
\end{blemma}
\begin{proof}
Induction on $n$ from $n = N$, where $A_{N,N}$ is everything and the sum is
empty. For the step: $A_{N,n+1}$ is $\mathcal{F}_n$-measurable (it
constrains coordinates $< n + 1$), so the tower property converts
$\int_{A_{N,n+1}} W_{n+1}$ into $\int_{A_{N,n+1}}
\mathbb{E}[W_{n+1} \mid \mathcal{F}_n]$. Split $A_{N,n+1}$ at time $n$ into
its request and non-request parts: on the request part the decay row pays
$\varepsilon \cdot \mu(A_{N,n+1} \cap \{X_n \in R\})$ --- exactly the new
summand --- and on the non-request part the nonincrease row holds (both
parts lie off $G$ at time $n$, since $A_{N,n+1}$ forbids a response at
$n$). Finally $A_{N,n+1} \subseteq A_{N,n}$ and $W_n \ge 0$ shrink the
integral domain monotonically, closing the induction. Two features deserve
remark. First, \emph{no sign condition on $\varepsilon$ is used here}: the
decrement paid at step $k$ is exactly the summand recorded at index $k$,
an exact cancellation --- positivity of $\varepsilon$ enters only in
Lemma~\ref{lem:qstb-bc}. Second, the two generator rows do genuinely different
work and neither subsumes the other: nonincrease carries the (possibly
long) non-request stretches; decay is charged only at requests.
\end{proof}

\begin{blemma}[summability and Borel--Cantelli]\label{lem:qstb-bc}
Under the hypotheses of Lemma~\ref{lem:qstb-window} with $\varepsilon > 0$
and $W$ bounded: almost surely, infinitely many requests imply infinitely
many responses.
\end{blemma}
\begin{proof}
Fix a window start $N$. Lemma~\ref{lem:qstb-window} bounds every
partial sum of the series below by $\mathbb{E}[W_N]$ uniformly (the
leading integral it retains is nonnegative), and a nonnegative series
is the supremum of its partial sums, so
\[
\varepsilon \sum_{j \ge 0} \mu\bigl(\{\text{no response in } [N, N+j]\}
\cap \{X_{N+j} \in R\}\bigr) \;\le\; \mathbb{E}[W_N] \;<\; \infty ,
\]
so the un-responded-request events of window $N$ are summable, and by the
Borel--Cantelli lemma their limsup is null. A path with infinitely many
requests but only finitely many responses has a finite (possibly empty)
set of response times; for any window start $N$ beyond all of them
($N = 0$ when there are none), the path lies in infinitely many
of window $N$'s un-responded-request events --- i.e.\ in that null limsup.
The union over the countably many window starts is null, which is the
claim. (Here $\varepsilon > 0$ is finally consumed, dividing the bound;
boundedness of $W$ supplies the finite right-hand side and, upstream,
integrability of every $W_n$.)
\end{proof}

\begin{blemma}[localization to the confinement region]\label{lem:qstb-loc}
The bundle's ranking rows (H4) hold only on $D \smallsetminus U$, while
Lemma~\ref{lem:qstb-window} wants global rows. Define the \emph{localized}
pair process: flag the state with the survival indicator of the event
``no visit to $U \cup D^{\mathrm{c}}$ so far'', zero the certificate at the
first exit, and count exited states as \emph{responded}. Then the flagged
process satisfies the global hypotheses of Lemmas
\ref{lem:qstb-window}--\ref{lem:qstb-bc}, and the conclusion, restricted to
paths confined forever, descends to the raw process: almost surely on
$\{\forall n,\ \omega_n \in D \smallsetminus U\}$, infinitely many
$R_i$-visits imply infinitely many $G_i$-visits.
\end{blemma}
\begin{proof}
Write it out. Let $\sigma := \inf\{n : \omega_n \in U \cup
D^{\mathrm{c}}\}$ (a stopping time --- the regions are measurable), and
localize by
\[
W'_n := W(\omega_n)\,\mathbf{1}[\sigma > n],
\qquad
G' \ \text{holds at time } n \ :\iff\ \omega_n \in G \ \text{or}\
\sigma \le n ,
\]
the zeroed certificate and the enlarged (exited-counts-as-responded)
response. Check the global hypotheses of
Lemma~\ref{lem:qstb-window} for $(W', G')$ by cases on the state class
at time $n$:

\emph{(surviving, non-request)} $\sigma > n$ and $\omega_n \notin R$:
off $G'$ means also $\omega_n \notin G$, so $\omega_n \in
(D \smallsetminus U) \smallsetminus G$; the mass row (H2) carries the
next step into $\bar{D}$, where $0 \le W \le C$ (H4), so
$W'_{n+1} = W(\omega_{n+1})\mathbf{1}[\sigma > n+1] \le W(\omega_{n+1})$
almost surely on this branch --- the masked certificate is dominated by
the unmasked one --- and
$\mathbb{E}[W'_{n+1} \mid \mathcal{F}_n] \le \int W \, dK(\omega_n)
\le W(\omega_n) = W'_n$ by the nonincrease row.

\emph{(surviving, request)} the same domination, closed by the decay
row: $\mathbb{E}[W'_{n+1} \mid \mathcal{F}_n] \le W'_n - \varepsilon$.

\emph{(exited)} $\sigma \le n$: the state is in the enlarged response
$G'$ by construction, so no generator row is demanded of it --- the global
hypotheses hold vacuously, and indeed $W'_n = 0$ there.

Bounds and adaptedness: $0 \le W'_n \le C$ everywhere (H4 on $\bar{D}$
while surviving --- reachable surviving states lie in $\bar D$ by (H2)
and the deterministic start --- zero after), and $W'$ is adapted since
$\sigma$ is a stopping time. (To consume
Lemma~\ref{lem:qstb-window} literally in its certificate reading, run
the chain on the flagged state space $S \times \{0,1\}$: the flag is a
running conjunction updated by membership of each new coordinate in
$U \cup D^{\mathrm{c}}$, so the flagged chain is again time-homogeneous
Markov, $W'$ and the enlarged response are state functions there ---
$W'(z, f) = f \cdot W(z)$, response $\iff$ $f = 0$ or $z \in G$ --- and
the flagged trajectory law projects to $\mu_x$.) Lemmas
\ref{lem:qstb-window}--\ref{lem:qstb-bc} then give: almost surely,
infinitely many requests-while-surviving imply infinitely many
$G'$-responses. Restricted to confined paths the flag never trips
($\sigma = \infty$), $W' = W$, the enlarged response set coincides with
$G_i$, and the conclusion is the raw one. The
exited-counts-as-responded device is not a weakening of the claim: the
theorem below asserts the pair implications only on the confinement event
of Lemma~\ref{lem:qstb-confine}, exactly where the device is invisible.
\end{proof}

\begin{bthm}[soundness of \rulename{QStreettBattery}, discrete profile,
conservative ranking scope]\label{thm:qstb-sound}
Under (H1)--(H5),
$\mu_x\{\omega \mid \forall i,\ (\exists^\infty n,\ \omega_n \in R_i)
\Rightarrow (\exists^\infty n,\ \omega_n \in G_i)\} \ge p$.
\end{bthm}
\begin{proof}
Per pair $i$, Lemmas \ref{lem:qstb-window}--\ref{lem:qstb-loc} give the pair
implication almost surely on the confinement event. The pairs are
countably many, so the almost-sure statements intersect (countable
intersection of co-null sets); the Streett event therefore contains the
confinement event up to a null set, and Lemma~\ref{lem:qstb-confine} bounds
its measure below by $p$. $\qedhere$
\end{proof}

\begin{bremark}[every hypothesis earns its seat]\label{rem:qstb-hyp}
\emph{(H1)}: at $p + y > 1$ the concluded atom admits start
states with $1 - V_0(x) < p$ and Lemma~\ref{lem:qstb-confine} yields less
than $p$ --- the sublevel atom and the probability bound are coupled by
exactly this row. \emph{(H2)}: drop the hull row and probability mass can
reach $U \smallsetminus \bar{D}$, where no floor row holds, voiding
Ville's barrier in Lemma~\ref{lem:qstb-confine}; drop the
operating-or-exit row and the stopped-process reading of the confinement
certificate is unsound (mass escapes the region where a generator row was ever
asserted) --- the two recomputed side-conditions are the fail-closed
pairing. \emph{(H3) nonnegativity}: load-bearing, not hygiene --- a signed
confinement certificate can go negative on one branch and thereby
subsidize exit mass on another; the recorded two-branch counterexample
lives in the \rulename{QSafetyBattery} section. \emph{(H4) strictness}:
at $\varepsilon_i = 0$, take $W_i \equiv 1$ on a one-state chain with
$K(s) = \delta_s$, $R_i = \{s\}$, $G_i = \varnothing$, $V_0 \equiv 0$:
every row holds, the path requests forever, and no response ever occurs
--- the conclusion fails at any $p > 0$. \emph{(H4) bound rows}: the
constant $C_i$ gives integrability of the localized certificate
(Lemma~\ref{lem:qstb-loc}) and the finite right-hand side of
Lemma~\ref{lem:qstb-bc}; nonnegativity of $W_i$ additionally powers the
monotone domain-shrink step inside Lemma~\ref{lem:qstb-window}.
\emph{(H5)}: initialization is what makes the region atom an
\emph{under}-approximation of the satisfaction set --- the workhorse
orientation of the satisfaction judgment.
\end{bremark}

\begin{bremark}[scope, loudly]\label{rem:qstb-scope}
Three ledger entries: the discrete-profile warrant, the cited
testimony, and the continuous-time profile. \emph{(i) Bound in
this edition}: the five excursion-cluster declarations (anchor table) ---
the windowed telescoping engine (Lemmas
\ref{lem:qstb-window}--\ref{lem:qstb-bc}), its additive-drift sibling for
plain reachability (the ranking hitting-time bound and almost-sure
hitting), and the nonnegative-supermartingale convergence theorem (part
of the cluster, available for the general excursion picture; the
presented proof path does not consume it) --- together with the
end-to-end assembly at the conservative scope
(\texttt{sound\_QStreettBattery} --- Lemmas \ref{lem:qstb-confine},
\ref{lem:qstb-loc} and Theorem~\ref{thm:qstb-sound} as one statement)
and the discharged carved-scope corollary
(\texttt{sound\_QStreettBattery\_carved\_discharged}), both in the
auxiliary module \texttt{BernsteinLogicAux.QStreett}. These warrant the
rule at the discrete-kernel profile. Also bound in this edition, under
the \rulename{QStreettCmplRequest} entry whose section carries its
event-level kernel: the formula tie
(\texttt{sound\_QStreettBattery\_tie}, with
\texttt{streett\_event\_mono}). \emph{(ii) Machine-checked and
cited, not bound}: the localization lemma under its
own name (\texttt{streett\_pair\_localized}), and the parameterized
route to the sharper \emph{carved} ranking scope, where decay is
demanded only off $G_i \cup J_i$ with
$J_i := \{V_i \le 1 - \gamma_i\} \cap
(I \smallsetminus [\![\rho_i]\!])$ --- $I$ is the display's confinement
region $\{V_0 < 1\}$; this section's exit set instances it as
$U = \{1 \le V_0\}$ --- and the absorbing certificates
$(V_i, \gamma_i)$ repay the carve-out: a L\'evy upward-convergence
argument shows frequent $J_i$-visits preclude frequent requests, the
absorption bound being discharged from kernel-level rows by a
time-shifted supermartingale bridge plus a conditional Ville inequality
(\texttt{sound\_QStreettBattery\_carved},
\texttt{absorptionBound\_of\_kernel}). Binding those in a later edition
is the named follow-up; until then this page claims them as
testimony, not as hash-bound warrant. \emph{(iii) The
martingale-problem profile, discharged whole}: the confinement half
and every zero-decrement ranking row ($\varepsilon_i = 0$:
nonincrease off the response region) are discharged by
\texttt{sound\_GSHS\_streett\_bundle\_mp}, and the strict-decrement
rows ($\varepsilon_i > 0$) by strict presentations ---
\texttt{MPFlowStrict} extends the martingale-problem presentation by
one pathwise progress row (along every path the accumulated generator
drift of the discounted certificate reaches $-\varepsilon$ by a
finite sample time; the intended certificate-side derivation is
$\varepsilon = \delta \cdot t_{\min}$ from a strict segment drift
bound and a minimum dwell time), and
\texttt{rankingSegmentDynkin\_of\_mp\_strict} derives from it the
ranking row family the $\varepsilon = 0$ route could not produce,
needing presentations only at requested-but-ungranted states. The
proof mechanism, for the record, is \emph{not} the lower-bound
optional-sampling inequality an earlier edition of this remark
predicted: the stopped certificate splits as martingale plus drift
accumulation, the existing upper-bound optional sampling serves the
martingale part, and dominated convergence on
$\varepsilon$-truncated drift accumulations supplies the strict term
(the lower-bound lemma was built anyway, completing the
optional-sampling module). The assembled
\texttt{sound\_GSHS\_streett\_bundle\_mp\_strict} therefore carries
no analytic per-model hypothesis --- all four statements live in
module \texttt{QtslAux.LayerD.MPDischarge} and are hash-bound in
this edition via the registered qtsl-aux package, so, like its
safety and deadline siblings, this rule carries no conditional
mark: each named profile is proven, and its theorems are bound.
\end{bremark}

\begin{bremark}[the formula tie]\label{rem:qstb-tie}
The theorem is stated over certified regions $R_i, G_i$; the concluded
judgment speaks of formulas $\rho_i, \varphi_i$. The sound directions are
asymmetric: certified \emph{request} regions may \emph{over}-approximate
the request denotation, certified \emph{response} regions must
\emph{under}-approximate the response denotation --- weakening the
hypothesis and strengthening the conclusion of each pair implication.
This is the content of the complement-request sub-rule
(\rulename{QStreettCmplRequest}, two sections on), whose event-level kernel
is two monotonicity steps.
\end{bremark}

\begin{bremark}[degenerate pairs, and why Streett is primitive]
A B\"uchi acceptance condition is the one-pair Streett condition with
trivially-true request; on that pair the bundle degenerates to the
classical recurrence certificate (the \rulename{QStreettDegenerate}
section). The converse embedding does not exist at parity of cost:
deterministic B\"uchi automata recognize only limit languages
\cite{landweber1969}, and re-encoding a $k$-pair Streett condition
B\"uchi-primitively costs an exponential-in-$k$ automaton product, where
the bundle pays $2k + 1$ certificates. Streett conditions enter the
literature as the complemented-pairs acceptance of \cite{streett1982},
the linear-time face of Rabin's pairs \cite{rabin1969}; the standard
survey is \cite{thomas1990}.
\end{bremark}

\begin{bremark}[relation to the literature]
The ranking row is a Foster--Lyapunov drift condition
\cite{foster1953,meyntweedie1993} in certificate form; its use to force
almost-sure hitting is the ranking-supermartingale argument of the
probabilistic-programs tradition \cite{chakarov2013}. The confinement
certificate is Ville's inequality \cite{ville1939}; the convergence item
of the bound cluster is Doob's supermartingale convergence theorem
\cite{doob1953}; the summability-to-limsup step is Borel--Cantelli (see
e.g.\ \cite[Thm.~2.3.1]{durrett2019}). As with the polynomial rules of
the base calculus, the certificate form moves all analytic content into
the \emph{choice} of $(V_0, W_i, \varepsilon_0)$: what the checker
replays is row arithmetic, and what this proof adds is only the
assembly. No completeness is claimed or needed here; the source
catalogue records a completeness theorem for the bundle relative to its
termination seam, outside this book's kernel-facing scope.
\end{bremark}

\begin{anchors}
Lemma \ref{lem:qstb-confine} & \texttt{sound\_QSafetyBattery\_failClosed} (bound under the \rulename{QSafetyBattery} entry; consumed whole)\\
Lemma \ref{lem:qstb-window} + no-sign-condition remark & \texttt{streett\_window\_invariant} (the engine's inner induction; subsumed by the bound \texttt{streett\_window\_sum\_le})\\
Lemma \ref{lem:qstb-bc}, summability form & \texttt{streett\_window\_sum\_le} (bound)\\
Lemma \ref{lem:qstb-bc}, Borel--Cantelli assembly & \texttt{streett\_ae\_frequently\_response} (bound)\\
additive-drift sibling (reachability) & \texttt{ranking\_hitting\_sum\_le}, \texttt{ae\_hits\_of\_ranking} (bound)\\
supermartingale convergence (cluster item) & \texttt{Supermartingale.exists\_ae\_tendsto\_of\_nonneg} (bound; not consumed by the presented path)\\
Lemma \ref{lem:qstb-loc} & \texttt{streett\_pair\_localized} (\texttt{BernsteinLogicAux.QStreett}; machine-checked, binding pending)\\
Theorem \ref{thm:qstb-sound} & \texttt{sound\_QStreettBattery} (same module, \textbf{bound in this edition})\\
Remark \ref{rem:qstb-tie} & \texttt{streett\_event\_mono}, \texttt{sound\_QStreettBattery\_tie} (same module; \textbf{bound in this edition} under the \rulename{QStreettCmplRequest} entry)\\
carved scope, discharged corollary & \texttt{sound\_QStreettBattery\_carved\_discharged} (same module, \textbf{bound in this edition})\\
Remark \ref{rem:qstb-scope}(ii), carved scope route & \texttt{AbsorptionBound}, \texttt{ae\_not\_frequently\_of\_absorption}, \texttt{supermartingale\_stopped\_shift}, \texttt{measure\_reach\_after\_le}, \texttt{absorptionBound\_of\_kernel}, \texttt{sound\_QStreettBattery\_carved} (same module, binding pending)\\
Remark \ref{rem:qstb-scope}(iii), martingale-problem profile & \texttt{sound\_GSHS\_streett\_bundle\_mp}, \texttt{MPFlowStrict}, \texttt{rankingSegmentDynkin\_of\_mp\_strict}, \texttt{sound\_GSHS\_streett\_bundle\_mp\_strict} (\texttt{QtslAux.LayerD.MPDischarge}; \textbf{bound in this edition} via the qtsl-aux package registration)\\
\end{anchors}


\rulesection{45}{QAvgBattery}{core rule}
\label{blcatrule-QAvgBattery}
\ruleblock{The rule (Steele metanotation)}
\sourcenote{as realized in the mechanized development; source row: \texttt{\texttt{QTSL-rules.toml}} --- \emph{QAvgBattery}. The metanotation statement is typeset by hand for notational quality; this rule cites no Lean statements yet (the absence is stated loudly below). The rule \emph{roster} is machine-checked (every mechanized rule tag round-trips to exactly one catalogue row); the statement \emph{text} is editorial.}
\begin{mathpar}
\inferrule*{
  \overline{\mathrm{drift}_m : R_m \vdash g - r_m - \mathcal{A}_m B_m \ \ge\ 0} \\
  \overline{\mathrm{guard}_j} \quad \overline{\mathrm{floor}_m} \quad \overline{\mathrm{nonneg}_m} \\
  \mathrm{tie} : \vdash g \bowtie c
}{
  \mathcal{M},\ x_0 \in R \vDash \mathsf{SReturn}{\bowtie}\bigl(c,\ \mathsf{poAvg}(o)\bigr)
}
\end{mathpar}

\par\medskip
\rulename{QAvgBattery} concludes long-run-average claims from the
certificate pair of the Poisson equation: a rational \emph{gain} $g$
and a \emph{bias} certificate family $B$, with the drift rows
carrying the constant term inside each row\textquotesingle s sequent
polynomial --- so the discharge species is the same sign-sequent
genre as every other battery, with no new checker architecture. On
the upper face the drifts certify gain minus reward minus
generator-of-bias nonnegative on scope, and the tie bounds the gain
above the claimed constant; on the lower face the orientations flip,
and the floors bound the bias above so the correction term dies in
the limit. The walk recomputes the required orientation from the
concluded root --- a producer never declares a face.

\cataloguenote{Entered at the roster census: the node was baked --- including its
tie proof field, which the safety and acceptance batteries still lack
--- but had never been catalogued. No Lean statement binds it yet;
the display is the schema-faithful schematic form, and the soundness
obligation (the Poisson-pair argument: telescoping the bias along the
trajectory, dividing by the horizon, the floors killing the
correction term) is stated in the proof section as an obligation,
not discharged.
}
\ruleblock{Theory rows (geolog)}
\sourcenote{The rule\textquotesingle s certificate content as it is declared in the geolog theory the checker loads --- payload schema, side-condition demands, and axiom rows, quoted from the theory source at every build of this catalogue (a stale quotation cannot exist). Axiom modes: \texttt{chase} = the engine derives; \texttt{report} = a violation is a loud judgment; \texttt{refuse} = a violation gates admission. Side-condition relations are discharged by bounded evaluation at the engine\textquotesingle s ground seam, never by derivation.}
\declhead{theory Qtsl}{QAvgBattery}{\badge{BadgeGreen}{quoted from the theory source at render}}
\begin{leansrc}
QAvgBattery : Set;
\end{leansrc}
\declhead{theory Qtsl}{QAvgBattery/*}{\badge{BadgeGreen}{quoted from the theory source at render}}
\begin{leansrc}
QAvgBattery/proof : QAvgBattery -> QProof;
QAvgBattery/gain : QAvgBattery -> Q;
QAvgBattery/arg : QAvgBattery -> StatePred;
QAvgBattery/cert : QAvgBattery -> QCertFam;
QAvgBattery/drifts : QAvgBattery -> QProofList;
QAvgBattery/guards : QAvgBattery -> QProofList;
QAvgBattery/lowbounds : QAvgBattery -> QProofList;
QAvgBattery/nonnegs : QAvgBattery -> QProofList;
QAvgBattery/tie : QAvgBattery -> Proof;
QAvgBattery/scope_warrant : QAvgBattery -> QInvScope;
\end{leansrc}
\ruleblock{Lean statement}
\badge{BadgeAmber}{not yet mechanized} this rule cites no Lean statements: nothing on this page is hash-bound for it, and its dashboard row shows the same absence across the statement, receipt, and correspondence columns.
\ruleblock{HOL4 statement}
\absentbox{NO HOL STATEMENT YET}{No HOL development of the extension theory exists.
}
\ruleblock{HOL verifier receipts}
\absentbox{NONE}{no HOL verification receipt exists for this rule.}
\ruleblock{Proof}

\sourcenote{Genre: core rule (tier 1), battery group. This section states the rule's soundness
obligation exactly, per the catalogue's completeness-over-polish
charter (a rule that can only be named-and-gapped is still a row);
the discharge is open debt on the mechanization column, and the
schema quoted in the entry is the baked ground truth.}

\textbf{Soundness obligation.}\enspace The Poisson-equation pair argument: with drift rows certifying
$g - r_m - \mathcal{A}_m B_m \ge 0$ on scope per mode (upper face),
guards at the seams, floors and nonnegativity on the bias family,
telescoping the compensated bias along any admitted trajectory bounds
the accumulated reward's time average by the gain plus a correction
term of order bias-over-horizon; the floors kill the correction in
the limit, and the tie transfers the gain bound to the claimed
constant. The lower face is the same argument with orientations
flipped. The obligation is stated over the same behavioural models as
the safety battery's, with the walk recomputing the required
orientation from the concluded root.

\begin{bremark}[status]
No Lean statement binds this node. The discharge species is the
landed sign-sequent genre --- the expected mechanization reuses the
battery kit's supermartingale machinery with the compensated bias in
the certificate slot --- and the node's tie field is already baked,
so no schema motion is expected at mechanization time.
\end{bremark}


\rulesection{46}{QRateFromClock}{core rule}
\label{blcatrule-QRateFromClock}
\ruleblock{The rule (Steele metanotation)}
\sourcenote{as realized in the mechanized development; source row: \texttt{\texttt{QTSL-rules.toml}} --- \emph{QRateFromClock}. The metanotation statement is typeset by hand for notational quality; this rule cites no Lean statements yet (the absence is stated loudly below). The rule \emph{roster} is machine-checked (every mechanized rule tag round-trips to exactly one catalogue row); the statement \emph{text} is editorial.}
\begin{mathpar}
\inferrule*{
  \mathsf{TrMeter}(t, d) \\
  \rho \cdot d \ge 1 \\
  \sigma \ge 1 \\
  \mathsf{qtsl\_rate\_clock}(n)
}{
  \mathcal{M} \vDash \mathsf{SRateLE}(\{t\},\ \rho,\ \sigma)
}
\end{mathpar}

\par\medskip
\rulename{QRateFromClock} discharges the rate-envelope specification
former from clock data alone: a transition whose clock carries a dwell
floor $d$ (demonic, forced, or carved --- never the instantaneous
race clock, which has no floor) satisfies
$\mathsf{SRateLE}(\{t\}, \rho, \sigma)$ for any $\rho \cdot d \ge 1$
and $\sigma \ge 1$. The rate lane\textquotesingle s other feed is the
supermartingale kit --- a different evidence family, deliberately not
conflated; this rule is the clock seam.

The premise is a premise \emph{class}, not a node:
$\mathsf{TrMeter}(t, d)$ names three things together --- the
transition\textquotesingle s clock-content conformance (exactly-one
content presentation, wing positivity; already checker-discharged
rows), a POSITIVE lower witness $0 < d \le d'$ below the declared
floor $d'$ (positivity is the class\textquotesingle s own formation:
a nonpositive witness would inhabit the class while its division law
is vacuous), and the meter law as a \emph{named obligation} --- the
conformance rows alone do not prove elapsed-time separation. The law:
under the declared meter semantics (dwell accrues per transition;
guard exit and own firing reset it), successive firings of $t$ are
separated by at least $d$ of elapsed time, hence
$N_t(T) \le T/d + 1$ on any window under every resolution. The
checker\textquotesingle s trusted contribution is two exact rational
inequalities ($\rho \cdot d \ge 1$; $\sigma \ge 1$) over conformance
rows it already checks --- with $d$ resolved deterministically as the
transition\textquotesingle s declared finite floor (no producer
freedom; the class is monotone in $d$, so the declared floor is the
tightest lawful witness).

The concluded envelope is an upper-prevision claim; the obligation
actually proved is pathwise and surely --- strictly stronger --- and
the entry records that strength so a future pathwise rate former can
strengthen the conclusion without touching the node. Singleton
rosters compose by the rate lane\textquotesingle s documented
disjoint-roster sum law (envelopes add on disjoint rosters); this
rule is that law\textquotesingle s first in-catalogue consumer.

\cataloguenote{The schema is baked in the extension theory: the node, its walk, the
race-refusal, and the per-wing meter-demand rows are quoted below
from the theory source; the per-transition conformance rows the
premise class reads were checker-discharged before the node landed.
What remains owed is the meter law\textquotesingle s first
formalization (the shared keystone of this rule family\textquotesingle s
continuous-time obligations --- the discrete-profile induction is
elementary). No Lean statement binds the node yet; the mechanization
column is honest open debt.
}
\ruleblock{Theory rows (geolog)}
\sourcenote{The rule\textquotesingle s certificate content as it is declared in the geolog theory the checker loads --- payload schema, side-condition demands, and axiom rows, quoted from the theory source at every build of this catalogue (a stale quotation cannot exist). Axiom modes: \texttt{chase} = the engine derives; \texttt{report} = a violation is a loud judgment; \texttt{refuse} = a violation gates admission. Side-condition relations are discharged by bounded evaluation at the engine\textquotesingle s ground seam, never by derivation.}
\par\medskip
The baked rows: the node with its two exact-rational envelope
reports, the race-refusal (the instantaneous wing carries no floor
and is refused as subject), the walk, and the per-wing meter-demand
rows that mint the premise class\textquotesingle s
conformance-plus-floor reading at the declared dwell floor.

\declhead{theory Qtsl}{QRateFromClock}{\badge{BadgeGreen}{quoted from the theory source at render}}
\begin{leansrc}
QRateFromClock : Set;
\end{leansrc}
\declhead{theory Qtsl}{QRateFromClock/*}{\badge{BadgeGreen}{quoted from the theory source at render}}
\begin{leansrc}
QRateFromClock/proof : QRateFromClock -> QProof;
QRateFromClock/tr    : QRateFromClock -> Transition;
QRateFromClock/rho   : QRateFromClock -> Q;
QRateFromClock/sigma : QRateFromClock -> Q;
\end{leansrc}
\declhead{theory Qtsl}{ax/wf/qrfc1\_rho\_pos}{\badge{BadgeGreen}{quoted from the theory source at render}}
\begin{leansrc}
report ax/wf/qrfc1_rho_pos : forall n : QRateFromClock, x : Q.
  n QRateFromClock/rho = x
  |- {l: 0, r: x} q_lt;
\end{leansrc}
\declhead{theory Qtsl}{ax/wf/qrfc1\_sigma\_ge\_one}{\badge{BadgeGreen}{quoted from the theory source at render}}
\begin{leansrc}
report ax/wf/qrfc1_sigma_ge_one : forall n : QRateFromClock, x : Q.
  n QRateFromClock/sigma = x
  |- {l: 1, r: x} q_le;
\end{leansrc}
\declhead{theory Qtsl}{ax/qtsl/rate\_clock\_race}{\badge{BadgeGreen}{quoted from the theory source at render}}
\begin{leansrc}
refuse ax/qtsl/rate_clock_race : forall n : QRateFromClock,
    t : Transition, e : EnnRat, k : TrClock.
  n QRateFromClock/tr = t,
  t Transition/clock = k,
  k = {majorant: e} tkRace
  |- false;
\end{leansrc}
\declhead{theory Qtsl}{qtsl\_rate\_clock}{\badge{BadgeGreen}{quoted from the theory source at render}}
\begin{leansrc}
qtsl_rate_clock : {node: QRateFromClock} -> Prop;
\end{leansrc}
\declhead{theory Qtsl}{ax/tf/qrfc1\_walk}{\badge{BadgeGreen}{quoted from the theory source at render}}
\begin{leansrc}
report ax/tf/qrfc1_walk : forall n : QRateFromClock.
  |- {node: n} qtsl_rate_clock;
\end{leansrc}
\declhead{theory Qtsl}{ax/tf/qrfc1\_meter\_forced}{\badge{BadgeGreen}{quoted from the theory source at render}}
\begin{leansrc}
report ax/tf/qrfc1_meter_forced : forall n : QRateFromClock,
    t : Transition, d : Q, k : TrClock.
  n QRateFromClock/tr = t,
  t Transition/clock = k,
  k = {dwell: d} tkForced
  |- {tr: t, d: d} qtsl_tr_meter;
\end{leansrc}
\declhead{theory Qtsl}{ax/tf/qrfc1\_meter\_demonic}{\badge{BadgeGreen}{quoted from the theory source at render}}
\begin{leansrc}
report ax/tf/qrfc1_meter_demonic : forall n : QRateFromClock,
    t : Transition, d : Q, e : EnnRat, k : TrClock.
  n QRateFromClock/tr = t,
  t Transition/clock = k,
  k = {dwellMin: e} tkDemonic,
  e = {q: d} enFin
  |- {tr: t, d: d} qtsl_tr_meter;
\end{leansrc}
\declhead{theory Qtsl}{ax/tf/qrfc1\_meter\_carved}{\badge{BadgeGreen}{quoted from the theory source at render}}
\begin{leansrc}
report ax/tf/qrfc1_meter_carved : forall n : QRateFromClock,
    t : Transition, fl : Q, ce : EnnRat, k : TrClock.
  n QRateFromClock/tr = t,
  t Transition/clock = k,
  k = {floor: fl, ceil: ce} tkCarved
  |- {tr: t, d: fl} qtsl_tr_meter;
\end{leansrc}
\declhead{theory Qtsl}{qtsl\_tr\_meter}{\badge{BadgeGreen}{quoted from the theory source at render}}
\begin{leansrc}
qtsl_tr_meter : {tr: Transition, d: Q} -> Prop;
\end{leansrc}
\declhead{theory Qtsl}{qtsl\_tr\_content}{\badge{BadgeGreen}{quoted from the theory source at render}}
\begin{leansrc}
qtsl_tr_content : {tr: Transition} -> Prop;             // side-condition
\end{leansrc}
\ruleblock{Lean statement}
\badge{BadgeAmber}{not yet mechanized} this rule cites no Lean statements: nothing on this page is hash-bound for it, and its dashboard row shows the same absence across the statement, receipt, and correspondence columns.
\ruleblock{HOL4 statement}
\absentbox{NO HOL STATEMENT YET}{No HOL development of the extension theory exists; this
node\textquotesingle s content is walk structure over baked rows.
}
\ruleblock{HOL verifier receipts}
\absentbox{NONE}{no HOL verification receipt exists for this rule.}
\ruleblock{Proof}

\sourcenote{Genre: core rule (tier 1), the timers family (the
rate-envelope former's clock seam). This section states the rule's
soundness obligation exactly; the rule's schema and walk are baked in
the extension theory (rows quoted in this section), and the meter
law's formalization is the named obligation of the mechanization
column.}

\textbf{Soundness obligation (meter law + pathwise domination).}\enspace
Under the declared meter semantics (dwell accrues per transition;
guard exit and own firing reset it), a transition $t$ whose clock
carries floor $d > 0$ fires with successive firings separated by at
least $d$ of elapsed time, hence $N_t(T) \le T/d + 1$ on any window
of length $T$, under every resolution. The witness $d$ is the
transition's declared finite floor, resolved deterministically by
the walk (the premise class demands a POSITIVE witness at most the
declared floor, and is monotone, so the declared floor is the
tightest lawful instance; no producer freedom enters). With the
checker-verified exact inequalities $\rho \cdot d \ge 1$ and
$\sigma \ge 1$,
$N_t(T) \le T/d + 1 \le \rho T + \sigma$ pathwise and surely; the
upper prevision of the roster firing count therefore satisfies the
envelope --- pathwise domination, strictly stronger than the
concluded prevision-level claim. The instantaneous race clock is
refused as subject: it carries no floor, and its rate story is a
different evidence family (the thinning majorant), deliberately not
conflated with this seam.

\begin{bremark}[status]
The schema is baked: the node, its walk, the race-refusal, and the
per-wing meter-demand rows are landed theory rows (quoted in this
section), and the premise class's per-transition conformance rows
were checker-discharged before the node landed. No Lean statement
binds the node yet; the meter LAW is the named obligation ---
elementary at the discrete profile, and the shared keystone of this
rule family's continuous-time lifts.
\end{bremark}


\rulesection{47}{QInvRelativize}{core rule}
\label{blcatrule-QInvRelativize}
\ruleblock{The rule (Steele metanotation)}
\sourcenote{as realized in the mechanized development; source row: \texttt{\texttt{QTSL-rules.toml}} --- \emph{QInvRelativize}. The metanotation statement is typeset by hand for notational quality; its mechanized ground truth is the Lean statement below, extracted from the sources at every build of this book and hash-bound. The rule \emph{roster} is machine-checked (every mechanized rule tag round-trips to exactly one catalogue row); the statement \emph{text} is editorial, checked by eye against the extracted Lean beside it.}
\begin{mathpar}
\inferrule*{
  \mathsf{BoxInv}(B, \nu, x)
}{
  \mathsf{AService}(B, T, \chi \cap \nu, x) \;\leftrightarrow\; \mathsf{AService}(B, T, \chi, x)
}
\end{mathpar}
$B$ an \emph{arbitrary} behaviour --- no closure field consumed; $\mathsf{AService}$ is the d\'ebut-as-infimum reading unpacked without infima (a visit at some $t < T + \varepsilon$ for every $\varepsilon > 0$), $\mathsf{BoxInv}$ the $\Box$ reading (every point of every solution from $x$).

\ruleflavor{certificate content out of scope of the current baked theory --- realized in Lean over the behavioural floor.}
\par\medskip
One of the three timed-core basis schemata proved non-admissible in
the gnand completeness development (machine independence witnesses;
see the proof section's anchors), entering as tier-1 core: its
soundness obligation is semantic, discharged at Putnam-exam-solution
detail in the proof section and mechanized in
BernsteinLogic.QtslTimedCore. Sound over \emph{all} behavioural models
--- no side condition; the proof consumes none of the closure laws.

\ruleblock{Lean statement}
\declhead{BernsteinLogic.QtslTimedCore}{sound\_QInvRelativize}{\badge{BadgeGreen}{statement bound}}
\begin{leansrc}
/-- **QInvRelativize** (obligation relativization through invariants), the
biconditional core: under `□ν`, the visit sets of `χ ∩ ν` and `χ`
coincide along every solution from `x`, so the two service claims
agree. Sound for ARBITRARY behaviors — none of B1–B3 is consumed.
Catalogue: tier-1 core; independence machine-checked — a 3-state
deadline-blind structure validates the remaining basis and every rule
while refuting this schema, so it is underivable. -/
theorem sound_QInvRelativize (B : Behavior X) (T : ℝ≥0) (ν χ : Set X) (x : X)
    (hbox : B.BoxInv ν x) :
    B.AService T (χ ∩ ν) x ↔ B.AService T χ x
\end{leansrc}
\ruleblock{Correspondence (Steele display $\leftrightarrow$ Lean)}
\absentbox{NO CORRESPONDENCE DICTIONARY YET}{no reviewer-certified dictionary maps the displayed rule onto the mechanized statements above; until one is written and certified, that mapping rests on this section\textquotesingle s editorial care alone.}
\ruleblock{Lean verifier receipts}
\sourcenote{module \texttt{BernsteinLogic.QtslTimedCore}; bound build \texttt{e6370ca0a2ac\ldots} (bound from receipt 6460a6dd4d227c9a, 2026-08-04T01:16:03Z). A receipt records that named checking layers accepted a build: (module, compiled-artifact SHA-256, environment SHA-256, checker layers), minted when the checking gate ran; the store receipts attesting the bound build are listed with the layers each run credited. Nothing here asks the reader to trust the store: every layer is an independently runnable program over the quoted sources --- rebuild the module and re-run the checkers to reproduce the verdict. Receipts attesting superseded builds credit no layer and are omitted from this rendering (the count of omissions, if any, is noted below).}
\begin{receipttable}
\receiptrow{6460a6dd4d227c9a}{2026-08-04T01:16:03Z}{build+decorrelated}{\badge{BadgeGreen}{current}}
\receiptrow{659fe3396a48b674}{2026-08-04T02:23:12Z}{build+decorrelated}{\badge{BadgeGreen}{current}}
\end{receipttable}
\sourcenote{4 further receipt(s) in the store attest superseded builds of this module and are omitted from this rendering.}
\layerline{\layeritem{build}{lake build (elaboration)}{\badge{BadgeGreen}{current}}\layeritem{replay}{leanchecker (official kernel replay)}{\badge{BadgeAmber}{stale only}}\layeritem{hardened}{safe\_verify (hardened screen)}{\badge{BadgeAmber}{stale only}}\layeritem{decorrelated}{lean4export + nanoda (decorrelated kernel)}{\badge{BadgeGreen}{current}}}
\sourcenote{\badge{BadgeGray}{comparator: none} no fixture in the three-implementation comparator harness (compiled / kernel-interpreted / reference-interpreted) exercises this rule.}
\ruleblock{HOL4 statement}
\sourcenote{theory \texttt{bern\_tr\_t5qdyn}, HOL4 development script \texttt{\texttt{bern\_tr\_t5qdynScript.sml}} \badge{BadgeGreen}{script bound}}
\holhead{Theorem sound\_QInvRelativize}
\begin{holsrc}
Theorem sound_QInvRelativize:
  ∀fwd bud nu chi (x:'x).
    qdBoxInv fwd nu x ⇒
    (qdAService fwd bud (chi INTER nu) x ⇔ qdAService fwd bud chi x)
\end{holsrc}
\ruleblock{HOL verifier receipts}
\badge{BadgeBlue}{documentary} \texttt{hol-strict} CLEAN: theory \texttt{bern\_tr\_t5qdyn}, 22
theorems, 0 cheat-tagged, 0 nonstandard axioms (whitelist:
\texttt{DISK\_THM}; the four \texttt{bool} axioms) --- verification
run of 2026-07-22 on an isolated build host, against a
content-addressed checkout of the exact sources cited here (git blob
\texttt{d64e1954\ldots}); an independent acceptance host ran the same
build green at the porting seating (2026-07-21). The timed core this
statement quantifies over is ported with an inventoried seam:
trajectories ride the full real line with $0 \le t$ guards at exactly
the positions where the nonnegative-time binders sit, and the
behaviour structure\textquotesingle s closure proof fields flatten to
an explicit wellformedness conjunction. This rule\textquotesingle s
twin is closure-free exactly as its Lean statement is --- stated over
a bare forward-solution set, strictly stronger than a
hypothesis-for-hypothesis transliteration.

\sourcenote{HOL receipts are documentary: quoted results of \texttt{hol-strict} runs
and census rows from the HOL4 development\textquotesingle s build
evidence, not minted receipt artefacts. A HOL receipt store (the Lean
store\textquotesingle s twin) is named future work.
}
\ruleblock{Proof}

\sourcenote{Genre: core rule (tier 1) --- an ordinary mathematical
proof at Putnam-exam-solution detail. This section also fixes the
semantic setting shared by the three timed-core sections
(\rulename{QInvRelativize}, \rulename{QFirstService},
\rulename{QFirstServiceSplit}); the later two abbreviate it as ``the
timed-core setting''.}

\textbf{Setting (shared by the timed-core sections).}\enspace A
\emph{behavioural model} is a state set $X$ with atom denotations
$\llbracket\cdot\rrbracket \subseteq X$, extended over $\neg, \wedge,
\vee$ as set operations (so $\llbracket\chi\wedge\nu\rrbracket =
\llbracket\chi\rrbracket \cap \llbracket\nu\rrbracket$, etc.), and,
for each $x \in X$, a set
$B_x$ of trajectories $\tau : \mathbb{R}_{\ge 0} \to X$, subject to
three closure conditions and nothing else: \textbf{(B1)}
$\tau \in B_x \Rightarrow \tau(0) = x$; \textbf{(B2)}
$\tau \in B_x \Rightarrow (s \mapsto \tau(t+s)) \in B_{\tau(t)}$;
\textbf{(B3)} $\tau \in B_x$ and $\tau' \in B_{\tau(t)}$ imply
$\tau \cdot_t \tau' \in B_x$, where $\tau \cdot_t \tau'$ follows
$\tau$ on $[0,t)$ and $\tau'$ (re-anchored) from $t$ on. Empty $B_x$
is permitted. The \emph{début} of $\tau$ into $E \subseteq X$ is
$\tau_E := \inf\{t : \tau(t) \in E\}$ ($\inf\emptyset = \infty$);
worst-case service and invariance are
$$x \vDash \mathsf{A}_T\varphi \iff \forall \tau \in B_x :
\tau_{\llbracket\varphi\rrbracket} \le T, \qquad
x \vDash \Box\varphi \iff \forall \tau \in B_x\ \forall t :
\tau(t) \in \llbracket\varphi\rrbracket,$$
where deadlines $T$ range over \emph{finite} nonnegative reals
(the mechanization types them so); in particular $\tau_E \le T$
forces the visit set nonempty.
For $S \subseteq \mathbb{R}_{\ge 0}$ and finite $T \ge 0$,
$\inf S \le T$ holds iff for
every $\varepsilon > 0$ some $s \in S$ has $s < T + \varepsilon$ ---
the infimum-free reading used throughout, and the form the mechanized
statements take. A début need \emph{not} be attained: nothing in B1--B3
forces a visit set to contain its infimum. The three sections split
exactly along that line: \rulename{QInvRelativize} is sound over all behavioural
models; the two first-service laws demand \emph{attained débuts} for
their obligation formulas (their sections state the attainment lemma
and the countermodel making the demand honest).

\textbf{Claim (\texttt{sound\_QInvRelativize}).}\enspace
$\Box\nu \wedge \mathsf{A}_T(\chi \wedge \nu) \;\leftrightarrow\;
\Box\nu \wedge \mathsf{A}_T\chi$ is valid in every behavioural model.

\begin{bthm}[soundness of \rulename{QInvRelativize}]\label{thm:QInvRelativize-sound}
At every state $x$ of every behavioural model:
$x \vDash \Box\nu \wedge \mathsf{A}_T(\chi\wedge\nu)$ iff
$x \vDash \Box\nu \wedge \mathsf{A}_T\chi$.
\end{bthm}
\begin{proof}
It suffices to show, for $x \vDash \Box\nu$, that
$x \vDash \mathsf{A}_T(\chi\wedge\nu) \iff x \vDash \mathsf{A}_T\chi$;
conjoining $\Box\nu$ on both sides then gives the displayed
biconditional (both sides carry it as a conjunct).

Let $\tau \in B_x$. By $\Box\nu$, every point of $\tau$ lies in
$\llbracket\nu\rrbracket$, so
$$\{t : \tau(t) \in \llbracket\chi \wedge \nu\rrbracket\}
= \{t : \tau(t) \in \llbracket\chi\rrbracket\} \cap
  \{t : \tau(t) \in \llbracket\nu\rrbracket\}
= \{t : \tau(t) \in \llbracket\chi\rrbracket\},$$
the second equality because the second intersecting set is all of
$\mathbb{R}_{\ge 0}$. Equal visit sets have equal infima:
$\tau_{\llbracket\chi\wedge\nu\rrbracket} =
\tau_{\llbracket\chi\rrbracket}$. So the two service conditions agree
trajectory by trajectory, hence at $x$.
\end{proof}

\begin{bremark}[why the invariant conjunct is everything]
Drop $\Box\nu$ and the right-to-left direction fails in a two-location
clocked model whose $\chi$-location lies outside $\nu$: the service
fact about $\chi$ says nothing about $\chi \wedge \nu$. (Left-to-right
needs no invariant --- every $\chi\wedge\nu$-visit is a $\chi$-visit,
so that direction is valid outright.) The proof uses
$\Box\nu$ once, exactly to collapse the two visit sets.
\end{bremark}

\begin{bremark}[no closure is consumed]
The proof never builds a trajectory: B1--B3 are not used, so the
schema is sound over arbitrary trajectory assignments. Its
independence witness must therefore deviate in the \emph{operator}
rather than the trajectory space, and does: the independence witness
is a machine-checked, deadline-blind three-state structure whose
service-operator interpretation deviates at an unreachable state; it
validates the seed basis --- the axiom set prior to the corrections
these sections record --- while refuting \rulename{QInvRelativize} directly.
Nothing else in the basis ties finite worst-case débuts to
reachability, which is precisely what this axiom asserts.
\end{bremark}

\begin{anchors}
Theorem \ref{thm:QInvRelativize-sound} & \texttt{sound\_QInvRelativize} (\texttt{BernsteinLogic/QtslTimedCore.lean}); semantics \texttt{Behavior.AService}, \texttt{Behavior.BoxInv}\\
independence witness & machine-checked three-state deviant structure (see the remark above)\\
\end{anchors}


\rulesection{48}{QFirstService}{core rule}
\label{blcatrule-QFirstService}
\ruleblock{The rule (Steele metanotation)}
\sourcenote{as realized in the mechanized development; source row: \texttt{\texttt{QTSL-rules.toml}} --- \emph{QFirstService}. The metanotation statement is typeset by hand for notational quality; its mechanized ground truth is the Lean statement below, extracted from the sources at every build of this book and hash-bound. The rule \emph{roster} is machine-checked (every mechanized rule tag round-trips to exactly one catalogue row); the statement \emph{text} is editorial, checked by eye against the extracted Lean beside it.}
\begin{mathpar}
\inferrule*{
  B.\mathsf{FusionClosed} \\ 0 < n \\
  \forall i.\ \mathsf{FirstVisits}(B, \alpha_i) \\
  \forall i.\ \mathsf{AService}(B, T_i, \alpha_i, x)
}{
  \mathsf{AService}\big(B,\ {\textstyle\min_i T_i},\ \mathsf{fsTarget}(\alpha, T),\ x\big)
}
\end{mathpar}
$\mathsf{fsTarget}(\alpha, T) = \bigcup_i \big(\alpha_i \cap \{y \mid \forall j \ne i.\ \mathsf{AService}(B, T_j, \alpha_j, y)\}\big)$ --- some obligation served \emph{right here} while every other still binds at its ORIGINAL budget; $\mathsf{FirstVisits}$ = nonempty visit sets have least elements (attained d\'ebuts).

\ruleflavor{certificate content out of scope of the current baked theory --- realized in Lean over fusion-closed behaviours with attained d\'ebuts.}
\par\medskip
One of the three timed-core basis schemata proved non-admissible in
the gnand completeness development (machine independence witnesses;
see the proof section's anchors), entering as tier-1 core: its
soundness obligation is semantic, discharged at Putnam-exam-solution
detail in the proof section and mechanized in
BernsteinLogic.QtslTimedCore. Sound over behavioural models with
\emph{attained d\'ebuts} for the obligation formulas --- a real side
condition: the proof section's closing remark carries the
countermodel (mechanized) witnessing failure without it.

\ruleblock{Lean statement}
\declhead{BernsteinLogic.QtslTimedCore}{sound\_QFirstService}{\badge{BadgeGreen}{statement bound}}
\begin{leansrc}
/-- **FS_n (first-service, EQ-FS)** — mechanized soundness. Over a
fusion-closed behavior with attained débuts for the obligations: if
every obligation binds at `x`, then within the tightest budget every
solution reaches a state where some obligation is served and every
other still binds at its ORIGINAL budget. Catalogue: tier-1 core
(schema family, n ≥ 2; the n = 1 instance is degenerate-derivable);
independence machine-checked — a blind 3-state structure validates the
FS-less basis over all definable sets while refuting a 2-obligation
instance of FS. Class boundary:
`qFirstService_attainment_necessary` below. -/
theorem sound_QFirstService (hB3 : B.FusionClosed) (hn : 0 < n)
    (α : Fin n → Set X) (T : Fin n → ℝ≥0) (x : X)
    (hatt : ∀ i, B.FirstVisits (α i))
    (hx : ∀ i, B.AService (T i) (α i) x) :
    B.AService (Finset.univ.inf' (Finset.univ_nonempty_iff.mpr ⟨⟨0, hn⟩⟩) T)
      (B.fsTarget α T) x
\end{leansrc}
\declhead{BernsteinLogic.QtslTimedCore}{qFirstService\_attainment\_necessary}{\badge{BadgeGreen}{statement bound}}
\begin{leansrc}
/-- **The theorem the master doc's countermodel E certifies**: there is
a fusion-closed behavior and a state satisfying the FS₂ premise at
budgets (2, 2) whose FS-target conclusion FAILS — so the attainment
hypothesis of `sound_QFirstService` is not removable, and FS_n is
invalid over raw C₀. Witness trajectory for the failure: `stay 0`;
along it every point of the FS-target would be a spine point at some
parameter in (1, 2) at which the OTHER atom still binds — but from a
rational point the escape `esc t t` never meets `Eb`, and from an
irrational point it never meets `Ea`. -/
theorem qFirstService_attainment_necessary :
    ∃ (S : Type) (B : Behavior S) (α : Fin 2 → Set S) (x : S),
      B.FusionClosed ∧
      (∀ i, B.AService 2 (α i) x) ∧
      ¬ B.AService 2 (B.fsTarget α (fun _ => 2)) x
\end{leansrc}
\ruleblock{Correspondence (Steele display $\leftrightarrow$ Lean)}
\absentbox{NO CORRESPONDENCE DICTIONARY YET}{no reviewer-certified dictionary maps the displayed rule onto the mechanized statements above; until one is written and certified, that mapping rests on this section\textquotesingle s editorial care alone.}
\ruleblock{Lean verifier receipts}
\sourcenote{module \texttt{BernsteinLogic.QtslTimedCore}; bound build \texttt{e6370ca0a2ac\ldots} (bound from receipt 6460a6dd4d227c9a, 2026-08-04T01:16:03Z). A receipt records that named checking layers accepted a build: (module, compiled-artifact SHA-256, environment SHA-256, checker layers), minted when the checking gate ran; the store receipts attesting the bound build are listed with the layers each run credited. Nothing here asks the reader to trust the store: every layer is an independently runnable program over the quoted sources --- rebuild the module and re-run the checkers to reproduce the verdict. Receipts attesting superseded builds credit no layer and are omitted from this rendering (the count of omissions, if any, is noted below).}
\begin{receipttable}
\receiptrow{6460a6dd4d227c9a}{2026-08-04T01:16:03Z}{build+decorrelated}{\badge{BadgeGreen}{current}}
\receiptrow{659fe3396a48b674}{2026-08-04T02:23:12Z}{build+decorrelated}{\badge{BadgeGreen}{current}}
\end{receipttable}
\sourcenote{4 further receipt(s) in the store attest superseded builds of this module and are omitted from this rendering.}
\layerline{\layeritem{build}{lake build (elaboration)}{\badge{BadgeGreen}{current}}\layeritem{replay}{leanchecker (official kernel replay)}{\badge{BadgeAmber}{stale only}}\layeritem{hardened}{safe\_verify (hardened screen)}{\badge{BadgeAmber}{stale only}}\layeritem{decorrelated}{lean4export + nanoda (decorrelated kernel)}{\badge{BadgeGreen}{current}}}
\sourcenote{\badge{BadgeGray}{comparator: none} no fixture in the three-implementation comparator harness (compiled / kernel-interpreted / reference-interpreted) exercises this rule.}
\ruleblock{HOL4 statement}
\sourcenote{theory \texttt{bern\_tr\_t5qdyn}, HOL4 development script \texttt{\texttt{bern\_tr\_t5qdynScript.sml}} \badge{BadgeGreen}{script bound}}
\holhead{Theorem sound\_QFirstService}
\begin{holsrc}
Theorem sound_QFirstService:
  ∀fwd al (Tb:num->real) n (x:'x).
    qdFusionClosed fwd ∧ 0 < n ∧
    (∀i. i < n ⇒ qdFirstVisits fwd (al i)) ∧
    (∀i. i < n ⇒ qdAService fwd (Tb i) (al i) x) ⇒
    qdAService fwd (rminN Tb n) (qdFsTarget fwd al Tb n) x
\end{holsrc}
\ruleblock{HOL verifier receipts}
\badge{BadgeBlue}{documentary} \texttt{hol-strict} CLEAN: theory \texttt{bern\_tr\_t5qdyn}, 22
theorems, 0 cheat-tagged, 0 nonstandard axioms (whitelist:
\texttt{DISK\_THM}; the four \texttt{bool} axioms) --- verification
run of 2026-07-22 on an isolated build host, against a
content-addressed checkout of the exact sources cited here (git blob
\texttt{d64e1954\ldots}); an independent acceptance host ran the same
build green at the porting seating (2026-07-21). The timed core this
statement quantifies over is ported with an inventoried seam:
trajectories ride the full real line with $0 \le t$ guards at exactly
the positions where the nonnegative-time binders sit, and the
behaviour structure\textquotesingle s closure proof fields flatten to
an explicit wellformedness conjunction; the fusion-closure hypothesis
is consumed exactly as displayed. The d\'ebut/attainment engine
lemmas behind the statement (first-visit and served-point existence)
are ported alongside as local instruments, mirroring their private
status on the Lean side.

\sourcenote{HOL receipts are documentary: quoted results of \texttt{hol-strict} runs
and census rows from the HOL4 development\textquotesingle s build
evidence, not minted receipt artefacts. A HOL receipt store (the Lean
store\textquotesingle s twin) is named future work.
}
\ruleblock{Proof}

\sourcenote{Genre: core rule (tier 1), in the timed-core setting fixed
at the \rulename{QInvRelativize} section. This section proves the attainment
lemma shared by the timed-core sections; the class boundary (attained
débuts are a genuine side condition) is certified by an explicit
countermodel, sketched in the final remark and mechanized alongside
the soundness lemma.}

\textbf{Claim (\texttt{sound\_QFirstService}).}\enspace For $n \ge 2$
obligations,
$$\bigwedge_{i \le n} \mathsf{A}_{T_i}\alpha_i \;\sqsubseteq\;
\mathsf{A}_{\min_i T_i}\Big(\underbrace{\bigvee_{i}\big(\alpha_i \wedge
\bigwedge_{i' \ne i}\mathsf{A}_{T_{i'}}\alpha_{i'}\big)}_{=:~\Phi,
\text{ the FS-target}}\Big)$$
is valid in every behavioural model with \emph{attained débuts} for the
$\alpha_i$: along every trajectory, a nonempty
$\alpha_i$-visit set has a least element.

\begin{blemma}[attainment]\label{lem:fs-attain}
Attained débuts hold: (i) for closed
$\llbracket\alpha_i\rrbracket$ in any model whose trajectories are
right-continuous into a topological state space; (ii) in every clocked
model, for every formula; (iii) in discrete time, for every formula.
\end{blemma}
\begin{proof}
(i) Let $E$ be closed, $\tau$ right-continuous, $V := \{t : \tau(t)
\in E\}$ nonempty, $t_0 := \inf V$. If $t_0 \notin V$, pick visits
$v_n \downarrow t_0$ with $v_n > t_0$; right-continuity at $t_0$ gives
$\tau(v_n) \to \tau(t_0)$, and since each $\tau(v_n) \in E$, every
neighbourhood of $\tau(t_0)$ meets $E$, so $\tau(t_0) \in \overline{E}
= E$ --- contradiction. (No Hausdorff hypothesis: only ``closed'' and
``converges'' are used.) (ii) A clocked trajectory is a
right-continuous step function through finitely many locations on any
bounded window, with the post-jump location taken \emph{at} the jump time,
and \emph{every} formula denotation is a union of locations --- a
subset of the finite discrete location set, hence closed; so (i)
applies with the discrete topology. Concretely: given any visit at
time $R$, the visit set restricted to $[0,R]$ is a finite union of
left-closed intervals, whose least point is the global least point.
(iii) Nonempty subsets
of $\mathbb{N}$ have least elements.
\end{proof}

\begin{bthm}[soundness of \rulename{QFirstService}]\label{thm:fs-sound}
In every behavioural model with attained débuts for
$\alpha_1,\dots,\alpha_n$, at every state $x$: if
$x \vDash \mathsf{A}_{T_i}\alpha_i$ for all $i$, then
$x \vDash \mathsf{A}_{\min_i T_i}\Phi$.
\end{bthm}
\begin{proof}
Let $\tau \in B_x$.

\emph{First services exist.} For each $i$,
$\tau_{\llbracket\alpha_i\rrbracket} \le T_i < \infty$ forces the
visit set nonempty; attainment gives its least element
$t_i = \tau_{\llbracket\alpha_i\rrbracket} \le T_i$.

\emph{The first service overall.} Put $t^* := \min_i t_i$, attained
at some $i_0$, and $y := \tau(t^*) \in
\llbracket\alpha_{i_0}\rrbracket$. Choosing $j^* \in \arg\min_i T_i$:
$t^* \le t_{j^*} \le T_{j^*} = \min_i T_i$.

\emph{Co-obligations persist at $y$, at original budgets.} Fix
$i' \ne i_0$, $\sigma \in B_y$, and $\varepsilon > 0$. The fusion
$\rho := \tau \cdot_{t^*} \sigma$ lies in $B_x$ by B3 (typed by B1:
$\sigma(0) = y = \tau(t^*)$), so
$\rho_{\llbracket\alpha_{i'}\rrbracket} \le T_{i'}$, giving a
$\rho$-visit of $\alpha_{i'}$ at some $s < T_{i'} + \varepsilon$. If
$s < t^*$ then $\rho(s) = \tau(s)$, so $t_{i'} \le s < t^* \le
t_{i'}$ --- impossible. Hence $s \ge t^*$ and $\rho(s) =
\sigma(s - t^*)$: a $\sigma$-visit at time $s - t^* \le s <
T_{i'} + \varepsilon$. As $\varepsilon$ was arbitrary,
$\sigma_{\llbracket\alpha_{i'}\rrbracket} \le T_{i'}$. (Indeed
$\le T_{i'} - t^*$; the schema states the weaker original-budget
bound, which is what survives grid transport --- the residual-split
law \rulename{QFirstServiceSplit} is where the true residual is
recovered.)

\emph{Assembly.} $\sigma$ was arbitrary, so $y \vDash
\mathsf{A}_{T_{i'}}\alpha_{i'}$ for every $i' \ne i_0$; with
$y \in \llbracket\alpha_{i_0}\rrbracket$, $y \in
\llbracket\Phi\rrbracket$, and $\tau$ visits it at
$t^* \le \min_i T_i$.
\end{proof}

\begin{bremark}[each hypothesis, remarked]
B1 types the fusion; B3 is the engine (facts about all of $B_y$ reach
back to the premise at $x$ only through fusion); B2 is not consumed.
Attainment is consumed exactly once, and essentially --- to produce
the service point $y = \tau(t^*)$: an infimum is not a visit. (The
placement of $s$ then runs on the ordinary lower-bound property of
$t_{i'}$.) It is not removable: see the next remark. At
$n = 1$ the instance degenerates and is derivable; the schema family
starts at $n = 2$.
\end{bremark}

\begin{bremark}[the class boundary is real]\label{rem:fs-cm}
Over bare B1--B3 models the schema is \emph{false}. Countermodel (mechanized
as \texttt{qFirstService\_attainment\_necessary}): a spine
$\{y_p : p \ge 0\}$ with $\llbracket a \rrbracket$ the rational
$p \in (1,2)$ and $\llbracket b \rrbracket$ the irrational ones; from
every $y_p$, the trajectories are the spine-ride and, for each
$s \ge p$ with $s > 1$, a ride-to-$s$-then-exit to an atom-free sink.
The family is fusion-closed; every trajectory from $y_0$ meets both
atoms before time $2$ (density), so the premise
$\mathsf{A}_2 a \wedge \mathsf{A}_2 b$ holds; but at every
$y_q \in \llbracket a \rrbracket$ (so $q \in (1,2)$, making the
immediate exit legal) that exit kills $\mathsf{A}_2 b$, at every
$y_\iota \in \llbracket b \rrbracket$ it kills $\mathsf{A}_2 a$, and
states outside the two atoms fail both disjuncts outright, so the
FS-target denotes
$\emptyset$ and the conclusion fails. Dense interleaving of disjoint
atoms along one trajectory is exactly what right-continuity into
closed denotations forbids (Lemma \ref{lem:fs-attain}(i)) --- the
countermodel lives strictly outside the attained class, as it must.
\end{bremark}

\begin{anchors}
Theorem \ref{thm:fs-sound} & \texttt{sound\_QFirstService}; class mixin \texttt{Behavior.FusionClosed}, \texttt{Behavior.FirstVisits}; target \texttt{Behavior.fsTarget}\\
Lemma \ref{lem:fs-attain} & prose (this section); clocked/discrete instances discharge \texttt{FirstVisits} per model\\
Remark \ref{rem:fs-cm} & \texttt{qFirstService\_attainment\_necessary} (the countermodel, mechanized)\\
independence witness & machine-checked structure validating the corrected basis with this schema removed while refuting the schema itself\\
\end{anchors}


\rulesection{49}{QFirstServiceSplit}{core rule}
\label{blcatrule-QFirstServiceSplit}
\ruleblock{The rule (Steele metanotation)}
\sourcenote{as realized in the mechanized development; source row: \texttt{\texttt{QTSL-rules.toml}} --- \emph{QFirstServiceSplit}. The metanotation statement is typeset by hand for notational quality; its mechanized ground truth is the Lean statement below, extracted from the sources at every build of this book and hash-bound. The rule \emph{roster} is machine-checked (every mechanized rule tag round-trips to exactly one catalogue row); the statement \emph{text} is editorial, checked by eye against the extracted Lean beside it.}
\begin{mathpar}
\inferrule*{
  B.\mathsf{FusionClosed} \\ 0 < n \\ \forall i.\ c < T_i \\
  \forall i.\ \mathsf{FirstVisits}(B, \alpha_i) \\
  \forall i.\ \mathsf{AService}(B, T_i, \alpha_i, x)
}{
  \mathsf{AService}\big(B,\ c,\ \mathsf{fsTarget}(\alpha, T) \cup \mathsf{residualCell}(\alpha, T, c),\ x\big)
}
\end{mathpar}
$\mathsf{residualCell}(\alpha, T, c) = \{y \mid \forall i.\ \mathsf{AService}(B, T_i - c, \alpha_i, y)\}$ --- the no-service-by-$c$ branch pins elapsed time by the cutoff, recovering TRUE residuals $T_i - c$ where FS$_n$ can only re-issue original budgets.

\ruleflavor{certificate content out of scope of the current baked theory --- realized in Lean over fusion-closed behaviours with attained d\'ebuts.}
\par\medskip
One of the three timed-core basis schemata proved non-admissible in
the gnand completeness development --- for this law the machine
independence witnesses cover integer and half-integer deadline menus,
with finer denominators open (see the proof section's closing remark
and anchors) --- entering as tier-1 core: its soundness obligation is
semantic, discharged at Putnam-exam-solution detail in the proof
section and mechanized in BernsteinLogic.QtslTimedCore. Sound over
behavioural models with \emph{attained d\'ebuts} for the obligation
formulas --- a real side condition; its witnessing countermodel is
prose-only (not yet mechanized; the QFirstService countermodel is the
mechanized one).

\ruleblock{Lean statement}
\declhead{BernsteinLogic.QtslTimedCore}{sound\_QFirstServiceSplit}{\badge{BadgeGreen}{statement bound}}
\begin{leansrc}
/-- **FS-R (residual-split first-service)** — mechanized soundness,
same class as FS_n. For any cutoff `c` strictly below every budget:
within `c`, every solution reaches EITHER a first-service state
(original co-budgets — the FS-target) OR the residual cell, where
every obligation holds at its true residual `T i - c`. Genuinely new
on top of FS_n; independence machine-checked — a minimal 6-state
witness settles every all-integer deadline menu, and two further
certified 6-state witnesses extend the underivability to all
half-integer menus. -/
theorem sound_QFirstServiceSplit (hB3 : B.FusionClosed) (hn : 0 < n)
    (α : Fin n → Set X) (T : Fin n → ℝ≥0) (c : ℝ≥0) (x : X)
    (hc : ∀ i, c < T i)
    (hatt : ∀ i, B.FirstVisits (α i))
    (hx : ∀ i, B.AService (T i) (α i) x) :
    B.AService c (B.fsTarget α T ∪ B.residualCell α T c) x
\end{leansrc}
\ruleblock{Correspondence (Steele display $\leftrightarrow$ Lean)}
\absentbox{NO CORRESPONDENCE DICTIONARY YET}{no reviewer-certified dictionary maps the displayed rule onto the mechanized statements above; until one is written and certified, that mapping rests on this section\textquotesingle s editorial care alone.}
\ruleblock{Lean verifier receipts}
\sourcenote{module \texttt{BernsteinLogic.QtslTimedCore}; bound build \texttt{e6370ca0a2ac\ldots} (bound from receipt 6460a6dd4d227c9a, 2026-08-04T01:16:03Z). A receipt records that named checking layers accepted a build: (module, compiled-artifact SHA-256, environment SHA-256, checker layers), minted when the checking gate ran; the store receipts attesting the bound build are listed with the layers each run credited. Nothing here asks the reader to trust the store: every layer is an independently runnable program over the quoted sources --- rebuild the module and re-run the checkers to reproduce the verdict. Receipts attesting superseded builds credit no layer and are omitted from this rendering (the count of omissions, if any, is noted below).}
\begin{receipttable}
\receiptrow{6460a6dd4d227c9a}{2026-08-04T01:16:03Z}{build+decorrelated}{\badge{BadgeGreen}{current}}
\receiptrow{659fe3396a48b674}{2026-08-04T02:23:12Z}{build+decorrelated}{\badge{BadgeGreen}{current}}
\end{receipttable}
\sourcenote{4 further receipt(s) in the store attest superseded builds of this module and are omitted from this rendering.}
\layerline{\layeritem{build}{lake build (elaboration)}{\badge{BadgeGreen}{current}}\layeritem{replay}{leanchecker (official kernel replay)}{\badge{BadgeAmber}{stale only}}\layeritem{hardened}{safe\_verify (hardened screen)}{\badge{BadgeAmber}{stale only}}\layeritem{decorrelated}{lean4export + nanoda (decorrelated kernel)}{\badge{BadgeGreen}{current}}}
\sourcenote{\badge{BadgeGray}{comparator: none} no fixture in the three-implementation comparator harness (compiled / kernel-interpreted / reference-interpreted) exercises this rule.}
\ruleblock{HOL4 statement}
\sourcenote{theory \texttt{bern\_tr\_t5qdyn}, HOL4 development script \texttt{\texttt{bern\_tr\_t5qdynScript.sml}} \badge{BadgeGreen}{script bound}}
\holhead{Theorem sound\_QFirstServiceSplit}
\begin{holsrc}
Theorem sound_QFirstServiceSplit:
  ∀fwd al (Tb:num->real) n c (x:'x).
    qdFusionClosed fwd ∧ 0 < n ∧ 0 ≤ c ∧
    (∀i. i < n ⇒ c < Tb i) ∧
    (∀i. i < n ⇒ qdFirstVisits fwd (al i)) ∧
    (∀i. i < n ⇒ qdAService fwd (Tb i) (al i) x) ⇒
    qdAService fwd c
      (qdFsTarget fwd al Tb n UNION qdResidualCell fwd al Tb n c) x
\end{holsrc}
\ruleblock{HOL verifier receipts}
\badge{BadgeBlue}{documentary} \texttt{hol-strict} CLEAN: theory \texttt{bern\_tr\_t5qdyn}, 22
theorems, 0 cheat-tagged, 0 nonstandard axioms (whitelist:
\texttt{DISK\_THM}; the four \texttt{bool} axioms) --- verification
run of 2026-07-22 on an isolated build host, against a
content-addressed checkout of the exact sources cited here (git blob
\texttt{d64e1954\ldots}); an independent acceptance host ran the same
build green at the porting seating (2026-07-21). The timed core this
statement quantifies over is ported with an inventoried seam:
trajectories ride the full real line with $0 \le t$ guards at exactly
the positions where the nonnegative-time binders sit, and the
behaviour structure\textquotesingle s closure proof fields flatten to
an explicit wellformedness conjunction. The residual
budget\textquotesingle s truncated subtraction is rendered as the
max-$0$ coercion --- agreeing with the nonnegative-reals reading on
the whole consulted domain, not only under the split
theorem\textquotesingle s guard.

\sourcenote{HOL receipts are documentary: quoted results of \texttt{hol-strict} runs
and census rows from the HOL4 development\textquotesingle s build
evidence, not minted receipt artefacts. A HOL receipt store (the Lean
store\textquotesingle s twin) is named future work.
}
\ruleblock{Proof}

\sourcenote{Genre: core rule (tier 1), in the timed-core setting fixed
at the \rulename{QInvRelativize} section, with the attainment lemma of the
\rulename{QFirstService} section. This is the fourth and final basis
correction from the gnand completeness analysis: the law that makes
the decision
procedure's deadline-tightening step derivable.}

\textbf{Claim (\texttt{sound\_QFirstServiceSplit}).}\enspace For $n \ge 1$
obligations and any cutoff $0 < c < \min_i T_i$,
$$\bigwedge_{i \le n} \mathsf{A}_{T_i}\alpha_i \;\sqsubseteq\;
\mathsf{A}_c\Big(\Phi \;\vee\; \underbrace{\bigwedge_i
\mathsf{A}_{T_i - c}\,\alpha_i}_{=:~P,\text{ the residual cell}}\Big)$$
is valid in every behavioural model with attained débuts for the
$\alpha_i$, where $\Phi$ is the FS-target of \rulename{QFirstService}
(read at $n = 1$ with the empty co-obligation conjunction as $\top$,
so $\Phi = \alpha_1$).
Reading: within the cutoff, either some first service has happened
(and the co-obligations persist at original budgets), or nothing has
been served and every obligation holds at its \emph{true} residual.

\begin{bthm}[soundness of \rulename{QFirstServiceSplit}]\label{thm:fsr-sound}
In every behavioural model with attained débuts for
$\alpha_1,\dots,\alpha_n$, at every state $x$: if
$x \vDash \mathsf{A}_{T_i}\alpha_i$ for all $i$ and
$0 < c < \min_i T_i$, then $x \vDash \mathsf{A}_c(\Phi \vee P)$.
\end{bthm}
\begin{proof}
Let $\tau \in B_x$. As in the \rulename{QFirstService} proof,
attainment gives each obligation a least visit time
$t_i \le T_i$ along $\tau$; put $t^* := \min_i t_i$, attained at
$i_0$. Split on the cutoff.

\emph{Case $t^* \le c$.} The \rulename{QFirstService} argument,
verbatim, puts $y := \tau(t^*)$ in $\llbracket\Phi\rrbracket$: it
serves $\alpha_{i_0}$, and for $i' \ne i_0$ every $\sigma \in B_y$
inherits $\sigma_{\llbracket\alpha_{i'}\rrbracket} \le T_{i'}$ through
the fusion $\tau \cdot_{t^*} \sigma \in B_x$, the prefix being
$\alpha_{i'}$-free below $t^*$ by minimality. So $\tau$ visits
$\llbracket\Phi \vee P\rrbracket$ at $t^* \le c$.

\emph{Case $t^* > c$.} Let $y := \tau(c)$; we show
$y \in \llbracket P \rrbracket$, so that $\tau$ visits the disjunction
at time $c$. Fix $i$, $\sigma \in B_y$, $\varepsilon > 0$. Fuse:
$\rho := \tau \cdot_c \sigma \in B_x$ by B3, so there is a
$\rho$-visit of $\alpha_i$ at some $s < T_i + \varepsilon$. The prefix
is service-free through $c$: for $s' < c$, $\rho(s') = \tau(s')$, and
a visit there would give $t_i \le s' < c < t^* \le t_i$ ---
impossible; and the point at $c$ itself is $\rho(c) = \sigma(0) =
\tau(c)$, likewise unvisited (a visit at $c$ would give
$t_i \le c < t^*$). So $s \ge c$ and $\rho(s) = \sigma(s - c)$: a
$\sigma$-visit at time $s - c < (T_i - c) + \varepsilon$ (subtract
$c$ across $s < T_i + \varepsilon$, using $c < T_i$ to keep the
residual budget a genuine nonnegative quantity). As $\varepsilon$ was
arbitrary, $\sigma_{\llbracket\alpha_i\rrbracket} \le T_i - c$.
\end{proof}

\begin{bremark}[each hypothesis, remarked]
$c < \min_i T_i$ keeps every residual budget strictly positive ---
the regime the law is stated for; at $c = \min_i T_i$ the split
degenerates (under attained débuts the served case always fires,
since $t^* \le \min_i T_i$) and the residual row would carry
$\mathsf{A}_0$. Attainment
is consumed only through $t^*$ --- and only the \emph{served} case consumes
it: the no-service case uses the visit $s$ directly, never its
infimum. That asymmetry is honest and sharp: over bare B1--B3 models
the residual branch alone survives, and the schema as a whole fails
--- the countermodel (prose; not mechanized) retunes the
\rulename{QFirstService} countermodel --- the dense cluster moved, the
escape-legality threshold lowered to match --- and adds ``lazy
wings'': branches
postponing every service to just inside the original budgets, legal
precisely because the premise pins only original budgets before the
cluster, defeating
target and residual simultaneously along the spine. Attained débuts
restore the dichotomy the case split runs on.
\end{bremark}

\begin{bremark}[what the law buys]
\rulename{QFirstService} can only re-issue original budgets at the
service point; \rulename{QFirstServiceSplit} recovers the true residual
$T_i - c$ on the no-service branch, where the elapsed time is pinned
by the cutoff rather than by an unknowable service instant. Iterated
at successive cutoffs this derives the decision procedure's
deadline-tightening
step exactly (its descent lemma), which is what closed the
derivation side of the gnand completeness theorem. Independence: a
six-state witness proves the law underivable from the
FS-corrected basis for integer deadline menus, extended to
half-integer menus by two certified witnesses; independence at finer
denominators is open, and citation is bounded accordingly.
\end{bremark}

\begin{anchors}
Theorem \ref{thm:fsr-sound} & \texttt{sound\_QFirstServiceSplit}; residual \texttt{Behavior.residualCell}\\
countermodel (class boundary) & prose only (not yet mechanized); sketched in the first remark\\
independence witnesses & six-state minimal witness + two half-integer-menu witnesses, machine-checked (see the closing remark for the citation bound)\\
\end{anchors}


\blcatsection{Streett Rules and the Specification Surface}
\rulesection{50}{QStreettDegenerate}{alias for \textcolor{DiffBlue}{\rulename{QStreettBattery}}}
\label{blcatrule-QStreettDegenerate}
\ruleblock{The rule (Steele metanotation)}
\sourcenote{as realized in the mechanized development; source row: \texttt{\texttt{QTSL-rules.toml}} --- \emph{QStreettDegenerate}. The metanotation statement is typeset by hand for notational quality; its mechanized ground truth is the Lean statement below, extracted from the sources at every build of this book and hash-bound. The rule \emph{roster} is machine-checked (every mechanized rule tag round-trips to exactly one catalogue row); the statement \emph{text} is editorial, checked by eye against the extracted Lean beside it.}
\begin{mathpar}
\inferrule*{
  \text{one \rulename{QStreettBattery} node at } \rho_i := \mathsf{STrue} \ \text{for every pair} \\
  \text{(absorbing families and } J_i \ \text{syntactically empty; ranking decay off } [\![\varphi_i]\!] \ \text{alone)}
}{
  \mathcal{M},\ x_0 \in \{V_0 \le y\} \vDash \mathsf{SProbGE}\bigl(p,\ \mathsf{PDwellStreett}(\mathsf{PairBase}(\mathsf{STrue}, \varphi_1), \ldots, \mathsf{PairBase}(\mathsf{STrue}, \varphi_k))\bigr)
}
\end{mathpar}

\ruleflavor{the pair at request $\top$ means $\mathrm{Inf}([\![\varphi_i]\!])$ --- the classical recurrence certificate; the degenerate-pair specialization erases the tautological rows before emission, so the emitted kit is the recurrence kit row for row.}
\par\medskip
\rulename{QStreettDegenerate} is the recurrence rule in its compact
spelling: to conclude, with probability at least $p$, that every response
region is visited infinitely often ($\mathsf{GF}$-shaped claims, the
classical B\"uchi/recurrence certificate), instantiate the Streett bundle
at the trivially-true request. It is an \emph{alias} (n\'e StreettDegenerate --- the first sub-rule of the battery row): one
\rulename{QStreettBattery} application, at request regions equal to the
whole state space, is the entire unpacking --- the emptiness it exploits
is syntactic, not semantic. When every state requests, the
finitely-many-requests disjunct of each pair is empty along the tail
filter, the absorbing-certificate family and its carved-out sets
$J_i$ are vacuous, and the ranking decay row is demanded off the
response region alone --- so the degenerate-pair specialization erases
exactly the tautological rows, and what the checker replays is the
recurrence certificate, row for row. A producer uses it wherever the
specification is recurrence rather than
request-response fairness; nothing about the certificate synthesis
changes except that the request machinery never materializes.

Scope, stated plainly: the
bound statement below is the degenerate-instance witness
\texttt{sound\_QStreettDegenerate} (module
\texttt{BernsteinLogicAux.QStreett} --- one application of
\texttt{sound\_QStreettBattery} at request set \texttt{univ}, whose
finitely-many-requests disjunct is empty along the tail filter; its
conclusion is the recurrence event
$\Pr[\forall i,\ \mathrm{Inf}(G_i)] \ge p$, exactly the amended
display\textquotesingle s claim). The witness covers the rule at the
\emph{discrete-kernel profile} (a time-homogeneous Markov chain
presented by a single step kernel); this entry once carried a
conditional mark inherited from the bundle\textquotesingle s
continuous-time residue, and it goes the way the
bundle\textquotesingle s went: at the martingale-problem profile the
bundle is discharged whole --- strict-decrement rows included,
hash-bound under the \rulename{QStreettBattery} entry --- and the
alias inherits that discharge the same way it inherits everything
else, as one bundle application at request set \texttt{univ}, per the
bundle section\textquotesingle s profile remark. The payload-level
unpack function is not among this edition\textquotesingle s bound
artefacts: the alias deliberately has no wire sort of its own
(producers elaborate before emission), and the transport table in the
proof is stated at the catalogued-field level.

\cataloguenote{The degenerate pair is also the honest cost anchor of the
Streett-primitive design: on recurrence-shaped workloads the bundle
emits the same rows a dedicated recurrence rule would, so the general
rule costs nothing where the special case suffices --- while the reverse
specialization (recurrence-primitive, Streett by automaton product)
would pay exponentially in the number of pairs.
}
\ruleblock{Lean statement}
\declhead{BernsteinLogicAux.QStreett}{sound\_QStreettDegenerate}{\badge{BadgeGreen}{statement bound}}
\begin{leansrc}
/-- **`QStreettDegenerate` (T9a)**: request `= ⊤` degenerates the pair to
`Inf(Gᵢ)` — the classical recurrence certificate, one `sound_QStreettBattery`
application with the request set `univ` (whose `Fin` is empty along `atTop`).
Alias tier: the statement is the bundle's, specialized and simplified. -/
theorem sound_QStreettDegenerate {G : ι → Set S} {W : ι → S → ℝ} {ε C : ι → ℝ}
    (hD : MeasurableSet D) (hDbar : MeasurableSet Dbar) (hU : MeasurableSet U)
    (hV₀ : Measurable V₀)
    (hG : ∀ i, MeasurableSet (G i)) (hW : ∀ i, Measurable (W i))
    (hΘ : p + thr ≤ 1)
    (hDsub : D ⊆ Dbar)
    (hR3 : ∀ y ∈ D, y ∉ U → K y Dbar = 1)
    (hR4 : ∀ y ∈ D, y ∉ U → K y (D ∪ U) = 1)
    (hV₀int : ∀ y ∈ D, y ∉ U → Integrable V₀ (K y))
    (hV₀drift : ∀ y ∈ D, y ∉ U → ∫ z, V₀ z ∂K y ≤ V₀ y)
    (hV₀floor : ∀ y ∈ U, y ∈ Dbar → 1 ≤ V₀ y)
    (hV₀nonneg : ∀ y ∈ Dbar, 0 ≤ V₀ y)
    (hε : ∀ i, 0 < ε i)
    (hW0 : ∀ i, ∀ y ∈ Dbar, 0 ≤ W i y) (hWC : ∀ i, ∀ y ∈ Dbar, W i y ≤ C i)
    (hWint : ∀ i, ∀ y ∈ D, y ∉ U → y ∉ G i → Integrable (W i) (K y))
    (hWdriftR : ∀ i, ∀ y ∈ D, y ∉ U → y ∉ G i → ∫ z, W i z ∂K y ≤ W i y - ε i)
    (hx : x ∈ D) (hthr : V₀ x ≤ thr) :
    ENNReal.ofReal p ≤ pathMeasure K x
      {ω | ∀ i, ∃ᶠ n in atTop, ω n ∈ G i}
\end{leansrc}
\par\medskip
\textbf{The general-stochastic instance.}\enspace The degenerate
reading needs no separate statement at the segment-model tier: it rides
the confined form of the bundle instance below with the vacuous premise
sets --- the same instantiation-at-the-skeleton the bundle
section\textquotesingle s general-stochastic block describes.

\declhead{BernsteinLogicAux.LayerD.Streett}{sound\_GSHS\_streett\_bundle\_confined}{\badge{BadgeGreen}{statement bound}}
\begin{leansrc}
/-- **`sound_QStreettBattery`, GSHS instance, confined form** (T9; assault
doc §2.7 Theorem G at the conservative ranking scope, mechanized
restart-free — module docstring for why G1/G2's epoch machinery is not
consumed). Hypothesis fields are segment-level and per-model: the three
containment rows, the θ = 0 confinement field for `V₀`, and per pair the
scoped ranking fields for `W i` (nonincrease off the response region,
`ε i`-decrease at request entries — per-step shape, module docstring for
the rate-form conversion). Conclusion: with probability at least `p` the
skeleton chain of `Rep(P)` stays alive in `D \ U` forever **and** every
certified pair satisfies the chain-level Streett condition. Semantic
transport of the chain events: module docstring (T9b monotonicity for
requests, entry reads for responses, W-Z for time-cofinality). -/
theorem sound_GSHS_streett_bundle_confined
    [IsMarkovKernel M.seg] [IsMarkovKernel M.jump]
    (hD : MeasurableSet D) (hDbar : MeasurableSet Dbar) (hU : MeasurableSet U)
    (hV₀ : Measurable V₀)
    (hRq : ∀ i, MeasurableSet (Rq i)) (hG : ∀ i, MeasurableSet (G i))
    (hW : ∀ i, Measurable (W i))
    (hΘ : p + thr ≤ 1) (hDsub : D ⊆ Dbar)
    (hLive : NoVanish M (D \ U)) (hHalt : HaltsCaught M (D \ U) (U ∩ Dbar))
    (hReset : ResetsInto M (D \ U) ((D ∪ U) ∩ Dbar))
    (hSegD₀ : RankingSegmentDynkin M (D \ U) V₀ 0)
    (hForce₀ : ForcedNonexpansion M V₀) (hNonneg₀ : NonnegV V₀)
    (hC₀0 : 0 ≤ C₀) (hV₀C : ∀ y, V₀ y ≤ C₀)
    (hFloor₀ : ∀ y ∈ U, y ∈ Dbar → 1 ≤ V₀ y)
    (hε : ∀ i, 0 < ε i)
    (hSegDW : ∀ i, RankingSegmentDynkin M ((D \ U) \ G i) (W i) 0)
    (hSegDWR : ∀ i, RankingSegmentDynkin M (((D \ U) ∩ Rq i) \ G i) (W i) (ε i))
    (hForceW : ∀ i, ForcedNonexpansion M (W i)) (hNonnegW : ∀ i, NonnegV (W i))
    (hC0 : ∀ i, 0 ≤ C i) (hWC : ∀ i, ∀ y, W i y ≤ C i)
    (hx : x ∈ D) (hthr : V₀ x ≤ thr) :
    ENNReal.ofReal p ≤ Battery.pathMeasure (skel M) (SkelState.alive x)
      {ω | (∀ n, ω n ∈ skelIn D \ skelBad U) ∧
        ∀ i, (∃ᶠ n in atTop, ω n ∈ skelIn (Rq i)) →
          ∃ᶠ n in atTop, ω n ∈ skelIn (G i)}
\end{leansrc}
\ruleblock{Correspondence (Steele display $\leftrightarrow$ Lean)}
\badge{BadgeGreen}{reviewer-certified against the displayed texts}
\sourcenote{The dictionary below is a reviewer-certified statement of how the mechanized statements above interpret the displayed rule: display metavariable $\mapsto$ Lean binder, each display hypothesis $\mapsto$ the premise or side condition carrying it, the display conclusion $\mapsto$ the goal --- with any deliberate difference in width between display and mechanization named outright. Certification is recorded as a hash triple over (the displayed cell, the extracted statement texts, this dictionary\textquotesingle s own text); if any of the three changes, this block turns amber until a reviewer re-certifies. The correspondence claim itself remains editorial --- reviewer-grade, not machine-checked; its machine-checked successor types each soundness statement against the rule as a structured object, and supersedes this block per rule as it lands.}
The displayed conclusion $\mathcal{M}
,\ x_0 \in \{V_0 \le y\} \vDash
\mathsf{SProbGE}\bigl(p,\ \mathsf{PDwellStreett}(\mathsf{PairBase}(\mathsf{STrue},
\varphi_1), \ldots, \mathsf{PairBase}(\mathsf{STrue}, \varphi_k))\bigr)$
$\mapsto$ \texttt{ofReal p ≤ pathMeasure K x \{ω | ∀ i, ∃ᶠ n in atTop,
ω n ∈ G i\}} for every start with \texttt{x ∈ D} and \texttt{V₀ x ≤
thr} (the Lean source\textquotesingle s identifier for the
display\textquotesingle s $y$) --- the recurrence event, with the
$\mathsf{STrue}$-requests \emph{erased
at the statement level}, not merely in the proof: no request binder
exists in \texttt{sound\_QStreettDegenerate} at all, which is the
display\textquotesingle s ``absorbing families and $J_i$ syntactically
empty'' carried into the type. The alias tier is visible in the proof
term: one \texttt{sound\_QStreettBattery} application at \texttt{Rq :=
fun \_ => Set.univ}, whose frequently-clause is discharged by
\texttt{univ}-membership.
\begin{itemize}
\item the inherited $V_0$ QSafetyBattery kit $\mapsto$ the four
  kernel-integral rows verbatim: \texttt{hV₀drift} (interior),
  \texttt{hV₀floor} (floor on $U \cap \bar{D}$), \texttt{hV₀nonneg}
  (on $\bar{D}$), with \texttt{hV₀int} the undisplayed integrability
  rider, plus the hull/fail-closed side conditions
  \texttt{hR3}/\texttt{hR4} and the threshold \texttt{hΘ : p + thr ≤
  1} --- the battery pattern, unchanged.
\item ``ranking decay off $[\![\varphi_i]\!]$ alone'' $\mapsto$
  \texttt{hWdriftR : ∀ i, ∀ y ∈ D, y ∉ U → y ∉ G i → ∫ z, W i z ∂K y
  ≤ W i y - ε i} --- one generator row where the bundle carries two: the
  request-conditioned decay and the off-response non-increase
  coincide when every state is a request, and the statement asks
  only the merged row (the proof re-splits it for the bundle).
\item nonneg + upper-bound rows $\mapsto$ \texttt{hW0}/\texttt{hWC}
  on $\bar{D}$; strictness $\mapsto$ \texttt{hε : ∀ i, 0 < ε i} ---
  note the specialized statement takes per-pair \texttt{ε i}, mildly
  wider than the display\textquotesingle s single $\varepsilon_0$.
\item the measurability pack
  (\texttt{hD}/\texttt{hDbar}/\texttt{hU}/\texttt{hV₀}/\texttt{hG}/
  \texttt{hW}) is undisplayed profile regularity, per the battery
  precedent.
\end{itemize}
Width note: this alias\textquotesingle s warrant is inherited from
the bundle at both profiles (hash-bound at the discrete-kernel
profile; the martingale-problem discharge, strict rows included,
hash-bound under the bundle entry); the producer-side erasure of
tautological rows before emission is compile behaviour with no Lean
counterpart --- correctly, since the certified statement never
receives the tautological rows to begin with.

\ruleblock{Lean verifier receipts}
\sourcenote{module \texttt{BernsteinLogicAux.QStreett}; bound build \texttt{469881beb770\ldots} (bound from receipt 9d1d848e0ba1718a, 2026-07-29T01:03:31Z). A receipt records that named checking layers accepted a build: (module, compiled-artifact SHA-256, environment SHA-256, checker layers), minted when the checking gate ran; the store receipts attesting the bound build are listed with the layers each run credited. Nothing here asks the reader to trust the store: every layer is an independently runnable program over the quoted sources --- rebuild the module and re-run the checkers to reproduce the verdict. Receipts attesting superseded builds credit no layer and are omitted from this rendering (the count of omissions, if any, is noted below).}
\begin{receipttable}
\receiptrow{e3990fe544d4708d}{2026-07-28T23:48:35Z}{build+decorrelated}{\badge{BadgeGreen}{current}}
\receiptrow{68e9374fa1b418fe}{2026-07-29T00:00:34Z}{build+decorrelated}{\badge{BadgeGreen}{current}}
\receiptrow{e5895c9706854cd3}{2026-07-29T00:49:04Z}{build+replay+hardened+decorrelated}{\badge{BadgeGreen}{current}}
\receiptrow{9d1d848e0ba1718a}{2026-07-29T01:03:31Z}{build+replay+hardened+decorrelated}{\badge{BadgeGreen}{current}}
\receiptrow{4a8d377b721b8523}{2026-07-31T22:00:09Z}{build+decorrelated}{\badge{BadgeGreen}{current}}
\receiptrow{d251a122540ab741}{2026-07-31T22:29:03Z}{build+decorrelated}{\badge{BadgeGreen}{current}}
\receiptrow{333635b8e047113a}{2026-08-02T23:31:14Z}{build+decorrelated}{\badge{BadgeGreen}{current}}
\receiptrow{9e96f80d63e4ad82}{2026-08-03T16:27:18Z}{build+decorrelated}{\badge{BadgeGreen}{current}}
\receiptrow{932b4b9e0b0fd323}{2026-08-03T22:46:02Z}{build+decorrelated}{\badge{BadgeGreen}{current}}
\end{receipttable}
\sourcenote{9 further receipt(s) in the store attest superseded builds of this module and are omitted from this rendering.}
\sourcenote{module \texttt{BernsteinLogicAux.LayerD.Streett}; bound build \texttt{7393065bbc7f\ldots} (bound from receipt abac471866cfcfd1, 2026-07-21T05:38:48Z). A receipt records that named checking layers accepted a build: (module, compiled-artifact SHA-256, environment SHA-256, checker layers), minted when the checking gate ran; the store receipts attesting the bound build are listed with the layers each run credited. Nothing here asks the reader to trust the store: every layer is an independently runnable program over the quoted sources --- rebuild the module and re-run the checkers to reproduce the verdict. Receipts attesting superseded builds credit no layer and are omitted from this rendering (the count of omissions, if any, is noted below).}
\begin{receipttable}
\receiptrow{0f3b50b4d3f0d33c}{2026-07-20T22:52:50Z}{build+decorrelated}{\badge{BadgeGreen}{current}}
\receiptrow{109669f848453ec2}{2026-07-21T00:03:19Z}{build+decorrelated}{\badge{BadgeGreen}{current}}
\receiptrow{7e17aa32b276623b}{2026-07-21T04:51:49Z}{build+decorrelated}{\badge{BadgeGreen}{current}}
\receiptrow{abac471866cfcfd1}{2026-07-21T05:38:48Z}{build+replay+hardened+decorrelated}{\badge{BadgeGreen}{current}}
\end{receipttable}
\sourcenote{6 further receipt(s) in the store attest superseded builds of this module and are omitted from this rendering.}
\layerline{\layeritem{build}{lake build (elaboration)}{\badge{BadgeGreen}{current}}\layeritem{replay}{leanchecker (official kernel replay)}{\badge{BadgeGreen}{current}}\layeritem{hardened}{safe\_verify (hardened screen)}{\badge{BadgeGreen}{current}}\layeritem{decorrelated}{lean4export + nanoda (decorrelated kernel)}{\badge{BadgeGreen}{current}}}
\sourcenote{\badge{BadgeGray}{comparator: none} no fixture in the three-implementation comparator harness (compiled / kernel-interpreted / reference-interpreted) exercises this rule.}
\ruleblock{HOL4 statement}
\absentbox{OUT OF SCOPE BY DESIGN}{Out of scope for the HOL side \emph{by design}, not by backlog: the rule
is a producer-side respelling that unpacks to a single
\rulename{QStreettBattery} node before any certificate is emitted, so no
wire object can ever contain a \rulename{QStreettDegenerate} node --- the
HOL statement that would verify a degenerate-request bundle is the
bundle\textquotesingle s own (itself absent today, for the reason given
in that section: the satisfaction stratum has no HOL artefact).
Checked against the HOL sources: the
\texttt{bern\_catalogue} theory is the base-calculus checker cone and
carries nothing Streett-shaped.
}
\ruleblock{HOL verifier receipts}
\absentbox{NONE}{no HOL verification receipt exists for this rule.}
\ruleblock{Proof}

\sourcenote{Genre: alias rule (tier 3) --- a respelling template. This
section additionally states, once for the whole batch of
satisfaction-level aliases, the general template those sections share;
the \rulename{QStreettCmplRequest}, \rulename{QStreettTargetCellCollapse},
\rulename{QNnfRewrite}, \rulename{QStreettRequestNnf},
\rulename{QAvoid}, \rulename{QStreettGF}, \rulename{QStreettFG},
and \rulename{QNegFAlias} sections reference it by
this section's rule name instead of restating it.}

\begin{brecipe}[the satisfaction-level alias template]\label{rec:sd-template}
An alias rule $A$ at the satisfaction stratum is warranted by a $1 \to 1$
unpack into a single application of a target rule $T$ (or, for the
spec-surface rewrites, into a single definitional denotation equality):
the $A$-node and its image conclude the same judgment from the same
premise data, payload transported, nothing derived. Soundness is then a
corollary, $\mathrm{sound}_A = \mathrm{unpack} \fatsemi \mathrm{sound}_T$,
and the certificate format need never carry $A$: producers elaborate
before emission. The five artefacts an alias owes this book are the same
as at the base stratum (the \rulename{SmallGain} section is the model):
the unpack, the filled-in target application, the payload-transport
table, the exactness identity, and the soundness corollary. One honest
divergence from the base-stratum model, stated once here: no alias of
this batch has a \emph{verified payload function} among this
edition\textquotesingle s bound artefacts. The
transport tables of this batch are therefore stated at the
catalogued-field level, and the unpack-function slot of each alias
says so plainly; where a discrete-profile
\emph{semantic} witness exists in the auxiliary Lean development it is
named (the soundness content is machine-checked even though the
payload-level function is not), with its statement pin
riding the bindings record wherever the per-section text says
hash-bound, and where none exists the status is stated as filed. In
every case it is the payload-level slot, never the semantic witness,
that remains outside the bound warrant; the per-section text says which
case holds.
\end{brecipe}

\begin{brecipe}[the filled-in application --- \rulename{QStreettDegenerate}
rewritten into its \rulename{QStreettBattery} instance]\label{rec:sd-filled}
A degenerate node claims, from the confinement certificate $V_0$ and
per-region ranking certificates $(W_i, \varepsilon_i)$, that with
probability $\ge p$ every response region $G_i$ is visited infinitely
often. Its image is the single bundle node at request regions
$R_i := S$ (every state requests):
\begin{mathpar}
\inferrule*[Right=\rulename{QStreettBattery}]{
  \text{confinement rows for } V_0 \\
  \bigl\{\vdash \text{nonincrease of } W_i \ \text{off } G_i\bigr\} \\
  \bigl\{\vdash \varepsilon_i\text{-decay of } W_i \ \text{off } G_i\bigr\}
  \ \text{(the request row at } R_i = S\text{)} \\\\
  p + y \le 1 \ \text{(threshold tie)}
}{
  \mathcal{M},\ x_0 \in \{V_0 \le y\} \vDash
  \mathsf{SProbGE}\bigl(p,\ \mathsf{PDwellStreett}(\mathsf{PairBase}(\mathsf{STrue},
  \varphi_1), \ldots, \mathsf{PairBase}(\mathsf{STrue}, \varphi_k))\bigr)
}
\end{mathpar}
Discharging the bundle's obligations at this instance: the request
regions are measurable trivially; the decay row \emph{subsumes} the
nonincrease row (off $G_i$, every state is a request state, so
nonincrease follows from decay with the kit\textquotesingle s
$\varepsilon_i > 0$ --- nonnegativity is all the subsumption spends); the
absorbing families $(V_i, \gamma_i)$ and carved sets $J_i$ are vacuous
--- each $J_i$ lives inside $I \smallsetminus [\![\rho_i]\!]$, with $I$
the confinement region of the bundle\textquotesingle s carved-scope
formula, and at $\rho_i = \top$ that base is
$I \smallsetminus S = \varnothing$, so
$J_i = \varnothing$ syntactically, the rows carved to $J_i$ and to
off-request regions are rows over the empty region, and the
specialization targets the bundle's conservative (absorber-free)
scope --- the scope its ranking-route soundness theorem consumes ---
so no request-region absorber obligation survives. The dead premise slots are honestly dead, and the
degenerate-pair specialization pass erases them before emission.
\end{brecipe}

\begin{brecipe}[payload transport, catalogued-field level]\label{rec:sd-transport}
\begin{center}\small
\begin{tabular}{@{}ll@{}}
\textbf{degenerate node} & \textbf{\rulename{QStreettBattery} node}\\
\midrule
$V_0$ rows, threshold, init & verbatim\\
response regions $G_i$, certificates $(W_i, \varepsilon_i)$ & verbatim (bound rows $0 \le W_i \le C_i$ and strict $\varepsilon_i > 0$ included)\\
shared side conditions (hull, fail-closed) + measurability pack & verbatim (the battery pattern\textquotesingle s)\\
--- & request regions $:= S$ (each pair; measurable trivially)\\
--- & absorbing families: empty; $J_i = \varnothing$ (syntactic)\\
decay rows off $G_i$ & request-decay rows (at $R_i = S$); nonincrease subsumed\\
\end{tabular}
\end{center}
Unpack-function slot: no verified payload function in this edition
(Recipe~\ref{rec:sd-template}); the semantic witness exists ---
\texttt{sound\_QStreettDegenerate}, one \texttt{sound\_QStreettBattery}
application at request set \texttt{univ} (auxiliary module; hash-bound
in this edition).
\end{brecipe}

\begin{bprop}[the exactness identity]\label{prop:sd-exact}
For every path $\omega$ and pair $i$, at $R_i = S$:
\[
\bigl(\exists^\infty n,\ \omega_n \in S\bigr) \Rightarrow
\bigl(\exists^\infty n,\ \omega_n \in G_i\bigr)
\quad\Longleftrightarrow\quad
\exists^\infty n,\ \omega_n \in G_i .
\]
\end{bprop}
\begin{proof}
The hypothesis $\exists^\infty n,\ \omega_n \in S$ holds of every path
(every time is a request time), so the implication and its conclusion are
equivalent --- one line, which is the point. Equivalently, in acceptance
vocabulary: the finitely-many-requests disjunct
$\mathrm{Fin}(S)$ is empty along the tail filter, so the pair condition
$\mathrm{Fin}(R_i) \vee \mathrm{Inf}(G_i)$ collapses to
$\mathrm{Inf}(G_i)$.
\end{proof}

\begin{bprop}[soundness corollary]\label{prop:sd-sound}
Under the degenerate node's rows,
$\mu_x\{\omega \mid \forall i,\ \exists^\infty n,\ \omega_n \in G_i\}
\ge p$.
\end{bprop}
\begin{proof}
Apply the bundle's soundness theorem (the \rulename{QStreettBattery}
section's --- at its proven scope, the discrete-kernel profile and
conservative ranking scope; the bundle\textquotesingle s
martingale-problem discharge transfers to this alias the same way,
one application at the trivially-true request) at $R_i = S$ and
rewrite the event by
Proposition~\ref{prop:sd-exact}. There is no other content --- the alias
genre working as intended. (Machine-checked as
\texttt{sound\_QStreettDegenerate}; hash-bound in this edition.)
\end{proof}

\begin{bremark}[this is the B\"uchi case]
A B\"uchi condition \emph{is} the one-pair Streett condition with
trivially-true request, so this alias is where the classical
recurrence certificate lives inside the Streett-primitive design.
Deterministic B\"uchi automata recognize only limit languages
\cite{landweber1969}, so the reverse direction --- expressing a general
Streett condition B\"uchi-primitively --- is not a respelling at all but
an automaton re-encoding, exponential in the number of pairs for the
classical translations; the asymmetry is why the bundle,
not the recurrence rule, is the primitive (\cite{streett1982,thomas1990};
see the bundle section's remark).
\end{bremark}

\begin{anchors}
Recipe \ref{rec:sd-filled} target & the \rulename{QStreettBattery} section, Theorem \& hypotheses (H1)--(H5)\\
Proposition \ref{prop:sd-exact} & syntactic; no formal anchor needed (one line of filter algebra)\\
Proposition \ref{prop:sd-sound} & \texttt{sound\_QStreettDegenerate} (\texttt{BernsteinLogicAux.QStreett}; machine-checked, hash-bound in this edition)\\
no wire tag, by design & producers elaborate before emission (Recipe~\ref{rec:sd-template})\\
\end{anchors}


\rulesection{51}{QStreettCmplRequest}{admissible rule}
\label{blcatrule-QStreettCmplRequest}
\ruleblock{The rule (Steele metanotation)}
\sourcenote{as realized in the mechanized development; source row: \texttt{\texttt{QTSL-rules.toml}} --- \emph{QStreettCmplRequest}. The metanotation statement is typeset by hand for notational quality; its mechanized ground truth is the Lean statement below, extracted from the sources at every build of this book and hash-bound. The rule \emph{roster} is machine-checked (every mechanized rule tag round-trips to exactly one catalogue row); the statement \emph{text} is editorial, checked by eye against the extracted Lean beside it.}
\begin{mathpar}
\inferrule*{
  \text{a \rulename{QStreettBattery} kit at certified regions } (R_i, G_i) \\
  \bigl\{ [\![\mathrm{cmpl}\,\rho_i]\!] \subseteq R_i \bigr\} \ \text{(requests over-approximated)} \\
  \bigl\{ G_i \subseteq [\![\varphi_i]\!] \bigr\} \ \text{(responses under-approximated)}
}{
  \mathcal{M},\ x_0 \in \{V_0 \le y\} \vDash \mathsf{SProbGE}\bigl(p,\ \mathsf{PDwellStreett}(\mathsf{PairCmpl}(\rho_1, \varphi_1), \ldots)\bigr)
}
\end{mathpar}

\ruleflavor{requests sit on the hypothesis side of each pair implication, so the orientation swap of the complement (the workhorse judgment certifies under-approximations; complements arrive as over-approximations) is exactly the sound direction, for free.}
\par\medskip
\rulename{QStreettCmplRequest} lets a Streett pair carry its request as a
\emph{complement}: the pair $(\mathrm{cmpl}\,\rho_i, \varphi_i)$ asks
that if states \emph{outside} $[\![\rho_i]\!]$ recur forever, responses
recur forever. The satisfaction judgment\textquotesingle s workhorse
orientation certifies regions that \emph{under}-approximate a
formula\textquotesingle s denotation --- one orientation, deliberately:
the complement of an under-approximation is an
\emph{over}-approximation, so a complement-typed request arrives with
the opposite orientation for free. The rule (n\'e StreettCmplRequest) is \emph{admissible}, not alias: its unpacking is one
\rulename{QStreettBattery} node \emph{plus} region-weakening premises ---
the certified request regions over-approximate the complement
denotations, the certified response regions under-approximate the
response denotations --- and the direction check that makes this the
sound weakening is genuine content, not transport. Requests sit on the
hypothesis side of each pair implication and responses on the conclusion
side, so over-approximating requests and under-approximating responses
each \emph{weaken the claimed event}: any path satisfying the certified
pair satisfies the semantic pair. What the checker recomputes: nothing
new in kind --- the weakening premises are ordinary region-containment
rows it already knows how to verify, and
the complement itself is consumed at the region-atom level (the atom
dualities of the \rulename{QNnfRewrite} section), never as a formula
node.

Scope, stated plainly (this entry is marked \emph{conditional}): the
bound statements below are the discrete-profile witnesses --- the
event-level direction lemma \texttt{streett\_event\_mono} (two
monotonicity steps, exactly the two lines the recipe below typesets) and
the assembled bundle-through-formula-tie
\texttt{sound\_QStreettBattery\_tie} (requests over-approximated,
responses under-approximated, concluded over the semantic satisfaction
predicates), both in module \texttt{BernsteinLogicAux.QStreett}, plus
the atom-level complement step \texttt{denot\_elimNeg} in
\texttt{BernsteinLogicAux.NNFDuality}. One named residue keeps the
flag --- and it is this entry\textquotesingle s own, not the
bundle\textquotesingle s (whose continuous-time residue is discharged
whole at the martingale-problem profile, strict rows included,
hash-bound under the \rulename{QStreettBattery} entry): the
continuous-time face of the complement step remains prose. At the
general-stochastic-hybrid profile, consuming a complemented
request requires fine-topology regularity of the region (that the
region is nearly Borel and contains its regular points), which is
argued in the design record.

\cataloguenote{Why one orientation of the workhorse judgment suffices, in one line: an
under-approximating region calculus is closed under the connectives the
rules consume, and where a complement appears, the swap to
over-approximation is exactly what the hypothesis side of an implication
wants --- demanding a second, over-approximating judgment kind would
duplicate every rule for zero content.
}
\ruleblock{Lean statement}
\declhead{BernsteinLogicAux.QStreett}{streett\_event\_mono}{\badge{BadgeGreen}{statement bound}}
\begin{leansrc}
/-- **`QStreettCmplRequest` (T9b), the event-level kernel** — the discrete face
of request/response event monotonicity: requests **over**-approximated and
responses **under**-approximated weaken the per-pair Streett event in the
sound direction. Two `Frequently.mono`s — the "two lines" the catalogue prices. -/
theorem streett_event_mono {Rq Rqhat G Ghat : Set S} (hR : Rq ⊆ Rqhat)
    (hG : Ghat ⊆ G) (ω : ℕ → S)
    (h : (∃ᶠ n in atTop, ω n ∈ Rqhat) → ∃ᶠ n in atTop, ω n ∈ Ghat) :
    (∃ᶠ n in atTop, ω n ∈ Rq) → ∃ᶠ n in atTop, ω n ∈ G
\end{leansrc}
\declhead{BernsteinLogicAux.QStreett}{sound\_QStreettBattery\_tie}{\badge{BadgeGreen}{statement bound}}
\begin{leansrc}
/-- **`QStreettCmplRequest` (T9b), assembled** — the bundle consumed through the
formula tie, direction-checked by request/response event monotonicity and the
antitone region-complement kernel: the certified request region
`Rq i` over-approximates the truth set of the semantic request `reqSat i` (for
`cmpl`-typed requests this is atom-level complement elimination's
`⟦cmpl φ⟧ ⊇ ⟦φ⟧ᶜ = ⟦¬φ⟧`), and the certified response region
under-approximates the semantic response. Conclusion over the semantic
satisfaction predicates; no measurability is asked of them. -/
theorem sound_QStreettBattery_tie {Rq G : ι → Set S} {W : ι → S → ℝ} {ε C : ι → ℝ}
    {reqSat respSat : ι → S → Prop}
    (hD : MeasurableSet D) (hDbar : MeasurableSet Dbar) (hU : MeasurableSet U)
    (hV₀ : Measurable V₀)
    (hRq : ∀ i, MeasurableSet (Rq i)) (hG : ∀ i, MeasurableSet (G i))
    (hW : ∀ i, Measurable (W i))
    (hreq : ∀ i y, reqSat i y → y ∈ Rq i)
    (hresp : ∀ i, ∀ y ∈ G i, respSat i y)
    (hΘ : p + thr ≤ 1)
    (hDsub : D ⊆ Dbar)
    (hR3 : ∀ y ∈ D, y ∉ U → K y Dbar = 1)
    (hR4 : ∀ y ∈ D, y ∉ U → K y (D ∪ U) = 1)
    (hV₀int : ∀ y ∈ D, y ∉ U → Integrable V₀ (K y))
    (hV₀drift : ∀ y ∈ D, y ∉ U → ∫ z, V₀ z ∂K y ≤ V₀ y)
    (hV₀floor : ∀ y ∈ U, y ∈ Dbar → 1 ≤ V₀ y)
    (hV₀nonneg : ∀ y ∈ Dbar, 0 ≤ V₀ y)
    (hε : ∀ i, 0 < ε i)
    (hW0 : ∀ i, ∀ y ∈ Dbar, 0 ≤ W i y) (hWC : ∀ i, ∀ y ∈ Dbar, W i y ≤ C i)
    (hWint : ∀ i, ∀ y ∈ D, y ∉ U → y ∉ G i → Integrable (W i) (K y))
    (hWdrift : ∀ i, ∀ y ∈ D, y ∉ U → y ∉ G i → ∫ z, W i z ∂K y ≤ W i y)
    (hWdriftR : ∀ i, ∀ y ∈ D, y ∉ U → y ∈ Rq i → y ∉ G i →
      ∫ z, W i z ∂K y ≤ W i y - ε i)
    (hx : x ∈ D) (hthr : V₀ x ≤ thr) :
    ENNReal.ofReal p ≤ pathMeasure K x
      {ω | ∀ i, (∃ᶠ n in atTop, reqSat i (ω n)) → ∃ᶠ n in atTop, respSat i (ω n)}
\end{leansrc}
\declhead{BernsteinLogicAux.NNFDuality}{denot\_elimNeg}{\badge{BadgeGreen}{statement bound}}
\begin{leansrc}
/-- **Denotation preservation of dual-atom elimination**: where each dual atom
denotes the complement (`hd` — the per-duality definitional axioms, e.g. the
`compl_le_atom` family above), elimination preserves the region. This is the
lawfulness of consuming a `cmpl`-typed request through swapped atoms
(atom-level complement elimination's atom step, the T9b input shape). -/
theorem denot_elimNeg {σ : Type*} (den : α → Set σ) (d : α → α)
    (hd : ∀ a, den (d a) = (den a)ᶜ) :
    ∀ φ : NForm α, (φ.elimNeg d).denot den = φ.denot den
\end{leansrc}
\par\medskip
\textbf{The general-stochastic instance, carved requests
discharged.}\enspace The complemented-request form at the segment-model
tier: the carved variants below instantiate the bundle on the skeleton
chain with the complement-carved request regions supplied and their
side conditions discharged --- plain and confined forms.

\declhead{BernsteinLogicAux.LayerD.Streett}{sound\_GSHS\_streett\_bundle\_carved\_discharged}{\badge{BadgeGreen}{statement bound}}
\begin{leansrc}
/-- **The full bundle datum, GSHS instance** (T9, carved scope, absorbing
certificates as segment-level rows) — the discrete-statement-shaped
corollary of the `_confined` form; matches
`QStreett.sound_QStreettBattery_carved_discharged` line for line. -/
theorem sound_GSHS_streett_bundle_carved_discharged
    [IsMarkovKernel M.seg] [IsMarkovKernel M.jump]
    (hD : MeasurableSet D) (hDbar : MeasurableSet Dbar) (hU : MeasurableSet U)
    (hV₀ : Measurable V₀)
    (hRq : ∀ i, MeasurableSet (Rq i)) (hG : ∀ i, MeasurableSet (G i))
    (hJ : ∀ i, MeasurableSet (J i)) (hW : ∀ i, Measurable (W i))
    (hΘ : p + thr ≤ 1) (hDsub : D ⊆ Dbar)
    (hLive : NoVanish M (D \ U)) (hHalt : HaltsCaught M (D \ U) (U ∩ Dbar))
    (hReset : ResetsInto M (D \ U) ((D ∪ U) ∩ Dbar))
    (hSegD₀ : RankingSegmentDynkin M (D \ U) V₀ 0)
    (hForce₀ : ForcedNonexpansion M V₀) (hNonneg₀ : NonnegV V₀)
    (hC₀0 : 0 ≤ C₀) (hV₀C : ∀ y, V₀ y ≤ C₀)
    (hFloor₀ : ∀ y ∈ U, y ∈ Dbar → 1 ≤ V₀ y)
    (hε : ∀ i, 0 < ε i) (hγ : ∀ i, 0 < γ i)
    (hSegDW : ∀ i, RankingSegmentDynkin M ((D \ U) \ (G i ∪ J i)) (W i) 0)
    (hSegDWR : ∀ i,
      RankingSegmentDynkin M (((D \ U) ∩ Rq i) \ (G i ∪ J i)) (W i) (ε i))
    (hForceW : ∀ i, ForcedNonexpansion M (W i)) (hNonnegW : ∀ i, NonnegV (W i))
    (hC0 : ∀ i, 0 ≤ C i) (hWC : ∀ i, ∀ y, W i y ≤ C i)
    (hVa : ∀ i, Measurable (Va i)) (hNonnegVa : ∀ i, NonnegV (Va i))
    (hCa0 : ∀ i, 0 ≤ Ca i) (hVaC : ∀ i, ∀ y, Va i y ≤ Ca i)
    (hVaFloor : ∀ i, ∀ y ∈ Rq i, 1 ≤ Va i y)
    (hSegDVa : ∀ i, RankingSegmentDynkin M ((Rq i)ᶜ) (Va i) 0)
    (hForceVa : ∀ i, ForcedNonexpansion M (Va i))
    (hVaJ : ∀ i, ∀ y ∈ J i, Va i y ≤ 1 - γ i)
    (hx : x ∈ D) (hthr : V₀ x ≤ thr) :
    ENNReal.ofReal p ≤ Battery.pathMeasure (skel M) (SkelState.alive x)
      {ω | ∀ i, (∃ᶠ n in atTop, ω n ∈ skelIn (Rq i)) →
        ∃ᶠ n in atTop, ω n ∈ skelIn (G i)}
\end{leansrc}
\declhead{BernsteinLogicAux.LayerD.Streett}{sound\_GSHS\_streett\_bundle\_carved\_discharged\_confined}{\badge{BadgeGreen}{statement bound}}
\begin{leansrc}
/-- **The full bundle datum, GSHS instance, confined form** (T9 with the
absorbing certificates; assault doc §2.7 Theorem G proper, at the
hypothesis-carrying ranking scope `I \ (⟦φᵢ⟧ ∪ Jᵢ)`): the ranking rows are
demanded only off the absorbed region `J i`, and in exchange each pair
carries an absorbing certificate `Va i` — nonnegative, bounded, `≥ 1` on
the request region, segment-nonincreasing off it, `≤ 1 − γ i` on `J i`.
The absorption bound is discharged on the skeleton by
`QStreett.absorptionBound_of_kernel` (the shifted supermartingale bridge +
conditional Ville — the mechanized replacement for prose-G2's epoch
argument); its off-request drift row is demanded at **every alive state**
(scope `(Rq i)ᶜ`), the global-sharp reading of the discrete discharge —
frozen states satisfy it for free. -/
theorem sound_GSHS_streett_bundle_carved_discharged_confined
    [IsMarkovKernel M.seg] [IsMarkovKernel M.jump]
    (hD : MeasurableSet D) (hDbar : MeasurableSet Dbar) (hU : MeasurableSet U)
    (hV₀ : Measurable V₀)
    (hRq : ∀ i, MeasurableSet (Rq i)) (hG : ∀ i, MeasurableSet (G i))
    (hJ : ∀ i, MeasurableSet (J i)) (hW : ∀ i, Measurable (W i))
    (hΘ : p + thr ≤ 1) (hDsub : D ⊆ Dbar)
    (hLive : NoVanish M (D \ U)) (hHalt : HaltsCaught M (D \ U) (U ∩ Dbar))
    (hReset : ResetsInto M (D \ U) ((D ∪ U) ∩ Dbar))
    (hSegD₀ : RankingSegmentDynkin M (D \ U) V₀ 0)
    (hForce₀ : ForcedNonexpansion M V₀) (hNonneg₀ : NonnegV V₀)
    (hC₀0 : 0 ≤ C₀) (hV₀C : ∀ y, V₀ y ≤ C₀)
    (hFloor₀ : ∀ y ∈ U, y ∈ Dbar → 1 ≤ V₀ y)
    (hε : ∀ i, 0 < ε i) (hγ : ∀ i, 0 < γ i)
    (hSegDW : ∀ i, RankingSegmentDynkin M ((D \ U) \ (G i ∪ J i)) (W i) 0)
    (hSegDWR : ∀ i,
      RankingSegmentDynkin M (((D \ U) ∩ Rq i) \ (G i ∪ J i)) (W i) (ε i))
    (hForceW : ∀ i, ForcedNonexpansion M (W i)) (hNonnegW : ∀ i, NonnegV (W i))
    (hC0 : ∀ i, 0 ≤ C i) (hWC : ∀ i, ∀ y, W i y ≤ C i)
    (hVa : ∀ i, Measurable (Va i)) (hNonnegVa : ∀ i, NonnegV (Va i))
    (hCa0 : ∀ i, 0 ≤ Ca i) (hVaC : ∀ i, ∀ y, Va i y ≤ Ca i)
    (hVaFloor : ∀ i, ∀ y ∈ Rq i, 1 ≤ Va i y)
    (hSegDVa : ∀ i, RankingSegmentDynkin M ((Rq i)ᶜ) (Va i) 0)
    (hForceVa : ∀ i, ForcedNonexpansion M (Va i))
    (hVaJ : ∀ i, ∀ y ∈ J i, Va i y ≤ 1 - γ i)
    (hx : x ∈ D) (hthr : V₀ x ≤ thr) :
    ENNReal.ofReal p ≤ Battery.pathMeasure (skel M) (SkelState.alive x)
      {ω | (∀ n, ω n ∈ skelIn D \ skelBad U) ∧
        ∀ i, (∃ᶠ n in atTop, ω n ∈ skelIn (Rq i)) →
          ∃ᶠ n in atTop, ω n ∈ skelIn (G i)}
\end{leansrc}
\ruleblock{Correspondence (Steele display $\leftrightarrow$ Lean)}
\badge{BadgeGreen}{reviewer-certified against the displayed texts}
\sourcenote{The dictionary below is a reviewer-certified statement of how the mechanized statements above interpret the displayed rule: display metavariable $\mapsto$ Lean binder, each display hypothesis $\mapsto$ the premise or side condition carrying it, the display conclusion $\mapsto$ the goal --- with any deliberate difference in width between display and mechanization named outright. Certification is recorded as a hash triple over (the displayed cell, the extracted statement texts, this dictionary\textquotesingle s own text); if any of the three changes, this block turns amber until a reviewer re-certifies. The correspondence claim itself remains editorial --- reviewer-grade, not machine-checked; its machine-checked successor types each soundness statement against the rule as a structured object, and supersedes this block per rule as it lands.}
Three bound statements, three display constituents. The
orientation-swap kernel --- the flavor text\textquotesingle s
``exactly the sound direction, for free'' --- $\mapsto$
\texttt{streett\_event\_mono}: \texttt{Rq ⊆ Rqhat} (requests
over-approximated) and \texttt{Ghat ⊆ G} (responses
under-approximated) transport the per-pair Streett implication
pathwise; two \texttt{Frequently.mono}s, the priced two lines,
axiom-light and \emph{deterministic} --- no measure appears.
\begin{itemize}
\item the assembled rule $\mapsto$
  \texttt{sound\_QStreettBattery\_tie}: the full bundle kit (the
  battery-pattern binder block, as on the degenerate-request page)
  plus the two tie hypotheses, and the polarity pair is carried
  exactly --- \texttt{hreq : ∀ i y, reqSat i y → y ∈ Rq i} (the
  certified region \emph{over}-approximates the semantic request
  truth set: for $\mathrm{cmpl}$-typed requests this is the
  atom-level complement-elimination inclusion
  $[\![\mathrm{cmpl}\,\varphi]\!] \supseteq
  [\![\varphi]\!]^{\mathsf{c}}$) and \texttt{hresp : ∀ i, ∀ y ∈ G i,
  respSat i y} (responses \emph{under}-approximated). The conclusion
  event quantifies over the semantic predicates
  \texttt{reqSat}/\texttt{respSat}, of which \emph{no measurability
  is asked} (the measurable objects are the certified regions ---
  the QSafetyBattery \texttt{sat} pattern), and both tie hypotheses
  are safe-direction widenings in the QClosure polarity sense.
\item the atom-level complement step $\mapsto$ \texttt{denot\_elimNeg}
  (the NNF module): where each dual atom denotes the complement,
  consuming a $\mathrm{cmpl}$-typed request through swapped atoms
  preserves the region --- the lawfulness clause cited at
  QStreettRequestNnf\textquotesingle s consumption boundary (its
  flavor text).
\end{itemize}
Width note: the conditional mark\textquotesingle s one remaining
residue, exactly as the roster carries it --- the fine-topology
content of the complement step at the
general-stochastic-hybrid profile (nearly-Borel regions,
$E^r \subseteq E$), which remains design-note prose with no bound
counterpart; the bundle side\textquotesingle s former continuous-time
residue is discharged and hash-bound under the bundle entry. \texttt{streett\_event\_mono} is
itself profile-free (pure path logic); the profile enters only with
the assembled tie.

\ruleblock{Lean verifier receipts}
\sourcenote{module \texttt{BernsteinLogicAux.QStreett}; bound build \texttt{469881beb770\ldots} (bound from receipt 9d1d848e0ba1718a, 2026-07-29T01:03:31Z). A receipt records that named checking layers accepted a build: (module, compiled-artifact SHA-256, environment SHA-256, checker layers), minted when the checking gate ran; the store receipts attesting the bound build are listed with the layers each run credited. Nothing here asks the reader to trust the store: every layer is an independently runnable program over the quoted sources --- rebuild the module and re-run the checkers to reproduce the verdict. Receipts attesting superseded builds credit no layer and are omitted from this rendering (the count of omissions, if any, is noted below).}
\begin{receipttable}
\receiptrow{e3990fe544d4708d}{2026-07-28T23:48:35Z}{build+decorrelated}{\badge{BadgeGreen}{current}}
\receiptrow{68e9374fa1b418fe}{2026-07-29T00:00:34Z}{build+decorrelated}{\badge{BadgeGreen}{current}}
\receiptrow{e5895c9706854cd3}{2026-07-29T00:49:04Z}{build+replay+hardened+decorrelated}{\badge{BadgeGreen}{current}}
\receiptrow{9d1d848e0ba1718a}{2026-07-29T01:03:31Z}{build+replay+hardened+decorrelated}{\badge{BadgeGreen}{current}}
\receiptrow{4a8d377b721b8523}{2026-07-31T22:00:09Z}{build+decorrelated}{\badge{BadgeGreen}{current}}
\receiptrow{d251a122540ab741}{2026-07-31T22:29:03Z}{build+decorrelated}{\badge{BadgeGreen}{current}}
\receiptrow{333635b8e047113a}{2026-08-02T23:31:14Z}{build+decorrelated}{\badge{BadgeGreen}{current}}
\receiptrow{9e96f80d63e4ad82}{2026-08-03T16:27:18Z}{build+decorrelated}{\badge{BadgeGreen}{current}}
\receiptrow{932b4b9e0b0fd323}{2026-08-03T22:46:02Z}{build+decorrelated}{\badge{BadgeGreen}{current}}
\end{receipttable}
\sourcenote{9 further receipt(s) in the store attest superseded builds of this module and are omitted from this rendering.}
\sourcenote{module \texttt{BernsteinLogicAux.NNFDuality}; bound build \texttt{97564f8e3950\ldots} (bound from receipt 21f04961748c0302, 2026-07-29T01:03:33Z). A receipt records that named checking layers accepted a build: (module, compiled-artifact SHA-256, environment SHA-256, checker layers), minted when the checking gate ran; the store receipts attesting the bound build are listed with the layers each run credited. Nothing here asks the reader to trust the store: every layer is an independently runnable program over the quoted sources --- rebuild the module and re-run the checkers to reproduce the verdict. Receipts attesting superseded builds credit no layer and are omitted from this rendering (the count of omissions, if any, is noted below).}
\begin{receipttable}
\receiptrow{1417cdf27f69842a}{2026-07-28T23:48:35Z}{build+decorrelated}{\badge{BadgeGreen}{current}}
\receiptrow{65fc1b51055d38bd}{2026-07-29T00:00:33Z}{build+decorrelated}{\badge{BadgeGreen}{current}}
\receiptrow{2abbbc61b6065811}{2026-07-29T00:49:05Z}{build+replay+hardened+decorrelated}{\badge{BadgeGreen}{current}}
\receiptrow{21f04961748c0302}{2026-07-29T01:03:33Z}{build+replay+hardened+decorrelated}{\badge{BadgeGreen}{current}}
\receiptrow{8879c5b9721acf0c}{2026-07-31T22:00:03Z}{build+decorrelated}{\badge{BadgeGreen}{current}}
\receiptrow{e47a2c2d7f152711}{2026-07-31T22:28:57Z}{build+decorrelated}{\badge{BadgeGreen}{current}}
\receiptrow{ce1bb97f7cbcaf4f}{2026-08-02T23:31:07Z}{build+decorrelated}{\badge{BadgeGreen}{current}}
\receiptrow{8755b28f77990ba0}{2026-08-03T16:27:11Z}{build+decorrelated}{\badge{BadgeGreen}{current}}
\receiptrow{64948e60a56ae3bc}{2026-08-03T22:45:55Z}{build+decorrelated}{\badge{BadgeGreen}{current}}
\end{receipttable}
\sourcenote{8 further receipt(s) in the store attest superseded builds of this module and are omitted from this rendering.}
\sourcenote{module \texttt{BernsteinLogicAux.LayerD.Streett}; bound build \texttt{7393065bbc7f\ldots} (bound from receipt abac471866cfcfd1, 2026-07-21T05:38:48Z). A receipt records that named checking layers accepted a build: (module, compiled-artifact SHA-256, environment SHA-256, checker layers), minted when the checking gate ran; the store receipts attesting the bound build are listed with the layers each run credited. Nothing here asks the reader to trust the store: every layer is an independently runnable program over the quoted sources --- rebuild the module and re-run the checkers to reproduce the verdict. Receipts attesting superseded builds credit no layer and are omitted from this rendering (the count of omissions, if any, is noted below).}
\begin{receipttable}
\receiptrow{0f3b50b4d3f0d33c}{2026-07-20T22:52:50Z}{build+decorrelated}{\badge{BadgeGreen}{current}}
\receiptrow{109669f848453ec2}{2026-07-21T00:03:19Z}{build+decorrelated}{\badge{BadgeGreen}{current}}
\receiptrow{7e17aa32b276623b}{2026-07-21T04:51:49Z}{build+decorrelated}{\badge{BadgeGreen}{current}}
\receiptrow{abac471866cfcfd1}{2026-07-21T05:38:48Z}{build+replay+hardened+decorrelated}{\badge{BadgeGreen}{current}}
\end{receipttable}
\sourcenote{6 further receipt(s) in the store attest superseded builds of this module and are omitted from this rendering.}
\layerline{\layeritem{build}{lake build (elaboration)}{\badge{BadgeGreen}{current}}\layeritem{replay}{leanchecker (official kernel replay)}{\badge{BadgeGreen}{current}}\layeritem{hardened}{safe\_verify (hardened screen)}{\badge{BadgeGreen}{current}}\layeritem{decorrelated}{lean4export + nanoda (decorrelated kernel)}{\badge{BadgeGreen}{current}}}
\sourcenote{\badge{BadgeGray}{comparator: none} no fixture in the three-implementation comparator harness (compiled / kernel-interpreted / reference-interpreted) exercises this rule.}
\ruleblock{HOL4 statement}
\absentbox{NO HOL STATEMENT YET}{Absent, precisely located: the rule\textquotesingle s realization rides
the typed request field of the Streett bundle --- and the HOL4
development (\texttt{bern\_catalogue}, the base-calculus checker cone)
carries nothing at the satisfaction stratum, so there is no HOL-side
wire object for a checker to consume. A HOL statement becomes
meaningful only after the extension stratum is ported; the weakening rows
themselves are ordinary containment sequents the base checker already
covers. Checked against the HOL sources this edition.
}
\ruleblock{HOL verifier receipts}
\absentbox{NONE}{no HOL verification receipt exists for this rule.}
\ruleblock{Proof}

\sourcenote{Genre: admissible rule (tier 2) --- a proof-tree-rewriting
recipe (the general admissible template is stated in the
\rulename{PConstExact} section; the batch-level honesty clause about
absent payload functions is Recipe~\ref{rec:sd-template} of the
\rulename{QStreettDegenerate}
section and applies here verbatim).}

\begin{brecipe}[the \rulename{QStreettCmplRequest} instance --- one node,
one bundle, per-pair weakening premises]\label{rec:scr-instance}
A node concluding
$\mathcal{M},\ x_0 \in \{V_0 \le y\} \vDash
\mathsf{SProbGE}\bigl(p,\ \mathsf{PDwellStreett}(\mathsf{PairCmpl}(\rho_1,
\varphi_1), \ldots)\bigr)$ rewrites to the derivation fragment
\begin{mathpar}
\inferrule*[Right=\rulename{QStreettBattery}]{
  \text{the bundle kit at certified regions } (R_i, G_i) \\\\
  \bigl\{\, [\![\mathrm{cmpl}\,\rho_i]\!] \subseteq R_i \,\bigr\}_i
  \ \text{(request weakening)} \\
  \bigl\{\, G_i \subseteq [\![\varphi_i]\!] \,\bigr\}_i
  \ \text{(response weakening)}
}{
  \mathcal{M},\ x_0 \in \{V_0 \le y\} \vDash
  \mathsf{SProbGE}\bigl(p,\ \mathsf{PDwellStreett}(\mathsf{PairCmpl}(\rho_1,
  \varphi_1), \ldots)\bigr)
}
\end{mathpar}
where the request-weakening premises are supplied for free by the
orientation swap: a certified under-approximation $A_i \subseteq
[\![\rho_i]\!]$ yields the over-approximation $A_i^{\mathrm{c}} \supseteq
[\![\rho_i]\!]^{\mathrm{c}}$, and $R_i := A_i^{\mathrm{c}}$ is a region
atom by the atom dualities (\rulename{QNnfRewrite} section) wherever
$A_i$\textquotesingle s atom class carries a dual --- the
ordered-comparison classes, which are the shapes normalized request
bodies produce; an over-approximation of any other shape is simply
supplied as its own certified region --- and the
complement is consumed at the atom level, never as a formula node. The
weakening premises are ordinary containment rows the checker already
knows how to verify; nothing new \emph{in kind} is
recomputed.
\end{brecipe}

\begin{bprop}[the direction lemma --- two lines]\label{prop:scr-mono}
Fix a path $\omega$ and one pair. If the certified regions satisfy
$\mathrm{Rq} \subseteq \widehat{\mathrm{Rq}}$ (requests
over-approximated when read semantic-to-certified) and
$\widehat{G} \subseteq G$ (responses under-approximated), then
\[
\Bigl[\bigl(\exists^\infty n,\ \omega_n \in \widehat{\mathrm{Rq}}\bigr)
\Rightarrow \bigl(\exists^\infty n,\ \omega_n \in \widehat{G}\bigr)\Bigr]
\;\Longrightarrow\;
\Bigl[\bigl(\exists^\infty n,\ \omega_n \in \mathrm{Rq}\bigr)
\Rightarrow \bigl(\exists^\infty n,\ \omega_n \in G\bigr)\Bigr] .
\]
\end{bprop}
\begin{proof}
Frequent membership is monotone in the set: frequent
$\mathrm{Rq}$-visits are frequent $\widehat{\mathrm{Rq}}$-visits
(hypothesis side, first monotonicity), and frequent $\widehat{G}$-visits
are frequent $G$-visits (conclusion side, second monotonicity). Two
applications of monotonicity --- which is the entire lemma, and why the
weakening premises are the only new obligations of the rewrite.
(Machine-checked as \texttt{streett\_event\_mono}; hash-bound in this
edition.)
\end{proof}

\begin{bprop}[semantic warrant of the rewrite]\label{prop:scr-sound}
Under the bundle kit's rows and the weakening premises of
Recipe~\ref{rec:scr-instance}, the concluded judgment holds: with
probability $\ge p$, every semantic pair implication
$(\exists^\infty\ \text{states outside } [\![\rho_i]\!]) \Rightarrow
(\exists^\infty\ \text{states in } [\![\varphi_i]\!])$ is satisfied.
\end{bprop}
\begin{proof}
The bundle theorem gives the certified-region Streett event probability
$\ge p$; Proposition~\ref{prop:scr-mono}, applied per pair with
$\mathrm{Rq} := [\![\mathrm{cmpl}\,\rho_i]\!] \subseteq R_i$ and
$G_i \subseteq [\![\varphi_i]\!]$, includes that event in the semantic
one; the certified event is measurable (its regions are the
kit\textquotesingle s), and the path measure, read as outer measure on
the possibly-nonmeasurable semantic event, is monotone --- the lower
bound transfers. (Machine-checked, assembled, as
\texttt{sound\_QStreettBattery\_tie}; hash-bound in this edition.)
\end{proof}

\begin{bremark}[why the orientations are not symmetric, and not optional]
Swap either direction and the rule is unsound. If a certified request
region \emph{under}-approximated the complement, a path could request
infinitely often in the semantic sense while the certificate sees only
finitely many requests --- the certified pair implication holds
vacuously, the semantic one fails. Dually, a response region
\emph{over}-approximating $[\![\varphi_i]\!]$ lets certified
``responses'' land outside the actual response set. The hypothesis-side
/ conclusion-side asymmetry of implication is the entire content of the
direction check, and it is why the complement orientation swap --- which
produces exactly over-approximations --- is fitted to requests and to
nothing else.
\end{bremark}

\begin{bremark}[profile residue]
At the discrete-kernel profile the lemma chain above is complete,
machine-checked, and hash-bound in this edition. At the
continuous-time martingale-problem profile the bundle side is
discharged whole per the bundle section\textquotesingle s profile
remark (confinement, zero-decrement, and strict-decrement rows;
hash-bound there); what keeps this entry\textquotesingle s
conditional mark is the complement consumption\textquotesingle s own
fine-topology content (regular points of the
complement region), which remains prose in the design record
--- a residue of this entry\textquotesingle s own, not of the
bundle\textquotesingle s.
\end{bremark}

\begin{anchors}
Recipe \ref{rec:scr-instance} target & the \rulename{QStreettBattery} section\\
atom-level complement step & \texttt{denot\_elimNeg} (\texttt{BernsteinLogicAux.NNFDuality}; machine-checked, hash-bound)\\
Proposition \ref{prop:scr-mono} & \texttt{streett\_event\_mono} (\texttt{BernsteinLogicAux.QStreett}; machine-checked, hash-bound)\\
Proposition \ref{prop:scr-sound} & \texttt{sound\_QStreettBattery\_tie} (same module, hash-bound)\\
no wire object in this edition & producers rewrite before emission (Recipe~\ref{rec:scr-instance})\\
\end{anchors}


\rulesection{52}{QStreettTargetCellCollapse}{alias for \textcolor{DiffBlue}{\rulename{QStreettBattery}}}
\label{blcatrule-QStreettTargetCellCollapse}
\ruleblock{The rule (Steele metanotation)}
\sourcenote{as realized in the mechanized development; source row: \texttt{\texttt{QTSL-rules.toml}} --- \emph{QStreettTargetCellCollapse}. The metanotation statement is typeset by hand for notational quality; its mechanized ground truth is the Lean statement below, extracted from the sources at every build of this book and hash-bound. The rule \emph{roster} is machine-checked (every mechanized rule tag round-trips to exactly one catalogue row); the statement \emph{text} is editorial, checked by eye against the extracted Lean beside it.}
\begin{mathpar}
\inferrule*[right={$g(\mathrm{slab}) \subseteq T$,\ $W_i \equiv 0$ on $T$\ (G5, recomputed)}]{
  \mathcal{K},\ \mathrm{slab} \vdash W_i \ge 0
}{
  \mathcal{K},\ \mathrm{slab} \vdash g \fatsemi W_i \le W_i
}
\end{mathpar}
$\mathcal{K}$ ranges over \rulename{QStreettBattery} kits; the slab region left of the turnstile scopes both sequents.

\ruleflavor{a deduction rule between sequents: the kits demand the nonnegativity premise anyway, and under the side conditions (checker-recomputed, never a trusted annotation) the reset lands where the ranking certificate is identically zero, so the conclusion follows pointwise from the premise. On slabs whose resets leave the zero region the side condition fails and the full guard row is owed as usual.}
\par\medskip
\rulename{QStreettTargetCellCollapse} is a deduction rule between
sequents of one \rulename{QStreettBattery} kit context (n\'e StreettTargetCellCollapse): with $\mathcal{K}$ the premise metavariable
ranging over bundle kits, it derives the ranking-guard row
$g \fatsemi W_i \le W_i$ on a guard slab from the nonnegativity row
$W_i \ge 0$ there --- a row the kit demands anyway --- under two side
conditions: the containment of the reset image in the zero-convention
region ($g(\mathrm{slab}) \subseteq T$, with $W_i \equiv 0$ on $T$).
The rule is sound because under the side conditions the
conclusion\textquotesingle s left side is literally zero on the slab.
Both side conditions are recomputed geometric checks, never trusted
annotations --- which is also why the rule cannot be abused: on a slab
whose resets leave the zero region, the side condition fails and the
full guard row is owed as usual.

Mechanization state, plainly: the bound statement below is the one-line
witness \texttt{ranking\_guard\_collapse} (module
\texttt{BernsteinLogicAux.QStreett}) --- under the zero-on-target
convention and the recomputed containment, the guard row follows from
the nonneg row on the slab, which is the entirety of the
rule\textquotesingle s semantic claim (pure state-space set algebra: no
kernel, no profile, no probability enters, so no profile residue
attaches and the entry carries no conditional flag). The payload-level
unpack owes no sort of its own --- the extension certificate
schema\textquotesingle s adjudication of the batch family
(\rulename{QStreettDegenerate} section, Recipe~\ref{rec:sd-template}):
the collapse is a citation (the guard row cites the nonneg
row\textquotesingle s proof), and the containment is a recomputed side
condition of the battery walk, never a stored row --- so there is
nothing to elaborate, by design.

\cataloguenote{The checker\textquotesingle s implementation is citation-shaped,
deliberately decoupled from the deduction-rule presentation: no wire
object ever carries a collapse node --- a producer that has already
discharged the nonneg row cites it where the guard row was owed and
emits nothing new, one row to one row. Tier flag, carried honestly from
the source catalogue: that one-to-one citation shape is the alias
tiering, but the licensing side condition is a recomputed containment
--- if side-condition weight is ever ruled heavier than the alias
criterion intends, the rule re-tiers admissible with machinery
unchanged. Flagged, not forced.
}
\ruleblock{Lean statement}
\declhead{BernsteinLogicAux.QStreett}{ranking\_guard\_collapse}{\badge{BadgeGreen}{statement bound}}
\begin{leansrc}
/-- **`QStreettTargetCellCollapse` (T9c)**: when the reset of a guard transition
lands inside a region where the ranking certificate is identically zero (the
`Wᵢ ≡ 0`-on-target-cells convention, G5), the ranking-guard row
`W (g y) ≤ W y` collapses to the nonneg row `0 ≤ W y` — one row to one row,
under the recomputed containment side-condition `g '' slab ⊆ T`. -/
theorem ranking_guard_collapse {W : S → ℝ} {T slab : Set S} {g : S → S}
    (hzero : ∀ y ∈ T, W y = 0) (himg : ∀ y ∈ slab, g y ∈ T) :
    ∀ y ∈ slab, 0 ≤ W y → W (g y) ≤ W y
\end{leansrc}
\ruleblock{Correspondence (Steele display $\leftrightarrow$ Lean)}
\badge{BadgeGreen}{reviewer-certified against the displayed texts}
\sourcenote{The dictionary below is a reviewer-certified statement of how the mechanized statements above interpret the displayed rule: display metavariable $\mapsto$ Lean binder, each display hypothesis $\mapsto$ the premise or side condition carrying it, the display conclusion $\mapsto$ the goal --- with any deliberate difference in width between display and mechanization named outright. Certification is recorded as a hash triple over (the displayed cell, the extracted statement texts, this dictionary\textquotesingle s own text); if any of the three changes, this block turns amber until a reviewer re-certifies. The correspondence claim itself remains editorial --- reviewer-grade, not machine-checked; its machine-checked successor types each soundness statement against the rule as a structured object, and supersedes this block per rule as it lands.}
One bound statement, and the displayed deduction rule is its literal
shape: \texttt{ranking\_guard\_collapse} concludes \texttt{∀ y ∈ slab,
0 ≤ W y → W (g y) ≤ W y} --- the premise sequent as the antecedent, the
concluded ranking-guard row as the consequent, per point of the slab. The two
licensing conditions map as: $W_i \equiv 0$ on target cells
$\mapsto$ \texttt{hzero : ∀ y ∈ T, W y = 0} (the G5 convention as a
hypothesis), and the recomputed containment $g(\mathrm{slab})
\subseteq T$ $\mapsto$ \texttt{himg : ∀ y ∈ slab, g y ∈ T} ---
consumed as a hypothesis, recomputed checker-side, which is the
tier-ambiguity flag\textquotesingle s exact locus rendered as a
binder. Width note: the statement is one line over arbitrary sets
and an arbitrary reset map \texttt{g : S → S} --- no slab geometry,
no polynomiality of $W$, no guard structure --- strictly wider than
the displayed instance, safe direction; the failure branch the
flavor text names (resets leaving the zero region owe the full guard
row) is representable precisely as \texttt{himg} failing, whereupon
the theorem offers nothing --- refusal-shaped, as it should be.

\ruleblock{Lean verifier receipts}
\sourcenote{module \texttt{BernsteinLogicAux.QStreett}; bound build \texttt{469881beb770\ldots} (bound from receipt 9d1d848e0ba1718a, 2026-07-29T01:03:31Z). A receipt records that named checking layers accepted a build: (module, compiled-artifact SHA-256, environment SHA-256, checker layers), minted when the checking gate ran; the store receipts attesting the bound build are listed with the layers each run credited. Nothing here asks the reader to trust the store: every layer is an independently runnable program over the quoted sources --- rebuild the module and re-run the checkers to reproduce the verdict. Receipts attesting superseded builds credit no layer and are omitted from this rendering (the count of omissions, if any, is noted below).}
\begin{receipttable}
\receiptrow{e3990fe544d4708d}{2026-07-28T23:48:35Z}{build+decorrelated}{\badge{BadgeGreen}{current}}
\receiptrow{68e9374fa1b418fe}{2026-07-29T00:00:34Z}{build+decorrelated}{\badge{BadgeGreen}{current}}
\receiptrow{e5895c9706854cd3}{2026-07-29T00:49:04Z}{build+replay+hardened+decorrelated}{\badge{BadgeGreen}{current}}
\receiptrow{9d1d848e0ba1718a}{2026-07-29T01:03:31Z}{build+replay+hardened+decorrelated}{\badge{BadgeGreen}{current}}
\receiptrow{4a8d377b721b8523}{2026-07-31T22:00:09Z}{build+decorrelated}{\badge{BadgeGreen}{current}}
\receiptrow{d251a122540ab741}{2026-07-31T22:29:03Z}{build+decorrelated}{\badge{BadgeGreen}{current}}
\receiptrow{333635b8e047113a}{2026-08-02T23:31:14Z}{build+decorrelated}{\badge{BadgeGreen}{current}}
\receiptrow{9e96f80d63e4ad82}{2026-08-03T16:27:18Z}{build+decorrelated}{\badge{BadgeGreen}{current}}
\receiptrow{932b4b9e0b0fd323}{2026-08-03T22:46:02Z}{build+decorrelated}{\badge{BadgeGreen}{current}}
\end{receipttable}
\sourcenote{9 further receipt(s) in the store attest superseded builds of this module and are omitted from this rendering.}
\layerline{\layeritem{build}{lake build (elaboration)}{\badge{BadgeGreen}{current}}\layeritem{replay}{leanchecker (official kernel replay)}{\badge{BadgeGreen}{current}}\layeritem{hardened}{safe\_verify (hardened screen)}{\badge{BadgeGreen}{current}}\layeritem{decorrelated}{lean4export + nanoda (decorrelated kernel)}{\badge{BadgeGreen}{current}}}
\sourcenote{\badge{BadgeGray}{comparator: none} no fixture in the three-implementation comparator harness (compiled / kernel-interpreted / reference-interpreted) exercises this rule.}
\ruleblock{HOL4 statement}
\absentbox{OUT OF SCOPE BY DESIGN}{Out of scope for the HOL side \emph{by design}: the rule rewrites one
premise row of a bundle kit into another \emph{before}
emission --- no wire object ever contains a collapse node, and the
surviving row is an ordinary base-calculus sequent the base checker
already covers. The HOL statement that would matter is the
bundle\textquotesingle s own, absent for the reasons given in that
section. Checked against the HOL sources: the \texttt{bern\_catalogue}
theory carries nothing Streett-shaped.
}
\ruleblock{HOL verifier receipts}
\absentbox{NONE}{no HOL verification receipt exists for this rule.}
\ruleblock{Proof}

\sourcenote{Genre: alias rule (tier 3); the batch's shared template,
including the honesty clause about absent payload functions, is
Recipe~\ref{rec:sd-template} of the
\rulename{QStreettDegenerate} section. This section's rule is stated in
deduction-rule form --- the respelling content is carried by the
payload-transport recipe and the checker's citation-shaped
implementation, noted there.}

\begin{brecipe}[the rule --- a deduction between kit sequents]
\label{rec:tcc-filled}
Fix a \rulename{QStreettBattery} kit context $\mathcal{K}$, a pair $i$, a
guard transition with reset map $g$ and guard slab $\mathrm{slab}$, and
let $T$ be a region on which the kit's convention $W_i \equiv 0$ holds
(the target cells). The rule derives the ranking-guard row from the
nonnegativity row the kit demands anyway --- the kit's bound rows live
on the hull $\bar{D}$, so on guard slabs inside the hull the premise
below is a row already owed, while a slab outside them simply owes its
premise sequent separately; the deduction as displayed takes the
on-slab sequent itself as premise either way:
\begin{mathpar}
\inferrule*[right={$g(\mathrm{slab}) \subseteq T$,\ $W_i \equiv 0$ on $T$\ (G5, recomputed)}]{
  \mathcal{K},\ \mathrm{slab} \vdash W_i \ge 0
}{
  \mathcal{K},\ \mathrm{slab} \vdash g \fatsemi W_i \le W_i
}
\end{mathpar}
Both side conditions are checker-recomputed (the G5-convention check),
never producer-supplied. The premise metavariable does the honest work:
premise and conclusion are sequents of the \emph{same} kit context, so
admitting the rule claims exactly that any kit admitting the premise
admits the conclusion --- Proposition~\ref{prop:tcc-sound}, and nothing
else.
\end{brecipe}

\begin{brecipe}[payload transport, catalogued-field level]\label{rec:tcc-transport}
\begin{center}\small
\begin{tabular}{@{}ll@{}}
\textbf{collapse node} & \textbf{bundle kit row}\\
\midrule
guard row $(W_i, g, \mathrm{slab})$ & the kit's nonneg row for $W_i$ on $\mathrm{slab}$, cited (licit where the slab lies in the cited row's region)\\
containment $g(\mathrm{slab}) \subseteq T$ & recomputed side-condition (never producer-supplied)\\
zero convention $W_i \equiv 0$ on $T$ & the kit's target-cell convention, verbatim (that it holds on $T$ is the recomputed G5 check, not a trusted annotation)\\
\end{tabular}
\end{center}
The checker's implementation is citation-shaped: a producer that has
already discharged the nonnegativity row cites it where the guard row was
owed, and emits nothing new --- one row to one row, which is the alias
tiering. Presentation and implementation are deliberately decoupled: the
rule \emph{is} the deduction above; the citation is how a certificate
spells an application of it. Unpack-function slot: nothing to
elaborate, by design --- the collapse is a citation (the guard row
cites the nonneg row\textquotesingle s proof) and the containment is a
recomputed side condition, never a stored row; the semantic witness
\texttt{ranking\_guard\_collapse} (auxiliary module) is hash-bound in
this edition.
\end{brecipe}

\begin{bprop}[the exactness identity --- one line]\label{prop:tcc-exact}
For every $y \in \mathrm{slab}$: if $g(y) \in T$ and $W_i \equiv 0$ on
$T$ and $0 \le W_i(y)$, then $W_i(g(y)) = 0 \le W_i(y)$.
\end{bprop}
\begin{proof}
Substitution --- the displayed chain is the proof. The guard row's left
side is identically zero on the slab, so the guard row and the nonneg
row assert the same inequality there.
\end{proof}

\begin{bprop}[soundness corollary --- admissibility of the rule]
\label{prop:tcc-sound}
Any bundle kit admitting the nonneg row on the slab admits the
ranking-guard row on the slab --- the displayed rule, read as a claim;
hence every conclusion drawn through \rulename{QStreettBattery} from the
collapsed kit is drawn from the full kit, with the bundle's own
soundness theorem as warrant.
\end{bprop}
\begin{proof}
Proposition~\ref{prop:tcc-exact} produces the guard row pointwise from
the nonneg row under the side conditions --- that collapse is the
machine-checked content (\texttt{ranking\_guard\_collapse}; hash-bound
in this edition), and it is the whole of what this section proves. The
guard row is consumed \emph{upstream} of the flattened kernel: for a
hybrid presentation it is one of the base-calculus sequents from which
the kernel-integral ranking rows --- the bundle kit's H4, which the
bundle theorem is what actually consumes --- are assembled by the
presentation-to-kernel compile bridge (the safety
battery\textquotesingle s staged form carries that bridge as its
explicit \texttt{hmix} hypothesis; the bundle theorem takes the
compiled kernel and its H4 rows directly), never re-proven here. The corollary therefore lives at the kit stratum: the
collapsed kit discharges every base-calculus obligation the full
kit\textquotesingle s guard row was owed for, and every conclusion
drawn through \rulename{QStreettBattery} whose compile bridge owed that
guard row is drawn from the full kit, the bundle\textquotesingle s own
soundness theorem --- its assembly hypothesis unchanged --- supplying
the rest.
\end{proof}

\begin{bremark}[why the side-condition is not decoration]
Drop the containment and the collapse is unsound: a reset landing at a
state with $W_i(g(y)) > W_i(y)$ (outside the zero region) breaks the
upstream derivation of the ranking generator rows that the bundle's
telescoping lemma consumes --- the certificate could be pumped back up
at each guard
crossing, and the summability of un-responded requests
(\rulename{QStreettBattery} section) fails. The containment is
checker-recomputed for precisely this reason; its weight is the tier
flag recorded in the commentary.
\end{bremark}

\begin{anchors}
Recipe \ref{rec:tcc-filled} premise context $\mathcal{K}$ & the \rulename{QStreettBattery} section's kit (H4 rows)\\
Proposition \ref{prop:tcc-exact} & \texttt{ranking\_guard\_collapse} (\texttt{BernsteinLogicAux.QStreett}; machine-checked, hash-bound)\\
side conditions & the zero-on-target containment, recomputed (the bundle's G5-convention check)\\
\end{anchors}


\rulesection{53}{QNnfRewrite}{alias rule}
\label{blcatrule-QNnfRewrite}
\ruleblock{The rule (Steele metanotation)}
\sourcenote{as realized in the mechanized development; source row: \texttt{\texttt{QTSL-rules.toml}} --- \emph{QNnfRewrite}. The metanotation statement is typeset by hand for notational quality; its mechanized ground truth is the Lean statement below, extracted from the sources at every build of this book and hash-bound. The rule \emph{roster} is machine-checked (every mechanized rule tag round-trips to exactly one catalogue row); the statement \emph{text} is editorial, checked by eye against the extracted Lean beside it.}
\begin{mathpar}
\inferrule*{ }{ [\![\neg\neg\varphi]\!] = [\![\varphi]\!] }
\and
\inferrule*{ }{ [\![\neg(\varphi \wedge \psi)]\!] = [\![\neg\varphi \vee \neg\psi]\!] }
\and
\inferrule*{ }{ [\![\neg(\varphi \vee \psi)]\!] = [\![\neg\varphi \wedge \neg\psi]\!] }
\and
\inferrule*{ }{ [\![\neg\{f \le c\}]\!] = \{f > c\} }
\and
\inferrule*{ }{ [\![\neg\{f < c\}]\!] = \{f \ge c\} }
\and
\inferrule*{ }{ [\![\neg\{f \ge c\}]\!] = \{f < c\} }
\and
\inferrule*{ }{ [\![\neg\{f > c\}]\!] = \{f \le c\} }
\end{mathpar}

\ruleflavor{seven axioms --- deduction rules with nothing above the line --- at the spec surface, each concluding its region-denotation equality; a satisfaction judgment rewrites across any of them in either direction with certified region and rows transported verbatim, and no negation node survives to the baked representation (surviving negative literals are the typing rejection\textquotesingle s business, not a rule\textquotesingle s).}
\par\medskip
\rulename{QNnfRewrite} is the spec-surface negation calculus (n\'e NNFDuality): double negation, De Morgan through
$\wedge/\vee$, and the four atom/strict swaps
$\neg\{f \le c\} \leftrightarrow \{f > c\}$,
$\neg\{f < c\} \leftrightarrow \{f \ge c\}$ --- both orientations of
both comparison directions. The seven dualities are \emph{axioms}:
deduction rules with nothing above the line, each concluding its
region-denotation equality, stated once per duality. Each licenses a
judgment-level $1 \to 1$ rewrite of a formula inside a satisfaction
judgment: the rewritten formula denotes the same set of states, so any
region under-approximating one denotation under-approximates the other,
and the judgment transports unchanged. A
producer uses the dualities to push negations down to atoms (the
negation-normal-form pass of the \rulename{QStreettRequestNnf} section)
so that every negative position is consumed either by an atom swap or by
the typing pass\textquotesingle s rejection --- never by a negation
node in the certificate. The checker recomputes nothing: the dualities
are definitional set identities over the state space, and the
strict-inequality sides interact with the closure discipline (rounding
of strict literals at positive positions) only through the separate
closure rules, deliberately out of this rule\textquotesingle s scope.

There is no single target rule, and hence no \emph{alias-for} header
here: each axiom stands on its own definitional equality.
Mechanization state, plainly: the bound
statements below are the four atom swaps, machine-checked as set-level
denotation equalities in both orientations of both comparison
directions (module \texttt{BernsteinLogicAux.NNFDuality}:
\texttt{compl\_le\_atom}, \texttt{compl\_lt\_atom},
\texttt{compl\_ge\_atom}, \texttt{compl\_gt\_atom}), together with
the normalization pass\textquotesingle s denotation-preservation
theorem \texttt{denot\_nnfAux}, inside which double negation and De
Morgan are consumed from the ambient set library --- each applied once
per duality, so the displayed equalities are covered jointly by the
swaps (atoms) and the pass theorem (connectives). These are definitional
set identities over the state space, warranted in full; the entry
carries no conditional flag.

\cataloguenote{Both orientations of both directions are stated deliberately: an early
one-sided draft of the atom swaps was caught in review, and the repaired
statement is the four-lemma family. The strict-atom sides are stated
here as \emph{denotation} facts only; whether a strict literal may
appear in a concluded region is the closure discipline\textquotesingle s
question, answered by polarity-directed rounding in its own catalogue
rows.
}
\ruleblock{Lean statement}
\declhead{BernsteinLogicAux.NNFDuality}{compl\_le\_atom}{\badge{BadgeGreen}{statement bound}}
\begin{leansrc}
/-- T12 atom swap: `¬{f ≤ c} ↔ {f > c}`. -/
theorem compl_le_atom : {x | f x ≤ c}ᶜ = {x | c < f x}
\end{leansrc}
\declhead{BernsteinLogicAux.NNFDuality}{compl\_lt\_atom}{\badge{BadgeGreen}{statement bound}}
\begin{leansrc}
/-- T12 atom swap: `¬{f < c} ↔ {f ≥ c}`. -/
theorem compl_lt_atom : {x | f x < c}ᶜ = {x | c ≤ f x}
\end{leansrc}
\declhead{BernsteinLogicAux.NNFDuality}{compl\_ge\_atom}{\badge{BadgeGreen}{statement bound}}
\begin{leansrc}
/-- T12 atom swap, the other orientation: `¬{f ≥ c} ↔ {f < c}`. -/
theorem compl_ge_atom : {x | c ≤ f x}ᶜ = {x | f x < c}
\end{leansrc}
\declhead{BernsteinLogicAux.NNFDuality}{compl\_gt\_atom}{\badge{BadgeGreen}{statement bound}}
\begin{leansrc}
/-- T12 atom swap, the other orientation: `¬{f > c} ↔ {f ≤ c}`. -/
theorem compl_gt_atom : {x | c < f x}ᶜ = {x | f x ≤ c}
\end{leansrc}
\declhead{BernsteinLogicAux.NNFDuality}{denot\_nnfAux}{\badge{BadgeGreen}{statement bound}}
\begin{leansrc}
/-- **Denotation preservation of the NNF pass, both polarities** — the T12
dualities (double negation `compl_compl`, De Morgan `Set.compl_inter`/
`Set.compl_union`) are exactly the cases of this induction; each is applied
once, which is the catalogue's "stated once as a definitional axiom per
duality" realized as proof structure. -/
theorem denot_nnfAux : ∀ (φ : Form α) (b : Bool),
    (φ.nnfAux b).denot den = if b then φ.denot den else (φ.denot den)ᶜ
\end{leansrc}
\ruleblock{Correspondence (Steele display $\leftrightarrow$ Lean)}
\badge{BadgeGreen}{reviewer-certified against the displayed texts}
\sourcenote{The dictionary below is a reviewer-certified statement of how the mechanized statements above interpret the displayed rule: display metavariable $\mapsto$ Lean binder, each display hypothesis $\mapsto$ the premise or side condition carrying it, the display conclusion $\mapsto$ the goal --- with any deliberate difference in width between display and mechanization named outright. Certification is recorded as a hash triple over (the displayed cell, the extracted statement texts, this dictionary\textquotesingle s own text); if any of the three changes, this block turns amber until a reviewer re-certifies. The correspondence claim itself remains editorial --- reviewer-grade, not machine-checked; its machine-checked successor types each soundness statement against the rule as a structured object, and supersedes this block per rule as it lands.}
The displayed atom/strict swaps map one-to-one onto the four bound
equalities, each a region-denotation identity:
$[\![\neg\{f \le c\}]\!] = \{f > c\}$ $\mapsto$
\texttt{compl\_le\_atom : \{x | f x ≤ c\}ᶜ = \{x | c < f x\}}, and
likewise \texttt{compl\_lt\_atom} ($\neg\{f < c\} = \{f \ge c\}$),
\texttt{compl\_ge\_atom}, \texttt{compl\_gt\_atom} (the reversed
orientations --- both directions bound separately, honouring the
recorded one-sidedness catch). Binder note, wider-safe: \texttt{f : σ
→ ℝ} is an arbitrary function on an arbitrary carrier --- the
polynomial atom structure of the certificate format is erased, so
the certified equalities cover every instance the rule will ever be
applied to. The propositional dualities ($\neg\neg\varphi$, De
Morgan both ways) $\mapsto$ the \texttt{neg}/\texttt{and}/\texttt{or}
cases of \texttt{denot\_nnfAux}\textquotesingle s both-polarity
induction, where \texttt{compl\_compl} /
\texttt{Set.compl\_inter} / \texttt{Set.compl\_union} are each
consumed from mathlib once --- the entry\textquotesingle s
``stated once per duality'' realized as proof structure rather than
as four more standalone statements. Width note: the entry is
unconditional --- set-level witnesses cover the catalogued claim in
full (axioms: \texttt{propext} alone) --- and the
grammar-level guarantee the flavor text mentions (no negation node
survives to the baked representation) is certified not here but by
the QStreettRequestNnf entry\textquotesingle s normal-form \emph{type}; this
page\textquotesingle s statements are the denotation half only, as
its note already rules.

\ruleblock{Lean verifier receipts}
\sourcenote{module \texttt{BernsteinLogicAux.NNFDuality}; bound build \texttt{97564f8e3950\ldots} (bound from receipt 21f04961748c0302, 2026-07-29T01:03:33Z). A receipt records that named checking layers accepted a build: (module, compiled-artifact SHA-256, environment SHA-256, checker layers), minted when the checking gate ran; the store receipts attesting the bound build are listed with the layers each run credited. Nothing here asks the reader to trust the store: every layer is an independently runnable program over the quoted sources --- rebuild the module and re-run the checkers to reproduce the verdict. Receipts attesting superseded builds credit no layer and are omitted from this rendering (the count of omissions, if any, is noted below).}
\begin{receipttable}
\receiptrow{1417cdf27f69842a}{2026-07-28T23:48:35Z}{build+decorrelated}{\badge{BadgeGreen}{current}}
\receiptrow{65fc1b51055d38bd}{2026-07-29T00:00:33Z}{build+decorrelated}{\badge{BadgeGreen}{current}}
\receiptrow{2abbbc61b6065811}{2026-07-29T00:49:05Z}{build+replay+hardened+decorrelated}{\badge{BadgeGreen}{current}}
\receiptrow{21f04961748c0302}{2026-07-29T01:03:33Z}{build+replay+hardened+decorrelated}{\badge{BadgeGreen}{current}}
\receiptrow{8879c5b9721acf0c}{2026-07-31T22:00:03Z}{build+decorrelated}{\badge{BadgeGreen}{current}}
\receiptrow{e47a2c2d7f152711}{2026-07-31T22:28:57Z}{build+decorrelated}{\badge{BadgeGreen}{current}}
\receiptrow{ce1bb97f7cbcaf4f}{2026-08-02T23:31:07Z}{build+decorrelated}{\badge{BadgeGreen}{current}}
\receiptrow{8755b28f77990ba0}{2026-08-03T16:27:11Z}{build+decorrelated}{\badge{BadgeGreen}{current}}
\receiptrow{64948e60a56ae3bc}{2026-08-03T22:45:55Z}{build+decorrelated}{\badge{BadgeGreen}{current}}
\end{receipttable}
\sourcenote{8 further receipt(s) in the store attest superseded builds of this module and are omitted from this rendering.}
\layerline{\layeritem{build}{lake build (elaboration)}{\badge{BadgeGreen}{current}}\layeritem{replay}{leanchecker (official kernel replay)}{\badge{BadgeGreen}{current}}\layeritem{hardened}{safe\_verify (hardened screen)}{\badge{BadgeGreen}{current}}\layeritem{decorrelated}{lean4export + nanoda (decorrelated kernel)}{\badge{BadgeGreen}{current}}}
\sourcenote{\badge{BadgeGray}{comparator: none} no fixture in the three-implementation comparator harness (compiled / kernel-interpreted / reference-interpreted) exercises this rule.}
\ruleblock{HOL4 statement}
\absentbox{OUT OF SCOPE BY DESIGN}{Out of scope for the HOL side \emph{by design}: the dualities are
producer-side formula rewrites that unpack to definitional denotation
equalities before any certificate is emitted --- no wire object carries
a negation node (the typing discipline rejects surviving negative
literals), so there is nothing for a wire checker, HOL or otherwise, to
verify. Checked against the HOL sources: the \texttt{bern\_catalogue}
theory is the base-calculus checker cone; formula-level negation of the
extension theory\textquotesingle s spec surface does not appear in it.
}
\ruleblock{HOL verifier receipts}
\absentbox{NONE}{no HOL verification receipt exists for this rule.}
\ruleblock{Proof}

\sourcenote{Genre: axioms at the spec surface (roster tier 3). The seven
dualities are deduction rules with nothing above the line, each
concluding a definitional region-denotation equality --- there is no
single target rule being respelled, and no premise to carry: the
equality is the entire content. The payload-transport honesty clause of
the batch is Recipe~\ref{rec:sd-template} of the
\rulename{QStreettDegenerate} section, and applies here vacuously (no
payload ever exists to transport --- the dualities are applied before
emission).}

\begin{brecipe}[the seven axioms]
\label{rec:nnf-filled}
Over a state space $S$, with $[\![\cdot]\!]$ the region denotation of
spec-surface formulas and $f : S \to \R$ ranging over atom left-hand
sides:
\begin{mathpar}
\inferrule*{ }{ [\![\neg\neg\varphi]\!] = [\![\varphi]\!] }
\and
\inferrule*{ }{ [\![\neg(\varphi \wedge \psi)]\!] = [\![\neg\varphi \vee \neg\psi]\!] }
\and
\inferrule*{ }{ [\![\neg(\varphi \vee \psi)]\!] = [\![\neg\varphi \wedge \neg\psi]\!] }
\and
\inferrule*{ }{ [\![\neg\{f \le c\}]\!] = \{f > c\} }
\and
\inferrule*{ }{ [\![\neg\{f < c\}]\!] = \{f \ge c\} }
\and
\inferrule*{ }{ [\![\neg\{f \ge c\}]\!] = \{f < c\} }
\and
\inferrule*{ }{ [\![\neg\{f > c\}]\!] = \{f \le c\} }
\end{mathpar}
Each axiom licenses a judgment-level $1 \to 1$ rewrite in either
direction: a satisfaction judgment whose formula contains one side
concludes identically with the other side substituted, the certified
region unchanged.
\end{brecipe}

\begin{brecipe}[payload transport]\label{rec:nnf-transport}
\begin{center}\small
\begin{tabular}{@{}ll@{}}
\textbf{rewritten node} & \textbf{image}\\
\midrule
formula (one duality applied at one position) & formula (the dual form at that position)\\
certified region, judgment kind, all rows & verbatim\\
\end{tabular}
\end{center}
Unpack-function slot: no verified payload function in this edition
(batch template) ---
and here the slot is moot by design: the typing
discipline rejects surviving negative literals, so the dualities exist
to be applied \emph{before} any payload exists.
\end{brecipe}

\begin{bprop}[soundness of the axioms --- two lines each]
\label{prop:nnf-exact}
All seven concluded equalities hold.
\end{bprop}
\begin{proof}
The propositional three are complement involution and De Morgan for
sets. The atom four: pointwise over $S$, the negation of $f(y) \le c$ is
$c < f(y)$ and the negation of $f(y) < c$ is $c \le f(y)$ --- the total
order on $\R$ --- and the remaining two are the same facts read in the
other orientation. (Machine-checked: the four atom swaps as set-level
equalities, and the propositional three consumed from the ambient set
library inside the normalization pass's denotation theorem; hash-bound
in this edition.)
\end{proof}

\begin{bprop}[soundness corollary]\label{prop:nnf-sound}
If $R \subseteq [\![\varphi]\!]$ and $\varphi'$ arises from $\varphi$ by
one duality, then $R \subseteq [\![\varphi']\!]$; hence every
satisfaction judgment transports across a duality unchanged, in both
directions.
\end{bprop}
\begin{proof}
Denotations are equal (Proposition~\ref{prop:nnf-exact}), and
substitution preserves that: region denotation is compositional ---
each connective's denotation is a fixed set operation on its
children's --- so by induction on the surrounding formula, replacing a
subformula by one of equal denotation leaves every enclosing
denotation unchanged. The machine-checked cousin of this one-position
congruence is the normalization pass's denotation-preservation
theorem, which pushes negation through \emph{all} eligible positions
in one traversal. The pass consumes the propositional facts --- the
complement of a conjunction or disjunction, once per constructor
case --- and treats atoms opaquely: a negatively-occurring atom
survives as a negated literal whose denotation is the complement, and
the four atom dualities are separate bound facts that swap that
complement for its closed-form partner. That division realizes what
the display means by ``stated once per duality'': each equality is
one definitional fact, cited where it is used, never re-derived.
\end{proof}

\begin{bremark}[what this rule deliberately does not say]
The strict-atom swaps are \emph{denotation} facts; whether a strict
inequality may appear in a concluded region atom is a different
question, answered by the closure discipline (polarity-directed
rounding) in its own catalogue rows. And the probability operator is
absent from the display on purpose: a negated probability bound at even
polarity is not rewritten but \emph{rejected} by the typing pass --- the
polarity theorem of the \rulename{QStreettRequestNnf} section is what
makes that rejection a surface-level check.
\end{bremark}

\begin{anchors}
atom swaps, both orientations & \texttt{compl\_le\_atom}, \texttt{compl\_lt\_atom}, \texttt{compl\_ge\_atom}, \texttt{compl\_gt\_atom} (\texttt{BernsteinLogicAux.NNFDuality}; machine-checked, hash-bound)\\
Proposition \ref{prop:nnf-sound} induction & \texttt{denot\_nnfAux} (same module, hash-bound; the whole-pass form --- the one-position congruence is the compositionality induction in the proof)\\
no wire object, by design & negation nodes are rejected by typing; producers rewrite before emission\\
\end{anchors}


\rulesection{54}{QStreettRequestNnf}{alias for \textcolor{DiffBlue}{\rulename{QStreettBattery}}}
\label{blcatrule-QStreettRequestNnf}
\ruleblock{The rule (Steele metanotation)}
\sourcenote{as realized in the mechanized development; source row: \texttt{\texttt{QTSL-rules.toml}} --- \emph{QStreettRequestNnf}. The metanotation statement is typeset by hand for notational quality; its mechanized ground truth is the Lean statement below, extracted from the sources at every build of this book and hash-bound. The rule \emph{roster} is machine-checked (every mechanized rule tag round-trips to exactly one catalogue row); the statement \emph{text} is editorial, checked by eye against the extracted Lean beside it.}
\begin{mathpar}
\inferrule*{
  \mathcal{M},\ x_0 \in R \vDash \mathsf{SProbGE}\bigl(p,\ \mathsf{PDwellStreett}(\ldots,\ \mathsf{PairCmpl}(\mathrm{elimNeg}(\mathrm{nnf}(\rho_i)),\ \varphi_i),\ \ldots)\bigr)
}{
  \mathcal{M},\ x_0 \in R \vDash \mathsf{SProbGE}\bigl(p,\ \mathsf{PDwellStreett}(\ldots,\ \mathsf{PairCmpl}(\rho_i,\ \varphi_i),\ \ldots)\bigr)
}
\and
\inferrule*{ }{ [\![\mathrm{elimNeg}(\mathrm{nnf}(\rho))]\!] = [\![\rho]\!] }
\end{mathpar}

\ruleflavor{per pair $(\rho_i, \varphi_i)$, applied at each $\mathrm{cmpl}$-typed pair while base-typed requests pass unchanged: the request \emph{body} is rewritten to literal normal form by the QNnfRewrite dualities --- negation-freedom of the output holds \emph{by construction} of the normal-form grammar (no negation node exists to build) --- with the $\mathrm{cmpl}$ tag RETAINED: the complement applies at the consumption boundary (QStreettCmplRequest), never as a literal. The displayed axiom is the warrant: complements of equal denotations are equal, so the retained tag consumes identically. Polarity theorem: a surviving negative literal is exactly an odd-polarity atom occurrence of the input --- so the neg-P rejection is a surface polarity check, stable under normalization.}
\par\medskip
\rulename{QStreettRequestNnf} normalizes each Streett pair into the
bundle\textquotesingle s expected input shape (n\'e StreettRequestNNF): base-typed requests pass through unchanged, and a
complement-typed request has its formula rewritten to negation normal
form --- negations pushed to the atoms by the \rulename{QNnfRewrite}
rewrites --- so that the complement is applied at the \emph{consumption
boundary} (the atom swaps, or the complement-request rule), never
carried as a formula literal. It is an alias: the rewrite is a $1 \to 1$
respelling of one pair inside one \rulename{QStreettBattery} node, and
its warrant is denotation preservation of the normalization pass ---
the region the pair constrains is unchanged, so the judgment transports
verbatim.

The load-bearing extra is the \emph{polarity theorem}: the surviving
negative literals of the normal form are exactly the odd-polarity atom
occurrences of the input formula. Its consequence is that the typing
pass\textquotesingle s rejection of negated probability literals ---
which guards against smuggling a probability \emph{upper} bound through
an even number of negations --- is equivalent to a surface-level
polarity check on the un-normalized formula: normalization can never
manufacture a rejected literal that was not already there at odd
polarity, and never loses one. What would break without it: a
normalization pass without the polarity invariant could present a
well-typed surface whose normal form hides a rejected literal, and the
consumption boundary would silently consume a bound in the wrong
direction.

Mechanization state, plainly: the pass and both theorems are
machine-checked structurally, and the bound statements below are their
statements (module \texttt{BernsteinLogicAux.NNFDuality}): the
normalizer \texttt{Form.nnf} lands in a literal normal form whose
\emph{type} has no negation node --- negation-freedom by construction,
the strongest form of the well-formedness claim --- with denotation
preservation \texttt{denot\_nnf} (the displayed
$[\![\mathrm{nnf}(\rho_i)]\!] = [\![\rho_i]\!]$), the polarity
theorem \texttt{negLits\_nnf} (surviving negative literals $=$
odd-polarity atom occurrences) and its rejection face
\texttt{nnf\_no\_negLit\_iff}, and the consumption-boundary
elimination \texttt{denot\_elimNeg}. The rule\textquotesingle s scope
is the propositional skeleton, deliberately --- the path-operator nodes
of the spec surface are handled by their own ruled aliases and typing
rejections, never by this pass --- and the witnesses cover that scope in
full, so the entry carries no conditional flag.

\cataloguenote{The pass is a reference implementation on the producer side; nothing
checker-facing depends on running it. What the certificate format sees
is only the normalized result, and the typing pass independently
rejects any surviving negative probability literal --- so a defective
producer normalization can produce rejections, never unsoundness.
}
\ruleblock{Lean statement}
\declhead{BernsteinLogicAux.NNFDuality}{denot\_nnf}{\badge{BadgeGreen}{statement bound}}
\begin{leansrc}
/-- Denotation preservation of `Form.nnf` (T22's soundness face). -/
theorem denot_nnf (φ : Form α) : (φ.nnf).denot den = φ.denot den
\end{leansrc}
\declhead{BernsteinLogicAux.NNFDuality}{negLits\_nnf}{\badge{BadgeGreen}{statement bound}}
\begin{leansrc}
/-- **The T22 polarity theorem** ("a surviving neg-P literal is ALWAYS the
rejected even-polarity form", catalogue T22 — the wf theorem): the surviving
negative literals of the normal form are exactly the **odd-polarity** atom
occurrences of the input. Consequently the typing pass's neg-literal rejection
(`qtsl_typing/v1`, catalogue R6) is equivalent to a surface-level polarity
check on the un-normalized formula. -/
theorem negLits_nnf (φ : Form α) : (φ.nnf).negLits = φ.occs false
\end{leansrc}
\declhead{BernsteinLogicAux.NNFDuality}{nnf\_no\_negLit\_iff}{\badge{BadgeGreen}{statement bound}}
\begin{leansrc}
/-- The rejection face of the polarity theorem: the normal form is free of
`P`-flagged negative literals iff no `P`-atom occurs at odd polarity. -/
theorem nnf_no_negLit_iff (φ : Form α) (P : α → Prop) :
    (∀ a ∈ (φ.nnf).negLits, ¬ P a) ↔ ∀ a ∈ φ.occs false, ¬ P a
\end{leansrc}
\declhead{BernsteinLogicAux.NNFDuality}{denot\_elimNeg}{\badge{BadgeGreen}{statement bound}}
\begin{leansrc}
/-- **Denotation preservation of dual-atom elimination**: where each dual atom
denotes the complement (`hd` — the per-duality definitional axioms, e.g. the
`compl_le_atom` family above), elimination preserves the region. This is the
lawfulness of consuming a `cmpl`-typed request through swapped atoms
(atom-level complement elimination's atom step, the T9b input shape). -/
theorem denot_elimNeg {σ : Type*} (den : α → Set σ) (d : α → α)
    (hd : ∀ a, den (d a) = (den a)ᶜ) :
    ∀ φ : NForm α, (φ.elimNeg d).denot den = φ.denot den
\end{leansrc}
\ruleblock{Correspondence (Steele display $\leftrightarrow$ Lean)}
\badge{BadgeGreen}{reviewer-certified against the displayed texts}
\sourcenote{The dictionary below is a reviewer-certified statement of how the mechanized statements above interpret the displayed rule: display metavariable $\mapsto$ Lean binder, each display hypothesis $\mapsto$ the premise or side condition carrying it, the display conclusion $\mapsto$ the goal --- with any deliberate difference in width between display and mechanization named outright. Certification is recorded as a hash triple over (the displayed cell, the extracted statement texts, this dictionary\textquotesingle s own text); if any of the three changes, this block turns amber until a reviewer re-certifies. The correspondence claim itself remains editorial --- reviewer-grade, not machine-checked; its machine-checked successor types each soundness statement against the rule as a structured object, and supersedes this block per rule as it lands.}
The flavor text\textquotesingle s ``negation-freedom holds by
construction of the normal-form grammar'' is certified \emph{by
type}: the target grammar \texttt{NForm} (positive/negative literals
under $\wedge$/$\vee$) has no negation constructor, so no theorem
about negation-freedom exists to state --- the inductive definition
is the certificate. The displayed constituents --- the judgment-transport rule and its
warranting axiom --- map onto the
four bound statements:
\begin{itemize}
\item $[\![\mathrm{elimNeg}(\mathrm{nnf}(\rho_i))]\!] =
  [\![\rho_i]\!]$ $\mapsto$ the composition of \texttt{denot\_nnf}
  (\texttt{(φ.nnf).denot den = φ.denot den} --- the normalizer
  preserves the region, QNnfRewrite\textquotesingle s dualities consumed as
  the induction) with \texttt{denot\_elimNeg}, whose hypothesis
  \texttt{hd : ∀ a, den (d a) = (den a)ᶜ} is the per-duality
  definitional-axiom pack: elimination is lawful exactly where a
  dual atom exists, which is the flavor text\textquotesingle s
  ``the complement applies at the consumption boundary, never as a
  literal'' --- the retained $\mathrm{cmpl}$ tag consumes at QStreettCmplRequest
  through this same lemma, and dropping the tag would change the
  pair\textquotesingle s truth value (a one-state counterexample
  witnesses it), which is why the rewrite retains it.
\item the polarity theorem $\mapsto$ \texttt{negLits\_nnf :
  (φ.nnf).negLits = φ.occs false} --- surviving negative literals
  are \emph{exactly} the odd-polarity atom occurrences --- with the
  rejection face \texttt{nnf\_no\_negLit\_iff} certifying the
  flavor text\textquotesingle s ``the neg-P rejection is a surface
  polarity check, stable under normalization.''
\end{itemize}
Width note: atoms and denotations are fully abstract (\texttt{α},
\texttt{den : α → Set σ}) --- the per-pair Streett context, the
base-typed-request pass-through, and the probability flavour of the
neg-P check are all consumers of these statements, not content of
them; unconditional, \texttt{propext} alone, set-level witnesses
covering the catalogued claim in full.

\ruleblock{Lean verifier receipts}
\sourcenote{module \texttt{BernsteinLogicAux.NNFDuality}; bound build \texttt{97564f8e3950\ldots} (bound from receipt 21f04961748c0302, 2026-07-29T01:03:33Z). A receipt records that named checking layers accepted a build: (module, compiled-artifact SHA-256, environment SHA-256, checker layers), minted when the checking gate ran; the store receipts attesting the bound build are listed with the layers each run credited. Nothing here asks the reader to trust the store: every layer is an independently runnable program over the quoted sources --- rebuild the module and re-run the checkers to reproduce the verdict. Receipts attesting superseded builds credit no layer and are omitted from this rendering (the count of omissions, if any, is noted below).}
\begin{receipttable}
\receiptrow{1417cdf27f69842a}{2026-07-28T23:48:35Z}{build+decorrelated}{\badge{BadgeGreen}{current}}
\receiptrow{65fc1b51055d38bd}{2026-07-29T00:00:33Z}{build+decorrelated}{\badge{BadgeGreen}{current}}
\receiptrow{2abbbc61b6065811}{2026-07-29T00:49:05Z}{build+replay+hardened+decorrelated}{\badge{BadgeGreen}{current}}
\receiptrow{21f04961748c0302}{2026-07-29T01:03:33Z}{build+replay+hardened+decorrelated}{\badge{BadgeGreen}{current}}
\receiptrow{8879c5b9721acf0c}{2026-07-31T22:00:03Z}{build+decorrelated}{\badge{BadgeGreen}{current}}
\receiptrow{e47a2c2d7f152711}{2026-07-31T22:28:57Z}{build+decorrelated}{\badge{BadgeGreen}{current}}
\receiptrow{ce1bb97f7cbcaf4f}{2026-08-02T23:31:07Z}{build+decorrelated}{\badge{BadgeGreen}{current}}
\receiptrow{8755b28f77990ba0}{2026-08-03T16:27:11Z}{build+decorrelated}{\badge{BadgeGreen}{current}}
\receiptrow{64948e60a56ae3bc}{2026-08-03T22:45:55Z}{build+decorrelated}{\badge{BadgeGreen}{current}}
\end{receipttable}
\sourcenote{8 further receipt(s) in the store attest superseded builds of this module and are omitted from this rendering.}
\layerline{\layeritem{build}{lake build (elaboration)}{\badge{BadgeGreen}{current}}\layeritem{replay}{leanchecker (official kernel replay)}{\badge{BadgeGreen}{current}}\layeritem{hardened}{safe\_verify (hardened screen)}{\badge{BadgeGreen}{current}}\layeritem{decorrelated}{lean4export + nanoda (decorrelated kernel)}{\badge{BadgeGreen}{current}}}
\sourcenote{\badge{BadgeGray}{comparator: none} no fixture in the three-implementation comparator harness (compiled / kernel-interpreted / reference-interpreted) exercises this rule.}
\ruleblock{HOL4 statement}
\absentbox{OUT OF SCOPE BY DESIGN}{Out of scope for the HOL side \emph{by design}: the normalization is a
producer-side rewrite completed before emission --- certificates carry
only its output, and the output\textquotesingle s well-formedness is
enforced by the typing pass, not by a checker replay of the rewrite. No
wire object exists for a HOL statement to be about. Checked against the
HOL sources: the \texttt{bern\_catalogue} theory carries no spec-surface
formula grammar.
}
\ruleblock{HOL verifier receipts}
\absentbox{NONE}{no HOL verification receipt exists for this rule.}
\ruleblock{Proof}

\sourcenote{Genre: alias rule (tier 3) --- a respelling template; the
batch's shared template is Recipe~\ref{rec:sd-template} of the
\rulename{QStreettDegenerate}
section.}

\begin{brecipe}[the filled-in application --- one pair, normalized]
\label{rec:rnnf-filled}
Fix a Streett pair $(\rho, \varphi)$ inside a \rulename{QStreettBattery}
node, with $\rho$ complement-typed. Write $\mathrm{nnf}$ for the
negation-normal-form pass on the propositional skeleton (atoms,
$\neg$, $\wedge$, $\vee$): negations pushed inward by the
\rulename{QNnfRewrite} rewrites until only literals remain. The rewrite
is
\[
\bigl(\mathrm{cmpl}\,\rho,\ \varphi\bigr)
\quad\rightsquigarrow\quad
\bigl(\mathrm{cmpl}\,\bigl(\mathrm{elimNeg}(\mathrm{nnf}(\rho))\bigr),\ \varphi\bigr)
\]
--- the body normalized, the complement tag \emph{retained}: per the
source discipline the complement applies at the consumption boundary
(the complement-request rule), never as a literal. After the rewrite
every literal of the request body is positive (atom-swapped where it
was negative), and the pair is in exactly the input shape the
bundle and the complement-request rule expect. A base-typed pair is
already in that shape: on it the rewrite is the identity, the transport
trivial, and the exactness identity reflexivity --- stated once here so
the corollary\textquotesingle s case split below is complete. One pair
to one pair,
inside one node: the alias shape.
\end{brecipe}

\begin{brecipe}[payload transport]\label{rec:rnnf-transport}
\begin{center}\small
\begin{tabular}{@{}ll@{}}
\textbf{un-normalized pair} & \textbf{normalized pair}\\
\midrule
request body $\rho$ ($\mathrm{cmpl}$ tag retained) & $\mathrm{elimNeg}(\mathrm{nnf}(\rho))$, negative literals eliminated through dual atoms\\
response formula, certificates, rows & verbatim\\
polarity bookkeeping & the polarity theorem (Proposition~\ref{prop:rnnf-polarity}); nothing stored\\
\end{tabular}
\end{center}
Unpack-function slot: no verified payload function in this edition
(batch template); the
structural mechanization exists (the normalizer and its theorems, named
below; the four theorems hash-bound in this edition).
\end{brecipe}

\begin{bprop}[the exactness identity --- denotation preservation]
\label{prop:rnnf-exact}
$[\![\mathrm{nnf}(\rho)]\!] = [\![\rho]\!]$, and eliminating negative
literals through dual atoms preserves the denotation.
\end{bprop}
\begin{proof}
Induction on $\rho$, with the \rulename{QNnfRewrite} equalities as the
cases --- double negation and De Morgan drive the recursion, the atom
swaps close the leaves. The elimination step substitutes, for each
negative literal, the dual atom denoting the complement; equality is
preserved literal by literal. (Machine-checked: \texttt{denot\_nnf} and
\texttt{denot\_elimNeg}, hash-bound in this edition; the eliminated form
has no negative literals at
all, \texttt{negLits\_elimNeg} --- an unpinned auxiliary of the same
module.) Complements of equal
sets are equal, so the retained $\mathrm{cmpl}$ tag's
consumption-boundary denotation $S \smallsetminus [\![\rho]\!]$ is
preserved as well --- the rewrite itself never discharges the
complement.
\end{proof}

\begin{bprop}[the polarity theorem]\label{prop:rnnf-polarity}
The surviving negative literals of $\mathrm{nnf}(\rho)$ are exactly the
atom occurrences of $\rho$ at odd polarity (an odd number of enclosing
negations). Consequently the normal form is free of negated occurrences
of a given atom class if and only if the input has no odd-polarity
occurrence of that class.
\end{bprop}
\begin{proof}
Strengthen to the polarity-indexed pass (normalizing $\rho$ or
$\neg\rho$ according to a polarity bit) and induct: at an atom the bit
decides the literal's sign; at a negation the bit flips; at $\wedge$ and
$\vee$ the bit distributes (De Morgan swapping the connective when
negative). The negative-literal set on one side and the odd-polarity
occurrence set on the other satisfy the same recursion. The rejection
face is a rewrite of the same equality. (Machine-checked:
\texttt{negLits\_nnf}, via the unpinned auxiliary
\texttt{negLits\_nnfAux}, and
\texttt{nnf\_no\_negLit\_iff}; both theorems hash-bound in this edition.
Terminology, once: a \emph{negative literal} is a signed-literal
constructor of the normal-form grammar, not an occurrence of the
formula negation node --- the grammar carries literals of both signs
and no negation constructor, which is the sense in which
negation-freedom of the output holds by \emph{type}; the elimination
pass then removes the negative-signed literals through dual atoms.)
\end{proof}

\begin{bprop}[soundness corollary]\label{prop:rnnf-sound}
A bundle node whose pair is normalized concludes identically to the
un-normalized node: through \rulename{QStreettBattery} directly for
base-typed requests, and through the \rulename{QStreettCmplRequest}
weakening for the consumed complement --- with
Proposition~\ref{prop:rnnf-exact} guaranteeing both nodes constrain the
same regions.
\end{bprop}
\begin{proof}
Immediate from Proposition~\ref{prop:rnnf-exact} and the target rules'
theorems, each at its own proven scope --- the rewrite itself is
propositional and profile-free; whatever conditional mark the
concluding bundle or weakening node carries is inherited there, not
here; the alias adds no semantic content of its own.
\end{proof}

\begin{bremark}[what the polarity theorem is protecting]
The typing discipline rejects probability-atom occurrences at odd
polarity --- an odd number of enclosing negations, whether worn as a
bare negated literal in even context or as a positive literal under an
odd stack --- because a negated lower bound is an upper bound, and the
certificate calculus deliberately has no rule concluding probability
upper bounds through this pass. The theorem makes the rejection
\emph{stable under normalization}: no sequence of dualities can
manufacture or hide a rejected literal, so the check may run on the
surface formula once, before normalization, with the same verdict. A
guard that had to chase literals through rewrites would be a moving
target; a polarity count is not.
\end{bremark}

\begin{anchors}
normalizer, negation-freedom by type & \texttt{Form.nnf}, \texttt{NForm} (\texttt{BernsteinLogicAux.NNFDuality}; definitions --- no statement pins, their theorems below carry the warrant)\\
Proposition \ref{prop:rnnf-exact} & \texttt{denot\_nnf}, \texttt{denot\_elimNeg} (same module, hash-bound); \texttt{negLits\_elimNeg} (unpinned auxiliary)\\
Proposition \ref{prop:rnnf-polarity} & \texttt{negLits\_nnf}, \texttt{nnf\_no\_negLit\_iff} (same module, hash-bound)\\
Proposition \ref{prop:rnnf-sound} targets & the \rulename{QStreettBattery} and \rulename{QStreettCmplRequest} sections\\
\end{anchors}


\rulesection{55}{QAvoid}{alias rule}
\label{blcatrule-QAvoid}
\ruleblock{The rule (Steele metanotation)}
\sourcenote{as realized in the mechanized development; source row: \texttt{\texttt{QTSL-rules.toml}} --- \emph{QAvoid}. The metanotation statement is typeset by hand for notational quality; its mechanized ground truth is the Lean statement below, extracted from the sources at every build of this book and hash-bound. The rule \emph{roster} is machine-checked (every mechanized rule tag round-trips to exactly one catalogue row); the statement \emph{text} is editorial, checked by eye against the extracted Lean beside it.}
\begin{mathpar}
\inferrule*{
  \mathcal{M},\ x_0 \in R \vDash \mathsf{SProbGE}\bigl(p,\ \mathsf{PAlways}(\mathsf{SNot}\,\beta)\bigr)
}{
  \mathcal{M},\ x_0 \in R \vDash \mathsf{P}_{\ge p}[\mathsf{avoid}\,\beta]
}
\end{mathpar}

\ruleflavor{$\mathsf{avoid}\,\beta := \mathsf{G}\neg\beta$ --- the unfolding is the rule (the premise is the target judgment at the expansion), so the alias is sound by definition: the equation is definitional at the set level (\texttt{rfl}-grade), and the negation atom is ruled in by the \rulename{QNnfRewrite} swaps. The surface constructor builds its expansion rather than a new node kind — witnessed by the producer-side concrete-syntax canonicality battery; no rule sort exists and none is needed.}
\par\medskip
\rulename{QAvoid} is the avoidance sugar of the spec surface:
$\mathsf{avoid}\,\beta := \mathsf{G}\neg\beta$ --- avoidance is safety
of the negation, exact, the negation being ruled in by the
\rulename{QNnfRewrite} atom swaps. An alias unpacking to \emph{one}
application of an existing rule (the safety battery at the avoid set),
and the strongest-fidelity member of the sugar trio (its siblings:
\rulename{QStreettGF}, \rulename{QStreettFG}): the identity the
definition rides is machine-checked as a path-event identity in the
target rule\textquotesingle s concluded shape (module
\texttt{BernsteinLogicAux.QtslAliases}), \texttt{sugarAvoid\_denot},
proof \texttt{rfl}, axiom-free --- definitional exactness, certified as
such by the kernel. A producer, and a reader of specifications, may use
the sugar constructor freely: it builds its expansion rather than a new
node kind, so nothing downstream --- typing, compilation, checking ---
ever sees a sugar node, and no checker recomputation exists for this
rule at all. There is no soundness lemma here because the equation is a
definition; the witness that the surface really does build the
expansion remains the reference implementation\textquotesingle s
concrete-syntax canonicality test battery --- an implementation
property, not a semantic claim, so it stays a test rather than a
theorem. No profile enters; the entry carries no conditional flag.

\cataloguenote{The avoid sugar is also the compatibility layer to plain safety
specification: avoidance properties state directly against the safety
kernel with no automaton machinery, which is the degenerate (and most
common) case of the classical-surface translation the trio provides.
}
\ruleblock{Lean statement}
\declhead{BernsteinLogicAux.QtslAliases}{sugarAvoid\_denot}{\badge{BadgeGreen}{statement bound}}
\begin{leansrc}
/-- **`SugarBatch` (T10), first sugar**: `avoid β := G ¬β` — avoidance *is*
safety of the negation, definitionally (the proof is `rfl`; the negation is
ruled in by the `NNFDuality.lean` atom swaps (rule QNnfRewrite)). Exact: no direction is lost. -/
theorem sugarAvoid_denot {S : Type*} (B : Set S) :
    {ω : ℕ → S | ∀ n, ω n ∉ B} = {ω : ℕ → S | ∀ n, ω n ∈ Bᶜ}
\end{leansrc}
\ruleblock{Correspondence (Steele display $\leftrightarrow$ Lean)}
\badge{BadgeGreen}{reviewer-certified against the displayed texts}
\sourcenote{The dictionary below is a reviewer-certified statement of how the mechanized statements above interpret the displayed rule: display metavariable $\mapsto$ Lean binder, each display hypothesis $\mapsto$ the premise or side condition carrying it, the display conclusion $\mapsto$ the goal --- with any deliberate difference in width between display and mechanization named outright. Certification is recorded as a hash triple over (the displayed cell, the extracted statement texts, this dictionary\textquotesingle s own text); if any of the three changes, this block turns amber until a reviewer re-certifies. The correspondence claim itself remains editorial --- reviewer-grade, not machine-checked; its machine-checked successor types each soundness statement against the rule as a structured object, and supersedes this block per rule as it lands.}
The displayed rule transports the target judgment across the
defining equation $\mathrm{avoid}\,\beta := \mathsf{G}\neg\beta$
(premise: the expansion\textquotesingle s judgment; conclusion: the
sugared form). \texttt{sugarAvoid\_denot} certifies that equation and
is \texttt{rfl} --- definitional at the set level, the
strongest possible fidelity, with the negation-atom step ruled in by
the \rulename{QNnfRewrite} swaps as the docstring notes. Width note: a
per-path equivalence over arbitrary state types --- no measure enters,
so the identity transports under every profile\textquotesingle s path
measure by congruence; the remaining accompanying claim (the entry\textquotesingle s flavor
text: the certificate
surface \emph{builds} the expansion rather than minting a node kind) is
a property of the concrete-syntax canonicality battery --- a test, not
a theorem, deliberately outside this bound set, exactly as the roster
witness says.

\ruleblock{Lean verifier receipts}
\sourcenote{module \texttt{BernsteinLogicAux.QtslAliases}; bound build \texttt{570599779c29\ldots} (bound from receipt 5a60f4d01fd81ac5, 2026-07-21T05:39:02Z). A receipt records that named checking layers accepted a build: (module, compiled-artifact SHA-256, environment SHA-256, checker layers), minted when the checking gate ran; the store receipts attesting the bound build are listed with the layers each run credited. Nothing here asks the reader to trust the store: every layer is an independently runnable program over the quoted sources --- rebuild the module and re-run the checkers to reproduce the verdict. Receipts attesting superseded builds credit no layer and are omitted from this rendering (the count of omissions, if any, is noted below).}
\begin{receipttable}
\receiptrow{0a3819a978e31127}{2026-07-21T04:51:36Z}{build+decorrelated}{\badge{BadgeGreen}{current}}
\receiptrow{0bb1074df705d863}{2026-07-21T05:21:58Z}{build+decorrelated}{\badge{BadgeGreen}{current}}
\receiptrow{5a60f4d01fd81ac5}{2026-07-21T05:39:02Z}{build+replay+hardened+decorrelated}{\badge{BadgeGreen}{current}}
\end{receipttable}
\sourcenote{11 further receipt(s) in the store attest superseded builds of this module and are omitted from this rendering.}
\layerline{\layeritem{build}{lake build (elaboration)}{\badge{BadgeGreen}{current}}\layeritem{replay}{leanchecker (official kernel replay)}{\badge{BadgeGreen}{current}}\layeritem{hardened}{safe\_verify (hardened screen)}{\badge{BadgeGreen}{current}}\layeritem{decorrelated}{lean4export + nanoda (decorrelated kernel)}{\badge{BadgeGreen}{current}}}
\sourcenote{\badge{BadgeGray}{comparator: none} no fixture in the three-implementation comparator harness (compiled / kernel-interpreted / reference-interpreted) exercises this rule.}
\ruleblock{HOL4 statement}
\absentbox{OUT OF SCOPE BY DESIGN}{Out of scope for the HOL side \emph{by design}, in the strongest form
an alias can have: the sugar is a definitional constructor that builds
its expansion at the surface --- no sugar node exists in any
representation downstream of parsing, wire format included, so there is
no artefact for any checker port to be about. Checked against the HOL
sources: the \texttt{bern\_catalogue} theory carries no spec-surface
grammar.
}
\ruleblock{HOL verifier receipts}
\absentbox{NONE}{no HOL verification receipt exists for this rule.}
\ruleblock{Proof}

\sourcenote{Genre: alias rule (tier 3) --- a respelling template; the
shared template is Recipe~\ref{rec:sd-template} of the
\rulename{QStreettDegenerate} section. This rule is the degenerate case
of even that template: the unpack is a \emph{definition}, so the five
artefacts collapse to the definitional equation, its reading, and the
target section.}

\begin{brecipe}[the filled-in application]\label{rec:qaa-filled}
\[
\mathsf{avoid}\,\beta \;:=\; \mathsf{G}\neg\beta
\]
The concluding judgment discharges through one existing rule:
avoidance through the safety battery (\rulename{QSafetyBattery}) at the
avoid set $[\![\beta]\!]$, with $\neg\beta$ read by the semantic
negation clause $[\![\neg\beta]\!] = S \smallsetminus [\![\beta]\!]$
(the \rulename{QNnfRewrite} atom dualities supply the syntactic normal
form --- the atom swap alone when $\beta$ is an atom, De Morgan
propagation when it is compound).
\end{brecipe}

\begin{brecipe}[payload transport]\label{rec:qaa-transport}
The expansion $\mathsf{G}\neg\beta$ carries the safety kit at
$U := [\![\beta]\!]$, verbatim. No unpack function is owed here,
because none will ever exist: the constructor \emph{is} its expansion,
so there is no sugar payload to transport from --- the strongest
possible form of producers-elaborate-before-emission.
\end{brecipe}

\begin{bprop}[the exactness identity]\label{prop:qaa-exact}
A path always avoids $[\![\beta]\!]$ iff it always satisfies
$\neg\beta$ --- definitional, by the negation clause
$[\![\neg\beta]\!] = S \smallsetminus [\![\beta]\!]$, with
\rulename{QNnfRewrite} normalizing the syntax where needed.
\end{bprop}
\begin{proof}
Definitional (the bound Lean identity is \texttt{rfl}, axiom-free ---
the kernel certifies the two sides as one term).
\end{proof}

\begin{bprop}[soundness corollary]\label{prop:qaa-sound}
The sugared judgment is concluded by the target rule's theorem, with
no additional obligation: the sugar constructor builds the expansion,
so the sugared judgment \emph{is} the target judgment.
\end{bprop}
\begin{proof}
Definitional unfolding plus the target section's theorem. The witness
that the surface implementation honours the definition --- that the
constructor builds its expansion rather than introducing a new node
kind --- is the reference implementation's concrete-syntax canonicality
test battery; no checker-facing artefact depends on it, since no
downstream representation contains a sugar node.
\end{proof}

\begin{anchors}
target & the \rulename{QSafetyBattery} section; negation via the \rulename{QNnfRewrite} atom swaps\\
the identity, event-level & \texttt{sugarAvoid\_denot} (\texttt{rfl}, axiom-free; bound in the entry)\\
siblings (the sugar trio) & \rulename{QStreettGF}, \rulename{QStreettFG}\\
constructor builds expansion & the reference implementation\textquotesingle s concrete-syntax canonicality battery (documentary; no wire artefact exists to check)\\
\end{anchors}


\rulesection{56}{QStreettGF}{alias rule}
\label{blcatrule-QStreettGF}
\ruleblock{The rule (Steele metanotation)}
\sourcenote{as realized in the mechanized development; source row: \texttt{\texttt{QTSL-rules.toml}} --- \emph{QStreettGF}. The metanotation statement is typeset by hand for notational quality; its mechanized ground truth is the Lean statement below, extracted from the sources at every build of this book and hash-bound. The rule \emph{roster} is machine-checked (every mechanized rule tag round-trips to exactly one catalogue row); the statement \emph{text} is editorial, checked by eye against the extracted Lean beside it.}
\begin{mathpar}
\inferrule*{
  \mathcal{M},\ x_0 \in R \vDash \mathsf{SProbGE}\bigl(p,\ \mathsf{PDwellStreett}(\mathsf{PairBase}(\mathsf{STrue}, \varphi))\bigr) \\
  \text{paths total over the state space}
}{
  \mathcal{M},\ x_0 \in R \vDash \mathsf{SProbGE}(p,\ \mathsf{GF}\,\varphi)
}
\end{mathpar}

\ruleflavor{$\mathsf{GF}\varphi := \mathsf{Streett}((\top, \varphi))$ --- recurrence is the one-pair Streett condition at the trivially-true request, \rulename{QStreettDegenerate}\textquotesingle s concluded shape, so the sugar lands on that rule\textquotesingle s bound statement without adjustment. The totality premise is the bound chart, exactly the certified statement\textquotesingle s scope. Killed-process classes are the named unbound residue: the identity extends whenever the cemetery is adjoined inside the region vocabulary \emph{and} killed paths are continued as total paths at \(\dagger\) (\(\top = V \cup \{\dagger\}\)) --- that \(\dagger\)-extended restatement is not yet part of the certified semantics, so it stays here rather than in a premise; with the cemetery outside the region vocabulary the rewrite direction over-claims, and that class is kept closed by the batteries\textquotesingle{} full-mass gates. The surface constructor builds its expansion rather than a new node kind — witnessed by the producer-side concrete-syntax canonicality battery; no rule sort exists and none is needed.}
\par\medskip
\rulename{QStreettGF} is the recurrence sugar of the spec surface:
$\mathsf{GF}\varphi := \mathsf{Streett}((\top, \varphi))$ --- recurrence
is the degenerate Streett pair, the \rulename{QStreettDegenerate}
section\textquotesingle s concluded shape, so the sugar lands on that
rule\textquotesingle s bound statement without adjustment. An alias
unpacking to \emph{one} application of an existing rule; with its
persistence twin \rulename{QStreettFG} it is the compatibility layer to
the classical temporal-logic surface ---
$\mathsf{GF}$/$\mathsf{FG}$ specifications translate to the
Streett-primitive kernel without automaton products, the practical
payoff of B\"uchi acceptance being a one-pair Streett condition. The
identity the definition rides is machine-checked (module
\texttt{BernsteinLogicAux.QtslAliases}): \texttt{sugarGF\_denot} ---
the per-pair Streett event at the trivially-true request \emph{is} the
recurrence event.

Scope, stated in the open because this member alone has one: the
mechanized identity lives on the total path type, so it consumes the
no-cemetery chart --- exact on total-path state spaces and on
killed-process classes with the cemetery adjoined inside the region
vocabulary ($\top = V \cup \{\dagger\}$); with the cemetery outside the
region vocabulary the rewrite direction over-claims, and that class is
kept closed by the batteries\textquotesingle{} full-mass gates. The
named repair (the entry\textquotesingle s residue): restate the
identity on the $\dagger$-extended state type, matching the promotion
lemma. The surface constructor builds its expansion rather than a new
node kind; the witness for that remains the reference
implementation\textquotesingle s concrete-syntax canonicality test
battery --- an implementation property, a test rather than a theorem.

\cataloguenote{Why the factoring put the honesty here: as a batch row, one PROVEN
badge covered three warrants of different semantic depths ---
\texttt{rfl}-grade, chart-consuming, and pure-logic --- and the
chart-consuming member\textquotesingle s scope restriction was a
memo-level footnote. As its own row it is display-level text: the
exactness chart is part of what this rule \emph{says}.
}
\ruleblock{Lean statement}
\declhead{BernsteinLogicAux.QtslAliases}{sugarGF\_denot}{\badge{BadgeGreen}{statement bound}}
\begin{leansrc}
/-- **`SugarBatch` (T10), second sugar**: `GF φ := Streett((⊤, φ))` — the
per-pair Streett event at the trivially-true request (every state requests,
and along `atTop` over `ℕ` the request is frequent) *is* the recurrence
event. The expansion's target is `QStreettDegenerate` (T9a), whose bound
statement `sound_QStreettDegenerate` concludes exactly the right-hand
event. -/
theorem sugarGF_denot {S : Type*} (G : Set S) (ω : ℕ → S) :
    ((∃ᶠ n in atTop, ω n ∈ (Set.univ : Set S)) → ∃ᶠ n in atTop, ω n ∈ G)
      ↔ ∃ᶠ n in atTop, ω n ∈ G
\end{leansrc}
\ruleblock{Correspondence (Steele display $\leftrightarrow$ Lean)}
\badge{BadgeGreen}{reviewer-certified against the displayed texts}
\sourcenote{The dictionary below is a reviewer-certified statement of how the mechanized statements above interpret the displayed rule: display metavariable $\mapsto$ Lean binder, each display hypothesis $\mapsto$ the premise or side condition carrying it, the display conclusion $\mapsto$ the goal --- with any deliberate difference in width between display and mechanization named outright. Certification is recorded as a hash triple over (the displayed cell, the extracted statement texts, this dictionary\textquotesingle s own text); if any of the three changes, this block turns amber until a reviewer re-certifies. The correspondence claim itself remains editorial --- reviewer-grade, not machine-checked; its machine-checked successor types each soundness statement against the rule as a structured object, and supersedes this block per rule as it lands.}
\texttt{sugarGF\_denot} ($\mathsf{GF}\varphi :=
\mathsf{Streett}((\top, \varphi))$) certifies the pair event at
request \texttt{Set.univ} $\leftrightarrow$ the bare recurrence event
--- exactly \texttt{sound\_QStreettDegenerate}\textquotesingle s
concluded form, so the sugar lands on
\rulename{QStreettDegenerate}\textquotesingle s bound statement without
adjustment. Fidelity caveat, bound honestly: the Lean statement lives
on the total path type, which IS the displayed rule\textquotesingle s
totality premise (``paths total over the state space'') --- the premise
and the bound scope now coincide exactly. The flavor
text\textquotesingle s $\dagger$-extended case (cemetery adjoined
inside the vocabulary with killed paths continued as total paths at
$\dagger$: exact; cemetery outside: the rewrite over-claims) is the
named unbound residue the bound statement does not and
cannot witness; it is carried by the roster note\textquotesingle s
residue line, and the $\dagger$-extended restatement named there would
move it into this bound set. No measure enters; the identity transports
under every profile\textquotesingle s path measure by congruence. The
builds-expansions claim rides the canonicality battery, deliberately
outside this bound set, exactly as the roster witness says.

\ruleblock{Lean verifier receipts}
\sourcenote{module \texttt{BernsteinLogicAux.QtslAliases}; bound build \texttt{570599779c29\ldots} (bound from receipt 5a60f4d01fd81ac5, 2026-07-21T05:39:02Z). A receipt records that named checking layers accepted a build: (module, compiled-artifact SHA-256, environment SHA-256, checker layers), minted when the checking gate ran; the store receipts attesting the bound build are listed with the layers each run credited. Nothing here asks the reader to trust the store: every layer is an independently runnable program over the quoted sources --- rebuild the module and re-run the checkers to reproduce the verdict. Receipts attesting superseded builds credit no layer and are omitted from this rendering (the count of omissions, if any, is noted below).}
\begin{receipttable}
\receiptrow{0a3819a978e31127}{2026-07-21T04:51:36Z}{build+decorrelated}{\badge{BadgeGreen}{current}}
\receiptrow{0bb1074df705d863}{2026-07-21T05:21:58Z}{build+decorrelated}{\badge{BadgeGreen}{current}}
\receiptrow{5a60f4d01fd81ac5}{2026-07-21T05:39:02Z}{build+replay+hardened+decorrelated}{\badge{BadgeGreen}{current}}
\end{receipttable}
\sourcenote{11 further receipt(s) in the store attest superseded builds of this module and are omitted from this rendering.}
\layerline{\layeritem{build}{lake build (elaboration)}{\badge{BadgeGreen}{current}}\layeritem{replay}{leanchecker (official kernel replay)}{\badge{BadgeGreen}{current}}\layeritem{hardened}{safe\_verify (hardened screen)}{\badge{BadgeGreen}{current}}\layeritem{decorrelated}{lean4export + nanoda (decorrelated kernel)}{\badge{BadgeGreen}{current}}}
\sourcenote{\badge{BadgeGray}{comparator: none} no fixture in the three-implementation comparator harness (compiled / kernel-interpreted / reference-interpreted) exercises this rule.}
\ruleblock{HOL4 statement}
\absentbox{OUT OF SCOPE BY DESIGN}{Out of scope for the HOL side \emph{by design}: the sugar is a
definitional constructor that builds its expansion at the surface ---
no sugar node exists in any representation downstream of parsing, wire
format included, so there is no artefact for any checker port to be
about. Checked against the HOL sources: the \texttt{bern\_catalogue}
theory carries no spec-surface grammar. (The chart restriction is a
property of the Lean statement\textquotesingle s ambient path type, not
of any wire artefact --- there is nothing HOL-side for it to attach to.)
}
\ruleblock{HOL verifier receipts}
\absentbox{NONE}{no HOL verification receipt exists for this rule.}
\ruleblock{Proof}

\sourcenote{Genre: alias rule (tier 3) --- a respelling template; the
shared template is Recipe~\ref{rec:sd-template} of the
\rulename{QStreettDegenerate} section. The unpack is a \emph{definition},
so the five artefacts collapse to the definitional equation, its
reading, its exactness chart, and the target section.}

\begin{brecipe}[the filled-in application]\label{rec:qgf-filled}
\[
\mathsf{GF}\varphi \;:=\; \mathsf{Streett}((\top, \varphi))
\]
The concluding judgment discharges through one existing rule:
recurrence through the degenerate bundle
(\rulename{QStreettDegenerate}).
\end{brecipe}

\begin{brecipe}[payload transport]\label{rec:qgf-transport}
The expansion $\mathsf{Streett}((\top, \varphi))$ carries the
degenerate kit, verbatim. No unpack function is owed here, because none
will ever exist: the constructor \emph{is} its expansion, so there is
no sugar payload to transport from --- the strongest possible form of
producers-elaborate-before-emission.
\end{brecipe}

\begin{bprop}[the exactness identity, and its chart]\label{prop:qgf-exact}
The pair $(\top, \varphi)$ holds iff $\varphi$-states occur infinitely
often (the degeneration identity, proved as the exactness identity,
Proposition~\ref{prop:sd-exact}, of the \rulename{QStreettDegenerate}
section). As mechanized, the identity is stated on the total path type;
it therefore holds on total-path state spaces, and on killed-process
classes whenever the cemetery is adjoined inside the region vocabulary
($\top = V \cup \{\dagger\}$). With the cemetery outside the region
vocabulary the rewrite direction over-claims --- a killed path visits no
state cofinally, yet the vocabulary cannot say so --- and that class is
kept closed by the batteries' full-mass gates. The named repair is to
restate the identity on the $\dagger$-extended state type, matching the
promotion lemma; until then the chart restriction is carried in the
rule's own display, in the open.
\end{bprop}
\begin{proof}
The identity: by the degeneration argument of the target section,
mechanized as the bound Lean statement. The chart restriction: read off
the bound statement's ambient path type, as displayed.
\end{proof}

\begin{bprop}[soundness corollary]\label{prop:qgf-sound}
On the stated chart, the sugared judgment is concluded by the target
rule's theorem, with no additional obligation: the sugar constructor
builds the expansion, so the sugared judgment \emph{is} the target
judgment.
\end{bprop}
\begin{proof}
Definitional unfolding plus the target section's theorem, on the chart
of Proposition~\ref{prop:qgf-exact}. The witness that the surface
implementation honours the definition --- that the constructor builds
its expansion rather than introducing a new node kind --- is the
reference implementation's concrete-syntax canonicality test battery;
no checker-facing artefact depends on it, since no downstream
representation contains a sugar node.
\end{proof}

\begin{anchors}
target & the \rulename{QStreettDegenerate} section (degeneration identity proved there)\\
the identity, event-level & \texttt{sugarGF\_denot} (bound in the entry; total path type = the chart)\\
the chart's named repair & the $\dagger$-extended restatement, $\top = V \cup \{\dagger\}$ (the promotion lemma's form)\\
twin & \rulename{QStreettFG} (the safety-polarized half; consumes no chart)\\
constructor builds expansion & the reference implementation\textquotesingle s concrete-syntax canonicality battery (documentary; no wire artefact exists to check)\\
\end{anchors}


\rulesection{57}{QStreettFG}{alias rule}
\label{blcatrule-QStreettFG}
\ruleblock{The rule (Steele metanotation)}
\sourcenote{as realized in the mechanized development; source row: \texttt{\texttt{QTSL-rules.toml}} --- \emph{QStreettFG}. The metanotation statement is typeset by hand for notational quality; its mechanized ground truth is the Lean statement below, extracted from the sources at every build of this book and hash-bound. The rule \emph{roster} is machine-checked (every mechanized rule tag round-trips to exactly one catalogue row); the statement \emph{text} is editorial, checked by eye against the extracted Lean beside it.}
\begin{mathpar}
\inferrule*{
  \mathcal{M},\ x_0 \in R \vDash \mathsf{SProbGE}\bigl(p,\ \mathsf{PDwellStreett}(\mathsf{PairBase}(\beta, \mathsf{SFalse}))\bigr)
}{
  \mathcal{M},\ x_0 \in R \vDash \mathsf{SProbGE}(p,\ \mathsf{FG}\,\neg\beta)
}
\end{mathpar}

\ruleflavor{$\mathsf{FG}\neg\beta := \mathsf{Streett}((\beta, \bot))$ --- eventual permanent avoidance is the one-pair Streett condition that forbids infinitely many \(\beta\)-visits: with an unsatisfiable response, the pair condition says exactly that \(\beta\)-states occur only finitely often. Pure per-path logic — the safety-polarized half of the recurrence/persistence pair, consuming no conservativity. The rule\textquotesingle s honest scope is this polarized form (\(\mathsf{FG}\) applied to \(\neg\beta\)): general persistence \(\mathsf{FG}\varphi\) is not this rule. The surface constructor builds its expansion rather than a new node kind — witnessed by the producer-side concrete-syntax canonicality battery; no rule sort exists and none is needed.}
\par\medskip
\rulename{QStreettFG} is the persistence-of-avoidance sugar of the spec
surface: $\mathsf{FG}\neg\beta := \mathsf{Streett}((\beta, \bot))$ ---
eventual permanent avoidance is the one-pair Streett condition that
forbids infinitely many $\beta$-visits: with an unsatisfiable response,
the pair condition says exactly that $\beta$-states occur only finitely
often. An alias unpacking to \emph{one} application of an existing rule
(one \rulename{QStreettBattery} node at the pair $(\beta, \bot)$); the
safety-polarized half of the recurrence/persistence pair --- its twin
is \rulename{QStreettGF}, and the two land at different points of the
same Streett rule. The honest scope is the polarized form
($\mathsf{FG}$ applied to $\neg\beta$): general persistence
$\mathsf{FG}\varphi$ is not this rule. The identity the definition
rides is machine-checked (module
\texttt{BernsteinLogicAux.QtslAliases}): \texttt{sugarFG\_denot} ---
the per-pair Streett event at an unsatisfiable response \emph{is}
eventual permanent avoidance, the finitely-often reading; pure per-path
logic, no conservativity and no chart hypothesis consumed. The surface
constructor builds its expansion rather than a new node kind; the
witness for that remains the reference implementation\textquotesingle s
concrete-syntax canonicality test battery --- an implementation
property, a test rather than a theorem. No profile enters; the entry
carries no conditional flag.

\cataloguenote{The two dual sugars land at different points of the same rule:
$\mathsf{GF}\varphi$ requests always and so erases the request
machinery, while $\mathsf{FG}\neg\beta$ makes $\beta$ itself the
request and forbids response --- the ranking certificate must then pay
out at every $\beta$-visit with nothing ever refunding it, which is
exactly a proof that $\beta$-visits are finitely many. That both
directions are single nodes of one rule --- where automaton-primitive
designs pay a product construction for one of them --- is the practical
content of the Streett-primitive choice.
}
\ruleblock{Lean statement}
\declhead{BernsteinLogicAux.QtslAliases}{sugarFG\_denot}{\badge{BadgeGreen}{statement bound}}
\begin{leansrc}
/-- **`SugarBatch` (T10), third sugar**: `FG ¬β := Streett((β, ⊥))` — the
per-pair Streett event at an unsatisfiable response says the request region
is visited at most finitely often, which over `ℕ` *is* eventual permanent
avoidance. The pair condition forbids infinitely many `β`-visits: with
response `⊥`, "infinitely often `β` implies infinitely often `False`" is
exactly "eventually always `¬β`". -/
theorem sugarFG_denot {S : Type*} (B : Set S) (ω : ℕ → S) :
    ((∃ᶠ n in atTop, ω n ∈ B) → ∃ᶠ n in atTop, ω n ∈ (∅ : Set S))
      ↔ ∀ᶠ n in atTop, ω n ∉ B
\end{leansrc}
\ruleblock{Correspondence (Steele display $\leftrightarrow$ Lean)}
\badge{BadgeGreen}{reviewer-certified against the displayed texts}
\sourcenote{The dictionary below is a reviewer-certified statement of how the mechanized statements above interpret the displayed rule: display metavariable $\mapsto$ Lean binder, each display hypothesis $\mapsto$ the premise or side condition carrying it, the display conclusion $\mapsto$ the goal --- with any deliberate difference in width between display and mechanization named outright. Certification is recorded as a hash triple over (the displayed cell, the extracted statement texts, this dictionary\textquotesingle s own text); if any of the three changes, this block turns amber until a reviewer re-certifies. The correspondence claim itself remains editorial --- reviewer-grade, not machine-checked; its machine-checked successor types each soundness statement against the rule as a structured object, and supersedes this block per rule as it lands.}
\texttt{sugarFG\_denot} ($\mathsf{FG}\neg\beta :=
\mathsf{Streett}((\beta, \bot))$) certifies the pair event at response
$\emptyset$ $\leftrightarrow$ eventual permanent avoidance
(\texttt{∀ᶠ}), the finitely-often reading --- stated, like its
siblings, in the target rule\textquotesingle s concluded shape, which
is what makes it consumable as a rewrite rather than a mere fact. Width
note: a per-path equivalence over arbitrary state types --- no measure
enters, so the identity transports under every profile\textquotesingle s
path measure by congruence; the remaining accompanying claim (the
entry\textquotesingle s flavor text: the certificate
surface \emph{builds} the expansion rather than minting a node kind) is
a property of the concrete-syntax canonicality battery --- a test, not
a theorem, deliberately outside this bound set, exactly as the roster
witness says.

\ruleblock{Lean verifier receipts}
\sourcenote{module \texttt{BernsteinLogicAux.QtslAliases}; bound build \texttt{570599779c29\ldots} (bound from receipt 5a60f4d01fd81ac5, 2026-07-21T05:39:02Z). A receipt records that named checking layers accepted a build: (module, compiled-artifact SHA-256, environment SHA-256, checker layers), minted when the checking gate ran; the store receipts attesting the bound build are listed with the layers each run credited. Nothing here asks the reader to trust the store: every layer is an independently runnable program over the quoted sources --- rebuild the module and re-run the checkers to reproduce the verdict. Receipts attesting superseded builds credit no layer and are omitted from this rendering (the count of omissions, if any, is noted below).}
\begin{receipttable}
\receiptrow{0a3819a978e31127}{2026-07-21T04:51:36Z}{build+decorrelated}{\badge{BadgeGreen}{current}}
\receiptrow{0bb1074df705d863}{2026-07-21T05:21:58Z}{build+decorrelated}{\badge{BadgeGreen}{current}}
\receiptrow{5a60f4d01fd81ac5}{2026-07-21T05:39:02Z}{build+replay+hardened+decorrelated}{\badge{BadgeGreen}{current}}
\end{receipttable}
\sourcenote{11 further receipt(s) in the store attest superseded builds of this module and are omitted from this rendering.}
\layerline{\layeritem{build}{lake build (elaboration)}{\badge{BadgeGreen}{current}}\layeritem{replay}{leanchecker (official kernel replay)}{\badge{BadgeGreen}{current}}\layeritem{hardened}{safe\_verify (hardened screen)}{\badge{BadgeGreen}{current}}\layeritem{decorrelated}{lean4export + nanoda (decorrelated kernel)}{\badge{BadgeGreen}{current}}}
\sourcenote{\badge{BadgeGray}{comparator: none} no fixture in the three-implementation comparator harness (compiled / kernel-interpreted / reference-interpreted) exercises this rule.}
\ruleblock{HOL4 statement}
\absentbox{OUT OF SCOPE BY DESIGN}{Out of scope for the HOL side \emph{by design}: the sugar is a
definitional constructor that builds its expansion at the surface ---
no sugar node exists in any representation downstream of parsing, wire
format included, so there is no artefact for any checker port to be
about. Checked against the HOL sources: the \texttt{bern\_catalogue}
theory carries no spec-surface grammar.
}
\ruleblock{HOL verifier receipts}
\absentbox{NONE}{no HOL verification receipt exists for this rule.}
\ruleblock{Proof}

\sourcenote{Genre: alias rule (tier 3) --- a respelling template; the
shared template is Recipe~\ref{rec:sd-template} of the
\rulename{QStreettDegenerate} section. The unpack is a \emph{definition},
so the five artefacts collapse to the definitional equation, its
reading, and the target section.}

\begin{brecipe}[the filled-in application]\label{rec:qfg-filled}
\[
\mathsf{FG}\neg\beta \;:=\; \mathsf{Streett}((\beta, \bot))
\]
The concluding judgment discharges through one existing rule: eventual
permanent avoidance through one \rulename{QStreettBattery} node at the
pair $(\beta, \bot)$. The honest scope is this polarized form
($\mathsf{FG}$ applied to $\neg\beta$): general persistence
$\mathsf{FG}\varphi$ is not this rule.
\end{brecipe}

\begin{brecipe}[payload transport]\label{rec:qfg-transport}
The expansion $\mathsf{Streett}((\beta, \bot))$ carries one bundle
pair: requests $[\![\beta]\!]$, response empty. No unpack function is
owed here, because none will ever exist: the constructor \emph{is} its
expansion, so there is no sugar payload to transport from --- the
strongest possible form of producers-elaborate-before-emission.
\end{brecipe}

\begin{bprop}[the exactness identity]\label{prop:qfg-exact}
The pair $(\beta, \bot)$ holds iff $\beta$-states occur only finitely
often, i.e.\ iff eventually the path permanently satisfies $\neg\beta$:
\[
\Bigl(
\bigl(\exists^\infty n,\ \omega_n \in [\![\beta]\!]\bigr) \Rightarrow
\bigl(\exists^\infty n,\ \omega_n \in \varnothing\bigr)
\Bigr)
\;\Longleftrightarrow\;
\neg\bigl(\exists^\infty n,\ \omega_n \in [\![\beta]\!]\bigr) ,
\]
since the consequent is false outright --- two lines: an implication
into falsehood is the negation of its hypothesis, and finitely many
$\beta$-visits along a discrete path is precisely a tail of
$\neg\beta$-states. Pure per-path logic; no conservativity and no chart
hypothesis is consumed.
\end{bprop}
\begin{proof}
As displayed; the two lines are shown.
\end{proof}

\begin{bprop}[soundness corollary]\label{prop:qfg-sound}
The sugared judgment is concluded by the target rule's theorem, with
no additional obligation: the sugar constructor builds the expansion,
so the sugared judgment \emph{is} the target judgment.
\end{bprop}
\begin{proof}
Definitional unfolding plus the target section's theorem. The witness
that the surface implementation honours the definition --- that the
constructor builds its expansion rather than introducing a new node
kind --- is the reference implementation's concrete-syntax canonicality
test battery; no checker-facing artefact depends on it, since no
downstream representation contains a sugar node.
\end{proof}

\begin{bremark}[why $\mathsf{FG}$ costs a bundle and $\mathsf{GF}$ costs
only its degenerate case]
The two dual sugars land at different points of the same rule.
$\mathsf{GF}\varphi$ (the twin, \rulename{QStreettGF}) requests always
and so erases the request machinery; $\mathsf{FG}\neg\beta$ makes
$\beta$ itself the request and forbids response --- the ranking
certificate must then pay out at every $\beta$-visit with nothing ever
refunding it, which is exactly a proof that $\beta$-visits are finitely
many. That both directions are single nodes of one rule --- where
automaton-primitive designs pay a product construction for one of them
--- is the practical content of the Streett-primitive choice (see the
\rulename{QStreettBattery} section's remark).
\end{bremark}

\begin{anchors}
target & the \rulename{QStreettBattery} section, one pair $(\beta, \bot)$\\
the identity, event-level & \texttt{sugarFG\_denot} (bound in the entry)\\
twin & \rulename{QStreettGF} (recurrence; the member that consumes a chart)\\
constructor builds expansion & the reference implementation\textquotesingle s concrete-syntax canonicality battery (documentary; no wire artefact exists to check)\\
\end{anchors}


\rulesection{58}{QNegFAlias}{alias for \textcolor{DiffBlue}{\rulename{QSafetyBattery}}}
\label{blcatrule-QNegFAlias}
\ruleblock{The rule (Steele metanotation)}
\sourcenote{as realized in the mechanized development; source row: \texttt{\texttt{QTSL-rules.toml}} --- \emph{QNegFAlias}. The metanotation statement is typeset by hand for notational quality; its mechanized ground truth is the Lean statement below, extracted from the sources at every build of this book and hash-bound. The rule \emph{roster} is machine-checked (every mechanized rule tag round-trips to exactly one catalogue row); the statement \emph{text} is editorial, checked by eye against the extracted Lean beside it.}
\begin{mathpar}
\inferrule*{
  \mathrm{child} : \mathsf{QRegionSat}\bigl(\mathsf{SProbGE}(p, \mathsf{PAlways}(\mathsf{SNot}\,\varphi)),\ R\bigr)
}{
  \mathcal{M},\ x_0 \in R \vDash \mathsf{SProbGE}(p, \mathsf{PNotPath}(\mathsf{PDeadline}(T, \varphi, \mathrm{false})))
}
\end{mathpar}

\par\medskip
\rulename{QNegFAlias} discharges a negated deadline: to certify, with
probability at least $p$, that $\varphi$ does \emph{not} occur within
horizon $T$ ($\neg\mathsf{F}_{\le T}\,\varphi$), it suffices to certify
the \emph{unbounded} safety claim --- $\varphi$ never occurs at all ---
over the same region, and the alias rewrites the one judgment into the
other: the child is a satisfaction judgment for
always-not-$\varphi$, the conclusion the negated-deadline judgment,
region verbatim. It is an alias to the safety battery
(\rulename{QSafetyBattery}): one node to one node, and the entire
mathematical content is a two-line denotation inclusion --- a path on
which $\varphi$ never holds certainly does not reach $\varphi$ within
$T$, and not-reaching-within-$T$ \emph{is} the bounded-always of the
negation. The rewrite is choice-free (no producer parameter beyond the
horizon already present in the goal), and the checker recomputes
nothing new: the child is an ordinary safety-battery conclusion.

The rule is deliberately lossy in one direction --- the unbounded claim
is strictly stronger than the bounded one --- which is the honest price
of reusing the safety kernel with no horizon bookkeeping. A producer
who needs the sharp bounded claim (because $\varphi$ \emph{does} occur
eventually, only late) uses the deadline machinery instead; this alias
serves the common case where the avoidance is permanent.

Mechanization state, plainly: the bound statements below are the two
lines, as priced. \texttt{negF\_denot\_subset} (module
\texttt{BernsteinLogicAux.QtslAliases}) is the denotation inclusion
--- a path on which $\varphi$ never holds does not reach $\varphi$
within $T$ --- pure path-set algebra, axiom-free (the kernel certifies
it depends on no axioms at all); and \texttt{sound\_QNegFAlias} is the
alias\textquotesingle s transport: a lower bound on the unbounded
safety event (the target\textquotesingle s conclusion shape, with
$\mathrm{sat} := \neg\varphi$) is a lower bound on the
negated-deadline event, for \emph{any} measure on path space ---
profile-generic, so the discrete-kernel and GSHS records read it
alike, and the entry carries no conditional flag.

\cataloguenote{Two flags travel with the row, both from the source catalogue. Tier
(flagged, not forced): alias as sorted, but if the folded denotation
inclusion must itself be a rule application, the tier moves to
admissible --- machinery unchanged either way. And the mirror is a
deliberate hole: there is no dual rule discharging a negated
\emph{always} by a bounded operator, because the conclusion it would
need ($\mathsf{F}\varphi$ unbounded) is structurally unrepresentable in
the spec grammar --- the unbounded-eventually constructor was excluded
by testable-negative design, unbounded eventuality surviving only as a
producer-choice provenance flag on deadlines. The pair must never be
ported as one duality entry: one side is a statable choice-free alias,
the other hides a horizon choice and a metatheory line inside the
deadline battery.
}
\ruleblock{Lean statement}
\declhead{BernsteinLogicAux.QtslAliases}{negF\_denot\_subset}{\badge{BadgeGreen}{statement bound}}
\begin{leansrc}
/-- **The denotation inclusion** (QTSL-rules.toml row `QNegFAlias` —
"the denotation inclusion `⟦G¬φ⟧ ⊆ ⟦G_{≤T}¬φ⟧ = ⟦¬F_{≤T}φ⟧` folds into the
rewrite"; the catalogue's "two lines"): a path on which `φ` never holds
certainly does not reach `φ` within `T` — and not-reaching-within-`T` *is*
the negated deadline event. Pure path-set algebra; no measure, no
measurability. -/
theorem negF_denot_subset {S : Type*} (sat : S → Prop) (T : ℕ) :
    {ω : ℕ → S | ∀ n, ¬ sat (ω n)} ⊆ {ω : ℕ → S | ¬ ∃ k ≤ T, sat (ω k)}
\end{leansrc}
\declhead{BernsteinLogicAux.QtslAliases}{sound\_QNegFAlias}{\badge{BadgeGreen}{statement bound}}
\begin{leansrc}
/-- **`sound_QNegFAlias`** (rule row `QNegFAlias`): the alias's
transport — a lower bound on the probability of the unbounded safety event
(`QSafetyBattery`'s conclusion shape, e.g. `sound_QSafetyBattery` /
`sound_QSafetyBattery_staged` with `sat := ¬φ`) is a lower bound on the
probability of the negated-deadline event. Profile-generic (any measure on
path space — discrete-kernel and GSHS records alike), by monotonicity
through `negF_denot_subset`. The rewrite is choice-free: no producer
parameter beyond the horizon `T` already present in the goal. -/
theorem sound_QNegFAlias {S : Type*} [MeasurableSpace S]
    (μ : Measure (ℕ → S)) {sat : S → Prop} {b : ENNReal} {T : ℕ}
    (h : b ≤ μ {ω | ∀ n, ¬ sat (ω n)}) :
    b ≤ μ {ω | ¬ ∃ k ≤ T, sat (ω k)}
\end{leansrc}
\ruleblock{Correspondence (Steele display $\leftrightarrow$ Lean)}
\badge{BadgeGreen}{reviewer-certified against the displayed texts}
\sourcenote{The dictionary below is a reviewer-certified statement of how the mechanized statements above interpret the displayed rule: display metavariable $\mapsto$ Lean binder, each display hypothesis $\mapsto$ the premise or side condition carrying it, the display conclusion $\mapsto$ the goal --- with any deliberate difference in width between display and mechanization named outright. Certification is recorded as a hash triple over (the displayed cell, the extracted statement texts, this dictionary\textquotesingle s own text); if any of the three changes, this block turns amber until a reviewer re-certifies. The correspondence claim itself remains editorial --- reviewer-grade, not machine-checked; its machine-checked successor types each soundness statement against the rule as a structured object, and supersedes this block per rule as it lands.}
The displayed rewrite --- child
$\mathsf{QRegionSat}(\mathsf{SProbGE}(p, \mathsf{PAlways}(\mathsf{SNot}\,
\varphi)), R)$ concluding
$\mathsf{SProbGE}(p, \mathsf{PNotPath}(\mathsf{PDeadline}(T, \varphi,
\mathrm{false})))$ --- $\mapsto$ the two-statement chain:
\texttt{negF\_denot\_subset} (the denotation inclusion \texttt{\{ω |
∀ n, ¬ sat (ω n)\} ⊆ \{ω | ¬ ∃ k ≤ T, sat (ω k)\}} --- pure path-set
algebra, axiom-free, no measure, no measurability) and
\texttt{sound\_QNegFAlias}, its transport: a lower bound \texttt{b ≤
μ} on the unbounded-safety event is a lower bound on the
negated-deadline event, for \emph{any} measure on path space ---
profile-generic by binder shape (\texttt{μ : Measure (ℕ → S)}
arbitrary), the \rulename{QUnionBound} pattern; discrete-kernel and
GSHS records alike are instances. The
display\textquotesingle s choice-freeness is a binder fact: no
producer parameter appears beyond the horizon \texttt{T} already in
the goal. Width note: the child\textquotesingle s
$\mathsf{SProbGE}$ shape enters only as the abstract bound \texttt{b ≤
μ \{…\}} --- the certified transport is wider than the displayed
probabilistic judgment (any \texttt{ENNReal} bound transports), safe
direction; paths are step-indexed (\texttt{ℕ → S}), the deadline
read discretely, as at \rulename{TaylorBound}.

\ruleblock{Lean verifier receipts}
\sourcenote{module \texttt{BernsteinLogicAux.QtslAliases}; bound build \texttt{570599779c29\ldots} (bound from receipt 5a60f4d01fd81ac5, 2026-07-21T05:39:02Z). A receipt records that named checking layers accepted a build: (module, compiled-artifact SHA-256, environment SHA-256, checker layers), minted when the checking gate ran; the store receipts attesting the bound build are listed with the layers each run credited. Nothing here asks the reader to trust the store: every layer is an independently runnable program over the quoted sources --- rebuild the module and re-run the checkers to reproduce the verdict. Receipts attesting superseded builds credit no layer and are omitted from this rendering (the count of omissions, if any, is noted below).}
\begin{receipttable}
\receiptrow{0a3819a978e31127}{2026-07-21T04:51:36Z}{build+decorrelated}{\badge{BadgeGreen}{current}}
\receiptrow{0bb1074df705d863}{2026-07-21T05:21:58Z}{build+decorrelated}{\badge{BadgeGreen}{current}}
\receiptrow{5a60f4d01fd81ac5}{2026-07-21T05:39:02Z}{build+replay+hardened+decorrelated}{\badge{BadgeGreen}{current}}
\end{receipttable}
\sourcenote{11 further receipt(s) in the store attest superseded builds of this module and are omitted from this rendering.}
\layerline{\layeritem{build}{lake build (elaboration)}{\badge{BadgeGreen}{current}}\layeritem{replay}{leanchecker (official kernel replay)}{\badge{BadgeGreen}{current}}\layeritem{hardened}{safe\_verify (hardened screen)}{\badge{BadgeGreen}{current}}\layeritem{decorrelated}{lean4export + nanoda (decorrelated kernel)}{\badge{BadgeGreen}{current}}}
\sourcenote{\badge{BadgeGray}{comparator: none} no fixture in the three-implementation comparator harness (compiled / kernel-interpreted / reference-interpreted) exercises this rule.}
\ruleblock{HOL4 statement}
\absentbox{OUT OF SCOPE BY DESIGN}{Out of scope for the HOL side \emph{by design}: the alias unpacks to a
single safety-battery node before emission --- the wire never carries a
negated-deadline node, so the HOL statement that would verify the
discharged certificate is the safety battery\textquotesingle s own
(absent today: the satisfaction stratum has no HOL artefact, per that
section). Checked against the HOL sources: the \texttt{bern\_catalogue}
theory is the base-calculus checker cone and carries no spec-surface
path operators.
}
\ruleblock{HOL verifier receipts}
\absentbox{NONE}{no HOL verification receipt exists for this rule.}
\ruleblock{Proof}

\sourcenote{Genre: alias rule (tier 3) --- a respelling template; the
batch's shared template is Recipe~\ref{rec:sd-template} of the
\rulename{QStreettDegenerate} section. The soundness content is two lines, recorded as such in the
source catalogue, and both are now mechanized and bound in the entry:
the inclusion (\texttt{negF\_denot\_subset}, axiom-free) and its
measure transport (\texttt{sound\_QNegFAlias}, profile-generic).}

\begin{brecipe}[the filled-in application]\label{rec:nf-filled}
\begin{mathpar}
\inferrule*[Right=\rulename{QNegFAlias}]{
  \mathrm{child} : \mathsf{QRegionSat}\bigl(\mathsf{SProbGE}(p,
  \mathsf{PAlways}(\mathsf{SNot}\,\varphi)),\ R\bigr)
  \ \text{(one \rulename{QSafetyBattery} conclusion)}
}{
  \mathcal{M},\ x_0 \in R \vDash
  \mathsf{SProbGE}\bigl(p,\ \mathsf{PNotPath}(\mathsf{PDeadline}(T,
  \varphi, \mathrm{false}))\bigr)
}
\end{mathpar}
One node to one node: the child is the safety battery's conclusion at
avoid set $[\![\varphi]\!]$, the concluded judgment re-reads the same
certified region against the negated deadline, horizon $T$ taken from
the goal. In the staged spec vocabulary the conclusion's formula is the
negated-path reading of a deadline node whose met-flag is false; the
child's is the always of the negation --- the rewrite touches the
formula only.
\end{brecipe}

\begin{brecipe}[payload transport]\label{rec:nf-transport}
\begin{center}\small
\begin{tabular}{@{}ll@{}}
\textbf{\rulename{QNegFAlias} node} & \textbf{\rulename{QSafetyBattery} node}\\
\midrule
certified region $R$, probability $p$ & verbatim\\
goal formula $\neg\mathsf{F}_{\le T}\,\varphi$ & child formula $\mathsf{G}\,\neg\varphi$ (horizon dropped)\\
horizon $T$ & not consumed (\autoref{rem:nf-lossy})\\
\end{tabular}
\end{center}
Unpack-function slot: no verified payload function in this edition
(batch template); the
semantic witness is bound in the entry (\texttt{negF\_denot\_subset} +
\texttt{sound\_QNegFAlias}).
\end{brecipe}

\begin{bprop}[the exactness inclusion --- two lines]\label{prop:nf-exact}
$[\![\mathsf{G}\,\neg\varphi]\!] \subseteq
[\![\mathsf{G}_{\le T}\,\neg\varphi]\!] =
[\![\neg\mathsf{F}_{\le T}\,\varphi]\!]$.
\end{bprop}
\begin{proof}
A path that never satisfies $\varphi$ in particular does not satisfy
$\varphi$ at any time $\le T$ --- the inclusion, by restriction of a
universal quantifier. And not satisfying $\varphi$ at any time $\le T$
is the definition of not reaching $\varphi$ within $T$ --- the equality,
definitional. Hence every path in the child's path set lies in the
conclusion's, and the probability bound transports monotonically.
\end{proof}

\begin{bprop}[soundness corollary]\label{prop:nf-sound}
If the child judgment holds --- from every state of $R$, with
probability $\ge p$ the path always avoids $[\![\varphi]\!]$ --- then
from every state of $R$, with probability $\ge p$ the path does not
reach $[\![\varphi]\!]$ within $T$.
\end{bprop}
\begin{proof}
Proposition~\ref{prop:nf-exact} plus monotonicity of measures, applied
under the safety battery's theorem. There is no other content --- and
this two-step corollary is precisely what the bound statements are:
the inclusion and the transport, one decl each.
\end{proof}

\begin{bremark}[the deliberate lossiness]\label{rem:nf-lossy}
The alias proves the bounded claim from the unbounded one, discarding
$T$. When $\varphi$ genuinely occurs eventually --- only after the
horizon --- the child judgment is unprovable and this alias is the wrong
tool; the deadline machinery (the deadline battery, with its geometric
decay certificates) proves the sharp bounded claim instead. The alias
exists because permanent avoidance is the common case, and for it the
safety kernel is both cheaper and already ported.
\end{bremark}

\begin{bremark}[the mirror that is not a rule]\label{rem:nf-mirror}
There is deliberately no dual entry discharging a negated
\emph{always}: its conclusion would need unbounded eventuality
($\mathsf{F}\varphi$), and the spec grammar has no unbounded-eventually
constructor --- an exclusion by testable-negative design, unbounded
eventuality surviving only as a provenance flag recording that a
deadline's horizon was producer-chosen. What replaces the missing rule
is three separate pieces: a negation-normal-form duality
(\rulename{QNnfRewrite}), a horizon instantiation (a producer choice,
flagged), and a weakening line \emph{inside} the deadline battery's
lemma where the flag is discharged. The asymmetry is worth this remark
because a port that treated the pair as one duality entry would smuggle
a choice and a metatheory line into what looks like a definitional
rewrite.
\end{bremark}

\begin{bremark}[tier flag]
Alias as sorted; if the folded inclusion of
Proposition~\ref{prop:nf-exact} must itself be a rule application
rather than a lemma of the unpack, the tier moves to admissible ---
flagged in the source catalogue, machinery unchanged either way.
\end{bremark}

\begin{anchors}
Recipe \ref{rec:nf-filled} target & the \rulename{QSafetyBattery} section\\
Proposition \ref{prop:nf-exact} + corollary & \texttt{negF\_denot\_subset} + \texttt{sound\_QNegFAlias} (two lines, as priced; bound in the entry)\\
Remark \ref{rem:nf-mirror} & the spec grammar\textquotesingle s constructor census (no unbounded-eventually node)\\
\end{anchors}


\blcatsection{Discharge Kits, Closure, and Arithmetic Leaves}
\rulesection{59}{HypMultiplier}{admissible rule}
\label{blcatrule-HypMultiplier}
\ruleblock{The rule (Steele metanotation)}
\sourcenote{as realized in the mechanized development; source row: \texttt{\texttt{QTSL-rules.toml}} --- \emph{HypMultiplier}. The metanotation statement is typeset by hand for notational quality; its mechanized ground truth is the Lean statement below, extracted from the sources at every build of this book and hash-bound. The rule \emph{roster} is machine-checked (every mechanized rule tag round-trips to exactly one catalogue row); the statement \emph{text} is editorial, checked by eye against the extracted Lean beside it.}
\begin{mathpar}
\inferrule*{
  \vdash \lambda \ge 0 \ \text{on } \mathrm{cell} \\
  \vdash N + \lambda\, g_q \le 0 \ \text{on } \mathrm{cell}
}{
  \vdash N \le 0 \ \text{on } \mathrm{cell} \cap \{g_q \ge 0\}
}
\end{mathpar}
the child atom $\{g_q \ge 0\}$ is a concluded child certificate\textquotesingle s region.

\ruleflavor{the child atom enters through the presentation bridge; both premises are \rulename{BernCone}-family rows, replayed by the base kernel — realized as the verified pointwise warrant \texttt{hypMultiplier\_sound}.}
\par\medskip
\rulename{HypMultiplier} is the discharge kit\textquotesingle s
multiplier move, and the only member of the kit that manufactures a
new row rather than reshaping the scope of an existing one. The
situation it serves: a premise row of a satisfaction-judgment rule
must hold not on a whole cell but on the part of the cell where a
child certificate\textquotesingle s region atom holds --- $N \le 0$ on
$\mathrm{cell} \cap \{g_q \ge 0\}$, with $\{g_q \ge 0\}$ the region
concluded by a nested certificate. The base
calculus\textquotesingle s rows are scoped to cells, not to
cell-with-polynomial-atom intersections, so the obligation cannot be
posted as a base row directly. The rule buys the restriction with one
multiplier: exhibit $\lambda$ and post two ordinary cone-family rows
--- $\lambda \ge 0$ on the cell, and $N + \lambda\, g_q \le 0$ on the
cell.

What the checker recomputes: nothing beyond the two premise rows,
which are cone-membership sequents the base kernel already replays;
the combination step is the degree-one corner of the
positivstellensatz family, the same multiplier discipline as the
calculus\textquotesingle s Farkas machinery. The producer typically
finds $\lambda$ by linear programming; only the two rows ship.

Both hypotheses are load-bearing, in opposite ways. Drop the sign row
and a negative multiplier flips the direction of the trade: with
$N \equiv 1$, $\lambda \equiv -1$, $g_q \equiv 1$, the combined row
$N + \lambda g_q \equiv 0 \le 0$ holds while the conclusion fails
everywhere on the child atom. Widen the conclusion to the whole cell
and the rule is unsound wherever the child atom fails: where
$g_q(x) < 0$ the term $\lambda(x)\, g_q(x)$ is a subsidy rather than
a tax, and the combined row bounds nothing. The proof section walks
both, counterexamples attached.

\cataloguenote{Before this rule\textquotesingle s mechanization, the move lived in
the producer toolkit as unchecked metatheory --- the discharge step
performed the combination and asserted the scoped row. The verified
lemma closes that trust boundary. The certificate row that would let
the wire format carry the move natively is out of scope of the current
baked theory, so today the lemma warrants the toolkit\textquotesingle s
elaboration rather than a kernel-replayed node.
}
\ruleblock{Lean statement}
\par\medskip
The auxiliary library states the extension theory\textquotesingle s
structural lemmas over the compiled semantic class: a bare state
type, scopes and cells as sets, observables as real-valued functions.
The presentation-level child atom enters through the named bridge
hypothesis (\texttt{hbridge}), which the presentation layer
discharges definitionally at compile time; nothing else about the
system enters at all.

\declhead{BernsteinLogicAux.Structural}{hypMultiplier\_sound}{\badge{BadgeGreen}{statement bound}}
\begin{leansrc}
/-- **T2 HypMultiplier**, soundness. `childAtom` is the presentation-level
child atom; its bridge names the compiled superlevel reading `{x | 0 ≤ g x}`.
The premises are the two compiled BernCone-family rows. -/
theorem hypMultiplier_sound (cell childAtom : Set X) (N lam g : X → ℝ)
    (hbridge : childAtom = {x | 0 ≤ g x})
    (hlam : ∀ x ∈ cell, 0 ≤ lam x)
    (hcomb : ∀ x ∈ cell, N x + lam x * g x ≤ 0) :
    ∀ x ∈ cell ∩ childAtom, N x ≤ 0
\end{leansrc}
\ruleblock{Correspondence (Steele display $\leftrightarrow$ Lean)}
\badge{BadgeGreen}{reviewer-certified against the displayed texts}
\sourcenote{The dictionary below is a reviewer-certified statement of how the mechanized statements above interpret the displayed rule: display metavariable $\mapsto$ Lean binder, each display hypothesis $\mapsto$ the premise or side condition carrying it, the display conclusion $\mapsto$ the goal --- with any deliberate difference in width between display and mechanization named outright. Certification is recorded as a hash triple over (the displayed cell, the extracted statement texts, this dictionary\textquotesingle s own text); if any of the three changes, this block turns amber until a reviewer re-certifies. The correspondence claim itself remains editorial --- reviewer-grade, not machine-checked; its machine-checked successor types each soundness statement against the rule as a structured object, and supersedes this block per rule as it lands.}
One bound statement, transcription-close. The displayed premises
$\vdash \lambda \ge 0$ on cell and $\vdash N + \lambda g_q \le 0$ on
cell $\mapsto$ \texttt{hlam}/\texttt{hcomb}, pointwise on
\texttt{cell}; the conclusion $\vdash N \le 0$ on $\mathrm{cell} \cap
\{g_q \ge 0\}$ $\mapsto$ the goal \texttt{∀ x ∈ cell ∩ childAtom, N x
≤ 0}, with the child atom entering through the System-bridge
hypothesis \texttt{hbridge : childAtom = \{x | 0 ≤ g x\}} --- the
presentation-level region named by an equality the presentation layer
will discharge definitionally; until then it is an honest binder, not
an assertion. Width note: \texttt{N}, \texttt{lam}, \texttt{g} are
arbitrary real-valued functions and \texttt{cell} an arbitrary set ---
the BernCone row structure and polynomiality are consumed one layer
down where the rows are replayed; the certified warrant is the
pointwise Farkas step alone, degree-1 positivstellensatz as displayed,
strictly wider than the polynomial instance, safe direction.

\ruleblock{Lean verifier receipts}
\sourcenote{module \texttt{BernsteinLogicAux.Structural}; bound build \texttt{1222cdcc326f\ldots} (bound from receipt cb0c3820c6877f22, 2026-07-29T01:03:36Z). A receipt records that named checking layers accepted a build: (module, compiled-artifact SHA-256, environment SHA-256, checker layers), minted when the checking gate ran; the store receipts attesting the bound build are listed with the layers each run credited. Nothing here asks the reader to trust the store: every layer is an independently runnable program over the quoted sources --- rebuild the module and re-run the checkers to reproduce the verdict. Receipts attesting superseded builds credit no layer and are omitted from this rendering (the count of omissions, if any, is noted below).}
\begin{receipttable}
\receiptrow{0b7e68907f45b27f}{2026-07-28T23:48:36Z}{build+decorrelated}{\badge{BadgeGreen}{current}}
\receiptrow{e97d5614d4f08c7e}{2026-07-29T00:00:35Z}{build+decorrelated}{\badge{BadgeGreen}{current}}
\receiptrow{27de21f047e6b194}{2026-07-29T00:49:05Z}{build+replay+hardened+decorrelated}{\badge{BadgeGreen}{current}}
\receiptrow{cb0c3820c6877f22}{2026-07-29T01:03:36Z}{build+replay+hardened+decorrelated}{\badge{BadgeGreen}{current}}
\receiptrow{3530841119a9a607}{2026-07-31T22:00:22Z}{build+decorrelated}{\badge{BadgeGreen}{current}}
\receiptrow{b6cc23eb3e5b7f97}{2026-07-31T22:29:16Z}{build+decorrelated}{\badge{BadgeGreen}{current}}
\receiptrow{bf3dacc3ab7f20ea}{2026-08-02T23:31:31Z}{build+decorrelated}{\badge{BadgeGreen}{current}}
\receiptrow{384a8569d2005df7}{2026-08-03T16:27:35Z}{build+decorrelated}{\badge{BadgeGreen}{current}}
\receiptrow{717e07acf9acac8a}{2026-08-03T22:46:19Z}{build+decorrelated}{\badge{BadgeGreen}{current}}
\end{receipttable}
\sourcenote{11 further receipt(s) in the store attest superseded builds of this module and are omitted from this rendering.}
\layerline{\layeritem{build}{lake build (elaboration)}{\badge{BadgeGreen}{current}}\layeritem{replay}{leanchecker (official kernel replay)}{\badge{BadgeGreen}{current}}\layeritem{hardened}{safe\_verify (hardened screen)}{\badge{BadgeGreen}{current}}\layeritem{decorrelated}{lean4export + nanoda (decorrelated kernel)}{\badge{BadgeGreen}{current}}}
\sourcenote{\badge{BadgeGray}{comparator: none} no fixture in the three-implementation comparator harness (compiled / kernel-interpreted / reference-interpreted) exercises this rule.}
\ruleblock{HOL4 statement}
\absentbox{NO HOL STATEMENT YET}{The HOL4 development\textquotesingle s object is the base
calculus\textquotesingle s wire-format checker (theory
\texttt{bern\_catalogue} ports the catalogue check-function cone).
This rule belongs to the extension theory\textquotesingle s discharge
kit: it postdates that port, and no row for the move exists in the
baked theory --- there is nothing for a HOL statement to be
about. The box becomes fillable when a revision carries the
move\textquotesingle s certificate rows and the corresponding check is
ported.
}
\ruleblock{HOL verifier receipts}
\absentbox{NONE}{no HOL verification receipt exists for this rule.}
\ruleblock{Proof}

\sourcenote{Genre: admissible rule (tier 2) --- a proof-tree-rewriting
recipe; the general admissible template, both species included, is
stated in the \rulename{Norm} section and instantiated here in its
boundary species (discharge against a named verified theorem). The rule
is sequent-side machinery: premises and conclusion are
all $\vdash$ rows, and what the extension theory contributes is the
\emph{scope} vocabulary --- the child atom $\{g_q \ge 0\}$ names a
region concluded by a nested certificate, a shape no base cell can
express. The rule closes a named trust boundary: before its
mechanization, the move was performed producer-side and asserted.}

\begin{brecipe}[the \rulename{HypMultiplier} instance]\label{rec:hm-instance}
\[
\inferrule*[Right=\rulename{HypMultiplier}]{
  \mathcal{T}_\lambda \rhd\ \vdash \lambda \ge 0 \ \text{on } \mathrm{cell} \\
  \mathcal{T}_c \rhd\ \vdash N + \lambda\, g_q \le 0 \ \text{on } \mathrm{cell}
}{
  \vdash N \le 0 \ \text{on } \mathrm{cell} \cap \{g_q \ge 0\}
}
\qquad\rightsquigarrow\qquad
\begin{array}{l}
\text{\upshape the premise subtrees stand unchanged;}\\
\text{\texttt{hypMultiplier\_sound} discharges the scoped row.}
\end{array}
\]
The rewrite is degenerate in the same honest sense as the base
catalogue\textquotesingle s ground-fact rule: the conclusion is not
derivable from catalogue nodes, because the base grammar has no row
form scoped to a cell-with-polynomial-atom intersection --- that scope
exists one layer up, where the child certificate that concluded
$\{g_q \ge 0\}$ lives. The two premise subtrees are ordinary
cone-family derivations, checked by the base kernel unchanged; the
node itself discharges against the verified pointwise warrant below,
which consumes the two premises plus the one bridge hypothesis naming
the child atom's compiled reading (the closing remark) --- and nothing
else.
\end{brecipe}

\begin{bprop}[semantic warrant]\label{prop:hm-sound}
If $\lambda \ge 0$ everywhere on the cell and $N + \lambda g_q \le 0$
everywhere on the cell, then $N \le 0$ everywhere on
$\mathrm{cell} \cap \{g_q \ge 0\}$.
\end{bprop}
\begin{proof}
Fix $x$ in the intersection. Then $g_q(x) \ge 0$ (membership in the
child atom) and $\lambda(x) \ge 0$ (the sign row), so
$\lambda(x)\, g_q(x) \ge 0$, and
\[
N(x) \;\le\; N(x) + \lambda(x)\, g_q(x) \;\le\; 0,
\]
the first inequality by nonnegativity of the added term, the second by
the combined row at $x$. That is the whole proof ---
\texttt{hypMultiplier\_sound} is the two displayed inequalities.
\end{proof}

\begin{bremark}[the sign row is half the rule]
Drop $\vdash \lambda \ge 0$ and the rule is unsound. Counterexample:
$N \equiv 1$, $\lambda \equiv -1$, $g_q \equiv 1$ on any nonempty cell
(a singleton suffices). The
combined row holds identically ($1 - 1 = 0 \le 0$), yet $N \equiv 1$
fails everywhere on $\mathrm{cell} \cap \{g_q \ge 0\} = \mathrm{cell}$.
A free-signed multiplier converts the tax the child atom levies into a
subsidy: the combination can cancel exactly the positivity the
conclusion denies.
\end{bremark}

\begin{bremark}[the scope restriction is the other half]
The conclusion cannot claim the whole cell. With $N \equiv 1$,
$\lambda \equiv 1$, $g_q \equiv -1$ on the same nonempty cell, both
premises hold
($\lambda \ge 0$; $1 - 1 = 0 \le 0$) while $N > 0$ everywhere --- and
the rule is nonetheless sound, because
$\mathrm{cell} \cap \{g_q \ge 0\} = \varnothing$ and the scoped claim
is vacuous. Where $g_q(x) < 0$ the term $\lambda(x) g_q(x)$ subsidizes
$N$ instead of taxing it, and the combined row bounds nothing: the
scoped conclusion is exactly as strong as the premises pay for, no
stronger.
\end{bremark}

\begin{bremark}[degree-one positivstellensatz, and where the search lives]
The move is the degree-one corner of certifying positivity over a
constrained set by nonnegative combinations: Farkas\textquotesingle{}
multiplier discipline \cite{farkas1902} (see \cite{schrijver1986} for
the sign conventions) in the affine case, sitting at the bottom of the
polynomial cone hierarchy of Krivine \cite{krivine1964} and Handelman
\cite{handelman1988} that the calculus\textquotesingle s cone rule
draws on. Nothing beyond degree one is consumed here: one multiplier,
one product, no certificate machinery beyond the two posted rows. The
\emph{search} for $\lambda$ (typically a linear program over the
cone-membership witnesses) is producer-side; only its result ships,
and the checker replays the two rows without ever seeing the search.
\end{bremark}

\begin{bremark}[the bridge hypothesis]
In the verified statement the child atom is an abstract set with one
bridge hypothesis naming its compiled superlevel reading
($\{x \mid 0 \le g_q(x)\}$). When the presentation layer of the
extension theory lands, the hypothesis discharges definitionally from
the compile --- the scoped conclusion itself is unchanged; only the
bridge binder becomes definitional. Until then it is the honest
record of exactly what the warrant assumes about the child
certificate\textquotesingle s conclusion --- its region is the
superlevel set it claims to be, and nothing more.
\end{bremark}

\begin{anchors}
Recipe \ref{rec:hm-instance} warrant & \texttt{hypMultiplier\_sound} (BernsteinLogicAux/Structural.lean)\\
premise rows & cone-family sequent machinery of the base calculus (the \rulename{BernCone} and \rulename{Farkas} sections)\\
child-atom reading & the \texttt{hbridge} hypothesis of \texttt{hypMultiplier\_sound}\\
trust boundary closed & the producer-side discharge step this lemma now warrants\\
\end{anchors}


\rulesection{60}{QFullDomainWeaken}{alias rule}
\label{blcatrule-QFullDomainWeaken}
\ruleblock{The rule (Steele metanotation)}
\sourcenote{as realized in the mechanized development; source row: \texttt{\texttt{QTSL-rules.toml}} --- \emph{QFullDomainWeaken}. The metanotation statement is typeset by hand for notational quality; its mechanized ground truth is the Lean statement below, extracted from the sources at every build of this book and hash-bound. The rule \emph{roster} is machine-checked (every mechanized rule tag round-trips to exactly one catalogue row); the statement \emph{text} is editorial, checked by eye against the extracted Lean beside it.}
\begin{mathpar}
\inferrule*{
  \vdash P \ \text{on } D_m \\
  \mathrm{scope} \subseteq D_m \ \text{(compile fact, not a posted premise)}
}{
  \vdash P \ \text{on scope}
}
\end{mathpar}
emitted shape: $\mathrm{scope} = D_m \smallsetminus \mathrm{Off}$ (the off-scoped row).

\ruleflavor{payload verbatim, scope replaced — one application at the region layer, realized as \texttt{fullDomainWeaken\_sound} (axiom-free) $+$ \texttt{\_offScope}.}
\par\medskip
\rulename{QFullDomainWeaken} is the cheapest move in the discharge
kit: an off-scoped row --- an obligation $P$ on
$D_m \smallsetminus \mathrm{Off}$, the part of a
mode\textquotesingle s domain outside an off-region --- is discharged
by proving $P$ on all of $D_m$ and forgetting the carve-out. It is an
alias: one scope-weakening application at the region layer,
transporting the row payload verbatim and replacing the scope by the
full domain.

A producer reaches for it first. When the row happens to hold
globally, the off-region was doing no work for this row, and the
discharge costs nothing beyond the containment of the scope in the
domain --- which is a compile guarantee, not a posted premise. When
the global row fails (the off-region was earning its keep), the move
is simply unavailable, and the kit\textquotesingle s working move,
\rulename{QVerifiedCover}, localizes the effort instead.

What the checker recomputes: nothing. Soundness is the monotonicity
of universal quantification --- a row asserted on a superset
restricts to any subset --- and the verified lemma is a one-line
restriction along an inclusion, provable without any classical
reasoning at all.

\cataloguenote{Two deliberate non-identifications. First, the alias\textquotesingle s
target is not a catalogue node: no region-weakening rule exists to
apply, and the unpacking discharges directly against the verified
monotonicity lemma --- the alias-tier twin of the admissible
tier\textquotesingle s recorded reduction-to-verified-machinery
boundary case. Second, the rule is deliberately NOT the base
calculus\textquotesingle s sequent weakening: that rule manages a
sequent\textquotesingle s premise list; this one manages the region a
row is quantified over, one extension layer up. Unifying them would
conflate two different quantifiers that merely rhyme.
}
\ruleblock{Lean statement}
\par\medskip
Stated over the compiled semantic class (sets over a bare state
type). The off-scoped shape the discharge plan actually emits is the
second declaration, an instance of the first; the scope-in-domain
containment enters as the bridge hypothesis the compile guarantees.

\declhead{BernsteinLogicAux.Structural}{fullDomainWeaken\_sound}{\badge{BadgeGreen}{statement bound}}
\begin{leansrc}
/-- **T3a FullDomainWeaken**, soundness. `scope` is the presentation-level
off-scoped row's scope; the bridge names the one fact the compile
guarantees about it: it lies inside the full domain. -/
theorem fullDomainWeaken_sound (Dm scope : Set X) (P : X → Prop)
    (hbridge : scope ⊆ Dm)
    (hrow : ∀ x ∈ Dm, P x) :
    ∀ x ∈ scope, P x
\end{leansrc}
\declhead{BernsteinLogicAux.Structural}{fullDomainWeaken\_offScope}{\badge{BadgeGreen}{statement bound}}
\begin{leansrc}
/-- The off-scope instance: the scope compiled as `D_m` minus an off-region.
(The shape the discharge plan actually emits.) -/
theorem fullDomainWeaken_offScope (Dm Off : Set X) (P : X → Prop)
    (hrow : ∀ x ∈ Dm, P x) :
    ∀ x ∈ Dm \ Off, P x
\end{leansrc}
\ruleblock{Correspondence (Steele display $\leftrightarrow$ Lean)}
\badge{BadgeGreen}{reviewer-certified against the displayed texts}
\sourcenote{The dictionary below is a reviewer-certified statement of how the mechanized statements above interpret the displayed rule: display metavariable $\mapsto$ Lean binder, each display hypothesis $\mapsto$ the premise or side condition carrying it, the display conclusion $\mapsto$ the goal --- with any deliberate difference in width between display and mechanization named outright. Certification is recorded as a hash triple over (the displayed cell, the extracted statement texts, this dictionary\textquotesingle s own text); if any of the three changes, this block turns amber until a reviewer re-certifies. The correspondence claim itself remains editorial --- reviewer-grade, not machine-checked; its machine-checked successor types each soundness statement against the rule as a structured object, and supersedes this block per rule as it lands.}
The displayed premise $\vdash P$ on $D_m$ $\mapsto$ \texttt{hrow : ∀ x
∈ Dm, P x}; the compile fact $\mathrm{scope} \subseteq D_m$ ---
displayed as ``not a posted premise'' --- $\mapsto$ the System-bridge
hypothesis \texttt{hbridge : scope ⊆ Dm}, an explicit binder here
because lean-aux lacks the compile layer that guarantees it; the
conclusion $\mapsto$ \texttt{∀ x ∈ scope, P x}. The emitted shape
$\mathrm{scope} = D_m \smallsetminus \mathrm{Off}$ is bound separately
as \texttt{fullDomainWeaken\_offScope}, where the bridge discharges by
\texttt{Set.diff\_subset} --- zero hypotheses beyond the row.
Width note: \texttt{P : X → Prop} arbitrary (payload verbatim, as
displayed --- the certified statement cannot even see the payload,
which is the strongest possible form of ``payload verbatim, scope
replaced''); the root statement is axiom-free pure
$\forall$-monotonicity, exactly the alias tier\textquotesingle s
price.

\ruleblock{Lean verifier receipts}
\sourcenote{module \texttt{BernsteinLogicAux.Structural}; bound build \texttt{1222cdcc326f\ldots} (bound from receipt cb0c3820c6877f22, 2026-07-29T01:03:36Z). A receipt records that named checking layers accepted a build: (module, compiled-artifact SHA-256, environment SHA-256, checker layers), minted when the checking gate ran; the store receipts attesting the bound build are listed with the layers each run credited. Nothing here asks the reader to trust the store: every layer is an independently runnable program over the quoted sources --- rebuild the module and re-run the checkers to reproduce the verdict. Receipts attesting superseded builds credit no layer and are omitted from this rendering (the count of omissions, if any, is noted below).}
\begin{receipttable}
\receiptrow{0b7e68907f45b27f}{2026-07-28T23:48:36Z}{build+decorrelated}{\badge{BadgeGreen}{current}}
\receiptrow{e97d5614d4f08c7e}{2026-07-29T00:00:35Z}{build+decorrelated}{\badge{BadgeGreen}{current}}
\receiptrow{27de21f047e6b194}{2026-07-29T00:49:05Z}{build+replay+hardened+decorrelated}{\badge{BadgeGreen}{current}}
\receiptrow{cb0c3820c6877f22}{2026-07-29T01:03:36Z}{build+replay+hardened+decorrelated}{\badge{BadgeGreen}{current}}
\receiptrow{3530841119a9a607}{2026-07-31T22:00:22Z}{build+decorrelated}{\badge{BadgeGreen}{current}}
\receiptrow{b6cc23eb3e5b7f97}{2026-07-31T22:29:16Z}{build+decorrelated}{\badge{BadgeGreen}{current}}
\receiptrow{bf3dacc3ab7f20ea}{2026-08-02T23:31:31Z}{build+decorrelated}{\badge{BadgeGreen}{current}}
\receiptrow{384a8569d2005df7}{2026-08-03T16:27:35Z}{build+decorrelated}{\badge{BadgeGreen}{current}}
\receiptrow{717e07acf9acac8a}{2026-08-03T22:46:19Z}{build+decorrelated}{\badge{BadgeGreen}{current}}
\end{receipttable}
\sourcenote{11 further receipt(s) in the store attest superseded builds of this module and are omitted from this rendering.}
\layerline{\layeritem{build}{lake build (elaboration)}{\badge{BadgeGreen}{current}}\layeritem{replay}{leanchecker (official kernel replay)}{\badge{BadgeGreen}{current}}\layeritem{hardened}{safe\_verify (hardened screen)}{\badge{BadgeGreen}{current}}\layeritem{decorrelated}{lean4export + nanoda (decorrelated kernel)}{\badge{BadgeGreen}{current}}}
\sourcenote{\badge{BadgeGray}{comparator: none} no fixture in the three-implementation comparator harness (compiled / kernel-interpreted / reference-interpreted) exercises this rule.}
\ruleblock{HOL4 statement}
\absentbox{NO HOL STATEMENT YET}{The HOL4 development\textquotesingle s object is the base
calculus\textquotesingle s wire-format checker (theory
\texttt{bern\_catalogue} ports the catalogue check-function cone).
This rule is scope machinery of the extension
theory\textquotesingle s discharge kit: it postdates that port, and
no row for the move exists in the baked theory --- there is
nothing for a HOL statement to be about. The box becomes fillable
when a revision carries the move\textquotesingle s certificate rows
and the corresponding check is ported.
}
\ruleblock{HOL verifier receipts}
\absentbox{NONE}{no HOL verification receipt exists for this rule.}
\ruleblock{Proof}

\sourcenote{Genre: alias rule (tier 3) --- a respelling template. An
alias owes the book five artefacts: the verified 1$\to$1 unpack, the
filled-in application of the target, the payload-transport table, the
exactness identity, and the soundness corollary through the
target\textquotesingle s lemma. This alias is the
tier\textquotesingle s degenerate case --- its target is a verified
lemma rather than a catalogue node, because no region-weakening rule
exists to apply --- and the five artefacts below say so honestly
rather than inventing a node.}

\begin{brecipe}[the alias template, degenerate-target form]\label{rec:fdw-template}
An alias rule $A$ is warranted by a verified function carrying its
payload to the payload of a \emph{single} target application --- one
node $\to$ one node, nothing derived. Where the target is not a
catalogue node but a verified theorem, the transport carries the
payload into the theorem\textquotesingle s hypotheses instead, and the
soundness corollary is the theorem itself: the alias-tier twin of the
admissible tier\textquotesingle s recorded
reduction-to-verified-machinery boundary case.
\end{brecipe}

\begin{brecipe}[the filled-in application]\label{rec:fdw-filled}
\[
\inferrule*[Right=\rulename{QFullDomainWeaken}]{
  \mathcal{T} \rhd\ \vdash P \ \text{on } D_m
}{
  \vdash P \ \text{on } D_m \smallsetminus \mathrm{Off}
}
\qquad\rightsquigarrow\qquad
\texttt{fullDomainWeaken\_offScope}\bigl[\mathcal{T}\bigr]
\;\rhd\; \vdash P \ \text{on } D_m \smallsetminus \mathrm{Off}
\]
with the general form --- any scope the compile certifies inside the
domain --- carried by \texttt{fullDomainWeaken\_sound}: from $P$ on
$D_m$ and $\mathrm{scope} \subseteq D_m$, conclude $P$ on
$\mathrm{scope}$.
\end{brecipe}

\begin{brecipe}[payload transport]\label{rec:fdw-transport}
\begin{center}\small
\begin{tabular}{@{}ll@{}}
\textbf{off-scoped row} & \textbf{full-domain row}\\
\midrule
row payload $P$ & $P$ (verbatim)\\
scope $D_m \smallsetminus \mathrm{Off}$ & $D_m$ (the off-region dropped)\\
$\mathrm{scope} \subseteq D_m$ & compile guarantee (\texttt{hbridge}), not a posted premise\\
\end{tabular}
\end{center}
\end{brecipe}

\begin{bprop}[exactness identity]\label{prop:fdw-identity}
$D_m \smallsetminus \mathrm{Off} \subseteq D_m$; more generally, every
off-scoped row\textquotesingle s scope is contained in its
mode\textquotesingle s domain.
\end{bprop}
\begin{proof}
A set difference is contained in its left operand --- one line
(\texttt{Set.diff\_subset}). The general containment is the compile
fact the bridge hypothesis names, supplied by the presentation layer,
never proven per-instance.
\end{proof}

\begin{bprop}[soundness corollary]\label{prop:fdw-sound}
If $P$ holds everywhere on $D_m$, it holds everywhere on any
$\mathrm{scope} \subseteq D_m$ --- in particular on
$D_m \smallsetminus \mathrm{Off}$.
\end{bprop}
\begin{proof}
Restriction of a universally quantified statement along an inclusion:
for $x \in \mathrm{scope}$, the inclusion gives $x \in D_m$, and the
full-domain row gives $P(x)$. \texttt{fullDomainWeaken\_sound} is this
one line, and it is axiom-free --- no classical reasoning, no choice:
pure monotonicity of $\forall$.
\end{proof}

\begin{bremark}[what is spent]
Nothing, soundness-wise --- and possibly everything,
feasibility-wise: if $P$ fails anywhere in $D_m \cap \mathrm{Off}$
(points of the off-region outside $D_m$ never enter the row), the
full-domain row is simply unprovable and the move is unavailable. The
off-region existing at all signals that some row somewhere needed it;
this move is the discharge for the rows that did not, and its
availability is itself a diagnosis: the carve-out was free for this
row. When it fails, the kit\textquotesingle s working move
(\rulename{QVerifiedCover}) localizes the work instead of abandoning
the carve-out.
\end{bremark}

\begin{bremark}[deliberately not the sequent weakening]
The base calculus has a weakening rule; this is deliberately not it.
Sequent weakening grows a premise \emph{list} --- what a sequent
assumes; this move shrinks a quantification \emph{region} --- where a
row is asserted. The two live one extension layer apart and neither
subsumes the other: a base sequent has no scope field to weaken, and
this move never touches a premise list. The kinship is real, and the
identification would be a type error.
\end{bremark}

\begin{anchors}
Recipe \ref{rec:fdw-filled} unpack & \texttt{fullDomainWeaken\_offScope} (the emitted shape)\\
general form & \texttt{fullDomainWeaken\_sound} (axiom-free)\\
Proposition \ref{prop:fdw-identity} & \texttt{Set.diff\_subset}; the \texttt{hbridge} inclusion\\
region-layer monotonicity, generic & \texttt{regionSat\_mono} (same module)\\
\end{anchors}


\rulesection{61}{QVerifiedCover}{admissible rule}
\label{blcatrule-QVerifiedCover}
\ruleblock{The rule (Steele metanotation)}
\sourcenote{as realized in the mechanized development; source row: \texttt{\texttt{QTSL-rules.toml}} --- \emph{QVerifiedCover}. The metanotation statement is typeset by hand for notational quality; its mechanized ground truth is the Lean statement below, extracted from the sources at every build of this book and hash-bound. The rule \emph{roster} is machine-checked (every mechanized rule tag round-trips to exactly one catalogue row); the statement \emph{text} is editorial, checked by eye against the extracted Lean beside it.}
\begin{mathpar}
\inferrule*{
  \vdash P \ \text{on } C \ \text{(the cover row)} \\
  \bigl\{\vdash K_i \subseteq \mathrm{Off}\bigr\}_i \ \text{(per-component containments)} \\
  D_m \smallsetminus C \subseteq \textstyle\bigcup_i K_i \ \text{(residual decomposition)}
}{
  \vdash P \ \text{on } D_m \smallsetminus \mathrm{Off}
}
\end{mathpar}

\ruleflavor{the per-component containments discharge by exact Bernstein upper bounds, and the residual decomposition is recomputed by the checker. Dimension-free as realized (\texttt{verifiedCover\_components\_sound}, arbitrary component index); the cover bookkeeping $C \subseteq D_m$ is checker-side well-formedness, not consumed by soundness.}
\par\medskip
\rulename{QVerifiedCover} is the discharge kit\textquotesingle s
working move: an off-scoped row --- $P$ on
$D_m \smallsetminus \mathrm{Off}$ --- is discharged by proving the
row on a cover $C$ and certifying that everything of $D_m$ the cover
misses lies inside the off-region. The un-covered residual is
decomposed into components $K_i$, and each component ships one
containment premise $K_i \subseteq \mathrm{Off}$, discharged by exact
Bernstein upper bounds --- ordinary sequent rows the base kernel
replays.

A producer reaches for it when \rulename{QFullDomainWeaken} fails ---
the row genuinely needs its carve-out --- but the on-region can be
covered cleanly. Cover \emph{selection} is producer-side search (the
current toolkit selects one-dimensional covers); the verified
statement is dimension-free, so richer covers need no new warrant.

What the checker recomputes: the residual decomposition
$D_m \smallsetminus C \subseteq \bigcup_i K_i$ (set bookkeeping over
the posted components) and the per-component containment rows. The
case analysis is exhaustive by construction: a point of the scope is
either covered, and the row-on-$C$ premise applies, or un-covered, in
which case it lies in some $K_i \subseteq \mathrm{Off}$ ---
contradicting membership in the scope. Drop the per-component
containments and an un-covered point outside the off-region escapes
both cases; drop the decomposition and the component list could
silently miss part of the residual.

\cataloguenote{One premise of the catalogued statement is deliberately absent from
the verified one: the cover bookkeeping $C \subseteq D_m$. The proof
never consumes it --- a cover reaching outside the domain proves more
rows than needed, harmlessly --- so the mechanized statement omits
it, which is strictly stronger in the sound direction; the checker
keeps the containment as cover well-formedness. The neighbouring
scope-preparation steps (hull closure, cell covers, off-scope
splitting, reset containment, floor boxes) are deliberately not rules
at all: they are side conditions the checker recomputes, and rows for
them would misstate that.
}
\ruleblock{Lean statement}
\par\medskip
Stated over the compiled semantic class (sets over a bare state
type), dimension-free. The residual form is the kernel statement; the
per-component form is the shape the discharge plan emits, with an
arbitrary component index --- nothing one-dimensional survives into
the warrant.

\declhead{BernsteinLogicAux.Structural}{verifiedCover\_sound}{\badge{BadgeGreen}{statement bound}}
\begin{leansrc}
/-- **T3b VerifiedCover**, soundness (residual form). `scope` is the
off-scoped row's compiled scope `D_m \ Off`; the row is proven on the cover
`C`, and the un-covered residual `D_m \ C` is certified inside the
off-region. -/
theorem verifiedCover_sound (Dm C Off scope : Set X) (P : X → Prop)
    (hbridge : scope = Dm \ Off)
    (hrowC : ∀ x ∈ C, P x)
    (hresidual : Dm \ C ⊆ Off) :
    ∀ x ∈ scope, P x
\end{leansrc}
\declhead{BernsteinLogicAux.Structural}{verifiedCover\_components\_sound}{\badge{BadgeGreen}{statement bound}}
\begin{leansrc}
/-- **T3b VerifiedCover**, per-component form: the residual is decomposed
into components `K i`, each separately certified inside the off-region (the
"exact Bernstein upper bounds" premises, one per component). -/
theorem verifiedCover_components_sound {ι : Type*} (Dm C Off scope : Set X)
    (K : ι → Set X) (P : X → Prop)
    (hbridge : scope = Dm \ Off)
    (hrowC : ∀ x ∈ C, P x)
    (hdecomp : Dm \ C ⊆ ⋃ i, K i)
    (hcomp : ∀ i, K i ⊆ Off) :
    ∀ x ∈ scope, P x
\end{leansrc}
\ruleblock{Correspondence (Steele display $\leftrightarrow$ Lean)}
\badge{BadgeGreen}{reviewer-certified against the displayed texts}
\sourcenote{The dictionary below is a reviewer-certified statement of how the mechanized statements above interpret the displayed rule: display metavariable $\mapsto$ Lean binder, each display hypothesis $\mapsto$ the premise or side condition carrying it, the display conclusion $\mapsto$ the goal --- with any deliberate difference in width between display and mechanization named outright. Certification is recorded as a hash triple over (the displayed cell, the extracted statement texts, this dictionary\textquotesingle s own text); if any of the three changes, this block turns amber until a reviewer re-certifies. The correspondence claim itself remains editorial --- reviewer-grade, not machine-checked; its machine-checked successor types each soundness statement against the rule as a structured object, and supersedes this block per rule as it lands.}
The displayed premises map onto \texttt{verifiedCover\_components\_sound}
one for one: the cover row $\vdash P$ on $C$ $\mapsto$ \texttt{hrowC};
the per-component containments $K_i \subseteq \mathrm{Off}$ $\mapsto$
\texttt{hcomp : ∀ i, K i ⊆ Off} (the exact-Bernstein-upper-bound
premises, one per component, over an \emph{arbitrary} index type ---
the dimension-freedom the flavor text claims is
a binder fact); the recomputed residual decomposition $D_m
\smallsetminus C \subseteq \bigcup_i K_i$ $\mapsto$ \texttt{hdecomp},
an explicit hypothesis (recompute rendered as binder, the
\rulename{QStreettTargetCellCollapse} pattern); the conclusion scope
$D_m \smallsetminus \mathrm{Off}$ $\mapsto$ the System-bridge equality
\texttt{hbridge : scope = Dm \textbackslash{} Off}. The residual form
\texttt{verifiedCover\_sound} is the same statement with the
decomposition pre-collapsed. Width note, divergence D1 as certified:
the catalogued well-formedness $C \subseteq D_m$ appears in NO binder
--- soundness demonstrably never consumes it, which the module header
banks as the divergence and the flavor text
carries; payload \texttt{P} arbitrary, per the QFullDomainWeaken pattern.

\ruleblock{Lean verifier receipts}
\sourcenote{module \texttt{BernsteinLogicAux.Structural}; bound build \texttt{1222cdcc326f\ldots} (bound from receipt cb0c3820c6877f22, 2026-07-29T01:03:36Z). A receipt records that named checking layers accepted a build: (module, compiled-artifact SHA-256, environment SHA-256, checker layers), minted when the checking gate ran; the store receipts attesting the bound build are listed with the layers each run credited. Nothing here asks the reader to trust the store: every layer is an independently runnable program over the quoted sources --- rebuild the module and re-run the checkers to reproduce the verdict. Receipts attesting superseded builds credit no layer and are omitted from this rendering (the count of omissions, if any, is noted below).}
\begin{receipttable}
\receiptrow{0b7e68907f45b27f}{2026-07-28T23:48:36Z}{build+decorrelated}{\badge{BadgeGreen}{current}}
\receiptrow{e97d5614d4f08c7e}{2026-07-29T00:00:35Z}{build+decorrelated}{\badge{BadgeGreen}{current}}
\receiptrow{27de21f047e6b194}{2026-07-29T00:49:05Z}{build+replay+hardened+decorrelated}{\badge{BadgeGreen}{current}}
\receiptrow{cb0c3820c6877f22}{2026-07-29T01:03:36Z}{build+replay+hardened+decorrelated}{\badge{BadgeGreen}{current}}
\receiptrow{3530841119a9a607}{2026-07-31T22:00:22Z}{build+decorrelated}{\badge{BadgeGreen}{current}}
\receiptrow{b6cc23eb3e5b7f97}{2026-07-31T22:29:16Z}{build+decorrelated}{\badge{BadgeGreen}{current}}
\receiptrow{bf3dacc3ab7f20ea}{2026-08-02T23:31:31Z}{build+decorrelated}{\badge{BadgeGreen}{current}}
\receiptrow{384a8569d2005df7}{2026-08-03T16:27:35Z}{build+decorrelated}{\badge{BadgeGreen}{current}}
\receiptrow{717e07acf9acac8a}{2026-08-03T22:46:19Z}{build+decorrelated}{\badge{BadgeGreen}{current}}
\end{receipttable}
\sourcenote{11 further receipt(s) in the store attest superseded builds of this module and are omitted from this rendering.}
\layerline{\layeritem{build}{lake build (elaboration)}{\badge{BadgeGreen}{current}}\layeritem{replay}{leanchecker (official kernel replay)}{\badge{BadgeGreen}{current}}\layeritem{hardened}{safe\_verify (hardened screen)}{\badge{BadgeGreen}{current}}\layeritem{decorrelated}{lean4export + nanoda (decorrelated kernel)}{\badge{BadgeGreen}{current}}}
\sourcenote{\badge{BadgeGray}{comparator: none} no fixture in the three-implementation comparator harness (compiled / kernel-interpreted / reference-interpreted) exercises this rule.}
\ruleblock{HOL4 statement}
\absentbox{NO HOL STATEMENT YET}{The HOL4 development\textquotesingle s object is the base
calculus\textquotesingle s wire-format checker (theory
\texttt{bern\_catalogue} ports the catalogue check-function cone).
This rule is scope machinery of the extension
theory\textquotesingle s discharge kit: it postdates that port, and
no row for the move exists in the baked theory --- there is
nothing for a HOL statement to be about. The box becomes fillable
when the extension\textquotesingle s certificate rows land in a wire
format and the corresponding check is ported.
}
\ruleblock{HOL verifier receipts}
\absentbox{NONE}{no HOL verification receipt exists for this rule.}
\ruleblock{Proof}

\sourcenote{Genre: admissible rule (tier 2) --- a proof-tree-rewriting
recipe. The rule discharges an off-scoped row by an exhaustive case
analysis: prove the row on a cover, certify everything un-covered into
the off-region. Its premises are ordinary sequent rows; the case
analysis is the verified content.}

The admissible template, general form, is stated once for the whole
catalogue as Recipe~\ref{rec:norm-template} of the \rulename{Norm}
section; this section instantiates it.

\begin{brecipe}[the \rulename{QVerifiedCover} instance]\label{rec:vc-instance}
\[
\inferrule*[Right=\rulename{QVerifiedCover}]{
  \mathcal{T}_C \rhd\ \vdash P \ \text{on } C \\
  \bigl\{\, \mathcal{T}_i \rhd\ \vdash K_i \subseteq \mathrm{Off} \,\bigr\}_i \\
  D_m \smallsetminus C \subseteq \textstyle\bigcup_i K_i
}{
  \vdash P \ \text{on } D_m \smallsetminus \mathrm{Off}
}
\qquad\rightsquigarrow\qquad
\begin{array}{l}
\text{\upshape premise subtrees stand;}\\
\text{\texttt{verifiedCover\_components\_sound}}\\
\text{\upshape discharges the scoped row.}
\end{array}
\]
The per-component containments $K_i \subseteq \mathrm{Off}$ are exact
Bernstein upper-bound rows --- ordinary sequents the base kernel
replays; the residual decomposition is checker-recomputed set
bookkeeping over the posted components, not a proof obligation. The
node discharges against the verified case analysis below.
\end{brecipe}

\begin{bprop}[semantic warrant, residual form]\label{prop:vc-sound}
If $P$ holds everywhere on $C$ and
$D_m \smallsetminus C \subseteq \mathrm{Off}$, then $P$ holds
everywhere on $D_m \smallsetminus \mathrm{Off}$.
\end{bprop}
\begin{proof}
Fix $x \in D_m \smallsetminus \mathrm{Off}$ and split on membership in
the cover. If $x \in C$, the cover row gives $P(x)$. If $x \notin C$,
then $x \in D_m \smallsetminus C \subseteq \mathrm{Off}$ ---
contradicting $x \notin \mathrm{Off}$. The two cases are exhaustive,
so the second is vacuous and the first decides
(\texttt{verifiedCover\_sound}).
\end{proof}

\begin{bprop}[semantic warrant, per-component form]\label{prop:vc-comp}
If $P$ holds everywhere on $C$, the residual decomposes as
$D_m \smallsetminus C \subseteq \bigcup_i K_i$, and each component
satisfies $K_i \subseteq \mathrm{Off}$, then $P$ holds everywhere on
$D_m \smallsetminus \mathrm{Off}$.
\end{bprop}
\begin{proof}
The union of the components is contained in $\mathrm{Off}$ (each is),
so $D_m \smallsetminus C \subseteq \mathrm{Off}$ and
Proposition~\ref{prop:vc-sound} applies
(\texttt{verifiedCover\_components\_sound}). The component index is an
arbitrary type: no finiteness, no dimension, no ordering of components
is consumed.
\end{proof}

\begin{bremark}[each premise family is load-bearing]
Drop the cover row and any $x \in C \cap (D_m \smallsetminus
\mathrm{Off})$ with $P(x)$ false refutes the conclusion while both
set-bookkeeping premises hold. Drop the per-component containments
and an un-covered point outside the off-region escapes both cases:
$x \in D_m \smallsetminus C$ with
$x \notin \mathrm{Off}$ and $P(x)$ false refutes the conclusion while
the remaining premises hold. Drop the residual decomposition and the
component list can silently miss part of the residual --- the same
point, now lying in no $K_i$, escapes again. The exhaustiveness of the
case analysis is exactly the conjunction of the last two; the cover
row is what it decides to.
\end{bremark}

\begin{bremark}[the omitted premise --- deliberate, and in the sound direction]
The catalogued statement carries a cover-bookkeeping premise
$C \subseteq D_m$; the verified statement omits it, because the proof
never consumes it: a cover reaching outside the domain proves more
rows than needed, harmlessly. Omitting an un-consumed hypothesis makes
the mechanized statement strictly stronger, in the sound direction ---
the checker retains the containment as cover \emph{well-formedness},
which is where it belongs.
\end{bremark}

\begin{bremark}[the kit as a spectrum]
The corner $C = D_m$ (empty residual, no components) is
\rulename{QFullDomainWeaken}; the opposite corner $C = \varnothing$
demands the entire domain certified into the off-region, at which
point the scope is empty and the row vacuous. Between the corners, the
cover is the dial trading global proof effort against per-component
containment certificates --- and cover \emph{selection} is
producer-side search (one-dimensional in the current toolkit), spent
before anything ships. The verified statement is dimension-free:
richer covers change the search, never the warrant.
\end{bremark}

\begin{anchors}
Recipe \ref{rec:vc-instance} warrant & \texttt{verifiedCover\_components\_sound}\\
Proposition \ref{prop:vc-sound} & \texttt{verifiedCover\_sound} (the case split)\\
containment premises & exact Bernstein upper-bound rows (the \rulename{BernCone} section)\\
scope reading & the \texttt{hbridge} hypothesis ($\mathrm{scope} = D_m \smallsetminus \mathrm{Off}$)\\
degenerate corner & the \rulename{QFullDomainWeaken} section\\
\end{anchors}


\rulesection{62}{QPositiveCombinationRoot}{admissible rule}
\label{blcatrule-QPositiveCombinationRoot}
\ruleblock{The rule (Steele metanotation)}
\sourcenote{as realized in the mechanized development; source row: \texttt{\texttt{QTSL-rules.toml}} --- \emph{QPositiveCombinationRoot}. The metanotation statement is typeset by hand for notational quality; its mechanized ground truth is the Lean statement below, extracted from the sources at every build of this book and hash-bound. The rule \emph{roster} is machine-checked (every mechanized rule tag round-trips to exactly one catalogue row); the statement \emph{text} is editorial, checked by eye against the extracted Lean beside it.}
\begin{mathpar}
\inferrule*{
  \mathrm{child}_1 : \mathsf{QRegionSat}(\varphi_1, R_1) \\
  \mathrm{child}_2 : \mathsf{QRegionSat}(\varphi_2, R_2)
}{
  \mathcal{M},\ x_0 \in R_1 \cap R_2 \vDash \varphi_1 \wedge \varphi_2
}
\and
\inferrule*{
  \mathrm{child}_1 : \mathsf{QRegionSat}(\varphi_1, R_1) \\
  \mathrm{child}_2 : \mathsf{QRegionSat}(\varphi_2, R_2)
}{
  \mathcal{M},\ x_0 \in R_1 \cup R_2 \vDash \varphi_1 \vee \varphi_2
}
\end{mathpar}

\ruleflavor{$n$-ary forms $\bigcap_i R_i$ / $\bigcup_i R_i$ over a block family (\texttt{positiveCombination\{And,Or,Inter,Union\}\_sound}); positive connectives only — negation rides the normal-form and closure machinery, never this rule.}
\par\medskip
\rulename{QPositiveCombinationRoot} assembles a certificate for a root
boolean combination out of certificates for its blocks, at zero new
row cost: if the children conclude $\mathsf{QRegionSat}(\varphi_1,
R_1)$ and $\mathsf{QRegionSat}(\varphi_2, R_2)$, the conjunction is
certified on $R_1 \cap R_2$ and the disjunction on $R_1 \cup R_2$ ---
and the $n$-ary forms take $\bigcap_i R_i$ and $\bigcup_i R_i$ over a
whole block family. A producer uses it at the top of every
specification whose root is a positive boolean shape: each block is
certified by its own sub-derivation (a battery, an invariance
certificate, a nested combination), and the root costs two set
operations.

What the checker recomputes: the region intersection or union itself.
The semantic clauses --- that the conjunction denotes the
intersection of the denotations, the disjunction the union --- enter
as definitional bridges supplied by the specification layer, never as
posted premises.

The word \emph{positive} in the name is load-bearing.
Under-approximation is preserved by intersection and union: a region
inside each denotation is inside their intersection; a union of
regions inside denotations is inside the union of the denotations. It
is REVERSED by complement: the complement of an under-approximation
over-approximates --- the empty region under-approximates every
denotation, and its complement, the whole space, over-approximates
every complement. A negation passed through this rule would silently
flip the judgment kind\textquotesingle s one orientation guarantee.
Negation therefore never passes through it; it rides the normal-form
and closure machinery instead, where the direction swap is handled
deliberately.

\cataloguenote{Admissible, not alias: several children and two lemma directions ---
more than a single application. The node shape (a pair of binary
connective nodes versus one $n$-ary combination node) is an enactment
decision still open on the schema side; the verified lemmas cover
both forms, so neither choice waits on new mathematics.
}
\ruleblock{Lean statement}
\par\medskip
\texttt{RegionSat} is the compiled reading of the workhorse judgment
kind --- the region under-approximates the denotation --- quoted in
full because the four lemmas\textquotesingle{} meaning turns on it.
The connective denotations enter through bridge hypotheses
(\texttt{denAnd} $=$ intersection of the children\textquotesingle s
denotations, \texttt{denOr} $=$ union), supplied by the specification
layer at compile time.

\declhead{BernsteinLogicAux.Structural}{positiveCombinationAnd\_sound}{\badge{BadgeGreen}{statement bound}}
\begin{leansrc}
/-- **T11**, conjunction: intersection of under-approximations
under-approximates the conjunction. -/
theorem positiveCombinationAnd_sound (R₁ R₂ den₁ den₂ denAnd : Set X)
    (hbridge : denAnd = den₁ ∩ den₂)
    (h₁ : RegionSat R₁ den₁) (h₂ : RegionSat R₂ den₂) :
    RegionSat (R₁ ∩ R₂) denAnd
\end{leansrc}
\declhead{BernsteinLogicAux.Structural}{positiveCombinationOr\_sound}{\badge{BadgeGreen}{statement bound}}
\begin{leansrc}
/-- **T11**, disjunction: union of under-approximations under-approximates
the disjunction. -/
theorem positiveCombinationOr_sound (R₁ R₂ den₁ den₂ denOr : Set X)
    (hbridge : denOr = den₁ ∪ den₂)
    (h₁ : RegionSat R₁ den₁) (h₂ : RegionSat R₂ den₂) :
    RegionSat (R₁ ∪ R₂) denOr
\end{leansrc}
\declhead{BernsteinLogicAux.Structural}{positiveCombinationInter\_sound}{\badge{BadgeGreen}{statement bound}}
\begin{leansrc}
/-- **T11**, n-ary conjunction (the staged n-ary node shape). -/
theorem positiveCombinationInter_sound {ι : Type*} (R den : ι → Set X)
    (denAll : Set X)
    (hbridge : denAll = ⋂ i, den i)
    (h : ∀ i, RegionSat (R i) (den i)) :
    RegionSat (⋂ i, R i) denAll
\end{leansrc}
\declhead{BernsteinLogicAux.Structural}{positiveCombinationUnion\_sound}{\badge{BadgeGreen}{statement bound}}
\begin{leansrc}
/-- **T11**, n-ary disjunction. -/
theorem positiveCombinationUnion_sound {ι : Type*} (R den : ι → Set X)
    (denAll : Set X)
    (hbridge : denAll = ⋃ i, den i)
    (h : ∀ i, RegionSat (R i) (den i)) :
    RegionSat (⋃ i, R i) denAll
\end{leansrc}
\declhead{BernsteinLogicAux.Structural}{RegionSat}{\badge{BadgeGreen}{definition bound}}
\begin{leansrc}
/-- The satisfaction judgment's workhorse kind, compiled: region `R`
under-approximates denotation `den`. -/
def RegionSat (R den : Set X) : Prop := R ⊆ den
\end{leansrc}
\ruleblock{Correspondence (Steele display $\leftrightarrow$ Lean)}
\badge{BadgeGreen}{reviewer-certified against the displayed texts}
\sourcenote{The dictionary below is a reviewer-certified statement of how the mechanized statements above interpret the displayed rule: display metavariable $\mapsto$ Lean binder, each display hypothesis $\mapsto$ the premise or side condition carrying it, the display conclusion $\mapsto$ the goal --- with any deliberate difference in width between display and mechanization named outright. Certification is recorded as a hash triple over (the displayed cell, the extracted statement texts, this dictionary\textquotesingle s own text); if any of the three changes, this block turns amber until a reviewer re-certifies. The correspondence claim itself remains editorial --- reviewer-grade, not machine-checked; its machine-checked successor types each soundness statement against the rule as a structured object, and supersedes this block per rule as it lands.}
The judgment kind itself is bound by display: \texttt{RegionSat R den
= R ⊆ den} --- $\mathsf{QRegionSat}$ compiled, the formula appearing
only through its denotation (the System bridge, module \S0). The two
displayed binary rules $\mapsto$
\texttt{positiveCombinationAnd\_sound} /
\texttt{Or\_sound}: children $\mathsf{QRegionSat}(\varphi_i, R_i)$
$\mapsto$ \texttt{h₁}/\texttt{h₂ : RegionSat Rᵢ denᵢ}; conclusions at
$R_1 \cap R_2$ / $R_1 \cup R_2$ verbatim; and the semantic clauses
$[\![\varphi_1 \wedge \varphi_2]\!] = [\![\varphi_1]\!] \cap
[\![\varphi_2]\!]$ (resp.\ $\cup$) $\mapsto$ the bridge hypotheses
\texttt{hbridge : denAnd = den₁ ∩ den₂} --- supplied by the future
spec layer, never asserted in lean-aux, exactly as the roster witness
says. The $n$-ary forms (\texttt{Inter\_sound}/\texttt{Union\_sound},
arbitrary index) cover the baked node shapes (\texttt{QAndRegion}/
\texttt{QOrRegion}, each n-ary over a children list). Width note: pure
lattice monotonicity over arbitrary sets --- no formula grammar, no
positivity syntax; the ``positive connectives only'' restriction is
enforced by which bridges exist to supply (negation has no
intersection/union clause), not by any binder here --- refusal lives
in the spec layer, the certified statements are the two monotone
directions the display prices.

\ruleblock{Lean verifier receipts}
\sourcenote{module \texttt{BernsteinLogicAux.Structural}; bound build \texttt{1222cdcc326f\ldots} (bound from receipt cb0c3820c6877f22, 2026-07-29T01:03:36Z). A receipt records that named checking layers accepted a build: (module, compiled-artifact SHA-256, environment SHA-256, checker layers), minted when the checking gate ran; the store receipts attesting the bound build are listed with the layers each run credited. Nothing here asks the reader to trust the store: every layer is an independently runnable program over the quoted sources --- rebuild the module and re-run the checkers to reproduce the verdict. Receipts attesting superseded builds credit no layer and are omitted from this rendering (the count of omissions, if any, is noted below).}
\begin{receipttable}
\receiptrow{0b7e68907f45b27f}{2026-07-28T23:48:36Z}{build+decorrelated}{\badge{BadgeGreen}{current}}
\receiptrow{e97d5614d4f08c7e}{2026-07-29T00:00:35Z}{build+decorrelated}{\badge{BadgeGreen}{current}}
\receiptrow{27de21f047e6b194}{2026-07-29T00:49:05Z}{build+replay+hardened+decorrelated}{\badge{BadgeGreen}{current}}
\receiptrow{cb0c3820c6877f22}{2026-07-29T01:03:36Z}{build+replay+hardened+decorrelated}{\badge{BadgeGreen}{current}}
\receiptrow{3530841119a9a607}{2026-07-31T22:00:22Z}{build+decorrelated}{\badge{BadgeGreen}{current}}
\receiptrow{b6cc23eb3e5b7f97}{2026-07-31T22:29:16Z}{build+decorrelated}{\badge{BadgeGreen}{current}}
\receiptrow{bf3dacc3ab7f20ea}{2026-08-02T23:31:31Z}{build+decorrelated}{\badge{BadgeGreen}{current}}
\receiptrow{384a8569d2005df7}{2026-08-03T16:27:35Z}{build+decorrelated}{\badge{BadgeGreen}{current}}
\receiptrow{717e07acf9acac8a}{2026-08-03T22:46:19Z}{build+decorrelated}{\badge{BadgeGreen}{current}}
\end{receipttable}
\sourcenote{11 further receipt(s) in the store attest superseded builds of this module and are omitted from this rendering.}
\layerline{\layeritem{build}{lake build (elaboration)}{\badge{BadgeGreen}{current}}\layeritem{replay}{leanchecker (official kernel replay)}{\badge{BadgeGreen}{current}}\layeritem{hardened}{safe\_verify (hardened screen)}{\badge{BadgeGreen}{current}}\layeritem{decorrelated}{lean4export + nanoda (decorrelated kernel)}{\badge{BadgeGreen}{current}}}
\sourcenote{\badge{BadgeGray}{comparator: none} no fixture in the three-implementation comparator harness (compiled / kernel-interpreted / reference-interpreted) exercises this rule.}
\ruleblock{HOL4 statement}
\absentbox{NO HOL STATEMENT YET}{The HOL4 development\textquotesingle s object is the base
calculus\textquotesingle s wire-format checker (theory
\texttt{bern\_catalogue} ports the catalogue check-function cone).
This rule concludes a satisfaction judgment of the extension theory:
it postdates that port, and although the combination nodes are baked
in the extension theory (\texttt{QAndRegion}/\texttt{QOrRegion}), no
HOL development consumes the extension theory\textquotesingle s rows
--- there is no ported check for a HOL statement to be about. The box
becomes fillable when a checker consuming the extension\textquotesingle s
certificate rows is ported.
}
\ruleblock{HOL verifier receipts}
\absentbox{NONE}{no HOL verification receipt exists for this rule.}
\ruleblock{Proof}

\sourcenote{Genre: admissible rule (tier 2) --- a proof-tree-rewriting
recipe. Admissible and not alias by count: several children, and two
lemma directions (one per connective). The rule adds zero sequent
rows; its entire content is two one-line lattice-monotonicity facts
about under-approximation.}

The admissible template, general form, is stated once for the whole
catalogue as Recipe~\ref{rec:norm-template} of the \rulename{Norm}
section; this section instantiates it.

\begin{brecipe}[the \rulename{QPositiveCombinationRoot} instance]\label{rec:pcr-instance}
\begin{mathpar}
\inferrule*[Right=$\wedge$-form]{
  \mathcal{T}_1 \rhd\ \mathsf{QRegionSat}(\varphi_1, R_1) \\
  \mathcal{T}_2 \rhd\ \mathsf{QRegionSat}(\varphi_2, R_2)
}{
  \mathcal{M},\ x_0 \in R_1 \cap R_2 \vDash \varphi_1 \wedge \varphi_2
}
\and
\inferrule*[Right=$\vee$-form]{
  \mathcal{T}_1 \rhd\ \mathsf{QRegionSat}(\varphi_1, R_1) \\
  \mathcal{T}_2 \rhd\ \mathsf{QRegionSat}(\varphi_2, R_2)
}{
  \mathcal{M},\ x_0 \in R_1 \cup R_2 \vDash \varphi_1 \vee \varphi_2
}
\end{mathpar}
Each node rewrites to its children --- which stand unchanged --- plus
one application of the corresponding lattice-monotonicity warrant
(Propositions \ref{prop:pcr-and} and \ref{prop:pcr-or}); the $n$-ary
forms replace the pair by a family and the binary operation by
$\bigcap_i$ or $\bigcup_i$. No sequent row is added anywhere: the
whole rewrite happens at the judgment layer.
\end{brecipe}

\begin{bprop}[conjunction: intersection preserves under-approximation]\label{prop:pcr-and}
If $R_1$ under-approximates the denotation of $\varphi_1$ and $R_2$
that of $\varphi_2$, then $R_1 \cap R_2$ under-approximates the
denotation of $\varphi_1 \wedge \varphi_2$ (which is the intersection
of the two denotations, by the specification layer\textquotesingle s
definitional clause).
\end{bprop}
\begin{proof}
$x \in R_1 \cap R_2$ gives $x \in R_1 \subseteq
[\![\varphi_1]\!]$ and $x \in R_2 \subseteq [\![\varphi_2]\!]$, so
$x \in [\![\varphi_1]\!] \cap [\![\varphi_2]\!] =
[\![\varphi_1 \wedge \varphi_2]\!]$. One line
(\texttt{positiveCombinationAnd\_sound}); the final equality is the
\texttt{denAnd} bridge, supplied at compile time, never posted.
\end{proof}

\begin{bprop}[disjunction: union preserves under-approximation]\label{prop:pcr-or}
Under the same hypotheses, $R_1 \cup R_2$ under-approximates the
denotation of $\varphi_1 \vee \varphi_2$ (the union of the
denotations).
\end{bprop}
\begin{proof}
$x \in R_1 \cup R_2$ lies in some $R_j \subseteq [\![\varphi_j]\!]
\subseteq [\![\varphi_1]\!] \cup [\![\varphi_2]\!] =
[\![\varphi_1 \vee \varphi_2]\!]$. One line
(\texttt{positiveCombinationOr\_sound}), with the \texttt{denOr}
bridge closing it.
\end{proof}

\begin{bremark}[the $n$-ary forms, and the empty corners]
The $n$-ary lemmas (\texttt{positiveCombinationInter\_sound},
\texttt{positiveCombinationUnion\_sound}) quantify over an arbitrary
index type --- no finiteness is consumed --- so a single $n$-ary
combination node brings no new mathematical idea beyond the same
monotonicity principle; the $n$-ary lemmas themselves are its warrant
(an arbitrary-index family is not a fold of binary steps, which is why
they are separate statements). The
empty corners are sound instances rather than exclusions: an empty
conjunction denotes everything and $\bigcap_{i \in \varnothing} R_i$
is the whole space, which under-approximates it; an empty disjunction
denotes nothing and the empty union is empty, which
under-approximates that. Neither degenerate certificate claims
anything false, and each is worthless in its own honest way: the
disjunction corner certifies the empty region, the conjunction corner
certifies a tautology.
\end{bremark}

\begin{bremark}[the polarity fence, with the two-line counterexample]
Complement REVERSES the direction this rule preserves: from
$R \subseteq [\![\varphi]\!]$ nothing at all follows about
$\mathrm{cmpl}\,R$ under-approximating
$\mathrm{cmpl}\,[\![\varphi]\!]$ --- the containment flips. Sharpest
instance: $R = \varnothing$ under-approximates every denotation, and
its complement is the whole space, which under-approximates only the
trivial one. A negation routed through this rule would silently
convert the judgment kind\textquotesingle s under-approximation
guarantee into its opposite. The name \emph{positive} is that fence:
negation rides the normal-form rewrites and the closure discipline
(the \rulename{QClosureExact} section and its siblings),
where the direction swap is handled deliberately, one atom at a time.
\end{bremark}

\begin{bremark}[the bridge hypotheses]
The semantic clauses $[\![\varphi_1 \wedge \varphi_2]\!] =
[\![\varphi_1]\!] \cap [\![\varphi_2]\!]$ and its $\vee$/union twin
enter the verified statements as named bridge hypotheses
(\texttt{denAnd}, \texttt{denOr}), to be discharged definitionally by
the specification layer\textquotesingle s compile. The warrant
deliberately does not assert them: what the connectives denote is the
specification layer\textquotesingle s to define, and this
rule\textquotesingle s content is only that under-approximation
respects whatever lattice structure that definition induces.
\end{bremark}

\begin{anchors}
Recipe \ref{rec:pcr-instance}, $\wedge$-form & \texttt{positiveCombinationAnd\_sound}\\
Recipe \ref{rec:pcr-instance}, $\vee$-form & \texttt{positiveCombinationOr\_sound}\\
$n$-ary forms & \texttt{positiveCombinationInter\_sound}, \texttt{positiveCombinationUnion\_sound}\\
the judgment reading & \texttt{RegionSat} (quoted in full above)\\
the polarity fence\textquotesingle s other side & \texttt{compl\_over\_approx} (the \rulename{QClosureOutward} section)\\
\end{anchors}


\rulesection{63}{QClosureExact}{alias rule}
\label{blcatrule-QClosureExact}
\ruleblock{The rule (Steele metanotation)}
\sourcenote{as realized in the mechanized development; source row: \texttt{\texttt{QTSL-rules.toml}} --- \emph{QClosureExact}. The metanotation statement is typeset by hand for notational quality; its mechanized ground truth is the Lean statement below, extracted from the sources at every build of this book and hash-bound. The rule \emph{roster} is machine-checked (every mechanized rule tag round-trips to exactly one catalogue row); the statement \emph{text} is editorial, checked by eye against the extracted Lean beside it.}
\begin{mathpar}
\inferrule*{ }{ \mathrm{cmpl}\{p < c\} = \{c \le p\} }
\and
\inferrule*{ }{ \mathrm{cmpl}\{c < p\} = \{p \le c\} }
\end{mathpar}

\ruleflavor{two axioms --- deduction rules with nothing above the line, each an exact set identity, both orientations stated (\texttt{closureExactLt\_sound}/\texttt{Gt}) — no margin spent, no side condition; the closed form is a closed set for continuous observables (\texttt{closureExact\_isClosed\_ge}/\texttt{le}).}
\par\medskip
The closure discipline is the calculus\textquotesingle s answer to
strict inequalities: rows quantify over closed cells and certificates
consume closed atoms, so strict literals --- $\{p < c\}$, $\{c < p\}$
--- never reach a row. Three moves resolve them, keyed to where the
literal sits. \rulename{QClosureExact}, this rule, is the
free one: the \emph{complement} of a strict literal is already a
closed atom, exactly --- $\mathrm{cmpl}\{p < c\} = \{c \le p\}$ ---
so a complemented strict literal is replaced by its closed form with
no loss, no margin, and no side condition. The sibling moves pay:
\rulename{QClosureOutward} closes the complement of a
\emph{closed} atom by an orientation swap that over-approximates (the
sound direction at negative positions), and
\rulename{QClosureInwardEps} closes a strict literal at a
positive position by rounding inward a producer-chosen
$\varepsilon > 0$ (the sound direction there). This rule is the alias
of the trio; the other two are admissible.

What the checker recomputes: nothing --- the rewrite is a
definitional set identity, stated in both orientations because both
strict forms occur in the atom grammar. The closedness that names the
move is a fact about continuous observables: the closed form is a
preimage of a closed ray, and the calculus\textquotesingle s
observables are polynomials.

\cataloguenote{The identity spends no margin --- contrast the sibling moves, which
each pay in one currency (a boundary cut, or an $\varepsilon$-band).
Like \rulename{QFullDomainWeaken}, the alias\textquotesingle s target
is a verified identity rather than a catalogue node. The
$\varepsilon$-accounting for the inward move travels as producer
annotation data on the certificate, never as a rule.
}
\ruleblock{Lean statement}
\par\medskip
The strict atom enters through the named bridge hypothesis
(\texttt{hbridge}) fixing its compiled sublevel reading; the
closedness lemmas ask only continuity of the observable.

\declhead{BernsteinLogicAux.Structural}{closureExactLt\_sound}{\badge{BadgeGreen}{statement bound}}
\begin{leansrc}
/-- **T13 EXACT**, `<`-orientation: the complement of `{p < c}` IS the
closed atom `{c ≤ p}` — no rounding, no margin spent. -/
theorem closureExactLt_sound (p : X → ℝ) (c : ℝ) (strictAtom : Set X)
    (hbridge : strictAtom = {x | p x < c}) :
    strictAtomᶜ = {x | c ≤ p x}
\end{leansrc}
\declhead{BernsteinLogicAux.Structural}{closureExactGt\_sound}{\badge{BadgeGreen}{statement bound}}
\begin{leansrc}
/-- **T13 EXACT**, `>`-orientation. -/
theorem closureExactGt_sound (p : X → ℝ) (c : ℝ) (strictAtom : Set X)
    (hbridge : strictAtom = {x | c < p x}) :
    strictAtomᶜ = {x | p x ≤ c}
\end{leansrc}
\declhead{BernsteinLogicAux.Structural}{closureExact\_isClosed\_ge}{\badge{BadgeGreen}{statement bound}}
\begin{leansrc}
/-- The closedness that names the move: `{c ≤ p}` is closed for continuous
`p`. -/
theorem closureExact_isClosed_ge [TopologicalSpace X] {p : X → ℝ}
    (hp : Continuous p) (c : ℝ) : IsClosed {x | c ≤ p x}
\end{leansrc}
\declhead{BernsteinLogicAux.Structural}{closureExact\_isClosed\_le}{\badge{BadgeGreen}{statement bound}}
\begin{leansrc}
/-- Closedness, the other orientation: `{p ≤ c}` is closed for continuous
`p`. -/
theorem closureExact_isClosed_le [TopologicalSpace X] {p : X → ℝ}
    (hp : Continuous p) (c : ℝ) : IsClosed {x | p x ≤ c}
\end{leansrc}
\ruleblock{Correspondence (Steele display $\leftrightarrow$ Lean)}
\badge{BadgeGreen}{reviewer-certified against the displayed texts}
\sourcenote{The dictionary below is a reviewer-certified statement of how the mechanized statements above interpret the displayed rule: display metavariable $\mapsto$ Lean binder, each display hypothesis $\mapsto$ the premise or side condition carrying it, the display conclusion $\mapsto$ the goal --- with any deliberate difference in width between display and mechanization named outright. Certification is recorded as a hash triple over (the displayed cell, the extracted statement texts, this dictionary\textquotesingle s own text); if any of the three changes, this block turns amber until a reviewer re-certifies. The correspondence claim itself remains editorial --- reviewer-grade, not machine-checked; its machine-checked successor types each soundness statement against the rule as a structured object, and supersedes this block per rule as it lands.}
Transcription-close, both orientations bound separately (the recorded
one-sidedness lesson as bind structure): $\mathrm{cmpl}\{p < c\} =
\{c \le p\}$ $\mapsto$ \texttt{closureExactLt\_sound}, whose goal is
the set identity \texttt{strictAtomᶜ = \{x | c ≤ p x\}} under the
System-bridge equality \texttt{hbridge : strictAtom = \{x | p x <
c\}}; the $>$-orientation $\mapsto$ \texttt{closureExactGt\_sound}.
The closedness that names the move $\mapsto$
\texttt{closureExact\_isClosed\_ge}/\texttt{\_le}, and their binder
row carries the one consumed regularity fact: \texttt{Continuous p}
--- the observable\textquotesingle s continuity, which for polynomial
atoms is free downstream but is named here, plus the
\texttt{[TopologicalSpace X]} instance those two lemmas alone need.
Width note: \texttt{p : X → ℝ} arbitrary --- the polynomial atom
grammar is erased, safe direction; the identity pair is margin-free
and side-condition-free exactly as displayed (no hypothesis beyond
the bridge), and no completeness claim exists to state --- an
identity loses nothing by construction.

\ruleblock{Lean verifier receipts}
\sourcenote{module \texttt{BernsteinLogicAux.Structural}; bound build \texttt{1222cdcc326f\ldots} (bound from receipt cb0c3820c6877f22, 2026-07-29T01:03:36Z). A receipt records that named checking layers accepted a build: (module, compiled-artifact SHA-256, environment SHA-256, checker layers), minted when the checking gate ran; the store receipts attesting the bound build are listed with the layers each run credited. Nothing here asks the reader to trust the store: every layer is an independently runnable program over the quoted sources --- rebuild the module and re-run the checkers to reproduce the verdict. Receipts attesting superseded builds credit no layer and are omitted from this rendering (the count of omissions, if any, is noted below).}
\begin{receipttable}
\receiptrow{0b7e68907f45b27f}{2026-07-28T23:48:36Z}{build+decorrelated}{\badge{BadgeGreen}{current}}
\receiptrow{e97d5614d4f08c7e}{2026-07-29T00:00:35Z}{build+decorrelated}{\badge{BadgeGreen}{current}}
\receiptrow{27de21f047e6b194}{2026-07-29T00:49:05Z}{build+replay+hardened+decorrelated}{\badge{BadgeGreen}{current}}
\receiptrow{cb0c3820c6877f22}{2026-07-29T01:03:36Z}{build+replay+hardened+decorrelated}{\badge{BadgeGreen}{current}}
\receiptrow{3530841119a9a607}{2026-07-31T22:00:22Z}{build+decorrelated}{\badge{BadgeGreen}{current}}
\receiptrow{b6cc23eb3e5b7f97}{2026-07-31T22:29:16Z}{build+decorrelated}{\badge{BadgeGreen}{current}}
\receiptrow{bf3dacc3ab7f20ea}{2026-08-02T23:31:31Z}{build+decorrelated}{\badge{BadgeGreen}{current}}
\receiptrow{384a8569d2005df7}{2026-08-03T16:27:35Z}{build+decorrelated}{\badge{BadgeGreen}{current}}
\receiptrow{717e07acf9acac8a}{2026-08-03T22:46:19Z}{build+decorrelated}{\badge{BadgeGreen}{current}}
\end{receipttable}
\sourcenote{11 further receipt(s) in the store attest superseded builds of this module and are omitted from this rendering.}
\layerline{\layeritem{build}{lake build (elaboration)}{\badge{BadgeGreen}{current}}\layeritem{replay}{leanchecker (official kernel replay)}{\badge{BadgeGreen}{current}}\layeritem{hardened}{safe\_verify (hardened screen)}{\badge{BadgeGreen}{current}}\layeritem{decorrelated}{lean4export + nanoda (decorrelated kernel)}{\badge{BadgeGreen}{current}}}
\sourcenote{\badge{BadgeGray}{comparator: none} no fixture in the three-implementation comparator harness (compiled / kernel-interpreted / reference-interpreted) exercises this rule.}
\ruleblock{HOL4 statement}
\absentbox{NO HOL STATEMENT YET}{The HOL4 development\textquotesingle s object is the base
calculus\textquotesingle s wire-format checker (theory
\texttt{bern\_catalogue} ports the catalogue check-function cone).
The closure discipline is a specification-surface rewrite layer of
the extension theory: it postdates that port, and no row for any
of the three moves exists in the baked theory --- there is nothing
for a HOL statement to be about. The box becomes fillable when a
revision carries the closure moves\textquotesingle{} certificate rows
and the corresponding check is ported.
}
\ruleblock{HOL verifier receipts}
\absentbox{NONE}{no HOL verification receipt exists for this rule.}
\ruleblock{Proof}

\sourcenote{Genre: alias rule (tier 3) --- a respelling template, here
in the degenerate-target form: the target is a verified set identity,
not a catalogue node. This section also carries the closure
discipline\textquotesingle s shared preamble; the
\rulename{QClosureOutward} and
\rulename{QClosureInwardEps} sections build on it by
reference.}

\begin{brecipe}[the closure discipline --- shared preamble for the trio]
\label{rec:cd-preamble}
The kernel\textquotesingle s rows quantify over closed cells; its
certificates consume closed atoms; its exact arithmetic decides
non-strict comparisons. Strict literals therefore never reach a row:
every strict atom is resolved, before rows are emitted, by one of
three moves keyed to the literal\textquotesingle s position.
\begin{center}\small
\begin{tabular}{@{}lll@{}}
\textbf{literal} & \textbf{move} & \textbf{cost}\\
\midrule
complement of a strict atom & \rulename{QClosureExact}: replace by the closed form & none (set identity)\\
complement of a closed atom & \rulename{QClosureOutward}: swap orientation, over-approximate & the cut $\{p = c\}$\\
strict atom, positive position & \rulename{QClosureInwardEps}: round inward by $\varepsilon > 0$ & the band $(c - \varepsilon, c)$; mirror $(c, c + \varepsilon)$\\
\end{tabular}
\end{center}
The polarity logic: at a \emph{positive} position --- the atom
contributes to what the certified region must lie inside ---
under-approximation is sound, so shrinking the atom (the inward move)
is the safe direction. At a \emph{negative} position --- the atom
arrives complemented --- the consuming side wants the complement
over-approximated, so growing it by the boundary cut (the outward
move) is the safe direction. The exact move is direction-free: an
identity is sound everywhere. Soundness-only is the
discipline\textquotesingle s contract: no move claims to preserve
provability, and a specification tight at a critical threshold fails
honestly --- the certificate is refused --- rather than being
silently widened.
\end{brecipe}

\begin{brecipe}[the \rulename{QClosureExact} instance]
\label{rec:ce-instance}
\[
\mathrm{cmpl}\{p < c\} \;\rightsquigarrow\; \{c \le p\}
\qquad\qquad
\mathrm{cmpl}\{c < p\} \;\rightsquigarrow\; \{p \le c\}
\]
Payload transport: observable $p$ and threshold $c$ verbatim;
orientation swapped; strictness dropped; the complement consumed. One
literal $\to$ one closed atom, and the respelling is \emph{exact}:
Proposition~\ref{prop:ce-identity} is an equality, not a containment,
so the rewrite is sound at every position, positive or negative ---
the only member of the trio with that property.
\end{brecipe}

\begin{bprop}[the exactness identity]\label{prop:ce-identity}
$\mathrm{cmpl}\{x \mid p(x) < c\} = \{x \mid c \le p(x)\}$, and
symmetrically
$\mathrm{cmpl}\{x \mid c < p(x)\} = \{x \mid p(x) \le c\}$.
\end{bprop}
\begin{proof}
Pointwise, $\neg\bigl(p(x) < c\bigr) \iff c \le p(x)$ --- totality of
the order on $\R$, nothing more.
\texttt{closureExactLt\_sound} and \texttt{closureExactGt\_sound} are
this equivalence under one set-extensionality step; there is no
analytic content anywhere in the move.
\end{proof}

\begin{bprop}[the closedness that names the move]\label{prop:ce-closed}
For continuous $p$, the atoms $\{c \le p\}$ and $\{p \le c\}$ are
closed sets.
\end{bprop}
\begin{proof}
Each is the preimage under $p$ of a closed ray
(\texttt{closureExact\_isClosed\_ge}, \texttt{\_le}). The
calculus\textquotesingle s observables are polynomials, hence
continuous; the lemmas ask continuity only, so any continuous
extension of the observable grammar inherits the move unchanged.
\end{proof}

\begin{bremark}[both orientations, deliberately]
The identity is stated twice because both strict forms occur in the
atom grammar, and a one-sided statement of a complement lemma is a
standing trap for author and porter alike. Stating both
orientations costs two lines and removes the trap.
\end{bremark}

\begin{bremark}[soundness corollary, and what the siblings inherit]
Substituting equals for equals needs no direction analysis, so the
soundness corollary is immediate: any judgment consuming
$\mathrm{cmpl}\{p < c\}$ may consume $\{c \le p\}$ with identical
denotation, and no margin is spent. The sibling moves are
containments, not identities, and their soundness is
position-dependent --- which is exactly why the discipline routes
each literal to the move whose direction its position wants
(Recipe~\ref{rec:cd-preamble}).
\end{bremark}

\begin{anchors}
Proposition \ref{prop:ce-identity} & \texttt{closureExactLt\_sound}, \texttt{closureExactGt\_sound}\\
Proposition \ref{prop:ce-closed} & \texttt{closureExact\_isClosed\_ge}, \texttt{closureExact\_isClosed\_le}\\
strict-atom reading & the \texttt{hbridge} hypothesis of the identity lemmas\\
the sibling moves & the \rulename{QClosureOutward} and \rulename{QClosureInwardEps} sections\\
\end{anchors}


\rulesection{64}{QClosureOutward}{admissible rule}
\label{blcatrule-QClosureOutward}
\ruleblock{The rule (Steele metanotation)}
\sourcenote{as realized in the mechanized development; source row: \texttt{\texttt{QTSL-rules.toml}} --- \emph{QClosureOutward}. The metanotation statement is typeset by hand for notational quality; its mechanized ground truth is the Lean statement below, extracted from the sources at every build of this book and hash-bound. The rule \emph{roster} is machine-checked (every mechanized rule tag round-trips to exactly one catalogue row); the statement \emph{text} is editorial, checked by eye against the extracted Lean beside it.}
\begin{mathpar}
\inferrule*{ }{ \mathrm{cmpl}\{p \le c\} \subseteq \{c \le p\} }
\and
\inferrule*{ }{ \mathrm{cmpl}\{c \le p\} \subseteq \{p \le c\} }
\and
\inferrule*{
  R \subseteq [\![\varphi]\!]
}{
  \mathrm{cmpl}\,[\![\varphi]\!] \subseteq \mathrm{cmpl}\,R
}
\end{mathpar}

\ruleflavor{the closed swap over-approximates the true complement, losing only the cut $\{p = c\}$ — the direction negative positions consume; atom face both orientations (\texttt{closureOutward\_sound}/\texttt{Ge}) $+$ the antitone region face (\texttt{compl\_over\_approx}).}
\par\medskip
\rulename{QClosureOutward} is the closure
discipline\textquotesingle s move for complements of \emph{closed}
atoms (the discipline itself is set out in the
\rulename{QClosureExact} section). The true complement of
$\{p \le c\}$ is the strict set $\{c < p\}$ --- not consumable by a
row --- and the move consumes the orientation swap $\{c \le p\}$
instead, which contains it: an over-approximation of the true
complement, losing only the cut $\{p = c\}$. Over-approximation is
exactly what the consuming side wants at negative positions: there
the atom describes a set to be excluded or avoided, and a certificate
against a superset is a fortiori a certificate against the set.

The warrant has two faces, both verified. The \emph{atom} face is the
pointwise containment above, in both orientations. The \emph{region}
face is the antitone kernel: a certified region under-approximates
its denotation, so its complement over-approximates the
denotation\textquotesingle s complement --- the direction transport
that lets a negative position consume a child
certificate\textquotesingle s complement soundly. What the checker
recomputes: the orientation swap itself, which is syntactic; the
containments are the verified lemmas\textquotesingle{} content, spent
once, globally.

The move is one-directional by design: at a \emph{positive} position
the same swap would overclaim --- take $p$ constant at $c$: the swap
$\{c \le p\}$ is the whole space while the true complement of
$\{p \le c\}$ is empty, the maximal over-claim --- so the
discipline\textquotesingle s routing table, not a side condition,
keeps this move on the negative side.

\cataloguenote{The conditional flag this entry once carried recorded the composite
normal-form induction that pushes complements through connectives ---
deliberately outside this rule\textquotesingle s own lemmas, since
that induction is stated once, for all consumers, with the
normal-form rewriting rules. It has since been mechanized where it
lives, so the row carries no conditional flag: \texttt{denot\_nnf}
(\texttt{BernsteinLogicAux.NNFDuality}, bound at the
\rulename{QStreettRequestNnf} entry) is the denotation-preserving
normal-form pass whose induction cases are exactly the
\rulename{QNnfRewrite} axioms, with \texttt{NForm.elimNeg} /
\texttt{denot\_elimNeg} the consumption-boundary elimination. A reader
tracing a complement through a full specification follows this rule
for the atom step and the normal-form sections for the traversal; the
composition-as-applied stays producer-side, which is why the traversal
is cited here, not bound.
}
\ruleblock{Lean statement}
\par\medskip
\texttt{RegionSat} (quoted in full) is the compiled reading of the
workhorse judgment kind, consumed by the region face; the atom faces
fix the consumed closed swap through the bridge hypothesis
(\texttt{hbridge}).

\declhead{BernsteinLogicAux.Structural}{closureOutward\_sound}{\badge{BadgeGreen}{statement bound}}
\begin{leansrc}
/-- **T13 OUTWARD**, `≤`-atom: the consumed closed swap `{c ≤ p}` contains
the true complement of `{p ≤ c}`. No margin is spent (the containment loses
only the measure-zero cut `{p = c}`). -/
theorem closureOutward_sound (p : X → ℝ) (c : ℝ) (cmplAtom : Set X)
    (hbridge : cmplAtom = {x | c ≤ p x}) :
    {x | p x ≤ c}ᶜ ⊆ cmplAtom
\end{leansrc}
\declhead{BernsteinLogicAux.Structural}{closureOutwardGe\_sound}{\badge{BadgeGreen}{statement bound}}
\begin{leansrc}
/-- **T13 OUTWARD**, `≥`-atom orientation. -/
theorem closureOutwardGe_sound (p : X → ℝ) (c : ℝ) (cmplAtom : Set X)
    (hbridge : cmplAtom = {x | p x ≤ c}) :
    {x | c ≤ p x}ᶜ ⊆ cmplAtom
\end{leansrc}
\declhead{BernsteinLogicAux.Structural}{compl\_over\_approx}{\badge{BadgeGreen}{statement bound}}
\begin{leansrc}
/-- The antitone region-complement kernel (region face): an
under-approximating certified region's complement OVER-approximates the
semantic complement — the direction negative positions consume. -/
theorem compl_over_approx {R den : Set X} (h : RegionSat R den) :
    denᶜ ⊆ Rᶜ
\end{leansrc}
\declhead{BernsteinLogicAux.Structural}{RegionSat}{\badge{BadgeGreen}{definition bound}}
\begin{leansrc}
/-- The satisfaction judgment's workhorse kind, compiled: region `R`
under-approximates denotation `den`. -/
def RegionSat (R den : Set X) : Prop := R ⊆ den
\end{leansrc}
\ruleblock{Correspondence (Steele display $\leftrightarrow$ Lean)}
\badge{BadgeGreen}{reviewer-certified against the displayed texts}
\sourcenote{The dictionary below is a reviewer-certified statement of how the mechanized statements above interpret the displayed rule: display metavariable $\mapsto$ Lean binder, each display hypothesis $\mapsto$ the premise or side condition carrying it, the display conclusion $\mapsto$ the goal --- with any deliberate difference in width between display and mechanization named outright. Certification is recorded as a hash triple over (the displayed cell, the extracted statement texts, this dictionary\textquotesingle s own text); if any of the three changes, this block turns amber until a reviewer re-certifies. The correspondence claim itself remains editorial --- reviewer-grade, not machine-checked; its machine-checked successor types each soundness statement against the rule as a structured object, and supersedes this block per rule as it lands.}
Three displayed clauses, three bound statements. The atom face, both
orientations: $\mathrm{cmpl}\{p \le c\} \subseteq \{c \le p\}$
$\mapsto$ \texttt{closureOutward\_sound} (goal \texttt{\{x | p x ≤
c\}ᶜ ⊆ cmplAtom} under the bridge \texttt{cmplAtom = \{x | c ≤ p
x\}}), mirrored by \texttt{closureOutwardGe\_sound} --- containments,
not identities: the flavor text\textquotesingle s ``losing only the
cut $\{p = c\}$'' is
commentary on the gap, and the certified statement claims exactly the
over-approximation direction, nothing finer. The region face $R
\subseteq [\![\varphi]\!] \Rightarrow \mathrm{cmpl}[\![\varphi]\!]
\subseteq \mathrm{cmpl}\,R$ $\mapsto$ \texttt{compl\_over\_approx :
denᶜ ⊆ Rᶜ} from \texttt{RegionSat R den} --- the antitone kernel,
the region face of the closure discipline, two tokens of proof. Width
note: this row once carried a conditional mark for a composite
residue --- the De Morgan \emph{traversal} of a full specification ---
which is discharged: the traversal lives with the
QNnfRewrite/QStreettRequestNnf normal-form machinery, deliberately not
here, and is machine-checked there (\texttt{denot\_nnf} +
\texttt{denot\_elimNeg}, cross-cited in the note), so the row carries
no conditional flag and what remains is only that division of binding
responsibility. The certified atom + region faces are unconditional
set-level facts over arbitrary observables and regions; no topology,
no continuity, no polynomiality is consumed anywhere in this bound
set.

\ruleblock{Lean verifier receipts}
\sourcenote{module \texttt{BernsteinLogicAux.Structural}; bound build \texttt{1222cdcc326f\ldots} (bound from receipt cb0c3820c6877f22, 2026-07-29T01:03:36Z). A receipt records that named checking layers accepted a build: (module, compiled-artifact SHA-256, environment SHA-256, checker layers), minted when the checking gate ran; the store receipts attesting the bound build are listed with the layers each run credited. Nothing here asks the reader to trust the store: every layer is an independently runnable program over the quoted sources --- rebuild the module and re-run the checkers to reproduce the verdict. Receipts attesting superseded builds credit no layer and are omitted from this rendering (the count of omissions, if any, is noted below).}
\begin{receipttable}
\receiptrow{0b7e68907f45b27f}{2026-07-28T23:48:36Z}{build+decorrelated}{\badge{BadgeGreen}{current}}
\receiptrow{e97d5614d4f08c7e}{2026-07-29T00:00:35Z}{build+decorrelated}{\badge{BadgeGreen}{current}}
\receiptrow{27de21f047e6b194}{2026-07-29T00:49:05Z}{build+replay+hardened+decorrelated}{\badge{BadgeGreen}{current}}
\receiptrow{cb0c3820c6877f22}{2026-07-29T01:03:36Z}{build+replay+hardened+decorrelated}{\badge{BadgeGreen}{current}}
\receiptrow{3530841119a9a607}{2026-07-31T22:00:22Z}{build+decorrelated}{\badge{BadgeGreen}{current}}
\receiptrow{b6cc23eb3e5b7f97}{2026-07-31T22:29:16Z}{build+decorrelated}{\badge{BadgeGreen}{current}}
\receiptrow{bf3dacc3ab7f20ea}{2026-08-02T23:31:31Z}{build+decorrelated}{\badge{BadgeGreen}{current}}
\receiptrow{384a8569d2005df7}{2026-08-03T16:27:35Z}{build+decorrelated}{\badge{BadgeGreen}{current}}
\receiptrow{717e07acf9acac8a}{2026-08-03T22:46:19Z}{build+decorrelated}{\badge{BadgeGreen}{current}}
\end{receipttable}
\sourcenote{11 further receipt(s) in the store attest superseded builds of this module and are omitted from this rendering.}
\layerline{\layeritem{build}{lake build (elaboration)}{\badge{BadgeGreen}{current}}\layeritem{replay}{leanchecker (official kernel replay)}{\badge{BadgeGreen}{current}}\layeritem{hardened}{safe\_verify (hardened screen)}{\badge{BadgeGreen}{current}}\layeritem{decorrelated}{lean4export + nanoda (decorrelated kernel)}{\badge{BadgeGreen}{current}}}
\sourcenote{\badge{BadgeGray}{comparator: none} no fixture in the three-implementation comparator harness (compiled / kernel-interpreted / reference-interpreted) exercises this rule.}
\ruleblock{HOL4 statement}
\absentbox{NO HOL STATEMENT YET}{The HOL4 development\textquotesingle s object is the base
calculus\textquotesingle s wire-format checker (theory
\texttt{bern\_catalogue} ports the catalogue check-function cone).
The closure discipline is a specification-surface rewrite layer of
the extension theory: it postdates that port, and no row for any
of the three moves exists in the baked theory --- there is nothing
for a HOL statement to be about. The box becomes fillable when a
revision carries the closure moves\textquotesingle{} certificate rows
and the corresponding check is ported.
}
\ruleblock{HOL verifier receipts}
\absentbox{NONE}{no HOL verification receipt exists for this rule.}
\ruleblock{Proof}

\sourcenote{Genre: admissible rule (tier 2) --- a proof-tree-rewriting
recipe. The closure discipline this move belongs to --- why strict
literals never reach a row, and how the three moves divide the
positions --- is set out in full in the
\rulename{QClosureExact} section; this section states only
what is native to the outward move: the two verified faces of its
warrant, and the direction analysis that confines it to negative
positions.}

The admissible template, general form, is stated once for the whole
catalogue as Recipe~\ref{rec:norm-template} of the \rulename{Norm}
section; this section instantiates it.

\begin{brecipe}[the \rulename{QClosureOutward} instance --- two faces]
\label{rec:co-instance}
\emph{Atom face.} At a negative position, the complement of a closed
atom is consumed as its orientation swap:
\[
\mathrm{cmpl}\{p \le c\} \;\rightsquigarrow\; \{c \le p\}
\qquad\qquad
\mathrm{cmpl}\{c \le p\} \;\rightsquigarrow\; \{p \le c\}
\]
warranted by the containment
$\mathrm{cmpl}\{p \le c\} \subseteq \{c \le p\}$
(Proposition~\ref{prop:co-atom}): the consumed atom
\emph{over-approximates} the true complement, losing only the cut
$\{p = c\}$.

\emph{Region face.} A negative position consuming the complement of a
\emph{certified region} $R$ (a child certificate\textquotesingle s
conclusion, not a syntactic atom) is warranted by the antitone kernel
(Proposition~\ref{prop:co-region}): certified under-approximation of
the denotation becomes over-approximation of the
denotation\textquotesingle s complement. Both faces discharge against
their verified containment; the orientation swap itself is syntactic
and recomputed by the checker.
\end{brecipe}

\begin{bprop}[atom face]\label{prop:co-atom}
$\mathrm{cmpl}\{x \mid p(x) \le c\} \subseteq \{x \mid c \le p(x)\}$,
and symmetrically
$\mathrm{cmpl}\{x \mid c \le p(x)\} \subseteq \{x \mid p(x) \le c\}$.
\end{bprop}
\begin{proof}
If $\neg(p(x) \le c)$ then $c < p(x)$, hence $c \le p(x)$ --- the
strict fact weakened to its non-strict form
(\texttt{closureOutward\_sound}; the mirror is
\texttt{closureOutwardGe\_sound}). The two sides differ exactly by
the cut $\{p = c\}$, which the swap adjoins.
\end{proof}

\begin{bprop}[region face --- the antitone kernel]\label{prop:co-region}
If a certified region under-approximates a denotation,
$R \subseteq [\![\varphi]\!]$, then
$\mathrm{cmpl}\,[\![\varphi]\!] \subseteq \mathrm{cmpl}\,R$: the
complement of the certificate over-approximates the complement of the
denotation.
\end{bprop}
\begin{proof}
Contrapositive, one line: $x \notin \mathrm{cmpl}\,R$ means
$x \in R \subseteq [\![\varphi]\!]$, so
$x \notin \mathrm{cmpl}\,[\![\varphi]\!]$
(\texttt{compl\_over\_approx}). This is the direction transport that
lets negative positions consume child certificates at all: the
judgment kind fixes one orientation, and complementation flips it in
precisely the way the consuming side wants.
\end{proof}

\begin{bremark}[what is lost, and why the direction is right]
Only the cut $\{p = c\}$ --- and in the sound direction. At a
negative position the atom names a set to be excluded or avoided;
certifying against a superset of it is a fortiori certifying against
it. The move never claims the cut is negligible --- for a nondegenerate
\emph{polynomial} observable ($p - c$ not identically zero on a
full-dimensional cell) it is a
measure-zero level set, but soundness nowhere consumes that: the
containment alone carries the warrant, so even a degenerate observable
whose cut is fat (at the extreme $p \equiv c$, the whole cell) only
makes the certificate more conservative, never wrong.
\end{bremark}

\begin{bremark}[one-directional by design]
At a \emph{positive} position the same swap would overclaim. Take $p$
constant at $c$: the swap $\{c \le p\}$ is the whole space, while the
true complement of $\{p \le c\}$ is empty --- the maximal over-claim.
Nothing in the move\textquotesingle s own premises excludes this; the
closure discipline\textquotesingle s routing table
(the preamble, \rulename{QClosureExact} section) is the
guard --- the recipe\textquotesingle s ``at a negative position''
names that routing assignment, not a premise the move checks for
itself --- sending positive-position strict literals to the inward
move instead. A porter implementing the discipline ports the routing with
it, or re-learns this remark the hard way.
\end{bremark}

\begin{bremark}[what is deliberately not warranted here]
The verified faces are the atom step and the region step. The
composite induction --- pushing complements through conjunctions,
disjunctions, and nested operators until they reach atoms --- is
stated once, for all consumers, with the normal-form rewriting rules
--- this rule\textquotesingle s warrant is deliberately partial in
exactly that sense: a complement traversing a full specification
rides the normal-form sections for the traversal and this rule for
the atoms it arrives at.
\end{bremark}

\begin{anchors}
Proposition \ref{prop:co-atom} & \texttt{closureOutward\_sound}, \texttt{closureOutwardGe\_sound}\\
Proposition \ref{prop:co-region} & \texttt{compl\_over\_approx}\\
the judgment reading & \texttt{RegionSat} (quoted in full above)\\
the routing guard & the closure-discipline preamble (\rulename{QClosureExact} section)\\
the composite traversal & the normal-form rewriting sections (deliberately not this rule)\\
\end{anchors}


\rulesection{65}{QClosureInwardEps}{admissible rule}
\label{blcatrule-QClosureInwardEps}
\ruleblock{The rule (Steele metanotation)}
\sourcenote{as realized in the mechanized development; source row: \texttt{\texttt{QTSL-rules.toml}} --- \emph{QClosureInwardEps}. The metanotation statement is typeset by hand for notational quality; its mechanized ground truth is the Lean statement below, extracted from the sources at every build of this book and hash-bound. The rule \emph{roster} is machine-checked (every mechanized rule tag round-trips to exactly one catalogue row); the statement \emph{text} is editorial, checked by eye against the extracted Lean beside it.}
\begin{mathpar}
\inferrule*{
  \varepsilon \in \Q_{>0}
}{
  \{p \le c - \varepsilon\} \subseteq \{p < c\}
}
\and
\inferrule*{
  \varepsilon \in \Q_{>0}
}{
  \{c + \varepsilon \le p\} \subseteq \{c < p\}
}
\end{mathpar}

\ruleflavor{the producer-chosen $\varepsilon$ elaborates a strict atom at a positive position to its rounded closed form ($\{p < c\} \rightsquigarrow \{p \le c - \varepsilon\}$, $\{c < p\} \rightsquigarrow \{c + \varepsilon \le p\}$), and the displayed containments are why that is sound: the rounded closed atom strengthens the claim (\texttt{closureInwardEpsLt\_sound}/\texttt{Gt}, any real $\varepsilon > 0$ — rationality subsumed by cast); soundness-only, no completeness preservation claimed — the honest failure at a critical threshold is refusal. Closedness rider: \texttt{closureInwardEps\_isClosed}.}
\par\medskip
\rulename{QClosureInwardEps} is the closure
discipline\textquotesingle s move for strict literals at positive
positions (the discipline is set out in the
\rulename{QClosureExact} section): the literal $\{p < c\}$
is replaced by the closed atom $\{p \le c - \varepsilon\}$ for a
producer-chosen rational $\varepsilon > 0$. The rounded atom is
contained in the literal, so at a position where the certified region
must lie \emph{inside} the denotation, the substitution strengthens
the obligation --- sound by construction --- and hands the kernel a
closed atom it can consume.

The $\varepsilon$ is real spend: the band between $c - \varepsilon$
and $c$ is surrendered (between $c$ and $c + \varepsilon$ in the
mirror orientation). The choice of $\varepsilon$ lives in the
producer toolkit --- searchable, like a multiplier or a cover --- and
its accounting travels as annotation data on the certificate, not as
a rule. What the checker recomputes: nothing beyond the rewritten
atom; the containment is the verified lemma, spent once, globally.

The contract is soundness-only, and the honest failure mode is
refusal: a specification exactly tight at its threshold --- the
intended witness satisfies $p < c$ everywhere, but $p$ approaches $c$
on the region --- admits no $\varepsilon$, and the certificate is
refused rather than silently widened. No completeness preservation is
claimed, deliberately.

\cataloguenote{The verified statement takes any real $\varepsilon > 0$; the
producer\textquotesingle s $\varepsilon$ being rational is a
representation fact about certificates, subsumed by the cast ---
soundness for every positive real is a fortiori soundness for every
positive rational. The closedness of the rounded atom is an instance
of the exact move\textquotesingle s closedness lemma at the shifted
threshold.
}
\ruleblock{Lean statement}
\par\medskip
Both orientations of the rounding, plus the closedness rider; the
rounded atom enters through the bridge hypothesis (\texttt{hbridge})
fixing its compiled sublevel reading, and $\varepsilon$ enters as a
real with $0 < \varepsilon$ (rationality of the
producer\textquotesingle s choice is subsumed by the cast).

\declhead{BernsteinLogicAux.Structural}{closureInwardEpsLt\_sound}{\badge{BadgeGreen}{statement bound}}
\begin{leansrc}
/-- **T13 INWARD-EPS**, `<`-orientation: `{p ≤ c - ε} ⊆ {p < c}` for any
`ε > 0`. Substituting the rounded closed atom at a positive position
strengthens the certified claim. -/
theorem closureInwardEpsLt_sound (p : X → ℝ) (c ε : ℝ) (hε : 0 < ε)
    (roundedAtom : Set X)
    (hbridge : roundedAtom = {x | p x ≤ c - ε}) :
    roundedAtom ⊆ {x | p x < c}
\end{leansrc}
\declhead{BernsteinLogicAux.Structural}{closureInwardEpsGt\_sound}{\badge{BadgeGreen}{statement bound}}
\begin{leansrc}
/-- **T13 INWARD-EPS**, `>`-orientation: `{c + ε ≤ p} ⊆ {c < p}`. -/
theorem closureInwardEpsGt_sound (p : X → ℝ) (c ε : ℝ) (hε : 0 < ε)
    (roundedAtom : Set X)
    (hbridge : roundedAtom = {x | c + ε ≤ p x}) :
    roundedAtom ⊆ {x | c < p x}
\end{leansrc}
\declhead{BernsteinLogicAux.Structural}{closureInwardEps\_isClosed}{\badge{BadgeGreen}{statement bound}}
\begin{leansrc}
/-- The inward-rounding closedness rider: the rounded atom is closed
(instance of the EXACT move's closedness at the shifted threshold). -/
theorem closureInwardEps_isClosed [TopologicalSpace X] {p : X → ℝ}
    (hp : Continuous p) (c ε : ℝ) : IsClosed {x | p x ≤ c - ε}
\end{leansrc}
\ruleblock{Correspondence (Steele display $\leftrightarrow$ Lean)}
\badge{BadgeGreen}{reviewer-certified against the displayed texts}
\sourcenote{The dictionary below is a reviewer-certified statement of how the mechanized statements above interpret the displayed rule: display metavariable $\mapsto$ Lean binder, each display hypothesis $\mapsto$ the premise or side condition carrying it, the display conclusion $\mapsto$ the goal --- with any deliberate difference in width between display and mechanization named outright. Certification is recorded as a hash triple over (the displayed cell, the extracted statement texts, this dictionary\textquotesingle s own text); if any of the three changes, this block turns amber until a reviewer re-certifies. The correspondence claim itself remains editorial --- reviewer-grade, not machine-checked; its machine-checked successor types each soundness statement against the rule as a structured object, and supersedes this block per rule as it lands.}
Both displayed containments bound, orientation for orientation:
$\{p \le c - \varepsilon\}
\subseteq \{p < c\}$ $\mapsto$ \texttt{closureInwardEpsLt\_sound}
(the rounded atom entering through the bridge \texttt{roundedAtom =
\{x | p x ≤ c − ε\}}, the strict margin as \texttt{hε : 0 < ε}), and
the $>$-orientation $\mapsto$ \texttt{closureInwardEpsGt\_sound};
the inward-rounding closedness rider (the rounded atom is closed)
$\mapsto$
\texttt{closureInwardEps\_isClosed} (continuity consumed there
alone, the QClosureExact pattern). Two width facts, one per direction of
the display: the certified statements take \emph{any real}
$\varepsilon > 0$ where the display asks $\varepsilon \in \Q_{>0}$
--- producer rationality subsumed by cast, divergence D3, statement
wider in the safe direction; and no completeness-preservation
statement exists in the module --- the flavor text\textquotesingle s
``soundness-only'' is certified by absence, with total loss at a
critical cut remaining the honest refusal mode, exactly as the
ledger inherits it.

\ruleblock{Lean verifier receipts}
\sourcenote{module \texttt{BernsteinLogicAux.Structural}; bound build \texttt{1222cdcc326f\ldots} (bound from receipt cb0c3820c6877f22, 2026-07-29T01:03:36Z). A receipt records that named checking layers accepted a build: (module, compiled-artifact SHA-256, environment SHA-256, checker layers), minted when the checking gate ran; the store receipts attesting the bound build are listed with the layers each run credited. Nothing here asks the reader to trust the store: every layer is an independently runnable program over the quoted sources --- rebuild the module and re-run the checkers to reproduce the verdict. Receipts attesting superseded builds credit no layer and are omitted from this rendering (the count of omissions, if any, is noted below).}
\begin{receipttable}
\receiptrow{0b7e68907f45b27f}{2026-07-28T23:48:36Z}{build+decorrelated}{\badge{BadgeGreen}{current}}
\receiptrow{e97d5614d4f08c7e}{2026-07-29T00:00:35Z}{build+decorrelated}{\badge{BadgeGreen}{current}}
\receiptrow{27de21f047e6b194}{2026-07-29T00:49:05Z}{build+replay+hardened+decorrelated}{\badge{BadgeGreen}{current}}
\receiptrow{cb0c3820c6877f22}{2026-07-29T01:03:36Z}{build+replay+hardened+decorrelated}{\badge{BadgeGreen}{current}}
\receiptrow{3530841119a9a607}{2026-07-31T22:00:22Z}{build+decorrelated}{\badge{BadgeGreen}{current}}
\receiptrow{b6cc23eb3e5b7f97}{2026-07-31T22:29:16Z}{build+decorrelated}{\badge{BadgeGreen}{current}}
\receiptrow{bf3dacc3ab7f20ea}{2026-08-02T23:31:31Z}{build+decorrelated}{\badge{BadgeGreen}{current}}
\receiptrow{384a8569d2005df7}{2026-08-03T16:27:35Z}{build+decorrelated}{\badge{BadgeGreen}{current}}
\receiptrow{717e07acf9acac8a}{2026-08-03T22:46:19Z}{build+decorrelated}{\badge{BadgeGreen}{current}}
\end{receipttable}
\sourcenote{11 further receipt(s) in the store attest superseded builds of this module and are omitted from this rendering.}
\layerline{\layeritem{build}{lake build (elaboration)}{\badge{BadgeGreen}{current}}\layeritem{replay}{leanchecker (official kernel replay)}{\badge{BadgeGreen}{current}}\layeritem{hardened}{safe\_verify (hardened screen)}{\badge{BadgeGreen}{current}}\layeritem{decorrelated}{lean4export + nanoda (decorrelated kernel)}{\badge{BadgeGreen}{current}}}
\sourcenote{\badge{BadgeGray}{comparator: none} no fixture in the three-implementation comparator harness (compiled / kernel-interpreted / reference-interpreted) exercises this rule.}
\ruleblock{HOL4 statement}
\absentbox{NO HOL STATEMENT YET}{The HOL4 development\textquotesingle s object is the base
calculus\textquotesingle s wire-format checker (theory
\texttt{bern\_catalogue} ports the catalogue check-function cone).
The closure discipline is a specification-surface rewrite layer of
the extension theory: it postdates that port, and no row for any
of the three moves exists in the baked theory --- there is nothing
for a HOL statement to be about. The box becomes fillable when a
revision carries the closure moves\textquotesingle{} certificate rows
and the corresponding check is ported.
}
\ruleblock{HOL verifier receipts}
\absentbox{NONE}{no HOL verification receipt exists for this rule.}
\ruleblock{Proof}

\sourcenote{Genre: admissible rule (tier 2) --- a proof-tree-rewriting
recipe. The closure discipline this move belongs to is set out in full
in the \rulename{QClosureExact} section; this section states
what is native to the inward move: the rounding containment in both
orientations, the closedness rider, and the two honesty clauses ---
strictness of $\varepsilon > 0$, and the deliberate absence of any
completeness claim.}

The admissible template, general form, is stated once for the whole
catalogue as Recipe~\ref{rec:norm-template} of the \rulename{Norm}
section; this section instantiates it.

\begin{brecipe}[the \rulename{QClosureInwardEps} instance]
\label{rec:cie-instance}
At a positive position, a strict literal rounds inward by a
producer-chosen $\varepsilon > 0$:
\[
\{p < c\} \;\rightsquigarrow\; \{p \le c - \varepsilon\}
\qquad\qquad
\{c < p\} \;\rightsquigarrow\; \{c + \varepsilon \le p\}
\]
warranted by the containment
$\{p \le c - \varepsilon\} \subseteq \{p < c\}$
(Proposition~\ref{prop:cie-sound}): the rounded closed atom
\emph{strengthens} the certified claim, which is the sound direction
where the certified region must lie inside the denotation. Payload
transport: observable and threshold verbatim, threshold shifted by
$\varepsilon$, strictness dropped; the $\varepsilon$ itself travels
as certificate annotation data, never as a rule premise.
\end{brecipe}

\begin{bprop}[the rounding containment]\label{prop:cie-sound}
For every $\varepsilon > 0$:
$\{x \mid p(x) \le c - \varepsilon\} \subseteq \{x \mid p(x) < c\}$,
and symmetrically
$\{x \mid c + \varepsilon \le p(x)\} \subseteq \{x \mid c < p(x)\}$.
\end{bprop}
\begin{proof}
$p(x) \le c - \varepsilon < c$, the second inequality because
$\varepsilon > 0$ --- one line of linear arithmetic
(\texttt{closureInwardEpsLt\_sound}; mirror
\texttt{closureInwardEpsGt\_sound}). The verified statements take any
\emph{real} $\varepsilon > 0$; the producer\textquotesingle s
$\varepsilon$ being rational is a representation fact about
certificates, subsumed by the cast.
\end{proof}

\begin{bprop}[closedness rider]\label{prop:cie-closed}
For continuous $p$, both rounded atoms are closed sets:
$\{p \le c - \varepsilon\}$ and $\{c + \varepsilon \le p\}$.
\end{bprop}
\begin{proof}
An instance of the exact move\textquotesingle s closedness at the
shifted threshold $c - \varepsilon$
(\texttt{closureInwardEps\_isClosed}, one application of
\texttt{closureExact\_isClosed\_le}); the mirrored orientation is
$\{c + \varepsilon \le p\} = p^{-1}\bigl([\,c + \varepsilon,
\infty)\bigr)$, closed by the same continuity --- the exact
move\textquotesingle s $\ge$-orientation closedness at threshold
$c + \varepsilon$ (the cited binding covers the $\le$ orientation;
the mirror discharges against the $\ge$ twin). The move thus
delivers exactly what the discipline asked of it: a closed atom the
kernel can consume.
\end{proof}

\begin{bremark}[$\varepsilon > 0$ strictly --- the boundary counterexample]
At $\varepsilon = 0$ the containment fails: $\{p \le c\}$ is not
contained in $\{p < c\}$, and the failure can be total --- take $p$
constant at $c$: the un-rounded atom is the whole space while the
literal is empty. The strict positivity of $\varepsilon$ is therefore
the whole rule, in the same sense that the sign row is half of
\rulename{HypMultiplier}: remove it and the move manufactures the
maximal over-claim from nothing.
\end{bremark}

\begin{bremark}[soundness-only, and the refusal posture]
No completeness preservation is claimed, deliberately. The
$\varepsilon$-band is surrendered --- $(c - \varepsilon, c)$ in the
first orientation, $(c, c + \varepsilon)$ in the mirror --- and a
specification exactly tight at its threshold --- the intended witness
satisfies $p < c$ everywhere while $p$ approaches $c$ on the region
--- admits no $\varepsilon$ at all: every rounding excludes points the
witness needs. The honest outcome is refusal (no certificate exists
through this move), never silent widening; a producer meeting refusal
here either finds slack elsewhere in the specification or reports the
tightness, and the discipline\textquotesingle s contract is that the
checker is never the party doing the forgiving.
\end{bremark}

\begin{bremark}[why not under a complement]
Rounding an atom that will be complemented is doubly wrong as a
discharge of \emph{this} obligation: the complement of the rounded
atom, $\{c - \varepsilon < p\}$, is strict again --- the closedness
goal is defeated, and the $\varepsilon$ was spent for nothing.
Complemented literals belong to the exact move (free, when the inner
literal is strict) or the outward move (when the inner atom is
closed); the routing table in the closure-discipline preamble
(\rulename{QClosureExact} section) assigns every case
exactly one move, and the assignment is what makes the trio a
discipline rather than three tricks.
\end{bremark}

\begin{anchors}
Proposition \ref{prop:cie-sound} & \texttt{closureInwardEpsLt\_sound}, \texttt{closureInwardEpsGt\_sound}\\
Proposition \ref{prop:cie-closed} & \texttt{closureInwardEps\_isClosed} (via \texttt{closureExact\_isClosed\_le})\\
rounded-atom reading & the \texttt{hbridge} hypothesis of the rounding lemmas\\
$\varepsilon$ accounting & certificate annotation data, producer-side (not a rule)\\
the routing guard & the closure-discipline preamble (\rulename{QClosureExact} section)\\
\end{anchors}


\rulesection{66}{TaylorBound}{alias for \textcolor{DiffBlue}{\rulename{PConstExact}}}
\label{blcatrule-TaylorBound}
\ruleblock{The rule (Steele metanotation)}
\sourcenote{as realized in the mechanized development; source row: \texttt{\texttt{QTSL-rules.toml}} --- \emph{TaylorBound}. The metanotation statement is typeset by hand for notational quality; its mechanized ground truth is the Lean statement below, extracted from the sources at every build of this book and hash-bound. The rule \emph{roster} is machine-checked (every mechanized rule tag round-trips to exactly one catalogue row); the statement \emph{text} is editorial, checked by eye against the extracted Lean beside it.}
\begin{mathpar}
\inferrule*[Right=\rulename{TaylorBound}]{
  q \in \Q; \quad \theta \in \Q_{>0}; \quad T \in \N; \quad n \in \N \\
  q \le \textstyle\sum_{k \le n} (\theta T)^k / k! \ \text{holds under exact evaluation in } \Q
}{
  \Phi \vdash q \le \textstyle\sum_{k \le n} (\theta T)^k / k!
}
\end{mathpar}

\ruleflavor{one \rulename{PConstExact} application, displayed at the deadline battery\textquotesingle s Taylor comparison — every term is wire-free (closed), so the comparison is decided by the verified closed evaluator; no QTSL sort exists and none is needed. The transcendental fact $\sum_{k \le n} (\theta T)^k / k! \le e^{\theta T}$ belongs to the consuming lemma, never the checker: the checker never needs to know $e$ exists.}
\par\medskip
\rulename{TaylorBound} supplies the deadline battery\textquotesingle s
rational surrogate for a transcendental bound. The deadline
rule\textquotesingle s soundness argument wants $q \le e^{\theta T}$
for a producer-chosen rational $q$, with decay rate
$\theta \in \Q_{>0}$ and horizon $T \in \Q_{\ge 0}$ from its payload:
its concluded sublevel region $\{V \le (1-p)\,q\}$ grows with $q$, so
the producer wants the sharpest $q$ the arithmetic affords. No checker
doing exact rational arithmetic can evaluate $e^{\theta T}$, and none
needs to: the certificate posts the single wire-free comparison
$q \le \sum_{k \le n} (\theta T)^k / k!$ --- a finite sum of
rationals, checked in exact $\Q$ --- and the consuming lemma carries
the enclosure $\sum_{k \le n} (\theta T)^k / k! \le e^{\theta T}$
(nonnegative Taylor remainder at nonnegative argument), so
$q \le e^{\theta T}$ follows without the checker ever knowing $e$
exists. A producer picks the truncation order $n$ and any convenient
$q$ at or below the partial sum.

What the checker recomputes is the target rule\textquotesingle s
entire check: exact closed evaluation of one comparison. No
certificate ever carries a \rulename{TaylorBound} node --- no sort
exists and none is needed; the producer emits the ground-fact row
directly, and this row is the catalogue-level name for that
elaboration.

What would break without each hypothesis: with rounded rather than
exact evaluation, $q$ could creep above the true partial sum and the
surrogate above $e^{\theta T}$; without wire-freedom the
``constant'' could depend on state and the enclosure is meaningless;
and the sign condition $\theta T \ge 0$ is load-bearing in the
consuming lemma --- at negative arguments the Taylor remainder
alternates in sign and the enclosure direction fails at alternate
truncation orders (the proof section carries the counterexample).

\cataloguenote{Deliberately \emph{not} aliased onto the interval-enclosure rule:
that would import witness-bracket machinery for a metatheory
one-liner, and the enclosure rule today lives only in a separate
reference-checker implementation besides. The seam this rule
instantiates is a design invariant of the whole extension page:
transcendental facts live in soundness lemmas (the Lean-side deadline
lemma consumes the standard library\textquotesingle s
partial-sum-below-exponential fact), arithmetic lives in the checker.
Mechanization state, plainly: the bound statements below are the
enclosure and its exact-$\Q$ face. \texttt{taylorLeaf\_le\_exp}
(module \texttt{BernsteinLogicAux.QtslAliases}) is the enclosure step
--- $q$ at or below the truncated sum is below $e^{x}$ at nonnegative
argument, the nonnegative Taylor remainder --- and
\texttt{taylorLeaf\_rat\_le\_exp} is the row end to end: the wire-free
comparison exactly as the checker replays it, in $\Q$, with the
$(\theta, T)$ payload shape, concluding $q \le e^{\theta T}$. The
consuming site is \texttt{sound\_QDeadlineBattery}\textquotesingle s
$hq$ premise (the target\textquotesingle s entry), whose internal
enclosure step these statements name as a rule-level fact. The sign
condition $0 \le \theta T$ is a stated hypothesis, exactly as
load-bearing as the commentary\textquotesingle s counterexample says
it is. No profile enters; the entry carries no conditional flag. The
target\textquotesingle s comparator exclusion ($\neq$ inadmissible)
travels with the alias to any port.
}
\ruleblock{Lean statement}
\declhead{BernsteinLogicAux.QtslAliases}{taylorLeaf\_le\_exp}{\badge{BadgeGreen}{statement bound}}
\begin{leansrc}
/-- **The enclosure step** (QTSL-rules.toml row `TaylorBound (T7)`): a
producer-chosen `q` at or below the order-`n` Taylor partial sum of
`exp` at a nonnegative argument is below `exp` itself — the nonnegative
Taylor remainder (`Real.sum_le_exp_of_nonneg`). The sign condition
`0 ≤ x` is load-bearing: at negative arguments the remainder alternates
in sign, so the enclosure fails at alternate truncation orders (at
`x = -1` it fails at `n = 0`, since `1 > e⁻¹`, and holds at `n = 1`) —
there is no all-`n` enclosure to be had (counter-witnesses in the
book's proof section). -/
theorem taylorLeaf_le_exp {q x : ℝ} {n : ℕ} (hx : 0 ≤ x)
    (hq : q ≤ ∑ k ∈ range (n + 1), x ^ k / (Nat.factorial k : ℝ)) :
    q ≤ Real.exp x
\end{leansrc}
\declhead{BernsteinLogicAux.QtslAliases}{taylorLeaf\_rat\_le\_exp}{\badge{BadgeGreen}{statement bound}}
\begin{leansrc}
/-- **`TaylorBound`, the exact-ℚ face, end to end** (rule row
`TaylorBound (T7)` — the single wire-free closed comparison, checked in
exact ℚ, that stands in for `q ≤ e^{θT}`): decay rate `θ ∈ ℚ≥0` and
horizon `T ∈ ℕ` from the deadline payload, truncation order `n`
producer-chosen. The checker replays the hypothesis; the conclusion —
which mentions `e` — lives only here (the catalogue's "the checker never
needs to know `e` exists"). Consuming site: `sound_QDeadlineBattery`'s
`hq` premise. -/
theorem taylorLeaf_rat_le_exp {qq θ : ℚ} {T n : ℕ} (hθ : 0 ≤ θ)
    (hleaf : qq ≤ ∑ k ∈ range (n + 1), (θ * T) ^ k / (Nat.factorial k : ℚ)) :
    (qq : ℝ) ≤ Real.exp (θ * T)
\end{leansrc}
\par\medskip
\textbf{The payload-level elaborator, displayed whole.}\enspace
$(q, \theta, T, n)$ to the one \rulename{PConstExact} payload, the
right side expanded producer-side in exact $\Q$; the checker replays
one closed comparison, and the transcendental conclusion appears only
in \texttt{taylorElab\_le\_exp}\textquotesingle s goal.

\declhead{BernsteinLogic.QtslElaborators}{taylorFormula\_true\_iff}{\badge{BadgeGreen}{statement bound}}
\begin{leansrc}
/-- The wire-level replay identity: the checker's `true` verdict is
exactly the ℚ comparison against the partial sum. -/
theorem taylorFormula_true_iff (q θ : Rat) (T n : Nat) :
    (taylorFormula q θ T n).evalClosed = some true ↔ q ≤ taylorPartialSum θ T n
\end{leansrc}
\declhead{BernsteinLogic.QtslElaborators}{taylorElab\_sound}{\badge{BadgeGreen}{statement bound}}
\begin{leansrc}
/-- The factoring, BL face: `sound_pConstExact` at the elaborated
payload (`prop:tb-sound`'s first half — validity of the posted
comparison at every valuation). -/
theorem taylorElab_sound {H : Hypernet} {seq : Sequent} {q θ : Rat} {T n : Nat}
    (hEval : (taylorFormula q θ T n).evalClosed = some true)
    (hpost : seq.post = taylorFormula q θ T n) :
    seq.semanticallyValid H
\end{leansrc}
\declhead{BernsteinLogic.QtslElaborators}{taylorElab\_le\_exp}{\badge{BadgeGreen}{statement bound}}
\begin{leansrc}
/-- The semantic face, end to end (`prop:tb-sound` composed with
`prop:tb-enclosure`): the checked wire comparison bounds `q` by
`e^{θT}` through the bound enclosure — the elaborator factoring through
the kernel replay into `taylorLeaf_rat_le_exp`. The sign condition
`0 ≤ θ` is load-bearing (the proof section's negative-argument
counterexample); the deadline rule's typing supplies it. -/
theorem taylorElab_le_exp {q θ : Rat} {T n : Nat} (hθ : 0 ≤ θ)
    (hEval : (taylorFormula q θ T n).evalClosed = some true) :
    (q : ℝ) ≤ Real.exp (θ * T)
\end{leansrc}
\declhead{BernsteinLogic.QtslElaborators}{taylorPartialSum}{\badge{BadgeGreen}{definition bound}}
\begin{leansrc}
/-- The order-`n` Taylor partial sum of `exp (θ·T)`, computed in exact
ℚ — the producer-side expansion of `rec:tb-transport`'s right side. -/
def taylorPartialSum (θ : Rat) (T n : Nat) : Rat :=
  ∑ k ∈ Finset.range (n + 1), (θ * (T : Rat)) ^ k / (Nat.factorial k : Rat)
\end{leansrc}
\declhead{BernsteinLogic.QtslElaborators}{taylorElab}{\badge{BadgeGreen}{definition bound}}
\begin{leansrc}
/-- **The TaylorBound payload elaborator** (`rec:tb-transport`): (q, θ,
T, n) plus the checked comparison, to the one `PConstExactData`
inhabitant. No `TaylorBound` node sort exists and none is needed
(`rec:tb-template`) — this function is the catalogue-level name for the
producer's elaboration, mechanized. -/
def taylorElab (H : Hypernet) (seq : Sequent) (q θ : Rat) (T n : Nat)
    (hEval : (taylorFormula q θ T n).evalClosed = some true)
    (hpost : seq.post = taylorFormula q θ T n) :
    PConstExactData H seq :=
  .mk H seq (.const q) (.const (taylorPartialSum θ T n)) .le (by decide)
    hEval hpost
\end{leansrc}
\ruleblock{Correspondence (Steele display $\leftrightarrow$ Lean)}
\badge{BadgeGreen}{reviewer-certified against the displayed texts}
\sourcenote{The dictionary below is a reviewer-certified statement of how the mechanized statements above interpret the displayed rule: display metavariable $\mapsto$ Lean binder, each display hypothesis $\mapsto$ the premise or side condition carrying it, the display conclusion $\mapsto$ the goal --- with any deliberate difference in width between display and mechanization named outright. Certification is recorded as a hash triple over (the displayed cell, the extracted statement texts, this dictionary\textquotesingle s own text); if any of the three changes, this block turns amber until a reviewer re-certifies. The correspondence claim itself remains editorial --- reviewer-grade, not machine-checked; its machine-checked successor types each soundness statement against the rule as a structured object, and supersedes this block per rule as it lands.}
The displayed wire-free comparison $q \le \sum_{k \le n}
(\theta T)^k / k!$ standing in for $q \le e^{\theta T}$ $\mapsto$ the
bound pair: \texttt{taylorLeaf\_le\_exp} (the enclosure step ---
partial sum below \texttt{Real.exp} at nonnegative argument, the
nonneg-remainder citation) and \texttt{taylorLeaf\_rat\_le\_exp}, the
end-to-end exact-$\Q$ face: \texttt{qq θ : ℚ}, \texttt{T n : ℕ},
hypothesis the exact-$\Q$ partial-sum comparison the checker replays,
conclusion \texttt{(qq : ℝ) ≤ Real.exp (θ * T)} --- the
flavor text\textquotesingle s ``the checker never needs to know $e$
exists'' as a binder fact: $e$ appears in the goal alone, never in a
hypothesis. The sign condition $\mapsto$ \texttt{hθ : 0 ≤ θ}
(load-bearing --- the alternating-remainder counterexample in the
proof section is why the binder exists); the truncation order
\texttt{n} is producer-chosen, unconstrained. Width note: the
consuming site (\texttt{sound\_QDeadlineBattery}\textquotesingle s
\texttt{hq} premise) is cited, not bound here --- correct
package-boundary behaviour, the deadline battery\textquotesingle s own
entry carries it; the displayed horizon $T \in \N$ matches the bound end-to-end
lemma\textquotesingle s \texttt{T : ℕ} exactly (the rational-horizon
generalization is mathematically available from
\texttt{taylorLeaf\_le\_exp} but is neither bound nor displayed); and
the displayed $\theta \in \Q_{>0}$ is a deliberately stronger surface
restriction than the bound \texttt{hθ : 0 ≤ θ} --- safe direction, the
deadline battery\textquotesingle s own decay-rate positivity.

\ruleblock{Lean verifier receipts}
\sourcenote{module \texttt{BernsteinLogicAux.QtslAliases}; bound build \texttt{570599779c29\ldots} (bound from receipt 5a60f4d01fd81ac5, 2026-07-21T05:39:02Z). A receipt records that named checking layers accepted a build: (module, compiled-artifact SHA-256, environment SHA-256, checker layers), minted when the checking gate ran; the store receipts attesting the bound build are listed with the layers each run credited. Nothing here asks the reader to trust the store: every layer is an independently runnable program over the quoted sources --- rebuild the module and re-run the checkers to reproduce the verdict. Receipts attesting superseded builds credit no layer and are omitted from this rendering (the count of omissions, if any, is noted below).}
\begin{receipttable}
\receiptrow{0a3819a978e31127}{2026-07-21T04:51:36Z}{build+decorrelated}{\badge{BadgeGreen}{current}}
\receiptrow{0bb1074df705d863}{2026-07-21T05:21:58Z}{build+decorrelated}{\badge{BadgeGreen}{current}}
\receiptrow{5a60f4d01fd81ac5}{2026-07-21T05:39:02Z}{build+replay+hardened+decorrelated}{\badge{BadgeGreen}{current}}
\end{receipttable}
\sourcenote{11 further receipt(s) in the store attest superseded builds of this module and are omitted from this rendering.}
\sourcenote{module \texttt{BernsteinLogic.QtslElaborators}; bound build \texttt{0205ec213b91\ldots} (hashed from the built olean at bind time). A receipt records that named checking layers accepted a build: (module, compiled-artifact SHA-256, environment SHA-256, checker layers), minted when the checking gate ran; the store receipts attesting the bound build are listed with the layers each run credited. Nothing here asks the reader to trust the store: every layer is an independently runnable program over the quoted sources --- rebuild the module and re-run the checkers to reproduce the verdict. Receipts attesting superseded builds credit no layer and are omitted from this rendering (the count of omissions, if any, is noted below).}
\begin{receipttable}
\receiptrow{a237f4cff27cb15a}{2026-07-22T10:06:47Z}{build+decorrelated}{\badge{BadgeGreen}{current}}
\end{receipttable}
\layerline{\layeritem{build}{lake build (elaboration)}{\badge{BadgeGreen}{current}}\layeritem{replay}{leanchecker (official kernel replay)}{\badge{BadgeGreen}{current}}\layeritem{hardened}{safe\_verify (hardened screen)}{\badge{BadgeGreen}{current}}\layeritem{decorrelated}{lean4export + nanoda (decorrelated kernel)}{\badge{BadgeGreen}{current}}}
\sourcenote{\badge{BadgeGray}{comparator: none} no fixture in the three-implementation comparator harness (compiled / kernel-interpreted / reference-interpreted) exercises this rule.}
\ruleblock{HOL4 statement}
\absentbox{OUT OF SCOPE BY DESIGN}{Out of scope for the HOL side \emph{by design}, not by backlog: the
wire format never carries a \rulename{TaylorBound} node (no sort
exists and none is needed --- the producer emits the ground-fact
comparison row directly), so no wire certificate a HOL checker reads
could ever contain one; HOL coverage for this rule \emph{is} the
target rule\textquotesingle s. Today even that statement is not
writable: the HOL4 development (theory \texttt{bern\_catalogue}, the
base catalogue\textquotesingle s check-function cone) predates
\rulename{PConstExact}, whose wire tag is unallocated --- see that
rule\textquotesingle s section for the honest state of the
target\textquotesingle s own box.
}
\ruleblock{HOL verifier receipts}
\absentbox{NONE}{no HOL verification receipt exists for this rule.}
\ruleblock{Proof}

\sourcenote{Genre: alias rule (tier 3) --- a respelling template. The
five artefacts an alias owes the book: the 1$\to$1 unpack, the filled-in
target application, the payload-transport table, the exactness identity,
and the soundness corollary through the target\textquotesingle s lemma.
All five follow, with one honest asymmetry marked where it occurs: this
rule\textquotesingle s exactness identity is not a ring identity but an
\emph{enclosure} --- the Taylor-remainder inequality --- and it lives in
the consuming lemma, never in the checker. Mechanization state: the
target\textquotesingle s lemma is kernel-verified; the alias-level
semantic content --- the enclosure and its exact-$\Q$ face --- is
mechanized and bound in the entry, and so --- now that the extension
certificate schema exists --- is the payload-level elaborator (the
remark below names the bound functions).}

\begin{brecipe}[the template instance --- no node sort, by design]
\label{rec:tb-template}
The deadline rule (\rulename{QDeadlineBattery}) wants the fact
$q \le e^{\theta T}$ for a producer-chosen rational surrogate $q$, with
decay rate $\theta \in \Q_{> 0}$ and horizon $T \in \Q_{\ge 0}$ from its
payload: its concluded sublevel region $\{V \le (1-p)\,q\}$ grows with
$q$, so the producer wants the sharpest $q$ the arithmetic affords. No
checker doing exact rational arithmetic can evaluate $e^{\theta T}$, and
none needs to: the certificate posts the single wire-free comparison
\[
q \;\le\; \textstyle\sum_{k \le n} (\theta T)^k / k!
\]
for a producer-chosen truncation order $n$ --- a finite sum of
rationals, one closed comparison --- and the consuming lemma carries the
enclosure to $e^{\theta T}$ (Proposition~\ref{prop:tb-enclosure}). The
unpack image is a single \rulename{PConstExact} node; indeed no
\rulename{TaylorBound} node sort exists and none is needed --- the
producer emits the ground-fact row directly, and this rule is the
catalogue-level name for that elaboration. One node $\to$ one node,
nothing derived.
\end{brecipe}

\begin{brecipe}[the filled-in application]\label{rec:tb-filled}
\begin{mathpar}
\inferrule*[Right=\rulename{PConstExact}]{
  q \le \textstyle\sum_{k \le n} (\theta T)^k / k! \ \text{holds under exact evaluation in } \Q
}{
  \Phi \vdash q \le \textstyle\sum_{k \le n} (\theta T)^k / k!
}
\end{mathpar}
with arbitrary premise $\Phi$, as the target licenses (the rule is
premise-free). Both sides are closed rational terms --- the sum is
finite, every summand exact --- so the target\textquotesingle s verified
evaluator decides the comparison in exact $\Q$, and the comparator
($\le$) lies inside the target\textquotesingle s admissible comparator
set.
\end{brecipe}

\begin{brecipe}[payload transport --- field for field]\label{rec:tb-transport}
\begin{center}\small
\begin{tabular}{@{}ll@{}}
\textbf{\rulename{TaylorBound} data} & \textbf{\rulename{PConstExact} payload}\\
\midrule
$q$ \textnormal{(the rational surrogate)} & left side of the comparison\\
$\theta, T, n$ \textnormal{(rate, horizon, order)} & the right side $\textstyle\sum_{k \le n} (\theta T)^k/k!$, expanded producer-side\\
comparator & $\le$ (inside the admissible set)\\
\textnormal{wire-freedom} & implied, not declared: the evaluator returns $\mathsf{none}$ on any wire\\
\end{tabular}
\end{center}
The checker sees one closed comparison and nothing else; $\theta$, $T$,
and $n$ survive only through the value of the right-hand term. The
transcendental target $e^{\theta T}$ appears in \emph{no} field --- that
is the point.
\end{brecipe}

\begin{bprop}[the enclosure --- the exactness identity, honest form]
\label{prop:tb-enclosure}
For $\theta T \ge 0$ and any $n$,
\[
\textstyle\sum_{k \le n} (\theta T)^k / k! \;\le\; e^{\theta T},
\qquad\text{hence}\qquad
q \le \textstyle\sum_{k \le n} (\theta T)^k / k! \;\Longrightarrow\; q \le e^{\theta T}.
\]
\end{bprop}
\begin{proof}
Two lines, though of a different species than a ring identity: at
$x := \theta T \ge 0$ every term of the exponential series is
nonnegative, so each partial sum is dominated by the limit --- in the
classical Taylor form, $e^{x} = \sum_{k \le n} x^k/k! + R_n$ with the
Lagrange remainder $R_n = e^{\xi} x^{n+1}/(n+1)! \ge 0$ for some
$\xi \in [0, x]$ \cite{rudin1976}. The sign condition is load-bearing:
at negative arguments the remainder alternates in sign and the
enclosure direction fails at alternate truncation orders ($n = 0$,
$x = -1$: the partial sum is $1 > e^{-1}$; at $n = 1$ it happens to
hold again --- there is no all-$n$ enclosure to be had) --- the
deadline rule\textquotesingle s typing ($\theta > 0$, $T \ge 0$) is
what entitles the lemma to the inequality, via $\theta T \ge 0$. Unlike the
base-calculus exemplar alias, whose two-line identity is checked by the
kernel as exact polynomial arithmetic, this identity is consumed
\emph{inside the deadline rule\textquotesingle s soundness lemma} (the
Lean side consumes the standard library\textquotesingle s
partial-sum-below-exponential fact,
\texttt{Real.sum\_le\_exp\_of\_nonneg}, there) and never reaches the
checker: the checker\textquotesingle s whole obligation is the exact
rational comparison of Recipe~\ref{rec:tb-filled}. The checker never
needs to know $e$ exists.
\end{proof}

\begin{bprop}[soundness corollary]\label{prop:tb-sound}
If the posted comparison checks and $\theta T \ge 0$ (the sign
hypothesis of Proposition~\ref{prop:tb-enclosure}), then
$q \le e^{\theta T}$, which is the fact the deadline
rule\textquotesingle s lemma consumes.
\end{bprop}
\begin{proof}
\texttt{sound\_pConstExact} (the target\textquotesingle s
kernel-verified lemma, quoted in the \rulename{PConstExact} section)
turns the $\mathsf{true}$ evaluation into the validity of
$q \le \sum_{k \le n} (\theta T)^k/k!$ at every valuation;
Proposition~\ref{prop:tb-enclosure} extends to $e^{\theta T}$ inside
the consuming lemma. There is no other content --- the alias genre
working as intended.
\end{proof}

\begin{bremark}[the payload-level elaborator, mechanized]
\label{rem:tb-elab}
The \emph{semantic} content of the respelling is bound in the entry:
\texttt{taylorLeaf\_le\_exp} is
Proposition~\ref{prop:tb-enclosure}\textquotesingle s enclosure step,
and \texttt{taylorLeaf\_rat\_le\_exp} is
Proposition~\ref{prop:tb-sound}\textquotesingle s exact-$\Q$ face end
to end. The payload-level elaborator --- previously this
page\textquotesingle s one absent artefact, held back until the
extension certificate schema existed --- is bound with it:
\texttt{taylorElab} (module \texttt{BernsteinLogic.QtslElaborators})
takes $(q, \theta, T, n)$ to exactly the one \rulename{PConstExact}
payload of Recipe~\ref{rec:tb-transport}, the right side expanded
producer-side in exact $\Q$ (\texttt{taylorPartialSum}), as an
inhabitant of the target\textquotesingle s own payload type --- the
comparator confinement discharged at elaboration, once --- with
\texttt{taylorElab\_sound} the factoring through
\texttt{sound\_pConstExact}, \texttt{taylorFormula\_true\_iff} the
wire-level replay identity, and \texttt{taylorElab\_le\_exp} the
semantic face end to end (the checked comparison bounding $q$ by
$e^{\theta T}$ through the bound enclosure --- the transcendental
conclusion still appears in no field, only in the goal). Stated at the
discrete horizon, $T : \mathbb{N}$, matching the bound exact-$\Q$ face
--- the entry\textquotesingle s width note carries the divergence from
the recipe\textquotesingle s rational-time payload; the enclosure of
Proposition~\ref{prop:tb-enclosure} itself needs only
$\theta T \ge 0$.
\end{bremark}

\begin{bremark}[deliberately not an enclosure-rule alias]
The calculus also has interval-enclosure machinery with witness
brackets. Respelling this leaf onto it was considered and declined:
that would import witness-bracket payload for a metatheory one-liner,
and the enclosure rule today lives only in a separate
reference-checker implementation besides. The seam this rule
instantiates is a design invariant of the whole extension page:
transcendental facts live in soundness lemmas, arithmetic lives in the
checker.
\end{bremark}

\begin{bremark}[the inherited comparator confinement]
\rulename{PConstExact} excludes one comparator ($\neq$) from its
admissible set, for reasons of a second soundness invariant unpacked in
that rule\textquotesingle s section. This leaf uses $\le$ only ---
comfortably inside --- but any port or mechanization of this rule must
\emph{preserve} the exclusion rather than re-derive the comparator set
from the semantic lemma alone; the flag travels with the alias.
\end{bremark}

\begin{anchors}
target lemma (Prop.~\ref{prop:tb-sound}) & \texttt{sound\_pConstExact}; \texttt{Formula.holds\_of\_evalClosed\_true} (quoted in the \rulename{PConstExact} section)\\
wire-freedom implied at the leaf & \texttt{Formula.evalClosed} ($\mathsf{none}$ on any wire; same section)\\
the enclosure, mechanized where consumed & \texttt{Real.sum\_le\_exp\_of\_nonneg} inside the deadline lemma \texttt{sound\_QDeadlineBattery} (see that rule\textquotesingle s entry)\\
the enclosure + exact-$\Q$ face, rule-level & \texttt{taylorLeaf\_le\_exp}, \texttt{taylorLeaf\_rat\_le\_exp} (bound in the entry)\\
payload-level elaborator (\autoref{rem:tb-elab}) & \texttt{taylorElab} + \texttt{taylorElab\_sound} + \texttt{taylorElab\_le\_exp}, over \texttt{taylorPartialSum} + \texttt{taylorFormula\_true\_iff} (module \texttt{BernsteinLogic.QtslElaborators}, bound in the entry)\\
\end{anchors}


\rulesection{67}{RootEval}{alias for \textcolor{DiffBlue}{\rulename{PConstExact}}}
\label{blcatrule-RootEval}
\ruleblock{The rule (Steele metanotation)}
\sourcenote{as realized in the mechanized development; source row: \texttt{\texttt{QTSL-rules.toml}} --- \emph{RootEval}. The metanotation statement is typeset by hand for notational quality; its mechanized ground truth is the Lean statement below, extracted from the sources at every build of this book and hash-bound. The rule \emph{roster} is machine-checked (every mechanized rule tag round-trips to exactly one catalogue row); the statement \emph{text} is editorial, checked by eye against the extracted Lean beside it.}
\begin{mathpar}
\inferrule*{
  R = \{x : V(x) \le y\} \ \text{(sublevel root atom)} \\
  \vdash V(x_0) \le y \ \text{(one \rulename{PConstExact} leaf: the substituted comparison is closed)}
}{
  \vdash x_0 \in R
}
\and
\inferrule*{
  R = \textstyle\prod_i [a_i, b_i] \ \text{(box root atom)} \\
  \forall i.\ a_i \le (x_0)_i \le b_i \ \text{(recomputed in exact } \Q\text{)}
}{
  \vdash x_0 \in R
}
\end{mathpar}

\ruleflavor{catalogued as the eval leaf of the region-to-satisfaction bridge (a BL \texttt{Proof} leaf; the bridge\textquotesingle s node sort is out of scope of the current baked theory — its baked warrant seam is \texttt{qsat\_warrant}); sublevel atoms discharge through \rulename{PConstExact} (the first form is that rule\textquotesingle s application at the substituted point), box atoms through the registry interval check (the second form: conformance recomputation, not a rule application).}
\par\medskip
\rulename{RootEval} is the membership leaf itself, catalogued as a
rule of its own because every initialized certificate ends with one:
membership of the specification\textquotesingle s initial point in the
\emph{root atom} --- the region atom concluded at the
certificate\textquotesingle s root (``root'' names the
certificate\textquotesingle s root claim, not roots of polynomials)
--- by exact rational evaluation. In the certificate schema it
is not a node sort at all: it is the eval leaf of the region-to-
satisfaction bridge (\rulename{QSatFromRegion} --- whose own node sort
is out of scope of the current baked theory; the baked warrant seam is
\texttt{qsat\_warrant}), a base-proof leaf. For a sublevel atom
$\{x : V(x) \le y\}$ the leaf is one \rulename{PConstExact}
application: substitute the rational point into the polynomial ---
exact, no rounding anywhere --- and decide one closed comparison. For
a box atom, the checker recomputes per-coordinate containment directly
(the side-condition registry: conformance recomputation, not a rule
application).

What the checker recomputes: everything --- the leaf \emph{is} a
recompute; there is no payload beyond the point and the atom. What
would break without exactness: under naive (non-directed) rounding a
point lying $\varepsilon$ outside the region could certify
membership, and the bridge rule would conclude satisfaction from a
false premise --- directed outward rounding keeps that direction
sound but surrenders the equivalence at the boundary --- the exactness identity in the
proof section is a literal identity of values, and it is the whole
rule. The comparator exclusion inherited from the target ($\neq$
inadmissible, a second-soundness-invariant exclusion unpacked in the
target\textquotesingle s section) must survive any port.

\cataloguenote{The judgment-level composition --- how the checked leaf feeds the
instantiation that concludes root satisfaction --- lives in the
\rulename{QSatFromRegion} section; this row owns the leaf.
Mechanization state, plainly: the bound statements below are the
exactness identity itself. \texttt{rootEval\_eval\_cast} (module
\texttt{BernsteinLogicAux.QtslAliases}) is the literal identity of
values --- evaluating the $\Q$-coefficient polynomial at the rational
point and casting equals evaluating the cast polynomial at the cast
point --- and \texttt{rootEval\_sublevel\_iff} /
\texttt{rootEval\_box\_iff} are its two atom shapes, each an
\emph{equivalence}: the exact-$\Q$ comparison decides the
real-semantics membership in both directions, which is precisely the
content a rounded evaluator would lose. The statements are over the
generic polynomial-evaluation vocabulary (\texttt{MvPolynomial} over
$\Q$, cast to $\R$); the wire-term evaluator that produces the
$\Q$-comparison is the target\textquotesingle s own kernel-verified
theorem, quoted in the \rulename{PConstExact} section, and the
box-atom branch\textquotesingle s registry recompute is checker-side
by design. No profile enters; the entry carries no conditional
flag.
}
\ruleblock{Lean statement}
\declhead{BernsteinLogicAux.QtslAliases}{rootEval\_eval\_cast}{\badge{BadgeGreen}{statement bound}}
\begin{leansrc}
/-- **The exactness identity** (QTSL-rules.toml row `RootEval (T19)` — "a
literal identity of values, and it is the whole rule"): evaluating a
ℚ-coefficient polynomial at a rational point and casting the result equals
evaluating the cast polynomial at the cast point. This is why the checker's
exact rational arithmetic (`PConstExact`'s replay) *is* the real-semantics
evaluation, with no rounding gap anywhere. -/
theorem rootEval_eval_cast {σ : Type*} (P : MvPolynomial σ ℚ) (xq : σ → ℚ) :
    ((MvPolynomial.eval xq P : ℚ) : ℝ)
      = MvPolynomial.eval (fun i => (xq i : ℝ)) (MvPolynomial.map (Rat.castHom ℝ) P)
\end{leansrc}
\declhead{BernsteinLogicAux.QtslAliases}{rootEval\_sublevel\_iff}{\badge{BadgeGreen}{statement bound}}
\begin{leansrc}
/-- **`RootEval`, sublevel atom, as an equivalence**: the exact ℚ comparison
at the initial point *decides* membership of the cast point in the cast
sublevel atom `{x : V(x) ≤ thr}` — both directions, which is the content a
rounded evaluator would lose (one direction would survive rounding; the
identity does not). The forward direction is the leaf as the checker
discharges it (rule row `RootEval (T19)`, sublevel branch →
`PConstExact`). -/
theorem rootEval_sublevel_iff {σ : Type*} (P : MvPolynomial σ ℚ) (xq : σ → ℚ)
    (thr : ℚ) :
    MvPolynomial.eval xq P ≤ thr
      ↔ MvPolynomial.eval (fun i => (xq i : ℝ))
          (MvPolynomial.map (Rat.castHom ℝ) P) ≤ (thr : ℝ)
\end{leansrc}
\declhead{BernsteinLogicAux.QtslAliases}{rootEval\_box\_iff}{\badge{BadgeGreen}{statement bound}}
\begin{leansrc}
/-- **`RootEval`, box atom, as an equivalence**: per-coordinate exact ℚ
containment decides containment of the cast point in the cast box (rule row
`RootEval (T19)`, box branch — the side-condition registry's per-coordinate
recompute). -/
theorem rootEval_box_iff {ι : Type*} (aq bq xq : ι → ℚ) :
    (∀ i, aq i ≤ xq i ∧ xq i ≤ bq i)
      ↔ ∀ i, ((aq i : ℝ) ≤ (xq i : ℝ) ∧ (xq i : ℝ) ≤ (bq i : ℝ))
\end{leansrc}
\par\medskip
\textbf{The payload-level elaborator, displayed whole.}\enspace The
closed left-hand term is built by verified substitution
(\texttt{polyClosedTerm}; its evaluation lemma says the
checker\textquotesingle s exact replay computes the value $V(x_0)$
itself), and the elaborated payload inhabits the
target\textquotesingle s own payload type --- comparator confinement
discharged at elaboration.

\declhead{BernsteinLogic.QtslElaborators}{polyClosedTerm\_evalQ}{\badge{BadgeGreen}{statement bound}}
\begin{leansrc}
@[simp] theorem polyClosedTerm_evalQ (x0 : Nat → Rat) :
    ∀ p : Poly, (polyClosedTerm x0 p).evalQ? = some (Poly.eval x0 p)
  | [] => by simp [polyClosedTerm, Poly.eval, Term.evalQ?]
  | mc :: rest => by
    have ih
\end{leansrc}
\declhead{BernsteinLogic.QtslElaborators}{rootEvalFormula\_true\_iff}{\badge{BadgeGreen}{statement bound}}
\begin{leansrc}
/-- The wire-level exactness identity, ℚ face (`prop:re-exact` at the
payload: the checker's `true` verdict IS the exact membership
inequality — both directions; the ℝ face stays
`Aliases.rootEval_sublevel_iff`). -/
theorem rootEvalFormula_true_iff (x0 : Nat → Rat) (V : Poly) (thr : Rat) :
    (rootEvalFormula x0 V thr).evalClosed = some true ↔ Poly.eval x0 V ≤ thr
\end{leansrc}
\declhead{BernsteinLogic.QtslElaborators}{rootEvalElab\_sound}{\badge{BadgeGreen}{statement bound}}
\begin{leansrc}
/-- The factoring, BL face: the elaborated payload is sound through the
target's kernel-verified lemma — one application of
`sound_pConstExact`, which is itself one step to
`Formula.holds_of_evalClosed_true` (`prop:re-sound`). -/
theorem rootEvalElab_sound {H : Hypernet} {seq : Sequent} {x0 : Nat → Rat}
    {V : Poly} {thr : Rat}
    (hEval : (rootEvalFormula x0 V thr).evalClosed = some true)
    (hpost : seq.post = rootEvalFormula x0 V thr) :
    seq.semanticallyValid H
\end{leansrc}
\declhead{BernsteinLogic.QtslElaborators}{polyClosedTerm}{\badge{BadgeGreen}{definition bound}}
\begin{leansrc}
/-- The closed polynomial term at a rational point. Wire-free by
construction, and its exact evaluation is `Poly.eval` at that point
(`polyClosedTerm_evalQ`) — the checker computes the value `V(x₀)`
itself, not an approximation to it. -/
def polyClosedTerm (x0 : Nat → Rat) (p : Poly) : Term :=
  p.foldr (fun mc acc => .add (.mul (.const mc.2) (monoClosedTerm x0 mc.1)) acc)
    (.const 0)
\end{leansrc}
\declhead{BernsteinLogic.QtslElaborators}{rootEvalFormula}{\badge{BadgeGreen}{definition bound}}
\begin{leansrc}
/-- The elaborated leaf formula: `V(x₀) ≤ thr` as one closed exact
comparison. -/
def rootEvalFormula (x0 : Nat → Rat) (V : Poly) (thr : Rat) : Formula :=
  .atom .le (polyClosedTerm x0 V) (.const thr)
\end{leansrc}
\declhead{BernsteinLogic.QtslElaborators}{rootEvalElab}{\badge{BadgeGreen}{definition bound}}
\begin{leansrc}
/-- **The RootEval payload elaborator** (`rec:re-transport`, sublevel
branch): the alias's data — point, atom — plus its own checked leaf, to
the one `PConstExactData` inhabitant. The comparator confinement
travels by construction (`.le ≠ .ne` discharged here, once). -/
def rootEvalElab (H : Hypernet) (seq : Sequent) (x0 : Nat → Rat) (V : Poly)
    (thr : Rat)
    (hEval : (rootEvalFormula x0 V thr).evalClosed = some true)
    (hpost : seq.post = rootEvalFormula x0 V thr) :
    PConstExactData H seq :=
  .mk H seq (polyClosedTerm x0 V) (.const thr) .le (by decide) hEval hpost
\end{leansrc}
\ruleblock{Correspondence (Steele display $\leftrightarrow$ Lean)}
\badge{BadgeGreen}{reviewer-certified against the displayed texts}
\sourcenote{The dictionary below is a reviewer-certified statement of how the mechanized statements above interpret the displayed rule: display metavariable $\mapsto$ Lean binder, each display hypothesis $\mapsto$ the premise or side condition carrying it, the display conclusion $\mapsto$ the goal --- with any deliberate difference in width between display and mechanization named outright. Certification is recorded as a hash triple over (the displayed cell, the extracted statement texts, this dictionary\textquotesingle s own text); if any of the three changes, this block turns amber until a reviewer re-certifies. The correspondence claim itself remains editorial --- reviewer-grade, not machine-checked; its machine-checked successor types each soundness statement against the rule as a structured object, and supersedes this block per rule as it lands.}
The displayed leaf rules --- exact rational evaluation concluding
root-atom membership, one form per atom shape --- $\mapsto$ three statements, and the
strength gradation is the dictionary\textquotesingle s content.
\texttt{rootEval\_eval\_cast} is the literal identity of values:
$\Q$-evaluation, cast, equals $\R$-evaluation of the cast polynomial
at the cast point (\texttt{MvPolynomial} over any variable type ---
dimension-free). The two atom shapes are then \emph{equivalences},
not implications: \texttt{rootEval\_sublevel\_iff} (exact-$\Q$
comparison $\leftrightarrow$ membership in the cast sublevel atom)
and \texttt{rootEval\_box\_iff} (per-coordinate $\Q$ containment
$\leftrightarrow$ cast-box membership) --- two-directional exactness,
which is precisely what a rounded evaluator would lose (under
directed outward rounding the sound forward direction survives and
the iff does not; undirected rounding loses even soundness): the entry\textquotesingle s
``no rounding gap anywhere'' as bind structure. Width note: the
checker consumes only the forward directions (leaf discharge via
\rulename{PConstExact} / the registry interval recompute); the
certified equivalences are strictly stronger, safe direction, and no
profile, measure, or topology enters any binder.

\ruleblock{Lean verifier receipts}
\sourcenote{module \texttt{BernsteinLogicAux.QtslAliases}; bound build \texttt{570599779c29\ldots} (bound from receipt 5a60f4d01fd81ac5, 2026-07-21T05:39:02Z). A receipt records that named checking layers accepted a build: (module, compiled-artifact SHA-256, environment SHA-256, checker layers), minted when the checking gate ran; the store receipts attesting the bound build are listed with the layers each run credited. Nothing here asks the reader to trust the store: every layer is an independently runnable program over the quoted sources --- rebuild the module and re-run the checkers to reproduce the verdict. Receipts attesting superseded builds credit no layer and are omitted from this rendering (the count of omissions, if any, is noted below).}
\begin{receipttable}
\receiptrow{0a3819a978e31127}{2026-07-21T04:51:36Z}{build+decorrelated}{\badge{BadgeGreen}{current}}
\receiptrow{0bb1074df705d863}{2026-07-21T05:21:58Z}{build+decorrelated}{\badge{BadgeGreen}{current}}
\receiptrow{5a60f4d01fd81ac5}{2026-07-21T05:39:02Z}{build+replay+hardened+decorrelated}{\badge{BadgeGreen}{current}}
\end{receipttable}
\sourcenote{11 further receipt(s) in the store attest superseded builds of this module and are omitted from this rendering.}
\sourcenote{module \texttt{BernsteinLogic.QtslElaborators}; bound build \texttt{0205ec213b91\ldots} (hashed from the built olean at bind time). A receipt records that named checking layers accepted a build: (module, compiled-artifact SHA-256, environment SHA-256, checker layers), minted when the checking gate ran; the store receipts attesting the bound build are listed with the layers each run credited. Nothing here asks the reader to trust the store: every layer is an independently runnable program over the quoted sources --- rebuild the module and re-run the checkers to reproduce the verdict. Receipts attesting superseded builds credit no layer and are omitted from this rendering (the count of omissions, if any, is noted below).}
\begin{receipttable}
\receiptrow{a237f4cff27cb15a}{2026-07-22T10:06:47Z}{build+decorrelated}{\badge{BadgeGreen}{current}}
\end{receipttable}
\layerline{\layeritem{build}{lake build (elaboration)}{\badge{BadgeGreen}{current}}\layeritem{replay}{leanchecker (official kernel replay)}{\badge{BadgeGreen}{current}}\layeritem{hardened}{safe\_verify (hardened screen)}{\badge{BadgeGreen}{current}}\layeritem{decorrelated}{lean4export + nanoda (decorrelated kernel)}{\badge{BadgeGreen}{current}}}
\sourcenote{\badge{BadgeGray}{comparator: none} no fixture in the three-implementation comparator harness (compiled / kernel-interpreted / reference-interpreted) exercises this rule.}
\ruleblock{HOL4 statement}
\absentbox{OUT OF SCOPE BY DESIGN}{Out of scope for the HOL side \emph{by design}, not by backlog: this
rule is a \emph{field} of the bridge rule
(\rulename{QSatFromRegion}\textquotesingle s eval slot), never a
node, so no wire certificate a HOL checker reads could ever contain a
\rulename{RootEval} node; HOL coverage for this rule \emph{is} the
target rule\textquotesingle s, and the box-atom branch is loader-side
registry recomputation, likewise not a wire-format rule. Today even
the target\textquotesingle s statement is not writable: the HOL4
development (theory \texttt{bern\_catalogue}, the base
catalogue\textquotesingle s check-function cone) predates
\rulename{PConstExact}, whose wire tag is unallocated --- see that
rule\textquotesingle s section for the honest state of the
target\textquotesingle s own box.
}
\ruleblock{HOL verifier receipts}
\absentbox{NONE}{no HOL verification receipt exists for this rule.}
\ruleblock{Proof}

\sourcenote{Genre: alias rule (tier 3) --- a respelling template: the
1$\to$1 unpack, the filled-in target application, the payload-transport
table, the exactness identity, and the soundness corollary through the
target\textquotesingle s lemma. This row is the membership leaf itself,
catalogued as a rule of its own because every initialized certificate
ends with one; in the extension certificate schema it is a \emph{field} of
the bridge rule (\rulename{QSatFromRegion}\textquotesingle s eval slot),
not a node sort. The judgment-level story lives in the bridge
rule\textquotesingle s section --- this section owns the leaf.
Mechanization state: the target\textquotesingle s lemma is
kernel-verified; the exactness identities are mechanized and bound in
the entry, and so --- now that the extension certificate schema exists
--- is the payload-level elaborator (the remark below names the bound
functions).}

\begin{brecipe}[the template instance]\label{rec:re-template}
One leaf $\to$ one node: membership of the initial point $x_0$ in the
certificate\textquotesingle s \emph{root atom} --- the region atom
concluded at the root of the certificate (naming note: ``root'' refers
to the certificate\textquotesingle s root claim, not to roots of
polynomials) --- by exact rational evaluation. For a \emph{sublevel}
atom $\{x : V(x) \le y\}$ the unpack image is a single
\rulename{PConstExact} application
(Recipe~\ref{rec:re-filled}). For a \emph{box} atom the leaf is not a
rule application at all: per-coordinate interval containment is
recomputed directly by the checker\textquotesingle s side-condition
registry (conformance recomputation) --- the alias reading covers the
sublevel branch, and the box branch is stated so the respelling never
overclaims. Nothing else is derived: the rule is a compact spelling of
one ground-fact row.
\end{brecipe}

\begin{brecipe}[the filled-in application]\label{rec:re-filled}
Sublevel atom $R = \{x : V(x) \le y\}$ with $V$ a polynomial
over $\Q$, threshold $y \in \Q$, and initial point
$x_0$ with rational coordinates:
\begin{mathpar}
\inferrule*[Right=\rulename{PConstExact}]{
  V(x_0) \le y \ \text{holds under exact evaluation in } \Q
}{
  \Phi \vdash V(x_0) \le y
}
\end{mathpar}
with arbitrary premise $\Phi$, as the target licenses (the rule is
premise-free). $V(x_0)$ denotes the closed term obtained by
substituting every coordinate of $x_0$ into $V$ --- closed by
construction, so the target\textquotesingle s verified evaluator
decides the comparison in exact $\Q$, and a $\mathsf{true}$ verdict
certifies closedness as a byproduct (the target section unpacks why).
\end{brecipe}

\begin{brecipe}[payload transport --- field for field]\label{rec:re-transport}
\begin{center}\small
\begin{tabular}{@{}ll@{}}
\textbf{\rulename{RootEval} data} & \textbf{\rulename{PConstExact} payload}\\
\midrule
$x_0$ \textnormal{(rational coordinates)} & substituted into the closed left-hand term\\
sublevel atom $(V, y)$ & the comparison\textquotesingle s two sides\\
comparator & $\le$ (inside the admissible set)\\
box atom \textnormal{(per-axis bounds)} & \emph{not} transported --- registry recompute (Recipe~\ref{rec:re-template})\\
\end{tabular}
\end{center}
\end{brecipe}

\begin{bprop}[the exactness identity]\label{prop:re-exact}
\[
x_0 \in \{x : V(x) \le y\}
\quad\Longleftrightarrow\quad
\mathrm{evalClosed}\bigl(V(x_0) \le y\bigr) = \mathsf{true}.
\]
\end{bprop}
\begin{proof}
Two lines, which is the point. Substituting a rational point into a
rational polynomial is exact arithmetic --- no rounding enters at any
step --- so the evaluator computes the value $V(x_0)$ itself, not an
approximation to it; and the decidable comparison of two exact
rationals is the membership predicate verbatim. This is the degenerate
--- and best --- case of the exactness identity an alias owes: not an
algebraic rearrangement but a literal identity of values. It is also
the whole rule: under naively rounded (non-directed) evaluation, a
point lying $\varepsilon$ outside the region could certify
membership, and the bridge rule\textquotesingle s instantiation (its
section, the extension page) would conclude satisfaction from a false
premise; directed outward rounding would keep that direction sound
but surrender the equivalence --- boundary points, where
$V(x_0) = y$ exactly, typically become undecidable (an
outward enclosure that straddles the threshold cannot decide; an
evaluation whose enclosure does not cross it --- an exact zero-width
value, or a one-sided enclosure touching the threshold --- still
decides).
\end{proof}

\begin{bprop}[soundness corollary]\label{prop:re-sound}
If the leaf checks, $x_0$ belongs to the root atom.
\end{bprop}
\begin{proof}
\texttt{Formula.holds\_of\_evalClosed\_true} (the target\textquotesingle s
kernel-verified evaluator theorem, quoted in the \rulename{PConstExact}
section): a closed comparison whose exact evaluation returns
$\mathsf{true}$ holds at every valuation. By
Proposition~\ref{prop:re-exact} that is membership: substitution
closed the formula, so $V(x_0) \le y$ mentions no wires,
its truth is valuation-independent, and that truth is literally the
sublevel atom\textquotesingle s membership predicate at $x_0$. There is no other
content --- the alias genre working as intended. (For the box branch
the recomputed per-coordinate containment plays the same role, warranted
by the recompute\textquotesingle s definition rather than a rule lemma.)
\end{proof}

\begin{bremark}[the payload-level elaborator, mechanized]
\label{rem:re-elab}
The \emph{semantic} content of the respelling is bound in the entry:
\texttt{rootEval\_eval\_cast} is
Proposition~\ref{prop:re-exact}\textquotesingle s literal identity of
values, and \texttt{rootEval\_sublevel\_iff} /
\texttt{rootEval\_box\_iff} are the two atom shapes as equivalences.
The payload-level elaborator --- previously this page\textquotesingle s
one absent artefact, held back until the extension certificate schema
existed --- is bound with it: \texttt{rootEvalElab} (module
\texttt{BernsteinLogic.QtslElaborators}) takes $(x_0, \text{atom})$ to
exactly the one \rulename{PConstExact} payload of
Recipe~\ref{rec:re-transport} (sublevel branch), as an inhabitant of
the target\textquotesingle s own payload type --- the comparator
confinement discharged at elaboration, once. The closed left-hand term
is built by verified substitution (\texttt{polyClosedTerm}, whose
evaluation lemma \texttt{polyClosedTerm\_evalQ} says the
checker\textquotesingle s exact replay computes
\emph{the value $V(x_0)$ itself}), \texttt{rootEvalFormula\_true\_iff}
is Proposition~\ref{prop:re-exact}\textquotesingle s identity at the
wire (the $\Q$ face --- both directions, as the proposition demands),
and \texttt{rootEvalElab\_sound} is the factoring through
\texttt{sound\_pConstExact}, itself one step to
\texttt{Formula.holds\_of\_evalClosed\_true}. The box branch is not
elaborated, by design (Recipe~\ref{rec:re-template}: registry
recompute, not a rule application).
\end{bremark}

\begin{bremark}[the inherited comparator confinement]
\rulename{PConstExact} excludes one comparator ($\neq$) from its
admissible set, for reasons of a second soundness invariant unpacked in
that rule\textquotesingle s section. The membership leaf uses $\le$
only --- comfortably inside --- but any port or mechanization of this
rule must \emph{preserve} the exclusion rather than re-derive the
comparator set from the semantic lemma alone; the flag travels with
the alias.
\end{bremark}

\begin{anchors}
target lemma (Prop.~\ref{prop:re-sound}) & \texttt{Formula.holds\_of\_evalClosed\_true}; \texttt{sound\_pConstExact} (quoted in the \rulename{PConstExact} section)\\
closedness implied at the leaf & \texttt{Formula.evalClosed} ($\mathsf{none}$ on any wire; same section)\\
exactness identities & \texttt{rootEval\_eval\_cast}, \texttt{rootEval\_sublevel\_iff}, \texttt{rootEval\_box\_iff} (bound in the entry)\\
payload-level elaborator (\autoref{rem:re-elab}) & \texttt{rootEvalElab} + \texttt{rootEvalElab\_sound}, over \texttt{polyClosedTerm} + \texttt{polyClosedTerm\_evalQ} + \texttt{rootEvalFormula\_true\_iff} (module \texttt{BernsteinLogic.QtslElaborators}, bound in the entry)\\
box-atom branch & checker-side registry recompute, by design --- no rule anchor\\
consumption site & the bridge rule\textquotesingle s eval field (\rulename{QSatFromRegion} section)\\
\end{anchors}


\rulesection{68}{BoxBern}{admissible rule}
\label{blcatrule-BoxBern}
\ruleblock{The rule (Steele metanotation)}
\sourcenote{as realized in the mechanized development; source row: \texttt{\texttt{QTSL-rules.toml}} --- \emph{BoxBern}. The metanotation statement is typeset by hand for notational quality; its mechanized ground truth is the Lean statement below, extracted from the sources at every build of this book and hash-bound. The rule \emph{roster} is machine-checked (every mechanized rule tag round-trips to exactly one catalogue row); the statement \emph{text} is editorial, checked by eye against the extracted Lean beside it.}
\begin{mathpar}
\inferrule*[Right=\rulename{BoxBern}]{
  N \in \Q[z_1, \ldots, z_n]; \quad D = \textstyle\prod_{v} [a_v, b_v] \subseteq \Q^n
  \\
  \varphi = (A,\ A^{-1},\ b): \quad A,\ A^{-1} : \Q^{[n] \times [n]},\ \ b : \Q^{[n]}; \qquad \varphi(w) = A\,w + b \ \text{(affine-rung witness; } \varphi \ \text{omitted reads as } \mathrm{id}\text{)}
  \\
  \vdash A^{-1} A = 1_{[n]} \ \text{(the one certified product, exact)}
  \\
  \hat N = \varphi \fatsemi N \ \text{(exact substitution)}
  \\
  \bigl\{b_\alpha \le 0\bigr\}_{\alpha} \ \text{(every tensor-product Bernstein coefficient of } \hat N \ \text{over } D\text{)}
}{
  D \vdash \hat N \le 0
}
\end{mathpar}

\ruleflavor{the witness is data with index structure --- two square rational tensors and an offset vector, the declared inverse carried rather than computed, its sole obligation the exact product (over a field, one side suffices) --- and the ladder above the affine rung is graded by each rung\textquotesingle s own grader: an unimplemented rung is a named refusing variant, never a partial check. The pullback $\hat N = \varphi \fatsemi N$ is recomputed by the checker and never rides the certificate, so a laundered pullback is unrepresentable rather than detected. The enacted rule is stated from its landed realization (\texttt{check\_box\_row}/\texttt{check\_wire\_family} in the workshop checker, over the wire-family grammar; refusal witness $(\alpha, b_\alpha)$) --- its certificate content is out of scope of the current baked theory: a box-cell row family is the named successor, and this cell moves onto the theory row when a revision carries it.}
\par\medskip
\rulename{BoxBern} is the tensor-product (box) sibling of
\rulename{BernCone}: where \rulename{BernCone} certifies a polynomial
predicate over a simplex-shaped region by exact Bernstein enclosure,
\rulename{BoxBern} certifies $\hat N \le 0$ over an axis-aligned
rational box, with the polynomial optionally pulled back through a
certified reparameterization witness first. A producer reaches for it
wherever certificate cells are boxes rather than simplices --- grid
covers, dyadic bisection cells, per-cell contraction rows for discrete
maps --- because the obligation count is the product of per-axis
degrees rather than the simplex count times a multinomial: at the
worked fixture\textquotesingle s shape (four variables, per-axis
degrees $(2,4,2,4)$, total degree $8$), one box costs $225$ exact
coefficients where the Kuhn-simplex decomposition of the same box
costs $24 \cdot 495 = 11{,}880$ --- a structural gap that compounds
with per-coefficient cost, measured at three orders of magnitude
wall-clock on the acceptance fixture.

What the checker recomputes, given the row $(N, \varphi?, D)$: that
the reparameterization witness certifies (for an affine witness,
shapes and the exact left-inverse identity $A^{-1}A = 1$; later ladder
rungs by their own graders; an unimplemented rung is a \emph{named
refusing variant}, never a partial check); the pullback
$\hat N = \varphi \fatsemi N$ by exact substitution --- no pullback ever
rides the certificate, so a laundered $\hat N$ is unrepresentable
rather than detected; and then the enclosure itself: substitute
$x_v = a_v + (b_v - a_v)\,t_v$ per axis, change basis per axis at the
polynomial\textquotesingle s actual degrees, and test every
tensor-product Bernstein coefficient against $\le 0$, refusing with
the first offending multi-index and its exact value. Everything is
exact rational arithmetic; nothing consults a float.

Each hypothesis is load-bearing. Drop the slot certification and an
uninvertible ``witness\textquotesingle\textquotesingle{} could smuggle
a coordinate change under which the claim means something else
entirely. Drop the recompute-the-pullback discipline and the checked
polynomial decouples from the claimed one. Drop the per-coefficient
sweep and the conclusion is unrelated to the box. The conclusion is
deliberately narrow --- one sign fact on one box --- because the row
\emph{species} that consumes it (coverage of the operative region,
S-procedure multipliers $\lambda \ge 0$, init/containment side rows)
owns every regional and compositional obligation; this rule is the
arithmetic leaf they all stand on.

\cataloguenote{The enclosure bound is conservative --- coefficients overshoot the
true range except at corners, where they interpolate --- and
completeness is recovered by subdivision exactly as for the simplex
rule: \rulename{Subdivide} remains the completeness partner. Degrees
enter the certified statement as \emph{data}: any degree family
admitting a valid tensor-Bernstein expansion encloses, so ``use the
actual degrees\textquotesingle\textquotesingle{} is a checker cost
policy, not a soundness premise. Degenerate axes ($a_v = b_v$) are
inside the statement --- the substitution sends the axis to a
constant and the enclosure is unaffected.
}
\ruleblock{Lean statement}
\declhead{BernsteinLogicAux.BoxBern}{boxBern\_enclosure}{\badge{BadgeGreen}{statement bound}}
\begin{leansrc}
/-- **The tensor-basis enclosure lemma** (soundness part (ii) of the
box-Bernstein rule): if the box-shifted claim polynomial
`N ∘ (lo + (hi − lo)·u)` has tensor-Bernstein coefficients `b`, then any
bounds on `b` bound `N` on the box. `hexp` is the checker's exact-`ℚ`
replay obligation (the shift and basis-conversion passes); the conclusion
quantifies over the landed arena's box set. -/
theorem boxBern_enclosure (lo hi : Fin n → ℚ) (N : MvPolynomial (Fin n) ℚ)
    (d : Fin n → ℕ) (b : (∀ v, Fin (d v + 1)) → ℚ) (m M : ℚ)
    (hexp : boxShift lo hi N = ∑ α, C (b α) * tensorBern d α)
    (hm : ∀ α, m ≤ b α) (hM : ∀ α, b α ≤ M) :
    ∀ x ∈ (ratBox lo hi).set, (m : ℝ) ≤ aeval x N ∧ aeval x N ≤ (M : ℝ)
\end{leansrc}
\declhead{BernsteinLogicAux.BoxBern}{boxBern\_enclosure\_minmax}{\badge{BadgeGreen}{statement bound}}
\begin{leansrc}
/-- The doc-verbatim min/max form (part (ii) as displayed):
`min_α b_α ≤ N(x) ≤ max_α b_α` on the box. -/
theorem boxBern_enclosure_minmax (lo hi : Fin n → ℚ)
    (N : MvPolynomial (Fin n) ℚ)
    (d : Fin n → ℕ) (b : (∀ v, Fin (d v + 1)) → ℚ)
    (hexp : boxShift lo hi N = ∑ α, C (b α) * tensorBern d α) :
    ∀ x ∈ (ratBox lo hi).set,
      ((Finset.univ.inf' Finset.univ_nonempty b : ℚ) : ℝ) ≤ aeval x N ∧
        aeval x N ≤ ((Finset.univ.sup' Finset.univ_nonempty b : ℚ) : ℝ)
\end{leansrc}
\declhead{BernsteinLogicAux.BoxBern}{boxBern\_nonpos}{\badge{BadgeGreen}{statement bound}}
\begin{leansrc}
/-- **The checker-facing sign corollary** (the rule's accept condition:
"accept iff every coefficient `b_α ≤ 0`"): nonpositive coefficients certify
`N ≤ 0` on the box — exactly one `hcells` premise of `Repar.t0_soundness` /
`t05_soundness_gridcx`, over the same `GridCx.Box` object. -/
theorem boxBern_nonpos (lo hi : Fin n → ℚ) (N : MvPolynomial (Fin n) ℚ)
    (d : Fin n → ℕ) (b : (∀ v, Fin (d v + 1)) → ℚ)
    (hexp : boxShift lo hi N = ∑ α, C (b α) * tensorBern d α)
    (hcoeff : ∀ α, b α ≤ 0) :
    ∀ x ∈ (ratBox lo hi).set, aeval x N ≤ 0
\end{leansrc}
\declhead{BernsteinLogicAux.BoxBern}{aeval\_boxShift}{\badge{BadgeGreen}{statement bound}}
\begin{leansrc}
/-- **Per-axis affine substitution exactness** (doc (ii) proof-shape item
"per-axis affine substitution exactness — a diagonal T0 map"): evaluating
the shifted polynomial at `u` is evaluating the original at the shifted
point. -/
theorem aeval_boxShift (lo hi : Fin n → ℚ) (u : Fin n → ℝ)
    (N : MvPolynomial (Fin n) ℚ) :
    aeval u (boxShift lo hi N)
      = aeval (fun v => (lo v : ℝ) + ((hi v : ℝ) - (lo v : ℝ)) * u v) N
\end{leansrc}
\declhead{BernsteinLogicAux.BoxBern}{sum\_aeval\_tensorBern}{\badge{BadgeGreen}{statement bound}}
\begin{leansrc}
/-- Sub-lemma 2c (partition of unity): the tensor basis sums to `1` on the
unit cube — "products of partitions of unity are partitions of unity"
(doc (ii)); univariate `bernsteinPolynomial.sum` distributed over the
product by `Fintype.prod_sum`. -/
theorem sum_aeval_tensorBern (d : Fin n → ℕ) (u : Fin n → ℝ) :
    ∑ α : ∀ v, Fin (d v + 1), aeval u (tensorBern d α) = 1
\end{leansrc}
\ruleblock{Correspondence (Steele display $\leftrightarrow$ Lean)}
\badge{BadgeGreen}{reviewer-certified against the displayed texts}
\sourcenote{The dictionary below is a reviewer-certified statement of how the mechanized statements above interpret the displayed rule: display metavariable $\mapsto$ Lean binder, each display hypothesis $\mapsto$ the premise or side condition carrying it, the display conclusion $\mapsto$ the goal --- with any deliberate difference in width between display and mechanization named outright. Certification is recorded as a hash triple over (the displayed cell, the extracted statement texts, this dictionary\textquotesingle s own text); if any of the three changes, this block turns amber until a reviewer re-certifies. The correspondence claim itself remains editorial --- reviewer-grade, not machine-checked; its machine-checked successor types each soundness statement against the rule as a structured object, and supersedes this block per rule as it lands.}
The displayed rule --- certified slot, exact pullback, coefficient
sweep, conclusion $\hat N \le 0$ on $D$ --- maps onto the bound
statements as follows. $D \mapsto$ \texttt{lo hi : Fin n → ℚ} through
\texttt{ratBox} (the same \texttt{GridCx.Box} object the transport
theorems quantify over); $\hat N \mapsto$ \texttt{N : MvPolynomial
(Fin n) ℚ} --- the \emph{pulled-back} polynomial, because the checker
recomputes the pullback and hands the enclosure lemma its result: the
slot legs (certification, replay) are consumed upstream by the
transport theorems (\texttt{t0\_soundness},
\texttt{t05\_soundness\_gridcx}, and the T1/TC rung lemmas), quoted
and bound in their own sections, never re-stated here. One
orientation divergence in the slot leg is named rather than left to
the reader: the enacted checker verifies the \emph{displayed} product
--- its single exact multiplication computes $A^{-1} A$ and compares
it with $1_{[n]}$ (\texttt{affine\_left\_certified}) --- whereas the
Lean witness structure stores the \emph{opposite} orientation
(\texttt{AffineWitness.certified} $: A \cdot A^{-1} = 1$) and derives
the displayed one as the theorem \texttt{left\_certified}
$: A^{-1} \cdot A = 1$ (via \texttt{mul\_eq\_one\_comm}; the matrices
are square over the commutative ring $\Q$, so either product
determines the other and the two orientations are equivalent). The
displayed coefficient table $\mapsto$ the expansion hypothesis of
\texttt{boxBern\_enclosure}: an explicit representation of the
\texttt{boxShift}-substituted polynomial as
$\sum_\alpha b_\alpha \prod_v B^{d_v}_{\alpha_v}$ --- the
u-space/x-space transport is a theorem (\texttt{aeval\_boxShift}),
not notation, and the degree family $d$ is a binder, not a fixed
choice. The displayed accept condition (every $b_\alpha \le 0$)
$\mapsto$ the hypothesis of \texttt{boxBern\_nonpos}, whose conclusion
is verbatim one \texttt{hcells} premise of the transport theorems ---
the two compose without glue, which is the certified content of
``exactly one premise of the landed transport
theorems.\textquotesingle\textquotesingle{} Width note, two-sided:
\texttt{boxBern\_enclosure\_minmax} is strictly stronger than the
accept condition needs (both range ends, where the rule consumes only
the max side), and nothing certified is wider than the display ---
the slot legs live upstream by construction, and no coverage or
multiplier obligation enters any binder here (the row species owns
them, as the display says).

\ruleblock{Lean verifier receipts}
\sourcenote{module \texttt{BernsteinLogicAux.BoxBern}; bound build \texttt{110ba4bda664\ldots} (bound from receipt 6935ac926dca961f, 2026-08-03T20:13:44Z). A receipt records that named checking layers accepted a build: (module, compiled-artifact SHA-256, environment SHA-256, checker layers), minted when the checking gate ran; the store receipts attesting the bound build are listed with the layers each run credited. Nothing here asks the reader to trust the store: every layer is an independently runnable program over the quoted sources --- rebuild the module and re-run the checkers to reproduce the verdict. Receipts attesting superseded builds credit no layer and are omitted from this rendering (the count of omissions, if any, is noted below).}
\begin{receipttable}
\receiptrow{63a3e12578528f22}{2026-08-02T02:06:52Z}{build+decorrelated}{\badge{BadgeGreen}{current}}
\receiptrow{c2fe5c957feb2507}{2026-08-02T23:30:40Z}{build+decorrelated}{\badge{BadgeGreen}{current}}
\receiptrow{665f63f0c0026da8}{2026-08-03T16:26:44Z}{build+decorrelated}{\badge{BadgeGreen}{current}}
\receiptrow{6935ac926dca961f}{2026-08-03T20:13:44Z}{build+decorrelated}{\badge{BadgeGreen}{current}}
\receiptrow{a3e7c09b848c5808}{2026-08-03T22:45:28Z}{build+decorrelated}{\badge{BadgeGreen}{current}}
\end{receipttable}
\sourcenote{8 further receipt(s) in the store attest superseded builds of this module and are omitted from this rendering.}
\layerline{\layeritem{build}{lake build (elaboration)}{\badge{BadgeGreen}{current}}\layeritem{replay}{leanchecker (official kernel replay)}{\badge{BadgeAmber}{stale only}}\layeritem{hardened}{safe\_verify (hardened screen)}{\badge{BadgeAmber}{stale only}}\layeritem{decorrelated}{lean4export + nanoda (decorrelated kernel)}{\badge{BadgeGreen}{current}}}
\sourcenote{\badge{BadgeGray}{comparator: none} no fixture in the three-implementation comparator harness (compiled / kernel-interpreted / reference-interpreted) exercises this rule.}
\ruleblock{HOL4 statement}
\absentbox{NO HOL STATEMENT YET}{The HOL family port (\texttt{checkFam\_boxbern} + a \texttt{famEq}
arm, the A1/A3 mechanism) is a named successor of the kernel leg that
landed this rule --- absence recorded at landing, with the mechanism
that will discharge it, so this reason cannot silently go stale the
way pre-wave absence reasons did.
}
\ruleblock{HOL verifier receipts}
\absentbox{NONE}{no HOL verification receipt exists for this rule.}
\ruleblock{Proof}

\sourcenote{Genre: admissible rule (tier 2) --- admissible by
\emph{reduction to verified checking machinery}, the taxonomy\textquotesingle s
recorded boundary case. The reduction shape is stated first; the
enclosure proposition it stands on is then proved in full, because it
is the section\textquotesingle s one piece of genuinely new
mathematics --- the tensor twin of the simplex grounding that
\rulename{BernCone} already carries.}

\begin{brecipe}[the reduction --- what the checker replays]
\label{rec:bb-reduction}
A \rulename{BoxBern} row is $(N,\ \varphi?,\ D)$: a claim polynomial
$N \in \Q[z_1, \dots, z_n]$, an optional reparameterization witness
$\varphi$ from the ladder ($\varnothing$ = identity), and an
axis-aligned box $D = \prod_v [a_v, b_v] \subseteq \Q^n$ in the
witness\textquotesingle s working coordinates. The discharge is three
replays, all exact $\Q$, each standing on a named verified theorem
rather than on catalogue structure:
\begin{enumerate}
\item \emph{slot certification} --- for an affine witness, shapes and
  the exact left-inverse identity $A^{-1}A = 1$ (a left inverse
  suffices over a field); later ladder rungs certify by their own
  graders, and a rung without a grader is a \emph{named refusing
  variant} --- refused loudly, never half-checked.
\item \emph{substitution replay} --- $\hat N = \varphi \fatsemi N$ by
  exact substitution. No pullback field exists in the row, so a forged
  $\hat N$ is unrepresentable rather than detected.
\item \emph{enclosure} --- the tensor-product Bernstein coefficients
  of $\hat N$ over $D$ (Proposition~\ref{prop:bb-enclosure}); accept
  iff every coefficient is $\le 0$, refusing with the witness
  $(\alpha, b_\alpha)$ otherwise.
\end{enumerate}
The conclusion, $\hat N \le 0$ on $D$, is deliberately exactly one
\texttt{hcells} premise of the landed transport theorems: coverage of
the operative region, S-procedure multipliers, and init/containment
side rows are owned by the consuming row species, never by this leaf.
\end{brecipe}

\begin{bprop}[tensor-product Bernstein enclosure]\label{prop:bb-enclosure}
Let $p \in \Q[x_1, \dots, x_n]$, let $D = \prod_v [a_v, b_v]$ with
rational endpoints ($a_v \le b_v$, degenerate axes allowed), and
suppose that after the per-axis affine substitution
$x_v = a_v + (b_v - a_v)\,u_v$ the resulting polynomial admits the
tensor-product Bernstein representation
\[
p\bigl(a + (b - a)\odot u\bigr) \;=\; \sum_{\alpha}\, b_\alpha
\prod_{v=1}^{n} B^{d_v}_{\alpha_v}(u_v),
\qquad
B^{d}_{i}(t) = \binom{d}{i}\, t^{i} (1-t)^{d-i},
\]
for some degree family $(d_v)$ and coefficients
$b_\alpha \in \Q$, $\alpha \in \prod_v \{0, \dots, d_v\}$. Then for
every $x \in D$,
\[
\min_\alpha b_\alpha \;\le\; p(x) \;\le\; \max_\alpha b_\alpha .
\]
In particular, if every $b_\alpha \le 0$ then $p \le 0$ on $D$.
\end{bprop}
\begin{proof}
Fix $x \in D$. Each coordinate lies in its interval, so there is a
$u \in [0,1]^n$ with $x_v = a_v + (b_v - a_v)u_v$ for every axis: on a
non-degenerate axis take $u_v = (x_v - a_v)/(b_v - a_v)$; on a
degenerate axis ($a_v = b_v$) the substitution is the constant $a_v$
regardless of $u_v$, and $u_v = 0$ serves. Substitution exactness ---
that evaluating $p$ at this $x$ equals evaluating the substituted
polynomial at $u$ --- is a per-axis affine replay, not a notational
identification (the bound statement keeps the two polynomials
distinct and transports along the substitution lemma).

On $[0,1]$, each Bernstein basis polynomial is nonnegative
($t^i \ge 0$, $(1-t)^{d-i} \ge 0$, binomial coefficients positive) and
each degree-$d$ family sums to $1$: $\sum_i B^d_i(t) =
\bigl(t + (1-t)\bigr)^d = 1$ by the binomial theorem. Products of
partitions of unity are partitions of unity, so the tensor family
$\bigl\{\prod_v B^{d_v}_{\alpha_v}(u_v)\bigr\}_\alpha$ is a family of
nonnegative weights summing to $1$ at every $u \in [0,1]^n$.
Evaluating the representation at $u$ therefore exhibits $p(x)$ as a
convex combination of the $b_\alpha$, and a convex combination lies
between the minimum and the maximum of its generators. The accept
condition is the corollary: if every generator is $\le 0$, so is
every convex combination.
\end{proof}

\begin{bremark}[degrees are data; degenerate axes are inside the statement]
Nothing in Proposition~\ref{prop:bb-enclosure} privileges the
polynomial\textquotesingle s actual per-axis degrees: \emph{any}
degree family admitting a valid representation encloses, so the
checker\textquotesingle s ``use the actual degrees after the
shift\textquotesingle\textquotesingle{} is a cost policy --- smallest
coefficient table --- and never a soundness premise. Likewise a
degenerate axis needs no exclusion: the substitution collapses it to
a constant, degrees drop, and the convex-combination argument is
untouched. Both points are formalization deltas made \emph{explicit}
in the bound statements, in the conservative direction --- the
expansion hypothesis speaks about the substituted polynomial (the
u-space/x-space transport is a theorem), and the degree family is a
binder.
\end{bremark}

\begin{bremark}[conservatism, and where completeness comes from]
The bound is not tight: coefficients overshoot the true range except
at the corners of $D$, where the representation interpolates
($B^d_i(0) = [i = 0]$, $B^d_i(1) = [i = d]$, so corner coefficients
\emph{equal} corner values). A true-but-refused row is repaired by
subdivision --- Bernstein enclosures converge under subdivision and
degree elevation --- so \rulename{Subdivide} is the completeness
partner here exactly as for the simplex rule, and refusal is always
the failure direction.
\end{bremark}

\begin{bremark}[why a box rule at all --- the obligation-count structure]
At per-axis degrees $(d_v)$, one box costs $\prod_v (d_v + 1)$
coefficients; the kernel-checkable simplex shape for the same box
costs $n!$ Kuhn simplices at $\binom{d + n}{n}$ coefficients each
(total degree $d$). At the acceptance fixture\textquotesingle s shape
--- $n = 4$, degrees $(2,4,2,4)$, $d = 8$ --- that is $3 \cdot 5
\cdot 3 \cdot 5 = 225$ against $24 \cdot \binom{12}{4} = 11{,}880$
per cell, and the per-coefficient work also drops (one-dimensional
passes against multinomial walks). The measured same-harness gap on
the fixture\textquotesingle s 26 cells is three orders of magnitude
wall-clock. None of this changes what is certified --- the rule is an
optimization of checker cost, not a widening of any claim.
\end{bremark}

\begin{anchors}
Recipe \ref{rec:bb-reduction} slot leg & \texttt{t0\_soundness} / \texttt{t05\_soundness\_gridcx} (their sections); later rungs per their graders, refusing variants named\\
Recipe \ref{rec:bb-reduction} replay leg & pullback recomputed, never carried --- laundering unrepresentable by row shape\\
Proposition \ref{prop:bb-enclosure} & \texttt{boxBern\_enclosure} (bound form), \texttt{boxBern\_enclosure\_minmax} (verbatim min/max form)\\
substitution exactness & \texttt{aeval\_boxShift} (per-axis affine replay; degenerate axes via the unit-scalar witness)\\
partition of unity & \texttt{sum\_aeval\_tensorBern} (tensor family sums to 1)\\
the accept corollary & \texttt{boxBern\_nonpos} --- conclusion verbatim one \texttt{hcells} premise of the transport theorems\\
\end{anchors}


\rulesection{69}{PExpTypeEnclosure}{admissible rule}
\label{blcatrule-PExpTypeEnclosure}
\ruleblock{The rule (Steele metanotation)}
\sourcenote{as realized in the mechanized development; source row: \texttt{\texttt{QTSL-rules.toml}} --- \emph{PExpTypeEnclosure}. The metanotation statement is typeset by hand for notational quality; its mechanized ground truth is the Lean statement below, extracted from the sources at every build of this book and hash-bound. The rule \emph{roster} is machine-checked (every mechanized rule tag round-trips to exactly one catalogue row); the statement \emph{text} is editorial, checked by eye against the extracted Lean beside it.}
\begin{mathpar}
\inferrule*[Right=\rulename{PExpTypeEnclosure}]{
  f \in F \quad n \in \mathbb{N} \quad r, R \in \Q
  \\
  a_k = a^f_k \ \text{for } k = 0..n \ \text{(exact coefficient recurrence)}
  \\
  0 \le r < n + 2
  \qquad
  R \;\ge\; \frac{r^{n+1}}{(n+1)!}\cdot\frac{n+2}{\,n+2-r\,}
}{
  \forall x,\ |x| \le r \;\Rightarrow\; \Bigl|f(x) - \textstyle\sum_{k \le n} a_k x^k\Bigr| \le R
}
\end{mathpar}

\ruleflavor{the coefficients are recomputed by the checker (never carried); the check is pure $\Q$-arithmetic and never evaluates $f$; the real functions enter only in the kernel-checked soundness theorem. Scope of $F$: entire functions with exact-$\Q$ Maclaurin coefficients dominated by $|a_k| \le 1/k!$ (entire of exponential type, normalized) --- $\exp$, $\sin$, $\cos$, and (qtsl tower) the reduced error-function kernel; a new in-class function costs one coefficient recurrence + one domination case, the tail argument untouched. Finite-radius functions ($\log$, rational powers) are certified by sibling routes, not by this rule; every downstream consumption decomposes into existing polynomial sequents (recentering, Taylor-model combination, endpoint substitution).}
\par\medskip
\rulename{PExpTypeEnclosure} is the calculus\textquotesingle s one
transcendental leaf: a remainder-bound rule that brackets a
non-polynomial function between two polynomials with all data rational,
so that every downstream obligation --- recentering, Taylor-model
combination, endpoint substitution into a generator inequality ---
decomposes into polynomial sequents the existing rules already check.
A producer reaches for it wherever certified dynamics carry
transcendental drift terms (pendulum-family $\sin/\cos$, rate and decay
$\exp$, Gaussian-kernel reductions): the rule certifies a Maclaurin
truncation once, at a radius, with an explicit remainder, and from
there the trusted arithmetic is exact rational computation exactly as
for every other leaf.

The witness is $(f, n, r, R)$: a function tag $f$ from the registered
library, a truncation degree $n$, a radius $r \in \Q_{\ge 0}$, and a
remainder budget $R \in \Q$. The check is pure rational arithmetic ---
the coefficients $a^f_k$ are recomputed from $(f, n)$ by exact per-tag
recurrence rather than carried (a laundered coefficient list is
unrepresentable on the wire), the radius must satisfy $r < n+2$, and
$R$ must dominate the closed-form tail bound
$\smash{\frac{r^{n+1}}{(n+1)!}\cdot\frac{n+2}{n+2-r}}$. The checker
never evaluates $f$; the real functions enter only in the
kernel-checked soundness theorem, which concludes the two-sided
bracket $|f(x) - \sum_{k \le n} a_k x^k| \le R$ on $[-r, r]$.

The scope is deliberately exact: \emph{entire functions with
exact-rational Maclaurin coefficients dominated by
$|a_k| \le 1/k!$} --- entire of exponential type, normalized so
$e^{|x|}$ dominates coefficientwise. All three conditions are
load-bearing: exact-$\Q$ recurrences make clause (i) checkable
term-by-term; the $1/k!$ domination is what lets ONE tail formula
serve every tag; entirety means any rational radius works, subject
only to the formula\textquotesingle s own guard $r < n+2$. This is
strictly narrower than ``anything with a closed-form Taylor
series\textquotesingle\textquotesingle: $\log$ and rational powers are
analytic but not entire (a finite-radius variant would need a
radius-ratio geometric tail and a domain side condition; they are
certified today by sibling routes built on their own bounds), and
$\tan$ is meromorphic (handled as $\sin/\cos$ plus division at the
consumer). Extending the library to a new in-class function costs one
coefficient recurrence and one domination case --- the tail argument
is untouched. The error-function kernel exercised exactly this clause:
under the $y = z^2$ reduction its coefficients are
$(-1)^k / (k!\,(2k+1))$, dominated by $1/k!$ verbatim, and the
soundness statement needed no edits beyond the per-tag cases.

\cataloguenote{Admissible by reduction to verified checking machinery (the same
boundary clause as \rulename{BoxBern}): the checker replays
\texttt{encCheck}, and the trust lands on the kernel-checked meaning
theorems, never on catalogue structure. Conservativity: no new sequent
theorems --- the rule concludes into a sibling judgment kind, and its
consumers re-enter the calculus as ordinary polynomial rows. Stated
outright rather than smoothed: this is the one rule where
transcendental meaning enters the trust story --- a deliberate,
priced, single-rule growth of the trusted base, carried by the
soundness theorems cited below. Slack is the producer\textquotesingle s
business (pick $n$, recenter, subdivide): subdivision shrinks bracket
width first-order while degree shrinks it factorially, and the
measured budgets on the acceptance systems ran at degree $9$ with
remainders near $10^{-7}$ --- orders below the certificate margins
they feed.
}
\ruleblock{Lean statement}
\declhead{BernsteinLogicAux.Enclosure}{encCheck\_sound}{\badge{BadgeGreen}{statement bound}}
\begin{leansrc}
/-- **The remainder-bound rule is sound** (note §1's boxed rule): a witness passing
`encCheck` certifies `|f(x) − P(x)| ≤ R` for every `|x| ≤ r`, where `f` is the tag's
real function and `P` the witness polynomial. -/
theorem encCheck_sound (w : EncWitness) (h : encCheck w = true) :
    ∀ x : ℝ, |x| ≤ (w.r : ℝ) →
      |w.tag.eval x - evalPoly w.coeffs x| ≤ (w.R : ℝ)
\end{leansrc}
\declhead{BernsteinLogicAux.Enclosure}{encCheck\_sound\_bracket}{\badge{BadgeGreen}{statement bound}}
\begin{leansrc}
/-- Two-sided bracket form: `P(x) − R ≤ f(x) ≤ P(x) + R` on `[−r, r]` — the pair of
polynomial brackets the certificate's sequents consume (note §3). -/
theorem encCheck_sound_bracket (w : EncWitness) (h : encCheck w = true) :
    ∀ x : ℝ, |x| ≤ (w.r : ℝ) →
      evalPoly w.coeffs x - (w.R : ℝ) ≤ w.tag.eval x ∧
      w.tag.eval x ≤ evalPoly w.coeffs x + (w.R : ℝ)
\end{leansrc}
\declhead{BernsteinLogicAux.Enclosure}{encCheck\_point}{\badge{BadgeGreen}{statement bound}}
\begin{leansrc}
/-- Point-bracket corollary (producer recentering, note §3(1)): evaluating the witness at
a rational point `c` with `|c| ≤ r` brackets the transcendental constant `f(c)`. -/
theorem encCheck_point (w : EncWitness) (h : encCheck w = true) (c : ℚ)
    (hc : |c| ≤ w.r) :
    evalPoly w.coeffs (c : ℝ) - (w.R : ℝ) ≤ w.tag.eval (c : ℝ) ∧
    w.tag.eval (c : ℝ) ≤ evalPoly w.coeffs (c : ℝ) + (w.R : ℝ)
\end{leansrc}
\declhead{BernsteinLogicAux.Enclosure}{linear\_interp\_le}{\badge{BadgeGreen}{statement bound}}
\begin{leansrc}
/-- Generator VCs are linear in each transcendental leaf occurrence: if `A + C·w ≤ 0`
holds at both endpoints of `[l, h]`, it holds for every `w ∈ [l, h]`. With this, per-cell
interval hypotheses for the enclosed leaf discharge as TWO ordinary polynomial sequents —
no new wire dimension. -/
theorem linear_interp_le {A C l h w : ℝ} (hlw : l ≤ w) (hwh : w ≤ h)
    (hl : A + C * l ≤ 0) (hh : A + C * h ≤ 0) : A + C * w ≤ 0
\end{leansrc}
\declhead{QtslAux.Enclosure}{encCheck\_sound}{\badge{BadgeGreen}{statement bound}}
\begin{leansrc}
theorem encCheck_sound (w : EncWitness) (h : encCheck w = true) :
    ∀ x : ℝ, |x| ≤ (w.r : ℝ) →
      |w.tag.eval x - evalPoly w.coeffs x| ≤ (w.R : ℝ)
\end{leansrc}
\declhead{QtslAux.Enclosure}{encCheck\_sound\_bracket}{\badge{BadgeGreen}{statement bound}}
\begin{leansrc}
theorem encCheck_sound_bracket (w : EncWitness) (h : encCheck w = true) :
    ∀ x : ℝ, |x| ≤ (w.r : ℝ) →
      evalPoly w.coeffs x - (w.R : ℝ) ≤ w.tag.eval x ∧
      w.tag.eval x ≤ evalPoly w.coeffs x + (w.R : ℝ)
\end{leansrc}
\declhead{QtslAux.Enclosure}{integral\_exp\_neg\_sq}{\badge{BadgeGreen}{statement bound}}
\begin{leansrc}
/-- The Gaussian integral reduces to the error-function kernel by the
substitution `t = z * s`: `∫ t in (0:ℝ)..z, exp (-(t ^ 2)) = z * T (z ^ 2)`.
Both sides are odd in `z`, and the oriented interval integral makes the
identity hold for every real `z` without a sign case split. Together with a
bracket for `2 / √π` (see `two_div_sqrt_pi_mem`) this carries polynomial
enclosures of `T` to enclosures of the error function. -/
theorem integral_exp_neg_sq (z : ℝ) :
    ∫ t in (0:ℝ)..z, Real.exp (-(t ^ 2)) = z * FnTag.erfT.eval (z ^ 2)
\end{leansrc}
\ruleblock{Correspondence (Steele display $\leftrightarrow$ Lean)}
\badge{BadgeGreen}{reviewer-certified against the displayed texts}
\sourcenote{The dictionary below is a reviewer-certified statement of how the mechanized statements above interpret the displayed rule: display metavariable $\mapsto$ Lean binder, each display hypothesis $\mapsto$ the premise or side condition carrying it, the display conclusion $\mapsto$ the goal --- with any deliberate difference in width between display and mechanization named outright. Certification is recorded as a hash triple over (the displayed cell, the extracted statement texts, this dictionary\textquotesingle s own text); if any of the three changes, this block turns amber until a reviewer re-certifies. The correspondence claim itself remains editorial --- reviewer-grade, not machine-checked; its machine-checked successor types each soundness statement against the rule as a structured object, and supersedes this block per rule as it lands.}
The displayed rule maps onto the bound statements as follows. The
witness $(f, n, r, R) \mapsto$ \texttt{EncWitness} (fields
\texttt{func}, \texttt{deg}, \texttt{radius}, \texttt{rem}, plus the
coefficient list governed by the displayed recurrence clause,
recompute semantics per the flavor text); the three displayed check clauses $\mapsto$ the definition
of \texttt{encCheck}, whose acceptance is the single hypothesis of
every soundness statement. The displayed conclusion $\mapsto$
\texttt{encCheck\_sound} (the absolute-value form) and
\texttt{encCheck\_sound\_bracket} (the same content as a two-sided
polynomial bracket --- equal width, different packaging);
\texttt{encCheck\_point} instantiates the bracket at a rational point
(the recentering step\textquotesingle s constant legs), and
\texttt{linear\_interp\_le} carries the endpoint-substitution step the
consumers use. One deliberate width difference, named outright: on the
wire the coefficient vector is never serialized --- the checker
recomputes it --- while the Lean witness carries the list explicitly
and clause (i) checks it against the recurrence; the wire grain is
therefore strictly narrower than the mechanized statement, and the
coefficient-mutation receipts live kernel-side precisely because those
mutations are unrepresentable on the wire. The two Lean packages prove
the same statement shape over different tag libraries (the qtsl
package adds the error-function kernel, with its meaning stated
against the Gaussian integral and bridged by
\texttt{integral\_exp\_neg\_sq}); the HOL statements are the
side-by-side port against \texttt{transcTheory}\textquotesingle s
natively per-index series --- reindexing-free exactly where the Lean
proof needed injection reindexing.

\ruleblock{Lean verifier receipts}
\sourcenote{module \texttt{BernsteinLogicAux.Enclosure}; bound build \texttt{aeef11c4ffdf\ldots} (bound from receipt c61f02086268c593, 2026-07-21T00:02:40Z). A receipt records that named checking layers accepted a build: (module, compiled-artifact SHA-256, environment SHA-256, checker layers), minted when the checking gate ran; the store receipts attesting the bound build are listed with the layers each run credited. Nothing here asks the reader to trust the store: every layer is an independently runnable program over the quoted sources --- rebuild the module and re-run the checkers to reproduce the verdict. Receipts attesting superseded builds credit no layer and are omitted from this rendering (the count of omissions, if any, is noted below).}
\begin{receipttable}
\receiptrow{ec9147f09c619209}{2026-07-17T21:43:40Z}{build+replay+hardened+decorrelated}{\badge{BadgeGreen}{current}}
\receiptrow{ef188baecdf2d5cc}{2026-07-19T18:23:51Z}{build+decorrelated}{\badge{BadgeGreen}{current}}
\receiptrow{74824f999e9c6c1a}{2026-07-20T01:21:03Z}{build+decorrelated}{\badge{BadgeGreen}{current}}
\receiptrow{c61f02086268c593}{2026-07-21T00:02:40Z}{build+decorrelated}{\badge{BadgeGreen}{current}}
\receiptrow{5f18c0fcd7885f0e}{2026-07-22T18:26:52Z}{build+decorrelated}{\badge{BadgeGreen}{current}}
\receiptrow{0b776138b6a8c8d7}{2026-07-31T21:59:44Z}{build+decorrelated}{\badge{BadgeGreen}{current}}
\receiptrow{2ee240defc1ac528}{2026-07-31T22:28:37Z}{build+decorrelated}{\badge{BadgeGreen}{current}}
\receiptrow{1fb846e38f09ba02}{2026-08-02T23:30:45Z}{build+decorrelated}{\badge{BadgeGreen}{current}}
\receiptrow{2f17504ab4610d57}{2026-08-03T16:26:49Z}{build+decorrelated}{\badge{BadgeGreen}{current}}
\receiptrow{c53cac58e2ed77d9}{2026-08-03T22:45:33Z}{build+decorrelated}{\badge{BadgeGreen}{current}}
\end{receipttable}
\sourcenote{module \texttt{QtslAux.Enclosure}; bound build \texttt{7de2b931c19f\ldots} (bound from receipt 840791fb258b3984, 2026-07-26T15:38:06Z). A receipt records that named checking layers accepted a build: (module, compiled-artifact SHA-256, environment SHA-256, checker layers), minted when the checking gate ran; the store receipts attesting the bound build are listed with the layers each run credited. Nothing here asks the reader to trust the store: every layer is an independently runnable program over the quoted sources --- rebuild the module and re-run the checkers to reproduce the verdict. Receipts attesting superseded builds credit no layer and are omitted from this rendering (the count of omissions, if any, is noted below).}
\begin{receipttable}
\receiptrow{969915618955be8f}{2026-07-25T11:27:08Z}{build+decorrelated}{\badge{BadgeGreen}{current}}
\receiptrow{840791fb258b3984}{2026-07-26T15:38:06Z}{build+replay+hardened+decorrelated}{\badge{BadgeGreen}{current}}
\end{receipttable}
\layerline{\layeritem{build}{lake build (elaboration)}{\badge{BadgeGreen}{current}}\layeritem{replay}{leanchecker (official kernel replay)}{\badge{BadgeGreen}{current}}\layeritem{hardened}{safe\_verify (hardened screen)}{\badge{BadgeGreen}{current}}\layeritem{decorrelated}{lean4export + nanoda (decorrelated kernel)}{\badge{BadgeGreen}{current}}}
\sourcenote{\badge{BadgeGray}{comparator: none} no fixture in the three-implementation comparator harness (compiled / kernel-interpreted / reference-interpreted) exercises this rule.}
\ruleblock{HOL4 statement}
\sourcenote{theory \texttt{bern\_enclosure}, HOL4 development script \texttt{\texttt{bern\_enclosureScript.sml}} \badge{BadgeGreen}{script bound}}
\holhead{Definition encCheck\_def}
\begin{holsrc}
Definition encCheck_def:
  encCheck w ⇔
    LENGTH w.coeffs = w.n + 1 ∧
    w.coeffs = GENLIST (maclaurin w.tag) (w.n + 1) ∧
    0 ≤ w.r ∧ w.r < &w.n + 2 ∧
    tailBound w.n w.r ≤ w.R
End
\end{holsrc}
\holhead{Theorem encCheck\_sound}
\begin{holsrc}
Theorem encCheck_sound:
  ∀w x.
    encCheck w ∧ abs x ≤ real_of_rat w.r ⇒
    abs (fnEval w.tag x - evalPolyR w.coeffs x) ≤ real_of_rat w.R
\end{holsrc}
\holhead{Theorem encCheck\_sound\_bracket}
\begin{holsrc}
Theorem encCheck_sound_bracket:
  ∀w x.
    encCheck w ∧ abs x ≤ real_of_rat w.r ⇒
    evalPolyR w.coeffs x - real_of_rat w.R ≤ fnEval w.tag x ∧
    fnEval w.tag x ≤ evalPolyR w.coeffs x + real_of_rat w.R
\end{holsrc}
\ruleblock{HOL verifier receipts}
\badge{BadgeBlue}{documentary} Honest partial, stated as such: the port deliberately carries the
MEANING theorem (\texttt{encCheck\_sound} against
\texttt{transcTheory}\textquotesingle s $\exp/\sin/\cos$), unlike the
neighbouring slices whose meaning stays Lean-side. Of the
theory\textquotesingle s fifteen theorems, the eleven foundation
theorems (datatypes through the factorial bound, including the full
rational-to-real cast bridge) were machine-verified in a batch build;
the four analysis theorems (\texttt{TERM\_LE},
\texttt{MAJORANT\_SUMS}, \texttt{TAIL\_MAIN},
\texttt{encCheck\_sound} with corollaries) are assembled from
interactively-verified tactic pieces, and a full in-file batch
re-verification currently exceeds the memory of every build host used
to date ($\sim$117\,GB resident observed at failure; a
$\ge$256\,GB host is the named path). Until that build lands, the HOL
half is a stated port with interactive verification, not a
batch-verified theory --- the badge here says exactly that. The
\texttt{encReceipts} table (seven rows, expected \texttt{TTTFFFF})
replays the same seven witnesses that appear byte-identically in the
Lean kernel receipts and the reference checker\textquotesingle s
battery: three accepts (including one with $R$ exactly at the tail
bound) and four single-mutation rejects.

\sourcenote{HOL receipts are documentary: quoted results of \texttt{hol-strict} runs
and census rows from the HOL4 development\textquotesingle s build
evidence, not minted receipt artefacts. A HOL receipt store (the Lean
store\textquotesingle s twin) is named future work.
}
\ruleblock{Proof}

\sourcenote{Genre: admissible rule (tier 2) --- admissible by
\emph{reduction to verified checking machinery}, the
taxonomy\textquotesingle s recorded boundary case. The reduction is
stated first; the tail bound it stands on is then proved in full,
because it is the section\textquotesingle s one piece of genuinely new
mathematics --- and the one place in the whole calculus where the
meaning of a transcendental function enters the trusted base.}

\begin{brecipe}[the reduction --- what the checker replays]
\label{rec:enc-reduction}
A \rulename{PExpTypeEnclosure} row is $(f,\ n,\ r,\ R)$: a function
tag $f$ from the registered library $F$, a truncation degree
$n \in \N$, a radius $r \in \Q_{\ge 0}$, and a remainder budget
$R \in \Q$. The discharge is three replays, all exact $\Q$:
\begin{enumerate}
\item \emph{coefficient recompute} --- the Maclaurin prefix
  $a_0, \dots, a_n$ is regenerated from $(f, n)$ by the
  tag\textquotesingle s exact rational recurrence; no coefficient ever
  rides the certificate, so a laundered coefficient list is
  unrepresentable rather than detected.
\item \emph{radius guard} --- $0 \le r < n + 2$, the side condition
  under which the closed-form tail bound below is finite and valid.
\item \emph{remainder comparison} --- $R \ge \mathrm{tail}(n, r)$
  where
  $\mathrm{tail}(n, r) = \dfrac{r^{n+1}}{(n+1)!}\cdot
   \dfrac{n+2}{\,n+2-r\,}$,
  a comparison of two explicit rationals.
\end{enumerate}
Nothing here evaluates $f$. The conclusion --- for all real $x$ with
$|x| \le r$, $\bigl|f(x) - \sum_{k \le n} a_k x^k\bigr| \le R$, i.e.\
the two-sided polynomial bracket $P - R \le f \le P + R$ on
$[-r, r]$ --- is the statement of the kernel-checked soundness
theorem, and that theorem is where the real $\exp$, $\sin$, $\cos$
(as the proof assistant\textquotesingle s own library defines them)
enter, once, for the whole calculus.
\end{brecipe}

\begin{bprop}[uniform Maclaurin tail bound]\label{prop:enc-tail}
Let $f$ be entire with Maclaurin coefficients satisfying
$|a_k| \le 1/k!$ for all $k$ (rationality of the $a_k$ is what makes
the \emph{check} exact; the bound itself needs only the domination).
For $n \in \N$ and $0 \le r < n + 2$,
\[
\sup_{|x| \le r}\ \Bigl| f(x) - \sum_{k \le n} a_k x^k \Bigr|
\;\le\;
\frac{r^{n+1}}{(n+1)!}\cdot\frac{n+2}{\,n+2-r\,}
\;=\; \mathrm{tail}(n, r).
\]
\end{bprop}

\begin{proof}
For $|x| \le r$ the omitted tail is dominated termwise:
\[
\Bigl|\sum_{k > n} a_k x^k\Bigr|
\;\le\; \sum_{k > n} \frac{r^k}{k!}.
\]
Write $N = n + 1$ and index the tail by $j \ge 0$ as $k = N + j$. The
factorial lower bound $(N + j)! \ge N!\,(N+1)^j$ --- induction on $j$:
the base is trivial, and the step multiplies by $N + j + 1 \ge N + 1$
--- gives
\[
\frac{r^{N+j}}{(N+j)!}
\;\le\; \frac{r^N}{N!}\Bigl(\frac{r}{N+1}\Bigr)^{j}.
\]
The guard $r < n + 2 = N + 1$ makes the ratio $q = r/(N+1)$ satisfy
$0 \le q < 1$, so the geometric series sums to
$1/(1 - q) = (N+1)/(N+1-r)$, and
\[
\sum_{k > n} \frac{r^k}{k!}
\;\le\; \frac{r^{N}}{N!}\cdot\frac{N+1}{\,N+1-r\,}
\;=\; \frac{r^{n+1}}{(n+1)!}\cdot\frac{n+2}{\,n+2-r\,}. \qedhere
\]
\end{proof}

Soundness of the rule is now the two-line composition the mechanized
statements carry: clause (1) certifies that the recomputed prefix
agrees with $f$\textquotesingle s true coefficients (a per-tag
recurrence identity, proven once per library member against the
library definition of $f$); every registered tag satisfies the
domination $|a_k| \le 1/k!$ (a per-tag lemma --- for $\exp$ an
identity, for $\sin/\cos$ the odd/even zero pattern, for the reduced
error-function kernel $|a_k| = 1/(k!\,(2k+1))$); and
Proposition~\ref{prop:enc-tail} turns clauses (2)--(3) into the
bracket. Accepting a witness with $R$ \emph{exactly} at
$\mathrm{tail}(n, r)$ is sound --- the receipt battery pins one such
row on purpose --- and every strengthening of the check (a smaller
$R$, a larger $r$, a shortened prefix, a flipped coefficient sign) is
refused by construction, which is what the four mutation receipts
witness.

\begin{bremark}[scope --- what the class is, and is not]
The hypotheses of Proposition~\ref{prop:enc-tail} are the exact scope
of the rule: entire, exact-$\Q$ coefficients, dominated by $1/k!$
(entire of exponential type, normalized so $e^{|x|}$ dominates
coefficientwise). Each condition excludes something familiar.
Finite-radius functions --- $\log$, rational powers --- fail
entirety: a variant for them would carry a radius-ratio geometric
tail \emph{and} a domain side condition, and they are certified today
by sibling routes built on their own bounds rather than by stretching
this rule. $\tan$ is meromorphic and out entirely; consumers build it
from $\sin/\cos$ and division. And $e^{x^2}$ is entire but not of
exponential type --- its coefficients escape the domination, so no
choice of $n$ rescues the one uniform tail formula. Conversely,
admitting a new in-class function costs exactly one coefficient
recurrence and one domination case; the tail argument, the check, and
the soundness statement shape are untouched --- the error-function
kernel entered this way, under a $y = z^2$ reduction chosen precisely
to land its coefficients inside the domination.
\end{bremark}

\begin{bremark}[slack is the producer\textquotesingle s business]
The rule certifies one truncation at one radius; tightness is policy,
not soundness. Recentering (the point brackets of the mechanized
\texttt{encCheck\_point}, combined by Taylor-model arithmetic over
$\Q$) and subdivision shrink bracket width first-order, while raising
the degree shrinks it factorially --- so producers buy accuracy with
$n$ and locality with cells, and every combination step re-enters the
calculus as ordinary polynomial rows. On the measured acceptance
systems, degree $9$ at unit radius gave remainders near $10^{-7}$,
orders below the certificate margins consuming them.
\end{bremark}


\blcatsection{Morse Topology}
\rulesection{70}{QMorseDischarge}{core rule}
\label{blcatrule-QMorseDischarge}
\ruleblock{The rule (Steele metanotation)}
\sourcenote{as realized in the mechanized development; source row: \texttt{\texttt{QTSL-rules.toml}} --- \emph{QMorseDischarge}. The metanotation statement is typeset by hand for notational quality; its mechanized ground truth is the Lean statement below, extracted from the sources at every build of this book and hash-bound. The rule \emph{roster} is machine-checked (every mechanized rule tag round-trips to exactly one catalogue row); the statement \emph{text} is editorial, checked by eye against the extracted Lean beside it.}
\begin{mathpar}
\inferrule*{
  L : \mathsf{QBlockLattice} \ \text{(blocks over one rational grid)} \\
  \bigl\{F_f \vdash \sigma_f\, \mathcal{A}\ell_f > 0\bigr\}_{f \in \mathrm{facets}(L)} \ \text{(facet-sign rows: } \ell_f \ \text{the facet\textquotesingle s normal coordinate, } \sigma_f \in \{\pm 1\} \ \text{the declared crossing sign, } F_f \ \text{the facet)} \\
  h : \mathsf{QCellMap} \ (\text{cell} \to \text{vertex} \sqcup \{\text{transit}\}) \\
  \bigl\{\vdash T(b) \subseteq \rho(h(b))\bigr\}_{h(b) \ne \text{transit}} \ \text{(containments, per vertex cell)} \\
  \bigl\{T(b) \vdash \mathcal{A} V_b + \varepsilon_b \le 0\bigr\}_{h(b) = \text{transit}}, \quad \varepsilon_b \in \Q_{>0} \ \text{(transits, per transit cell)} \\
  \bigl\{\vdash \rho(u) \cap \rho(v) = \emptyset\bigr\}_{u \ne v \in G}, \quad \bigl\{\vdash \rho(v) \subseteq D\bigr\}_{v \in G} \ \text{(} \rho\text{-wf: rows where not box-aligned, recomputed where box-aligned)} \\
  \mathrm{CR}(\mathcal{M}|_D) \subseteq \textstyle\bigcup_{h(b) \ne \text{transit}} T(b) \ \text{(cell-level chain-recurrence cover)} \\
  \mathsf{MorseAssembly}(L, h, G, \rho) \ \text{(the assembly ledger, defined below)}
}{
  \mathcal{M} \vDash \mathsf{Morse}(D, G, \rho)
}
\end{mathpar}

\ruleflavor{the lattice payload carries name, grid, arena, and blocks, and enters as certificate data (a shared reference; the rows impose the claims); its per-facet sign rows are displayed inline as the second premise --- each certifies a uniform strict sign of the flow across one grid facet ($\mathcal{A}\ell_f$ is the vector field\textquotesingle s facet-normal component on the facet, frozen at the gridline), and the transit certificates $V_b$ ride the same payload. The $\rho$-wf obligations split by alignment: box-aligned disjointness/containment is bounded interval arithmetic, recomputed; non-box-aligned pairs are per-pair BL rows. The displayed $\mathsf{MorseAssembly}$ package comprises the assembly ledger --- arena compactness, per-cell-pair crossing certificates on shared facets, orbit-interval order rows --- plus profile identification and the ambient metric structure chain recurrence needs (model-class); the chain-recurrence cover is displayed as its own premise. Checker-recomputed: lattice/cell-map presentation conformance (closure recompute, join-irreducible readout, the acyclicity check on the spec graph, the arena inside the top member\textquotesingle s realization --- the nonemptiness row).}
\par\medskip
\rulename{QMorseDischarge} concludes the global Morse judgment: all
recurrence of the model inside the arena lives in the declared regions
$\rho(v)$, and transitions between them run only down the declared
order. It is reduction-to-base in its cleanest form: the certificate is
a lattice of blocks over one rational grid whose every internal facet
carries a strictly-inward sign row of the base calculus, a map $h$
sending the lattice\textquotesingle s cells to spec vertices or to
``transit'', containment rows placing each vertex cell inside its
region, and per-transit-cell decrease rows $\mathcal{A}V + \varepsilon
\le 0$ --- every leaf a sequent the base kernel already checks, and
zero new kernel rules.

A producer reaches for it from the refinement loop: synthesize the
block family (splitting cells where a facet sign fails to certify,
merging where refinement is not needed), pick transit observables from
low-degree templates, and match cells to the spec graph. Failure mode
is refusal with a pointer: a spec whose resolution exceeds what the
certificate can contain reports exactly which cells obstruct, and that
report is itself the useful artefact.

What the checker recomputes rather than trusts: the lattice closure
under union and grid-level intersection; the readout --- the
join-irreducible members, their unique lower covers, the cells
$T(b) = \mathrm{cl}(b \smallsetminus b^{-})$, and the induced cell
order; convexity of $h$\textquotesingle s fibres and monotonicity of
$h$ on vertex cells; acyclicity of the spec graph; and the box-aligned
well-formedness checks on $\rho$. The witness the rule assembles is
$M(v) := \mathrm{Inv}\bigl(\bigcup_{h(b)=v} T(b)\bigr)$ --- the
invariant part of the fibre\textquotesingle s cell-union, not the
union of per-cell invariant parts; the proof section shows the
distinction is load-bearing.

Each hypothesis earns its place. Drop the containment rows and the
location clause fails. Drop the transit rows and a transit cell could
harbour recurrence the $\rho$-cover misses. Drop fibre convexity and
the judgment could assert a strict order the certificate does not
enforce. Drop facet strictness and three arguments die at once:
no-facet-riding, fibre disjointness, and the chain-margin bound
(``strictness pays a second time'').

The mechanization is conditional, honestly: the bound theorems prove
the witness assembly with a ledger of named hypothesis families ---
G3 (compactness of the invariant parts), discharged outright from
model-class facts; G4 (disjointness of distinct fibres), reduced to
transit certificates on fibre overlaps; G5 (the chain-recurrence
cover), reduced to the cell-level cover; and G6 (the order clause),
which remains a hypothesis of the statements bound here, its content
being the block lattice\textquotesingle s orbit combinatorics. The
proof section renders each as mathematics, together with G2, the
pointwise-to-integrated bridge that every certificate-shaped
hypothesis consumes.

\cataloguenote{The judgment is deliberately an ``at most'' claim: empty Morse sets
are permitted, so the rule certifies confinement of recurrence and of
transition order, never existence of structure --- the existence half
(Conley indices, connection witnesses) is a separately priced
annotation tier, absent by design. End-to-end certificate receipts
exist in the development\textquotesingle s exact-arithmetic
certificate-checker lane, which accepts both direct and
product-assembled Morse-family certificates against a mutant battery;
those receipts ride the checker\textquotesingle s own test suite, not
this book\textquotesingle s receipt store.
}
\ruleblock{Theory rows (geolog)}
\sourcenote{The rule\textquotesingle s certificate content as it is declared in the geolog theory the checker loads --- payload schema, side-condition demands, and axiom rows, quoted from the theory source at every build of this catalogue (a stale quotation cannot exist). Axiom modes: \texttt{chase} = the engine derives; \texttt{report} = a violation is a loud judgment; \texttt{refuse} = a violation gates admission. Side-condition relations are discharged by bounded evaluation at the engine\textquotesingle s ground seam, never by derivation.}
\par\medskip
The rule\textquotesingle s official certificate content in the geolog
theory the checker loads. \texttt{QMorse} is the concluded judgment
kind --- the graph-over-regions claim, with its arena --- and
\texttt{QMorseDischarge} is this rule\textquotesingle s payload
entity, field for field the display\textquotesingle s premise kit
(lattice, cell map, containment/transit/well-formedness proof-leaf
lists, and the per-transit-cell certificate polynomials in
\texttt{transit\_certs}). The arena tie is the quoted axiom: a
discharged claim\textquotesingle s arena \emph{is} the
lattice\textquotesingle s, as a stated equation between hash-consed
values. The R1 side condition is the demand relation
\texttt{qtsl\_morse\_readout} --- lattice closure and readout
recompute, discharged by bounded evaluation at the
engine\textquotesingle s ground seam. And the transit rows\textquotesingle
$\mathcal{A}V$ is not an ambient operator: \texttt{qtsl\_gen\_form} is
the derived generator-form seam tying the subject system and the
certificate component $V$ to the hypernet that says how $V$ can change
with duration --- recomputed, never trusted as stored data.

\declhead{theory Qtsl}{QMorse}{\badge{BadgeGreen}{quoted from the theory source at render}}
\begin{leansrc}
QMorse          : Set;   // entity: graph-over-regions claim
\end{leansrc}
\declhead{theory Qtsl}{QMorse/*}{\badge{BadgeGreen}{quoted from the theory source at render}}
\begin{leansrc}
QMorse/judgment : QMorse -> QJudgment;
QMorse/graph : QMorse -> QDag;
QMorse/rho   : QMorse -> QVertexRegionList;
QMorse/arena : QMorse -> QArena;
\end{leansrc}
\declhead{theory Qtsl}{QMorseDischarge}{\badge{BadgeGreen}{quoted from the theory source at render}}
\begin{leansrc}
QMorseDischarge : Set;   // entity: the lattice-reading discharge
\end{leansrc}
\declhead{theory Qtsl}{QMorseDischarge/*}{\badge{BadgeGreen}{quoted from the theory source at render}}
\begin{leansrc}
QMorseDischarge/proof : QMorseDischarge -> QProof;
QMorseDischarge/lattice : QMorseDischarge -> QBlockLattice;
QMorseDischarge/h : QMorseDischarge -> QCellMap;
QMorseDischarge/containments : QMorseDischarge -> QProofList;
QMorseDischarge/transits : QMorseDischarge -> QProofList;
QMorseDischarge/transit_certs : QMorseDischarge -> QCellPolyList;
QMorseDischarge/rho_wf : QMorseDischarge -> QProofList;
\end{leansrc}
\declhead{theory Qtsl}{ax/qmorsedischarge/arena}{\badge{BadgeGreen}{quoted from the theory source at render}}
\begin{leansrc}
chase ax/qmorsedischarge/arena : forall r : QMorseDischarge,
    p : QProof, j : QMorse, lat : QBlockLattice, k : QJudgment,
    ac : Shape, gc : QGridCx.
  r QMorseDischarge/proof = p,
  p QProof/concludes = k,
  j QMorse/judgment = k,
  r QMorseDischarge/lattice = lat,
  lat QBlockLattice/arena = ac,
  lat QBlockLattice/cx = gc
  |- j QMorse/arena = {cells: ac, cx: gc} qarMk;
\end{leansrc}
\declhead{theory Qtsl}{qtsl\_morse\_readout}{\badge{BadgeGreen}{quoted from the theory source at render}}
\begin{leansrc}
qtsl_morse_readout : {dag: QDag, h: QCellMap, lat: QBlockLattice,
                      rho: QVertexRegionList} -> Prop;
\end{leansrc}
\declhead{theory Qtsl}{qtsl\_gen\_form}{\badge{BadgeGreen}{quoted from the theory source at render}}
\begin{leansrc}
qtsl_gen_form : {sys: System, cert: HypernetRef, gen: HypernetRef}
  -> Prop;   // side-condition (oracle discharge / checker
\end{leansrc}
\ruleblock{Lean statement}
\declhead{BernsteinLogic.QtslSoundness}{qMorseDischarge\_sound}{\badge{BadgeGreen}{statement bound}}
\begin{leansrc}
/-- QMorseDischarge soundness — Lemma M(v), transcribed as the witness
assembly `M(v) := MaxInv(⋃_{h(b)=v} T(b))` (the note's own definition,
"the invariant part of the fiber's cell-union, NOT the union of per-cell
invariant parts" — its self-caught first-draft defect). PROVED outright:
the location clause (from the containment rows) and the invariance
clause (`maxInv_isInvariant`). INTERFACE (gaps G3–G6, each the named
prose paragraph of the note's M(v)/M(ii)/M(iii) proofs): compactness,
fiber disjointness, the chain-recurrence cover — into which the transit
rows enter via `transit_maxInv_empty` (M(iv), mechanized above) — and
the order clause. A successor leg mechanizing a gap discharges the
corresponding hypothesis; the statement is already the roster's. -/
theorem qMorseDischarge_sound (φ : Semiflow X) (D : Set X)
    (ρ : Vtx → Set X) (cells : ι → Set X) (hmap : ι → Option Vtx)
    (hContain : ∀ i v, hmap i = some v → cells i ⊆ ρ v)
    (hCompact : ∀ v, IsCompact (φ.MaxInv (fiberUnion cells hmap v)))
    (hDisj : ∀ v w, v ≠ w →
      Disjoint (φ.MaxInv (fiberUnion cells hmap v))
        (φ.MaxInv (fiberUnion cells hmap w)))
    (hCR : ∀ x ∈ D, ChainRecurrentIn φ D x →
      ∃ v, x ∈ φ.MaxInv (fiberUnion cells hmap v))
    (hOrder : ∀ γ : ℝ → X, φ.CompleteOrbit γ → Set.range γ ⊆ D →
      (∀ v, ¬ Set.range γ ⊆ φ.MaxInv (fiberUnion cells hmap v)) →
      ∃ v w, w < v ∧ alphaSet γ ⊆ φ.MaxInv (fiberUnion cells hmap v) ∧
        omegaSet γ ⊆ φ.MaxInv (fiberUnion cells hmap w)) :
    MorseSat φ D ρ
\end{leansrc}
\declhead{BernsteinLogic.QtslSoundness}{MorseWitness}{\badge{BadgeGreen}{statement bound}}
\begin{leansrc}
/-- A witness for `M ⊨_D Morse(G, ρ)` — the graded Morse-decomposition
semantics, clause for clause. `Vtx` carries the QDag's reachability
order (the edge set is presentation; v1 semantics consumes only the
reachability order); empty `M v` are permitted (the deliberate v1
"at most this structure" polarity). Order convention: source-to-sink,
`w < v` = connecting orbits may run from `M v` down to `M w`. -/
structure MorseWitness (φ : Semiflow X) (D : Set X) [PartialOrder Vtx]
    (ρ : Vtx → Set X) where
  /-- the family `{M(v)}` of Morse sets. -/
  M : Vtx → Set X
  /-- each compact… -/
  isCompact : ∀ v, IsCompact (M v)
  /-- …invariant… -/
  invariant : ∀ v, φ.IsInvariantIn (M v)
  /-- clause 1 (location): `M(v) ⊆ ρ(v)`. -/
  located : ∀ v, M v ⊆ ρ v
  /-- …pairwise disjoint. -/
  disjoint : ∀ v w, v ≠ w → Disjoint (M v) (M w)
  /-- clause 2 (completeness of the cover): every chain-recurrent point
  of `φ|_D` lies in some `M(v)`. -/
  crCover : ∀ x ∈ D, ChainRecurrentIn φ D x → ∃ v, x ∈ M v
  /-- clause 3 (order): a complete orbit in `D` contained in no `M(v)`
  connects a strictly higher Morse set to a strictly lower one. -/
  order : ∀ γ : ℝ → X, φ.CompleteOrbit γ → Set.range γ ⊆ D →
    (∀ v, ¬ Set.range γ ⊆ M v) →
    ∃ v w, w < v ∧ alphaSet γ ⊆ M v ∧ omegaSet γ ⊆ M w
\end{leansrc}
\declhead{BernsteinLogic.QtslSoundness}{MorseSat}{\badge{BadgeGreen}{statement bound}}
\begin{leansrc}
/-- The `QMorse` judgment's semantics: some family witnesses it. -/
def MorseSat (φ : Semiflow X) (D : Set X) [PartialOrder Vtx]
    (ρ : Vtx → Set X) : Prop
\end{leansrc}
\declhead{BernsteinLogic.QtslDischarge}{qMorseDischarge\_sound\_modelClass}{\badge{BadgeGreen}{statement bound}}
\begin{leansrc}
/-- `qMorseDischarge_sound` with the gap ledger at its post-discharge
state. Same conclusion, same rows; the interface hypotheses are now:

* model-class facts (the ODE-doctrine profile's warrants, per the
  System notion): spacewise-continuous flow, compact cells —
  G3 is GONE, these replace it;
* `hOverlap` — G4's residue: shared-facet overlaps carry transit
  certificates ("strictness pays a second time");
* `hCell` — G5's residue: the CELL-level chain-recurrence cover
  (M(iii), KMV-III, literature strength);
* `hTransitCert` — M(iv)'s input, now in certificate form;
* `hOrder` — G6, unchanged (wall named in the module header).

The witness assembly and the proved clauses (location, invariance) are
inherited from `qMorseDischarge_sound` untouched. -/
theorem qMorseDischarge_sound_modelClass (φ : Semiflow X) (D : Set X)
    (ρ : Vtx → Set X) (cells : ι → Set X) (hmap : ι → Option Vtx)
    (hflow : ∀ t, Continuous (φ.flow t))
    (hcells : ∀ i, IsCompact (cells i))
    (hContain : ∀ i v, hmap i = some v → cells i ⊆ ρ v)
    (hOverlap : ∀ v w, v ≠ w →
      φ.HasTransitCert
        (fiberUnion cells hmap v ∩ fiberUnion cells hmap w))
    (hCell : ∀ x ∈ D, ChainRecurrentIn φ D x → ∃ i, x ∈ φ.MaxInv (cells i))
    (hTransitCert : ∀ i, hmap i = none → φ.HasTransitCert (cells i))
    (hOrder : ∀ γ : ℝ → X, φ.CompleteOrbit γ → Set.range γ ⊆ D →
      (∀ v, ¬ Set.range γ ⊆ φ.MaxInv (fiberUnion cells hmap v)) →
      ∃ v w, w < v ∧ alphaSet γ ⊆ φ.MaxInv (fiberUnion cells hmap v) ∧
        omegaSet γ ⊆ φ.MaxInv (fiberUnion cells hmap w)) :
    MorseSat φ D ρ
\end{leansrc}
\declhead{BernsteinLogic.QtslDischarge}{maxInv\_isCompact}{\badge{BadgeGreen}{statement bound}}
\begin{leansrc}
/-- G3 discharged: `MaxInv` of a compact region is compact, given a
Hausdorff state space and spacewise-continuous flow maps — the "flow
continuity" the gap ledger names, and nothing else. Route: the set of
complete orbits inside `C` is a closed subset of the compact product
`∀ _ : ℝ, C` (Tychonoff; orbit constraints are equalizers of
continuous maps, closed in T2), and `MaxInv C` is its continuous image
under evaluation at 0. -/
theorem maxInv_isCompact [T2Space X] {φ : Semiflow X} {C : Set X}
    (hflow : ∀ t, Continuous (φ.flow t)) (hC : IsCompact C) :
    IsCompact (φ.MaxInv C)
\end{leansrc}
\declhead{BernsteinLogic.QtslDischarge}{fiberUnion\_isCompact}{\badge{BadgeGreen}{statement bound}}
\begin{leansrc}
/-- Finitely many compact cells have compact fiber unions. -/
theorem fiberUnion_isCompact {Vtx ι : Type*} [Finite ι]
    {cells : ι → Set X} {hmap : ι → Option Vtx}
    (hcells : ∀ i, IsCompact (cells i)) (v : Vtx) :
    IsCompact (fiberUnion cells hmap v)
\end{leansrc}
\declhead{BernsteinLogic.QtslDischarge}{HasTransitCert.disjoint\_maxInv}{\badge{BadgeGreen}{statement bound}}
\begin{leansrc}
/-- The G4 reduction ("strictness pays a second time"): if the overlap of two regions carries a transit
certificate, their maximal invariant sets are disjoint — a common
point would be forward-confined to the overlap by
`maxInv_forward_confined` applied on both sides. No backward
uniqueness, no injectivity of the flow maps. -/
theorem HasTransitCert.disjoint_maxInv {φ : Semiflow X} {C₁ C₂ : Set X}
    (h : φ.HasTransitCert (C₁ ∩ C₂)) :
    Disjoint (φ.MaxInv C₁) (φ.MaxInv C₂)
\end{leansrc}
\declhead{BernsteinLogic.QtslDischarge}{HasTransitCert.maxInv\_empty}{\badge{BadgeGreen}{statement bound}}
\begin{leansrc}
/-- Lemma M(iv) through the certificate interface: a transit-certified
region has empty maximal invariant set. (Shorter than the original
route: forward confinement suffices.) -/
theorem HasTransitCert.maxInv_empty {φ : Semiflow X} {C : Set X}
    (h : φ.HasTransitCert C) : φ.MaxInv C = ∅
\end{leansrc}
\declhead{BernsteinLogic.QtslDischarge}{crCover\_assemble}{\badge{BadgeGreen}{statement bound}}
\begin{leansrc}
/-- The cover-completeness paragraph of Lemma M(v), mechanized: the
fiber-level chain-recurrence cover follows from the CELL-level cover
(M(iii), the KMV-III margin argument — the residue, deliberately at
literature strength) plus transit certificates on transit cells
(M(iv), through `HasTransitCert.maxInv_empty`). -/
theorem crCover_assemble {φ : Semiflow X} {D : Set X}
    {cells : ι → Set X} {hmap : ι → Option Vtx}
    (hCell : ∀ x ∈ D, ChainRecurrentIn φ D x → ∃ i, x ∈ φ.MaxInv (cells i))
    (hTransit : ∀ i, hmap i = none → φ.HasTransitCert (cells i)) :
    ∀ x ∈ D, ChainRecurrentIn φ D x →
      ∃ v, x ∈ φ.MaxInv (fiberUnion cells hmap v)
\end{leansrc}
\declhead{BernsteinLogic.QtslOrderAssembly}{qMorseDischarge\_sound\_orderAssembled}{\badge{BadgeGreen}{statement bound}}
\begin{leansrc}
/-- `qMorseDischarge_sound_modelClass` with G6 at its residue: `hOrder`
is DERIVED (`orderClause_assemble`), and the interface carries instead
arena compactness (`hD`, a model-class fact: the arena is a compact
box-union), the two R1 shadows (`hConv`, `hMono`), and **G6a**
(`hInterval` — M(ii)'s orbit-interval conclusion, prose tier A; see
`orderClause_assemble`). No statement in `QtslSoundness.lean` or the
leg-5 module changes shape. -/
theorem qMorseDischarge_sound_orderAssembled (φ : Semiflow X) (D : Set X)
    (ρ : Vtx → Set X) (cells : ι → Set X) (hmap : ι → Option Vtx)
    (hD : IsCompact D)
    (hflow : ∀ t, Continuous (φ.flow t))
    (hcells : ∀ i, IsCompact (cells i))
    (hContain : ∀ i v, hmap i = some v → cells i ⊆ ρ v)
    (hOverlap : ∀ v w, v ≠ w →
      φ.HasTransitCert
        (fiberUnion cells hmap v ∩ fiberUnion cells hmap w))
    (hCell : ∀ x ∈ D, ChainRecurrentIn φ D x → ∃ i, x ∈ φ.MaxInv (cells i))
    (hTransitCert : ∀ i, hmap i = none → φ.HasTransitCert (cells i))
    (hConv : ∀ i j k v, i ≤ j → j ≤ k →
      hmap i = some v → hmap k = some v → hmap j = some v)
    (hMono : ∀ i j v w, i ≤ j →
      hmap i = some v → hmap j = some w → v ≤ w)
    (hInterval : ∀ γ : ℝ → X, φ.CompleteOrbit γ → Set.range γ ⊆ D →
      ∃ bα bω : ι, bω ≤ bα ∧
        alphaSet γ ⊆ cells bα ∧ omegaSet γ ⊆ cells bω ∧
        ∀ t : ℝ, ∃ i, bω ≤ i ∧ i ≤ bα ∧ γ t ∈ cells i) :
    MorseSat φ D ρ
\end{leansrc}
\declhead{BernsteinLogic.QtslEventualMember}{qMorseDischarge\_sound\_eventualMember}{\badge{BadgeGreen}{statement bound}}
\begin{leansrc}
/-- `qMorseDischarge_sound_orderAssembled` with G6a DERIVED: the
`hInterval` hypothesis is replaced by the block-family interface — the
five R1-recompute combinatorial rows, per-member forward invariance
(G1's class), closedness of join-irreducible realizations, and the
strict-inwardness shadows (F1's meet-confinement clause and the two
F2 absorption facts). No statement in `QtslSoundness.lean` or the
leg-5/leg-6 modules changes shape. What remains of G6 after this leg:
nothing under the name G6a — the residue classes are those already on
the ledger (G1b facet-row bridges per member, G2-class
pointwise-to-integrated for the shadows, R1 recompute conformance). -/
theorem qMorseDischarge_sound_eventualMember (φ : Semiflow X) (D : Set X)
    (ρ : Vtx → Set X) (cells : ι → Set X) (hmap : ι → Option Vtx)
    (blocks : β → Set X) (emb : ι → β) (cover : ι → β) (bD : β)
    (hD : IsCompact D)
    (hflow : ∀ t, Continuous (φ.flow t))
    (hcells : ∀ i, IsCompact (cells i))
    (hContain : ∀ i v, hmap i = some v → cells i ⊆ ρ v)
    (hOverlap : ∀ v w, v ≠ w →
      φ.HasTransitCert
        (fiberUnion cells hmap v ∩ fiberUnion cells hmap w))
    (hCell : ∀ x ∈ D, ChainRecurrentIn φ D x → ∃ i, x ∈ φ.MaxInv (cells i))
    (hTransitCert : ∀ i, hmap i = none → φ.HasTransitCert (cells i))
    (hConv : ∀ i j k v, i ≤ j → j ≤ k →
      hmap i = some v → hmap k = some v → hmap j = some v)
    (hMono : ∀ i j v w, i ≤ j →
      hmap i = some v → hmap j = some w → v ≤ w)
    (hbD : D ⊆ blocks bD)
    (hJoin : ∀ b : β, ∀ x ∈ blocks b,
      ∃ i : ι, x ∈ blocks (emb i) ∧ blocks (emb i) ⊆ blocks b)
    (hcover : ∀ i : ι, blocks (cover i) ⊂ blocks (emb i))
    (hcellsSup : ∀ i : ι, blocks (emb i) \ blocks (cover i) ⊆ cells i)
    (horder : ∀ i j : ι, blocks (emb i) ⊆ blocks (emb j) → i ≤ j)
    (hclosed : ∀ i : ι, IsClosed (blocks (emb i)))
    (hInv : ∀ b : β, φ.Invariant (blocks b))
    (hMeet : ∀ a b : β, ∃ c : β, blocks c ⊆ blocks a ∩ blocks b ∧
      ∀ γ : ℝ → X, φ.CompleteOrbit γ → Set.range γ ⊆ D →
        ∀ s t : ℝ, (∀ u, s ≤ u → u ≤ t → γ u ∈ blocks a ∩ blocks b) →
          ∀ u, s ≤ u → u ≤ t → γ u ∈ blocks c)
    (hAbsorbOmega : ∀ γ : ℝ → X, φ.CompleteOrbit γ → Set.range γ ⊆ D →
      ∀ b : β, (omegaSet γ ∩ blocks b).Nonempty →
        ∃ T : ℝ, ∀ t, T ≤ t → γ t ∈ blocks b)
    (hAbsorbAlpha : ∀ γ : ℝ → X, φ.CompleteOrbit γ → Set.range γ ⊆ D →
      ∀ b : β, (alphaSet γ ∩ blocks b).Nonempty →
        ∃ T : ℝ, ∀ t, t ≤ T → γ t ∈ blocks b) :
    MorseSat φ D ρ
\end{leansrc}
\declhead{BernsteinLogic.QtslCrossing}{qMorseDischarge\_sound\_cellCover}{\badge{BadgeGreen}{statement bound}}
\begin{leansrc}
/-- `qMorseDischarge_sound_modelClass` with G4 at its sharpened residue:
the fiber-overlap hypothesis `hOverlap` (one transit certificate on each
whole fiber-intersection — not the row shape) is replaced by per-cell-
pair crossing certificates `hCross` (the row shape: one certified facet
functional per shared facet). New model-class fact consumed: trajectory
continuity `htraj` — the same hypothesis G1's
`invariant_of_boundaryDwell` already takes. Everything else is
inherited unchanged; G5's and G6's residues are exactly as in
`qMorseDischarge_sound_modelClass`. -/
theorem qMorseDischarge_sound_cellCover (φ : Semiflow X) (D : Set X)
    (ρ : Vtx → Set X) (cells : ι → Set X) (hmap : ι → Option Vtx)
    (hflow : ∀ t, Continuous (φ.flow t))
    (htraj : ∀ x, Continuous fun t : ℝ≥0 => φ.flow t x)
    (hcells : ∀ i, IsCompact (cells i))
    (hContain : ∀ i v, hmap i = some v → cells i ⊆ ρ v)
    (hCross : ∀ v w, v ≠ w → ∀ i j, hmap i = some v → hmap j = some w →
      φ.HasCrossingCert (cells i ∩ cells j))
    (hCell : ∀ x ∈ D, ChainRecurrentIn φ D x → ∃ i, x ∈ φ.MaxInv (cells i))
    (hTransitCert : ∀ i, hmap i = none → φ.HasTransitCert (cells i))
    (hOrder : ∀ γ : ℝ → X, φ.CompleteOrbit γ → Set.range γ ⊆ D →
      (∀ v, ¬ Set.range γ ⊆ φ.MaxInv (fiberUnion cells hmap v)) →
      ∃ v w, w < v ∧ alphaSet γ ⊆ φ.MaxInv (fiberUnion cells hmap v) ∧
        omegaSet γ ⊆ φ.MaxInv (fiberUnion cells hmap w)) :
    MorseSat φ D ρ
\end{leansrc}
\ruleblock{Correspondence (Steele display $\leftrightarrow$ Lean)}
\badge{BadgeAmber}{a bound Lean statement CHANGED since certification --- re-review and re-certify}
\sourcenote{The dictionary below is a reviewer-certified statement of how the mechanized statements above interpret the displayed rule: display metavariable $\mapsto$ Lean binder, each display hypothesis $\mapsto$ the premise or side condition carrying it, the display conclusion $\mapsto$ the goal --- with any deliberate difference in width between display and mechanization named outright. Certification is recorded as a hash triple over (the displayed cell, the extracted statement texts, this dictionary\textquotesingle s own text); if any of the three changes, this block turns amber until a reviewer re-certifies. The correspondence claim itself remains editorial --- reviewer-grade, not machine-checked; its machine-checked successor types each soundness statement against the rule as a structured object, and supersedes this block per rule as it lands.}
The displayed judgment $\mathcal{M} \vDash \mathsf{Morse}(D, G, \rho)$
$\mapsto$ the goal \texttt{MorseSat φ D ρ} = \texttt{Nonempty
(MorseWitness φ D ρ)} --- the Morse note \S2.2 clauses (compactness,
invariance, location, disjointness, chain-recurrence cover, order)
packaged as a structure, so the certified conclusion is existential
exactly as the judgment reads; the witness assembly $M(v) :=
\mathrm{MaxInv}\bigl(\bigcup_{h(b)=v} T(b)\bigr)$ lives inside the
proof term of \texttt{qMorseDischarge\_sound}, never in its statement.
The display\textquotesingle s spec graph $G$ enters as the binder pair
\texttt{Vtx} + \texttt{[PartialOrder Vtx]}: the semantics consumes
only the reachability order (the \texttt{MorseWitness} docstring says
so in words), and the dag check on the presentation\textquotesingle s
edge set is R1 recompute with no Lean counterpart --- correctly.
Order polarity carried exactly: \texttt{hOrder}\textquotesingle s
conclusion places $\alpha$-limits in the strictly \emph{higher} set
and $\omega$-limits in the lower (\texttt{w < v ∧ alphaSet γ ⊆ \ldots
v ∧ omegaSet γ ⊆ \ldots w}), the display\textquotesingle s
``transitions run only down the declared order.''
\begin{itemize}
\item the E2 lattice $L$ and cell map $h$ $\mapsto$ the abstract
  binders \texttt{cells : ι → Set X} and \texttt{hmap : ι → Option
  Vtx} (\texttt{none} = transit, the display\textquotesingle s
  $\sqcup\,\{\text{transit}\}$): the whole lattice stratum ---
  closure recompute, join-irreducible readout, $T(b)$, fibre
  convexity and monotonicity of $h$, box-aligned $\rho$-wf --- is
  presentation the checker recomputes (side condition R1), and the
  bound statements never see it; they quantify over arbitrary cell
  families, strictly wider than the display, in the safe direction.
\item containment rows $\mapsto$ \texttt{hContain} (\texttt{hmap i =
  some v → cells i ⊆ ρ v}), one-to-one at the semantic stratum.
\item transit rows $\mathcal{A}V + \varepsilon \le 0$ $\mapsto$
  \texttt{hTransitCert i} (per \texttt{hmap i = none}) in the
  semantic-shadow form \texttt{HasTransitCert}: some observable $V$,
  bounded below on the cell, with explicit margin $\varepsilon > 0$
  and the \emph{integrated} decrease clause $V(\Phi_t x) \le V(x) -
  \varepsilon t$ along confined prefixes. The pointwise-to-integrated
  bridge (gap G2; the \texttt{FlowDerivOn} profile identification) is
  outside this bound set --- the roster witness cites its
  mechanization (\texttt{QtslMeanValue.lean}), not bound here.
\item $\rho$-wf rows $\mapsto$ nothing: no bound hypothesis consumes
  them (fibre disjointness arrives through \texttt{hOverlap}
  instead) --- R1 territory, like the lattice.
\item the conditional ledger, exactly as bound:
  \texttt{qMorseDischarge\_sound} carries G3--G6 as hypotheses
  (\texttt{hCompact}/\texttt{hDisj}/\texttt{hCR}/\texttt{hOrder});
  \texttt{\_modelClass} discharges G3 outright
  (\texttt{maxInv\_isCompact}: Tychonoff, consuming \texttt{[T2Space
  X]}, spacewise flow continuity \texttt{hflow}, \texttt{[Finite ι]}
  + compact cells via \texttt{fiberUnion\_isCompact} --- model-class
  facts all), reduces G4 to overlap transit certificates
  (\texttt{HasTransitCert.disjoint\_maxInv}: forward confinement on
  both sides, no backward uniqueness), reduces G5 to the cell-level
  cover \texttt{hCell} (M(iii), KMV-III, literature strength) through
  \texttt{crCover\_assemble} + \texttt{HasTransitCert.maxInv\_empty}
  (M(iv)), and carries G6 (\texttt{hOrder}) untouched.
\end{itemize}
Width note, both directions stated. Three assembly statements extend
the bound set (the three assembly blocks below): the
order assembly (\texttt{qMorseDischarge\_sound\_orderAssembled},
module \texttt{QtslOrderAssembly}) derives \texttt{hOrder} --- G6
proper --- carrying instead arena compactness, the two R1 shadows, and
G6a, the orbit-interval row; the eventual-member assembly
(\texttt{qMorseDischarge\_sound\_eventualMember}, module
\texttt{QtslEventualMember}) derives G6a in turn from the block-family
interface (the five R1-recompute combinatorial rows, per-member
forward invariance, closedness, strict-inwardness shadows), leaving
nothing of G6 under its own name; and the crossing assembly
(\texttt{qMorseDischarge\_sound\_cellCover}, module
\texttt{QtslCrossing}) sharpens G4 to the row shape --- per-cell-pair
crossing certificates on shared facets in place of whole-fibre-overlap
transit certificates. All three are hash-bound, so the certified
reach extends past \texttt{\_modelClass}\textquotesingle s ledger and
the conditional residue shrinks to the ledger\textquotesingle s
G1b/G2-class profile identifications (the
$\mathrm{FlowDerivOn}\,V\,(\mathcal{A}V)$ chain-rule bridges --- the
mean-value and Dini profile strata remain cited theorems, not bound),
the cell-level chain-recurrence cover --- now a displayed premise
($\mathrm{CR}(\mathcal{M}|_D) \subseteq \bigcup T(b)$ over vertex
cells) --- $\mapsto$ \texttt{hCell} (the fibre-to-cell margin argument
is the proof\textquotesingle s stated residue), and the displayed
$\mathsf{MorseAssembly}$ package $\mapsto$ the remaining
crossing/order/interface hypotheses plus recompute
conformance. Binder cargo: \texttt{[PseudoMetricSpace
X]} is consumed by chain recurrence itself --- $(\varepsilon,T)$-chains
need distances --- ambient structure the assembly package carries.

\ruleblock{Lean verifier receipts}
\sourcenote{module \texttt{BernsteinLogic.QtslSoundness}; bound build \texttt{77005991d61d\ldots} (bound from receipt 8130a7a7a679cef4, 2026-08-04T01:16:06Z). A receipt records that named checking layers accepted a build: (module, compiled-artifact SHA-256, environment SHA-256, checker layers), minted when the checking gate ran; the store receipts attesting the bound build are listed with the layers each run credited. Nothing here asks the reader to trust the store: every layer is an independently runnable program over the quoted sources --- rebuild the module and re-run the checkers to reproduce the verdict. Receipts attesting superseded builds credit no layer and are omitted from this rendering (the count of omissions, if any, is noted below).}
\begin{receipttable}
\receiptrow{8130a7a7a679cef4}{2026-08-04T01:16:06Z}{build+decorrelated}{\badge{BadgeGreen}{current}}
\receiptrow{9eb3706354082144}{2026-08-04T02:23:12Z}{build+decorrelated}{\badge{BadgeGreen}{current}}
\end{receipttable}
\sourcenote{5 further receipt(s) in the store attest superseded builds of this module and are omitted from this rendering.}
\sourcenote{module \texttt{BernsteinLogic.QtslDischarge}; bound build \texttt{4aa46910a600\ldots} (bound from receipt dd2946805f790f99, 2026-08-04T01:16:09Z). A receipt records that named checking layers accepted a build: (module, compiled-artifact SHA-256, environment SHA-256, checker layers), minted when the checking gate ran; the store receipts attesting the bound build are listed with the layers each run credited. Nothing here asks the reader to trust the store: every layer is an independently runnable program over the quoted sources --- rebuild the module and re-run the checkers to reproduce the verdict. Receipts attesting superseded builds credit no layer and are omitted from this rendering (the count of omissions, if any, is noted below).}
\begin{receipttable}
\receiptrow{dd2946805f790f99}{2026-08-04T01:16:09Z}{build+decorrelated}{\badge{BadgeGreen}{current}}
\receiptrow{e1bc68890e2aebf1}{2026-08-04T02:23:11Z}{build+decorrelated}{\badge{BadgeGreen}{current}}
\end{receipttable}
\sourcenote{14 further receipt(s) in the store attest superseded builds of this module and are omitted from this rendering.}
\sourcenote{module \texttt{BernsteinLogic.QtslOrderAssembly}; bound build \texttt{1fef28f78168\ldots} (bound from receipt f6e3628b18a7ed89, 2026-07-20T08:27:03Z). A receipt records that named checking layers accepted a build: (module, compiled-artifact SHA-256, environment SHA-256, checker layers), minted when the checking gate ran; the store receipts attesting the bound build are listed with the layers each run credited. Nothing here asks the reader to trust the store: every layer is an independently runnable program over the quoted sources --- rebuild the module and re-run the checkers to reproduce the verdict. Receipts attesting superseded builds credit no layer and are omitted from this rendering (the count of omissions, if any, is noted below).}
\begin{receipttable}
\receiptrow{b066c7d1be7303f1}{2026-07-19T04:29:30Z}{build+decorrelated}{\badge{BadgeGreen}{current}}
\receiptrow{ef9ccb70d0449d3b}{2026-07-19T09:07:23Z}{build+decorrelated}{\badge{BadgeGreen}{current}}
\receiptrow{f6e3628b18a7ed89}{2026-07-20T08:27:03Z}{build+decorrelated}{\badge{BadgeGreen}{current}}
\receiptrow{08fb682d6ccdcbc7}{2026-07-20T20:44:50Z}{build+decorrelated}{\badge{BadgeGreen}{current}}
\receiptrow{62d115665ebe5578}{2026-07-20T21:01:19Z}{build+replay+hardened+decorrelated}{\badge{BadgeGreen}{current}}
\receiptrow{b2fd112afe983043}{2026-07-21T05:28:46Z}{build+replay+hardened+decorrelated}{\badge{BadgeGreen}{current}}
\end{receipttable}
\sourcenote{module \texttt{BernsteinLogic.QtslEventualMember}; bound build \texttt{9282e7682307\ldots} (bound from receipt 839bcd47d53458bf, 2026-07-20T08:27:05Z). A receipt records that named checking layers accepted a build: (module, compiled-artifact SHA-256, environment SHA-256, checker layers), minted when the checking gate ran; the store receipts attesting the bound build are listed with the layers each run credited. Nothing here asks the reader to trust the store: every layer is an independently runnable program over the quoted sources --- rebuild the module and re-run the checkers to reproduce the verdict. Receipts attesting superseded builds credit no layer and are omitted from this rendering (the count of omissions, if any, is noted below).}
\begin{receipttable}
\receiptrow{d9e8d15537b8c931}{2026-07-19T11:09:44Z}{build+decorrelated}{\badge{BadgeGreen}{current}}
\receiptrow{dcd494c897af133b}{2026-07-19T11:29:34Z}{build+decorrelated}{\badge{BadgeGreen}{current}}
\receiptrow{839bcd47d53458bf}{2026-07-20T08:27:05Z}{build+decorrelated}{\badge{BadgeGreen}{current}}
\receiptrow{0b348bff26bc835c}{2026-07-20T20:44:43Z}{build+decorrelated}{\badge{BadgeGreen}{current}}
\receiptrow{fbf234c5b555fbd6}{2026-07-20T21:01:12Z}{build+replay+hardened+decorrelated}{\badge{BadgeGreen}{current}}
\receiptrow{a67d63ba90d91fb5}{2026-07-21T05:28:37Z}{build+replay+hardened+decorrelated}{\badge{BadgeGreen}{current}}
\end{receipttable}
\sourcenote{2 further receipt(s) in the store attest superseded builds of this module and are omitted from this rendering.}
\sourcenote{module \texttt{BernsteinLogic.QtslCrossing}; bound build \texttt{29a1be816112\ldots} (bound from receipt 3dec8934e00e8465, 2026-07-20T08:27:06Z). A receipt records that named checking layers accepted a build: (module, compiled-artifact SHA-256, environment SHA-256, checker layers), minted when the checking gate ran; the store receipts attesting the bound build are listed with the layers each run credited. Nothing here asks the reader to trust the store: every layer is an independently runnable program over the quoted sources --- rebuild the module and re-run the checkers to reproduce the verdict. Receipts attesting superseded builds credit no layer and are omitted from this rendering (the count of omissions, if any, is noted below).}
\begin{receipttable}
\receiptrow{58226ade60f5d15c}{2026-07-19T04:33:17Z}{build+decorrelated}{\badge{BadgeGreen}{current}}
\receiptrow{8fa273e016e7de3b}{2026-07-19T09:07:20Z}{build+decorrelated}{\badge{BadgeGreen}{current}}
\receiptrow{3dec8934e00e8465}{2026-07-20T08:27:06Z}{build+decorrelated}{\badge{BadgeGreen}{current}}
\receiptrow{204c60d7cf1cf318}{2026-07-20T20:44:39Z}{build+decorrelated}{\badge{BadgeGreen}{current}}
\receiptrow{574f67e1c43258ae}{2026-07-20T21:01:08Z}{build+replay+hardened+decorrelated}{\badge{BadgeGreen}{current}}
\receiptrow{5e4d769e0abef0a1}{2026-07-21T05:28:33Z}{build+replay+hardened+decorrelated}{\badge{BadgeGreen}{current}}
\end{receipttable}
\layerline{\layeritem{build}{lake build (elaboration)}{\badge{BadgeGreen}{current}}\layeritem{replay}{leanchecker (official kernel replay)}{\badge{BadgeGreen}{current}}\layeritem{hardened}{safe\_verify (hardened screen)}{\badge{BadgeGreen}{current}}\layeritem{decorrelated}{lean4export + nanoda (decorrelated kernel)}{\badge{BadgeGreen}{current}}}
\sourcenote{\badge{BadgeGray}{comparator: none} no fixture in the three-implementation comparator harness (compiled / kernel-interpreted / reference-interpreted) exercises this rule.}
\ruleblock{HOL4 statement}
\absentbox{NO HOL STATEMENT YET}{Absent, and honestly so at two levels. The HOL4
development\textquotesingle s object is the base-kernel checker cone
--- the first-order computational surface of the sequent-stratum
certificate contract (dtype factorization, exact rational arithmetic,
the catalogue check functions) --- and the satisfaction-judgment
stratum this rule concludes in has no HOL twin: no Morse-family or
judgment-stratum definition exists in the ported scripts (verified
against the port\textquotesingle s definition inventory). The base
leaves this rule consumes --- facet-sign, containment, and decrease
rows --- ARE covered by the ported checker; what has no HOL
counterpart is the judgment-stratum soundness development itself,
which today is Lean-only. A port would be new assurance surface over
the semiflow stratum, not a translation of existing checker code.
}
\ruleblock{HOL verifier receipts}
\absentbox{NONE}{no HOL verification receipt exists for this rule.}
\ruleblock{Proof}

\sourcenote{Genre: core rule (tier 1) --- an ordinary mathematical proof
at Putnam-exam-solution detail; the largest mechanized lemma set in this
book. The rule's mechanization is \emph{conditional}: the soundness
theorem is proved with a ledger of four named hypothesis families, of
which one is discharged outright below, two are reduced to
transit-certificate form, and the fourth --- the order clause --- is
carried as an explicit hypothesis by the discharged-floor theorem and
discharged by the separately bound assembly legs (ledger remark). The
ledger remark near the end states each as mathematics. Every hypothesis is used and remarked on; the formal-anchor
table maps each step to the Lean declaration that carries it.}

\textbf{Setting.}\enspace Fix a forward semiflow $\varphi$ on a metric
state space $X$ ($\Phi_0 = \mathrm{id}$,
$\Phi_{s+t} = \Phi_t \circ \Phi_s$), and a compact forward-invariant
arena $D \subseteq X$. Four standard notions, exactly as the mechanized
statements carry them:

\emph{(complete orbits)} $\gamma : \R \to X$ is a complete orbit when
$\gamma(s + t) = \Phi_t(\gamma(s))$ for all $s \in \R$, $t \geq 0$ ---
a two-sided trajectory compatible with the forward flow.

\emph{(maximal invariant part)} for $C \subseteq X$,
\[
\mathrm{Inv}(C) \;:=\; \{\, x \ :\ \exists\, \gamma \ \text{complete
orbit},\ \gamma(0) = x,\ \mathrm{range}\,\gamma \subseteq C \,\}.
\]
An \emph{invariant set} $S$ (the sense fit for a semiflow) is one that
is forward-invariant and satisfies $S \subseteq \mathrm{Inv}(S)$ ---
every point rides a complete orbit inside.

\emph{(limit sets)} $\omega(\gamma) := \bigcap_{T}
\mathrm{cl}\,\gamma([T, \infty))$ and $\alpha(\gamma) := \bigcap_{T}
\mathrm{cl}\,\gamma((-\infty, T])$.

\emph{(chains)} an $(\varepsilon, T)$-chain from $x$ to $y$ inside $D$
is a finite concatenation of orbit segments of duration $\geq T$ with
jumps of size $< \varepsilon$ at points of $D$ (generated by: a single
step $x, y \in D$, $t \geq T$,
$\mathrm{dist}(\Phi_t x, y) < \varepsilon$; and transitivity). A point
is \emph{chain-recurrent in $D$} when for every $\varepsilon > 0$ and
$T$ there is an $(\varepsilon, T)$-chain from $x$ to itself.

\textbf{The judgment.}\enspace
$\mathcal{M} \vDash \mathsf{Morse}(D, G, \rho)$ holds when some family
$\{M(v)\}_{v \in V(G)}$ --- indexed by the poset the graph $G$
presents by reachability --- satisfies:
each $M(v)$ compact and invariant; the family pairwise disjoint;
\emph{(location)} $M(v) \subseteq \rho(v)$;
\emph{(completeness of the cover)} every chain-recurrent point of
$\varphi|_D$ lies in some $M(v)$; and
\emph{(order)} every complete orbit in $D$ contained in no single
$M(v)$ has $\alpha(\gamma) \subseteq M(v)$ and
$\omega(\gamma) \subseteq M(w)$ for some $w < v$. The order runs
source-to-sink: $w < v$ means connecting orbits may run from $M(v)$
down to $M(w)$, so the $\alpha$-side vertex is the higher one. Empty
$M(v)$ are permitted --- the judgment is an ``at most this structure''
claim (recurrence confined to the $\rho(v)$, transitions confined to
the order), never an existence claim; the existence half is a
separately priced annotation tier (Conley indices, connection
matrices), deliberately absent here.

\textbf{The rule's data and the checker's recompute.}\enspace The
certificate supplies: the family object --- a finite lattice of
box-union blocks over one declared rational grid, each internal facet
carrying a strictly-inward sign row (the \rulename{QBlockInv} section's
row species; the same rows serve both rules); a cell map $h$ from the
lattice's cells to $V(G) \sqcup \{\mathsf{transit}\}$; per vertex cell
a containment row $T(b) \subseteq \rho(h(b))$; per transit cell a
polynomial $V$ with decrease rows $\mathcal{A}V + \varepsilon \leq 0$;
and well-formedness rows for $\rho$. The checker \emph{recomputes},
rather than trusts, the lattice readout: closure under union and
grid-level intersection, the join-irreducible members $b$ with unique
lower covers $b^-$, the cells $T(b) := \mathrm{cl}(b \smallsetminus
b^-)$, the induced cell order, convexity of $h$'s fibres and
monotonicity of $h$ on vertex cells, and acyclicity of $G$. Writing
$\mathrm{fib}(v) := \bigcup_{h(b) = v} T(b)$ for the fibre's
cell-union, the witness the rule assembles is
\[
M(v) \;:=\; \mathrm{Inv}\big(\mathrm{fib}(v)\big)
\]
--- the invariant part \emph{of the fibre's cell-union}, not the union
of per-cell invariant parts. The distinction is load-bearing
(\autoref{rem:qmd-mv}).

\begin{blemma}[$\mathrm{Inv}$ is located and invariant]
\label{lem:qmd-inv}
\emph{(i)} $\mathrm{Inv}(C) \subseteq C$.
\emph{(ii)} $\mathrm{Inv}(C)$ is an invariant set.
\end{blemma}
\begin{proof}
(i) The witnessing orbit passes through $x$ at time $0$ with range in
$C$. (ii) Two facts, both by time-shifting: shifting a complete orbit
by any $s$ ($u \mapsto \gamma(u + s)$) is again a complete orbit, and
every point $\gamma(s)$ of an orbit with range in $C$ therefore lies in
$\mathrm{Inv}(C)$ with the shifted orbit as witness. Forward
invariance: for $x \in \mathrm{Inv}(C)$ with orbit $\gamma$,
$\Phi_t x = \gamma(0 + t)$ is such a point. The complete-orbit-through
clause: the same statement at every shift shows
$\mathrm{range}\,\gamma \subseteq \mathrm{Inv}(C)$, so the witnessing
orbit itself runs inside $\mathrm{Inv}(C)$.
\end{proof}

\begin{blemma}[location]\label{lem:qmd-loc}
Under the containment rows, $\mathrm{fib}(v) \subseteq \rho(v)$, hence
$M(v) \subseteq \rho(v)$ by Lemma~\ref{lem:qmd-inv}(i).
\end{blemma}
\begin{proof}
A union of cells each inside $\rho(v)$ is inside $\rho(v)$.
\end{proof}

\begin{blemma}[forward confinement]\label{lem:qmd-conf}
If $x \in \mathrm{Inv}(C)$ then $\Phi_t x \in C$ for every
$t \geq 0$.
\end{blemma}
\begin{proof}
The witnessing orbit's forward half \emph{is} the forward orbit:
$\Phi_t x = \Phi_t(\gamma(0)) = \gamma(t) \in C$. This one-line lemma
is the section's quiet workhorse: a semiflow is deterministic
\emph{forward}, so membership in an $\mathrm{Inv}$ set constrains the
actual forward trajectory --- and nothing here assumes backward
uniqueness (a semiflow may branch backward; complete orbits through a
point need not be unique). The disjointness reduction below survives
precisely because it only ever uses forward halves.
\end{proof}

\begin{blemma}[transit certificates kill invariance]\label{lem:qmd-transit}
Call $C$ \emph{transit-certified} when there are $V : X \to \R$ and
$\varepsilon > 0$ with $V$ bounded below on $C$ and, along every
forward trajectory confined to $C$,
$V(\Phi_t x) \leq V(x) - \varepsilon\, t$. Then no point is
forward-confined to $C$; in particular
$\mathrm{Inv}(C) = \emptyset$.
\end{blemma}
\begin{proof}
Suppose $\Phi_t x \in C$ for all $t$, and let $m$ bound $V$ below on
$C$. By the Archimedean property choose $n$ with
$n\,\varepsilon > V(x) - m$; the decrease clause at horizon $n$ gives
$V(\Phi_n x) \leq V(x) - \varepsilon n < m$, contradicting
$\Phi_n x \in C$. The $\mathrm{Inv}$ statement follows by
Lemma~\ref{lem:qmd-conf}. (The certified rows are pointwise ---
$\mathcal{A}V + \varepsilon \leq 0$ on the cell, for the deterministic
generator $\mathcal{A} = f \cdot \nabla$ --- and reach the integrated
clause used here by integrating the derivative along the trajectory;
that bridge is ledger item G2 in \autoref{rem:qmd-ledger}, stated
once and consumed by every certificate-shaped hypothesis in this
section.)
\end{proof}

\begin{blemma}[disjointness from overlap certificates]
\label{lem:qmd-disj}
If $C_1 \cap C_2$ is transit-certified, then
$\mathrm{Inv}(C_1)$ and $\mathrm{Inv}(C_2)$ are disjoint.
\end{blemma}
\begin{proof}
A common point is forward-confined to $C_1$ and to $C_2$
(Lemma~\ref{lem:qmd-conf} twice), hence to $C_1 \cap C_2$ --- which
Lemma~\ref{lem:qmd-transit} forbids. No backward uniqueness, no
injectivity of the flow maps.
\end{proof}

\begin{bremark}[why the certificate sits on the overlap ---
``strictness pays a second time'']\label{rem:qmd-overlap}
Distinct fibres' cell-unions intersect only in shared boundary faces
(cell interiors are disjoint by the readout's construction ---
generically facets, but possibly lower-dimensional: two cells may meet
in an edge or a corner alone). On a certified facet incident to a
shared face, the facet polynomial supplies the observable ---
\emph{oriented to decrease}: the rows certify inward growth of $\ell$,
so the transit observable Lemma~\ref{lem:qmd-transit} consumes is
$V = -\ell$, bounded below on the compact overlap, with the certified
strict rate as its $\varepsilon$; the shared face lies in $\ell$'s
zero set, and confinement to the face would force $\ell \equiv 0$.
Where the overlap is a union of several faces, one signed observable
per face does not by itself assemble into a single global certificate
--- the row-shaped assembly is per-cell-pair crossing certificates on
shared facets, the development's G4 sharpening, bound in this edition
(anchor table) --- and the hypothesis
Theorem~\ref{thm:qmd-sound} takes is the assembled overlap
certificate. The strict margin of the \rulename{QBlockInv} rows is
thus consumed a second time, as the transit raw material for facet
regions. One might
hope to weaken the hypothesis to per-piece certificates --- each
$C_i$ separately transit-certified --- and conclude something about
the union; that is false. On the circle rotation, two overlapping
closed arcs each carry a monotone angular observable (a transit
certificate for the arc), yet the union is the whole circle, every
point recurrent: certificates on pieces do not compose across an
overlapping union, because no global observable exists. The overlap
form is the correct, and certifiable, hypothesis.
\end{bremark}

\begin{blemma}[compactness of the invariant parts]\label{lem:qmd-cpt}
Suppose $X$ is Hausdorff, each flow map $\Phi_t : X \to X$ is
continuous, and $C$ is compact. Then $\mathrm{Inv}(C)$ is compact.
Consequently each $\mathrm{fib}(v)$, a finite union of compact cells,
has compact invariant part.
\end{blemma}
\begin{proof}
Let $K \subseteq X^{\R}$ (product topology) be the set of complete
orbits with range in $C$:
\[
K \;=\; \Big(\bigcap_{s \in \R,\ t \geq 0}
\{\gamma : \gamma(s+t) = \Phi_t(\gamma(s))\}\Big)
\;\cap\; \prod_{r \in \R} C .
\]
Each orbit constraint is the equalizer of two continuous maps ---
evaluation at $s + t$, and $\Phi_t \circ (\text{evaluation at } s)$
--- hence closed, $X$ being Hausdorff; the product
$\prod_{\R} C$ is compact by Tychonoff (the one genuinely infinitary
step of this section); a closed subset of a compact set is compact. And
$\mathrm{Inv}(C)$ is exactly the image of $K$ under evaluation at $0$,
a continuous map, hence compact. The fibre statement: a finite union
of compact cells is compact, and the above applies to it.
\end{proof}

\begin{blemma}[the cover assembles from the cell level]
\label{lem:qmd-cover}
Suppose \emph{(cell-level cover)} every chain-recurrent point of
$\varphi|_D$ lies in $\mathrm{Inv}(T(b))$ for some cell $b$, and
\emph{(transit rows)} every transit cell is transit-certified. Then
every chain-recurrent point of $\varphi|_D$ lies in some
$M(v) = \mathrm{Inv}(\mathrm{fib}(v))$.
\end{blemma}
\begin{proof}
Let $x$ be chain-recurrent, $x \in \mathrm{Inv}(T(b))$. If $h(b) =
\mathsf{transit}$, Lemma~\ref{lem:qmd-transit} makes
$\mathrm{Inv}(T(b))$ empty --- contradiction; so $h(b)$ is a vertex
$v$. Then $T(b) \subseteq \mathrm{fib}(v)$, and $\mathrm{Inv}$ is
monotone in its region (a complete orbit inside the cell is inside the
fibre's union), so $x \in \mathrm{Inv}(\mathrm{fib}(v))$.
\end{proof}

\begin{bthm}[soundness of \rulename{QMorseDischarge}, at the
discharged floor]\label{thm:qmd-sound}
Suppose:
\emph{(model class)} $X$ is a Hausdorff metric space, each flow map
$\Phi_t$ is continuous, and every cell is compact;
\emph{(containment rows)} $T(b) \subseteq \rho(h(b))$ for every vertex
cell;
\emph{(overlap certificates)} for $v \neq w$, the overlap
$\mathrm{fib}(v) \cap \mathrm{fib}(w)$ is transit-certified;
\emph{(cell-level cover)} every chain-recurrent point of $\varphi|_D$
lies in some cell's invariant part;
\emph{(transit rows)} every transit cell is transit-certified; and
\emph{(order clause)} every complete orbit in $D$ contained in no
single $M(v)$ connects a strictly higher fibre's invariant part to a
strictly lower one, in the sense of the judgment's order clause with
$M(v) = \mathrm{Inv}(\mathrm{fib}(v))$.
Then $\mathcal{M} \vDash \mathsf{Morse}(D, G, \rho)$.
\end{bthm}
\begin{proof}
Take the family $M(v) := \mathrm{Inv}(\mathrm{fib}(v))$ and check the
witness clauses. \emph{Location}: Lemma~\ref{lem:qmd-loc}.
\emph{Invariance}: Lemma~\ref{lem:qmd-inv}(ii). \emph{Compactness}:
Lemma~\ref{lem:qmd-cpt} under the model-class facts.
\emph{Disjointness}: Lemma~\ref{lem:qmd-disj} applied to the overlap
certificates. \emph{Cover}: Lemma~\ref{lem:qmd-cover} from the
cell-level cover and the transit rows. \emph{Order}: the hypothesis,
verbatim. All six clauses hold, so the family witnesses the judgment.
\end{proof}

\begin{bremark}[the gap ledger, as mathematics]\label{rem:qmd-ledger}
The rule was first mechanized as the bare witness assembly: location
and invariance proved outright, and four clauses --- labelled G3--G6 in
the development, the labels this book adopts as the rule's hypothesis
names --- taken as opaque interface hypotheses. The discharge stratum
bound by this section moves the ledger to the state
Theorem~\ref{thm:qmd-sound} exhibits:

\emph{G3 (compactness of the invariant parts): proven.}
Lemma~\ref{lem:qmd-cpt} derives it from model-class facts alone ---
Hausdorff state space, spacewise-continuous flow, compact cells. No
compactness hypothesis survives at the fibre level.

\emph{G4 (disjointness of distinct fibres' invariant parts): reduced,
and sharpened to rows.} Lemma~\ref{lem:qmd-disj} trades it for transit
certificates on pairwise overlaps; the residue is the overlap
observable of \autoref{rem:qmd-overlap} --- shared-facet
functionals whose certified strict signs are the raw material ---
together with the G2 bridge below; and the row-shaped assembly
(per-cell-pair crossing certificates on shared facets in place of
whole-overlap certificates) is mechanized and bound
(\texttt{qMorseDischarge\_sound\_cellCover}, anchor table).

\emph{G5 (completeness of the chain-recurrence cover): reduced.}
Lemma~\ref{lem:qmd-cover} trades the fibre-level cover for the
cell-level one. The residue is the margin argument: a chain-recurrent
point's chains stay within its own cell, because the strict facet
margin is a uniform bound on admissible chain jumps near facets --- an
$\varepsilon$-chain with $\varepsilon$ below the margin can no more
exit a block than an orbit can. The classical home of this
containment is Conley's theory \cite{conley1978}; its algorithmic
form is the chain-recurrence algorithm of
Kalies--Mischaikow--van der Vorst \cite{kmv2005}.

\emph{G6 (the order clause): a hypothesis of
Theorem~\ref{thm:qmd-sound}, discharged by the bound assembly legs.}
Its mathematical content is the orbit
combinatorics of the block lattice: the \emph{eventual member} of a
forward orbit --- the minimal lattice member containing a forward
tail, well-defined because the tail-containing family is finite,
closed under the lattice meet, and nonempty: the arena sits inside the
top member's realization (one of the combinatorial interface rows the
assembly legs consume), so the top member contains every forward
tail --- descends the cell
order and stabilizes, giving every complete orbit in $D$ extreme
cells $b_\omega \leq b_\alpha$ that bracket every visited cell; the
fibre-convexity case analysis then converts bracketing into the
witness's order clause (the extremes are not transit, by
Lemma~\ref{lem:qmd-transit} against their nonempty invariant limit
sets; a same-fibre pair traps the whole orbit in its fibre by
convexity, contradicting the ``no single $M(v)$'' hypothesis; distinct
fibres yield the strict order by monotonicity of $h$). The
development's assembly legs mechanize exactly this --- the cell order
entering as theorem-interface data, the eventual-member construction
proved against combinatorial row hypotheses --- and are bound in this
edition (\texttt{qMorseDischarge\_sound\_orderAssembled},
\texttt{\_eventualMember}; anchor table), so nothing of G6 survives
under its own name: the residue is the R1-recompute combinatorial rows
and strict-inwardness shadows those legs consume, plus the G2-class
profile identifications below. Theorem~\ref{thm:qmd-sound} is stated
at the discharged floor and takes the clause as an explicit hypothesis
there; the rule's conditional badge is judged against the full bound
set.

\emph{G2 (the pointwise-to-integrated bridge), named for
completeness:} every ``transit-certified'' hypothesis above is the
integrated semantic shadow of pointwise rows
$\mathcal{A}V + \varepsilon \leq 0$; the bridge is one-variable
calculus along trajectories, and its surviving analytic content is
the identification of the certified polynomial rate with the
trajectory-wise derivative on the region --- the same model-class
profile warrant as ledger item (a) of the \rulename{QBlockInv}
section, whose G1 (block invariance) completes the shared ledger.
\end{bremark}

\begin{bremark}[the witness is $\mathrm{Inv}$ of the union, not the
union of $\mathrm{Inv}$s]\label{rem:qmd-mv}
Define instead $M'(v) := \bigcup_{h(b) = v} \mathrm{Inv}(T(b))$ and
the order clause fails: a complete orbit crossing between two cells of
the \emph{same} fibre lies in no per-cell invariant part (its range is
in neither cell alone), so it belongs to no $M'(v)$ while also
connecting nothing strictly ordered --- the ``orbit contained in no
Morse set'' case triggers with no higher/lower pair to name. The
correct definition absorbs such orbits into their fibre's $M(v)$,
which is exactly the role of fibre convexity: every cell between the
orbit's extremes maps to the same vertex, so the whole orbit lies in
the fibre's union and hence in its invariant part. Without convexity
the judgment could assert a strict order the certificate does not
enforce --- convexity of $h$'s fibres is well-formedness, not a
nicety, and the checker recomputes it.
\end{bremark}

\begin{bremark}[relation to the literature]
The judgment is Conley's Morse decomposition \cite{conley1978} in
certified form: a decomposition indexed by a two-element order is an
attractor--repeller pair, and the general form is the lattice-refined
iteration of that idea; the completeness clause is the containment of
chain recurrence in the Morse sets, and the certified staircase of
lattice members is a finite shard of a complete Lyapunov function
(Conley's fundamental theorem). The block-lattice route --- attracting
blocks closed under union and intersection, join-irreducibles, the
tessellation into cells $T(b)$, and Birkhoff duality
\cite{birkhoff1937} between the attractor lattice and the Morse-graph
poset --- is the program of Kalies--Mischaikow--van der Vorst's lattice
series \cite{kmv2016,kmv2022}; the chain-recurrence cover's
algorithmic form is their earlier chain-recurrence algorithm
\cite{kmv2005}; and the computational tradition of certifying global
dynamics by index machinery is surveyed in Mischaikow--Mrozek
\cite{mischaikowmrozek2002}. What the rule adds to that tradition is
the artefact discipline: those pipelines end in trusted computation,
while here every analytic claim rides certified sign rows --- base
kernel sequents in exact $\Q$ --- plus combinatorics the checker
recomputes, with the semantic warrant being the lemma chain above.
\end{bremark}

\begin{anchors}
witness assembly, location + invariance proved & \texttt{qMorseDischarge\_sound} (QtslSoundness.lean); \texttt{maxInv\_subset}, \texttt{maxInv\_isInvariant}\\
Lemma \ref{lem:qmd-inv} & \texttt{maxInv\_subset}, \texttt{maxInv\_invariant}, \texttt{maxInv\_isInvariant}\\
Lemma \ref{lem:qmd-loc} & \texttt{fiberUnion\_subset}; rows \texttt{hContain}\\
Lemma \ref{lem:qmd-conf} & \texttt{maxInv\_forward\_confined} (QtslDischarge.lean)\\
Lemma \ref{lem:qmd-transit} & \texttt{HasTransitCert}, \texttt{HasTransitCert.not\_forward\_confined}, \texttt{HasTransitCert.maxInv\_empty}; the bare form \texttt{transit\_maxInv\_empty} (QtslSoundness.lean)\\
Lemma \ref{lem:qmd-disj} & \texttt{HasTransitCert.disjoint\_maxInv}; hypothesis \texttt{hOverlap}\\
Lemma \ref{lem:qmd-cpt} & \texttt{maxInv\_isCompact} (Tychonoff route), \texttt{fiberUnion\_isCompact}\\
Lemma \ref{lem:qmd-cover} & \texttt{crCover\_assemble}; \texttt{maxInv\_mono}, \texttt{subset\_fiberUnion}; hypotheses \texttt{hCell}, \texttt{hTransitCert}\\
Theorem \ref{thm:qmd-sound} & \texttt{qMorseDischarge\_sound\_modelClass}; order hypothesis \texttt{hOrder} (ledger G6)\\
judgment semantics & \texttt{MorseWitness}, \texttt{MorseSat}, \texttt{EpsChainIn}, \texttt{ChainRecurrentIn}, \texttt{omegaSet}, \texttt{alphaSet}\\
G6's assembly legs (\textbf{bound in this edition}) & \texttt{orderClause\_assemble}, \texttt{eventualMember\_interval}, \texttt{qMorseDischarge\_sound\_orderAssembled}, \texttt{\_eventualMember} (QtslOrderAssembly, QtslEventualMember)\\
G4's row sharpening (\textbf{bound in this edition}) & \texttt{qMorseDischarge\_sound\_cellCover} (QtslCrossing)\\
\end{anchors}


\rulesection{71}{QMorseProduct}{core rule}
\label{blcatrule-QMorseProduct}
\ruleblock{The rule (Steele metanotation)}
\sourcenote{as realized in the mechanized development; source row: \texttt{\texttt{QTSL-rules.toml}} --- \emph{QMorseProduct}. The metanotation statement is typeset by hand for notational quality; its mechanized ground truth is the Lean statement below, extracted from the sources at every build of this book and hash-bound. The rule \emph{roster} is machine-checked (every mechanized rule tag round-trips to exactly one catalogue row); the statement \emph{text} is editorial, checked by eye against the extracted Lean beside it.}
\begin{mathpar}
\inferrule*{
  \mathcal{M}_1 \vDash \mathsf{Morse}(D_1, G_1, \rho_1) \\
  \mathcal{M}_2 \vDash \mathsf{Morse}(D_2, G_2, \rho_2) \\
  L_1, L_2 : \mathsf{QBlockLattice} \ \text{(the factor lattices)}
}{
  \mathcal{M}_1 \times \mathcal{M}_2 \vDash \mathsf{Morse}(D_1 \times D_2,\ G_1 \times G_2,\ \rho_1 \times \rho_2)
}
\end{mathpar}

\ruleflavor{the two factor judgments enter as derivation premises (judgment-kind edges); the factor lattices enter as certificate data (shared references --- the premise judgments\textquotesingle{} own certificates impose their claims), consulted by the checker\textquotesingle s product tie.}
\par\medskip
\rulename{QMorseProduct} composes two global judgments along the
state-space axis: disjoint variable sets, product arena, product spec
graph, product regions. A producer reaches for it exactly where global
structure is engineered compositionally --- certify each factor
subsystem\textquotesingle s Morse graph, then take the product ---
because the alternative (re-discharging the composite through
\rulename{QMorseDischarge} on the composite lattice) recomputes an
exponentially larger readout and re-derives the cell map.

The rule is fenced to the deterministic-flow class, and the fence is
the rule\textquotesingle s own mathematical content: the composite
certificate\textquotesingle s readout --- join-irreducibles, cells,
cell order --- must be the product of the factor readouts, and for
flows that is a theorem (the attractor lattice of a product flow is
the coproduct of the factor attractor lattices), while for maps it is
false: products of maps exhibit emergent, phase-locked attractors
outside the coproduct, so a composite readout computed from factor
certificates would misrepresent the product\textquotesingle s global
structure. For maps the rule must not be derived at all.

What the checker recomputes: wire-disjointness of the factors, the
product surgery on graphs and regions, and --- on the certificate
route --- the composite lattice\textquotesingle s readout. Facet and
transit rows ride along by citation: every composite facet is normal
to one factor\textquotesingle s axis and its inwardness polynomial is
that factor\textquotesingle s, the other factor\textquotesingle s
coordinates entering as free box variables; a factor transit
observable still decreases on cell $\times$ anything, the other
factor\textquotesingle s drift differentiating it to zero.

The witness-level lemma chain is mechanized in full: the product
construction on the semiflow stratum (\texttt{Semiflow.prod}), the
factoring lemmas for orbits, limit sets, and chains, the six-clause
witness transport (\texttt{MorseWitness.prod}), and the soundness
theorem (\texttt{qMorseProduct\_sound}) --- unconditional, with no
interface hypotheses: from the two factor judgments, the product
judgment, on the product arena, product order, and product regions.
The bound statements consume neither compactness nor invariance of
the arenas (the proof never needed them), and Hausdorffness exactly
where the transport says ``compact, hence closed.''

\cataloguenote{The proof section separates two strata deliberately: the
witness-transport argument (compactness, invariance, location,
disjointness, cover, order --- all coordinate-wise) is elementary and
does not consume continuity of time, while the readout-transport
theorem that warrants the certificate route is flow-specific. The
fence binds the rule as catalogued; the decomposition is recorded so
the mechanization and any map-side successor judgment know where the
flow hypothesis actually bites. The mechanization respects the split
exactly: the bound statements carry the witness-level entailment in
full and import nothing of Kalies--Wehbe --- the fence's warrant
remains at literature strength, where it belongs, since it certifies
the canonicity of the certificate-route composite readout rather than
the satisfaction-level soundness bound here.
}
\ruleblock{Lean statement}
\declhead{BernsteinLogic.QtslProduct}{qMorseProduct\_sound}{\badge{BadgeGreen}{statement bound}}
\begin{leansrc}
/-- QMorseProduct soundness, witness level (thm:qmp-sound; the bind
target). From the two factor judgments — the rule's E1 premises `left,
right : QMorse` — conclude the product judgment on the product arena,
product reachability order, and product regions. UNCONDITIONAL: no
interface hypotheses (contrast the `\dcond` family); the analytic
content is the witness transport above. The E2 lattice premises and the
`flows_only` fence's readout-transport warrant (Kalies–Wehbe) live at
the certificate stratum and are deliberately not consumed — see the
module header and the book page's fence remark. -/
theorem qMorseProduct_sound
    (h₁ : MorseSat φ₁ D₁ ρ₁) (h₂ : MorseSat φ₂ D₂ ρ₂) :
    MorseSat (φ₁.prod φ₂) (D₁ ×ˢ D₂)
      (fun p : Vtx₁ × Vtx₂ => ρ₁ p.1 ×ˢ ρ₂ p.2)
\end{leansrc}
\declhead{BernsteinLogic.QtslProduct}{MorseWitness.prod}{\badge{BadgeGreen}{statement bound}}
\begin{leansrc}
/-- Witness transport (lem:qmp-transport): the product family
`M(v, w) := M₁(v) × M₂(w)`, indexed by the product order, witnesses the
product judgment. Six clauses, each the book paragraph:

* compactness — a product of compacts is compact;
* invariance — coordinate-wise forward invariance, and a product point
  rides the pairing of the factor witnesses' complete orbits;
* location — the product of the factor locations;
* disjointness — distinct product vertices differ in some coordinate,
  and that factor's disjointness kills the intersection;
* cover — a chain-recurrent product point has chain-recurrent
  coordinates (§4), each covered by its factor family;
* order — a product orbit contained in no `M(v, w)` has an uncontained
  coordinate; three cases (both uncontained / first only / second
  only), the contained coordinate handled by compact-hence-closed and
  the closure bound (§3), the α-side vertex strictly higher in the
  product order exactly as the source-to-sink convention demands. -/
def MorseWitness.prod (W₁ : MorseWitness φ₁ D₁ ρ₁)
    (W₂ : MorseWitness φ₂ D₂ ρ₂) :
    MorseWitness (φ₁.prod φ₂) (D₁ ×ˢ D₂)
      (fun p : Vtx₁ × Vtx₂ => ρ₁ p.1 ×ˢ ρ₂ p.2)
\end{leansrc}
\declhead{BernsteinLogic.QtslProduct}{completeOrbit\_prod\_iff}{\badge{BadgeGreen}{statement bound}}
\begin{leansrc}
/-- Orbits factor (lem:qmp-factor(i)): `γ` is a complete orbit of the
product iff each coordinate is a complete orbit of its factor. -/
theorem completeOrbit_prod_iff {γ : ℝ → X₁ × X₂} :
    (φ₁.prod φ₂).CompleteOrbit γ ↔
      φ₁.CompleteOrbit (fun s => (γ s).1) ∧
      φ₂.CompleteOrbit (fun s => (γ s).2)
\end{leansrc}
\declhead{BernsteinLogic.QtslProduct}{omegaSet\_prod\_subset}{\badge{BadgeGreen}{statement bound}}
\begin{leansrc}
/-- ω-limit sets factor (lem:qmp-factor(ii)):
`ω(γ) ⊆ ω(γ₁) × ω(γ₂)`. Book proof verbatim: the image of `[T, ∞)`
under `γ` sits inside the product of the coordinate images; closure of
a product is the product of closures. -/
theorem omegaSet_prod_subset (γ : ℝ → X₁ × X₂) :
    omegaSet γ ⊆
      (omegaSet fun s => (γ s).1) ×ˢ (omegaSet fun s => (γ s).2)
\end{leansrc}
\declhead{BernsteinLogic.QtslProduct}{alphaSet\_prod\_subset}{\badge{BadgeGreen}{statement bound}}
\begin{leansrc}
/-- α-limit sets factor (lem:qmp-factor(ii), the `likewise for α`). -/
theorem alphaSet_prod_subset (γ : ℝ → X₁ × X₂) :
    alphaSet γ ⊆
      (alphaSet fun s => (γ s).1) ×ˢ (alphaSet fun s => (γ s).2)
\end{leansrc}
\declhead{BernsteinLogic.QtslProduct}{EpsChainIn.fst}{\badge{BadgeGreen}{statement bound}}
\begin{leansrc}
/-- Chains project, first coordinate (lem:qmp-factor(iii)): under the
max metric (mathlib's product metric, `Prod.dist_eq`) factor distance
is ≤ product distance, so a chain step projects to a chain step — jump
points project into `D₁`, durations unchanged — and concatenations
project to concatenations. -/
theorem EpsChainIn.fst {D₁ : Set X₁} {D₂ : Set X₂} {ε : ℝ} {T : ℝ≥0}
    {x y : X₁ × X₂}
    (h : EpsChainIn (φ₁.prod φ₂) (D₁ ×ˢ D₂) ε T x y) :
    EpsChainIn φ₁ D₁ ε T x.1 y.1
\end{leansrc}
\declhead{BernsteinLogic.QtslProduct}{EpsChainIn.snd}{\badge{BadgeGreen}{statement bound}}
\begin{leansrc}
/-- Chains project, second coordinate. -/
theorem EpsChainIn.snd {D₁ : Set X₁} {D₂ : Set X₂} {ε : ℝ} {T : ℝ≥0}
    {x y : X₁ × X₂}
    (h : EpsChainIn (φ₁.prod φ₂) (D₁ ×ˢ D₂) ε T x y) :
    EpsChainIn φ₂ D₂ ε T x.2 y.2
\end{leansrc}
\declhead{BernsteinLogic.QtslProduct}{ChainRecurrentIn.fst}{\badge{BadgeGreen}{statement bound}}
\begin{leansrc}
/-- Chain recurrence projects, first coordinate (lem:qmp-factor(iii),
"hence a point chain-recurrent in `D₁ × D₂` has both coordinates
chain-recurrent in their arenas"). -/
theorem ChainRecurrentIn.fst {D₁ : Set X₁} {D₂ : Set X₂} {x : X₁ × X₂}
    (h : ChainRecurrentIn (φ₁.prod φ₂) (D₁ ×ˢ D₂) x) :
    ChainRecurrentIn φ₁ D₁ x.1
\end{leansrc}
\declhead{BernsteinLogic.QtslProduct}{ChainRecurrentIn.snd}{\badge{BadgeGreen}{statement bound}}
\begin{leansrc}
/-- Chain recurrence projects, second coordinate. -/
theorem ChainRecurrentIn.snd {D₁ : Set X₁} {D₂ : Set X₂} {x : X₁ × X₂}
    (h : ChainRecurrentIn (φ₁.prod φ₂) (D₁ ×ˢ D₂) x) :
    ChainRecurrentIn φ₂ D₂ x.2
\end{leansrc}
\declhead{BernsteinLogic.QtslProduct}{CompleteOrbit.prodMk}{\badge{BadgeGreen}{statement bound}}
\begin{leansrc}
/-- Pairing complete orbits (lem:qmp-factor(i), ⇐): factor orbits ride
together as a product orbit. -/
theorem CompleteOrbit.prodMk {γ₁ : ℝ → X₁} {γ₂ : ℝ → X₂}
    (h₁ : φ₁.CompleteOrbit γ₁) (h₂ : φ₂.CompleteOrbit γ₂) :
    (φ₁.prod φ₂).CompleteOrbit fun s => (γ₁ s, γ₂ s)
\end{leansrc}
\declhead{BernsteinLogic.QtslProduct}{omegaSet\_subset\_closure\_range}{\badge{BadgeGreen}{statement bound}}
\begin{leansrc}
/-- `ω(γ) ⊆ cl(range γ)` — one intersect and one image bound; feeds the
order clause's contained cases (book:
`α(γ₂), ω(γ₂) ⊆ cl(range γ₂) ⊆ M₂(w₀)`). -/
theorem omegaSet_subset_closure_range (γ : ℝ → X) :
    omegaSet γ ⊆ closure (Set.range γ)
\end{leansrc}
\declhead{BernsteinLogic.QtslProduct}{alphaSet\_subset\_closure\_range}{\badge{BadgeGreen}{statement bound}}
\begin{leansrc}
/-- `α(γ) ⊆ cl(range γ)` — the α twin. -/
theorem alphaSet_subset_closure_range (γ : ℝ → X) :
    alphaSet γ ⊆ closure (Set.range γ)
\end{leansrc}
\ruleblock{Correspondence (Steele display $\leftrightarrow$ Lean)}
\absentbox{NO CORRESPONDENCE DICTIONARY YET}{no reviewer-certified dictionary maps the displayed rule onto the mechanized statements above; until one is written and certified, that mapping rests on this section\textquotesingle s editorial care alone.}
\ruleblock{Lean verifier receipts}
\sourcenote{module \texttt{BernsteinLogic.QtslProduct}; bound build \texttt{1b85f883dfdc\ldots} (bound from receipt d5b482c9d6ed9b6c, 2026-07-20T03:30:32Z). A receipt records that named checking layers accepted a build: (module, compiled-artifact SHA-256, environment SHA-256, checker layers), minted when the checking gate ran; the store receipts attesting the bound build are listed with the layers each run credited. Nothing here asks the reader to trust the store: every layer is an independently runnable program over the quoted sources --- rebuild the module and re-run the checkers to reproduce the verdict. Receipts attesting superseded builds credit no layer and are omitted from this rendering (the count of omissions, if any, is noted below).}
\begin{receipttable}
\receiptrow{d5b482c9d6ed9b6c}{2026-07-20T03:30:32Z}{build+decorrelated}{\badge{BadgeGreen}{current}}
\receiptrow{15fa0566d44e84fd}{2026-07-20T20:44:52Z}{build+decorrelated}{\badge{BadgeGreen}{current}}
\receiptrow{e8214fce26d91980}{2026-07-20T21:01:21Z}{build+replay+hardened+decorrelated}{\badge{BadgeGreen}{current}}
\receiptrow{b87a989a1475042c}{2026-07-21T05:28:48Z}{build+replay+hardened+decorrelated}{\badge{BadgeGreen}{current}}
\end{receipttable}
\layerline{\layeritem{build}{lake build (elaboration)}{\badge{BadgeGreen}{current}}\layeritem{replay}{leanchecker (official kernel replay)}{\badge{BadgeGreen}{current}}\layeritem{hardened}{safe\_verify (hardened screen)}{\badge{BadgeGreen}{current}}\layeritem{decorrelated}{lean4export + nanoda (decorrelated kernel)}{\badge{BadgeGreen}{current}}}
\sourcenote{\badge{BadgeGray}{comparator: none} no fixture in the three-implementation comparator harness (compiled / kernel-interpreted / reference-interpreted) exercises this rule.}
\ruleblock{HOL4 statement}
\absentbox{NO HOL STATEMENT YET}{Absent at the judgment stratum: the HOL4 development\textquotesingle s
object is the base-kernel checker cone, and no Morse-family definition
exists in the ported scripts (verified against the
port\textquotesingle s definition inventory) --- so the
satisfaction-judgment semantics this rule\textquotesingle s soundness
is stated in has no HOL twin to mirror the now-complete Lean
mechanization into. A future port has its interface fixed twice over:
the witness clauses quoted in the \rulename{QMorseDischarge} section
and the product lemma chain bound here.
}
\ruleblock{HOL verifier receipts}
\absentbox{NONE}{no HOL verification receipt exists for this rule.}
\ruleblock{Proof}

\sourcenote{Genre: core rule (tier 1) --- an ordinary mathematical proof
at Putnam-exam-solution detail, mechanized \emph{in full}: the
accompanying development proves the product construction, the factoring
lemmas, the witness transport, and the soundness theorem with no
interface hypotheses (contrast the \rulename{QMorseDischarge} ledger;
within the Morse family only \rulename{QBasin} shares this fully-proved
status). The lemma chain below fixed the mechanization's interface in
advance, and the formal-anchor table maps each step to the declaration
that carries it. The readout-transport theorem of the fence remark is
deliberately \emph{not} mechanized: it warrants the canonicity of the
certificate-route composite readout, not the witness-level entailment
proved here, and remains at literature strength.}

\textbf{Setting and claim.}\enspace Let $\varphi_1, \varphi_2$ be
forward semiflows on Hausdorff metric state spaces $X_1, X_2$, with
compact forward-invariant arenas $D_1, D_2$. The \emph{product
semiflow} $\varphi_1 \times \varphi_2$ on $X_1 \times X_2$ (with the
max metric) flows coordinate-wise:
$\Phi_t(x_1, x_2) = (\Phi^1_t x_1, \Phi^2_t x_2)$; the semiflow laws
hold coordinate-wise. The rule: from
$\mathcal{M}_1 \vDash \mathsf{Morse}(D_1, G_1, \rho_1)$ and
$\mathcal{M}_2 \vDash \mathsf{Morse}(D_2, G_2, \rho_2)$, with
disjoint variable sets, conclude
\[
\mathcal{M}_1 \times \mathcal{M}_2 \ \vDash\
\mathsf{Morse}\big(D_1 \times D_2,\ G_1 \times G_2,\ \rho_1 \times \rho_2\big),
\]
where $G_1 \times G_2$ presents the product order
($(v_1, v_2) \geq (w_1, w_2)$ iff $v_i \geq w_i$ in both coordinates)
and $(\rho_1 \times \rho_2)(v_1, v_2) := \rho_1(v_1) \times
\rho_2(v_2)$.

\begin{blemma}[orbits, limit sets, and chains factor]
\label{lem:qmp-factor}
\emph{(i)} $\gamma = (\gamma_1, \gamma_2)$ is a complete orbit of
$\varphi_1 \times \varphi_2$ iff each $\gamma_i$ is a complete orbit
of $\varphi_i$.
\emph{(ii)} $\omega(\gamma) \subseteq \omega(\gamma_1) \times
\omega(\gamma_2)$, and likewise for $\alpha$.
\emph{(iii)} the projection of an $(\varepsilon, T)$-chain inside
$D_1 \times D_2$ is an $(\varepsilon, T)$-chain inside $D_i$; hence a
point chain-recurrent in $D_1 \times D_2$ has both coordinates
chain-recurrent in their arenas.
\end{blemma}
\begin{proof}
(i) is the coordinate-wise flow law read in both directions. (ii):
$\gamma([T, \infty)) \subseteq \gamma_1([T,\infty)) \times
\gamma_2([T,\infty))$; closure of a product is the product of
closures; and an intersection of product sets is the product of the
intersections --- pointwise, $(x_1, x_2)$ lies in every
$A_T \times B_T$ iff each $x_i$ lies in every factor set. (iii): under
the max metric, factor distance is $\leq$ product distance, so a
single chain step projects to a chain step (jump points project into
$D_i$, durations are unchanged), and concatenations project to
concatenations.
\end{proof}

\begin{blemma}[witness transport]\label{lem:qmp-transport}
Let $\{M_1(v)\}_{v \in V(G_1)}$ and $\{M_2(w)\}_{w \in V(G_2)}$
witness the factor judgments. Then the product family
\[
M(v, w) \;:=\; M_1(v) \times M_2(w),
\]
indexed by the product order, witnesses the product judgment.
\end{blemma}
\begin{proof}
Six clauses.

\emph{Compactness}: a product of two compact sets is compact.

\emph{Invariance}: forward invariance holds coordinate-wise; and a
point $(x_1, x_2)$ of the product rides the pairing of the factor
witnesses' complete orbits through $x_1$ and $x_2$
(Lemma~\ref{lem:qmp-factor}(i)), which stays inside the product set.

\emph{Location}: $M_1(v) \times M_2(w) \subseteq \rho_1(v) \times
\rho_2(w)$, the product of the factor locations.

\emph{Disjointness}: if $(v, w) \neq (v', w')$ they differ in some
coordinate, say $v \neq v'$; then $M_1(v) \cap M_1(v') = \emptyset$
makes the products disjoint regardless of the other coordinate.

\emph{Cover}: let $(x_1, x_2)$ be chain-recurrent in $D_1 \times D_2$.
By Lemma~\ref{lem:qmp-factor}(iii) each $x_i$ is chain-recurrent in
$D_i$, hence covered by its factor family: $x_1 \in M_1(v)$,
$x_2 \in M_2(w)$ for some $v, w$; so
$(x_1, x_2) \in M(v, w)$.

\emph{Order}: let $\gamma = (\gamma_1, \gamma_2)$ be a complete orbit
in $D_1 \times D_2$ contained in no single $M(v, w)$. Containment in a
product set is coordinate-wise, so
--- negating ``some $(v,w)$ contains $\gamma$'' ---
either $\gamma_1$ is contained in no $M_1(v)$, or $\gamma_2$ in no
$M_2(w)$. Three cases.
\emph{(both uncontained)}: the factor order clauses give
$v > v'$ with $\alpha(\gamma_1) \subseteq M_1(v)$,
$\omega(\gamma_1) \subseteq M_1(v')$, and $w > w'$ likewise for
$\gamma_2$; by Lemma~\ref{lem:qmp-factor}(ii),
$\alpha(\gamma) \subseteq M(v, w)$ and
$\omega(\gamma) \subseteq M(v', w')$, with
$(v, w) > (v', w')$ strictly in the product order.
\emph{($\gamma_1$ uncontained, $\gamma_2 \subseteq M_2(w_0)$)}: the
first factor gives $v > v'$ as above; for the second,
$\mathrm{range}\,\gamma_2 \subseteq M_2(w_0)$ with $M_2(w_0)$ compact
in a Hausdorff space, hence closed, gives
$\alpha(\gamma_2), \omega(\gamma_2) \subseteq
\mathrm{cl}(\mathrm{range}\,\gamma_2) \subseteq M_2(w_0)$; so
$\alpha(\gamma) \subseteq M(v, w_0)$,
$\omega(\gamma) \subseteq M(v', w_0)$, and
$(v, w_0) > (v', w_0)$ strictly.
\emph{(the symmetric case)}: exchange the roles of the factors.
In each case the $\alpha$-side vertex is the strictly higher one, as
the source-to-sink convention demands.
\end{proof}

\begin{bthm}[soundness of \rulename{QMorseProduct}, witness level]
\label{thm:qmp-sound}
If both factor judgments hold, the product judgment holds.
\end{bthm}
\begin{proof}
Choose factor witnesses and apply Lemma~\ref{lem:qmp-transport}.
\end{proof}

\begin{bremark}[where the flow hypothesis actually binds --- the rule's
fence]\label{rem:qmp-fence}
Nothing in Lemmas~\ref{lem:qmp-factor}--\ref{lem:qmp-transport}
consumes continuity of time: read over $\mathbb{N}$-time systems, the
same six clauses transport verbatim, at the same ``at most'' polarity.
The rule is nonetheless fenced to the deterministic-\emph{flow} class,
for reasons that live above the witness stratum. The certificate-level
discharge route does not consume semantic witnesses: it composes the
factor \emph{lattices} (the join-closure of the products of factor
blocks, every composite facet's inwardness polynomial being a factor's
with the other factor's coordinates riding along as free box
variables) and must know that the composite lattice's recomputed
readout --- its join-irreducibles, cells, and cell order --- is the
product of the factor readouts, so that the composite tessellation
represents the product dynamics at the resolution the factors
certified. For flows this is a theorem: the attractor lattice of a
product flow is the coproduct of the factor attractor lattices
(Kalies--Wehbe \cite{kalieswehbe2026}), whose Birkhoff dual
\cite{birkhoff1937} is exactly ``composite Morse graph $=$ product of
factor graphs.'' For maps it is false: products of maps exhibit
emergent, phase-locked attractors outside the coproduct
(Kalies--Wehbe's counterexample --- already visible in the product of
two period-two orbits, whose diagonal and anti-diagonal are attractors
of the product generated by no product of factor attractors), so a
composite readout computed from factor certificates would silently
misrepresent the product's global structure, and the judgment's
ecosystem --- which reads a certified Morse graph as the canonical
decomposition at the declared resolution --- would be poisoned at
exactly the canonicity it relies on. Hence the fence: for maps the
rule must not be derived at all; for flows the readout-transport
theorem is the rule's semantic import, and it is the reason this core
rule carries its own lemma rather than reducing to lattice
combinatorics.
\end{bremark}

\begin{bremark}[what rides along for free]\label{rem:qmp-rides}
Two certificate-level transports cost nothing beyond citation, and are
worth recording with the rule even though the witness-level proof
above does not touch them. \emph{Facet rows}: every internal facet of
a composite block is normal to one factor's axis, and its inwardness
polynomial \emph{is} that factor's facet polynomial --- the other
factor's coordinates enter as free box variables, so the composite
row's enclosure data is the factor's, replicated along the free
coordinates. \emph{Transit rows}: if $\mathcal{A}_1 V \leq
-\varepsilon$ on a factor cell, the same $V$ read on the product
satisfies $\mathcal{A}_{1 \times 2} V = \mathcal{A}_1 V \leq
-\varepsilon$ on cell $\times$ anything --- the other factor's drift
differentiates $V$ to zero. Both ride the same disjointness of the variable sets --- the facet
rows by weakening alone, the transit rows by weakening plus the
product-generator identity
$\mathcal{A}_{1\times 2}(V \circ \pi_i) = (\mathcal{A}_i V) \circ \pi_i$,
which is the precise content of ``differentiates $V$ to zero.''
\end{bremark}

\begin{bremark}[the mechanization consumes no arena hypotheses]
\label{rem:qmp-arena}
The setting above posits compact forward-invariant arenas because that
is what arenas are in this catalogue's ecosystem; the transport itself
never uses either property, and the mechanized statements quantify over
bare sets $D_1, D_2$ --- strictly more general, and faithful to the
witness clauses, which constrain the arena only through the judgments.
Hausdorffness is consumed exactly once, where the proof of
Lemma~\ref{lem:qmp-transport} says ``compact in a Hausdorff space,
hence closed''; the mechanization takes it as a topological hypothesis
over a pseudometric base, which every metric state space satisfies.
\end{bremark}

\begin{bremark}[relation to the literature]
The rule is the certified face of compositionality results for global
dynamics: Kalies--Wehbe \cite{kalieswehbe2026} prove the
attractor-lattice coproduct isomorphism for product flows and delimit
it (maps fail; skew products degrade to one-way results), inside the
lattice-theoretic program of Kalies--Mischaikow--van der Vorst
\cite{kmv2016,kmv2022}. The witness-level transport
(Lemma~\ref{lem:qmp-transport}) is elementary Conley theory
\cite{conley1978} over products; the observation that its ``at most''
clauses survive even where the canonical readout does not
(\autoref{rem:qmp-fence}) is exactly why the judgment's fence is
stated at the rule, not at the witness.
\end{bremark}

\begin{anchors}
product semiflow (Setting) & \texttt{Semiflow.prod}, \texttt{prod\_flow} (QtslProduct.lean); the max metric is the product metric, \texttt{Prod.dist\_eq}\\
Lemma \ref{lem:qmp-factor}(i) & \texttt{completeOrbit\_prod\_iff}, \texttt{CompleteOrbit.prodMk}\\
Lemma \ref{lem:qmp-factor}(ii) & \texttt{omegaSet\_prod\_subset}, \texttt{alphaSet\_prod\_subset}\\
Lemma \ref{lem:qmp-factor}(iii) & \texttt{EpsChainIn.fst}/\texttt{.snd}, \texttt{ChainRecurrentIn.fst}/\texttt{.snd}\\
Lemma \ref{lem:qmp-transport} & \texttt{MorseWitness.prod} (all six clauses; the contained cases of the order clause via \texttt{alphaSet\_subset\_closure\_range}, \texttt{omegaSet\_subset\_closure\_range})\\
Theorem \ref{thm:qmp-sound} & \texttt{qMorseProduct\_sound} --- unconditional, no interface hypotheses\\
witness clauses (the interface targeted) & \texttt{MorseWitness}, \texttt{MorseSat} (QtslSoundness.lean, bound in the \rulename{QMorseDischarge} section)\\
Remark \ref{rem:qmp-fence} (readout transport, flows) & \emph{deliberately not mechanized}: Kalies--Wehbe \cite{kalieswehbe2026}, attractor-lattice coproduct isomorphism; Birkhoff duality \cite{birkhoff1937}\\
interim discharge path & \rulename{QMorseDischarge} on the composite lattice (that section's theorem chain) --- no longer needed as warrant; retained as the certificate-route alternative\\
\end{anchors}


\rulesection{72}{QBlockInv}{core rule}
\label{blcatrule-QBlockInv}
\ruleblock{The rule (Steele metanotation)}
\sourcenote{as realized in the mechanized development; source row: \texttt{\texttt{QTSL-rules.toml}} --- \emph{QBlockInv}. The metanotation statement is typeset by hand for notational quality; its mechanized ground truth is the Lean statement below, extracted from the sources at every build of this book and hash-bound. The rule \emph{roster} is machine-checked (every mechanized rule tag round-trips to exactly one catalogue row); the statement \emph{text} is editorial, checked by eye against the extracted Lean beside it.}
\begin{mathpar}
\inferrule*{
  N : \mathsf{QRegionAtom} \ \text{(the block)} \\
  \bigl\{F_f \vdash \mathcal{A}\ell_f \ge \varepsilon_f\bigr\}_{f \in \mathrm{facets}(N)}, \quad \varepsilon_f \in \Q_{>0} \ \text{(} \ell_f \ \text{the facet\textquotesingle s inward affine coordinate; } F_f \ \text{its facet region)} \\
  \mathsf{FacetBridge}(\mathcal{M}, D, N, (F_f, \varepsilon_f)_f) \ \text{(the bridge package, defined below)} \\
  \vdash X_0 \subseteq N \ \text{(\texttt{init\_containment}, a single BL leaf)}
}{
  \mathcal{M} \vDash \mathsf{QInvScope}(N, X_0)
}
\end{mathpar}

\ruleflavor{each facet-inward row bounds the flow derivative of the facet\textquotesingle s inward coordinate ($\ell_f$: affine, vanishing on the facet hyperplane, nonnegative on the block\textquotesingle s side) below by a strictly positive rational margin, on the facet region $F_f$. The displayed $\mathsf{FacetBridge}$ package comprises: the block realized closed in its arena; each facet region a neighbourhood of its entire facet box, corners included (tangential margin, not merely normal thickening); the per-facet wedge test (one-directional by design); arena invariance ($N \subseteq D$, $D$ invariant); trajectory continuity and the pointwise-to-integrated row bridge --- checker-recomputed or interface-carried per profile (model-class).}
\par\medskip
\rulename{QBlockInv} concludes the invariance-scope judgment: the block
$N$ is forward-invariant and encloses the declared initial region. It
is the Morse family\textquotesingle s workhorse producer --- the
judgment it mints is consumed by \rulename{QBasin}\textquotesingle s
safety factory, by scope-restriction of reach obligations, and as the
arena warrant of nested certificates --- and its premise kit is
minimal: per internal facet of the box-union block, one strictly-inward
sign row with explicit positive rational margin (facets on the arena
boundary owe nothing), plus one containment row $X_0 \subseteq N$.

What the checker recomputes: the block\textquotesingle s grid geometry
--- that the bitmap realizes a closed box-union inside its arena, and
the facet-list well-formedness --- plus the replay of every row through
the base kernel. What would break without each hypothesis: without the
sign rows, exit through a facet; without strictness-by-margin, a strict
field sign could not be certified through the non-strict enclosure
predicates of exact-$\Q$ arithmetic (and the Morse rules that cite the
same rows would lose the arguments where strictness is genuinely
load-bearing); without the containment row, a start outside the block;
without closedness of the realization, the first-exit
argument\textquotesingle s sup-time step fails.

The mechanization is conditional at a named interface: the bound
theorems descend the invariance obligation through a first-exit
skeleton (closed block, continuous trajectories, boundary dwell) and a
facet-row bridge (integrated growth of the facet observables through a
one-directional wedge test), leaving exactly three hypothesis families
--- the integrated row clause (reachable from the pointwise rows
through a mean-value bridge whose surviving content is the
profile-level identification of the certified rate with the
trajectory-wise derivative), the grid facts of the box-union, and
arena invariance (nested citation or model-class warrant, terminating
at the whole space). The proof section states the ledger, and why the
wedge test must be one-directional (reentrant corners of box-unions).

\cataloguenote{The same facet rows serve three consumers --- this rule, the lattice
membership of the family object, and the Morse
discharge\textquotesingle s facet arguments --- which is the design
reason strict margins are mandatory here even though this
rule\textquotesingle s own conclusion would survive non-strict signs
(the proof section\textquotesingle s strictness remark).
}
\ruleblock{Lean statement}
\declhead{BernsteinLogic.QtslSoundness}{qBlockInv\_sound}{\badge{BadgeGreen}{statement bound}}
\begin{leansrc}
/-- QBlockInv soundness, the M(i) corollary (roster: "2 lines"). The
facet rows deliver `hInv` through Lemma M(i) — strict facet-inwardness
of the compact box-union block implies forward invariance (gap G1, tier
A prose); `init_containment` delivers `hInit`. The rule-level content is
the assembly of the `QInvScope` claim. -/
theorem qBlockInv_sound {φ : Semiflow X} {N X0 : Set X}
    (hInv : φ.Invariant N) (hInit : X0 ⊆ N) :
    φ.InvScopeSat N X0
\end{leansrc}
\declhead{BernsteinLogic.QtslDischarge}{qBlockInv\_sound\_of\_dwell}{\badge{BadgeGreen}{statement bound}}
\begin{leansrc}
/-- QBlockInv soundness with G1 at its residue: the interface fact
drops from "the block is invariant" to "boundary points dwell"
(G1b — the facet-row bridge). -/
theorem qBlockInv_sound_of_dwell {φ : Semiflow X} {N X0 : Set X}
    (hN : IsClosed N)
    (htraj : ∀ x, Continuous fun t : ℝ≥0 => φ.flow t x)
    (hdwell : ∀ x ∈ N, x ∈ frontier N →
      ∃ δ : ℝ≥0, 0 < δ ∧ ∀ s ≤ δ, φ.flow s x ∈ N)
    (hInit : X0 ⊆ N) : φ.InvScopeSat N X0
\end{leansrc}
\declhead{BernsteinLogic.QtslDischarge}{invariant\_of\_boundaryDwell}{\badge{BadgeGreen}{statement bound}}
\begin{leansrc}
/-- G1's tier-A "first-exit argument", mechanized. A closed region with
continuous trajectories is forward-invariant as soon as every BOUNDARY
point dwells (stays in the region for some positive time) — interior
points dwell by continuity alone. The residue G1b (certified strict
facet-inward sign rows ⇒ the boundary dwell) is the polynomial-field
bridge, prose tier A, same class as G2.

Proof shape = the note's first-exit: the set of times whose whole
prefix stays in `N` is nonempty, closed under the sup (closedness +
continuity), and extends past its sup by the dwell at the sup point —
so it is unbounded, i.e. all of `ℝ≥0`. -/
theorem invariant_of_boundaryDwell {φ : Semiflow X} {N : Set X}
    (hN : IsClosed N)
    (htraj : ∀ x, Continuous fun t : ℝ≥0 => φ.flow t x)
    (hdwell : ∀ x ∈ N, x ∈ frontier N →
      ∃ δ : ℝ≥0, 0 < δ ∧ ∀ s ≤ δ, φ.flow s x ∈ N) :
    φ.Invariant N
\end{leansrc}
\declhead{BernsteinLogic.QtslDischarge}{HasFacetRows.boundaryDwell}{\badge{BadgeGreen}{statement bound}}
\begin{leansrc}
/-- G1b mechanized at its honest stratum: facet rows + block-in-arena +
a forward-invariant arena ⇒ the boundary dwell that
`invariant_of_boundaryDwell` consumes. Proof shape = the entry half of
the interior-entry argument ("interior entry is the same sign fact
read forward"): at a boundary point, trajectory continuity confines a short
prefix to the wedge neighborhood and every active certified region at
once (finitely many); on that prefix each active observable grows from
`0` at rate `ε i` — nonnegative — arena invariance supplies the `D`
conjunct, and the wedge test re-admits every prefix point. Inactive
facets never appear: the wedge test quotes only active ones. -/
theorem HasFacetRows.boundaryDwell {φ : Semiflow X} {D N : Set X}
    (h : φ.HasFacetRows D N) (hND : N ⊆ D) (hD : φ.Invariant D)
    (htraj : ∀ x, Continuous fun t : ℝ≥0 => φ.flow t x) :
    ∀ x ∈ N, x ∈ frontier N →
      ∃ δ : ℝ≥0, 0 < δ ∧ ∀ s ≤ δ, φ.flow s x ∈ N
\end{leansrc}
\declhead{BernsteinLogic.QtslDischarge}{qBlockInv\_sound\_of\_facetRows}{\badge{BadgeGreen}{statement bound}}
\begin{leansrc}
/-- QBlockInv soundness from the certificate's facet rows — the
composed G1 corollary at the current floor. Standing interface facts,
each at its named honest stratum (module header, G1 entry): `hN` and
`hND` are grid facts of the block (compact box-union in its arena),
`hD` is the model-class arena fact (`invariant_univ` when the block
has no ∂D-faces), `hrows` carries the G2-class integrated row residue
and the R1-class wedge geometry, `htraj` is flow continuity (G3's
model-class clause). -/
theorem qBlockInv_sound_of_facetRows {φ : Semiflow X} {D N X0 : Set X}
    (hN : IsClosed N)
    (htraj : ∀ x, Continuous fun t : ℝ≥0 => φ.flow t x)
    (hND : N ⊆ D) (hD : φ.Invariant D) (hrows : φ.HasFacetRows D N)
    (hInit : X0 ⊆ N) : φ.InvScopeSat N X0
\end{leansrc}
\declhead{BernsteinLogic.QtslMeanValue}{qBlockInv\_sound\_of\_rows}{\badge{BadgeGreen}{statement bound}}
\begin{leansrc}
/-- QBlockInv soundness from pointwise rows — G1's chain at its new
floor: per-facet row bounds + profile + wedge geometry + arena
invariance + model-class regularity ⇒ `InvScopeSat`. Composes
`hasFacetRows_of_rows` into `qBlockInv_sound_of_facetRows`; every
hypothesis is now either R1-recomputable grid geometry, a model-class
fact, exact row arithmetic, or the profile identification. -/
theorem qBlockInv_sound_of_rows {φ : Semiflow X} {D N X0 : Set X}
    {n : ℕ} {ℓ : Fin n → X → ℝ} {ε : Fin n → ℝ} {F : Fin n → Set X}
    (W : Fin n → X → ℝ)
    (hN : IsClosed N)
    (htraj : ∀ x, Continuous fun t : ℝ≥0 => φ.flow t x)
    (hND : N ⊆ D) (hD : φ.Invariant D)
    (hε : ∀ i, 0 < ε i)
    (hℓ : ∀ i, ContinuousOn (ℓ i) (F i))
    (hW : ∀ i, φ.FlowDerivOn (ℓ i) (W i) (F i))
    (hrow : ∀ i, ∀ y ∈ F i, ε i ≤ W i y)
    (hlocal : ∀ x ∈ frontier N, ∃ (A : Set (Fin n)) (U : Set X), U ∈ 𝓝 x ∧
      (∀ i ∈ A, ℓ i x = 0 ∧ F i ∈ 𝓝 x) ∧
      ∀ y ∈ U, y ∈ D → (∀ i ∈ A, 0 ≤ ℓ i y) → y ∈ N)
    (hInit : X0 ⊆ N) : φ.InvScopeSat N X0
\end{leansrc}
\declhead{BernsteinLogic.QtslGridGeometry}{qBlockInv\_sound\_of\_grid}{\badge{BadgeGreen}{statement bound}}
\begin{leansrc}
/-- QBlockInv soundness on the grid model, end to end: grid geometry
discharged outright; the remaining hypotheses are exactly the named
residues — the arena invariance interface `hArena` (residue (c), see
the two reduction theorems below) and the integrated row clause
(residue (a), G2's class) — plus flow continuity (G3's model-class
clause). -/
theorem qBlockInv_sound_of_grid (φ : Semiflow (Fin d → ℝ)) (B : Finset G.CellIdx)
    {X0 : Set (Fin d → ℝ)}
    (htraj : ∀ x, Continuous fun t : ℝ≥0 => φ.flow t x)
    (hArena : φ.Invariant G.arena)
    (ε : G.FacetData → ℝ) (F : G.FacetData → Set (Fin d → ℝ))
    (hε : ∀ p, G.IsFacet B p → 0 < ε p)
    (hF : ∀ p, G.IsFacet B p → ∀ x ∈ G.facetBox p, F p ∈ 𝓝 x)
    (hrow : ∀ p, G.IsFacet B p → ∀ x ∈ F p, ∀ t : ℝ≥0,
      (∀ s ≤ t, φ.flow s x ∈ F p) →
      G.obs p x + ε p * (t : ℝ) ≤ G.obs p (φ.flow t x))
    (hInit : X0 ⊆ G.blockRegion B) :
    φ.InvScopeSat (G.blockRegion B) X0
\end{leansrc}
\declhead{BernsteinLogic.QtslGridGeometry}{qBlockInv\_sound\_of\_pointwiseRows}{\badge{BadgeGreen}{statement bound}}
\begin{leansrc}
/-- CAPSTONE of the same-day gap-fill pair (this file × the interval
mean-value bridge, `QtslMeanValue.lean`): grid block + POINTWISE row
bounds, consumed through the profile predicate `FlowDerivOn`, ⇒
`InvScopeSat`. Every remaining hypothesis is a model-class regularity
fact (`htraj`, `hV`), exact row arithmetic (`hε`, `hrow`), the arena
interface (`hArena`, residue (c)'s reductions), the certified-region
neighborhood fact (`hF`), or the profile identification (`hW` —
`W p = 𝒜(obs p)` for the ODE doctrine, the sole analytic remainder of
the whole G1 chain). -/
theorem qBlockInv_sound_of_pointwiseRows (φ : Semiflow (Fin d → ℝ))
    (B : Finset G.CellIdx) {X0 : Set (Fin d → ℝ)}
    (htraj : ∀ x, Continuous fun t : ℝ≥0 => φ.flow t x)
    (hArena : φ.Invariant G.arena)
    (ε : G.FacetData → ℝ) (F : G.FacetData → Set (Fin d → ℝ))
    (W : G.FacetData → (Fin d → ℝ) → ℝ)
    (hε : ∀ p, G.IsFacet B p → 0 < ε p)
    (hF : ∀ p, G.IsFacet B p → ∀ x ∈ G.facetBox p, F p ∈ 𝓝 x)
    (hV : ∀ p, G.IsFacet B p → ContinuousOn (G.obs p) (F p))
    (hW : ∀ p, G.IsFacet B p → φ.FlowDerivOn (G.obs p) (W p) (F p))
    (hrow : ∀ p, G.IsFacet B p → ∀ y ∈ F p, ε p ≤ W p y)
    (hInit : X0 ⊆ G.blockRegion B) :
    φ.InvScopeSat (G.blockRegion B) X0
\end{leansrc}
\ruleblock{Correspondence (Steele display $\leftrightarrow$ Lean)}
\badge{BadgeAmber}{a bound Lean statement CHANGED since certification --- re-review and re-certify}
\sourcenote{The dictionary below is a reviewer-certified statement of how the mechanized statements above interpret the displayed rule: display metavariable $\mapsto$ Lean binder, each display hypothesis $\mapsto$ the premise or side condition carrying it, the display conclusion $\mapsto$ the goal --- with any deliberate difference in width between display and mechanization named outright. Certification is recorded as a hash triple over (the displayed cell, the extracted statement texts, this dictionary\textquotesingle s own text); if any of the three changes, this block turns amber until a reviewer re-certifies. The correspondence claim itself remains editorial --- reviewer-grade, not machine-checked; its machine-checked successor types each soundness statement against the rule as a structured object, and supersedes this block per rule as it lands.}
The displayed judgment $\mathcal{M} \vDash \mathsf{QInvScope}(N, X_0)$
$\mapsto$ the goal \texttt{InvScopeSat N X0} = \texttt{φ.Invariant N ∧
X0 ⊆ N} --- the claim record\textquotesingle s denotation verbatim.
The five bound statements form one discharge chain:
\texttt{qBlockInv\_sound} (the roster\textquotesingle s ``two lines'':
invariance + containment $\Rightarrow$ the claim, M(i) as hypothesis),
\texttt{\_of\_dwell} riding \texttt{invariant\_of\_boundaryDwell} (G1
mechanized: closed block + trajectory continuity + boundary dwell
$\Rightarrow$ invariance, the first-exit sup-time argument), and
\texttt{\_of\_facetRows} riding \texttt{HasFacetRows.boundaryDwell},
the deepest bound form, whose kit is the honest ledger:
\begin{itemize}
\item the block $N : \mathsf{QRegionAtom}$ $\mapsto$ an abstract
  \texttt{N : Set X} with \texttt{hN : IsClosed N}: the grid stratum
  (bitmap realizes a closed box-union in its arena; facet-list
  well-formedness) is checker-side recompute, and closedness is the
  one grid fact the bound statements consume --- as an explicit
  hypothesis, not a derived fact.
\item facet-inward sign rows $\mapsto$ \texttt{hrows :
  φ.HasFacetRows D N}: finitely many facet observables $\ell_i$ with
  explicit positive margins (\texttt{∀ i, 0 < ε i} --- the
  display\textquotesingle s strictness-by-margin, carried exactly),
  each growing at \emph{integrated} rate $\ge \varepsilon_i$ along
  confined prefixes --- the pointwise row $\mathcal{A}\ell \ge
  \varepsilon$ reaches that form through the G2-class mean-value
  bridge, outside this bound set --- plus the local wedge test,
  one-directional by design (reentrant corners of box-unions defeat
  the $\leftrightarrow$ form), with ``$\partial D$-faces owe
  nothing'' landing as the wedge test\textquotesingle s \texttt{y ∈
  D} conjunct.
\item \texttt{init\_containment} $\mapsto$ \texttt{hInit : X0 ⊆ N},
  one-to-one, unchanged at every link of the chain.
\item the displayed \textsf{FacetBridge} package (its contents
  defined in the flavor text) $\mapsto$
  \texttt{htraj} (trajectory continuity --- model-class), and the
  arena pair \texttt{hND : N ⊆ D} + \texttt{hD : φ.Invariant D} ---
  the certified statement makes the invariance regress explicit (the
  roster witness names its termination at \texttt{invariant\_univ},
  which is in the module but not this bound set).
\end{itemize}
Width note: the conditional flag\textquotesingle s three named
families are visible verbatim in the facet-row
statement\textquotesingle s binders --- the integrated row clause
(inside \texttt{HasFacetRows}), the grid facts (\texttt{hN}, plus the
wedge-test geometry carried \emph{inside} \texttt{hrows} rather than
discharged), and arena invariance (\texttt{hD}). The bound set
extends past the two-line core (the two blocks below): the
row-to-shadow discharge \texttt{qBlockInv\_sound\_of\_rows} (module
\texttt{QtslMeanValue} --- pointwise rows reach the integrated shadow
form through the interval mean-value bridge, leaving the profile
identification \texttt{hW} as the sole analytic remainder of the G1
chain), the grid-model discharge \texttt{qBlockInv\_sound\_of\_grid}
(module \texttt{QtslGridGeometry} --- grid geometry discharged
outright, residues exactly the arena interface and the integrated row
clause), and the capstone \texttt{qBlockInv\_sound\_of\_pointwiseRows}
(grid block + pointwise rows through \texttt{FlowDerivOn} to the
judgment) are hash-bound, so the conditional residue shrinks to
the profile identification plus the arena-invariance interface
(its reduction theorems terminate at \texttt{invariant\_univ}), the
facet-neighbourhood interface clause (each certified region $F_i$ a
neighbourhood of every point of its facet box, corners included ---
the capstone consumes it as an explicit hypothesis) and
model-class regularity. One deliberate width inversion: the
display\textquotesingle s conclusion tie is an unconditional judgment
about the model, while \texttt{qBlockInv\_sound} alone is wider than
the rule (any invariant set qualifies); the facet-row form narrows
back to certificate-shaped input. No hypothesis in the bound set
demands the block be a box-union --- the geometry enters only through
what the wedge test can actually be discharged for.

\ruleblock{Lean verifier receipts}
\sourcenote{module \texttt{BernsteinLogic.QtslSoundness}; bound build \texttt{77005991d61d\ldots} (bound from receipt 8130a7a7a679cef4, 2026-08-04T01:16:06Z). A receipt records that named checking layers accepted a build: (module, compiled-artifact SHA-256, environment SHA-256, checker layers), minted when the checking gate ran; the store receipts attesting the bound build are listed with the layers each run credited. Nothing here asks the reader to trust the store: every layer is an independently runnable program over the quoted sources --- rebuild the module and re-run the checkers to reproduce the verdict. Receipts attesting superseded builds credit no layer and are omitted from this rendering (the count of omissions, if any, is noted below).}
\begin{receipttable}
\receiptrow{8130a7a7a679cef4}{2026-08-04T01:16:06Z}{build+decorrelated}{\badge{BadgeGreen}{current}}
\receiptrow{9eb3706354082144}{2026-08-04T02:23:12Z}{build+decorrelated}{\badge{BadgeGreen}{current}}
\end{receipttable}
\sourcenote{5 further receipt(s) in the store attest superseded builds of this module and are omitted from this rendering.}
\sourcenote{module \texttt{BernsteinLogic.QtslDischarge}; bound build \texttt{4aa46910a600\ldots} (bound from receipt dd2946805f790f99, 2026-08-04T01:16:09Z). A receipt records that named checking layers accepted a build: (module, compiled-artifact SHA-256, environment SHA-256, checker layers), minted when the checking gate ran; the store receipts attesting the bound build are listed with the layers each run credited. Nothing here asks the reader to trust the store: every layer is an independently runnable program over the quoted sources --- rebuild the module and re-run the checkers to reproduce the verdict. Receipts attesting superseded builds credit no layer and are omitted from this rendering (the count of omissions, if any, is noted below).}
\begin{receipttable}
\receiptrow{dd2946805f790f99}{2026-08-04T01:16:09Z}{build+decorrelated}{\badge{BadgeGreen}{current}}
\receiptrow{e1bc68890e2aebf1}{2026-08-04T02:23:11Z}{build+decorrelated}{\badge{BadgeGreen}{current}}
\end{receipttable}
\sourcenote{14 further receipt(s) in the store attest superseded builds of this module and are omitted from this rendering.}
\sourcenote{module \texttt{BernsteinLogic.QtslMeanValue}; bound build \texttt{ed04b606934b\ldots} (bound from receipt c2f57d21d0b1c270, 2026-07-20T08:27:08Z). A receipt records that named checking layers accepted a build: (module, compiled-artifact SHA-256, environment SHA-256, checker layers), minted when the checking gate ran; the store receipts attesting the bound build are listed with the layers each run credited. Nothing here asks the reader to trust the store: every layer is an independently runnable program over the quoted sources --- rebuild the module and re-run the checkers to reproduce the verdict. Receipts attesting superseded builds credit no layer and are omitted from this rendering (the count of omissions, if any, is noted below).}
\begin{receipttable}
\receiptrow{160082fd7894faf1}{2026-07-19T07:39:26Z}{build+decorrelated}{\badge{BadgeGreen}{current}}
\receiptrow{260eab5039951515}{2026-07-19T09:07:22Z}{build+decorrelated}{\badge{BadgeGreen}{current}}
\receiptrow{c2f57d21d0b1c270}{2026-07-20T08:27:08Z}{build+decorrelated}{\badge{BadgeGreen}{current}}
\receiptrow{f75d120f9128f80c}{2026-07-20T20:44:49Z}{build+decorrelated}{\badge{BadgeGreen}{current}}
\receiptrow{891ea7ab9c4f79ff}{2026-07-20T21:01:17Z}{build+replay+hardened+decorrelated}{\badge{BadgeGreen}{current}}
\receiptrow{af025f8810d8a6cc}{2026-07-21T05:28:44Z}{build+replay+hardened+decorrelated}{\badge{BadgeGreen}{current}}
\end{receipttable}
\sourcenote{module \texttt{BernsteinLogic.QtslGridGeometry}; bound build \texttt{9060a661920e\ldots} (bound from receipt 64d6a21041b4042f, 2026-07-20T08:27:10Z). A receipt records that named checking layers accepted a build: (module, compiled-artifact SHA-256, environment SHA-256, checker layers), minted when the checking gate ran; the store receipts attesting the bound build are listed with the layers each run credited. Nothing here asks the reader to trust the store: every layer is an independently runnable program over the quoted sources --- rebuild the module and re-run the checkers to reproduce the verdict. Receipts attesting superseded builds credit no layer and are omitted from this rendering (the count of omissions, if any, is noted below).}
\begin{receipttable}
\receiptrow{bf51898b80b37884}{2026-07-19T12:35:27Z}{build+decorrelated}{\badge{BadgeGreen}{current}}
\receiptrow{64d6a21041b4042f}{2026-07-20T08:27:10Z}{build+decorrelated}{\badge{BadgeGreen}{current}}
\receiptrow{c7deda94ed05f8e5}{2026-07-20T20:44:45Z}{build+decorrelated}{\badge{BadgeGreen}{current}}
\receiptrow{6bab7a61aa41db61}{2026-07-20T21:01:14Z}{build+replay+hardened+decorrelated}{\badge{BadgeGreen}{current}}
\receiptrow{5e72e577b5f436e5}{2026-07-21T05:28:39Z}{build+replay+hardened+decorrelated}{\badge{BadgeGreen}{current}}
\end{receipttable}
\sourcenote{2 further receipt(s) in the store attest superseded builds of this module and are omitted from this rendering.}
\layerline{\layeritem{build}{lake build (elaboration)}{\badge{BadgeGreen}{current}}\layeritem{replay}{leanchecker (official kernel replay)}{\badge{BadgeGreen}{current}}\layeritem{hardened}{safe\_verify (hardened screen)}{\badge{BadgeGreen}{current}}\layeritem{decorrelated}{lean4export + nanoda (decorrelated kernel)}{\badge{BadgeGreen}{current}}}
\sourcenote{\badge{BadgeGray}{comparator: none} no fixture in the three-implementation comparator harness (compiled / kernel-interpreted / reference-interpreted) exercises this rule.}
\ruleblock{HOL4 statement}
\sourcenote{theory \texttt{bern\_tr\_t5qdyn}, HOL4 development script \texttt{\texttt{bern\_tr\_t5qdynScript.sml}} \badge{BadgeGreen}{script bound}}
\holhead{Definition sfInvariant\_def}
\begin{holsrc}
Definition sfInvariant_def:
  sfInvariant (fl:real->'x->'x) (sv:'x->bool) ⇔
    ∀x. x IN sv ⇒ ∀t. 0 ≤ t ⇒ fl t x IN sv
End
\end{holsrc}
\holhead{Definition sfInvScopeSat\_def}
\begin{holsrc}
Definition sfInvScopeSat_def:
  sfInvScopeSat (fl:real->'x->'x) (sv:'x->bool) (X0:'x->bool) ⇔
    sfInvariant fl sv ∧ X0 SUBSET sv
End
\end{holsrc}
\holhead{Theorem sfQBlockInv\_sound}
\begin{holsrc}
Theorem sfQBlockInv_sound:
  ∀fl (N:'x->bool) X0.
    sfInvariant fl N ∧ X0 SUBSET N ⇒ sfInvScopeSat fl N X0
\end{holsrc}
\ruleblock{HOL verifier receipts}
\badge{BadgeBlue}{documentary} \texttt{hol-strict} CLEAN: theory \texttt{bern\_tr\_t5qdyn}, 22
theorems, 0 cheat-tagged, 0 nonstandard axioms (whitelist:
\texttt{DISK\_THM}; the four \texttt{bool} axioms) --- verification
run of 2026-07-22 on an isolated build host, against a
content-addressed checkout of the exact sources cited here (git blob
\texttt{d64e1954\ldots}); an independent acceptance host ran the same
build green at the porting seating (2026-07-21). What is ported is
the semiflow-stratum root: the flow law, the invariance and
invariance-scope-satisfaction predicates, and the root soundness
theorem concluding the scope judgment from invariance plus
containment --- the exact root the Lean development\textquotesingle s
satisfaction stratum supplies for this rule. The facet-row bridge
(the first-exit skeleton and integrated-growth descent that discharge
invariance from the certificate\textquotesingle s sign rows) and the
grid-geometry recomputation remain Lean-side; the conditional reading
of this section\textquotesingle s mechanization carries over to the
HOL twin unchanged. The facet-sign and containment rows the rule
consumes are covered by the ported base checker as before.

\sourcenote{HOL receipts are documentary: quoted results of \texttt{hol-strict} runs
and census rows from the HOL4 development\textquotesingle s build
evidence, not minted receipt artefacts. A HOL receipt store (the Lean
store\textquotesingle s twin) is named future work.
}
\ruleblock{Proof}

\sourcenote{Genre: core rule (tier 1) --- an ordinary mathematical proof
at Putnam-exam-solution detail. The rule's mechanization is
\emph{conditional}: the section descends the invariance obligation from
an opaque interface fact to three named hypothesis families ---
integrated facet-row clauses, grid geometry of the block, arena
invariance --- and the ledger remark at the end states exactly what
remains and why. Every hypothesis is used and remarked on; the
formal-anchor table maps each step to the Lean declaration that carries
it.}

\textbf{Setting and claim.}\enspace Fix a forward semiflow $\varphi$ on
a state space $X$ ($\Phi_0 = \mathrm{id}$,
$\Phi_{s+t} = \Phi_t \circ \Phi_s$; the \rulename{QBasin} section's
setting, which also fixes the $\mathsf{QInvScope}$ denotation: $S$
forward-invariant and $X_0 \subseteq S$). The rule's data: a
\emph{block} $N$ --- the realization of a box-union over the
certificate's declared rational grid, a compact closed set --- sitting
in an \emph{arena} $D \supseteq N$; per internal facet of $N$, one
strictly-inward sign row of the base calculus (a
\rulename{BernCone}-family sequent asserting the field's normal
component points into the block, with an explicit positive rational
margin; facets on $\partial D$ owe nothing); and one containment row
$X_0 \subseteq N$. The claim: these data warrant
$\mathsf{QInvScope}(N, X_0)$ --- established below at the discharged
floor, conditional on the named ledger of
\autoref{rem:qbi-ledger}.

\begin{blemma}[rule-level assembly]\label{lem:qbi-assembly}
If $N$ is forward-invariant and $X_0 \subseteq N$, then
$\mathsf{QInvScope}(N, X_0)$ holds.
\end{blemma}
\begin{proof}
The judgment \emph{is} that conjunction. The lemma is stated to fix
where the rule-level content ends: everything below descends the
invariance conjunct toward the certificate's rows.
\end{proof}

\begin{blemma}[the first-exit skeleton: closed $+$ continuous $+$
boundary dwell $\Rightarrow$ invariant]\label{lem:qbi-dwell}
Suppose $N$ is closed, every trajectory $t \mapsto \Phi_t x$ is
continuous, and every boundary point dwells: for each
$x \in N \cap \partial N$ there is $\delta > 0$ with
$\Phi_s x \in N$ for all $0 \leq s \leq \delta$. Then $N$ is
forward-invariant.
\end{blemma}
\begin{proof}
First, \emph{every} point of $N$ dwells: an interior point $x$ has
$\Phi_0 x = x \in \mathrm{int}\,N$, so continuity of the trajectory at
time $0$ confines a short prefix to $\mathrm{int}\,N$ --- the dwell
hypothesis is needed only at the frontier, which is where the
certificate's rows live.

Now fix $x \in N$ and $t \geq 0$, and suppose toward a contradiction
that $\Phi_t x \notin N$. Consider the set of \emph{prefix-confined
times}
\[
S \;:=\; \{\, u \geq 0 \ :\ \Phi_s x \in N \ \text{for all } s \leq u
\,\}.
\]
$S$ contains $0$, and every element of $S$ is $\leq t$ (an element
$> t$ would confine time $t$ itself); throughout, times $s, u$ range
over $[0, \infty)$. Let $T := \sup S \leq t$. Every
time $s < T$ is confined (it is below some element of $S$). The sup
itself is confined: for $T > 0$ the preimage of $N$ under the
continuous trajectory is closed and contains $[0, T)$, hence contains
$T \in \mathrm{cl}\,[0,T)$; for $T = 0$ it is the hypothesis
$x \in N$. So $T \in S$. But now the dwell at the point $\Phi_T x \in N$
supplies $\delta > 0$ with $\Phi_s(\Phi_T x) \in N$ for
$0 \leq s \leq \delta$,
i.e. (semiflow law) $\Phi_u x \in N$ for $T \leq u \leq T + \delta$;
combined with $T \in S$ this puts $T + \delta \in S$, contradicting
$T = \sup S$ with $\delta > 0$.
\end{proof}

\begin{bremark}[each skeleton hypothesis is load-bearing]
\emph{Closedness}: on $X = \R$ with $\Phi_t x = x + t$ and the
half-open $N = [0, 1)$, every point of $N$ dwells (interior points by
continuity; the one frontier point of $N$, namely $0$, dwells since
$\Phi_s 0 = s \in [0,1)$ for small $s$ --- the other frontier point
$1$ is not in $N$), yet the trajectory from $0$ exits: the proof fails
exactly at
``the sup itself is confined'', which is where closedness is consumed.
\emph{Trajectory continuity}: with a discontinuous time-dependence the
preimage argument dies and a trajectory may jump across the frontier
without ever presenting a boundary point to the dwell hypothesis; in
this stratum continuity is a model-class fact (polynomial-drift flows),
not a theorem. \emph{Dwell at the frontier}: dropping it re-opens exit
through the boundary --- it is precisely the obligation the facet rows
discharge (next lemma).
\end{bremark}

\begin{blemma}[the facet-row bridge: certified rows $\Rightarrow$
boundary dwell]\label{lem:qbi-bridge}
Suppose the certificate supplies, for block $N$ in arena $D$: finitely
many \emph{facet observables} $\ell_1, \dots, \ell_n : X \to \R$ (the
affine facet polynomials), margins $\varepsilon_i > 0$, and certified
regions $F_i$ (boxes around the facets, thickened to full dimension
\emph{with tangential margin}: each $F_i$ must be a neighbourhood of
every point of its facet, corners included --- a box thickened only in
the normal direction is a neighbourhood of no corner point, and the
wedge test below demands one), such that

\emph{(row, integrated form)} along any trajectory segment confined to
$F_i$, $\ell_i$ grows at rate at least $\varepsilon_i$:
if $\Phi_s x \in F_i$ for all $s \leq t$, then
$\ell_i(\Phi_t x) \geq \ell_i(x) + \varepsilon_i\, t$; and

\emph{(local wedge test)} every frontier point $x$ of $N$ has an
\emph{active set} $A \subseteq \{1, \dots, n\}$ --- each active
observable vanishing at $x$, each active $F_i$ a neighbourhood of $x$
--- and a neighbourhood $U$ of $x$ such that every $y \in U$ with
$y \in D$ and $\ell_i(y) \geq 0$ for all $i \in A$ lies in $N$.

If moreover $N \subseteq D$, $D$ is forward-invariant, and trajectories
are continuous, then every boundary point of $N$ dwells.
\end{blemma}
\begin{proof}
Fix $x \in N \cap \partial N$ with active set $A$ and wedge
neighbourhood $U$. The set $U \cap \bigcap_{i \in A} F_i$ is a finite
intersection of neighbourhoods of $x$, hence a neighbourhood; continuity
of the trajectory at time $0$ (where it sits at $x$) confines a prefix
$[0, \delta]$ to it, some $\delta > 0$. Now take any
$0 \leq s \leq \delta$
and check the wedge test's three conditions at $y = \Phi_s x$:
membership in $U$ is the confinement; membership in $D$ is arena
invariance applied to $x \in N \subseteq D$; and for each active $i$,
the whole prefix up to $s$ is confined to $F_i$, so the integrated row
gives $\ell_i(\Phi_s x) \geq \ell_i(x) + \varepsilon_i s
= 0 + \varepsilon_i s \geq 0$, using exactly the activation fact
$\ell_i(x) = 0$. The wedge test re-admits $\Phi_s x \in N$. Inactive
facets never appear: the test quotes only active ones.
\end{proof}

\begin{bthm}[soundness of \rulename{QBlockInv}, at the discharged
floor]\label{thm:qbi-sound}
Suppose: $N$ is closed with $N \subseteq D$; $D$ is forward-invariant;
trajectories are continuous; the certificate's facet rows satisfy
Lemma~\ref{lem:qbi-bridge}'s two clauses; and $X_0 \subseteq N$. Then
$\mathsf{QInvScope}(N, X_0)$ holds.
\end{bthm}
\begin{proof}
Lemma~\ref{lem:qbi-bridge} delivers the boundary dwell;
Lemma~\ref{lem:qbi-dwell} converts it to forward invariance of $N$;
Lemma~\ref{lem:qbi-assembly} assembles the judgment with the
containment row.
\end{proof}

\begin{bremark}[the wedge test is one-directional by design]
\label{rem:qbi-wedge}
At a \emph{reentrant} corner of a box-union block --- an inner corner
of an L-shape --- the block is locally the \emph{union} of two
half-spaces $\{\ell_1 \geq 0\} \cup \{\ell_2 \geq 0\}$, while the
wedge test's hypothesis set is their \emph{intersection}: the
implication ``both $\ell_i \geq 0$ $\Rightarrow$ $y \in N$'' holds
(the intersection sits inside the union), but its converse is false
there --- points of $N$ with one active observable negative exist on
the other branch. The biconditional form of the test would therefore
be unsatisfiable at reentrant corners, while the one-directional form
is exactly what the first-exit argument consumes
(Lemma~\ref{lem:qbi-bridge} never needs to \emph{characterize} the
frontier, only to re-admit prefix points). A certificate whose rows
are jointly unsatisfiable at a boundary point --- checkerboard blocks
touching at a corner and demanding opposite strict signs there ---
simply makes the bridge's hypotheses unfulfillable: soundness only
consumes what is given.
\end{bremark}

\begin{bremark}[$\partial D$-faces owe nothing; the arena regress
terminates]
Facets of $N$ on the arena boundary carry no rows: in
Lemma~\ref{lem:qbi-bridge} they enter only through the conjunct
$y \in D$, which arena invariance supplies. The arena's own invariance
is either a nested $\mathsf{QInvScope}$ citation (a chart sub-block
certified by its own facet rows --- this same rule, one level up) or
the model-class fact that a chart's denotation confines trajectories to
its domain; the regress terminates at the whole space, which is
invariant vacuously.
\end{bremark}

\begin{bremark}[strictness of the margins: where it is and is not
consumed]\label{rem:qbi-strict}
The invariance conclusion of this section survives
$\varepsilon_i = 0$: the integrated clause then says $\ell_i$ is
non-decreasing along confined segments, and growth-from-zero still
yields $\ell_i \geq 0$ --- consistent with the classical invariance
theorem (Nagumo \cite{nagumo1942}), which needs only subtangentiality.
The explicit positive rational margin is nonetheless mandatory in the
certificate convention, for two reasons that live outside this
section. First, exactness: a \emph{strict} field sign cannot be
certified directly by the non-strict enclosure predicates of the base
calculus; ``non-strict row on $\mathcal{A}\ell - \varepsilon$ with
explicit $\varepsilon > 0$'' is how strictness becomes checkable in
exact $\Q$. Second, the same rows are consumed by the Morse-family
rules, where strictness is genuinely load-bearing --- no orbit rides a
facet, distinct fibres' invariant parts are disjoint, and the facet
margin bounds admissible chain jumps (``strictness pays a second
time''; see the \rulename{QMorseDischarge} section's ledger). A margin
quietly dropped to zero would leave this rule's conclusion intact and
break those.
\end{bremark}

\begin{bremark}[the hypothesis ledger --- what remains, named]
\label{rem:qbi-ledger}
The chain above discharges the rule's invariance obligation down to
three hypothesis families, and the rule's Lean column renders with the
hypothesis-flavoured check because of exactly these:

\emph{(a) the integrated row clause: discharged, and bound.} The
certificate's rows are pointwise: per facet box, a base-calculus
sequent bounding the field's certified rate,
$\mathcal{A}\ell_i \geq \varepsilon_i$ (the generator
$\mathcal{A} = f \cdot \nabla$ of the deterministic model class).
Lemma~\ref{lem:qbi-bridge} consumes the \emph{integrated} clause along
trajectory segments. The bridge between them is one-variable calculus
--- a mean-value argument along the trajectory, mechanized and bound
in this edition (anchor table) --- and its surviving analytic content
is the identification of the certified polynomial rate with the
trajectory-wise derivative on the region: a model-class profile
warrant (for smooth flows the identification is literal; switched and
hybrid profiles replace it by one-sided derivative bounds, which is
their own seam, deliberately not this section's).

\emph{(b) grid geometry of the block: closed for grid blocks, and
bound --- except one interface clause.} Closedness of $N$,
$N \subseteq D$, the vanishing of each facet observable on its facet
box, and the wedge re-admission implication are
finite geometry of the box-union over its declared rational grid ---
recomputable combinatorics, proved outright for grid blocks in the
grid-geometry module and bound in this edition (anchor table),
end-to-end through the capstone that takes grid block plus pointwise
row bounds to the concluded judgment. The one geometric clause that
stays interface-level even there: each certified region $F_i$ must be
a neighbourhood of every point of its facet box, corners included ---
the tangential-margin requirement of Lemma~\ref{lem:qbi-bridge} ---
which the capstone consumes as a hypothesis on the certificate's
regions rather than proving from the boxes. The interface-level
theorems above remain stated with (b) as hypotheses, which is what the
general (non-grid) reading consumes.

\emph{(c) arena invariance.} As in the regress remark: nested
citation, or model-class warrant, terminating at the whole space.

None of the three is an opaque dynamical assumption; each names a
specific fact with a specific owner --- (a) the profile stratum, (b)
the checker's recompute (and the grid module), (c) the arena's own
certificate. That is the sense in which this core rule is
\emph{conditionally} mechanized: the cited theorems are sorry-free,
and their hypotheses are the ledger above.
\end{bremark}

\begin{bremark}[relation to the literature]
The block-with-inward-facets certificate is the exact-arithmetic
descendant of Conley theory's isolating blocks
\cite{conleyeaston1971}: compact sets whose boundary behaviour is
transverse, so that dynamical claims reduce to boundary sign checks.
The invariance argument itself is Nagumo's theorem
\cite{nagumo1942} in first-exit form, specialized to polyhedral sets,
where the subtangentiality condition becomes finitely many polynomial
sign rows --- the form set-invariance control has long exploited
\cite{blanchini1999}. What the certificate stratum adds is the
discipline: strict signs via non-strict predicates with exact rational
margins (\autoref{rem:qbi-strict}), and a one-directional wedge
test chosen so that box-union (nonconvex) blocks are admissible
without any frontier characterization
(\autoref{rem:qbi-wedge}).
\end{bremark}

\begin{anchors}
Lemma \ref{lem:qbi-assembly} & \texttt{qBlockInv\_sound} (QtslSoundness.lean)\\
Lemma \ref{lem:qbi-dwell}, first-exit skeleton & \texttt{invariant\_of\_boundaryDwell} (QtslDischarge.lean); composed as \texttt{qBlockInv\_sound\_of\_dwell}\\
Lemma \ref{lem:qbi-bridge}, row shadow + wedge test & \texttt{HasFacetRows} (the definition); \texttt{HasFacetRows.boundaryDwell}\\
Theorem \ref{thm:qbi-sound}, end-to-end & \texttt{qBlockInv\_sound\_of\_facetRows}\\
arena regress terminus & \texttt{invariant\_univ}; nested-citation path \texttt{invariant\_of\_invScopeSat} (development)\\
ledger item (a), the mean-value bridge & \texttt{integrated\_ge\_of\_flowDeriv\_ge}, composed as \texttt{qBlockInv\_sound\_of\_rows} (QtslMeanValue; \textbf{bound in this edition})\\
ledger item (b), grid facts & \texttt{Grid.wedge\_containment}, \texttt{Grid.qBlockInv\_sound\_of\_grid}; capstone \texttt{Grid.qBlockInv\_sound\_of\_pointwiseRows} (QtslGridGeometry; \textbf{bound in this edition})\\
facet rows & strictly-inward \rulename{BernCone}-family sequents with explicit rational margin (base kernel)\\
\end{anchors}


\rulesection{73}{QLatticeInv}{alias for \textcolor{DiffBlue}{\rulename{QBlockInv}}}
\label{blcatrule-QLatticeInv}
\ruleblock{The rule (Steele metanotation)}
\sourcenote{as realized in the mechanized development; source row: \texttt{\texttt{QTSL-rules.toml}} --- \emph{QLatticeInv}. The metanotation statement is typeset by hand for notational quality; its mechanized ground truth is the Lean statement below, extracted from the sources at every build of this book and hash-bound. The rule \emph{roster} is machine-checked (every mechanized rule tag round-trips to exactly one catalogue row); the statement \emph{text} is editorial, checked by eye against the extracted Lean beside it.}
\begin{mathpar}
\inferrule*{
  L : \mathsf{QBlockLattice} \\
  C \ \text{(\texttt{downset}, a cell-set; membership and down-closedness recomputed)} \\
  \mathcal{M} \vDash \mathsf{QInvScope}\bigl(\textstyle\bigcup_{i \in C} \mathrm{blocks}(\mathrm{emb}\ i),\ X_0\bigr) \ \text{(the cited \rulename{QBlockInv} conclusion)} \\
  \vdash X_0 \subseteq \textstyle\bigcup\downarrow C \ \text{(\texttt{init\_containment})}
}{
  \mathcal{M} \vDash \mathsf{QInvScope}\bigl(\textstyle\bigcup\downarrow C,\ X_0\bigr)
}
\end{mathpar}

\par\medskip
\rulename{QLatticeInv} is the family object\textquotesingle s
invariance alias: any down-closed cell-set $C$ of the recomputed
readout names an invariant region $\bigcup\downarrow C$, and the rule
concludes its invariance-scope judgment by citation --- the unpacking
is a single \rulename{QBlockInv} application whose block is the
lattice member realizing the down-set (existence by join-closure) and
whose facet rows are that member\textquotesingle s own, already
present in the family object. No row is re-proven; the producer merely
selects which down-set encloses the initial region away from the avoid
sets.

What the checker recomputes: membership of the realization in the
lattice, and down-closedness of $C$ in the recomputed cell order ---
readout facts, not dynamics. What would break without them: a cell-set
that is not down-closed names a region orbits can leave (an orbit
descends the cell order out of it), and the exactness identity between
the down-set\textquotesingle s cell union and the
member\textquotesingle s realization fails --- already on a two-member
chain: the top member alone is not down-closed, and its cell omits
everything the bottom member contributes.

Mechanization state, plainly: the bound statements below are the
exactness identity and the soundness corollaries, at both semantic
strata (module \texttt{BernsteinLogic.QtslLatticeInv}).
\texttt{qLatticeInv\_exact} proves the identity that makes the
respelling exact --- for down-closed $C$, the union of $C$\textquotesingle
s cells equals the union of its join-irreducible realizations --- by
the discharge chain\textquotesingle s own well-founded descent through
lower covers (\texttt{exists\_cellMem\_subset},
over the same descent rows the eventual-member interface consumes ---
rows that assume no distributivity,
plus one row this alias adds: cell inside member, the readout closure
half\textquotesingle s set-level shadow). One sharp point, because the
page\textquotesingle s classical route consumes distributivity: the
mechanized route consumes none, and the row pack does not encode
join-primality --- absorption-positive families satisfy every row, and
the bound (union-form) identity holds in them regardless --- because
the descent runs inside a single witnessed member of $C$ and never
forms the join. The join-form reading (the union as the member
$N_C$\textquotesingle s realization) costs exactly
join-realizes-as-union, the checker\textquotesingle s recomputed
membership of the down-set\textquotesingle s realization, strictly
weaker than distributivity; that recompute is also what refuses the
diamond-lattice witness when its top is realized past the
atoms\textquotesingle{} union. The corollaries
\texttt{Semiflow.qLatticeInv\_sound} /
\texttt{Behavior.qLatticeInv\_sound} (and their
\texttt{\_of\_blockInv} end-to-end forms) are
$\mathrm{sound}_A = \mathrm{sound}_T \circ \mathrm{unpack}$ through
the cited \rulename{QBlockInv} conclusion, delivered as judgment-level
cargo. Unconditional: the lattice rows are the checker\textquotesingle
s recomputed readout facts entering as theorem parameters (R1 shadows),
and join-closure membership of the down-set\textquotesingle s
realization is deliberately \emph{not} consumed --- with
down-closedness consumed and the cited member judgment supplied, the
entailment holds for the bare cell union, so the rule is sound a
fortiori with the checker\textquotesingle s extra side condition. The
payload-level unpack elaborator is bound beside the corollaries (as
with the sibling extension rules; module
\texttt{BernsteinLogic.QtslElaborators}):
\texttt{LatticeInvPayload} reifies the wire payload --- the leg-7 row
pack, the down-set, the containment row ---
\texttt{LatticeInvPayload.unpack} is the verified transport onto the
\rulename{QBlockInv} payload (block $:=$ the down-set\textquotesingle
s realization; the containment row rewritten along the exactness
identity; facets a citation, never payload data), and the
\texttt{qLatticeInv\_sound\_elab} pair is $\mathrm{sound}_A =
\mathrm{sound}_T \circ \mathrm{unpack}$ with
\texttt{qBlockInv\_sound} as the literal $\mathrm{sound}_T$, at both
semantic strata. The tier is flagged, not forced:
if registry-backed citation of the member\textquotesingle s rows is
ruled heavier than the alias criterion allows, the rule reclassifies
as admissible with all machinery unchanged.

\cataloguenote{The warrant for the reading --- once an orbit enters a lattice member
it never leaves, so visited cells descend the cell order and down-sets
are forward-invariant --- is the same orbit combinatorics the Morse
discharge\textquotesingle s order clause rests on, consumed here at
its easy end.
}
\ruleblock{Lean statement}
\declhead{BernsteinLogic.QtslLatticeInv}{qLatticeInv\_exact}{\badge{BadgeGreen}{statement bound}}
\begin{leansrc}
/-- **prop:qli-exact, the exactness identity.** For a down-closed cell-set
`C`, the union of the cells of `C` equals the union of the realizations of
`C`'s join-irreducible members — the realization of the display's member
`N_C = ⋁ C` (join = union in the union-closed family object). The ⊇ half
is leg 7's well-founded descent through lower covers
(`Semiflow.exists_cellMem_subset`): a member point lands in the cell of a
join-irreducible below it, and down-closedness (through `horder`) pins
that cell inside `C`. The ⊆ half is `hcellSub`, cell inside member. This
is what makes the alias's respelling exact rather than merely `⊇`. -/
theorem qLatticeInv_exact
    (hdc : IsLowerSet C)
    (hJoin : ∀ b : β, ∀ x ∈ blocks b,
      ∃ i : ι, x ∈ blocks (emb i) ∧ blocks (emb i) ⊆ blocks b)
    (hcover : ∀ i : ι, blocks (cover i) ⊂ blocks (emb i))
    (hcellsSup : ∀ i : ι, blocks (emb i) \ blocks (cover i) ⊆ cells i)
    (horder : ∀ i j : ι, blocks (emb i) ⊆ blocks (emb j) → i ≤ j)
    (hcellSub : ∀ i : ι, cells i ⊆ blocks (emb i)) :
    (⋃ i ∈ C, cells i) = ⋃ i ∈ C, blocks (emb i)
\end{leansrc}
\declhead{BernsteinLogic.QtslLatticeInv}{member\_subset\_cells\_le}{\badge{BadgeGreen}{statement bound}}
\begin{leansrc}
/-- **"cells tile members"** (the page proof's first readout fact), at the
principal down-set: every point of a join-irreducible's realization lies
in the cell of some join-irreducible at-or-below it in the cell order.
The identity's ⊇ half, exposed as its own statement because the book page
names it. -/
theorem member_subset_cells_le
    (hJoin : ∀ b : β, ∀ x ∈ blocks b,
      ∃ i : ι, x ∈ blocks (emb i) ∧ blocks (emb i) ⊆ blocks b)
    (hcover : ∀ i : ι, blocks (cover i) ⊂ blocks (emb i))
    (hcellsSup : ∀ i : ι, blocks (emb i) \ blocks (cover i) ⊆ cells i)
    (horder : ∀ i j : ι, blocks (emb i) ⊆ blocks (emb j) → i ≤ j)
    (i : ι) :
    blocks (emb i) ⊆ ⋃ j ∈ Set.Iic i, cells j
\end{leansrc}
\declhead{BernsteinLogic.QtslLatticeInv}{Semiflow.qLatticeInv\_sound}{\badge{BadgeGreen}{statement bound}}
\begin{leansrc}
/-- **prop:qli-sound, through the target** — the alias soundness corollary
in its genre-native shape `sound_A = sound_T ∘ unpack`: the premise is the
cited `QBlockInv` node's conclusion on the member `N_C` (its realization:
the union of the down-set's join-irreducible realizations), delivered as
judgment-level cargo (the `qBasin_sound` precedent — the facet rows live
in the cited node, not here); the conclusion is the alias's displayed
judgment on `⋃↓C`. The bridge is `qLatticeInv_exact`, and there is no
other content — the alias genre working as intended. -/
theorem Semiflow.qLatticeInv_sound {φ : Semiflow X} {X0 : Set X}
    (hdc : IsLowerSet C)
    (hJoin : ∀ b : β, ∀ x ∈ blocks b,
      ∃ i : ι, x ∈ blocks (emb i) ∧ blocks (emb i) ⊆ blocks b)
    (hcover : ∀ i : ι, blocks (cover i) ⊂ blocks (emb i))
    (hcellsSup : ∀ i : ι, blocks (emb i) \ blocks (cover i) ⊆ cells i)
    (horder : ∀ i j : ι, blocks (emb i) ⊆ blocks (emb j) → i ≤ j)
    (hcellSub : ∀ i : ι, cells i ⊆ blocks (emb i))
    (hMember : φ.InvScopeSat (⋃ i ∈ C, blocks (emb i)) X0) :
    φ.InvScopeSat (⋃ i ∈ C, cells i) X0
\end{leansrc}
\declhead{BernsteinLogic.QtslLatticeInv}{Semiflow.qLatticeInv\_sound\_of\_blockInv}{\badge{BadgeGreen}{statement bound}}
\begin{leansrc}
/-- **rec:qli-filled, composed end-to-end**: the filled-in `QBlockInv`
application. Premises: forward invariance of the member's realization
(the cited facet-row chain's output — dischargeable per member via
`qBlockInv_sound_of_facetRows` and kin) and the rule's own containment
row, stated on the alias's displayed region and transported verbatim
across the identity. -/
theorem Semiflow.qLatticeInv_sound_of_blockInv {φ : Semiflow X} {X0 : Set X}
    (hdc : IsLowerSet C)
    (hJoin : ∀ b : β, ∀ x ∈ blocks b,
      ∃ i : ι, x ∈ blocks (emb i) ∧ blocks (emb i) ⊆ blocks b)
    (hcover : ∀ i : ι, blocks (cover i) ⊂ blocks (emb i))
    (hcellsSup : ∀ i : ι, blocks (emb i) \ blocks (cover i) ⊆ cells i)
    (horder : ∀ i j : ι, blocks (emb i) ⊆ blocks (emb j) → i ≤ j)
    (hcellSub : ∀ i : ι, cells i ⊆ blocks (emb i))
    (hInv : φ.Invariant (⋃ i ∈ C, blocks (emb i)))
    (hInit : X0 ⊆ ⋃ i ∈ C, cells i) :
    φ.InvScopeSat (⋃ i ∈ C, cells i) X0
\end{leansrc}
\declhead{BernsteinLogic.QtslLatticeInv}{Behavior.qLatticeInv\_sound}{\badge{BadgeGreen}{statement bound}}
\begin{leansrc}
/-- `Semiflow.qLatticeInv_sound` at the behavioral floor — the identity is
dynamics-free, so the alias transports over solution-trajectory-set
semantics verbatim (the two-directions carve: `QInvScope` is
safety-directional, sound here without existence hypotheses). -/
theorem Behavior.qLatticeInv_sound {B : Behavior X} {X0 : Set X}
    (hdc : IsLowerSet C)
    (hJoin : ∀ b : β, ∀ x ∈ blocks b,
      ∃ i : ι, x ∈ blocks (emb i) ∧ blocks (emb i) ⊆ blocks b)
    (hcover : ∀ i : ι, blocks (cover i) ⊂ blocks (emb i))
    (hcellsSup : ∀ i : ι, blocks (emb i) \ blocks (cover i) ⊆ cells i)
    (horder : ∀ i j : ι, blocks (emb i) ⊆ blocks (emb j) → i ≤ j)
    (hcellSub : ∀ i : ι, cells i ⊆ blocks (emb i))
    (hMember : B.InvScopeSat (⋃ i ∈ C, blocks (emb i)) X0) :
    B.InvScopeSat (⋃ i ∈ C, cells i) X0
\end{leansrc}
\declhead{BernsteinLogic.QtslLatticeInv}{Behavior.qLatticeInv\_sound\_of\_blockInv}{\badge{BadgeGreen}{statement bound}}
\begin{leansrc}
/-- `Semiflow.qLatticeInv_sound_of_blockInv` at the behavioral floor. -/
theorem Behavior.qLatticeInv_sound_of_blockInv {B : Behavior X} {X0 : Set X}
    (hdc : IsLowerSet C)
    (hJoin : ∀ b : β, ∀ x ∈ blocks b,
      ∃ i : ι, x ∈ blocks (emb i) ∧ blocks (emb i) ⊆ blocks b)
    (hcover : ∀ i : ι, blocks (cover i) ⊂ blocks (emb i))
    (hcellsSup : ∀ i : ι, blocks (emb i) \ blocks (cover i) ⊆ cells i)
    (horder : ∀ i j : ι, blocks (emb i) ⊆ blocks (emb j) → i ≤ j)
    (hcellSub : ∀ i : ι, cells i ⊆ blocks (emb i))
    (hInv : B.Invariant (⋃ i ∈ C, blocks (emb i)))
    (hInit : X0 ⊆ ⋃ i ∈ C, cells i) :
    B.InvScopeSat (⋃ i ∈ C, cells i) X0
\end{leansrc}
\par\medskip
\textbf{The payload-level unpack, displayed whole.}\enspace
Recipe~\ref{rec:qli-filled} as a function: the wire payload reified
field-for-field against the register, the containment row transported
along the exactness identity, the target\textquotesingle s soundness
consumed literally in the factoring pair.

\declhead{BernsteinLogic.QtslElaborators}{LatticeInvPayload}{\badge{BadgeGreen}{statement bound}}
\begin{leansrc}
/-- **The QLatticeInv wire payload, Lean face** — fields mirror the
geolog register (section J's adjudication: "registers already staged;
the elaborator's domains are complete"): the family object with its
recomputed readout (`QLatticeInv/lattice`, as the leg-7 R1-shadow row
pack — no lattice structure minted, the leg-7 discipline), the
down-set (`QLatticeInv/downset`, with down-closedness the checker's
recomputed side condition), and the rule's own containment row
(`QLatticeInv/init_containment`). Membership of the down-set's
realization in the lattice is deliberately NOT a field — checker-side,
not consumed (soundness holds for the bare union, a fortiori with the
extra condition). -/
structure LatticeInvPayload (X ι β : Type*) [Finite ι] [PartialOrder ι] where
  /-- Member realizations (the family object's blocks). -/
  blocks : β → Set X
  /-- Join-irreducible embedding over the cell index. -/
  emb : ι → β
  /-- Lower covers. -/
  cover : ι → β
  /-- Cells `T(b) = cl(b ∖ b⁻)`, as realized sets. -/
  cells : ι → Set X
  /-- The down-set `C` (wire field `downset`). -/
  downset : Set ι
  /-- The initial region `X₀`. -/
  initRegion : Set X
  /-- Down-closedness in the recomputed cell order (R1 recompute). -/
  hdc : IsLowerSet downset
  /-- Birkhoff join-decomposition, realization form (R1 shadow). -/
  hJoin : ∀ b : β, ∀ x ∈ blocks b,
    ∃ i : ι, x ∈ blocks (emb i) ∧ blocks (emb i) ⊆ blocks b
  /-- Strict lower covers (R1 shadow). -/
  hcover : ∀ i : ι, blocks (cover i) ⊂ blocks (emb i)
  /-- Cell ⊇ member-minus-cover (R1 shadow). -/
  hcellsSup : ∀ i : ι, blocks (emb i) \ blocks (cover i) ⊆ cells i
  /-- Order reflects realization inclusion (R1 shadow). -/
  horder : ∀ i j : ι, blocks (emb i) ⊆ blocks (emb j) → i ≤ j
  /-- Cell ⊆ its member's realization (the row this alias adds). -/
  hcellSub : ∀ i : ι, cells i ⊆ blocks (emb i)
  /-- The rule's own containment row (wire field `init_containment`),
  stated on the alias's displayed region. -/
  hInit : initRegion ⊆ ⋃ i ∈ downset, cells i
\end{leansrc}
\declhead{BernsteinLogic.QtslElaborators}{BlockInvPayload}{\badge{BadgeGreen}{statement bound}}
\begin{leansrc}
/-- **The QBlockInv payload, Lean face** (`rec:qli-transport`'s
codomain): the block region and the containment row. The facet rows are
a CITATION (judgement-level cargo, the `qBasin_sound` precedent) — they
enter the factoring theorems as the cited node's invariance, never as
payload data. -/
structure BlockInvPayload (X : Type*) where
  /-- The block `N_C` — here its realization, the bare union (the
  join realizes as the union in the union-closed family object; that
  the union is a lattice MEMBER is the checker's recomputed side
  condition, deliberately not consumed). -/
  block : Set X
  /-- The initial region `X₀`, transported verbatim. -/
  initRegion : Set X
  /-- `init_containment`, rewritten along `qLatticeInv_exact`. -/
  hInit : initRegion ⊆ block
\end{leansrc}
\declhead{BernsteinLogic.QtslElaborators}{Semiflow.qLatticeInv\_sound\_elab}{\badge{BadgeGreen}{statement bound}}
\begin{leansrc}
/-- **sound_A = sound_T ∘ unpack, semiflow stratum** — the target's
soundness lemma (`qBlockInv_sound`, literally sound_T) applied to the
UNPACKED payload, concluding the alias's displayed judgement. The cited
member's facet-row chain enters as `hInv` (judgement-level cargo). -/
theorem Semiflow.qLatticeInv_sound_elab {φ : Semiflow X}
    (p : LatticeInvPayload X ι β) (hInv : φ.Invariant p.unpack.block) :
    φ.InvScopeSat (⋃ i ∈ p.downset, p.cells i) p.initRegion
\end{leansrc}
\declhead{BernsteinLogic.QtslElaborators}{Behavior.qLatticeInv\_sound\_elab}{\badge{BadgeGreen}{statement bound}}
\begin{leansrc}
/-- **sound_A = sound_T ∘ unpack, behavioural floor** — the identity is
dynamics-free, so the unpack transports at every semantics for free. -/
theorem Behavior.qLatticeInv_sound_elab {B : Behavior X}
    (p : LatticeInvPayload X ι β) (hInv : B.Invariant p.unpack.block) :
    B.InvScopeSat (⋃ i ∈ p.downset, p.cells i) p.initRegion
\end{leansrc}
\declhead{BernsteinLogic.QtslElaborators}{LatticeInvPayload.unpack}{\badge{BadgeGreen}{definition bound}}
\begin{leansrc}
/-- **The verified unpack** — `rec:qli-filled` as a function, the first
of the five alias artefacts (`unpack : Step_A → Step_T`): block :=
the down-set's realization `N_C`, facets cited (not payload),
`init_containment` transported along the exactness identity
`qLatticeInv_exact` (`⋃↓C = N_C` as regions — the alias's conclusion is
the target's on the nose). -/
def LatticeInvPayload.unpack (p : LatticeInvPayload X ι β) : BlockInvPayload X where
  block := ⋃ i ∈ p.downset, p.blocks (p.emb i)
  initRegion := p.initRegion
  hInit := by
    rw [← qLatticeInv_exact p.hdc p.hJoin p.hcover p.hcellsSup p.horder p.hcellSub]
    exact p.hInit
\end{leansrc}
\ruleblock{Correspondence (Steele display $\leftrightarrow$ Lean)}
\badge{BadgeGreen}{reviewer-certified against the displayed texts}
\sourcenote{The dictionary below is a reviewer-certified statement of how the mechanized statements above interpret the displayed rule: display metavariable $\mapsto$ Lean binder, each display hypothesis $\mapsto$ the premise or side condition carrying it, the display conclusion $\mapsto$ the goal --- with any deliberate difference in width between display and mechanization named outright. Certification is recorded as a hash triple over (the displayed cell, the extracted statement texts, this dictionary\textquotesingle s own text); if any of the three changes, this block turns amber until a reviewer re-certifies. The correspondence claim itself remains editorial --- reviewer-grade, not machine-checked; its machine-checked successor types each soundness statement against the rule as a structured object, and supersedes this block per rule as it lands.}
The displayed judgment $\mathcal{M} \vDash
\mathsf{QInvScope}(\bigcup\downarrow C, X_0)$ $\mapsto$
\texttt{InvScopeSat (⋃ i ∈ C, cells i) X0}, at BOTH semantic strata
(\texttt{Semiflow.}/\texttt{Behavior.qLatticeInv\_sound} --- the
exactness identity is dynamics-free, so the behavioural transport is
verbatim). The alias\textquotesingle s entire content is
\texttt{qLatticeInv\_exact}: for down-closed \texttt{C} (\texttt{hdc
: IsLowerSet C} --- the R1 down-closedness recompute as a binder),
the union of $C$\textquotesingle s cells EQUALS the union of its
join-irreducible realizations --- the member $N_C$\textquotesingle s
realization exactly where the checker\textquotesingle s membership
recompute holds (join realizes as union; the bound statement is on the
bare union and consumes no distributivity) --- so the respelling is
exact, not merely one-sided.
\begin{itemize}
\item the E2 lattice $L$ $\mapsto$ the Birkhoff row pack, abstract:
  \texttt{hJoin} (every block point lands in a join-irreducible below
  it), \texttt{hcover} (strict lower covers), \texttt{hcellsSup}
  (cell $\supseteq$ member minus cover), \texttt{horder} (order
  reflects realization inclusion) --- byte-identical to the Morse
  discharge\textquotesingle s eventual-member interface: no lattice
  structure is minted, the rows ARE the interface --- plus the one row
  new to this alias,
  \texttt{hcellSub} (cell $\subseteq$ member: the readout closure
  half\textquotesingle s set-level shadow).
\item the cited \rulename{QBlockInv} premise $\mapsto$
  \texttt{hMember : InvScopeSat (⋃ i ∈ C, blocks (emb i)) X0} ---
  judgment-level cargo on the member\textquotesingle s realization
  (the \texttt{qBasin\_sound} precedent: facet rows live in the cited
  node, never here); \texttt{\_of\_blockInv} composes the target in
  end-to-end (its \texttt{hInv}/\texttt{hInit} pair is
  \texttt{qBlockInv\_sound}\textquotesingle s kit, the containment
  row transported across the identity).
\item \texttt{init\_containment} $\mapsto$ \texttt{hInit : X0 ⊆ ⋃ i ∈
  C, cells i}, stated on the alias\textquotesingle s displayed region.
\item \texttt{member\_subset\_cells\_le} exposes the
  page\textquotesingle s ``cells tile members'' fact at the principal
  down-set --- bound because the proof section names it.
\end{itemize}
Width note: the module is topology-free and dynamics-free except
through the cited judgment --- \texttt{cells}/\texttt{blocks}
arbitrary set families over an abstract ordered index; the grid
realization of the lattice is R1 territory, strictly wider in the
safe direction. Module row bound to an attested build (the receipt
provenance is recorded in the bindings).

\ruleblock{Lean verifier receipts}
\sourcenote{module \texttt{BernsteinLogic.QtslLatticeInv}; bound build \texttt{467dc57e1c30\ldots} (bound from receipt 98755f4d8dbefb10, 2026-07-21T04:48:40Z). A receipt records that named checking layers accepted a build: (module, compiled-artifact SHA-256, environment SHA-256, checker layers), minted when the checking gate ran; the store receipts attesting the bound build are listed with the layers each run credited. Nothing here asks the reader to trust the store: every layer is an independently runnable program over the quoted sources --- rebuild the module and re-run the checkers to reproduce the verdict. Receipts attesting superseded builds credit no layer and are omitted from this rendering (the count of omissions, if any, is noted below).}
\begin{receipttable}
\receiptrow{98755f4d8dbefb10}{2026-07-21T04:48:40Z}{build+decorrelated}{\badge{BadgeGreen}{current}}
\receiptrow{210205ac3ea031f4}{2026-07-21T05:17:21Z}{build+decorrelated}{\badge{BadgeGreen}{current}}
\receiptrow{9a870584fbb58e05}{2026-07-21T05:28:41Z}{build+replay+hardened+decorrelated}{\badge{BadgeGreen}{current}}
\end{receipttable}
\sourcenote{3 further receipt(s) in the store attest superseded builds of this module and are omitted from this rendering.}
\sourcenote{module \texttt{BernsteinLogic.QtslElaborators}; bound build \texttt{0205ec213b91\ldots} (hashed from the built olean at bind time). A receipt records that named checking layers accepted a build: (module, compiled-artifact SHA-256, environment SHA-256, checker layers), minted when the checking gate ran; the store receipts attesting the bound build are listed with the layers each run credited. Nothing here asks the reader to trust the store: every layer is an independently runnable program over the quoted sources --- rebuild the module and re-run the checkers to reproduce the verdict. Receipts attesting superseded builds credit no layer and are omitted from this rendering (the count of omissions, if any, is noted below).}
\begin{receipttable}
\receiptrow{a237f4cff27cb15a}{2026-07-22T10:06:47Z}{build+decorrelated}{\badge{BadgeGreen}{current}}
\end{receipttable}
\layerline{\layeritem{build}{lake build (elaboration)}{\badge{BadgeGreen}{current}}\layeritem{replay}{leanchecker (official kernel replay)}{\badge{BadgeGreen}{current}}\layeritem{hardened}{safe\_verify (hardened screen)}{\badge{BadgeGreen}{current}}\layeritem{decorrelated}{lean4export + nanoda (decorrelated kernel)}{\badge{BadgeGreen}{current}}}
\sourcenote{\badge{BadgeGray}{comparator: none} no fixture in the three-implementation comparator harness (compiled / kernel-interpreted / reference-interpreted) exercises this rule.}
\ruleblock{HOL4 statement}
\absentbox{NO HOL STATEMENT YET}{Absent at the judgment stratum: no judgment-stratum definition exists
in the ported base-kernel checker scripts (verified against the
port\textquotesingle s definition inventory), so the satisfaction-level
witnesses --- now mechanized on the Lean side (module
\texttt{QtslLatticeInv}, the entry\textquotesingle s bound statements)
--- have no HOL twin to mirror them. The target rule\textquotesingle s
rows (facet signs, containment) ARE covered by the ported checker.
}
\ruleblock{HOL verifier receipts}
\absentbox{NONE}{no HOL verification receipt exists for this rule.}
\ruleblock{Proof}

\sourcenote{Genre: alias rule (tier 3) --- a respelling template. An
alias owes the book five artefacts: the verified 1$\to$1 unpack
function, the filled-in application of the target rule, the
payload-transport table, the exact algebraic identity that makes the
respelling correct, and the soundness corollary through the target
rule's lemma. All five are present --- the filled-in application, the
transport table, the exactness identity and the soundness corollary
mechanized and bound in the entry
(\texttt{qLatticeInv\_exact}, Proposition~\ref{prop:qli-exact}; the
\texttt{qLatticeInv\_sound} family, Proposition~\ref{prop:qli-sound},
at the semiflow \emph{and} behavioural strata), and --- now that the
extension certificate schema exists --- the verified payload-level
unpack function itself, \texttt{LatticeInvPayload.unpack} with its
factoring pair (the remark below names the bound functions). (Tier
flag, carried honestly: the
unpack cites facet rows through the checker's recomputed lattice
readout --- if registry-backed citation is ruled heavier than the alias
criterion allows, the rule reclassifies as admissible with all
machinery unchanged.)}

\begin{bremark}[the payload-level elaborator, mechanized]
\label{rem:qli-elab}
The \emph{semantic} content of this page is bound in the entry:
\texttt{qLatticeInv\_exact} is Proposition~\ref{prop:qli-exact}
(over the discharge chain's own Birkhoff rows --- no lattice structure
is minted), and the \texttt{qLatticeInv\_sound}
family is Proposition~\ref{prop:qli-sound} in both its
through-the-target and end-to-end forms, at both semantic strata. The
payload-level elaborator --- previously this page's one absent
artefact, held back until the extension certificate schema existed ---
is bound with it, as with the sibling extension rules:
\texttt{LatticeInvPayload} (module
\texttt{BernsteinLogic.QtslElaborators}) reifies the rule's wire
payload --- family object as the leg-7 row pack, down-set, containment
row --- and \texttt{LatticeInvPayload.unpack} takes it to the single
\rulename{QBlockInv} payload of Recipe~\ref{rec:qli-filled}
(\texttt{BlockInvPayload}: block $:=$ the down-set's realization, the
containment row transported along Proposition~\ref{prop:qli-exact};
facets are a citation, entering the factoring theorems as the cited
node's invariance, never as payload data). The factoring pair
\texttt{Semiflow.qLatticeInv\_sound\_elab} /
\texttt{Behavior.qLatticeInv\_sound\_elab} is
$\mathrm{sound}_A = \mathrm{sound}_T \circ \mathrm{unpack}$ with
\texttt{qBlockInv\_sound} as the literal $\mathrm{sound}_T$, applied to
the unpacked payload. The checker-side ingredients (membership of the
down-set's realization in the lattice, down-closedness of the cell-set)
are recomputed side conditions, already specified with the rule; the
mechanization deliberately consumes the down-closedness and leaves the
membership fact checker-side (with down-closedness consumed and the
cited member judgment supplied, soundness holds for the bare cell
union, a fortiori with the extra condition).
\end{bremark}

\begin{brecipe}[the alias template]\label{rec:qli-template}
An alias rule $A$ is warranted by a verified function
$\mathrm{unpack} : \mathrm{Step}_A \to \mathrm{Step}_T$ into the
payload of a \emph{single} target node $T$, such that the $A$-node and
its image conclude the same judgment from the same premises: one node
$\to$ one node, payload transported, nothing derived. Soundness is then
a corollary, $\mathrm{sound}_A = \mathrm{sound}_T \circ
\mathrm{unpack}$.
\end{brecipe}

\textbf{The rule.}\enspace The certificate supplies the family object
(the block lattice of the \rulename{QMorseDischarge} section, with its
recomputed readout: join-irreducible members $b$, lower covers $b^-$,
cells $T(b) := \mathrm{cl}(b \smallsetminus b^-)$, the induced cell
order), a cell-set $C$ of join-irreducibles that is \emph{down-closed}
in the cell order, and one containment row
$X_0 \subseteq \bigcup\downarrow C$. It concludes
$\mathsf{QInvScope}\big(\bigcup\downarrow C,\ X_0\big)$, writing
$\bigcup\downarrow C$ for the union of the cells of $C$.

\begin{brecipe}[the filled-in application --- \rulename{QLatticeInv}
rewritten \emph{into} the \rulename{QBlockInv} instance]
\label{rec:qli-filled}
Let $N_C := \bigvee_{b \in C} b$ --- the lattice join of the down-set's
members, itself a member of the family object by join-closure. The
\rulename{QLatticeInv} node rewrites to the single \rulename{QBlockInv}
application
\begin{mathpar}
\inferrule*[Right=\rulename{QBlockInv}]{
  N_C \ \text{(the member: block data from the family object)} \\
  \bigl\{\vdash \text{facet-inward sign rows of } N_C\bigr\}
  \ \text{(cited from the family object, byte-identical)} \\
  \vdash X_0 \subseteq N_C \ \text{(the rule's own containment row,
  rewritten along Proposition~\ref{prop:qli-exact})}
}{
  \mathcal{M} \vDash \mathsf{QInvScope}(N_C,\ X_0)
}
\end{mathpar}
and the conclusion is the alias's conclusion \emph{on the nose} by
Proposition~\ref{prop:qli-exact}: $N_C = \bigcup\downarrow C$ as
regions. Every internal facet of $N_C$ already carries its sign row in
the family object --- membership of $N_C$ in the lattice is exactly
the checker's recomputed side condition, and no row is re-proven.
\end{brecipe}

\begin{bprop}[the exactness identity]\label{prop:qli-exact}
For a down-closed set $C$ of join-irreducibles,
\[
\bigcup_{b \in C} T(b) \;=\; \bigvee_{b \in C} b
\qquad \text{(equality of realizations).}
\]
\end{bprop}
\begin{proof}
Two readout facts, one lattice fact. \emph{(cells tile members)} Every
member $N$ of the lattice is the union of the cells of the
join-irreducibles below it, $N = \bigcup \{T(b) : b \in \mathsf{J},\
b \subseteq N\}$: downward induction along lower covers --- $T(b)$ is
by definition what $b$ adds over $b^-$, so peeling cells descends
through the lattice and exhausts $N$. \emph{(Birkhoff exactness)} In a
finite \emph{distributive} lattice (a lattice of regions closed under
union and intersection --- a ring of sets --- qualifies) the map
$N \mapsto \{b \in \mathsf{J} : b
\subseteq N\}$ is an isomorphism onto down-closed sets of
join-irreducibles \cite{birkhoff1937}; for down-closed
$C$ the join $N_C$ then satisfies $\{b \in \mathsf{J} : b \subseteq N_C\} =
C$. Distributivity --- join-primality of the irreducibles --- is what
forbids the join from absorbing new join-irreducibles, and it is not
decoration: in the diamond lattice $M_3$, two of the three atoms form
a down-closed set of join-irreducibles whose join absorbs the third. Substituting the second fact into the
first: $N_C = \bigcup_{b \in C} T(b)$. (Where each half lives: the
tiling is the readout's own construction, recomputed by the checker
per certificate; the isomorphism is the one classical fact, and
down-closedness of $C$ is the recomputed side condition that makes the
identity exact rather than merely $\supseteq$.) \emph{Mechanization
note, the honest disanalogy, adjudicated against the descent:}
\texttt{qLatticeInv\_exact} binds the identity in \emph{union form} ---
$\bigcup_{b \in C} T(b) = \bigcup_{b \in C} b$ as realizations, the
right side the bare union over the down-set --- and proves it from the
row pack alone: \emph{no distributivity hypothesis is consumed}, and
none is encoded, since the rows are satisfiable in absorption-positive
families. The route never forms the join: each point of the union
arrives with a witnessing member $b \in C$; the well-founded descent
through lower covers (the discharge chain's
\texttt{exists\_cellMem\_subset}) runs \emph{inside that one member},
landing the point in the cell of a join-irreducible $b' \subseteq b$;
and the pin-step is two rows --- the induced-order binding
\texttt{horder} turns $b' \subseteq b$ into $b' \leq b$ in the cell
order, and down-closedness closes $b' \in C$. Join-primality forbids
absorption below the \emph{join}; the descent never visits the join,
so the prohibition is sidestepped, not re-proven. What distributivity
carries is therefore only the display seam --- reading the bare union
as the realization of the member $N_C = \bigvee C$ --- and the seam
consumes less than the full hypothesis: it needs exactly
join-realizes-as-union, which is the checker's recomputed membership
of the down-set's realization (a member that is a union of members is
their least upper bound), a consequence of the ring-of-sets structure
strictly weaker than distributivity (union-closed absorption-positive
families exist: there the identity holds in \emph{both} forms while
the classifying map still fails). The $M_3$ witness above, realized as
sets --- atom realizations $\{1\}, \{2\}, \{3\}$, top $\{1,2,3\}$,
$C$ the first two atoms --- separates the strata cleanly: every row
holds, the union form holds ($\{1,2\} = \{1,2\}$), the join form
fails ($\{1,2\} \neq \{1,2,3\}$), and it is the membership recompute
that refuses. Nothing mechanized rests on the classical citation; it
remains as warrant genealogy for this proof's route.
\end{proof}

\begin{brecipe}[payload transport --- field for field]
\label{rec:qli-transport}
\begin{center}\small
\begin{tabular}{@{}ll@{}}
\textbf{\rulename{QLatticeInv} payload} & \textbf{\rulename{QBlockInv} payload}\\
\midrule
family object $+$ down-set $C$ & block $N_C = \bigvee C$ (member, by join-closure)\\
\emph{(no own rows)} & \texttt{facets} $:=$ $N_C$'s rows, cited from the family object\\
\texttt{init\_containment} & \texttt{init\_containment} (rewritten along Proposition~\ref{prop:qli-exact}: $\bigcup\downarrow C = N_C$)\\
conclusion region $\bigcup\downarrow C$ & conclusion region $N_C$ (Proposition~\ref{prop:qli-exact})\\
\end{tabular}
\end{center}
The checker's recomputed side conditions at the alias: $C$ down-closed
in the recomputed cell order, and $N_C$ a member whose facet rows are
all present --- both facts about the readout, not about the dynamics.
\end{brecipe}

\begin{bprop}[soundness corollary, through the target]
\label{prop:qli-sound}
If the family object's rows are valid and
$X_0 \subseteq \bigcup\downarrow C$, then --- under the
\rulename{QBlockInv} section's standing hypotheses (its conditional
ledger, applied to the member $N_C$) ---
$\mathsf{QInvScope}\big(\bigcup\downarrow C,\ X_0\big)$ holds.
\end{bprop}
\begin{proof}
The member $N_C$ is a strictly facet-inward compact block of the
family object, so the \rulename{QBlockInv} section's theorem chain
applies to it verbatim (its rows \emph{are} rows of the family
object); the conclusion transports along
Proposition~\ref{prop:qli-exact}. Mechanized, this proposition is one
composition, $\mathrm{sound}_{\textsc{QLatticeInv}} =
\mathrm{sound}_{\textsc{QBlockInv}} \circ \mathrm{unpack}$ --- there
is no other content, which is the alias genre working as intended: the
through-the-target form (\texttt{Semiflow.qLatticeInv\_sound} /
\texttt{Behavior.qLatticeInv\_sound}) takes the cited node's
$\mathsf{QInvScope}$ conclusion as judgment-level cargo and rewrites
along the identity; the end-to-end form (\texttt{\_of\_blockInv})
composes through \texttt{qBlockInv\_sound} verbatim. Both strata are
covered because the identity is dynamics-free --- the alias transports
at every semantics for free.
\end{proof}

\begin{bremark}[why the rule earns its keep]
Down-closed cell-sets are the invariant enclosures the Morse readout
offers for free: once an orbit enters a member it never leaves, so the
sequence of cells it visits is non-increasing in the cell order, and
the down-set of any initial cell-set is forward-invariant and encloses
the reachable set. The alias turns each such enclosure into a one-node
citation --- the member's rows already exist in the family object ---
where a direct proof would re-derive facet rows the certificate
already carries. Downstream, the minted $\mathsf{QInvScope}$ judgment
feeds \rulename{QBasin}'s safety factory and scope-restriction of
reach obligations, which is where the row-count payoff compounds.
\end{bremark}

\begin{anchors}
target soundness & \texttt{qBlockInv\_sound} (QtslSoundness.lean) and the \rulename{QBlockInv} section's discharge chain\\
Proposition \ref{prop:qli-exact} & \texttt{qLatticeInv\_exact} (QtslLatticeInv.lean, bound in the entry; union form, no distributivity consumed) via the \texttt{exists\_cellMem\_subset} descent; Birkhoff representation \cite{birkhoff1937} as warrant genealogy (nothing mechanized rests on it)\\
``cells tile members'' & \texttt{member\_subset\_cells\_le} (principal down-set instance)\\
alias-level corollary & \texttt{Semiflow.qLatticeInv\_sound} $+$ \texttt{\_of\_blockInv} and the behavioural-floor pair (bound in the entry)\\
payload-level elaborator (\autoref{rem:qli-elab}) & \texttt{LatticeInvPayload.unpack} + \texttt{Semiflow.qLatticeInv\_sound\_elab} + \texttt{Behavior.qLatticeInv\_sound\_elab} (module \texttt{BernsteinLogic.QtslElaborators}, bound in the entry)\\
side conditions & down-closedness: consumed (\texttt{IsLowerSet} $+$ \texttt{horder}); membership of $N_C$ in the lattice: checker-side, deliberately not consumed\\
\end{anchors}


\rulesection{74}{QBasin}{core rule}
\label{blcatrule-QBasin}
\ruleblock{The rule (Steele metanotation)}
\sourcenote{as realized in the mechanized development; source row: \texttt{\texttt{QTSL-rules.toml}} --- \emph{QBasin}. The metanotation statement is typeset by hand for notational quality; its mechanized ground truth is the Lean statement below, extracted from the sources at every build of this book and hash-bound. The rule \emph{roster} is machine-checked (every mechanized rule tag round-trips to exactly one catalogue row); the statement \emph{text} is editorial, checked by eye against the extracted Lean beside it.}
\begin{mathpar}
\inferrule*{
  \mathcal{M} \vDash \mathsf{QInvScope}(S, X_0) \\
  \bigl\{\vdash S \cap [\![\beta_j]\!] = \emptyset\bigr\}_{j} \ \text{(\texttt{avoid\_disjoint} rows)}
}{
  \mathcal{M},\ x_0 \in X_0 \vDash \mathsf{SProbGE}\bigl(1,\ \mathsf{PAlways}(\mathsf{SNot}\,\beta_j)\bigr)
}
\end{mathpar}

\ruleflavor{the invariance premise is deliberately the \\textsf{QInvScope} judgment and not the Morse-graph judgment: basin avoidance consumes invariance of $S$ from $X_0$ alone, and asking the stronger judgment would couple the rule to machinery its soundness never uses.}
\par\medskip
\rulename{QBasin} is the campaign rule: safety atoms by citation. From
one invariance-scope judgment and one disjointness fact per avoid
region, it concludes the safety-shaped region judgment
$\mathsf{SProbGE}\bigl(1,\ \mathsf{PAlways}(\mathsf{SNot}\,\beta)\bigr)$
(surface syntax $\mathsf{P}_{\ge 1}[\mathsf{G}\,\neg\beta]$) on the
scope\textquotesingle s init region --- the deterministic reading, with
no probabilistic machinery: for the deterministic model class the
certified claim is universal avoidance, and no supermartingale
certificate is synthesized or needed. The economics are the point: the
invariance judgment is proved once (typically by \rulename{QBlockInv}
or \rulename{QLatticeInv}, whose facet rows already live in the family
object), and each additional avoid region costs exactly its own
disjointness rows and nothing else --- the invariance side is cited,
never re-certified, even when the regions are taken jointly (the batch
theorem).

What the checker recomputes: the conclusion tie (the cited
judgment\textquotesingle s region and init fields equal the
rule\textquotesingle s own), plus the per-box disjointness sign rows,
or the grid bookkeeping when the avoid region is cell-aligned. What
would break without each hypothesis: without invariance, a trajectory
exits the scope and reaches the avoid set (the proof
section\textquotesingle s drift counterexample); without the
containment, failure at time zero; without disjointness, failure at
time zero inside the scope.

The premise is deliberately the invariance judgment and NOT the global
Morse judgment: recurrence structure alone implies no enclosure of the
init region --- an orbit may transit the avoid set en route to its
Morse set. Basin-minted and battery-minted safety atoms are
type-identical conclusions, so downstream citations consume either
producer indifferently; keying on provenance is forbidden by design.

This is the one Morse-family rule whose mechanization is
unconditional: the cited theorems are proved outright, with no
interface hypotheses.

\cataloguenote{The probability-one costume ($\mathsf{SProbGE}(1, \cdot)$) matters for
composition: the
conclusion lands in the same judgment form the probabilistic batteries
mint, so tower citations, boolean combination at the root, and the
assume--guarantee spine consume deterministic and stochastic safety
atoms through one interface.
}
\ruleblock{Lean statement}
\declhead{BernsteinLogic.QtslSoundness}{qBasin\_sound}{\badge{BadgeGreen}{statement bound}}
\begin{leansrc}
/-- QBasin soundness (roster: "3 lines against the Morse note's lemma
corpus"; axioms note §3.5) — FULLY PROVED, no interface facts. Premises:
the cited `QInvScope` judgment (the E1 edge `QBasin/inv`) and the
avoid-disjointness fact (`QBasin/avoid_disjoint`'s rows). Conclusion:
the safety-shaped `QRegionSat` on the scope's init region. Note the
invariance premise arrives as judgment-level cargo — the certificate-
level fact `QInvScope` carries — which is why no dynamics beyond the
semiflow algebra is consumed here (the roster's §4.6 correction: the
premise is NOT the `QMorse` judgment). -/
theorem qBasin_sound {φ : Semiflow X} {S X0 B : Set X}
    (hscope : φ.InvScopeSat S X0) (hAvoid : S ∩ B = ∅) :
    φ.SafetySat X0 B
\end{leansrc}
\declhead{BernsteinLogic.QtslSoundness}{qBasin\_batch}{\badge{BadgeGreen}{statement bound}}
\begin{leansrc}
/-- The batch-factory payoff, mechanized ("one family object, many
G-judgments, no new rows" — zero new rows per additional `β` after the
first containment): one invariant scope
discharges every avoid-region at the price of its own disjointness fact
alone — even jointly. -/
theorem qBasin_batch {φ : Semiflow X} {S X0 : Set X} {J : Type*}
    {B : J → Set X} (hscope : φ.InvScopeSat S X0)
    (hAvoid : ∀ j, S ∩ B j = ∅) :
    φ.SafetySat X0 (⋃ j, B j)
\end{leansrc}
\ruleblock{Correspondence (Steele display $\leftrightarrow$ Lean)}
\badge{BadgeAmber}{a bound Lean statement CHANGED since certification --- re-review and re-certify}
\sourcenote{The dictionary below is a reviewer-certified statement of how the mechanized statements above interpret the displayed rule: display metavariable $\mapsto$ Lean binder, each display hypothesis $\mapsto$ the premise or side condition carrying it, the display conclusion $\mapsto$ the goal --- with any deliberate difference in width between display and mechanization named outright. Certification is recorded as a hash triple over (the displayed cell, the extracted statement texts, this dictionary\textquotesingle s own text); if any of the three changes, this block turns amber until a reviewer re-certifies. The correspondence claim itself remains editorial --- reviewer-grade, not machine-checked; its machine-checked successor types each soundness statement against the rule as a structured object, and supersedes this block per rule as it lands.}
Both statements are read against the abstract model class: the
display\textquotesingle s model $\mathcal{M}$ is any forward semiflow
$\varphi$ (\texttt{Semiflow X}) on a bare state type; the v1 model
class realizes it, and nothing in either statement consumes the
realization.
\begin{itemize}
\item $S$, $X_0$ (the scope\textquotesingle s region and init region)
  $\mapsto$ the binders \texttt{S X0 : Set X}; the avoid
  predicate\textquotesingle s extent $[\![\beta_j]\!]$ $\mapsto$
  \texttt{B : Set X} in \texttt{qBasin\_sound}, and the indexed family
  \texttt{B : J $\to$ Set X} in \texttt{qBasin\_batch} (the
  display\textquotesingle s $j$-indexed premise row, taken jointly).
\item premise $\mathrm{inv} : \mathsf{QInvScope}(S, X_0)$ $\mapsto$
  \texttt{hscope} $: \varphi.\texttt{InvScopeSat}\ S\ X_0$ ---
  definitionally the conjunction of forward invariance of $S$ and
  $X_0 \subseteq S$, the judgment\textquotesingle s two clauses.
\item each \texttt{avoid\_disjoint} row $\vdash S \cap [\![\beta_j]\!]
  = \emptyset$ $\mapsto$ \texttt{hAvoid} $: S \cap B = \emptyset$
  (batch: $\forall j,\ S \cap B\,j = \emptyset$).
\item conclusion $\mathcal{M},\ x_0 \in X_0 \vDash
  \mathsf{SProbGE}\bigl(1,\ \mathsf{PAlways}(\mathsf{SNot}\,\beta_j)\bigr)$
  $\mapsto$ the goal
  $\varphi.\texttt{SafetySat}\ X_0\ B$: every trajectory from $X_0$
  avoids $B$ at every forward time --- the deterministic
  probability-one reading of the displayed judgment.
  \texttt{qBasin\_batch} concludes
  $\varphi.\texttt{SafetySat}\ X_0\ (\bigcup_j B\,j)$, which yields the
  displayed per-$j$ family of conclusions by monotonicity.
\end{itemize}
Width note: the display\textquotesingle s E1 annotation, glossed by
the flavor text (the premise is
the invariance judgment, \emph{not} the Morse judgment), is enforced by
type rather than by hypothesis --- \texttt{qBasin\_sound} consumes an
\texttt{InvScopeSat} fact and nothing of Morse shape fits that premise
slot. No hypothesis of either statement is unused, and neither
statement is wider than the displayed rule.

\ruleblock{Lean verifier receipts}
\sourcenote{module \texttt{BernsteinLogic.QtslSoundness}; bound build \texttt{77005991d61d\ldots} (bound from receipt 8130a7a7a679cef4, 2026-08-04T01:16:06Z). A receipt records that named checking layers accepted a build: (module, compiled-artifact SHA-256, environment SHA-256, checker layers), minted when the checking gate ran; the store receipts attesting the bound build are listed with the layers each run credited. Nothing here asks the reader to trust the store: every layer is an independently runnable program over the quoted sources --- rebuild the module and re-run the checkers to reproduce the verdict. Receipts attesting superseded builds credit no layer and are omitted from this rendering (the count of omissions, if any, is noted below).}
\begin{receipttable}
\receiptrow{8130a7a7a679cef4}{2026-08-04T01:16:06Z}{build+decorrelated}{\badge{BadgeGreen}{current}}
\receiptrow{9eb3706354082144}{2026-08-04T02:23:12Z}{build+decorrelated}{\badge{BadgeGreen}{current}}
\end{receipttable}
\sourcenote{5 further receipt(s) in the store attest superseded builds of this module and are omitted from this rendering.}
\layerline{\layeritem{build}{lake build (elaboration)}{\badge{BadgeGreen}{current}}\layeritem{replay}{leanchecker (official kernel replay)}{\badge{BadgeAmber}{stale only}}\layeritem{hardened}{safe\_verify (hardened screen)}{\badge{BadgeAmber}{stale only}}\layeritem{decorrelated}{lean4export + nanoda (decorrelated kernel)}{\badge{BadgeGreen}{current}}}
\sourcenote{\badge{BadgeGray}{comparator: none} no fixture in the three-implementation comparator harness (compiled / kernel-interpreted / reference-interpreted) exercises this rule.}
\ruleblock{HOL4 statement}
\absentbox{NO HOL STATEMENT YET}{Absent: the HOL4 development\textquotesingle s object is the
base-kernel checker cone, and the satisfaction-judgment stratum this
rule concludes in has no HOL twin --- no judgment-stratum definition
exists in the ported scripts (verified against the
port\textquotesingle s definition inventory). The disjointness rows
this rule consumes ARE covered by the ported checker; the three-line
judgment-stratum soundness lemma is Lean-only today.
}
\ruleblock{HOL verifier receipts}
\absentbox{NONE}{no HOL verification receipt exists for this rule.}
\ruleblock{Proof}

\sourcenote{Genre: core rule (tier 1) --- an ordinary mathematical proof at
Putnam-exam-solution detail. Every hypothesis of the rule is used below
and remarked on; the formal-anchor table maps each step to the Lean
declaration that carries it. This rule concludes a judgment one stratum
above the sequent calculus: premises typed $\vdash$ are ordinary sequents
the base kernel replays, and the rule's own soundness lives at the
judgment stratum, over a dynamical model. Uniquely among the Morse-family
rules in this book, its mechanization is \emph{unconditional}: the cited
theorems are proved outright, with no interface hypotheses.}

\textbf{Setting and claim.}\enspace Fix a \emph{forward semiflow} $\varphi$
on a state space $X$: a map $\Phi : \R_{\geq 0} \times X \to X$, written
$\Phi_t$, with $\Phi_0 = \mathrm{id}$ and $\Phi_{s+t} = \Phi_s \fatsemi
\Phi_t$ (the abstract shadow of the model class: a deterministic
polynomial-drift flow on its chart domains; the two semiflow laws are
model-frame data --- the proofs below consume only the rule premises
(a)--(c)). Regions denote subsets of
$X$. Two judgment forms participate:

\emph{(invariance scope)}\enspace $\mathsf{QInvScope}(S, X_0)$ asserts the
conjunction: $S$ is forward-invariant ($x \in S \Rightarrow \Phi_t x \in S$
for all $t \geq 0$) and $X_0 \subseteq S$.

\emph{(safety)}\enspace the region judgment
$\mathcal{M},\ x_0 \in X_0 \vDash
\mathsf{SProbGE}\bigl(1,\ \mathsf{PAlways}(\mathsf{SNot}\,\beta)\bigr)$
asserts, in the deterministic reading, that every point of $X_0$ avoids
$B := \llbracket\beta\rrbracket$ at every forward time:
$\forall x \in X_0,\ \forall t \geq 0,\ \Phi_t x \notin B$. The
probability-one former $\mathsf{SProbGE}(1, \cdot)$
is the degenerate probabilistic costume of this
universal statement --- no measure, no stochastic machinery, and no
supermartingale certificate participates anywhere in this section (the
model-class fence: deterministic flows only).

The rule's premises: the cited invariance judgment (an internal edge ---
the claim is consumed, not its derivation), and one disjointness fact
$S \cap \llbracket\beta_j\rrbracket = \emptyset$ per avoid region,
discharged by per-box sign rows of the base calculus (or by cell
bookkeeping the checker recomputes, when $\beta_j$ is aligned to the
certificate's grid).

\begin{bthm}[soundness of \rulename{QBasin}]\label{thm:qbasin-sound}
Let $S, X_0, B \subseteq X$ with
\emph{(a)} $S$ forward-invariant,
\emph{(b)} $X_0 \subseteq S$
(jointly: the cited $\mathsf{QInvScope}(S, X_0)$ judgment), and
\emph{(c)} $S \cap B = \emptyset$.
Then every $x \in X_0$ satisfies $\Phi_t x \notin B$ for all
$t \geq 0$.
\end{bthm}
\begin{proof}
Three sentences. By (b), $x \in S$; by (a), $\Phi_t x \in S$ for every
$t$. Were $\Phi_t x \in B$, it would lie in $S \cap B$, which (c)
declares empty. Hence $\Phi_t x \notin B$.
\end{proof}

\begin{bprop}[the batch form]\label{prop:qbasin-batch}
Let $\{B_j\}_{j \in J}$ be any family of avoid regions with, for each
$j$, the single fact $S \cap B_j = \emptyset$. Under (a) and (b) of
Theorem~\ref{thm:qbasin-sound}, every $x \in X_0$ avoids
$\bigcup_j B_j$ at every forward time --- equivalently, each $B_j$
separately, and all of them jointly.
\end{bprop}
\begin{proof}
A point of the union lies in some $B_j$; apply
Theorem~\ref{thm:qbasin-sound} to that $j$. Nothing about $J$ is used ---
the index family may be infinite, though a certificate of course carries
finitely many rows.
\end{proof}

\begin{bremark}[each hypothesis is load-bearing]\label{rem:qbasin-hyp}
Drop (a), keeping (b) and (c): on $X = \R$ with the unit-speed drift
$\Phi_t x = x + t$, take $S = X_0 = [0,1]$ and $B = [2,3]$. Then
$X_0 \subseteq S$ and $S \cap B = \emptyset$, yet the trajectory from
$x = 0$ enters $B$ at $t = 2$ --- invariance is the entire dynamical
content of the rule, and it arrives wholesale through the cited
judgment. Drop (b): a point $x \in X_0 \smallsetminus S$ may sit in $B$
already, violating the conclusion at $t = 0$ --- concretely, the
identity flow with $S = \{0\}$ and $X_0 = B = \{1\}$: (a) and (c) hold,
and the start state is the violation. Drop (c): a point of
$S \cap B$ reachable from $X_0$ (at worst a point of
$X_0 \cap B$ itself, again at $t = 0$) violates it --- concretely, the
identity flow with $X_0 = S = B = \{0\}$: (a) and (b) hold, and the
start state sits in the avoid set. The two set-theoretic
hypotheses fail at time zero when dropped; the dynamical one fails at
positive time --- which is why (a) is a judgment with its own producing
rules while (b) and (c) are single containment/disjointness rows.
\end{bremark}

\begin{bremark}[the premise is the invariance judgment, not the Morse
judgment]\label{rem:qbasin-premise}
A natural mis-design would take the global Morse judgment as the
premise: it is the stronger-looking claim. But
$\mathsf{Morse}(G, \rho)$ alone implies no particular invariant
enclosure of $X_0$ --- an orbit from $X_0$ may transit
$\llbracket\beta\rrbracket$ \emph{en route} to its Morse set
(concretely: $\dot{x} = 3 - x$ on $\R$, Morse structure the attractor
$\{3\}$, $X_0 = \{0\}$, $B = (1, 2)$ --- the orbit crosses $B$ on its
way to the attractor), so
recurrence structure without an enclosure proves no safety at all. The
sound premise is the certificate-level invariance fact, and the Morse
machinery contributes through the \emph{factory} instead: every lattice
member of the family object (and every down-closed cell-set, via
\rulename{QLatticeInv}) yields a $\mathsf{QInvScope}$ judgment whose
facet rows already exist in the family object. This rule then mints one
safety atom per avoid region at the price of that region\textquotesingle s
own disjointness rows and \emph{nothing else} --- the invariance side is
cited, never re-certified, so no invariance or facet row is ever added
however many avoid regions arrive --- which is the row-count payoff the
batch form
(Proposition~\ref{prop:qbasin-batch}) states as mathematics.
\end{bremark}

\begin{bremark}[what the checker recomputes]
At this rule the checker verifies: that the cited judgment's region and
init fields equal the rule's own scope data (the conclusion tie), and
the disjointness rows --- per-box, per-atom sign sequents of the base
calculus, or grid bookkeeping when $\beta_j$ is cell-aligned. Nothing
else: soundness is pure logic over the judgment's semantic content,
which is why the mechanized statement is three lines and
hypothesis-free. Basin-minted and battery-minted safety atoms are
\emph{type-identical} conclusions, deliberately: downstream citations
consume the judgment and may not key on its producer.
\end{bremark}

\begin{bremark}[relation to the literature]
Safety-by-invariance is the classical frame of set-invariance control
--- Nagumo's subtangentiality theorem \cite{nagumo1942} is the
ancestral invariance certificate, and Blanchini's survey
\cite{blanchini1999} the modern map of the field --- and of barrier
certificates \cite{prajnajadbabaie2004}, which in their usual packaging
synthesize one function
per specification, coupling the invariance argument to the safety
argument. The rule's point is the \emph{decoupling}: invariance is
certified once, as a judgment (by strictly inward facet rows, at
\rulename{QBlockInv}), and each safety claim is thereafter a citation
plus disjointness rows. The batch corollary is the certificate-side
face of the elementary observation that one invariant set separates its
contents from every disjoint avoid set at once.
\end{bremark}

\begin{anchors}
Theorem \ref{thm:qbasin-sound} & \texttt{qBasin\_sound} (QtslSoundness.lean); premise cargo \texttt{hscope}, \texttt{hAvoid}\\
Proposition \ref{prop:qbasin-batch} & \texttt{qBasin\_batch}\\
invariance-scope denotation & \texttt{Semiflow.InvScopeSat}\\
safety denotation ($\mathsf{SProbGE}(1, \cdot)$, deterministic reading) & \texttt{Semiflow.SafetySat}\\
disjointness rows & per-box \rulename{BernCone}-family sequents (base kernel); cell-aligned case = recomputed bookkeeping\\
\end{anchors}


\blcatsection{Existence Certificates}
\rulesection{75}{QZeroByMiranda}{core rule}
\label{blcatrule-QZeroByMiranda}
\ruleblock{The rule (Steele metanotation)}
\sourcenote{as realized in the mechanized development; source row: \texttt{\texttt{QTSL-rules.toml}} --- \emph{QZeroByMiranda}. The metanotation statement is typeset by hand for notational quality; this rule cites no Lean statements yet (the absence is stated loudly below). The rule \emph{roster} is machine-checked (every mechanized rule tag round-trips to exactly one catalogue row); the statement \emph{text} is editorial.}
\begin{mathpar}
\inferrule*{
  b : \mathsf{QMirandaBattery} \\
  \mathsf{qtsl\_miranda\_battery}(b,\ C,\ g,\ \Pi)
}{
  \vDash \mathsf{QZeroExists}(g,\ C)
}
\end{mathpar}

\par\medskip
\rulename{QZeroByMiranda} is Miranda\textquotesingle s theorem as a
certificate: $2m$ face rows over an equidimensional polynomial view,
each row a sign sequent claiming the signed component restricted to
one face of the cell nonpositive on the lower face and nonnegative on
the upper --- and the theorem concludes a zero of the vector field
inside the cell. Existence only: nothing in this family can conclude
uniqueness, count, or non-existence. Face geometry is derived from
each row\textquotesingle s indices, never stored --- a
producer-supplied face would be a trusted row --- and the
face-restriction recompute (exact-rational substitution) is the
walk\textquotesingle s derived-sequent authority.

\cataloguenote{The battery walk checks completeness (each (axis, face) exactly
once), per-axis component agreement, component bijectivity, and index
ranges as decided-false violations; non-box atoms, non-canonical or
over-wide polynomials, list ambiguity, and budget overruns refuse
with no verdict. Entered at the roster census (baked, never
catalogued); no Lean statement binds it yet.
}
\ruleblock{Theory rows (geolog)}
\sourcenote{The rule\textquotesingle s certificate content as it is declared in the geolog theory the checker loads --- payload schema, side-condition demands, and axiom rows, quoted from the theory source at every build of this catalogue (a stale quotation cannot exist). Axiom modes: \texttt{chase} = the engine derives; \texttt{report} = a violation is a loud judgment; \texttt{refuse} = a violation gates admission. Side-condition relations are discharged by bounded evaluation at the engine\textquotesingle s ground seam, never by derivation.}
\declhead{theory Qtsl}{QZeroByMiranda}{\badge{BadgeGreen}{quoted from the theory source at render}}
\begin{leansrc}
QZeroByMiranda : Set;
\end{leansrc}
\declhead{theory Qtsl}{QZeroByMiranda/*}{\badge{BadgeGreen}{quoted from the theory source at render}}
\begin{leansrc}
QZeroByMiranda/proof : QZeroByMiranda -> QProof;
QZeroByMiranda/g : QZeroByMiranda -> PolyVec;
QZeroByMiranda/cell : QZeroByMiranda -> QRegionAtom;
QZeroByMiranda/battery : QZeroByMiranda -> QMirandaBattery;
\end{leansrc}
\declhead{theory Qtsl}{qtsl\_miranda\_battery}{\badge{BadgeGreen}{quoted from the theory source at render}}
\begin{leansrc}
qtsl_miranda_battery : {battery: QMirandaBattery,
                        cell: QRegionAtom, g: PolyVec,
                        params: QRegionAtom} -> Prop;
\end{leansrc}
\ruleblock{Lean statement}
\badge{BadgeAmber}{not yet mechanized} this rule cites no Lean statements: nothing on this page is hash-bound for it, and its dashboard row shows the same absence across the statement, receipt, and correspondence columns.
\ruleblock{HOL4 statement}
\absentbox{NO HOL STATEMENT YET}{No HOL development of the extension theory exists.
}
\ruleblock{HOL verifier receipts}
\absentbox{NONE}{no HOL verification receipt exists for this rule.}
\ruleblock{Proof}

\sourcenote{Genre: core rule (tier 1), existence group. This section states the rule's soundness
obligation exactly, per the catalogue's completeness-over-polish
charter (a rule that can only be named-and-gapped is still a row);
the discharge is open debt on the mechanization column, and the
schema quoted in the entry is the baked ground truth.}

\textbf{Soundness obligation.}\enspace Miranda's theorem, certified: if each signed component of the
equidimensional polynomial view is nonpositive on its lower face and
nonnegative on its upper face of the cell (the $2m$ face rows, with
face geometry derived from row indices and each restriction
recomputed by exact-rational substitution), then the view has a zero
in the cell --- for every parameter point of the parameter cell.
Existence only; uniqueness, count, and non-existence are outside the
family by design.

\begin{bremark}[status]
No Lean statement binds this node. Miranda's theorem
(Poincar\'e--Miranda) is classical; the mechanization route would
state it over the boxed cell with the walk's derived face
restrictions as hypotheses.
\end{bremark}


\rulesection{76}{QEquilibriumFromInward}{admissible rule}
\label{blcatrule-QEquilibriumFromInward}
\ruleblock{The rule (Steele metanotation)}
\sourcenote{as realized in the mechanized development; source row: \texttt{\texttt{QTSL-rules.toml}} --- \emph{QEquilibriumFromInward}. The metanotation statement is typeset by hand for notational quality; this rule cites no Lean statements yet (the absence is stated loudly below). The rule \emph{roster} is machine-checked (every mechanized rule tag round-trips to exactly one catalogue row); the statement \emph{text} is editorial.}
\begin{mathpar}
\inferrule*{
  \mathrm{inv} : \mathsf{QBlockInv}\ \text{(strictly-inward facets, box block)}
}{
  \vDash \mathsf{QEquilibriumExists}(f,\ C)
}
\end{mathpar}

\par\medskip
\rulename{QEquilibriumFromInward} is a derived rule: a strictly-inward
facet battery on a box block --- every facet\textquotesingle s sign
sequent points the field inward --- is already the existence battery
for the negated field with identity pairing, so an equilibrium
inside the block follows by citation, with zero new oracle demand.
The premises are the block-invariant evidence rows themselves; the
unpacking instantiates the zero-existence obligation.

\cataloguenote{For fixed points the battery\textquotesingle s subject is the map
minus the identity; the elaboration recompute tying the two is a
named deferral at the ground seam. Entered at the roster census
(baked, never catalogued); no Lean statement binds it yet.
}
\ruleblock{Theory rows (geolog)}
\sourcenote{The rule\textquotesingle s certificate content as it is declared in the geolog theory the checker loads --- payload schema, side-condition demands, and axiom rows, quoted from the theory source at every build of this catalogue (a stale quotation cannot exist). Axiom modes: \texttt{chase} = the engine derives; \texttt{report} = a violation is a loud judgment; \texttt{refuse} = a violation gates admission. Side-condition relations are discharged by bounded evaluation at the engine\textquotesingle s ground seam, never by derivation.}
\declhead{theory Qtsl}{QEquilibriumFromInward}{\badge{BadgeGreen}{quoted from the theory source at render}}
\begin{leansrc}
QEquilibriumFromInward : Set;
\end{leansrc}
\declhead{theory Qtsl}{QEquilibriumFromInward/*}{\badge{BadgeGreen}{quoted from the theory source at render}}
\begin{leansrc}
QEquilibriumFromInward/proof : QEquilibriumFromInward -> QProof;
QEquilibriumFromInward/inv : QEquilibriumFromInward -> QBlockInv;
\end{leansrc}
\ruleblock{Lean statement}
\badge{BadgeAmber}{not yet mechanized} this rule cites no Lean statements: nothing on this page is hash-bound for it, and its dashboard row shows the same absence across the statement, receipt, and correspondence columns.
\ruleblock{HOL4 statement}
\absentbox{NO HOL STATEMENT YET}{No HOL development of the extension theory exists.
}
\ruleblock{HOL verifier receipts}
\absentbox{NONE}{no HOL verification receipt exists for this rule.}
\ruleblock{Proof}

\sourcenote{Genre: admissible rule (tier 2), existence group. This section states the rule's soundness
obligation exactly, per the catalogue's completeness-over-polish
charter (a rule that can only be named-and-gapped is still a row);
the discharge is open debt on the mechanization column, and the
schema quoted in the entry is the baked ground truth.}

\textbf{Soundness obligation.}\enspace A strictly-inward facet battery on a box block is the existence
battery for the negated field with identity pairing: inwardness on
every facet gives exactly the sign pattern Miranda's theorem demands
for $-f$, so the unpacking cites the block-invariant evidence rows
and instantiates the zero-existence obligation --- premises by
citation, zero new oracle demand. Soundness is therefore inherited
from the zero-existence rule's obligation.

\begin{bremark}[status]
No Lean statement binds this node; its obligation factors
entirely through the zero-existence rule's. For fixed points the
battery's subject is the map minus the identity, and the elaboration
recompute tying the two is a named deferral at the ground seam.
\end{bremark}


\rulesection{77}{QCoverByMiranda}{core rule}
\label{blcatrule-QCoverByMiranda}
\ruleblock{The rule (Steele metanotation)}
\sourcenote{as realized in the mechanized development; source row: \texttt{\texttt{QTSL-rules.toml}} --- \emph{QCoverByMiranda}. The metanotation statement is typeset by hand for notational quality; this rule cites no Lean statements yet (the absence is stated loudly below). The rule \emph{roster} is machine-checked (every mechanized rule tag round-trips to exactly one catalogue row); the statement \emph{text} is editorial.}
\begin{mathpar}
\inferrule*{
  b : \mathsf{QMirandaBattery} \\
  \mathsf{qtsl\_miranda\_battery}\bigl(b,\ W,\ \sigma - \mathrm{id}_{Z},\ Z\bigr)
}{
  \vDash Z \subseteq \sigma(W)
}
\end{mathpar}

\par\medskip
\rulename{QCoverByMiranda} is the chart-image discharge at
per-(patch cell, target cell) granularity: the battery\textquotesingle
s subject is chart minus target in joint variables, with the target
cell as parameters --- a certified zero of $\sigma(w) - z$ in the
patch cell for every $z$ in the target cell is exactly the
containment of the target cell in the chart image. Union to a region
cover rides the existing cover machinery; a cover is never inferred
from an image enclosure --- enclosures bound images from outside,
and only the zero-existence argument bounds them from inside.

\cataloguenote{Entered at the roster census (baked, never catalogued); no Lean
statement binds it yet. The battery walk is shared with the
zero-existence rule.
}
\ruleblock{Theory rows (geolog)}
\sourcenote{The rule\textquotesingle s certificate content as it is declared in the geolog theory the checker loads --- payload schema, side-condition demands, and axiom rows, quoted from the theory source at every build of this catalogue (a stale quotation cannot exist). Axiom modes: \texttt{chase} = the engine derives; \texttt{report} = a violation is a loud judgment; \texttt{refuse} = a violation gates admission. Side-condition relations are discharged by bounded evaluation at the engine\textquotesingle s ground seam, never by derivation.}
\declhead{theory Qtsl}{QCoverByMiranda}{\badge{BadgeGreen}{quoted from the theory source at render}}
\begin{leansrc}
QCoverByMiranda : Set;
\end{leansrc}
\declhead{theory Qtsl}{QCoverByMiranda/*}{\badge{BadgeGreen}{quoted from the theory source at render}}
\begin{leansrc}
QCoverByMiranda/proof : QCoverByMiranda -> QProof;
QCoverByMiranda/chart : QCoverByMiranda -> PolyVec;
QCoverByMiranda/wcell : QCoverByMiranda -> QRegionAtom;
QCoverByMiranda/zcell : QCoverByMiranda -> QRegionAtom;
QCoverByMiranda/battery : QCoverByMiranda -> QMirandaBattery;
\end{leansrc}
\ruleblock{Lean statement}
\badge{BadgeAmber}{not yet mechanized} this rule cites no Lean statements: nothing on this page is hash-bound for it, and its dashboard row shows the same absence across the statement, receipt, and correspondence columns.
\ruleblock{HOL4 statement}
\absentbox{NO HOL STATEMENT YET}{No HOL development of the extension theory exists.
}
\ruleblock{HOL verifier receipts}
\absentbox{NONE}{no HOL verification receipt exists for this rule.}
\ruleblock{Proof}

\sourcenote{Genre: core rule (tier 1), existence group. This section states the rule's soundness
obligation exactly, per the catalogue's completeness-over-polish
charter (a rule that can only be named-and-gapped is still a row);
the discharge is open debt on the mechanization column, and the
schema quoted in the entry is the baked ground truth.}

\textbf{Soundness obligation.}\enspace For each target point $z$ of the target cell, the battery (subject:
chart minus target in joint variables, parameters the target cell)
certifies a zero of $\sigma(w) - z$ in the patch cell --- i.e.\ a
preimage of $z$ under the chart --- so the target cell is contained
in the chart image. Union to a region cover rides the existing cover
machinery; a cover is never inferred from an image enclosure.

\begin{bremark}[status]
No Lean statement binds this node; the obligation factors through
the zero-existence rule's, applied parametrically in the target
point.
\end{bremark}


\blcatsection{Composition and Aggregation}
\rulesection{78}{QWireCompose}{core rule}
\label{blcatrule-QWireCompose}
\ruleblock{The rule (Steele metanotation)}
\sourcenote{as realized in the mechanized development; source row: \texttt{\texttt{QTSL-rules.toml}} --- \emph{QWireCompose}. The metanotation statement is typeset by hand for notational quality; its mechanized ground truth is the Lean statement below, extracted from the sources at every build of this book and hash-bound. The rule \emph{roster} is machine-checked (every mechanized rule tag round-trips to exactly one catalogue row); the statement \emph{text} is editorial, checked by eye against the extracted Lean beside it.}
\begin{mathpar}
\inferrule*[Right=\rulename{QWireCompose}]{
  \bigl\{\Sigma_i \vDash (I_i, \alpha_i, \gamma_i[, \kappa_i])\bigr\}_{i=1}^{n} \\
  \mathsf{QWireCert}\bigl(w,\ (\mathit{spec}_i)_{i} \Rightarrow \mathit{spec}_{\mathrm{out}},\ \kappa\bigr) \\
  \mathrm{rank} : \{1, \ldots, n\} \to \N; \quad \bigl\{\mathrm{rank}_i < \mathrm{rank}_j\bigr\}_{(i \to j) \in w} \ \text{(the causality certificate: one rank per component, strictly increasing along every wire)} \\
  \mathrm{W0}\!: \ \text{the composite martingale problem admits a c\`adl\`ag strong-Markov solution, unique in law, from every initial state (stochastic mechanisms; the martingale-problem sense --- distinct from the aggregation language\textquotesingle s own W0)} \\
  \mathrm{W1}\!: \ \text{transition data factor through per-place colour multisets (SDCPN mechanisms; typing-discharged in the aggregation language)} \\
  \mathrm{W2}\!: \ \text{the within-instant immediate-jump relation is well-founded, of uniformly bounded depth} \\
  \mathrm{W3}\!: \ \text{simultaneous enablings carry a declared serialization or a certified reset commutation (a.s.-null collision row for stochastic timers)} \\
  \kappa_{\mathrm{out}} := \textstyle\bigoplus_{\mathrm{legs}} \kappa_{\mathrm{leg}}, \quad \kappa_a \oplus \kappa_b := \kappa_a + \kappa_b \ \text{(slack addition along lens composition; checker-recomputed, never stored)}
}{
  \mathsf{SysWire}(w, [\Sigma_1, \ldots, \Sigma_n]) \vDash \mathit{spec}_{\mathrm{out}}
}
\end{mathpar}

\ruleflavor{satisfaction of an interface-bearing spec is plain \texttt{QSat} --- behaviours whose inputs conform to the assumption deliver the guarantee; closed $\mathcal{M} \vDash s$ is the empty-interface instance (defined once at the judgment-kinds home; the former subscripted display $\vDash_{\mathrm{I}}$ and the \texttt{QSatI} sort are both dissolved). The first two premises are derivation edges (judgment-kind premises); the well-posedness families enter as proof leaves on the node\textquotesingle s own well-posedness list, each owed only by the mechanism profiles that name it --- W0 by stochastic profiles, W1 by SDCPN profiles, W2/W3 the syntactic class (their four-way recompute relation is not yet enacted; the leaves enter as posted proof rows) --- and the causality ranks ride the node as one natural per component. The causality certificate and the slack composite are recomputed by the checker. Tier split derived from mechanism field-family population, never a stored tag: \textbf{admissible} where every mechanism derives ODE/discrete --- the composite certificate $\sum_i \lambda_i V_i$ is syntactically constructible and the node unpacks through the sequent-layer transport (\rulename{Promise}/\rulename{Syzygy}) --- \textbf{core} where any mechanism is stochastic (SDE/GSHS): the warrant is the per-profile semantic theorems. This display is the acyclic $n$-ary form (the causality certificate makes a directed cycle unrepresentable); feedback rides the binary small-gain cells, sum-form and gauge-form.}
\par\medskip
\rulename{QWireCompose} is the satisfaction-judgment
stratum\textquotesingle s composition rule --- assume--guarantee
transport over a structured system presentation. From
interface-indexed satisfaction of each component
($\mathsf{QSatI}(\Sigma_i, \mathit{spec}_i)$: every behaviour of
$\Sigma_i$ whose inputs satisfy the assumption $\alpha_i$ delivers the
guarantee $\gamma_i$, with quantitative slack where declared) and a
wiring certificate (the lens square transporting delivered guarantees
into downstream assumptions, with slack
$\kappa \in \mathcal{K}^0_\infty$), it concludes satisfaction of the
wired composite at the output spec. A producer uses it to certify a
system of interacting components --- each certified once, against its
own datasheet --- without ever building the monolithic certificate.

The soundness warrant factors, and the factorization organizes
everything on this page: $\mathrm{sound} = $ (transport arithmetic)
$\circ$ (semantic reading). The transport half is profile-generic
exact polynomial arithmetic --- the $(1, \lambda, 1, 1)$ cancellation
of cross-supplies, the sequent-layer machinery of \rulename{Promise}
and \rulename{Syzygy} --- and consumes only linearity of the
obligation operator over the wiring. The semantic half is
per-profile, and it is where this entry\textquotesingle s conditional
check lives. What the cited theorems prove, per profile and shape:
for general stochastic hybrid systems (GSHS), \emph{series}
composition (the conditioning tower, qualitative and product-bound
forms, plus the certified pipeline through the maximal-inequality
kernel), \emph{parallel} composition (the weighted-sum
supermartingale --- the semantic face of the syzygy weights), and
\emph{feedback in sum-form} (the cancellation run under conditional
expectation along the joint jump skeleton, with non-expansive
boundary rows) are proven; the $n$-ary \emph{acyclic} assembly is
proven at any $n$ by strong induction on the causality certificate,
with union-bound slack composition; and the per-profile tier split is
itself a proven, derived function of the presentation. \emph{Feedback
in gauge-form} (jumps that expand the certificate, repaid by flow
decay) is partial: its reduction theorem and the max-form
piecewise-deterministic instance are machine-checked end-to-end, while
the finite-index price of the diffusion term and the admissible-window
pinch remain open --- with the diffusion cell\textquotesingle s
parking reason retired: within-segment decay for gauged certificates
is now available at the martingale-problem profile, so the open part
is the accounting, not the stochastic calculus. The
deterministic-ODE row is at statement level;
the discrete-kernel row\textquotesingle s assembly is proven --- B1
$n$-ary acyclic, B2 parallel, and B3 sum-form small-gain, end-to-end
over the trajectory kernel --- while the sampled-data row remains
stated, not yet assembled. The discrete row carries a substrate
asymmetry worth stating with the grid: it consumes \emph{no}
interface hypothesis fields (contrast the GSHS cells\textquotesingle{}
D0--D4 data --- Dynkin, d\'ebut sampling, skeleton adaptedness),
because its execution semantics is mathlib\textquotesingle s own
\texttt{Kernel.traj}. The commentary states
this grid plainly because the dashboard\textquotesingle s conditional
mark warrants exactly the proven cells, never the whole catalogued
rule.

Two fences are part of the statement, not omissions. Undirected
stochastic wire sharing (common noise) is out of scope: stochastic
components sharing a wire need exactly one determining side, and the
proven parallel cell covers directed coupling and disjoint tensor
only. And the $n$-ary theorems are acyclic by construction --- the
presentation\textquotesingle s causality certificate makes a directed
cycle unrepresentable --- so feedback rides the binary small-gain
cells, where the coupling of jump times between the two components is
handled by the joint-skeleton argument.

What a checker recomputes, the rule\textquotesingle s row sorts now
baked: the causality certificate (a rank function strictly
increasing along wires), the slack composite (lens composition plus
slack addition, with the class membership rows discharged by
\rulename{QWireCertDischarge}), and the well-posedness row families
for the stochastic profiles. The component satisfactions and the
wiring squares are cited judgments and inherited sequent leaves ---
the rule adds no kernel arithmetic of its own.

\cataloguenote{The tier split is derived from mechanism field-family population,
never stored as a tag: admissible where every mechanism in the
presentation derives the ODE or discrete profile (the composite
certificate $\sum_i \lambda_i V_i$ is syntactically constructible and
the node unpacks through the sequent-layer transport), core where any
mechanism is stochastic (no unpacking exists; the semantic theorems
are the warrant). One rule node either way. The judgment sorts the
rule concludes into are baked in the extension theory; the checker
chain that would consume them is not yet ported --- the rule is
proven ahead of being wireable, and the dashboard row
shows exactly that state.
}
\ruleblock{Lean statement}
\declhead{BernsteinLogicAux.AGCompose.GshsSeries}{certified}{\badge{BadgeGreen}{statement bound}}
\begin{leansrc}
/-- **C1, certified form**: the full conditional-on-interface pipeline. Component
certificates enter as I-bridge outputs on the composite law — nonnegative
supermartingales along the stopped joint skeleton [D1][D2], with barrier values on
their bad sets (floor-row transport [D2]) and violation-witness inclusions [D3]
(the stopped skeleton samples the début of a violation; for component 2 the
certificate process is understood α₂-conformance-stopped, which is why its witness
hypothesis is relative to `conform₂`). The landed [K1] (`reach_avoid_div`,
`Ville.lean`) converts each certificate to a failure bound; `fail_bound` composes
them. Conclusion: the composite QSatI failure bound with certificate-valued ε's,

`μ (conform₁ ∩ satOutᶜ) ≤ ofReal (∫ Y₁ 0) / ofReal a₁ + ofReal (∫ Y₂ 0) / ofReal a₂`. -/
theorem certified [IsFiniteMeasure μ]
    {Y₁ Y₂ : ℕ → Ω → ℝ} {X₁ : ℕ → Ω → S₁} {X₂ : ℕ → Ω → S₂}
    {B₁ : Set S₁} {B₂ : Set S₂} {a₁ a₂ : ℝ}
    -- [D1]+[D2]: I-bridge outputs on the composite law
    (hY₁ : Supermartingale Y₁ 𝒢 μ) (hY₁0 : 0 ≤ Y₁)
    (hY₂ : Supermartingale Y₂ 𝒢 μ) (hY₂0 : 0 ≤ Y₂)
    -- floor-row barriers at skeleton samples [D2]
    (ha₁ : 0 < a₁) (ha₂ : 0 < a₂)
    (hbar₁ : ∀ n ω, X₁ n ω ∈ B₁ → a₁ ≤ Y₁ n ω)
    (hbar₂ : ∀ n ω, X₂ n ω ∈ B₂ → a₂ ≤ Y₂ n ω)
    -- violation witnesses at skeleton samples [D3]
    (hwit₁ : conform₁ ∩ sat₁ᶜ ⊆ {ω | ∃ n, X₁ n ω ∈ B₁})
    (hwit₂ : conform₂ ∩ sat₂ᶜ ⊆ {ω | ∃ n, X₂ n ω ∈ B₂})
    -- the wiring square (κ-slack folded) and the output reading
    (hw : conform₁ ∩ sat₁ ⊆ conform₂)
    (hout : conform₁ ∩ sat₁ ∩ sat₂ ⊆ satOut) :
    μ (conform₁ ∩ satOutᶜ) ≤
      ENNReal.ofReal (∫ ω, Y₁ 0 ω ∂μ) / ENNReal.ofReal a₁ +
      ENNReal.ofReal (∫ ω, Y₂ 0 ω ∂μ) / ENNReal.ofReal a₂
\end{leansrc}
\declhead{BernsteinLogicAux.AGCompose.GshsParallel}{certified}{\badge{BadgeGreen}{statement bound}}
\begin{leansrc}
/-- **C2, certified form** (n-ary, sum-form): component certificates — nonnegative
supermartingales `Y i` w.r.t. the joint skeleton filtration [D1][D2], syzygy
weights `λᵢ ≥ 0`, and per-component floor rows giving a common barrier level `a` on
each bad set (boundary rows transport pointwise through sums, inventory §3.C
claim 1: on `Bad i` the term `λᵢ·Yᵢ` alone clears `a`, the other terms being
nonnegative) — yield the composite avoid bound via the landed [K1]:

`ofReal a * μ {∃ n, X n ∈ ⋃ i ∈ s, Bad i} ≤ ofReal (∑ i ∈ s, λᵢ · ∫ Yᵢ 0)`.

The right side is the syzygy-weighted initial certificate mass — the semantic
reading of the composite certificate `∑ λᵢVᵢ` at the initial skeleton time. To
recover per-component barrier levels `aᵢ`, instantiate `a := min over s of λᵢ·aᵢ`. -/
theorem certified (s : Finset ι) [IsFiniteMeasure μ]
    {Y : ι → ℕ → Ω → ℝ} {lam : ι → ℝ} {X : ℕ → Ω → St} {Bad : ι → Set St} {a : ℝ}
    (hY : ∀ i ∈ s, Supermartingale (Y i) 𝒢 μ)
    (hY0 : ∀ i ∈ s, 0 ≤ Y i)
    (hlam : ∀ i ∈ s, 0 ≤ lam i)
    (hbar : ∀ i ∈ s, ∀ n ω, X n ω ∈ Bad i → a ≤ lam i * Y i n ω) :
    ENNReal.ofReal a * μ {ω | ∃ n, X n ω ∈ ⋃ i ∈ s, Bad i} ≤
      ENNReal.ofReal (∑ i ∈ s, lam i * ∫ ω, Y i 0 ω ∂μ)
\end{leansrc}
\declhead{BernsteinLogicAux.AGCompose.GshsParallel}{certified\_pair}{\badge{BadgeGreen}{statement bound}}
\begin{leansrc}
/-- **C2, certified form, two components** — the inventory's `V₁ + λV₂` shape,
directly. Barrier reading: component 1's floor clears `a` on `B₁`, component 2's
λ-weighted floor clears `a` on `B₂`; the composite avoid bound is

`ofReal a * μ {∃ n, X n ∈ B₁ ∪ B₂} ≤ ofReal (∫ Y₁ 0 + λ · ∫ Y₂ 0)`. -/
theorem certified_pair [IsFiniteMeasure μ]
    {Y₁ Y₂ : ℕ → Ω → ℝ} {X : ℕ → Ω → St} {B₁ B₂ : Set St} {lam a : ℝ}
    (hY₁ : Supermartingale Y₁ 𝒢 μ) (hY₁0 : 0 ≤ Y₁)
    (hY₂ : Supermartingale Y₂ 𝒢 μ) (hY₂0 : 0 ≤ Y₂) (hlam : 0 ≤ lam)
    (hbar₁ : ∀ n ω, X n ω ∈ B₁ → a ≤ Y₁ n ω)
    (hbar₂ : ∀ n ω, X n ω ∈ B₂ → a ≤ lam * Y₂ n ω) :
    ENNReal.ofReal a * μ {ω | ∃ n, X n ω ∈ B₁ ∪ B₂} ≤
      ENNReal.ofReal (∫ ω, Y₁ 0 ω ∂μ + lam * ∫ ω, Y₂ 0 ω ∂μ)
\end{leansrc}
\declhead{BernsteinLogicAux.GshsFeedbackSum}{supermartingale}{\badge{BadgeGreen}{statement bound}}
\begin{leansrc}
/-- **The composite certificate is a supermartingale** for the joint-skeleton
filtration — the C3a conclusion's engine. Follows from `drift` by discarding the
nonnegative θ-margin. -/
theorem supermartingale : Supermartingale F.W 𝒢 μ
\end{leansrc}
\declhead{BernsteinLogicAux.GshsFeedbackSum}{avoid\_bound}{\badge{BadgeGreen}{statement bound}}
\begin{leansrc}
/-- **The composite avoid bound** (K1, consumed as the landed
`BernsteinLogic.Ville.reach_avoid`): a barrier `a` for `W` on a sampled bad set `B`
bounds the probability that any joint-skeleton sample reaches `B`. -/
theorem avoid_bound {S : Type*} {X : ℕ → Ω → S} {B : Set S} {a : ℝ}
    (hbarrier : ∀ k ω, X k ω ∈ B → a ≤ F.W k ω) :
    ENNReal.ofReal a * μ {ω | ∃ k, X k ω ∈ B} ≤ ENNReal.ofReal (∫ ω, F.W 0 ω ∂μ)
\end{leansrc}
\declhead{BernsteinLogicAux.MaxFormPDMP}{sound\_GSHS\_safety\_barrier}{\badge{BadgeGreen}{statement bound}}
\begin{leansrc}
/-- **The T5 consumer at a general positive barrier** (Statement II's
Step 7 barrier `min(1, λ/λ')` needs this): with floor `m > 0` on the avoid
region, the reach probability is bounded by `V x / m`. Proved by scaling
the certificate to barrier `1` and consuming the landed
`sound_GSHS_safety`; no new engine plumbing. -/
theorem sound_GSHS_safety_barrier [IsMarkovKernel M.seg] [IsMarkovKernel M.jump]
    {W : State → ℝ} {U : Set State} {m : ℝ} (hm : 0 < m)
    (hW : Measurable W) (hU : MeasurableSet U)
    (hSegD : SegmentDynkin M 0 W) (hForce : ForcedNonexpansion M W)
    (hNonneg : NonnegV W) {C : ℝ} (hC0 : 0 ≤ C) (hWC : ∀ y, W y ≤ C)
    (hFloor : ∀ y ∈ U, m ≤ W y) (x : State) :
    Battery.pathMeasure (skel M) (SkelState.alive x) {ω | ∃ n, ω n ∈ skelBad U}
      ≤ ENNReal.ofReal (W x / m)
\end{leansrc}
\par\medskip
\textbf{The presentation layer, bound by display.}\enspace The
\texttt{System} structure is the abstraction the composition theorems
consume: the field-family population pattern of each mechanism (from
which the doctrine profile is \emph{derived}, never stored) and the
wiring order with its causality certificate carried as data --- a rank
function strictly increasing along wires, which is precisely what the
checker-side causality condition recomputes. \texttt{System.qagTier}
derives the rule\textquotesingle s tier from the same data, with both
unfoldings proven; \texttt{admissible\_constructs} makes the
admissible face\textquotesingle s \textquotedblleft syntactically
constructible\textquotedblright{} clause a checkable statement rather
than prose.

\declhead{BernsteinLogicAux.QtslCompose.Presentation}{System}{\badge{BadgeGreen}{statement bound}}
\begin{leansrc}
/-- A System presentation, at the composition-relevant abstraction: `n` charts
with their field-family population patterns, the wiring DAG (`upstream j` =
the charts whose outputs wire into `j`'s inputs), and the causality
certificate carried as **data** — a rank function strictly increasing along
wires, exactly what the checker's causality-DAG side condition recomputes.
Feedback (a directed cycle) is *unrepresentable* in this structure; the
feedback transport shape rides the landed binary cells
(`AGCompose/GshsFeedbackSum.lean`, `GaugeSmallGain.lean`, `MaxFormPDMP.lean`)
— see the AGCompose module docstring. -/
structure System (n : ℕ) : Type where
  charts : Fin n → ChartShape
  upstream : Fin n → Finset (Fin n)
  rank : Fin n → ℕ
  causality : ∀ j, ∀ i ∈ upstream j, rank i < rank j
\end{leansrc}
\declhead{BernsteinLogicAux.QtslCompose.Presentation}{System.qagTier}{\badge{BadgeGreen}{statement bound}}
\begin{leansrc}
/-- The `QWireCompose` tier at this System, **derived from the derived
profiles** (never a stored tag): admissible iff every chart's profile admits a
syntactically constructible composite certificate. -/
def System.qagTier {n : ℕ} (S : System n) : Tier
\end{leansrc}
\declhead{BernsteinLogicAux.QtslCompose.Presentation}{System.qagTier\_admissible\_iff}{\badge{BadgeGreen}{statement bound}}
\begin{leansrc}
/-- The split, unfolded to the profile characterization: admissible iff every
chart derives ODE or discrete. -/
theorem System.qagTier_admissible_iff {n : ℕ} (S : System n) :
    S.qagTier = .admissible ↔
      ∀ i, (S.charts i).profile = .ode ∨ (S.charts i).profile = .discrete
\end{leansrc}
\declhead{BernsteinLogicAux.QtslCompose.Presentation}{System.qagTier\_core\_iff}{\badge{BadgeGreen}{statement bound}}
\begin{leansrc}
/-- The core face of the split: core iff some chart is stochastic (SDE or
GSHS). -/
theorem System.qagTier_core_iff {n : ℕ} (S : System n) :
    S.qagTier = .core ↔
      ∃ i, (S.charts i).profile = .sde ∨ (S.charts i).profile = .gshs
\end{leansrc}
\declhead{BernsteinLogicAux.QtslCompose.Presentation}{admissible\_constructs}{\badge{BadgeGreen}{statement bound}}
\begin{leansrc}
/-- The core face of the split: core iff some chart is stochastic (SDE or
GSHS). -/
theorem System.qagTier_core_iff {n : ℕ} (S : System n) :
    S.qagTier = .core ↔
      ∃ i, (S.charts i).profile = .sde ∨ (S.charts i).profile = .gshs := by
  have hsplit : S.qagTier = .admissible ∨ S.qagTier = .core := by
    unfold System.qagTier
    split <;> simp
  constructor
  · intro h
    by_contra hno
    push_neg at hno
    have : S.qagTier = .admissible := by
      rw [S.qagTier_admissible_iff]
      intro i
      have := hno i
      cases hp : (S.charts i).profile <;> simp_all
    rw [this] at h
    exact absurd h (by simp)
  · intro ⟨i, hi⟩
    rcases hsplit with h | h
    · rw [S.qagTier_admissible_iff] at h
      rcases h i with ho | hd <;> rcases hi with hs | hs <;> simp_all
    · exact h

/-! ## The admissible face's constructibility, pointwise

The composite certificate for the admissible tier is the weighted sum
`∑ i, λ i • V i`; its composite obligation row follows *pointwise* from the
component dissipation rows (`e i ≤ 0` on the operating region) and the wiring
identification rows (`g k = 0` there), at nonnegative certificate weights and
free-sign wiring multipliers. This is the pointwise shadow of the landed
sequent-layer theorem (`sound_syzygy`, main package — the (T) layer of the
inventory §0 factorization); it lives here so that the tier definition's
"syntactically constructible" clause is a statement, not prose. -/
theorem admissible_constructs {ι κIdx X : Type*} [Fintype ι] [Fintype κIdx]
    {R : Set X} (lam : ι → ℝ) (e : ι → X → ℝ) (nu : κIdx → ℝ) (g : κIdx → X → ℝ)
    (hlam : ∀ i, 0 ≤ lam i)
    (he : ∀ i, ∀ x ∈ R, e i x ≤ 0)
    (hg : ∀ k, ∀ x ∈ R, g k x = 0) :
    ∀ x ∈ R, (∑ i, lam i * e i x) + (∑ k, nu k * g k x) ≤ 0
\end{leansrc}
\declhead{BernsteinLogicAux.QtslCompose.AGCompose}{sat\_of\_conform}{\badge{BadgeGreen}{statement bound}}
\begin{leansrc}
/-- **Conformance propagation** (pointwise core). If the wiring squares
transport external conformance + upstream deliveries into each component's
conformance, and each component delivers on conforming inputs (at this `ω`),
then every component delivers. Strong induction on the causality rank — the
presentation's causality certificate is exactly what makes this well-founded. -/
theorem sat_of_conform
    (hw : ∀ j, conformExt ∩ (⋂ i ∈ S.upstream j, sat i) ⊆ conform j)
    {ω : Ω} (h : ∀ j, ω ∈ conform j → ω ∈ sat j)
    (hext : ω ∈ conformExt) : ∀ j, ω ∈ sat j
\end{leansrc}
\declhead{BernsteinLogicAux.QtslCompose.AGCompose}{qualitative}{\badge{BadgeGreen}{statement bound}}
\begin{leansrc}
/-- **QWireCompose, qualitative form** (a.e.; the n-ary generalization of
`AGCompose.Series.qualitative`): component QSatI's as a.e. implications +
wiring squares yield the composite a.e. delivery of every guarantee, hence of
any output spec read from them (`hout`). -/
theorem qualitative [MeasurableSpace Ω] {μ : Measure Ω}
    (hw : ∀ j, conformExt ∩ (⋂ i ∈ S.upstream j, sat i) ⊆ conform j)
    (h : ∀ j, ∀ᵐ ω ∂μ, ω ∈ conform j → ω ∈ sat j)
    {satOut : Set Ω} (hout : conformExt ∩ (⋂ j, sat j) ⊆ satOut) :
    ∀ᵐ ω ∂μ, ω ∈ conformExt → ω ∈ satOut
\end{leansrc}
\declhead{BernsteinLogicAux.QtslCompose.AGCompose}{fail\_subset}{\badge{BadgeGreen}{statement bound}}
\begin{leansrc}
/-- **The n-ary failure decomposition** (pure set algebra; the rule's
union-bound side of slack composition): a composite failure under conforming
external inputs localizes at some component that failed *while conforming* —
namely a rank-minimal failing one, whose upstream all delivered. -/
theorem fail_subset
    (hw : ∀ j, conformExt ∩ (⋂ i ∈ S.upstream j, sat i) ⊆ conform j) :
    conformExt ∩ (⋂ j, sat j)ᶜ ⊆ ⋃ j, conform j ∩ (sat j)ᶜ
\end{leansrc}
\declhead{BernsteinLogicAux.QtslCompose.AGCompose}{fail\_bound}{\badge{BadgeGreen}{statement bound}}
\begin{leansrc}
/-- **QWireCompose, probabilistic form: n-ary union-bound slack composition** —
the inventory §3.B1 display `ℙ(composite fails) ≤ ∑ i (1 − p_i)` at any `n`,
in failure-measure form (`ε j` plays `1 − p_j`). -/
theorem fail_bound [MeasurableSpace Ω] {μ : Measure Ω} {ε : Fin n → ℝ≥0∞}
    (hw : ∀ j, conformExt ∩ (⋂ i ∈ S.upstream j, sat i) ⊆ conform j)
    (hε : ∀ j, μ (conform j ∩ (sat j)ᶜ) ≤ ε j) :
    μ (conformExt ∩ (⋂ j, sat j)ᶜ) ≤ ∑ j, ε j
\end{leansrc}
\declhead{BernsteinLogicAux.QtslCompose.AGCompose}{fail\_bound\_out}{\badge{BadgeGreen}{statement bound}}
\begin{leansrc}
/-- `fail_bound` read at an output spec (`hout` = the output reading of the
delivered guarantees, as in the binary cells). -/
theorem fail_bound_out [MeasurableSpace Ω] {μ : Measure Ω} {ε : Fin n → ℝ≥0∞}
    {satOut : Set Ω}
    (hw : ∀ j, conformExt ∩ (⋂ i ∈ S.upstream j, sat i) ⊆ conform j)
    (hout : conformExt ∩ (⋂ j, sat j) ⊆ satOut)
    (hε : ∀ j, μ (conform j ∩ (sat j)ᶜ) ≤ ε j) :
    μ (conformExt ∩ satOutᶜ) ≤ ∑ j, ε j
\end{leansrc}
\par\medskip
\textbf{The B-row (discrete-kernel) assembly.}\enspace The row the
commentary\textquotesingle s grid records as proven: B1
$n$-ary acyclic union-bound composition through the first-failure
device, B2 parallel weighted-sum on the product kernel, B3 sum-form
small-gain cancellation --- machine-checked end-to-end over
\texttt{Kernel.traj} through the battery substrate, with no semantic
hypothesis fields (the profile identification enters as the choice of
substrate, which is the conditional mark\textquotesingle s honest
residue for this row).

\declhead{BernsteinLogicAux.QtslCompose.AGComposeDiscrete}{firstFail}{\badge{BadgeGreen}{statement bound}}
\begin{leansrc}
/-- The event that component `j` is a **first** violator: at some time the path
is in `j`'s violation set, no component's violation set having been visited
before. Distinct components' first-failure events may overlap (simultaneous
first violations); the union bound is indifferent. -/
def firstFail (bad : Fin m → Set S) (j : Fin m) : Set (ℕ → S)
\end{leansrc}
\declhead{BernsteinLogicAux.QtslCompose.AGComposeDiscrete}{reach\_subset\_iUnion\_firstFail}{\badge{BadgeGreen}{statement bound}}
\begin{leansrc}
/-- **The first-failure decomposition** (the discrete-profile face of K4's event
decomposition; kin of the event-stratum `AGCompose.fail_subset`, with
time-minimality replacing rank-minimality): a path that ever visits any
violation set has a first violation, and that violation belongs to some
component. Pure path-set algebra. -/
theorem reach_subset_iUnion_firstFail (bad : Fin m → Set S) :
    {ω : ℕ → S | ∃ n, ω n ∈ ⋃ i, bad i} ⊆ ⋃ j, firstFail bad j
\end{leansrc}
\declhead{BernsteinLogicAux.QtslCompose.AGComposeDiscrete}{measure\_firstFail\_le}{\badge{BadgeGreen}{statement bound}}
\begin{leansrc}
/-- **The stopped-process device** (inventory §3.B1), isolated: stopping the
composite chain at the first violation by *any* component, component `j`'s
certificate — with drift live off the **joint** violation set and barrier only
on `j`'s **own** violation set — bounds the probability that `j` fails first
by its initial value. K1 (`Supermartingale.ville_ineq`) at barrier `1`;
contrast `Battery.measure_reach_le`, whose barrier must cover the whole stop
set and which therefore cannot see a multi-component decomposition. -/
theorem measure_firstFail_le {bad : Fin m → Set S} {V : S → ℝ} {C : ℝ} {j : Fin m}
    (hA : MeasurableSet (⋃ i, bad i)) (hV : Measurable V)
    (hV0 : ∀ y, 0 ≤ V y) (hVC : ∀ y, V y ≤ C)
    (hbar : ∀ y ∈ bad j, 1 ≤ V y)
    (hdrift : ∀ y, y ∉ (⋃ i, bad i) → ∫ z, V z ∂K y ≤ V y) (x : S) :
    pathMeasure K x (firstFail bad j) ≤ ENNReal.ofReal (V x)
\end{leansrc}
\declhead{BernsteinLogicAux.QtslCompose.AGComposeDiscrete}{fail\_bound}{\badge{BadgeGreen}{statement bound}}
\begin{leansrc}
/-- **QWireCompose B-row, probabilistic form** (engine form; inventory §3.B1's
display `ℙ(composite fails) ≤ ∑ᵢ (1−pᵢ)` at any `n`, semantically discharged
over `Kernel.traj`): per-component bounded certificates with drift off the
joint violation set and barriers on their own violation sets bound the
composite failure probability by the sum of initial certificate values. The
ε-shaped quantities the event-stratum `AGCompose.fail_bound` consumes as
hypotheses are here established from kernel-level data. -/
theorem fail_bound {bad : Fin m → Set S} {V : Fin m → S → ℝ} {C : Fin m → ℝ}
    (hA : MeasurableSet (⋃ i, bad i))
    (hV : ∀ j, Measurable (V j)) (hV0 : ∀ j y, 0 ≤ V j y)
    (hVC : ∀ j y, V j y ≤ C j)
    (hbar : ∀ j, ∀ y ∈ bad j, 1 ≤ V j y)
    (hdrift : ∀ j, ∀ y, y ∉ (⋃ i, bad i) → ∫ z, V j z ∂K y ≤ V j y) (x : S) :
    pathMeasure K x {ω | ∃ n, ω n ∈ ⋃ i, bad i}
      ≤ ∑ j, ENNReal.ofReal (V j x)
\end{leansrc}
\declhead{BernsteinLogicAux.QtslCompose.AGComposeDiscrete}{fail\_bound\_wired}{\badge{BadgeGreen}{statement bound}}
\begin{leansrc}
/-- **QWireCompose B-row, wired form** (the catalogued rule's hypothesis shape):
per-component **datasheet** rows — drift valid only on the component's own
conformance region — plus the state-level wiring square (external conformance
folded; joint upstream delivery implies conformance, the I-bridge reading of
`QWireCert` at this stratum) yield the same union-bound conclusion. The
`System`'s wiring order supplies the square's index set; its rank/causality
data is consumed by the event-stratum assembly, not here (time-minimality
replaces the rank induction — see the module header, and `fail_bound_upstream`
for the honest-width form this now specializes). -/
theorem fail_bound_wired (Sys : System m) {conf bad : Fin m → Set S}
    {V : Fin m → S → ℝ} {C : Fin m → ℝ}
    (hA : MeasurableSet (⋃ i, bad i))
    (hw : ∀ j, (⋂ i ∈ Sys.upstream j, (bad i)ᶜ) ⊆ conf j)
    (hV : ∀ j, Measurable (V j)) (hV0 : ∀ j y, 0 ≤ V j y)
    (hVC : ∀ j y, V j y ≤ C j)
    (hbar : ∀ j, ∀ y ∈ bad j, 1 ≤ V j y)
    (hdrift : ∀ j, ∀ y ∈ conf j, y ∉ bad j → ∫ z, V j z ∂K y ≤ V j y) (x : S) :
    pathMeasure K x {ω | ∃ n, ω n ∈ ⋃ i, bad i}
      ≤ ∑ j, ENNReal.ofReal (V j x)
\end{leansrc}
\declhead{BernsteinLogicAux.QtslCompose.AGComposeDiscrete}{ae\_avoid\_of\_zero}{\badge{BadgeGreen}{statement bound}}
\begin{leansrc}
/-- **The qualitative reading** (inventory §3.B1's "a.s.-invariance-shaped"
half): started strictly inside every certificate's zero set, the composite
chain a.s. never violates any component's guarantee — the conditioning-tower
display degenerates, at the discrete profile, to the union bound at total
mass zero. -/
theorem ae_avoid_of_zero {bad : Fin m → Set S} {V : Fin m → S → ℝ}
    {C : Fin m → ℝ}
    (hA : MeasurableSet (⋃ i, bad i))
    (hV : ∀ j, Measurable (V j)) (hV0 : ∀ j y, 0 ≤ V j y)
    (hVC : ∀ j y, V j y ≤ C j)
    (hbar : ∀ j, ∀ y ∈ bad j, 1 ≤ V j y)
    (hdrift : ∀ j, ∀ y, y ∉ (⋃ i, bad i) → ∫ z, V j z ∂K y ≤ V j y) {x : S}
    (hzero : ∀ j, V j x = 0) :
    ∀ᵐ ω ∂pathMeasure K x, ∀ n, ω n ∉ ⋃ i, bad i
\end{leansrc}
\declhead{BernsteinLogicAux.QtslCompose.AGComposeDiscrete}{drift\_parallelComp}{\badge{BadgeGreen}{statement bound}}
\begin{leansrc}
/-- **B2, the linearity row** (inventory §3.B2's one-paragraph argument,
machine-checked): component drift rows transport to the weighted-sum
certificate on the synchronous product — the syzygy combination at weights
`(1, λ)`, verbatim. Off the lifted joint violation set both component rows
are live; linearity of the integral over `K₁ ∥ₖ K₂` does the rest. -/
theorem drift_parallelComp {bad₁ : Set S₁} {bad₂ : Set S₂}
    {V₁ : S₁ → ℝ} {V₂ : S₂ → ℝ} {lam : ℝ} (hlam : 0 ≤ lam)
    (hint₁ : ∀ y, y ∉ bad₁ → Integrable V₁ (K₁ y))
    (hint₂ : ∀ y, y ∉ bad₂ → Integrable V₂ (K₂ y))
    (hd₁ : ∀ y, y ∉ bad₁ → ∫ z, V₁ z ∂K₁ y ≤ V₁ y)
    (hd₂ : ∀ y, y ∉ bad₂ → ∫ z, V₂ z ∂K₂ y ≤ V₂ y) :
    ∀ p : S₁ × S₂, p ∉ (Prod.fst ⁻¹' bad₁ ∪ Prod.snd ⁻¹' bad₂) →
      ∫ q, (V₁ q.1 + lam * V₂ q.2) ∂(K₁ ∥ₖ K₂) p
        ≤ V₁ p.1 + lam * V₂ p.2
\end{leansrc}
\declhead{BernsteinLogicAux.QtslCompose.AGComposeDiscrete}{avoid\_parallelComp}{\badge{BadgeGreen}{statement bound}}
\begin{leansrc}
/-- **B2, assembled** (the discrete face of `GshsParallel.certified`'s
weighted-sum supermartingale claim, end-to-end): on the synchronous product
chain, the weighted-sum certificate bounds the probability that *either*
component ever violates by its initial value. The weight's role is visible in
the barrier hypotheses: component 2's barrier is carried at level `1/λ`
(stated as `1 ≤ λ·V₂` — how the syzygy weights meet the common barrier). -/
theorem avoid_parallelComp {bad₁ : Set S₁} {bad₂ : Set S₂}
    {V₁ : S₁ → ℝ} {V₂ : S₂ → ℝ} {C₁ C₂ lam : ℝ} (hlam : 0 ≤ lam)
    (hbad₁ : MeasurableSet bad₁) (hbad₂ : MeasurableSet bad₂)
    (hV₁ : Measurable V₁) (hV₂ : Measurable V₂)
    (hV₁0 : ∀ y, 0 ≤ V₁ y) (hV₂0 : ∀ y, 0 ≤ V₂ y)
    (hV₁C : ∀ y, V₁ y ≤ C₁) (hV₂C : ∀ y, V₂ y ≤ C₂)
    (hbar₁ : ∀ y ∈ bad₁, 1 ≤ V₁ y) (hbar₂ : ∀ y ∈ bad₂, 1 ≤ lam * V₂ y)
    (hd₁ : ∀ y, y ∉ bad₁ → ∫ z, V₁ z ∂K₁ y ≤ V₁ y)
    (hd₂ : ∀ y, y ∉ bad₂ → ∫ z, V₂ z ∂K₂ y ≤ V₂ y) (x : S₁ × S₂) :
    pathMeasure (K₁ ∥ₖ K₂) x
        {ω | ∃ n, ω n ∈ (Prod.fst ⁻¹' bad₁ ∪ Prod.snd ⁻¹' bad₂ : Set (S₁ × S₂))}
      ≤ ENNReal.ofReal (V₁ x.1 + lam * V₂ x.2)
\end{leansrc}
\declhead{BernsteinLogicAux.QtslCompose.AGComposeDiscrete}{smallGain\_drift}{\badge{BadgeGreen}{statement bound}}
\begin{leansrc}
/-- **B3, the cancellation row** (inventory §3.B3; the semantic face of the
checker-level `(1, λ, 1, 1)` syzygy fold `sound_smallGain`, at the kernel
level): per-component dissipation rows with cross-supplies on the joint
closed-loop kernel, under the small-gain membership rows, yield the
`θ`-margin drift row for `W = V₁ + λV₂`. Only `0 ≤ λ`, `0 ≤ θ` are consumed;
the rule-level ranges `λ > 0`, `θ ∈ (0,1)` are well-formedness, never
hypotheses (the width find of the C3a cell, reproduced at this profile). -/
theorem smallGain_drift (hlam : 0 ≤ lam)
    (hint₁ : ∀ y, y ∉ A → Integrable V₁ (K y))
    (hint₂ : ∀ y, y ∉ A → Integrable V₂ (K y))
    (hd₁ : ∀ y, y ∉ A → ∫ z, V₁ z ∂K y ≤ V₁ y - alpha₁ y + sigma₁ y)
    (hd₂ : ∀ y, y ∉ A → ∫ z, V₂ z ∂K y ≤ V₂ y - alpha₂ y + sigma₂ y)
    (hm₁ : ∀ y, y ∉ A → lam * sigma₂ y ≤ (1 - theta) * alpha₁ y)
    (hm₂ : ∀ y, y ∉ A → sigma₁ y ≤ lam * ((1 - theta) * alpha₂ y)) :
    ∀ y, y ∉ A → ∫ z, (V₁ z + lam * V₂ z) ∂K y
      ≤ (V₁ y + lam * V₂ y) - theta * (alpha₁ y + lam * alpha₂ y)
\end{leansrc}
\declhead{BernsteinLogicAux.QtslCompose.AGComposeDiscrete}{avoid\_smallGain}{\badge{BadgeGreen}{statement bound}}
\begin{leansrc}
/-- **B3, assembled** (the discrete face of the C3a conclusion, end-to-end):
on the closed loop, the small-gain data yields the composite avoid bound for
the weighted-sum certificate — Foster–Lyapunov metatheory
(`ville_ineq`/optional stopping, through `Battery.measure_reach_le`) closing
the cell exactly as the inventory's §3.B3 statement names it. The `α`-sign
hypotheses convert the `θ`-margin into a plain drift row; the barrier is the
caller's (how `1 ≤ W` distributes between components is model data). -/
theorem avoid_smallGain (hlam : 0 ≤ lam) (htheta : 0 ≤ theta)
    {C₁ C₂ : ℝ} (hA : MeasurableSet A)
    (hV₁ : Measurable V₁) (hV₂ : Measurable V₂)
    (hV₁0 : ∀ y, 0 ≤ V₁ y) (hV₂0 : ∀ y, 0 ≤ V₂ y)
    (hV₁C : ∀ y, V₁ y ≤ C₁) (hV₂C : ∀ y, V₂ y ≤ C₂)
    (hbar : ∀ y ∈ A, 1 ≤ V₁ y + lam * V₂ y)
    (halpha₁ : ∀ y, y ∉ A → 0 ≤ alpha₁ y)
    (halpha₂ : ∀ y, y ∉ A → 0 ≤ alpha₂ y)
    (hd₁ : ∀ y, y ∉ A → ∫ z, V₁ z ∂K y ≤ V₁ y - alpha₁ y + sigma₁ y)
    (hd₂ : ∀ y, y ∉ A → ∫ z, V₂ z ∂K y ≤ V₂ y - alpha₂ y + sigma₂ y)
    (hm₁ : ∀ y, y ∉ A → lam * sigma₂ y ≤ (1 - theta) * alpha₁ y)
    (hm₂ : ∀ y, y ∉ A → sigma₁ y ≤ lam * ((1 - theta) * alpha₂ y)) (x : S) :
    pathMeasure K x {ω | ∃ n, ω n ∈ A}
      ≤ ENNReal.ofReal (V₁ x + lam * V₂ x)
\end{leansrc}
\ruleblock{Correspondence (Steele display $\leftrightarrow$ Lean)}
\badge{BadgeAmber}{a bound Lean statement CHANGED since certification --- re-review and re-certify}
\sourcenote{The dictionary below is a reviewer-certified statement of how the mechanized statements above interpret the displayed rule: display metavariable $\mapsto$ Lean binder, each display hypothesis $\mapsto$ the premise or side condition carrying it, the display conclusion $\mapsto$ the goal --- with any deliberate difference in width between display and mechanization named outright. Certification is recorded as a hash triple over (the displayed cell, the extracted statement texts, this dictionary\textquotesingle s own text); if any of the three changes, this block turns amber until a reviewer re-certifies. The correspondence claim itself remains editorial --- reviewer-grade, not machine-checked; its machine-checked successor types each soundness statement against the rule as a structured object, and supersedes this block per rule as it lands.}
The display is one $n$-ary rule; the bound set is a \emph{grid}, and
the dictionary\textquotesingle s first duty is to say which display
constituent lands in which theorem. All statements live over the
compiled behaviour space (\texttt{Ω}, the composite execution
law\textquotesingle s carrier), reached from the presentation through
bridge hypotheses --- the events \texttt{conform j} / \texttt{sat j}
(component $j$\textquotesingle s input-conformance and
guarantee-delivery, interpreted over composite behaviours: the
I-bridge of the kernel-semantics interface) --- so the certified
statements quantify over exactly the Systems whose compilation
discharges those bridges.
\begin{itemize}
\item the E1 premises $\Sigma_i \vDash (\ldots)$ (the
  \texttt{QSatI} judgments) $\mapsto$,
  $n$-ary face: \texttt{h : ∀ j, ∀ᵐ ω ∂μ, ω ∈ conform j → ω ∈ sat j}
  (\texttt{qualitative}) and quantitatively \texttt{hε : μ (conform j
  ∩ (sat j)ᶜ) ≤ ε j} (\texttt{fail\_bound}) --- per-component
  satisfaction as a.e.\ implication / failure-measure bound.
\item the $\mathsf{QWireCert}$ premise $\mapsto$ \texttt{hw : ∀ j,
  conformExt ∩ (⋂ i ∈ S.upstream j, sat i) ⊆ conform j}: the wiring
  squares with the $\kappa$-slack \emph{folded into the predicates}.
  Consequence, stated plainly: the displayed $\kappa_{\mathrm{out}}$
  recompute row has no counterpart in the $n$-ary bound statements
  --- quantitative slack composition appears instead as the
  union-bound $\sum_j \varepsilon_j$ (\texttt{fail\_bound\_out}, the
  inventory\textquotesingle s $\sum_i (1 - p_i)$ display at any $n$).
\item the causality certificate $\mapsto$ \texttt{System.rank} +
  \texttt{System.causality} as \emph{data} (the structure is bound by
  display): rank strictly increasing along wires is a field, the
  induction of \texttt{sat\_of\_conform} runs on it, and a directed
  cycle is unrepresentable --- the Lean face of ``recomputed.''
\item W0--W3 well-posedness rows $\mapsto$ nothing bound ---
  schema-stratum families, honestly absent.
\item the conclusion $\mathsf{SysWire}(w, \vec{\Sigma}) \vDash
  \mathit{spec}_{\mathrm{out}}$ $\mapsto$ \texttt{qualitative}\textquotesingle s
  goal (a.e.\ delivery of \texttt{satOut} on conforming inputs) and
  \texttt{fail\_bound\_out}\textquotesingle s (failure measure at the
  output spec), with \texttt{hout} the output reading of the
  delivered guarantees; \texttt{fail\_subset} is the pure
  set-algebra localization engine (a composite failure localizes at a
  rank-minimal failing component, conforming there).
\item the tier split (the display\textquotesingle s
  admissible/core tier split, stated in the flavor text) $\mapsto$ \texttt{System.qagTier}, a
  \emph{derived function of the presentation} with both unfoldings
  proven (\texttt{qagTier\_admissible\_iff}: every chart ODE or
  discrete; \texttt{qagTier\_core\_iff}: some chart stochastic), and
  \texttt{admissible\_constructs} making ``syntactically
  constructible'' a statement: composite row of $\sum_i \lambda_i
  V_i$ from component dissipation rows + wiring identification rows,
  pointwise, at nonnegative certificate weights and free-sign wiring
  multipliers --- the pointwise shadow of the main
  package\textquotesingle s \texttt{sound\_syzygy} (witness prose,
  package boundary).
\item the binary GSHS cells, premise shapes as bound:
  \texttt{Series.certified} consumes two supermartingale certificates
  (\texttt{Supermartingale Yᵢ 𝒢 μ}, nonneg, positive barriers,
  floor-row transport \texttt{hbarᵢ}, violation-witness inclusions
  \texttt{hwitᵢ}) + one wiring square \texttt{conform₁ ∩ sat₁ ⊆
  conform₂}, concluding the composite failure bound with
  certificate-valued $\varepsilon$\textquotesingle s (initial mass
  over barrier, per component); \texttt{Parallel.certified} /
  \texttt{certified\_pair} conclude the avoid bound for the
  weighted-sum certificate (the semantic face of the syzygy
  weights); \texttt{FeedbackSum.supermartingale} /
  \texttt{avoid\_bound} operate on the bundled
  \texttt{SumFormFeedback} structure --- and its binder rows carry a
  width find: soundness consumes only \texttt{0 ≤ lambda}, \texttt{0
  ≤ theta}; the rule-level ranges $\lambda > 0$, $\theta \in (0,1)$
  are well-formedness, never hypotheses --- with the $(1, \lambda, 1,
  1)$ cancellation entering as the two integrated a.e.\ membership
  rows \texttt{mem₁}/\texttt{mem₂};
  \texttt{sound\_GSHS\_safety\_barrier} is the $\kappa = 0$ max-form
  PDMP instance at a general positive floor \texttt{m}, consuming
  measurability, kernel instances, Dynkin/non-expansion data, and a
  global bound \texttt{C} on the certificate.
\item the B-row (discrete-kernel) assembly
  (\texttt{AGComposeDiscrete}; end-to-end over \texttt{Kernel.traj}
  through the battery substrate, with \emph{no} D-layer hypothesis
  fields --- the semantic hypotheses of the GSHS cells above have no
  counterpart here because the discrete substrate is in Lean): the E1
  premises $\mapsto$ per-component bounded certificates \texttt{V j}
  with datasheet drift on \texttt{conf j} and barrier \texttt{1 ≤ V j}
  on \texttt{bad j} (state-level); the $\mathsf{QWireCert}$ premise
  $\mapsto$ \texttt{hw : ⋂ i ∈ upstream j, (bad i)ᶜ ⊆ conf j} --- the
  square pre-folded at the state stratum, while rank and causality
  stay event-stratum: time-minimality (the first-failure device,
  \texttt{firstFail} / \texttt{measure\_firstFail\_le}, the
  K1-consumption primitive the multi-component setting needs) replaces
  the rank induction, so \texttt{System.rank}/\texttt{causality} are
  honestly unused in \texttt{fail\_bound\_wired}; the conclusion
  $\mapsto$ the failure-measure bound $\sum_j
  \mathrm{ofReal}(V_j\,x)$, the
  $\varepsilon$\textquotesingle s of the event-stratum
  \texttt{fail\_bound} \emph{established} rather than assumed. K2
  (optional stopping) is consumed transitively --- it is bridge step 1
  inside \texttt{ville\_ineq} --- and no direct optional-stopping
  consumption is needed for avoid-shaped conclusions. B2/B3 transpose
  the binary GSHS cells to the kernel level:
  \texttt{drift\_parallelComp} / \texttt{avoid\_parallelComp} (the
  weighted-sum row on \texttt{K₁ ∥ₖ K₂}, the discrete face of the
  parallel cell) and \texttt{smallGain\_drift} /
  \texttt{avoid\_smallGain} (the $(1, \lambda, 1, 1)$ cancellation,
  consuming only \texttt{0 ≤ lam} and \texttt{0 ≤ theta} --- the
  sum-form cell\textquotesingle s width find reproduced at this
  stratum).
\end{itemize}
Width note: the conditional mark\textquotesingle s content is the
grid-vs-rule gap the commentary already states (the deterministic-ODE row at statement level, the sampled-data row stated but not yet assembled, feedback in gauge-form partial; the discrete-kernel row's assembly is discharged, so the gap is one row lighter) --- the dictionary adds
that within the \emph{proven} cells the mappings above are exact, that
the $n$-ary statements are acyclic by type (not by side condition),
and that no bound statement consumes the $\mathcal{K}^0_\infty$
structure of $\kappa$ itself (that lives in
\rulename{QWireCertDischarge}\textquotesingle s bound set; here slack
is pre-folded or union-bounded).

\ruleblock{Lean verifier receipts}
\sourcenote{module \texttt{BernsteinLogicAux.AGCompose.GshsSeries}; bound build \texttt{ff5cb858ef68\ldots} (bound from receipt 911818331ed966fc, 2026-07-20T01:42:13Z). A receipt records that named checking layers accepted a build: (module, compiled-artifact SHA-256, environment SHA-256, checker layers), minted when the checking gate ran; the store receipts attesting the bound build are listed with the layers each run credited. Nothing here asks the reader to trust the store: every layer is an independently runnable program over the quoted sources --- rebuild the module and re-run the checkers to reproduce the verdict. Receipts attesting superseded builds credit no layer and are omitted from this rendering (the count of omissions, if any, is noted below).}
\begin{receipttable}
\receiptrow{cb0f28f9ca6b1c44}{2026-07-19T18:23:37Z}{build+decorrelated}{\badge{BadgeGreen}{current}}
\receiptrow{c973aa90108c94d9}{2026-07-20T01:20:47Z}{build+decorrelated}{\badge{BadgeGreen}{current}}
\receiptrow{911818331ed966fc}{2026-07-20T01:42:13Z}{build+decorrelated}{\badge{BadgeGreen}{current}}
\receiptrow{eada292d784c9848}{2026-07-20T20:46:24Z}{build+decorrelated}{\badge{BadgeGreen}{current}}
\receiptrow{1a58b3dfc83a9cb5}{2026-07-20T20:57:57Z}{build+replay+hardened+decorrelated}{\badge{BadgeGreen}{current}}
\receiptrow{3358b4d03fa77d7c}{2026-07-21T00:02:20Z}{build+decorrelated}{\badge{BadgeGreen}{current}}
\receiptrow{d4f9f9c67009d52f}{2026-07-21T05:38:28Z}{build+replay+hardened+decorrelated}{\badge{BadgeGreen}{current}}
\receiptrow{f86a2ccbe493f911}{2026-07-22T18:26:47Z}{build+decorrelated}{\badge{BadgeGreen}{current}}
\receiptrow{6d17867e1f4f7cdc}{2026-07-31T21:59:37Z}{build+decorrelated}{\badge{BadgeGreen}{current}}
\receiptrow{f4d14ab85f9d9554}{2026-07-31T22:28:31Z}{build+decorrelated}{\badge{BadgeGreen}{current}}
\receiptrow{90881889f02ba9fe}{2026-08-02T23:30:38Z}{build+decorrelated}{\badge{BadgeGreen}{current}}
\receiptrow{684b3dcbc3c2e336}{2026-08-03T16:26:42Z}{build+decorrelated}{\badge{BadgeGreen}{current}}
\receiptrow{782c96e76de07ffa}{2026-08-03T22:45:26Z}{build+decorrelated}{\badge{BadgeGreen}{current}}
\end{receipttable}
\sourcenote{2 further receipt(s) in the store attest superseded builds of this module and are omitted from this rendering.}
\sourcenote{module \texttt{BernsteinLogicAux.AGCompose.GshsParallel}; bound build \texttt{8b8d0ee45b1a\ldots} (hashed from the built olean at bind time). A receipt records that named checking layers accepted a build: (module, compiled-artifact SHA-256, environment SHA-256, checker layers), minted when the checking gate ran; the store receipts attesting the bound build are listed with the layers each run credited. Nothing here asks the reader to trust the store: every layer is an independently runnable program over the quoted sources --- rebuild the module and re-run the checkers to reproduce the verdict. Receipts attesting superseded builds credit no layer and are omitted from this rendering (the count of omissions, if any, is noted below).}
\begin{receipttable}
\receiptrow{f455dcef7549cd27}{2026-07-18T13:22:39Z}{build+decorrelated}{\badge{BadgeGreen}{current}}
\end{receipttable}
\sourcenote{12 further receipt(s) in the store attest superseded builds of this module and are omitted from this rendering.}
\sourcenote{module \texttt{BernsteinLogicAux.GshsFeedbackSum}; bound build \texttt{a0ce53c03e8d\ldots} (bound from receipt 45ff739d5d3e0306, 2026-07-20T07:45:49Z). A receipt records that named checking layers accepted a build: (module, compiled-artifact SHA-256, environment SHA-256, checker layers), minted when the checking gate ran; the store receipts attesting the bound build are listed with the layers each run credited. Nothing here asks the reader to trust the store: every layer is an independently runnable program over the quoted sources --- rebuild the module and re-run the checkers to reproduce the verdict. Receipts attesting superseded builds credit no layer and are omitted from this rendering (the count of omissions, if any, is noted below).}
\begin{receipttable}
\receiptrow{45ff739d5d3e0306}{2026-07-20T07:45:49Z}{build+decorrelated}{\badge{BadgeGreen}{current}}
\receiptrow{eb766d96123685ba}{2026-07-20T20:46:31Z}{build+decorrelated}{\badge{BadgeGreen}{current}}
\receiptrow{f0ef6207375847ca}{2026-07-20T20:58:04Z}{build+replay+hardened+decorrelated}{\badge{BadgeGreen}{current}}
\receiptrow{cf3cd505d8c77450}{2026-07-21T00:02:49Z}{build+decorrelated}{\badge{BadgeGreen}{current}}
\receiptrow{51f1db6c22c527e0}{2026-07-21T05:38:37Z}{build+replay+hardened+decorrelated}{\badge{BadgeGreen}{current}}
\receiptrow{474b9edb3b9dbacd}{2026-07-22T18:26:54Z}{build+decorrelated}{\badge{BadgeGreen}{current}}
\receiptrow{b0859917a7e80e56}{2026-07-31T21:59:48Z}{build+decorrelated}{\badge{BadgeGreen}{current}}
\receiptrow{e0b2e5fdbdb04064}{2026-07-31T22:28:42Z}{build+decorrelated}{\badge{BadgeGreen}{current}}
\receiptrow{57b040e29c83ab7a}{2026-08-02T23:30:50Z}{build+decorrelated}{\badge{BadgeGreen}{current}}
\receiptrow{d425bb06e18d9d85}{2026-08-03T16:26:54Z}{build+decorrelated}{\badge{BadgeGreen}{current}}
\receiptrow{30583333e955b957}{2026-08-03T22:45:37Z}{build+decorrelated}{\badge{BadgeGreen}{current}}
\end{receipttable}
\sourcenote{3 further receipt(s) in the store attest superseded builds of this module and are omitted from this rendering.}
\sourcenote{module \texttt{BernsteinLogicAux.MaxFormPDMP}; bound build \texttt{426c38790122\ldots} (bound from receipt 0a8c057e8cbf9aab, 2026-07-20T01:42:14Z). A receipt records that named checking layers accepted a build: (module, compiled-artifact SHA-256, environment SHA-256, checker layers), minted when the checking gate ran; the store receipts attesting the bound build are listed with the layers each run credited. Nothing here asks the reader to trust the store: every layer is an independently runnable program over the quoted sources --- rebuild the module and re-run the checkers to reproduce the verdict. Receipts attesting superseded builds credit no layer and are omitted from this rendering (the count of omissions, if any, is noted below).}
\begin{receipttable}
\receiptrow{42fadeaf4c4d0d5b}{2026-07-19T18:24:26Z}{build+decorrelated}{\badge{BadgeGreen}{current}}
\receiptrow{81de1969c895bf9f}{2026-07-20T01:21:39Z}{build+decorrelated}{\badge{BadgeGreen}{current}}
\receiptrow{0a8c057e8cbf9aab}{2026-07-20T01:42:14Z}{build+decorrelated}{\badge{BadgeGreen}{current}}
\receiptrow{12151b1f2bcd3a91}{2026-07-20T20:46:35Z}{build+decorrelated}{\badge{BadgeGreen}{current}}
\receiptrow{5821a7ee1f92823e}{2026-07-20T20:58:08Z}{build+replay+hardened+decorrelated}{\badge{BadgeGreen}{current}}
\receiptrow{c1facca742c8a57f}{2026-07-21T00:03:24Z}{build+decorrelated}{\badge{BadgeGreen}{current}}
\receiptrow{b37c4c71b35d3ed7}{2026-07-21T05:38:51Z}{build+replay+hardened+decorrelated}{\badge{BadgeGreen}{current}}
\receiptrow{e32a92fac2e848c6}{2026-07-22T18:27:07Z}{build+decorrelated}{\badge{BadgeGreen}{current}}
\receiptrow{2ed08aae31ac4f48}{2026-07-31T22:00:02Z}{build+decorrelated}{\badge{BadgeGreen}{current}}
\receiptrow{1c0dd63c4cc8d828}{2026-07-31T22:28:56Z}{build+decorrelated}{\badge{BadgeGreen}{current}}
\receiptrow{47a3e55f8aa07fc7}{2026-08-02T23:31:07Z}{build+decorrelated}{\badge{BadgeGreen}{current}}
\receiptrow{5ada5a799af57e12}{2026-08-03T16:27:10Z}{build+decorrelated}{\badge{BadgeGreen}{current}}
\receiptrow{c2de6b6e1b9e676c}{2026-08-03T22:45:54Z}{build+decorrelated}{\badge{BadgeGreen}{current}}
\end{receipttable}
\sourcenote{2 further receipt(s) in the store attest superseded builds of this module and are omitted from this rendering.}
\sourcenote{module \texttt{BernsteinLogicAux.QtslCompose.Presentation}; bound build \texttt{efd8b570583d\ldots} (bound from receipt c2304cb9372c9e10, 2026-07-21T05:39:15Z). A receipt records that named checking layers accepted a build: (module, compiled-artifact SHA-256, environment SHA-256, checker layers), minted when the checking gate ran; the store receipts attesting the bound build are listed with the layers each run credited. Nothing here asks the reader to trust the store: every layer is an independently runnable program over the quoted sources --- rebuild the module and re-run the checkers to reproduce the verdict. Receipts attesting superseded builds credit no layer and are omitted from this rendering (the count of omissions, if any, is noted below).}
\begin{receipttable}
\receiptrow{0d24201fcc020c14}{2026-07-20T07:45:50Z}{build+decorrelated}{\badge{BadgeGreen}{current}}
\receiptrow{c6d195b51ddfe44f}{2026-07-21T00:04:00Z}{build+decorrelated}{\badge{BadgeGreen}{current}}
\receiptrow{c2304cb9372c9e10}{2026-07-21T05:39:15Z}{build+replay+hardened+decorrelated}{\badge{BadgeGreen}{current}}
\end{receipttable}
\sourcenote{11 further receipt(s) in the store attest superseded builds of this module and are omitted from this rendering.}
\sourcenote{module \texttt{BernsteinLogicAux.QtslCompose.AGCompose}; bound build \texttt{88a891d2d053\ldots} (bound from receipt b967c74819cd1516, 2026-07-21T05:39:06Z). A receipt records that named checking layers accepted a build: (module, compiled-artifact SHA-256, environment SHA-256, checker layers), minted when the checking gate ran; the store receipts attesting the bound build are listed with the layers each run credited. Nothing here asks the reader to trust the store: every layer is an independently runnable program over the quoted sources --- rebuild the module and re-run the checkers to reproduce the verdict. Receipts attesting superseded builds credit no layer and are omitted from this rendering (the count of omissions, if any, is noted below).}
\begin{receipttable}
\receiptrow{91f9e1cce3d69343}{2026-07-19T03:38:55Z}{build+decorrelated}{\badge{BadgeGreen}{current}}
\receiptrow{ccc58392dbd207a9}{2026-07-19T18:24:44Z}{build+decorrelated}{\badge{BadgeGreen}{current}}
\receiptrow{875fc85d5bcb76ca}{2026-07-20T01:21:56Z}{build+decorrelated}{\badge{BadgeGreen}{current}}
\receiptrow{23e1df54aafc0165}{2026-07-20T20:46:49Z}{build+decorrelated}{\badge{BadgeGreen}{current}}
\receiptrow{02e33e58a762ea11}{2026-07-20T20:58:22Z}{build+replay+hardened+decorrelated}{\badge{BadgeGreen}{current}}
\receiptrow{9690cee413e67035}{2026-07-21T00:03:48Z}{build+decorrelated}{\badge{BadgeGreen}{current}}
\receiptrow{b967c74819cd1516}{2026-07-21T05:39:06Z}{build+replay+hardened+decorrelated}{\badge{BadgeGreen}{current}}
\end{receipttable}
\sourcenote{6 further receipt(s) in the store attest superseded builds of this module and are omitted from this rendering.}
\sourcenote{module \texttt{BernsteinLogicAux.QtslCompose.AGComposeDiscrete}; bound build \texttt{d5233632f95a\ldots} (bound from receipt 88ff74180adf5b95, 2026-08-03T20:13:53Z). A receipt records that named checking layers accepted a build: (module, compiled-artifact SHA-256, environment SHA-256, checker layers), minted when the checking gate ran; the store receipts attesting the bound build are listed with the layers each run credited. Nothing here asks the reader to trust the store: every layer is an independently runnable program over the quoted sources --- rebuild the module and re-run the checkers to reproduce the verdict. Receipts attesting superseded builds credit no layer and are omitted from this rendering (the count of omissions, if any, is noted below).}
\begin{receipttable}
\receiptrow{88ff74180adf5b95}{2026-08-03T20:13:53Z}{build+decorrelated}{\badge{BadgeGreen}{current}}
\end{receipttable}
\sourcenote{11 further receipt(s) in the store attest superseded builds of this module and are omitted from this rendering.}
\layerline{\layeritem{build}{lake build (elaboration)}{\badge{BadgeGreen}{current}}\layeritem{replay}{leanchecker (official kernel replay)}{\badge{BadgeGreen}{current}}\layeritem{hardened}{safe\_verify (hardened screen)}{\badge{BadgeGreen}{current}}\layeritem{decorrelated}{lean4export + nanoda (decorrelated kernel)}{\badge{BadgeGreen}{current}}}
\sourcenote{\badge{BadgeGray}{comparator: none} no fixture in the three-implementation comparator harness (compiled / kernel-interpreted / reference-interpreted) exercises this rule.}
\ruleblock{HOL4 statement}
\absentbox{NO HOL STATEMENT YET}{The satisfaction-judgment extension postdates the HOL4 port entirely:
the ported cone is the base calculus\textquotesingle s first-order
check-function surface (theory \texttt{bern\_catalogue}, stopping
before even the base proof-object layer), and no extension-theory
definition exists anywhere in the HOL development. A HOL statement of
this rule presupposes a HOL port of the satisfaction stratum, which is
its own campaign; until then this box is the honest state of the
edifice.
}
\ruleblock{HOL verifier receipts}
\absentbox{NONE}{no HOL verification receipt exists for this rule.}
\ruleblock{Proof}

\sourcenote{Genre: core rule (tier 1), with a per-profile tier split
derived from the presentation. The coverage is deliberately
conditional: this section proves the cells the mechanization proves,
names the bridge hypotheses under which they apply, and states what
remains open --- never claiming past the mechanized grid.}

\textbf{Setting and claim.}\enspace A \emph{system} is a structured
presentation: finitely many \emph{mechanisms} (primitive dynamics
declarations --- drift, diffusion, jump, and guard--reset field
families over typed wires), plus \emph{wirings} (variable
identifications) carrying a \emph{causality certificate} --- a rank
function strictly increasing along wires, so a directed cycle is
unrepresentable in the assembled structure. Each mechanism's
\emph{doctrine profile} is derived from which field families are
populated: drift only $=$ ODE; $+$diffusion $=$ SDE; $+$real-time
jumps $=$ general stochastic hybrid (GSHS); counter-time jumps only
$=$ discrete kernel. An interface-indexed satisfaction judgment
$\mathsf{QSatI}(\Sigma, (I, \alpha, \gamma[, \kappa]))$ asserts: every
behaviour of $\Sigma$ whose input trajectory satisfies the assumption
$\alpha$ pointwise delivers the guarantee $\gamma$ (with quantitative
margin per the slack $\kappa$ where declared). A wiring certificate
$\mathsf{QWireCert}(w, (\mathit{spec}_i) \Rightarrow
\mathit{spec}_{\mathrm{out}}, \kappa)$ asserts the lens square
\cite{capuccimyers2026}: delivered upstream guarantees, degraded by at
most $\kappa$, imply downstream assumptions through $w$, and the
output spec reads off the delivered guarantees (the square's discharge
is \rulename{QWireCertDischarge}'s business). The rule: from
$\Sigma_i \vDash \mathit{spec}_i$ for $i = 1, \dots, n$ and
the wiring certificate, conclude
$\mathsf{SysWire}(w, [\Sigma_1, \ldots, \Sigma_n]) \vDash
\mathit{spec}_{\mathrm{out}}$.

The warrant factors as
$\mathrm{sound} = \text{(T)} \circ \text{(S)}$: \emph{transport
arithmetic} --- that composite obligation \emph{rows} follow from
component rows plus wiring rows --- which is profile-generic exact
polynomial arithmetic; and the \emph{semantic reading} --- that the
checked row family implies the behavioural claim --- which is
per-profile, and is where the tier split and every open cell live.

\begin{blemma}[transport is profile-silent]\label{lem:qag-transport}
With residuals $e_i := \mathcal{L}_i V_i + \alpha_i - \sigma_i$ and
membership rows $m_1 := \lambda\sigma_2 - (1{-}\theta)\alpha_1$,
$m_2 := \sigma_1 - \lambda(1{-}\theta)\alpha_2$,
\[
e_1 + \lambda e_2 + m_1 + m_2
\;=\;
\mathcal{L}(V_1 + \lambda V_2) + \theta\,(\alpha_1 + \lambda\alpha_2),
\]
an identity of polynomials --- the cross-supplies cancel by
construction of the multipliers. More generally, for nonnegative
certificate weights $\lambda_i$ and free-sign wiring multipliers
$\nu_k$: if each $e_i \leq 0$ and each wiring identity $g_k = 0$ holds
on the operating region, then
$\sum_i \lambda_i e_i + \sum_k \nu_k g_k \leq 0$ there.
\end{blemma}
\begin{proof}
The displayed identity is a two-line expansion; the only property of
the obligation operator $\mathcal{L}$ it consumes is \emph{linearity
over the wiring} --- which every profile's operator has:
$\nabla(\cdot) \cdot f$ for ODE, one-step conditional expectation
minus identity for kernels, and the extended generator's interior
terms for GSHS. The pointwise fold is
\texttt{admissible\_constructs}: each $\lambda_i e_i \leq 0$
(nonnegative times nonpositive), each $\nu_k g_k = 0$ regardless of
sign, finite sums of nonpositives are nonpositive. At the sequent
layer this is the \rulename{Promise}/\rulename{Syzygy} machinery, and
the small-gain instance is exactly the $(1, \lambda, 1, 1)$-weighted
\rulename{Syzygy} node --- which is why the split of
\autoref{rem:qag-tier} can define ``admissible'' as
\emph{syntactically constructible}: the composite certificate for the
non-stochastic profiles \emph{is} $\sum_i \lambda_i V_i$, and its row
is this lemma. The break in uniformity, wherever it appears below, is
never here: GSHS forced jumps contribute boundary rows, not
$\mathcal{L}$-terms, so part of the composite obligation family
bypasses linearity entirely and must be handled semantically.
\end{proof}

\begin{blemma}[conformance propagation --- the acyclic engine]
\label{lem:qag-conform}
Let a system presentation $S$ have $n$ components, and work in the
composite behaviour space. Write $\mathrm{conform}_j$ for the event
that component $j$'s inputs conform to $\alpha_j$, $\mathrm{sat}_j$
for delivery of $\gamma_j$, and $\mathrm{ext}$ for external-input
conformance. Suppose the wiring squares transport, pointwise:
$\mathrm{ext} \cap \bigcap_{i \in \mathrm{upstream}(j)} \mathrm{sat}_i
\subseteq \mathrm{conform}_j$ for every $j$. If at a behaviour $\omega
\in \mathrm{ext}$ every component delivers on conformance
($\omega \in \mathrm{conform}_j \Rightarrow \omega \in
\mathrm{sat}_j$), then every component delivers at $\omega$.
\end{blemma}
\begin{proof}
Strong induction on the causality rank (\texttt{sat\_of\_conform}):
if every component of rank $< N$ delivers, then for $j$ of rank $< N{+}1$
all of $j$'s upstream components have strictly smaller rank (the
causality certificate), hence deliver; the wiring square then puts
$\omega \in \mathrm{conform}_j$, and the per-component implication
closes. The causality certificate is precisely what makes the
induction well-founded, and its necessity is a two-component
counterexample: components in a cycle, each delivering \emph{only if}
the other does, satisfy every wiring square while both fail --- the
implication chain has no base case. That failure is the classical
unsoundness of naive circular assume--guarantee
\cite{abadilamport1993, mcmillan1999}; this calculus fences it
structurally (cycles unrepresentable here) and readmits feedback only
through the guarded small-gain cells of
Lemma~\ref{lem:qag-feedback}, where the well-foundedness device is the
$\theta$-margin.
\end{proof}

\begin{blemma}[failure decomposition and slack composition]
\label{lem:qag-union}
Under Lemma~\ref{lem:qag-conform}'s wiring hypothesis, a composite
failure under conforming external inputs localizes at some component
that failed \emph{while conforming}:
\[
\mathrm{ext} \cap \Bigl(\bigcap_j \mathrm{sat}_j\Bigr)^{\mathsf{c}}
\;\subseteq\;
\bigcup_j \bigl(\mathrm{conform}_j \cap \mathrm{sat}_j^{\mathsf{c}}\bigr),
\]
hence for any measure with per-component conforming-failure bounds
$\mu(\mathrm{conform}_j \cap \mathrm{sat}_j^{\mathsf{c}}) \leq
\varepsilon_j$, the composite fails with mass at most
$\sum_j \varepsilon_j$ --- the $\sum_j (1 - p_j)$ display, at any $n$.
\end{blemma}
\begin{proof}
Pick a \emph{rank-minimal} failing component $j_m$ (finitely many
components; the failing set is nonempty). Every upstream component of
$j_m$ has strictly smaller rank, hence delivers (minimality), so the
wiring square puts $\omega \in \mathrm{conform}_{j_m}$: a conforming
failure (\texttt{fail\_subset}). Monotonicity and finite subadditivity
of the measure close (\texttt{fail\_bound}; read at an output spec,
\texttt{fail\_bound\_out}). The per-component event
$\mathrm{conform}_j \cap \mathrm{sat}_j^{\mathsf{c}}$ --- not bare
failure --- is forced: component $j$'s certificate priced failure
\emph{only under conforming inputs}; charging it for
non-conforming failures would demand a guarantee nobody certified,
and the rank-minimality argument is exactly what shows the composite
never needs one.
\end{proof}

\begin{blemma}[the stochastic-hybrid cells: series, parallel,
feedback sum-form]\label{lem:qag-feedback}
Over the composite execution law of a GSHS composite, with component
certificates entering as nonnegative supermartingales along the
(stopped) joint jump skeleton, with barrier values on their bad sets
and violation-witness inclusions:
\emph{(i, series)} the conditioning tower composes --- qualitatively,
delivery a.e.\ on conforming inputs chains through the wiring square;
quantitatively, an upstream success mass $\geq p_1$ and a downstream
conditional success probability $\geq p_2$ --- the conditional bound
holding almost everywhere \emph{on the upstream success event}, itself
measurable in the conditioning $\sigma$-algebra --- compose to
composite success mass $\geq p_2 \cdot p_1$ (stronger than the
union-bound shape by exactly $(1-p_1)(1-p_2) \ge 0$);
and the certified pipeline bounds the composite conforming-failure
mass by $\int Y_1(0) / a_1 + \int Y_2(0) / a_2$ through the maximal
inequality.
\emph{(ii, parallel)} for nonnegative weights, the weighted sum
$\sum_i \lambda_i V_i$ of supermartingales w.r.t.\ one (joint)
filtration is a supermartingale --- the semantic face of the syzygy
weights --- and the composite avoid bound follows through the same
maximal inequality.
\emph{(iii, feedback, sum-form)} for two components in feedback whose
certificates satisfy, along the \emph{joint} skeleton, the
between-jumps conditional decrease (pre-jump value plus integrated
assumption-minus-supply increment conditionally below the current
value) and \emph{non-expansive} jump accounting (post-jump value a.e.\
at most pre-jump value), with the integrated membership rows of the
small-gain pair at $\lambda, \theta \geq 0$: the composite
$W = V_1 + \lambda V_2$ is a nonnegative supermartingale for the
joint-skeleton filtration, with the $\theta$-margin banked by the
compensated process, and the composite avoid bound follows.
\end{blemma}
\begin{proof}
(i) is \texttt{AGCompose.Series}: the qualitative form is an a.e.\
implication chain through the pointwise wiring square
(\texttt{qualitative}); the product form integrates the downstream
conditional bound over the upstream event --- the set-integral of the
conditional expectation of the indicator over an upstream-measurable
set is the joint mass (\texttt{tower\_step}, \texttt{tower}); the
certified form (\texttt{certified}) converts each certificate to a
failure bound via the maximal inequality for nonnegative
supermartingales (Ville's inequality, mechanized in the accompanying
development) and composes by the failure decomposition. Exogeneity is
where acyclicity earns its keep semantically: in a cascade each
component's input process is measurable w.r.t.\ upstream trajectories,
so the component's certificate argument runs pathwise in its input.

(ii) is \texttt{AGCompose.Parallel}: induction on the finite index
set, scalar-multiple and sum closure of supermartingales
(\texttt{supermartingale\_weighted\_sum}), then the barrier argument
as in (i) (\texttt{certified}, \texttt{certified\_pair}).

(iii) is the file quoted above (\texttt{SumFormFeedback}): folding
jump non-expansion into the between-jumps decrease gives the
per-component skeleton step; the weighted combination of the two steps
cancels the cross-supplies \emph{under conditional expectation} ---
Lemma~\ref{lem:qag-transport} run pointwise inside the conditional ---
leaving the $\theta$-weighted assumption margin (\texttt{drift});
nonnegativity and the supermartingale property assemble
(\texttt{supermartingale}), and the avoid bound is the maximal
inequality again (\texttt{avoid\_bound}, with the continuous-time
reading conditional on the d\'ebut-sampling inclusion,
\texttt{debut\_avoid\_bound}). Two hypotheses deserve their remark.
\emph{The joint skeleton is not a convenience}: in feedback no
component's input is exogenous --- the jump chain interleaves both
components' guard hits, each component's between-jumps argument is
valid only between \emph{joint} jumps, and forced jumps, having no
intensity, enter only through the jump-accounting field at those
times. \emph{Non-expansion is the sum-form's load-bearing wall}: only
$\lambda, \theta \geq 0$ and the non-expansive accounting are consumed
--- no dwell-time or schedule hypothesis appears anywhere in the
descent, which is a checkable structural fact about the statements ---
and the rule-level ranges $\lambda > 0$, $\theta \in (0,1)$ are
well-formedness backing the published guard-modulus reading, not
soundness needs.
\end{proof}

\begin{bthm}[the $n$-ary acyclic assembly]\label{thm:qag-nary}
Let $S$ be an acyclic system presentation with wiring squares
transporting pointwise as in Lemma~\ref{lem:qag-conform}. If each
component delivers a.e.\ on conformance, the composite delivers its
output spec a.e.\ on external conformance
(\texttt{QtslCompose.qualitative}); and with per-component
conforming-failure bounds $\varepsilon_j$, the composite output-spec
failure mass is at most $\sum_j \varepsilon_j$
(\texttt{fail\_bound\_out}) --- at any $n$, by
Lemmas~\ref{lem:qag-conform} and~\ref{lem:qag-union}.
\end{bthm}
\begin{proof}
Immediate from the two lemmas: the a.e.\ form collects the countable
family of per-component a.e.\ implications, applies the pointwise
propagation on the common full-measure set, and reads the output spec
off the delivered guarantees; the probabilistic form is the failure
decomposition plus subadditivity. The per-profile content enters
entirely through the \emph{bridge hypotheses} --- that the
component-certificate events and their bounds exist over the composite
law is the per-profile discharge of Lemma~\ref{lem:qag-feedback} (GSHS
acyclic cells) or its lighter analogues; this theorem is the
profile-generic event algebra on top.
\end{proof}

\begin{bremark}[the tier split, derived --- never a stored tag]
\label{rem:qag-tier}
The rule node's tier is a function of the presentation
(\texttt{System.qagTier}): \emph{admissible} iff every mechanism's
derived profile admits a syntactically constructible composite
certificate --- proven equivalent to ``every mechanism derives ODE or
discrete'' (\texttt{qagTier\_admissible\_iff}) --- and \emph{core} iff
some mechanism is stochastic (\texttt{qagTier\_core\_iff}). On the
admissible face the node unpacks: the composite certificate is
$\sum_i \lambda_i V_i$ (Lemma~\ref{lem:qag-transport}), transported at
the sequent layer by the \rulename{Promise}/\rulename{Syzygy}
machinery, and the rule adds no semantic content. On the core face no
unpacking exists --- the warrant \emph{is}
Lemmas~\ref{lem:qag-conform}--\ref{lem:qag-feedback} --- which is the
taxonomy's definition of core applied honestly, one rule node either
way. Deriving the tier from field-family population (rather than
storing a system-level tag) is forced by mixed composites: a
sampled-data plant-plus-controller would be ill-typed across its own
components under any stored per-system tag.
\end{bremark}

\begin{bremark}[what is \emph{not} proven --- the open cells, loud]
\label{rem:qag-gaps}
The conditional mark on this rule's dashboard row warrants exactly the
cells above, and the honest remainder is:
\emph{(a)} the deterministic-ODE row's semantic reading is at
statement level (its literature anchor is the certified-lens
composition theorem \cite{capuccimyers2026}; the classical comparison
argument is standard but not yet written in-house);
\emph{(b)} the discrete-kernel row's assembly is discharged ---
$n$-ary acyclic union-bound, parallel weighted-sum, and sum-form
small-gain, end-to-end over the trajectory kernel (the width note's
dictionary names the statements) --- with the \emph{sampled-data} row
the one still stated, not yet assembled (item (d));
\emph{(c)} feedback in \emph{gauge-form} --- jumps that expand the
certificate ($V(g(x)) \leq c\,V(x)$, $c > 1$), repaid by flow decay
under schedule hypotheses --- is partial: the gauge-then-sum reduction
and the machine-checked max-form piecewise-deterministic instance are
landed sorry-free, while the finite-index price of the diffusion term
and the admissible-window pinch remain open research --- though the
reason the diffusion cell was \emph{parked} is retired: within-segment
supermartingale decay for gauged certificates is now available at the
continuous-time martingale-problem profile (the battery
family\textquotesingle s discharge), so what remains is the finite-index
accounting and the window pinch themselves, not a missing stochastic
calculus; sum-form
composites provably do not cover this regime, which is exactly why the
core tier exists. Non-expansive-reset models live entirely in
Lemma~\ref{lem:qag-feedback}(iii), so nothing below gates on (c);
\emph{(d)} the mixed sampled-data row is stated, not proven ---
deterministic throughout, the natural first boundary-row exercise;
\emph{(e)} undirected stochastic wire sharing (common noise) is
\emph{fenced}, not missing: stochastic components sharing a wire need
one determining side, and for fenced models no discharge of the
joint-filtration hypothesis exists --- the lemmas are correctly
unusable there;
\emph{(f)} the $n$-ary \emph{product-form} (tower) composition is
binary only --- its $n$-ary form needs a filtration chain along the
wiring order and is real further work;
\emph{(g)} the judgment sorts the rule concludes into are not yet
staged in the certificate schema, so no end-to-end checkable
certificate exercises the rule yet. Each bridge hypothesis
(the interface layers carrying a component certificate onto the
composite law) is a named hypothesis discharged per model class ---
parameters, never silent assumptions.
\end{bremark}

\begin{bremark}[relation to the literature]
The acyclic discipline is the composing-specifications tradition
\cite{abadilamport1993}, with circularity readmitted only under a
well-foundedness device \cite{mcmillan1999} --- here the
$\theta$-margin of the small-gain pair, whose Lyapunov formulation
descends from the nonlinear small-gain theorems
\cite{jiangteelpraly1994, jiangmareelswang1996}. The parallel cell is
compositional dissipativity \cite{willems1972, moylanhill1978} run at
the satisfaction layer --- the weighted-sum supermartingale is the
stochastic face of the storage-function sum. The wiring square, its
slack algebra, and the certified-lens composition law are
\cite{capuccimyers2026}. What the mechanized grid adds over all of
these is the certificate form: exact-arithmetic transport rows plus
per-profile semantic lemmas with their hypotheses reified as checkable
data, rather than a metatheorem quantified over models nobody
replays.
\end{bremark}

\begin{anchors}
Lemma \ref{lem:qag-transport}: cancellation, pointwise fold & \texttt{admissible\_constructs}; sequent layer \texttt{sound\_promise}, \texttt{sound\_syzygy}, \texttt{sound\_smallGain} (main package)\\
Lemma \ref{lem:qag-conform}: rank induction & \texttt{sat\_of\_conform}; \texttt{System} (fields \texttt{rank}, \texttt{causality})\\
Lemma \ref{lem:qag-union}: rank-minimal localization, union bound & \texttt{fail\_subset}, \texttt{fail\_bound}, \texttt{fail\_bound\_out}\\
Lemma \ref{lem:qag-feedback}(i): series & \texttt{AGCompose.Series.qualitative}, \texttt{tower\_step}, \texttt{tower}, \texttt{fail\_bound}, \texttt{certified}\\
Lemma \ref{lem:qag-feedback}(ii): parallel & \texttt{AGCompose.Parallel.supermartingale\_weighted\_sum}, \texttt{certified}, \texttt{certified\_pair}\\
Lemma \ref{lem:qag-feedback}(iii): feedback sum-form & \texttt{SkeletonCertificate} (fields \texttt{dynkin}, \texttt{jump\_nonexpansive}), \texttt{SkeletonCertificate.step}, \texttt{SumFormFeedback.drift}, \texttt{supermartingale}, \texttt{avoid\_bound}, \texttt{debut\_avoid\_bound}, \texttt{compensated\_supermartingale}, \texttt{margin\_budget}\\
gauge-form partial (max-form instance) & \texttt{sound\_GSHS\_safety\_barrier} (MaxFormPDMP)\\
Theorem \ref{thm:qag-nary} & \texttt{QtslCompose.qualitative}, \texttt{fail\_bound\_out}\\
Remark \ref{rem:qag-tier}: derived split & \texttt{System.qagTier}, \texttt{qagTier\_admissible\_iff}, \texttt{qagTier\_core\_iff}; \texttt{ChartShape.profile}\\
\end{anchors}


\rulesection{79}{QWireCertDischarge}{core rule}
\label{blcatrule-QWireCertDischarge}
\ruleblock{The rule (Steele metanotation)}
\sourcenote{as realized in the mechanized development; source row: \texttt{\texttt{QTSL-rules.toml}} --- \emph{QWireCertDischarge}. The metanotation statement is typeset by hand for notational quality; its mechanized ground truth is the Lean statement below, extracted from the sources at every build of this book and hash-bound. The rule \emph{roster} is machine-checked (every mechanized rule tag round-trips to exactly one catalogue row); the statement \emph{text} is editorial, checked by eye against the extracted Lean beside it.}
\begin{mathpar}
\inferrule*[Right=\rulename{QWireCertDischarge}]{
  \bigl\{R_\ell,\ \varphi(s) \wedge \alpha(v(s), a) \vdash \mathrm{Lift}_s\, \varphi(u(s, a))\bigr\}_{\ell},\ \ \bigl\{R_\ell,\ \varphi(s) \vdash \gamma(v(s))\bigr\}_{\ell} \ \text{(mechanism-certification legs; } s \ \text{a mechanism state, } a \ \text{an admissible input)} \\
  \bigl\{R_\ell \vdash \alpha_2(w(o_1), a_2) + \kappa(\gamma_1(o_1)) \ge \alpha_1(w^{\sharp}(o_1, a_2))\bigr\}_{\ell} \ \text{(assumption legs)} \\
  \bigl\{R_\ell \vdash \gamma_1(o_1) \ge \gamma_2(w(o_1))\bigr\}_{\ell} \ \text{(guarantee legs)} \\
  \vdash \kappa(0) = 0 \ \text{(closed evaluation)} \\
  [0, \infty) \vdash \kappa' \ge 0 \ \text{(derivative-nonnegativity row)} \\
  \vdash \deg \kappa > 0,\ \mathrm{lead}(\kappa) \ge 0 \ \text{(leading-coefficient row)}
}{
  \mathsf{QWireCert}\bigl(w,\ (\mathit{spec}_i) \Rightarrow \mathit{spec}_{\mathrm{out}},\ \kappa\bigr)
}
\end{mathpar}
every leg family indexes over its posted proof leaves $\ell$, and each leaf binds its own region $R_\ell$ as the antecedent of its sequent; the three arithmetic rows certify $\kappa \in \mathcal{K}^0_\infty$ (zero at zero, monotone on the ray, unbounded).

\ruleflavor{$\mathcal{K}^0_\infty$ is the class under which slacks add and lenses compose; the premises are pointwise implications over polynomial data, one inherited BL sequent per posted leaf, the leaf\textquotesingle s antecedent carrying its region --- no region cover is stored on the node, and the obligation that the posted regions jointly cover the consuming composition\textquotesingle s interface domain lives one row up, at the consumption site. The two Capucci--Myers condition families are the mechanism-certification legs (n\'e chart certification) and the wiring-lens legs. No new checker machinery: the rule is leaf assembly, and its semantic content lives one row up in the composition rule.}
\par\medskip
\rulename{QWireCertDischarge} is the wiring-certificate
judgment\textquotesingle s own discharge: it assembles a
$\mathsf{QWireCert}$ --- the license \rulename{QWireCompose} consumes
--- entirely from inherited BL sequent leaves. Four condition families
are checked, each indexing over its posted proof leaves $\ell$, and
each leaf binds its own region $R_\ell$ as the antecedent of its own
sequent. The two \emph{mechanism-certification legs} certify the
component\textquotesingle s contract at its mechanism:
$R_\ell,\ \varphi(s) \wedge \alpha(v(s), a) \vdash
\mathrm{Lift}_s\, \varphi(u(s, a))$ --- the certified invariant,
together with the assumption evaluated at the mechanism\textquotesingle
s readout and an admissible input, is preserved by the update --- and
$R_\ell,\ \varphi(s) \vdash \gamma(v(s))$ --- the invariant delivers
the guarantee at the readout --- where $s$ ranges over the
mechanism\textquotesingle s state carrier and $a$ over its admissible
input carrier. The two \emph{lens legs} certify the wiring square: the
\emph{assumption leg} (delivered upstream guarantees, degraded by at
most the slack $\kappa$, imply the downstream assumption through the
wiring map) and the \emph{guarantee leg} (the upstream guarantee
dominates the composite\textquotesingle s output reading). For
polynomial data all four are pointwise implications --- ordinary
sequents the base calculus
already checks. Three further arithmetic rows certify the
slack\textquotesingle s class membership,
$\kappa \in \mathcal{K}^0_\infty$: the constant coefficient is zero
(closed evaluation), the derivative is nonnegative on the ray (a
derivative-nonnegativity row), and the leading coefficient is positive
at positive degree.

What the cited theorems warrant --- and the scope of this
entry\textquotesingle s conditional check: that the three row families
genuinely certify class membership (a polynomial passing them is zero
at zero, monotone on $[0, \infty)$, and unbounded), and that the class
supports the composition rule\textquotesingle s slack algebra ---
slacks are closed under pointwise sum, closed under composition, the
two-slack contract chain composes to the summed slack, and the
guarantee legs chain along a cascade. That algebra is the
quantitative content of \textquotedblleft compose by lens composition
and by addition of the slacks\textquotedblright.

Each membership row is load-bearing for the algebra. Drop zero-at-zero
and a constant slack certifies a square that degrades even exact
guarantee delivery --- and the nonnegativity that anchors the additive
algebra (a class slack is nonnegative on the ray \emph{because} it is
zero at zero and monotone) fails with it. Drop monotonicity and
composition of slacks fails outright: the outer slack of a two-lens
cascade need no longer bound the chained degradation, since a
non-monotone $\kappa_1$ can reward a worse input. Drop unboundedness
and bounded candidates enter ($\kappa \equiv 0$ passes the other two
rows): the discharge then certifies membership in a class the slack
does not belong to, and the published reading of $\kappa$ as a
reusable comparison-function-style gain is lost. The proof section
walks each closure with the field it consumes.

\cataloguenote{Core-as-glue: the rule is leaf assembly, and the semantic content of
what a wiring certificate \emph{buys} lives one row up, in
\rulename{QWireCompose}\textquotesingle s semantic reading. The
non-strict monotone reading of the class follows the
catalogue\textquotesingle s own row roster (a derivative-nonnegativity
row certifies monotone, not strictly increasing); the zero slack is
thereby not a member --- a slack-free square is a square with the
$\kappa$-term absent, not with $\kappa = 0$. The judgment sort itself
is baked in the extension theory; the mechanized warrant was
proven ahead of the bake, and the square leaves are landed
main-package sequent machinery, cited rather than re-bound here.
}
\ruleblock{Lean statement}
\par\medskip
\textbf{The slack class, bound by display.}\enspace
\texttt{IsClassK0Inf} is the catalogue\textquotesingle s
$\mathcal{K}^0_\infty$: zero at zero, monotone on the nonnegative ray,
unbounded at infinity. Its closure theorems (\texttt{add} --- slacks
add; \texttt{comp} --- lenses compose) are proven in the same module
inside the class namespace; \texttt{polynomial\_discharge} maps the
three checked row families onto the three fields, and
\texttt{slack\_triangle}/\texttt{guarantee\_leg\_comp} are the
quantitative and guarantee-side composition faces.

\declhead{BernsteinLogicAux.QtslCompose.WireCertDischarge}{polynomial\_discharge}{\badge{BadgeGreen}{statement bound}}
\begin{leansrc}
/-- **The polynomial discharge** — the rule note's three membership rows, each
mapped to the class field it certifies:
1. zero-at-zero = closed evaluation (`p.coeff 0 = 0`, checked by evaluating);
2. monotone = the derivative-nonnegativity row (`p' ≥ 0` on the ray);
3. unbounded = positive leading coefficient (with positive degree).
A polynomial passing the three BL row families is genuinely a 𝒦⁰∞ slack. -/
theorem polynomial_discharge (p : Polynomial ℝ)
    (h0 : p.coeff 0 = 0)
    (hderiv : ∀ x : ℝ, 0 ≤ x → 0 ≤ p.derivative.eval x)
    (hdeg : 0 < p.degree) (hlead : 0 ≤ p.leadingCoeff) :
    IsClassK0Inf (fun x => p.eval x)
\end{leansrc}
\declhead{BernsteinLogicAux.QtslCompose.WireCertDischarge}{slack\_triangle}{\badge{BadgeGreen}{statement bound}}
\begin{leansrc}
/-- **The additive slack triangle** (the quantitative face of "compose by
lens composition and by addition of the slacks"): two contract inequalities
with slack compose to one with the summed slack. Stated pointwise over an
arbitrary interface carrier — the per-region reading is `X` = the region. -/
theorem slack_triangle {X : Type*} {f g h : X → ℝ} {c₁ c₂ : ℝ}
    (h₁ : ∀ x, f x ≤ g x + c₁) (h₂ : ∀ x, g x ≤ h x + c₂) :
    ∀ x, f x ≤ h x + (c₁ + c₂)
\end{leansrc}
\declhead{BernsteinLogicAux.QtslCompose.WireCertDischarge}{guarantee\_leg\_comp}{\badge{BadgeGreen}{statement bound}}
\begin{leansrc}
/-- The guarantee-side lens leg composes (the κ-free leg of the CM square:
`γ₁ ≥ γ₂ ∘ w` chains along a cascade). -/
theorem guarantee_leg_comp {A B C : Type*} {w₁ : A → B} {w₂ : B → C}
    {γ₁ : A → ℝ} {γ₂ : B → ℝ} {γ₃ : C → ℝ}
    (h₁ : ∀ a, γ₂ (w₁ a) ≤ γ₁ a) (h₂ : ∀ b, γ₃ (w₂ b) ≤ γ₂ b) :
    ∀ a, γ₃ (w₂ (w₁ a)) ≤ γ₁ a
\end{leansrc}
\declhead{BernsteinLogicAux.QtslCompose.WireCertDischarge}{nonneg}{\badge{BadgeGreen}{statement bound}}
\begin{leansrc}
/-- A 𝒦⁰∞ slack is nonnegative on the nonnegative ray. -/
theorem nonneg (h : IsClassK0Inf κ) {x : ℝ} (hx : 0 ≤ x) : 0 ≤ κ x
\end{leansrc}
\declhead{BernsteinLogicAux.QtslCompose.WireCertDischarge}{add}{\badge{BadgeGreen}{statement bound}}
\begin{leansrc}
/-- **Slacks add** ("by addition of the slacks"): the class is closed under
pointwise sum. -/
theorem add (h₁ : IsClassK0Inf κ₁) (h₂ : IsClassK0Inf κ₂) :
    IsClassK0Inf (fun x => κ₁ x + κ₂ x)
\end{leansrc}
\declhead{BernsteinLogicAux.QtslCompose.WireCertDischarge}{comp}{\badge{BadgeGreen}{statement bound}}
\begin{leansrc}
/-- **Lenses compose**: the class is closed under composition (the slack of a
two-lens cascade). -/
theorem comp (h₁ : IsClassK0Inf κ₁) (h₂ : IsClassK0Inf κ₂) :
    IsClassK0Inf (κ₁ ∘ κ₂)
\end{leansrc}
\declhead{BernsteinLogicAux.QtslCompose.WireCertDischarge}{IsClassK0Inf}{\badge{BadgeGreen}{definition bound}}
\begin{leansrc}
/-- Class 𝒦⁰∞ membership (the `QWireCert` slack class): zero at zero,
monotone on the nonnegative ray, unbounded at infinity. -/
structure IsClassK0Inf (κ : ℝ → ℝ) : Prop where
  zero_at_zero : κ 0 = 0
  mono : MonotoneOn κ (Ici 0)
  unbounded : Tendsto κ atTop atTop
\end{leansrc}
\ruleblock{Correspondence (Steele display $\leftrightarrow$ Lean)}
\badge{BadgeGreen}{reviewer-certified against the displayed texts}
\sourcenote{The dictionary below is a reviewer-certified statement of how the mechanized statements above interpret the displayed rule: display metavariable $\mapsto$ Lean binder, each display hypothesis $\mapsto$ the premise or side condition carrying it, the display conclusion $\mapsto$ the goal --- with any deliberate difference in width between display and mechanization named outright. Certification is recorded as a hash triple over (the displayed cell, the extracted statement texts, this dictionary\textquotesingle s own text); if any of the three changes, this block turns amber until a reviewer re-certifies. The correspondence claim itself remains editorial --- reviewer-grade, not machine-checked; its machine-checked successor types each soundness statement against the rule as a structured object, and supersedes this block per rule as it lands.}
The displayed conclusion $\mathsf{QWireCert}(w, (\mathit{spec}_i)
\Rightarrow \mathit{spec}_{\mathrm{out}}, \kappa)$ is concluded by
\emph{no} bound statement --- deliberately: the judgment sort lives
theory-side (the note), the square leaves are landed sequent machinery
cited as witness prose, and what this bound set certifies is the
\emph{semantic warrant behind the arithmetic rows} plus the class
algebra the rule\textquotesingle s composition claims lean on. The
class itself is bound by display: \texttt{IsClassK0Inf} =
zero-at-zero + \texttt{MonotoneOn κ (Ici 0)} + \texttt{Tendsto κ
atTop atTop} --- $\mathcal{K}^0_\infty$ as the catalogue reads it,
non-strict monotone (the note\textquotesingle s ruling, visible in
the field type).
\begin{itemize}
\item the three membership rows $\mapsto$
  \texttt{polynomial\_discharge}, one hypothesis per row per field:
  closed evaluation $\mapsto$ \texttt{p.coeff 0 = 0} $\Rightarrow$
  \texttt{zero\_at\_zero}; derivative-nonnegativity on the ray
  $\mapsto$ \texttt{hderiv} $\Rightarrow$ \texttt{mono} (through the
  mean-value machinery \texttt{monotoneOn\_of\_deriv\_nonneg});
  leading-coefficient row $\mapsto$ \texttt{0 < p.degree} +
  \texttt{0 ≤ p.leadingCoeff} $\Rightarrow$ \texttt{unbounded}. The
  displayed row asks only $\mathrm{lead}(\kappa) \ge 0$ and so does
  the statement --- nonstrict suffices because a degree-witnessing
  leading coefficient is nonzero by definition; the
  commentary\textquotesingle s ``positive'' is the derived fact, not
  an extra demand.
\item ``slacks add, lenses compose'' $\mapsto$
  \texttt{slack\_triangle} (two contract inequalities chain to the
  summed slack, stated pointwise over an arbitrary carrier --- the
  per-region reading is the carrier choice) and
  \texttt{guarantee\_leg\_comp} (the $\kappa$-free guarantee leg
  chains along a cascade). The class \emph{closure} lemmas
  \texttt{IsClassK0Inf.add} and \texttt{IsClassK0Inf.comp} (with the
  nonnegativity anchor \texttt{IsClassK0Inf.nonneg}) are bound with
  the entry and carry the commentary\textquotesingle s closure claims
  (sum, composition), alongside the triangle and the guarantee leg.
\item polynomiality of $\kappa$ $\mapsto$ \texttt{p : Polynomial ℝ}
  in the discharge theorem only; the class and its algebra are stated
  for arbitrary \texttt{ℝ → ℝ} functions --- wider than the
  display\textquotesingle s polynomial slacks, safe direction.
\end{itemize}
Width note: the four leaf families of the display map onto the baked
node\textquotesingle s two proof lists --- the mechanism-certification
legs (both shapes) ride \texttt{QWireCertDischarge/mech\_rows}, and
the assumption and guarantee lens legs ride
\texttt{QWireCertDischarge/lens\_rows} --- with every leaf\textquotesingle
s own $R_\ell$ as its sequent\textquotesingle s antecedent. In the
bound set the lens legs appear
only through their compositional shadows (\texttt{slack\_triangle},
\texttt{guarantee\_leg\_comp}); the per-leaf discharge of all four
families is
inherited BL sequents, main package, outside this page\textquotesingle s
bindings --- the package-boundary rule, stated in the module header.

\ruleblock{Lean verifier receipts}
\sourcenote{module \texttt{BernsteinLogicAux.QtslCompose.WireCertDischarge}; bound build \texttt{cec5cc1749e1\ldots} (bound from receipt 90dcbae3ace744d7, 2026-07-19T18:24:58Z). A receipt records that named checking layers accepted a build: (module, compiled-artifact SHA-256, environment SHA-256, checker layers), minted when the checking gate ran; the store receipts attesting the bound build are listed with the layers each run credited. Nothing here asks the reader to trust the store: every layer is an independently runnable program over the quoted sources --- rebuild the module and re-run the checkers to reproduce the verdict. Receipts attesting superseded builds credit no layer and are omitted from this rendering (the count of omissions, if any, is noted below).}
\begin{receipttable}
\receiptrow{7aeb2249883cf96c}{2026-07-19T03:38:57Z}{build+decorrelated}{\badge{BadgeGreen}{current}}
\receiptrow{90dcbae3ace744d7}{2026-07-19T18:24:58Z}{build+decorrelated}{\badge{BadgeGreen}{current}}
\receiptrow{164a0a543f6637cb}{2026-07-20T01:22:11Z}{build+decorrelated}{\badge{BadgeGreen}{current}}
\receiptrow{100eb93aa5f473af}{2026-07-20T20:47:01Z}{build+decorrelated}{\badge{BadgeGreen}{current}}
\receiptrow{1f74787842459411}{2026-07-20T20:58:34Z}{build+replay+hardened+decorrelated}{\badge{BadgeGreen}{current}}
\receiptrow{807e980b9db27d7a}{2026-07-21T00:04:07Z}{build+decorrelated}{\badge{BadgeGreen}{current}}
\receiptrow{efa933bf3665661d}{2026-07-21T05:39:20Z}{build+replay+hardened+decorrelated}{\badge{BadgeGreen}{current}}
\receiptrow{64dad6f82e052f30}{2026-07-22T18:27:19Z}{build+decorrelated}{\badge{BadgeGreen}{current}}
\receiptrow{2a1434c39627845c}{2026-07-31T22:00:17Z}{build+decorrelated}{\badge{BadgeGreen}{current}}
\receiptrow{d12854bca5ce90c3}{2026-07-31T22:29:11Z}{build+decorrelated}{\badge{BadgeGreen}{current}}
\receiptrow{b0d9534dd5d7039b}{2026-08-02T23:31:22Z}{build+decorrelated}{\badge{BadgeGreen}{current}}
\receiptrow{bf1f56225eddc2cf}{2026-08-03T16:27:26Z}{build+decorrelated}{\badge{BadgeGreen}{current}}
\receiptrow{e839dbc9244ca3d2}{2026-08-03T22:46:10Z}{build+decorrelated}{\badge{BadgeGreen}{current}}
\end{receipttable}
\layerline{\layeritem{build}{lake build (elaboration)}{\badge{BadgeGreen}{current}}\layeritem{replay}{leanchecker (official kernel replay)}{\badge{BadgeGreen}{current}}\layeritem{hardened}{safe\_verify (hardened screen)}{\badge{BadgeGreen}{current}}\layeritem{decorrelated}{lean4export + nanoda (decorrelated kernel)}{\badge{BadgeGreen}{current}}}
\sourcenote{\badge{BadgeGray}{comparator: none} no fixture in the three-implementation comparator harness (compiled / kernel-interpreted / reference-interpreted) exercises this rule.}
\ruleblock{HOL4 statement}
\absentbox{NO HOL STATEMENT YET}{The satisfaction-judgment extension postdates the HOL4 port entirely:
the ported cone is the base calculus\textquotesingle s check-function
surface (theory \texttt{bern\_catalogue}), and no extension-theory
definition exists in the HOL development --- in particular nothing
shaped like a wiring-certificate discharge. The rule\textquotesingle s
sequent \emph{leaves} would already be expressible in the ported
vocabulary, but the judgment they assemble would not; the box waits on
a HOL port of the satisfaction stratum.
}
\ruleblock{HOL verifier receipts}
\absentbox{NONE}{no HOL verification receipt exists for this rule.}
\ruleblock{Proof}

\sourcenote{Genre: core rule (tier 1), core-as-glue --- the rule
assembles inherited sequent leaves into the wiring-certificate
judgment, and its own mathematical content is deliberately thin: that
the checked arithmetic rows genuinely certify slack-class membership,
and that the class supports the composition algebra the composition
rule recomputes. The section proves exactly that and marks the
assembly gap loudly.}

\textbf{Setting and claim.}\enspace A wiring certificate
$\mathsf{QWireCert}(w, (\mathit{spec}_i) \Rightarrow
\mathit{spec}_{\mathrm{out}}, \kappa)$ carries two per-region families
of pointwise implications --- the lens square's assumption leg
$\alpha_2(w(o_1), a_2) + \kappa(\gamma_1(o_1)) \geq
\alpha_1(w^{\sharp}(o_1, a_2))$ and guarantee leg
$\gamma_1(o_1) \geq \gamma_2(w(o_1))$
\cite{capuccimyers2026} --- which for polynomial data are ordinary
sequents the base calculus checks, plus a slack
$\kappa \in \mathcal{K}^0_\infty$: zero at zero, monotone on
$[0, \infty)$, unbounded at infinity (the \texttt{IsClassK0Inf}
structure quoted above; the non-strict monotone reading is the
catalogued one). The rule's own obligations, proven here: the three
checked rows certify membership for polynomial slacks
(Lemma~\ref{lem:qwc-poly}), and the class supports ``compose by lens
composition and by addition of the slacks''
(Lemma~\ref{lem:qwc-algebra}).

\begin{blemma}[the polynomial discharge]\label{lem:qwc-poly}
Let $p$ be a real polynomial with \emph{(R1)} constant coefficient
zero, \emph{(R2)} derivative nonnegative on $[0, \infty)$, and
\emph{(R3)} positive degree with nonnegative leading coefficient. Then
$x \mapsto p(x)$ is in $\mathcal{K}^0_\infty$.
\end{blemma}
\begin{proof}
(\texttt{polynomial\_discharge}.) Zero at zero: $p(0)$ \emph{is} the
constant coefficient, so (R1) is the field it certifies --- and note
the checker's form of (R1) is one closed exact evaluation, the
ground-fact pattern. Monotone: $[0, \infty)$ is convex and $p$ is
everywhere continuous and differentiable, so nonnegativity of $p'$ on
the interior gives monotonicity on the ray (the mean-value
inequality); (R2) is a single-variable nonnegativity row of exactly
the calculus's cone-checkable shape. Unbounded: a polynomial of
positive degree with nonnegative leading coefficient tends to
$+\infty$ --- (R3) is one syntactic degree check plus one sign row.
Each row maps onto exactly one class field; none is spare.
\end{proof}

\begin{blemma}[the slack algebra]\label{lem:qwc-algebra}
$\mathcal{K}^0_\infty$ is closed under pointwise sum and under
composition; two contract inequalities with constant slacks chain to
one with the summed slack
($f \leq g + c_1$ and $g \leq h + c_2$ pointwise give
$f \leq h + (c_1 + c_2)$); and guarantee legs chain along a cascade
($\gamma_2 \circ w_1 \leq \gamma_1$ and $\gamma_3 \circ w_2 \leq
\gamma_2$ give $\gamma_3 \circ w_2 \circ w_1 \leq \gamma_1$).
\end{blemma}
\begin{proof}
Sum (\texttt{IsClassK0Inf.add}): zeros add; monotonicity is termwise;
for unboundedness, eventually the second summand is nonnegative (it
tends to infinity), so the sum eventually dominates the first, which
tends to infinity. Composition (\texttt{IsClassK0Inf.comp}):
$\kappa_1(\kappa_2(0)) = \kappa_1(0) = 0$; monotone-composed-monotone
is monotone \emph{provided the inner map lands in the outer's domain
of monotonicity} --- which is where nonnegativity of a class slack on
the ray (zero at zero plus monotone) is consumed, not decoration;
unbounded-composed-unbounded is unbounded. The triangle
(\texttt{slack\_triangle}) and the guarantee chain
(\texttt{guarantee\_leg\_comp}) are one-line pointwise transitivity,
stated over an arbitrary interface carrier so the per-region reading
is instantiation, not generalization. In application the
``constants'' are the slack values at the delivered guarantees --- the
triangle is applied pointwise per region --- while composing the slack
\emph{functions} of a cascade is the closure pair's job.
\end{proof}

\begin{bremark}[why each membership row is load-bearing]
\label{rem:qwc-rows}
Drop (R1) and $\kappa(x) = x + 1$ enters the class candidates: a
wiring square then certifies with degradation even on \emph{exact}
upstream delivery ($\kappa(0) = 1$); and the anchor of the algebra ---
nonnegativity of a class slack on the ray, proved from zero-at-zero
plus monotonicity --- fails too, witness $\kappa(x) = x - 1$. Drop (R2) and
$\kappa(x) = x^2 - x$ enters (zero at zero, positive degree, positive
lead): negative on $(0, 1)$, so the nonnegativity anchor fails, and
non-monotone, so composition closure fails at the monotonicity step
(Lemma~\ref{lem:qwc-algebra}'s proof consumes it twice: once as the
outer map's field, once through inner-map nonnegativity), and the
quantitative reading --- a weaker delivered guarantee never
\emph{shrinks} the licensed degradation --- inverts. Drop (R3) and
bounded slacks enter ($\kappa \equiv 0$ passes (R1)--(R2) alone); the
discharge then no longer certifies unboundedness, and $\kappa$ loses
its comparison-function-style gain reading. The zero polynomial is excluded by (R3) deliberately: a
slack-free square is a square with the $\kappa$-term \emph{absent},
never one with $\kappa = 0$, so the additive algebra needs no zero
element.
\end{bremark}

\begin{bremark}[the assembly gap, loud]\label{rem:qwc-gap}
What this section does \emph{not} contain: the square leaves
themselves (landed sequent machinery of the base calculus, cited by
the certificate --- pointwise polynomial implications per region,
checked by the cone machinery), and the end-to-end assembly of the
judgment object, which awaits the staging of the judgment sort in the
certificate schema. The dashboard's conditional mark scopes the
warrant to exactly Lemmas~\ref{lem:qwc-poly}
and~\ref{lem:qwc-algebra}: the arithmetic rows mean what they claim,
and the class composes. The semantic content of what a wiring
certificate \emph{buys} a composite is deliberately not here --- it is
the composition rule's semantic reading, one section over.
\end{bremark}

\begin{anchors}
class definition (bound by display) & \texttt{IsClassK0Inf} (fields \texttt{zero\_at\_zero}, \texttt{mono}, \texttt{unbounded})\\
Lemma \ref{lem:qwc-poly} & \texttt{polynomial\_discharge}\\
Lemma \ref{lem:qwc-algebra}: sum, composition & \texttt{IsClassK0Inf.add}, \texttt{IsClassK0Inf.comp}; nonnegativity anchor \texttt{IsClassK0Inf.nonneg}\\
Lemma \ref{lem:qwc-algebra}: triangle, guarantee chain & \texttt{slack\_triangle}, \texttt{guarantee\_leg\_comp}\\
square legs (cited, not re-proved) & inherited \rulename{BernCone}-family sequents (main package)\\
\end{anchors}


\rulesection{80}{QWireComposeMarks}{core rule}
\label{blcatrule-QWireComposeMarks}
\ruleblock{The rule (Steele metanotation)}
\sourcenote{as realized in the mechanized development; source row: \texttt{\texttt{QTSL-rules.toml}} --- \emph{QWireComposeMarks}. The metanotation statement is typeset by hand for notational quality; its mechanized ground truth is the Lean statement below, extracted from the sources at every build of this book and hash-bound. The rule \emph{roster} is machine-checked (every mechanized rule tag round-trips to exactly one catalogue row); the statement \emph{text} is editorial, checked by eye against the extracted Lean beside it.}
\begin{mathpar}
\inferrule*[Right=\rulename{QWireComposeMarks}]{
  \bigl\{\Sigma_i \vDash (I_i,\ A_i,\ E_i)\bigr\}_{i=1}^{n}
  \quad \text{where } A_i \equiv \textstyle\bigwedge_{j \in \mathrm{in}(i)} \mathrm{ev}(g_j)\ \mathsf{PathPred},
  \ \ E_i \equiv g_i\ \mathsf{ExistsSpec}
  \ \ \text{(conditional delivery: } \forall\text{-conforming-}\exists\text{-delivering, properness presupposed)} \\
  \bigl\{\ell_i\ \text{latency rows}\bigr\}_{i=1}^{n}
  \ \text{(each premise ships its delivery-stage arithmetic + seriality, at its own profile tier)} \\
  \mathsf{QWireCert}\bigl(w,\ (\mathit{spec}_i)_i \Rightarrow \mathit{spec}_{\mathrm{out}},\ -\bigr)
  \qquad
  \bigwedge_i g_i \vdash \mathrm{marks}(\mathit{spec}_{\mathrm{out}})
  \ \text{(mark-entailment cut)} \\
  \mathsf{MarkFlow}\bigl(w,\ (t_k)_k,\ (P_k)_{k \in \mathrm{cyc}_0(w)}\bigr)
  \ \text{(the schedule walk, two columns: stage rows --- difference-constraint feasibility at the labels; carve rows --- containment $P_k \subseteq \mathrm{mark}_k$ + carved-delivery citations + one species-tagged diagonal certificate per zero-lag cycle, exact values as the degenerate species)} \\
  \mathrm{W0}\ \text{(composite properness --- the existing family; the liveness construction\textquotesingle s tier discharge is stepwise, schedule-relative)}
}{
  \mathsf{SysWire}(w, [\Sigma_1, \ldots, \Sigma_n]) \vDash \mathit{spec}_{\mathrm{out}}
}
\end{mathpar}

\ruleflavor{the $\exists$-face sibling of \rulename{QWireCompose}: same wiring premise vocabulary, opposite polarity economics. \textbf{No causality certificate, no ranks, no guard column} --- a dip under a mark survives every extension, so a cyclic wiring conspires \emph{for} the conclusion, and the well-foundedness kit appears precisely in the rule whose conclusion face is $\mathsf{ForallSpec}$, never here. In its place, one composite-level ingredient in two columns: the $\mathsf{MarkFlow}$ walk --- the $\exists$-dual of the causality walk, same arithmetic species (difference constraints over the declaration closure, bounded recompute), opposite verdict role: the $\forall$-walk demands strict increase along every wire, refusing all cycles; this walk demands \emph{consistency}, admitting every cycle whose schedule closes and refusing time-travel. Its stage arithmetic forces every dependency cycle to total lag zero with one shared label --- so the cycles the rule admits are exactly the instantaneous self-supporting loops, and there the walk\textquotesingle s carve column takes over: producer-supplied \emph{carves} (one predicate per cycle mark, presented as a scalar-codomain gate hypernet cited by digest --- its sign-fiber the carve; a plated aggregate gate for token-population members; never a trace, and never a bare point of a continuous state space) checked by containment arithmetic, carved-delivery citations, and one diagonal certificate per cycle asserting the instant loop closes jointly \emph{inside} the carve product --- the discriminant carve-invariance alone provably misses. Exact value pins survive as the degenerate diagonal species, honest exactly where states are finitely presentable. Off the cycles no carve is needed: the $\forall\exists$ premises chain deliveries in topological order, producer-side --- and the premises themselves are the maximal-carve rows, so the whole column is the rule\textquotesingle s own premise form run at tighter predicates. The conclusion is root satisfaction with the demanded marks in the record\textquotesingle s $\mathsf{ExistsSpec}$ list --- one turnstile, no subscript; the face lives in the argument. The joint witness (record, joint law, PDE solution) is constructed once, inside the soundness metatheorem, and never appears on the wire. Guarded premises are not exempt: the composition consumer eliminates each component guard at consumption --- initial membership discharged before the premise is read as a guarantee --- exactly as at \rulename{QWireCompose} (the mirrored elimination row).}
\par\medskip
\rulename{QWireComposeMarks} is the satisfaction stratum\textquotesingle s
composition rule for cyclic wiring at the liveness face --- the polarity
complement of \rulename{QWireCompose}, whose causality certificate makes
a directed cycle unrepresentable. The polarity fact underneath: a safety
violation and a mark realization are both finite-stage events with
opposite persistence. A violation poisons every extension, so circular
reasoning about guarantees-as-floors needs a well-foundedness device; a
dip under a mark survives every extension, so a cycle\textquotesingle s
self-support is productive --- the two-relay loop that refutes the
unguarded circular safety rule is, read at the mark face, exactly the
witness that realizes ``both outputs hit $1$.'' Accordingly this rule
carries \emph{no causality certificate, no ranks, no guards}: from
per-component conditional-delivery judgments (every admitted resolution
whose interface conforms to the assumed in-neighbour marks admits some
resolution through it realizing the own mark --- with a properness
presupposition on the conforming sub-body), their latency rows, the
wiring plumbing, one mark-entailment sequent, the $\mathsf{MarkFlow}$
walk, and the properness family, it concludes satisfaction of the wired
composite at an output spec whose demanded content is its
record\textquotesingle s $\mathsf{ExistsSpec}$ list.

The one composite-level ingredient is the walk, and it runs in two
columns of the same arithmetic species --- bounded recompute over the
declaration closure, never chase derivation, no well-foundedness
anywhere. The \emph{stage column} is the causality walk\textquotesingle s
$\exists$-dual: producer-supplied stage labels per mark, checked against
the difference-constraint system the latency rows induce --- where the
$\forall$-walk demands strict increase along every wire (refusing all
cycles), this walk demands consistency, admitting every cycle whose
schedule closes and refusing time-travel (a relay against an initialized
unit delay demands $t_x = t_y$ and $t_y = t_x + 1$ at once). Its own
arithmetic forces every admitted cycle to total lag zero and one shared
label --- so the admitted cycles are precisely the instantaneous
self-supporting loops. There the \emph{carve column} takes over:
producer-supplied carves --- one \emph{predicate} per cycle mark,
presented as a scalar-codomain gate hypernet cited by digest, its
sign-fiber the carve; a plated aggregate gate where the members are
token-population systems; never a trace, and never a bare point of a
continuous state space --- checked by three row families: containment
(the carve refines its mark), carved-delivery citations of the very
premises the rule already holds (instantiated at the carve events),
and one \emph{diagonal} certificate per cycle, asserting the instant
loop closes jointly inside the carve product. The diagonal is the
load-bearing novelty: carve-invariance alone provably misses it, and
exact value pins survive as its degenerate species --- honest exactly
where states are finitely presentable, which is the reason the column
generalized. The premises themselves are the maximal-carve rows
(machine-checked as an equivalence), so the whole column is the
rule\textquotesingle s own premise form run at tighter predicates ---
a dial from the sketch\textquotesingle s judgments down to points. The
proof section derives both columns from countermodels: delete the
stage column and the relay--delay loop concludes falsely; delete the
diagonal and the swap loop with a swap-closed mark --- the relay
loop\textquotesingle s evil twin, distinguishable by delivered content
alone --- concludes falsely while every premise holds, \emph{including
carve-invariance at every carve}. Both refutations, the refuse/admit
pairs, the discriminant pair, and the subsumption of pins by carves
are machine-checked.

The soundness warrant is a metatheorem whose joint object --- record,
joint law, or solution, per profile --- is constructed once in the
proof and never appears on the wire. The mechanized instance is the
discrete-relational tier: the assembly engine in carved form (carve
rows + schedule + seriality + diagonal closure produce a joint run
realizing every mark at its label, by stage recursion --- the
walk\textquotesingle s rows align every in-neighbour of a marked
component at its own label, where the carve-at-label invariant holds,
so the carved-delivery rows discharge the marked constraints; the
pinned engine is re-derived through it at singleton carves), and the
chaining bridge at both ends of the dial (on acyclic-instant networks
the conditional-delivery premises alone conclude --- via derived
points, and equally via the maximal carve with no diagonal presented
--- so the carve column enters the wire only where a cycle genuinely
closed). The diagonal\textquotesingle s per-species discharge splits
the labor the way the stochastic tier already does: the checker
verifies finite species rows (a witness point, a bounded enumeration,
Lipschitz/gain rows with cycle product below one, convexity or
order-interval flags), and the fixpoint theorem that converts them
into joint solvability --- exhaustion, Banach, Kakutani, Tarski ---
lives in the metatheorem, per profile, proved once. Properness earns a tier
note: the liveness construction consumes \emph{stepwise} properness
(the composite stays live at every stage, given the schedule\textquotesingle s
pins), because bare nonemptiness of the behaviour set is witnessed by
runs the construction cannot use --- blocking components and
pin-sensitive instant subsystems are the countermodels, named in the
proof section.

Two fences are part of the statement. Marks are single-dip
graded-reach demands --- finitely realizable, extension-persistent;
recurrence-shaped liveness (Streett-class acceptance around cycles)
stays outside both cyclic variants. And the rule deliberately trades
completeness for finite presentation: a component delivering only on
its favourite conforming input can witness a joint behaviour that no
conditional-delivery premise certifies --- certifying such composites
without presenting behaviours is the strategy form, future work.

\cataloguenote{The classical anchors sit on the safety side: circular
assume--guarantee is unsound for safety without a well-foundedness
device \cite{abadilamport1993}, and the standard repairs induct over
time \cite{mcmillan1999}. This rule is the observation that the
liveness-mark fragment of the specification language needs no such
device at all --- greatest-fixpoint territory, where the walk replaces
the ranking: schedules must close around cycles (and deliveries must
agree around the instantaneous ones --- agreement certified over
carved predicates, with exact values as the degenerate species) rather
than cycles being forbidden to close. The certificate carries
judgments, stage labels, gate-hypernet digests with their containment,
delivery, and diagonal rows, and one sequent; nothing behaviour-sized.
Mechanization posture, stated loudly: the discrete-relational tier is
proven end-to-end in carved form with the pinned form subsumed (the
Lean cells below, sorry-free, standard axiom closures); stochastic and
PDE instances of the metatheorem, the per-species diagonal devices at
continuous tiers (Banach for the contraction species, Kakutani for the
convex, Tarski for the order-theoretic), the history-dependent
mechanism tier, and the general difference-constraint walk are open,
and the schema realization (an additive sibling node kind beside the
acyclic rule\textquotesingle s, carrying premise citations, latency
rows, the two-column walk as a side condition, and the reused
properness family) is staged for the extension theory\textquotesingle s
next revision, which is why this entry\textquotesingle s dashboard row
shows a catalogued rule ahead of its wire format.
}
\ruleblock{Lean statement}
\par\medskip
\textbf{The discrete-relational tier, bound whole.}\enspace The engine,
the chaining bridge, and the battery --- the hypotheses-necessity
refutations are first-class bound statements, not test scaffolding: the
rule\textquotesingle s two walk columns each carry their forcing
countermodel in the same module that proves the repaired kit sound.

\declhead{BernsteinLogicAux.QtslCompose.AGComposeMarks}{marks\_of\_pinned\_delivery}{\badge{BadgeGreen}{statement bound}}
\begin{leansrc}
/-- **The assembly engine** (the soundness metatheorem's discrete-relational
instance, pinned form): pin rows + schedule walk + seriality + schedule-relative
instant closure produce a joint run realizing every mark at its label. The
joint behaviour exists only under this conclusion's binder. -/
theorem marks_of_pinned_delivery (hser : Serial N) (hfeas : MarkFlowFeasible N t)
    (hpin : PinnedDelivery N mark a) (hclose : InstantClosure N a t) :
    ∃ ω : ℕ → Fin n → V, JointRun N ω ∧ ∀ i, ω (t i) i ∈ mark i
\end{leansrc}
\declhead{BernsteinLogicAux.QtslCompose.AGComposeMarks}{pins\_of\_conditional\_delivery}{\badge{BadgeGreen}{statement bound}}
\begin{leansrc}
/-- **Chaining**: on an acyclic network (any topological order of the
dependency graph), the ∀∃ conditional-delivery premises alone produce a pin —
each component's delivered value is chosen against its in-neighbours' already-
chosen deliveries, in topological order. This is why the premise-shaped rule
needs no value column on the wire in the acyclic case: the pin is derivable,
producer-side, from the premises it already cites. -/
theorem pins_of_conditional_delivery (hloc : DepLocal N) (ord : Fin n → ℕ)
    (hord : ∀ i, ∀ j ∈ N.dep i, ord j < ord i)
    (hdel : ∀ i, ConditionalDelivery N mark i) :
    ∃ a : Fin n → V, PinnedDelivery N mark a
\end{leansrc}
\declhead{BernsteinLogicAux.QtslCompose.AGComposeMarks}{instantClosure\_of\_acyclic}{\badge{BadgeGreen}{statement bound}}
\begin{leansrc}
/-- Acyclicity dissolves the joint instant solve: with a topological order,
seriality alone closes every instant, pins included — coordinate by
coordinate, in order. (Holds for ANY pin vector `a`; the pinned coordinates'
own admission is the pointed rows' business, not this closure's.) -/
theorem instantClosure_of_acyclic (hser : Serial N) (hloc : DepLocal N)
    (ord : Fin n → ℕ) (hord : ∀ i, ∀ j ∈ N.dep i, ord j < ord i)
    (a : Fin n → V) (t : Fin n → ℕ) :
    InstantClosure N a t
\end{leansrc}
\declhead{BernsteinLogicAux.QtslCompose.AGComposeMarks}{marks\_of\_conditional\_delivery\_acyclic}{\badge{BadgeGreen}{statement bound}}
\begin{leansrc}
/-- **The premise-shaped rule, acyclic instance**: on a network with a
topological order, the ∀∃ conditional-delivery premises + the schedule walk +
seriality conclude — no pin appears among the hypotheses (it is derived).
This is the discrete face of the design's claim that the value column is
forced only by instant cycles. -/
theorem marks_of_conditional_delivery_acyclic (hser : Serial N) (hloc : DepLocal N)
    (ord : Fin n → ℕ) (hord : ∀ i, ∀ j ∈ N.dep i, ord j < ord i)
    (hfeas : MarkFlowFeasible N t) (hdel : ∀ i, ConditionalDelivery N mark i) :
    ∃ ω : ℕ → Fin n → V, JointRun N ω ∧ ∀ i, ω (t i) i ∈ mark i
\end{leansrc}
\declhead{BernsteinLogicAux.QtslCompose.AGComposeMarks}{z1\_walk\_infeasible}{\badge{BadgeGreen}{statement bound}}
\begin{leansrc}
/-- The schedule walk REFUSES the relay–delay loop: its rows demand
`t 0 = t 1 ∧ t 1 = t 0 + 1`. -/
theorem z1_walk_infeasible : ¬ ∃ t : Fin 2 → ℕ, MarkFlowFeasible z1Net t
\end{leansrc}
\declhead{BernsteinLogicAux.QtslCompose.AGComposeMarks}{allOnes\_walk\_feasible}{\badge{BadgeGreen}{statement bound}}
\begin{leansrc}
/-- The schedule walk ADMITS the all-ones loop at the shared label `0`. -/
theorem allOnes_walk_feasible : MarkFlowFeasible allOnesNet ![0, 0]
\end{leansrc}
\declhead{BernsteinLogicAux.QtslCompose.AGComposeMarks}{z1\_refutes\_walkless}{\badge{BadgeGreen}{statement bound}}
\begin{leansrc}
/-- **M1 (the walk deleted): the pin-rowed rule WITHOUT the schedule walk is
unsound** — the relay–delay loop satisfies seriality, the pin rows, and the
instant closure, yet realizes no mark on any run at any stage. (The design
memo's forcing argument for MarkFlow, machine-checked; `z1_walk_infeasible`
shows the walk indeed refuses this instance.) -/
theorem z1_refutes_walkless :
    ¬ (∀ (m : ℕ) (W : Type) (_ : Nonempty W) (N : MarkNet m W)
        (mark : Fin m → Set W) (a : Fin m → W) (t : Fin m → ℕ),
        Serial N → PinnedDelivery N mark a → InstantClosure N a t →
        ∃ ω : ℕ → Fin m → W, JointRun N ω ∧ ∀ i, ∃ k, ω k i ∈ mark i)
\end{leansrc}
\declhead{BernsteinLogicAux.QtslCompose.AGComposeMarks}{z5\_refutes\_pinless}{\badge{BadgeGreen}{statement bound}}
\begin{leansrc}
/-- **M2 (the pins deleted): the premise-shaped rule without a value column is
unsound** — conditional delivery (∀∃) at every component, dep-locality,
a feasible schedule, unpinned joint instant closure, AND a nonempty composite
all hold for the swap loop with swap-closed marks, yet no joint run realizes
any mark: each component maps the mark set into the mark set (the premises
are true), but the loop's instant fixpoint lands outside it. The one
composite-level ingredient the stage arithmetic cannot see is the VALUE
consistency — the pin rows. -/
theorem z5_refutes_pinless :
    ¬ (∀ (m : ℕ) (W : Type) (_ : Nonempty W) (N : MarkNet m W)
        (mark : Fin m → Set W) (t : Fin m → ℕ),
        Serial N → DepLocal N → (∀ i, ConditionalDelivery N mark i) →
        MarkFlowFeasible N t → W0Serial N → (∃ ω, JointRun N ω) →
        ∃ ω : ℕ → Fin m → W, JointRun N ω ∧ ∀ i, ∃ k, ω k i ∈ mark i)
\end{leansrc}
\declhead{BernsteinLogicAux.QtslCompose.AGComposeMarks}{z5\_no\_pin}{\badge{BadgeGreen}{statement bound}}
\begin{leansrc}
/-- **The repair refuses Z5**: no pin exists — the pin rows demand a value
assignment with `a 0 = swap3 (a 1)` and `a 1 = a 0` inside the marks, and the
swap has no fixpoint there. (With `z5_refutes_pinless`: the pin rows are
exactly the ingredient that separates the swap loop from the all-ones loop,
which `allOnes_marks_realized` below certifies through the same engine.) -/
theorem z5_no_pin : ¬ ∃ a : Fin 2 → Fin 3, PinnedDelivery z5Net z5Mark a
\end{leansrc}
\declhead{BernsteinLogicAux.QtslCompose.AGComposeMarks}{allOnes\_marks\_realized}{\badge{BadgeGreen}{statement bound}}
\begin{leansrc}
/-- **Positive control**: the assembly engine's full hypothesis kit is
satisfiable on the all-ones relay loop and concludes — both marks realized at
the shared label `0` (the design's Role-2 example, end-to-end: self-support
productive under the ∃-face). -/
theorem allOnes_marks_realized :
    ∃ ω : ℕ → Fin 2 → Bool, JointRun allOnesNet ω ∧
      ∀ i, ω ((![0, 0] : Fin 2 → ℕ) i) i ∈ boolMark i
\end{leansrc}
\par\medskip
\textbf{The carve generalization (r5), bound whole.}\enspace Pins as
predicates: the carved engine at general lags, the subsumption of value
pins as singleton carves (degeneracy as a theorem), the dial identity
(the rule\textquotesingle s $\forall\exists$ premise IS the
maximal-carve row), the acyclic bridge at the maximal carve, the
discriminant pair (carve-invariance holds on the swap loop; no carve
inside the marks closes its diagonal), the carve-level mutants, and the
fat-carve positive control.

\declhead{BernsteinLogicAux.QtslCompose.AGComposeCarves}{marks\_of\_carved\_delivery}{\badge{BadgeGreen}{statement bound}}
\begin{leansrc}
/-- **The generalized assembly engine** (C1): carve rows + schedule walk +
seriality + diagonal closure produce a joint run realizing every mark at its
label. No hypothesis names a value; the joint behaviour — and every delivered
point — exists only under this conclusion's binder. -/
theorem marks_of_carved_delivery (hser : Serial N) (hfeas : MarkFlowFeasible N t)
    (hpin : CarvedDelivery N mark P) (hclose : CarveClosure N P t) :
    ∃ ω : ℕ → Fin n → V, JointRun N ω ∧ ∀ i, ω (t i) i ∈ mark i
\end{leansrc}
\declhead{BernsteinLogicAux.QtslCompose.AGComposeCarves}{carvedDelivery\_of\_pinned}{\badge{BadgeGreen}{statement bound}}
\begin{leansrc}
/-- The r4 pin rows are the singleton instance of the carve rows. -/
theorem carvedDelivery_of_pinned (hpin : PinnedDelivery N mark a) :
    CarvedDelivery N mark (fun i => {a i})
\end{leansrc}
\declhead{BernsteinLogicAux.QtslCompose.AGComposeCarves}{carveClosure\_of\_pinned}{\badge{BadgeGreen}{statement bound}}
\begin{leansrc}
/-- The r4 closure discharges the singleton diagonal. The walk is
load-bearing: every in-neighbour of a pinned instant coordinate is itself
labeled `k` (zero-lag row), so it is either an instant coordinate held at its
pin by the closure or a pipelined coordinate held at its pin by conformance —
and then the pointed row fires. -/
theorem carveClosure_of_pinned (hfeas : MarkFlowFeasible N t)
    (hpin : PinnedDelivery N mark a) (hclose : InstantClosure N a t) :
    CarveClosure N (fun i => {a i}) t
\end{leansrc}
\declhead{BernsteinLogicAux.QtslCompose.AGComposeCarves}{marks\_of\_pinned\_delivery\_via\_carve}{\badge{BadgeGreen}{statement bound}}
\begin{leansrc}
/-- **Degeneracy is a theorem, not a coincidence**: the r4 engine's exact
statement, re-derived through the generalized engine at singleton carves. -/
theorem marks_of_pinned_delivery_via_carve (hser : Serial N)
    (hfeas : MarkFlowFeasible N t) (hpin : PinnedDelivery N mark a)
    (hclose : InstantClosure N a t) :
    ∃ ω : ℕ → Fin n → V, JointRun N ω ∧ ∀ i, ω (t i) i ∈ mark i
\end{leansrc}
\declhead{BernsteinLogicAux.QtslCompose.AGComposeCarves}{carvedDelivery\_mark\_iff\_conditional}{\badge{BadgeGreen}{statement bound}}
\begin{leansrc}
/-- At `P = mark`, the carve rows ARE the conditional-delivery premises: his
sketch's judgments and the r4 pin rows are the two ends of one dial. -/
theorem carvedDelivery_mark_iff_conditional :
    CarvedDelivery N mark mark ↔ ∀ i, ConditionalDelivery N mark i
\end{leansrc}
\declhead{BernsteinLogicAux.QtslCompose.AGComposeCarves}{carveClosure\_of\_acyclic}{\badge{BadgeGreen}{statement bound}}
\begin{leansrc}
/-- Acyclicity dissolves the diagonal at ANY delivery-closed carve: with a
topological order, pinned instant coordinates choose from `F ∩ P` one
`ord`-level at a time (their in-neighbours are all labeled `k` by the walk:
pipelined ones held in-carve by conformance, instant ones settled at earlier
levels), and unpinned coordinates close by seriality as in r4. -/
theorem carveClosure_of_acyclic (hser : Serial N) (hloc : DepLocal N)
    (ord : Fin n → ℕ) (hord : ∀ i, ∀ j ∈ N.dep i, ord j < ord i)
    (hfeas : MarkFlowFeasible N t) (hpin : CarvedDelivery N mark P) :
    CarveClosure N P t
\end{leansrc}
\declhead{BernsteinLogicAux.QtslCompose.AGComposeCarves}{marks\_acyclic\_via\_maximal\_carve}{\badge{BadgeGreen}{statement bound}}
\begin{leansrc}
/-- **The dial's other end, machine-checked** (the second derivation of r4's
acyclic corollary): on acyclic networks, his sketch's ∀∃ premises conclude
THROUGH the carve engine at the MAXIMAL carve `P = mark` — no point derived,
no diagonal presented. With `marks_of_conditional_delivery_acyclic` (r4,
via derived points) this brackets the dial from both ends. -/
theorem marks_acyclic_via_maximal_carve (hser : Serial N) (hloc : DepLocal N)
    (ord : Fin n → ℕ) (hord : ∀ i, ∀ j ∈ N.dep i, ord j < ord i)
    (hfeas : MarkFlowFeasible N t)
    (hdel : ∀ i, ConditionalDelivery N mark i) :
    ∃ ω : ℕ → Fin n → V, JointRun N ω ∧ ∀ i, ω (t i) i ∈ mark i
\end{leansrc}
\declhead{BernsteinLogicAux.QtslCompose.AGComposeCarves}{z5\_carvedDelivery\_mark}{\badge{BadgeGreen}{statement bound}}
\begin{leansrc}
/-- **C5 — invariance is NOT the discriminant**: the swap loop's carve rows
HOLD at the maximal carve (swap maps the mark set into the mark set; the
relay copies it). What Z5 loses is the diagonal, not the delivery. -/
theorem z5_carvedDelivery_mark : CarvedDelivery z5Net z5Mark z5Mark
\end{leansrc}
\declhead{BernsteinLogicAux.QtslCompose.AGComposeCarves}{z5\_no\_carve\_closure}{\badge{BadgeGreen}{statement bound}}
\begin{leansrc}
/-- **C6 — the diagonal IS the discriminant**: NO carve system inside the
marks closes the swap loop's instant diagonal — containment and closure alone
contradict (the delivery rows are not even needed). The r4 refusal
`z5_no_pin` is this fact's singleton shadow. -/
theorem z5_no_carve_closure :
    ¬ ∃ P : Fin 2 → Set (Fin 3), (∀ i, P i ⊆ z5Mark i) ∧
      CarveClosure z5Net P ![0, 0]
\end{leansrc}
\declhead{BernsteinLogicAux.QtslCompose.AGComposeCarves}{z5\_refutes\_closureless}{\badge{BadgeGreen}{statement bound}}
\begin{leansrc}
/-- **M2′ (the diagonal deleted): the carve rule without `CarveClosure` is
unsound** — carve rows at the maximal carve, dep-locality, a feasible
schedule, unpinned joint instant closure, and a nonempty composite all hold
on the swap loop, yet no joint run realizes any mark. The r4 mutant
`z5_refutes_pinless` restated in r5 vocabulary: what the stage arithmetic
cannot see, carve-invariance cannot see either. -/
theorem z5_refutes_closureless :
    ¬ (∀ (m : ℕ) (W : Type) (_ : Nonempty W) (N : MarkNet m W)
        (mark P : Fin m → Set W) (t : Fin m → ℕ),
        Serial N → DepLocal N → CarvedDelivery N mark P →
        MarkFlowFeasible N t → W0Serial N → (∃ ω, JointRun N ω) →
        ∃ ω : ℕ → Fin m → W, JointRun N ω ∧ ∀ i, ∃ k, ω k i ∈ mark i)
\end{leansrc}
\declhead{BernsteinLogicAux.QtslCompose.AGComposeCarves}{z1\_refutes\_walkless\_carve}{\badge{BadgeGreen}{statement bound}}
\begin{leansrc}
/-- **M1′ (the walk deleted): the carve rule without `MarkFlowFeasible` is
unsound** — the relay–delay loop carries carve rows AND a discharging
diagonal at the singleton carve `{true}` (the closure's conformance
hypothesis hands the delay's base value in-carve, and the relay copies it),
yet realizes no mark on any run. The walk is forced at carve level exactly
as at value level. -/
theorem z1_refutes_walkless_carve :
    ¬ (∀ (m : ℕ) (W : Type) (_ : Nonempty W) (N : MarkNet m W)
        (mark P : Fin m → Set W) (t : Fin m → ℕ),
        Serial N → CarvedDelivery N mark P → CarveClosure N P t →
        ∃ ω : ℕ → Fin m → W, JointRun N ω ∧ ∀ i, ∃ k, ω k i ∈ mark i)
\end{leansrc}
\declhead{BernsteinLogicAux.QtslCompose.AGComposeCarves}{fatCarve\_marks\_realized}{\badge{BadgeGreen}{statement bound}}
\begin{leansrc}
/-- **The measure-zero objection, answered at theorem level**: the relay loop
realizes its marks through the engine with carve = mark = `{v ≠ 0}` — a
NON-singleton carve; no hypothesis names a value. (The diagonal closes at a
point the PROOF chooses; the kit never presents one. Z5's swap loop, with the
same carve and the same walk, has no such point — `z5_no_carve_closure`.) -/
theorem fatCarve_marks_realized :
    ∃ ω : ℕ → Fin 2 → Fin 3, JointRun relay3Net ω ∧
      ∀ i, ω ((![0, 0] : Fin 2 → ℕ) i) i ∈ relay3Mark i
\end{leansrc}
\ruleblock{Correspondence (Steele display $\leftrightarrow$ Lean)}
\badge{BadgeAmber}{the Steele display and a bound Lean statement and the dictionary text CHANGED since certification --- re-review and re-certify}
\sourcenote{The dictionary below is a reviewer-certified statement of how the mechanized statements above interpret the displayed rule: display metavariable $\mapsto$ Lean binder, each display hypothesis $\mapsto$ the premise or side condition carrying it, the display conclusion $\mapsto$ the goal --- with any deliberate difference in width between display and mechanization named outright. Certification is recorded as a hash triple over (the displayed cell, the extracted statement texts, this dictionary\textquotesingle s own text); if any of the three changes, this block turns amber until a reviewer re-certifies. The correspondence claim itself remains editorial --- reviewer-grade, not machine-checked; its machine-checked successor types each soundness statement against the rule as a structured object, and supersedes this block per rule as it lands.}
The display is the rule at the satisfaction stratum; the bound set is
its soundness metatheorem\textquotesingle s discrete-relational
instance, over an explicit network model (`MarkNet`: in-neighbour sets,
per-component delivery lags, memoryless admitted-output relations over
wire snapshots, initial segments; `JointRun` the per-stage constraint
satisfaction). The dictionary, display constituent by bound statement:
\begin{itemize}
\item the conditional-delivery premises $\Sigma_i \vDash (I_i, A_i,
  E_i)$ $\mapsto$ \texttt{ConditionalDelivery} (every mark-conforming
  snapshot admits an own-mark output through it) --- the memoryless
  instantiation; the latency rows\textquotesingle{} seriality
  obligation $\mapsto$ \texttt{Serial}; dep-locality (a component
  reads only its declared in-neighbours) $\mapsto$ \texttt{DepLocal}.
\item the $\mathsf{MarkFlow}$ stage column $\mapsto$
  \texttt{MarkFlowFeasible} --- the walk\textquotesingle s Lean
  \emph{specification} at the exact-lag row shape ($t_i = t_j +
  \mathrm{lag}_i$ per in-edge, $\mathrm{lag}_i \leq t_i$); the general
  difference-constraint species (windows, rational stages) is the wire
  format\textquotesingle s and is honestly absent here.
\item the carve column\textquotesingle s row families $\mapsto$
  \texttt{CarvedDelivery} (containment + carved delivery: every
  snapshot with in-neighbour reads inside the carves admits an output
  in the own carve) and \texttt{CarveClosure} (the diagonal: at every
  stage, for every carve-conforming base, the instant subsystem closes
  jointly with labeled coordinates in-carve and admitted) --- the
  latter consumed \emph{abstractly}: its per-species certificate
  dischargers are the wire format\textquotesingle s, its fixpoint
  devices the metatheorem\textquotesingle s. The r4 value column
  $\mapsto$ \texttt{PinnedDelivery} + \texttt{InstantClosure},
  subsumed as the singleton case by
  \texttt{carvedDelivery\_of\_pinned} / \texttt{carveClosure\_of\_pinned} /
  \texttt{marks\_of\_pinned\_delivery\_via\_carve}; the premises
  themselves are the maximal-carve rows,
  \texttt{carvedDelivery\_mark\_iff\_conditional}.
\item the W0-species rows $\mapsto$ \texttt{InstantClosure} /
  \texttt{CarveClosure} (schedule-relative stepwise properness) ---
  deliberately STRONGER than bare nonemptiness, the tier discharge the
  construction consumes; the unpinned form \texttt{W0Serial} appears
  only in the refuted premise-only statements, to make those
  refutations as strong as possible.
\item the conclusion $\mapsto$ the existential in
  \texttt{marks\_of\_carved\_delivery} (engine) /
  \texttt{marks\_of\_conditional\_delivery\_acyclic} +
  \texttt{marks\_acyclic\_via\_maximal\_carve} (both dial ends of the
  acyclic case): some \texttt{JointRun} realizes every mark at its
  label (the joint witness under the binder, never a hypothesis).
\item the countermodel battery, as prereg\textquotesingle d mutants:
  \texttt{z1\_refutes\_walkless} / \texttt{z1\_refutes\_walkless\_carve}
  (stage column deleted, at value and carve level),
  \texttt{z5\_refutes\_pinless} / \texttt{z5\_refutes\_closureless}
  (pin/diagonal column deleted --- against kits with \texttt{W0Serial}
  and composite nonemptiness granted), the discriminant pair
  \texttt{z5\_carvedDelivery\_mark} (invariance holds on the swap
  loop) / \texttt{z5\_no\_carve\_closure} (no carve inside the marks
  closes its diagonal), \texttt{z5\_no\_pin} (its singleton shadow),
  \texttt{z1\_walk\_infeasible} / \texttt{allOnes\_walk\_feasible}
  (the walk\textquotesingle s refuse/admit pair),
  \texttt{allOnes\_marks\_realized} (positive control through the
  pinned engine) + \texttt{allOnes\_direct} (independent witness) +
  \texttt{fatCarve\_marks\_realized} (positive control through the
  carved engine at a non-singleton carve --- no hypothesis names a
  value).
\end{itemize}
Width note: the conditional mark\textquotesingle s content is the
tier gap stated in the commentary --- within the mechanized tier the
mappings above are exact, the engine is proven at general lags, and
the mark-entailment sequent is consumed at its identity instance (the
conclusion realizes the premises\textquotesingle{} own marks; the cut
against a distinct output-spec mark list is sequent-layer machinery,
not this module\textquotesingle s).

\ruleblock{Lean verifier receipts}
\sourcenote{module \texttt{BernsteinLogicAux.QtslCompose.AGComposeMarks}; no attested build yet; source-bound only. The statements above are hash-bound to the source file, but no build attestation exists for this module yet, so no receipt can be judged to attest the quoted build.}
\badge{BadgeAmber}{NO ATTESTED BUILD BOUND for this module}
\sourcenote{7 receipt(s) in the store attest superseded builds of this module and are omitted from this rendering.}
\sourcenote{module \texttt{BernsteinLogicAux.QtslCompose.AGComposeCarves}; no attested build yet; source-bound only. The statements above are hash-bound to the source file, but no build attestation exists for this module yet, so no receipt can be judged to attest the quoted build.}
\badge{BadgeAmber}{NO ATTESTED BUILD BOUND for this module}
\sourcenote{6 receipt(s) in the store attest superseded builds of this module and are omitted from this rendering.}
\layerline{\layeritem{build}{lake build (elaboration)}{\badge{BadgeAmber}{stale only}}\layeritem{replay}{leanchecker (official kernel replay)}{\badge{BadgeGray}{none}}\layeritem{hardened}{safe\_verify (hardened screen)}{\badge{BadgeGray}{none}}\layeritem{decorrelated}{lean4export + nanoda (decorrelated kernel)}{\badge{BadgeAmber}{stale only}}}
\sourcenote{\badge{BadgeGray}{comparator: none} no fixture in the three-implementation comparator harness (compiled / kernel-interpreted / reference-interpreted) exercises this rule.}
\ruleblock{HOL4 statement}
\absentbox{NO HOL STATEMENT YET}{The satisfaction-judgment extension postdates the HOL4 port entirely
(the ported cone is the base calculus\textquotesingle s first-order
check-function surface), and this rule\textquotesingle s wire format is
itself staged ahead of the schema — there is no wire checker for a HOL
statement to mirror yet. The box inherits \rulename{QWireCompose}\textquotesingle s
honest state one rule further out.
}
\ruleblock{HOL verifier receipts}
\absentbox{NONE}{no HOL verification receipt exists for this rule.}
\ruleblock{Proof}

\sourcenote{Genre: core rule (tier 1). The soundness warrant is a
metatheorem whose joint witness never appears on the wire; this section
proves it at the discrete-relational tier --- the tier the mechanization
proves --- and derives every certificate row from a countermodel that
breaks the rule without it. Per-profile instances beyond this tier are
stated as open, never claimed.}

\textbf{Setting and claim.}\enspace The acyclic composition rule's
causality certificate makes a directed cycle unrepresentable; this rule
is its polarity complement, existing precisely to admit cyclic wiring
--- for \emph{liveness marks only}. The polarity fact that makes this
sound where the safety analogue fails: a safety violation and a mark
realization are both finite-stage events with opposite persistence. A
violation of a prefix-closed guarantee poisons every extension, so under
a $\forall$-reading a cyclic wiring can conspire to admit a bad
behaviour and a well-foundedness device is mandatory. A dip under a
mark survives every extension, so under an $\exists$-reading the same
conspiracy works \emph{for} the conclusion: circular self-support is a
post-fixpoint, productive rather than vicious. The rule concludes
$\mathsf{SysWire}(w, [\Sigma_1, \ldots, \Sigma_n]) \vDash
\mathit{spec}_{\mathrm{out}}$ where the output spec's demanded content
is its record's $\mathsf{ExistsSpec}$ list --- one turnstile, no
subscript: the face lives in the argument. Premises: per-component
\emph{conditional delivery} (for every admitted resolution whose
interface conforms to the assumed in-neighbour marks, some admitted
resolution through it realizes the own mark --- the
$\forall$-conforming-$\exists$-delivering reading, with a properness
presupposition on the conforming sub-body), each shipping its
\emph{latency row} (delivery-stage arithmetic at its own profile tier);
the wiring certificate's plumbing; one mark-entailment sequent
$\bigwedge_i g_i \vdash \mathrm{marks}(\mathit{spec}_{\mathrm{out}})$;
the $\mathsf{MarkFlow}$ walk in two columns (below); and the
W0-species properness rows.

\begin{bremark}[the walk, stage column: what breaks without it]
\label{rem:qwcm-z1}
Take a relay ($\mathrm{out} \equiv \mathrm{in}$, instantaneous) closed
against a unit delay initialized to $0$, marks ``hits $1$'' on both
wires. Both conditional-delivery premises are true --- a relay fed a
trace hitting $1$ emits one; a delay shifts a hit by one stage --- and
the composite is proper. But its unique behaviour is all-zeros: the
initialization propagates around the cycle forever, and the conclusion
is false. Premise shape alone is therefore unsound; something
composite-level is forced. The forced ingredient is \emph{schedule
arithmetic}: the producer supplies one stage label $t_k$ per mark, and
the checker verifies the finite difference-constraint system assembled
from the latency rows --- label equality across instantaneous wires,
shifts through delay media. The relay--delay loop is refused exactly
here: its rows demand $t_x = t_y$ and $t_y = t_x + 1$, infeasible ---
while the loop of two relays passes at the shared label $0$. This is
the $\exists$-dual of the causality walk: the same arithmetic species
(difference constraints over the declaration closure, bounded
recompute, never chase derivation), the opposite verdict role --- the
$\forall$-walk demands strict increase along every wire, refusing all
cycles; this walk demands consistency, admitting every cycle whose
schedule closes. Feasibility of difference constraints is a
Bellman--Ford-class check, not an induction device: no
well-foundedness enters. Mechanized: \texttt{z1\_refutes\_walkless}
(every other hypothesis of the pinned engine true, conclusion false),
with the refuse/admit pair \texttt{z1\_walk\_infeasible} /
\texttt{allOnes\_walk\_feasible}.
\end{bremark}

\begin{bremark}[the walk, value column: what breaks without it]
\label{rem:qwcm-z5}
The stage column has a blind spot, and it is exactly where this rule
lives. Around any dependency cycle the latency rows sum, so a feasible
schedule forces every cycle to total lag zero and one shared label:
the admitted cycles are the instantaneous self-supporting loops, whose
marks all land at one instant. There the $\exists$-chosen delivery
\emph{values} must agree around the loop --- and no stage arithmetic
can see values. Witness: a pointwise swap gate over $\{0, 1, 1'\}$
($1 \mapsto 1'$, $1' \mapsto 1$, $0 \mapsto 0$) closed against a relay,
with mark ``hits $\{1, 1'\}$'' on both wires --- a mark set closed
under the swap. Every remaining hypothesis holds: both conditional
deliveries (the swap maps the mark set into itself; the relay copies
it), the schedule ($t_x = t_y = 0$), properness (all-zeros --- $0$ is
the swap's fixpoint), even joint instant seriality. But every joint
run satisfies $x = \mathrm{swap}(x)$ around the loop, hence is
all-zeros, and the conclusion is false: the swap loop is the relay
loop's evil twin, distinguished by values alone. The forced repair is
the walk's \emph{value column}: on each zero-lag cycle the producer
supplies a \emph{pin} $\alpha_k$ per mark --- the delivered value,
spec-representable and label-sized, the value analogue of $t_k$, never
a trace --- and the checker verifies mark membership per pinned wire
plus \emph{pointed delivery}: each pinned component's
conditional-delivery judgment cited at the pins (through every
conforming input agreeing with the neighbouring pins, the own pin is
admitted). The relay loop passes ($\alpha_x = \alpha_y = 1$); the swap
loop is refused ($\alpha_x = \mathrm{swap}(\alpha_y)$,
$\alpha_y = \alpha_x$ has no solution inside the marks). Mechanized:
\texttt{z5\_refutes\_pinless} (the premise-only kit accepts a false
conclusion) and \texttt{z5\_no\_pin} (the pin rows refuse the
witness). Off the cycles no pin is needed --- see the chaining lemma.
An exact value is, however, the wrong general datum --- the next
remark generalizes the column and keeps this one as its degenerate
species.
\end{bremark}

\begin{bremark}[the carve column: pins generalized to predicates]
\label{rem:qwcm-carve}
An exact pin is measure-zero data: honest for a finite alphabet,
unpresentable for members whose wire states are continuous or
infinite-dimensional. The column therefore generalizes: per zero-lag-
cycle mark the producer supplies a \emph{carve} --- a predicate on the
wire\textquotesingle s state carrier, presented as a scalar-codomain
gate hypernet cited by digest (its nonnegativity locus the carve), or
a plated aggregate gate where the member is a token-population system.
The soundness question this raises is the design\textquotesingle s
crux: a point composes trivially around an instant cycle (the joint
assignment \emph{is} the pin vector), while a carve leaves each
delivery free within a set, so joint existence becomes a nonemptiness
obligation over the carve product --- the \emph{diagonal}. Carve-
invariance does not settle it: on the swap loop the carve rows hold at
the maximal carve --- the swap maps the mark set into itself
(\texttt{z5\_carvedDelivery\_mark}) --- yet \emph{no} carve inside the
marks closes the diagonal (\texttt{z5\_no\_carve\_closure}): around
the loop any jointly-admitted value is a swap fixpoint, and the only
one lies outside the marks. So the carve rows come in three families:
\emph{containment} (carve refines mark), \emph{carved delivery} (the
rule\textquotesingle s own conditional-delivery judgment, cited at the
carve events --- the premises are literally the maximal-carve rows,
\texttt{carvedDelivery\_mark\_iff\_conditional}, so the column is one
dial from the sketch\textquotesingle s judgments down to r4\textquotesingle s
points), and one \emph{diagonal certificate} per cycle, species-tagged:
a witness point or bounded enumeration where states are finite (the
degenerate species --- and the swap loop is refused mechanically
there); contraction rows with exact-rational cycle gain below one
(Banach closes the diagonal --- small-gain arithmetic reused at the
$\exists$-face); convexity or order-interval flags (Kakutani, Tarski).
The checker verifies only the species\textquotesingle{} finite rows;
each fixpoint theorem lives in the soundness metatheorem, per profile
--- exactly where Ionescu--Tulcea already lives for the stochastic
assembly. Mechanized at this tier: the carved engine
(\texttt{marks\_of\_carved\_delivery}, the diagonal consumed
abstractly as \texttt{CarveClosure}), the subsumption of pins as
singleton carves (\texttt{marks\_of\_pinned\_delivery\_via\_carve} ---
degeneracy as a theorem), the carve-level mutants
(\texttt{z5\_refutes\_closureless},
\texttt{z1\_refutes\_walkless\_carve}), and a fat-carve positive
control in which no hypothesis names a value
(\texttt{fatCarve\_marks\_realized}).
\end{bremark}

\begin{bremark}[the reading of the premise, the properness
presupposition, and W0: what breaks without each]
\label{rem:qwcm-rest}
Three further hypotheses earn their places by countermodels at the
design tier. (i) Reading the component premise as co-realization
(``some behaviour realizes assumptions and guarantee together'') is
unsound: a component can co-realize only in a mode the closed loop
never enters (a gate that cooperates only on inputs the wiring cannot
produce), so the conditioning must do real work ---
$\forall$-conforming-$\exists$-delivering is the reading under which
the metatheorem is true. (ii) The properness presupposition on each
premise's conforming sub-body blocks vacuous conditioning: a component
admitting no conforming input satisfies the conditional vacuously
while delivering nothing. (iii) The W0-species rows block the empty
composite (an inconsistent loop --- a negation gate against a relay
--- satisfies every safety spec and no liveness spec), and for the
liveness \emph{construction} their tier discharge is stepwise: the
composite must stay live at every stage given the schedule's pins,
because bare nonemptiness of the behaviour set is witnessed by runs
the construction cannot use (a component may block after delivering;
an instant subsystem's solvability may flip on a pinned coordinate's
value). Stepwise properness is what the mechanized engine consumes.
\end{bremark}

\begin{blemma}[the assembly engine --- soundness at the
discrete-relational tier]\label{lem:qwcm-engine}
Fix $n$ components over a wire alphabet, each with in-neighbours, a
delivery lag ($0$ = instantaneous), a memoryless admitted-output
relation reading its lagged snapshot, and an initial segment. Suppose:
(a) seriality (initial sets and output relations nonempty --- the
mechanism-presentation obligation of the latency rows); (b) a feasible
schedule; (c) pins with their membership and pointed-delivery rows;
(d) schedule-relative instant closure (at every stage, with pinned
instantaneous coordinates held at their pins and pipelined coordinates
at any given values, the remaining instantaneous coordinates close
jointly). Then some joint run realizes every mark at its label.

\emph{Proof (as mechanized, \texttt{marks\_of\_pinned\_delivery}).}
Build the run by stage recursion. At stage $k$, pipelined coordinates
choose: their initial value before the pipe fills; \emph{the pin} at
their own label (legal because $t_i \geq \mathrm{lag}_i$ is a walk
row); any admitted output otherwise (seriality). Instantaneous
coordinates close jointly by (d), with the schedule's pinned ones held
at their pins. Two facts discharge the pinned constraints. First, the
\emph{pin invariant}: at its label every coordinate carries its pin ---
by construction for pipelined coordinates, by (d)'s base clause for
instantaneous ones. Second, the walk's alignment: for a marked
component $i$, every in-neighbour $j$ satisfies $t_i = t_j +
\mathrm{lag}_i$, so the snapshot $i$ reads at delivery time is taken at
stage $t_j$ --- where, by the invariant, $j$ carries its pin. The
pointed-delivery row therefore applies verbatim and admits the pin.
Extension past the last label is seriality plus (d) at unpinned
stages. The realized-at-label conclusion is the invariant plus the
membership rows. \emph{Carved form (as mechanized,
\texttt{marks\_of\_carved\_delivery}):} the same recursion with the
invariant weakened from ``carries its pin'' to ``lands inside its
carve'' --- pipelined labeled coordinates now \emph{choose} from the
intersection of their admitted set with their carve (nonempty by the
carved-delivery row, since the walk aligns their in-neighbours at
labels where the invariant holds), and instantaneous labeled
coordinates close through the diagonal certificate, in-carve and
admitted at once. The pinned statement above is re-derived from the
carved engine at singleton carves
(\texttt{marks\_of\_pinned\_delivery\_via\_carve}). \qedhere
\end{blemma}

\begin{blemma}[the chaining bridge --- no pins where no cycle bites]
\label{lem:qwcm-chain}
On a network whose dependency graph admits a topological order, the
conditional-delivery premises alone produce pins satisfying the rows
of Lemma~\ref{lem:qwcm-engine}: choose each component's delivered
value by its premise, applied to the already-chosen deliveries of its
in-neighbours, in topological order; dep-locality (a component reads
only its declared in-neighbours) turns the chosen witness into the
pointed row. Moreover acyclicity dissolves the joint instant solve:
seriality alone furnishes the instant closure, coordinate by
coordinate in order. Hence the \emph{premise-shaped rule with no value
column is sound on acyclic-instant networks}
(\texttt{marks\_of\_conditional\_delivery\_acyclic}) --- the pins are
producer-derivable search data there, entering the wire only where the
walk's own arithmetic says a cycle closed. \emph{Proof (as
mechanized, \texttt{pins\_of\_conditional\_delivery} +
\texttt{instantClosure\_of\_acyclic}).} Induction on topological
levels: at each level, apply the premise at the snapshot assembling
the lower levels' pins (conforming by their membership conjuncts);
transport the chosen admission along dep-locality to every
dep-agreeing snapshot. The instant closure by the same induction, with
seriality supplying unpinned choices. The bridge holds at the dial's
other end too: the same premises \emph{are} the maximal-carve rows,
and acyclicity dissolves the diagonal at any delivery-closed carve
(\texttt{carveClosure\_of\_acyclic}), so the carved engine also
concludes with nothing pinned and nothing carved on the wire
(\texttt{marks\_acyclic\_via\_maximal\_carve}). \qedhere
\end{blemma}

\begin{bremark}[the metatheorem beyond this tier]\label{rem:qwcm-tiers}
The joint object lives in the soundness proof and only there: at the
discrete-relational tier it is the recursion above; at stochastic
tiers the joint law is assembled by Ionescu--Tulcea/Kolmogorov
extension along the wiring's stagewise dependency structure, with the
mark-escaping resolutions selected per the premises; for PDE profiles,
whatever solution object the profile's semantic theorem provides. The
rule's premise kit does not change per profile --- only the bridge
lemmas do, exactly as in the acyclic rule's grid. Those instances, the
per-species diagonal devices at continuous tiers (Banach for the
contraction species, Kakutani/Brouwer for the convex, Knaster--Tarski
for the order-theoretic --- each converting its species' finite rows
into the joint instant solve the carved engine consumes abstractly),
history-dependent mechanisms, and the general difference-constraint
walk (windows, rational stages) are open; the mechanized tier is the
one this section proves, and the conditional mark warrants exactly it.
Two scope fences are part of the statement: marks are single-dip
graded-reach demands (finitely realizable, extension-persistent) ---
recurrence-shaped liveness stays outside; and the rule trades
completeness for finite presentation --- a component delivering only
on its favourite conforming input can witness a joint behaviour no
conditional-delivery premise certifies.
\end{bremark}


\rulesection{81}{QSimTransport}{core rule}
\label{blcatrule-QSimTransport}
\ruleblock{The rule (Steele metanotation)}
\sourcenote{as realized in the mechanized development; source row: \texttt{\texttt{QTSL-rules.toml}} --- \emph{QSimTransport}. The metanotation statement is typeset by hand for notational quality; its mechanized ground truth is the Lean statement below, extracted from the sources at every build of this book and hash-bound. The rule \emph{roster} is machine-checked (every mechanized rule tag round-trips to exactly one catalogue row); the statement \emph{text} is editorial, checked by eye against the extracted Lean beside it.}
\begin{mathpar}
\inferrule*[Right=\rulename{QSimTransport}]{
  S' \vDash \mathit{spec}' \\
  \mathsf{QSim}(\sigma, S, S')
}{
  S \vDash \sigma^{*}\mathit{spec}'
}
\end{mathpar}

\ruleflavor{certification transport along a coarse-graining: $\sigma$ is an exact simulation witness, $S'$ coarse-grains $S$ --- if $S'$ is certified and $\mathsf{QSim}(\sigma, S, S')$ holds ($\sigma$ maps every behaviour of $S$ to a behaviour of $S'$), then $S$ is certified for the pulled-back spec ($\sigma^{*}$ pulls back the whole interface spec, assumptions and guarantees). Expectation specs pull back with the bound intact; region/invariance specs pull back to the preimage region. The former side condition (per-mode polynomial map + compatibility squares as BL leaves) is the kind\textquotesingle s discharge, one node down: \rulename{QSimDischarge} consumes exactly those squares ($\sigma$-strict tier; the kernel tier joins with its own bounded coverage walk), \rulename{QSimCompose} composes witnesses. Exactness is what $\mathsf{QSim}$ asserts --- the approximate $\sigma$-$\varepsilon$ class concludes $\mathsf{QDist}$, never this premise. One-sided functional refinement --- a simulation witness is never a two-sided coupling. satisfaction of an interface-bearing spec is plain \texttt{QSat}, as at \rulename{QWireCompose}.}
\par\medskip
\rulename{QSimTransport} transports certification along a
coarse-graining: if the abstract system $S'$ is certified, and the
simulation judgment $\mathsf{QSim}(\sigma, S, S')$ holds --- its
discharge, one node down, checks the per-mode polynomial
map\textquotesingle s compatibility squares as inherited BL leaves ---
then the concrete system $S$ is certified for the pulled-back spec
$\sigma^{*}\mathit{spec}'$. A producer uses it in the direction the
rule states: certify the abstraction \emph{once} (small state space,
cheap rows), then reuse that certificate under every refinement that
simulates it, paying only the simulation judgment\textquotesingle s
discharge per refinement.

What the cited theorems warrant is three strata. The
\emph{behavioural} form: over compiled behaviour classes, if $\sigma$
maps every behaviour of $S$ to a behaviour of $S'$ (the semantic content
of the squares) and $S'$\textquotesingle s behaviours all satisfy the
spec, then $S$\textquotesingle s behaviours satisfy the pullback --- two
lines, which is the honest size of the rule\textquotesingle s semantic
content; everything substantive in an application is the sequent-level
discharge of the squares. The \emph{finite-kernel} instance, where the
development has machine-checked simulation machinery: a functional simulation
witness at the finite Markov tier is exactly a pushforward-kernel
identity (per state, the pushed concrete kernel row equals the abstract
kernel row at the image state), the marginal-process lemma then
transports laws at every time, and both spec-transport faces follow ---
expectation specs pull back along $\sigma$ with the bound intact, and
region/invariance specs pull back to the preimage region. And the
\emph{general measurable-kernel} tier, mirroring the finite instance
name-for-name over Markov kernels and the Giry monad: the same
pushforward identity (now a \texttt{Measure.map} equation per state),
the same law transport at every horizon, both spec-transport faces ---
with the finite tier\textquotesingle s nonnegativity hypotheses gone,
true by type.

The witness is one-sided by design: a simulation map is functional
refinement, not a two-sided relational coupling. A functional
simulation induces a degenerate coupling, but the transport theorems
never route through it --- they are the pushforward identity plus
summation identities, and the relational tier\textquotesingle s
coupling machinery (joint kernels, absorption hypotheses) is
deliberately not shared. What breaks without the hypotheses: without
the compatibility squares the pushed law and the abstract law decouple
at the first step, and the pullback claims anything; without the
nonnegativity hypotheses of the region form, signed
\textquotedblleft mass\textquotedblright{} could cancel and hide
probability outside the preimage region.

\cataloguenote{The general measurable-kernel tier is discharged (third stratum,
bound below). The simulation premise is the judgment kind
$\mathsf{QSim}$, baked in the extension theory together with its
discharge and composition nodes: the former side condition (map data +
compatibility squares) is that kind\textquotesingle s discharge, so
nothing here renders as a hypothesis mark. Exactness is what the kind
asserts --- the approximate class concludes distance judgments
instead. What remains beyond the three proven strata: the
categorical strengthening under which the composition rule could cite
the lens-composition theorem as a lemma \emph{template} rather than
proving its cells directly --- the direct proofs never pull back along
strict weakenings, which is why the strengthening is not a dependency
of anything on this page. No end-to-end certificate exercises the
rule yet. Related reading, not
dependency: an exchangeability quotient is a simulation map, so
aggregate-space reduction can be \emph{read} as an instance of this
rule --- it stands on its own proof.
}
\ruleblock{Lean statement}
\declhead{BernsteinLogicAux.QtslCompose.SimTransport}{pullback}{\badge{BadgeGreen}{statement bound}}
\begin{leansrc}
/-- **QSimTransport, behavioral form.** `Beh`/`Beh'` are the compiled
behavior classes of the concrete and abstract Systems; `hsim` is the
semantic content of the simulation witness (σ maps behaviors to behaviors);
`hsat` is `QSatI(S', spec')`. Conclusion: `QSatI(S, σ* spec')`. -/
theorem pullback {Ω Ω' : Type*} {σ : Ω → Ω'} {Beh : Set Ω} {Beh' : Set Ω'}
    {spec' : Set Ω'} (hsim : σ '' Beh ⊆ Beh') (hsat : Beh' ⊆ spec') :
    Beh ⊆ σ ⁻¹' spec'
\end{leansrc}
\declhead{BernsteinLogicAux.QtslCompose.SimTransport}{expect\_push}{\badge{BadgeGreen}{statement bound}}
\begin{leansrc}
/-- Observable pullback: the expectation of `g` under a pushed law is the
expectation of the pulled-back observable `g ∘ σ` under the original law.
(The `σ* spec` of the rule display, at expectation-spec granularity.) -/
theorem expect_push (σ : α → β) (μ : α → ℝ) (g : β → ℝ) :
    ∑ b, push σ μ b * g b = ∑ s, μ s * g (σ s)
\end{leansrc}
\declhead{BernsteinLogicAux.QtslCompose.SimTransport}{sim\_expect}{\badge{BadgeGreen}{statement bound}}
\begin{leansrc}
/-- **QSimTransport, expectation form** (finite tier). Along a simulation
witness `PushKernel σ K Q`, the time-`n` expectation of a pulled-back
observable on the concrete chain *equals* the abstract chain's expectation —
`push_law` composed with `expect_push`. -/
theorem sim_expect {σ : α → β} {K : α → α → ℝ} {Q : β → β → ℝ}
    (hσ : PushKernel σ K Q) (μ₀ : α → ℝ) (g : β → ℝ) (n : ℕ) :
    ∑ s, law K μ₀ n s * g (σ s) = ∑ b, law Q (push σ μ₀) n b * g b
\end{leansrc}
\declhead{BernsteinLogicAux.QtslCompose.SimTransport}{sim\_expect\_bound}{\badge{BadgeGreen}{statement bound}}
\begin{leansrc}
/-- Certified expectation bounds transport: a bound proven against the
abstract (coarse) chain holds for the concrete chain at the pulled-back
observable. This is CM's rule 2 at the finite tier, in the direction the
rule states it (certify the abstraction once, reuse under refinement). -/
theorem sim_expect_bound {σ : α → β} {K : α → α → ℝ} {Q : β → β → ℝ}
    (hσ : PushKernel σ K Q) (μ₀ : α → ℝ) (g : β → ℝ) {n : ℕ} {b₀ : ℝ}
    (hb : ∑ b, law Q (push σ μ₀) n b * g b ≤ b₀) :
    ∑ s, law K μ₀ n s * g (σ s) ≤ b₀
\end{leansrc}
\declhead{BernsteinLogicAux.QtslCompose.SimTransport}{sim\_region}{\badge{BadgeGreen}{statement bound}}
\begin{leansrc}
/-- **QSimTransport, region form** (finite tier): if the abstract law at time
`n` puts no mass outside `R'`, the concrete law puts no mass outside
`σ ⁻¹' R'` — invariance specs pull back to the preimage region. -/
theorem sim_region {σ : α → β} {K : α → α → ℝ} {Q : β → β → ℝ}
    (hσ : PushKernel σ K Q) {μ₀ : α → ℝ}
    (hK : ∀ s t, 0 ≤ K s t) (hμ₀ : ∀ s, 0 ≤ μ₀ s) {n : ℕ} {R' : Set β}
    (hR' : ∀ b, b ∉ R' → law Q (push σ μ₀) n b = 0) :
    ∀ s, σ s ∉ R' → law K μ₀ n s = 0
\end{leansrc}
\declhead{BernsteinLogicAux.QtslCompose.SimTransportMeas}{push\_law}{\badge{BadgeGreen}{statement bound}}
\begin{leansrc}
/-- **QSimTransport, measurable-kernel tier — law transport.** Along a
simulation witness, the pushforward of the time-`n` concrete law is the
time-`n` abstract law of the pushed start. (The measurable-tier face of
`QRelChain.push_law`; both spec-transport faces below are corollaries.) -/
theorem push_law {σ : α → β} {K : Kernel α α} {Q : Kernel β β}
    (hσ : Measurable σ) (h : PushKernel σ K Q) (μ : Measure α) :
    ∀ n, (law K μ n).map σ = law Q (μ.map σ) n
  | 0 => rfl
  | n + 1 => by
    rw [law_succ, map_bind _ _ hσ, funext h, ← bind_map _ _ hσ,
      push_law hσ h μ n, law_succ]

/-! ## The two spec-transport faces -/

/-- Expectation specs pull back along `σ` with the bound intact —
`ℝ≥0∞`-observable face, an *equality* as at the finite tier. -/
theorem sim_expect_lintegral {σ : α → β} {K : Kernel α α} {Q : Kernel β β}
    (hσ : Measurable σ) (h : PushKernel σ K Q) (μ : Measure α)
    {g : β → ℝ≥0∞} (hg : Measurable g) (n : ℕ) :
    ∫⁻ s, g (σ s) ∂law K μ n = ∫⁻ b, g b ∂law Q (μ.map σ) n
\end{leansrc}
\declhead{BernsteinLogicAux.QtslCompose.SimTransportMeas}{sim\_expect\_bound}{\badge{BadgeGreen}{statement bound}}
\begin{leansrc}
/-- Certified expectation bounds transport in the direction the rule states:
certify the abstraction once, reuse under refinement. -/
theorem sim_expect_bound {σ : α → β} {K : Kernel α α} {Q : Kernel β β}
    (hσ : Measurable σ) (h : PushKernel σ K Q) (μ : Measure α)
    {g : β → ℝ≥0∞} (hg : Measurable g) {n : ℕ} {b₀ : ℝ≥0∞}
    (hb : ∫⁻ b, g b ∂law Q (μ.map σ) n ≤ b₀) :
    ∫⁻ s, g (σ s) ∂law K μ n ≤ b₀
\end{leansrc}
\declhead{BernsteinLogicAux.QtslCompose.SimTransportMeas}{sim\_region}{\badge{BadgeGreen}{statement bound}}
\begin{leansrc}
/-- Region/invariance specs pull back to the preimage region. The finite
tier's nonnegativity hypotheses (`hK`, `hμ₀` — "signed mass could cancel")
are absent: at this tier they hold by type. -/
theorem sim_region {σ : α → β} {K : Kernel α α} {Q : Kernel β β}
    (hσ : Measurable σ) (h : PushKernel σ K Q) (μ : Measure α)
    {R' : Set β} (hR' : MeasurableSet R') {n : ℕ}
    (hz : law Q (μ.map σ) n R'ᶜ = 0) :
    law K μ n (σ ⁻¹' R')ᶜ = 0
\end{leansrc}
\declhead{BernsteinLogicAux.QtslCompose.SimTransportMeas}{pushKernel\_iff\_comp}{\badge{BadgeGreen}{statement bound}}
\begin{leansrc}
/-- A functional simulation witness is exactly commutation of the step
kernels with the deterministic kernel of `σ`: the square the lens/fibred
strengthening would functor over. -/
theorem pushKernel_iff_comp {σ : α → β} (hσ : Measurable σ)
    {K : Kernel α α} {Q : Kernel β β} :
    PushKernel σ K Q ↔
      Kernel.deterministic σ hσ ∘ₖ K = Q ∘ₖ Kernel.deterministic σ hσ
\end{leansrc}
\ruleblock{Correspondence (Steele display $\leftrightarrow$ Lean)}
\badge{BadgeGreen}{reviewer-certified against the displayed texts}
\sourcenote{The dictionary below is a reviewer-certified statement of how the mechanized statements above interpret the displayed rule: display metavariable $\mapsto$ Lean binder, each display hypothesis $\mapsto$ the premise or side condition carrying it, the display conclusion $\mapsto$ the goal --- with any deliberate difference in width between display and mechanization named outright. Certification is recorded as a hash triple over (the displayed cell, the extracted statement texts, this dictionary\textquotesingle s own text); if any of the three changes, this block turns amber until a reviewer re-certifies. The correspondence claim itself remains editorial --- reviewer-grade, not machine-checked; its machine-checked successor types each soundness statement against the rule as a structured object, and supersedes this block per rule as it lands.}
Three strata, as the commentary\textquotesingle s conditional scope
states; the display maps onto each at its own granularity.
\emph{Behavioural}: the displayed conclusion $S \vDash
\sigma^{*}\mathit{spec}'$ $\mapsto$ \texttt{pullback}\textquotesingle s
goal \texttt{Beh ⊆ σ ⁻¹\textquotesingle{} spec\textquotesingle} over
compiled behaviour classes; the E1 premise $S' \vDash
\mathit{spec}'$ $\mapsto$ \texttt{hsat : Beh\textquotesingle{} ⊆
spec\textquotesingle}; the judgment premise
$\mathsf{QSim}(\sigma, S, S')$ $\mapsto$
\texttt{hsim : σ \textquotesingle\textquotesingle{} Beh ⊆
Beh\textquotesingle} --- the kind\textquotesingle s \emph{semantic
content} exactly: the per-mode polynomial structure of $\sigma$ is
erased (\texttt{σ} is any function), and the BL discharge of the
compatibility squares lives in the kind\textquotesingle s own
discharge node, outside this bound set. The
statement is axiom-free set algebra --- two lines, the honest size of
the rule\textquotesingle s semantic content.
\begin{itemize}
\item \emph{finite-kernel instance}: the simulation witness $\mapsto$
  \texttt{PushKernel σ K Q} --- the per-state marginal identity
  \texttt{push σ (K s) = Q (σ s)}, which is what a functional
  simulation IS at this tier.
\item expectation specs pull back ``with the bound intact'' $\mapsto$
  \texttt{sim\_expect}, and the certified statement is
  \emph{stronger} than the display\textquotesingle s reading: an
  \emph{equality} of expectations at every horizon (\texttt{∑ s, law
  K μ₀ n s * g (σ s) = ∑ b, law Q (push σ μ₀) n b * g b}), with
  \texttt{sim\_expect\_bound} the displayed inequality direction as
  its corollary (certify the abstraction once, reuse under
  refinement).
\item region/invariance specs pull back to the preimage $\mapsto$
  \texttt{sim\_region}: abstract law puts no mass outside
  $R'$ $\Rightarrow$ concrete law puts none outside
  $\sigma^{-1} R'$ --- and this face alone consumes
  nonnegativity (\texttt{hK}, \texttt{hμ₀}): the commentary\textquotesingle s
  ``signed mass could cancel'' hypothesis, visible as exactly two
  binder rows.
\item \texttt{expect\_push} is the observable-pullback engine
  ($\sigma^{*}$ at expectation granularity), fibre-summation algebra.
\end{itemize}
\emph{Measurable-kernel tier} (second bound block, module
\texttt{SimTransportMeas}, namespace \texttt{…QtslCompose.Meas}): the
simulation witness $\mapsto$
\texttt{Meas.PushKernel σ K Q} --- per state, \texttt{(K s).map σ = Q
(σ s)} over \texttt{ProbabilityTheory.Kernel}s; \texttt{push\_law} is
the law transport at every horizon over the \texttt{n}-fold Giry bind;
\texttt{sim\_expect\_bound} and \texttt{sim\_region} are the two
spec-transport faces in measure vocabulary (the addressability
equalities \texttt{sim\_expect\_lintegral}/\texttt{\_integral} live in
the module, cited not bound); and \texttt{pushKernel\_iff\_comp} is
the intertwining square --- \texttt{Kernel.deterministic σ hσ ∘ₖ K = Q
∘ₖ Kernel.deterministic σ hσ}, push-after-step equals step-after-push
--- the shape the lens/fibred strengthening would functor over, so
that residue names a bound \emph{object}, not a blank. The finite
tier\textquotesingle s nonnegativity binders (\texttt{hK},
\texttt{hμ₀}) vanish at this tier by type, which is the ``signed mass
could cancel'' clause discharged by construction.
Width note: one-sidedness is structural in the bound statements ---
nothing consumes or constructs a joint kernel; \texttt{push\_law} (the
shared marginal-process lemma) is genuinely the same lemma the
relational tier rides, at both tiers, while the coupling machinery is
deliberately absent from this bound set. The conditional mark is
dropped: the former side condition is the baked $\mathsf{QSim}$
premise; what remains named --- the lex/fibred (categorical)
strengthening --- is a dependency of nothing on this page.

\ruleblock{Lean verifier receipts}
\sourcenote{module \texttt{BernsteinLogicAux.QtslCompose.SimTransport}; bound build \texttt{b2a09651f9a0\ldots} (bound from receipt dc0250279a3a2d0d, 2026-07-19T18:24:55Z). A receipt records that named checking layers accepted a build: (module, compiled-artifact SHA-256, environment SHA-256, checker layers), minted when the checking gate ran; the store receipts attesting the bound build are listed with the layers each run credited. Nothing here asks the reader to trust the store: every layer is an independently runnable program over the quoted sources --- rebuild the module and re-run the checkers to reproduce the verdict. Receipts attesting superseded builds credit no layer and are omitted from this rendering (the count of omissions, if any, is noted below).}
\begin{receipttable}
\receiptrow{edfcab6a34d8774b}{2026-07-19T03:38:57Z}{build+decorrelated}{\badge{BadgeGreen}{current}}
\receiptrow{dc0250279a3a2d0d}{2026-07-19T18:24:55Z}{build+decorrelated}{\badge{BadgeGreen}{current}}
\receiptrow{714abbd46c61ec9c}{2026-07-20T01:22:07Z}{build+decorrelated}{\badge{BadgeGreen}{current}}
\receiptrow{28aeb5d1779c6cdc}{2026-07-20T20:46:58Z}{build+decorrelated}{\badge{BadgeGreen}{current}}
\receiptrow{1bb39a35738bbfe3}{2026-07-20T20:58:31Z}{build+replay+hardened+decorrelated}{\badge{BadgeGreen}{current}}
\receiptrow{1ae45ad709a68e43}{2026-07-21T00:04:02Z}{build+decorrelated}{\badge{BadgeGreen}{current}}
\receiptrow{5b96bfe1e468ce59}{2026-07-21T05:39:17Z}{build+replay+hardened+decorrelated}{\badge{BadgeGreen}{current}}
\receiptrow{016e87c2921cd31d}{2026-07-22T18:27:17Z}{build+decorrelated}{\badge{BadgeGreen}{current}}
\receiptrow{e776a27adcccb860}{2026-07-31T22:00:16Z}{build+decorrelated}{\badge{BadgeGreen}{current}}
\receiptrow{cf62752ca2ea53ab}{2026-07-31T22:29:10Z}{build+decorrelated}{\badge{BadgeGreen}{current}}
\receiptrow{3c2ed3fc0dce1ac5}{2026-08-02T23:31:21Z}{build+decorrelated}{\badge{BadgeGreen}{current}}
\receiptrow{6121f42e978829a4}{2026-08-03T16:27:25Z}{build+decorrelated}{\badge{BadgeGreen}{current}}
\receiptrow{3a4a61220d562456}{2026-08-03T22:46:09Z}{build+decorrelated}{\badge{BadgeGreen}{current}}
\end{receipttable}
\sourcenote{module \texttt{BernsteinLogicAux.QtslCompose.SimTransportMeas}; bound build \texttt{9530b0b3c465\ldots} (bound from receipt 3c92f1e3a2986972, 2026-07-20T08:24:45Z). A receipt records that named checking layers accepted a build: (module, compiled-artifact SHA-256, environment SHA-256, checker layers), minted when the checking gate ran; the store receipts attesting the bound build are listed with the layers each run credited. Nothing here asks the reader to trust the store: every layer is an independently runnable program over the quoted sources --- rebuild the module and re-run the checkers to reproduce the verdict. Receipts attesting superseded builds credit no layer and are omitted from this rendering (the count of omissions, if any, is noted below).}
\begin{receipttable}
\receiptrow{3c92f1e3a2986972}{2026-07-20T08:24:45Z}{build+decorrelated}{\badge{BadgeGreen}{current}}
\receiptrow{06e7f179212d370d}{2026-07-20T20:47:00Z}{build+decorrelated}{\badge{BadgeGreen}{current}}
\receiptrow{c314290e315215f2}{2026-07-20T20:58:33Z}{build+replay+hardened+decorrelated}{\badge{BadgeGreen}{current}}
\receiptrow{c471feb792abf799}{2026-07-21T00:04:04Z}{build+decorrelated}{\badge{BadgeGreen}{current}}
\receiptrow{b18bbdfab8a60463}{2026-07-21T05:39:18Z}{build+replay+hardened+decorrelated}{\badge{BadgeGreen}{current}}
\receiptrow{16821f18670ae2eb}{2026-07-22T18:27:18Z}{build+decorrelated}{\badge{BadgeGreen}{current}}
\receiptrow{fb024d6489c5da75}{2026-07-31T22:00:16Z}{build+decorrelated}{\badge{BadgeGreen}{current}}
\receiptrow{ed5047f4fc559336}{2026-07-31T22:29:10Z}{build+decorrelated}{\badge{BadgeGreen}{current}}
\receiptrow{e618e69852b924d8}{2026-08-02T23:31:21Z}{build+decorrelated}{\badge{BadgeGreen}{current}}
\receiptrow{19820df5b08bd97a}{2026-08-03T16:27:25Z}{build+decorrelated}{\badge{BadgeGreen}{current}}
\receiptrow{682c99d152518997}{2026-08-03T22:46:09Z}{build+decorrelated}{\badge{BadgeGreen}{current}}
\end{receipttable}
\layerline{\layeritem{build}{lake build (elaboration)}{\badge{BadgeGreen}{current}}\layeritem{replay}{leanchecker (official kernel replay)}{\badge{BadgeGreen}{current}}\layeritem{hardened}{safe\_verify (hardened screen)}{\badge{BadgeGreen}{current}}\layeritem{decorrelated}{lean4export + nanoda (decorrelated kernel)}{\badge{BadgeGreen}{current}}}
\sourcenote{\badge{BadgeGray}{comparator: none} no fixture in the three-implementation comparator harness (compiled / kernel-interpreted / reference-interpreted) exercises this rule.}
\ruleblock{HOL4 statement}
\absentbox{NO HOL STATEMENT YET}{The satisfaction-judgment extension postdates the HOL4 port entirely:
the ported cone is the base calculus\textquotesingle s check-function
surface (theory \texttt{bern\_catalogue}), and no extension-theory
definition --- nothing shaped like simulation transport --- exists in
the HOL development. The box waits on a HOL port of the satisfaction
stratum, which is its own campaign.
}
\ruleblock{HOL verifier receipts}
\absentbox{NONE}{no HOL verification receipt exists for this rule.}
\ruleblock{Proof}

\sourcenote{Genre: core rule (tier 1). The rule's semantic content is
honestly small --- the behavioural pullback is two lines --- and the
section says so rather than inflating it; the substance of any
application is the sequent-level discharge of the compatibility
squares, plus the finite-kernel instance where the transport is proven
against actual chain machinery. The strata beyond the two walked here,
and the genuine gaps, are marked loudly.}

\textbf{Setting and claim.}\enspace A simulation witness
$\sigma : S \to S'$ is a per-mode polynomial map whose
\emph{compatibility squares} --- commutation of $\sigma$ with the two
systems' dynamics declarations, mode by mode --- are checked as
inherited sequent leaves; the rule consumes the witness as the
judgment $\mathsf{QSim}(\sigma, S, S')$, whose own discharge node
carries those squares. The rule: from
$S' \vDash \mathit{spec}'$ and the judgment, conclude
$S \vDash \sigma^{*}\mathit{spec}'$ --- certify the
abstraction once, reuse under refinement \cite{lynchvaandrager1995,
capuccimyers2026}.

\begin{bthm}[behavioral pullback]\label{thm:qst-pullback}
Let $\mathrm{Beh}$, $\mathrm{Beh}'$ be the compiled behaviour classes
of $S$, $S'$, and suppose $\sigma$ maps behaviours to behaviours
($\sigma[\mathrm{Beh}] \subseteq \mathrm{Beh}'$ --- the semantic
content of the compatibility squares). If every behaviour of $S'$
satisfies $\mathit{spec}'$, then every behaviour of $S$ satisfies
$\sigma^{-1}[\mathit{spec}']$.
\end{bthm}
\begin{proof}
$\omega \in \mathrm{Beh}$ gives
$\sigma(\omega) \in \mathrm{Beh}' \subseteq \mathit{spec}'$, i.e.\
$\omega \in \sigma^{-1}[\mathit{spec}']$
(\texttt{pullback}). Two lines is the honest size: the rule's
\emph{semantic} content is inclusion-chasing, and everything
substantive in an application is establishing
$\sigma[\mathrm{Beh}] \subseteq \mathrm{Beh}'$ from the checked
squares --- per-mode polynomial identities, discharged by the base
calculus. The hypothesis is not droppable to
``$\sigma$ is any map'': without the squares the image of a concrete
behaviour need not be an abstract behaviour, and the pullback claims
anything at all.
\end{proof}

\textbf{The finite-kernel instance.}\enspace At the finite Markov
tier the witness has teeth the behavioural form leaves abstract: a
\emph{functional simulation} is per-state marginal identity ---
pushing the concrete kernel row at $s$ along $\sigma$ yields the
abstract kernel row at $\sigma(s)$ (a pushforward-kernel witness) ---
and the transport theorems are proven against the chain machinery.

\begin{blemma}[observable pullback]\label{lem:qst-expect}
For any law $\mu$ on the concrete states and observable $g$ on the
abstract states,
$\sum_b (\sigma_{*}\mu)(b) \, g(b) = \sum_s \mu(s) \, g(\sigma(s))$;
and along a pushforward-kernel witness, the time-$n$ expectation of a
pulled-back observable on the concrete chain \emph{equals} the
abstract chain's expectation of $g$ at time $n$ (started from the
pushed initial law). Hence any certified expectation bound proven
against the abstract chain holds for the concrete chain at the
pulled-back observable.
\end{blemma}
\begin{proof}
The first identity is summation over fibres of $\sigma$
(\texttt{expect\_push}). The second composes it with the
marginal-process lemma --- the per-state marginal identities
propagate, by induction on time, to equality of pushed laws at every
time --- giving \texttt{sim\_expect}; the bound transports by
rewriting (\texttt{sim\_expect\_bound}). Note the direction: the
producer proves one bound on the small chain and every refinement
inherits it; nothing is re-proven per refinement beyond the witness
identities themselves. The interface reading rides the same
equalities: an assumption event $A'$ pulls back with its mass intact
(the pushed law at time $n$ \emph{equals} the abstract law, so
$\nu_n(\sigma^{-1}A') = \lambda_n(A')$), hence assumption-weighted
bounds transport, and conditional expectations transport whenever the
conditioning mass is positive.
\end{proof}

\begin{blemma}[region pullback]\label{lem:qst-region}
Along a pushforward-kernel witness with entrywise-nonnegative kernel
and initial law: if the abstract law at time $n$ puts no mass outside
$R'$, the concrete law at time $n$ puts no mass outside
$\sigma^{-1}[R']$.
\end{blemma}
\begin{proof}
For $s$ with $\sigma(s) \notin R'$: the concrete mass at $s$ is at
most the pushed mass at $\sigma(s)$ (one nonnegative term of a
fibre sum), which equals the abstract mass at $\sigma(s)$ (the
marginal-process lemma), which is zero; nonnegativity closes the
squeeze (\texttt{sim\_region}). The nonnegativity hypotheses are
load-bearing, not hygiene: with signed ``mass'' the fibre sum could
cancel --- a negative sibling hiding positive mass at $s$ outside the
preimage region --- and the single-term bound fails. Nonnegativity is
what makes ``no mass outside $R'$'' transfer pointwise through a
many-to-one map.
\end{proof}

\begin{bremark}[refinement is one-sided; a simulation is not a
coupling]\label{rem:qst-onesided}
The witness machinery genuinely shared with the relational
(two-chain) tier is exactly the pushforward-kernel identity and the
marginal-process lemma --- a functional simulation \emph{is} such a
witness. Nothing else transfers: the relational tier's certificates
are joint-kernel data over a product space with absorption and drift
hypotheses, and none of that is simulation transport. A functional
simulation induces a degenerate coupling, but the theorems above
never route through it --- they are pushforward plus summation
identities, full stop. Conflating the two would import hypotheses
(coupling emptiness, absorption) that this rule neither needs nor
checks.
\end{bremark}

\begin{bremark}[the gaps, loud]\label{rem:qst-gaps}
The section walks two strata --- behavioural, and finite-kernel. The
general measurable-kernel tier is mechanized and bound in this
rule's entry, mirroring the finite instance name-for-name over
Markov kernels (the same pushforward identity per state, now a
\texttt{Measure.map} equation; law transport at every horizon; both
spec-transport faces --- with the finite tier's nonnegativity
hypotheses gone, true by type); it is cited there rather than
re-walked here. Not yet mechanized: the categorical strengthening
under which the composition rule could cite the lens-composition
theorem of
\cite{capuccimyers2026} as a lemma \emph{template} --- the direct
composition proofs never pull back along strict weakenings, so
nothing in this book depends on that strengthening; it is named so
that its absence is a recorded choice, not an oversight. The
judgment kinds the rule consumes and concludes into are staged in
the certificate schema (the simulation judgment together with its
discharge and composition nodes), but no end-to-end certificate yet
exercises the rule --- loader and end-to-end work, not mathematics. Related
reading, not dependency: an exchangeability quotient is a simulation
map, so aggregate-space reduction can be \emph{read} as an instance
of this rule; it stands on its own proof and cites nothing here.
\end{bremark}

\begin{anchors}
Theorem \ref{thm:qst-pullback} & \texttt{pullback}\\
Lemma \ref{lem:qst-expect} & \texttt{expect\_push}, \texttt{sim\_expect}, \texttt{sim\_expect\_bound}; marginal-process lemma \texttt{push\_law} (chain substrate module)\\
Lemma \ref{lem:qst-region} & \texttt{sim\_region}; law nonnegativity \texttt{law\_nonneg} (chain substrate module)\\
witness form & \texttt{PushKernel} (chain substrate module) --- the per-state marginal identity\\
compatibility squares (cited, not re-proved) & inherited sequent leaves, per mode (main package)\\
\end{anchors}


\rulesection{82}{QFlatten}{admissible rule}
\label{blcatrule-QFlatten}
\ruleblock{The rule (Steele metanotation)}
\sourcenote{as realized in the mechanized development; source row: \texttt{\texttt{QTSL-rules.toml}} --- \emph{QFlatten}. The metanotation statement is typeset by hand for notational quality; its mechanized ground truth is the Lean statement below, extracted from the sources at every build of this book and hash-bound. The rule \emph{roster} is machine-checked (every mechanized rule tag round-trips to exactly one catalogue row); the statement \emph{text} is editorial, checked by eye against the extracted Lean beside it.}
\begin{mathpar}
\inferrule*[Right=\rulename{QFlatten}]{
  \mathsf{QFlatEq}(\mathrm{sys}, \mathit{flat}) \\
  \mathcal{M}_{\mathit{flat}},\ x_0 \in R \vDash \varphi \ \text{(the closed-system pipeline\textquotesingle s judgment)}
}{
  \mathrm{sys},\ x_0 \in R \vDash \varphi
}
\end{mathpar}

\ruleflavor{$\mathit{flat}$ is the flattening of the composite $\mathrm{sys}$; the substitution rule it always was: satisfaction transports across the flattening equation. The former side condition (polynomial substitution + guard-product bookkeeping) is the kind\textquotesingle s discharge --- oracle-shaped by construction (\texttt{qtsl\_flatten}, a bounded deterministic recompute over declaration structure), with the refusal boundary stated positively as the judgment\textquotesingle s scope: shared-carrier confinement, and binary wiring with $n$-ary/directed cases reducing by currying = the checker\textquotesingle s substitution recompute (refusal-shaped completeness cost, never a soundness hypothesis). The guarantee the rule was minted to state --- $\mathsf{SysWire}$ adds structure, never expressiveness --- is carried by the kind\textquotesingle s existence: the flattened manifest\textquotesingle s compiled semantics is the composite\textquotesingle s (theorem at both kernel tiers on the synchronous-product face), and the projections remain free \rulename{QSimTransport} witnesses under sibling well-formedness (unit row sums). The clock rows (schedule hypotheses, simultaneity, non-Zeno) are premise classes cited by every composition rule, deliberately not rows here.}
\par\medskip
\rulename{QFlatten} transports satisfaction across the flattening
equation: given the judgment $\mathsf{QFlatEq}(\mathrm{sys},
\mathit{flat})$ --- $\mathit{flat}$ is the flattening of the composite
presentation, discharged by the checker\textquotesingle s substitution
recompute --- a verdict on the monolithic manifest is a verdict on the
composite. A
producer reaches for it as the bridge to every closed-system pipeline:
compile nodes, certificate synthesis, and the whole battery machinery
operate on one flat manifest, while modelling happens compositionally.
The rule\textquotesingle s two guarantees pull in opposite directions
and are both load-bearing: wiring adds \emph{structure}, never
\emph{expressiveness}. Never expressiveness: the flattened
manifest\textquotesingle s compiled semantics is exactly the
composite\textquotesingle s --- at the compiled class this is
definitional (a presentation\textquotesingle s semantics IS its
flattened object\textquotesingle s), and at the finite-kernel tier it
is a theorem: the monolithic chain\textquotesingle s law is, at every
horizon, the product of the component laws. Added structure: the
flattened manifest remembers its components --- the projections are
exact simulation witnesses for free, so every component-level
specification remains addressable on the flattened object through
\rulename{QSimTransport}.

What the judgment\textquotesingle s discharge recomputes: the
substitution itself --- polynomial
substitution plus guard-product bookkeeping --- which is why the rule is
admissible rather than core: it discharges against verified
recomputation and the semantic witness pair, adding no new kernel
content. Directed substitution wirings do \emph{not} reduce to the
independent synchronous product of the witness pair --- wiring an
upstream output into a downstream input creates conditional dependence
the product law does not carry; they rest on the checker\textquotesingle
s substitution recompute itself (currying the upstream trajectory into
a dependent downstream kernel), definitional in the compiled class.

What would break: component well-formedness (unit row sums) is what
makes the projections exact witnesses --- against a defective second
factor, the first projection\textquotesingle s pushed row picks up the
leftover mass as a spurious scale factor, and the component-spec
addressability claim fails. The scheduling hypotheses (clock rows,
simultaneity, non-Zeno families) are deliberately premise classes cited
by every composition rule rather than rows of this one.

\cataloguenote{The flattening premise is the judgment kind $\mathsf{QFlatEq}$,
baked in the extension theory; its discharge is the flatten
recompute behind a registry seam, with the shared-carrier and
binary-wiring confinements as the judgment\textquotesingle s stated
scope --- refusal-shaped, never a hypothesis of this rule.
The proven fragment spans two tiers: the finite-kernel tier, on the
same substrate as the relational chain machinery --- the tier where
\textquotedblleft law\textquotedblright, \textquotedblleft
product\textquotedblright, and \textquotedblleft
pushforward\textquotedblright{} are finite exact-arithmetic objects ---
and the general measurable-kernel tier, where they
are the Giry-monad objects (second bound block). The remaining
confinement is the \emph{independent binary product} witness pair --- a
design fence, not a gap: $n$-ary independent wirings iterate it, while
directed wirings rest on the checker\textquotesingle s substitution
recompute (a dependent downstream kernel, not the independent product).
The general compiled-class statement is substitution correctness,
checker recomputation by design.
}
\ruleblock{Lean statement}
\declhead{BernsteinLogicAux.QtslCompose.Flatten}{law\_prodKernel}{\badge{BadgeGreen}{statement bound}}
\begin{leansrc}
/-- **The flatten witness** (`QFlatten`'s "never expressiveness" face): the
flattened monolithic chain's law is, at every time, the product of the
component chains' laws — the composite presentation and its monolithic
manifest compile to the same semantics. -/
theorem law_prodKernel (K₁ : α → α → ℝ) (K₂ : β → β → ℝ)
    (μ₁ : α → ℝ) (μ₂ : β → ℝ) :
    ∀ n, law (prodKernel K₁ K₂) (prodDist μ₁ μ₂) n
      = prodDist (law K₁ μ₁ n) (law K₂ μ₂ n)
  | 0 => rfl
  | n + 1 => by
    rw [law_succ, law_prodKernel K₁ K₂ μ₁ μ₂ n, step_prodKernel, law_succ,
      law_succ]

section Projections

variable [DecidableEq α] [DecidableEq β]

/-- The added *structure*: the first projection is a simulation witness
(`QSimTransport.PushKernel`) from the flattened manifest onto component 1 —
for free, whenever component 2 is well-formed. -/
theorem pushKernel_fst {K₁ : α → α → ℝ} {K₂ : β → β → ℝ}
    (h₂ : ∀ b, ∑ b', K₂ b b' = 1) :
    PushKernel Prod.fst (prodKernel K₁ K₂) K₁
\end{leansrc}
\declhead{BernsteinLogicAux.QtslCompose.Flatten}{pushKernel\_fst}{\badge{BadgeGreen}{statement bound}}
\begin{leansrc}
/-- The added *structure*: the first projection is a simulation witness
(`QSimTransport.PushKernel`) from the flattened manifest onto component 1 —
for free, whenever component 2 is well-formed. -/
theorem pushKernel_fst {K₁ : α → α → ℝ} {K₂ : β → β → ℝ}
    (h₂ : ∀ b, ∑ b', K₂ b b' = 1) :
    PushKernel Prod.fst (prodKernel K₁ K₂) K₁
\end{leansrc}
\declhead{BernsteinLogicAux.QtslCompose.Flatten}{pushKernel\_snd}{\badge{BadgeGreen}{statement bound}}
\begin{leansrc}
/-- Mirror: the second projection simulates onto component 2. -/
theorem pushKernel_snd {K₁ : α → α → ℝ} {K₂ : β → β → ℝ}
    (h₁ : ∀ a, ∑ a', K₁ a a' = 1) :
    PushKernel Prod.snd (prodKernel K₁ K₂) K₂
\end{leansrc}
\declhead{BernsteinLogicAux.QtslCompose.FlattenMeas}{law\_parallelComp}{\badge{BadgeGreen}{statement bound}}
\begin{leansrc}
/-- **QFlatten, measurable-kernel tier — the flatten witness** ("never
expressiveness" face): the flattened monolithic chain's law is, at every
horizon, the product of the component chains' laws — the composite
presentation and its monolithic manifest compile to the same semantics.
(The measurable-tier face of `law_prodKernel`.) -/
theorem law_parallelComp (μ : Measure α) (ν : Measure β)
    (K : Kernel α α) (Q : Kernel β β)
    [SFinite μ] [SFinite ν] [IsSFiniteKernel K] [IsSFiniteKernel Q] :
    ∀ n, law (K ∥ₖ Q) (μ.prod ν) n = (law K μ n).prod (law Q ν n)
  | 0 => rfl
  | n + 1 => by
    rw [law_succ, law_parallelComp μ ν K Q n, bind_parallelComp,
      ← law_succ, ← law_succ]

/-! ## The added structure: projections are simulation witnesses -/

/-- The added *structure* ("adds structure" face): the first projection is a
simulation witness (`SimTransportMeas.PushKernel`) from the flattened
manifest onto component 1 — for free, consuming exactly the *sibling's*
Markov property (`[IsMarkovKernel Q]`: unit total mass per row) and nothing
of `K`'s own well-formedness. Margin-zero boundary, displayed by
`Measure.map_fst_prod`: against a defective sibling the pushed row is
`(Q p.2 univ) • K p.1` — the leak coefficient is the sibling's total mass.
(The measurable-tier face of the finite tier's `pushKernel_fst`.) -/
theorem pushKernel_fst {K : Kernel α α} {Q : Kernel β β}
    [IsSFiniteKernel K] [IsMarkovKernel Q] :
    PushKernel Prod.fst (K ∥ₖ Q) K
\end{leansrc}
\declhead{BernsteinLogicAux.QtslCompose.FlattenMeas}{pushKernel\_fst}{\badge{BadgeGreen}{statement bound}}
\begin{leansrc}
/-- The added *structure* ("adds structure" face): the first projection is a
simulation witness (`SimTransportMeas.PushKernel`) from the flattened
manifest onto component 1 — for free, consuming exactly the *sibling's*
Markov property (`[IsMarkovKernel Q]`: unit total mass per row) and nothing
of `K`'s own well-formedness. Margin-zero boundary, displayed by
`Measure.map_fst_prod`: against a defective sibling the pushed row is
`(Q p.2 univ) • K p.1` — the leak coefficient is the sibling's total mass.
(The measurable-tier face of the finite tier's `pushKernel_fst`.) -/
theorem pushKernel_fst {K : Kernel α α} {Q : Kernel β β}
    [IsSFiniteKernel K] [IsMarkovKernel Q] :
    PushKernel Prod.fst (K ∥ₖ Q) K
\end{leansrc}
\declhead{BernsteinLogicAux.QtslCompose.FlattenMeas}{pushKernel\_snd}{\badge{BadgeGreen}{statement bound}}
\begin{leansrc}
/-- Mirror: the second projection simulates onto component 2, consuming
exactly the first component's Markov property. -/
theorem pushKernel_snd {K : Kernel α α} {Q : Kernel β β}
    [IsMarkovKernel K] [IsSFiniteKernel Q] :
    PushKernel Prod.snd (K ∥ₖ Q) Q
\end{leansrc}
\declhead{BernsteinLogicAux.QtslCompose.FlattenMeas}{bind\_parallelComp}{\badge{BadgeGreen}{statement bound}}
\begin{leansrc}
/-- One step of the flattened (synchronous-product) chain factors as the
product of the component steps. (The measurable-tier face of
`step_prodKernel`.) Proof is pure `∘ₘ`-algebra:
`Measure.prod_comp_right`/`prod_comp_left` are the two one-sided
measure-level product-bind identities, fused by
`Kernel.parallelComp_comp_parallelComp`. -/
theorem bind_parallelComp (μ : Measure α) (ν : Measure β)
    (K : Kernel α α) (Q : Kernel β β)
    [SFinite μ] [SFinite ν] [IsSFiniteKernel K] [IsSFiniteKernel Q] :
    (μ.prod ν).bind (K ∥ₖ Q) = (μ.bind K).prod (ν.bind Q)
\end{leansrc}
\ruleblock{Correspondence (Steele display $\leftrightarrow$ Lean)}
\badge{BadgeGreen}{reviewer-certified against the displayed texts}
\sourcenote{The dictionary below is a reviewer-certified statement of how the mechanized statements above interpret the displayed rule: display metavariable $\mapsto$ Lean binder, each display hypothesis $\mapsto$ the premise or side condition carrying it, the display conclusion $\mapsto$ the goal --- with any deliberate difference in width between display and mechanization named outright. Certification is recorded as a hash triple over (the displayed cell, the extracted statement texts, this dictionary\textquotesingle s own text); if any of the three changes, this block turns amber until a reviewer re-certifies. The correspondence claim itself remains editorial --- reviewer-grade, not machine-checked; its machine-checked successor types each soundness statement against the rule as a structured object, and supersedes this block per rule as it lands.}
The displayed judgment premise $\mathsf{QFlatEq}(\mathrm{sys},
\mathit{flat})$ --- the flattened manifest has the
\emph{same} compiled semantics, each projection $\pi_i$ a free
\rulename{QSimTransport} witness --- $\mapsto$ the theorem pair at the
finite-kernel tier: \texttt{law\_prodKernel} (never expressiveness)
and \texttt{pushKernel\_fst}/\texttt{\_snd} (added structure). The
kind\textquotesingle s discharge (the substitution recompute) has
\emph{no} Lean
counterpart by design: at the compiled class the
presentation\textquotesingle s semantics IS its flattened
object\textquotesingle s, so that discharge is checker recomputation,
and what is theorem-shaped is exactly what is bound. The
rule\textquotesingle s own transport step (a verdict on
$\mathcal{M}_{\mathit{flat}}$ is a verdict on $\mathrm{sys}$) is
substitution of equals along the judgment --- no separate statement
to bind.
\begin{itemize}
\item $\mathsf{SysWire}(w, \mathit{Sys}_1, \ldots, \mathit{Sys}_n)$
  $\mapsto$ the \emph{binary synchronous product}: kernels
  \texttt{K₁ : α → α → ℝ}, \texttt{K₂ : β → β → ℝ} over finite types,
  flattened as \texttt{prodKernel} on \texttt{α × β}. Two confinements
  at this block, both design fences: binary-independent (directed and
  $n$-ary substitution wirings rest on the checker\textquotesingle s
  substitution recompute --- currying the upstream trajectory into a
  \emph{dependent} downstream kernel, which is not the independent
  product bound here; per the module header), and
  finite-kernel tier (\texttt{[Fintype α] [Fintype β]}) --- the latter
  discharged by the measurable-tier block below, leaving the
  independent binary product as this bound set\textquotesingle s
  remaining confinement.
\item same compiled semantics $\mapsto$ \texttt{law\_prodKernel}:
  \texttt{law (prodKernel K₁ K₂) (prodDist μ₁ μ₂) n = prodDist (law K₁
  μ₁ n) (law K₂ μ₂ n)} at every \texttt{n} --- and its binder row is
  the width fact: the identity is proven for \emph{arbitrary} kernel
  and law data, no Markov/distribution well-formedness consumed ---
  pure bilinear algebra, strictly wider than the display\textquotesingle s
  well-formed-system reading, safe direction.
\item projections as free \rulename{QSimTransport} witnesses $\mapsto$
  \texttt{pushKernel\_fst : PushKernel Prod.fst (prodKernel K₁ K₂) K₁}
  --- consuming only \texttt{∀ b, ∑ b', K₂ b b' = 1}: the
  \emph{sibling} component\textquotesingle s unit row sums, nothing of
  its own, and not even nonnegativity --- the certified statement
  isolates exactly which well-formedness fragment the
  commentary\textquotesingle s ``what would break'' paragraph names
  (a defective second factor leaks scale into the pushed row).
  \texttt{\_snd} mirrors. \texttt{[DecidableEq]} binders enter only
  here (the pushforward\textquotesingle s fibre sums).
\end{itemize}
\emph{Measurable-kernel tier} (second bound block, module
\texttt{FlattenMeas}, namespace \texttt{…QtslCompose.Meas}): the
product is mathlib\textquotesingle s own pair
--- \texttt{K ∥ₖ Q} (\texttt{Kernel.parallelComp}) over
\texttt{Measure.prod} --- and the theorem pair transposes
name-for-name:
\texttt{law\_parallelComp} (never expressiveness: \texttt{law (K ∥ₖ Q)
(μ.prod ν) n = (law K μ n).prod (law Q ν n)} at every horizon, over
\texttt{SimTransportMeas}\textquotesingle s Giry-bind \texttt{law}) and
\texttt{pushKernel\_fst}/\texttt{\_snd} (added structure: projections
as \texttt{Meas.PushKernel} witnesses), with \texttt{bind\_parallelComp}
the one-step glue. Binder floor, the honest analogue of the finite
tier\textquotesingle s ``arbitrary data'' width fact:
\texttt{SFinite}/\texttt{IsSFiniteKernel} --- product-measure theory is
s-finite theory --- still strictly wider than the
display\textquotesingle s Markov/probability reading, same safe
direction. And the ``what would break'' margin-zero witness is
displayed by mathlib itself: \texttt{Measure.map\_fst\_prod} gives the
pushed row as \texttt{(ν univ) • μ} --- leak coefficient the
sibling\textquotesingle s total mass --- with \texttt{pushKernel\_fst}
consuming exactly \texttt{[IsMarkovKernel Q]} and nothing of
\texttt{K}\textquotesingle s own well-formedness, not even finiteness.
Width note: the clock rows (schedule/simultaneity/non-Zeno) are
premise classes deliberately absent from both display and bound
statements; the closure facts at both tiers live in their modules but
outside this bound set --- cited, not certified here: finite
\texttt{isMarkov\_prodKernel}, \texttt{isDist\_prodDist},
\texttt{push\_fst\_prodDist}; measurable \texttt{sfinite\_law},
\texttt{isProbabilityMeasure\_law}, the addressability compositions
\texttt{sim\_expect\_fst}/\texttt{\_snd} (component specs consumed
end-to-end through \rulename{QSimTransport}\textquotesingle s
measurable engine), the bridge \texttt{pushKernel\_iff\_stochSim}, and
the mathlib instances \texttt{IsMarkovKernel (K ∥ₖ Q)} /
\texttt{prod.instIsProbabilityMeasure}.

\ruleblock{Lean verifier receipts}
\sourcenote{module \texttt{BernsteinLogicAux.QtslCompose.Flatten}; bound build \texttt{30ed43793772\ldots} (bound from receipt 8cd9a31cf76c1595, 2026-07-19T18:24:49Z). A receipt records that named checking layers accepted a build: (module, compiled-artifact SHA-256, environment SHA-256, checker layers), minted when the checking gate ran; the store receipts attesting the bound build are listed with the layers each run credited. Nothing here asks the reader to trust the store: every layer is an independently runnable program over the quoted sources --- rebuild the module and re-run the checkers to reproduce the verdict. Receipts attesting superseded builds credit no layer and are omitted from this rendering (the count of omissions, if any, is noted below).}
\begin{receipttable}
\receiptrow{e4496ce5107a3441}{2026-07-19T03:38:56Z}{build+decorrelated}{\badge{BadgeGreen}{current}}
\receiptrow{8cd9a31cf76c1595}{2026-07-19T18:24:49Z}{build+decorrelated}{\badge{BadgeGreen}{current}}
\receiptrow{62e3b3d615f78f37}{2026-07-20T01:22:02Z}{build+decorrelated}{\badge{BadgeGreen}{current}}
\receiptrow{b45a913927ac5488}{2026-07-20T20:46:52Z}{build+decorrelated}{\badge{BadgeGreen}{current}}
\receiptrow{d52960f0b18dc76a}{2026-07-20T20:58:26Z}{build+replay+hardened+decorrelated}{\badge{BadgeGreen}{current}}
\receiptrow{0eabb0d2fdfbc00c}{2026-07-21T00:03:55Z}{build+decorrelated}{\badge{BadgeGreen}{current}}
\receiptrow{1f0f373238ef2c2b}{2026-07-21T05:39:11Z}{build+replay+hardened+decorrelated}{\badge{BadgeGreen}{current}}
\receiptrow{35d0f3c7c71856fc}{2026-07-22T18:27:16Z}{build+decorrelated}{\badge{BadgeGreen}{current}}
\receiptrow{53dde5ffb738ab3a}{2026-07-31T22:00:14Z}{build+decorrelated}{\badge{BadgeGreen}{current}}
\receiptrow{50f8cfda2d25e116}{2026-07-31T22:29:08Z}{build+decorrelated}{\badge{BadgeGreen}{current}}
\receiptrow{edc0e93a8721d195}{2026-08-02T23:31:19Z}{build+decorrelated}{\badge{BadgeGreen}{current}}
\receiptrow{a6fae671c06c5ebb}{2026-08-03T16:27:23Z}{build+decorrelated}{\badge{BadgeGreen}{current}}
\receiptrow{0427fd20db259dcd}{2026-08-03T22:46:07Z}{build+decorrelated}{\badge{BadgeGreen}{current}}
\end{receipttable}
\sourcenote{module \texttt{BernsteinLogicAux.QtslCompose.FlattenMeas}; bound build \texttt{40d14da6aacd\ldots} (bound from receipt 4747724a289473cd, 2026-07-20T08:24:46Z). A receipt records that named checking layers accepted a build: (module, compiled-artifact SHA-256, environment SHA-256, checker layers), minted when the checking gate ran; the store receipts attesting the bound build are listed with the layers each run credited. Nothing here asks the reader to trust the store: every layer is an independently runnable program over the quoted sources --- rebuild the module and re-run the checkers to reproduce the verdict. Receipts attesting superseded builds credit no layer and are omitted from this rendering (the count of omissions, if any, is noted below).}
\begin{receipttable}
\receiptrow{4747724a289473cd}{2026-07-20T08:24:46Z}{build+decorrelated}{\badge{BadgeGreen}{current}}
\receiptrow{2a11a368eea6adc5}{2026-07-20T20:46:54Z}{build+decorrelated}{\badge{BadgeGreen}{current}}
\receiptrow{d11ebe966df6a113}{2026-07-20T20:58:27Z}{build+replay+hardened+decorrelated}{\badge{BadgeGreen}{current}}
\receiptrow{ba9ffebde141f18b}{2026-07-21T00:03:56Z}{build+decorrelated}{\badge{BadgeGreen}{current}}
\receiptrow{13a1681b1015f8c2}{2026-07-21T05:39:13Z}{build+replay+hardened+decorrelated}{\badge{BadgeGreen}{current}}
\receiptrow{89eb1271b74a3f48}{2026-07-22T18:27:16Z}{build+decorrelated}{\badge{BadgeGreen}{current}}
\receiptrow{d86334f5a6f3694d}{2026-07-31T22:00:15Z}{build+decorrelated}{\badge{BadgeGreen}{current}}
\receiptrow{4aceb786dcde76f2}{2026-07-31T22:29:09Z}{build+decorrelated}{\badge{BadgeGreen}{current}}
\receiptrow{e040a1347e8cd2a0}{2026-08-02T23:31:20Z}{build+decorrelated}{\badge{BadgeGreen}{current}}
\receiptrow{52c94a6a75909deb}{2026-08-03T16:27:23Z}{build+decorrelated}{\badge{BadgeGreen}{current}}
\receiptrow{0a269083e637c33e}{2026-08-03T22:46:07Z}{build+decorrelated}{\badge{BadgeGreen}{current}}
\end{receipttable}
\layerline{\layeritem{build}{lake build (elaboration)}{\badge{BadgeGreen}{current}}\layeritem{replay}{leanchecker (official kernel replay)}{\badge{BadgeGreen}{current}}\layeritem{hardened}{safe\_verify (hardened screen)}{\badge{BadgeGreen}{current}}\layeritem{decorrelated}{lean4export + nanoda (decorrelated kernel)}{\badge{BadgeGreen}{current}}}
\sourcenote{\badge{BadgeGray}{comparator: none} no fixture in the three-implementation comparator harness (compiled / kernel-interpreted / reference-interpreted) exercises this rule.}
\ruleblock{HOL4 statement}
\absentbox{NO HOL STATEMENT YET}{The HOL4 development checks the base calculus\textquotesingle s wire
format (theory \texttt{bern\_catalogue}). \rulename{QFlatten} operates
on system presentations of the extension theory --- baked, System
stratum and all; its semantic witnesses live in the
auxiliary Lean package. No HOL development consumes the extension
theory\textquotesingle s rows; the box becomes writable when a
flatten recompute joins a ported checker chain.
}
\ruleblock{HOL verifier receipts}
\absentbox{NONE}{no HOL verification receipt exists for this rule.}
\ruleblock{Proof}

\sourcenote{Genre: admissible rule (tier 2) --- a proof-tree-rewriting
recipe of the reduction-to-verified-machinery kind: the node discharges
against the checker's substitution recompute plus a semantic witness
pair, not against catalogue nodes. Scope, stated loudly: the witness
pair is proven at the finite-kernel tier, on the synchronous-product
wiring; the general compiled-class statement is substitution
correctness, definitional by the presentation semantics.}

\begin{brecipe}[the \rulename{QFlatten} rewrite]\label{rec:qfl-rewrite}
A \rulename{QFlatten} node transports satisfaction across the
flattening judgment --- $\mathit{flat}$ the monolithic manifest of a
composite presentation:
\[
\inferrule*[Right=\rulename{QFlatten}]{\;\mathsf{QFlatEq}(\mathrm{sys}, \mathit{flat}) \\ \mathcal{M}_{\mathit{flat}},\ x_0 \in R \vDash \varphi\;}{\mathrm{sys},\ x_0 \in R \vDash \varphi}
\qquad\rightsquigarrow\qquad
\text{\upshape the checker replays the substitution.}
\]
The rewrite target contains zero catalogue nodes: a presentation's
compiled semantics \emph{is} its flattened object's, so at the compiled
class the judgment's discharge is a recomputation --- polynomial
substitution
plus guard-product bookkeeping --- and the rule adds no kernel semantic
content (the admissibility criterion). What deserves proof, and gets it
below, is the slogan the rule is trusted for: \textbf{wiring adds
structure, never expressiveness}. Directed substitution wirings do
\emph{not} reduce to the independent synchronous product treated below
--- wiring an upstream output into a downstream input creates
conditional dependence the product law does not carry --- so the
flatten witness below attests the independent-product case, and the
directed case rests on the substitution recompute (the judgment's
discharge above) alone.
\end{brecipe}

\begin{bprop}[never expressiveness: the flatten witness]\label{prop:qfl-law}
Let $K_1, K_2$ be finite kernels with initial laws $\mu_1, \mu_2$, and
let the flattened manifest be the product kernel
$K_{\otimes}\bigl((a,b), (a',b')\bigr) = K_1(a, a')\, K_2(b, b')$ with
initial law $(\mu_1 \otimes \mu_2)(a, b) = \mu_1(a)\, \mu_2(b)$. Then
at every horizon $n$,
\[
\mathrm{law}_{K_\otimes}\bigl(\mu_1 \otimes \mu_2,\ n\bigr)
\;=\; \mathrm{law}_{K_1}(\mu_1, n) \,\otimes\, \mathrm{law}_{K_2}(\mu_2, n):
\]
the monolithic chain's law is, at every time, exactly the product of
the component laws --- the composite presentation and its flattened
manifest compile to the same semantics.
\end{bprop}
\begin{proof}
One step factors, by Fubini for finite sums and bilinearity:
\[
\mathrm{step}_{K_\otimes}(\mu_1 \otimes \mu_2)(a', b')
= \sum_{a, b} \mu_1(a) \mu_2(b)\, K_1(a, a') K_2(b, b')
= \Bigl(\sum_a \mu_1(a) K_1(a, a')\Bigr)\Bigl(\sum_b \mu_2(b) K_2(b, b')\Bigr),
\]
i.e.\ $\mathrm{step}_{K_\otimes}(\mu_1 \otimes \mu_2) =
\mathrm{step}_{K_1}\mu_1 \otimes \mathrm{step}_{K_2}\mu_2$. Induction
on $n$ closes. Closure rides along: the product of kernels is a kernel
and the product of laws is a law (nonnegativity termwise; unit masses
by the same bilinear factorization ---
$\sum_{a,b} \mu_1(a)\,\mu_2(b) = 1 \cdot 1$ for the product law,
$\sum_{a',b'} K_1(a, a')\,K_2(b, b') = 1 \cdot 1$ per product row).
\end{proof}

\begin{bprop}[added structure: the projections are free simulation
witnesses]\label{prop:qfl-proj}
If $K_2$ is row-stochastic, the first projection is an \emph{exact}
intertwining witness from $K_\otimes$ onto $K_1$:
$\pi_{1\#} K_\otimes\bigl((a,b), \cdot\bigr) = K_1(a, \cdot)$ for every
product state $(a, b)$; symmetrically for $\pi_2$ under
row-stochastic $K_1$; and $\pi_{1\#}(\mu_1 \otimes \mu_2) = \mu_1$
whenever $\mu_2$ has unit mass.
\end{bprop}
\begin{proof}
$\pi_{1\#} K_\otimes\bigl((a,b), \cdot\bigr)(a') =
\sum_{b'} K_1(a, a') K_2(b, b') = K_1(a, a') \sum_{b'} K_2(b, b') =
K_1(a, a')$, the last step by the unit row sum. The initial-law
statement is the same computation at $\mu$. Consequence: the flattened
manifest \emph{remembers} its components --- each projection is a free
exact-intertwining (\texttt{PushKernel}) witness, the
\rulename{QSimTransport} witness form, instantiating the strict
(zero-degradation) column of the \rulename{QDistTransport} matrix ---
so every component-level specification remains addressable on the
flattened object with zero degradation. $\mathsf{SysWire}$ adds these
handles and nothing else: Proposition~\ref{prop:qfl-law} caps
expressiveness, this proposition delivers the structure.
\end{proof}

\begin{bremark}[why the well-formedness hypotheses are not decoration]
Both directions consume them. In Proposition~\ref{prop:qfl-proj}, drop
the unit row sum of the \emph{other} factor and the projection witness
fails: the pushed row picks up $\sum_{b'} K_2(b, b')$ as a spurious
scale factor, so a sub-stochastic (leaky) second component would make
every component-1 judgment transport against a deflated law --- the
component-spec addressability claim silently weakens. In
Proposition~\ref{prop:qfl-law}, unit masses are what make the product
law a probability law at every time. The scheduling hypotheses of
composite dynamics --- clock rows, simultaneity, non-Zeno families ---
are deliberately \emph{premise classes} cited by every composition rule
rather than rows of this one: flattening is about state and kernel
structure, and importing the clock discipline here would duplicate its
single authoritative statement.
\end{bremark}

\begin{bremark}[the tier boundary, honestly]
The propositions above are the synchronous-product instance at the
finite-kernel tier --- the tier where law, product, and pushforward are
finite exact-arithmetic objects, shared with the relational-chain
machinery. General wirings (directed substitution, guard products)
reduce by the checker's own recompute; the reduction is definitional in
the compiled class and is not re-proven as a theorem here. Nothing
above claims the measure/kernel tier; its flatten-commutes statements
--- the same theorem pair over Markov kernels and product measures ---
are mechanized and bound in this rule's entry, cited there rather than
re-walked here.
\end{bremark}

\begin{anchors}
Recipe \ref{rec:qfl-rewrite} judgment discharge & the substitution recompute (polynomial substitution + guard-product bookkeeping; checker-side, behind the flattening judgment's registry seam)\\
Proposition \ref{prop:qfl-law} & \texttt{prodKernel}; \texttt{prodDist}; \texttt{step\_prodKernel}; \texttt{law\_prodKernel}; closure \texttt{isMarkov\_prodKernel}, \texttt{isDist\_prodDist}\\
Proposition \ref{prop:qfl-proj} & \texttt{pushKernel\_fst}; \texttt{pushKernel\_snd}; \texttt{push\_fst\_prodDist}\\
consumed by & \rulename{QSimTransport} (the \texttt{PushKernel} witness form); the strict column of \rulename{QDistTransport}'s matrix (\texttt{strict\_cell\_exact})\\
\end{anchors}


\rulesection{83}{QLinGrowthDischarge}{core rule}
\label{blcatrule-QLinGrowthDischarge}
\ruleblock{The rule (Steele metanotation)}
\sourcenote{as realized in the mechanized development; source row: \texttt{\texttt{QTSL-rules.toml}} --- \emph{QLinGrowthDischarge}. The metanotation statement is typeset by hand for notational quality; this rule cites no Lean statements yet (the absence is stated loudly below). The rule \emph{roster} is machine-checked (every mechanized rule tag round-trips to exactly one catalogue row); the statement \emph{text} is editorial.}
\begin{mathpar}
\inferrule*{
  \overline{\mathrm{env}_i : B \vdash\ \pm f_i(x) \ \le\ a + b\,\lVert x \rVert}
}{
  \vDash \mathsf{QLinGrowth}(f,\ a,\ b,\ B)
}
\end{mathpar}

\par\medskip
\rulename{QLinGrowthDischarge} concludes the linear-growth-envelope
kind with the constants carried: membership of a polynomial vector
field in the envelope $(a, b)$ over a scope box, from per-coordinate
affine-envelope sign sequents --- exact-rational, checker-replayed.
Downstream cells consume the pair $(a, b)$ arithmetically, never a
membership bit: the kind exists precisely so that growth constants
compose.

\cataloguenote{Entered at the roster census (baked, never catalogued); no Lean
statement binds it yet. The displayed norm form is schematic --- the
envelope rows\textquotesingle{} own sequent shape is the authority.
}
\ruleblock{Theory rows (geolog)}
\sourcenote{The rule\textquotesingle s certificate content as it is declared in the geolog theory the checker loads --- payload schema, side-condition demands, and axiom rows, quoted from the theory source at every build of this catalogue (a stale quotation cannot exist). Axiom modes: \texttt{chase} = the engine derives; \texttt{report} = a violation is a loud judgment; \texttt{refuse} = a violation gates admission. Side-condition relations are discharged by bounded evaluation at the engine\textquotesingle s ground seam, never by derivation.}
\declhead{theory Qtsl}{QLinGrowth}{\badge{BadgeGreen}{quoted from the theory source at render}}
\begin{leansrc}
QLinGrowth    : Set;
\end{leansrc}
\declhead{theory Qtsl}{QLinGrowth/*}{\badge{BadgeGreen}{quoted from the theory source at render}}
\begin{leansrc}
QLinGrowth/judgment : QLinGrowth -> QJudgment;
QLinGrowth/f : QLinGrowth -> PolyVec;
QLinGrowth/a : QLinGrowth -> Q;
QLinGrowth/b : QLinGrowth -> Q;
QLinGrowth/scope : QLinGrowth -> Box;
\end{leansrc}
\declhead{theory Qtsl}{QLinGrowthDischarge}{\badge{BadgeGreen}{quoted from the theory source at render}}
\begin{leansrc}
QLinGrowthDischarge : Set;
\end{leansrc}
\declhead{theory Qtsl}{QLinGrowthDischarge/*}{\badge{BadgeGreen}{quoted from the theory source at render}}
\begin{leansrc}
QLinGrowthDischarge/proof : QLinGrowthDischarge -> QProof;
QLinGrowthDischarge/envelopes : QLinGrowthDischarge -> QProofList;
\end{leansrc}
\ruleblock{Lean statement}
\badge{BadgeAmber}{not yet mechanized} this rule cites no Lean statements: nothing on this page is hash-bound for it, and its dashboard row shows the same absence across the statement, receipt, and correspondence columns.
\ruleblock{HOL4 statement}
\absentbox{NO HOL STATEMENT YET}{No HOL development of the extension theory exists.
}
\ruleblock{HOL verifier receipts}
\absentbox{NONE}{no HOL verification receipt exists for this rule.}
\ruleblock{Proof}

\sourcenote{Genre: core rule (tier 1), the discharge species. This section states the rule's soundness
obligation exactly, per the catalogue's completeness-over-polish
charter (a rule that can only be named-and-gapped is still a row);
the discharge is open debt on the mechanization column, and the
schema quoted in the entry is the baked ground truth.}

\textbf{Soundness obligation.}\enspace The per-coordinate affine-envelope sign sequents on the scope box
jointly certify membership of the field in the linear-growth envelope
with constants $(a, b)$; the kind carries the constants so that
downstream consumers compose them arithmetically. The obligation is
that the envelope rows' sequent semantics entail the envelope
inequality coordinatewise on the scope box --- a per-row sign-sequent
instantiation.

\begin{bremark}[status]
No Lean statement binds this node. The displayed norm form is
schematic; the envelope rows' own sequent shape (exact-rational,
checker-replayed) is the authority.
\end{bremark}


\rulesection{84}{QMarginalCompat}{core rule}
\label{blcatrule-QMarginalCompat}
\ruleblock{The rule (Steele metanotation)}
\sourcenote{as realized in the mechanized development; source row: \texttt{\texttt{QTSL-rules.toml}} --- \emph{QMarginalCompat}. The metanotation statement is typeset by hand for notational quality; this rule cites no Lean statements yet (the absence is stated loudly below). The rule \emph{roster} is machine-checked (every mechanized rule tag round-trips to exactly one catalogue row); the statement \emph{text} is editorial.}
\begin{mathpar}
\inferrule*{
  \overline{\mathrm{dom}_k}\ \ \text{(per-contributor dominance rows)} \\
  \mathsf{qtsl\_stochwire\_shape}(\mathit{cl},\ \mathit{sp})
}{
  \vDash \mathsf{QMarginalCompat}(\mathit{sp},\ \mathit{cl})
}
\end{mathpar}

\par\medskip
\rulename{QMarginalCompat} packages marginal compatibility for
stochastic wires: a belief names convex-analytic content
(the indicator of a named cell set, or a registered function class),
a stochastic wire carries one with its convexity,
lower-semicontinuity, and properness row kits, and this node collects
the per-contributor dominance rows over a shared sample space. The
contributor list\textquotesingle s shape agreement against the space
is an admission demand, discharged by the shape walk quoted below.

\cataloguenote{Entered at the roster census (baked, never catalogued); no Lean
statement binds it yet; display is schema-faithful schematic.
}
\ruleblock{Theory rows (geolog)}
\sourcenote{The rule\textquotesingle s certificate content as it is declared in the geolog theory the checker loads --- payload schema, side-condition demands, and axiom rows, quoted from the theory source at every build of this catalogue (a stale quotation cannot exist). Axiom modes: \texttt{chase} = the engine derives; \texttt{report} = a violation is a loud judgment; \texttt{refuse} = a violation gates admission. Side-condition relations are discharged by bounded evaluation at the engine\textquotesingle s ground seam, never by derivation.}
\declhead{theory Qtsl}{QMarginalCompat}{\badge{BadgeGreen}{quoted from the theory source at render}}
\begin{leansrc}
QMarginalCompat : Set;
\end{leansrc}
\declhead{theory Qtsl}{QMarginalCompat/*}{\badge{BadgeGreen}{quoted from the theory source at render}}
\begin{leansrc}
QMarginalCompat/proof : QMarginalCompat -> QProof;
QMarginalCompat/space : QMarginalCompat -> Shape;
QMarginalCompat/contributors : QMarginalCompat -> QStochContribList;
QMarginalCompat/dominance_rows : QMarginalCompat -> QProofList;
\end{leansrc}
\declhead{theory Qtsl}{qtsl\_stochwire\_shape}{\badge{BadgeGreen}{quoted from the theory source at render}}
\begin{leansrc}
qtsl_stochwire_shape : {contribs: QStochContribList, space: Shape}
  -> Prop;   // side-condition
\end{leansrc}
\declhead{theory Qtsl}{ax/wf/stochwire\_shape}{\badge{BadgeGreen}{quoted from the theory source at render}}
\begin{leansrc}
report ax/wf/stochwire_shape : forall m : QMarginalCompat,
    sp : Shape, cl : QStochContribList.
  m QMarginalCompat/space = sp,
  m QMarginalCompat/contributors = cl
  |- {contribs: cl, space: sp} qtsl_stochwire_shape;
\end{leansrc}
\ruleblock{Lean statement}
\badge{BadgeAmber}{not yet mechanized} this rule cites no Lean statements: nothing on this page is hash-bound for it, and its dashboard row shows the same absence across the statement, receipt, and correspondence columns.
\ruleblock{HOL4 statement}
\absentbox{NO HOL STATEMENT YET}{No HOL development of the extension theory exists.
}
\ruleblock{HOL verifier receipts}
\absentbox{NONE}{no HOL verification receipt exists for this rule.}
\ruleblock{Proof}

\sourcenote{Genre: core rule (tier 1), stochastic wires. This section states the rule's soundness
obligation exactly, per the catalogue's completeness-over-polish
charter (a rule that can only be named-and-gapped is still a row);
the discharge is open debt on the mechanization column, and the
schema quoted in the entry is the baked ground truth.}

\textbf{Soundness obligation.}\enspace The per-contributor dominance rows, over a shared sample space whose
shape agreement with the contributor list is an admission demand,
jointly certify marginal compatibility of the contributors'
beliefs --- the packaging judgment consumed by stochastic wire
composition. The obligation is that dominance per contributor entails
compatibility of the marginals on the shared space.

\begin{bremark}[status]
No Lean statement binds this node; the convex-analytic kits
(convexity, lower semicontinuity, properness) ride the stochastic
wire's own rows and are consumed here by citation.
\end{bremark}


\rulesection{85}{QAggReduce}{core rule}
\label{blcatrule-QAggReduce}
\ruleblock{The rule (Steele metanotation)}
\sourcenote{as realized in the mechanized development; source row: \texttt{\texttt{QTSL-rules.toml}} --- \emph{QAggReduce}. The metanotation statement is typeset by hand for notational quality; its mechanized ground truth is the Lean statement below, extracted from the sources at every build of this book and hash-bound. The rule \emph{roster} is machine-checked (every mechanized rule tag round-trips to exactly one catalogue row); the statement \emph{text} is editorial, checked by eye against the extracted Lean beside it.}
\begin{mathpar}
\inferrule*{
  (\vec{s}, \vec{k}) \in K \vdash G(\vec{s}, \vec{k}) \le 0 \ \text{(one \rulename{BernCone}-family sequent over the aggregate coordinates)} \\
  s^P_j := \textstyle\sum_{t:P} \phi_j(c_t), \qquad k_P := \#P \ \text{(the aggregate coordinates, defined)} \\
  \bigl\{\vdash l_j\, k_P \le s^P_j \le u_j\, k_P\bigr\}_{j, P} \ \text{(coupling rows; } [l_j, u_j] \ \text{a Bernstein enclosure of } \phi_j \ \text{on } B_P\text{)} \\
  \bigl\{\vdash 0 \le k_P \le M_P\bigr\}_{P} \ \text{(count rows)} \\
  \bigl\{\vdash y^{\top} \vec{k} = y^{\top} \vec{k}(x_0)\bigr\}_{y \in Y} \ \text{(P-invariant conservation rows: each cited witness } y \ge 0,\ y^{\top} C = 0 \ \text{verified against the uncoloured skeleton\textquotesingle s incidence matrix } C\text{; } x_0 \ \text{the initial marking)} \\
  \bigl\{\vdash (s^P_\alpha)^2 \le k_P\, s^P_{2\alpha}\bigr\}_{\alpha, P} \ \text{(moment rows, optional tightening)} \\
  \forall x \in \mathrm{scope},\ \forall t{:}P \ \text{in } x:\ c_t \in B_P \ \text{(colour-box scope)}
}{
  x \in \mathrm{scope} \vdash G\bigl(\vec{s}(x), \vec{k}(x)\bigr) \le 0
}
\end{mathpar}

\ruleflavor{the defined aggregate coordinates replace the certificate\textquotesingle s feature dictionary in the display: the dictionary is certificate data (a shared reference) whose entire displayed content is the two defining equations, with $\phi_j$ cited by the coupling rows. The enclosure box discharge is a side condition. The checker never sees a marking: the multiset quotient is compiled away producer-side, aggregate dimension independent of the marking bound $M$; every row family is valid at every reachable in-scope marking under its stated premises (the moment and count-lower rows unconditionally; the count-upper row under the marking bound $M$; conservation rows by their verified witness vectors, the conserved value pinned at the initial marking (the rows ride the count-row proof list --- no dedicated conservation field exists on the baked node, a named schema divergence); coupling rows under their per-token bounds), so adding in-scope rows only shrinks $K$ and never cuts a reachable in-scope aggregate point.}
\par\medskip
\rulename{QAggReduce} is the aggregation rule for nets whose state is a
\emph{marking} --- a multiset of coloured tokens distributed over places,
with variable dimension and a mandatory exchangeability quotient
(tokens within a place are interchangeable, so any legitimate
observable is a symmetric function of each place\textquotesingle s
colour multiset). It replaces an obligation quantified over all
in-scope markings by a single polynomial obligation over finitely many
\emph{aggregate} coordinates: token counts $k_P$ and pooled features
$s^P_j = \sum_{t:P} \phi_j(c_t)$ for a finite feature dictionary. A
producer reaches for it because the alternative --- flattening marking
shapes into modes --- explodes combinatorially in the token bound and
breaks the symmetry; here obligation count and dimension are
independent of the marking bound, which enters only as a constant in
count rows.

What the checker recomputes: one hypothesis-carrying polynomial sequent
$G \le 0$ over the aggregate box, whose hypothesis rows assemble the
containment certificate for a polytope $K$ covering every in-scope
marking\textquotesingle s aggregate image --- coupling rows
$l_j k_P \le s^P_j \le u_j k_P$ from enclosures of each feature on its
place\textquotesingle s colour box, count rows $0 \le k_P \le M_P$ plus
linear conservation laws $y^{\top} \vec{k} = y^{\top} \vec{k}(x_0)$
(place invariants \cite{murata1989}, the conserved value pinned at the
initial marking; each cited witness $y \ge 0$, $y^{\top} C = 0$
checker-verified against the uncoloured skeleton\textquotesingle s
incidence matrix), and
optional moment rows $(s^P_\alpha)^2 \le k_P\, s^P_{2\alpha}$
(Cauchy--Schwarz over the bag). Soundness of the reduction is
evaluation: if the aggregate image lands in $K$ and $G \le 0$ on $K$,
the marking-quantified obligation holds pointwise. The whole content is
the containment certificate, with each row family proven valid at its
stated scope: nonnegativity and moment rows for \emph{every} bag
outright; coupling rows for every bag under their colour-scope premise
(the box discharge); count bounds and conservation rows as certified
reachability facts. Under those side conditions, adding a row shrinks
$K$ without ever cutting a reachable aggregate point.

What would break: a row NOT valid for every reachable marking makes $K$
fail containment and the conclusion simply false --- the central
instance being the coupling rows\textquotesingle{} scope premise, since
an enclosure of $\phi_j$ on the colour box says nothing about tokens
outside it; the box must be \emph{discharged} (exit-guard transitions,
an invariance proof, or stopped-process semantics), never assumed. The
converse failure is honest: not every polytope point is a realizable
marking, so $K$ over-approximates --- a completeness cost, refusal-shaped
and never soundness-shaped, with the row menu as its dial.

\cataloguenote{The operative aggregate is the measurable-query observable
$\langle \mathrm{marking}, f \rangle$ over bags of colours; counts and
pools are the $f \equiv 1$ and $f = \mathrm{id}$ degenerate queries,
and a query is a semiring contraction of the count tensor against $f$
along the colour axis. Semantically the pools are commutative-monoid
folds along the tokens-to-places classifying map, with the multiset as
the universal aggregate \cite{spivak2021agg}; the dictionary\textquotesingle s
expressive ceiling per shape is the classical generation of
multisymmetric invariants by polarized power sums \cite{weyl1939}. The
coupling rows here are deterministic pool--count arithmetic within one
marking --- \emph{not} probabilistic couplings; the identically-named
rows of \rulename{QCoupleDrift} are joint-kernel marginal data over a
product of chains. Same word, disjoint mathematics, neither cites the
other.
}
\ruleblock{Lean statement}
\declhead{BernsteinLogicAux.QtslCompose.AggReduce}{aggReduce}{\badge{BadgeGreen}{statement bound}}
\begin{leansrc}
/-- **QAggReduce, semantic form** ("proof: evaluation"): if `G ≤ 0` on `K`
and every in-scope point's aggregate vector lands in `K`, then the
marking-quantified obligation holds. `agg` is the compiled aggregate map
`x ↦ (s⃗(x), k⃗(x))` — its pointwise containment (`hin`) is what the
coupling/count/moment rows + the P-invariant conservation laws certify. -/
theorem aggReduce {X A : Type*} (agg : X → A) (scope : Set X) (K : Set A)
    (G : A → ℝ) (hG : ∀ a ∈ K, G a ≤ 0) (hin : ∀ x ∈ scope, agg x ∈ K) :
    ∀ x ∈ scope, G (agg x) ≤ 0
\end{leansrc}
\declhead{BernsteinLogicAux.QtslCompose.AggReduce}{coupling\_row}{\badge{BadgeGreen}{statement bound}}
\begin{leansrc}
/-- **The coupling row is valid for every marking** (NOT a probabilistic
coupling — see the module docstring's fence): if every token's `f`-value
lies in `[l, u]`, the pool–count coupling `l · k ≤ s ≤ u · k` holds. Adding
these rows to the polytope `K` never cuts a reachable aggregate point. -/
theorem coupling_row {m : Multiset C} {f : C → ℝ} {l u : ℝ}
    (hl : ∀ c ∈ m, l ≤ f c) (hu : ∀ c ∈ m, f c ≤ u) :
    l * m.card ≤ query m f ∧ query m f ≤ u * m.card
\end{leansrc}
\declhead{BernsteinLogicAux.QtslCompose.AggReduce}{count\_row\_le}{\badge{BadgeGreen}{statement bound}}
\begin{leansrc}
/-- The count row's upper half, from any marking bound (the `M_P` bound is a
reachability/P-invariant fact, consumed as a hypothesis — the box-discharge
side condition [D10], not proven here). -/
theorem count_row_le {m : Multiset C} {M : ℕ} (h : m.card ≤ M) :
    (m.card : ℝ) ≤ M
\end{leansrc}
\declhead{BernsteinLogicAux.QtslCompose.AggReduce}{count\_row\_nonneg}{\badge{BadgeGreen}{statement bound}}
\begin{leansrc}
/-- The count row's lower half: counts are nonnegative. -/
theorem count_row_nonneg (m : Multiset C) : (0 : ℝ) ≤ m.card
\end{leansrc}
\declhead{BernsteinLogicAux.QtslCompose.AggReduce}{moment\_row}{\badge{BadgeGreen}{statement bound}}
\begin{leansrc}
/-- **The moment row is valid for every marking**: `⟨m, f⟩² ≤ k · ⟨m, f²⟩` —
Cauchy–Schwarz over the bag (the rule display's `(s_α)² ≤ k · s_{2α}`,
with `s_α = ⟨m, x^α⟩` instantiated at `f = x^α`). -/
theorem moment_row (m : Multiset C) (f : C → ℝ) :
    (query m f) ^ 2 ≤ m.card * query m (fun c => f c ^ 2)
\end{leansrc}
\declhead{BernsteinLogicAux.QtslCompose.AggReduce}{query\_const\_one}{\badge{BadgeGreen}{statement bound}}
\begin{leansrc}
/-- Counts are the `f ≡ 1` degenerate query (the 14:20 correction's
normal form: `k_P = ⟨m, 1⟩`). -/
theorem query_const_one (m : Multiset C) :
    query m (fun _ => (1 : ℝ)) = m.card
\end{leansrc}
\declhead{BernsteinLogicAux.QtslCompose.AggReduce}{query\_count}{\badge{BadgeGreen}{statement bound}}
\begin{leansrc}
/-- **The contraction-row form** (the γ-trace): the query observable is the
count tensor contracted against `f` along the color axis. -/
theorem query_count [DecidableEq C] (m : Multiset C) (f : C → ℝ) :
    query m f = ∑ c ∈ m.toFinset, (m.count c : ℝ) * f c
\end{leansrc}
\ruleblock{Correspondence (Steele display $\leftrightarrow$ Lean)}
\badge{BadgeGreen}{reviewer-certified against the displayed texts}
\sourcenote{The dictionary below is a reviewer-certified statement of how the mechanized statements above interpret the displayed rule: display metavariable $\mapsto$ Lean binder, each display hypothesis $\mapsto$ the premise or side condition carrying it, the display conclusion $\mapsto$ the goal --- with any deliberate difference in width between display and mechanization named outright. Certification is recorded as a hash triple over (the displayed cell, the extracted statement texts, this dictionary\textquotesingle s own text); if any of the three changes, this block turns amber until a reviewer re-certifies. The correspondence claim itself remains editorial --- reviewer-grade, not machine-checked; its machine-checked successor types each soundness statement against the rule as a structured object, and supersedes this block per rule as it lands.}
The displayed conclusion $\forall x$ in scope: $G(\vec{s}(x),
\vec{k}(x)) \le 0$ $\mapsto$ \texttt{aggReduce}\textquotesingle s goal
\texttt{∀ x ∈ scope, G (agg x) ≤ 0}, with the displayed premise
$\vdash G \le 0$ on $K$ $\mapsto$ \texttt{hG} and the entire row
apparatus $\mapsto$ \texttt{hin : ∀ x ∈ scope, agg x ∈ K} --- the
containment certificate, which the coupling/count/moment/conservation
rows exist to assemble. The binder row is the biggest width fact on
this page: \texttt{aggReduce} is stated for \emph{any} types \texttt{X
A}, any \texttt{agg : X → A}, any \texttt{K : Set A}, any \texttt{G :
A → ℝ} --- no polynomiality of $G$, no polytopality of $K$, no
aggregate box --- so the certified semantic core is strictly wider
than the displayed polynomial instance, safe direction (proof by
evaluation needs none of the structure the checker needs).
\begin{itemize}
\item coupling rows $l_j k_P \le s^P_j \le u_j k_P$ $\mapsto$
  \texttt{coupling\_row}: hypotheses \texttt{∀ c ∈ m, l ≤ f c} and
  \texttt{f c ≤ u} --- the enclosure quantified over tokens
  \emph{in the bag}, which is precisely where the box-discharge side
  condition lands in the certified statement (an enclosure on
  the colour box reaches the hypothesis only for markings confined to
  the box; the flavor text carries this as the box-discharge side
  condition, the statement as the membership guard).
\item count rows $0 \le k_P \le M_P$ $\mapsto$
  \texttt{count\_row\_nonneg} (unconditional) +
  \texttt{count\_row\_le}, whose bound \texttt{m.card ≤ M} is a
  \emph{hypothesis} --- the $M_P$ bound is a reachability/P-invariant
  fact, consumed, never proven. The displayed conservation rows
  $y^{\top} \vec{k} = y^{\top} \vec{k}(x_0)$ (P-invariants, the
  conserved value pinned at the initial marking; witnesses $y \ge 0$,
  $y^{\top} C = 0$ checker-verified) have no bound Lean counterpart at
  all: they are
  reachability facts of the net, outside this module\textquotesingle s
  every-bag warrants. On the baked node they ride the count-row proof
  list --- no dedicated conservation field exists, a named schema
  divergence.
\item moment rows $(s^P_\alpha)^2 \le k_P\, s^P_{2\alpha}$ $\mapsto$
  \texttt{moment\_row}, unconditional for every bag (Cauchy--Schwarz
  by induction), instantiated at \texttt{f = x\^{}α}.
\item the defined aggregate coordinates $s^P_j := \sum_{t:P}
  \phi_j(c_t)$ and $k_P := \#P$ $\mapsto$ per-query functions
  \texttt{f : C → ℝ}, arbitrary --- polynomiality of the $\phi_j$ is
  producer-side; \texttt{query\_count} certifies the
  contraction-row form $\langle m, f\rangle = \sum_c
  \mathrm{count}(c) \cdot f(c)$ (the $\gamma$-trace) under
  \texttt{[DecidableEq C]}, and \texttt{query\_const\_one} pins counts
  as the $f \equiv 1$ degenerate query.
\end{itemize}
Width note: every bound row warrant is per-bag arithmetic --- the
multiset quotient is already taken (\texttt{Multiset C} IS the
exchangeability-quotiented object), so the display\textquotesingle s
``the checker never sees a marking'' has as its Lean face that no
statement here mentions places, nets, or reachability; scope
containment \texttt{hin} is where all of that enters, as one
hypothesis.
The conditional mark\textquotesingle s content, at its exact extent:
the count-bound and conservation rows are reachability facts of the
net --- consumed as hypotheses, never proven here --- and a coupling
row\textquotesingle s colour-box confinement must be discharged by
the producer (exit-guard transitions, an invariance certificate, or
stopped-process semantics), never assumed. Every other row family is
warranted for every bag outright.

\ruleblock{Lean verifier receipts}
\sourcenote{module \texttt{BernsteinLogicAux.QtslCompose.AggReduce}; bound build \texttt{be9dbd259553\ldots} (bound from receipt bd4d8aca0400b822, 2026-07-21T04:51:40Z). A receipt records that named checking layers accepted a build: (module, compiled-artifact SHA-256, environment SHA-256, checker layers), minted when the checking gate ran; the store receipts attesting the bound build are listed with the layers each run credited. Nothing here asks the reader to trust the store: every layer is an independently runnable program over the quoted sources --- rebuild the module and re-run the checkers to reproduce the verdict. Receipts attesting superseded builds credit no layer and are omitted from this rendering (the count of omissions, if any, is noted below).}
\begin{receipttable}
\receiptrow{5f6de8c6a1ff5236}{2026-07-21T01:51:28Z}{build+decorrelated}{\badge{BadgeGreen}{current}}
\receiptrow{bd4d8aca0400b822}{2026-07-21T04:51:40Z}{build+decorrelated}{\badge{BadgeGreen}{current}}
\receiptrow{c37acfaf586210ee}{2026-07-21T05:39:09Z}{build+replay+hardened+decorrelated}{\badge{BadgeGreen}{current}}
\receiptrow{24c5f02176c0ee6a}{2026-07-22T18:27:15Z}{build+decorrelated}{\badge{BadgeGreen}{current}}
\receiptrow{1eaf40d323232bf8}{2026-07-31T22:00:14Z}{build+decorrelated}{\badge{BadgeGreen}{current}}
\receiptrow{47734200303b43f0}{2026-07-31T22:29:07Z}{build+decorrelated}{\badge{BadgeGreen}{current}}
\receiptrow{712878ff235d6c76}{2026-08-02T23:31:19Z}{build+decorrelated}{\badge{BadgeGreen}{current}}
\receiptrow{36b853e33e482b54}{2026-08-03T16:27:22Z}{build+decorrelated}{\badge{BadgeGreen}{current}}
\receiptrow{cb8f537f8267cf3e}{2026-08-03T22:46:06Z}{build+decorrelated}{\badge{BadgeGreen}{current}}
\end{receipttable}
\sourcenote{6 further receipt(s) in the store attest superseded builds of this module and are omitted from this rendering.}
\layerline{\layeritem{build}{lake build (elaboration)}{\badge{BadgeGreen}{current}}\layeritem{replay}{leanchecker (official kernel replay)}{\badge{BadgeGreen}{current}}\layeritem{hardened}{safe\_verify (hardened screen)}{\badge{BadgeGreen}{current}}\layeritem{decorrelated}{lean4export + nanoda (decorrelated kernel)}{\badge{BadgeGreen}{current}}}
\sourcenote{\badge{BadgeGray}{comparator: none} no fixture in the three-implementation comparator harness (compiled / kernel-interpreted / reference-interpreted) exercises this rule.}
\ruleblock{HOL4 statement}
\absentbox{NO HOL STATEMENT YET}{The HOL4 development checks the base calculus\textquotesingle s wire
format (theory \texttt{bern\_catalogue}). \rulename{QAggReduce} is an
extension-edge semantic rule whose node is baked in the extension
theory; what reaches the base checker after the
compile is ordinary hypothesis-carrying \rulename{BernCone} sequents,
which the HOL side already covers as sequents without knowing their
aggregation provenance. A HOL statement \emph{about the rule} becomes
writable when a checker consuming the extension theory\textquotesingle s
rows is ported.
}
\ruleblock{HOL verifier receipts}
\absentbox{NONE}{no HOL verification receipt exists for this rule.}
\ruleblock{Proof}

\sourcenote{Genre: core rule (tier 1). Scope, stated loudly: the
semantic core (the aggregate-space reduction) is proven, and the
row-warrant families are proven at their stated scopes ---
nonnegativity and moment rows for every bag, coupling rows under their
colour-scope premise, while count upper bounds and place-invariant
rows enter as certified in-scope reachability facts, not bagwise
identities; the marking-space compile itself
--- producing the aggregate map and its rows from a net --- is checker
and compiler recomputation, outside these statements. Every hypothesis
is used and remarked on.}

\textbf{Setting and claim.}\enspace The state of a coloured net is a
\emph{marking}: for each place $P$, a finite multiset of colours drawn
from a compact rational box $B_P \subset \R^{n_P}$ --- so the state
space is a disjoint sum, over count vectors, of symmetric powers, with
variable dimension and a mandatory exchangeability quotient (a
``function of the state'' that reads token identities is not a function
on the quotient at all). The certificate language works in
\emph{aggregate coordinates}: counts $k_P = \#P$ and pooled features
$s^P_j = \sum_{t : P} \phi_j(c_t)$ for a finite dictionary
$\Phi = \{\phi_j\}$ of polynomials. An obligation arrives as
``for every in-scope marking $x$:
$G\bigl(\vec{s}(x), \vec{k}(x)\bigr) \leq 0$'' with $G$ polynomial. The
rule replaces marking quantification by one polynomial sequent over a
polytope $K$ in aggregate space, whose hypothesis rows are the
containment certificate. The claim (\texttt{aggReduce} + the row
warrants): the reduction is sound, and every row family below --- each
under its stated scope premise --- may be added to $K$'s description
without ever cutting a reachable aggregate point.

\begin{blemma}[queries, counts, and the contraction form]
\label{lem:qar-query}
For a bag $m$ of colours and $f : \text{colour} \to \R$, write
$\langle m, f \rangle = \sum_{c \in m} f(c)$ (with multiplicity). Then
$\langle m, \mathbf{1} \rangle = \#m$, and
$\langle m, f \rangle = \sum_{c \in \mathrm{supp}(m)}
\mathrm{count}_m(c) \cdot f(c)$ --- the count tensor contracted against
$f$ along the colour axis.
\end{blemma}
\begin{proof}
Both are fibre decompositions of a finite sum: group the terms of
$\sum_{c \in m} f(c)$ by value of $c$. Counts and pools are thus the
$f \equiv 1$ and coordinate-projection instances of one observable
class, and every axis of the contraction is either summed out (colour)
or an output (the query index) --- there is no un-read axis, which is
why, for the chosen finite dictionary, aggregate dimension is the
dictionary's index set: set by the dictionary, never by the marking.
\end{proof}

\begin{blemma}[coupling rows are valid for every bag]\label{lem:qar-couple}
If $l \leq f(c) \leq u$ for every colour $c$ occurring in $m$, then
\[
l \cdot \#m \;\leq\; \langle m, f \rangle \;\leq\; u \cdot \#m .
\]
\end{blemma}
\begin{proof}
A sum of $\#m$ terms, each in $[l, u]$ (induction on the bag). In use,
$[l, u]$ is a certified enclosure of $\phi_j$ on the colour box $B_P$
--- polynomial-enclosure machinery the calculus already owns --- so the
row couples the pool wire to the count wire. Two remarks. The bounds
hypothesis quantifies over colours \emph{in the bag}: the row is
therefore only as good as the guarantee that tokens stay in $B_P$ ---
the scope premise of \autoref{rem:qar-box}. And the degenerate case
is load-bearing: at $k_P = 0$ the two inequalities pin $s^P_j = 0$ ---
emptiness is polytope geometry, not a case split.
\end{proof}

\begin{blemma}[count rows]\label{lem:qar-count}
$0 \leq \#m$ always; and $\#m \leq M$ for any certified marking bound
$M$. Linear conservation laws of the uncoloured skeleton (place
invariants \cite{murata1989}) are further valid count rows.
\end{blemma}
\begin{proof}
Cardinalities are nonnegative integers; the bound and the invariants
are reachability facts of the net, consumed as hypotheses (they are
certified upstream as reachability/place-invariant facts, never
re-proven per obligation). Integer-valuedness of $k_P$ is what
the counting closure discipline (\rulename{QCountAtomRounding})
exploits at complemented count atoms.
\end{proof}

\begin{blemma}[moment rows are valid for every bag]\label{lem:qar-moment}
For every bag $m$ and every $f$:
$\langle m, f \rangle^2 \leq \#m \cdot \langle m, f^2 \rangle$.
\end{blemma}
\begin{proof}
Cauchy--Schwarz over the counting measure on the bag:
$\bigl(\sum_{c \in m} 1 \cdot f(c)\bigr)^2 \leq
\bigl(\sum_{c \in m} 1^2\bigr)\bigl(\sum_{c \in m} f(c)^2\bigr)$.
Instantiated at monomial features this is the family
$(s^P_\alpha)^2 \leq k_P\, s^P_{2\alpha}$ --- optional tightenings:
each one added shrinks the relaxation $K$, and by this lemma never
below the reachable image.
\end{proof}

\begin{bthm}[soundness of \rulename{QAggReduce}]\label{thm:qar-sound}
Let $\mathrm{agg}$ be the aggregate map
$x \mapsto (\vec{s}(x), \vec{k}(x))$, let
$K$ satisfy $\mathrm{agg}(x) \in K$ for every in-scope $x$, and suppose
the sequent $G \leq 0$ holds on $K$. Then
$G(\mathrm{agg}(x)) \leq 0$ for every in-scope $x$.
\end{bthm}
\begin{proof}
Evaluation: $\mathrm{agg}(x) \in K$ and $G \leq 0$ on $K$. That the
proof is one line \emph{is} the honest size of the semantic core ---
the entire content of a \rulename{QAggReduce} node is the containment
certificate for $K$, which Lemmas
\ref{lem:qar-couple}--\ref{lem:qar-moment} assemble from rows each
valid at every reachable marking: a conjunction of per-marking-valid
constraints contains the whole reachable aggregate image, so the single
polynomial sequent over $K$ (discharged by the hypothesis-carrying
cone machinery of the base calculus) covers every in-scope marking at
once. Obligation count and dimension are independent of the marking
bound, which enters only as a constant in Lemma~\ref{lem:qar-count}'s
rows --- the reason this rule exists instead of a mode explosion over
marking shapes.
\end{proof}

\begin{bremark}[the box-discharge scope premise]\label{rem:qar-box}
Lemma~\ref{lem:qar-couple}'s hypothesis --- every occurring colour in
$B_P$ --- is a claim about the \emph{process}, and with nondegenerate
diffusion a process exits every compact box with positive probability.
An enclosure row quantified over states the process leaves is
unsoundness, not incompleteness; so the box must be \emph{discharged},
by one of three named routes: model box exit as a forced transition
(the boundary becomes part of the model, so no in-scope state ever
lies outside the box and the exit's successor region supplies its own
rows), prove invariance of the box for every reachable state
(vanishing diffusion at the boundary plus inward drift --- polynomial,
box-supported conditions landing in the existing row machinery), or
certify the process stopped at exit --- the stopped state kept inside
the box, with the exit event carried in the specification's
probability budget. Which route is a per-model
declaration; that \emph{some} route is declared is a side condition of
the rule. Dropping it admits the counterexample: a generator row verified
on $B_P$ says nothing at the first token that leaves, and the
concluded obligation can be false on reachable states.
\end{bremark}

\begin{bremark}[what the relaxation costs, and what it never costs]
\label{rem:qar-gap}
$K$ over-approximates the reachable aggregate image --- not every
polytope point is a realizable marking. The cost is
\emph{completeness}: a true obligation may fail on a spurious point of
$K$ and the sequent then refuses; it is never soundness, by
Theorem~\ref{thm:qar-sound}. The row menu is the dial: integer
rounding at count atoms, conservation laws, and moment rows each cut
spurious regions. The drain example, worked to its numbers in
\autoref{rem:qar-drain}, is the sharpest small witness: its
interior row is \emph{infeasible} over the real-relaxed
off-target scope $k > 0$ (it evaluates to $+\tfrac14$ at
$k = \tfrac12$, $s_2 = 0$, a point no reachable marking occupies) and
discharges, with an exact linear-multiplier witness, once
\rulename{QCountAtomRounding} tightens the scope to $k \geq 1$ ---
the tightening menu earning its keep on the first toy. Conversely a
row that fails at some reachable in-scope marking (a naively-cleared
mean at $k = 0$, say) would cut reachable points and break
Theorem~\ref{thm:qar-sound}'s containment hypothesis --- validity at
every reachable in-scope marking, under the row family's stated
premises, is the admission criterion for the menu: the nonnegativity
and moment rows meet it for every bag outright, the coupling rows
meet it once their colour-scope premise is discharged, and the count
bounds and invariants meet it as certified reachability facts
consumed as hypotheses.
\end{bremark}

\begin{bremark}[the drain example, worked to its numbers]
\label{rem:qar-drain}
One place, colour $x \in [-2, 2]$, one completion transition consuming
a single token at rate $\lambda(x) = 1 + x^2$, no source; the deadline
certificate is $V = 1 + k$ with $\theta = \tfrac12$ (counts alone
suffice for $V$; the pooled feature $s_2 = \sum_t x_t^2$ enters the
interior row through the rate, so its coupling row joins the scope).
The generator: the continuous part vanishes ($V$ reads no colours),
and each completion removes one token, so
$\mathcal{L}V = \sum_t \lambda(x_t)\bigl(V(k{-}1) - V(k)\bigr)
= -\sum_t (1 + x_t^2) = -(k + s_2)$. The interior row is
\[
\theta V + \mathcal{L}V
\;=\; \tfrac12(1 + k) - (k + s_2)
\;=\; -\tfrac12(k - 1) - s_2 .
\]
Over the real-relaxed off-target scope
$\{\, k > 0,\ 0 \le s_2 \le 4k \,\}$ --- the coupling row
(Lemma~\ref{lem:qar-couple}) at the enclosure $x^2 \in [0, 4]$ on the
colour box --- the row obligation fails: $(k, s_2) = (\tfrac12, 0)$ is
a scope point but no marking (counts are integers), and there the row
evaluates to $\tfrac12 \cdot \tfrac32 - \tfrac12 = +\tfrac14$. Once
\rulename{QCountAtomRounding} rounds the off-target scope
$\mathrm{cmpl}\{k \le 0\}$ to $k \ge 1$, the rewriting displayed above
\emph{is} the exact linear-multiplier witness: multiplier $\tfrac12$
on the scope row $k - 1 \ge 0$ and $1$ on the coupling floor
$s_2 \ge 0$ reproduce the row identically --- a polynomial identity,
no residue, no numerics --- and the certificate is tight at
$(k, s_2) = (1, 0)$: the single-cold-token state, the true bottleneck,
since the slowest completion rate is $\lambda(0) = 1$. Every number
cited in \autoref{rem:qar-gap} recomputes from this display; the
example's other rows (floor, atom, root) are not the story here and
are not staged.
\end{bremark}

\begin{bremark}[relation to the literature]
The aggregate coordinates are the classical multisymmetric power sums:
per fixed shape, polarized power sums generate the full invariant
algebra in characteristic $0$ \cite{weyl1939}, which locates the
dictionary's expressiveness ceiling (generation degree grows with the
shape --- a resolution bound, priced as completeness, never claimed
away). Degree-one certificates in counts alone recover classical
place-invariant analysis \cite{murata1989}. Semantically, pools are
commutative-monoid folds along the tokens-to-places classifying map,
with the multiset as the universal aggregate every other one factors
through \cite{spivak2021agg} --- exchangeability and
factoring-through-the-free-commutative-monoid are the same statement
from two sides. The identically-named ``coupling rows'' of the
relational tier (\rulename{QCoupleDrift}) are probability couplings
over a product of chains; the rows here are deterministic pool--count
arithmetic within one marking. Same word, disjoint mathematics,
deliberately uncited in both directions.
\end{bremark}

\begin{anchors}
Lemma \ref{lem:qar-query} & \texttt{query}; \texttt{query\_const\_one}; \texttt{query\_count}\\
Lemma \ref{lem:qar-couple} & \texttt{coupling\_row} (degenerate pin at $k = 0$: its two halves)\\
Lemma \ref{lem:qar-count} & \texttt{count\_row\_nonneg}; \texttt{count\_row\_le}\\
Lemma \ref{lem:qar-moment} & \texttt{moment\_row}\\
Theorem \ref{thm:qar-sound} & \texttt{aggReduce}\\
integrality at count atoms & \rulename{QCountAtomRounding} section (\texttt{countAtomRounding\_sound}/\texttt{\_eq})\\
\end{anchors}


\rulesection{86}{QCountAtomRounding}{alias rule}
\label{blcatrule-QCountAtomRounding}
\ruleblock{The rule (Steele metanotation)}
\sourcenote{as realized in the mechanized development; source row: \texttt{\texttt{QTSL-rules.toml}} --- \emph{QCountAtomRounding}. The metanotation statement is typeset by hand for notational quality; its mechanized ground truth is the Lean statement below, extracted from the sources at every build of this book and hash-bound. The rule \emph{roster} is machine-checked (every mechanized rule tag round-trips to exactly one catalogue row); the statement \emph{text} is editorial, checked by eye against the extracted Lean beside it.}
\begin{mathpar}
\inferrule*{
  k_P \ \text{a count observable, integer-valued at every reachable marking} \\
  j \in \mathbb{Z}
}{
  \mathrm{cmpl}\,\{k_P \le j\} \ \rightsquigarrow \ \{k_P \ge j + 1\} \quad \text{(hypothesis positions, reachable scope)}
}
\end{mathpar}

\ruleflavor{integrality is the compile\textquotesingle s guarantee. Exact over the reachable locus (the complement over reachable markings \emph{is} $\{k_P \ge j + 1\}$, not $\{k_P > j\}$), and unconditionally never wider ($\{k_P \ge j+1\} \subseteq \{k_P > j\}$) --- rounding to the next integer is free and strictly tighter than the real complement, the counting analogue of the strict-atom closure discipline.}
\par\medskip
\rulename{QCountAtomRounding} is the counting closure discipline: when a
count observable $k_P$ is complemented in a hypothesis position, the
real-arithmetic complement of $\{k_P \le j\}$ --- the open condition
$\{k_P > j\}$ --- is replaced by the closed integer-rounded atom
$\{k_P \ge j + 1\}$. Over reachable markings the two are \emph{equal}
(counts are integers there), and everywhere the rounded atom is
contained in the open one, so the rewrite is free and never wider
than the real complement --- strictly tighter wherever the slab is
realized: substituting it in hypothesis position never widens the
constraint set, and the gained slab $j < k_P < j + 1$
(unreachable, but visible to real-arithmetic relaxations) is often
exactly what makes a downstream combination certificate feasible. It is
the counting analogue of the strict-atom closure discipline for
continuous observables, with integrality replacing topology.

There is nothing for the checker to recompute beyond the syntactic
rewrite itself: producers apply it before emission, and what the wire
carries is the already-tightened atom inside an aggregation
rule\textquotesingle s hypothesis rows. The one genuine hypothesis is
integrality --- that $k_P$ takes integer values at every reachable
marking. That is a guarantee of the marking-space compile (counts are
cardinalities), not of the rewrite; in the mechanization it enters as a
named hypothesis, discharged definitionally when the net layer lands.
Without it the exactness identity fails (a fractional
\textquotedblleft count\textquotedblright{} at $k_P = j + \tfrac12$
satisfies the true complement but not the rounded atom), and only the
containment direction survives --- which is not enough in hypothesis
position: the rounded scope is \emph{smaller}, so a row verified only
on it silently skips exactly such fractional points, and a misapplied
rounding can silently accept rather than refuse (the proof
section\textquotesingle s integrality remark carries the witness). The
integrality warrant is load-bearing for soundness, not a convenience.

\cataloguenote{Load-bearing on the first worked aggregation example: a drain
net\textquotesingle s interior row is infeasible over the real-relaxed
off-target scope $k > 0$ (it fails at $k = \tfrac12$) and discharges
with the rounded scope $k \ge 1$, via an exact linear-multiplier
witness. The rewrite rides with \rulename{QAggReduce}\textquotesingle s
family: its target is an atom-denotation identity consumed by count
rows, not a separate rule node --- one atom in, one atom out.
}
\ruleblock{Lean statement}
\declhead{BernsteinLogicAux.Structural}{countAtomRounding\_sound}{\badge{BadgeGreen}{statement bound}}
\begin{leansrc}
/-- **P3.3 CountAtomRounding**, the soundness side (hypothesis position):
every reachable point of the true complement of `{k ≤ j}` satisfies the
tightened atom `{j + 1 ≤ k}`. This is the side the drain example's interior
row consumes — the aggregate polytope `K` must contain every in-scope
reachable point, and it does. -/
theorem countAtomRounding_sound (k : X → ℝ) (j : ℤ) (Reach : Set X)
    (hbridge : ∀ x ∈ Reach, ∃ n : ℤ, k x = (n : ℝ)) :
    ∀ x ∈ Reach, x ∈ {x | k x ≤ (j : ℝ)}ᶜ → (j : ℝ) + 1 ≤ k x
\end{leansrc}
\declhead{BernsteinLogicAux.Structural}{countAtomRounding\_tighter}{\badge{BadgeGreen}{statement bound}}
\begin{leansrc}
/-- **P3.3**, the free-tightening direction: the rounded atom is contained
in the real complement unconditionally — substituting it never widens the
constraint set, so the rewrite is strictly tighter (by the cut
`(j, j+1)`), which is what buys the drain example its Farkas
feasibility. Stated at any real threshold `j`: the containment needs no
integrality — only exactness (`countAtomRounding_eq`) and soundness
(`countAtomRounding_sound`) ride the integer-threshold bridge. -/
theorem countAtomRounding_tighter (k : X → ℝ) (j : ℝ) :
    {x | j + 1 ≤ k x} ⊆ {x | j < k x}
\end{leansrc}
\declhead{BernsteinLogicAux.Structural}{countAtomRounding\_eq}{\badge{BadgeGreen}{statement bound}}
\begin{leansrc}
/-- **P3.3**, exactness on the reachable locus: over integer-valued points
the rewrite loses nothing — the catalogue's "the complement ... over
reachable markings IS `{k_P ≥ j + 1}`", verbatim as a set identity. -/
theorem countAtomRounding_eq (k : X → ℝ) (j : ℤ) (Reach : Set X)
    (hbridge : ∀ x ∈ Reach, ∃ n : ℤ, k x = (n : ℝ)) :
    Reach ∩ {x | k x ≤ (j : ℝ)}ᶜ = Reach ∩ {x | (j : ℝ) + 1 ≤ k x}
\end{leansrc}
\ruleblock{Correspondence (Steele display $\leftrightarrow$ Lean)}
\badge{BadgeGreen}{reviewer-certified against the displayed texts}
\sourcenote{The dictionary below is a reviewer-certified statement of how the mechanized statements above interpret the displayed rule: display metavariable $\mapsto$ Lean binder, each display hypothesis $\mapsto$ the premise or side condition carrying it, the display conclusion $\mapsto$ the goal --- with any deliberate difference in width between display and mechanization named outright. Certification is recorded as a hash triple over (the displayed cell, the extracted statement texts, this dictionary\textquotesingle s own text); if any of the three changes, this block turns amber until a reviewer re-certifies. The correspondence claim itself remains editorial --- reviewer-grade, not machine-checked; its machine-checked successor types each soundness statement against the rule as a structured object, and supersedes this block per rule as it lands.}
The entry\textquotesingle s two claims (exactness and never-wider,
now the flavor text) and the displayed integrality premise
map onto the three bound statements. The compile\textquotesingle s
integrality guarantee --- the rule\textquotesingle s System bridge,
as the roster says in words --- $\mapsto$ \texttt{hbridge : ∀ x ∈
Reach, ∃ n : ℤ, k x = (n : ℝ)}: counts integer-valued at every
reachable point, a named hypothesis because lean-aux lacks the net
layer that guarantees it. Then: the hypothesis-position rewrite
$\mapsto$ \texttt{countAtomRounding\_sound} (every reachable point of
the true complement of $\{k \le j\}$ satisfies $\{j + 1 \le k\}$ ---
the side the drain example\textquotesingle s interior row consumes);
``unconditionally never wider'' $\mapsto$
\texttt{countAtomRounding\_tighter}, whose binder row is the width
find: NO integrality, NO reachable scope --- the containment $\{j+1
\le k\} \subseteq \{j < k\}$ holds for arbitrary real observables,
so the substitution never \emph{widens} anywhere; soundness of the
hypothesis-position substitution itself rides the exactness side ---
the integrality hypothesis --- certified separately as
\texttt{countAtomRounding\_eq} --- the
catalogue\textquotesingle s ``the complement over reachable markings
IS $\{k_P \ge j+1\}$'' verbatim as a set identity on the reachable
locus. Width note: \texttt{k : X → ℝ} throughout;
\texttt{\_tighter} now takes $j : \mathbb{R}$ outright (the
containment needs no integrality, so the anchor is stated at the
proposition\textquotesingle s full width), while \texttt{\_sound} and
\texttt{\_eq} keep the $j : \mathbb{Z}$ cast their integer-threshold
bridge demands ---
the count-observable grammar (the $f \equiv 1$ query) is upstream;
divergence D2\textquotesingle s question (which side is load-bearing)
is answered by the bind structure itself: soundness rides
\texttt{\_sound}, the free tightening rides \texttt{\_tighter}, and
each is consumable independently.

\ruleblock{Lean verifier receipts}
\sourcenote{module \texttt{BernsteinLogicAux.Structural}; bound build \texttt{1222cdcc326f\ldots} (bound from receipt cb0c3820c6877f22, 2026-07-29T01:03:36Z). A receipt records that named checking layers accepted a build: (module, compiled-artifact SHA-256, environment SHA-256, checker layers), minted when the checking gate ran; the store receipts attesting the bound build are listed with the layers each run credited. Nothing here asks the reader to trust the store: every layer is an independently runnable program over the quoted sources --- rebuild the module and re-run the checkers to reproduce the verdict. Receipts attesting superseded builds credit no layer and are omitted from this rendering (the count of omissions, if any, is noted below).}
\begin{receipttable}
\receiptrow{0b7e68907f45b27f}{2026-07-28T23:48:36Z}{build+decorrelated}{\badge{BadgeGreen}{current}}
\receiptrow{e97d5614d4f08c7e}{2026-07-29T00:00:35Z}{build+decorrelated}{\badge{BadgeGreen}{current}}
\receiptrow{27de21f047e6b194}{2026-07-29T00:49:05Z}{build+replay+hardened+decorrelated}{\badge{BadgeGreen}{current}}
\receiptrow{cb0c3820c6877f22}{2026-07-29T01:03:36Z}{build+replay+hardened+decorrelated}{\badge{BadgeGreen}{current}}
\receiptrow{3530841119a9a607}{2026-07-31T22:00:22Z}{build+decorrelated}{\badge{BadgeGreen}{current}}
\receiptrow{b6cc23eb3e5b7f97}{2026-07-31T22:29:16Z}{build+decorrelated}{\badge{BadgeGreen}{current}}
\receiptrow{bf3dacc3ab7f20ea}{2026-08-02T23:31:31Z}{build+decorrelated}{\badge{BadgeGreen}{current}}
\receiptrow{384a8569d2005df7}{2026-08-03T16:27:35Z}{build+decorrelated}{\badge{BadgeGreen}{current}}
\receiptrow{717e07acf9acac8a}{2026-08-03T22:46:19Z}{build+decorrelated}{\badge{BadgeGreen}{current}}
\end{receipttable}
\sourcenote{11 further receipt(s) in the store attest superseded builds of this module and are omitted from this rendering.}
\layerline{\layeritem{build}{lake build (elaboration)}{\badge{BadgeGreen}{current}}\layeritem{replay}{leanchecker (official kernel replay)}{\badge{BadgeGreen}{current}}\layeritem{hardened}{safe\_verify (hardened screen)}{\badge{BadgeGreen}{current}}\layeritem{decorrelated}{lean4export + nanoda (decorrelated kernel)}{\badge{BadgeGreen}{current}}}
\sourcenote{\badge{BadgeGray}{comparator: none} no fixture in the three-implementation comparator harness (compiled / kernel-interpreted / reference-interpreted) exercises this rule.}
\ruleblock{HOL4 statement}
\absentbox{NO HOL STATEMENT YET}{The HOL4 development checks the base calculus\textquotesingle s wire
format (theory \texttt{bern\_catalogue}). This rewrite is applied
producer-side, before emission: even with the aggregation node baked,
the wire is expected to carry the already-tightened atom, so the HOL
statement that would matter is the aggregation
sequent\textquotesingle s --- and no checker consuming the extension
theory\textquotesingle s rows is ported yet.
Nothing exists for a HOL statement to be about; if a future format ever
carries the un-rounded atom, the rewrite\textquotesingle s exactness
identity becomes the portable statement.
}
\ruleblock{HOL verifier receipts}
\absentbox{NONE}{no HOL verification receipt exists for this rule.}
\ruleblock{Proof}

\sourcenote{Genre: alias rule (tier 3) --- a respelling template. One
variation from the alias norm, stated up front: the unpack's target is
an \emph{atom-denotation identity} consumed by
\rulename{QAggReduce}'s count rows (the pattern of the strict-atom
closure moves), not a separate rule node --- one atom in, one atom out,
producers rewrite before emission. The five alias artefacts follow in
that key: the unpack, the filled-in application, the payload transport,
the exactness identity, and the soundness corollary.}

\begin{brecipe}[the respelling]\label{rec:car-unpack}
The unpack acts on a complemented count atom in hypothesis position:
\[
\mathrm{cmpl}\,\{\, k_P \leq j \,\}
\qquad\rightsquigarrow\qquad
\{\, k_P \geq j + 1 \,\},
\qquad j \in \mathbb{Z},
\]
where $k_P$ is a count observable --- integer-valued at every reachable
marking, because counts are cardinalities (the compile's integrality
guarantee, entering the mechanization as a named hypothesis until the
net layer discharges it definitionally). Payload transport, field for
field:
\begin{center}\small
\begin{tabular}{@{}lll@{}}
\textbf{field} & \textbf{before} & \textbf{after}\\
\midrule
observable & $k_P$ & $k_P$ (verbatim)\\
threshold & $j$ & $j + 1$ (integer bump)\\
orientation & complemented $\leq$ (open $>$) & closed $\geq$\\
scope & reachable markings & reachable markings (verbatim)\\
\end{tabular}
\end{center}
The filled-in application is a scope row of an aggregation obligation:
where a compile node scopes a row ``off the target $\{k_P \leq j\}$'',
the row's hypothesis list carries $k_P \geq j + 1$ in place of the
real-arithmetic complement.
\end{brecipe}

\begin{bprop}[exactness over the reachable locus]\label{prop:car-eq}
On integer-valued points, for $j \in \mathbb{Z}$,
$\mathrm{Reach} \cap \mathrm{cmpl}\{k_P \leq j\}
= \mathrm{Reach} \cap \{k_P \geq j + 1\}$. (Both restrictions are
load-bearing: at a non-integer threshold the integer $n = \lceil j
\rceil$ would satisfy $n > j$ but not $n \geq j + 1$.)
\end{bprop}
\begin{proof}
At a reachable point, $k_P = n$ for some integer $n$; and for integers,
$n > j \iff n \geq j + 1$. Two lines, which is the point: the
complement of $\{k_P \leq j\}$ \emph{over reachable markings} is
$\{k_P \geq j+1\}$, not $\{k_P > j\}$ --- nothing is lost.
\end{proof}

\begin{bprop}[free tightening, unconditionally]\label{prop:car-tight}
$\{k_P \geq j + 1\} \subseteq \{k_P > j\}$ on \emph{all} points,
integer-valued or not.
\end{bprop}
\begin{proof}
$j + 1 \leq k_P$ implies $j < k_P$. So substituting the rounded atom in
hypothesis position never widens the constraint set; what it removes
is exactly the slab $j < k_P < j + 1$ --- strictly, whenever that
slab's preimage is nonempty: unreachable by any marking, but fully
visible to a real-arithmetic relaxation, which is where the gain
lands.
\end{proof}

\begin{bprop}[soundness corollary]\label{prop:car-sound}
Every reachable point of the true complement of $\{k_P \leq j\}$
satisfies the tightened atom; hence any row valid on
$\mathrm{Reach} \cap \{k_P \geq j + 1\}$ is valid on the reachable part
of the true complement --- the direction hypothesis positions consume,
and by Proposition~\ref{prop:car-eq} nothing more is being assumed than
the complement itself, reachably.
\end{bprop}
\begin{proof}
Immediate from Proposition~\ref{prop:car-eq}, left-to-right. The
asymmetry of positions is worth one sentence: in \emph{hypothesis}
position a tighter scope is free (fewer points to verify, no reachable
point dropped, by exactness); in \emph{conclusion} position a tightened
atom would be a stronger claim and is not licensed by this rule ---
the alias is a hypothesis-position discipline, exactly like its
continuous sibling.
\end{proof}

\begin{bremark}[why integrality is not decoration]\label{rem:car-int}
Drop it and exactness fails: a fractional ``count'' at
$k_P = j + \tfrac12$ satisfies the true complement but not the rounded
atom, so a row verified only on the rounded scope would silently skip
it --- if such points were reachable, that is unsoundness in hypothesis
position. The integrality hypothesis is precisely the statement that
they are not: counts are cardinalities of token multisets. Note there is \emph{no} graceful failure mode when the rewrite is
misapplied to a non-count observable: in hypothesis (scope) position
the rounded atom \emph{shrinks} the region a row must be verified on,
so misapplication can silently accept --- the half-integer point above
is exactly the witness. The integrality warrant is load-bearing, not a
convenience.
\end{bremark}

\begin{bremark}[load-bearing on the first worked example]
The drain net's deadline certificate needs its interior row off the
target $\{k \leq 0\}$. Over the real-relaxed scope $k > 0$ the row is
infeasible --- it evaluates to $+\tfrac14$ at $(k, s_2) =
(\tfrac12, 0)$, a point no marking occupies --- while over the rounded
scope $k \geq 1$ it discharges with an exact linear-multiplier witness
(both numbers recompute in the \rulename{QAggReduce} section's
drain-example remark, where the row and its multipliers are displayed).
The rounding is the counting analogue of the strict-atom closure
discipline (polarity-directed rounding of strict literals), with
integrality playing the role topology plays there: both convert an
open condition into the closed atom the row machinery wants, at zero
semantic cost on the states that matter.
\end{bremark}

\begin{anchors}
Proposition \ref{prop:car-eq} & \texttt{countAtomRounding\_eq}\\
Proposition \ref{prop:car-tight} & \texttt{countAtomRounding\_tighter}\\
Proposition \ref{prop:car-sound} & \texttt{countAtomRounding\_sound}\\
integrality bridge & the named hypothesis in \texttt{countAtomRounding\_sound} and \texttt{\_eq} (compile-side guarantee); deliberately absent from \texttt{\_tighter}, whose containment is unconditional\\
consumed by & \rulename{QAggReduce} scope rows (count atoms); the drain example's interior row\\
\end{anchors}


\rulesection{87}{QPAlwaysFromFold}{core rule}
\label{blcatrule-QPAlwaysFromFold}
\ruleblock{The rule (Steele metanotation)}
\sourcenote{as realized in the mechanized development; source row: \texttt{\texttt{QTSL-rules.toml}} --- \emph{QPAlwaysFromFold}. The metanotation statement is typeset by hand for notational quality; this rule cites no Lean statements yet (the absence is stated loudly below). The rule \emph{roster} is machine-checked (every mechanized rule tag round-trips to exactly one catalogue row); the statement \emph{text} is editorial.}
\begin{mathpar}
\inferrule*{
  \mathrm{upper} : \mathsf{qpaf\_upper}(n,\ c_{\mathrm{u}}) \\
  \mathrm{lower} : \mathsf{qpaf\_lower}(n,\ c_{\mathrm{l}}) \\
  \mathsf{qtsl\_palways\_fold}(n)
}{
  \mathcal{M} \vDash \mathsf{SProbGE}\bigl(p,\ \mathsf{PAlways}(\mathsf{SAnd}[\mathsf{SAtomLE}(c, V),\ \mathsf{SAtomGE}(c, V)])\bigr)
}
\end{mathpar}

\par\medskip
\rulename{QPAlwaysFromFold} ties plated-fold safety evidence to the
judgment stratum: it concludes a surely-invariant claim --- a
conserved-charge invariant, held with probability one --- over the
demonic composite of a token-population subject, from certificates
about a plated fold of the population. Coverage telescopes footprint
rows \emph{pathwise} across every resolution of every clock wing: the
telescoping argument reads firings, not clocks, so race, demonic, and
forced transitions are covered uniformly and no scheduler resolution
escapes. Initial membership is checked inside the evaluator\textquotesingle
s walk (the deadline-table precedent): there is no warrant row, no
region-level child, and no boxed region --- the subject\textquotesingle
s marking is unbounded, and a boxed region for it would be a type
error, not a convenience.

A producer cites one certificate per direction: an upper-bound
citation for the $\le$ face, a lower-bound citation for the $\ge$
face --- at least one, at most one row each; the lower\textquotesingle
s features are recomputed by the walk, never trusted. The accepted
conclusion shapes are one atom or the displayed two-sided conjunction
with equal polynomials; anything wider refuses loudly (a narrow walk
mis-verifies nothing).

\cataloguenote{The verification content lives in the evaluator\textquotesingle s
walk, whose semantics are normative in the plated-transfer design:
conclusion resolution, spec shape, atom decomposition, cert
citations, subject agreement including the plate-to-population token
bindings, fold scale, polynomial correspondence against the
count-wire enumeration, init membership from the bag marking,
sibling-claim presence, and the polarity roster --- with the
pathwise-species demand on covered transitions: an expectation-species
footprint row never warrants a surely-invariant. No Lean statement
binds the node yet; the schema and walk are baked, the walk mutant is
red in the theory battery, and the mechanization column is honest
open debt.
}
\ruleblock{Theory rows (geolog)}
\sourcenote{The rule\textquotesingle s certificate content as it is declared in the geolog theory the checker loads --- payload schema, side-condition demands, and axiom rows, quoted from the theory source at every build of this catalogue (a stale quotation cannot exist). Axiom modes: \texttt{chase} = the engine derives; \texttt{report} = a violation is a loud judgment; \texttt{refuse} = a violation gates admission. Side-condition relations are discharged by bounded evaluation at the engine\textquotesingle s ground seam, never by derivation.}
\par\medskip
The baked rows: the node, its direction citations, the walk, and the
plate-population binding the walk\textquotesingle s subject-agreement
step reads.

\declhead{theory Qtsl}{QPAlwaysFromFold}{\badge{BadgeGreen}{quoted from the theory source at render}}
\begin{leansrc}
QPAlwaysFromFold : Set;
\end{leansrc}
\declhead{theory Qtsl}{QPAlwaysFromFold/*}{\badge{BadgeGreen}{quoted from the theory source at render}}
\begin{leansrc}
QPAlwaysFromFold/proof : QPAlwaysFromFold -> QProof;
\end{leansrc}
\declhead{theory Qtsl}{qpaf\_upper}{\badge{BadgeGreen}{quoted from the theory source at render}}
\begin{leansrc}
qpaf_upper : {node: QPAlwaysFromFold, cert: HCert} -> Prop;
\end{leansrc}
\declhead{theory Qtsl}{qpaf\_lower}{\badge{BadgeGreen}{quoted from the theory source at render}}
\begin{leansrc}
qpaf_lower : {node: QPAlwaysFromFold, cert: HCert} -> Prop;
\end{leansrc}
\declhead{theory Qtsl}{ax/wf/qpalwaysfold\_walk}{\badge{BadgeGreen}{quoted from the theory source at render}}
\begin{leansrc}
report ax/wf/qpalwaysfold_walk : forall n : QPAlwaysFromFold.
  |- {node: n} qtsl_palways_fold;
\end{leansrc}
\declhead{theory Qtsl}{hplate\_token}{\badge{BadgeGreen}{quoted from the theory source at render}}
\begin{leansrc}
hplate_token : {plate: HPlate, token: Token} -> Prop;
\end{leansrc}
\ruleblock{Lean statement}
\badge{BadgeAmber}{not yet mechanized} this rule cites no Lean statements: nothing on this page is hash-bound for it, and its dashboard row shows the same absence across the statement, receipt, and correspondence columns.
\ruleblock{HOL4 statement}
\absentbox{NO HOL STATEMENT YET}{No HOL development of the extension theory exists; this
node\textquotesingle s content is walk structure over baked rows.
}
\ruleblock{HOL verifier receipts}
\absentbox{NONE}{no HOL verification receipt exists for this rule.}
\ruleblock{Proof}

\sourcenote{Genre: core rule (tier 1), the plated-transfer genus (safety face). This section states the rule's soundness
obligation exactly, per the catalogue's completeness-over-polish
charter (a rule that can only be named-and-gapped is still a row);
the discharge is open debt on the mechanization column, and the
schema quoted in the entry is the baked ground truth.}

\textbf{Soundness obligation.}\enspace With the walk's ten steps discharged --- conclusion resolution, spec
shape, atom decomposition, direction citations, subject agreement
including plate-to-population token bindings, fold scale, polynomial
correspondence against the count-wire enumeration, initial membership
from the bag marking, sibling-claim presence, and the polarity roster
with its pathwise-species demand --- the cited fold certificates'
footprint rows telescope pathwise across every resolution of every
clock wing, so the conserved-charge invariant holds along every run:
the concluded claim holds surely, and any probability threshold in
$(0,1]$ is sound.

\begin{bremark}[status]
No Lean statement binds this node. The walk semantics are
normative in the plated-transfer design; the walk mutant is red in
the theory battery, which certifies the walk's presence, not the
transfer theorem --- the transfer theorem's mechanization is the open
debt.
\end{bremark}


\rulesection{88}{QFireBoundFromFold}{core rule}
\label{blcatrule-QFireBoundFromFold}
\ruleblock{The rule (Steele metanotation)}
\sourcenote{as realized in the mechanized development; source row: \texttt{\texttt{QTSL-rules.toml}} --- \emph{QFireBoundFromFold}. The metanotation statement is typeset by hand for notational quality; this rule cites no Lean statements yet (the absence is stated loudly below). The rule \emph{roster} is machine-checked (every mechanized rule tag round-trips to exactly one catalogue row); the statement \emph{text} is editorial.}
\begin{mathpar}
\inferrule*{
  \mathrm{cert} : \mathsf{HCert} \\
  d_{\min} > 0 \\
  \mathsf{qtsl\_firebound\_fold}(n)
}{
  \mathcal{M} \vDash \mathsf{SReturnLE}\bigl(B,\ \mathsf{retFireCount}(\mathit{trs})\bigr)
}
\end{mathpar}

\par\medskip
\rulename{QFireBoundFromFold} is the ranking face of the same genus:
from \emph{one} gap-rowed plated certificate it concludes a roster
firing-count ceiling --- the specification\textquotesingle s warranted
component is the count bound over the roster firing-count run
functional. The transfer theorem\textquotesingle s pathwise face,
applied: with footprint, coverage, floor, and flow discharged, every
roster firing drops the fold by at least $d_{\min}$ in the reading
scale, the fold never rises anywhere else, and the fold is floored at
zero --- so roster firings number at most
$\mathrm{fold}_\varphi(m_0)/d_{\min}$ along \emph{every} run,
uniformly in token count and interleaving.

What the checker recomputes: two exact inequalities --- each roster
row\textquotesingle s gap at least $d_{\min}$, and $B \cdot d_{\min}$
at least the fold\textquotesingle s initial value --- and never an
argmin: $d_{\min}$ is pinned by the producer, the same witness
species as a pinned tolerance. Counting is firing-metered per roster
member --- deliberately not region-entry counting (the crossing-count
acceptance device) and not the instant-crossing count (the
crossing-budget family\textquotesingle s compensator-mass species):
three counting species, three devices, none interchangeable.

\cataloguenote{Closed-judgment route, deliberately: no warrant row, no region-level
child, no boxed certificate-family table for an unbounded marking;
init arithmetic lives inside the walk. The cert citation is a
function, not a relation --- a ranking consequence cites exactly one
certificate, and exactly-one is semantic here (totality is the
sort\textquotesingle s own obligation); contrast the safety
sibling\textquotesingle s two-direction relation shape. No Lean
statement binds the node yet; the schema, the positivity report row,
and the walk are baked with red mutants in the theory battery, and
the mechanization column is honest open debt.
}
\ruleblock{Theory rows (geolog)}
\sourcenote{The rule\textquotesingle s certificate content as it is declared in the geolog theory the checker loads --- payload schema, side-condition demands, and axiom rows, quoted from the theory source at every build of this catalogue (a stale quotation cannot exist). Axiom modes: \texttt{chase} = the engine derives; \texttt{report} = a violation is a loud judgment; \texttt{refuse} = a violation gates admission. Side-condition relations are discharged by bounded evaluation at the engine\textquotesingle s ground seam, never by derivation.}
\par\medskip
The baked rows: the node with its pinned decrement, the positivity
report, the walk, and the roster relation the concluded run
functional reads.

\declhead{theory Qtsl}{QFireBoundFromFold}{\badge{BadgeGreen}{quoted from the theory source at render}}
\begin{leansrc}
QFireBoundFromFold : Set;
\end{leansrc}
\declhead{theory Qtsl}{QFireBoundFromFold/*}{\badge{BadgeGreen}{quoted from the theory source at render}}
\begin{leansrc}
QFireBoundFromFold/proof : QFireBoundFromFold -> QProof;
QFireBoundFromFold/cert  : QFireBoundFromFold -> HCert;
QFireBoundFromFold/dmin  : QFireBoundFromFold -> Q;
\end{leansrc}
\declhead{theory Qtsl}{ax/wf/qfirebound\_dmin\_pos}{\badge{BadgeGreen}{quoted from the theory source at render}}
\begin{leansrc}
report ax/wf/qfirebound_dmin_pos : forall n : QFireBoundFromFold,
    d : Q.
  n QFireBoundFromFold/dmin = d
  |- {l: 0, r: d} q_lt;
\end{leansrc}
\declhead{theory Qtsl}{ax/wf/qfirebound\_walk}{\badge{BadgeGreen}{quoted from the theory source at render}}
\begin{leansrc}
report ax/wf/qfirebound_walk : forall n : QFireBoundFromFold.
  |- {node: n} qtsl_firebound_fold;
\end{leansrc}
\declhead{theory Qtsl}{retFireCount}{\badge{BadgeGreen}{quoted from the theory source at render}}
\begin{leansrc}
retFireCount : {trs: TrList} -> PathObservable;
\end{leansrc}
\declhead{theory Qtsl}{qtsl\_firecount\_roster}{\badge{BadgeGreen}{quoted from the theory source at render}}
\begin{leansrc}
qtsl_firecount_roster : {sys: System, trs: TrList} -> Prop;
\end{leansrc}
\ruleblock{Lean statement}
\badge{BadgeAmber}{not yet mechanized} this rule cites no Lean statements: nothing on this page is hash-bound for it, and its dashboard row shows the same absence across the statement, receipt, and correspondence columns.
\ruleblock{HOL4 statement}
\absentbox{NO HOL STATEMENT YET}{No HOL development of the extension theory exists; this
node\textquotesingle s content is walk structure over baked rows.
}
\ruleblock{HOL verifier receipts}
\absentbox{NONE}{no HOL verification receipt exists for this rule.}
\ruleblock{Proof}

\sourcenote{Genre: core rule (tier 1), the plated-transfer genus (ranking face). This section states the rule's soundness
obligation exactly, per the catalogue's completeness-over-polish
charter (a rule that can only be named-and-gapped is still a row);
the discharge is open debt on the mechanization column, and the
schema quoted in the entry is the baked ground truth.}

\textbf{Soundness obligation.}\enspace With footprint, coverage, floor, and flow discharged for the one
gap-rowed certificate, every roster firing drops the fold by at least
$d_{\min}$ in the reading scale, the fold never rises anywhere else,
and the fold is floored at zero; hence roster firings number at most
$\mathrm{fold}_\varphi(m_0)/d_{\min}$ along every run, uniformly in
token count and interleaving --- the concluded firing-count ceiling.
The checker's contribution is two exact inequalities (each roster
row's gap at least $d_{\min}$; the ceiling times $d_{\min}$ at least
the fold's initial value); $d_{\min}$ is pinned by the producer and
never derived.

\begin{bremark}[status]
No Lean statement binds this node. The walk semantics are
normative in the transfer-former design (twelve steps); the theory
battery's mutants are red at the walk and positivity loci.
\end{bremark}


\rulesection{89}{QPAlwaysTFromFold}{core rule}
\label{blcatrule-QPAlwaysTFromFold}
\ruleblock{The rule (Steele metanotation)}
\sourcenote{as realized in the mechanized development; source row: \texttt{\texttt{QTSL-rules.toml}} --- \emph{QPAlwaysTFromFold}. The metanotation statement is typeset by hand for notational quality; this rule cites no Lean statements yet (the absence is stated loudly below). The rule \emph{roster} is machine-checked (every mechanized rule tag round-trips to exactly one catalogue row); the statement \emph{text} is editorial.}
\begin{mathpar}
\inferrule*{
  n\,\mathsf{cert} = c \\ n\,\mathsf{goalfold} = f_g \\ n\,\mathsf{horizon} = T \\
  \mathsf{coverage\_lb}(c) \\ \mathsf{flow\_lb}(c) \\
  \forall r.\ u_r > 0 \Rightarrow \mathsf{TrMeter}(\mathit{tr}(r), d_r) \\
  \mathsf{qtsl\_palwayst\_fold}(n)
}{
  \mathcal{M} \vDash \mathsf{SProbGE}\bigl(p,\ \mathsf{PAlwaysT}(T,\ \mathrm{cmpl}(\mathit{goal}))\bigr)
}
\end{mathpar}

\par\medskip
\rulename{QPAlwaysTFromFold} is the timed-floor face of the
plated-transfer genus --- the third sibling beside
\rulename{QPAlwaysFromFold} (safety) and \rulename{QFireBoundFromFold}
(ranking): from one certificate whose footprint rows carry the
\emph{drop-at-most} polarity (per firing, the fold decreases by at
most $u_r$ --- including $u_r = 0$, never decreases), together with
per-row clock floors, it concludes that no run of the system reaches
a designated goal region within a horizon $T$ --- a surely-satisfied
time-bounded invariance claim over the goal\textquotesingle s
complement.

The soundness assembly is the transfer telescope\textquotesingle s
mirror face plus a meter law: with drop-at-most coverage of every
may-touch transition, mirror flow, per-row footprint sequents, and
--- for conjugated reading kinds --- the conjugation-parameter
conformance discharged (integrality of conjugated powers and the
order-preservation domain; the transformed-space telescope is
unjustified without it), every run satisfies
$\mathrm{fold}_\varphi(m_t) \ge \mathrm{fold}_\varphi(m_0) - \sum_r
u_r N_r(t)$; with clock floors $d_r > 0$ on every counted row,
$N_r(t) \le \max(0, (t - o_r)/d_r) + 1$ under every resolution. A run
reaching the goal region (where the fold is at most $f_g$) by time
$T$ would force $\mathrm{fold}_\varphi(m_0) - f_g \le \sum_{r : u_r >
0} u_r \cdot (\max(0, (T - o_r)/d_r) + 1)$ --- and the
checker\textquotesingle s whole trusted contribution is ONE exact
rational verification of the strict reverse inequality. $T$, $f_g$,
and every $u_r$ are pinned by the producer; the checker never
computes an argmin or an optimum.

Counting from below stays deliberately \emph{inside} the walk: the
specification language refuses a standing ``at least $n$ events
occur\textquotesingle\textquotesingle{} former (a liveness claim
consuming survival mass), and this rule honors that refusal --- the
conclusion is safety-polarity, violated by finite prefixes. The
conclusion is Zeno-robust without any crossing budget: rows with
$u_r = 0$ (instantaneous rows included) may fire unboundedly, since
they cannot drop the fold; an instantaneous row with $u_r > 0$
refuses outright (it has no time floor to meter).

\cataloguenote{The schema is baked in the extension theory: the node, the
fourteen-step walk (normative in the judgment-tie design), the two
mirror side-conditions, the per-wing meter demands, and the
drop-at-most footprint polarity are all landed rows, quoted below
from the theory source. The polarity\textquotesingle s formation
carried its guard obligations into the bake --- most critically that
the \emph{existing} upper-coverage side condition counts exactly
rows of the two elder polarities (safety, gap), stated as a closed
positive membership rather than a prohibition: a drop-at-most row
bounds nothing above, and an evaluator that counted one would unsound
the two elder siblings.
The mechanization ladder names the discrete-unfolding profile first
(round-indexed runs); the continuous lift waits on a first
formalization of the clock meter semantics, shared with
\rulename{QRateFromClock} and \rulename{QPAlwaysTFromLineage}. No
Lean statement binds the node yet; the mechanization column is
honest open debt.
}
\ruleblock{Theory rows (geolog)}
\sourcenote{The rule\textquotesingle s certificate content as it is declared in the geolog theory the checker loads --- payload schema, side-condition demands, and axiom rows, quoted from the theory source at every build of this catalogue (a stale quotation cannot exist). Axiom modes: \texttt{chase} = the engine derives; \texttt{report} = a violation is a loud judgment; \texttt{refuse} = a violation gates admission. Side-condition relations are discharged by bounded evaluation at the engine\textquotesingle s ground seam, never by derivation.}
\par\medskip
The baked rows: the node with its pinned horizon and goal-fold
fields, the optional lineage-offsets citation, the walk, the two
mirror side-conditions on the cited certificate, the per-wing meter
demands, and the drop-at-most polarity with its displacement bound.

\declhead{theory Qtsl}{QPAlwaysTFromFold}{\badge{BadgeGreen}{quoted from the theory source at render}}
\begin{leansrc}
QPAlwaysTFromFold : Set;
\end{leansrc}
\declhead{theory Qtsl}{QPAlwaysTFromFold/*}{\badge{BadgeGreen}{quoted from the theory source at render}}
\begin{leansrc}
QPAlwaysTFromFold/proof : QPAlwaysTFromFold -> QProof;
QPAlwaysTFromFold/cert  : QPAlwaysTFromFold -> HCert;
QPAlwaysTFromFold/goalfold : QPAlwaysTFromFold -> Q;   // f_g, pinned
QPAlwaysTFromFold/horizon  : QPAlwaysTFromFold -> Q;   // T, pinned
\end{leansrc}
\declhead{theory Qtsl}{ax/wf/qpatf\_horizon\_nonneg}{\badge{BadgeGreen}{quoted from the theory source at render}}
\begin{leansrc}
report ax/wf/qpatf_horizon_nonneg : forall n : QPAlwaysTFromFold,
    x : Q.
  n QPAlwaysTFromFold/horizon = x
  |- {l: 0, r: x} q_le;
\end{leansrc}
\declhead{theory Qtsl}{qpatf\_offsets}{\badge{BadgeGreen}{quoted from the theory source at render}}
\begin{leansrc}
qpatf_offsets : {node: QPAlwaysTFromFold,
                 lineage: QPAlwaysTFromLineage} -> Prop;
\end{leansrc}
\declhead{theory Qtsl}{qtsl\_palwayst\_fold}{\badge{BadgeGreen}{quoted from the theory source at render}}
\begin{leansrc}
qtsl_palwayst_fold : {node: QPAlwaysTFromFold} -> Prop;
\end{leansrc}
\declhead{theory Qtsl}{ax/tf/qpatf\_walk}{\badge{BadgeGreen}{quoted from the theory source at render}}
\begin{leansrc}
report ax/tf/qpatf_walk : forall n : QPAlwaysTFromFold.
  |- {node: n} qtsl_palwayst_fold;
\end{leansrc}
\declhead{theory Qtsl}{ax/tf/qpatf\_coverage\_lb}{\badge{BadgeGreen}{quoted from the theory source at render}}
\begin{leansrc}
report ax/tf/qpatf_coverage_lb : forall n : QPAlwaysTFromFold,
    c : HCert.
  n QPAlwaysTFromFold/cert = c
  |- {cert: c} qtsl_hplate_coverage_lb;
\end{leansrc}
\declhead{theory Qtsl}{ax/tf/qpatf\_flow\_lb}{\badge{BadgeGreen}{quoted from the theory source at render}}
\begin{leansrc}
report ax/tf/qpatf_flow_lb : forall n : QPAlwaysTFromFold,
    c : HCert.
  n QPAlwaysTFromFold/cert = c
  |- {cert: c} qtsl_hplate_flow_lb;
\end{leansrc}
\declhead{theory Qtsl}{ax/tf/qpatf\_meter\_forced}{\badge{BadgeGreen}{quoted from the theory source at render}}
\begin{leansrc}
report ax/tf/qpatf_meter_forced : forall n : QPAlwaysTFromFold,
    c : HCert, r : HFootprintRow, t : Transition, d : Q,
    k : TrClock.
  n QPAlwaysTFromFold/cert = c,
  r HFootprintRow/cert = c,
  r HFootprintRow/polarity = hrpDropAtMost,
  r HFootprintRow/tr = t,
  t Transition/clock = k,
  k = {dwell: d} tkForced
  |- {tr: t, d: d} qtsl_tr_meter;
\end{leansrc}
\declhead{theory Qtsl}{ax/tf/qpatf\_meter\_demonic}{\badge{BadgeGreen}{quoted from the theory source at render}}
\begin{leansrc}
report ax/tf/qpatf_meter_demonic : forall n : QPAlwaysTFromFold,
    c : HCert, r : HFootprintRow, t : Transition, d : Q,
    e : EnnRat, k : TrClock.
  n QPAlwaysTFromFold/cert = c,
  r HFootprintRow/cert = c,
  r HFootprintRow/polarity = hrpDropAtMost,
  r HFootprintRow/tr = t,
  t Transition/clock = k,
  k = {dwellMin: e} tkDemonic,
  e = {q: d} enFin
  |- {tr: t, d: d} qtsl_tr_meter;
\end{leansrc}
\declhead{theory Qtsl}{ax/tf/qpatf\_meter\_carved}{\badge{BadgeGreen}{quoted from the theory source at render}}
\begin{leansrc}
report ax/tf/qpatf_meter_carved : forall n : QPAlwaysTFromFold,
    c : HCert, r : HFootprintRow, t : Transition, fl : Q,
    ce : EnnRat, k : TrClock.
  n QPAlwaysTFromFold/cert = c,
  r HFootprintRow/cert = c,
  r HFootprintRow/polarity = hrpDropAtMost,
  r HFootprintRow/tr = t,
  t Transition/clock = k,
  k = {floor: fl, ceil: ce} tkCarved
  |- {tr: t, d: fl} qtsl_tr_meter;
\end{leansrc}
\declhead{theory Qtsl}{hrpDropAtMost}{\badge{BadgeGreen}{quoted from the theory source at render}}
\begin{leansrc}
hrpDropAtMost : HRowPolarity;
\end{leansrc}
\declhead{theory Qtsl}{qtsl\_hplate\_coverage\_lb}{\badge{BadgeGreen}{quoted from the theory source at render}}
\begin{leansrc}
qtsl_hplate_coverage_lb : {cert: HCert} -> Prop;
\end{leansrc}
\declhead{theory Qtsl}{qtsl\_hplate\_flow\_lb}{\badge{BadgeGreen}{quoted from the theory source at render}}
\begin{leansrc}
qtsl_hplate_flow_lb : {cert: HCert} -> Prop;
\end{leansrc}
\ruleblock{Lean statement}
\badge{BadgeAmber}{not yet mechanized} this rule cites no Lean statements: nothing on this page is hash-bound for it, and its dashboard row shows the same absence across the statement, receipt, and correspondence columns.
\ruleblock{HOL4 statement}
\absentbox{NO HOL STATEMENT YET}{No HOL development of the extension theory exists; this
node\textquotesingle s content is walk structure over baked rows.
}
\ruleblock{HOL verifier receipts}
\absentbox{NONE}{no HOL verification receipt exists for this rule.}
\ruleblock{Proof}

\sourcenote{Genre: core rule (tier 1), the plated-transfer genus
(timed-floor face). This section states the rule's soundness
obligation exactly, per the catalogue's completeness-over-polish
charter; the rule's schema and walk are baked in the extension
theory (rows quoted in this section), and the discharge below is the
named obligation of the mechanization column.}

\textbf{Soundness obligation (mirror telescope + meter law).}\enspace
Let $F := \mathrm{fold}_\varphi(m_0)$, recomputed by the walk from the
declared initial marking. (Telescope.) With drop-at-most coverage of
every may-touch transition, per-row footprint sequents (produced-side
fold contribution at least consumed side minus $u_r$, in the reading
scale, under the guard hypotheses, for all colors in the boxes),
mirror flow (the fold does not decrease between firings), and --- for
conjugated reading kinds --- the conjugation-parameter conformance
(integrality of conjugated powers; the order-preservation domain,
even powers requiring certified nonnegative feature ranges: without
it the transformed-space reading of the telescope is unjustified)
discharged,
every run satisfies
$\mathrm{fold}_\varphi(m_t) \ge F - \sum_r u_r N_r(t)$, where
$N_r(t)$ counts firings of row $r$ by time $t$. (Meter.) With every
counted row ($u_r > 0$) carrying a clock floor $d_r > 0$ and
first-firing offset $o_r$ (from a cited lineage node, else the
one-dwell default $o_r := d_r$),
$N_r(t) \le \max(0, (t - o_r)/d_r) + 1$ under every resolution.
(Assembly.) The goal region carries a certified fold ceiling $f_g$;
a run reaching it at $t \le T$ forces
$F - f_g \;\le\; \sum_r u_r N_r(t) \;\le\; \sum_{r : u_r > 0} u_r
\cdot \bigl(\max(0, \tfrac{T - o_r}{d_r}) + 1\bigr)$,
using $u_r \ge 0$ throughout and monotonicity in $t$; the
checker-verified strict inequality
$\sum_{r : u_r > 0} u_r \cdot (\max(0, \tfrac{T - o_r}{d_r}) + 1) <
F - f_g$ refutes it, so no admitted run reaches the goal region
within $T$ --- the concluded time-bounded invariance, surely.

The conclusion is safety-polarity and consumes no crossing budget:
rows with $u_r = 0$ (instantaneous rows included) may fire
unboundedly without touching the telescope's right side; an
instantaneous row with $u_r > 0$ refuses at the walk's meter step.

\begin{bremark}[status]
The schema is baked: the node, its fourteen-step walk, the mirror
side-conditions, the per-wing meter demands, and the drop-at-most
polarity are landed theory rows (quoted in this section). No Lean
statement binds the node yet: the named obligation is the
discrete-unfolding profile first (round-indexed runs), with the
continuous lift behind the clock meter semantics' first
formalization --- the keystone shared across this rule family.
\end{bremark}


\rulesection{90}{QPAlwaysTFromLineage}{core rule}
\label{blcatrule-QPAlwaysTFromLineage}
\ruleblock{The rule (Steele metanotation)}
\sourcenote{as realized in the mechanized development; source row: \texttt{\texttt{QTSL-rules.toml}} --- \emph{QPAlwaysTFromLineage}. The metanotation statement is typeset by hand for notational quality; this rule cites no Lean statements yet (the absence is stated loudly below). The rule \emph{roster} is machine-checked (every mechanized rule tag round-trips to exactly one catalogue row); the statement \emph{text} is editorial.}
\begin{mathpar}
\inferrule*{
  \mathsf{unit\_lineage}(c) \\
  \{\mathsf{hrpGap}(1), \mathsf{hrpDropAtMost}(1)\}\ \text{per advancing row on } \chi = -\psi\ \text{(pathwise species)} \\
  \mathsf{TrMeter}(\mathit{tr}(r), d)\ \text{per advancing row} \\
  \forall \tau \in m_0.\ \psi(\tau) \le 0 \\
  T < k \cdot d \\
  \mathsf{qtsl\_palwayst\_lineage}(n)
}{
  \mathcal{M} \vDash \mathsf{SProbGE}\bigl(p,\ \mathsf{PAlwaysT}(T,\ \textstyle\bigwedge_{q\,:\,\mathrm{depth}(q) \ge k} \mathrm{count}_q = 0)\bigr)
}
\end{mathpar}

\par\medskip
\rulename{QPAlwaysTFromLineage} is a causal-depth meter: tokens whose
production chain has depth $k$ cannot exist strictly before time
$k \cdot d$. It concludes plate-emptiness --- no token of declared
causal depth at least $k$ exists at any time $t \le T$, for a pinned
horizon $T$ checked STRICTLY below $k \cdot d$ (the induction proves
emptiness through round $k - 1$ and permits a depth-$k$ token in the
round-$k$ marking, so the inclusive endpoint is deliberately not
concluded; on the round grid $T = (k-1) \cdot d$ is the tightest
warranted choice, and $k$ is a positive integer) --- from four
ingredients: \emph{unit lineage} (every covered transition consumes
at most one and produces at most one token across the cited plates,
so per-token lineage is a path: for a row consuming a cited token and
producing a cited successor, the fold-level displacement of the depth
dictionary is the successor\textquotesingle s depth minus the
consumed token\textquotesingle s; birth rows root at depth at most
one, and a death row on a positive-depth token satisfies neither
bracket and is refused --- a completeness cost, never an
unsoundness); a
\emph{bracket} of footprint rows on the negated depth dictionary
pinning each chain firing\textquotesingle s advance to exactly one
(the drop-at-most row and the gap row at equal witness --- both sides
load-bearing, since a decrease bound alone would admit one firing
advancing depth by ten --- and every bracket row of the pathwise
deterministic-reset species: a stochastic moment row could satisfy
both sequents at expected advance one while individual outcomes jump
farther, defeating the surely conclusion); clock floors $d$ on every
advancing row; and
initial depth at most zero, recomputed by the walk from the declared
initial marking (universally over its tokens --- vacuous when empty).

The rule has a second face: a \rulename{QPAlwaysTFromFold} node may
cite it as the source of its per-row first-firing offsets ($o_r$ :=
the row\textquotesingle s consumed-tuple minimum declared depth over
the guard box, times $d$ --- recomputed by the CONSUMING
walk, never read from stored producer data, because an inflated
offset errs on the unsound side; a row whose guard box admits no
consumed cited token takes the consumer\textquotesingle s one-dwell
default instead). This face is unaffected by the strict endpoint: a
firing that consumes a depth-$k$ token happens at round $k$ or
later, time at least $k \cdot d$, which is exactly the offset claim.

The rule mints \emph{profile-gated}, and the gate is stated honestly
rather than smoothed over: per-token chain spacing is NOT derivable
from the per-transition meter law when several tokens share a guard
--- a transition\textquotesingle s accrued dwell can be spent on a
freshly produced token. This is the classical timed-Petri-net
memory-policy fork (single-server versus per-token timers). The
profile, exactly: every covered row carries the one period $d$ as its
dwell parameter (demonic and forced wings admitted; a carved wing at
floor equal to ceiling $d$, so firings cannot fall off the round
grid); the subject is declared at the round-indexed step reading ---
firings at round $r$ read the round-$r$ marking and produce into
round $r + 1$, rounds at times $r \cdot d$, the theory\textquotesingle s
own discrete-time-unfolding reading, whose canonical case is the
all-demonic uniform-dwell net; and the cited plates\textquotesingle{}
token colors do not evolve between firings (the brackets constrain
firing displacement only, so between-round color drift would evade
them). Under that profile the per-token spacing IS the step
law\textquotesingle s own content and the soundness induction is
clean; the continuous-time lift is a genuine semantics decision the
proposal prices without baking.

\cataloguenote{The schema is baked in the extension theory: the node, its walk, and
its structural side-conditions are quoted below from the theory
source. Profile-gated from birth, per the catalogue\textquotesingle s
precedent of rules proven at a named profile and filed beyond it; the
concluded specification former itself is profile-independent --- the
gate lives in the walk\textquotesingle s side conditions, never in
the conclusion\textquotesingle s meaning. A worked instance of the
discrete induction (a fabrication-line model with a 370-deep
production chain) is already kernel-checked; the oracle obligation
that generalizes it to the profile rule remains owed, and no Lean
statement binds the node yet --- the mechanization column is honest
open debt.
}
\ruleblock{Theory rows (geolog)}
\sourcenote{The rule\textquotesingle s certificate content as it is declared in the geolog theory the checker loads --- payload schema, side-condition demands, and axiom rows, quoted from the theory source at every build of this catalogue (a stale quotation cannot exist). Axiom modes: \texttt{chase} = the engine derives; \texttt{report} = a violation is a loud judgment; \texttt{refuse} = a violation gates admission. Side-condition relations are discharged by bounded evaluation at the engine\textquotesingle s ground seam, never by derivation.}
\par\medskip
The baked rows: the node with its positive depth, positive period,
and nonnegative horizon reports, the unit-lineage structural
side-condition on the cited certificate, and the causal-depth walk.

\declhead{theory Qtsl}{QPAlwaysTFromLineage}{\badge{BadgeGreen}{quoted from the theory source at render}}
\begin{leansrc}
QPAlwaysTFromLineage : Set;
\end{leansrc}
\declhead{theory Qtsl}{QPAlwaysTFromLineage/*}{\badge{BadgeGreen}{quoted from the theory source at render}}
\begin{leansrc}
QPAlwaysTFromLineage/proof : QPAlwaysTFromLineage -> QProof;
QPAlwaysTFromLineage/cert  : QPAlwaysTFromLineage -> HCert;
QPAlwaysTFromLineage/depth   : QPAlwaysTFromLineage -> Nat;
QPAlwaysTFromLineage/period  : QPAlwaysTFromLineage -> Q;
QPAlwaysTFromLineage/horizon : QPAlwaysTFromLineage -> Q;
\end{leansrc}
\declhead{theory Qtsl}{ax/wf/qpatl\_depth\_pos}{\badge{BadgeGreen}{quoted from the theory source at render}}
\begin{leansrc}
report ax/wf/qpatl_depth_pos : forall n : QPAlwaysTFromLineage,
    k : Nat.
  n QPAlwaysTFromLineage/depth = k
  |- {l: 0, r: k} nat_lt;
\end{leansrc}
\declhead{theory Qtsl}{ax/wf/qpatl\_period\_pos}{\badge{BadgeGreen}{quoted from the theory source at render}}
\begin{leansrc}
report ax/wf/qpatl_period_pos : forall n : QPAlwaysTFromLineage,
    d : Q.
  n QPAlwaysTFromLineage/period = d
  |- {l: 0, r: d} q_lt;
\end{leansrc}
\declhead{theory Qtsl}{ax/wf/qpatl\_horizon\_nonneg}{\badge{BadgeGreen}{quoted from the theory source at render}}
\begin{leansrc}
report ax/wf/qpatl_horizon_nonneg : forall n : QPAlwaysTFromLineage,
    x : Q.
  n QPAlwaysTFromLineage/horizon = x
  |- {l: 0, r: x} q_le;
\end{leansrc}
\declhead{theory Qtsl}{qtsl\_unit\_lineage}{\badge{BadgeGreen}{quoted from the theory source at render}}
\begin{leansrc}
qtsl_unit_lineage : {cert: HCert} -> Prop;
\end{leansrc}
\declhead{theory Qtsl}{qtsl\_palwayst\_lineage}{\badge{BadgeGreen}{quoted from the theory source at render}}
\begin{leansrc}
qtsl_palwayst_lineage : {node: QPAlwaysTFromLineage} -> Prop;
\end{leansrc}
\declhead{theory Qtsl}{ax/tf/qpatl\_walk}{\badge{BadgeGreen}{quoted from the theory source at render}}
\begin{leansrc}
report ax/tf/qpatl_walk : forall n : QPAlwaysTFromLineage.
  |- {node: n} qtsl_palwayst_lineage;
\end{leansrc}
\declhead{theory Qtsl}{ax/tf/qpatl\_unit\_lineage}{\badge{BadgeGreen}{quoted from the theory source at render}}
\begin{leansrc}
report ax/tf/qpatl_unit_lineage : forall n : QPAlwaysTFromLineage,
    c : HCert.
  n QPAlwaysTFromLineage/cert = c
  |- {cert: c} qtsl_unit_lineage;
\end{leansrc}
\ruleblock{Lean statement}
\badge{BadgeAmber}{not yet mechanized} this rule cites no Lean statements: nothing on this page is hash-bound for it, and its dashboard row shows the same absence across the statement, receipt, and correspondence columns.
\ruleblock{HOL4 statement}
\absentbox{NO HOL STATEMENT YET}{No HOL development of the extension theory exists; this
node\textquotesingle s content is walk structure over baked rows.
}
\ruleblock{HOL verifier receipts}
\absentbox{NONE}{no HOL verification receipt exists for this rule.}
\ruleblock{Proof}

\sourcenote{Genre: core rule (tier 1), the plated-transfer genus
(causal-depth face), profile-gated from birth. This section states
the rule's soundness obligation exactly, gate included; the rule's
schema and walk are baked in the extension theory (rows quoted in
this section), and the discharge below is the named obligation of
the mechanization column.}

\textbf{Soundness obligation (depth induction at the
discrete-unfolding profile).}\enspace At the gated profile (every
covered row carrying the one period $d$ as its dwell parameter ---
demonic and forced wings admitted, carved wings at floor equal to
ceiling $d$; the subject declared at the round-indexed step reading,
firings at round $r$ reading the round-$r$ marking and producing
into round $r + 1$, rounds at times $r \cdot d$; cited
plates' token colors static between firings), by
induction over rounds: unit lineage makes each token's production
history a path through covered firings; the bracket rows on the
negated depth dictionary --- every one of the pathwise
deterministic-reset species --- pin each chain firing's depth advance
to exactly one (the gap row bounds the advance from below, the
drop-at-most row from above --- both sides load-bearing); initial
depth is at most zero, universally over the declared initial
marking's tokens (vacuous when empty), by the walk's recomputation.
Hence a token present at round $m$ arose through at
most $m$ chain firings in strictly increasing rounds and has declared
depth at most $m$ --- so plates of declared depth at least $k$
($k \ge 1$) are empty at rounds $0, \ldots, k-1$: emptiness at every
time $t \le T$ for the pinned horizon $T < k \cdot d$, surely. The
endpoint is strict by construction --- the induction permits a
depth-$k$ token in the round-$k$ marking itself, so the inclusive
horizon $k \cdot d$ is deliberately not concluded.

\textbf{The gate, exactly.}\enspace Beyond the profile, per-token
chain spacing is NOT a consequence of the per-transition meter law
when tokens share a guard: a transition's accrued dwell can be spent
on a freshly produced token, collapsing the chain's temporal spacing
to one dwell for arbitrarily many chain steps. This is the classical
timed-Petri-net memory-policy fork (single-server versus per-token
timers); the continuous-time lift is a semantics decision the
proposal prices without baking, and the rule refuses outside the
profile rather than concluding from an unproven spacing premise.

\begin{bremark}[status]
The schema is baked: the node, its walk, and its structural
side-conditions are landed theory rows (quoted in this section). No
Lean statement binds the node yet. A worked instance of the discrete
induction (a fabrication-line model with a 370-deep production
chain) is already kernel-checked; the named obligation generalizes
it to the profile rule.
\end{bremark}


\blcatsection{Distance Transport}
\rulesection{91}{QCoupleDrift}{core rule}
\label{blcatrule-QCoupleDrift}
\ruleblock{The rule (Steele metanotation)}
\sourcenote{as realized in the mechanized development; source row: \texttt{\texttt{QTSL-rules.toml}} --- \emph{QCoupleDrift}. The metanotation statement is typeset by hand for notational quality; its mechanized ground truth is the Lean statement below, extracted from the sources at every build of this book and hash-bound. The rule \emph{roster} is machine-checked (every mechanized rule tag round-trips to exactly one catalogue row); the statement \emph{text} is editorial, checked by eye against the extracted Lean beside it.}
\begin{mathpar}
\inferrule*{
  \bigl\{\vdash K(p, \cdot) \ge 0,\ \textstyle\sum_q K(p, q) = 1,\ \pi_{i\#} K(p, \cdot) = \widehat{K}_i(p_i, \cdot)\bigr\} \ \text{(coupling rows, exact } \Q\text{; } \widehat{K}_i \ \text{the completed component kernels)} \\
  \bigl\{\vdash K(p, p) = 1 \ \text{for both-absorbed } p\bigr\} \ \text{(freeze rows)} \\
  \bigl\{\vdash \textstyle\sum_q K(p, q)\, V(q) \le V(p)\bigr\} \ \text{(relational supermartingale generator row, every product state)} \\
  \vdash V \ge 0 \\
  \bigl\{\vdash V \ge 1 \ \text{(disagreeing), resp.\ } V \ge |v_1 - v_2|, \ \text{on both-absorbed states}\bigr\} \ \text{(terminal domination)} \\
  \text{absorption: per-component reach criterion}; \quad \mu_0 \ \text{a probability law with } \pi_{i\#}\mu_0 = \text{the component initial laws (initial marginal rows)}; \quad \vdash \textstyle\sum_p \mu_0(p)\, V(p) \le b
}{
  \mathsf{QDist}[d](\mathcal{M}_1, \mathcal{M}_2) \le b \qquad d \in \{\mathrm{TV},\ \text{gap-pseudometrics}\}
}
\end{mathpar}

\ruleflavor{the coupling marginals are tied to the \emph{completed} component kernels $\widehat{K}_i$ (absorbing states frozen --- the completion carries the component freezes), and the absorption criterion is recomputed by the checker (a per-component graph reach). Coupling emptiness refuses, never lies; the $\mathbb{W}_1$ registry entry and the mixture-polynomial tier are designed, not yet realized --- the finite-chain tier above is the proven fragment.}
\par\medskip
\rulename{QCoupleDrift} is the coupling-witnessed discharge of the
distance judgment: from a certified coupling of two absorbing Markov
chains and a relational supermartingale on the product, it concludes a
bound on how far the two systems\textquotesingle{} final readouts can
differ --- total-variation distance of the final readout laws, or the
gap of final expected values. A producer reaches for it to compare two
models of the same system (a nominal chain against a faulted one, a
fine model against a simplification): the producer chooses a coupling
--- a product-space kernel whose marginals are the two chains --- that
keeps the pair together where the chains agree, and a nonnegative
function $V$ that decreases along the coupled dynamics and dominates
the disagreement indicator at absorption. The certified bound is
$V$\textquotesingle s initial expectation.

What the checker recomputes, all in exact rational arithmetic: that the
coupling rows form a genuine kernel (entrywise nonnegative, unit row
sums) whose row marginals are exactly the two completed component
kernels; that both-absorbed states are frozen; the drift inequality
$\sum_q K(p,q)\, V(q) \le V(p)$ at every product state; nonnegativity
of $V$; the terminal-domination rows ($V \ge 1$ on disagreeing
both-absorbed states for the total-variation form, $V \ge |v_1 - v_2|$
for the expectation-gap form); the per-component absorption criterion
as a finite-graph reachability check; and the initial bound
$\sum_p \mu_0(p)\, V(p) \le b$.

Each hypothesis is load-bearing, and the proof section walks the
omission counterexamples: broken marginals sever the product chain from
the systems being compared; without nonnegativity the drift induction
cannot be localized to the absorbing set; without terminal domination
the coupling inequality has nothing to consume; without absorption the
final pair law is not a coupling of the final component laws (mass that
never both-absorbs escapes the accounting). A coupling that cannot be
completed refuses --- the coupling is an input premise, so when none
exists no certificate can be formed: the rule never lies by emptiness.

\cataloguenote{The judgment-primary form ($\mathsf{QDist}[d] \le \varepsilon$, one
conclusion per certificate) deliberately differs from pre/post-style
relational program logics \cite{hsu2017}: it composes at judgment level
(triangle, weaken, symmetry where flagged), and expectation-gap
readouts enter as gap-pseudometric registry entries rather than as a
second judgment form. The proven fragment is the finite-chain
certificate tier, and its conclusion now reads in general vocabulary:
the certified bound dominates the Jordan-decomposition total-variation
distance of the final readout laws as measures (see the width note);
the mixture-polynomial tier and
the $\mathbb{W}_1$ registry entry are designed, not yet realized. A
pilot study ran four exact-rational certificates against an independent
checker with a mutation battery; on correlated faulting the relational
route certified a total-variation bound roughly three times tighter
than the best single-chain route --- the reason the tier exists.
}
\ruleblock{Lean statement}
\declhead{BernsteinLogicAux.TvDist}{exists\_hahn}{\badge{BadgeGreen}{statement bound}}
\begin{leansrc}
/-- **The Hahn-set bridge**: for finite measures of equal total mass there
is a measurable set `i` on which the difference attains `tvDist`, and which
dominates the difference on every measurable set. This is the single lemma
through which all Jordan-decomposition bookkeeping flows; everything
downstream (metric axioms, coupling inequality, finite-tier seam) consumes
only its statement. -/
theorem exists_hahn (μ ν : Measure X) [IsFiniteMeasure μ] [IsFiniteMeasure ν]
    (hmass : μ.real univ = ν.real univ) :
    ∃ i : Set X, MeasurableSet i ∧ μ.real i - ν.real i = tvDist μ ν ∧
      ∀ A : Set X, MeasurableSet A → μ.real A - ν.real A ≤ μ.real i - ν.real i
\end{leansrc}
\declhead{BernsteinLogicAux.TvDist}{tvDist\_le\_coupling}{\badge{BadgeGreen}{statement bound}}
\begin{leansrc}
/-- **The coupling inequality**: if `π` on `X × X` has marginals `μ` and
`ν`, then `d_TV(μ,ν) ≤ π {r | r.1 ≠ r.2}`. The off-diagonal is evaluated by
the *outer* measure, so no measurability of the diagonal (and no separation
hypothesis on `X`) is needed; the marginal conditions themselves force the
total masses to agree. -/
theorem tvDist_le_coupling {μ ν : Measure X} [IsFiniteMeasure μ] [IsFiniteMeasure ν]
    (π : Measure (X × X)) (hπ₁ : π.map Prod.fst = μ) (hπ₂ : π.map Prod.snd = ν) :
    tvDist μ ν ≤ (π {r : X × X | r.1 ≠ r.2}).toReal
\end{leansrc}
\declhead{BernsteinLogicAux.TvDist}{tvDist\_ofFn}{\badge{BadgeGreen}{statement bound}}
\begin{leansrc}
/-- **The finite-tier compatibility lemma**: on a finite discrete space the
general TV distance of the induced measures is the checker's own
`(1/2)·L¹` total variation. Requires only nonnegativity and equal totals
(sub-probability pairs of equal mass included). -/
theorem tvDist_ofFn {p q : L → ℝ} (hp0 : ∀ l, 0 ≤ p l) (hq0 : ∀ l, 0 ≤ q l)
    (htot : ∑ l, p l = ∑ l, q l) :
    tvDist (ofFn p) (ofFn q) = QRelChain.tvDist p q
\end{leansrc}
\declhead{BernsteinLogicAux.TvDist}{sound\_QRelTV\_measure}{\badge{BadgeGreen}{statement bound}}
\begin{leansrc}
/-- **`sound_QRelTV`, conclusion re-scoped to the general object**: under
exactly the hypotheses of `QRelChain.sound_QRelTV` — the checker's
obligations C1–C9, nothing added — the certified bound dominates the
*general* total-variation distance (`TV.tvDist`, the Jordan/coupling
object of this module) between the final readout laws, viewed as measures
on the readout space. `Fintype` survives only where the certificate
arithmetic lives (the chains and the readout space); the conclusion is the
unrestricted object. -/
theorem sound_QRelTV_measure {L : Type*} [Fintype L] [DecidableEq L]
    [MeasurableSpace L] [MeasurableSingletonClass L] {f : σ₁ → L} {g : σ₂ → L}
    {V : σ₁ × σ₂ → ℝ} {b : ℝ}
    (hK : QRelChain.IsMarkov K) (hμ₀ : QRelChain.IsDist μ₀)
    (hK₁ : QRelChain.IsMarkov K₁) (hK₂ : QRelChain.IsMarkov K₂)
    (hm₁ : QRelChain.PushKernel Prod.fst K K₁)
    (hm₂ : QRelChain.PushKernel Prod.snd K K₂)
    (hfreeze : ∀ p ∈ A₁ ×ˢ A₂, K p p = 1)
    (hfr₁ : ∀ a ∈ A₁, K₁ a a = 1) (hfr₂ : ∀ a ∈ A₂, K₂ a a = 1)
    (habs₁ : Filter.Tendsto
      (fun n => ∑ a ∈ A₁, QRelChain.law K₁ (QRelChain.push Prod.fst μ₀) n a)
      Filter.atTop (nhds 1))
    (habs₂ : Filter.Tendsto
      (fun n => ∑ a ∈ A₂, QRelChain.law K₂ (QRelChain.push Prod.snd μ₀) n a)
      Filter.atTop (nhds 1))
    (hV0 : ∀ p, 0 ≤ V p) (hdrift : ∀ p, ∑ r, K p r * V r ≤ V p)
    (hterm : ∀ p ∈ A₁ ×ˢ A₂, f p.1 ≠ g p.2 → 1 ≤ V p)
    (hb : ∑ p, μ₀ p * V p ≤ b) :
    tvDist
      (ofFn (QRelChain.push f (QRelChain.finalLaw K₁ (QRelChain.push Prod.fst μ₀) A₁)))
      (ofFn (QRelChain.push g (QRelChain.finalLaw K₂ (QRelChain.push Prod.snd μ₀) A₂)))
      ≤ b
\end{leansrc}
\declhead{BernsteinLogicAux.QRelChain}{sound\_QRelTV}{\badge{BadgeGreen}{statement bound}}
\begin{leansrc}
/-- **Soundness of `qrel-explicit-v0` `tv_le` certificates** (README
§"Soundness chain", assembled; the `sound_QCoupleDrift` obligation of the
doctrine sketch at the explicit finite tier, TV readout): under the checker's
obligations — C1/C3 (`IsMarkov`), C4 (`PushKernel`, both faces), the
completed-kernel freeze, C5 (`0 ≤ V`), C6 (drift), C7 (`V ≥ 1` on
disagreeing both-absorbed states), C8 (absorption, semantic form, per
component), C9 (`E_{μ₀}[V] ≤ b`) — the certified bound dominates the true
total-variation distance of the two final readout laws. -/
theorem sound_QRelTV [Fintype L] [DecidableEq L] {f : σ₁ → L} {g : σ₂ → L}
    {V : σ₁ × σ₂ → ℝ} {b : ℝ}
    (hK : IsMarkov K) (hμ₀ : IsDist μ₀) (hK₁ : IsMarkov K₁) (hK₂ : IsMarkov K₂)
    (hm₁ : PushKernel Prod.fst K K₁) (hm₂ : PushKernel Prod.snd K K₂)
    (hfreeze : ∀ p ∈ A₁ ×ˢ A₂, K p p = 1)
    (hfr₁ : ∀ a ∈ A₁, K₁ a a = 1) (hfr₂ : ∀ a ∈ A₂, K₂ a a = 1)
    (habs₁ : Tendsto (fun n => ∑ a ∈ A₁, law K₁ (push Prod.fst μ₀) n a)
      atTop (𝓝 1))
    (habs₂ : Tendsto (fun n => ∑ a ∈ A₂, law K₂ (push Prod.snd μ₀) n a)
      atTop (𝓝 1))
    (hV0 : ∀ p, 0 ≤ V p) (hdrift : ∀ p, ∑ q, K p q * V q ≤ V p)
    (hterm : ∀ p ∈ A₁ ×ˢ A₂, f p.1 ≠ g p.2 → 1 ≤ V p)
    (hb : ∑ p, μ₀ p * V p ≤ b) :
    tvDist (push f (finalLaw K₁ (push Prod.fst μ₀) A₁))
      (push g (finalLaw K₂ (push Prod.snd μ₀) A₂)) ≤ b
\end{leansrc}
\declhead{BernsteinLogicAux.QRelChain}{sound\_QRelExpGap}{\badge{BadgeGreen}{statement bound}}
\begin{leansrc}
/-- **Soundness of `qrel-explicit-v0` `expgap_le` certificates** (README
§"Soundness chain", assembled; Jensen replaces the coupling lemma): under the
same obligations with C7 in its expectation-gap form
(`|v₁ − v₂| ≤ V` at both-absorbed states), the certified bound dominates the
true gap of final expected readout values. -/
theorem sound_QRelExpGap {v₁ : σ₁ → ℝ} {v₂ : σ₂ → ℝ} {V : σ₁ × σ₂ → ℝ} {b : ℝ}
    (hK : IsMarkov K) (hμ₀ : IsDist μ₀) (hK₁ : IsMarkov K₁) (hK₂ : IsMarkov K₂)
    (hm₁ : PushKernel Prod.fst K K₁) (hm₂ : PushKernel Prod.snd K K₂)
    (hfreeze : ∀ p ∈ A₁ ×ˢ A₂, K p p = 1)
    (hfr₁ : ∀ a ∈ A₁, K₁ a a = 1) (hfr₂ : ∀ a ∈ A₂, K₂ a a = 1)
    (habs₁ : Tendsto (fun n => ∑ a ∈ A₁, law K₁ (push Prod.fst μ₀) n a)
      atTop (𝓝 1))
    (habs₂ : Tendsto (fun n => ∑ a ∈ A₂, law K₂ (push Prod.snd μ₀) n a)
      atTop (𝓝 1))
    (hV0 : ∀ p, 0 ≤ V p) (hdrift : ∀ p, ∑ q, K p q * V q ≤ V p)
    (hterm : ∀ p ∈ A₁ ×ˢ A₂, |v₁ p.1 - v₂ p.2| ≤ V p)
    (hb : ∑ p, μ₀ p * V p ≤ b) :
    |(∑ a₁ ∈ A₁, finalMass K₁ (push Prod.fst μ₀) a₁ * v₁ a₁) -
      ∑ a₂ ∈ A₂, finalMass K₂ (push Prod.snd μ₀) a₂ * v₂ a₂| ≤ b
\end{leansrc}
\ruleblock{Correspondence (Steele display $\leftrightarrow$ Lean)}
\badge{BadgeGreen}{reviewer-certified against the displayed texts}
\sourcenote{The dictionary below is a reviewer-certified statement of how the mechanized statements above interpret the displayed rule: display metavariable $\mapsto$ Lean binder, each display hypothesis $\mapsto$ the premise or side condition carrying it, the display conclusion $\mapsto$ the goal --- with any deliberate difference in width between display and mechanization named outright. Certification is recorded as a hash triple over (the displayed cell, the extracted statement texts, this dictionary\textquotesingle s own text); if any of the three changes, this block turns amber until a reviewer re-certifies. The correspondence claim itself remains editorial --- reviewer-grade, not machine-checked; its machine-checked successor types each soundness statement against the rule as a structured object, and supersedes this block per rule as it lands.}
The displayed $d \in \{\mathrm{TV}, \text{gap-pseudometrics}\}$ is two
theorems: $\mathsf{QDist}[\mathrm{TV}](\mathcal{M}_1, \mathcal{M}_2)
\le b$ $\mapsto$ \texttt{sound\_QRelTV}\textquotesingle s goal
\texttt{tvDist (push f (finalLaw K₁ …)) (push g (finalLaw K₂ …)) ≤ b}
--- the true total-variation distance of the two final \emph{readout}
laws (the readout maps \texttt{f}/\texttt{g} are the displayed
pseudometric\textquotesingle s observation structure, explicit in the
binder) --- and the gap form $\mapsto$
\texttt{sound\_QRelExpGap}\textquotesingle s absolute-difference goal.
The obligation roster C1--C9 maps row family by row family:
\begin{itemize}
\item coupling rows $\mapsto$ \texttt{hK : IsMarkov K} (entrywise
  nonneg, unit sums) + \texttt{hm₁}/\texttt{hm₂} (\texttt{PushKernel
  Prod.fst/snd K Kᵢ} --- marginals equal the completed component
  kernels, both faces).
\item freeze rows $\mapsto$ \texttt{hfreeze} (both-absorbed diagonal
  $= 1$) + \texttt{hfr₁}/\texttt{hfr₂} (component freezes).
\item generator row $\mapsto$ \texttt{hdrift : ∀ p, ∑ q, K p q * V q ≤ V
  p}, every product state, displayed vocabulary carried exactly;
  $V \ge 0$ $\mapsto$ \texttt{hV0}.
\item terminal domination $\mapsto$ \texttt{hterm}, per distance flavour, exactly as displayed: the
  indicator form (\texttt{f p.1 ≠ g p.2 → 1 ≤ V p}) and the gap form
  (\texttt{|v₁ p.1 − v₂ p.2| ≤ V p}), both on both-absorbed states
  only.
\item the absorption criterion carries this
  entry\textquotesingle s one representation delta: the display asks the
  per-component reach criterion (recomputed, per the flavor text); the BOUND
  statements consume absorption in \emph{semantic} form ---
  \texttt{habs₁}/\texttt{habs₂}, absorbed mass tending to one. The
  \texttt{\_of\_reach} variants consuming the checker\textquotesingle s
  own graph criterion (via \texttt{tendsto\_absorbedMass\_one})
  exist in the module and are cited by the roster witness, but are
  not in this bound set --- machine-checked and cited, not bound, in
  this edition.
\item initial marginal rows $\mapsto$ carried \emph{by construction}:
  the component initial laws in the goal ARE \texttt{push Prod.fst/snd
  μ₀} --- no separate hypothesis exists because the statement never
  admits mismatched initials; $\vdash \sum_p \mu_0(p) V(p) \le b$
  $\mapsto$ \texttt{hb}.
\end{itemize}
Width note: the \emph{conclusion} is the general object ---
\texttt{TV.sound\_QRelTV\_measure} (module
\texttt{BernsteinLogicAux.TvDist}) restates
\texttt{sound\_QRelTV}\textquotesingle s bound against the
Jordan-decomposition total-variation distance of the final readout
laws viewed as measures, with zero new hypotheses, through the general
coupling inequality (\texttt{tvDist\_le\_coupling}: any coupling,
off-diagonal weighed by outer measure, no diagonal-measurability
assumption) and the finite-tier seam \texttt{tvDist\_ofFn}
(\texttt{QRelChain.tvDist} \emph{is} the general distance restricted);
TV is packaged as a metric on probability measures (\texttt{TVProb}).
The conditional mark therefore re-scopes to the certificate-arithmetic
tier: \texttt{Fintype} chains with exact real arithmetic standing in
for exact $\Q$ --- a fragment fence, by design. The countable tier now carries kernel-level twins
(\texttt{BernsteinLogicAux.CountableTier}: the drift-induction reach
bound, the optional-stopping reach lower bound, and geometric
absorption, over \texttt{Countable} state spaces with row hypotheses as
\texttt{tsum} series, summability supplied by the unit-row-sum and
boundedness rows), machine-checked and cited, not bound, like the
\texttt{\_of\_reach} variants --- the $\mathbb{W}_1$
registry entry and the mixture-polynomial tier remain designed, not
realized, exactly as the flavor text says;
coupling emptiness has no Lean face because the theorems consume a
supplied coupling --- refusal-shaped completeness cost lives
producer-side.

\ruleblock{Lean verifier receipts}
\sourcenote{module \texttt{BernsteinLogicAux.TvDist}; bound build \texttt{54f7e48568aa\ldots} (bound from receipt 1c592f5cc199fad5, 2026-07-24T09:28:46Z). A receipt records that named checking layers accepted a build: (module, compiled-artifact SHA-256, environment SHA-256, checker layers), minted when the checking gate ran; the store receipts attesting the bound build are listed with the layers each run credited. Nothing here asks the reader to trust the store: every layer is an independently runnable program over the quoted sources --- rebuild the module and re-run the checkers to reproduce the verdict. Receipts attesting superseded builds credit no layer and are omitted from this rendering (the count of omissions, if any, is noted below).}
\begin{receipttable}
\receiptrow{a908c952d3907ab4}{2026-07-22T16:31:55Z}{build+decorrelated}{\badge{BadgeGreen}{current}}
\receiptrow{4f8b7d7adcc8f5d2}{2026-07-22T18:27:35Z}{build+decorrelated}{\badge{BadgeGreen}{current}}
\receiptrow{1c592f5cc199fad5}{2026-07-24T09:28:46Z}{build+decorrelated}{\badge{BadgeGreen}{current}}
\receiptrow{2cc8a0435e486a23}{2026-07-31T22:00:34Z}{build+decorrelated}{\badge{BadgeGreen}{current}}
\receiptrow{b78489856a2622df}{2026-07-31T22:29:28Z}{build+decorrelated}{\badge{BadgeGreen}{current}}
\receiptrow{10be3d04a0098b77}{2026-08-02T23:31:44Z}{build+decorrelated}{\badge{BadgeGreen}{current}}
\receiptrow{9bcf8f97503e6168}{2026-08-03T16:27:47Z}{build+decorrelated}{\badge{BadgeGreen}{current}}
\receiptrow{c41aeb2d063d9b00}{2026-08-03T22:46:32Z}{build+decorrelated}{\badge{BadgeGreen}{current}}
\end{receipttable}
\sourcenote{module \texttt{BernsteinLogicAux.QRelChain}; bound build \texttt{64ae4b299db4\ldots} (bound from receipt efcd1abd35a1b7f0, 2026-07-19T18:24:37Z). A receipt records that named checking layers accepted a build: (module, compiled-artifact SHA-256, environment SHA-256, checker layers), minted when the checking gate ran; the store receipts attesting the bound build are listed with the layers each run credited. Nothing here asks the reader to trust the store: every layer is an independently runnable program over the quoted sources --- rebuild the module and re-run the checkers to reproduce the verdict. Receipts attesting superseded builds credit no layer and are omitted from this rendering (the count of omissions, if any, is noted below).}
\begin{receipttable}
\receiptrow{54e0e5134c7a203f}{2026-07-19T01:11:59Z}{build+decorrelated}{\badge{BadgeGreen}{current}}
\receiptrow{1ac5f362df2230fb}{2026-07-19T01:58:14Z}{build+decorrelated}{\badge{BadgeGreen}{current}}
\receiptrow{efcd1abd35a1b7f0}{2026-07-19T18:24:37Z}{build+decorrelated}{\badge{BadgeGreen}{current}}
\receiptrow{5a18ef8a1d5d481d}{2026-07-20T01:21:49Z}{build+decorrelated}{\badge{BadgeGreen}{current}}
\receiptrow{1119d63191cbca38}{2026-07-20T20:46:42Z}{build+decorrelated}{\badge{BadgeGreen}{current}}
\receiptrow{53ac120ea74296cc}{2026-07-20T20:58:15Z}{build+replay+hardened+decorrelated}{\badge{BadgeGreen}{current}}
\receiptrow{1189ecab6f8e092f}{2026-07-21T00:03:39Z}{build+decorrelated}{\badge{BadgeGreen}{current}}
\receiptrow{4a715ffa478d6acd}{2026-07-21T05:38:59Z}{build+replay+hardened+decorrelated}{\badge{BadgeGreen}{current}}
\receiptrow{ab2a676a4de29f86}{2026-07-22T18:27:12Z}{build+decorrelated}{\badge{BadgeGreen}{current}}
\receiptrow{a1b101422ff8770b}{2026-07-31T22:00:09Z}{build+decorrelated}{\badge{BadgeGreen}{current}}
\receiptrow{6aed590f9ba0b241}{2026-07-31T22:29:03Z}{build+decorrelated}{\badge{BadgeGreen}{current}}
\receiptrow{b83a01eca4bc0784}{2026-08-02T23:31:14Z}{build+decorrelated}{\badge{BadgeGreen}{current}}
\receiptrow{2bbbc458342cfb9e}{2026-08-03T16:27:17Z}{build+decorrelated}{\badge{BadgeGreen}{current}}
\receiptrow{04d2eae82194b82e}{2026-08-03T22:46:01Z}{build+decorrelated}{\badge{BadgeGreen}{current}}
\end{receipttable}
\layerline{\layeritem{build}{lake build (elaboration)}{\badge{BadgeGreen}{current}}\layeritem{replay}{leanchecker (official kernel replay)}{\badge{BadgeGreen}{current}}\layeritem{hardened}{safe\_verify (hardened screen)}{\badge{BadgeGreen}{current}}\layeritem{decorrelated}{lean4export + nanoda (decorrelated kernel)}{\badge{BadgeGreen}{current}}}
\sourcenote{\badge{BadgeGray}{comparator: none} no fixture in the three-implementation comparator harness (compiled / kernel-interpreted / reference-interpreted) exercises this rule.}
\ruleblock{HOL4 statement}
\absentbox{NO HOL STATEMENT YET}{The HOL4 development checks the base calculus\textquotesingle s wire
format (theory \texttt{bern\_catalogue}, the first-order check-function
cone of the core catalogue). \rulename{QCoupleDrift} concludes a
distance judgment of the extension theory: its rule sort is out of
scope of the current baked theory, and the soundness spine lives in the
auxiliary Lean package. Nothing exists yet for a HOL statement to be
about; the box becomes writable when a revision carries the relational
certificate rows and its checker chain is ported.
}
\ruleblock{HOL verifier receipts}
\absentbox{NONE}{no HOL verification receipt exists for this rule.}
\ruleblock{Proof}

\sourcenote{Genre: core rule (tier 1). Scope, stated loudly: the
finite-chain tier --- finite state spaces, distribution-evolution
semantics, exact-rational certificate data instantiated over $\R$ by
cast. The mixture-polynomial tier and the $\mathbb{W}_1$ conclusion
family are designed, not proven, and nothing below claims them. Every
hypothesis of the rule is used and remarked on.}

\textbf{Setting and claim.}\enspace Two finite chains are compared:
state spaces $\sigma_1, \sigma_2$, row-stochastic kernels $K_1, K_2$,
absorbing target sets $A_i$ (each $a \in A_i$ frozen: $K_i(a, a) = 1$),
and readouts $f : \sigma_1 \to L$, $g : \sigma_2 \to L$ (resp.\ values
$v_1, v_2 : \sigma_i \to \R$). The certificate supplies a
\emph{coupling} kernel $K$ on $\sigma_1 \times \sigma_2$ with initial
law $\mu_0$, and a function $V \geq 0$ with bound $b$. The chain's law
evolves by $\mu_{n+1}(t) = \sum_s \mu_n(s) K(s, t)$; the \emph{final
law} of an absorbing state is its monotone limit mass
$\mathrm{fin}(a) = \sup_n \mu_n(a)$. Hypotheses (the checker's
obligation list, in the order the proof consumes them): the coupling
rows make $K$ a kernel whose row marginals are exactly $K_1, K_2$
(\emph{marginal identities}); both-absorbed states
$A_1 \times A_2$ are frozen; $V$ is a supermartingale for $K$
($\sum_q K(p, q) V(q) \leq V(p)$, every $p$); $V \geq 0$; terminal
domination ($V(p) \geq 1$ on both-absorbed $p$ with
$f(p_1) \neq g(p_2)$, resp.\ $V(p) \geq |v_1(p_1) - v_2(p_2)|$); each
component chain absorbs almost surely (certified by a finite-graph
reach criterion); and $\sum_p \mu_0(p) V(p) \leq b$. The claims
(\texttt{sound\_QRelTV}, \texttt{sound\_QRelExpGap}): the
total-variation distance of the two final readout laws, resp.\ the gap
of final expected values, is at most $b$.

\begin{blemma}[marginal-process lemma]\label{lem:qcd-marginal}
If every row of $K$ pushes along $\pi_i$ to the corresponding row of
$K_i$, then at every time $n$ the pushforward of the coupled law is the
component law: $\pi_{i\#}\mu_n = \mathrm{law}_i(n)$, where
$\mathrm{law}_i$ evolves $\pi_{i\#}\mu_0$ under $K_i$.
\end{blemma}
\begin{proof}
Induction on $n$; the base is definitional. For the step it suffices
that the row identities intertwine the two Markov operators:
\[
(\pi_{\#}\,\mathrm{step}_K\,\mu)(c)
= \sum_{\pi(t) = c} \sum_{s} \mu(s)\, K(s, t)
= \sum_{s} \mu(s) \sum_{\pi(t) = c} K(s, t)
= \sum_{s} \mu(s)\, K_i(\pi(s), c)
= (\mathrm{step}_{K_i}\, \pi_{\#}\mu)(c),
\]
where the third equality is the pushed-row identity at $s$ and the last
regroups the outer sum over the fibres of $\pi$. This is the coupling's
first duty: the product chain really \emph{is} a joint presentation of
the two systems being compared.
\end{proof}

\begin{blemma}[supermartingale induction]\label{lem:qcd-drift}
$\sum_p \mu_n(p)\, V(p) \leq \sum_p \mu_0(p)\, V(p)$ for all $n$.
\end{blemma}
\begin{proof}
One step: $\sum_t (\mathrm{step}\,\mu)(t) V(t) = \sum_s \mu(s)
\sum_t K(s, t) V(t) \leq \sum_s \mu(s) V(s)$, using the generator row at
each $s$ and $\mu(s) \geq 0$. Iterate.
\end{proof}

\begin{blemma}[the finite optional-stopping bound]\label{lem:qcd-stop}
Let $A \subseteq$ states be frozen, and let $0 \leq \gamma \leq V$ on
$A$. Then $\sum_{a \in A} \mathrm{fin}(a)\, \gamma(a) \leq b$.
\end{blemma}
\begin{proof}
At every finite time,
\[
\sum_{a \in A} \mu_n(a)\, \gamma(a)
\;\leq\; \sum_{a \in A} \mu_n(a)\, V(a)
\;\leq\; \sum_{\text{all } s} \mu_n(s)\, V(s)
\;\leq\; b,
\]
the middle step because $V \geq 0$ lets the sum extend off $A$, the
last by Lemma~\ref{lem:qcd-drift} and the initial bound. Frozen states
have monotone mass ($\mu_{n+1}(a) \geq \mu_n(a) K(a,a) = \mu_n(a)$), so
each term converges upward to $\mathrm{fin}(a)\gamma(a)$; the limit of
the finite sum obeys the uniform bound. This is optional stopping with
monotone convergence of finite sums doing Fatou's work --- no measure
theory enters at this tier.
\end{proof}

\begin{blemma}[coupling inequality]\label{lem:qcd-couple}
For any nonnegative weight $\gamma$ on $\sigma_1 \times \sigma_2$,
\[
\mathrm{TV}\bigl((f \circ \pi_1)_{\#}\gamma,\ (g \circ \pi_2)_{\#}\gamma\bigr)
\;\leq\; \gamma\{\, p : f(p_1) \neq g(p_2) \,\}.
\]
\end{blemma}
\begin{proof}
Fix a label $l$. Splitting each pushed mass over the agreement region
$\{f = g\}$ and its complement, the agreement contributions to the two
sides are \emph{equal} (there, $f(p_1) = l \iff g(p_2) = l$), so
\[
\bigl| (f\pi_1)_{\#}\gamma(l) - (g\pi_2)_{\#}\gamma(l) \bigr|
\;\leq\; \gamma\{f \neq g,\ f(p_1) = l\} + \gamma\{f \neq g,\ g(p_2) = l\}.
\]
Summing over $l$ counts the disagreement mass twice, and
$\mathrm{TV} = \tfrac12 \sum_l |\cdot|$ halves it. This is the
classical coupling lemma \cite[\S I.2]{lindvall1992}, in Doeblin's
tradition \cite{doeblin1938}; stated for arbitrary nonnegative
$\gamma$ because normalization is never consumed --- Step 2 of the
theorem applies it to the final pair law, whose unit mass
Lemma~\ref{lem:qcd-final} supplies separately.
\end{proof}

\begin{blemma}[final marginal identities]\label{lem:qcd-final}
Suppose additionally each component chain absorbs almost surely
(its absorbed mass tends to $1$). Then for $a_1 \in A_1$:
$\sum_{a_2 \in A_2} \mathrm{fin}_K(a_1, a_2) =
\mathrm{fin}_{K_1}(a_1)$, and symmetrically. Consequently the final
pair law on $A_1 \times A_2$ is a genuine \emph{coupling of the two
final component laws}, of total mass $1$.
\end{blemma}
\begin{proof}
Let $S_n = \sum_{a_2 \in A_2} \mu_n(a_1, a_2)$ and
$T_n = \sum_{s_2} \mu_n(a_1, s_2)$; by Lemma~\ref{lem:qcd-marginal},
$T_n$ is the component-1 law at $a_1$. Then $S_n \leq T_n$, and
\[
T_n - S_n \;=\; \textstyle\sum_{s_2 \notin A_2} \mu_n(a_1, s_2)
\;\leq\; \textstyle\sum_{s_2 \notin A_2} (\pi_{2\#}\mu_n)(s_2)
\;=\; 1 - (\text{component-2 absorbed mass at } n) \;\longrightarrow\; 0 :
\]
the \emph{one-absorbed-loner} mass --- pairs whose first coordinate has
frozen while the second still wanders --- is dominated by the second
chain's unabsorbed mass, which its absorption hypothesis kills.
Squeeze; both-absorbed pairs are frozen, so $S_n$ also converges to the
final pair mass. Total mass $1$ follows by summing over $A_1$ and
applying component-1 absorption.
\end{proof}

\begin{bthm}[soundness, total-variation form]\label{thm:qcd-tv}
Under all hypotheses,
$\mathrm{TV}\bigl(f_{\#}\mathrm{fin}_{K_1},\ g_{\#}\mathrm{fin}_{K_2}\bigr)
\leq b$.
\end{bthm}
\begin{proof}
\emph{Step 1.} By Lemma~\ref{lem:qcd-final} the final pair law $\gamma$
on $A_1 \times A_2$ has marginals $\mathrm{fin}_{K_1}$,
$\mathrm{fin}_{K_2}$, so the two final readout laws are the two
pushforwards of the single weight $\gamma$.
\emph{Step 2.} Lemma~\ref{lem:qcd-couple} bounds their total-variation
distance by $\gamma$'s disagreement mass.
\emph{Step 3.} Terminal domination says the disagreement indicator is
$\leq V$ on the frozen set $A_1 \times A_2$ (and $\geq 0$ everywhere),
so Lemma~\ref{lem:qcd-stop} with $\gamma = \mathbf{1}[f \neq g]$ bounds
the disagreement mass by $b$.
\end{proof}

\begin{bthm}[soundness, expectation-gap form]\label{thm:qcd-gap}
With terminal domination in the form
$V(p) \geq |v_1(p_1) - v_2(p_2)|$ on $A_1 \times A_2$,
\[
\Bigl| \textstyle\sum_{a_1} \mathrm{fin}_{K_1}(a_1)\, v_1(a_1)
- \sum_{a_2} \mathrm{fin}_{K_2}(a_2)\, v_2(a_2) \Bigr| \;\leq\; b .
\]
\end{bthm}
\begin{proof}
Rewrite both final expectations over the final pair law via
Lemma~\ref{lem:qcd-final}; the difference becomes
$\sum_p \mathrm{fin}_K(p)\,(v_1(p_1) - v_2(p_2))$ over
$A_1 \times A_2$. The triangle inequality moves the absolute value
inside, and Lemma~\ref{lem:qcd-stop} with
$\gamma = |v_1 - v_2|$ closes. The coupling lemma is not consumed here
--- domination of the value gap replaces domination of the disagreement
indicator; two conclusion families, one engine.
\end{proof}

\begin{bremark}[the absorption criterion is discharged, not assumed]
\label{rem:qcd-reach}
The checker certifies absorption per component as a finite-graph fact:
from every state, the frozen set is reachable through
positive-probability edges. Finiteness uniformizes this into one
horizon $N$ and one margin $\varepsilon > 0$ valid for all states; each
$N$-step block then absorbs at least an $\varepsilon$-fraction of the
still-loose mass (frozen mass never returns), so the unabsorbed mass
decays geometrically and the absorbed mass tends to $1$
(\texttt{tendsto\_absorbedMass\_one}). The \texttt{\_of\_reach}
theorem forms consume the graph criterion directly.
\end{bremark}

\begin{bremark}[why each hypothesis is needed]\label{rem:qcd-omit}
\emph{Marginal identities}: drop them and $\gamma$ in
Theorem~\ref{thm:qcd-tv} no longer has the two systems' final laws as
marginals --- the bound then constrains an unrelated pair (couple
$\mathcal{M}_1$ with a copy of \emph{itself} instead of
$\mathcal{M}_2$: every absorbed pair agrees, so $V = 0$ on the
reachable diagonal --- with $V = 1$ on the unreachable disagreeing
absorbed pairs --- passes every other row at $b = 0$, whatever
$\mathcal{M}_2$ is). \emph{$V \geq 0$}: the extension step in
Lemma~\ref{lem:qcd-stop} ($\sum_A \leq \sum_{\text{all}}$) fails, and
negative values can subsidize the generator row: from an initial state of
value $0$, branch half-half into a frozen \emph{disagreeing} state of
value $1$ and a frozen \emph{agreeing} state of value $-1$ (domination
constrains only disagreeing states). The generator row holds at every state
($\tfrac12 \cdot 1 + \tfrac12 \cdot (-1) = 0 \leq 0$), the initial
bound passes at $b = 0$ --- and the disagreement mass is $\tfrac12$. \emph{Terminal domination}: with
$V \equiv 0$ and disagreeing finals, generator and initial rows pass at
$b = 0$ while the true distance is $1$. \emph{The freeze}: monotone
convergence of frozen mass (Lemmas \ref{lem:qcd-stop},
\ref{lem:qcd-final}) is what makes ``final law'' well-defined;
unfrozen, the two-state swap chain puts mass $1$ on the target at even
times and $0$ at odd --- $\sup_n \mu_n(a)$ is no longer a limit, and
certified mass drains back out.
\emph{Absorption}: without it the one-absorbed-loner mass survives ---
take $K_2$ the identity off $A_2$: the pair never both-absorbs, the
final pair law is $0$, and Lemma~\ref{lem:qcd-couple} holds vacuously
about sub-laws that no longer represent the systems; the marginal
identity of Lemma~\ref{lem:qcd-final} is exactly what fails --- and
the bound itself fails with it: component 1 absorbing at once,
component 2 frozen nowhere it reaches, $V = 0$ on the reachable ray
and $1$ on the (unreachable) disagreeing both-absorbed pairs passes
every other row at $b = 0$, while the final readout laws sit at
total-variation distance $\tfrac12$.
\emph{Initial bound}: the conclusion's $b$ is
$V$'s initial expectation; without it nothing quantitative survives.
A coupling that cannot satisfy the marginal rows simply does not
exist as a certificate --- the coupling is an input premise, so its
emptiness refuses (a completeness cost), never lies.
\end{bremark}

\begin{bremark}[relation to the literature]
The engine is the coupling method for Markov chains
\cite{doeblin1938,lindvall1992}: total variation bounded by
disagreement probability of a coupling, here in distribution-evolution
form over finite spaces \cite[Ch.~4--5]{levinperes2017}, with a
relational supermartingale supplying the quantitative disagreement
bound --- the certificate-side sibling of coupling-based relational
program logics \cite{hsu2017}, transposed from pre/post judgments to a
single distance judgment that composes by metric axioms (triangle,
weakening) at judgment level. The expectation-gap form replaces the
coupling lemma by direct domination --- the gap-pseudometric readout.
On correlated perturbations the relational route is strictly sharper
than composing two one-sided certificates, which is the rule's reason
to exist: the coupling sees the correlation, unary certificates
cannot.
\end{bremark}

\begin{anchors}
Lemma \ref{lem:qcd-marginal} & \texttt{PushKernel}; \texttt{push\_step}; \texttt{push\_law}\\
Lemma \ref{lem:qcd-drift} & \texttt{expect\_step\_le}; \texttt{expect\_law\_le}\\
Lemma \ref{lem:qcd-stop} & \texttt{law\_mono\_absorbing}; \texttt{finalMass}; \texttt{finalExpect\_le}\\
Lemma \ref{lem:qcd-couple} & \texttt{tvDist}; \texttt{tvDist\_push\_le\_disagree}\\
Lemma \ref{lem:qcd-final} & \texttt{sum\_finalMass\_snd\_eq}; \texttt{sum\_finalMass\_fst\_eq}; \texttt{sum\_finalMass\_prod\_eq\_one}; \texttt{push\_finalLaw\_fst}/\texttt{\_snd}\\
Theorem \ref{thm:qcd-tv} & \texttt{sound\_QRelTV} (graph-criterion form \texttt{sound\_QRelTV\_of\_reach})\\
Theorem \ref{thm:qcd-gap} & \texttt{sound\_QRelExpGap} (graph-criterion form \texttt{sound\_QRelExpGap\_of\_reach})\\
Remark \ref{rem:qcd-reach} & \texttt{exists\_uniform\_reach}; \texttt{absorbed\_decay}; \texttt{tendsto\_absorbedMass\_one}\\
\end{anchors}


\rulesection{92}{QFixBall}{core rule}
\label{blcatrule-QFixBall}
\ruleblock{The rule (Steele metanotation)}
\sourcenote{as realized in the mechanized development; source row: \texttt{\texttt{QTSL-rules.toml}} --- \emph{QFixBall}. The metanotation statement is typeset by hand for notational quality; its mechanized ground truth is the Lean statement below, extracted from the sources at every build of this book and hash-bound. The rule \emph{roster} is machine-checked (every mechanized rule tag round-trips to exactly one catalogue row); the statement \emph{text} is editorial, checked by eye against the extracted Lean beside it.}
\begin{mathpar}
\inferrule*{
  d : \mathsf{QDistRegistry}(X_d); \qquad (X_d, d) \ \text{complete} \\
  \vdash \rho < 1 \ \text{two-point on } \overline{B}(\hat\mu, r) \ \text{(contraction row)} \\
  \vdash d(T\hat\mu, \hat\mu) \le \delta \ \text{(residual row)} \\
  \vdash \delta + \rho\, r \le r \ \text{(radius row, division-free)}
}{
  \vdash \mathsf{QDist}[d](\hat\mu,\ \mathrm{fix}\,T) \le r
}
\end{mathpar}

\ruleflavor{every hypothesis anchored at the known snapshot $\hat\mu$, never at the unknown fixed point: the contraction row discharges by operator-norm/psdCert rows at the affine tier, the residual row is one concrete sweep at the known point, and completeness is mechanized at the affine/Euclidean instantiation ($\mathbb{W}_2$ completeness by instance on $(\mathcal{P}_2(X), W_2)$). The conclusion\textquotesingle s definite description needs no separate license premise: this rule is the only derivation that makes $\mathrm{fix}\,T$ speakable, and it discharges the description\textquotesingle s well-formedness itself --- exactly one fixed point exists in $\overline{B}(\hat\mu, r)$, manufactured by the contraction (the $\exists!$ of the soundness statement; no existence stratum consumed). The registry typing has no baked sort yet (the promoted $\mathsf{QDist}$ kind awaits its row species); completeness is the instance hypothesis of the Lean statement. Refusals: $\rho \ge 1$, and registry entries the subject sort does not carry.}
\par\medskip
\rulename{QFixBall} is the contraction-witnessed discharge of the
distance judgment: it certifies that a computed snapshot $\hat\mu$ lies
within distance $r$ of the fixed point of an update map $T$ --- and, in
the same breath, that \textquotedblleft the fixed point of
$T$\textquotedblright{} is a well-defined description at all (exactly
one fixed point exists in the certified ball). A producer reaches for
it wherever an iterative solver has been run to approximate a
fixed-point object --- a stationary law, a variational posterior, a
self-consistent field --- and the deliverable is not the iterate but a
certified \emph{ball} around it. Its three rows realize the anchoring
inversion: every hypothesis is stated at the \emph{known} snapshot,
never at the unknown fixed point --- the contraction row bounds
$T$\textquotesingle s two-point ratio on the closed ball around
$\hat\mu$, the residual row measures one concrete application
$d(T\hat\mu, \hat\mu) \le \delta$, and the radius row
$\delta + \rho\, r \le r$ is division-free arithmetic.

What the checker recomputes: the contraction bound $\rho < 1$ (at the
affine tier, an operator-norm row discharged by positive-semidefinite
certificate machinery on $\rho^2 I - A^{\top}\!A$-shaped data), the residual
evaluation (one sweep of the map at the snapshot, exact rational
arithmetic), and the radius inequality. Everything else --- ball
self-mapping, existence, uniqueness, the anchored bound
$d(\hat\mu, \mathrm{fix}\,T) \le \delta/(1-\rho) \le r$ --- is
manufactured by the mathematics from those three rows; no existence
stratum is consumed, so the judgment is existential without vacuity.

What would break: with $\rho \ge 1$ the fixed point need not be unique
or exist and the description license fails --- the rule refuses (the
identity map has residual $0$ at every point and fixes all of them).
Without the radius row the ball need not be self-mapped and can miss
the fixed point entirely, even with a genuine contraction. Without the
residual row nothing ties the ball to the dynamics. The completeness of
the subject space is a genuine hypothesis, carried honestly --- and both
named instantiations now discharge it by instance: the affine/Euclidean
case directly, and the Wasserstein case on $(\mathcal{P}_2(X), W_2)$
over a complete separable base (the instantiation block below) --- never
an axiom smuggled in.

\cataloguenote{The rule concludes onto a definite-description subject
($\mathrm{fix}\,T$) and therefore carries the description
well-formedness discharge inside the rule --- the one new typing
obligation of the distance-judgment grammar. An independent
exact-arithmetic checker replays a pilot certificate\textquotesingle s
eleven obligations; its mutation battery refuses each mutilated
certificate at a named obligation. Producer strategy (proposing
$\rho$ from iterate data, pinning the affine reduction) is deliberately
outside the kernel.
}
\ruleblock{Lean statement}
\declhead{BernsteinLogicAux.QFixBall}{sound\_QFixBall}{\badge{BadgeGreen}{statement bound}}
\begin{leansrc}
/-- **`sound_QFixBall`** — the rule's full conclusion: the certified
rows produce EXACTLY ONE fixed point of `f` in `closedBall x₀ r`.
Existence is manufactured by the contraction itself (the third case of
the two-direction carve: no existence stratum is consumed, and the
judgment is existential without vacuity). `[CompleteSpace α]` is the
honest boundary — see the header. -/
theorem sound_QFixBall [CompleteSpace α] (hK : K < 1)
    (hlip : LipschitzOnWith K f (closedBall x₀ r))
    (hres : dist (f x₀) x₀ ≤ δ) (hrad : δ + K * r ≤ r) :
    ∃! y, y ∈ closedBall x₀ r ∧ IsFixedPt f y
\end{leansrc}
\declhead{BernsteinLogicAux.QFixBall}{sound\_QFixBall\_affine}{\badge{BadgeGreen}{statement bound}}
\begin{leansrc}
/-- **`sound_QFixBall_affine`** — the caviwball-v0 instantiation shape:
snapshot mean vector in `EuclideanSpace ℝ (Fin M)` (completeness is an
instance — the `[CompleteSpace]` boundary discharges for free here),
affine sweep map, operator-norm + residual + radius rows ⇒ unique
fixed point in the certified ball. -/
theorem sound_QFixBall_affine {M : ℕ}
    (A : EuclideanSpace ℝ (Fin M) →L[ℝ] EuclideanSpace ℝ (Fin M))
    (c x₀ : EuclideanSpace ℝ (Fin M)) {K : ℝ≥0} {δ r : ℝ}
    (hK : K < 1) (hA : ‖A‖₊ ≤ K)
    (hres : dist (A x₀ + c) x₀ ≤ δ) (hrad : δ + K * r ≤ r) :
    ∃! y, y ∈ closedBall x₀ r ∧ IsFixedPt (fun x => A x + c) y
\end{leansrc}
\declhead{BernsteinLogicAux.QFixBall}{dist\_fixedPoint\_le}{\badge{BadgeGreen}{statement bound}}
\begin{leansrc}
/-- **The anchored a-priori bound** — for ANY fixed point in the ball
(uniqueness not consumed): `dist x₀ μ⋆ ≤ δ / (1 - K)`. Proved by one
triangle inequality through the snapshot's image; the literature's
one-point-at-the-limit anchoring never appears. -/
theorem dist_fixedPoint_le (hK : K < 1)
    (hlip : LipschitzOnWith K f (closedBall x₀ r))
    (hres : dist (f x₀) x₀ ≤ δ) (hrad : δ + K * r ≤ r)
    {y : α} (hy : y ∈ closedBall x₀ r) (hfy : IsFixedPt f y) :
    dist x₀ y ≤ δ / (1 - K)
\end{leansrc}
\declhead{BernsteinLogicAux.QFixBall}{anchored\_le\_radius}{\badge{BadgeGreen}{statement bound}}
\begin{leansrc}
/-- The radius row also caps the anchored bound by `r` itself
(`δ/(1-K) ≤ r`) — the form the checker's ball-radius obligation uses. -/
theorem anchored_le_radius (hK : K < 1)
    (hrad : δ + K * r ≤ r) : δ / (1 - K) ≤ r
\end{leansrc}
\par\medskip
\textbf{The Wasserstein instantiation, discharged by instance.}\enspace
The completeness hypothesis named above is closed on the space where
this rule\textquotesingle s stochastic clients live: the statement
below runs \texttt{sound\_QFixBall} on $(\mathcal{P}_2(X), W_2)$ ---
second-moment laws under quadratic Wasserstein distance over a
complete separable metric space --- with the \texttt{CompleteSpace}
side condition supplied by instance resolution rather than hypothesis.
No completeness assumption survives in the statement; the affine and
Wasserstein instantiations now stand on the same footing.

\declhead{BernsteinLogicAux.Wasserstein.Complete}{sound\_QFixBall\_wasserstein}{\badge{BadgeGreen}{statement bound}}
\begin{leansrc}
/-- **`sound_QFixBall`, instantiated at the Wasserstein-2 space.** The
`[CompleteSpace]` hypothesis — the honest boundary of
`BernsteinLogicAux/QFixBall.lean` (header item 1), and the side
condition that kept the rule's coverage marked conditional — is
discharged here **by instance**: for any complete separable metric `X`,
certified contraction/residual/radius rows on `(𝒫₂(X), W₂)` produce
the unique fixed point in the certified ball. -/
theorem sound_QFixBall_wasserstein {T : P2 X → P2 X} {μ₀ : P2 X} {K : ℝ≥0}
    {δ r : ℝ} (hK : K < 1)
    (hlip : LipschitzOnWith K T (closedBall μ₀ r))
    (hres : dist (T μ₀) μ₀ ≤ δ) (hrad : δ + K * r ≤ r) :
    ∃! ν, ν ∈ closedBall μ₀ r ∧ IsFixedPt T ν
\end{leansrc}
\ruleblock{Correspondence (Steele display $\leftrightarrow$ Lean)}
\badge{BadgeGreen}{reviewer-certified against the displayed texts}
\sourcenote{The dictionary below is a reviewer-certified statement of how the mechanized statements above interpret the displayed rule: display metavariable $\mapsto$ Lean binder, each display hypothesis $\mapsto$ the premise or side condition carrying it, the display conclusion $\mapsto$ the goal --- with any deliberate difference in width between display and mechanization named outright. Certification is recorded as a hash triple over (the displayed cell, the extracted statement texts, this dictionary\textquotesingle s own text); if any of the three changes, this block turns amber until a reviewer re-certifies. The correspondence claim itself remains editorial --- reviewer-grade, not machine-checked; its machine-checked successor types each soundness statement against the rule as a structured object, and supersedes this block per rule as it lands.}
The displayed conclusion splits, and so do the bound statements: the
description license --- $\mathrm{fix}\,T$ well-defined (unique) in
$\overline{B}(\hat\mu, r)$ --- $\mapsto$ \texttt{sound\_QFixBall :
∃! y, y ∈ closedBall x₀ r ∧ IsFixedPt f y} (existence manufactured by
the contraction itself; no existence stratum consumed, exactly as the
roster says); the distance clause $\mathsf{QDist}[d](\hat\mu,
\mathrm{fix}\,T) \le r$ is then ball membership, with the sharper
anchored bound $\delta/(1-\rho)$ carried by
\texttt{dist\_fixedPoint\_le} --- whose binder row shows uniqueness is
NOT consumed (any fixed point in the ball obeys it) --- and capped by
the radius row via \texttt{anchored\_le\_radius}.
\begin{itemize}
\item contraction row, two-point on the ball $\mapsto$ \texttt{hlip :
  LipschitzOnWith K f (closedBall x₀ r)} with \texttt{K : ℝ≥0},
  \texttt{hK : K < 1} --- the $\rho \ge 1$ refusal is the strict
  binder; at the affine tier the operator-norm row $\mapsto$
  \texttt{hA : ‖A‖₊ ≤ K} in \texttt{sound\_QFixBall\_affine}
  (EuclideanSpace, the caviwball-v0 pinned mean map).
\item residual row $\mapsto$ \texttt{hres : dist (f x₀) x₀ ≤ δ} ---
  anchored at the snapshot, and no hypothesis anywhere mentions the
  fixed point: the flavor text\textquotesingle s ``never at the unknown
  fixed point'' is a checkable property of the binder list.
\item radius row $\mapsto$ \texttt{hrad : δ + K * r ≤ r},
  division-free verbatim.
\item registry completeness $\mapsto$ the instance binder
  \texttt{[CompleteSpace α]} --- a named hypothesis of the general
  statement, never an axiom, and discharged by instance resolution in
  both named instantiations: directly for the affine/Euclidean case,
  and for $\mathbb{W}_2$ by the Wasserstein block
  (\texttt{sound\_QFixBall\_wasserstein}) below.
\end{itemize}
Width note: the subject space is an arbitrary metric space --- the
registry-entry structure of $d$ is erased, so the certified reach
covers every registry instantiation whose completeness discharge
arrives (both named instantiations\textquotesingle{} discharges are now
in hand), strictly wider than the displayed registry framing, safe
direction; the means-only (variance-conclusion) refusal and the
Gaussian-pinning producer strategy are kernel-external, no Lean
counterpart, correctly.

\ruleblock{Lean verifier receipts}
\sourcenote{module \texttt{BernsteinLogicAux.QFixBall}; bound build \texttt{68796bc2f9db\ldots} (bound from receipt 060c7093ff4899bf, 2026-07-21T04:51:42Z). A receipt records that named checking layers accepted a build: (module, compiled-artifact SHA-256, environment SHA-256, checker layers), minted when the checking gate ran; the store receipts attesting the bound build are listed with the layers each run credited. Nothing here asks the reader to trust the store: every layer is an independently runnable program over the quoted sources --- rebuild the module and re-run the checkers to reproduce the verdict. Receipts attesting superseded builds credit no layer and are omitted from this rendering (the count of omissions, if any, is noted below).}
\begin{receipttable}
\receiptrow{a88bac06d271526e}{2026-07-21T00:27:07Z}{build+decorrelated}{\badge{BadgeGreen}{current}}
\receiptrow{060c7093ff4899bf}{2026-07-21T04:51:42Z}{build+decorrelated}{\badge{BadgeGreen}{current}}
\receiptrow{ae5851604e101130}{2026-07-21T05:38:57Z}{build+replay+hardened+decorrelated}{\badge{BadgeGreen}{current}}
\receiptrow{85c6bac865d222ba}{2026-07-22T18:27:12Z}{build+decorrelated}{\badge{BadgeGreen}{current}}
\receiptrow{bf722ef76b576c08}{2026-07-31T22:00:08Z}{build+decorrelated}{\badge{BadgeGreen}{current}}
\receiptrow{0caae5e60f19815a}{2026-07-31T22:29:02Z}{build+decorrelated}{\badge{BadgeGreen}{current}}
\receiptrow{9780a725a44c924b}{2026-08-02T23:31:13Z}{build+decorrelated}{\badge{BadgeGreen}{current}}
\receiptrow{490335613ee08847}{2026-08-03T16:27:17Z}{build+decorrelated}{\badge{BadgeGreen}{current}}
\receiptrow{e68b74ba7bf3e8de}{2026-08-03T22:46:01Z}{build+decorrelated}{\badge{BadgeGreen}{current}}
\end{receipttable}
\sourcenote{5 further receipt(s) in the store attest superseded builds of this module and are omitted from this rendering.}
\sourcenote{module \texttt{BernsteinLogicAux.Wasserstein.Complete}; bound build \texttt{06292c010c42\ldots} (bound from receipt fc2b784b4f0c347e, 2026-07-29T08:09:21Z). A receipt records that named checking layers accepted a build: (module, compiled-artifact SHA-256, environment SHA-256, checker layers), minted when the checking gate ran; the store receipts attesting the bound build are listed with the layers each run credited. Nothing here asks the reader to trust the store: every layer is an independently runnable program over the quoted sources --- rebuild the module and re-run the checkers to reproduce the verdict. Receipts attesting superseded builds credit no layer and are omitted from this rendering (the count of omissions, if any, is noted below).}
\begin{receipttable}
\receiptrow{f660300e93359aad}{2026-07-28T23:54:07Z}{build+decorrelated}{\badge{BadgeGreen}{current}}
\receiptrow{ddcacf582a03fe61}{2026-07-29T05:43:06Z}{build+decorrelated}{\badge{BadgeGreen}{current}}
\receiptrow{fc2b784b4f0c347e}{2026-07-29T08:09:21Z}{build+replay+hardened+decorrelated}{\badge{BadgeGreen}{current}}
\receiptrow{439956f3f0160d8a}{2026-07-31T22:00:35Z}{build+decorrelated}{\badge{BadgeGreen}{current}}
\receiptrow{d52523d643b3462b}{2026-07-31T22:29:29Z}{build+decorrelated}{\badge{BadgeGreen}{current}}
\receiptrow{a5c8597e9187735d}{2026-08-02T23:31:45Z}{build+decorrelated}{\badge{BadgeGreen}{current}}
\receiptrow{9f4e8bbd7042154f}{2026-08-03T16:27:49Z}{build+decorrelated}{\badge{BadgeGreen}{current}}
\receiptrow{c0e6026aa1939c74}{2026-08-03T22:46:33Z}{build+decorrelated}{\badge{BadgeGreen}{current}}
\end{receipttable}
\sourcenote{4 further receipt(s) in the store attest superseded builds of this module and are omitted from this rendering.}
\layerline{\layeritem{build}{lake build (elaboration)}{\badge{BadgeGreen}{current}}\layeritem{replay}{leanchecker (official kernel replay)}{\badge{BadgeGreen}{current}}\layeritem{hardened}{safe\_verify (hardened screen)}{\badge{BadgeGreen}{current}}\layeritem{decorrelated}{lean4export + nanoda (decorrelated kernel)}{\badge{BadgeGreen}{current}}}
\sourcenote{\badge{BadgeGray}{comparator: none} no fixture in the three-implementation comparator harness (compiled / kernel-interpreted / reference-interpreted) exercises this rule.}
\ruleblock{HOL4 statement}
\absentbox{NO HOL STATEMENT YET}{The HOL4 development\textquotesingle s object is the base
calculus\textquotesingle s wire-format checker (theory
\texttt{bern\_catalogue} ports the first-order check-function cone of
the core catalogue). \rulename{QFixBall} lives one extension edge above
that stratum --- it concludes a distance judgment of the extension
theory, its rule sort is out of scope of the current baked theory, and
its mechanization lives in the auxiliary Lean package. There is nothing
checker-shaped for a HOL statement to be about yet; the box becomes
writable when a revision of the theory carries the certificate
objects and the corresponding checker
chain is ported.
}
\ruleblock{HOL verifier receipts}
\absentbox{NONE}{no HOL verification receipt exists for this rule.}
\ruleblock{Proof}

\sourcenote{Genre: core rule (tier 1) --- Banach's fixed-point theorem
\cite{banach1922} in certificate form. Scope, stated loudly: the
theorems are parametric in a complete metric space; the affine/Euclidean
instantiation is mechanized, and completeness of the Wasserstein space
is a \emph{named hypothesis}, not a discharged one. Every hypothesis of
the rule is used below and remarked on.}

\textbf{Setting and claim.}\enspace Fix a metric space $(X, d)$, a map
$T : X \to X$, a snapshot $\hat\mu \in X$ (the object the producer
actually computed), and rational certificate data
$\rho \in \Q \cap [0, 1)$, $\delta, r \in \Q$. The three certified rows:
\begin{align*}
&\text{(H-\textsc{anchor})}\quad d(Tx, Ty) \leq \rho\, d(x, y)
  \ \text{ for all } x, y \in \overline{B}(\hat\mu, r)
  && \text{(two-point contraction \emph{on the ball})},\\
&\text{(H-\textsc{resid})}\quad d(T\hat\mu, \hat\mu) \leq \delta
  && \text{(one concrete sweep at the snapshot)},\\
&\text{(H-\textsc{radius})}\quad \delta + \rho\, r \leq r
  && \text{(division-free; equivalent to } r \geq \delta/(1-\rho)\text{)}.
\end{align*}
The claim (\texttt{sound\_QFixBall}): if $X$ is complete, there is
\emph{exactly one} fixed point of $T$ in $\overline{B}(\hat\mu, r)$ ---
so the description $\mathrm{fix}\,T$ is licensed on the ball, and the
distance judgment $d(\hat\mu, \mathrm{fix}\,T) \leq r$ holds; moreover
the sharper anchored bound $d(\hat\mu, \mathrm{fix}\,T) \leq
\delta/(1-\rho)$ holds (\texttt{dist\_fixedPoint\_le},
\texttt{anchored\_le\_radius}). Every hypothesis is anchored at the
\emph{known} point $\hat\mu$; the unknown fixed point appears only in
conclusions --- the anchoring inversion that makes the statement
checkable at all.

\begin{blemma}[the rows are sign-coherent]\label{lem:qfb-signs}
H-\textsc{resid} forces $\delta \geq 0$; together with $\rho < 1$,
H-\textsc{radius} forces $r \geq 0$, so the ball is honest.
\end{blemma}
\begin{proof}
$0 \leq d(T\hat\mu, \hat\mu) \leq \delta$. From H-\textsc{radius},
$(1 - \rho)\, r \geq \delta \geq 0$, and $1 - \rho > 0$.
\end{proof}

\begin{blemma}[self-mapping is derived, never assumed]\label{lem:qfb-maps}
$T$ maps $\overline{B}(\hat\mu, r)$ into itself.
\end{blemma}
\begin{proof}
Let $d(x, \hat\mu) \leq r$. Both $x$ and $\hat\mu$ lie in the ball
(Lemma~\ref{lem:qfb-signs} for $\hat\mu$), so H-\textsc{anchor} applies
to the pair:
\[
d(Tx, \hat\mu) \;\leq\; d(Tx, T\hat\mu) + d(T\hat\mu, \hat\mu)
\;\leq\; \rho\, d(x, \hat\mu) + \delta
\;\leq\; \rho\, r + \delta \;\leq\; r . \qedhere
\]
\end{proof}

\begin{blemma}[in-ball uniqueness]\label{lem:qfb-unique}
If $y, z \in \overline{B}(\hat\mu, r)$ are both fixed by $T$, then
$y = z$.
\end{blemma}
\begin{proof}
$d(y, z) = d(Ty, Tz) \leq \rho\, d(y, z)$ with $\rho < 1$ forces
$d(y, z) \leq 0$, hence $= 0$. Note the proof consumes only the
\emph{two-point} row, on the ball --- no iteration, no completeness.
\end{proof}

\begin{blemma}[the anchored a-priori bound]\label{lem:qfb-anchor}
Any fixed point $y \in \overline{B}(\hat\mu, r)$ satisfies
$d(\hat\mu, y) \leq \delta/(1 - \rho) \leq r$.
\end{blemma}
\begin{proof}
One triangle inequality through the snapshot's image:
\[
d(\hat\mu, y) \;\leq\; d(\hat\mu, T\hat\mu) + d(T\hat\mu, Ty)
\;\leq\; \delta + \rho\, d(\hat\mu, y),
\]
so $(1 - \rho)\, d(\hat\mu, y) \leq \delta$. The cap by $r$ is
H-\textsc{radius} rearranged: $\delta \leq (1 - \rho)\, r$
(\texttt{anchored\_le\_radius}). Uniqueness is not consumed --- the
bound holds for \emph{any} fixed point in the ball, which is what lets
the rule state it before existence is settled.
\end{proof}

\begin{bthm}[soundness of \rulename{QFixBall}]\label{thm:qfb-sound}
Let $X$ be complete and let the three rows hold. Then $T$ has exactly
one fixed point $\mu^{*}$ in $\overline{B}(\hat\mu, r)$, and
$d(\hat\mu, \mu^{*}) \leq \delta/(1-\rho) \leq r$.
\end{bthm}
\begin{proof}
\emph{Existence.} By Lemma~\ref{lem:qfb-maps} the closed ball is
$T$-invariant; it is a closed subset of a complete space, hence
complete, and $T$ restricted to it is a $\rho$-contraction
(H-\textsc{anchor}). Banach's theorem \cite{banach1922} applies to the
restriction: the Picard iterates $T^n \hat\mu$ form a Cauchy sequence
($d(T^{n+1}\hat\mu, T^n\hat\mu) \leq \rho^n \delta$, summable) whose
limit $\mu^{*}$ lies in the ball and is fixed. Existence is thus
\emph{manufactured} by the contraction --- the rule consumes no
separate existence certificate, so the concluded description is
existential without vacuity.

\emph{Uniqueness} is Lemma~\ref{lem:qfb-unique}; the \emph{bound} is
Lemma~\ref{lem:qfb-anchor}. Together: $\mathrm{fix}\,T$ is well-defined
on the ball (the description license the judgment carries) and the
distance judgment holds at radius $r$, with $\delta/(1-\rho)$ as the
sharper certified value.
\end{proof}

\begin{bprop}[the affine instantiation]\label{prop:qfb-affine}
On a finite-dimensional Euclidean space, let $T x = A x + c$ with $A$
linear and an operator-norm row $\lVert A \rVert \leq \rho < 1$. Then
H-\textsc{anchor} holds globally (a fortiori on the ball), completeness
holds by instance, and Theorem~\ref{thm:qfb-sound} applies verbatim
(\texttt{sound\_QFixBall\_affine}).
\end{bprop}
\begin{proof}
$d(Tx, Ty) = \lVert A(x - y) \rVert \leq \lVert A \rVert\,
\lVert x - y \rVert \leq \rho\, d(x, y)$. The operator-norm row is the
affine tier's discharge of H-\textsc{anchor}: it is certified by
positive-semidefiniteness of $\rho^2 I - A^{\top} A$, which is
\rulename{PSDCert}-family machinery on rational data.
\end{proof}

\begin{bremark}[why each hypothesis is needed]\label{rem:qfb-omit}
\emph{Drop $\rho < 1$}: the identity map has residual $0$ at every
snapshot and every radius row holds with $\rho = 1$, yet every point is
fixed --- the description license (uniqueness) fails, so the rule
\emph{refuses} $\rho \geq 1$ rather than concluding.
\emph{Drop H-\textsc{radius}}: on $\R$, take $Tx = \tfrac{1}{2}x + 2$,
$\hat\mu = 0$, so $\rho = \tfrac12$, $\delta = 2$; the unique fixed
point is $4$. With $r = 1$ the ball $[-1, 1]$ contains no fixed point
--- and indeed $\delta + \rho r = \tfrac52 > 1$: the radius row fails,
correctly. \emph{Drop H-\textsc{resid}}: nothing ties the ball to the
dynamics --- any $r$ would pass on any map with small $\rho$, concluding
a distance bound to a fixed point that can sit arbitrarily far away.
\emph{Weaken completeness}: on $X = \Q \cap [1, 2]$ (incomplete),
$Tx = \tfrac{x}{2} + \tfrac{1}{x}$ maps $X$ into itself with
$|T'| \le \tfrac12$ and contracts toward the irrational $\sqrt{2}$ ---
iterates Cauchy, no fixed point in $X$, and the certificate rows even
pass ($\hat\mu = \tfrac32$, $r = \tfrac12$, $\rho = \tfrac12$,
$\delta = \tfrac{1}{12}$, with $\delta + \rho r = \tfrac13 \le
\tfrac12$); the conclusion's existential half is exactly what
completeness buys. The two-point row being
required only \emph{on the ball} is deliberate: locality is what a
finite certificate can attest (globally, $\lVert A \rVert < 1$ may even
be false while the ball-restricted bound holds).
\end{bremark}

\begin{bremark}[relation to the literature]
The theorem is Banach's \cite{banach1922}; the closed-ball variant with
the condition $d(T\hat\mu, \hat\mu) \leq (1 - \rho)\, r$ and the
a-priori bound $\delta/(1-\rho)$ is the textbook local form
\cite[\S I.1]{granasdugundji2003}. The certificate-form contribution is
the \emph{anchoring inversion}: classical statements estimate the
distance from an iterate to the limit; the rule re-expresses every
hypothesis as a row checkable at the known snapshot in exact rational
arithmetic --- contraction as an operator-norm (or two-point) row,
residual as one evaluated sweep, radius as division-free arithmetic ---
so the analytic content survives while everything the checker touches
is decidable. The completeness gap for measure-space subjects
(Wasserstein distance over a Euclidean base) is carried as a named
hypothesis; the pilot certificate discharges it by pinning the subject
to the affine/Euclidean reduction of
Proposition~\ref{prop:qfb-affine}.
\end{bremark}

\begin{anchors}
Lemma \ref{lem:qfb-signs} & \texttt{delta\_nonneg}, \texttt{radius\_nonneg}\\
Lemma \ref{lem:qfb-maps} & \texttt{mapsTo\_closedBall}\\
Lemma \ref{lem:qfb-unique} & \texttt{fixedPoint\_unique}\\
Lemma \ref{lem:qfb-anchor} & \texttt{dist\_fixedPoint\_le}; \texttt{anchored\_le\_radius}\\
Theorem \ref{thm:qfb-sound} & \texttt{sound\_QFixBall} (existence via the contraction on the complete invariant ball)\\
Proposition \ref{prop:qfb-affine} & \texttt{lipschitzWith\_affine}; \texttt{sound\_QFixBall\_affine}\\
\end{anchors}


\rulesection{93}{QDistTransport}{core rule}
\label{blcatrule-QDistTransport}
\ruleblock{The rule (Steele metanotation)}
\sourcenote{as realized in the mechanized development; source row: \texttt{\texttt{QTSL-rules.toml}} --- \emph{QDistTransport}. The metanotation statement is typeset by hand for notational quality; its mechanized ground truth is the Lean statement below, extracted from the sources at every build of this book and hash-bound. The rule \emph{roster} is machine-checked (every mechanized rule tag round-trips to exactly one catalogue row); the statement \emph{text} is editorial, checked by eye against the extracted Lean beside it.}
\begin{mathpar}
\inferrule*{
  J = \bigl(\mathcal{M}_c \vDash s\bigr) \ \text{(strict, law-level)} \quad \text{or} \quad J = \bigl(\mathsf{QDist}[\mathrm{TV}](\mu_c, \nu) \le \delta\bigr), \ \text{on the coarse side} \\
  \bigl\{\vdash \mathrm{tv}\bigl(\pi_{\#}\, K(s, \cdot),\ Q(\pi(s), \cdot)\bigr) \le \varepsilon\bigr\}_{s \in S_{\mathrm{f}}} \ \text{(} \sigma\text{-witness rows, one per fine state; } \varepsilon = 0 \ \text{is exact intertwining)} \\
  \vdash \mu^{0}_{c} = \pi_{\#}\, \mu^{0}_{\mathrm{f}} \ \text{(initial row)} \\
  n \in \N \ \text{(declared horizon; } \varepsilon\text{-cell only)}
}{
  \text{exact cell: } J \ \text{transports unchanged} \qquad \varepsilon\text{-cell: } \mathsf{QDist}[\mathrm{TV}]\bigl(\pi_{\#}(\text{fine law at } n),\ \text{coarse law at } n\bigr) \le n\,\varepsilon
}
\end{mathpar}

\ruleflavor{$\pi : S_{\mathrm{f}} \to S_c$ is the state projection, $K$ and $Q$ the fine and coarse kernels, $\mu^0$ the initial laws; each $\sigma$-witness row asserts that one row of the fine kernel pushes along $\pi$ to within total-variation $\varepsilon$ of the coarse kernel\textquotesingle s row at the projected state. Unbounded horizons refuse without a contraction row. The displayed $J$ is restricted to the built judgment classes (strict law-level specs and $\mathsf{QDist}[\mathrm{TV}]$); the successor cells widen it as they land. Transport is a matrix (judgment class $\times$ witness class), each built cell carrying its own lemma and refusal condition; sequential witnesses compose under the certified budget $\varepsilon_1 + \varepsilon_2$ (an upper bound --- the actual composite defect may be smaller), and the strict column is the $\varepsilon = 0$ instantiation of the family. Built at the finite-kernel tier; the Lipschitz-pushforward ($\mathbb{W}_p$) and description-subject cells are named successors.}
\par\medskip
\rulename{QDistTransport} moves judgments across a simulation map. It
is a \emph{matrix} of transport cells --- judgment class (strict spec
vs.\ distance judgment) $\times$ witness class (exact intertwining
vs.\ $\varepsilon$-approximate intertwining) --- each cell carrying its
own conclusion form, its own lemma, and its own refusal condition. A
producer reaches for it to certify a fine model through a coarse one:
prove the spec on the abstraction, exhibit the simulation witness, and
transport. The two built kernel-tier cells: under an \emph{exact}
witness (every row of the fine kernel pushes along the readout map
exactly onto the coarse kernel\textquotesingle s row) laws transport
unchanged at every horizon, so any judgment about the coarse law pulls
back with zero degradation; under an $\varepsilon$-witness (every row
pushes to within total-variation distance $\varepsilon$) a strict
judgment does NOT transport --- what transports is a \emph{distance}
judgment, with $\varepsilon$ inherited and horizon-compounded: the
pushed fine law tracks the coarse law to $n\,\varepsilon$ at declared
horizon $n$. The degrading cell is what makes approximate simulation a
distance-judgment \emph{producer}.

What the checker recomputes: the per-row intertwining identities (exact
case) or per-row total-variation defects against the declared
$\varepsilon$ (approximate case), well-formedness of both kernels, and
the horizon arithmetic. The compounding law $n\,\varepsilon$ is proved,
never guessed: one-step defect by mixture convexity, then the coarse
Markov operator\textquotesingle s total-variation nonexpansiveness
(data processing --- the coarse chain cannot amplify tracking error),
then telescoping. Sequential witnesses compose with certified budget
$\varepsilon_1 + \varepsilon_2$ --- an upper bound with no additional
bookkeeping slack (plain addition is the honest lattice at this tier),
though the composite\textquotesingle s actual defect may be smaller.

What would break: without the coarse kernel\textquotesingle s row
stochasticity the operator could amplify differences and the telescoping
collapses; without a genuine probability law at the start, mixture
convexity fails; and without a contraction row the linear growth in $n$
is sharp, which is why the $\varepsilon$-cell refuses unbounded
horizons rather than extrapolating.

\cataloguenote{The strict column is the $\varepsilon = 0$ instantiation of the family
--- one biconditional, so the exact cell re-exports the marginal-process
lemma rather than restating it. Analytic lineage: approximate
(bi)simulation with quantitative degradation \cite{girardpappas2007}
and behavioural pseudometrics for probabilistic systems
\cite{desharnais2004}. Built at the finite-kernel tier (the scope the
Lean statements carry); the Lipschitz-pushforward cells for
$\mathbb{W}_p$ registry entries, the description-subject cells, and the
per-cell composition-coherence squares are named successors.
}
\ruleblock{Lean statement}
\declhead{BernsteinLogicAux.QDistTransport}{sound\_QDistTransport\_eps}{\badge{BadgeGreen}{statement bound}}
\begin{leansrc}
/-- **`sound_QDistTransport_eps` — the degrading cell's transport
theorem** (matrix cell: strict judgment × ε-witness ⇒ distance
judgment): at horizon `n`, the pushed fine law and the coarse law are
within TV-distance `n·ε`. Telescoping: one-step defect (`ε`, above) +
TV-nonexpansiveness of the coarse operator (§3) + triangle. This is
the theorem that makes "ε inherited, horizon-compounded" a proved
compounding law rather than a guess; with no contraction row the
linear growth is the honest rate (the refusal condition for unbounded
horizons lives spec-side). -/
theorem sound_QDistTransport_eps [DecidableEq β] {π : α → β}
    {K : α → α → ℝ} {Q : β → β → ℝ} {ε : ℝ} {μ₀ : α → ℝ}
    (hK : IsMarkov K) (hQ : IsMarkov Q) (hμ₀ : IsDist μ₀)
    (h : EpsPushKernel π K Q ε) (n : ℕ) :
    tvDist (push π (law K μ₀ n)) (law Q (push π μ₀) n) ≤ n * ε
\end{leansrc}
\declhead{BernsteinLogicAux.QDistTransport}{epsPushKernel\_comp}{\badge{BadgeGreen}{statement bound}}
\begin{leansrc}
/-- **ε-arithmetic composes** (derived, not guessed): transport along
`σ₁` then `σ₂` carries the certified budget `ε₁ + ε₂` — an upper bound
(triangle through the mid-tier pushforward + data processing; the actual
composite defect may be smaller). Plain addition is the right
composition law at the TV tier. -/
theorem epsPushKernel_comp [DecidableEq β] [DecidableEq γ]
    {π₁ : α → β} {π₂ : β → γ} {K₁ : α → α → ℝ} {K₂ : β → β → ℝ}
    {K₃ : γ → γ → ℝ} {ε₁ ε₂ : ℝ}
    (h₁ : EpsPushKernel π₁ K₁ K₂ ε₁) (h₂ : EpsPushKernel π₂ K₂ K₃ ε₂) :
    EpsPushKernel (π₂ ∘ π₁) K₁ K₃ (ε₁ + ε₂)
\end{leansrc}
\declhead{BernsteinLogicAux.QDistTransport}{sound\_QDistSelfTransport}{\badge{BadgeGreen}{statement bound}}
\begin{leansrc}
/-- **Matrix cell (distance judgment × strict witness), TV entry**: a
TV distance judgment transports along any simulation/readout map
unchanged (registry reading: TV → TV at Lipschitz constant 1). -/
theorem sound_QDistSelfTransport [DecidableEq β] (π : α → β)
    {μ ν : α → ℝ} {ε : ℝ} (h : tvDist μ ν ≤ ε) :
    tvDist (push π μ) (push π ν) ≤ ε
\end{leansrc}
\declhead{BernsteinLogicAux.QDistTransport}{strict\_cell\_exact}{\badge{BadgeGreen}{statement bound}}
\begin{leansrc}
/-- **Matrix cell (strict judgment × strict/kernel witness)**: under
the exact witness, laws transport EXACTLY at every horizon
(`QRelChain.push_law` re-exported as the matrix's strict kernel cell)
— so any judgment about the coarse law pulls back with zero
degradation, and the degrading cell's `n·ε` bound specializes to `0`. -/
theorem strict_cell_exact [DecidableEq β] {π : α → β} {K : α → α → ℝ}
    {Q : β → β → ℝ} (h : PushKernel π K Q) (μ₀ : α → ℝ) (n : ℕ) :
    push π (law K μ₀ n) = law Q (push π μ₀) n
\end{leansrc}
\declhead{BernsteinLogicAux.QDistTransport}{pushKernel\_iff\_eps\_zero}{\badge{BadgeGreen}{statement bound}}
\begin{leansrc}
/-- The exact intertwining is precisely the `ε = 0` witness. -/
theorem pushKernel_iff_eps_zero [DecidableEq β] {π : α → β}
    {K : α → α → ℝ} {Q : β → β → ℝ} :
    PushKernel π K Q ↔ EpsPushKernel π K Q 0
\end{leansrc}
\ruleblock{Correspondence (Steele display $\leftrightarrow$ Lean)}
\badge{BadgeGreen}{reviewer-certified against the displayed texts}
\sourcenote{The dictionary below is a reviewer-certified statement of how the mechanized statements above interpret the displayed rule: display metavariable $\mapsto$ Lean binder, each display hypothesis $\mapsto$ the premise or side condition carrying it, the display conclusion $\mapsto$ the goal --- with any deliberate difference in width between display and mechanization named outright. Certification is recorded as a hash triple over (the displayed cell, the extracted statement texts, this dictionary\textquotesingle s own text); if any of the three changes, this block turns amber until a reviewer re-certifies. The correspondence claim itself remains editorial --- reviewer-grade, not machine-checked; its machine-checked successor types each soundness statement against the rule as a structured object, and supersedes this block per rule as it lands.}
The display is a matrix and the bound set is its built column, cell
by cell. The $\varepsilon$-cell\textquotesingle s conclusion
$\mathsf{QDist}[\mathrm{TV}](\text{fine pushed law}, \text{coarse
law}) \le n\varepsilon$ at horizon $n$ $\mapsto$
\texttt{sound\_QDistTransport\_eps : tvDist (push π (law K μ₀ n))
(law Q (push π μ₀) n) ≤ n * ε} --- the compounding law as a proved
induction (one-step defect + TV-nonexpansiveness of the coarse
operator + triangle), never a guess.
\begin{itemize}
\item the $\sigma$-witness rows $\mapsto$ \texttt{EpsPushKernel π K Q
  ε} (every fine row pushes to within TV $\varepsilon$ of the coarse
  row at the projected state), with the exact-intertwining column
  certified as its $\varepsilon = 0$ face:
  \texttt{pushKernel\_iff\_eps\_zero}, and the strict cell
  \texttt{strict\_cell\_exact} re-exporting \texttt{push\_law} ---
  laws transport \emph{exactly} at every horizon, zero degradation.
\item the initial row (coarse initial law $= \pi_{\#}$ fine initial)
  $\mapsto$ carried by construction: the coarse arguments in every
  goal ARE \texttt{push π μ₀} --- mismatched initials are
  unrepresentable in the statement, so no hypothesis row exists.
\item the declared horizon $\mapsto$ the explicit \texttt{n : ℕ}
  binder; the refusal condition for unbounded horizons without a
  contraction row is spec-side with no Lean counterpart --- the
  honest linear rate IS the certified statement.
\item sequential accumulation $\mapsto$ \texttt{epsPushKernel\_comp}:
  certified composite budget $\varepsilon_1 + \varepsilon_2$, plain
  addition as the displayed lattice --- an upper bound (the actual
  composite defect may be smaller); the distance-judgment $\times$
  strict-witness cell $\mapsto$ \texttt{sound\_QDistSelfTransport}
  (TV data processing at $L = 1$).
\item consumed regularity: the $\varepsilon$-cell alone demands
  well-formedness (\texttt{IsMarkov} both kernels, \texttt{IsDist
  μ₀}) --- the mixture-convexity and nonexpansiveness steps use it
  --- while the strict cell and the composition law are
  hypothesis-free beyond their witnesses: visible in the binder
  rows.
\end{itemize}
Width note: finite-kernel tier (\texttt{Fintype},
\texttt{DecidableEq} for fibre sums), as the flavor text
confines it; the unbuilt cells ($\mathbb{W}_p$ Lipschitz
pushforward, description subjects, per-cell coherence squares) are
the conditional residue, named in the note, bound by nothing --- and
the judgment-class abstraction of the display\textquotesingle s $J$
(``or a strict spec\textquotesingle\textquotesingle) is realized
concretely: strict specs ride \texttt{strict\_cell\_exact}, distance
judgments ride the TV theorems; no single parameterized $J$ exists in
Lean, the matrix\textquotesingle s rows are separate theorems.

\ruleblock{Lean verifier receipts}
\sourcenote{module \texttt{BernsteinLogicAux.QDistTransport}; bound build \texttt{4bd1c5103f38\ldots} (bound from receipt 6ff1345414960d8f, 2026-07-21T05:38:55Z). A receipt records that named checking layers accepted a build: (module, compiled-artifact SHA-256, environment SHA-256, checker layers), minted when the checking gate ran; the store receipts attesting the bound build are listed with the layers each run credited. Nothing here asks the reader to trust the store: every layer is an independently runnable program over the quoted sources --- rebuild the module and re-run the checkers to reproduce the verdict. Receipts attesting superseded builds credit no layer and are omitted from this rendering (the count of omissions, if any, is noted below).}
\begin{receipttable}
\receiptrow{6ff1345414960d8f}{2026-07-21T05:38:55Z}{build+replay+hardened+decorrelated}{\badge{BadgeGreen}{current}}
\end{receipttable}
\sourcenote{13 further receipt(s) in the store attest superseded builds of this module and are omitted from this rendering.}
\layerline{\layeritem{build}{lake build (elaboration)}{\badge{BadgeGreen}{current}}\layeritem{replay}{leanchecker (official kernel replay)}{\badge{BadgeGreen}{current}}\layeritem{hardened}{safe\_verify (hardened screen)}{\badge{BadgeGreen}{current}}\layeritem{decorrelated}{lean4export + nanoda (decorrelated kernel)}{\badge{BadgeGreen}{current}}}
\sourcenote{\badge{BadgeGray}{comparator: none} no fixture in the three-implementation comparator harness (compiled / kernel-interpreted / reference-interpreted) exercises this rule.}
\ruleblock{HOL4 statement}
\absentbox{NO HOL STATEMENT YET}{The HOL4 development checks the base calculus\textquotesingle s wire
format (theory \texttt{bern\_catalogue}). The transport
matrix\textquotesingle s cells conclude judgments of the extension
theory, whose rule sorts are out of scope of the current baked theory;
the mechanized cells live in the auxiliary Lean package. Nothing exists
yet for a HOL statement to be about; the box becomes writable when a
revision carries the distance-judgment rows and its checker chain is
ported.
}
\ruleblock{HOL verifier receipts}
\absentbox{NONE}{no HOL verification receipt exists for this rule.}
\ruleblock{Proof}

\sourcenote{Genre: core rule (tier 1). Scope, stated loudly: the
finite-kernel tier --- finite state spaces, distribution evolution,
total-variation entries of the registry. The Lipschitz-pushforward
cells for Wasserstein entries, the description-subject cells, and the
per-cell composition-coherence squares are named successors, not
claimed. Every hypothesis is used and remarked on.}

\textbf{Setting and claim.}\enspace A \emph{fine} chain (states
$\alpha$, kernel $K$, initial law $\mu_0$), a \emph{coarse} chain
(states $\beta$, kernel $Q$), and a readout map
$\pi : \alpha \to \beta$. Write $\pi_{\#}$ for pushforward,
$\mathrm{law}_K(n)$ for the time-$n$ law, and
$\mathrm{TV}(p, q) = \tfrac12 \sum_l |p(l) - q(l)|$. The witness comes
in two grades: \emph{exact intertwining} --- every kernel row pushes
exactly, $\pi_{\#} K(s, \cdot) = Q(\pi(s), \cdot)$ --- and
\emph{$\varepsilon$-intertwining} ---
$\mathrm{TV}\bigl(\pi_{\#} K(s, \cdot),\ Q(\pi(s), \cdot)\bigr) \leq
\varepsilon$ for every $s$. The transport matrix's built cells, which
are the claims proved below: \textbf{(strict $\times$ exact)} laws
transport unchanged at every horizon, so any judgment about the coarse
law pulls back with zero degradation
(\texttt{strict\_cell\_exact}); \textbf{(distance $\times$ strict
witness)} a total-variation distance judgment moves along any readout
map unchanged (\texttt{sound\_QDistSelfTransport});
\textbf{(strict $\times$ $\varepsilon$-witness --- the degrading
cell)} what transports is a \emph{distance} judgment:
$\mathrm{TV}\bigl(\pi_{\#}\mathrm{law}_K(n),\
\mathrm{law}_Q(n)\bigr) \leq n\,\varepsilon$ at declared horizon $n$
(\texttt{sound\_QDistTransport\_eps}); and witnesses compose with
$\varepsilon_1 + \varepsilon_2$ (\texttt{epsPushKernel\_comp}).

\begin{blemma}[data processing: pushforward is TV-nonexpansive]
\label{lem:qdt-dpi}
For any map $\pi$ and any $p, q$:
$\mathrm{TV}(\pi_{\#} p, \pi_{\#} q) \leq \mathrm{TV}(p, q)$.
\end{blemma}
\begin{proof}
Per label $b$, $|\sum_{\pi(s) = b} (p - q)(s)| \leq
\sum_{\pi(s) = b} |p - q|(s)$; summing over $b$ regroups the right side
into $\sum_s |p - q|(s)$. Merging states can only cancel differences,
never create them \cite[Ch.~4]{levinperes2017}. This lemma \emph{is}
the (distance $\times$ strict-witness) cell: a certified
$\mathrm{TV} \leq \varepsilon$ on the fine side transports to the
coarse side unchanged --- the registry reading of total variation as a
Lipschitz-$1$ entry under pushforward.
\end{proof}

\begin{blemma}[the coarse operator cannot amplify]\label{lem:qdt-markov}
If $Q$ is row-stochastic with nonnegative entries, then
$\mathrm{TV}(\mathrm{step}_Q\, p,\ \mathrm{step}_Q\, q) \leq
\mathrm{TV}(p, q)$.
\end{blemma}
\begin{proof}
$|\mathrm{step}_Q(p - q)(t)| \leq \sum_s |p - q|(s)\, Q(s, t)$ by the
triangle inequality and $Q \geq 0$; summing over $t$ and swapping sums,
the inner $\sum_t Q(s, t) = 1$ leaves $\sum_s |p - q|(s)$. Both
hypotheses on $Q$ are consumed: nonnegativity to pull the absolute
value through, unit row sums to collapse the mass.
\end{proof}

\begin{blemma}[one-step defect]\label{lem:qdt-defect}
Let $\mu$ be a probability law and let the $\varepsilon$-intertwining
rows hold. Then
$\mathrm{TV}\bigl(\pi_{\#}(\mathrm{step}_K \mu),\
\mathrm{step}_Q(\pi_{\#}\mu)\bigr) \leq \varepsilon$.
\end{blemma}
\begin{proof}
Both sides are $\mu$-mixtures of per-state rows:
$\pi_{\#}(\mathrm{step}_K \mu) = \sum_s \mu(s)\,
\pi_{\#}K(s, \cdot)$ and $\mathrm{step}_Q(\pi_{\#}\mu) =
\sum_s \mu(s)\, Q(\pi(s), \cdot)$ (the second by regrouping the coarse
step over the fibres of $\pi$). Total variation is jointly convex ---
a mixture of pairs is bounded by the mixture of pairwise distances ---
so the defect is at most $\sum_s \mu(s) \cdot \varepsilon =
\varepsilon$, using that $\mu$ is nonnegative with total mass $1$.
\end{proof}

\begin{bthm}[the degrading cell: horizon-compounded transport]
\label{thm:qdt-eps}
With $K, Q$ row-stochastic, $\mu_0$ a probability law, and the
$\varepsilon$-intertwining rows: for every $n$,
\[
\mathrm{TV}\bigl(\pi_{\#}\,\mathrm{law}_K(n),\ \mathrm{law}_Q(n)\bigr)
\;\leq\; n\, \varepsilon,
\]
where the coarse chain starts from $\pi_{\#}\mu_0$.
\end{bthm}
\begin{proof}
Induction on $n$; the base is $\mathrm{TV}(x, x) = 0$. Step, by the
triangle inequality through the hybrid term
$\mathrm{step}_Q\bigl(\pi_{\#}\,\mathrm{law}_K(n)\bigr)$:
\[
\mathrm{TV}\bigl(\pi_{\#}\,\mathrm{law}_K(n{+}1),\ \mathrm{law}_Q(n{+}1)\bigr)
\;\leq\;
\underbrace{\mathrm{TV}\bigl(\pi_{\#}\mathrm{step}_K\,\mathrm{law}_K(n),\
\mathrm{step}_Q\,\pi_{\#}\mathrm{law}_K(n)\bigr)}_{\leq\ \varepsilon
\ \text{(Lemma \ref{lem:qdt-defect}, at the time-}n\text{ law)}}
\;+\;
\underbrace{\mathrm{TV}\bigl(\mathrm{step}_Q\,\pi_{\#}\mathrm{law}_K(n),\
\mathrm{step}_Q\,\mathrm{law}_Q(n)\bigr)}_{\leq\ n \varepsilon
\ \text{(Lemma \ref{lem:qdt-markov} + induction)}} .
\]
The compounding law is thus \emph{proved}, not guessed: one fresh
defect per step, prior error never amplified. A strict judgment does
NOT transport along an $\varepsilon$-witness; what transports is a
distance judgment with $\varepsilon$ inherited and horizon-compounded
--- precisely this cell.
\end{proof}

\begin{bprop}[witness composition]\label{prop:qdt-comp}
If $\sigma_1 = (\pi_1, \varepsilon_1)$ intertwines $K_1$ with $K_2$ and
$\sigma_2 = (\pi_2, \varepsilon_2)$ intertwines $K_2$ with $K_3$, then
$\pi_2 \circ \pi_1$ is an $(\varepsilon_1 + \varepsilon_2)$-witness
from $K_1$ to $K_3$.
\end{bprop}
\begin{proof}
At any state $s$, triangle through the mid-tier row
$\pi_{2\#} K_2(\pi_1(s), \cdot)$: the first leg is
$\mathrm{TV}\bigl(\pi_{2\#}\pi_{1\#}K_1(s,\cdot),\
\pi_{2\#}K_2(\pi_1 s, \cdot)\bigr) \leq \varepsilon_1$ by
Lemma~\ref{lem:qdt-dpi} applied to $\sigma_1$'s row defect; the second
leg is $\sigma_2$'s row at $\pi_1(s)$. Plain addition records the honest
\emph{upper bound} at this tier --- $\varepsilon_1 + \varepsilon_2$ is
what the triangle gives, with no minimality claim: data processing can
gain slack (a constant second projection has composite defect $0$
whatever $\varepsilon_1$ was).
\end{proof}

\begin{bprop}[the strict column is the $\varepsilon = 0$ instantiation]
\label{prop:qdt-strict}
Exact intertwining holds iff the rows are $0$-intertwining; under it,
$\pi_{\#}\,\mathrm{law}_K(n) = \mathrm{law}_Q(n)$ for every $n$, and
every judgment about the coarse law pulls back with zero degradation.
\end{bprop}
\begin{proof}
($\Rightarrow$) equal rows have distance $0$. ($\Leftarrow$)
$\mathrm{TV}(p, q) \leq 0$ forces $p = q$ entrywise (a sum of absolute
values is zero only if each vanishes). The exact law transport is then
the marginal-process induction (the same pushed-row-identity induction
as the \rulename{QCoupleDrift} section's marginal-process lemma, there
instantiated at product projections, here at an arbitrary readout
map), or equivalently Theorem~\ref{thm:qdt-eps} at $\varepsilon = 0$.
\end{proof}

\begin{bremark}[why each hypothesis is needed]\label{rem:qdt-omit}
\emph{Row-stochasticity of $Q$} (Lemma~\ref{lem:qdt-markov}): a
non-stochastic coarse operator amplifies --- with $Q = 2\,\mathrm{Id}$
each step doubles any existing difference and the telescoping
collapses; with signed entries the absolute value cannot be pulled
through. \emph{$\mu$ a probability law}
(Lemma~\ref{lem:qdt-defect}): mixture convexity needs nonnegative
weights of total mass one --- weights summing to $2$ double the
defect. \emph{The declared horizon}: the linear rate is sharp to first order
without further rows --- on the two-state fine chain that leaks mass
$\varepsilon$ toward an absorbing label each step, against a static
coarse chain, the distance at horizon $n$ is
$1 - (1 - \varepsilon)^n = n\varepsilon - O(n^2\varepsilon^2)$, so no
$o(n\varepsilon)$ bound holds, and over unbounded horizons the error
climbs to order one --- the $\varepsilon$-cell therefore
\emph{refuses} unbounded horizons rather than extrapolating; a
contraction row is the named mechanism that would buy a
uniform-in-$n$ bound, and it is a different certificate. \emph{Per-row quantification} of the intertwining: a
defect bound only at states charged by $\mu_0$ fails at later times,
since the law spreads --- the rows are quantified over all states
precisely so Lemma~\ref{lem:qdt-defect} applies at \emph{every}
time-$n$ law.
\end{bremark}

\begin{bremark}[relation to the literature]
The degrading cell is the kernel-tier, exact-arithmetic face of
approximate simulation with quantitative degradation
\cite{girardpappas2007}: there, an $\varepsilon$-approximate
(bi)simulation transports properties up to a metric slack; here the
slack is total variation, the systems are finite kernels, and the
horizon-compounding law is proved from mixture convexity plus data
processing. Behavioural pseudometrics for probabilistic systems
\cite{desharnais2004} are the semantic tradition the distance-judgment
registry discretizes; the coupling toolkit is classical
\cite{lindvall1992,levinperes2017}. The matrix framing --- judgment
class $\times$ witness class, one lemma and one refusal condition per
cell --- is the rule's own contribution: it makes ``what survives
transport along a defective witness'' a typed question with a proved
answer per cell, instead of a convention.
\end{bremark}

\begin{anchors}
Lemma \ref{lem:qdt-dpi} & \texttt{tvDist\_push\_le}; cell form \texttt{sound\_QDistSelfTransport}\\
Lemma \ref{lem:qdt-markov} & \texttt{tvDist\_step\_le}\\
Lemma \ref{lem:qdt-defect} & \texttt{EpsPushKernel}; \texttt{push\_step\_defect}\\
Theorem \ref{thm:qdt-eps} & \texttt{sound\_QDistTransport\_eps} (triangle: \texttt{tvDist\_triangle})\\
Proposition \ref{prop:qdt-comp} & \texttt{epsPushKernel\_comp}; \texttt{push\_comp}\\
Proposition \ref{prop:qdt-strict} & \texttt{pushKernel\_iff\_eps\_zero}; \texttt{eq\_of\_tvDist\_le\_zero}; \texttt{strict\_cell\_exact}\\
\end{anchors}


\blcatsection{Stored-Chain Certificate Species}
\rulesection{94}{DtmcAbsorbBracket}{core rule}
\label{blcatrule-DtmcAbsorbBracket}
\ruleblock{The rule (Steele metanotation)}
\sourcenote{as realized in the mechanized development; source row: \texttt{\texttt{QTSL-rules.toml}} --- \emph{DtmcAbsorbBracket}. The metanotation statement is typeset by hand for notational quality; this rule cites no Lean statements yet (the absence is stated loudly below). The rule \emph{roster} is machine-checked (every mechanized rule tag round-trips to exactly one catalogue row); the statement \emph{text} is editorial.}
\begin{mathpar}
\inferrule*{
  q : \mathsf{DAbsQuery} \\
  \mathsf{qtsl\_dabs\_wf}(q) \\
  \text{stored rows of } q.\mathsf{chain}\ (\mathsf{dprob}, \mathsf{dinit}, \mathsf{dclass}, \mathsf{dreject})
}{
  \mathsf{qtsl\_dabs\_absorb\_bracket}(q,\ c,\ lo,\ hi)
}
\end{mathpar}

\par\medskip
\rulename{DtmcAbsorbBracket} is a certificate species rather than a
proof-node rule: a checker-recomputed claim over the stored rows of a
finite chain, introducing zero new judgment sorts and zero new
specification formers. It certifies guaranteed two-sided brackets on
the eventual-absorption measure of one class of a finite stored-row
chain --- and, through its normalized sibling
\rulename{DtmcPosteriorBracket}, on normalized conditional measures
(posteriors of discrete probabilistic programs, where the compilation
of a program to a chain is an untrusted translation outside the trust
path). The species is stated at full generality: any finite chain, any
class labeling, any horizon.

The evaluation semantics, once: $m_0$ is the chain\textquotesingle s
stored initial measure; $m_{t+1} = m_t K$, where $K$ is the stored
transition-row matrix COMPLETED by a self-loop at every state with no
stored outgoing row (\emph{rowless-is-absorbing}). For a query with
horizon $T$, class labeling, and rejection labeling: $N_c$ is the
class-$c$ mass of $m_T$; $S = \sum_c N_c$; $R = 1 - S -
(\text{rejected mass of } m_T)$; and $R_c$ is the
\emph{reachability-sharpened} per-class residual --- the part of $R$
sitting at states from which some class-$c$ state is reachable
through strictly positive stored rows (a backward walk seeded at the
class, run under the same work budget; budget exhaustion during the
walk refuses the whole claim, never verdicts). The claim holds iff
$lo \le N_c$ and $N_c + R_c \le hi$. Mass that can never arrive in
the class is never charged to it: a class unreachable from every
alive state has $R_c = 0$ and a width-zero bracket at any $T \ge 1$,
with no termination assumption; and $R_c \le R$ always, so any claim
admissible under the uniform-residual reading stays admissible.

The completion convention is the species\textquotesingle{}
load-bearing design move: a truncated row set can only WIDEN brackets
(unclassified mass accrues to the residual, never to any class), so
row-set completeness is a tightness concern rather than a soundness
obligation --- there is no closure side condition, deliberately, and
under-expansion is unrepresentable as an attack. Derived quantities
($N_c$, $S$, $R$, $R_c$, ratio denominators) are never stored: the
evaluator recomputes them from the chain rows, so forging an
intermediate is unrepresentable; claims may be looser than canonical,
never tighter.

Soundness has two sides. Below, absorption is monotone: classed states
are demanded stored-absorbing (the well-formedness side condition),
so class mass never decreases past the horizon. Above, the
class\textquotesingle s co-reachable set $B_C$ (states with a
strictly-positive-support path into the class; $C \subseteq B_C$) has
a forward-closed complement --- no positive-support row leads from
outside $B_C$ into it --- so the mass of $B_C$ is nonincreasing and
later class gain is bounded by $\mathrm{mass}_T(B_C \setminus C) =
R_c$. The uniform-residual bound ($R_c$ relaxed to $1 -
\mathrm{mass}_T(\mathrm{Abs})$) is the pre-sharpening corollary. The
argument is mechanized in the accompanying Lean development
(\texttt{SuFunctor/Adequacy/DtmcAbsorb.lean}; sorry-free, standard
axiom closure, axiom-audited on every build): the sharpened upper
side \texttt{dtmcAbsorb\_mass\_le\_reach} is notably HYPOTHESIS-FREE
on the class --- it holds for any class of any finite stochastic
chain, the well-formedness demands protecting only the monotone lower
side; beside it sit \texttt{dtmcAbsorbBracket\_reach\_sound} (the
bracket, both sides), \texttt{dtmcAbsorb\_reach\_residual\_le} (the
never-loosens warrant $R_c \le 1 - \mathrm{mass}_T(\mathrm{Abs})$,
itself needing only that absorbing states absorb),
\texttt{dtmcAbsorb\_eventual\_bracket\_reach} (the claim shape), the
forward-closed generalization (\texttt{FinDTMC.ForwardClosed},
\texttt{dtmc\_mass\_mono\_of\_forwardClosed}) from which the
incumbent monotonicity lemmas re-derive in one line, and a
kernel-checked strict-improvement witness
(\texttt{idBool\_strict\_improvement}: an unreachable class where the
uniform-residual bracket has width one and the sharpened bracket
width zero --- a standing regression tooth, re-checked every build).
The mechanized statements are deliberately universally quantified
($\forall\, t' \ge T$) rather than supremum-shaped: over $\Q$ the
supremum of a monotone bounded mass sequence need not exist as a
rational, and the universal form is the citable content --- any
eventual-absorption reading (sup or limit, in any completion of
$\Q$) inherits the claimed bracket.

One clause of the species is graded prose by design: the
\emph{truncation-transfer} reading --- stored rows agreeing with an
external program\textquotesingle s chain on rowed states implies the
bracket transfers to the program --- is a coupling argument about a
translation that lives OUTSIDE the trust path, and its assurance
instrument is a battery (vendored source, Monte-Carlo differential,
closed-form containment), not a kernel theorem. The certificate
itself never claims more than the stored rows warrant.

\cataloguenote{The schema is baked in the certificate theory as a species sibling
extending the chain grain, purely additive --- the query sort with
anchor-mediated chain binding (the query\textquotesingle s chain IS
its anchor\textquotesingle s chain, labels bound by chase rows), the
horizon field, the class and rejection labelings, and the three
side-condition Props with their report, refusal, and chase rows, all
quoted below from the theory source. The doubly-classed state is
unrepresentable at the attribute grain (the class labeling is a
keyed function row). The mechanized soundness lemmas are named in
the commentary; no Lean statement is bound into this book yet ---
the binding column is honest open debt, and joins when the
statement-binding machinery registers the development\textquotesingle s
package.
}
\ruleblock{Theory rows (geolog)}
\sourcenote{The rule\textquotesingle s certificate content as it is declared in the geolog theory the checker loads --- payload schema, side-condition demands, and axiom rows, quoted from the theory source at every build of this catalogue (a stale quotation cannot exist). Axiom modes: \texttt{chase} = the engine derives; \texttt{report} = a violation is a loud judgment; \texttt{refuse} = a violation gates admission. Side-condition relations are discharged by bounded evaluation at the engine\textquotesingle s ground seam, never by derivation.}
\par\medskip
The chain grain this species reads: states, stored transition rows,
stored initial rows, and the row-mass admission demand. The species
adds no rows here; it consumes them.

\declhead{theory Dtmc}{DState}{\badge{BadgeGreen}{quoted from the theory source at render}}
\begin{leansrc}
DState : Set;   // one state of one chain
\end{leansrc}
\declhead{theory Dtmc}{dprob}{\badge{BadgeGreen}{quoted from the theory source at render}}
\begin{leansrc}
dprob : {src: DState, dst: DState} -> Q;
\end{leansrc}
\declhead{theory Dtmc}{dinit}{\badge{BadgeGreen}{quoted from the theory source at render}}
\begin{leansrc}
dinit : {chain: Dtmc, state: DState} -> Q;
\end{leansrc}
\declhead{theory Dtmc}{qtsl\_dtmc\_rowmass}{\badge{BadgeGreen}{quoted from the theory source at render}}
\begin{leansrc}
qtsl_dtmc_rowmass  : {chain: Dtmc} -> Prop;   // side-condition
\end{leansrc}
\par\medskip
The baked rows of the species: the query sort with its anchor and
horizon, the class and rejection labelings, the disjointness
refusal, the labels-bind-own-chain chase rows, the well-formedness
demand, and the two bracket claims with their interval-order
reports.

\declhead{theory DtmcAbsorb}{DAbsQuery}{\badge{BadgeGreen}{quoted from the theory source at render}}
\begin{leansrc}
DAbsQuery : Set;
\end{leansrc}
\declhead{theory DtmcAbsorb}{DAbsQuery/*}{\badge{BadgeGreen}{quoted from the theory source at render}}
\begin{leansrc}
DAbsQuery/anchor  : DAbsQuery -> DState;
DAbsQuery/horizon : DAbsQuery -> Nat;
\end{leansrc}
\declhead{theory DtmcAbsorb}{dclass}{\badge{BadgeGreen}{quoted from the theory source at render}}
\begin{leansrc}
dclass : {q: DAbsQuery, state: DState} -> Nat;
\end{leansrc}
\declhead{theory DtmcAbsorb}{dreject}{\badge{BadgeGreen}{quoted from the theory source at render}}
\begin{leansrc}
dreject : {q: DAbsQuery, state: DState} -> Prop;
\end{leansrc}
\declhead{theory DtmcAbsorb}{ax/wf/dabs\_class\_reject\_disjoint}{\badge{BadgeGreen}{quoted from the theory source at render}}
\begin{leansrc}
refuse ax/wf/dabs_class_reject_disjoint : forall q : DAbsQuery,
    s : DState, c : Nat.
  c = {q: q, state: s} dclass,
  {q: q, state: s} dreject
  |- false;
\end{leansrc}
\declhead{theory DtmcAbsorb}{ax/wf/dabs\_class\_same\_dtmc}{\badge{BadgeGreen}{quoted from the theory source at render}}
\begin{leansrc}
chase ax/wf/dabs_class_same_dtmc : forall q : DAbsQuery, s : DState,
    c : Nat.
  c = {q: q, state: s} dclass
  |- s DState/dtmc = q DAbsQuery/anchor DState/dtmc;
\end{leansrc}
\declhead{theory DtmcAbsorb}{ax/wf/dabs\_reject\_same\_dtmc}{\badge{BadgeGreen}{quoted from the theory source at render}}
\begin{leansrc}
chase ax/wf/dabs_reject_same_dtmc : forall q : DAbsQuery, s : DState.
  {q: q, state: s} dreject
  |- s DState/dtmc = q DAbsQuery/anchor DState/dtmc;
\end{leansrc}
\declhead{theory DtmcAbsorb}{qtsl\_dabs\_wf}{\badge{BadgeGreen}{quoted from the theory source at render}}
\begin{leansrc}
qtsl_dabs_wf : {q: DAbsQuery} -> Prop;   // side-condition
\end{leansrc}
\declhead{theory DtmcAbsorb}{ax/wf/dabs\_wf}{\badge{BadgeGreen}{quoted from the theory source at render}}
\begin{leansrc}
report ax/wf/dabs_wf : forall q : DAbsQuery.
  |- {q: q} qtsl_dabs_wf;
\end{leansrc}
\declhead{theory DtmcAbsorb}{qtsl\_dabs\_absorb\_bracket}{\badge{BadgeGreen}{quoted from the theory source at render}}
\begin{leansrc}
qtsl_dabs_absorb_bracket : {q: DAbsQuery, class: Nat, lo: Q, hi: Q}
  -> Prop;   // side-condition
\end{leansrc}
\declhead{theory DtmcAbsorb}{ax/wf/dabs\_absorb\_order}{\badge{BadgeGreen}{quoted from the theory source at render}}
\begin{leansrc}
report ax/wf/dabs_absorb_order : forall q : DAbsQuery, c : Nat,
    l : Q, h : Q.
  {q: q, class: c, lo: l, hi: h} qtsl_dabs_absorb_bracket
  |- {l: l, r: h} q_le;
\end{leansrc}
\ruleblock{Lean statement}
\badge{BadgeAmber}{not yet mechanized} this rule cites no Lean statements: nothing on this page is hash-bound for it, and its dashboard row shows the same absence across the statement, receipt, and correspondence columns.
\ruleblock{HOL4 statement}
\absentbox{NO HOL STATEMENT YET}{No HOL development of the extension theory exists; this
species\textquotesingle{} content is an exact-rational evaluator
recompute over baked rows, and its soundness theorems live in the
Lean development named in the commentary.
}
\ruleblock{HOL verifier receipts}
\absentbox{NONE}{no HOL verification receipt exists for this rule.}
\ruleblock{Proof}

\sourcenote{Genre: certificate species (tier 1 semantic soundness;
checker recomputation, no proof node). The proof below is the
mechanized argument restated at textbook detail; the named Lean
declarations carry it, sorry-free at standard axiom closure.}

\textbf{Setting.}\enspace A finite stored-row chain: state set $X$,
stored transition rows $K_0 \subseteq X \times X \to \Q_{\ge 0}$
satisfying the row-mass admission demand (each rowed source's stored
masses total exactly one), stored initial measure $m_0$, horizon
$T \in \N$, class labeling and rejection labeling over $X$. Let $K$
be $K_0$ completed by a unit self-loop at every state with no stored
outgoing row (\emph{rowless-is-absorbing}); $K$ is then a stochastic
matrix on $X$, and $m_{t+1} := m_t K$ defines the exact-rational
evolution. Write $C$ for the set of class-$c$ states,
$\mathrm{Abs}$ for the labeled (classed or rejected) states --- each
demanded stored-absorbing by the well-formedness side condition:
rowless, or exactly the unit self-loop --- and
$N_c := \mathrm{mass}_T(C)$, $S := \sum_{c'} N_{c'}$,
$R := 1 - S - (\text{rejected mass of } m_T)$.

\textbf{Lemma 1 (monotone absorption --- the lower side).}\enspace
For $t' \ge T$: $\mathrm{mass}_{t'}(C) \ge \mathrm{mass}_T(C)$.
\emph{Proof.} Each state of $C$ is stored-absorbing, so under $K$ its
mass never leaves; class mass is nondecreasing step by step.
(Removing the well-formedness hypothesis admits the counterexample
the mutant battery pins: a classed state with a stored mass-leaking
outgoing row loses class mass at the very next step.)\hfill$\square$

\textbf{Lemma 2 (forward-closed complement --- the upper side).}\enspace
Let $B_C$ be the co-reachable set of $C$: states admitting a path
into $C$ through strictly positive stored rows ($C \subseteq B_C$).
Then $X \setminus B_C$ is forward-closed --- no positive-support row
of $K$ leads from outside $B_C$ into $B_C$ (such a row would itself
witness co-reachability of its source) --- hence
$\mathrm{mass}_{t'}(B_C)$ is nonincreasing in $t'$.
\emph{Proof.} For a forward-closed set $F$, $\mathbf{1}_F K^\top
\ge \mathbf{1}_F$ pointwise on $F$'s complement contributes nothing,
and mass on $F$ can only stay or grow; apply to
$F = X \setminus B_C$ and complement. This is the mechanized
generalization (\texttt{FinDTMC.ForwardClosed},
\texttt{dtmc\_mass\_mono\_of\_forwardClosed}), from which Lemma~1
re-derives in one line: a stored-absorbing set is forward-closed
around each of its points.\hfill$\square$

\textbf{Theorem (the bracket).}\enspace Define
$R_c := \mathrm{mass}_T(B_C \setminus C)$ (the
reachability-sharpened per-class residual, recomputed by a backward
walk over strictly positive stored rows, seeded at $C$, under the
same work budget --- exhaustion refuses the whole claim). For every
$t' \ge T$:
$$\mathrm{mass}_{t'}(C) \in [\,\mathrm{mass}_T(C),\
\mathrm{mass}_T(C) + R_c\,].$$
\emph{Proof.} Lower: Lemma 1. Upper: class mass at $t'$ sits inside
$B_C$ (mass outside $B_C$ can never arrive in $C$, by Lemma~2
applied along the trajectory), and $C$'s gain since $T$ is at most
$B_C$'s non-$C$ mass at $T$, which is $R_c$, since $B_C$'s total is
nonincreasing while its $C$-part is nondecreasing.\hfill$\square$

The mechanized upper side (\texttt{dtmcAbsorb\_mass\_le\_reach}) is
HYPOTHESIS-FREE on the class: it holds for any class of any finite
stochastic chain --- the well-formedness demands protect only
Lemma~1. The never-loosens warrant
(\texttt{dtmcAbsorb\_reach\_residual\_le}) gives $R_c \le
1 - \mathrm{mass}_T(\mathrm{Abs})$, needing only that labeled states
absorb, so every claim admissible under the uniform residual remains
admissible under the sharpened one; and the strict-improvement tooth
(\texttt{idBool\_strict\_improvement}) exhibits an unreachable class
where the uniform bracket has width one and the sharpened bracket
width zero --- a standing kernel-checked regression witness. A class
unreachable from every alive state has $R_c = 0$: a width-zero
bracket at any $T \ge 1$, with no termination assumption anywhere.

\textbf{Corollary (eventual readings).}\enspace The statements are
universally quantified over $t' \ge T$ deliberately: over $\Q$ the
supremum of the monotone bounded sequence $\mathrm{mass}_{t'}(C)$
need not exist as a rational. Any eventual-absorption reading ---
supremum or limit, in any completion of $\Q$ --- inherits the
claimed bracket from the universal form
(\texttt{dtmcAbsorb\_eventual\_bracket\_reach}).

\textbf{Why each hypothesis.}\enspace Row-mass admission: without
it, doubled stored masses inflate $N_c$ (a battery locus). The
well-formedness demand: Lemma 1's counterexample above. The budget
refusal on the backward walk: a partial walk could under-build
$B_C$, dropping states whose mass belongs in $R_c$ --- so exhaustion
refuses rather than verdicts. Interval order $lo \le hi$ is a report
row; inversion refuses. And the completion convention
(rowless-is-absorbing) makes truncation \emph{widen} brackets:
omitted rows leave their states absorbing and their mass in the
residual, never in a class --- there is no closure side condition to
attack, and claims may be looser than canonical, never tighter.

\begin{bremark}[status]
The schema is baked (species rows quoted in this section); the
argument above is mechanized sorry-free at standard axiom closure,
axiom-audited on every build, in the module named in the commentary.
No Lean statement is bound into this book yet: the binding column is
honest open debt until the statement-binding machinery registers the
development's package.
\end{bremark}


\rulesection{95}{DtmcPosteriorBracket}{core rule}
\label{blcatrule-DtmcPosteriorBracket}
\ruleblock{The rule (Steele metanotation)}
\sourcenote{as realized in the mechanized development; source row: \texttt{\texttt{QTSL-rules.toml}} --- \emph{DtmcPosteriorBracket}. The metanotation statement is typeset by hand for notational quality; this rule cites no Lean statements yet (the absence is stated loudly below). The rule \emph{roster} is machine-checked (every mechanized rule tag round-trips to exactly one catalogue row); the statement \emph{text} is editorial.}
\begin{mathpar}
\inferrule*{
  q : \mathsf{DAbsQuery} \\
  \mathsf{qtsl\_dabs\_wf}(q) \\
  \text{stored rows of } q.\mathsf{chain}\ (\mathsf{dprob}, \mathsf{dinit}, \mathsf{dclass}, \mathsf{dreject}) \\
  S > 0\ \text{(recomputed; refusal otherwise)}
}{
  \mathsf{qtsl\_dabs\_posterior\_bracket}(q,\ c,\ lo,\ hi)
}
\end{mathpar}

\par\medskip
\rulename{DtmcPosteriorBracket} is the normalized face of
\rulename{DtmcAbsorbBracket}: a bracket on the \emph{posterior}
eventual-absorption measure of one class, with normalizer the total
classed mass. Hard conditioning is a CLASS, not a score: rejected
mass (failed conditioning) and unlabeled or nonterminating mass are
excluded from the normalizer by the same labeling mechanism, so every
stored row remains a probability row and the chain grain\textquotesingle s
row-mass admission demand is untouched. This is what lets posteriors
of discrete probabilistic programs ride the species without any
weighted-row machinery.

The claim holds iff $lo \le N_c/(S+R)$ and $(N_c+R)/(S+R) \le hi$,
with $N_c$, $S$, $R$ recomputed by the same $T$-step exact-rational
evaluation as the sibling; it REFUSES at $S = 0$ (no
positive-evidence witness). The mathematical content beyond the
sibling\textquotesingle s conservation bracket is the correlated-ratio
lemma --- pure algebra, no chain in sight: from $N_c \le M_c$
classwise, $\sum_c (M_c - N_c) \le R$, and $S > 0$, the true
posterior lies in $[\,N_c/(S+R),\ (N_c+R)/(S+R)\,]$. The correlation
is the point: numerator and denominator draw their errors from ONE
budget, so both endpoints use the single normalizer $S + R$ rather
than the looser interval quotient --- never worse, often sharper
(measured $3.0\times$ tighter than the interval-quotient baseline at
matched truncation on the species\textquotesingle{} first worked
example). Mechanized as \texttt{dtmcPosteriorBracket\_sound}, over
any finite class alphabet, in the same module and with the same
receipts as the sibling.

An honesty note carried from the design: the bracket width is exactly
$R/(S+R)$, class-uniform --- the tightness of the whole judgment is
one number, which is what keeps comparison-table claim shapes honest.
The posterior face deliberately keeps the UNIFORM residual $R$ even
though the unnormalized sibling mints with the reachability-sharpened
$R_c$: the single-normalizer argument requires one shared budget
across all classes, and the class-uniform width claim stays honest
only under it. An asymmetric refinement (uniform lower endpoint,
class-specific upper endpoint $(N_c+R_c)/(S+R_c)$) is sound and is
deliberately NOT minted --- it boards a future revision of the
theory if and when a use case demands it.

\cataloguenote{The schema is baked in the certificate theory alongside the sibling:
same query sort, labelings, and well-formedness demand; its own claim
relation with its own interval-order report; the $S > 0$ refusal is
evaluator-side (a recompute refusal, not a stored row). Quoted below
from the theory source. No Lean statement is bound into this book
yet --- the binding column is honest open debt, as the sibling.
}
\ruleblock{Theory rows (geolog)}
\sourcenote{The rule\textquotesingle s certificate content as it is declared in the geolog theory the checker loads --- payload schema, side-condition demands, and axiom rows, quoted from the theory source at every build of this catalogue (a stale quotation cannot exist). Axiom modes: \texttt{chase} = the engine derives; \texttt{report} = a violation is a loud judgment; \texttt{refuse} = a violation gates admission. Side-condition relations are discharged by bounded evaluation at the engine\textquotesingle s ground seam, never by derivation.}
\par\medskip
The baked rows beyond the sibling\textquotesingle s: the posterior
claim relation and its interval-order report. The query sort,
labelings, and well-formedness demand are the sibling\textquotesingle s
own rows, quoted in its section.

\declhead{theory DtmcAbsorb}{qtsl\_dabs\_posterior\_bracket}{\badge{BadgeGreen}{quoted from the theory source at render}}
\begin{leansrc}
qtsl_dabs_posterior_bracket : {q: DAbsQuery, class: Nat, lo: Q,
                               hi: Q} -> Prop;   // side-condition
\end{leansrc}
\declhead{theory DtmcAbsorb}{ax/wf/dabs\_posterior\_order}{\badge{BadgeGreen}{quoted from the theory source at render}}
\begin{leansrc}
report ax/wf/dabs_posterior_order : forall q : DAbsQuery, c : Nat,
    l : Q, h : Q.
  {q: q, class: c, lo: l, hi: h} qtsl_dabs_posterior_bracket
  |- {l: l, r: h} q_le;
\end{leansrc}
\ruleblock{Lean statement}
\badge{BadgeAmber}{not yet mechanized} this rule cites no Lean statements: nothing on this page is hash-bound for it, and its dashboard row shows the same absence across the statement, receipt, and correspondence columns.
\ruleblock{HOL4 statement}
\absentbox{NO HOL STATEMENT YET}{No HOL development of the extension theory exists; this
species\textquotesingle{} content is an exact-rational evaluator
recompute over baked rows plus one algebraic lemma, mechanized in
the Lean development named in the commentary.
}
\ruleblock{HOL verifier receipts}
\absentbox{NONE}{no HOL verification receipt exists for this rule.}
\ruleblock{Proof}

\sourcenote{Genre: certificate species (tier 1 semantic soundness;
checker recomputation, no proof node). The content beyond the
unnormalized sibling is one algebraic lemma, proved here in full;
the mechanized form is \texttt{dtmcPosteriorBracket\_sound}, over
any finite class alphabet.}

\textbf{Setting.}\enspace As in \rulename{DtmcAbsorbBracket}: from
the $T$-step evaluation, per-class truncation masses $N_c \ge 0$,
their total $S = \sum_c N_c$, and the residual $R \ge 0$. Write
$M_c$ for the true eventual class-$c$ mass (any completion-valued
reading). The sibling's conservation bracket gives, classwise,
$N_c \le M_c$ and $\sum_c (M_c - N_c) \le R$.

\textbf{Lemma (correlated ratio).}\enspace If additionally $S > 0$,
then for every class $c$:
$$\frac{M_c}{\sum_{c'} M_{c'}} \;\in\;
\Bigl[\,\frac{N_c}{S+R},\ \frac{N_c+R}{S+R}\,\Bigr].$$
\emph{Proof.} Let $\delta_{c'} := M_{c'} - N_{c'} \ge 0$ and
$\Delta := \sum_{c'} \delta_{c'} \le R$. The denominator is
$\sum_{c'} M_{c'} = S + \Delta \in [S,\ S+R]$, and $S > 0$ makes
every quotient well-defined. Upper endpoint:
$$\frac{M_c}{S+\Delta} = \frac{N_c + \delta_c}{S + \Delta}
\le \frac{N_c + \Delta}{S + \Delta} \le \frac{N_c + R}{S + R},$$
the first step since $\delta_c \le \Delta$, the second since
$x \mapsto (N_c + x)/(S + x)$ is nondecreasing for $x \ge 0$ (as
$N_c \le S$) and $\Delta \le R$. Lower endpoint:
$$\frac{M_c}{S+\Delta} \ge \frac{N_c}{S+\Delta} \ge
\frac{N_c}{S+R}.$$
\hfill$\square$

The correlation is the content: numerator and denominator draw their
errors from the ONE budget $\Delta \le R$, so both endpoints use the
single normalizer $S + R$. The uncorrelated alternative --- bracket
numerator and denominator separately and take the interval quotient
$[\,N_c/(S+R),\ (N_c+R)/S\,]$ --- is never better and usually worse;
at matched truncation on the species' first worked example the
correlated form measured $3.0\times$ tighter.

\textbf{Why each hypothesis.}\enspace $S > 0$: at $S = 0$ the
posterior is $0/0$-shaped and the species REFUSES (no
positive-evidence witness) --- a refusal, never a verdict.
Nonnegativity of the $\delta_{c'}$ is exactly the sibling's monotone
lower side; the one-budget bound $\Delta \le R$ is its conservation
upper side, summed. Hard conditioning enters as a CLASS (rejected
mass leaves the normalizer by labeling), so every stored row remains
a probability row and no weighted-row machinery is consumed.

\textbf{Width honesty.}\enspace Both endpoints differ by exactly
$R/(S+R)$, class-uniformly: the tightness of the whole judgment is
one number. The posterior face deliberately keeps the uniform $R$
(the one-budget argument above requires a single normalizer shared
by all classes); the asymmetric per-class refinement with upper
endpoint $(N_c+R_c)/(S+R_c)$ is sound and deliberately not part of
this revision --- it boards a future one if a use case demands it.

\begin{bremark}[status]
The schema is baked (claim relation and interval-order report quoted
in this section; query sort and labelings in the sibling's). The
lemma above is mechanized sorry-free at standard axiom closure over
any finite class alphabet. No Lean statement is bound into this book
yet: the binding column is honest open debt, as the sibling's.
\end{bremark}


\rulesection{96}{QGenFormContraction}{core rule}
\label{blcatrule-QGenFormContraction}
\ruleblock{The rule (Steele metanotation)}
\sourcenote{as realized in the mechanized development; source row: \texttt{\texttt{QTSL-rules.toml}} --- \emph{QGenFormContraction}. The metanotation statement is typeset by hand for notational quality; this rule cites no Lean statements yet (the absence is stated loudly below). The rule \emph{roster} is machine-checked (every mechanized rule tag round-trips to exactly one catalogue row); the statement \emph{text} is editorial.}
\begin{mathpar}
\inferrule*{
  w : \mathsf{DGenContraction} \\
  \text{stored rows: } \mathsf{dweight}\ (A = \mathrm{supp}\,V,\ \text{nonempty}),\ \mathsf{dprob},\ \mathsf{dinit} \\
  \textstyle\sum_{y \in A} K(x,y)\,V(y) \le c\,V(x)\ \ \forall x \in A \\
  K(x,y) = 0\ \text{for stored rows } x \notin A,\ y \in A \\
  \min_A V > 0 \\
  \mathsf{qtsl\_dtmc\_genform\_contraction}(w)
}{
  \text{each claimed row of } w:\ \mathsf{qtsl\_dgc\_partial\_lo},\ \mathsf{qtsl\_dgc\_ev\_bracket},\ \mathsf{qtsl\_dgc\_survival\_tail},\ \mathsf{qtsl\_dgc\_z\_one}
}
\end{mathpar}

\par\medskip
\rulename{QGenFormContraction} strengthens the chain-content
generator-form seam from the incumbent\textquotesingle s boolean
claim --- transfer of the certificate is nonpositive-drift --- to a
score-carrying, regional, \emph{quantitative} contraction whose
conclusions include the incumbent\textquotesingle s and add geometric
tails, first-moment tails, and a termination-mass-one claim. It is a
certificate species, not a proof-node rule: checker recomputation
over stored rows, zero new judgment sorts, zero new formers. The
mechanized carrier is an arbitrary nonnegative weight matrix (the
score factorization $s(x)K(x,y)$ is presentation, not mathematics),
so stochastic, sub-probability, and soft-score row sets ride ONE
induction.

This is the fold-don\textquotesingle t-sibling reading of the change
discipline enacted: the would-be separate weighted-tail species folds
into the incumbent seam, and the incumbent survives as the
degenerate face (score $\equiv 1$, $c = 1$, region = all states) ---
its rows, wire shape, and consuming lemmas untouched. Subsumption is
mechanized twice over, as the strengthening discipline demands: the
instance chain (the weighted core at $s \equiv 1$ on a stochastic
kernel re-derives the previously proven stochastic core\textquotesingle s
exact statement:
\texttt{massAt\_expect\_le\_of\_contraction\_of\_weighted}) and the
composed theorem
(\texttt{qGenFormContraction\_subsumes\_genform}: from the
incumbent\textquotesingle s recomputed nonpositive generator form,
the strengthened engine derives the incumbent
consumption\textquotesingle s exact path-law conclusion, never by
citing the incumbent\textquotesingle s own proof) --- so every
certificate admitted under the incumbent discharges under the
strengthening with an unchanged conclusion.

The evaluator recomputes, per stored row: the region-restricted
weighted contraction $\sum_{y \in A} K(x,y)V(y) \le c\,V(x)$ on the
region $A$; no re-entry ($K(x,y) = 0$ for stored rows from outside
$A$ into $A$ --- equivalently, the complement is closed); and the
strictly positive region minimum $\mu = \min_A V > 0$. Conclusions,
each bound exactly where its claim row is spent: weighted alive-mass
decay $\langle m_t|_A, V\rangle \le c^t \langle m_0|_A, V\rangle$;
for $c < 1$, window exit-weight truncation and the time-weighted
first-moment tail (the expected-value rows\textquotesingle{}
engine); the exact partial survival sum (the lower claim); and ---
through the DERIVED conversion $1/\mu$ rescaling the mechanized
unit-floor lemma by scale-invariance --- the bare alive-mass tail
$\mathrm{mass}_t(A) \le c^t \langle m_0|_A, V\rangle/\mu$, the
stored-row form of a hitting-time tail
$\Pr[\text{hitting time} > t] \le M c^t$ (a self-stabilization
time\textquotesingle s region is the not-yet-stabilized
configurations, whose complement\textquotesingle s closedness is
exactly the no-re-entry recompute).

The enacted schema\textquotesingle s deviations from the drafted
proposal are each load-bearing. There is NO stored exit or
conversion constant: the constant is derived ($1/\min_A V$, with the
minimum demanded strictly positive) because a stored constant would
be a forgeable derived quantity --- the species\textquotesingle{}
own never-store-what-you-can-recompute discipline. The region IS the
weight relation\textquotesingle s support (a stored weight row is
the stored membership; no separate region relation to fall out of
sync). Four NAMED claim relations replace a flavor datum --- partial
lower bound, expected-exit-time bracket, survival tail, and
termination-mass-one --- making the wire self-documenting and
letting each side condition bind exactly where its conclusion is
spent (the $c < 1$ demand binds exactly at tail-claiming rows; at
$c = 1$ the incumbent decay-only conclusion remains). A PRESENT
score field REFUSES at this revision (deferred-safe-fail): the
mechanized core is already score-carrying, and the evaluator face
for present scores boards its own future revision rather than
shipping unexercised. Higher moments ($k \ge 2$) are deliberately
absent --- the same per-term reduction with $\sum t^k c^t$ closed
forms, filed to the first leg that needs them.

The soundness engine is the weighted-transfer induction (one-step
transfer, induction on $t$, geometric series --- no stochastic
analysis anywhere), mechanized sorry-free at standard axiom closure
and axiom-audited on every build
(\texttt{SuFunctor/Adequacy/DtmcContraction.lean}): \texttt{wmassAt}
with the instance tie \texttt{massAt\_eq\_wmassAt}; the core
\texttt{wmassAt\_expect\_le\_of\_contraction} (certificate
nonnegativity DROPPED --- strictly stronger than the drafted
specification); the score face; the regional packaging via the
region-supported observable; exit domination; the exact
arithmetico-geometric identity and series tails; and the two tail
assemblies, in universally-quantified partial-sum form (over $\Q$
the limit need not exist; any completion-valued reading inherits the
bound). The alive-mass face lives on its own module
(\texttt{SuFunctor/Adequacy/DtmcAliveTail.lean}):
\texttt{massAt\_expect\_on\_le\_of\_regionContraction} and
\texttt{dtmcAliveTail\_of\_regionContraction}, whose two added
hypotheses --- the floor and the closed complement --- are each
witnessed load-bearing by kernel-checked countermodels COMMITTED
beside the lemmas (\texttt{aliveTail\_needs\_floor},
\texttt{aliveTail\_needs\_no\_reentry}), re-checked on every build.
The species\textquotesingle{} mutant battery was run as committed
throwaway edits, each red at its predicted load-bearing proof step.

Two further faces are deliberately NOT part of this revision. The
TEMPLATE face (geometric certificate families $V(n) = M\gamma^n$
over offset families, collapsing region instances to single rational
inequalities) and the PATTERN face (finitely-presented infinite
chains and product-profile templates) are deferred to their own
design work: the pattern face is underdetermined at this grain (no
canonical state-coordinate binding; a pattern\textquotesingle s
generated kernel and the stored completion\textquotesingle s
rowless-is-absorbing reading are two different kernels and must not
share one chain reading; and the finite-state induction does not by
itself warrant infinite-sum exchange --- a carrier theorem is owed),
and no admitted certificate demands the one-dimensional template at
the stored-weights grain, where the finite face stores every weight.
Their collapse warrants are already mechanized and banked for that
future work (\texttt{genFormTemplate\_collapse} for one-dimensional
offset families; \texttt{aliveStep\_factor},
\texttt{transferRatio\_of\_ge}, \texttt{productProfile\_collapse}
and its window form,
\texttt{SuFunctor/Adequacy/DtmcProductProfile.lean}, proven at full
profile generality), together with a kernel-checked necessity
witness (\texttt{rampProfile\_needs\_window\_reach}) showing the
class-check family alone cannot cover ramp-resident exterior states
--- so the face\textquotesingle s eventual schema must carry a
profile-length demand beside window coverage. What remains for that
face is the infinite-carrier bridge, the kernel-object tie, and the
stored/template composition glue; until then, template- and
pattern-shaped claims live outside the certified lane, labeled as
such.

\cataloguenote{The FINITE FACE is baked in the certificate theory as a species
sibling extending the chain grain, purely additive: the witness
entity (anchor-mediated chain binding, absolute horizon, the
deterministic-prologue offset datum, the contraction constant, the
optional score reference whose presence refuses at this revision),
the keyed weight relation whose support is the region (nonempty
demanded), the engine-decided demand minted per witness by a report
row, and the four named claim relations with their range and
interval-order reports --- all quoted below from the theory source.
The incumbent generator-form rows are untouched (additive delta
only), and the sibling species\textquotesingle{} deliberately-absent
notes carry their pointer flips from the same revision. The
mechanized soundness lemmas are named in the commentary; no Lean
statement is bound into this book yet --- the binding column is
honest open debt.
}
\ruleblock{Theory rows (geolog)}
\sourcenote{The rule\textquotesingle s certificate content as it is declared in the geolog theory the checker loads --- payload schema, side-condition demands, and axiom rows, quoted from the theory source at every build of this catalogue (a stale quotation cannot exist). Axiom modes: \texttt{chase} = the engine derives; \texttt{report} = a violation is a loud judgment; \texttt{refuse} = a violation gates admission. Side-condition relations are discharged by bounded evaluation at the engine\textquotesingle s ground seam, never by derivation.}
\par\medskip
The incumbent seam this species strengthens --- and survives as the
degenerate face of: the chain-content generator-form demand.

\declhead{theory Dtmc}{qtsl\_dtmc\_genform}{\badge{BadgeGreen}{quoted from the theory source at render}}
\begin{leansrc}
qtsl_dtmc_genform : {sys: System, cert: HypernetRef,
                     gen: HypernetRef} -> Prop;
\end{leansrc}
\par\medskip
The baked rows of the finite face: the witness entity, the weight
relation with its same-chain chase row, the engine-decided demand,
and the four named claim relations with their wf reports (constant
range, interval order on expected-value brackets, tail-coefficient
ranges, and the strict-contraction demand exactly where an
expected-value bracket is claimed).

\declhead{theory DtmcContraction}{DGenContraction}{\badge{BadgeGreen}{quoted from the theory source at render}}
\begin{leansrc}
DGenContraction : Set;
\end{leansrc}
\declhead{theory DtmcContraction}{DGenContraction/*}{\badge{BadgeGreen}{quoted from the theory source at render}}
\begin{leansrc}
DGenContraction/anchor  : DGenContraction -> DState;
DGenContraction/horizon : DGenContraction -> Nat;
DGenContraction/offset  : DGenContraction -> Nat;
DGenContraction/c       : DGenContraction -> Q;
DGenContraction/score   : DGenContraction -> HypernetRef;
\end{leansrc}
\declhead{theory DtmcContraction}{dweight}{\badge{BadgeGreen}{quoted from the theory source at render}}
\begin{leansrc}
dweight : {cw: DGenContraction, state: DState} -> Q;
\end{leansrc}
\declhead{theory DtmcContraction}{ax/wf/dgc\_weight\_same\_dtmc}{\badge{BadgeGreen}{quoted from the theory source at render}}
\begin{leansrc}
chase ax/wf/dgc_weight_same_dtmc : forall w : DGenContraction,
    s : DState, v : Q.
  v = {cw: w, state: s} dweight
  |- s DState/dtmc = w DGenContraction/anchor DState/dtmc;
\end{leansrc}
\declhead{theory DtmcContraction}{qtsl\_dtmc\_genform\_contraction}{\badge{BadgeGreen}{quoted from the theory source at render}}
\begin{leansrc}
qtsl_dtmc_genform_contraction : {cw: DGenContraction} -> Prop;
\end{leansrc}
\declhead{theory DtmcContraction}{ax/wf/dgc\_demand}{\badge{BadgeGreen}{quoted from the theory source at render}}
\begin{leansrc}
report ax/wf/dgc_demand : forall w : DGenContraction.
  |- {cw: w} qtsl_dtmc_genform_contraction;
\end{leansrc}
\declhead{theory DtmcContraction}{qtsl\_dgc\_partial\_lo}{\badge{BadgeGreen}{quoted from the theory source at render}}
\begin{leansrc}
qtsl_dgc_partial_lo : {cw: DGenContraction, lo: Q} -> Prop;
\end{leansrc}
\declhead{theory DtmcContraction}{qtsl\_dgc\_ev\_bracket}{\badge{BadgeGreen}{quoted from the theory source at render}}
\begin{leansrc}
qtsl_dgc_ev_bracket : {cw: DGenContraction, lo: Q, hi: Q} -> Prop;
\end{leansrc}
\declhead{theory DtmcContraction}{qtsl\_dgc\_survival\_tail}{\badge{BadgeGreen}{quoted from the theory source at render}}
\begin{leansrc}
qtsl_dgc_survival_tail : {cw: DGenContraction, coeff: Q, rate: Q}
  -> Prop;   // side-condition (R21)
\end{leansrc}
\declhead{theory DtmcContraction}{qtsl\_dgc\_z\_one}{\badge{BadgeGreen}{quoted from the theory source at render}}
\begin{leansrc}
qtsl_dgc_z_one : {cw: DGenContraction} -> Prop;
\end{leansrc}
\declhead{theory DtmcContraction}{ax/wf/dgc\_c\_range}{\badge{BadgeGreen}{quoted from the theory source at render}}
\begin{leansrc}
report ax/wf/dgc_c_range : forall w : DGenContraction, c : Q.
  c = w DGenContraction/c
  |- {l: 0, r: c} q_le, {l: c, r: 1} q_le;
\end{leansrc}
\declhead{theory DtmcContraction}{ax/wf/dgc\_ev\_order}{\badge{BadgeGreen}{quoted from the theory source at render}}
\begin{leansrc}
report ax/wf/dgc_ev_order : forall w : DGenContraction, l : Q, h : Q.
  {cw: w, lo: l, hi: h} qtsl_dgc_ev_bracket
  |- {l: l, r: h} q_le;
\end{leansrc}
\declhead{theory DtmcContraction}{ax/wf/dgc\_tail\_range}{\badge{BadgeGreen}{quoted from the theory source at render}}
\begin{leansrc}
report ax/wf/dgc_tail_range : forall w : DGenContraction, b : Q,
    r : Q.
  {cw: w, coeff: b, rate: r} qtsl_dgc_survival_tail
  |- {l: 0, r: b} q_le, {l: 0, r: r} q_le, {l: r, r: 1} q_lt;
\end{leansrc}
\declhead{theory DtmcContraction}{ax/wf/dgc\_ev\_strict}{\badge{BadgeGreen}{quoted from the theory source at render}}
\begin{leansrc}
report ax/wf/dgc_ev_strict : forall w : DGenContraction, l : Q,
    h : Q, c : Q.
  {cw: w, lo: l, hi: h} qtsl_dgc_ev_bracket,
  c = w DGenContraction/c
  |- {l: c, r: 1} q_lt;
\end{leansrc}
\ruleblock{Lean statement}
\badge{BadgeAmber}{not yet mechanized} this rule cites no Lean statements: nothing on this page is hash-bound for it, and its dashboard row shows the same absence across the statement, receipt, and correspondence columns.
\ruleblock{HOL4 statement}
\absentbox{NO HOL STATEMENT YET}{No HOL development of the extension theory exists; this
species\textquotesingle{} content is an exact-rational evaluator
recompute over baked rows, and its soundness theorems live in the
Lean development named in the commentary.
}
\ruleblock{HOL verifier receipts}
\absentbox{NONE}{no HOL verification receipt exists for this rule.}
\ruleblock{Proof}

\sourcenote{Genre: certificate species (tier 1 semantic soundness;
checker recomputation strengthening an existing seam). The proof is
the weighted-transfer induction plus exact series arithmetic --- no
stochastic analysis anywhere; the mechanized lemma roster is named
inline.}

\textbf{Setting.}\enspace A finite stored-row chain with completed
kernel $K$ (rowless-is-absorbing, as in the sibling species); a
stored weight function $V$ on the region $A := \mathrm{supp}\,V$
(the region IS the support of the stored weight rows --- nonempty
demanded); a contraction constant $c \in [0,1]$; and the recomputed
side conditions: regional contraction
$\sum_{y \in A} K(x,y)\,V(y) \le c\,V(x)$ for every $x \in A$;
no re-entry, $K(x,y) = 0$ for stored rows with $x \notin A$,
$y \in A$ (equivalently, the complement of $A$ is closed); and
strictly positive region minimum $\mu := \min_A V > 0$. The carrier
of the mechanized core is an arbitrary nonnegative weight matrix ---
the score factorization $s(x)K(x,y)$ is presentation, not
mathematics --- so stochastic, sub-probability, and soft-score row
sets ride one induction; certificate nonnegativity is NOT assumed by
the core (\texttt{wmassAt\_expect\_le\_of\_contraction}), which is
strictly stronger than the drafted specification.

\textbf{Lemma 1 (one-step transfer).}\enspace With $w_t :=
\langle m_t|_A,\ V\rangle$ the region-supported weighted mass:
$w_{t+1} \le c\, w_t$. \emph{Proof.} Expand $m_{t+1} = m_t K$;
mass entering $A$ from outside contributes nothing (no re-entry),
and each $x \in A$ contributes $m_t(x) \sum_{y \in A} K(x,y)V(y)
\le c\, m_t(x) V(x)$ by the recomputed contraction, using
nonnegativity of kernel entries and of the evolving mass (the two
mutant loci at this step). Sum over $x \in A$.\hfill$\square$

\textbf{Theorem (decay, tails, termination).}\enspace By induction,
$w_t \le c^t w_0$ (weighted alive-mass decay --- at $c = 1$ this is
exactly the incumbent seam's nonpositive-drift conclusion). Per-step
exit weight is dominated through the region values (exit
domination), so for $c < 1$ and any window $[T, T')$, the geometric
series gives the truncation bound
$\kappa\, c^T \langle m_0, V\rangle/(1-c)$ and --- via the exact
arithmetico-geometric identity
(\texttt{sum\_range\_mul\_pow\_mul\_eq}) and its tail estimates
(\texttt{sum\_Ico\_pow\_le}, \texttt{sum\_Ico\_mul\_pow\_le}) --- the
time-weighted first-moment tail
$\kappa\, \langle m_0, V\rangle\, c^T (T(1-c)+c)/(1-c)^2$, both in
universally-quantified partial-sum form: over $\Q$ the limit need
not exist, and any completion-valued normalizer, expected-value
numerator, or moment row inherits the bound
(\texttt{wmassAt\_exit\_tail\_le},
\texttt{wmassAt\_exit\_moment\_tail\_le}). Dropping $c < 1$ from a
tail claim is unsound, not vacuous --- at $c = 1$ the rational
division convention would make the assembled bound false (the
battery's witness) --- which is why the strict-contraction demand
binds exactly where a tail-shaped claim row is spent.

\textbf{Corollary (bare alive-mass tail --- the derived
conversion).}\enspace Where $\mu = \min_A V > 0$, apply the
unit-floor lemma to the rescaled certificate $V/\mu$ (contraction
and no-re-entry are scale-invariant; the floor $V/\mu \ge 1$ on $A$
holds by definition of $\mu$):
$$\mathrm{mass}_t(A) \;\le\; c^t\, \langle m_0|_A,\ V\rangle/\mu$$
(\texttt{massAt\_expect\_on\_le\_of\_regionContraction},
\texttt{dtmcAliveTail\_of\_regionContraction}) --- the stored-row
form of the hitting-time tail $\Pr[\text{hitting time} > t] \le
M c^t$. The conversion constant $1/\mu$ is DERIVED, never stored: a
stored constant would be a forgeable derived quantity. Both added
hypotheses are witnessed load-bearing by kernel-checked
countermodels committed beside the lemmas and re-checked every
build: without the floor, one state at $V = 1/2$, $c = 1$ makes the
bound read $1 \le 1/2$ (\texttt{aliveTail\_needs\_floor}); without
the closed complement, two states with the exterior feeding the
region at probability one make it read $1 \le 0$ at $t = 1$
(\texttt{aliveTail\_needs\_no\_reentry}). With stochastic rows at
every region state, an initial probability measure, $0 \le c < 1$,
no re-entry, and the positive minimum, weighted decay converts to
bare survival decay and forces termination mass exactly one --- the
fourth claim relation; a sign-free weighted contraction alone would
NOT suffice (the negative-weight self-loop countermodel is a battery
tooth).

\textbf{Subsumption (the strengthening's extra obligation).}\enspace
Mechanized twice over, as a strengthening of a landed seam must be:
the instance chain --- the weighted core at $s \equiv 1$ on a
stochastic kernel re-derives the previously proven stochastic core's
exact statement
(\texttt{massAt\_expect\_le\_of\_contraction\_of\_weighted}) --- and
the composed theorem
(\texttt{qGenFormContraction\_subsumes\_genform}): from the
incumbent's recomputed nonpositive generator form, the strengthened
engine at the degenerate face (score $\equiv 1$, $c = 1$, region =
all states) derives the incumbent consumption's exact path-law
conclusion through the duality bridge, never by citing the
incumbent's own proof. Every certificate admitted under the
incumbent therefore discharges under the strengthening with an
unchanged conclusion; the incumbent's rows, wire shape, and
consuming lemmas are untouched.

\textbf{Prologue offsets.}\enspace The offset datum $k$ (absent = 0)
excludes a deterministic prologue from the claim window: partial
survival sums index from $k$, and where a partial or expected-value
claim spends it, the well-formedness demand recomputes prologue
determinism --- killed-on-exit evolution from the full stored
initial measure with total surviving mass exactly one at step
$k - 1$ ($k = 0$ vacuous). An unearned offset would shift the tail
window onto mass the induction never controlled; the recompute
refuses it.

\begin{bremark}[status]
The finite face's schema is baked (entity, weight relation, demand,
and the four named claim relations quoted in this section); the
whole species' soundness, the subsumption pair, and the alive-mass
face are mechanized sorry-free at standard axiom closure,
axiom-audited on every build, with the species' mutant battery run
as committed throwaway edits, each red at its predicted load-bearing
step. No Lean statement is bound into this book yet: the binding
column is honest open debt. The template and pattern faces are NOT
part of this revision --- their collapse warrants and the necessity
witness for the eventual profile-length demand are mechanized and
banked, and the faces board the revision that follows their own
design work; until then template- and pattern-shaped claims live
outside the certified lane, labeled as such.
\end{bremark}


\blcatsection{The Soundness Theorem}
{\itshape Its content is three conformance obligations on derivation edges, and together they are what turns the page\textquotesingle s per-rule \emph{relative} soundness lemmas --- each proved on the hypothesis that the base calculus is sound --- into a single verdict about a whole certificate, unconditional once those per-rule lemmas, base soundness, and the loader\textquotesingle s conformance passes are supplied.}

\rulesection{97}{ProofWellFormedness}{edge discipline (glue)}
\label{blcatrule-ProofWellFormedness}
\ruleflavor{baked as \emph{typing}, not as a node sort: E1 — derivation premises are judgment kinds XOR inherited $\mathsf{Proof}$, never \texttt{QProof}, never a rule sort (the baked proof-node sorts\textquotesingle{} premise fields carry exactly these types); E2 — certificate-data references are not derivation edges; E3 — the admittedness + acyclicity pass over the certificate\textquotesingle s own row graph (loader conformance, deliberately a check rather than a theory axiom: a cyclic certificate would chase-satisfy every axiom while proving nothing).}
\ruleblock{Lean statement}
\par\medskip
The discipline, mechanized at exactly the certificate-object level
the recipes state it. The certificate shape binds by display --- the four definitions
ARE the E1/E2 typing claims --- and the theorems bind as statements:
the E3 bridge, the induction template in well-founded and
finite-acyclic forms with root corollary, and the consumption-site
transport pair. The induction theorem proper
(\texttt{towerCitation\_sound\_of\_wf}) closes under no axioms at all.

\declhead{BernsteinLogicAux.Acceptance}{acyclic\_wellFounded}{\badge{BadgeGreen}{statement bound}}
\begin{leansrc}
/-- **The E3 bridge**: on a FINITE certificate, the loader's checkable
acyclicity pass yields exactly the well-foundedness the induction needs.
This is why E3 can be loader conformance rather than an axiom of the
theory: well-foundedness is an induction principle the logical framework
deliberately lacks, but on finite row sets it is equivalent to the
decidable absence of citation cycles. -/
theorem acyclic_wellFounded [Finite ι] (C : ι → Row σ J ι)
    (hacyc : Acyclic C) : WellFounded (cites C)
\end{leansrc}
\declhead{BernsteinLogicAux.Acceptance}{towerCitation\_sound\_of\_wf}{\badge{BadgeGreen}{statement bound}}
\begin{leansrc}
/-- **`rec:tc-template`, well-founded form** — the extension-edge
induction, in the shape the reflection theorem's induction consumes.
Hypotheses are the template's clauses: clause 1 = `hbase` (base
soundness; in deployment, the BL checker-soundness frontier — residue
R4); clause 2 = `hside` + `hstep` (the acceptance relation's local-check
clause as its own parameter, then the per-rule conditional-soundness
step in the cookie-cutter shape — the per-rule lemmas are CONSUMED here,
never proven: residue R1); clause 3 = `hwf` (the well-foundedness the E3
pass buys via `acyclic_wellFounded`). Conclusion: every admitted row's
conclusion-denotation holds — at each row the induction hypothesis
supplies exactly what the per-rule lemma consumes, and clause 1 grounds
the leaves. -/
theorem towerCitation_sound_of_wf
    (C : ι → Row σ J ι) (baseTrue : σ → Prop) (sem : J → Prop)
    (SideOK : ι → Prop)
    (hbase : ∀ i s, C i = Row.base s → baseTrue s)
    (hside : ∀ i, SideOK i)
    (hstep : ∀ i prems concl, C i = Row.ext prems concl → SideOK i →
      (∀ j ∈ prems, Row.holds baseTrue sem (C j)) → sem concl)
    (hwf : WellFounded (cites C)) :
    ∀ i, Row.holds baseTrue sem (C i)
\end{leansrc}
\declhead{BernsteinLogicAux.Acceptance}{towerCitation\_sound}{\badge{BadgeGreen}{statement bound}}
\begin{leansrc}
/-- **`rec:tc-template`, headline form** — the discipline as the checker
consumes it: a FINITE certificate whose citation graph passes the E3
acyclicity check, base soundness, local checks, and the per-rule
conditional lemmas jointly warrant every admitted row's
conclusion-denotation. This is the theorem that turns the book's
per-rule RELATIVE soundness lemmas into an unconditional verdict about
a whole certificate. -/
theorem towerCitation_sound [Finite ι]
    (C : ι → Row σ J ι) (baseTrue : σ → Prop) (sem : J → Prop)
    (SideOK : ι → Prop)
    (hbase : ∀ i s, C i = Row.base s → baseTrue s)
    (hside : ∀ i, SideOK i)
    (hstep : ∀ i prems concl, C i = Row.ext prems concl → SideOK i →
      (∀ j ∈ prems, Row.holds baseTrue sem (C j)) → sem concl)
    (hacyc : Acyclic C) :
    ∀ i, Row.holds baseTrue sem (C i)
\end{leansrc}
\declhead{BernsteinLogicAux.Acceptance}{towerCitation\_root}{\badge{BadgeGreen}{statement bound}}
\begin{leansrc}
/-- **Root soundness** (the acceptance relation's clause-4 shape): if
some admitted row of an accepted certificate concludes the root
judgment, that judgment's denotation holds. The abstract skeleton of
the reflection theorem's statement. -/
theorem towerCitation_root [Finite ι]
    (C : ι → Row σ J ι) (baseTrue : σ → Prop) (sem : J → Prop)
    (SideOK : ι → Prop)
    (hbase : ∀ i s, C i = Row.base s → baseTrue s)
    (hside : ∀ i, SideOK i)
    (hstep : ∀ i prems concl, C i = Row.ext prems concl → SideOK i →
      (∀ j ∈ prems, Row.holds baseTrue sem (C j)) → sem concl)
    (hacyc : Acyclic C)
    {Jroot : J} (hroot : ∃ i prems, C i = Row.ext prems Jroot) :
    sem Jroot
\end{leansrc}
\declhead{BernsteinLogicAux.Acceptance}{citePos\_sound}{\badge{BadgeGreen}{statement bound}}
\begin{leansrc}
/-- **`prop:tc-consume` (a), positive citations**: the child's concluded
region `R` under-approximates its formula's denotation `D`, and any
MONOTONE consuming claim proved over `R` transfers to `D`. Which
operators are monotone in which positions is each consuming rule's own
lemma's business; this lemma fixes only the direction. A certificate can
lose completeness at a citation, never soundness. -/
theorem citePos_sound {X : Type*} {R D : Set X} {P : Set X → Prop}
    (hmono : Monotone P) (hsub : R ⊆ D) (hR : P R) : P D
\end{leansrc}
\declhead{BernsteinLogicAux.Acceptance}{citeCmpl\_subset}{\badge{BadgeGreen}{statement bound}}
\begin{leansrc}
/-- **`prop:tc-consume` (b), the orientation swap**: an
under-approximating certified region's complement OVER-approximates the
denotation's complement — the conservative direction where regions
appear negatively. (Set-inclusion twin of
`BernsteinLogic.Structural.compl_over_approx`, which states the same
direction through the compiled region-judgment reading.) -/
theorem citeCmpl_subset {X : Type*} {R D : Set X} (hsub : R ⊆ D) :
    Dᶜ ⊆ Rᶜ
\end{leansrc}
\declhead{BernsteinLogicAux.Acceptance}{citeCmpl\_sound}{\badge{BadgeGreen}{statement bound}}
\begin{leansrc}
/-- **`prop:tc-consume` (b), consumer form**: an ANTITONE claim proved on
the citable complement `Rᶜ` covers the true violation set `Dᶜ` — a
premise row proved on, or a probability bound proved for reaching, the
over-approximation covers the truth. -/
theorem citeCmpl_sound {X : Type*} {R D : Set X} {P : Set X → Prop}
    (hanti : Antitone P) (hsub : R ⊆ D) (hR : P Rᶜ) : P Dᶜ
\end{leansrc}
\declhead{BernsteinLogicAux.Acceptance}{Row}{\badge{BadgeGreen}{definition bound}}
\begin{leansrc}
/-- **E1, as a type**: one row of a decoded certificate. A `base` row is a
base-calculus proof leaf concluding the sequent `s` (the `⊢` side); an
`ext` row is an extension row carrying its premise slots — references to
other rows, the derivation edges, and nothing else — and its concluded
satisfaction judgment (the `⊨` side). A premise slot can land on nothing
but another admitted row, consumed as that row's conclusion. -/
inductive Row (σ J ι : Type*) where
  | base (s : σ)
  | ext (prems : List ι) (concl : J)
\end{leansrc}
\declhead{BernsteinLogicAux.Acceptance}{Row.cited}{\badge{BadgeGreen}{definition bound}}
\begin{leansrc}
/-- **E2, in one definition**: the derivation edges out of a row are its
premise slots and nothing else — the induction metric counts these only.
Certificate-data references (shared certificate families, lattice
objects) are unrepresented here by design; their erasure is the loader's
conformance obligation, not this module's. -/
def Row.cited : Row σ J ι → List ι
  | .base _ => []
  | .ext prems _ => prems
\end{leansrc}
\declhead{BernsteinLogicAux.Acceptance}{Row.holds}{\badge{BadgeGreen}{definition bound}}
\begin{leansrc}
/-- The denotation of a row's CONCLUSION, under a kernel semantics
`baseTrue` for base sequents and a judgment semantics `sem` (the
judgment-denotation interface, kept abstract — zero measure theory in
this module by design). The two-case shape is E1's XOR: an edge lands on
kernel truth or on a judgment denotation; there is no third thing to
cite. -/
def Row.holds (baseTrue : σ → Prop) (sem : J → Prop) : Row σ J ι → Prop
  | .base s => baseTrue s
  | .ext _ concl => sem concl
\end{leansrc}
\declhead{BernsteinLogicAux.Acceptance}{cites}{\badge{BadgeGreen}{definition bound}}
\begin{leansrc}
/-- Row `j` is cited by row `i`: a derivation edge from `i` down to `j`.
The well-founded induction descends along this relation. -/
def cites (C : ι → Row σ J ι) (j i : ι) : Prop :=
  j ∈ (C i).cited
\end{leansrc}
\declhead{BernsteinLogicAux.Acceptance}{Acyclic}{\badge{BadgeGreen}{definition bound}}
\begin{leansrc}
/-- **E3's checkable half**, as the induction consumes it: no row
transitively cites itself. On a finite certificate this is precisely
acyclicity of the citation graph — the decidable pass the loader runs. -/
def Acyclic (C : ι → Row σ J ι) : Prop :=
  ∀ i, ¬ Relation.TransGen (cites C) i i
\end{leansrc}
\ruleblock{Correspondence (Steele display $\leftrightarrow$ Lean)}
\absentbox{NO CORRESPONDENCE DICTIONARY YET}{no reviewer-certified dictionary maps the displayed rule onto the mechanized statements above; until one is written and certified, that mapping rests on this section\textquotesingle s editorial care alone.}
\ruleblock{Lean verifier receipts}
\sourcenote{module \texttt{BernsteinLogicAux.Acceptance}; bound build \texttt{e78746b9f98a\ldots} (bound from receipt 3d1fca564d6e11dc, 2026-07-20T03:41:07Z). A receipt records that named checking layers accepted a build: (module, compiled-artifact SHA-256, environment SHA-256, checker layers), minted when the checking gate ran; the store receipts attesting the bound build are listed with the layers each run credited. Nothing here asks the reader to trust the store: every layer is an independently runnable program over the quoted sources --- rebuild the module and re-run the checkers to reproduce the verdict. Receipts attesting superseded builds credit no layer and are omitted from this rendering (the count of omissions, if any, is noted below).}
\begin{receipttable}
\receiptrow{3d1fca564d6e11dc}{2026-07-20T03:41:07Z}{build+decorrelated}{\badge{BadgeGreen}{current}}
\receiptrow{3546fbecf9c242a2}{2026-07-20T20:46:26Z}{build+decorrelated}{\badge{BadgeGreen}{current}}
\receiptrow{e684b09bdf01a640}{2026-07-20T20:57:59Z}{build+replay+hardened+decorrelated}{\badge{BadgeGreen}{current}}
\receiptrow{fe82a04ec14e7cee}{2026-07-21T00:02:22Z}{build+decorrelated}{\badge{BadgeGreen}{current}}
\receiptrow{f0f540e41c59285c}{2026-07-21T05:38:30Z}{build+replay+hardened+decorrelated}{\badge{BadgeGreen}{current}}
\receiptrow{69a2515828782208}{2026-07-22T18:26:47Z}{build+decorrelated}{\badge{BadgeGreen}{current}}
\receiptrow{bd7c9443b7f8649a}{2026-07-31T21:59:38Z}{build+decorrelated}{\badge{BadgeGreen}{current}}
\receiptrow{652943c163b17f25}{2026-07-31T22:28:31Z}{build+decorrelated}{\badge{BadgeGreen}{current}}
\receiptrow{dcd3934b54c0db35}{2026-08-02T23:30:39Z}{build+decorrelated}{\badge{BadgeGreen}{current}}
\receiptrow{5d9475b9d8bbe3c2}{2026-08-03T16:26:43Z}{build+decorrelated}{\badge{BadgeGreen}{current}}
\receiptrow{cbbad707861e8740}{2026-08-03T22:45:26Z}{build+decorrelated}{\badge{BadgeGreen}{current}}
\end{receipttable}
\layerline{\layeritem{build}{lake build (elaboration)}{\badge{BadgeGreen}{current}}\layeritem{replay}{leanchecker (official kernel replay)}{\badge{BadgeGreen}{current}}\layeritem{hardened}{safe\_verify (hardened screen)}{\badge{BadgeGreen}{current}}\layeritem{decorrelated}{lean4export + nanoda (decorrelated kernel)}{\badge{BadgeGreen}{current}}}
\sourcenote{\badge{BadgeGray}{comparator: none} no fixture in the three-implementation comparator harness (compiled / kernel-interpreted / reference-interpreted) exercises this rule.}
\ruleblock{HOL4 statement}
\absentbox{OUT OF SCOPE BY DESIGN}{Out of scope by structure: there is no rule statement to port --- the
discipline concludes nothing and has no check function of its own. Its
one checkable obligation, the admittedness + acyclicity pass, is
certificate-loader conformance, which would arrive with a HOL-side
extension-theory loader as conformance code, not as a rule check.
Independently, the HOL4 development\textquotesingle s object today is
the base calculus\textquotesingle s wire-format checker (theory
\texttt{bern\_catalogue}, the base catalogue\textquotesingle s
check-function cone); no extension-theory HOL development exists for
such a pass to live in.
}
\ruleblock{HOL verifier receipts}
\absentbox{NONE}{no HOL verification receipt exists for this rule.}
\ruleblock{Proof}

\sourcenote{Genre: edge discipline (glue) --- the extension
page\textquotesingle s fourth category, deliberately outside the three
tiers. This entry is not a rule: it has no conclusion judgment, so
``semantic soundness of the conclusion'' does not type-check against
it, and it has nothing to unpack, so neither rewriting genre applies
--- it \emph{is} the unpacking mechanism. What it owes the book
instead is the warrant every other section of the extension page
silently leans on: each rule\textquotesingle s soundness lemma is
proved \emph{relative} to base soundness, and this discipline is how
those conditional lemmas compose into a single verdict about a whole
certificate --- unconditional once the per-rule lemmas, base
soundness, and the loader\textquotesingle s conformance passes are in
hand; the closing remark names what remains conditional. The presentation apes the recipe form: the general
induction template, then the discipline\textquotesingle s obligations
as the instance that makes the template apply.}

\begin{brecipe}[what the discipline licenses]\label{rec:tc-license}
A concluded child judgment
$\mathcal{M},\ x_0 \in R \vDash \varphi_{\mathrm{child}}$ (the region-level
satisfaction kind: every point of $R$ satisfies
$\varphi_{\mathrm{child}}$, with $R$ under-approximating the
formula\textquotesingle s denotation) makes $R$ \emph{citable} as a
region atom in any parent row\textquotesingle s scope or target. The
working example: a nested probability former
$\mathsf{SProbGE}(p, \psi)$ (surface syntax $\mathsf{P}_{\ge p}[\psi]$)
inside a state formula compiles its child
first and cites the child\textquotesingle s concluded region as the
sublevel atom $\{V_{\mathrm{child}} \le 1 - p\}$. The license is an
\emph{edge} between rows of one certificate, never a node: the
citation itself concludes nothing.
\end{brecipe}

\begin{brecipe}[the induction template]\label{rec:tc-template}
Let a certificate be a finite set of admitted rows: base-calculus
proof rows (leaves, concluding sequents on the $\vdash$ side) and
extension rows (concluding satisfaction judgments on the $\vDash$
side), each extension row\textquotesingle s premise fields referencing
other rows; call a row \emph{admitted} when its local checks pass ---
its premise fields are well typed and its side conditions recompute
true. Suppose:
\begin{enumerate}
\item \emph{base soundness} --- every admitted base row\textquotesingle s
  sequent is true in kernel semantics;
\item \emph{per-rule conditional soundness} --- for every extension
  rule sort $R$ a lemma $\mathrm{sound}_R$: if the referenced child
  judgments\textquotesingle{} denotations hold, the base leaves are
  true, and the side conditions recompute true, then the concluded
  judgment\textquotesingle s denotation holds;
\item \emph{acyclicity} --- the certificate\textquotesingle s citation
  graph is acyclic.
\end{enumerate}
Then every admitted row\textquotesingle s judgment denotation holds:
induct over the citation order, which clause 3 (with finiteness)
makes well-founded; at each row the induction hypothesis supplies the
child denotations, admittedness supplies the recomputed side
conditions, and clause 1 grounds the leaves --- together exactly what
clause 2\textquotesingle s lemma consumes.
The conditional lemmas are the per-rule sections of this book; the
induction is this section.
\end{brecipe}

\begin{brecipe}[the instance --- the discipline\textquotesingle s three
obligations]\label{rec:tc-instance}
The discipline is staged as \emph{typing and conformance}, not as a
node sort, in three obligations, each of which is one premise of
Recipe~\ref{rec:tc-template} made checkable:
\begin{itemize}
\item \textbf{E1 (edge typing).} Every premise field of an extension
  row is typed as a satisfaction-judgment kind or as an inherited
  base-proof leaf, exclusively --- never as a raw rule row of either
  theory. This makes the induction\textquotesingle s case analysis
  exhaustive: every edge lands on something clause 1 or the induction
  hypothesis covers; no edge can cite an object that has no soundness
  account.
\item \textbf{E2 (data/edge separation).} Certificate-\emph{data}
  references --- a certificate family reused across rows, a lattice
  object --- are not derivation edges. The induction metric counts
  derivation edges only, so shared data can neither smuggle a cycle
  past the acyclicity pass nor impose ordering the derivation
  structure does not have.
\item \textbf{E3 (admittedness + acyclicity).} A loader pass over the
  certificate\textquotesingle s own row graph checks that every
  referenced row is admitted and that the citation graph is acyclic
  --- clause 3 verbatim. This is \emph{conformance}, deliberately not
  an axiom of the theory: an existence-of-concluding-proof axiom would
  let the theory\textquotesingle s own saturation fabricate proof
  rows, and well-foundedness is an induction principle, which the
  logical framework deliberately lacks. A cyclic certificate would
  satisfy every local check while proving nothing --- or worse: take
  two rows, each citing the other\textquotesingle s concluded
  judgment as its sole premise, with per-rule lemmas of the shape
  $J_B \Rightarrow J_A$ and $J_A \Rightarrow J_B$. In an instance
  where both judgments are false, both implications hold, every
  per-row obligation and the admittedness pass discharge, and only
  acyclicity objects; accepting the pair would not merely warrant
  nothing --- it would assert falsehoods.
\end{itemize}
\end{brecipe}

\begin{bprop}[consumption-site soundness]\label{prop:tc-consume}
Let the child conclude $\mathcal{M},\ x_0 \in R \vDash \varphi$ with
$R \subseteq [\![\varphi]\!]$. Then at every consumption site the
citation is used in the sound direction:
\begin{enumerate}
\item[(a)] a \emph{positive} citation passes $R$ exactly, and any
  monotone claim over $R$ (``the state reaches / remains in $R$'')
  entails the same claim for $[\![\varphi]\!]$;
\item[(b)] a \emph{complement} citation swaps orientation:
  $\mathrm{cmpl}(R) \supseteq \mathrm{cmpl}([\![\varphi]\!])$, so a
  premise row proved on $\mathrm{cmpl}(R)$ --- or an \emph{upper}
  probability bound proved for reaching it --- covers the true
  violation set.
\end{enumerate}
\end{bprop}
\begin{proof}
(a) is monotonicity of the consuming operator in
$R \subseteq [\![\varphi]\!]$; (b) is the same inclusion complemented.
The discipline contributes exactly this inclusion and its swap;
\emph{which} operators are monotone in which positions is each
consuming rule\textquotesingle s own lemma\textquotesingle s business,
stated and proved in that rule\textquotesingle s section.
Under-approximation on the positive side and over-approximation on the
complement side are both conservative: a certificate can lose
completeness at a citation, never soundness.
\end{proof}

\begin{bremark}[why this is not a rule, and what the mechanization is]
The three tiers each predicate on a conclusion: core is the concluded
judgment\textquotesingle s semantic soundness; admissible and alias
are verified unpackings of a concluded node. A discipline with no
conclusion satisfies none of the three definitions --- it is the
induction step the tiers\textquotesingle{} lemmas plug into, which is
why it renders under its own category head rather than wearing a tier
that does not type-check. Its one \emph{checkable} obligation (E3)
lands as certificate-loader conformance code. The warrant itself is
now mechanized at exactly the certificate-object level this
section\textquotesingle s recipes state it
(\texttt{BernsteinLogicAux/Acceptance.lean}), anchor by anchor:
Recipe~\ref{rec:tc-template} is \texttt{towerCitation\_sound} ---
well-founded induction over admitted rows, consuming the per-rule
lemmas as its step hypothesis --- in finite-acyclic (headline) and
well-founded forms, with root corollary
\texttt{towerCitation\_root}; the E1 typing of
Recipe~\ref{rec:tc-instance} is the row type itself, so the
exhaustiveness of the induction\textquotesingle s case analysis is
type-checked (two consumptions exist: kernel truth of a base leaf,
denotation of a concluded judgment --- no third case to cite); E2 is
the edge metric \texttt{Row.cited} (premise slots only); E3\textquotesingle s
sufficiency is the bridge \texttt{acyclic\_wellFounded} --- on a
finite certificate the checkable acyclicity pass yields exactly the
well-foundedness the induction needs, which is why conformance
suffices where the logical framework deliberately lacks an induction
axiom; and Proposition~\ref{prop:tc-consume} is the transport pair
\texttt{citePos\_sound} / \texttt{citeCmpl\_subset} /
\texttt{citeCmpl\_sound}. The induction theorem proper
(\texttt{towerCitation\_sound\_of\_wf}) depends on no axioms at all;
the bridge and transport lemmas close under the standard three.
\end{bremark}

\begin{bremark}[what remains conditional, named]
This discipline\textquotesingle s mechanization is conditional on
named residues rather than a vague gap.
\textbf{R1} --- the per-rule lemmas enter the induction as a
hypothesis: that is the discipline\textquotesingle s literal content
(it composes conditional lemmas; it does not prove them), but the
type-level identification of that hypothesis\textquotesingle s
instances with this book\textquotesingle s bound per-rule statement
family awaits the deep embedding and the per-rule schema wrappers;
until then the identification is editorial, rule by rule.
\textbf{R2} --- E2 is discharged by \emph{representation}: the
abstract row type carries derivation edges only, and the faithfulness
of the concrete certificate format\textquotesingle s erasure onto it
(no derivation edge lost, no data reference kept) is loader/decode
conformance, outside the module. \textbf{R3} --- admittedness
(E3\textquotesingle s first half) is absorbed by the totality of the
abstract certificate\textquotesingle s indexing: every reference IS an
admitted row, by type; the loader\textquotesingle s admittedness check
is what establishes that totality. (Clause~1, base soundness, is the
base calculus\textquotesingle s own soundness frontier and enters as
the named hypothesis it always was --- a premise of the template, not
a residue of the mechanization.)
\end{bremark}



\begingroup
\renewcommand{\bibname}{References}
\renewcommand{\chapter}{\section}
\pdfbookmark[1]{References}{blcat-refs}%
\clearpage
\begin{thebibliography}{99}
\footnotesize

\bibitem{farkas1902}
J.~Farkas.
\emph{Theorie der einfachen Ungleichungen}.
Journal f\"ur die reine und angewandte Mathematik 124 (1902), 1--27.

\bibitem{schrijver1986}
A.~Schrijver.
\emph{Theory of Linear and Integer Programming}.
Wiley, 1986. (Farkas variants and multiplier sign conventions: \S7.)

\bibitem{krivine1964}
J.-L.~Krivine.
\emph{Anneaux pr\'eordonn\'es}.
Journal d'Analyse Math\'ematique 12 (1964), 307--326.

\bibitem{handelman1988}
D.~Handelman.
\emph{Representing polynomials by positive linear functions on compact
convex polyhedra}.
Pacific Journal of Mathematics 132 (1988), no.~1, 35--62.

\bibitem{tellegen1952}
B.~D.~H.~Tellegen.
\emph{A general network theorem, with applications}.
Philips Research Reports 7 (1952), 259--269.

\bibitem{willems1972}
J.~C.~Willems.
\emph{Dissipative dynamical systems, Part I: General theory}.
Archive for Rational Mechanics and Analysis 45 (1972), 321--351.

\bibitem{moylanhill1978}
P.~J.~Moylan and D.~J.~Hill.
\emph{Stability criteria for large-scale systems}.
IEEE Transactions on Automatic Control 23 (1978), no.~2, 143--149.

\bibitem{vanderschaftjeltsema2014}
A.~van~der~Schaft and D.~Jeltsema.
\emph{Port-Hamiltonian Systems Theory: An Introductory Overview}.
Foundations and Trends in Systems and Control 1 (2014), no.~2--3, 173--378.

\bibitem{jiangteelpraly1994}
Z.-P.~Jiang, A.~R.~Teel, and L.~Praly.
\emph{Small-gain theorem for ISS systems and applications}.
Mathematics of Control, Signals and Systems 7 (1994), 95--120.

\bibitem{jiangmareelswang1996}
Z.-P.~Jiang, I.~M.~Y.~Mareels, and Y.~Wang.
\emph{A Lyapunov formulation of the nonlinear small-gain theorem for
interconnected ISS systems}.
Automatica 32 (1996), no.~8, 1211--1215.

\bibitem{astromfuruta2000}
K.~J.~{\AA}str\"om and K.~Furuta.
\emph{Swinging up a pendulum by energy control}.
Automatica 36 (2000), no.~2, 287--295.

\bibitem{steele2017}
G.~L.~Steele~Jr.
\emph{It's Time for a New Old Language}.
Keynote, ACM SIGPLAN Symposium on Principles and Practice of Parallel
Programming (PPoPP), 2017. (The ``Computer Science Metanotation'' this
book's rule cells are set in.)


\bibitem{gentzen1935}
G.~Gentzen.
\emph{Untersuchungen \"uber das logische Schlie{\ss}en I, II}.
Mathematische Zeitschrift 39 (1935), 176--210 and 405--431. English
translation in M.~E.~Szabo (ed.), \emph{The Collected Papers of Gerhard
Gentzen}, North-Holland, 1969.

\bibitem{girard1987}
J.-Y.~Girard.
\emph{Linear logic}.
Theoretical Computer Science 50 (1987), no.~1, 1--101. (Source of the
additive/multiplicative vocabulary for context-sharing vs.\
context-splitting rule formats.)

\bibitem{birkhoff1967}
G.~Birkhoff.
\emph{Lattice Theory}.
3rd ed., American Mathematical Society Colloquium Publications 25, 1967.
(M$_3$, the diamond, as the canonical modular non-distributive lattice.)

\bibitem{rudin1976}
W.~Rudin.
\emph{Principles of Mathematical Analysis}.
3rd ed., McGraw--Hill, 1976. (Taylor's theorem with the Lagrange form
of the remainder: Theorem 5.15; the nonnegative-remainder enclosure at
nonnegative argument used by the \rulename{TaylorBound} section.)

\bibitem{banach1922}
S.~Banach.
\emph{Sur les op\'erations dans les ensembles abstraits et leur
application aux \'equations int\'egrales}.
Fundamenta Mathematicae 3 (1922), 133--181.

\bibitem{granasdugundji2003}
A.~Granas and J.~Dugundji.
\emph{Fixed Point Theory}.
Springer Monographs in Mathematics, Springer, 2003. (The closed-ball
condition $d(Tx_0, x_0) \leq (1 - \rho)\,r$ and the a-priori bound
$d(x_0, x^*) \leq \delta/(1-\rho)$: \S I.1.)

\bibitem{doeblin1938}
W.~Doeblin.
\emph{Expos\'e de la th\'eorie des cha\^ines simples constantes de
Markov \`a un nombre fini d'\'etats}.
Revue Math\'ematique de l'Union Interbalkanique 2 (1938), 77--105.
(The origin of the coupling method for Markov chains.)

\bibitem{lindvall1992}
T.~Lindvall.
\emph{Lectures on the Coupling Method}.
Wiley, 1992; corrected reprint, Dover, 2002. (The coupling inequality
$\mathrm{TV} \leq \Pr[X \neq Y]$: \S I.2.)

\bibitem{levinperes2017}
D.~A.~Levin and Y.~Peres, with E.~L.~Wilmer.
\emph{Markov Chains and Mixing Times}.
2nd edition, American Mathematical Society, 2017. (Total-variation
distance, couplings, and the data-processing/projection inequality:
Ch.~4--5.)

\bibitem{hsu2017}
J.~Hsu.
\emph{Probabilistic Couplings for Probabilistic Reasoning}.
PhD thesis, University of Pennsylvania, 2017. arXiv:1710.09951.
(Couplings as certificates in relational program logics --- the
pre/post-style sibling of the judgment form used here.)

\bibitem{girardpappas2007}
A.~Girard and G.~J.~Pappas.
\emph{Approximation metrics for discrete and continuous systems}.
IEEE Transactions on Automatic Control 52 (2007), no.~5, 782--798.
(Approximate simulation with quantitative degradation --- the analytic
ancestor of the $\varepsilon$-witness transport cell.)

\bibitem{desharnais2004}
J.~Desharnais, V.~Gupta, R.~Jagadeesan, and P.~Panangaden.
\emph{Metrics for labelled Markov processes}.
Theoretical Computer Science 318 (2004), no.~3, 323--354.
(Behavioral pseudometrics for probabilistic systems.)

\bibitem{weyl1939}
H.~Weyl.
\emph{The Classical Groups: Their Invariants and Representations}.
Princeton University Press, 1939. (Polarized power sums generate the
multisymmetric invariants in characteristic 0.)

\bibitem{murata1989}
T.~Murata.
\emph{Petri nets: Properties, analysis and applications}.
Proceedings of the IEEE 77 (1989), no.~4, 541--580. (Place invariants
--- the linear conservation laws entering as free count rows.)

\bibitem{spivak2021agg}
D.~I.~Spivak, R.~Garner, and A.~D.~Fairbanks.
\emph{Functorial aggregation}.
arXiv:2111.10968. (Aggregation as commutative-monoid folds along a
classifying map; the multiset as the universal aggregate.)

\bibitem{dashkovskiyrufferwirth2007}
S.~Dashkovskiy, B.~S.~R\"uffer, F.~R.~Wirth.
An ISS small gain theorem for general networks.
\emph{Mathematics of Control, Signals, and Systems} 19(2):93--122,
2007. DOI 10.1007/s00498-007-0014-8. (The $n$-component ISS small-gain
theorem over arbitrary interconnection graphs whose certificate form
the \rulename{SmallGainN} section catalogues; its monotone-operator
fixpoint machinery lives producer-side as existence theory for the
$(\lambda, B, \theta)$ side data.)

\bibitem{nelsonoppen1980}
G.~Nelson and D.~C.~Oppen.
\emph{Fast decision procedures based on congruence closure}.
Journal of the ACM 27 (1980), no.~2, 356--364.

\bibitem{nieuwenhuisoliveras2007}
R.~Nieuwenhuis and A.~Oliveras.
\emph{Fast congruence closure and extensions}.
Information and Computation 205 (2007), no.~4, 557--574.

\bibitem{bernstein1912}
S.~N.~Bernstein.
\emph{D\'emonstration du th\'eor\`eme de Weierstrass fond\'ee sur le
calcul des probabilit\'es}.
Communications de la Soci\'et\'e Math\'ematique de Kharkov (2) 13
(1912), 1--2.

\bibitem{cargoshisha1966}
G.~T.~Cargo and O.~Shisha.
\emph{The Bernstein form of a polynomial}.
Journal of Research of the National Bureau of Standards 70B (1966),
79--81. (Bernstein coefficients enclose the range on the simplex ---
the sound direction this calculus's kernel uses.)

\bibitem{farouki2012}
R.~T.~Farouki.
\emph{The Bernstein polynomial basis: a centennial retrospective}.
Computer Aided Geometric Design 29 (2012), no.~6, 379--419.

\bibitem{decasteljau1959}
P.~de~Casteljau.
\emph{Outillages m\'ethodes calcul}.
Technical report, A.~Citro\"en, Paris, 1959. (The founding
subdivision/evaluation algorithm for polynomials in Bernstein form.)

\bibitem{laneriesenfeld1980}
J.~M.~Lane and R.~F.~Riesenfeld.
\emph{A theoretical development for the computer generation and display
of piecewise polynomial surfaces}.
IEEE Transactions on Pattern Analysis and Machine Intelligence 2
(1980), no.~1, 35--46.

\bibitem{garloff1986}
J.~Garloff.
\emph{Convergent bounds for the range of multivariate polynomials}.
In: Interval Mathematics 1985 (K.~Nickel, ed.), Lecture Notes in
Computer Science 212, Springer, 1986, 37--56. (Subdivision sharpens
Bernstein enclosures --- the completeness partner of the enclosure
soundness used here.)

\bibitem{ehlers2017}
R.~Ehlers.
\emph{Formal verification of piece-wise linear feed-forward neural
networks}.
Proceedings of ATVA 2017 (Automated Technology for Verification and
Analysis), 2017. (The ``triangle'' linear relaxation of ReLU.)

\bibitem{wongkolter2018}
E.~Wong and J.~Z.~Kolter.
\emph{Provable defenses against adversarial examples via the convex
outer adversarial polytope}.
Proceedings of ICML 2018, PMLR 80, 2018.

\bibitem{zhang2018crown}
H.~Zhang, T.-W.~Weng, P.-Y.~Chen, C.-J.~Hsieh, and L.~Daniel.
\emph{Efficient neural network robustness certification with general
activation functions}.
Advances in Neural Information Processing Systems 31 (NeurIPS 2018).
(CROWN: adaptive two-sided linear bounds.)

\bibitem{singh2019deeppoly}
G.~Singh, T.~Gehr, M.~P\"uschel, and M.~Vechev.
\emph{An abstract domain for certifying neural networks}.
Proceedings of the ACM on Programming Languages 3 (POPL 2019),
Article 41. (DeepPoly.)

\bibitem{xu2020lirpa}
K.~Xu, Z.~Shi, H.~Zhang, Y.~Wang, K.-W.~Chang, M.~Huang,
B.~Kailkhura, X.~Lin, and C.-J.~Hsieh.
\emph{Automatic perturbation analysis for scalable certified
robustness and beyond}.
Advances in Neural Information Processing Systems 33 (NeurIPS 2020).
(The LiRPA framework and name.)

\bibitem{katz2017reluplex}
G.~Katz, C.~Barrett, D.~L.~Dill, K.~Julian, and M.~J.~Kochenderfer.
\emph{Reluplex: an efficient SMT solver for verifying deep neural
networks}.
Proceedings of CAV 2017 (Computer Aided Verification), 2017.

\bibitem{bunel2020branch}
R.~Bunel, J.~Lu, I.~Turkaslan, P.~H.~S.~Torr, P.~Kohli, and
M.~P.~Kumar.
\emph{Branch and bound for piecewise linear neural network
verification}.
Journal of Machine Learning Research 21 (2020).

\bibitem{dalrymple2026psd}
David ``davidad'' Dalrymple.
\emph{Certifying strict positive definiteness via an approximate
triangular inverse} (post). X (formerly Twitter), 2026, status
2022316806094913669. (The triangular-inverse method:
sign-and-pattern checks on $A \cdot Z$ for an untrusted triangular
witness $Z$; no divisions, witness accuracy never load-bearing.)

\bibitem{irving2026certificate}
G.~Irving.
\emph{Answer to: Certificate for strict positive definiteness of a
sparse matrix}.
Computational Science Stack Exchange, 2026 (question 45365, answer
45400). (The rounding-tolerant integer variant and its certificate-size
analysis.)

\bibitem{irvingaks}
G.~Irving.
\emph{aks} (Lean formalization repository): the \texttt{Random}
library's diagonal-dominance PSD bridge and a certified
$65536$-node spectral expander.
\texttt{github.com/girving/aks}, 2026.

\bibitem{gershgorin1931}
S.~Gerschgorin.
\emph{\"Uber die Abgrenzung der Eigenwerte einer Matrix}.
Bulletin de l'Acad\'emie des Sciences de l'URSS, Classe des sciences
math\'ematiques et naturelles 6 (1931), 749--754.

\bibitem{taussky1949}
O.~Taussky.
\emph{A recurring theorem on determinants}.
American Mathematical Monthly 56 (1949), no.~10, 672--676. (The
Levy--Desplanques strict-dominance nonsingularity theorem and its
history.)

\bibitem{boman2005}
E.~G.~Boman, D.~Chen, O.~Parekh, and S.~Toledo.
\emph{On factor width and symmetric H-matrices}.
Linear Algebra and its Applications 405 (2005), 239--248.
(Factor-width-2 = scaled diagonal dominance --- the class the /DD
tier's congruence certificates inhabit.)

\bibitem{hornjohnson2013}
R.~A.~Horn and C.~R.~Johnson.
\emph{Matrix Analysis}, 2nd edition.
Cambridge University Press, 2013. (Sylvester's criterion, congruence,
and Gershgorin discs, as cited in the certificate proofs.)

\bibitem{abadilamport1993}
M.~Abadi and L.~Lamport.
\emph{Composing Specifications}.
ACM Transactions on Programming Languages and Systems 15 (1993),
no.~1, 73--132.

\bibitem{mcmillan1999}
K.~L.~McMillan.
\emph{Circular Compositional Reasoning about Liveness}.
In: Correct Hardware Design and Verification Methods (CHARME 1999),
Lecture Notes in Computer Science 1703, Springer, 1999, 342--346.

\bibitem{lynchvaandrager1995}
N.~Lynch and F.~Vaandrager.
\emph{Forward and Backward Simulations, Part I: Untimed Systems}.
Information and Computation 121 (1995), no.~2, 214--233.

\bibitem{capuccimyers2026}
M.~Capucci and D.~J.~Myers.
Certified lenses and compositional system certification.
arXiv:2604.03017 (2026). (The wiring-square conditions, the
$\mathcal{K}^0_\infty$ slack discipline, and the
compose-by-lens-composition-and-slack-addition law cited in the
composition sections.)


\bibitem{ville1939}
J.~Ville.
\emph{\'Etude critique de la notion de collectif}.
Monographies des Probabilit\'es, Gauthier-Villars, Paris, 1939.
(The maximal inequality for nonnegative supermartingales --- the
probability engine of every battery rule.)

\bibitem{doob1953}
J.~L.~Doob.
\emph{Stochastic Processes}.
Wiley, New York, 1953. (Systematic supermartingale theory; maximal
inequalities and optional sampling.)

\bibitem{kushner1967}
H.~J.~Kushner.
\emph{Stochastic Stability and Control}.
Mathematics in Science and Engineering 33, Academic Press, New York,
1967. (Supermartingale safety certificates for stochastic systems.)

\bibitem{meyntweedie2009}
S.~P.~Meyn and R.~L.~Tweedie.
\emph{Markov Chains and Stochastic Stability}.
2nd edition, Cambridge University Press, 2009. (Foster--Lyapunov drift
conditions; the geometric drift form behind the deadline battery.)

\bibitem{prajnajadbabaiepappas2007}
S.~Prajna, A.~Jadbabaie, and G.~J.~Pappas.
\emph{A framework for worst-case and stochastic safety verification
using barrier certificates}.
IEEE Transactions on Automatic Control 52 (2007), no.~8, 1415--1428.

\bibitem{streett1982}
R.~S.~Streett.
\emph{Propositional dynamic logic of looping and converse is elementarily
decidable}.
Information and Control 54 (1982), no.~1--2, 121--141. (The
complemented-pairs acceptance condition.)

\bibitem{rabin1969}
M.~O.~Rabin.
\emph{Decidability of second-order theories and automata on infinite
trees}.
Transactions of the American Mathematical Society 141 (1969), 1--35.

\bibitem{landweber1969}
L.~H.~Landweber.
\emph{Decision problems for $\omega$-automata}.
Mathematical Systems Theory 3 (1969), no.~4, 376--384. (Deterministic
B\"uchi automata recognize exactly the limit languages.)

\bibitem{thomas1990}
W.~Thomas.
\emph{Automata on infinite objects}.
In: Handbook of Theoretical Computer Science, vol.~B, 133--191.
Elsevier, 1990.

\bibitem{foster1953}
F.~G.~Foster.
\emph{On the stochastic matrices associated with certain queuing
processes}.
Annals of Mathematical Statistics 24 (1953), no.~3, 355--360. (The
additive-drift criterion.)

\bibitem{meyntweedie1993}
S.~P.~Meyn and R.~L.~Tweedie.
\emph{Markov Chains and Stochastic Stability}.
Springer, 1993. (Foster--Lyapunov drift conditions, systematized.)

\bibitem{chakarov2013}
A.~Chakarov and S.~Sankaranarayanan.
\emph{Probabilistic program analysis with martingales}.
In: Computer Aided Verification (CAV 2013), Lecture Notes in Computer
Science 8044, 511--526. Springer, 2013. (Ranking supermartingales.)

\bibitem{durrett2019}
R.~Durrett.
\emph{Probability: Theory and Examples}.
5th edition, Cambridge University Press, 2019. (Borel--Cantelli.)

\bibitem{conley1978}
C.~Conley.
\emph{Isolated Invariant Sets and the Morse Index}.
CBMS Regional Conference Series in Mathematics 38, American Mathematical
Society, 1978.

\bibitem{conleyeaston1971}
C.~Conley and R.~Easton.
\emph{Isolated invariant sets and isolating blocks}.
Transactions of the American Mathematical Society 158 (1971), 35--61.

\bibitem{birkhoff1937}
G.~Birkhoff.
\emph{Rings of sets}.
Duke Mathematical Journal 3 (1937), no.~3, 443--454. (The representation
theorem for finite distributive lattices via join-irreducibles.)

\bibitem{nagumo1942}
M.~Nagumo.
\emph{\"Uber die Lage der Integralkurven gew\"ohnlicher
Differentialgleichungen}.
Proceedings of the Physico-Mathematical Society of Japan (3) 24 (1942),
551--559.

\bibitem{blanchini1999}
F.~Blanchini.
\emph{Set invariance in control}.
Automatica 35 (1999), no.~11, 1747--1767.

\bibitem{prajnajadbabaie2004}
S.~Prajna and A.~Jadbabaie.
\emph{Safety verification of hybrid systems using barrier certificates}.
Hybrid Systems: Computation and Control (HSCC 2004), LNCS 2993,
Springer, 2004, 477--492.

\bibitem{kmv2005}
W.~D.~Kalies, K.~Mischaikow, and R.~C.~A.~M.~van~der~Vorst.
\emph{An algorithmic approach to chain recurrence}.
Foundations of Computational Mathematics 5 (2005), no.~4, 409--449.

\bibitem{kmv2016}
W.~D.~Kalies, K.~Mischaikow, and R.~C.~A.~M.~van~der~Vorst.
\emph{Lattice structures for attractors II}.
Foundations of Computational Mathematics 16 (2016), 1151--1191.

\bibitem{kmv2022}
W.~D.~Kalies, K.~Mischaikow, and R.~C.~A.~M.~van~der~Vorst.
\emph{Lattice structures for attractors III}.
Journal of Dynamics and Differential Equations 34 (2022), 1729--1768.
arXiv:1911.09382.

\bibitem{mischaikowmrozek2002}
K.~Mischaikow and M.~Mrozek.
\emph{Conley index}.
Handbook of Dynamical Systems, Vol.~2, North-Holland, 2002, 393--460.

\bibitem{kalieswehbe2026}
W.~D.~Kalies and T.~Wehbe.
\emph{Compositionality of Global Dynamics in Product and Skew-Product
Systems}.
arXiv:2606.30925, 2026.

\bibitem{statman1979}
R.~Statman.
\emph{Lower bounds on Herbrand's theorem}.
Proceedings of the American Mathematical Society 75 (1979), no.~1,
104--107. (Non-elementary lower bounds on quantifier-inference
elimination; with \cite{orevkov1979}, the classical witness for the
cut-elimination blow-up.)

\bibitem{orevkov1979}
V.~P.~Orevkov.
\emph{Lower bounds for lengthening of proofs after cut-elimination}.
Zapiski Nauchnykh Seminarov LOMI 88 (1979), 137--162; English
translation in Journal of Soviet Mathematics (1982).
\end{thebibliography}

\endgroup
\endgroup
